% Tensor-Field Simulation: How the Synthesis Surface Simulates Its Targets
% Build: latexmk -pdf main.tex   (pdflatex + bibtex; TeX Live 2023)
\documentclass[11pt,letterpaper,oneside,openany]{memoir}

\usepackage{style/simbook}
\usepackage{style/simtikz}

% ── memoir page layout ───────────────────────────────────────────────────────
\setlrmarginsandblock{1.15in}{1.15in}{*}
\setulmarginsandblock{1.1in}{1.25in}{*}
\checkandfixthelayout
\setsecnumdepth{subsection}
\maxtocdepth{section}

% memoir chapter style
\chapterstyle{veelo}
\setlength{\beforechapskip}{0pt}

% running heads
\makepagestyle{simhead}
\makeevenhead{simhead}{\thepage}{}{\small\itshape\leftmark}
\makeoddhead{simhead}{\small\itshape\rightmark}{}{\thepage}
\makeevenfoot{simhead}{}{}{}
\makeoddfoot{simhead}{}{}{}
\pagestyle{simhead}

\title{\Huge\bfseries Tensor-Field Simulation\\[2mm]
  \Large How the Synthesis Surface Simulates Its Targets}
\author{The Palace Whiteroom}
\date{Draft --- \today}

\includeonly{
  frontmatter/preface,
  frontmatter/howtoread,
  frontmatter/notation,
  part1-orientation/ch01-simulator-and-targets,
  part1-orientation/ch02-lifecycle,
  part1-orientation/ch03-big-idea,
  part1-orientation/ch04-situating,
  part2-fieldtheory/ch05-maxwell,
  part2-fieldtheory/ch06-reductions,
  part2-fieldtheory/ch07-weak,
  part2-fieldtheory/ch08-galerkin,
  part2-fieldtheory/ch09-functionspaces,
  part2-fieldtheory/ch10-derham,
  part2-fieldtheory/ch11-refelement,
  part2-fieldtheory/ch12-assembly,
  part2-fieldtheory/ch13-matrixfree,
  part2-fieldtheory/ch14-bc,
  part3-calculus/ch15-language,
  part3-calculus/ch16-records,
  part3-calculus/ch17-tensors,
  part3-calculus/ch18-shapes,
  part3-calculus/ch19-fold,
  part3-calculus/ch20-combinators,
  part3-calculus/ch21-monad,
  part3-calculus/ch22-strata,
  part3-calculus/ch23-operators,
  part3-calculus/ch24-rotations,
  part4-machinery/ch25-krylov,
  part4-machinery/ch26-cg,
  part4-machinery/ch27-gmres,
  part4-machinery/ch28-orthog,
  part4-machinery/ch29-givens,
  part4-machinery/ch30-precond,
  part4-machinery/ch31-smoothers,
  part4-machinery/ch32-multigrid,
  part4-machinery/ch33-eigenvalues,
  part4-machinery/ch34-krylovschur,
  part4-machinery/ch35-deflation,
  part4-machinery/ch36-time,
  part4-machinery/ch37-amr,
  part4-machinery/ch38-reductions,
  part5-modalities/ch39-electrostatics,
  part5-modalities/ch40-magnetostatics,
  part5-modalities/ch41-eigenmodes,
  part5-modalities/ch42-driven,
  part5-modalities/ch43-transient,
  part5-modalities/ch44-waveports,
  part5-modalities/ch45-fulltrace,
  part6-synthesis/ch46-synthlib,
  part6-synthesis/ch47-fivelibs,
  part6-synthesis/ch48-composing,
  part6-synthesis/ch49-wholemachine,
  appendices/appA-calculus,
  appendices/appB-shapes,
  appendices/appC-laws,
  appendices/appD-numerics,
  appendices/appE-kernel,
  appendices/appF-stack,
  appendices/appG-glossary,
}

\begin{document}

\frontmatter
\begin{titlingpage}
  \centering
  \vspace*{2cm}
  {\Huge\bfseries Tensor-Field Simulation\par}
  \vspace{6mm}
  {\Large How the Synthesis Surface Simulates Its Targets\par}
  \vspace{2cm}
  {\large An illustrated introduction to simulating Maxwell's equations\\
   as a calculus of immutable tensors and combinators\par}
  \vfill
  {\large The Palace Whiteroom\par}
  \vspace{4mm}
  {\itshape Draft --- \today\par}
\end{titlingpage}

\cleardoublepage
\tableofcontents*
\clearpage
\listoffigures*

\chapter*{Preface}
\addcontentsline{toc}{chapter}{Preface}
\markboth{Preface}{Preface}

This is a book about a single, surprisingly small idea: that a large physics
simulator---the kind that engineers use to design antennas, microwave filters,
superconducting qubits, and electric machines---can be understood, in its
entirety, as the repeated composition of a handful of mathematical building
blocks. The simulator we use as our running example is
\emph{Palace}, an open-source three-dimensional electromagnetic field solver.
But almost nothing in this book is specific to Palace. The shapes we will
draw---assemble, solve, reduce; fold, map, iterate---are the shapes of nearly
every field simulator ever written.

\section*{Why this book exists}

Most simulation software is written in the idiom of its era's hardware. The
classic finite-element code is a sea of nested loops that march over arrays,
overwriting memory in place: \texttt{for each element, for each quadrature
point, accumulate into the global matrix}. That idiom is fast, it is familiar,
and it is almost perfectly designed to hide what the program is actually
computing. To read such a code is to reverse-engineer a physical theory from a
memory-access pattern.

This book takes the opposite stance. We present the same algorithms, but in a
form where the \emph{mathematics is in the foreground and the memory management
has disappeared.} A field is a \emph{tensor}---an immutable array of numbers
with named axes. A simulation step is a \emph{pure function} from tensors to
tensors. A loop is a \emph{combinator}---a higher-order function like ``fold''
or ``map'' that captures the \emph{shape} of the iteration without committing to
how it is scheduled on a particular machine. In this language, a whole
electrostatic capacitance extraction is three lines; an eigenmode analysis is
one black-box call wrapped in a readout; a frequency sweep is a map over a list.

Why bother? Three reasons.

\begin{itemize}
  \item \textbf{Clarity.} When mutation is gone, the data flow of an algorithm
  is exactly its mathematical dependency graph. You can \emph{see} that
  conjugate-gradient touches the operator once per step, that the only thing a
  Krylov solver writes back each cycle is the current iterate, that the
  preconditioner is just an approximate inverse applied as a function.
  \item \textbf{Composability.} Because every piece is a pure function, pieces
  snap together. The same \texttt{solve} routine serves electrostatics,
  magnetostatics, and the inner loop of an eigensolver. We will spend the whole
  book discovering that there are really only \emph{four} ways the simulator
  ever solves anything.
  \item \textbf{Hardware independence.} Pure tensor functions are exactly what
  modern accelerator libraries (GPUs, tensor processors, automatic
  differentiation frameworks) want to consume. The functional rendering is not
  an academic nicety; it is the form in which the algorithm can be \emph{handed
  to a backend} and run, unchanged, on a graphics card.
\end{itemize}

\section*{Three ways to look at a simulator}

There are three honest ways to describe a program this size, and a good engineer
learns to switch between them fluently. This book deliberately uses all three.

\begin{description}
  \item[Top-down (what does it do?).] Start from the user's command---``compute
  the capacitance of this geometry''---and trace the work downward into smaller
  and smaller pieces. This is the view of \cref{ch:targets} and
  \cref{ch:lifecycle}: the six analyses Palace offers and the lifecycle that
  runs them.
  \item[Bottom-up (what is this piece?).] Start from a primitive---``what
  exactly does \texttt{fe\_assemble} compute?''---and build understanding
  upward. This is the view of \cref{part:fieldtheory} and
  \cref{part:machinery}: the finite-element theory and the numerical
  algorithms, each studied on its own terms.
  \item[As a synthesized library (what is the code?).] Look at the whole thing
  as one small, well-organized library of composed functions. This is the view
  of \cref{part:collection}: the \emph{synthesis surface}, where every analysis
  is reassembled from the parts as a few lines of pseudocode.
\end{description}

The book's title points at the third view because it is the one that ties the
other two together. ``The synthesis surface'' is the name we give to the
simulator rendered as a library; ``simulates its targets'' is the claim that
this small library, composed in the right way, reproduces every physical
analysis the original C\texttt{++} program performs. The book is the argument
for that claim, told slowly enough that a reader who has just finished
sophomore-year mathematics can follow every step.

\section*{What you will be able to do}

By the end you should be able to: read the simulator's algorithms in the
pseudo-language we use throughout; explain, for each of the six analyses, what
physical quantity it computes and which equation it discretizes; recognize the
four ``solve corners'' and the three ``reduction shapes'' that organize the
entire system; follow a finite-element assembly from a weak form to a global
operator; describe how conjugate gradients, GMRES, multigrid, and a
shift-and-invert eigensolver work and why each appears where it does; and read
the synthesized library as a coherent whole. None of this requires prior
exposure to functional programming, monads, or finite-element analysis---only
the willingness to meet a few new ideas one at a time.

\section*{A note on the diagrams}

This is an illustrated book on purpose. Nearly every idea here has a picture
that is easier to remember than its formula: the lifecycle pipeline, the de Rham
complex, the multigrid V-cycle, the matrix-free contraction chain, the
shift-and-invert spectrum. Where a diagram and an equation say the same thing,
we give both, and we recommend you trust the diagram first and use the equation
to make it precise.

\vspace{1em}
\noindent\hfill\textit{The Palace Whiteroom}\hfill\null

\chapter*{How to read this book}
\addcontentsline{toc}{chapter}{How to read this book}
\markboth{How to read this book}{How to read this book}

\section*{Who this book is for}

This book assumes you have completed roughly two years of an undergraduate
engineering or applied-mathematics curriculum. Concretely, we assume you are
comfortable with:

\begin{itemize}[nosep]
  \item \textbf{Vector calculus}---gradient, divergence, and curl; line, surface,
  and volume integrals; the divergence and Stokes theorems; integration by parts.
  \item \textbf{Linear algebra}---matrices and vectors, inner products and norms,
  orthogonality and projections, symmetric/positive-definite matrices,
  eigenvalues and eigenvectors, and the idea of solving $\mat{A}\vx=\vb$.
  \item \textbf{Ordinary differential equations}---linear systems, and the idea
  of marching a solution forward in time.
  \item \textbf{A first course in electromagnetism}---you have seen Maxwell's
  equations at least once, even if they are hazy.
\end{itemize}

We do \emph{not} assume any background in functional programming, type theory,
the finite-element method, or numerical iterative solvers. Those are precisely
what the book teaches. Where a topic needs more mathematics than the list above,
we build it up from the list above.

\section*{The shape of the book}

The book has six parts and a set of appendices, arranged so that each part rests
on the ones before it. \Cref{fig:roadmap} shows how they depend on one another;
you can read it as a suggested route.

\begin{figure}[t]
\centering
\begin{tikzpicture}[node distance=6mm and 10mm,
  pt/.style={stage, minimum width=34mm, minimum height=9mm},
  ap/.style={opbox, minimum width=34mm, minimum height=8mm, fill=simBgViolet, draw=simViolet}]
  \node[pt] (p1) {\textbf{I.}~Orientation};
  \node[pt, below left=of p1, xshift=-6mm] (p2) {\textbf{II.}~Field Theory};
  \node[pt, below right=of p1, xshift=6mm] (p3) {\textbf{III.}~The Calculus};
  \node[pt, below=18mm of p1] (p4) {\textbf{IV.}~Numerical Machinery};
  \node[pt, below=of p4] (p5) {\textbf{V.}~The Modalities};
  \node[pt, below=of p5] (p6) {\textbf{VI.}~The Collection};
  \node[ap, right=22mm of p4] (app) {Appendices:\\\footnotesize deep semantics};
  \draw[depedge] (p1) -- (p2);
  \draw[depedge] (p1) -- (p3);
  \draw[depedge] (p2) -- (p4);
  \draw[depedge] (p3) -- (p4);
  \draw[depedge] (p4) -- (p5);
  \draw[depedge] (p5) -- (p6);
  \draw[refedge] (app) -- (p3);
  \draw[refedge] (app) -- (p4);
  \node[cornertag, right=1mm of p1] {start here};
\end{tikzpicture}
\caption[Reading roadmap]{The parts and their dependencies. Solid arrows mean
``rests on''; the dotted arrows mark the appendices, which deepen the calculus
(Part~III) and the machinery (Part~IV) and can be read on demand.}
\label{fig:roadmap}
\end{figure}

\textbf{Part~I (Orientation)} is the map of the whole territory. It introduces
the six analyses, the program that runs them, and the single idea---fields as
tensors, steps as pure functions---that the rest of the book elaborates. Read it
first and in full; everything later refers back to it.

\textbf{Parts~II and~III} are the two halves of the foundation, and they are
independent of each other. \textbf{Part~II (The Field Theory)} builds the
physics and the finite-element method from your vector calculus.
\textbf{Part~III (The Abstract Framework)} builds the small functional calculus
from your linear algebra. A reader who wants the physics first can read
II~then~III; a reader who wants the computational language first can read
III~then~II. Both are needed before Part~IV.

\textbf{Part~IV (The Numerical Machinery)} is the algorithmic heart: how linear
systems, eigenproblems, time evolution, mesh adaptation, and physical-quantity
extraction actually run. \textbf{Part~V (The Constituent Modalities)} then walks
through the six analyses one at a time, each as a specific composition of the
machinery. \textbf{Part~VI (The Collection)} reassembles everything into the
synthesized library and steps back to view the whole machine.

The \textbf{appendices} hold the precise version of the computational
semantics---the formal grammar, the reduction and typing rules, the algebraic
laws, the floating-point conventions, and the inner details of the matrix-free
kernel. They are written with the same diagrams and discussion as the main text,
but they are reference material: reach for them when you want the exact rule
behind a picture, not on a first pass.

\section*{The order in which ideas are introduced}

Part~III teaches the functional calculus in a deliberate sequence, each idea
built only from the ones before it. If you have never programmed in this style,
follow this order and resist the urge to skip:

\begin{enumerate}[nosep]
  \item \emph{Values and immutability}---a tensor is a value you never modify.
  \item \emph{Higher-order functions}---functions that take and return functions.
  \item \emph{The fold}---one loop combinator, \texttt{iterate\_while}, from
  which every iteration is built.
  \item \emph{Records}---bundles of named fields, updated by making fresh copies.
  \item \emph{The state monad}---a disciplined way to thread one evolving piece
  of state so that its single point of change stays visible.
  \item \emph{Operators as values}---built once, applied many times.
  \item \emph{Shapes}---symbolic contracts on a tensor's axes.
  \item \emph{Rotations}---the idea that ties the layered representations together.
\end{enumerate}

\section*{The boxes}

Four kinds of boxed aside recur throughout, each with its own color so you can
find them at a glance:

\begin{keyidea}
A \textbf{Key idea} box distills the single most important takeaway of a section
into one or two sentences. If you remember nothing else from a section, remember
the box.
\end{keyidea}

\begin{pitfall}
A \textbf{Pitfall} box flags a tempting mistake---an algorithm that looks
correct but quietly computes the wrong thing, or a notation that means something
other than it appears to.
\end{pitfall}

\begin{physicsbox}
A \textbf{physical picture} box anchors an abstract construction to the concrete
device or experiment it models, so the mathematics never floats free of the
engineering.
\end{physicsbox}

\begin{notationbox}
A \textbf{Notation} box introduces or clarifies a piece of the pseudo-language
or the mathematical notation at the moment you first need it.
\end{notationbox}

\noindent Numbered \textbf{worked examples} carry a calculation all the way
through, by hand, on a small enough problem that you can check every step. They
are the rungs of the ladder; do them with pencil in hand.

\chapter*{A notation primer}
\addcontentsline{toc}{chapter}{A notation primer}
\markboth{A notation primer}{A notation primer}

Two notations run through this book: ordinary mathematics, set in the usual way,
and a \emph{pseudo-language} in which we write the simulator's algorithms. This
short chapter teaches you to read both. You do not need to memorize anything
here; treat it as a reference and come back when a symbol is unfamiliar.

\section*{Reading the pseudo-language}

The algorithms in this book are written in a small, deliberately readable
notation borrowed from typed functional programming. It is not a language you
will run; it is a language you will \emph{read}, the way a physicist reads a
Feynman diagram. It always appears in a shaded, monospaced box, like this:

\begin{lstlisting}[style=l4style]
-- a function's TYPE: "given a scalar and two tensors, return a tensor"
axpy :: Scalar -> Tensor[N] -> Tensor[N] -> Tensor[N]

-- its DEFINITION: y' = a*x + y  (computed as a fresh value, nothing overwritten)
axpy a x y = a .* x .+ y
\end{lstlisting}

A handful of conventions cover almost everything:

\begin{description}[style=nextline,leftmargin=1.5em]
  \item[\code{f :: A -> B}] is a \emph{type signature}: ``\code{f} is a function
  taking an \code{A} and producing a \code{B}.'' The arrow \code{->} reads
  ``maps to.''
  \item[Currying.] A function of several arguments is written
  \code{f :: A -> B -> C}, read ``given an \code{A}, then a \code{B}, produce a
  \code{C}.'' Arguments are supplied one at a time, separated by spaces:
  \code{f x y}. This lets us supply some arguments now and the rest later---the
  basis of ``build once, apply many times.''
  \item[Records] are bundles of named fields, written in braces:
  \lstinline[style=l4style]|{ x: Tensor[N], it: Int, converged: Bool }|. A field
  is read with a dot, \code{s.x}. A record is \emph{updated} by making a fresh
  copy with one field changed: \lstinline[style=l4style]|{ ...s, it: s.it + 1 }|
  means ``\code{s} again, but with \code{it} increased by one.'' The original
  \code{s} is untouched.
  \item[Lists] are written \code{[a, b, c]}; a \emph{list comprehension}
  \code{[ f x | x <- xs ]} reads ``the list of \code{f x} for each \code{x} in
  \code{xs}.''
  \item[\code{foldl}, \code{foldr}, \code{map}] are the standard
  \emph{combinators}: \code{map f xs} applies \code{f} to every element;
  \code{foldl f z xs} threads an accumulator left to right through the list.
  These three, and one loop combinator we introduce in \cref{ch:fold}, express
  every iteration in the book.
  \item[\code{do} blocks] sequence a computation top to bottom, like imperative
  code:
\begin{lstlisting}[style=l4style]
do  Ax  <- apply A x        -- name an intermediate result
    a   <- dot Ax x
    modify (\s -> { ...s, it: s.it + 1 })   -- the one place state changes
    return a
\end{lstlisting}
  We use \code{do} blocks only where a computation threads evolving state; the
  rest is plain function application. The full meaning of \code{do} and
  \code{modify} is built carefully in \cref{ch:monad}---for now, read a
  \code{do} block as ``these steps, in order.''
\end{description}

\begin{notationbox}
The pseudo-language is set in a \emph{box}, never inside a mathematical formula.
This is a real convention, not a stylistic one: the notation reuses symbols
(such as the \texttt{\$} sign in shape expressions) that would clash with
mathematics typesetting. When you see a shaded box, you are reading code; when
you see an inline formula like $\mat{K}\vx=\vb$, you are reading mathematics.
\end{notationbox}

\section*{Shapes, in brief}

Every tensor carries a \emph{shape}---the sizes and meanings of its axes---written
in square brackets after \code{Tensor}. We will treat shapes carefully in
\cref{ch:tensors}; for now you need only three forms:

\begin{itemize}[nosep]
  \item \code{Tensor[N]} is a flat vector of length $N$---a single column of
  degrees of freedom, the workhorse shape of every linear solve.
  \item \code{Tensor[m, n]} is a two-axis array (think: a small dense matrix).
  \item \lstinline[style=l4style]|Tensor[(S: ...)]| means ``a tensor whose axes
  form a run we will call \code{S}, of unspecified rank''; a later
  \lstinline[style=l4style]|Tensor[$S]| asserts ``the same shape \code{S}
  again.'' This lets a single function describe an operation that works on
  vectors, matrices, or higher arrays at once.
\end{itemize}

\section*{Mathematical conventions}

Vector fields are set upright and bold: $\vE$ (electric field), $\vB$ (magnetic
flux density), $\vJ$ (current density), $\vA$ (vector potential). The
differential operators are $\Grad$ (gradient), $\Div$ (divergence), and
$\Curl$ (curl). Discrete operators---matrices that act on vectors of degrees of
freedom---are set in a sans-serif font: $\mat{K}$ (a stiffness operator),
$\mat{M}$ (a mass operator), $\mat{C}$ (a damping operator). Inner products are
$\inner{\vx}{\vy}$ and norms are $\norm{\vx}$. The real and complex numbers are
$\Real$ and $\Complex$; $\Re$ and $\Im$ take real and imaginary parts. The four
finite-element function spaces, introduced in \cref{ch:derham}, are $\Hone$,
$\Hcurl$, $\Hdiv$, and $\Ltwo$.

\Cref{tab:symbols} collects the symbols used most often.

\begin{table}[t]
\centering
\caption[Frequently used symbols]{Frequently used symbols and the page or chapter
where each is introduced.}
\label{tab:symbols}
\small
\begin{tabular}{@{}lll@{}}
\toprule
Symbol & Meaning & Introduced \\
\midrule
$\vE,\vB,\vA$ & electric field, magnetic flux density, vector potential & \cref{ch:maxwell} \\
$\varepsilon,\mu,\sigma$ & permittivity, permeability, conductivity & \cref{ch:maxwell} \\
$\mat{K},\mat{M},\mat{C}$ & stiffness, mass, damping operators & \cref{ch:weak} \\
$a(u,v)$ & a bilinear (weak) form & \cref{ch:weak} \\
$\Hone,\Hcurl,\Hdiv,\Ltwo$ & the de Rham function spaces & \cref{ch:derham} \\
$N$ & number of (true) degrees of freedom & \cref{ch:fem} \\
$\omega,\lambda$ & angular frequency, eigenvalue & \cref{ch:maxwell} \\
$Q$ & quality factor of a resonant mode & \cref{ch:eigenmodes} \\
$S_{ij}$ & scattering-matrix (S-parameter) entry & \cref{ch:driven} \\
\code{Tensor[...]} & a typed, immutable array & \cref{ch:tensors} \\
\code{iterate\_while} & the loop combinator & \cref{ch:fold} \\
\code{Solve} & the state monad threading \code{SimState} & \cref{ch:monad} \\
\bottomrule
\end{tabular}
\end{table}


\mainmatter

\part{Orientation}\label{part:orientation}
\chapter{The Simulator and Its Targets}
\label{ch:targets}

\chapterorient{Before we open the machine, we should know what it makes. This
chapter is a guided tour of the six questions our simulator can answer, the
physical situations each one models, and the single number or table each one
produces. By the end you will be able to look at an engineering problem and say
which of the six analyses it calls for---and you will have met, in physical
form, every target the rest of the book learns to compute.}

\section{What an electromagnetic field simulator does}

An electromagnetic field simulator takes a description of a physical
device---its geometry, its materials, and how it is excited---and computes the
electromagnetic fields inside and around it. From those fields it extracts the
quantities an engineer actually wants: the capacitance between two conductors,
the inductance of a coil, the resonant frequencies of a cavity, the fraction of
a signal a microwave component reflects.

The device we use throughout this book is \emph{Palace}, a three-dimensional
finite-element solver for Maxwell's equations. Its internals are the subject of
the entire book; in this first chapter we stay outside the machine and look only
at its interface. From the outside, every run has the same shape
(\cref{fig:blackbox}): a \emph{configuration} goes in---a mesh of the geometry,
the material properties, the boundary excitations, and a choice of analysis---and
a \emph{physical product} comes out.

\begin{figure}[t]
\centering
\begin{tikzpicture}[node distance=14mm]
  \node[data, align=left, text width=24mm] (cfg)
    {geometry\\materials\\excitation\\analysis type};
  \node[stage, minimum width=34mm, minimum height=16mm, right=of cfg] (sim)
    {\textbf{the simulator}\\[1pt]\footnotesize assemble $\to$ solve $\to$ reduce};
  \node[product, align=left, text width=24mm, right=of sim] (out)
    {capacitance,\\inductance,\\frequencies,\\S-parameters,\\\dots};
  \draw[flow] (cfg) -- node[above,font=\scriptsize]{config} (sim);
  \draw[flow] (sim) -- node[above,font=\scriptsize]{product} (out);
\end{tikzpicture}
\caption[The simulator as a black box]{From the outside, every analysis is the
same black box: a configuration enters on the left and a physical product leaves
on the right. The three-word caption inside the box---\emph{assemble, solve,
reduce}---is the skeleton we will open up in \cref{ch:situating} and spend the
rest of the book filling in.}
\label{fig:blackbox}
\end{figure}

What changes from run to run is the \emph{analysis type}---the choice of which
physical question to ask. Palace offers six. They look, at first, like six
different programs. A central claim of this book, made precise in
\cref{ch:situating} and proved component by component thereafter, is that they
are really one program with a few parts swapped. But to make that claim land, we
first need the six in their own right. We take them in turn, each with the
physical picture it models and the quantity it yields.

\section{The six analyses}

\Cref{fig:sixanalyses} shows all six at a glance; the subsections that follow
walk through them one by one. Keep one eye on the figure as you read---the point
of the section is that these six very different-looking physical situations will,
by the end of the book, turn out to be six settings of the same dial.

\begin{figure}[p]
\centering
\setlength{\tabcolsep}{6pt}
\renewcommand{\arraystretch}{1.2}
\begin{tabular}{@{}cc@{}}
% --- electrostatic: capacitor ---
\begin{tikzpicture}[scale=0.95]
  \fill[simBlue!12] (-1.1,-0.7) rectangle (1.1,0.7);
  \draw[very thick,simInk] (-1.1,0.7)--(1.1,0.7);
  \draw[very thick,simInk] (-1.1,-0.7)--(1.1,-0.7);
  \node[font=\scriptsize] at (-1.4,0.7) {$+V$};
  \node[font=\scriptsize] at (-1.4,-0.7) {$0$};
  \foreach \x in {-0.8,-0.4,0,0.4,0.8}{\draw[flow,simRust,shorten >=1pt] (\x,0.62)--(\x,-0.62);}
  \node[font=\scriptsize\itshape] at (0.7,0.0) {$\vE$};
\end{tikzpicture}
&
% --- magnetostatic: coil ---
\begin{tikzpicture}[scale=0.95]
  \draw[very thick,simRust] (0,0) ellipse (0.45 and 0.9);
  \draw[-{Stealth[length=2mm]},simRust,thick] (0.45,0.2)--(0.45,-0.1);
  \node[font=\scriptsize] at (0.85,-0.4) {$I$};
  \foreach \y in {-0.5,0,0.5}{\draw[flow,simBlue] (-1.1,\y) .. controls (-0.2,\y+0.15) and (0.2,\y+0.15) .. (1.1,\y);}
  \node[font=\scriptsize\itshape] at (1.0,0.75) {$\vB$};
\end{tikzpicture}
\\[-1mm]
{\footnotesize\textbf{Electrostatic} $\to$ capacitance}
&
{\footnotesize\textbf{Magnetostatic} $\to$ inductance}
\\[3mm]
% --- eigenmode: cavity ---
\begin{tikzpicture}[scale=0.95]
  \draw[very thick,simInk] (-1.1,-0.7) rectangle (1.1,0.7);
  \draw[simTeal,thick,domain=-1.05:1.05,samples=60] plot (\x,{0.45*sin(180*\x+90)});
  \draw[simTeal!55,thick,domain=-1.05:1.05,samples=60] plot (\x,{-0.45*sin(180*\x+90)});
  \node[font=\scriptsize] at (0.0,-1.05) {$f,\,Q$};
\end{tikzpicture}
&
% --- driven: 2-port network ---
\begin{tikzpicture}[scale=0.95]
  \node[draw,thick,simInk,fill=simBgGray,minimum width=12mm,minimum height=12mm] (b) at (0,0) {\scriptsize DUT};
  \draw[flow,simGreen] (-1.5,0.25)--(b.west|-0,0.25) coordinate[midway](in);
  \draw[-{Stealth[length=2mm]},simRust] (b.west|-0,-0.25)--(-1.5,-0.25);
  \draw[flow,simGreen] (b.east)--(1.5,0);
  \node[font=\scriptsize] at (-1.2,0.55) {in};
  \node[font=\scriptsize] at (-1.25,-0.55) {refl.};
  \node[font=\scriptsize] at (1.05,-0.35) {$S_{ij}$};
\end{tikzpicture}
\\[-1mm]
{\footnotesize\textbf{Eigenmode} $\to$ frequencies \& $Q$}
&
{\footnotesize\textbf{Driven} $\to$ S-parameters}
\\[3mm]
% --- transient: pulse over time ---
\begin{tikzpicture}[scale=0.95]
  \draw[->,simGray] (-1.2,-0.55)--(1.3,-0.55) node[right,font=\scriptsize]{$t$};
  \draw[simRust,thick,domain=-1.1:1.1,samples=60]
    plot (\x,{0.7*exp(-6*\x*\x)*cos(360*\x)-0.0});
  \node[font=\scriptsize\itshape] at (-1.0,0.55) {$\vJ(t)$};
\end{tikzpicture}
&
% --- waveguide cross-section ---
\begin{tikzpicture}[scale=0.95]
  \draw[very thick,simInk] (-0.9,-0.7) rectangle (0.9,0.7);
  \foreach \x in {-0.55,-0.18,0.18,0.55}{
    \draw[-{Stealth[length=1.6mm]},simViolet] (\x,-0.5)--(\x,0.5);}
  \draw[-{Stealth[length=2.4mm]},simBlue,thick] (1.05,-0.5)--(1.05,0.5)
    node[right,font=\scriptsize]{$k_n$};
\end{tikzpicture}
\\[-1mm]
{\footnotesize\textbf{Transient} $\to$ time-domain fields}
&
{\footnotesize\textbf{Boundary mode} $\to$ waveguide modes}
\end{tabular}
\caption[The six analyses]{The six physical situations the simulator models.
Each is a different question about the same Maxwell physics; each yields a
different product. The rest of Part~I shows that all six share one skeleton, and
\cref{part:modalities} studies each in depth.}
\label{fig:sixanalyses}
\end{figure}

\subsection{Electrostatic: capacitance}

Hold two conductors at fixed voltages with an insulator (a \emph{dielectric})
between them. Charge collects on their surfaces, and an electric field $\vE$
fills the gap. The ratio of stored charge to applied voltage is the
\emph{capacitance}, the quantity that governs how a circuit stores energy in its
electric field and how quickly signals can charge a node.

The electrostatic analysis computes the capacitance of a system of conductors.
It does so by solving for the electric potential $V$ when each conductor in turn
is held at unit voltage and the others at zero, then reading the capacitance off
the resulting fields. The governing equation is the simplest in the book---a
single scalar Poisson equation, $-\Div(\varepsilon\Grad V)=0$, the same
equation that describes steady heat flow or an ideal fluid's potential. The
product is an $n\times n$ \emph{capacitance matrix} relating the charges on $n$
conductors to their voltages.

\subsection{Magnetostatic: inductance}

Run a steady current around a loop of wire and a magnetic field $\vB$ threads
through it. The ratio of magnetic flux to current is the \emph{inductance}, the
magnetic counterpart of capacitance and the quantity that governs how a circuit
stores energy in its magnetic field and resists changes in current.

The magnetostatic analysis computes inductance the same way electrostatics
computes capacitance: drive each current source in turn, solve for the field,
and reduce. The equation is the \emph{curl--curl} equation
$\Curl(\nu\,\Curl\vA)=\vJ$ for a magnetic vector potential $\vA$, where
$\nu=1/\mu$ is the reluctivity. We will see in \cref{ch:magnetostatics} that
this analysis is, almost line for line, the electrostatic analysis with the
function space changed and one weight in the final reduction swapped---our first
concrete instance of ``the same machine, one part different.''

\subsection{Eigenmode: resonant frequencies and quality factors}

Strike a bell and it rings at particular pitches; close off a region of space
with conducting walls and the electromagnetic field inside will likewise
``ring'' at particular frequencies, its \emph{resonant modes}. Each mode has a
frequency $f$ and a \emph{quality factor} $Q$ measuring how slowly it loses
energy. Cavity resonators, superconducting qubits, and filter design all live
or die by these numbers.

The eigenmode analysis finds the resonances with no external drive at all. It
assembles the operators once and hands them to a specialized eigenvalue solver,
which returns the modes as eigenpairs: each eigenvalue gives a frequency, each
eigenvector a mode field. Mathematically this is a \emph{generalized eigenvalue
problem} $\mat{K}\vx=\lambda\mat{M}\vx$. It is the one analysis whose ``solve''
is a single call to an opaque library routine rather than a loop we write
ourselves---a distinction that will turn out to matter a great deal.

\subsection{Driven: scattering parameters}

Feed a sinusoidal signal into one port of a microwave component and some of it
passes through, some reflects, and the proportions depend on frequency. The
\emph{scattering parameters}, or S-parameters, are the complex numbers $S_{ij}$
that record how much of a wave entering port $j$ leaves port $i$. They are the
universal currency of microwave and RF engineering; every filter, coupler, and
antenna is specified by its S-parameters across a frequency band.

The driven analysis computes S-parameters by solving the time-harmonic Maxwell
system at each frequency in a swept band and projecting the resulting fields onto
the ports. The operator now depends on the frequency $\omega$:
$\mat{A}(\omega)=\mat{K}+i\omega\mat{C}-\omega^2\mat{M}$. Because the operator
changes with every frequency, this analysis rebuilds and re-solves inside the
sweep---the opposite of electrostatics, which assembles once and reuses. That
contrast is the heart of \cref{ch:driven}.

\subsection{Transient: time-domain fields}

Inject a pulse of current and watch the fields evolve, step by step, in time.
The transient analysis marches the time-domain Maxwell equations forward from an
initial state, producing the full history of the fields: a movie rather than a
snapshot. It is the analysis of choice when the excitation is broadband, when
nonlinear effects matter, or when one simply wants to watch a wave propagate.

The governing equation is now a second-order system in time,
$\mat{M}\,\ddot{\vs}+\mat{C}\,\dot{\vs}+\mat{K}\,\vs=\vJ(t)$, integrated by an
ordinary-differential-equation solver. Where the driven analysis solved
independent problems at independent frequencies, the transient analysis solves a
\emph{chain}: each time step begins from the state the previous step left
behind. This is our first genuinely sequential computation, and it will be the
motivating example for the ``fold'' in \cref{ch:fold}.

\subsection{Boundary mode: waveguide propagation modes}

A waveguide---a hollow metal pipe, a coaxial line, a microstrip---carries
electromagnetic energy along its length in a set of characteristic patterns
called \emph{propagation modes}, each with a \emph{propagation constant} $k_n$
that says how fast it travels and whether it travels at all. Before the driven
analysis can excite a device through a port, it must know the modes of the
waveguide feeding that port.

The boundary-mode analysis computes those modes. It is distinguished from the
other five by a preliminary step: it extracts a two-dimensional cross-section
from the three-dimensional mesh and solves an eigenvalue problem \emph{on the
cross-section}. The solve itself is, once again, the opaque eigenvalue call of
the eigenmode analysis---now posed on a surface rather than a volume. The product
is a table of modes, each with its propagation constant and its transverse field
pattern.

\section{Six questions, one machine}

\Cref{tab:sixanalyses} lays the six side by side. Read down the columns and a
pattern jumps out that no single analysis reveals: the products differ, the
equations differ, the function spaces differ---but the \emph{structure} of every
row is identical. Each analysis assembles one or more operators from the mesh and
materials, solves something, and reduces the solution to a physical quantity.

\begin{table}[t]
\centering
\caption[The six analyses, compared]{The six analyses share a single structure:
assemble, solve, reduce. The columns that change are the physics (equation,
space) and the \emph{kind} of solve. \Cref{ch:situating} names the four kinds of
solve and the three kinds of reduction that this table is quietly tabulating.}
\label{tab:sixanalyses}
\small
\setlength{\tabcolsep}{4.5pt}
\begin{tabular}{@{}llll@{}}
\toprule
Analysis & Physical question & Governing equation & Product \\
\midrule
Electrostatic & capacitance of conductors & $-\Div(\varepsilon\Grad V)=0$ & capacitance matrix \\
Magnetostatic & inductance of current loops & $\Curl(\nu\Curl\vA)=\vJ$ & inductance matrix \\
Eigenmode & resonances of a cavity & $\mat{K}\vx=\lambda\mat{M}\vx$ & frequencies $f$, factors $Q$ \\
Driven & frequency response & $\mat{A}(\omega)\vE=\vb(\omega)$ & S-parameters $S_{ij}(\omega)$ \\
Transient & time-domain response & $\mat{M}\ddot{\vs}+\mat{C}\dot{\vs}+\mat{K}\vs=\vJ(t)$ & field history \\
Boundary mode & waveguide modes & $\mat{A}(\omega)\vx=\sigma\mat{B}\vx$ (2-D) & modes $k_n$, fields \\
\bottomrule
\end{tabular}
\end{table}

\begin{keyidea}
The simulator offers six analyses, but they are six settings of one machine.
Every one of them \emph{assembles} operators from the geometry and materials,
\emph{solves} a problem built from those operators, and \emph{reduces} the
solution to a physical quantity. What distinguishes the analyses is which
equation they discretize, on which function space, and---above all---\emph{what
kind of solve} the middle stage performs. There are only four kinds, and naming
them is the business of \cref{ch:situating}.
\end{keyidea}

This is the thesis the book exists to develop. We have it now only as a
promise---a table whose last-but-one column we have not yet learned to read. The
next chapter opens the machine far enough to see the stages \emph{assemble,
solve, reduce} as actual steps in an actual program: the lifecycle that every
one of these six analyses runs through, from the moment the user presses
\textsf{enter} to the moment the product is written out.

\summarysection
A field simulator turns a configuration---geometry, materials, excitation,
analysis type---into a physical product. Palace offers six analyses:
electrostatic (capacitance), magnetostatic (inductance), eigenmode (resonant
frequencies and quality factors), driven (S-parameters), transient (time-domain
fields), and boundary mode (waveguide propagation modes). Despite their physical
variety, all six share one skeleton: \emph{assemble, solve, reduce}. The book's
central project is to show that the variety lives almost entirely in the
``solve'' stage, of which there are only four kinds.

\furtherreading
For the physics behind the six analyses, any undergraduate electromagnetics
text---Griffiths~\citep{griffiths2017} for the fields, Pozar~\citep{pozar2011}
for the microwave quantities (S-parameters, waveguide modes, $Q$)---is ideal
background. The simulator itself is documented at the Palace
project~\citep{palace2023}. No functional-programming or finite-element
background is needed yet; we build both from scratch beginning in
\cref{part:fieldtheory} and \cref{part:framework}.

\exercisessection
\begin{enumerate}[nosep]
  \item For each of the six analyses, name one engineering device whose design
  would require it. (Example: the eigenmode analysis is essential to designing a
  superconducting-qubit resonator.)
  \item The electrostatic and magnetostatic rows of \cref{tab:sixanalyses} look
  the most alike. List every column in which they differ. (We will return to
  exactly this comparison in \cref{ch:magnetostatics}.)
  \item Which two analyses produce a result that depends on frequency? Which one
  produces a result that depends on time? What does this suggest about how their
  ``solve'' stages must differ?
  \item The driven and boundary-mode analyses both involve frequency, but only
  one of them ``rings'' with no input. Which, and why does that make its solve a
  search for eigenvalues rather than a response to a drive?
\end{enumerate}

\chapter{A Simulation's Life Cycle}
\label{ch:lifecycle}

\chapterorient{In the last chapter we stood outside the machine and watched a
configuration go in and a physical product come out. Now we open the lid just
far enough to see the stages the run passes through---parse, configure,
dispatch, build the mesh, solve, refine, output---and what travels between them.
One stage will turn out to be the hinge on which all six analyses swing, and one
will be our first glimpse of the single most important structure in the book:
the \emph{fold}. By the end you will be able to trace any Palace run from the
press of \textsf{enter} to the file it writes, and to point to the exact line
where ``one program'' becomes ``six.''}

\section{From black box to pipeline}

\Cref{ch:targets} showed every analysis as one black box: configuration in,
product out. That picture is true but coarse. If we slice the box open along the
direction the data flows, we find not one step but a short \emph{pipeline} of
them, run in sequence, the same sequence every time, for every analysis. This
chapter is a tour of that pipeline---the \emph{lifecycle} of a simulation.

The lifecycle is worth a chapter of its own for two reasons. First, it is the
\emph{shared skeleton}: all six analyses pass through exactly these stages in
exactly this order, so understanding the lifecycle once means understanding the
top-level shape of all six at once. Second, two of its stages carry the book's
central ideas in seed form---the \emph{dispatch} that turns one program into six,
and the \emph{adaptive loop} that gives us our first fold. We meet both here, in
their natural habitat, before we study either in depth.

In real source, the lifecycle lives in two files. A short driver, \code{main.cpp},
sets the run up and chooses which analysis to perform; a shared base class,
\code{basesolver.cpp}, runs the chosen analysis under an adaptive outer loop.
We will not read C++ in this book---we read the lifecycle as a sequence of pure
transformations on values---but it helps to know that the two halves we are
about to describe are two real, separate pieces of the program.

\section{The six stages}

\Cref{fig:lifecycle} is the whole chapter in one picture. A run proceeds left to
right through six stages. Most of the arrows are ordinary forward flow; one is
dashed and points \emph{backward}---the adaptive loop, the only place the
pipeline ever revisits a stage. Keep the figure in view; the subsections below
walk the stages one at a time.

\begin{figure}[t]
\centering
\begin{tikzpicture}[node distance=7mm and 9mm]
  \node[stage] (parse) {parse\\config};
  \node[stage, right=of parse] (dev) {configure\\device};
  \node[stage, right=of dev] (disp) {dispatch on\\problem type};
  \node[stage, right=of disp] (mesh) {build\\mesh};
  % the solve / estimate / refine band
  \node[opbox, below=12mm of disp, xshift=10mm] (solve) {Solve};
  \node[opbox, right=of solve] (est) {estimate};
  \node[opbox, right=of est] (mark) {mark};
  \node[opbox, right=of mark] (ref) {refine};
  \node[product, below=12mm of mesh, xshift=22mm] (out) {output\\product};
  % top row forward flow
  \draw[flow] (parse) -- (dev);
  \draw[flow] (dev) -- (disp);
  \draw[flow] (disp) -- (mesh);
  % into the solve band
  \draw[flow] (mesh.south) |- (solve.west);
  % the adaptive band, left to right
  \draw[flow] (solve) -- node[above,font=\scriptsize]{soln.} (est);
  \draw[flow] (est) -- node[above,font=\scriptsize]{$\eta_K$} (mark);
  \draw[flow] (mark) -- (ref);
  % the back-edge: refine -> solve (the AMR loop)
  \draw[backflow] (ref.south) -- ++(0,-6mm) -| node[pos=0.25,below,font=\scriptsize]{refined mesh}
    (solve.south);
  % exit to output
  \draw[flow] (est.south) -- ++(0,-3mm) -| (out.north);
  % a faint band behind the loop
  \begin{scope}[on background layer]
    \node[band, fill=simBgGold, fit=(solve)(est)(mark)(ref),
      inner sep=6pt] (loopband) {};
  \end{scope}
  \node[cornertag, above=0.5pt of loopband.north] {the adaptive estimate--mark--refine loop};
\end{tikzpicture}
\caption[The simulation lifecycle]{The lifecycle of a single run. Four set-up
stages along the top---\emph{parse}, \emph{configure}, \emph{dispatch},
\emph{build mesh}---feed into the shaded \emph{solve--estimate--mark--refine}
loop, whose dashed back-edge is the only place the pipeline revisits a stage.
When adaptive refinement is switched off, the loop runs its body exactly once
and the back-edge is never taken; what remains is a straight line from config to
product. The hinge stage is \emph{dispatch}, expanded in \cref{fig:dispatch}.}
\label{fig:lifecycle}
\end{figure}

\subsection{Parse the configuration}

The run begins with a single file: a configuration written by the user. Parsing
it produces one in-memory value---call it the \code{config}---that holds
everything the run needs to know: the mesh file and polynomial order, the
material properties, the boundary excitations, the refinement settings, and,
crucially, the \emph{problem type} naming which of the six analyses to perform.
In Palace this value is an \code{IoData} record built in one line near the top of
\code{main.cpp}.

The single most important fact about the \code{config} is that it is
\emph{read-only}. It is built once, at the start, and then \emph{never modified}.
Every later stage consults it but none writes to it. We will lean on this fact
repeatedly: because the configuration cannot change mid-run, every stage that
depends only on the configuration computes the same thing no matter when we run
it, which is exactly the property that lets us reorder, hoist, and parallelize
the later stages with confidence.

\begin{notationbox}
We use \code{config} for the parsed configuration value throughout the book. It
is a \emph{record}---a bundle of named fields, like
\lstinline[style=l4style]|{problem_type, mesh, materials}|---and we read its
fields with a dot, as in \lstinline[style=l4style]|config.problem_type|. Records
get their full treatment in \cref{ch:bigidea}; for now, think of \code{config} as
a read-only struct.
\end{notationbox}

\subsection{Configure the device}

Before any arithmetic happens, the run brings up the hardware and numerical
libraries it will compute on: it selects a \emph{device} (a CPU or a GPU),
initializes the tensor backend that will run the element kernels, and starts the
linear-algebra and eigenvalue libraries. All of these choices are read from the
\code{config}; none of them touches the physics.

This stage is mostly plumbing, and we will say little more about it---but its
\emph{existence} is the seed of one of the book's themes. The reason Palace can
ask, in one line of configuration, to run the very same simulation on a laptop
CPU or a datacenter GPU is that the computation downstream is expressed
\emph{against an abstract device}, not against any particular one. The functional,
immutable-tensor presentation this book develops is precisely the shape that
makes a single description run on many devices. We will return to this point in
\cref{ch:bigidea}; here, simply note that ``configure the device'' is where the
backend is chosen, and that everything after it is device-agnostic.

\subsection{Dispatch on the problem type}

Now comes the hinge. The run reads one field of the \code{config}---the problem
type---and on the strength of that single value selects which of the six
analyses to perform. In source it is a \code{switch} on
\lstinline[style=l4style]|config.problem_type| with six branches; in our reading
it is a function that maps a problem type to a \emph{driver}, the object that
knows how to run one analysis end to end:
\begin{lstlisting}[style=l4style]
dispatch :: ProblemType -> Driver
dispatch ELECTROSTATIC = electrostaticDriver
dispatch MAGNETOSTATIC = magnetostaticDriver
dispatch EIGENMODE     = eigenmodeDriver
dispatch DRIVEN        = drivenDriver
dispatch TRANSIENT     = transientDriver
dispatch BOUNDARYMODE  = boundaryModeDriver
\end{lstlisting}

This is the place where ``one program'' becomes ``six.'' Everything before
dispatch is shared by all analyses; everything after is, in part, specialized to
the one we chose. We call this point the \emph{specialization seam}, and it is so
central that it gets its own figure (\cref{fig:dispatch}) and its own section
(\cref{sec:seam}) below.

\begin{figure}[t]
\centering
\begin{tikzpicture}[node distance=4mm]
  \node[stage, minimum width=24mm] (sw) {\code{switch}\\\footnotesize\itshape problem type};
  % six driver boxes in a fan
  \node[opbox, above right=10mm and 22mm of sw] (es) {electrostatic};
  \node[opbox, above right=3mm  and 22mm of sw] (ms) {magnetostatic};
  \node[opbox, right=22mm of sw, yshift=2mm]    (em) {eigenmode};
  \node[opbox, below right=3mm  and 22mm of sw] (dr) {driven};
  \node[opbox, below right=10mm and 22mm of sw] (tr) {transient};
  \node[opbox, below right=17mm and 22mm of sw] (bm) {boundary-mode};
  \draw[flow] (sw.east) -- (es.west);
  \draw[flow] (sw.east) -- (ms.west);
  \draw[flow] (sw.east) -- (em.west);
  \draw[flow] (sw.east) -- (dr.west);
  \draw[flow] (sw.east) -- (tr.west);
  \draw[flow] (sw.east) -- (bm.west);
  % one-word hints
  \node[cornertag, right=2mm of es] {capacitance};
  \node[cornertag, right=2mm of ms] {inductance};
  \node[cornertag, right=2mm of em] {resonance};
  \node[cornertag, right=2mm of dr] {S-parameters};
  \node[cornertag, right=2mm of tr] {time march};
  \node[cornertag, right=2mm of bm] {waveguide};
\end{tikzpicture}
\caption[The dispatch fan-out]{The specialization seam. A single
\code{switch} on the problem type fans out to six drivers, one per analysis,
each tagged with the product it computes. This is the only branch in the entire
lifecycle that depends on \emph{which} analysis we are running; every other stage
is shared. The six drivers are the subject of \cref{part:modalities}; that they
share one skeleton is the claim \cref{ch:situating} makes precise.}
\label{fig:dispatch}
\end{figure}

\subsection{Build the mesh}

With a driver chosen, the run builds the computational mesh: it loads the
geometry file named in the \code{config}, applies any driver-specific
preprocessing, partitions it, and performs any \emph{a-priori} refinement the
user requested. The output is a \code{mesh}---the discretized geometry on which
the fields will live.

Two things about this stage matter for the chapters ahead. First, the mesh is
the first value the lifecycle \emph{produces} (as opposed to merely reads): the
\code{config} came from the user, but the \code{mesh} is computed. Second, mesh
construction is \emph{driver-agnostic in shape}: every analysis builds a mesh the
same way, even though the function spaces later laid over it differ. We treat the
mesh as the second of the values that flow between stages, after the config.

\subsection{Solve}

Here the actual physics happens. The chosen driver takes the mesh and the
configuration, assembles the operators the analysis needs, solves the resulting
system, and reduces the solution to the analysis's physical product. This is the
\emph{assemble--solve--reduce} skeleton we met in \cref{ch:targets}, now occupying
a single stage of the lifecycle. Its internals are different for each driver---an
electrostatic solve is a handful of linear solves, an eigenmode solve is one
opaque eigenvalue call, a transient solve is a march through time---and unpacking
those differences is the work of \cref{part:modalities}.

For the lifecycle, the solve stage produces two outputs, and it is worth being
precise about both. The obvious one is the \emph{product}: the capacitance
matrix, the S-parameters, the field history---the thing the user ran the
simulation to obtain. The less obvious one is a set of \emph{error
indicators}: one nonnegative number $\eta_K$ per mesh element, estimating how
much of the solution's error is concentrated there. The product is what the user
wants; the error indicators are what the \emph{next} stage wants.

\subsection{Estimate, mark, refine---the adaptive loop}
\label{sec:amrloop}

The final structural stage is the one that makes \cref{fig:lifecycle} a loop
rather than a line. Having solved once and obtained error indicators, the run
asks a question: is the solution accurate enough? If the total error is below the
user's tolerance, it stops and outputs the product. If not, it improves the mesh
where the error is worst and solves again. Concretely:

\begin{enumerate}[nosep]
  \item \textbf{Estimate.} Turn the solution into the per-element error
  indicators $\eta_K$ (this is folded into the solve stage above; we name it
  separately because it is conceptually its own step).
  \item \textbf{Mark.} Select the subset of elements carrying the bulk of the
  error---a small set responsible for a large fraction of $\sum_K \eta_K^2$.
  \item \textbf{Refine.} Subdivide the marked elements, producing a finer mesh
  concentrated exactly where the solution needs the most resolution.
  \item \textbf{Repeat.} Solve again on the refined mesh, re-estimate, and loop,
  until the error falls below tolerance or a resource budget is exhausted.
\end{enumerate}

This is \emph{adaptive mesh refinement} (AMR), and it is studied in full in
\cref{ch:amr}. What concerns us here is its \emph{shape}, because that shape is
the most important structure in this book.

\subsection{Finalize and output}

When the loop terminates, the run writes the product to disk, records run
metadata (timings, memory, the configuration used), and tears down the device
and libraries. The product itself---the capacitance, the modes, the field
history---is written by the driver, not by the top-level \code{main}; \code{main}
writes only the metadata. This matters less for the physics than for the
bookkeeping: it tells us that the ``output'' of the lifecycle is exactly the
driver's product, and that the lifecycle ROOT contributes the \emph{scaffold}
around it, not the product itself.

\section{What flows between the stages}

We have named the stages; now read the figure the other way and ask what travels
the arrows. Four kinds of value flow through the lifecycle, and keeping them
straight is half of understanding it.

\begin{description}[style=nextline, leftmargin=0pt, font=\normalfont\itshape]
  \item[The configuration] enters at parse and is \emph{read-only} thereafter.
  Every stage may consult it; none may change it. It is the one value present
  from the first stage to the last.
  \item[The mesh] is produced by the build stage and consumed by every solve.
  It is the one value the \emph{loop} modifies: each refinement replaces the mesh
  with a finer one. Outside the loop it is built once and read.
  \item[The solution and its error indicators] are produced by each solve. The
  solution becomes (after reduction) the product; the error indicators drive the
  mark--refine decision. They are the values that \emph{close the loop}---the
  output of one iteration that becomes (via the refined mesh) the input of the
  next.
  \item[The product] is the final output: the physical quantity the user asked
  for. It is produced by the driver and written out at finalize.
\end{description}

\begin{keyidea}
The configuration is read-only and lives for the whole run; the mesh is built
once and then \emph{refined in the loop}; the solution and error indicators are
recomputed every iteration; the product is written at the end. The only value
that genuinely \emph{threads through} the adaptive loop---changing from one
iteration to the next---is the mesh (carrying its indicators). Everything else is
either fixed for the run or recomputed from scratch each pass. That single
threaded value is what makes the loop a \emph{fold}.
\end{keyidea}

\section{The specialization seam: one program, six analyses}
\label{sec:seam}

We promised in \cref{ch:targets} that the six analyses are one machine with parts
swapped. The lifecycle lets us say exactly \emph{where} the swap happens: at the
dispatch stage, and almost nowhere else.

Read \cref{fig:lifecycle} again with this question in mind: which stages depend on
\emph{which} analysis we are running? Parsing the config does not---it reads the
same file format for all six. Configuring the device does not. Building the mesh
does not, in shape. The adaptive loop does not---estimate, mark, and refine are
the same procedure regardless of what was solved. Even \emph{finalize} is
uniform. The \emph{only} stage that branches on the problem type is dispatch, and
the only stage whose internals differ between analyses is the solve it selects.

This is why we call dispatch the \emph{specialization seam}. Above it, the program
is one shared skeleton; below it, a driver fills in the one stage---solve---that
genuinely differs. The drivers are not six unrelated programs sharing a launcher;
they are six fillings of one mould. \Cref{ch:situating} makes this precise by
showing that even the solve stage varies in a tightly constrained
way---there are only \emph{four kinds} of solve among the six analyses---and that
the reductions that follow come in only three shapes.

\begin{pitfall}
It is tempting to read ``six drivers'' as ``six programs'' and conclude that the
shared lifecycle is just a convenience---a common \code{main} that happens to
launch any of them. That gets the emphasis backwards. The lifecycle is not
incidental scaffolding around six independent solvers; it is the \emph{primary
structure}, and the drivers are specializations of its one variable stage. The
whole book is organized around this inversion: we study the shared skeleton first
(in \cref{ch:situating} and the parts that follow) and the per-analysis
differences last (\cref{part:modalities}), because the shared part is larger and
more fundamental than the differences.
\end{pitfall}

\section{First sight of the fold}
\label{sec:fold}

Return to the adaptive loop of \cref{sec:amrloop}, and look past what it
\emph{does} to the \emph{shape} of how it does it. We have a starting value (the
initial mesh). We repeatedly apply a step that takes the current value and
produces the next one (solve, estimate, mark, refine $\to$ a finer mesh). We stop
when a condition on the accumulated state holds (the error is small enough, or
the budget is spent). And we report a result derived along the way (the product
of the final solve).

A starting value, a step that carries state forward, a stopping condition, a
result threaded through: this pattern is a \emph{fold}. It is the single most
reused structure in this book, and the adaptive loop is your first encounter with
it. In the framework's vocabulary the combinator is called \code{fold\_solve},
and the adaptive loop is written, schematically,
\begin{lstlisting}[style=l4style]
-- run the driver under the adaptive estimate-mark-refine loop
adaptive cfg drv mesh0 = fold_solve step mesh0
  where step m = let (product, etas) = drv cfg m   -- solve + estimate
                     m'               = refine (mark etas) m
                 in  (product, m')                 -- carry the refined mesh forward
\end{lstlisting}

The crucial property of a fold---the one that distinguishes it from a simple
loop over independent items---is that each step \emph{depends on the previous}.
The mesh fed to iteration $n+1$ is exactly the mesh that iteration $n$ produced;
the iterations do not commute and cannot be run in parallel or reordered. This is
in sharp contrast to, say, the electrostatic analysis's inner solves, which run
the \emph{same} operator against several independent right-hand sides and could
be done in any order. The distinction between an independent family of solves (a
\emph{map}) and a state-threaded chain of solves (a \emph{fold}) is one of the
four ``solve corners'' \cref{ch:situating} will name, and the fold is studied as a
combinator in its own right in \cref{ch:fold}.

\begin{keyidea}
The adaptive loop is a \emph{fold}: a step applied repeatedly, each pass
consuming the value the previous pass produced, terminating on a condition over
the accumulated state. The threaded value is the mesh. This same fold structure
will reappear as the engine of time-stepping, of Krylov solvers, and of
eigenvalue iterations---once you can see the fold here, you will see it
everywhere.
\end{keyidea}

\subsection{When the loop degenerates to a line}

There is one more observation, small but clarifying. Adaptive refinement is
\emph{optional}: a user who wants a single solve on a fixed mesh simply sets the
maximum refinement count to zero. When they do, the loop's stopping condition is
true before the first refinement, so the body runs exactly once: solve, estimate,
done. The back-edge in \cref{fig:lifecycle} is never taken, and the lifecycle
collapses to a straight line---config, device, dispatch, mesh, one solve, output.

This is worth dwelling on because it means the \emph{non-adaptive} lifecycle is
not a different program; it is the adaptive lifecycle with its fold run for a
single iteration. A fold of one step is just that step. So the electrostatic
analysis of \cref{ch:targets}, which most readers will picture as a simple
straight-through computation, is in truth the very same adaptive lifecycle---it
has just degenerated to a single pass. We lose nothing by always picturing the
loop; the line is its special case.

\begin{example}
A capacitance extraction with refinement disabled runs: parse the config;
configure the CPU; dispatch to the electrostatic driver; build the mesh; solve
once (yielding the capacitance matrix and a set of error indicators that nobody
will consult); find that the refinement budget is zero, so skip refinement; write
the capacitance matrix; finalize. Six stages, the loop run once. Re-enable
refinement and the only change is that the solve--estimate--mark--refine band
repeats until the indicators fall below tolerance---the same program, its fold
unrolled further.
\end{example}

\section{Where this leaves us}

We have opened the black box of \cref{ch:targets} into a pipeline of six stages
and identified its two load-bearing features: the \emph{dispatch} that
specializes one shared skeleton into six analyses, and the \emph{adaptive loop}
whose shape is a fold. We have catalogued the four kinds of value that flow
between stages, and we have seen that the familiar straight-through simulation is
just the adaptive lifecycle with its fold run once.

Two threads now lead forward. The \emph{specialization} thread asks: given that
the analyses differ only at the solve stage, exactly how do they differ?
\Cref{ch:situating} answers with the four solve corners and three reduction
shapes---the map that the rest of the book fills in. The \emph{representation}
thread asks: what is this ``value-threading,'' ``read-only configuration,''
``fold'' language we have started speaking, and why is it the right one for a
simulator? \Cref{ch:bigidea} answers that, taking up the shift from in-place
array mutation to pure functions over immutable tensors that the whole book is
built on.

\summarysection
Every Palace run follows the same lifecycle: parse the read-only configuration;
configure the compute device; \emph{dispatch} on the problem type to one of six
drivers; build the mesh; run the driver's \emph{solve} (assemble--solve--reduce),
which yields both a product and per-element error indicators; and run the
\emph{adaptive estimate--mark--refine loop}, refining the mesh where the error is
worst and re-solving until the error is small enough, before writing the product.
Four kinds of value flow through: the read-only config (whole run), the mesh
(built once, refined in the loop), the solution and its error indicators
(recomputed each pass), and the product (written at the end). Dispatch is the
\emph{specialization seam}---the one stage that branches on the analysis---so the
six analyses are one shared skeleton with a single swapped stage. The adaptive
loop is a \emph{fold}: a step that threads the mesh forward, our first sight of
the structure that organizes the entire book; when refinement is off, the fold
runs once and the lifecycle is a straight line.

\furtherreading
The lifecycle described here is implemented in the Palace solver's top-level
driver and its shared solver base class~\citep{palace2023}; readers comfortable
with C++ can profitably read \code{main.cpp} (the set-up and dispatch) alongside
this chapter. The adaptive-refinement loop rests on a-posteriori error estimation
and bulk marking, due to Zienkiewicz and Zhu~\citep{zienkiewiczzhu1987} and
D\"orfler~\citep{dorfler1996} respectively; we develop both in \cref{ch:amr}. The
fold, the structural heart of \cref{sec:fold}, is a staple of functional
programming; a reader curious to get ahead may consult any introduction to the
subject~\citep{peytonjones2003}, though we build the idea from nothing in
\cref{ch:fold}.

\exercisessection
\begin{enumerate}[nosep]
  \item List the six stages of the lifecycle in order. For each, state whether
  its behaviour depends on which of the six analyses is being run. (Exactly one
  should depend on it in its \emph{branching}, and one more in its
  \emph{internals}.)
  \item The configuration is read-only for the entire run. Name two concrete
  benefits this gives---one about reasoning, one about hardware. (Hint: re-read
  the parse and configure-device subsections.)
  \item A user sets the maximum refinement count to zero. Trace the lifecycle
  step by step and explain, in terms of the fold, why the result is a single
  straight-through solve. How many times does the step run?
  \item Of the four kinds of value that flow between stages, exactly one is
  threaded through the adaptive loop (changed from one iteration to the next).
  Which? Why does that single threaded value make the loop a fold rather than a
  map?
  \item The dispatch stage is the only one that branches on the problem type.
  Suppose someone proposed adding a seventh analysis to Palace. Using
  \cref{fig:lifecycle}, predict which stages they would have to touch and which
  they could leave untouched. (We test this prediction against reality in
  \cref{part:modalities}.)
\end{enumerate}

\chapter{The Big Idea: Tensor-Field Simulation}
\label{ch:bigidea}

\chapterorient{We have seen what the simulator computes (\cref{ch:targets}) and
the stages every run passes through (\cref{ch:lifecycle}). This chapter is about
the single idea that organizes the whole book---the lens through which we will
read every operator, every solver, every driver. It is a small idea, and a
familiar one once you see it: \emph{a simulation is a sequence of transformations
of immutable arrays of numbers.} Not a program that scribbles over a block of
memory, but a chain of pure functions, each taking a field and returning a new
one. By the end you will understand why we choose this lens, what it buys us, and
the three questions it teaches us to ask of any computation.}

\section{A field is a tensor}

Every analysis in \cref{ch:targets} solves for a \emph{field}: a quantity that
takes a value at every point of space. The electrostatic analysis solves for the
potential $V$, one number at each point. The magnetostatic analysis solves for
the vector potential $\vA$, three numbers at each point. The eigenmode analysis
solves for a mode field $\vE$, again a vector at each point.

A computer cannot store a value at \emph{every} point---there are infinitely
many. So the first thing any simulator does is \emph{discretize}: it replaces the
continuous field with a finite list of numbers, the \emph{degrees of freedom}
(DoFs), from which the field everywhere can be reconstructed. How that list is
chosen---which points, which basis functions---is the business of the finite
element method, and we build it carefully in \cref{part:fieldtheory}. For now,
only the shape of the answer matters: \emph{a discretized field is a finite,
ordered collection of numbers}. If a mesh has $N$ degrees of freedom, the
potential becomes a list of $N$ real numbers; the $\vE$-field becomes a list of
$N$ numbers too (one per edge degree of freedom).

Such an ordered, rectangular collection of numbers is exactly what we call a
\emph{tensor}. A list of $N$ numbers is a tensor of shape $[N]$. A grid of $H$
rows and $W$ columns is a tensor of shape $[H, W]$. A field that stores a
3-vector at each cell of an $H \times W$ grid is a tensor of shape $[H, W, 3]$.
The numbers $N$, or $H, W, 3$, are the \emph{axes} of the tensor; each axis has a
length, and we will increasingly give each axis a \emph{name} so we never lose
track of what it counts.

\begin{keyidea}
A discretized field \emph{is} a tensor---a finite, named-axis array of numbers.
The whole simulator, from this height, is a machine that consumes tensors
(the mesh, the materials, the excitation) and produces tensors (the solution
field) and finally small tensors (the capacitance matrix, the table of
frequencies). To understand the simulator is to understand the transformations
that carry one tensor to the next.
\end{keyidea}

\begin{figure}[t]
\centering
\begin{tikzpicture}[scale=1.0]
  % --- left: a field over a mesh, as Tensor[N] ---
  \begin{scope}[shift={(-3.4,0)}]
    % little triangulated patch
    \draw[simGray, semithick] (-0.9,-0.7) -- (0.9,-0.7) -- (0.6,0.7) -- (-0.7,0.6) -- cycle;
    \draw[simGray, semithick] (-0.9,-0.7) -- (0.6,0.7);
    \draw[simGray, semithick] (0.9,-0.7) -- (-0.7,0.6);
    \foreach \p in {(-0.9,-0.7),(0.9,-0.7),(0.6,0.7),(-0.7,0.6),(-0.05,-0.0)}{
      \fill[simRust] \p circle (1.4pt);}
    \node[font=\scriptsize\itshape, simGray] at (0,-1.15) {a field on a mesh};
  \end{scope}
  \node[font=\Large] at (-1.7,0) {$\rightsquigarrow$};
  % the Tensor[N] box
  \node[opbox, minimum width=30mm, minimum height=9mm] (tN) at (0.1,0)
    {\lstinline[style=l4style]|Tensor[N]|};
  \node[font=\scriptsize\itshape, simGray, below=1mm of tN]
    {$N$ degrees of freedom, one axis};
  % --- right: a multi-axis tensor Tensor[H, W, 3] ---
  \begin{scope}[shift={(4.6,0)}]
    \def\dx{0.16}\def\dy{0.10}
    % a small 3-deep stack of HxW grids to suggest axes
    \foreach \k/\c in {0/simBgGray,1/simBgBlue,2/simBgRust}{
      \begin{scope}[shift={(\k*\dx,\k*\dy)}]
        \fill[\c, draw=simGray] (-0.8,-0.55) rectangle (0.8,0.55);
        \draw[simGray!60] (-0.8,-0.18) -- (0.8,-0.18);
        \draw[simGray!60] (-0.8,0.18) -- (0.8,0.18);
        \draw[simGray!60] (-0.27,-0.55) -- (-0.27,0.55);
        \draw[simGray!60] (0.27,-0.55) -- (0.27,0.55);
      \end{scope}}
    \node[font=\scriptsize, simInk] at (-1.15,0.0) {$H$};
    \node[font=\scriptsize, simInk] at (0.0,-0.85) {$W$};
    \node[font=\scriptsize, simRust] at (1.15,0.75) {$3$};
    \node[font=\scriptsize\ttfamily, simInk, below=4mm of {(0,-0.55)}] (lab)
      at (0.1,-1.0) {Tensor[H, W, 3]};
  \end{scope}
\end{tikzpicture}
\caption[A tensor as a labeled box of axes]{Two tensors, drawn as labeled boxes
of axes. \emph{Left:} a scalar field on a mesh of $N$ degrees of freedom is a
one-axis tensor, written \lstinline[style=l4style]|Tensor[N]|. \emph{Right:} a
field storing a 3-vector at every cell of an $H\times W$ grid is a three-axis
tensor \lstinline[style=l4style]|Tensor[H, W, 3]|, with named axes $H$, $W$, and
the length-3 component axis. Reading and reasoning about these shapes is a skill
we develop in full in \cref{ch:tensors}; here we need only the picture.}
\label{fig:tensorbox}
\end{figure}

\Cref{fig:tensorbox} shows the two pictures we will lean on. On the left, a
scalar field over a mesh, collapsed into a single-axis tensor
\lstinline[style=l4style]|Tensor[N]|; on the right, a small multi-axis tensor
\lstinline[style=l4style]|Tensor[H, W, 3]| whose three named axes remind us that
a tensor is not just ``a big array'' but an array \emph{with structure we can
name}. We will treat shapes seriously---they are a kind of type, and they catch
mistakes the way ordinary types do---beginning in \cref{ch:tensors}. In this
chapter we keep them at arm's length and let the picture do the work.

\section{Two ways to evolve a state}

Now suppose we have our field as a tensor and we want to advance it: collide and
stream a fluid, take a time step, apply one iteration of a solver. There are two
very different mental models for ``advancing,'' and the choice between them is the
true subject of this chapter.

\subsection{The picture you probably have: in-place mutation}

If you have written numerical code in C or Fortran---or read any classical
simulation textbook---you carry a particular picture of how a field evolves. You
allocate an array. You loop over its entries. You \emph{overwrite} them in place.
The canonical inner loop looks like this update of a vector $\vy$ by a scaled
vector $\vx$ (the operation BLAS calls \texttt{axpy}, for ``$a$ times $x$ plus
$y$''):

\begin{lstlisting}[style=l4style, language=]
for (int i = 0; i < n; i++)
    y[i] = a * x[i] + y[i];     // overwrite y[i] in place
\end{lstlisting}

\noindent After the loop, the array named \texttt{y} holds different numbers than
before. The \emph{name} \texttt{y} survived; the \emph{values} it points to were
destroyed and replaced. There is exactly one block of memory, and it is mutated.
A whole simulation, in this picture, is a long sequence of such overwrites: the
same arrays, scribbled over again and again, step after step, until the final
state is whatever happens to be left in memory when the loop ends.

This picture is not wrong---it is, more or less, what the silicon does, and it is
why classical simulators are fast. But it has a quiet cost: the meaning of the
program is tangled up with \emph{when} each overwrite happens. To know what
\texttt{y} holds, you must know exactly which loop iterations have run and in what
order. The data and its history are the same thing.

\subsection{The picture we will use: pure functions over immutable values}

Here is the other model. A tensor is a \emph{value}, like the number $7$. You
never change a value; you only compute new ones from it. The expression $3 + 4$
does not ``overwrite'' the $3$---it produces a fresh $7$ and leaves $3$ untouched.
We treat tensors the same way. The \texttt{axpy} update becomes a \emph{pure
function}: it takes $a$, $\vx$, and $\vy$, and \emph{returns a brand-new tensor},
leaving its inputs exactly as they were.

\begin{lstlisting}[style=l4style]
axpy :: Scalar -> Tensor[N] -> Tensor[N] -> Tensor[N]
axpy a x y = a * x + y      -- returns a FRESH tensor; x and y unchanged
\end{lstlisting}

\noindent The arrow signature reads ``\texttt{axpy} takes a scalar and two
length-$N$ tensors and gives back a length-$N$ tensor.'' Nothing is overwritten.
If we want to advance a state $\vs$ repeatedly, we do not loop over memory; we
\emph{thread a sequence of distinct values} through a step function:
\[
  \vs_0 \;\xrightarrow{\;\mathrm{step}\;}\; \vs_1
        \;\xrightarrow{\;\mathrm{step}\;}\; \vs_2
        \;\xrightarrow{\;\mathrm{step}\;}\; \cdots
\]
Each $\vs_{k+1} = \mathrm{step}(\vs_k)$ is a fresh tensor produced from the one
before it. The old states still ``exist'' as values---we simply stop referring to
them, and the machine is free to reclaim their memory. \Cref{fig:mutvalue}
draws the two models side by side.

\begin{figure}[t]
\centering
\begin{tikzpicture}[node distance=9mm]
  % ============ LEFT: in-place mutation ============
  \begin{scope}
    \node[font=\bfseries\small, simRust] at (0,2.55) {In-place mutation};
    % a single memory cell, mutated across steps
    \node[draw=simRust, fill=simBgRust, thick, rounded corners=2pt,
          minimum width=20mm, minimum height=9mm] (cell) at (0,1.3)
      {\ttfamily y};
    \node[font=\scriptsize\itshape, simGray, right=1mm of cell]
      {one block of memory};
    % self-loop "overwrite" arrows
    \draw[backflow] (cell.north west) to[out=140,in=40,looseness=5]
      node[above,font=\scriptsize]{\texttt{y := a*x + y}} (cell.north east);
    \node[font=\scriptsize, simRust] at (0,-0.05) {step 1};
    \node[font=\scriptsize, simRust] at (0,-0.55) {step 2};
    \node[font=\scriptsize, simRust] at (0,-1.05) {step 3};
    \draw[backflow] (-1.0,-0.05) -- (cell.south west);
    \node[align=center, font=\scriptsize\itshape, simGray, text width=34mm]
      at (0,-1.85) {the same cell, scribbled over again and again};
  \end{scope}
  % ============ RIGHT: value threading ============
  \begin{scope}[shift={(6.6,0)}]
    \node[font=\bfseries\small, simBlue] at (0,2.55) {Value threading};
    \node[draw=simBlue, fill=simBgBlue, thick, rounded corners=2pt,
          minimum width=13mm, minimum height=8mm] (s0) at (0,1.4) {$\vs_0$};
    \node[draw=simBlue, fill=simBgBlue, thick, rounded corners=2pt,
          minimum width=13mm, minimum height=8mm] (s1) at (0,0.0) {$\vs_1$};
    \node[draw=simBlue, fill=simBgBlue, thick, rounded corners=2pt,
          minimum width=13mm, minimum height=8mm] (s2) at (0,-1.4) {$\vs_2$};
    \draw[flow] (s0) -- node[right,font=\scriptsize]{\texttt{step}} (s1);
    \draw[flow] (s1) -- node[right,font=\scriptsize]{\texttt{step}} (s2);
    \node[align=center, font=\scriptsize\itshape, simGray, text width=36mm]
      at (0,-2.25) {distinct immutable tensors, each produced from the last};
  \end{scope}
\end{tikzpicture}
\caption[Array mutation versus value threading]{The same computation under two
memory models. \emph{Left:} one block of memory named \texttt{y} is overwritten
in place, step after step; to know what it holds you must know how many steps have
run. \emph{Right:} each step is a pure function producing a \emph{fresh}
immutable tensor $\vs_{k+1}$ from $\vs_k$; the states are distinct values and the
chain of arrows \emph{is} the computation's history.}
\label{fig:mutvalue}
\end{figure}

\begin{pitfall}
The value-threading picture does \emph{not} mean we keep every intermediate
tensor in memory forever, nor that we copy a giant array on every step. It is a
statement about \emph{meaning}, not about machine resources. An implementation is
free to reuse the storage of $\vs_k$ for $\vs_{k+1}$ once nothing else refers to
$\vs_k$---exactly the in-place overwrite, recovered as an optimization. We are
choosing how to \emph{describe} the computation, not forbidding the compiler from
being clever. \Cref{ch:monad} makes this reuse precise.
\end{pitfall}

\subsection{Same algorithm, two memory models}

It is worth seeing, plainly, that these two pictures compute \emph{the same
thing}. Whether we overwrite \texttt{y} in place or thread a chain of fresh
tensors $\vs_0 \to \vs_1 \to \vs_2$, the arithmetic performed is identical and so
is the final answer. \Cref{fig:twomodels} makes this concrete: one dependency
graph, the bare sequence of arithmetic operations, annotated twice---once as
overwrites of a single cell, once as a chain of values.

\begin{figure}[t]
\centering
\begin{tikzpicture}[node distance=11mm]
  % the shared op chain
  \node[opbox, minimum width=14mm] (a) {step};
  \node[opbox, minimum width=14mm, right=14mm of a] (b) {step};
  \node[opbox, minimum width=14mm, right=14mm of b] (c) {step};
  \node[data, left=10mm of a] (in) {field};
  \node[data, right=10mm of c] (out) {result};
  \draw[flow] (in) -- (a);
  \draw[flow] (a) -- (b);
  \draw[flow] (b) -- (c);
  \draw[flow] (c) -- (out);
  % mutation annotation (above, rust)
  \node[font=\scriptsize, simRust] at ($(in)+(0,0.62)$) {\texttt{y}};
  \node[font=\scriptsize, simRust] at ($(a)!0.5!(b)+(0,0.62)$) {\texttt{y}};
  \node[font=\scriptsize, simRust] at ($(b)!0.5!(c)+(0,0.62)$) {\texttt{y}};
  \node[font=\scriptsize, simRust] at ($(c)!0.5!(out)+(0,0.62)$) {\texttt{y}};
  \node[font=\scriptsize\itshape, simRust, anchor=west] at (-3.0,0.62)
    {as mutation:};
  % value annotation (below, blue)
  \node[font=\scriptsize, simBlue] at ($(in)+(0,-0.62)$) {$\vs_0$};
  \node[font=\scriptsize, simBlue] at ($(a)!0.5!(b)+(0,-0.62)$) {$\vs_1$};
  \node[font=\scriptsize, simBlue] at ($(b)!0.5!(c)+(0,-0.62)$) {$\vs_2$};
  \node[font=\scriptsize, simBlue] at ($(c)!0.5!(out)+(0,-0.62)$) {$\vs_3$};
  \node[font=\scriptsize\itshape, simBlue, anchor=west] at (-3.0,-0.62)
    {as values:};
\end{tikzpicture}
\caption[One algorithm, two annotations]{The same dependency graph, annotated
twice. The boxes are the actual arithmetic; the chain of arrows is fixed. Read it
along the top (rust) and every arrow carries the \emph{same name} \texttt{y}---a
single cell, overwritten three times. Read it along the bottom (blue) and every
arrow carries a \emph{distinct value} $\vs_k$---a chain of immutable tensors. The
two readings are the same computation; they differ only in what the names mean.}
\label{fig:twomodels}
\end{figure}

This is the crucial observation, and it is liberating: \emph{choosing the
value-threading reading costs us nothing in the final answer.} It is a change of
\emph{description}, not of algorithm. So if it costs nothing, why prefer it? The
next section is the answer.

\section{Why the pure picture pays}

Three payoffs, in increasing order of how much they shape the rest of the book.

\paragraph{The data flow becomes the dependency graph.}
In the mutation picture, the arrows in \cref{fig:twomodels} are implicit---you
recover them by tracing which loop reads which array after which write. In the
value picture, the arrows are \emph{explicit and literal}: $\vs_1$ is the input to
the step that makes $\vs_2$, full stop. A pure function's output depends on
exactly its inputs and nothing else---no hidden global it quietly read, no array
some other loop overwrote behind its back. So the program text \emph{is} the
mathematical dependency graph. This is not a metaphor; it is what lets us, later,
read off an algorithm's structure---what it iterates, what it folds, what it can
do in parallel---directly from its written form (\cref{ch:fold}).

\paragraph{Pure pieces compose without surprises.}
Because a pure step depends only on its inputs, two steps that do not share an
input cannot interfere, and a step behaves the same no matter what ran before it.
This makes the pieces \emph{composable}: we can snap them together, reason about
each in isolation, substitute one implementation for another, and trust that the
whole behaves as the sum of its parts. Half of this book is the discovery that
six very different-looking simulators are built from one small set of such pieces,
combined in slightly different ways. That discovery is only possible because the
pieces compose cleanly---which is only true because they are pure.

\paragraph{The pure picture is the one the accelerator wants.}
Our eventual target is not a single CPU walking an array one entry at a time; it
is a GPU (or another massively parallel backend) applying an operation to
\emph{all} entries of a tensor at once. Such hardware is happiest when told ``here
is an immutable input tensor; produce an output tensor''---a pure, whole-tensor
transformation with no aliasing and no ordering hazards to police. The
value-threading description is, almost word for word, the form these backends
consume. We call this the \emph{implementation view}: the same spec, read as a
library of pure tensor functions ready to lower onto an accelerator. It is the
destination of the whole book, assembled explicitly in \cref{part:collection}.

\begin{keyidea}
Mutation and value-threading compute the same answer, so the choice is free---and
we spend that freedom on the pure picture because it makes the \emph{data flow
explicit} (the program text is the dependency graph), makes the \emph{pieces
compose} (pure functions don't interfere), and matches the \emph{accelerator
backend} (whole-tensor, no aliasing). Clarity, composability, and the GPU: three
reasons, one decision.
\end{keyidea}

\section{A physical analogy: stream, then collide}

A small physical example makes the ``make a new state from the old one'' rhythm
concrete. The lattice-Boltzmann method simulates a fluid by tracking, at each cell
of a grid, how much fluid is moving in each of a few discrete directions---a
little tensor of \emph{populations} per cell. One time step has two sub-steps. In
\emph{streaming}, every population slides to the neighboring cell in its
direction. In \emph{collision}, the populations within each cell relax toward a
local equilibrium. Stream, then collide; that is one step.

In a classical code this is a nest of loops overwriting a population array in
place. In our picture it is a pure function from the old state to a fresh one:

\begin{lstlisting}[style=l4style]
step :: FluidState -> FluidState
step s =
  let streamed = stream  s.populations in   -- slide to neighbors
  let collided = collide streamed       in   -- relax toward equilibrium
  { ...s, populations: collided, t: s.t + 1 }
\end{lstlisting}

\noindent The step reads the current state \texttt{s}, computes two intermediate
tensors (\texttt{streamed}, then \texttt{collided}), and returns a \emph{new}
record with the advanced populations and an incremented step count. Nothing is
overwritten. A whole simulation is just this step, applied again and again to its
own output---exactly the $\vs_0 \to \vs_1 \to \vs_2$ chain of
\cref{fig:mutvalue}. Our solvers will have the same shape: a step that turns one
field-state into the next, threaded through an iteration. We have borrowed the
fluid's rhythm; we will reuse it for time integration in \cref{ch:transient} and,
combinator by combinator, throughout \cref{part:framework}.

\begin{physicsbox}
``Stream then collide'' is the lattice-Boltzmann heartbeat, but the
\emph{shape}---read the current state, compute a fresh next state, advance a
counter---is universal. A time-domain Maxwell step does it. One iteration of a
linear solver does it. One round of mesh refinement does it. Every sequential
computation in this book is a step function threaded through an iteration; the
physics only changes what happens \emph{inside} the step.
\end{physicsbox}

\section{Three questions the calculus asks}

We have made one decision---describe simulations as pure transformations of
immutable tensors. To turn that decision into a working language, we will, in
\cref{part:framework}, build a small \emph{calculus}: a tiny, precise notation for
exactly these transformations. We close this chapter with a first look at the
three questions that calculus is organized to answer. You have already met all
three informally in this chapter; naming them now plants the signposts for
Part~III.

\begin{enumerate}[label=\textbf{Q\arabic*.}, leftmargin=*]
  \item \textbf{What operations happen?} The verbs---the primitive
  transformations and the named combinations of them. \texttt{axpy} is a verb;
  \texttt{stream} and \texttt{collide} are verbs; ``apply this operator,''
  ``take this inner product,'' ``assemble this matrix'' are verbs. Pinning down
  the vocabulary of operations, each with its shape signature like
  \lstinline[style=l4style]|axpy :: Scalar -> Tensor[N] -> Tensor[N] -> Tensor[N]|,
  is the work of \cref{ch:tensors} and \cref{ch:operators}.

  \item \textbf{Who owns the state?} Not every value plays the same role. The
  evolving field-state is threaded from step to step; a solver's matrix entries
  are fixed once and only read; a step's scratch intermediates live and die
  inside one step. Keeping these straight---which values evolve, which are
  read-only parameters, which are ephemeral---is what lets us thread state purely
  without losing track of it. \Cref{ch:monad} gives this its own discipline,
  the \emph{state monad}, and shows that the entire ``mutation'' of a solve is one
  carefully placed write.

  \item \textbf{How is state evolution coordinated?} A list of verbs is not yet an
  algorithm; something must say \emph{do this, then this, repeat until converged,
  collect the results}. That coordination---sequencing, iteration, convergence
  tests, mapping over a family---is itself made of a few reusable shapes we call
  \emph{combinators}. The most important is the \emph{fold}: apply a step to its
  own output, over and over, threading a state through. \Cref{ch:fold} is built
  around it.
\end{enumerate}

\begin{notationbox}
These three questions---\emph{what operations}, \emph{who owns state}, \emph{how
coordinated}---are exactly the three axes the formal calculus separates, and they
are the spine of \cref{part:framework}: one question per cluster of chapters
(verbs and shapes; the state monad; the fold and its relatives). Whenever a later
chapter feels intricate, return here and ask which of the three it is answering.
\end{notationbox}

We now have the lens. \Cref{ch:situating} uses it immediately: it lays the six
analyses of \cref{ch:targets} onto this framework, showing that the ``solve''
stage of every one is one of just four coordination shapes---four ways of
threading the step. With that, Part~I is complete, and we can begin building the
field theory and the calculus the rest of the book rests on.

\summarysection
A discretized field is a \emph{tensor}: a finite, named-axis array of numbers, the
degrees of freedom collected into a rectangular shape like
\lstinline[style=l4style]|Tensor[N]| or \lstinline[style=l4style]|Tensor[H, W, 3]|.
A simulation is a sequence of transformations of such tensors. There are two ways
to describe ``transformation.'' The classical one---\emph{in-place mutation}---
overwrites a single block of memory step after step, tangling a value with the
history of writes that produced it. The one we adopt---\emph{pure functions over
immutable values}---threads a chain of distinct tensors $\vs_0 \to \vs_1 \to
\cdots$, each freshly produced from the last. The two compute the same answer, so
the choice is free; we spend it on the pure picture because it makes the data flow
into an explicit dependency graph, makes the pieces compose without interference,
and matches the whole-tensor form a GPU/accelerator backend wants (the
implementation view). The lattice-Boltzmann ``stream then collide'' step is the
prototype of the step functions our solvers thread. Finally, the calculus we build
in Part~III is organized around three questions: \emph{what operations happen},
\emph{who owns the state}, and \emph{how state evolution is coordinated}.

\furtherreading
The shift from mutable arrays to pure functions over immutable values is the
central idea of functional programming; the classic, readable entry points are
Wadler's exposition of how monads tame state and effects in a pure
language~\citep{wadler1995monads}, and the standard reference for the language
that crystallized these ideas~\citep{peytonjones2003}. No functional-programming
background is assumed---\cref{ch:monad} builds the state monad we will need from
scratch. The tensor-and-shape model previewed here is developed fully in
\cref{ch:tensors}; the fold, our central combinator, in \cref{ch:fold}. Readers
who want the physical analogy in depth will find the lattice-Boltzmann method in
any computational fluid dynamics text; we use only its rhythm, not its details.

\exercisessection
\begin{enumerate}[nosep]
  \item A scalar potential on a mesh with $N=10^6$ degrees of freedom, and an
  $\vE$-field on the same mesh, are both stored as
  \lstinline[style=l4style]|Tensor[N]|. What does the single axis $N$ count in
  each case, and why is it the \emph{same} $N$? (Hint: revisit which degrees of
  freedom each field lives on; the full story is \cref{ch:functionspaces}.)
  \item Write the \texttt{axpy} update $y \leftarrow a x + y$ once as a C-style
  in-place loop and once as a pure function returning a fresh tensor. For each,
  state precisely what is true of the input \texttt{x} and the input \texttt{y}
  \emph{after} the operation completes.
  \item In \cref{fig:twomodels} the two readings differ only in what the names on
  the arrows mean. Explain why a compiler is nonetheless free to \emph{implement}
  the value-threading reading as in-place mutation. What condition on $\vs_k$ must
  hold for it to reuse that tensor's storage for $\vs_{k+1}$?
  \item Three of this chapter's payoffs were: explicit dependency graph,
  composability, accelerator-friendliness. For each, give a one-sentence reason it
  is a \emph{consequence of purity}---that is, of a step depending on nothing but
  its declared inputs.
  \item The lattice-Boltzmann \texttt{step} returns
  \lstinline[style=l4style]|{ ...s, populations: collided, t: s.t + 1 }|, a fresh
  record. Rewrite ``stream then collide'' as a description that mutates a single
  population array in place, and identify the one place where the pure version's
  meaning would become ambiguous if two such steps were run on overlapping memory.
\end{enumerate}

\chapter{Situating the Drivers on the Framework}
\label{ch:situating}

\chapterorient{We have met the six analyses (\cref{ch:targets}) and watched a
single run move from configuration to product through one lifecycle
(\cref{ch:lifecycle}). This chapter is where those threads tie together. We make
the shared \emph{skeleton} explicit, then ask the question the whole book turns
on: of the six analyses, how many genuinely different things happen in the
middle ``solve'' stage? The answer---\emph{four}---is the most important
taxonomy in the book. We name the four \emph{solve corners} and the three
\emph{reduction shapes}, and we lay out the grid that the rest of the book
fills in, one cell at a time. By the end you will be able to place any of the
six drivers on the framework and say exactly which stage it swaps.}

\section{The skeleton, made explicit}

Every analysis Palace runs, we claimed in \cref{ch:targets}, is the same
machine with a few parts swapped. \Cref{ch:lifecycle} showed the machine as an
actual program: parse the configuration, build a mesh, dispatch on the analysis
type, run a solve, and write out the product---with an optional adaptive
refinement loop wrapped around the solve. Strip that lifecycle down to its
load-bearing core and four words remain:
\[
  \texttt{config} \;\longrightarrow\; \texttt{assemble} \;\longrightarrow\;
  \texttt{solve} \;\longrightarrow\; \texttt{reduce} \;\longrightarrow\;
  \texttt{product}.
\]
This is \emph{the skeleton}. It is worth saying slowly, because everything in
the book hangs on it (\cref{fig:skeleton}).

\begin{itemize}[nosep]
  \item \textbf{config.} The user's description of the problem---geometry,
  materials, excitations, and the choice of analysis---parsed into a single
  read-only record. Nothing downstream ever mutates it.
  \item \textbf{assemble.} Turn the continuous physics on the mesh into one or
  more discrete \emph{operators}: large (conceptually) matrices like the
  stiffness $\Kop$, the mass $\Mop$, and the damping $\Cop$. Assembly is a
  \emph{fold} over the weak-form terms (\cref{ch:fem}); for now, treat it as the
  box that produces the operators.
  \item \textbf{solve.} Use those operators to compute the field, or fields, the
  analysis needs. This is the stage with all the variety, and the subject of
  the next section.
  \item \textbf{reduce.} Collapse the computed field(s) into the small physical
  quantity the user actually asked for---a capacitance matrix, a list of
  resonant frequencies, a table of S-parameters.
  \item \textbf{product.} The reduced quantity, written out. The reason the run
  existed.
\end{itemize}

\begin{figure}[t]
\centering
\begin{tikzpicture}[node distance=9mm and 11mm]
  \node[data] (cfg) {config};
  \node[stage, right=of cfg] (asm) {assemble};
  \node[stage, right=of asm] (solve) {\textbf{solve}};
  \node[stage, right=of solve] (red) {reduce};
  \node[product, right=of red] (prod) {product};
  \draw[flow] (cfg) -- (asm);
  \draw[flow] (asm) -- node[above,font=\scriptsize]{$\Kop,\Mop,\dots$} (solve);
  \draw[flow] (solve) -- node[above,font=\scriptsize]{field(s)} (red);
  \draw[flow] (red) -- (prod);
  % AMR back-edge
  \draw[backflow] (red.south) ++(0,-1mm)
    .. controls +(0,-7mm) and +(0,-7mm) .. (asm.south)
    node[midway, below, font=\scriptsize, text=simRust]
      {refine mesh (AMR outer fold)};
\end{tikzpicture}
\caption[The skeleton]{The skeleton every analysis shares:
\texttt{config} $\to$ \texttt{assemble} $\to$ \texttt{solve} $\to$
\texttt{reduce} $\to$ \texttt{product}.
The dashed back-edge is the adaptive-mesh-refinement outer loop
(\cref{ch:amr}): when refinement is on, the reduced error estimate feeds a mesh
update and the whole assemble--solve--reduce body runs again; when it is off,
the body runs exactly once. The bold \texttt{solve} box is where the six
analyses diverge---and they diverge into only four shapes.}
\label{fig:skeleton}
\end{figure}

One feature of the lifecycle does not fit inside a single left-to-right pass:
adaptive mesh refinement. When refinement is enabled, the run does not stop
after one reduce. Instead the error it estimates feeds a mesh update, and the
entire assemble--solve--reduce body runs again on the finer mesh---and again,
until the error is small enough or a step budget is spent. That is the dashed
back-edge in \cref{fig:skeleton}. It is an \emph{outer fold} wrapped around the
skeleton: each pass begins from the mesh the previous pass refined. When
refinement is off, the fold degenerates to a single pass, and the skeleton is
exactly the straight line. We will meet this outer fold again as a genuine
computation in \cref{ch:amr}; for now, note only that it leaves the
five-stage skeleton intact and merely runs it more than once.

\begin{keyidea}
The skeleton is \texttt{config} $\to$ \texttt{assemble} $\to$ \texttt{solve}
$\to$ \texttt{reduce} $\to$ \texttt{product},
optionally wrapped in an adaptive-refinement outer fold. The configuration is
read-only; assembly produces operators; the solve produces fields; the
reduction produces the physical product. Of the five stages, four are nearly
identical across all six analyses. \emph{Almost the entire difference between
the six analyses lives in the \texttt{solve} stage}---and, to a lesser degree,
in the \texttt{reduce} stage. This chapter classifies both.
\end{keyidea}

\section{The four solve corners}

Here is the central observation of the book. Look at what the ``solve'' stage
actually does in each of the six analyses, and ignore the physics: ignore which
equation, which function space, which materials. Ask only about the
\emph{shape} of the computation---how many solves happen, whether they depend on
one another, whether the operator stays fixed or changes, and whether we write
the loop or hand the whole problem to a library. When you strip the physics
away like this, the six analyses collapse onto just \emph{four} distinct shapes.
We call them the four \emph{solve corners}.

Two questions separate them (\cref{fig:corners}). First: \emph{does the operator
stay fixed, or does it change as we go?} An analysis might assemble its operator
once and reuse it for every solve, or it might have to rebuild the operator
before each solve. Second: \emph{is the computation a map or a fold?} A
\emph{map} runs a collection of \emph{independent} solves---each one could be
done on a separate computer, in any order, because none depends on another's
result. A \emph{fold} runs a \emph{chain} of solves---each one starts from the
previous one's output, so they are intrinsically sequential. Crossing these two
questions gives a $2\times 2$ grid; three of its cells are occupied by real
corners, and a fourth corner sits outside the grid entirely, because in it we do
not write a loop at all. We take the four in turn.

\begin{figure}[p]
\centering
\begin{tikzpicture}[font=\footnotesize]
  % axis labels
  \node[cornertag, anchor=south] at (0,4.55) {operator \emph{fixed}};
  \node[cornertag, anchor=south] at (4.6,4.55) {operator \emph{varying}};
  \node[cornertag, rotate=90, anchor=south] at (-2.55,3.3) {a \textbf{map} (independent)};
  \node[cornertag, rotate=90, anchor=south] at (-2.55,0.9) {a \textbf{fold} (chained)};

  % --- top-left: fixed-operator map ---
  \begin{scope}[shift={(-1.3,2.5)}]
    \node[stage, minimum width=24mm, minimum height=22mm, anchor=south west,
      fill=simBgBlue, align=center] (q1) at (0,0) {};
    \node[anchor=north, font=\scriptsize\bfseries] at (1.2,2.0)
      {Fixed-operator map};
    % one operator box, three independent solves off it
    \node[opbox, minimum width=7mm] (op1) at (0.35,1.0) {$\Kop$};
    \foreach \i/\y in {1/1.45,2/1.0,3/0.55}{
      \node[product, minimum width=4mm, minimum height=3.5mm, scale=0.6]
        (s1\i) at (1.95,\y) {};
      \draw[flow, shorten >=1pt] (op1.east) -- (s1\i.west);}
    \node[anchor=north, font=\scriptsize\ttfamily] at (1.2,0.42)
      {solve\_family};
    \node[anchor=north, font=\scriptsize\itshape] at (1.2,0.08)
      {electrostatic, magnetostatic};
  \end{scope}

  % --- top-right: operator-varying map ---
  \begin{scope}[shift={(3.3,2.5)}]
    \node[stage, minimum width=24mm, minimum height=22mm, anchor=south west,
      fill=simBgBlue, align=center] (q2) at (0,0) {};
    \node[anchor=north, font=\scriptsize\bfseries] at (1.2,2.0)
      {Operator-varying map};
    \foreach \i/\y in {1/1.5,2/1.0,3/0.5}{
      \node[opbox, minimum width=6mm, scale=0.75] (op2\i) at (0.5,\y)
        {$\Kop(\omega)$};
      \node[product, minimum width=4mm, minimum height=3.5mm, scale=0.6]
        (t2\i) at (1.95,\y) {};
      \draw[flow, shorten >=1pt] (op2\i.east) -- (t2\i.west);}
    \node[anchor=north, font=\scriptsize\ttfamily] at (1.2,0.42)
      {frequency\_sweep};
    \node[anchor=north, font=\scriptsize\itshape] at (1.2,0.08)
      {driven};
  \end{scope}

  % --- bottom-left: state-threaded fold ---
  \begin{scope}[shift={(-1.3,0.0)}]
    \node[stage, minimum width=24mm, minimum height=22mm, anchor=south west,
      fill=simBgGreen, align=center] (q3) at (0,0) {};
    \node[anchor=north, font=\scriptsize\bfseries] at (1.2,2.0)
      {State-threaded fold};
    \foreach \i/\x in {1/0.3,2/0.95,3/1.6}{
      \node[opbox, minimum width=4mm, minimum height=4mm, scale=0.7]
        (f\i) at (\x,1.0) {};}
    \draw[flow, shorten >=1pt] (f1.east) -- (f2.west);
    \draw[flow, shorten >=1pt] (f2.east) -- (f3.west);
    \node[anchor=north, font=\scriptsize\ttfamily] at (1.2,0.42)
      {fold\_solve};
    \node[anchor=north, font=\scriptsize\itshape] at (1.2,0.08)
      {transient, AMR loop};
  \end{scope}

  % --- bottom-right: black-box (outside the grid) ---
  \begin{scope}[shift={(3.3,0.0)}]
    \node[extern, minimum width=24mm, minimum height=22mm, anchor=south west,
      align=center] (q4) at (0,0) {};
    \node[anchor=north, font=\scriptsize\bfseries] at (1.2,2.0)
      {Opaque black-box};
    \node[extern, minimum width=12mm, minimum height=8mm] (bb) at (1.2,1.0)
      {\scriptsize library};
    \node[anchor=north, font=\scriptsize\ttfamily] at (1.2,0.42)
      {eigsolve};
    \node[anchor=north, font=\scriptsize\itshape] at (1.2,0.08)
      {eigenmode, boundary-mode};
  \end{scope}

  \node[font=\scriptsize\itshape, text=simViolet, align=center]
    at (4.5,-0.65) {(no loop we write)};
\end{tikzpicture}
\caption[The four solve corners]{The four solve corners. Two questions
classify a solve: is the operator \emph{fixed} or \emph{varying} as we proceed
(columns), and is the computation a \emph{map} of independent solves or a
\emph{fold} of chained solves (rows)? Three corners occupy the resulting grid.
A fixed-operator map (top-left) reuses one operator across independent solves; an
operator-varying map (top-right) rebuilds the operator before each independent
solve; a state-threaded fold (bottom-left) chains solves so each starts from the
last one's output. The fourth corner (bottom-right, dashed) sits \emph{outside}
the grid: we write no loop at all---we hand the whole problem to a library
routine. Each corner names its combinator and the drivers that use it.}
\label{fig:corners}
\end{figure}

\subsection{Corner 1: the fixed-operator map}

The simplest corner. Assemble the operator \emph{once}, then run a family of
independent solves that all reuse it. Each solve has the same operator on the
left-hand side and a different right-hand side; the solves do not depend on one
another, so they form a \emph{map} over the family of right-hand sides.

The exemplar is electrostatics (\cref{ch:electrostatics}). To compute the
capacitance of $n$ conductors, we assemble the stiffness operator $\Kop$ once,
then solve $\Kop\,\vx_i = \vb_i$ for $n$ different right-hand sides---one per
conductor held at unit voltage. The operator $\Kop$ never changes; only the
right-hand side does. We capture $\Kop$ once, outside the loop, and map the
solver over the family. Magnetostatics works identically, with a different
operator and a different family. In the framework's vocabulary this corner is
the combinator \code{solve\_family}, written
\[
  \texttt{solve\_family}\;\Kop\;[\vb_1,\dots,\vb_n]
    \;=\; [\,\texttt{solve}\;\Kop\;\vb_1,\;\dots,\;\texttt{solve}\;\Kop\;\vb_n\,].
\]
The single load-bearing fact is the \emph{hoist}: the operator is set up once,
before the map, not once per element. Everything else about this corner follows
from that.

\subsection{Corner 2: the operator-varying map}

Now keep the map---a family of independent solves---but change one thing: the
operator is no longer fixed. It must be rebuilt before each solve.

The exemplar is the driven analysis (\cref{ch:driven}). To compute S-parameters
across a frequency band, we solve the time-harmonic system at each frequency
$\omega$ in the sweep. But the system operator depends on frequency:
$\mat{A}(\omega) = \Kop + i\omega\,\Cop - \omega^2\Mop$. So inside the sweep,
\emph{before each solve}, we must reassemble $\mat{A}(\omega)$ at the current
frequency. The solves are still independent---frequency $\omega_3$ does not need
the answer at $\omega_2$---so this is still a map, not a fold. What distinguishes
it from corner~1 is precisely that the operator is rebuilt \emph{inside} the
loop rather than hoisted outside it. The combinator is \code{frequency\_sweep}.

This contrast---hoist the operator versus rebuild it---is small to state and
large in consequence, and it is the entire structural difference between
electrostatics and the driven analysis. We will return to it in detail in
\cref{ch:driven}; for now, hold onto the picture in \cref{fig:corners}: one
operator feeding many solves (corner~1) versus a fresh operator for each solve
(corner~2).

\begin{pitfall}
It is tempting to call corner~2 a ``harder version of corner~1'' and leave it
at that. It is harder, but the difference is not a matter of degree---it is a
matter of \emph{kind}. In corner~1 the expensive operator setup happens once and
is amortized over the whole family; in corner~2 it happens every single
iteration. An implementation that ``optimizes'' corner~2 by accidentally
hoisting the operator out of the loop does not run faster---it computes the
wrong answer, because it solves every frequency against the operator of the
first one. The position of one setup call relative to one loop is load-bearing.
\end{pitfall}

\subsection{Corner 3: the state-threaded fold}

Drop the independence. In corners~1 and~2 the solves were a map: independent,
reorderable, parallelizable. In corner~3 they are a \emph{fold}: each solve
begins from the state the previous solve produced. The computation is a chain,
and the chain cannot be reordered or parallelized, because step $k$ literally
needs the output of step $k-1$ as its input.

The exemplar is the transient analysis (\cref{ch:transient}). To compute the
time-domain response, we march forward in time: the field state at time
$t_{k+1}$ is computed from the field state at time $t_k$ by taking one time
step. The state is \emph{threaded} through the computation, each step
transforming it and handing it to the next. This is a \emph{fold}, the
fundamental sequential combinator we will study at length in \cref{ch:fold}; in
the framework it is \code{fold\_solve}. The adaptive-refinement outer loop from
\cref{fig:skeleton} is also a fold of this shape---each refinement pass begins
from the mesh the last pass produced---which is why the same combinator appears
for both.

The distinction between corner~3 and the two map corners is the distinction
between a \emph{movie} and a \emph{contact sheet}. A map produces a set of
independent snapshots that happen to be collected together; a fold produces a
single evolving history in which each frame depends on the last. We will see in
\cref{ch:fold} that the map is in fact the degenerate special case of the fold
in which each step ignores the running state---but the operational difference,
parallelizable versus strictly sequential, is real and central.

\subsection{Corner 4: the opaque black-box}

The fourth corner is different in nature from the other three. In corners~1--3
\emph{we} write the loop: we control the map or the fold, and we can see every
solve it performs. In corner~4 we write no loop at all. We assemble the
operators, hand the entire problem to a specialized library routine, and receive
the answer. The iteration that produces it happens \emph{inside} the library,
where our framework cannot see it.

The exemplar is the eigenmode analysis (\cref{ch:eigenvalues}). To find the
resonant frequencies of a cavity, we assemble the operator pencil
$(\Kop,\Cop,\Mop)$ once and pass it to an eigenvalue solver---SLEPc or ARPACK,
in Palace's case---with a single call. The solver runs its own sophisticated
iteration (a Krylov--Schur or Arnoldi process; \cref{ch:eigenvalues}) and
returns the converged eigenpairs. We do not see the iteration; we see one call
and its result. The boundary-mode analysis uses the same corner, on a
two-dimensional cross-section rather than a volume. In the framework this corner
is the combinator \code{eigsolve}, and it is deliberately \emph{opaque}: it
names the library boundary and the contract across it, and stops there.

This opacity is not a gap in our understanding; it is an honest record of where
Palace's authors chose to call out to someone else's code. \Cref{ch:eigenvalues}
both documents that boundary and---separately---reconstructs what the library
does internally, in our own vocabulary, so that the black box becomes a glass
box for the reader who wants to look inside.

\begin{keyidea}
The four solve corners are the most important taxonomy in this book. They are
distinguished by two questions plus one exception:
\begin{itemize}[nosep]
  \item \textbf{Fixed-operator map} (\code{solve\_family}): assemble once, then
  map an independent family of solves over the fixed operator. \emph{Hoist the
  operator.} Exemplars: electrostatic, magnetostatic.
  \item \textbf{Operator-varying map} (\code{frequency\_sweep}): a map of
  independent solves, but rebuild the operator \emph{inside} the loop because it
  changes each step. \emph{Do not hoist.} Exemplar: driven.
  \item \textbf{State-threaded fold} (\code{fold\_solve}): a chain, not a map---
  each step's input is the previous step's output. \emph{Sequential.}
  Exemplars: transient, and the AMR outer loop.
  \item \textbf{Opaque black-box} (\code{eigsolve}): no loop we write at all;
  one call to a library solver returns the answer. Exemplars: eigenmode,
  boundary-mode.
\end{itemize}
Every later chapter that studies a driver is, in part, the study of which corner
that driver sits in and why.
\end{keyidea}

\section{The three reduction shapes}

The solve stage has four shapes; the reduce stage has three. After a driver has
produced its field or family of fields, the \texttt{reduce} stage collapses
them into the small physical quantity the user wants. As with the solve stage,
the physics differs from analysis to analysis but the \emph{shape} of the
reduction does not---and there are only three shapes (\cref{fig:reductions}).

\begin{figure}[t]
\centering
\begin{tikzpicture}[font=\footnotesize]
  % --- (a) symmetric Gram matrix ---
  \begin{scope}[shift={(0,0)}]
    \foreach \i in {0,1,2}{
      \foreach \j in {0,1,2}{
        \pgfmathtruncatemacro{\sym}{(\i==\j)?1:0}
        \node[draw=simGray, semithick,
          fill={\ifnum\i=\j simBgRust\else simBgBlue\fi},
          minimum size=6.5mm, inner sep=0pt]
          (g\i\j) at (\j*0.66, -\i*0.66) {};}}
    % symmetry mirror line
    \draw[densely dashed, simRust, semithick]
      (-0.45,0.45) -- (1.77,-1.77);
    \node[font=\scriptsize\bfseries, align=center] at (0.66,-2.45)
      {Rank-2 Gram};
    \node[font=\scriptsize, align=center] at (0.66,-2.9)
      {symmetric: $G_{ij}=G_{ji}$};
    \node[font=\scriptsize\itshape, align=center] at (0.66,-3.4)
      {capacitance,\\inductance};
  \end{scope}

  % --- (b) non-symmetric port-projection ---
  \begin{scope}[shift={(4.2,0)}]
    \foreach \i in {0,1,2}{
      \foreach \j in {0,1,2}{
        \node[draw=simGray, semithick, fill=simBgGreen,
          minimum size=6.5mm, inner sep=0pt]
          (p\i\j) at (\j*0.66, -\i*0.66) {};}}
    \node[font=\scriptsize, text=simGreen] at (0.66,0.55)
      {ports out};
    \node[font=\scriptsize, rotate=90, text=simGreen] at (-0.62,-0.66)
      {ports in};
    \node[font=\scriptsize\bfseries, align=center] at (0.66,-2.45)
      {Rank-2 projection};
    \node[font=\scriptsize, align=center] at (0.66,-2.9)
      {\emph{not} symmetric};
    \node[font=\scriptsize\itshape, align=center] at (0.66,-3.4)
      {S-parameters};
  \end{scope}

  % --- (c) rank-1 scalar / mode table ---
  \begin{scope}[shift={(8.4,0)}]
    \foreach \i in {0,1,2,3}{
      \node[draw=simGray, semithick, fill=simBgGold,
        minimum width=18mm, minimum height=5mm, inner sep=0pt]
        (r\i) at (0, -\i*0.6) {};}
    \node[font=\scriptsize] at (0,0) {mode 1: $f_1, Q_1$};
    \node[font=\scriptsize] at (0,-0.6) {mode 2: $f_2, Q_2$};
    \node[font=\scriptsize] at (0,-1.2) {mode 3: $f_3, Q_3$};
    \node[font=\scriptsize] at (0,-1.8) {$\vdots$};
    \node[font=\scriptsize\bfseries, align=center] at (0,-2.45)
      {Rank-1 table};
    \node[font=\scriptsize, align=center] at (0,-2.9)
      {one row per element};
    \node[font=\scriptsize\itshape, align=center] at (0,-3.4)
      {eigenfreq.\,$+\,Q$,\\energy, modes};
  \end{scope}
\end{tikzpicture}
\caption[The three reduction shapes]{The three reduction shapes. A
\emph{rank-2 Gram} (left) is a symmetric matrix of energy-like inner products
over a family of solutions, $G_{ij}=w(i,j)\,\vx_j^{\!\top}\Kop\,\vx_i$; its
symmetry (dashed mirror) is structural. A \emph{rank-2 projection} (centre) is
also a matrix, but it projects an ``in'' family onto an ``out'' family and is
\emph{deliberately not} a Gram---it is generally non-symmetric. A \emph{rank-1
table} (right) is a one-dimensional list: one row of scalars per mode, per
domain, or per propagation mode. The distinction between the left and centre
shapes is the single most over-unified point in the design---see the pitfall.}
\label{fig:reductions}
\end{figure}

\subsection{The rank-2 symmetric Gram}

The first shape is a \emph{Gram matrix}: a symmetric matrix of energy-like inner
products formed over a family of solutions. Given a solution family
$\vx_1,\dots,\vx_n$ and the assembled operator $\Kop$, the Gram reduction
computes
\[
  G_{ij} \;=\; w(i,j)\;\vx_j^{\!\top}\,\Kop\,\vx_i,
\]
a matrix whose $(i,j)$ entry is an operator-weighted inner product of solutions
$i$ and $j$. Because $\Kop$ is symmetric, $G_{ij}=G_{ji}$: the result is a
symmetric matrix, structurally, not by coincidence.

Both capacitance and inductance are Gram reductions. They use the \emph{same
verb}---\code{gram\_reduce} in the framework---and differ only in the weight
$w(i,j)$. Capacitance uses the unit weight $w(i,j)=1$ (the energy form
$C_{ii}=2U/V^2$ with unit applied voltage); inductance uses the
current-normalized weight $w(i,j)=1/(I_i I_j)$. One verb, two weights, two
physical products. This is the cleanest instance of the book's recurring
refrain: the same machine with one part---here, a scalar weight function---
swapped. We study it fully in \cref{ch:reductions}.

\subsection{The rank-2 port-projection}

The second shape is also a matrix, and it is tempting to file it under the
first. Resist. The S-parameter matrix is a \emph{projection}, not a Gram. Its
entry $S_{ij}(\omega)$ records how much of a wave entering port $j$ leaves port
$i$: it projects the solved field for each driven port onto each output port's
mode. It is a matrix indexed by (out-port, in-port), but the two index families
play different roles, the operation is a linear projection rather than an
energy-weighted self-inner-product, and the result is in general
\emph{not symmetric}. The framework gives it its own verb,
\code{sparameter\_reduce}, deliberately distinct from \code{gram\_reduce}.

\begin{pitfall}
Two reductions producing matrices is not enough to make them the same reduction.
The Gram products (capacitance, inductance) and the port-projection
(S-parameters) both yield rank-2 results, and a tidy-minded designer is sorely
tempted to unify them under one ``matrix reduction'' verb. This is a mistake.
The Gram is a \emph{symmetric self-pairing} of a family with itself through an
operator; the port-projection is an \emph{asymmetric projection} of one family
onto a different basis. Forcing them into one verb either smuggles a false
symmetry into the S-matrix or burdens the Gram with machinery it does not need.
The right move---which the framework makes---is two verbs that happen to share
an output rank. \emph{Sharing a shape is not sharing a meaning.}
\end{pitfall}

\subsection{The rank-1 scalar or mode table}

The third shape is a \emph{table}: a one-dimensional list with one row per
element, each row holding a few scalars (or, in one case, a few small fields).
There is no matrix structure, no pairing of a family with itself---just a map
over a collection, producing one tuple per item.

Three products live here. Eigenfrequencies and quality factors form a table with
one row per mode: $(f_m, Q_m)$ for each resonant mode. Energy and participation
factors form a table with one row per material domain: how much of the field
energy sits in each region. And waveguide propagation modes form a table with
one row per mode, each carrying a propagation constant $k_n$ and a small
transverse field pattern (the one case where the row holds fields, not just
scalars). All three are rank-1 reductions---a \code{map} over the items the
driver produced, with no cross-coupling between rows.

\begin{keyidea}
The three reduction shapes are: the \textbf{rank-2 symmetric Gram} (capacitance,
inductance---one verb, the weight is the only difference); the \textbf{rank-2
port-projection} (S-parameters---a projection, deliberately not a Gram); and the
\textbf{rank-1 scalar or mode table} (eigenfrequency$+Q$, energy and
participation, waveguide modes---one row per item). A matrix-shaped result does
not by itself mean a Gram: the projection shares the Gram's rank but not its
symmetry or its meaning.
\end{keyidea}

\section{The driver \texorpdfstring{$\times$}{x} framework grid}

We can now lay the six drivers on the framework all at once. \Cref{tab:grid}
is the grid the rest of the book fills in. Each row is a driver; each column is
a stage of the skeleton---the function space it assembles on, what it assembles,
which solve corner its solve sits in, which reduction shape its reduce uses, and
the product it yields. Every cell is a forward reference: the chapters of
Parts~II through~V supply the contents of these cells in full, one stage at a
time.

\begin{table}[t]
\centering
\caption[The driver $\times$ framework grid]{The driver $\times$ framework grid:
the six drivers laid on the skeleton. Reading across a row tells you how that
analysis is built; reading down a column tells you how that stage varies across
analyses. The \emph{solve corner} and \emph{reduction} columns are the two that
this chapter named; almost every difference between rows lives there. The
electrostatic row is the simplest; every other row is the electrostatic row
with one or two stages swapped (last paragraph).}
\label{tab:grid}
\footnotesize
\setlength{\tabcolsep}{3.5pt}
\renewcommand{\arraystretch}{1.3}
\begin{tabular}{@{}llp{27mm}ll@{}}
\toprule
Driver & FE space & \multicolumn{1}{l}{Assemble} & Solve corner & Reduction $\to$ product \\
\midrule
Electrostatic & $\Hone$ & $\Kop$ once
  & fixed-operator map & Gram ($w{=}1$) $\to$ capacitance \\
Magnetostatic & $\Hcurl$ & $\Kop$ once
  & fixed-operator map & Gram ($w{=}\tfrac{1}{I_iI_j}$) $\to$ inductance \\
Driven & $\Hcurl$ & $\Kop,\Cop,\Mop$ once; $\mat{A}(\omega)$ rebuilt in loop
  & operator-varying map & projection $\to$ S-parameters \\
Transient & $\Hcurl$ & $\Kop,\Cop,\Mop$ once
  & state-threaded fold & table $\to$ field history \\
Eigenmode & $\Hcurl$ & $\Kop,\Cop,\Mop$ once
  & opaque black-box & table $\to$ frequencies, $Q$ \\
Boundary mode & $\Hcurl{\oplus}\Hone$ & block pencil (2-D)
  & opaque black-box & table $\to$ waveguide modes \\
\bottomrule
\end{tabular}
\end{table}

Read down the \emph{Solve corner} column and the chapter's central claim becomes
a glance: two drivers map over a fixed operator, one maps over a varying
operator, one folds a threaded state, and two hand the whole problem to a black
box. Read down the \emph{Reduction} column and the three shapes appear: two
Grams, one projection, three tables. The physics columns (space, assemble)
supply variety too, but they are the subject of Part~II; the two columns this
chapter named are the ones that organize the book.

There is a punchline, and it is worth stating plainly. \emph{Electrostatics is
the universal entry point, and every other analysis is electrostatics with one
stage swapped.} Magnetostatics swaps the function space ($\Hone\to\Hcurl$) and
the Gram weight, and changes nothing else. The driven analysis keeps the map but
swaps the solve corner from fixed to operator-varying, and swaps the reduction
from Gram to projection. The transient analysis swaps the map for a fold.
Eigenmode swaps the loop-we-write for a black box. Boundary mode is eigenmode on
a two-dimensional cross-section. Six analyses, one anchor, a handful of swaps:
that is the whole architecture of the simulator, and it is why the book can
teach electrostatics once, in depth (\cref{ch:electrostatics}), and then present
each remaining analysis as a small, well-localized set of differences from it.

\summarysection
The skeleton is
\texttt{config} $\to$ \texttt{assemble} $\to$ \texttt{solve} $\to$
\texttt{reduce} $\to$ \texttt{product},
optionally wrapped in an adaptive-refinement outer fold. Almost all the variety
between the six analyses lives in two stages. The \texttt{solve} stage has four
\emph{corners}: the fixed-operator map (\code{solve\_family}; electrostatic,
magnetostatic), the operator-varying map (\code{frequency\_sweep}; driven), the
state-threaded fold (\code{fold\_solve}; transient and the AMR loop), and the
opaque black-box (\code{eigsolve}; eigenmode, boundary-mode). The \texttt{reduce}
stage has three \emph{shapes}: the rank-2 symmetric Gram (capacitance and
inductance, one verb differing only by weight), the rank-2 port-projection
(S-parameters, a projection and deliberately not a Gram), and the rank-1
scalar-or-mode table (eigenfrequency$+Q$, energy, waveguide modes). The driver
$\times$ framework grid (\cref{tab:grid}) places all six on these axes, and the
punchline organizes the rest of the book: electrostatics is the universal entry
point, and every other analysis swaps one stage.

\furtherreading
The four-corners and three-shapes taxonomy is this book's own organizing
device, distilled from the structure of the Palace solver
itself~\citep{palace2023}; reading the source's six driver entry points side by
side is the best way to see the shared skeleton in the wild. The distinction
between an independent map and a state-threaded fold is the functional
programmer's daily bread; we build it from scratch in \cref{ch:fold}, assuming
nothing. For the numerical content behind the corners---why the fixed-operator
map can amortize a factorization, why the eigenvalue corner is best left to a
specialized library---Saad's account of iterative methods~\citep{saad2003} is
the standard reference, and we draw on it throughout Part~IV.

\exercisessection
\begin{enumerate}[nosep]
  \item For each of the four solve corners, state the two-question
  classification that places it: is the operator fixed or varying, and is the
  computation a map or a fold? (The black-box corner answers ``neither row''---
  explain why.)
  \item The driven analysis (corner~2) and electrostatics (corner~1) both run a
  family of independent solves. Pinpoint the single line of pseudocode whose
  position---inside the loop versus outside it---is the entire difference between
  them. Why does moving it change the \emph{answer}, not merely the speed?
  \item Capacitance and inductance share one reduction verb and differ only in a
  weight. Write down each weight, and explain in one sentence why inductance
  needs the current normalization that capacitance does not.
  \item S-parameters and capacitance both produce a matrix, yet the book insists
  they are different reduction shapes. Give two structural reasons the
  S-parameter matrix is not a Gram. (Hint: consider symmetry, and consider what
  each index ranges over.)
  \item Using \cref{tab:grid}, describe the transient analysis as ``the
  electrostatic analysis with these stages swapped.'' Exactly which cells change,
  and which stay the same?
\end{enumerate}


\part{The Field Theory}\label{part:fieldtheory}
\chapter{Maxwell's Equations: A Working Review}
\label{ch:maxwell}

% Chapter-local vector macros not in style/simbook.sty (\vs is the time-domain
% state vector; the Poynting vector, wavevector, and z-hat are local here).
\newcommand{\vS}{\mathbf{S}}
\providecommand{\vk}{\mathbf{k}}
\providecommand{\vz}{\mathbf{z}}

\chapterorient{You have met Maxwell's equations before---probably as four lines
on a slide, admired and then set aside. This chapter brings them back as
\emph{working} equations: the ones the simulator actually discretizes. We will
read each of the four laws for its physical content, attach the materials
through the constitutive relations, pin down what happens at a wall or an
interface, account for the energy, and---most important for everything that
follows---split the physics into the three regimes (static, time-harmonic,
time-domain) that become the analyses of Part~II onward. By the end you will
recognize, in raw continuous form, every operator the rest of the book learns
to assemble and solve.}

\section{The four laws, read for meaning}

Electromagnetism is the statement that two vector fields---the electric field
$\vE$ and the magnetic field $\vB$---fill all of space, push on charges, and
generate each other. Maxwell's four equations say exactly \emph{how} the fields
arrange themselves around their sources and how a change in one provokes the
other. In differential (point-by-point) form, in a region of space with charge
density $\rho$ and current density $\vJ$, they are

\begin{align}
  \Div \vD &= \rho, & &\text{(Gauss's law for $\vE$)} \label{eq:gauss-e}\\
  \Div \vB &= 0, & &\text{(Gauss's law for $\vB$)} \label{eq:gauss-b}\\
  \Curl \vE &= -\,\partial_t \vB, & &\text{(Faraday's law)} \label{eq:faraday}\\
  \Curl \vH &= \vJ + \partial_t \vD. & &\text{(Amp\`ere--Maxwell law)} \label{eq:ampere}
\end{align}

Here $\vD$ is the electric displacement and $\vH$ the magnetic field intensity;
they are the ``material-aware'' partners of $\vE$ and $\vB$, related to them by
the constitutive relations of \cref{sec:constitutive}. For the moment read
$\vD$ as ``$\vE$, scaled by how the material responds'' and $\vH$ as
``$\vB$, scaled the same way.'' The operators $\Div$ and $\Curl$ are the
ordinary divergence and curl of vector calculus; if they feel rusty,
\cref{tab:vectorcalc} is a one-line reminder of what each measures.

\begin{table}[t]
\centering
\caption[The two derivative operators]{The two first-order vector-calculus
operators that build every Maxwell equation. \emph{Divergence} measures
net outflow from a point (a source); \emph{curl} measures local circulation
(a swirl). Maxwell's equations are entirely statements about the divergence
and curl of $\vE$ and $\vB$.}
\label{tab:vectorcalc}
\small
\setlength{\tabcolsep}{6pt}
\renewcommand{\arraystretch}{1.25}
\begin{tabular}{@{}lll@{}}
\toprule
Operator & Reads & Measures \\
\midrule
$\Div \vF$ & ``div of $\vF$'' & net flux out of an infinitesimal box (a source/sink) \\
$\Curl \vF$ & ``curl of $\vF$'' & circulation around an infinitesimal loop (a swirl) \\
\bottomrule
\end{tabular}
\end{table}

A field theory is much easier to feel than to memorize. Take the four laws one
at a time, each with the picture it draws (\cref{fig:fourlaws}).

\begin{figure}[t]
\centering
\begin{tikzpicture}[scale=1.0]
% --- Gauss for E: charge with radial field lines ---
\begin{scope}[shift={(0,0)}]
  \fill[simRust] (0,0) circle (2.2pt);
  \node[font=\scriptsize] at (0.0,-0.28) {$+q$};
  \foreach \a in {0,45,...,315}{
    \draw[flow,simBlue,shorten <=3pt]
      (0,0) -- (\a:0.95);}
  \node[font=\scriptsize\itshape,simBlue] at (0.85,0.85) {$\vE$};
  \node[font=\footnotesize] at (0,-1.35) {$\Div\vD=\rho$};
  \node[font=\scriptsize,text=simGray,align=center] at (0,-1.75)
    {charge is\\ a source};
\end{scope}
% --- Gauss for B: dipole loop, no source ---
\begin{scope}[shift={(3.6,0)}]
  \draw[flow,simTeal] (0,0) circle (0.75);
  \draw[flow,simTeal] (0,0) circle (0.42);
  \node[font=\scriptsize\itshape,simTeal] at (0.95,0.0) {$\vB$};
  \fill[simInk] (0,0) circle (1.2pt);
  \node[font=\footnotesize] at (0,-1.35) {$\Div\vB=0$};
  \node[font=\scriptsize,text=simGray,align=center] at (0,-1.75)
    {lines close;\\ no source};
\end{scope}
% --- Faraday: changing B makes circulating E ---
\begin{scope}[shift={(7.2,0)}]
  \draw[-{Stealth[length=2mm]},simRust,thick] (0,-0.55)--(0,0.55)
    node[right,font=\scriptsize]{$\partial_t\vB$};
  \draw[flow,simBlue] (0.85,0) arc (0:300:0.85);
  \node[font=\scriptsize\itshape,simBlue] at (-0.95,0.5) {$\vE$};
  \node[font=\footnotesize] at (0,-1.35) {$\Curl\vE=-\partial_t\vB$};
  \node[font=\scriptsize,text=simGray,align=center] at (0,-1.75)
    {changing $\vB$\\ swirls $\vE$};
\end{scope}
% --- Ampere-Maxwell: current makes circulating H ---
\begin{scope}[shift={(10.8,0)}]
  \draw[-{Stealth[length=2mm]},simGreen,thick] (0,-0.55)--(0,0.55)
    node[right,font=\scriptsize]{$\vJ$};
  \draw[flow,simViolet] (0.85,0) arc (0:300:0.85);
  \node[font=\scriptsize\itshape,simViolet] at (-0.95,0.5) {$\vH$};
  \node[font=\footnotesize] at (0,-1.35) {$\Curl\vH=\vJ+\partial_t\vD$};
  \node[font=\scriptsize,text=simGray,align=center] at (0,-1.75)
    {current swirls\\ $\vH$};
\end{scope}
\end{tikzpicture}
\caption[The four laws as pictures]{Each Maxwell equation as a field-line
sketch. \emph{Gauss for $\vE$}: field lines begin and end on charge.
\emph{Gauss for $\vB$}: field lines form closed loops---there is no magnetic
charge. \emph{Faraday}: a time-varying magnetic field wraps an electric field
around it. \emph{Amp\`ere--Maxwell}: a current (or a time-varying $\vD$) wraps a
magnetic field around it. The last two are mirror images with a crucial sign
difference.}
\label{fig:fourlaws}
\end{figure}

\subsection{Gauss's law: charge is the source of \texorpdfstring{$\vE$}{E}}

Equation~\eqref{eq:gauss-e}, $\Div\vD=\rho$, says that electric field lines
\emph{begin and end on charge}. The divergence of $\vD$ at a point is the net
number of field lines streaming out of an infinitesimal box around that
point; the equation pins that outflow to the charge density enclosed. Where
there is positive charge, lines emanate; where negative, they terminate; in
empty space the divergence is zero and lines merely pass through. Integrate
over a closed surface and you recover the familiar integral form,
$\oint \vD\cdot\vn\,\mathrm{d}S = Q_{\text{enc}}$: the flux through any closed
surface equals the charge inside. This is the law that makes a capacitor a
capacitor, and it is the entire content of the electrostatic analysis.

\subsection{Gauss's law for \texorpdfstring{$\vB$}{B}: no magnetic charge}

Equation~\eqref{eq:gauss-b}, $\Div\vB=0$, is the same statement with the
right-hand side set permanently to zero: there is no magnetic charge. Magnetic
field lines never begin or end; they always close on themselves
(\cref{fig:fourlaws}, second panel). This innocent-looking zero is
load-bearing for the discretization. It forces $\vB$ to be \emph{divergence
free}, and a finite-element method that fails to honour that property produces
spurious, unphysical modes. We will spend real effort in \cref{ch:bc} building
function spaces that respect $\Div\vB=0$ exactly rather than approximately.

\subsection{Faraday's law: a changing \texorpdfstring{$\vB$}{B} drives \texorpdfstring{$\vE$}{E}}

Equation~\eqref{eq:faraday}, $\Curl\vE=-\partial_t\vB$, is the first of the two
\emph{coupling} laws---the ones that knit $\vE$ and $\vB$ together into a single
electromagnetic field. It says a magnetic field that changes in time wraps an
electric field around itself. The curl of $\vE$---its local circulation---is
fed by the rate of change of $\vB$ threading through. This is the law of the
electric generator and the transformer: move a magnet, and the changing flux
induces a circulating $\vE$ that drives current. The minus sign is Lenz's law
in disguise; it ensures the induced effect \emph{opposes} the change that
caused it, which is why energy is conserved rather than spontaneously created.

\subsection{Amp\`ere--Maxwell: current (and changing \texorpdfstring{$\vD$}{D}) drives \texorpdfstring{$\vH$}{H}}

Equation~\eqref{eq:ampere}, $\Curl\vH=\vJ+\partial_t\vD$, is Faraday's mirror.
A current $\vJ$ wraps a magnetic field around itself---the law of the
electromagnet. Maxwell's own contribution was the second term, the
\emph{displacement current} $\partial_t\vD$: even with no moving charge, a
\emph{changing} electric field sources $\vH$ exactly as a real current would.
That term is what lets the two coupling laws close a loop---$\vE$ feeds $\vB$
through Faraday, $\vB$ feeds $\vE$ through Amp\`ere--Maxwell---and that loop,
propagating through space, \emph{is} light. Without the displacement current,
electromagnetic waves do not exist.

We can make ``that loop is light'' precise in three lines, and it is worth
doing once. Take the curl of Faraday's law~\eqref{eq:faraday} and feed in
Amp\`ere--Maxwell~\eqref{eq:ampere} (in vacuum, $\vJ=\bm0$,
$\vD=\varepsilon_0\vE$, $\vB=\mu_0\vH$):
\[
  \Curl(\Curl\vE)
    = -\,\partial_t(\Curl\vB)
    = -\,\mu_0\varepsilon_0\,\partial_t^2\vE.
\]
The vector identity $\Curl\Curl\vE=\Grad(\Div\vE)-\Lap\vE$, together with
$\Div\vE=0$ in charge-free vacuum, collapses the left side to $-\Lap\vE$, and we
are left with the \emph{wave equation}
\begin{equation}
  \Lap\vE - \mu_0\varepsilon_0\,\partial_t^2\vE = \bm0. \label{eq:waveeq}
\end{equation}
This is the equation of a disturbance propagating at speed
$c=1/\sqrt{\mu_0\varepsilon_0}$---and the only thing it took to produce it was
the displacement current, which supplied the second time derivative. Maxwell's
quiet extra term turned a pair of static field laws into a dynamical theory of
waves. The same operator $\Curl\Curl-\,(\text{constant})\,\partial_t^2$ is, in
the frequency domain, the curl--curl operator $\mat K-\omega^2\mat M$ of
Worked example~\ref{wex:phasor}; we will meet it again and again.

\begin{keyidea}
The four laws split cleanly into two pairs. The two \emph{Gauss} laws are
constraints on the field at one instant---they say where lines can begin and
end. The two \emph{curl} laws are the dynamics---they say how $\vE$ and $\vB$
generate each other in time. Static analyses keep only the constraints; wave
and resonance analyses turn on the coupling. Every analysis in this book is a
choice of which terms of \eqref{eq:gauss-e}--\eqref{eq:ampere} survive.
\end{keyidea}

\subsection{The continuity equation: charge is conserved}

The four laws are not independent; they already contain the conservation of
charge. Take the divergence of the Amp\`ere--Maxwell law~\eqref{eq:ampere}. The
divergence of any curl is identically zero, so the left side vanishes:
\[
  0 = \Div(\Curl\vH) = \Div\vJ + \partial_t(\Div\vD).
\]
Substitute Gauss's law $\Div\vD=\rho$ into the last term and you obtain the
\emph{continuity equation}
\begin{equation}
  \Div\vJ + \partial_t\rho = 0. \label{eq:continuity}
\end{equation}
In words: the net current flowing out of any region equals the rate at which
the charge inside it drops. Charge is never created or destroyed---it only
flows. This is not an extra assumption; it is a \emph{consequence} of Maxwell's
equations, and the displacement current is precisely the term that makes it
come out right. (Had Maxwell omitted $\partial_t\vD$, the divergence of
\eqref{eq:ampere} would have forced $\Div\vJ=0$, false whenever charge piles
up.) We will meet the identity $\Div(\Curl(\cdot))=0$ again as one half of the
\emph{de Rham complex} in \cref{ch:derham}; it is the same algebraic fact wearing
different clothes.

\begin{pitfall}
It is tempting to treat $\Div\vJ+\partial_t\rho=0$ as a fifth law to be imposed
separately. It is not---it follows from the other four. In the discrete setting
the analogue is subtle: a finite-element scheme that discretizes the curl laws
\emph{without} a compatible discretization of the divergence will generally
\emph{violate} discrete continuity, leaking charge. Honouring these identities
\emph{discretely}, not just continuously, is the entire reason for the
structured function spaces of \cref{ch:functionspaces}.
\end{pitfall}

\subsection{The integral forms, and why we keep the differential ones}
\label{sec:integral}

Each differential law has an \emph{integral} twin, obtained by integrating over
a region and applying the two great theorems of vector calculus---the
divergence theorem (for the Gauss laws) and Stokes' theorem (for the curl
laws). The divergence theorem turns a volume integral of a divergence into a
flux through the bounding surface; Stokes' theorem turns a surface integral of a
curl into a circulation around the bounding loop (\cref{fig:theorems}). Applied
to the four laws they give

\begin{align}
  \oint_{\partial \mathcal V} \vD\cdot\vn\,\mathrm{d}S &= \int_{\mathcal V}\rho\dV
    = Q_{\text{enc}}, & &\text{(Gauss, integral)}\\
  \oint_{\partial \mathcal V} \vB\cdot\vn\,\mathrm{d}S &= 0, & &\text{(no monopole)}\\
  \oint_{\partial \mathcal S} \vE\cdot\mathrm{d}\bm\ell
    &= -\frac{\mathrm d}{\mathrm dt}\int_{\mathcal S}\vB\cdot\vn\,\mathrm{d}S,
    & &\text{(Faraday: emf $=-\dot\Phi$)}\\
  \oint_{\partial \mathcal S} \vH\cdot\mathrm{d}\bm\ell
    &= \int_{\mathcal S}\Big(\vJ+\partial_t\vD\Big)\cdot\vn\,\mathrm{d}S.
    & &\text{(Amp\`ere--Maxwell)}
\end{align}

These are the forms a physicist reaches for first: ``flux through a closed
surface equals enclosed charge,'' ``the emf around a loop equals minus the rate
of change of flux through it.'' They are equivalent to the differential laws
wherever the fields are smooth---the differential and integral forms are two
readings of the same content, related by the integral theorems. So why does a
simulator prefer the \emph{differential} forms? Because the finite-element
method discretizes a \emph{local} statement: it asks the equation to hold,
weakly, on each little element, and a differential law is already local. The
integral forms remain the right tool for \emph{deriving} the interface
conditions of \cref{sec:bc} (shrink a loop or a pillbox across the interface)
and for \emph{checking} a computed field (does the flux through this surface
equal the charge we put inside?), and we will use them for exactly those two
purposes.

\begin{figure}[t]
\centering
\begin{tikzpicture}[scale=1.0]
% --- divergence theorem: volume with outward flux ---
\begin{scope}[shift={(0,0)}]
  \draw[simInk,thick] (0,0) circle (1.0);
  \fill[simBgBlue,opacity=0.5] (0,0) circle (1.0);
  \node[font=\scriptsize] at (0,0) {$\mathcal V$};
  \foreach \a in {20,90,160,230,300}{
    \draw[-{Stealth[length=1.8mm]},simBlue,thick] (\a:1.0)--(\a:1.45);}
  \node[font=\scriptsize,align=center] at (0,-1.55)
    {$\displaystyle\int_{\mathcal V}\!\Div\vF\,\mathrm{d}V
      =\!\oint_{\partial\mathcal V}\!\vF\cdot\vn\,\mathrm{d}S$};
  \node[font=\footnotesize,text=simGray] at (0,1.7) {divergence theorem};
\end{scope}
% --- Stokes theorem: surface with circulation on boundary ---
\begin{scope}[shift={(5.4,0)}]
  \fill[simBgGold,opacity=0.6] (0,0) ellipse (1.05 and 0.7);
  \draw[flow,simRust] (1.05,0) arc (0:355:1.05 and 0.7);
  \node[font=\scriptsize] at (0,0) {$\mathcal S$};
  \draw[-{Stealth[length=1.8mm]},simGreen,thick] (0,0)--(0,0.55)
    node[right,font=\scriptsize]{$\Curl\vF$};
  \node[font=\scriptsize,align=center] at (0,-1.55)
    {$\displaystyle\int_{\mathcal S}\!\Curl\vF\cdot\vn\,\mathrm{d}S
      =\!\oint_{\partial\mathcal S}\!\vF\cdot\mathrm d\bm\ell$};
  \node[font=\footnotesize,text=simGray] at (0,1.7) {Stokes' theorem};
\end{scope}
\end{tikzpicture}
\caption[The two integral theorems]{The divergence theorem (left) converts the
volume integral of a divergence into the outward flux through the bounding
surface; it turns the differential Gauss laws into their integral forms.
Stokes' theorem (right) converts the surface integral of a curl into the
circulation around the bounding loop; it turns Faraday's and Amp\`ere's laws
into their integral forms. The same two theorems, applied to a test function
instead of the field itself, produce the \emph{weak form} of \cref{ch:weak}.}
\label{fig:theorems}
\end{figure}

These two theorems are not a one-time convenience. The integration by parts that
produces the finite-element \emph{weak form} in \cref{ch:weak} is the
divergence theorem applied to a product of the field and a test function---the
very same identity, used to move a derivative off the field and onto the test
function. The integral forms are the first appearance of the machinery the whole
discretization rests on.

\section{Materials: the constitutive relations}
\label{sec:constitutive}

Maxwell's four laws contain four fields---$\vE,\vD,\vB,\vH$---but only two are
truly independent. The other two are tied to them by the \emph{material}: the
\emph{constitutive relations} encode how a given substance responds to a field.
For a linear, isotropic medium they are three scalar multiplications,

\begin{align}
  \vD &= \varepsilon\,\vE, & &\text{(permittivity: electric response)} \label{eq:const-eps}\\
  \vB &= \mu\,\vH, & &\text{(permeability: magnetic response)} \label{eq:const-mu}\\
  \vJ &= \sigma\,\vE. & &\text{(conductivity: Ohm's law)} \label{eq:const-sigma}
\end{align}

The three coefficients are the only place the \emph{stuff} of the device enters
the physics. Everything else is geometry and excitation.

\begin{itemize}
  \item The \textbf{permittivity} $\varepsilon$ measures how much a material
  polarizes---how its bound charges shift---in response to $\vE$, and so how
  much electric flux $\vD$ a given field produces. In vacuum
  $\varepsilon=\varepsilon_0\approx\SI{8.854e-12}{\farad\per\meter}$. For a
  dielectric one writes $\varepsilon=\varepsilon_r\varepsilon_0$ with the
  \emph{relative permittivity} $\varepsilon_r$ (the ``dielectric constant'')
  ranging from $\approx 1$ (air) through $\approx 3.9$ (silicon dioxide) and
  $\approx 11.7$ (silicon) to $\approx 80$ (water).

  \item The \textbf{permeability} $\mu$ plays the same role magnetically. In
  vacuum $\mu=\mu_0=4\pi\times10^{-7}\,\si{\henry\per\meter}$. Most materials of
  interest in microwave engineering are non-magnetic, $\mu_r\approx 1$;
  ferromagnets ($\mu_r$ in the hundreds or thousands) are the exception. In the
  magnetostatic analysis it is the \emph{reluctivity} $\nu=1/\mu$ that appears
  in the operator, exactly as $\varepsilon$ appears in the electrostatic one.

  \item The \textbf{conductivity} $\sigma$ measures how freely charge flows:
  Ohm's law~\eqref{eq:const-sigma} says the current is proportional to the
  field. A perfect insulator has $\sigma=0$; a perfect conductor has
  $\sigma\to\infty$; copper sits at $\approx\SI{6e7}{\siemens\per\meter}$. The
  conductivity is what makes a material \emph{lossy}---it is the term that
  drains electromagnetic energy into heat, and (as we will see in
  \cref{sec:regimes}) it is what gives a resonant mode a finite quality factor
  $Q$ rather than ringing forever.
\end{itemize}

\begin{table}[t]
\centering
\caption[Typical material parameters]{Representative values of the three
constitutive coefficients. Relative permittivity $\varepsilon_r$ and relative
permeability $\mu_r$ are dimensionless ratios to the vacuum values
$\varepsilon_0\approx\SI{8.854e-12}{\farad\per\meter}$ and
$\mu_0=4\pi\times10^{-7}\,\si{\henry\per\meter}$; conductivity $\sigma$ is in
$\si{\siemens\per\meter}$. The enormous range of $\sigma$---from an insulator
to a metal---is what separates a dielectric from a conductor, and what makes a
boundary a PEC wall in the limit.}
\label{tab:materials}
\small
\setlength{\tabcolsep}{7pt}
\renewcommand{\arraystretch}{1.2}
\begin{tabular}{@{}lccc@{}}
\toprule
Material & $\varepsilon_r$ & $\mu_r$ & $\sigma\;(\si{\siemens\per\meter})$ \\
\midrule
Vacuum             & $1$     & $1$    & $0$ \\
Air                & $1.0006$& $1$    & $\approx 0$ \\
Silicon dioxide (\(\mathrm{SiO_2}\)) & $3.9$ & $1$ & $\sim 10^{-12}$ \\
Silicon            & $11.7$  & $1$    & varies (doping) \\
Water (\SI{20}{\celsius}) & $\approx 80$ & $1$ & $\sim 5\times10^{-4}$ \\
Copper             & ---     & $1$    & $5.96\times10^{7}$ \\
Iron (soft)        & ---     & $\sim 5000$ & $1.0\times10^{7}$ \\
\bottomrule
\end{tabular}
\end{table}

\begin{notationbox}
Three symbols, three responses: $\varepsilon$ couples $\vE\!\to\!\vD$ (electric),
$\mu$ couples $\vH\!\to\!\vB$ (magnetic), $\sigma$ couples $\vE\!\to\!\vJ$
(conduction). Their reciprocals appear in the operators: reluctivity
$\nu=1/\mu$ weights the curl--curl term, and $1/\varepsilon$ never appears
alone---$\varepsilon$ itself is the weight in the electrostatic stiffness. Keep
the pairing $(\vD,\vE,\varepsilon)$ and $(\vB,\vH,\mu)$ straight: $\vD$ and
$\vB$ are the ``flux'' fields (the ones whose divergence the Gauss laws
constrain); $\vE$ and $\vH$ are the ``intensity'' fields (the ones whose curl
the dynamic laws constrain).
\end{notationbox}

\subsection{Isotropic, anisotropic, and the material tensor}

Equations~\eqref{eq:const-eps}--\eqref{eq:const-sigma} assume the material
responds the \emph{same in every direction}---it is \emph{isotropic}, and the
coefficients are scalars. Real materials need not be so accommodating. A
crystal, a layered substrate, or a magnetized ferrite responds differently
along different axes, and then $\varepsilon$ becomes a $3\times3$ symmetric
\emph{tensor} $\bm{\varepsilon}$:
\[
  \vD = \bm{\varepsilon}\,\vE,
  \qquad
  \bm{\varepsilon} =
  \begin{pmatrix}
    \varepsilon_{xx} & \varepsilon_{xy} & \varepsilon_{xz}\\
    \varepsilon_{xy} & \varepsilon_{yy} & \varepsilon_{yz}\\
    \varepsilon_{xz} & \varepsilon_{yz} & \varepsilon_{zz}
  \end{pmatrix}.
\]
Now $\vD$ need not point along $\vE$. The isotropic case is the diagonal one
with all three entries equal. The simulator treats the tensor case as the
general one and the scalar case as a special diagonal---a pattern we will see
repeatedly: handle the general object, specialize for free. The material
coefficient becomes a per-point weight $Q$ inside the bilinear form
$a(u,v)=(Q\,\diffop u,\diffop v)_\dom$ of \cref{ch:galerkin}, and whether $Q$
is a scalar or a $3\times3$ tensor changes nothing in the structure of the
assembly---only the arithmetic at each quadrature point.

\section{Interfaces and walls: boundary conditions}
\label{sec:bc}

A simulation lives in a bounded box, and the box is made of materials that
meet at interfaces and terminate at walls. What the fields do at those
surfaces is dictated by Maxwell's equations themselves, applied across the
discontinuity. Two conditions matter, and they are the continuous ancestors of
the \emph{essential} and \emph{natural} boundary conditions that the
finite-element method will later impose (\cref{ch:weak}, \cref{ch:bc}).

\begin{figure}[t]
\centering
\begin{tikzpicture}[scale=1.0]
  % two media separated by an interface
  \fill[simBgBlue]  (-3,0) rectangle (3,1.8);
  \fill[simBgGold]  (-3,-1.8) rectangle (3,0);
  \draw[very thick,simInk] (-3,0)--(3,0);
  \node[font=\footnotesize] at (-2.5,1.4) {medium 1};
  \node[font=\scriptsize,text=simGray] at (-2.5,1.0) {$\varepsilon_1,\mu_1$};
  \node[font=\footnotesize] at (-2.5,-1.4) {medium 2};
  \node[font=\scriptsize,text=simGray] at (-2.5,-1.0) {$\varepsilon_2,\mu_2$};
  % interface normal
  \draw[-{Stealth[length=2.2mm]},simRust,thick] (0,0)--(0,0.9)
     node[right,font=\scriptsize]{$\vn$};
  % tangential E continuous: same arrow both sides
  \draw[-{Stealth[length=2mm]},simBlue,thick] (1.2,0.0)--(2.2,0.0)
     node[above,font=\scriptsize]{$\vE_{\parallel}$};
  \draw[densely dashed,simBlue] (1.7,0.06)--(1.7,-0.5);
  \draw[-{Stealth[length=2mm]},simBlue,thick] (1.2,-0.0)--(2.2,-0.0);
  \node[font=\scriptsize,text=simBlue,align=center] at (1.7,-0.85)
     {$\vE_\parallel$ continuous};
  % normal B continuous: vertical arrow crossing
  \draw[-{Stealth[length=2.4mm]},simTeal,thick] (-1.7,-0.85)--(-1.7,0.85);
  \node[font=\scriptsize,text=simTeal] at (-1.15,0.55) {$\vB_\perp$};
  \node[font=\scriptsize,text=simTeal,align=center] at (-1.6,-1.15)
     {$\vB_\perp$ continuous};
\end{tikzpicture}
\caption[Interface boundary conditions]{At an interface between two media the
\emph{tangential} electric field $\vE_\parallel$ and the \emph{normal} magnetic
flux $\vB_\perp$ pass through unchanged. These two continuities are what the
finite-element spaces $\Hcurl$ (tangentially continuous, for $\vE$) and $\Hdiv$
(normally continuous, for $\vB$) are built to honour exactly. The complementary
components ($\vD_\perp$, $\vH_\parallel$) may jump by surface charge or current.}
\label{fig:interface}
\end{figure}

\subsection{Tangential \texorpdfstring{$\vE$}{E} and normal \texorpdfstring{$\vB$}{B} are continuous}

Apply Faraday's law to a thin loop straddling the interface, then shrink the
loop's height to zero. The flux through the vanishing loop goes to zero, and
what survives is the statement that the \emph{tangential component of $\vE$ is
continuous across the interface}:
\begin{equation}
  \vn\times(\vE_1-\vE_2) = \bm 0. \label{eq:bc-Etan}
\end{equation}
Symmetrically, apply Gauss's law for $\vB$ to a thin pillbox and shrink its
height: the \emph{normal component of $\vB$ is continuous},
\begin{equation}
  \vn\cdot(\vB_1-\vB_2) = 0. \label{eq:bc-Bnorm}
\end{equation}
The complementary components are \emph{not} continuous in general: the normal
$\vD$ can jump by a surface charge density, and the tangential $\vH$ can jump by
a surface current. But \eqref{eq:bc-Etan} and \eqref{eq:bc-Bnorm}---tangential
$\vE$, normal $\vB$---always hold. They are the reason the discrete spaces of
\cref{ch:functionspaces} are designed as they are: $\Hcurl$ enforces exactly
tangential continuity (the right home for $\vE$), and $\Hdiv$ enforces exactly
normal continuity (the right home for $\vB$). The physics chooses the function
space.

\subsection{PEC and PMC walls}

Two idealized walls bound nearly every simulation.

A \emph{perfect electric conductor} (PEC) is the $\sigma\to\infty$ limit of a
metal. Inside it the field is zero, so by tangential continuity
\eqref{eq:bc-Etan} the tangential electric field must vanish \emph{at the
surface}:
\begin{equation}
  \vn\times\vE = \bm 0 \quad\text{on a PEC wall.} \label{eq:pec}
\end{equation}
A PEC wall is a mirror that shorts out any tangential $\vE$; it is the
boundary condition at the metal plates of a capacitor and the walls of a
cavity. It will become an \emph{essential} (Dirichlet) condition: imposed
directly on the solution by removing the constrained degrees of freedom.

A \emph{perfect magnetic conductor} (PMC) is the dual idealization,
$\vn\times\vH=\bm0$. No real material is a PMC, but the condition arises
naturally as a \emph{symmetry plane}---a surface across which the problem is
mirror-symmetric---and it is the condition the finite-element method imposes by
doing \emph{nothing} (a \emph{natural} condition). The contrast is worth
holding onto: PEC must be enforced; PMC happens for free. That asymmetry,
traced to the integration by parts of \cref{ch:weak}, is why one of them costs
work to impose and the other costs nothing.

\begin{notationbox}
\textbf{Essential} (Dirichlet) conditions constrain the solution value and must
be \emph{built into} the function space---the simulator eliminates the
constrained degrees of freedom (\code{eliminate\_essential\_bc} in the
machinery of \cref{ch:bc}). \textbf{Natural} (Neumann-like) conditions
constrain a derivative and fall out of the weak form automatically; the
simulator imposes them by adding nothing. PEC $\to$ essential; PMC and
symmetry $\to$ natural. This single distinction reappears in every driver.
\end{notationbox}

\section{Energy and Poynting's theorem}
\label{sec:energy}

A simulation that finds the fields is only halfway done; the engineer wants a
\emph{number}---a capacitance, an inductance, a quality factor---and almost
every such number is an \emph{energy}. The electromagnetic field stores energy
in space at a density
\begin{equation}
  u = \tfrac12\,\vE\cdot\vD + \tfrac12\,\vH\cdot\vB
    = \tfrac12\,\varepsilon\,\abs{\vE}^2 + \tfrac{1}{2\mu}\,\abs{\vB}^2,
  \label{eq:energy-density}
\end{equation}
the first term electric, the second magnetic. Integrate $u$ over the device and
you have the total stored energy $U$---and from $U$ come the reductions:
capacitance is $C = 2U/V^2$ at unit voltage, inductance is $L = 2U/I^2$ at unit
current. These are the very quantities the \emph{energy reductions} of
\cref{ch:reductions} compute. Equation~\eqref{eq:energy-density} is their
integrand.

\subsection{Poynting's theorem: where the energy goes}

Energy density alone is a snapshot; we also want to know how energy
\emph{moves}. That is the content of Poynting's theorem, and it follows from
Maxwell's equations by one inspired manipulation. Define the \emph{Poynting
vector}
\begin{equation}
  \vS = \vE\times\vH, \label{eq:poynting-vec}
\end{equation}
the directional flow of electromagnetic power per unit area (\cref{fig:poynting}).
Now form the combination $\vH\cdot(\Curl\vE) - \vE\cdot(\Curl\vH)$. A standard
vector identity rewrites it as a divergence,
\[
  \vH\cdot(\Curl\vE) - \vE\cdot(\Curl\vH) = \Div(\vE\times\vH) = \Div\vS.
\]
Substitute Faraday's law~\eqref{eq:faraday} for $\Curl\vE$ and the
Amp\`ere--Maxwell law~\eqref{eq:ampere} for $\Curl\vH$:
\[
  \Div\vS
  = \vH\cdot(-\partial_t\vB) - \vE\cdot(\vJ + \partial_t\vD)
  = -\,\partial_t\!\left(\tfrac12\vH\cdot\vB + \tfrac12\vE\cdot\vD\right)
    - \vE\cdot\vJ.
\]
Recognizing the bracket as the energy density $u$ of
\eqref{eq:energy-density}, we arrive at \emph{Poynting's theorem}:
\begin{equation}
  \partial_t u + \Div\vS = -\,\vE\cdot\vJ. \label{eq:poynting}
\end{equation}
Read it as a balance sheet for energy at every point. The first term is the
rate the stored energy changes; the second is the net power flowing \emph{out}
(the divergence of the energy flux $\vS$); the right side is the power
delivered to the currents. When $\vJ=\sigma\vE$ (Ohmic material), the source
term is $-\sigma\abs{\vE}^2\le 0$: energy is \emph{dissipated}, drained
irreversibly into heat. That loss term is what gives a resonator its finite
$Q$.

\begin{figure}[t]
\centering
\begin{tikzpicture}[scale=1.0]
  % E field (blue, right), H field (green, up/out), S = E x H (rust, into-page-ish)
  \draw[-{Stealth[length=2.6mm]},simBlue,very thick] (0,0)--(2.0,0)
     node[right,font=\small]{$\vE$};
  \draw[-{Stealth[length=2.6mm]},simGreen,very thick] (0,0)--(0,2.0)
     node[above,font=\small]{$\vH$};
  \draw[-{Stealth[length=2.6mm]},simRust,very thick] (0,0)--(1.45,1.45)
     node[above right,font=\small]{$\vS=\vE\times\vH$};
  % little flux disk
  \fill[simRust!18] (0.0,0.0) -- (2.0,0.0) arc (0:90:2.0) -- cycle;
  \node[font=\scriptsize,text=simGray,align=center] at (3.6,1.0)
     {power flows\\ along $\vS$};
  % right angle ticks
  \draw[simInk] (0.18,0)--(0.18,0.18)--(0,0.18);
\end{tikzpicture}
\caption[The Poynting vector]{The Poynting vector $\vS=\vE\times\vH$ points
perpendicular to both fields and gives the direction and density of
electromagnetic power flow. In a propagating wave, $\vE$, $\vH$, and $\vS$ form
a right-handed triple and $\vS$ points the way the wave travels. Poynting's
theorem \eqref{eq:poynting} is the local accounting that $\partial_t u+\Div\vS$
equals the power delivered to charges.}
\label{fig:poynting}
\end{figure}

\begin{keyidea}
Poynting's theorem is the conservation of energy for the field, derived
\emph{from} Maxwell's equations rather than assumed alongside them---just as
continuity \eqref{eq:continuity} was conservation of charge. The pattern is
worth noticing: the deep facts of the theory (energy conservation, charge
conservation) are not extra axioms but algebraic consequences of the four laws,
extracted by a vector identity. The energy density \eqref{eq:energy-density} is
the integrand the output reductions sum; the loss term $\vE\cdot\vJ$ is what
makes a quality factor finite.
\end{keyidea}

\subsection{Stored energy, the quality factor, and participation}
\label{sec:Q}

The energy bookkeeping is not an afterthought; it is what most analyses
\emph{report}. Three quantities, all built from the energy density
\eqref{eq:energy-density}, recur throughout the book.

First, in a resonant mode the electric and magnetic energies are equal on
average and exchange completely twice per cycle---energy sloshes from the
electric field to the magnetic field and back, just as a pendulum trades
kinetic for potential energy. The total $U=U_E+U_B$ is what the eigenmode
analysis integrates.

Second, when a small loss is present, the mode does not ring forever; its
amplitude decays, and the \emph{quality factor}
\[
  Q = \omega\,\frac{U_{\text{stored}}}{P_{\text{loss}}}
    = \frac{\omega}{\kappa}
\]
measures how many radians of oscillation it takes to drain a factor $1/e$ of the
energy. A high-$Q$ cavity ($Q\sim10^6$ or more for a superconducting resonator)
rings for millions of cycles; a lossy one barely rings at all. The numerator is
the stored field energy of \eqref{eq:energy-density}; the denominator is the
power-loss integral of $\sigma\abs{\vE}^2$ from Poynting's
theorem~\eqref{eq:poynting}. This is the $Q$ the eigenmode reduction of
\cref{ch:reductions} computes, and the reason that chapter cares about the loss
term at all.

Third, engineers want to know \emph{where} a mode keeps its energy---in which
dielectric, in which gap. The \emph{participation ratio} of a region is the
fraction of the total energy stored there,
\[
  p_{\text{region}} = \frac{\int_{\text{region}} u\dV}{\int_{\text{all}} u\dV},
\]
a number between $0$ and $1$ that sums to one over a partition of the device. It
is the headline output of a great deal of quantum-device design: the
participation of a lossy interface sets the limit on a qubit's coherence. Like
$Q$ and like capacitance, it is just the energy density
\eqref{eq:energy-density} integrated and divided---which is why a single
\emph{energy reduction} (\cref{ch:reductions}) computes all three.

\section{Three regimes: static, time-harmonic, time-domain}
\label{sec:regimes}

The full Maxwell system is a coupled set of time-dependent partial differential
equations. Almost no engineering question needs it in that generality. Instead
each analysis lives in one of three \emph{regimes}, distinguished entirely by
\emph{how time enters}, and the choice of regime is what turns the one physics
into the several analyses of \cref{ch:targets}. \Cref{fig:regimes} is the map;
the rest of this section walks it.

\begin{figure}[t]
\centering
\begin{tikzpicture}[node distance=8mm]
  \node[stage, minimum width=30mm, minimum height=14mm] (full)
    {\textbf{Maxwell}\\[1pt]\footnotesize $\partial_t$ present};
  \node[stage, minimum width=30mm, below left=14mm and 8mm of full] (static)
    {\textbf{Static}\\[1pt]\footnotesize $\partial_t\to 0$};
  \node[stage, minimum width=30mm, below=22mm of full] (harm)
    {\textbf{Time-harmonic}\\[1pt]\footnotesize $\partial_t\to i\omega$};
  \node[stage, minimum width=30mm, below right=14mm and 8mm of full] (td)
    {\textbf{Time-domain}\\[1pt]\footnotesize $\partial_t$ kept};
  \draw[flow] (full) -- (static);
  \draw[flow] (full) -- (harm);
  \draw[flow] (full) -- (td);
  % products
  \node[product, below=7mm of static, minimum width=30mm, font=\scriptsize]
    {electro- / magnetostatic\\ $\to$ $C$, $L$};
  \node[product, below=7mm of harm, minimum width=34mm, font=\scriptsize]
    {eigenmode, driven\\ $\to$ $f,Q$, S-params};
  \node[product, below=7mm of td, minimum width=30mm, font=\scriptsize]
    {transient\\ $\to$ field history};
  % the key substitution callout
  \node[font=\scriptsize\itshape, text=simRust, right=10mm of harm, align=center]
    (sub) {the phasor\\ substitution};
  \draw[backflow] (sub) -- (harm);
\end{tikzpicture}
\caption[The three regimes]{The three regimes of Maxwell's equations, set by
how the time derivative is treated. Dropping $\partial_t$ gives the
\emph{static} problems (electrostatic, magnetostatic). Replacing
$\partial_t\to i\omega$ gives the \emph{time-harmonic} problems (eigenmode,
driven) and produces the complex curl--curl operators of
\cref{ch:reductions-pde}. Keeping $\partial_t$ gives the \emph{time-domain}
transient problem. The single substitution $\partial_t\to i\omega$ is the
bridge from the time-domain to the frequency-domain analyses.}
\label{fig:regimes}
\end{figure}

\subsection{Static: time stands still}

When nothing changes in time, every $\partial_t$ vanishes and the two coupling
laws \emph{decouple}. Faraday's law becomes $\Curl\vE=\bm0$, so $\vE$ is a
gradient (\cref{sec:potentials}); Amp\`ere's law loses its displacement current
and becomes $\Curl\vH=\vJ$. The electric and magnetic problems no longer talk
to each other, and we are left with two independent, self-contained problems:

\begin{itemize}
  \item \textbf{Electrostatics.} With $\vE=-\Grad V$ and $\vD=\varepsilon\vE$,
  Gauss's law $\Div\vD=\rho$ becomes a single scalar second-order equation for
  the potential,
  \begin{equation}
    -\Div(\varepsilon\,\Grad V) = \rho, \label{eq:poisson}
  \end{equation}
  the \emph{Poisson equation} (Laplace's equation $-\Div(\varepsilon\Grad
  V)=0$ in charge-free regions). This is the electrostatic analysis, and it is
  the simplest governing equation in the book.

  \item \textbf{Magnetostatics.} With $\vB=\Curl\vA$ and $\vH=\nu\vB$,
  Amp\`ere's law $\Curl\vH=\vJ$ becomes the \emph{curl--curl} equation for the
  vector potential,
  \begin{equation}
    \Curl(\nu\,\Curl\vA) = \vJ. \label{eq:curlcurl-static}
  \end{equation}
\end{itemize}

These are the equations behind the electrostatic and magnetostatic analyses of
\cref{ch:targets}, and \cref{ch:reductions-pde} derives them in full. Notice
already how alike they are in shape---a weighted second-order operator equal to
a source---differing only in the function space ($V$ scalar, $\vA$ vector) and
the differential operator ($\Grad$ versus $\Curl$). That kinship is the seed of
``the same machine, one part swapped.''

\subsection{Time-harmonic: the phasor substitution}

Most microwave engineering asks not ``what happens in time?'' but ``what is the
steady response at frequency $\omega$?'' If every field oscillates at a single
angular frequency $\omega$, we can write each as the real part of a complex
\emph{phasor} times $e^{i\omega t}$:
\[
  \vE(\vx,t) = \Re\!\big\{\hat{\vE}(\vx)\,e^{i\omega t}\big\},
\]
and similarly for $\vB,\vD,\vH,\vJ$. The complex amplitude $\hat{\vE}(\vx)$
carries both magnitude and phase at every point; time has been factored out.
The payoff is a single, beautiful simplification. Differentiating
$e^{i\omega t}$ in time just multiplies by $i\omega$:
\begin{equation}
  \partial_t \;\longrightarrow\; i\omega. \label{eq:phasor-sub}
\end{equation}
Every time derivative in Maxwell's equations becomes an algebraic factor
$i\omega$, and the partial differential equations \emph{in space and time}
collapse to partial differential equations \emph{in space alone}, with complex
coefficients. The phasor curl laws read
\begin{align}
  \Curl\hat{\vE} &= -\,i\omega\,\hat{\vB}, \label{eq:faraday-phasor}\\
  \Curl\hat{\vH} &= \hat{\vJ} + i\omega\,\hat{\vD}. \label{eq:ampere-phasor}
\end{align}
Worked example~\ref{wex:phasor} carries this substitution all the way to the curl--curl
operator that the driven and eigenmode analyses solve. The substitution
\eqref{eq:phasor-sub} is the single most important bridge in Part~II: it is
what turns the time-domain Maxwell system into the frequency-domain operators
$\mat K + i\omega\mat C - \omega^2\mat M$ that \cref{ch:reductions-pde} derives
and that \cref{part:machinery} learns to solve.

It is worth seeing in advance where those three operators come from, because the
pattern recurs in every harmonic analysis. After the weak form and
discretization of \cref{ch:weak,ch:galerkin}, each \emph{term} of the harmonic
curl--curl equation contributes one matrix:
\begin{itemize}
  \item the curl--curl term $\Curl(\nu\,\Curl\hat{\vE})$ becomes the
  \emph{stiffness} $\mat K$ (no $\omega$ in it---it is the static operator);
  \item a conduction/loss term $i\omega\sigma\hat{\vE}$ (present when the medium
  conducts, or a port radiates) becomes $i\omega\mat C$, the \emph{damping}
  operator---first power of $\omega$, and the only one carrying the imaginary
  unit;
  \item the displacement term $-\omega^2\varepsilon\hat{\vE}$ becomes
  $-\omega^2\mat M$, with the \emph{mass} $\mat M$---second power of $\omega$,
  the term that makes the problem a resonance.
\end{itemize}
The full operator $\mat A(\omega)=\mat K + i\omega\mat C - \omega^2\mat M$ is a
\emph{quadratic} polynomial in $\omega$ whose three coefficient matrices are
assembled once and reused at every frequency. Drop the drive and you have the
eigenmode problem $\mat K\hat{\vE}=\omega^2\mat M\hat{\vE}$ (a generalized
eigenproblem when $\mat C=0$, a quadratic one when it does not); keep the drive
and you sweep $\omega$, re-forming $\mat A(\omega)$ and solving
$\mat A(\omega)\hat{\vE}=\vb(\omega)$ at each step. The single phasor
substitution has produced the entire frequency-domain operator algebra the
later chapters live in.

\subsection{Time-domain: keep the clock running}

Sometimes the excitation is a broadband pulse, or the materials are nonlinear,
or one simply wants the movie. Then no single frequency suffices and the time
derivatives must be \emph{kept}. After spatial discretization the result is a
second-order system of ordinary differential equations in time,
\[
  \mat M\,\ddot{\vs}(t) + \mat C\,\dot{\vs}(t) + \mat K\,\vs(t) = \vJ(t),
\]
the same $\mat M,\mat C,\mat K$ operators as the harmonic case but now marched
forward step by step rather than solved frequency by frequency. This is the
transient analysis. Its defining feature---each step begins from the state the
previous step produced---makes it the book's first genuinely \emph{sequential}
computation, and the motivating example for the fold of \cref{ch:fold}.

\begin{keyidea}
One physics, three regimes, set by the fate of $\partial_t$. \textbf{Drop} it
($\partial_t\to0$): static, real operators, two decoupled problems (electro-
and magnetostatic). \textbf{Algebraize} it ($\partial_t\to i\omega$):
time-harmonic, complex operators, the curl--curl/Helmholtz family (eigenmode,
driven). \textbf{Keep} it: time-domain, a marched ODE system (transient). The
substitution $\partial_t\to i\omega$ is the hinge between the second and third,
and the source of every complex operator in the book.
\end{keyidea}

\begin{table}[t]
\centering
\caption[The three regimes, compared]{The three regimes side by side. The
column that changes everything is the fate of the time derivative; it dictates
whether the operator is real or complex, whether the problem couples $\vE$ and
$\vB$, and which analyses of \cref{ch:targets} it produces. The discrete
operators $\mat K,\mat C,\mat M$ are the same matrices in every column---only
their combination differs.}
\label{tab:regimes}
\small
\setlength{\tabcolsep}{5pt}
\renewcommand{\arraystretch}{1.3}
\begin{tabular}{@{}llll@{}}
\toprule
 & Static & Time-harmonic & Time-domain \\
\midrule
$\partial_t$ becomes & $0$ & $i\omega$ (algebraic) & kept (a derivative) \\
Unknown & real field & complex phasor & real field history \\
Coupling & decoupled & coupled & coupled \\
Discrete operator & $\mat K$ & $\mat K+i\omega\mat C-\omega^2\mat M$
  & $\mat M\ddot{\vs}+\mat C\dot{\vs}+\mat K\vs$ \\
Solve shape & one linear solve & solve per $\omega$ / eigenproblem & marched fold \\
Analyses & electro-/magnetostatic & eigenmode, driven & transient \\
\bottomrule
\end{tabular}
\end{table}

\section{Potentials: solving the constraints automatically}
\label{sec:potentials}

The static analyses were posed not for the fields $\vE,\vB$ directly but for
\emph{potentials}---a scalar $V$ and a vector $\vA$. The reason is elegant: a
potential \emph{automatically satisfies} one of the Gauss laws, removing it
from the list of equations to solve. Potentials trade a constrained field for an
unconstrained one.

\subsection{The scalar potential \texorpdfstring{$V$}{V}}

In the static (or curl-free) case Faraday's law gives $\Curl\vE=\bm0$. A vector
identity from vector calculus states that the curl of \emph{any} gradient
vanishes, $\Curl(\Grad f)=\bm0$. So if we write
\begin{equation}
  \vE = -\Grad V, \label{eq:E-from-V}
\end{equation}
then $\Curl\vE = -\Curl(\Grad V)=\bm 0$ is satisfied \emph{identically}, for
\emph{any} scalar function $V$ (\cref{fig:potentials}, left). The curl law is no
longer an equation to enforce; it is built into the representation. We are left
with only Gauss's law, which \eqref{eq:E-from-V} turns into the Poisson
equation~\eqref{eq:poisson}. The minus sign is convention: it makes $\vE$ point
``downhill,'' from high potential to low, matching the picture of a positive
charge rolling down a voltage slope.

\subsection{The vector potential \texorpdfstring{$\vA$}{A}}

The magnetic side runs in parallel. Gauss's law for $\vB$ says $\Div\vB=0$.
The companion identity is that the divergence of \emph{any} curl vanishes,
$\Div(\Curl\vF)=0$ (the same identity that gave us continuity in
\eqref{eq:continuity}). So if we write
\begin{equation}
  \vB = \Curl\vA, \label{eq:B-from-A}
\end{equation}
then $\Div\vB=\Div(\Curl\vA)=0$ holds \emph{identically} for any vector field
$\vA$ (\cref{fig:potentials}, right). Gauss's law for $\vB$ disappears, and
Amp\`ere's law becomes the curl--curl equation
\eqref{eq:curlcurl-static}. The two potentials are the two halves of the same
trick: each rides on a vanishing-derivative identity ($\Curl\Grad=0$ and
$\Div\Curl=0$) to dissolve one Gauss law and leave one equation behind. Those
two identities are exactly the two composites of the de Rham complex
$\Hone\xrightarrow{\Grad}\Hcurl\xrightarrow{\Curl}\Hdiv\xrightarrow{\Div}\Ltwo$
that \cref{ch:derham} builds the discrete function spaces to honour.

\begin{figure}[t]
\centering
\begin{tikzpicture}[scale=1.0]
% left: scalar potential, E = -grad V over a hill
\begin{scope}[shift={(0,0)}]
  % contour "hill"
  \foreach \r/\op in {1.05/12, 0.75/22, 0.45/36}{
    \draw[simGray!60] (0,0) circle (\r);}
  \node[font=\scriptsize] at (0,0) {high $V$};
  % E arrows pointing downhill (outward, but E=-grad V so toward low V outside)
  \foreach \a in {30,120,210,300}{
    \draw[-{Stealth[length=2mm]},simBlue,thick]
      (\a:0.5)--(\a:1.25);}
  \node[font=\scriptsize\itshape,simBlue] at (1.35,1.0) {$\vE=-\Grad V$};
  \node[font=\footnotesize] at (0,-1.7) {scalar potential};
\end{scope}
% right: vector potential, B = curl A
\begin{scope}[shift={(5.0,0)}]
  % A as a circulation
  \draw[flow,simViolet] (0,0) circle (0.8);
  \node[font=\scriptsize\itshape,simViolet] at (-1.05,0.6) {$\vA$};
  % B coming out of the loop
  \draw[-{Stealth[length=2.4mm]},simTeal,very thick] (0,0)--(0,0.0);
  \fill[simTeal] (0,0) circle (1.6pt);
  \node[font=\scriptsize\itshape,simTeal] at (0.45,0.25) {$\vB$};
  \node[font=\footnotesize] at (0,-1.7) {vector potential};
  \node[font=\scriptsize,text=simGray] at (0,-1.4) {$\vB=\Curl\vA$};
\end{scope}
\end{tikzpicture}
\caption[The potential representations]{The two potential representations.
\emph{Left}: writing $\vE=-\Grad V$ makes $\vE$ point downhill on the potential
landscape and satisfies $\Curl\vE=\bm0$ identically. \emph{Right}: writing
$\vB=\Curl\vA$ makes $\vB$ the circulation of $\vA$ and satisfies $\Div\vB=0$
identically. Each potential dissolves one Gauss law for free, riding on a
vanishing-derivative identity of vector calculus.}
\label{fig:potentials}
\end{figure}

\subsection{Gauge freedom, briefly}

There is a price, gentle but real. The potentials are not unique. Adding any
constant to $V$ leaves $\vE=-\Grad V$ unchanged. Adding any gradient $\Grad
\chi$ to $\vA$ leaves $\vB=\Curl\vA$ unchanged, because $\Curl(\Grad\chi)=\bm0$.
This freedom---the \emph{gauge freedom}---means the magnetostatic curl--curl
operator \eqref{eq:curlcurl-static} has a nontrivial null space: every gradient
field is invisible to it. One fixes the freedom by an extra condition (a
\emph{gauge}, such as $\Div\vA=0$), or, as the simulator does, by projecting
the solution onto the gradient-free part with a divergence-free projector
(\cref{ch:bc}). For now simply note that the null space is not a bug to be
feared but the direct fingerprint of gauge freedom---the same vanishing-curl
identity that made the potential useful, seen from the other side.

\section{Worked examples}

The examples below are deliberately small enough to check with a pencil. The
first exercises Gauss's law and the energy integral end to end---it is the
electrostatic analysis in miniature. The second carries the phasor substitution
to the operator the time-harmonic analyses actually solve. The third puts the
interface conditions of \cref{sec:bc} to work on a concrete refraction of field
lines.

\begin{workedexample}{Parallel-plate capacitor: field, charge, energy}{capacitor}
Two parallel conducting plates of area $A$ sit a distance $d$ apart, the gap
filled with a dielectric of permittivity $\varepsilon=\varepsilon_r\varepsilon_0$.
The top plate is held at potential $V_0$, the bottom at $0$ (\cref{fig:capacitor}).
Find the field, the charge, the capacitance, and the stored energy---ignoring
fringing, so the field is uniform and vertical between the plates.

\emph{Field from the potential.} With the plates close compared to their size,
the potential varies linearly across the gap: $V(z)=V_0\,z/d$ for $0\le z\le
d$. Then
\[
  \vE = -\Grad V = -\frac{dV}{dz}\,\hat{\vz} = -\frac{V_0}{d}\,\hat{\vz},
\]
a uniform field of magnitude $E=V_0/d$ pointing from the high plate to the low.
With $V_0=\SI{1}{\volt}$ and $d=\SI{1}{\micro\meter}$ this is
$E=\SI{e6}{\volt\per\meter}$---fields inside small devices are large.

\emph{Charge from Gauss's law.} The displacement is $\vD=\varepsilon\vE$, of
magnitude $D=\varepsilon V_0/d$. A pillbox Gauss surface straddling the top
plate sees flux $D$ leave through its face of area $A$, so the enclosed surface
charge is
\[
  Q = D\,A = \frac{\varepsilon A}{d}\,V_0.
\]

\emph{Capacitance.} By definition $C=Q/V_0$, so
\begin{equation}
  C = \frac{\varepsilon A}{d}. \label{eq:cap}
\end{equation}
The classic result: capacitance grows with plate area and dielectric constant,
and shrinks with separation.

\emph{Energy, two ways.} The stored energy from the energy
density~\eqref{eq:energy-density}, integrated over the gap volume $V_{\text{gap}}=Ad$,
is
\[
  U = \int \tfrac12\varepsilon\abs{\vE}^2\dV
    = \tfrac12\varepsilon\Big(\frac{V_0}{d}\Big)^{\!2}\,(Ad)
    = \tfrac12\,\frac{\varepsilon A}{d}\,V_0^2
    = \tfrac12\,C\,V_0^2,
\]
recovering the familiar $U=\tfrac12 C V_0^2$. Read backwards, this says
$C=2U/V_0^2$: \emph{the capacitance is twice the field energy at unit
voltage}---exactly the energy reduction the electrostatic analysis performs in
\cref{ch:reductions}. A concrete number: $\varepsilon_r=3.9$ (silicon dioxide),
$A=\SI{100}{\micro\meter\squared}=\SI{e-10}{\meter\squared}$,
$d=\SI{100}{\nano\meter}$ gives
$C=3.9\times\SI{8.854e-12}{}\times\SI{e-10}{}/\SI{e-7}{}\approx\SI{3.5e-14}{\farad}=\SI{35}{\femto\farad}$,
a typical on-chip value.
\end{workedexample}

\begin{figure}[t]
\centering
\begin{physicsbox}
\physicsframe{
  \fill[simBgBlue] (-1.6,-0.9) rectangle (1.6,0.9);
  \draw[very thick,simInk] (-1.6,0.9)--(1.6,0.9);
  \draw[very thick,simInk] (-1.6,-0.9)--(1.6,-0.9);
  \node[font=\scriptsize] at (-2.0,0.9) {$V_0$};
  \node[font=\scriptsize] at (-2.0,-0.9) {$0$};
  \node[font=\scriptsize,text=simGray] at (0,1.25) {area $A$};
  \draw[{Stealth[length=1.5mm]}-{Stealth[length=1.5mm]},simGray]
    (1.95,0.9)--(1.95,-0.9);
  \node[font=\scriptsize,text=simGray] at (2.25,0) {$d$};
  \foreach \x in {-1.2,-0.6,0,0.6,1.2}{
    \draw[-{Stealth[length=1.8mm]},simBlue,thick] (\x,0.82)--(\x,-0.82);}
  \node[font=\scriptsize\itshape,simBlue] at (1.0,0.45) {$\vE$};
  \node[font=\scriptsize,text=simGreen] at (0,0) {$\varepsilon$};
}
\end{physicsbox}
\caption[Parallel-plate capacitor]{The parallel-plate capacitor of
Worked example~\ref{wex:capacitor}. A uniform vertical field $\vE=-(V_0/d)\hat{\vz}$ fills the
dielectric-filled gap; Gauss's law gives the surface charge, and the field
energy gives the capacitance $C=\varepsilon A/d = 2U/V_0^2$.}
\label{fig:capacitor}
\end{figure}

\begin{workedexample}{Phasor reduction to the curl--curl operator}{phasor}
Take a source-free region ($\vJ=\bm0$, $\rho=0$) of a linear medium
($\vD=\varepsilon\vE$, $\vB=\mu\vH$). Show that the time-harmonic Maxwell
equations reduce to a single second-order operator equation for $\hat{\vE}$,
the \emph{curl--curl} (vector Helmholtz) equation that the driven and eigenmode
analyses solve.

\emph{Start from the phasor curl laws.} With $\vJ=\bm0$, equations
\eqref{eq:faraday-phasor}--\eqref{eq:ampere-phasor} read
\[
  \Curl\hat{\vE} = -\,i\omega\,\hat{\vB} = -\,i\omega\mu\,\hat{\vH},
  \qquad
  \Curl\hat{\vH} = i\omega\,\hat{\vD} = i\omega\varepsilon\,\hat{\vE}.
\]

\emph{Eliminate $\hat{\vH}$.} Solve the first for $\hat{\vH}=
-\,\dfrac{1}{i\omega\mu}\Curl\hat{\vE}$ and substitute into the second:
\[
  \Curl\!\left(-\frac{1}{i\omega\mu}\Curl\hat{\vE}\right)
    = i\omega\varepsilon\,\hat{\vE}.
\]
Multiply through by $i\omega$ and write $\nu=1/\mu$. The factor $i\omega$ cancels
on the left to leave $-\Curl(\nu\Curl\hat{\vE})$ on one side, and $i\omega\cdot
i\omega = -\omega^2$ on the other:
\begin{equation}
  \Curl(\nu\,\Curl\hat{\vE}) - \omega^2\varepsilon\,\hat{\vE} = \bm0.
  \label{eq:curlcurl-harmonic}
\end{equation}

\emph{Read the result.} Equation~\eqref{eq:curlcurl-harmonic} is the
\emph{curl--curl} operator equation. Its two terms are exactly the discrete
operators of \cref{ch:targets}: the curl--curl term becomes the stiffness
$\mat K$ (weighted by reluctivity $\nu$), and the $\omega^2$ term becomes the
mass $\mat M$ (weighted by permittivity $\varepsilon$). With no drive it is the
generalized eigenvalue problem $\mat K\hat{\vE}=\omega^2\mat M\hat{\vE}$---the
\emph{eigenmode} analysis, whose eigenvalues are the squared resonant
frequencies. Restore a current source $\hat{\vJ}$ and add a loss term
$i\omega\sigma$ (from $\vJ=\sigma\vE$) and the operator becomes
$\mat K+i\omega\mat C-\omega^2\mat M$---the \emph{driven} analysis. Both
analyses are this one phasor reduction, with and without a right-hand side.
This is the single calculation that turns the four time-domain laws into the
frequency-domain machinery the rest of the book solves.
\end{workedexample}

\begin{workedexample}{Refraction of electric field lines at an interface}{refract}
A static electric field crosses the planar interface between two dielectrics,
$\varepsilon_1$ above and $\varepsilon_2$ below, with no free surface charge.
On the upper side the field makes an angle $\theta_1$ with the normal. Find the
angle $\theta_2$ on the lower side, and show that field lines \emph{bend} at the
interface like light through glass.

\emph{Set up the two conditions.} Resolve each field into a normal component
(along $\vn$) and a tangential component (in the surface). The interface
conditions of \cref{sec:bc} give two equations. Tangential $\vE$ is continuous
\eqref{eq:bc-Etan}:
\[
  E_1\sin\theta_1 = E_2\sin\theta_2.
\]
With no surface charge, normal $\vD=\varepsilon\vE$ is continuous:
\[
  \varepsilon_1 E_1\cos\theta_1 = \varepsilon_2 E_2\cos\theta_2.
\]

\emph{Divide to eliminate the magnitudes.} Dividing the first equation by the
second cancels $E_1$ and $E_2$ and leaves a clean relation between the angles:
\begin{equation}
  \frac{\tan\theta_1}{\varepsilon_1} = \frac{\tan\theta_2}{\varepsilon_2},
  \qquad\text{i.e.}\qquad
  \tan\theta_2 = \frac{\varepsilon_2}{\varepsilon_1}\,\tan\theta_1.
  \label{eq:refract}
\end{equation}
This is the electrostatic ``law of refraction.''

\emph{Put numbers in.} Take air above ($\varepsilon_1=\varepsilon_0$) and a
dielectric below with $\varepsilon_2=4\varepsilon_0$, and let the field enter at
$\theta_1=\SI{30}{\degree}$. Then $\tan\theta_2 = 4\tan\SI{30}{\degree} =
4\times 0.577 = 2.31$, so $\theta_2 = \arctan(2.31)\approx\SI{66.6}{\degree}$.
The field line bends \emph{away} from the normal on entering the denser
dielectric---the opposite sense from a light ray, because here it is the normal
\emph{flux} $\vD$, not the field $\vE$, that is conserved. The lesson for the
simulator is concrete: the two interface conditions of \cref{sec:bc} are not
decorative; they are algebraic constraints the discrete field must satisfy at
every material boundary, and they are exactly what the $\Hcurl$ and $\Hdiv$
function spaces of \cref{ch:functionspaces} are constructed to enforce.
\end{workedexample}

\begin{pitfall}
The cancellation in Worked example~\ref{wex:phasor} relied on $\omega\neq0$. At $\omega=0$ the
phasor substitution \eqref{eq:phasor-sub} sends $\partial_t\to0$ and we are back
in the static regime, where the mass term vanishes and the curl--curl equation
degenerates to the magnetostatic $\Curl(\nu\Curl\vA)=\vJ$. The three regimes are
not separate physics; the static case is the $\omega\to0$ limit of the harmonic
one. Watch the limit: the $\omega^2\mat M$ term that defines a resonance simply
switches off.
\end{pitfall}

\section{A third look: the plane wave}
\label{sec:planewave}

It is worth seeing the solution that the source-free curl--curl equation
\eqref{eq:curlcurl-harmonic} hides: the \emph{plane wave}, the simplest
nontrivial field in all of electromagnetism and the one that makes ``Maxwell's
equations are light'' concrete. In a uniform medium try
$\hat{\vE}=\vE_0\,e^{-i\vk\cdot\vx}$, a field of fixed amplitude $\vE_0$ and a
single propagation direction $\vk$ (\cref{fig:planewave}). The curl of such a
field turns into $-i\vk\times$, and \eqref{eq:curlcurl-harmonic} becomes the
algebraic \emph{dispersion relation}
\[
  \abs{\vk}^2 = \omega^2\mu\varepsilon,
  \qquad\text{i.e.}\qquad
  \abs{\vk} = \frac{\omega}{c}, \quad c=\frac{1}{\sqrt{\mu\varepsilon}},
\]
the statement that the wave travels at speed $c=1/\sqrt{\mu\varepsilon}$. In
vacuum $c=1/\sqrt{\mu_0\varepsilon_0}\approx\SI{3e8}{\meter\per\second}$---the
speed of light, recovered from $\mu_0$ and $\varepsilon_0$ alone, with no light
ever mentioned. Faraday's law \eqref{eq:faraday-phasor} then forces
$\vk\cdot\vE_0=0$ (the field is transverse to the direction of travel) and ties
$\vH_0$ perpendicular to both, so $\vE$, $\vH$, and the Poynting vector
$\vS=\vE\times\vH$ form the right-handed triple of \cref{fig:poynting}, with
$\vS$ pointing along $\vk$.

\begin{figure}[t]
\centering
\begin{tikzpicture}[scale=1.0]
  % propagation axis
  \draw[->,simGray] (-0.3,0)--(6.4,0) node[right,font=\scriptsize]{$\vk$};
  % E field: blue sinusoid in the vertical plane
  \draw[simBlue,thick,domain=0:5.9,samples=100,smooth]
    plot (\x,{0.85*sin(180*\x)});
  \node[font=\scriptsize\itshape,simBlue] at (0.55,1.05) {$\vE$};
  % a few E vectors riding the crest
  \foreach \x in {0.5,1.5,2.5,3.5,4.5,5.5}{
    \draw[-{Stealth[length=1.3mm]},simBlue,semithick]
      (\x,0)--(\x,{0.85*sin(180*\x)});}
  % H field: smaller transverse green sinusoid in phase
  \draw[simGreen,thick,domain=0:5.9,samples=100,smooth]
    plot (\x,{0.40*sin(180*\x)});
  \node[font=\scriptsize\itshape,simGreen] at (1.55,-0.62) {$\vH$};
  \node[font=\scriptsize,text=simRust] at (5.0,0.95) {$\vS\parallel\vk$};
\end{tikzpicture}
\caption[A plane wave]{A plane electromagnetic wave. The electric field $\vE$
(blue) and magnetic field $\vH$ (green, drawn here as the smaller transverse
component) oscillate in phase, perpendicular to each other and to the
propagation direction $\vk$. The wave moves at speed $c=1/\sqrt{\mu\varepsilon}$
set by the dispersion relation $\abs{\vk}=\omega/c$, carrying power along
$\vS=\vE\times\vH$. This solution falls out of the same source-free curl--curl
operator \eqref{eq:curlcurl-harmonic} the eigenmode and driven analyses
assemble.}
\label{fig:planewave}
\end{figure}

\subsection{Lossy media and the skin depth}

The dispersion relation already carries the materials inside it, and turning on
conductivity makes the wave \emph{decay}. In a conductor the
Amp\`ere--Maxwell law gains the conduction current $\vJ=\sigma\vE$, and the
phasor curl law becomes $\Curl\hat{\vH}=(\sigma+i\omega\varepsilon)\hat{\vE}$.
Carrying this through the elimination of Worked example~\ref{wex:phasor} replaces
$\omega^2\mu\varepsilon$ by $\omega^2\mu\varepsilon - i\omega\mu\sigma$, so the
wavenumber $\abs{\vk}$ acquires an imaginary part. A complex $\vk=\beta - i/\delta$
means the field $e^{-i\vk\cdot\vx}$ carries a real exponential decay
$e^{-x/\delta}$: the wave dies away as it penetrates, over a characteristic
length---the \emph{skin depth}
\[
  \delta = \sqrt{\frac{2}{\omega\mu\sigma}}\qquad(\text{good conductor}).
\]
At microwave frequencies $\delta$ in copper is well under a micron, which is why
fields do not penetrate metal and why the PEC idealization
$\vn\times\vE=\bm0$ of \cref{sec:bc} is so accurate: a real conductor confines
the field to a vanishingly thin surface layer. The same imaginary part is the
microscopic origin of the loss term $-\sigma\abs{\vE}^2$ in Poynting's theorem
\eqref{eq:poynting}, and of the finite quality factor $Q$ a resonant cavity
exhibits.

Light is Maxwell's equations solving themselves in empty space; a damped wave
in a conductor is the same equations solving themselves in a lossy one. That
both fall out of the single curl--curl operator the simulator assembles---the
operator Worked example~\ref{wex:phasor} derived---is the unity the whole book is about.

\summarysection
Maxwell's four equations split into two \emph{Gauss} constraints
($\Div\vD=\rho$, $\Div\vB=0$) and two \emph{curl} dynamics (Faraday
$\Curl\vE=-\partial_t\vB$, Amp\`ere--Maxwell $\Curl\vH=\vJ+\partial_t\vD$); the
displacement current closes the loop that becomes light, and forces the
\emph{continuity} of charge $\Div\vJ+\partial_t\rho=0$ as a consequence.
Materials enter only through the three constitutive coefficients
($\vD=\varepsilon\vE$, $\vB=\mu\vH$, $\vJ=\sigma\vE$), scalar for isotropic
media and tensor for anisotropic ones. At interfaces, tangential $\vE$ and
normal $\vB$ are continuous; PEC walls force $\vn\times\vE=\bm0$ (an essential
condition), PMC and symmetry walls impose $\vn\times\vH=\bm0$ for free (a
natural condition). The field stores energy at density
$\tfrac12\varepsilon\abs{\vE}^2+\tfrac1{2\mu}\abs{\vB}^2$, and Poynting's
theorem $\partial_t u+\Div\vS=-\vE\cdot\vJ$ accounts for its flow and loss---the
integrand of the output reductions. The physics divides into three regimes by
the fate of $\partial_t$: dropping it gives the static Poisson and curl--curl
problems posed via the potentials $\vE=-\Grad V$ and $\vB=\Curl\vA$;
substituting $\partial_t\to i\omega$ gives the time-harmonic curl--curl operator
$\mat K+i\omega\mat C-\omega^2\mat M$ behind the eigenmode and driven analyses;
keeping it gives the marched transient system. Every operator the rest of
Part~II assembles is one of these reductions.

\furtherreading
For the physics at a sophomore's level, Griffiths~\citep{griffiths2017} is the
standard and the clearest---its treatment of the four laws, the constitutive
relations, boundary conditions, and Poynting's theorem is exactly the
background assumed here. For the engineering-oriented frequency-domain picture
(phasors, the curl--curl/Helmholtz equation, waveguides) and the finite-element
formulation that \cref{ch:weak} begins, Jin~\citep{jin2014} is the canonical
reference. The mathematically careful treatment of the function spaces
$\Hcurl,\Hdiv$ and the boundary conditions that motivate them---material we
will need in earnest in \cref{ch:functionspaces}---is developed rigorously by
Monk~\citep{monk2003}.

\exercisessection
\begin{enumerate}[nosep]
  \item \textbf{Reading the laws.} For each of the four Maxwell equations, state
  in one sentence what physical quantity is the ``source'' on the right-hand
  side and what field property (divergence or curl) it controls. Which two
  equations couple $\vE$ and $\vB$, and which two are instantaneous
  constraints?

  \item \textbf{Continuity by hand.} Starting from Faraday's law
  $\Curl\vE=-\partial_t\vB$, take the divergence of both sides and use
  $\Div(\Curl\,\cdot)=0$ to show that $\partial_t(\Div\vB)=0$---so once
  $\Div\vB=0$ holds it holds forever. Explain why this is the magnetic analogue
  of the continuity equation~\eqref{eq:continuity}.

  \item \textbf{A different dielectric.} Repeat the capacitance calculation of
  Worked example~\ref{wex:capacitor} for two plates of area $A=\SI{1}{\milli\meter\squared}$
  separated by $d=\SI{10}{\micro\meter}$ of a material with $\varepsilon_r=10$.
  Find $C$, then the energy stored at $V_0=\SI{5}{\volt}$, and verify
  $C=2U/V_0^2$.

  \item \textbf{The displacement current matters.} Suppose Maxwell had written
  Amp\`ere's law without the $\partial_t\vD$ term. Take the divergence and show
  that the resulting equation forces $\Div\vJ=0$ everywhere. Give a physical
  situation (e.g.\ a charging capacitor) where $\Div\vJ\neq0$, and explain why
  the missing term is therefore essential.

  \item \textbf{Refraction, reversed.} Worked example~\ref{wex:refract} sent a
  field from air into a denser dielectric. Reverse it: a field inside the
  dielectric ($\varepsilon_2=4\varepsilon_0$) makes $\SI{66.6}{\degree}$ with the
  normal and exits into air ($\varepsilon_1=\varepsilon_0$). Use the refraction
  law~\eqref{eq:refract} to find the exit angle, and confirm it is the
  $\SI{30}{\degree}$ you started with---refraction is reversible. Then explain,
  in one sentence, why a field line crossing into a \emph{higher}-permittivity
  region always bends away from the normal.

  \item \textbf{The phasor factor.} Show explicitly that if
  $f(t)=\Re\{\hat f\,e^{i\omega t}\}$ then $\partial_t f=\Re\{i\omega\hat
  f\,e^{i\omega t}\}$, justifying the substitution
  $\partial_t\to i\omega$~\eqref{eq:phasor-sub}. What does the factor $i$ do to
  the \emph{phase} of the oscillation, and why does that make ``a derivative''
  into ``a $\SI{90}{\degree}$ phase shift''?

  \item \textbf{From harmonic to static.} Take the harmonic curl--curl equation
  \eqref{eq:curlcurl-harmonic} and let $\omega\to0$. Which term vanishes? Show
  that with a current source restored you recover the magnetostatic equation
  \eqref{eq:curlcurl-static}, confirming that the static regime is the
  zero-frequency limit of the harmonic one.

  \item \textbf{Energy in a resonant mode.} In a lossless cavity a mode
  oscillates with the electric and magnetic energies exchanging back and forth,
  $U_E(t)+U_B(t)=U_{\text{total}}=\text{const}$. Using Poynting's theorem
  \eqref{eq:poynting} with $\vJ=\bm0$ and no net outflow through the PEC walls,
  argue that the total stored energy is conserved. Then add a small
  conductivity $\sigma$ and explain, from the loss term $-\sigma\abs{\vE}^2$,
  why the mode now decays---and why that decay rate is what the quality factor
  $Q$ measures.
\end{enumerate}

\chapter{The Five Reductions of Maxwell}
\label{ch:reductions-pde}

% local: velocity unknown in the first-order transient recast (no \vv in simbook)
\providecommand{\vv}{\mathbf{v}}

\chapterorient{Maxwell's four equations describe every electromagnetic
phenomenon at once---which is precisely why, taken whole, they are too rich to
compute against directly. Every practical simulation first \emph{specializes}
them: it fixes a regime (static, time-harmonic, or time-domain), discards the
terms that vanish there, and lands on a single, tractable partial differential
equation. This chapter performs that specialization five times, once for each
of the analyses of \cref{ch:targets}. We will watch the same four equations
collapse into a Poisson problem, a curl--curl problem, an eigenproblem, a
complex-valued Helmholtz system, and a second-order evolution---and we will end
by noticing that, beneath their different faces, all five are assembled from
the \emph{same three operators}.}

\section{From one theory to five problems}

In \cref{ch:maxwell} we reviewed Maxwell's equations in their full, coupled,
time-dependent glory. They are a beautiful and complete description of the
electromagnetic field, and for exactly that reason they are not yet a
\emph{computation}. A computation needs a finite, well-posed question: given
\emph{this} geometry, \emph{these} materials, and \emph{this} excitation, solve
for \emph{that} field. The art of turning physics into simulation is the art of
\emph{reduction}---of asking a question narrow enough that Maxwell's equations,
restricted to it, become a single solvable problem.

The six analyses of \cref{ch:targets} are six such questions. Five of them
(every one but the boundary-mode analysis, which is an eigenproblem posed on a
two-dimensional slice and which we postpone to \cref{ch:waveports}) reduce
Maxwell's equations to a governing PDE of one of three kinds:

\begin{itemize}[nosep]
  \item a \emph{boundary-value problem}---solve once for a static field
  (electrostatic, magnetostatic);
  \item an \emph{eigenvalue problem}---find the fields that ring with no drive
  (eigenmode);
  \item an \emph{evolution problem}---march a field forward in a parameter,
  either frequency (driven) or time (transient).
\end{itemize}

This chapter derives all five. Each derivation follows the same three-step
recipe, and it is worth naming the recipe before we run it, because the
repetition is the point.

\begin{keyidea}
Every reduction is three moves: \textbf{(1) choose a regime}---decide whether
$\partial/\partial t$ is zero (static), is replaced by $i\omega$
(time-harmonic), or is kept (time-domain); \textbf{(2) introduce a
potential or pick the primary field}---so that one of Maxwell's curl or
divergence constraints is satisfied identically; \textbf{(3) substitute the
constitutive law} ($\vD=\varepsilon\vE$, $\vB=\mu\vH$, $\vJ=\sigma\vE$) and
collect terms into a single PDE for the chosen unknown. The regime decides the
\emph{kind} of problem; the potential decides the \emph{function space}; the
constitutive law decides the \emph{material weights}.
\end{keyidea}

\Cref{fig:reductiontree} previews where the recipe lands. One box---Maxwell's
equations---fans out into five leaves, each a different PDE. Keep it in view;
the sections that follow are a guided walk down its five branches.

\begin{figure}[t]
\centering
\begin{tikzpicture}[
    every node/.style={font=\footnotesize},
    leaf/.style={product, text width=30mm, minimum height=13mm, font=\scriptsize},
    regime/.style={band, fill=simBgGray, draw=simGray, font=\scriptsize\itshape},
    node distance=6mm]
  % root
  \node[stage, minimum width=46mm, minimum height=12mm] (root)
    {\textbf{Maxwell's equations}\\[1pt]\scriptsize $\Curl\vE=-\partial_t\vB,\ \Curl\vH=\vJ+\partial_t\vD$};
  % regime band
  \node[regime, below left=10mm and 22mm of root] (rstat) {static $\partial_t=0$};
  \node[regime, below=10mm of root] (rharm) {harmonic $\partial_t\!\to i\omega$};
  \node[regime, below right=10mm and 22mm of root] (rtime) {time-domain};
  \draw[depedge] (root) -- (rstat);
  \draw[depedge] (root) -- (rharm);
  \draw[depedge] (root) -- (rtime);
  % leaves
  \node[leaf, below=7mm of rstat, xshift=-12mm] (es)
    {\textbf{Electrostatic}\\[1pt]$-\Div(\varepsilon\Grad V)=\rho$};
  \node[leaf, below=7mm of rstat, xshift=14mm] (ms)
    {\textbf{Magnetostatic}\\[1pt]$\Curl(\nu\Curl\vA)=\vJ$};
  \node[leaf, below=7mm of rharm, xshift=-9mm] (eig)
    {\textbf{Eigenmode}\\[1pt]$\Curl(\nu\Curl\vE)=\omega^2\varepsilon\vE$};
  \node[leaf, below=7mm of rharm, xshift=14mm] (dr)
    {\textbf{Driven}\\[1pt]$(\Kop+i\omega\Cop-\omega^2\Mop)\vE=\vb$};
  \node[leaf, below=7mm of rtime] (tr)
    {\textbf{Transient}\\[1pt]$\Mop\ddot\vs+\Cop\dot\vs+\Kop\vs=\vJ(t)$};
  \draw[flow] (rstat) -- (es);
  \draw[flow] (rstat) -- (ms);
  \draw[flow] (rharm) -- (eig);
  \draw[flow] (rharm) -- (dr);
  \draw[flow] (rtime) -- (tr);
\end{tikzpicture}
\caption[Maxwell's equations and their five reductions]{One theory, five
problems. The choice of \emph{regime} (static / time-harmonic / time-domain)
splits Maxwell's equations into three families; within each family the choice
of unknown and excitation fixes the individual PDE. The static branch yields two
boundary-value problems, the harmonic branch an eigenproblem and a driven
system, the time-domain branch an evolution. By the end of the chapter every
leaf will be re-expressed in the three operators $\Kop,\Mop,\Cop$.}
\label{fig:reductiontree}
\end{figure}

\begin{notationbox}
Throughout, $\varepsilon$ is the electric permittivity, $\mu$ the magnetic
permeability, and $\sigma$ the conductivity; all may vary with position and may
be tensors (anisotropic materials). We write $\nu=\mu^{-1}$ for the
\emph{reluctivity}, the magnetic counterpart of $\varepsilon$. Capital
$V$ is a scalar potential, $\vA$ a vector potential, $\vE,\vB,\vH,\vD,\vJ$ the
usual fields, and $\rho$ the free charge density. Sans-serif $\Kop,\Mop,\Cop$
denote the three \emph{assembled operators}---the discrete stiffness, mass, and
damping matrices---introduced in \cref{sec:threeoperators}.
\end{notationbox}

\section{Electrostatic: Maxwell collapses to Poisson}
\label{sec:electrostatic}

We begin with the simplest reduction, and the one the rest of the book treats as
its anchor. The physical question is the one from \cref{ch:targets}: hold a set
of conductors at fixed voltages and ask for the field between them.

\begin{physicsbox}
Two or more conductors sit in a dielectric. Each is an \emph{equipotential}: a
perfect conductor cannot sustain a tangential electric field, so its surface is
a surface of constant $V$. We prescribe those voltages and ask for the potential
$V$---hence the field $\vE=-\Grad V$---in the dielectric between them. Nothing
moves, nothing oscillates: this is a \emph{static} problem,
$\partial/\partial t = 0$.
\end{physicsbox}

\subsection{The derivation}

Apply the recipe. \emph{Step 1, regime:} set every time derivative to zero.
Faraday's law $\Curl\vE = -\partial_t\vB$ then loses its right-hand side:
\begin{equation}
  \Curl\vE = \mathbf{0}.
  \label{eq:es-curlfree}
\end{equation}
\emph{Step 2, potential:} a curl-free field on a simply connected domain is the
gradient of a scalar. (This is the converse of $\Curl\Grad\equiv\mathbf 0$, the
first exactness identity of the de Rham complex we meet in
\cref{ch:functionspaces}.) So there exists a scalar potential $V$ with
\begin{equation}
  \vE = -\Grad V,
  \label{eq:es-potential}
\end{equation}
the minus sign a convention making $\vE$ point ``downhill,'' from high to low
potential. Introducing $V$ \emph{automatically} satisfies
\cref{eq:es-curlfree}; we have traded a vector field obeying a constraint for an
unconstrained scalar field.

\emph{Step 3, the remaining law and the constitutive relation.} The equation we
have not yet used is Gauss's law, $\Div\vD = \rho$, together with the dielectric
constitutive relation $\vD = \varepsilon\vE$. Substituting
\cref{eq:es-potential},
\begin{equation}
  \rho \;=\; \Div\vD \;=\; \Div(\varepsilon\vE)
        \;=\; \Div\bigl(\varepsilon(-\Grad V)\bigr)
        \;=\; -\Div(\varepsilon\Grad V).
\end{equation}
Rearranging gives the governing equation of the electrostatic analysis:
\begin{equation}
  \boxed{\;-\Div(\varepsilon\Grad V) = \rho\;}
  \label{eq:es-poisson}
\end{equation}
the \emph{generalized Poisson equation}. In the charge-free dielectric between
the conductors $\rho=0$, and it degenerates to the homogeneous
\emph{Laplace equation} $-\Div(\varepsilon\Grad V)=0$---the equation Palace
actually solves, since the free charge lives only on the conductor surfaces and
enters through the boundary conditions, not the interior. If $\varepsilon$ is a
spatial constant it factors out, leaving the textbook $\Lap V = -\rho/\varepsilon$.

\begin{remark}
The same operator $-\Div(\varepsilon\Grad\,\cdot\,)$ governs steady heat
conduction (with $\varepsilon$ the thermal conductivity and $V$ the
temperature), steady incompressible potential flow, and the equilibrium shape of
a soap film. The electrostatic solver is, mathematically, a general second-order
elliptic solver wearing an electromagnetic hat. This is the first hint of the
book's thesis: the \emph{physics} is in the coefficients and the boundary data,
not in the operator.
\end{remark}

\subsection{Boundary conditions}

A PDE is only half a problem; the boundary conditions are the other half, and in
electrostatics they are where the physics enters. There are two kinds, and the
distinction---\emph{essential} versus \emph{natural}---runs through the whole of
\cref{part:fieldtheory}.

\begin{description}[style=nextline,leftmargin=1em,font=\normalfont\itshape]
  \item[Dirichlet (essential) conditions] prescribe the value of $V$ on a
  boundary. On each conductor surface $\Gamma_i$ we set $V=V_i$, the applied
  terminal voltage. These conditions are called \emph{essential} because they
  constrain the solution \emph{space} directly: an admissible $V$ must equal
  $V_i$ on $\Gamma_i$ before we even ask it to satisfy the PDE.
  \item[Neumann (natural) conditions] prescribe the \emph{flux}
  $\varepsilon\,\partial V/\partial n = -\vn\cdot\vD$ across a boundary. A
  homogeneous Neumann condition $\partial V/\partial n = 0$ means no normal
  field---a symmetry plane, or a far boundary. These are called \emph{natural}
  because, as we will see in \cref{ch:weak}, they fall out of the variational
  formulation \emph{for free}, requiring no constraint on the space.
\end{description}

Palace runs the electrostatic analysis as a \emph{family} of solves: it holds
one conductor at unit voltage ($V_i=1$) and all the others grounded ($V_j=0$,
$j\neq i$), solves \cref{eq:es-poisson}, and repeats for each conductor in turn.
The driver source makes this concrete---it loops over the terminal set,
\code{laplace\_op.GetExcitationVector(idx, ...)} forming each unit-voltage
right-hand side, and solves the \emph{same} assembled operator each
time.\footnote{\code{palace/drivers/electrostaticsolver.cpp:59-69}---the
per-terminal loop forming each RHS and solving the once-assembled stiffness
operator. We trace this driver end to end in \cref{ch:electrostatics}.} The
operator is assembled \emph{once}; only the right-hand side changes. That
``assemble once, solve a family'' shape is the fixed-operator solve corner of
\cref{ch:situating}.

\subsection{The operator that will be assembled}

When we discretize \cref{eq:es-poisson} by the finite-element method
(\cref{ch:galerkin,ch:assembly}), the differential operator
$-\Div(\varepsilon\Grad\,\cdot\,)$ becomes a matrix. Multiplying
\cref{eq:es-poisson} by a test function $v$ and integrating by parts (the subject
of \cref{ch:weak}) turns it into the \emph{bilinear form}
\begin{equation}
  a(u,v) = \int_\Omega \varepsilon\,\Grad u\cdot\Grad v \dV,
  \label{eq:es-bilinear}
\end{equation}
the $\varepsilon$-weighted \emph{stiffness} or \emph{diffusion} form. Assembled
over the mesh it becomes a sparse symmetric positive-(semi)definite matrix we
call $\Kop$, the \emph{stiffness matrix}. In the specification's vocabulary this
is one \code{weak\_form\_term}: the pair (coefficient $\varepsilon$, differential
operator \emph{gradient}), realizing $a(u,v)=(\varepsilon\,\Grad u,\Grad
v)$.\footnote{The differential-operator/coefficient factorization is the
\code{weak\_form\_term} value; the electrostatic stiffness is its
\emph{gradient} variant, witnessed at
\code{palace/models/laplaceoperator.cpp:191-194}. All four reductions in this
chapter are built from this one family of terms; see
\cref{tab:weakforms}.} The unknown $V$ lives in the scalar space $\Hone$
(continuous, nodal) of \cref{ch:functionspaces}; the matrix $\Kop$ is the
electrostatic instance of the stiffness operator that will recur, in different
clothes, in every reduction to follow.

\section{Magnetostatic: the curl--curl equation}
\label{sec:magnetostatic}

The magnetic dual of electrostatics is nearly its mirror image---which is the
whole pedagogical point. We run the identical recipe, change the field that the
potential satisfies, and a different (but structurally parallel) PDE drops out.

\begin{physicsbox}
A steady current $\vJ$ flows in conductors and coils; it produces a static
magnetic field $\vB$ threading the surrounding space. Nothing changes in
time, so again $\partial/\partial t = 0$, but now the \emph{magnetic} side of
Maxwell is active. We ask for the field $\vB$, hence the flux through circuits,
hence inductance.
\end{physicsbox}

\subsection{The derivation}

\emph{Step 1, regime:} static again, time derivatives zero. Amp\`ere's law
$\Curl\vH = \vJ + \partial_t\vD$ loses its displacement-current term, leaving
the magnetostatic Amp\`ere law
\begin{equation}
  \Curl\vH = \vJ.
  \label{eq:ms-ampere}
\end{equation}
\emph{Step 2, potential:} the constraint to satisfy identically is now Gauss's
law for magnetism, $\Div\vB = 0$. A divergence-free field is the curl of a
vector field (the converse of $\Div\Curl\equiv 0$, the \emph{second} exactness
identity of the de Rham complex). So introduce the \emph{magnetic vector
potential} $\vA$ with
\begin{equation}
  \vB = \Curl\vA.
  \label{eq:ms-potential}
\end{equation}
This automatically satisfies $\Div\vB=0$, just as $\vE=-\Grad V$ automatically
satisfied $\Curl\vE=\mathbf0$. The parallel is exact: where electrostatics
introduced a \emph{scalar} potential to kill a \emph{curl} constraint,
magnetostatics introduces a \emph{vector} potential to kill a \emph{divergence}
constraint. They sit one slot apart in the de Rham complex
(\cref{ch:functionspaces}).

\emph{Step 3, constitutive relation.} Use $\vH = \nu\vB$ with reluctivity
$\nu = \mu^{-1}$, and substitute \cref{eq:ms-potential} into
\cref{eq:ms-ampere}:
\begin{equation}
  \vJ \;=\; \Curl\vH \;=\; \Curl(\nu\vB) \;=\; \Curl(\nu\,\Curl\vA).
\end{equation}
This is the governing equation of the magnetostatic analysis, the
\emph{curl--curl equation}:
\begin{equation}
  \boxed{\;\Curl(\nu\,\Curl\vA) = \vJ\;}
  \label{eq:curlcurl}
\end{equation}
Compare it term for term with Poisson \cref{eq:es-poisson}: $\Grad\mapsto\Curl$,
$\varepsilon\mapsto\nu$, scalar $V\mapsto$ vector $\vA$, charge
$\rho\mapsto$ current $\vJ$. It is, almost literally, the electrostatic equation
with the gradient replaced by the curl and the scalar by a vector. This is the
``same machine, one part swapped'' promise of \cref{ch:targets} made into an
equation.

\subsection{The gauge non-uniqueness, and why it is a feature of the operator}

There is a subtlety here that does not arise in electrostatics, and it will cost
us real attention later, so we name it now. The vector potential $\vA$ is
\emph{not unique}. If $\vA$ solves \cref{eq:curlcurl}, so does $\vA + \Grad\phi$
for \emph{any} scalar field $\phi$, because
\begin{equation}
  \Curl\bigl(\vA + \Grad\phi\bigr) = \Curl\vA + \Curl\Grad\phi
    = \Curl\vA + \mathbf 0 = \Curl\vA = \vB,
\end{equation}
using the exactness identity $\Curl\Grad\equiv\mathbf0$ a second time. Adding any
gradient leaves $\vB$---the physically meaningful field---unchanged. This freedom
is called \emph{gauge freedom}, and it means the curl--curl operator
$\Curl(\nu\,\Curl\,\cdot\,)$ has a large \emph{null space}: it annihilates every
gradient field $\Grad\phi$. The assembled matrix $\Kop$ is therefore singular
(positive \emph{semi}definite, not definite).

\begin{pitfall}
A naive iterative solver applied to the singular curl--curl system can wander in
the gradient null space and fail to converge, even though the physical field
$\vB$ is perfectly well-defined. Two standard fixes---\emph{fixing the gauge}
(adding a constraint that selects one $\vA$) and \emph{projecting onto the
divergence-free subspace} (the discrete Helmholtz decomposition)---are the
business of \cref{ch:bc}. For now, simply note that the null space is real, that
it comes directly from the exactness identity $\Curl\Grad=\mathbf0$, and that it
is the price magnetostatics pays for its vector potential.
\end{pitfall}

\subsection{The operator}

Discretizing \cref{eq:curlcurl} yields the bilinear form
\begin{equation}
  a(u,v) = \int_\Omega \nu\,\Curl u\cdot\Curl v \dV,
  \label{eq:ms-bilinear}
\end{equation}
the $\nu$-weighted \emph{curl--curl} form---the \code{weak\_form\_term} with
coefficient $\nu$ and differential operator \emph{curl}.\footnote{The
\emph{curl} variant of \code{weak\_form\_term}, witnessed at
\code{palace/models/curlcurloperator.cpp:178-181}; the \emph{same}
\code{BilinearForm} build-up as the electrostatic stiffness, differing only in
the integrator (curl instead of gradient).} Assembled, it is again a stiffness
matrix $\Kop$---the \emph{magnetostatic} stiffness. The unknown $\vA$ now lives
in the vector space $\Hcurl$ (tangentially continuous, edge-based) rather than
the scalar $\Hone$; that change of function space, and the swap
$\varepsilon\to\nu$ in the coefficient, are \emph{the entire difference} between
the electrostatic and magnetostatic analyses. The solve corner, the assembly
machinery, the reduction-to-a-matrix---all are shared.

\section{Eigenmode: a generalized eigenproblem}
\label{sec:eigenmode}

We now leave the static world. Switch on time, but let it run \emph{without an
external source}: ask what fields the cavity can sustain on its own. Those are
its resonant modes, and finding them is an eigenvalue problem.

\begin{physicsbox}
Seal a region of space with conducting walls and put nothing inside---no charge,
no driving current. The empty cavity is not, electromagnetically, silent: it can
hold standing-wave fields that oscillate forever (in the lossless idealization)
at particular frequencies, its \emph{resonant modes}. Each mode is a field
pattern $\vE(\vx)$ paired with a frequency $\omega$. We seek the pairs.
\end{physicsbox}

\subsection{The time-harmonic reduction}

When a field oscillates sinusoidally at a single angular frequency $\omega$, we
write it as $\vE(\vx,t) = \Re\{\vE(\vx)\,e^{i\omega t}\}$ and work with the
complex \emph{phasor} amplitude $\vE(\vx)$. The payoff is that time derivatives
become algebra:
\begin{equation}
  \frac{\partial}{\partial t}\,e^{i\omega t} = i\omega\,e^{i\omega t},
  \qquad\text{so}\qquad
  \partial_t \;\longmapsto\; i\omega,\quad
  \partial_t^2 \;\longmapsto\; -\omega^2 .
  \label{eq:harmonic-rule}
\end{equation}
This substitution---the \emph{time-harmonic} or \emph{frequency-domain}
reduction---is the single most useful trick in computational electromagnetics,
and it powers both this analysis and the next. Apply it to source-free Maxwell.
Faraday and Amp\`ere, with no current and no charge, become
\begin{equation}
  \Curl\vE = -i\omega\vB = -i\omega\mu\vH,
  \qquad
  \Curl\vH = i\omega\vD = i\omega\varepsilon\vE.
\end{equation}
Take the curl of the first, divide by $\mu$ (use $\nu=\mu^{-1}$), and substitute
the second:
\begin{align}
  \Curl(\nu\,\Curl\vE)
    &= \Curl(-i\omega\vH)
     = -i\omega\,\Curl\vH \notag\\
    &= -i\omega\,(i\omega\varepsilon\vE)
     = \omega^2\varepsilon\vE.
\end{align}
The result is the source-free time-harmonic curl--curl equation:
\begin{equation}
  \boxed{\;\Curl(\nu\,\Curl\vE) = \omega^2\,\varepsilon\,\vE\;}
  \label{eq:eigen-pde}
\end{equation}
Look at what this is. The left side is the \emph{same} curl--curl operator that
governed magnetostatics, now acting on $\vE$. The right side is $\omega^2$ times
$\vE$ weighted by $\varepsilon$. This is an \emph{eigenvalue equation}: we seek
fields $\vE$ that the curl--curl operator merely \emph{rescales}, the scale
factor being $\omega^2\varepsilon$. Only special $\omega$ admit a nonzero
solution---the resonant frequencies---and each comes with its mode shape $\vE$.

\subsection{In operator form}
\label{sec:eigen-operatorform}

Discretize. The left side $\Curl(\nu\,\Curl\,\cdot\,)$ assembles, exactly as in
magnetostatics, into the stiffness matrix $\Kop$. The right-side weighting
$\varepsilon\,\cdot$ is the \code{weak\_form\_term} with coefficient
$\varepsilon$ and differential operator \emph{identity}---it assembles into the
$\varepsilon$-weighted \emph{mass matrix} $\Mop$.\footnote{The \emph{identity}
variant of \code{weak\_form\_term}, the $L^2$ pairing $a(u,v)=(Q\,u,v)$; the
mass matrix is assembled by the very same \code{BilinearForm} fold, witnessed at
\code{palace/models/spaceoperator.cpp:278}. Three reductions in (electrostatic,
magnetostatic, eigenmode) and we have already met all three of the operators
$\Kop,\Mop,\Cop$ except the damping $\Cop$.} With $\lambda=\omega^2$ the
discrete form of \cref{eq:eigen-pde} is the \emph{generalized eigenvalue
problem}
\begin{equation}
  \boxed{\;\Kop\,\vx = \lambda\,\Mop\,\vx,\qquad \lambda=\omega^2\;}
  \label{eq:gevp}
\end{equation}
Each eigenvalue $\lambda$ gives a resonant frequency $\omega=\sqrt\lambda$; each
eigenvector $\vx$ gives the mode field. ``Generalized'' refers to the matrix
$\Mop$ on the right: an \emph{ordinary} eigenproblem would read
$\Kop\vx=\lambda\vx$, but the $\varepsilon$-weighting puts $\Mop$ there instead.
Solving \cref{eq:gevp} is the one analysis whose ``solve'' is a single call to a
specialized eigenvalue library rather than a loop we write ourselves---the
opaque black-box solve corner of \cref{ch:situating}, treated in depth in
\cref{ch:eigenmodes}.\footnote{The driver hands the assembled pencil
$(\Kop,\Mop)$ to one library call, \code{eigen->Solve()}; there is no
Palace-authored iteration around it (\code{palace/drivers/eigensolver.cpp:367}).}

\subsection{The lossy case: a quadratic eigenproblem}

The derivation above assumed a lossless cavity. Real cavities lose energy---to
finite wall conductivity, to dielectric absorption, to radiation through ports.
Loss introduces a conduction current $\vJ=\sigma\vE$ into Amp\`ere's law, which
under the harmonic rule \cref{eq:harmonic-rule} contributes a term linear in
$i\omega$. Carrying it through the derivation adds a middle operator, the
\emph{damping} matrix $\Cop$ (the $\sigma$-weighted mass term), and the
eigenproblem becomes \emph{quadratic} in the eigenvalue:
\begin{equation}
  \bigl(\Kop + \lambda\,\Cop + \lambda^2\,\Mop\bigr)\,\vx = \mathbf 0,
  \qquad \lambda = i\omega.
  \label{eq:qevp}
\end{equation}
Now $\lambda$ is complex; its imaginary part is the oscillation frequency and its
real part the decay rate, from which the quality factor $Q$ is read off (a mode
that decays slowly has high $Q$). The lossless problem \cref{eq:gevp} is the
special case $\Cop=0$, with $\lambda=\omega^2$ real. The polynomial structure
$\Kop+\lambda\Cop+\lambda^2\Mop$---three operators, combined with powers of
$\lambda$---is worth fixing in mind, because the \emph{very same} polynomial is
about to reappear in the driven analysis, with $\lambda$ replaced by a fixed
$i\omega$.

\section{Driven: a complex-valued Helmholtz system}
\label{sec:driven}

The eigenmode analysis asked which frequencies the cavity \emph{chooses}. The
driven analysis asks the complementary question: \emph{we} choose the frequency,
push the system at it, and ask how it responds.

\begin{physicsbox}
Feed a sinusoidal signal into a microwave device at a frequency $\omega$ of our
choosing. The device responds with a steady oscillation at that same frequency:
some of the signal passes through, some reflects. We want the response---the
fields, and from them the scattering parameters---at each frequency across a
band. Unlike the eigenmode problem, there \emph{is} a source, and $\omega$ is an
input, not an unknown.
\end{physicsbox}

\subsection{The derivation}

Take the full time-harmonic Maxwell system---now \emph{with} sources and
\emph{with} loss---and apply the harmonic rule \cref{eq:harmonic-rule}. The same
curl-of-Faraday manipulation as before, but retaining the conduction current
$\sigma\vE$ and an impressed source current $\vJ_{\text{src}}$, gives
\begin{equation}
  \Curl(\nu\,\Curl\vE)
    + i\omega\,\sigma\,\vE
    - \omega^2\,\varepsilon\,\vE
    = -i\omega\,\vJ_{\text{src}}.
  \label{eq:driven-pde}
\end{equation}
This is a \emph{Helmholtz-type} equation: a curl--curl operator plus a
zeroth-order term proportional to $-\omega^2$, the vector analogue of the scalar
Helmholtz equation $u'' + k^2 u = 0$ we will meet in
\cref{wex:cavity1d}. Read its three pieces against the operators we have already
named:
\begin{itemize}[nosep]
  \item $\Curl(\nu\,\Curl\,\cdot\,)$ is the curl--curl \emph{stiffness} $\Kop$
  (\cref{sec:magnetostatic});
  \item $\sigma\,\cdot$ is the conduction/loss term---the $\sigma$-weighted
  \emph{damping} $\Cop$, entering linearly in $i\omega$;
  \item $\varepsilon\,\cdot$ is the displacement-current \emph{mass} $\Mop$
  (\cref{sec:eigen-operatorform}), entering as $-\omega^2$.
\end{itemize}

\subsection{In operator form}

Assembling each term, \cref{eq:driven-pde} becomes the discrete \emph{driven
system}
\begin{equation}
  \boxed{\;\bigl(\Kop + i\omega\,\Cop - \omega^2\,\Mop\bigr)\,\vE = \vb(\omega)\;}
  \label{eq:driven-system}
\end{equation}
with right-hand side $\vb(\omega)=-i\omega\,\vJ_{\text{src}}$ the discretized
source. Define $\mat{A}(\omega) = \Kop + i\omega\,\Cop - \omega^2\,\Mop$: this is the
\emph{frequency-dependent system operator}. The structural resemblance to the
quadratic eigenproblem \cref{eq:qevp} is not a coincidence---both are the
polynomial $\Kop + \lambda\Cop + \lambda^2\Mop$ in disguise. The eigenproblem
\emph{solves for} the $\lambda=i\omega$ that make it singular; the driven problem
\emph{fixes} $\omega$ and solves the (generally nonsingular) linear system
\cref{eq:driven-system} for $\vE$.

\begin{keyidea}
The driven and eigenmode analyses share one operator polynomial,
$\Kop+\lambda\Cop+\lambda^2\Mop$, and differ only in what they hold fixed.
\emph{Eigenmode}: the right-hand side is zero and $\lambda$ is the unknown---a
\emph{search for the $\omega$ that resonate}. \emph{Driven}: $\omega$ is given,
the right-hand side is the source, and the field is the unknown---a
\emph{response at a chosen $\omega$}. Resonance and response are two questions
about the same three matrices.
\end{keyidea}

\subsection{The solve corner: assemble \emph{inside} the loop}

There is one operational fact about \cref{eq:driven-system} that the rest of the
book leans on heavily. The operator $\mat{A}(\omega)$ \emph{depends on the
frequency}. A driven analysis sweeps a whole band of frequencies, and at each one
it must \emph{rebuild} $\mat{A}(\omega)$ from the fixed pieces $\Kop,\Cop,\Mop$ and
re-solve. Contrast electrostatics, which assembled its operator \emph{once} and
reused it across the whole terminal family: there the operator was constant and
the assembly hoisted out of the loop. Here the operator varies, and the assembly
lives \emph{inside} the loop.\footnote{The driver builds the solver once but
captures the operator per frequency: \code{ksp.SetOperators(*A, *P)} sits
\emph{inside} the frequency loop (\code{palace/drivers/drivensolver.cpp:176-180}),
the exact negation of the electrostatic hoist. This operator-varying-map versus
fixed-operator-map distinction is a load-bearing axis of \cref{ch:situating} and
\cref{ch:driven}.} The three matrices $\Kop,\Cop,\Mop$ are assembled once and
held; only their \emph{combination} $\mat{A}(\omega)$ is rebuilt per step. We are
beginning to see why isolating the three operators matters so much: they are the
fixed, reusable raw material, and the analysis is the recipe for recombining
them.

\section{Transient: a second-order evolution}
\label{sec:transient}

The final reduction keeps time in full. Rather than assume a single frequency, we
inject an arbitrary pulse and watch the field evolve, step by step, in time.

\begin{physicsbox}
Inject a current pulse $\vJ(t)$---a short burst, a step, a modulated wave packet,
anything---into a device and follow the fields as they propagate, reflect, and
ring down. The output is not a number or a spectrum but a \emph{movie}: the
entire time history of the field. This is the analysis of choice when the
excitation is broadband, or when one simply wants to watch a wave move.
\end{physicsbox}

\subsection{The derivation}

We do \emph{not} apply the harmonic rule; we keep $\partial_t$ as a genuine time
derivative. The same curl-of-Faraday manipulation that produced the driven
equation, but with $i\omega\mapsto\partial_t$ and $-\omega^2\mapsto\partial_t^2$
(reversing the substitution \cref{eq:harmonic-rule}), gives the time-domain
\emph{second-order wave equation}
\begin{equation}
  \varepsilon\,\frac{\partial^2\vE}{\partial t^2}
    + \sigma\,\frac{\partial\vE}{\partial t}
    + \Curl(\nu\,\Curl\vE)
    = -\frac{\partial\vJ_{\text{src}}}{\partial t}.
  \label{eq:transient-pde}
\end{equation}
Each term is the time-domain image of a term in the driven equation
\cref{eq:driven-pde}: the displacement-current $\varepsilon\,\partial_t^2$, the
conduction $\sigma\,\partial_t$, the curl--curl $\Curl(\nu\Curl\,\cdot\,)$. The
harmonic and transient equations are the \emph{same physics} expressed in
frequency and in time; the substitution \cref{eq:harmonic-rule} is the bridge
between them.

\subsection{Semi-discretization in space}

Discretize \emph{only the space} (leave time continuous for the moment)---this is
the \emph{method of lines}. Each spatial term assembles into the matrix we
already know, and the field becomes a vector of time-dependent degrees of freedom
$\vs(t)$. The result is a system of \emph{ordinary} differential equations in
time:
\begin{equation}
  \boxed{\;\Mop\,\ddot\vs + \Cop\,\dot\vs + \Kop\,\vs = \vJ(t)\;}
  \label{eq:transient-ode}
\end{equation}
where the dots are time derivatives, $\Mop$ is the $\varepsilon$-mass, $\Cop$ the
$\sigma$-damping, $\Kop$ the curl--curl stiffness, and $\vJ(t)$ the discretized
forcing $-\partial_t\vJ_{\text{src}}$. This is precisely the equation of a damped
mechanical oscillator---mass times acceleration, plus damping times velocity,
plus stiffness times displacement, equals the applied force---and the analogy is
exact: $\Mop$ resists change in the field (inertia), $\Kop$ pulls it back toward
equilibrium (restoring stiffness), $\Cop$ bleeds energy away (damping). The names
$\Mop$ (mass), $\Kop$ (stiffness), $\Cop$ (damping) are borrowed wholesale from
structural dynamics, and they earn their keep here.

\subsection{Recast as a first-order IVP}

Standard time integrators expect a \emph{first-order} system $\dot\vy =
\vF(\vy,t)$. We reduce the order by carrying the velocity as an extra unknown.
Set $\vv = \dot\vs$ and stack $(\vs,\vv)$ into a doubled state vector; then
\cref{eq:transient-ode} becomes
\begin{equation}
  \frac{d}{dt}\begin{pmatrix}\vs\\\vv\end{pmatrix}
  =
  \begin{pmatrix}
    \vv\\
    \Mop^{-1}\bigl(\vJ(t) - \Cop\,\vv - \Kop\,\vs\bigr)
  \end{pmatrix},
  \label{eq:transient-firstorder}
\end{equation}
a first-order initial-value problem in the doubled state, ready for an
off-the-shelf ODE integrator.\footnote{Palace builds exactly this first-order
operator (\code{TimeDependentFirstOrderOperator}) and hands it to a library
integrator---Generalized-$\alpha$, SDIRK, or others---constructed once outside
the loop (\code{palace/models/timeoperator.cpp:311-335}). The detailed time
integration is \cref{ch:time}.} The crucial feature, the one that sets transient
apart from every other analysis in this chapter, is its \emph{sequential}
character: each time step starts from the state the previous step produced. The
electrostatic terminal family and the driven frequency sweep were collections of
\emph{independent} solves; the transient march is a \emph{chain}, where step
$n{+}1$ cannot begin until step $n$ is done. This carry-the-state-forward
structure is the \emph{fold}, the keystone of \cref{ch:fold} and the solve corner
of \cref{ch:situating}, and it is why the transient analysis is the book's
canonical example of a genuinely stateful computation.

\section{Two worked reductions}
\label{sec:worked}

The five derivations above are general. To make them concrete---and checkable by
hand---we now solve two one-dimensional instances in full. The first reduces
electrostatics to an ODE and solves it; the second reduces a wave to a Helmholtz
equation and extracts a spectrum of cavity modes, previewing the eigenmode
analysis with real numbers.

\begin{workedexample}{Parallel plates: electrostatics in one dimension}{plates1d}
Two large parallel conducting plates sit at $x=0$ and $x=L$, the lower held at
$V=0$, the upper at $V=V_0$, with a uniform dielectric and \emph{no free charge}
between them. Find the potential and field in the gap.

\medskip
\emph{Reduce.} By symmetry the field depends only on $x$, so $V=V(x)$ and the
Poisson equation \cref{eq:es-poisson} with constant $\varepsilon$ and $\rho=0$
collapses to the ordinary differential equation
\begin{equation*}
  -\varepsilon\,V''(x) = 0
  \qquad\Longrightarrow\qquad
  V''(x) = 0,\quad 0<x<L,
\end{equation*}
the one-dimensional Laplace equation, with Dirichlet data $V(0)=0$, $V(L)=V_0$.

\emph{Solve.} The general solution of $V''=0$ is the affine function
$V(x)=ax+b$. The boundary conditions fix the constants: $V(0)=b=0$ and
$V(L)=aL=V_0$, so $a=V_0/L$. Hence
\begin{equation*}
  V(x) = \frac{V_0}{L}\,x,
  \qquad
  E(x) = -V'(x) = -\frac{V_0}{L}.
\end{equation*}
The potential ramps \emph{linearly} across the gap and the field is
\emph{uniform}---the textbook parallel-plate result, here derived as the
one-dimensional reduction of the electrostatic PDE. The surface charge on the
upper plate follows from Gauss's law as
$\sigma_s = \vn\cdot\vD = \varepsilon|E| = \varepsilon V_0/L$, and since the
total charge is $\sigma_s$ times the plate area $A$, the capacitance is
\begin{equation*}
  C = \frac{Q}{V_0} = \frac{\varepsilon A V_0/L}{V_0}
    = \frac{\varepsilon A}{L},
\end{equation*}
recovering $C=\varepsilon A/L$. The whole electrostatic pipeline---solve for $V$,
differentiate to $\vE$, integrate the flux, divide by voltage---runs here in four
lines of algebra; \cref{ch:electrostatics} runs the identical pipeline in three
dimensions on a finite-element mesh.
\end{workedexample}

\Cref{fig:plates1d} draws the result: the linear potential ramp and the uniform
field it produces.

\begin{figure}[t]
\centering
\begin{tikzpicture}
\begin{axis}[
    width=0.62\linewidth, height=46mm,
    xlabel={position $x/L$}, ylabel={},
    xmin=0, xmax=1, ymin=-0.15, ymax=1.15,
    axis lines=left, xtick={0,0.5,1}, ytick={0,0.5,1},
    yticklabel style={font=\scriptsize}, xticklabel style={font=\scriptsize},
    legend style={font=\scriptsize, draw=simGray!50, at={(0.02,0.98)},
      anchor=north west},
    every axis plot/.append style={very thick},
  ]
  \addplot[simBlue, domain=0:1, samples=2] {x};
  \addlegendentry{$V(x)/V_0 = x/L$}
  \addplot[simRust, domain=0:1, samples=2] {0.5};
  \addlegendentry{$-E\,L/V_0 = 1$ (uniform)}
\end{axis}
\end{tikzpicture}
\caption[The one-dimensional parallel-plate solution]{The parallel-plate
reduction of \cref{wex:plates1d}. The potential $V$ (blue) ramps linearly from
$0$ to $V_0$ across the gap; the field $E=-V'$ (rust) is constant. This is the
simplest non-trivial instance of every electrostatic solve in the book---and the
exact answer the finite-element method must reproduce on a one-element mesh.}
\label{fig:plates1d}
\end{figure}

\begin{workedexample}{The one-dimensional cavity: from wave to spectrum}{cavity1d}
A field $u(x,t)$ on the interval $0\le x\le L$ satisfies the lossless wave
equation $\partial_t^2 u = c^2\,\partial_x^2 u$, with the field pinned to zero at
the ends, $u(0,t)=u(L,t)=0$ (perfectly conducting walls). Find the frequencies at
which the cavity can ring, and the mode shapes.

\medskip
\emph{Reduce (time-harmonic).} Seek a single-frequency solution
$u(x,t)=\Re\{u(x)\,e^{i\omega t}\}$. The harmonic rule \cref{eq:harmonic-rule}
sends $\partial_t^2\mapsto-\omega^2$, so the wave equation becomes the
one-dimensional \emph{Helmholtz equation}
\begin{equation*}
  -\omega^2 u(x) = c^2\,u''(x)
  \qquad\Longrightarrow\qquad
  u''(x) + k^2\,u(x) = 0,
  \qquad k \equiv \frac{\omega}{c},
\end{equation*}
with $u(0)=u(L)=0$. This is the one-dimensional scalar shadow of the
source-free curl--curl eigenproblem \cref{eq:eigen-pde}: a second-order operator
on the left, $k^2$ times the field on the right, no source---an
\emph{eigenproblem for $k$}.

\emph{Solve.} The general solution of $u''+k^2u=0$ is
$u(x)=\alpha\sin(kx)+\beta\cos(kx)$. The condition $u(0)=0$ kills the cosine,
$\beta=0$. The condition $u(L)=0$ then demands $\alpha\sin(kL)=0$; a nonzero mode
($\alpha\neq 0$) requires
\begin{equation*}
  \sin(kL) = 0
  \qquad\Longrightarrow\qquad
  kL = n\pi,\quad n=1,2,3,\dots
\end{equation*}
So the cavity rings only at the \emph{discrete} wavenumbers and frequencies
\begin{equation*}
  \boxed{\;k_n = \frac{n\pi}{L},\qquad \omega_n = c\,k_n = \frac{n\pi c}{L}\;}
\end{equation*}
with mode shapes $u_n(x)=\sin(n\pi x/L)$. The lowest mode $n=1$ fits a half
wavelength in the cavity; higher modes fit $n$ half-wavelengths
(\cref{fig:cavitymodes}).

\emph{Connect to the operator form.} Discretize $-u''$ on a grid and it becomes a
stiffness matrix $\Kop$; discretize the $u$ on the right and it becomes a mass
matrix $\Mop$. The continuous eigenproblem $-u''=k^2u$ is then exactly the
discrete generalized eigenproblem $\Kop\vx=\lambda\Mop\vx$ of \cref{eq:gevp},
with $\lambda=k^2$. The countable spectrum $k_n=n\pi/L$ is the one-dimensional
preview of the resonant frequencies the eigenmode solver returns in three
dimensions---now you have computed them by hand.
\end{workedexample}

\begin{figure}[t]
\centering
\begin{tikzpicture}
\begin{axis}[
    width=0.8\linewidth, height=50mm,
    xlabel={position $x/L$}, ylabel={mode shape $u_n(x)$},
    xmin=0, xmax=1, ymin=-1.25, ymax=1.25,
    axis lines=middle, xtick={0,0.25,0.5,0.75,1}, ytick={-1,0,1},
    yticklabel style={font=\scriptsize}, xticklabel style={font=\scriptsize},
    xlabel style={font=\scriptsize, at={(axis description cs:0.5,-0.04)}},
    ylabel style={font=\scriptsize},
    legend style={font=\scriptsize, draw=simGray!50, at={(1.0,1.18)},
      anchor=north east, legend columns=3, column sep=4pt},
    every axis plot/.append style={very thick, samples=80},
  ]
  \addplot[simBlue, domain=0:1] {sin(deg(pi*x))};
  \addlegendentry{$n=1$}
  \addplot[simTeal, domain=0:1] {sin(deg(2*pi*x))};
  \addlegendentry{$n=2$}
  \addplot[simRust, domain=0:1] {sin(deg(3*pi*x))};
  \addlegendentry{$n=3$}
\end{axis}
\end{tikzpicture}
\caption[Standing-wave modes of the one-dimensional cavity]{The first three
cavity modes $u_n(x)=\sin(n\pi x/L)$ of \cref{wex:cavity1d}, the eigenfunctions
of $-u''=k^2u$ with $u(0)=u(L)=0$. Mode $n$ fits $n$ half-wavelengths between the
walls and rings at $\omega_n=n\pi c/L$. These standing waves are the
one-dimensional ancestors of the three-dimensional cavity modes the eigenmode
analysis computes in \cref{ch:eigenmodes}.}
\label{fig:cavitymodes}
\end{figure}

\section{The three operators, recombined five ways}
\label{sec:threeoperators}

Step back and look at the five boxed governing equations together. They were
derived from different regimes, live on different function spaces, and produce
different products---yet, written in operator form, they are built from a
\emph{single vocabulary of three matrices}:
\begin{itemize}[nosep]
  \item the \emph{stiffness} $\Kop$---the second-order spatial operator
  ($\varepsilon$-weighted $\Grad$ for the scalar problems, $\nu$-weighted
  $\Curl$ for the vector ones);
  \item the \emph{mass} $\Mop$---the zeroth-order $\varepsilon$-weighted $L^2$
  pairing, the displacement-current term;
  \item the \emph{damping} $\Cop$---the $\sigma$-weighted $L^2$ pairing, the
  conduction/loss term.
\end{itemize}
Every reduction in this chapter is a way of \emph{combining} these three, as
\cref{fig:kmc} shows. The static problems use $\Kop$ alone. The eigenproblem
pairs $\Kop$ against $\Mop$. The driven problem combines all three with the
complex weights $1,\,i\omega,\,-\omega^2$. The transient problem combines all
three again, but as coefficients of $\vs,\dot\vs,\ddot\vs$. The \emph{matrices}
are shared; the \emph{recipe for combining them} is the analysis.

\begin{figure}[t]
\centering
\begin{tikzpicture}[
    every node/.style={font=\footnotesize},
    op/.style={stage, minimum width=15mm, minimum height=10mm,
      font=\small\bfseries},
    use/.style={product, text width=33mm, minimum height=9mm, font=\scriptsize,
      align=left},
    node distance=5mm]
  % the three shared operators, left column
  \node[op] (K) {$\Kop$};
  \node[op, below=8mm of K] (M) {$\Mop$};
  \node[op, below=8mm of M] (C) {$\Cop$};
  \node[font=\scriptsize\itshape, above=1mm of K] {three shared operators};
  % the five recombinations, right column
  \node[use, right=34mm of K, yshift=12mm] (es) {\textbf{Electrostatic:} $\ \Kop\vx=\vb$};
  \node[use, below=3mm of es] (ms) {\textbf{Magnetostatic:} $\ \Kop\vx=\vb$};
  \node[use, below=3mm of ms] (eig) {\textbf{Eigenmode:} $\ \Kop\vx=\lambda\Mop\vx$};
  \node[use, below=3mm of eig] (dr) {\textbf{Driven:} $\ (\Kop{+}i\omega\Cop{-}\omega^2\Mop)\vx=\vb$};
  \node[use, below=3mm of dr] (tr) {\textbf{Transient:} $\ \Mop\ddot\vx{+}\Cop\dot\vx{+}\Kop\vx=\vb$};
  % edges K -> all (K appears in all five)
  \draw[depedge] (K.east) -- (es.west);
  \draw[depedge] (K.east) -- (ms.west);
  \draw[depedge] (K.east) -- (eig.west);
  \draw[depedge] (K.east) -- (dr.west);
  \draw[depedge] (K.east) -- (tr.west);
  % edges M -> eigen, driven, transient
  \draw[refedge] (M.east) -- (eig.west);
  \draw[refedge] (M.east) -- (dr.west);
  \draw[refedge] (M.east) -- (tr.west);
  % edges C -> driven, transient
  \draw[backflow] (C.east) -- (dr.west);
  \draw[backflow] (C.east) -- (tr.west);
\end{tikzpicture}
\caption[The $\Kop/\Mop/\Cop$ recurrence]{Three operators, five recombinations.
The stiffness $\Kop$ (solid blue) appears in every analysis; the mass $\Mop$
(dotted) joins for the eigenmode, driven, and transient problems; the damping
$\Cop$ (dashed rust) appears only where loss or conduction matters (driven,
transient). The static analyses are $\Kop$ alone. The figure is the chapter's
thesis in one picture: \emph{the analysis is not new physics, it is a new way of
combining the same three matrices.}}
\label{fig:kmc}
\end{figure}

This recurrence is the structural payoff of the chapter, and it is worth stating
as the central idea.

\begin{keyidea}
The five reductions of Maxwell are built from \emph{three} assembled
operators---stiffness $\Kop$, mass $\Mop$, damping $\Cop$---recombined five ways.
What distinguishes the analyses is not the operators (those are shared) but
\emph{how they are combined and what is solved}: $\Kop$ alone for a static
boundary-value problem; the pencil $(\Kop,\Mop)$ for an eigenproblem; the
frequency-weighted sum $\Kop+i\omega\Cop-\omega^2\Mop$ for a driven linear
system; the time-derivative-weighted sum for a transient evolution. Build the
three operators well, and you have built every analysis.
\end{keyidea}

\subsection{The regimes on a timeline}

One more lens before the summary table. The five reductions partition cleanly by
how they treat time, and \cref{fig:timeline} arranges them on that axis. The
static problems freeze it; the harmonic problems replace it with a single complex
frequency; the transient problem keeps it and marches. Reading left to right is
reading the harmonic rule \cref{eq:harmonic-rule} forward and backward:
\emph{static} is the $\omega\to 0$ limit, \emph{harmonic} is a single $\omega$,
\emph{transient} is all frequencies at once (a pulse is a superposition of
harmonics).

\begin{figure}[t]
\centering
\begin{tikzpicture}[every node/.style={font=\footnotesize}]
  % the time axis
  \draw[flow, simGray] (-0.3,0) -- (12.3,0)
    node[right, font=\scriptsize\itshape, text=simGray] {treatment of $\partial_t$};
  % three regime bands
  \node[band, fill=simBgBlue, draw=simBlue, minimum width=30mm] (stat) at (1.8,1)
    {\textbf{Static}\\[-1pt]\scriptsize $\partial_t = 0$};
  \node[band, fill=simBgGold, draw=simGold, minimum width=38mm] (harm) at (6.2,1)
    {\textbf{Time-harmonic}\\[-1pt]\scriptsize $\partial_t \to i\omega$};
  \node[band, fill=simBgRust, draw=simRust, minimum width=30mm] (time) at (10.5,1)
    {\textbf{Time-domain}\\[-1pt]\scriptsize $\partial_t$ kept};
  \foreach \n in {stat,harm,time}{\draw[refedge] (\n.south) -- (\n.south|-0,0.06);}
  % leaves under each band
  \node[font=\scriptsize, align=center, below=2mm of {(1.8,0)}] (le1)
    {electrostatic\\ magnetostatic\\[1pt]\itshape boundary-value};
  \node[font=\scriptsize, align=center, below=2mm of {(6.2,0)}] (le2)
    {eigenmode \quad driven\\[1pt]\itshape eigen \quad linear-solve};
  \node[font=\scriptsize, align=center, below=2mm of {(10.5,0)}] (le3)
    {transient\\[1pt]\itshape evolution (fold)};
  \draw[refedge] (1.8,0) -- (le1.north);
  \draw[refedge] (6.2,0) -- (le2.north);
  \draw[refedge] (10.5,0) -- (le3.north);
\end{tikzpicture}
\caption[The reductions arranged by regime]{The five reductions along the
time axis. Freezing time ($\partial_t=0$) gives the two static boundary-value
problems; replacing it with a single complex frequency ($\partial_t\to i\omega$)
gives the eigenmode (no source) and driven (with source) problems; keeping it
gives the transient evolution. The static problems are the $\omega\to0$ corner of
the harmonic ones; the transient problem is, by Fourier superposition, all
frequencies at once.}
\label{fig:timeline}
\end{figure}

\section{The five reductions, side by side}

\Cref{tab:reductions} collects the chapter. Read across each row to see one
analysis whole---its PDE, the operators it assembles, the function space its
unknown lives in, and the solve corner that handles it. Read \emph{down} the
``operators'' column and the recurrence of \cref{fig:kmc} reappears in symbols:
$\Kop$ in every row, $\Mop$ from the eigenmode down, $\Cop$ only where there is
loss.

\begin{table}[t]
\centering
\caption[The five reductions of Maxwell, compared]{The five PDE reductions of
this chapter, with the operators each assembles, the function space its unknown
inhabits (\cref{ch:functionspaces}), and the solve corner that processes it
(\cref{ch:situating}). The boundary-mode analysis---a sixth, eigen-type problem
posed on a 2-D cross-section---is deferred to \cref{ch:waveports}.}
\label{tab:reductions}
\small
\setlength{\tabcolsep}{4pt}
\renewcommand{\arraystretch}{1.25}
\begin{tabular}{@{}lllll@{}}
\toprule
Analysis & Governing PDE & Operators & Space & Solve corner \\
\midrule
Electrostatic
  & $-\Div(\varepsilon\Grad V)=\rho$
  & $\Kop$
  & $\Hone$
  & fixed-operator map \\
Magnetostatic
  & $\Curl(\nu\Curl\vA)=\vJ$
  & $\Kop$
  & $\Hcurl$
  & fixed-operator map \\
Eigenmode
  & $\Curl(\nu\Curl\vE)=\omega^2\varepsilon\vE$
  & $\Kop,\Mop\ (\Cop)$
  & $\Hcurl$
  & black-box eigen \\
Driven
  & $(\Kop{+}i\omega\Cop{-}\omega^2\Mop)\vE=\vb$
  & $\Kop,\Mop,\Cop$
  & $\Hcurl$
  & operator-varying map \\
Transient
  & $\Mop\ddot\vs{+}\Cop\dot\vs{+}\Kop\vs=\vJ(t)$
  & $\Kop,\Mop,\Cop$
  & $\Hcurl$
  & state-threaded fold \\
\bottomrule
\end{tabular}
\end{table}

Three columns of \cref{tab:reductions} are promissory notes the rest of
\cref{part:fieldtheory} and \cref{part:framework} redeem. The \emph{operators}
column---how each $\Kop,\Mop,\Cop$ is built from a mesh and materials by
integrating a bilinear form---is the assembly story of
\cref{ch:galerkin,ch:assembly}, and the bilinear-form family that unifies all of
them is previewed in \cref{tab:weakforms}. The \emph{space} column---why $V$
belongs in $\Hone$ and $\vE,\vA$ in $\Hcurl$, and what the de Rham complex has to
do with it---is \cref{ch:functionspaces,ch:derham}. The \emph{solve corner}
column---the four kinds of solve and why each analysis lands where it does---is
the framework of \cref{part:framework}. This chapter has supplied the rows; the
book supplies the columns.

\begin{table}[t]
\centering
\caption[The bilinear-form family behind the operators]{Every operator in
\cref{tab:reductions} is assembled from one family of \emph{weak-form terms}
$a(u,v)=(Q\,\diffop u,\diffop v)$, differing only in the coefficient $Q$ and the
differential operator $\diffop$. Three of the four variants are exercised by the
five reductions; the divergence variant awaits an analysis that needs it. This is
the \code{weak\_form\_term} family of \cref{ch:galerkin}.}
\label{tab:weakforms}
\small
\setlength{\tabcolsep}{6pt}
\renewcommand{\arraystretch}{1.25}
\begin{tabular}{@{}llll@{}}
\toprule
$\diffop$ & Bilinear form $a(u,v)$ & Operator & Appears in \\
\midrule
$\Grad$ (gradient)
  & $(\varepsilon\,\Grad u,\Grad v)$
  & stiffness $\Kop$ (scalar)
  & electrostatic \\
$\Curl$ (curl)
  & $(\nu\,\Curl u,\Curl v)$
  & stiffness $\Kop$ (vector)
  & magneto-, eigen, driven, transient \\
$I$ (identity)
  & $(\varepsilon\,u,v)$ or $(\sigma\,u,v)$
  & mass $\Mop$, damping $\Cop$
  & eigen, driven, transient \\
$\Div$ (divergence)
  & $(Q\,\Div u,\Div v)$
  & (div--div stiffness)
  & ---\,(no reduction here) \\
\bottomrule
\end{tabular}
\end{table}

\summarysection
Maxwell's equations are too rich to compute against whole; every analysis first
\emph{reduces} them to a single governing PDE by a three-step recipe---fix the
regime, introduce a potential or primary field, substitute the constitutive law.
Run five times, the recipe yields: the \emph{electrostatic} Poisson equation
$-\Div(\varepsilon\Grad V)=\rho$ (curl-free $\vE=-\Grad V$ plus Gauss), assembled
as the scalar stiffness $\Kop$ on $\Hone$; the \emph{magnetostatic} curl--curl
equation $\Curl(\nu\Curl\vA)=\vJ$ (divergence-free $\vB=\Curl\vA$ plus
Amp\`ere), assembled as the vector stiffness $\Kop$ on $\Hcurl$ with a gauge null
space; the \emph{eigenmode} generalized eigenproblem $\Kop\vx=\lambda\Mop\vx$
($\lambda=\omega^2$) from source-free time-harmonic curl--curl, with the lossy
case quadratic in $\lambda$; the \emph{driven} Helmholtz system
$(\Kop+i\omega\Cop-\omega^2\Mop)\vE=\vb(\omega)$ from time-harmonic Maxwell with
sources and loss, re-assembled per frequency; and the \emph{transient}
second-order ODE $\Mop\ddot\vs+\Cop\dot\vs+\Kop\vs=\vJ(t)$, recast as a
first-order IVP and marched as a sequential fold. Beneath the variety, all five
are built from the \emph{same three operators} $\Kop,\Mop,\Cop$---a stiffness, a
mass, a damping---recombined five ways. The analysis is the recipe for combining
them, not new physics.

\furtherreading
The reductions in this chapter are standard computational-electromagnetics
material. Jin's \emph{The Finite Element Method in Electromagnetics}
\citep{jin2014} derives the static, time-harmonic, and time-domain formulations
at book length and is the natural next text for a reader who wants every step
filled in; its treatment of the curl--curl equation and its gauge subtleties
(\cref{ch:bc}) is especially thorough. Monk's \emph{Finite Element Methods for
Maxwell's Equations} \citep{monk2003} supplies the rigorous functional-analytic
backing---the $\Hcurl$ spaces, the de Rham complex, well-posedness---that
\cref{ch:functionspaces,ch:derham} develop more gently. For the engineering
side---what the driven and eigenmode analyses \emph{compute} (S-parameters,
quality factors, cavity and waveguide modes) and why they matter---Pozar's
\emph{Microwave Engineering} \citep{pozar2011} is the standard reference and the
source of the resonant-cavity picture behind \cref{wex:cavity1d}.

\exercisessection
\begin{enumerate}[nosep]
  \item \textbf{The recipe, applied.} For each of the five reductions, name (a)
  the regime (value of $\partial_t$), (b) the potential introduced or primary
  field chosen, and (c) the Maxwell equation used in the final substitution.
  Present your answer as a five-row table.

  \item \textbf{Poisson with charge.} Repeat \cref{wex:plates1d} but with a
  uniform free charge density $\rho_0$ filling the gap, so the equation is
  $-\varepsilon V''=\rho_0$ with $V(0)=0$, $V(L)=V_0$. Solve for $V(x)$ and
  $E(x)$, and verify that the field is no longer uniform. What is the
  field at each plate?

  \item \textbf{Cavity overtones.} For the one-dimensional cavity of
  \cref{wex:cavity1d} with $L=\SI{5}{\centi\meter}$ and wave speed $c$ the speed
  of light, compute the three lowest resonant frequencies $f_n=\omega_n/2\pi$ in
  gigahertz. By what factor do successive modes differ in frequency?

  \item \textbf{Static as a limit.} Show that the driven system operator
  $\mat{A}(\omega)=\Kop+i\omega\Cop-\omega^2\Mop$ reduces to the magnetostatic
  stiffness $\Kop$ as $\omega\to 0$. Which physical effects do the dropped terms
  $i\omega\Cop$ and $\omega^2\Mop$ represent, and why is it consistent that they
  vanish in the static limit?

  \item \textbf{Eigenproblem versus driven, structurally.} The quadratic
  eigenproblem \cref{eq:qevp} and the driven system \cref{eq:driven-system} both
  contain the polynomial $\Kop+\lambda\Cop+\lambda^2\Mop$ (with
  $\lambda=i\omega$). Explain precisely what each problem holds fixed and what it
  solves for, and why one is an eigenvalue search while the other is a linear
  solve.

  \item \textbf{The damping term's origin.} In the lossless eigenmode derivation
  (\cref{sec:eigenmode}) the damping $\Cop$ was absent; in the driven and
  transient derivations it appeared. Trace exactly which physical effect (and
  which term in Maxwell's equations) introduces $\Cop$, and explain why a
  \emph{lossless} cavity has $\Cop=0$ but a real one does not.

  \item \textbf{The mechanical analogy.} The transient system
  $\Mop\ddot\vs+\Cop\dot\vs+\Kop\vs=\vJ(t)$ is formally identical to a damped
  mechanical oscillator. Match each of $\Mop,\Cop,\Kop$ to its mechanical
  counterpart (inertia, friction, spring), and predict qualitatively how the
  field rings down when $\Cop$ is small versus large.

  \item \textbf{A reduction of your own.} Suppose a problem is \emph{lossy and
  static}: a steady current flows through a conductor of finite conductivity
  $\sigma$ (so $\vJ=\sigma\vE$) with no time dependence. Starting from
  $\Curl\vE=\mathbf0$ and $\Div\vJ=0$, introduce a scalar potential and derive the
  governing PDE. Which of the chapter's five reductions does your result most
  resemble, and what plays the role of the coefficient $\varepsilon$?
\end{enumerate}

\chapter{From Strong to Weak: Integration by Parts}
\label{ch:weak}

% local vector-field macros not in simbook.sty (trial/test fields)
\providecommand{\vu}{\mathbf{u}}
\providecommand{\vv}{\mathbf{v}}
\providecommand{\vp}{\mathbf{p}}
\providecommand{\vz}{\mathbf{z}}

\chapterorient{The previous chapter handed us five partial differential
equations, each written in \emph{strong form}: a differential operator set
equal to a source, holding at every point of the domain. A computer cannot
discretize that equation directly---it asks too much of the unknown, and it
asks it pointwise. This chapter performs the single move that makes the finite
element method possible: multiply by a test function, integrate over the
domain, and integrate by parts. The derivative is ``spread around,'' the
smoothness demand is halved, and the equation reorganizes into a
\emph{bilinear form} on the left and a \emph{linear functional} on the right.
We do this slowly in one dimension, carefully in three, and then for the curl
operator the vector problems need. The payoff is a single template,
$a(u,v)=(Q\,\diffop u,\diffop v)$, that covers every term every Palace driver
assembles.}

\section{Why weak forms}

A partial differential equation in strong form is a demanding object. Take the
electrostatic equation from \cref{ch:reductions-pde},
\begin{equation}
  -\Div(\varepsilon\,\Grad u) = f \quad\text{in } \dom,
  \label{eq:strong-es}
\end{equation}
and read it literally. It asserts that a particular second-order differential
operator, applied to the unknown potential $u$, equals the source $f$ at
\emph{every single point} of the domain $\dom$. For that assertion even to make
sense, $u$ must be twice differentiable: we have to be able to take its
gradient, multiply by the material coefficient $\varepsilon$, take the
divergence of the result, and have a genuine function come out. If the
permittivity $\varepsilon$ jumps across a material interface---as it does in
every real capacitor, where a dielectric meets a vacuum---then $\varepsilon\Grad
u$ is discontinuous, its divergence is not an ordinary function at the
interface, and \cref{eq:strong-es} fails to hold there in any classical sense.
The strong form has painted us into a corner: the physically correct solution
is not smooth enough to satisfy the equation we wrote down.

\begin{physicsbox}
Across the boundary between two dielectrics, the physics does \emph{not} ask
the field to be smooth. It asks for two interface conditions: the tangential
electric field is continuous, and the \emph{normal} component of the flux
$\varepsilon\Grad u$ is continuous---while the normal derivative of $u$ itself
jumps. A formulation that bakes in twice-differentiability cannot represent this
kink. The weak form, by contrast, treats the interface conditions as
\emph{consequences} rather than assumptions, and asks only that $u$ have one
finite derivative. It is the formulation that matches the physics.
\end{physicsbox}

The weak form escapes the corner with a classical trick. Rather than demand that
a differential equation hold pointwise, we test it against a family of probe
functions and demand only that the \emph{averages} match. Concretely: pick a
\emph{test function} $v$, multiply the equation by it, and integrate over the
domain. If \cref{eq:strong-es} holds pointwise then certainly
\begin{equation}
  \int_\dom \bigl(-\Div(\varepsilon\Grad u)\bigr)\,v \dV
  = \int_\dom f\,v \dV
  \label{eq:weighted-residual}
\end{equation}
holds for every $v$. So far we have only lost information---we replaced a
pointwise statement with a family of integral statements. The decisive step
comes next: integrate the left side by parts. One derivative moves off $u$ and
onto $v$. The second derivative on $u$ disappears, the requirement that $u$ be
twice differentiable evaporates, and what remains asks $u$ to have only
\emph{one} derivative---a demand the discontinuous-flux solution can meet.

\begin{figure}[t]
\centering
\begin{tikzpicture}[node distance=10mm]
  \node[stage, text width=30mm, align=center] (strong)
    {\textbf{strong form}\\[2pt]\footnotesize $-\Div(\varepsilon\Grad u)=f$\\[1pt]\scriptsize needs $u\in$ ``$C^2$''};
  \node[opbox, right=14mm of strong, text width=24mm] (test)
    {multiply by $v$,\\ integrate};
  \node[opbox, right=12mm of test, text width=26mm] (ibp)
    {integrate\\ by parts};
  \node[product, right=14mm of ibp, text width=30mm, align=center] (weak)
    {\textbf{weak form}\\[2pt]\footnotesize $a(u,v)=\ell(v)$\\[1pt]\scriptsize needs $u\in$ ``$C^1$''};
  \draw[flow] (strong) -- (test);
  \draw[flow] (test) -- (ibp);
  \draw[flow] (ibp) -- (weak);
  % derivative-count annotation underneath
  \node[font=\scriptsize\itshape, simRust, below=7mm of test.south] (d2) {2 derivatives on $u$};
  \node[font=\scriptsize\itshape, simGreen, below=7mm of weak.south] (d1) {1 on $u$, 1 on $v$};
  \draw[refedge] (test.south) -- (d2);
  \draw[refedge] (weak.south) -- (d1);
\end{tikzpicture}
\caption[Strong to weak: spreading the derivative]{The strong-to-weak pipeline.
Multiplying by a test function $v$ and integrating turns a pointwise equation
into an averaged one; integrating by parts then moves one derivative off the
unknown $u$ and onto the test function $v$. Two derivatives concentrated on $u$
become one derivative on each---the smoothness demand on the unknown is halved,
which is exactly what lets the solution have the kinks the physics requires.}
\label{fig:strongweak}
\end{figure}

There is a second reason to prefer the weak form, and it is not merely
technical: it is the form in which the physics was \emph{originally} stated.
Long before anyone wrote Poisson's equation as a differential equation,
physicists characterized equilibrium by a variational principle. A soap film
settles into the shape of least area; a hanging chain into the shape of least
potential energy; an electrostatic field into the configuration that minimizes
the stored energy $\tfrac12\int_\dom \varepsilon\,\abs{\Grad u}^2\dV$ subject to
the imposed voltages. The strong-form PDE is what you get by differentiating
that energy and setting the result to zero---the Euler--Lagrange equation of the
variational problem. The weak form sits one step \emph{before} the
differentiation: it is the statement that the energy is stationary, that no
small ``virtual'' perturbation $v$ of the field changes the energy to first
order. We will make this precise in \cref{sec:virtual-work}; for now, hold onto
the slogan.

\begin{keyidea}
A weak form replaces ``the equation holds at every point'' with ``the equation,
averaged against every test function, holds.'' Integration by parts then trades
one derivative on the unknown for one derivative on the test function. Two gains
follow at once: the solution may be \emph{less smooth} (one derivative instead
of two), and the equation takes the form of a \emph{variational} (energy)
principle---the form in which the physics naturally lives.
\end{keyidea}

The rest of this chapter executes this move three times, with increasing
generality, and then extracts the single template they share.

\section{The one-dimensional archetype}
\label{sec:1d}

We begin where every student of the subject begins: the one-dimensional
Poisson equation. It is small enough to do completely by hand, and every
structural feature of the general theory is already present in it.

Let the domain be the interval $\dom = (0,1)$, and consider the boundary-value
problem
\begin{equation}
  -u''(x) = f(x) \quad\text{for } x\in(0,1), \qquad u(0)=u(1)=0.
  \label{eq:1d-strong}
\end{equation}
Here $u$ is the unknown (think of it as a temperature along an insulated rod, or
the deflection of a stretched string under a transverse load $f$), and the
boundary conditions pin both ends to zero. This is the strong form: it asks for
a function $u$, twice differentiable, whose second derivative equals $-f$
everywhere.

\subsection{Multiply, integrate, integrate by parts}

Choose a \emph{test function} $v$. For now we require only two things of it: it
should be once-differentiable (so that $v'$ makes sense), and---this is the
crucial choice---it should \emph{vanish at the boundary}, $v(0)=v(1)=0$. We will
see in a moment exactly why. Multiply \cref{eq:1d-strong} by $v$ and integrate
over $(0,1)$:
\begin{equation}
  -\int_0^1 u''(x)\,v(x)\dx = \int_0^1 f(x)\,v(x)\dx.
  \label{eq:1d-weighted}
\end{equation}
Now integrate the left-hand side by parts. Recall the one-dimensional formula
$\int_0^1 p'\,q = \bigl[p\,q\bigr]_0^1 - \int_0^1 p\,q'$. Apply it with $p=u'$
(so $p'=u''$) and $q=v$:
\begin{equation}
  -\int_0^1 u''\,v\dx
  = -\Bigl[\,u'(x)\,v(x)\,\Bigr]_0^1 + \int_0^1 u'(x)\,v'(x)\dx.
  \label{eq:1d-ibp}
\end{equation}
Look at the boundary term $\bigl[u'\,v\bigr]_0^1 = u'(1)v(1) - u'(0)v(0)$. This
is exactly where our requirement $v(0)=v(1)=0$ earns its keep: \emph{both}
$v(1)$ and $v(0)$ are zero, so the entire boundary term vanishes. We chose the
test functions to die at the boundary precisely so that this term would
disappear. What survives is clean:
\begin{equation}
  \int_0^1 u'(x)\,v'(x)\dx = \int_0^1 f(x)\,v(x)\dx
  \qquad\text{for all admissible } v.
  \label{eq:1d-weak}
\end{equation}
This is the \emph{weak form} of \cref{eq:1d-strong}. Read it carefully and
notice what has changed. On the left, $u$ now appears only through its
\emph{first} derivative $u'$, never $u''$. The unknown no longer needs to be
twice differentiable; it needs one derivative, and that derivative may be
merely square-integrable. The price was symmetry: where the strong form treated
$u$ specially, the weak form treats $u$ and $v$ even-handedly---each contributes
exactly one derivative.

\begin{figure}[t]
\centering
\begin{tikzpicture}[scale=1.0]
  % axis
  \draw[->,simGray] (-0.3,0) -- (5.3,0) node[right,font=\scriptsize]{$x$};
  \draw[->,simGray] (0,-0.3) -- (0,2.4);
  \node[font=\scriptsize,simGray,below] at (0,0) {$0$};
  \node[font=\scriptsize,simGray,below] at (5,0) {$1$};
  \draw[simGray,densely dotted] (5,0) -- (5,2.2);
  % the solution u (a smooth bump)
  \draw[simBlue,very thick,domain=0:5,samples=80]
    plot (\x,{1.9*sin(180*\x/5)});
  \node[simBlue,font=\small] at (4.0,1.95) {$u(x)$};
  % a hat test function v
  \draw[simRust,thick] (0,0) -- (2.4,1.25) -- (5,0);
  \node[simRust,font=\small] at (1.9,1.45) {$v(x)$};
  % endpoints clamped
  \fill[simInk] (0,0) circle (1.6pt);
  \fill[simInk] (5,0) circle (1.6pt);
  \node[font=\scriptsize,simRust] at (3.6,0.55) {$v(0)=v(1)=0$};
\end{tikzpicture}
\caption[A test function against the unknown]{The unknown $u$ (blue) and one
admissible test function $v$ (rust), here a piecewise-linear ``hat.'' The test
function is required to vanish at both ends, $v(0)=v(1)=0$; this is what kills
the boundary term in the integration by parts. A hat like this is the prototype
of the finite-element basis functions of \cref{ch:functionspaces}: the weak form
only ever pairs $u'$ against $v'$, so $v$ needs just one (possibly jumpy)
derivative.}
\label{fig:hat}
\end{figure}

\subsection{Naming the parts: bilinear form and linear functional}

The weak form \cref{eq:1d-weak} has a structure worth naming, because the names
will persist unchanged through every generalization to come. Define
\begin{align}
  a(u,v) &\;=\; \int_0^1 u'(x)\,v'(x)\dx, \label{eq:1d-bilinear}\\
  \ell(v) &\;=\; \int_0^1 f(x)\,v(x)\dx. \label{eq:1d-linear}
\end{align}
The first, $a(u,v)$, is a \emph{bilinear form}: it takes two functions and
returns a number, and it is linear in each argument separately. Linearity in
the first slot means $a(\alpha u_1 + \beta u_2, v) = \alpha\,a(u_1,v) +
\beta\,a(u_2,v)$, and similarly in the second---both are immediate from the
linearity of the integral and the derivative. The second, $\ell(v)$, is a
\emph{linear functional}: it takes one function and returns a number, linearly.
With these names, the entire weak problem collapses to one line:
\begin{equation}
  \boxed{\;\text{find } u \text{ with } u(0)=u(1)=0 \text{ such that } a(u,v)=\ell(v) \text{ for all admissible } v.\;}
  \label{eq:1d-abstract}
\end{equation}
This abstract shape---$a(u,v) = \ell(v)$, a symmetric bilinear form equated to a
linear functional---is the form the finite element method discretizes. When in
\cref{ch:galerkin} we replace the infinite-dimensional space of admissible
functions with a finite-dimensional one spanned by hat functions like the one in
\cref{fig:hat}, the bilinear form $a$ becomes a matrix, the functional $\ell$
becomes a vector, and \cref{eq:1d-abstract} becomes a linear system $\mat{K}\vx
= \vb$. Everything downstream is bookkeeping on $a$ and $\ell$.

\begin{notationbox}
We will write the bilinear form as an inner product whenever it is convenient:
$a(u,v) = \inner{u'}{v'}$, where $\inner{\cdot}{\cdot}$ is the $\Ltwo(\dom)$
inner product $\inner{p}{q} = \int_\dom p\,q$. The two notations
$\int_0^1 u'v'\dx$ and $\inner{u'}{v'}$ mean the same thing; the inner-product
notation makes the structure ``differential operator applied to each argument,
paired in $\Ltwo$'' visible at a glance, and it is the notation in which the
unifying family of \cref{sec:family} is cleanest.
\end{notationbox}

\subsection{A complete worked example}

Let us carry the 1-D problem all the way to a number, so that nothing is
abstract.

\begin{workedexample}{The 1-D weak form with a concrete load}{1dfull}
Solve, in weak form, the problem
\[
  -u''(x) = \pi^2\sin(\pi x) \quad\text{on } (0,1), \qquad u(0)=u(1)=0,
\]
and verify the weak statement on a specific test function.

\medskip\noindent\textbf{Strong solution (for checking).}
The strong problem has the exact solution $u(x)=\sin(\pi x)$: indeed
$u''=-\pi^2\sin(\pi x)$, so $-u'' = \pi^2\sin(\pi x)=f$, and $u(0)=u(1)=0$. We
keep this in our pocket to check the weak form against.

\medskip\noindent\textbf{The bilinear form and functional.}
With $f(x)=\pi^2\sin(\pi x)$ the weak problem is: find $u$ with $u(0)=u(1)=0$ such
that, for every admissible $v$,
\[
  a(u,v)=\int_0^1 u'(x)\,v'(x)\dx
  \;=\;
  \ell(v)=\int_0^1 \pi^2\sin(\pi x)\,v(x)\dx.
\]

\medskip\noindent\textbf{Test it on $v(x)=\sin(\pi x)$.}
This $v$ is admissible: it is smooth and vanishes at both ends. With the candidate
solution $u(x)=\sin(\pi x)$ we have $u'(x)=\pi\cos(\pi x)$ and
$v'(x)=\pi\cos(\pi x)$, so the left side is
\[
  a(u,v)=\int_0^1 \pi^2\cos^2(\pi x)\dx
  = \pi^2\cdot\tfrac12 = \tfrac{\pi^2}{2},
\]
using $\int_0^1\cos^2(\pi x)\dx=\tfrac12$. The right side is
\[
  \ell(v)=\int_0^1 \pi^2\sin^2(\pi x)\dx
  = \pi^2\cdot\tfrac12 = \tfrac{\pi^2}{2},
\]
using $\int_0^1\sin^2(\pi x)\dx=\tfrac12$. The two agree: $a(u,v)=\ell(v)=\pi^2/2$.
The weak equation is satisfied on this test function, as it must be for the true
solution.

\medskip\noindent\textbf{Where the boundary term went.}
The integration by parts produced the term $-[u'v]_0^1 = -\bigl(u'(1)v(1) -
u'(0)v(0)\bigr)$. Here $u'(x)=\pi\cos(\pi x)$, so $u'(1)=-\pi$ and $u'(0)=\pi$ are
both \emph{nonzero}; the boundary term vanished entirely because $v(1)=v(0)=0$,
not because the derivative of $u$ was small. This is the general mechanism: the
homogeneous Dirichlet condition on the \emph{test} space, not any property of the
solution, is what discards the boundary contribution.
\end{workedexample}

\subsection{What changes at a Neumann end}
\label{sec:1d-neumann}

The Dirichlet condition $u(0)=u(1)=0$ in \cref{eq:1d-strong} entered the weak
form in two distinct ways, and it is worth separating them, because the
distinction is the seed of the whole essential-versus-natural story of
\cref{sec:essential-natural}. First, the condition on $u$ itself restricts the
\emph{trial space}: we look for the solution among functions that already vanish
at the ends. Second, the matching condition on $v$ restricts the \emph{test
space}, and that is what annihilated the boundary term. The boundary condition on
$u$ is enforced by \emph{constraining the space}; the boundary term, meanwhile,
is a separate object that we made disappear by a choice on $v$.

Now change one end. Suppose instead of $u(1)=0$ we are told the \emph{derivative}
there: $u'(1) = g$ for some given flux $g$ (a Neumann, or natural, condition---
physically, a prescribed heat flux or surface charge), while keeping $u(0)=0$.
The strong problem is now
\[
  -u'' = f \text{ on } (0,1), \qquad u(0)=0, \quad u'(1)=g.
\]
Repeat the derivation. We multiply by $v$ and integrate by parts exactly as
before, arriving again at \cref{eq:1d-ibp}. But now we can no longer ask $v(1)=0$:
the value of $u$ at $x=1$ is \emph{not} prescribed, so the trial space does not
constrain it, and correspondingly we must leave the test functions free there.
We retain $v(0)=0$ (the left end is still Dirichlet). The boundary term is
therefore
\[
  -\Bigl[u'v\Bigr]_0^1 = -\bigl(u'(1)v(1) - u'(0)v(0)\bigr)
  = -u'(1)\,v(1) = -g\,v(1),
\]
because $v(0)=0$ kills the left contribution while $u'(1)=g$ is \emph{known}. The
weak form becomes
\begin{equation}
  \underbrace{\int_0^1 u'v'\dx}_{a(u,v)}
  \;=\;
  \underbrace{\int_0^1 f\,v\dx + g\,v(1)}_{\ell(v)}.
  \label{eq:1d-neumann-weak}
\end{equation}
Two things happened, and both are general. The Dirichlet end ($x=0$) was
enforced by constraining the space and produced nothing in the equation---it is
an \emph{essential} condition. The Neumann end ($x=1$) was \emph{not} imposed on
the space at all; instead the prescribed flux $g$ rode in through the surviving
boundary term and landed in the functional $\ell$ as a source---it is a
\emph{natural} condition. A homogeneous Neumann condition ($g=0$) would land
nothing and need no enforcement: it is the condition you get ``for free'' by
doing nothing, which is why the weak form is sometimes said to satisfy it
naturally. Keep these two roles in view; \cref{sec:essential-natural} is nothing
but their three-dimensional restatement.

\section{The three-dimensional scalar case: electrostatics}
\label{sec:3d-scalar}

We now lift the 1-D archetype to three dimensions. The structure will be
identical; only the integration-by-parts identity changes, from the
one-dimensional formula to the divergence theorem. This is the case that
governs the electrostatic driver---the capacitance analysis of
\cref{ch:targets}---so we work it concretely.

The strong form is \cref{eq:strong-es},
\[
  -\Div(\varepsilon\,\Grad u) = f \quad\text{in } \dom,
\]
with $\dom\subset\Real^3$ a bounded domain, $\varepsilon$ the (possibly
position-dependent, possibly matrix-valued) permittivity, $u$ the electrostatic
potential, and $f$ a charge density. Following the 1-D recipe to the letter:
choose a scalar test function $v$, multiply, and integrate over $\dom$:
\begin{equation}
  -\int_\dom \Div(\varepsilon\Grad u)\,v\dV = \int_\dom f\,v\dV.
  \label{eq:3d-weighted}
\end{equation}

\subsection{Green's first identity}

The three-dimensional analogue of integration by parts is \emph{Green's first
identity}, and it is itself a one-line consequence of the divergence theorem. We
derive it so that nothing is taken on faith. Start from the product rule for the
divergence of a scalar times a vector field. For a scalar field $v$ and a vector
field $\vF$,
\begin{equation}
  \Div(v\,\vF) = \Grad v\cdot\vF + v\,\Div\vF.
  \label{eq:prodrule}
\end{equation}
Integrate this over $\dom$ and apply the divergence theorem
$\int_\dom\Div(\,\cdot\,)\dV=\oint_\bdry(\,\cdot\,)\cdot\vn\,\mathrm{d}S$ to the
left side, where $\vn$ is the outward unit normal on the boundary $\bdry$:
\[
  \oint_\bdry v\,\vF\cdot\vn\,\mathrm{d}S
  = \int_\dom \Grad v\cdot\vF\dV + \int_\dom v\,\Div\vF\dV.
\]
Rearranging,
\begin{equation}
  \int_\dom v\,\Div\vF\dV
  = -\int_\dom \Grad v\cdot\vF\dV
  + \oint_\bdry v\,(\vF\cdot\vn)\,\mathrm{d}S.
  \label{eq:green1}
\end{equation}
This is Green's first identity. Compare it term-for-term with the 1-D formula
\cref{eq:1d-ibp}: $\int v\,\Div\vF$ plays the role of $\int u''v$, the
volume term $\int\Grad v\cdot\vF$ is the new $\int u'v'$, and the surface
integral $\oint_\bdry$ is the boundary term $[u'v]_0^1$. The divergence theorem
\emph{is} integration by parts in three dimensions; the only genuine change is
that the ``boundary'' of an interval (two points) has become the boundary of a
solid (a surface).

\subsection{Applying it to the electrostatic problem}

Set $\vF = \varepsilon\Grad u$ in Green's identity \cref{eq:green1}. Then
$\Div\vF = \Div(\varepsilon\Grad u)$, and the identity reads
\[
  \int_\dom v\,\Div(\varepsilon\Grad u)\dV
  = -\int_\dom \Grad v\cdot(\varepsilon\Grad u)\dV
  + \oint_\bdry v\,(\varepsilon\Grad u)\cdot\vn\,\mathrm{d}S.
\]
Substitute this into the left side of the weighted residual
\cref{eq:3d-weighted} (which carries an overall minus sign):
\begin{equation}
  \int_\dom \varepsilon\,\Grad u\cdot\Grad v\dV
  \;-\;
  \oint_\bdry v\,(\varepsilon\Grad u)\cdot\vn\,\mathrm{d}S
  \;=\;
  \int_\dom f\,v\dV.
  \label{eq:3d-with-boundary}
\end{equation}
There is the volume bilinear form we were after, $\int_\dom\varepsilon\Grad
u\cdot\Grad v$, sitting on the left exactly as $\int u'v'$ did in one dimension.
And there is the boundary integral, the three-dimensional descendant of the
$[u'v]_0^1$ term, carrying the normal flux $(\varepsilon\Grad u)\cdot\vn$ across
$\bdry$. What we do with that surface integral---make it vanish, or keep it as a
source---is precisely the essential/natural dichotomy, and it is important
enough to get its own section.

\subsection{Splitting the boundary: essential and natural}
\label{sec:essential-natural}

Partition the boundary into two disjoint pieces,
\[
  \bdry = \Gamma_D \cup \Gamma_N, \qquad \Gamma_D\cap\Gamma_N=\emptyset,
\]
where on $\Gamma_D$ (the \emph{Dirichlet} part) the potential is prescribed,
$u = u_0$, and on $\Gamma_N$ (the \emph{Neumann} part) the normal flux is
prescribed, $(\varepsilon\Grad u)\cdot\vn = g$. Physically, $\Gamma_D$ is a
conductor held at a fixed voltage and $\Gamma_N$ is a surface with a given
surface charge (with $g=0$ the natural ``no-flux'' wall).

\begin{figure}[t]
\centering
\begin{tikzpicture}[scale=1.0]
  % a blobby domain
  \fill[simBgBlue]
    (0,0) .. controls (0.4,1.6) and (2.2,2.0) .. (3.6,1.4)
          .. controls (4.8,0.95) and (4.6,-0.6) .. (3.2,-1.0)
          .. controls (1.8,-1.4) and (-0.4,-1.0) .. (0,0) -- cycle;
  % Dirichlet portion (thick blue arc, lower-left)
  \draw[simBlue,line width=2.2pt]
    (0,0) .. controls (1.8,-1.4) and (3.2,-1.0) .. (3.2,-1.0);
  % Neumann portion (thick green arc, upper)
  \draw[simGreen,line width=2.2pt]
    (0,0) .. controls (0.4,1.6) and (2.2,2.0) .. (3.6,1.4)
          .. controls (4.8,0.95) and (4.6,-0.6) .. (3.2,-1.0);
  \node[simInk,font=\large] at (2.1,0.3) {$\dom$};
  \node[simBlue,font=\small] at (1.4,-1.5) {$\Gamma_D:\ u=u_0$};
  \node[simGreen,font=\small] at (4.7,1.55) {$\Gamma_N:\ (\varepsilon\Grad u)\cdot\vn=g$};
  % outward normal on Neumann part
  \draw[-{Stealth[length=2mm]},simGreen,thick] (3.55,1.42) -- (4.25,1.85)
     node[right,font=\scriptsize]{$\vn$};
  % test function vanishes on Dirichlet
  \node[simBlue,font=\scriptsize\itshape] at (1.5,-0.75) {$v=0$ here};
\end{tikzpicture}
\caption[The boundary splits into Dirichlet and Neumann parts]{Green's first
identity produces a boundary integral over all of $\bdry$. Splitting
$\bdry=\Gamma_D\cup\Gamma_N$ disposes of it in two ways. On the Dirichlet part
$\Gamma_D$ (blue) the potential is prescribed and we constrain the test space to
vanish, $v=0$, so the integral there drops---the condition is \emph{essential}.
On the Neumann part $\Gamma_N$ (green) the normal flux $(\varepsilon\Grad
u)\cdot\vn=g$ is prescribed; the integral survives and carries $g$ into the
right-hand side as a surface source---the condition is \emph{natural}.}
\label{fig:greenbdry}
\end{figure}

Now dispose of the boundary integral $\oint_\bdry v(\varepsilon\Grad
u)\cdot\vn\,\mathrm{d}S$ piece by piece, exactly as in the 1-D Neumann
discussion of \cref{sec:1d-neumann}.

\emph{On $\Gamma_D$}, where $u$ is prescribed, we have no information about the
flux $(\varepsilon\Grad u)\cdot\vn$---it is part of the unknown response. We
remove the difficulty the same way we did in one dimension: we demand that the
test functions vanish on $\Gamma_D$, $v|_{\Gamma_D}=0$. Then the integral over
$\Gamma_D$ is zero, regardless of the unknown flux. The Dirichlet datum $u=u_0$
is enforced separately, by constraining the trial space to functions that equal
$u_0$ on $\Gamma_D$. This is the \emph{essential} boundary condition: it lives in
the \emph{space}, not in the equation, and it leaves no trace in $a$ or $\ell$.

\emph{On $\Gamma_N$}, where the flux is prescribed, the integrand
$(\varepsilon\Grad u)\cdot\vn = g$ is \emph{known}. We keep the test functions
free there, substitute $g$, and move the resulting surface integral to the
right-hand side. It becomes a source term in the functional. This is the
\emph{natural} boundary condition: it lives in the \emph{equation}, contributing
to $\ell$.

Putting it together, the electrostatic weak form is
\begin{equation}
  \boxed{\;
  \underbrace{\int_\dom \varepsilon\,\Grad u\cdot\Grad v\dV}_{a(u,v)}
  \;=\;
  \underbrace{\int_\dom f\,v\dV \;+\; \oint_{\Gamma_N} g\,v\,\mathrm{d}S}_{\ell(v)}
  \;}
  \label{eq:3d-weak}
\end{equation}
for all test functions $v$ with $v|_{\Gamma_D}=0$, sought among trial functions
$u$ with $u|_{\Gamma_D}=u_0$. The homogeneous-Neumann case $g=0$---the most
common wall condition in a capacitance problem, where no charge sits on the
insulating outer boundary---drops the surface integral entirely and is enforced
by doing nothing at all. That is the sense in which Neumann conditions are
``natural.''

\begin{pitfall}
It is tempting to think the essential (Dirichlet) condition is the ``easy'' one
because it looks simpler in the strong form. In the weak/discrete setting it is
the \emph{harder} one to implement: it must be built into the function space, by
eliminating or constraining the degrees of freedom that sit on $\Gamma_D$
(\cref{ch:bc}). The natural (Neumann) condition, which looks like the
complicated one, requires almost nothing---a homogeneous Neumann wall is handled
by simply not assembling a boundary term. ``Essential'' means ``you must enforce
it yourself''; ``natural'' means ``the weak form enforces it for you.''
\end{pitfall}

\begin{workedexample}{The electrostatic weak form, term by term}{esweak}
Take a parallel-plate capacitor: $\dom$ the slab between two conducting plates,
the top plate held at $u=V$ and the bottom at $u=0$ (both Dirichlet, so both part
of $\Gamma_D$), the side walls insulating (homogeneous Neumann, $g=0$, so part of
$\Gamma_N$), no free charge in the dielectric ($f=0$), and a uniform scalar
permittivity $\varepsilon$.

\medskip\noindent\textbf{Identify the pieces of the template.}
Reading \cref{eq:3d-weak} against $a(u,v)=(Q\,\diffop u,\diffop v)$:
\begin{itemize}\setlength\itemsep{1pt}
  \item the coefficient is $Q=\varepsilon$, the permittivity;
  \item the differential operator is $\diffop=\Grad$, the gradient;
  \item the bilinear form is $a(u,v)=\inner{\varepsilon\Grad u}{\Grad v}_{\Ltwo(\dom)}=\int_\dom\varepsilon\,\Grad u\cdot\Grad v\dV$.
\end{itemize}

\medskip\noindent\textbf{Assemble the functional.}
With $f=0$ in the volume and $g=0$ on the insulating walls, the functional is
\[
  \ell(v) = \int_\dom 0\cdot v\dV + \oint_{\Gamma_N} 0\cdot v\,\mathrm{d}S = 0.
\]
The right-hand side is \emph{empty}. The problem is not, however, trivial: the
nonzero plate voltage $u=V$ is an \emph{inhomogeneous essential} condition, and
it enters not through $\ell$ but through the trial space (the solution must equal
$V$ on the top plate). In the discrete setting this inhomogeneous Dirichlet datum
is ``lifted'' into the right-hand side by the elimination procedure of
\cref{ch:bc}---which is exactly why an apparently source-free problem
($f=g=0$) still produces a nonzero load vector. The physics confirms it: a
charged capacitor with no free charge in the dielectric still stores energy,
driven entirely by the imposed plate voltages.

\medskip\noindent\textbf{The product.}
Solving the weak problem yields the potential $u$; the capacitance follows from
the stored energy $U=\tfrac12\int_\dom\varepsilon\abs{\Grad u}^2\dV =
\tfrac12\,a(u,u)$ via $C = 2U/V^2$. The bilinear form $a$ evaluated on the
solution against \emph{itself} is the energy---a connection we exploit in
\cref{sec:virtual-work} and again in the capacitance reduction of
\cref{ch:electrostatics}.
\end{workedexample}

\section{The vector case: curl--curl}
\label{sec:curlcurl}

The electrostatic problem lives on a \emph{scalar} field, the potential $u$, and
its natural derivative is the gradient. The magnetostatic and eigenmode problems
of \cref{ch:reductions-pde} live on \emph{vector} fields---the magnetic vector
potential $\vA$, the electric field $\vE$---and their natural derivative is the
\emph{curl}. The integration-by-parts move is structurally identical, but the
identity it rests on is the curl analogue of Green's formula, and the function
space is different. We sketch the derivation and read off the same template.

The strong form of the magnetostatic problem is the curl--curl equation
\begin{equation}
  \Curl(\nu\,\Curl\vu) = \vJ \quad\text{in }\dom,
  \label{eq:curlcurl-strong}
\end{equation}
where $\vu$ is the unknown vector field (the vector potential $\vA$), $\nu=1/\mu$
is the reluctivity (the magnetic coefficient, the analogue of $\varepsilon$), and
$\vJ$ is the source current density. Multiply by a vector test function $\vv$ and
integrate---now the multiplication is a dot product, since both factors are
vectors:
\begin{equation}
  \int_\dom \Curl(\nu\Curl\vu)\cdot\vv\dV = \int_\dom \vJ\cdot\vv\dV.
  \label{eq:curlcurl-weighted}
\end{equation}

\subsection{The curl integration-by-parts identity}

The identity we need plays the exact role Green's first identity played for the
gradient. For two vector fields $\vp$ and $\vv$,
\begin{equation}
  \int_\dom (\Curl\vp)\cdot\vv\dV
  = \int_\dom \vp\cdot(\Curl\vv)\dV
  + \oint_\bdry (\vp\times\vn)\cdot\vv\,\mathrm{d}S.
  \label{eq:curlibp}
\end{equation}
Like Green's identity, this follows from a vector product rule---$\Div(\vp\times
\vv) = (\Curl\vp)\cdot\vv - \vp\cdot(\Curl\vv)$---integrated over $\dom$ with the
divergence theorem applied to the left side; the surface term
$\oint_\bdry(\vp\times\vv)\cdot\vn\,\mathrm{d}S$ is rewritten as
$\oint_\bdry(\vp\times\vn)\cdot\vv\,\mathrm{d}S$ by the scalar triple-product
identity. The shape is the same as before: a derivative on $\vp$ moves to a
derivative on $\vv$, at the cost of a boundary integral. Apply it with
$\vp = \nu\Curl\vu$:
\[
  \int_\dom \Curl(\nu\Curl\vu)\cdot\vv\dV
  = \int_\dom (\nu\Curl\vu)\cdot(\Curl\vv)\dV
  + \oint_\bdry \bigl((\nu\Curl\vu)\times\vn\bigr)\cdot\vv\,\mathrm{d}S.
\]

\begin{figure}[t]
\centering
\begin{tikzpicture}[
    row/.style={font=\small},
    lbl/.style={font=\footnotesize\bfseries, text=simInk},
  ]
  % column headers
  \node[lbl] at (0,2.0) {dimension / operator};
  \node[lbl] at (3.7,2.0) {move a derivative\dots};
  \node[lbl] at (8.4,2.0) {\dots leaving a boundary term};
  \draw[simGray,semithick] (-2.0,1.6) -- (10.6,1.6);
  % 1-D row
  \node[row,simBlue,anchor=west] at (-2.0,1.0)
    {1-D: \ $\displaystyle\int u''\,v$};
  \node[row,anchor=west] at (2.3,1.0)
    {$\displaystyle =-\int u'\,v'$};
  \node[row,simRust,anchor=west] at (6.6,1.0)
    {$+\,\bigl[u'\,v\bigr]_0^1$};
  % 3-D gradient row
  \node[row,simBlue,anchor=west] at (-2.0,0.1)
    {$\Grad$: \ $\displaystyle\int v\,\Div\vF$};
  \node[row,anchor=west] at (2.3,0.1)
    {$\displaystyle =-\int \Grad v\cdot\vF$};
  \node[row,simRust,anchor=west] at (6.6,0.1)
    {$+\displaystyle\oint_\bdry v\,(\vF\cdot\vn)$};
  % curl row
  \node[row,simBlue,anchor=west] at (-2.0,-0.8)
    {$\Curl$: \ $\displaystyle\int(\Curl\vp)\cdot\vv$};
  \node[row,anchor=west] at (2.3,-0.8)
    {$\displaystyle =\int \vp\cdot(\Curl\vv)$};
  \node[row,simRust,anchor=west] at (6.6,-0.8)
    {$+\displaystyle\oint_\bdry(\vp\times\vn)\cdot\vv$};
  \draw[simGray,semithick] (-2.0,-1.3) -- (10.6,-1.3);
  \node[font=\scriptsize\itshape,simGray,anchor=west] at (-2.0,-1.75)
    {one move, three settings: a derivative crosses from one factor to the other, and a boundary term is born.};
\end{tikzpicture}
\caption[The three integration-by-parts identities, aligned]{The same move in
three settings. One-dimensional integration by parts, Green's first identity (for
the gradient), and the curl integration-by-parts identity all have the identical
shape: a derivative is transferred from one factor to the other (middle column),
at the cost of a boundary integral (rust). The 1-D boundary ``surface'' is two
endpoints; in three dimensions it is a genuine surface $\bdry$. Choosing test
functions that vanish (gradient case) or whose tangential part vanishes (curl
case) on the Dirichlet boundary kills that term---the one mechanism that recurs
throughout the chapter.}
\label{fig:threeibp}
\end{figure}

Substituting into the weighted residual \cref{eq:curlcurl-weighted} and, as
before, choosing test functions whose \emph{tangential} component vanishes on the
Dirichlet boundary (so the surface term drops there) gives the weak form
\begin{equation}
  \boxed{\;
  \underbrace{\int_\dom \nu\,(\Curl\vu)\cdot(\Curl\vv)\dV}_{a(\vu,\vv)}
  \;=\;
  \underbrace{\int_\dom \vJ\cdot\vv\dV \;+\;(\text{natural-BC surface term})}_{\ell(\vv)}
  \;}
  \label{eq:curlcurl-weak}
\end{equation}
for all admissible $\vv$. The natural boundary condition here prescribes the
\emph{tangential} field $\nu\Curl\vu\times\vn$ (physically, a tangential
magnetic field $\vH$); the essential condition prescribes the tangential
component of $\vu$ itself. The pattern is the gradient story with ``value on the
boundary'' replaced by ``tangential component on the boundary''---a change forced
by the geometry of the curl, treated fully in \cref{ch:functionspaces} and
\cref{ch:bc}.

\subsection{The mass term}
\label{sec:mass}

Eigenmode and driven problems pair the curl--curl stiffness term with a second,
\emph{lower-order} term that carries no derivative at all. In the eigenmode
problem $\Curl(\nu\Curl\vE)=\omega^2\varepsilon\vE$, the right side
$\omega^2\varepsilon\vE$ contributes, after multiplying by $\vv$ and integrating,
the term
\begin{equation}
  m(\vu,\vv) = \int_\dom \varepsilon\,\vu\cdot\vv\dV,
  \label{eq:mass}
\end{equation}
with \emph{no} integration by parts---there is no derivative to move. This is the
\emph{mass} term (a name inherited from the analogous term in structural
dynamics, where the coefficient is literally a mass density). It is the
``identity'' member of the family: the differential operator applied to $\vu$ and
$\vv$ is the identity, $\diffop\vu=\vu$, and the pairing is just the
$Q$-weighted $\Ltwo$ inner product. The eigenmode weak form is then the
\emph{pencil} $a(\vu,\vv)=\omega^2\,m(\vu,\vv)$, which discretizes to the
generalized eigenvalue problem $\mat{K}\vx=\lambda\mat{M}\vx$ of
\cref{ch:targets}---the stiffness $\mat{K}$ from the curl--curl term, the mass
$\mat{M}$ from \cref{eq:mass}.

That the curl--curl term and the mass term are \emph{the same kind of object}---a
coefficient-weighted $\Ltwo$ pairing of a differential operator applied to both
arguments, differing only in \emph{which} operator---is the observation the next
section turns into the chapter's payoff.

\section{The unifying bilinear-form family}
\label{sec:family}

Lay the three bilinear forms we have derived side by side, written in the
inner-product notation of the notation box:
\begin{align*}
  \text{electrostatic (gradient):}\quad & a(u,v) = \inner{\varepsilon\,\Grad u}{\Grad v}_{\Ltwo(\dom)},\\
  \text{magnetostatic (curl):}\quad & a(\vu,\vv) = \inner{\nu\,\Curl\vu}{\Curl\vv}_{\Ltwo(\dom)},\\
  \text{mass (identity):}\quad & m(\vu,\vv) = \inner{\varepsilon\,\vu}{\vv}_{\Ltwo(\dom)}.
\end{align*}
The three are visibly instances of one shape. In each, a differential operator
is applied to both the trial and the test function; the two results are paired in
the $\Ltwo(\dom)$ inner product; and a material coefficient $Q$ weights the
pairing. Only \emph{which differential operator} changes: $\Grad$, then $\Curl$,
then the identity. Naming the differential operator $\diffop$ and the coefficient
$Q$, all three collapse into a single template.

\begin{keyidea}
\textbf{The bilinear-form family.}
Every domain bilinear form the simulator assembles has the form
\begin{equation}
  a(u,v) \;=\; \bigl(Q\,\diffop u,\ \diffop v\bigr)_{\Ltwo(\dom)}
            \;=\; \int_\dom (Q\,\diffop u)\cdot(\diffop v)\dV,
  \label{eq:family}
\end{equation}
where $\diffop\in\{\,\text{value},\ \Grad,\ \Curl,\ \Div\,\}$ is a differential
operator applied to both trial and test functions, and $Q$ is a material
coefficient (scalar or matrix-valued). A weak-form problem is built by choosing,
for each term, a pair $(Q,\diffop)$. The four choices of $\diffop$ are:
\[
\begin{array}{llll}
  \diffop=\text{value} & a(u,v)=(Q\,u,\,v) & \text{mass / }\Ltwo\text{ pairing} & \text{any space}\\[2pt]
  \diffop=\Grad        & a(u,v)=(Q\,\Grad u,\,\Grad v) & \text{electrostatic (diffusion)} & \Hone\\[2pt]
  \diffop=\Curl        & a(u,v)=(Q\,\Curl u,\,\Curl v) & \text{magnetostatic / eigenmode} & \Hcurl\\[2pt]
  \diffop=\Div         & a(u,v)=(Q\,\Div u,\,\Div v) & \text{div--div} & \Hdiv
\end{array}
\]
\end{keyidea}

This template is the load-bearing abstraction of the entire field-theory part of
the book, and it is worth pausing on what it buys us. A weak-form term is now a
\emph{piece of data}: a pair $(\text{coefficient } Q,\ \text{differential
operator } \diffop)$. To specify the physics of a problem is to write down a
short list of such pairs. The electrostatic stiffness is the single term
$(\varepsilon,\Grad)$; the magnetostatic stiffness is $(\nu,\Curl)$; the
eigenmode problem is two terms, $(\nu,\Curl)$ for the stiffness and
$(\varepsilon,\text{value})$ for the mass. The machinery that turns such a pair
into an actual matrix---element quadrature, basis evaluation, assembly---is
\emph{completely indifferent} to which pair it was handed; it is the same fold
over terms in every case, the subject of \cref{ch:galerkin} and
\cref{ch:assembly}.

\begin{figure}[t]
\centering
\setlength{\tabcolsep}{0pt}
\renewcommand{\arraystretch}{1.0}
\begin{tikzpicture}[
    cell/.style={draw=simGray, semithick, minimum width=33mm, minimum height=20mm,
      align=center, font=\footnotesize, inner sep=2pt},
    hdr/.style={font=\small\bfseries, text=simInk},
  ]
  % four cells in a row
  \node[cell, fill=simBgGray] (val) at (0,0)
    {$\diffop=\text{value}$\\[3pt]$(Q\,u,\;v)$\\[3pt]\textit{mass}\\[1pt]\scriptsize any space};
  \node[cell, fill=simBgBlue, right=2mm of val] (grad)
    {$\diffop=\Grad$\\[3pt]$(Q\,\Grad u,\;\Grad v)$\\[3pt]\textit{electrostatic}\\[1pt]\scriptsize $\Hone$};
  \node[cell, fill=simBgRust, right=2mm of grad] (curl)
    {$\diffop=\Curl$\\[3pt]$(Q\,\Curl u,\;\Curl v)$\\[3pt]\textit{magnetostatic}\\[1pt]\scriptsize $\Hcurl$};
  \node[cell, fill=simBgGreen, right=2mm of curl] (div)
    {$\diffop=\Div$\\[3pt]$(Q\,\Div u,\;\Div v)$\\[3pt]\textit{div--div}\\[1pt]\scriptsize $\Hdiv$};
  % unifying brace above
  \draw[decorate,decoration={brace,amplitude=5pt},simInk]
    ($(val.north west)+(0,3mm)$) -- ($(div.north east)+(0,3mm)$)
    node[midway,above=5pt,font=\small\bfseries]{$a(u,v)=(Q\,\diffop u,\;\diffop v)_{\Ltwo(\dom)}$};
  % de Rham hint below
  \node[font=\scriptsize\itshape,simGray] at ($(val.south)!0.5!(div.south)+(0,-6mm)$)
    {one template, four differential operators --- the de Rham ladder of \cref{ch:derham}};
\end{tikzpicture}
\caption[The four-operator bilinear-form family]{The bilinear-form family as a
labeled grid. A single template $a(u,v)=(Q\,\diffop u,\diffop v)_{\Ltwo(\dom)}$
(brace) specializes to four terms by the choice of differential operator
$\diffop$. The value (mass) term carries no derivative; gradient, curl, and
divergence each carry one, and each is well-formed on a different
finite-element space---$\Hone$, $\Hcurl$, $\Hdiv$---the slots of the de Rham
complex of \cref{ch:derham}. The coefficient $Q$ rides every cell uniformly.}
\label{fig:family}
\end{figure}

\begin{remark}[How Palace records a term]
This is not an abstraction imposed from outside; it is the structure Palace's own
code already has. The simulator builds an operator by pushing
\emph{integrator} objects onto a list, each integrator carrying a coefficient and
a differential-operator type: the diffusion integrator for $\diffop=\Grad$, the
curl--curl integrator for $\diffop=\Curl$, the mass integrator for the value
term. Every integrator inherits the same coefficient slot from a shared base
class, which is exactly the statement that $Q$ is variant-invariant while
$\diffop$ is the variant axis. The pair $(Q,\diffop)$ is, almost literally, a row
in Palace's integrator list. We pick up this thread, and the fold that consumes
the list, in \cref{ch:assembly}.
\end{remark}

A fourth operator, $\diffop=\Div$, completes the family on grounds of symmetry:
it is the div--div term, well-formed on the $\Hdiv$ space, and it fills the last
slot of the de Rham ladder $\Hone\xrightarrow{\Grad}\Hcurl\xrightarrow{\Curl}
\Hdiv\xrightarrow{\Div}\Ltwo$ that organizes \cref{ch:derham}. Among the five
analyses of this book it happens that no driver assembles a div--div stiffness
term, so it appears here as the symmetric completion of the family rather than as
a witnessed workhorse---but the template is the more compelling for being closed
under the full ladder of differential operators.

\section{The weak form as virtual work}
\label{sec:virtual-work}

We close the conceptual arc by returning to the energy principle promised in
\cref{sec:virtual-work}'s forward reference at the start of the chapter. The
bilinear-form family is not just a computational convenience; when $a$ is
symmetric and positive---as the electrostatic and magnetostatic stiffness forms
are---the weak problem $a(u,v)=\ell(v)$ is exactly the condition for a certain
energy to be minimized.

Define the energy functional
\begin{equation}
  J(w) \;=\; \tfrac12\,a(w,w) \;-\; \ell(w).
  \label{eq:energy}
\end{equation}
For the electrostatic problem, $\tfrac12 a(w,w)=\tfrac12\int_\dom
\varepsilon\abs{\Grad w}^2\dV$ is the electrostatic field energy stored by the
trial potential $w$, and $\ell(w)$ is the work done by the sources; $J$ is the
total potential energy. The claim is that the solution $u$ of the weak problem is
precisely the minimizer of $J$.

\begin{proposition}[Variational equivalence]
\label{prop:varequiv}
Let $a$ be a symmetric bilinear form ($a(w,z)=a(z,w)$) that is positive
($a(w,w)>0$ for $w\neq0$), and let $\ell$ be linear. Then $u$ solves the weak
problem $a(u,v)=\ell(v)$ for all $v$ \emph{if and only if} $u$ minimizes
$J(w)=\tfrac12 a(w,w)-\ell(w)$.
\end{proposition}

\begin{proof}
Perturb the candidate $u$ by an arbitrary admissible variation $v$, scaled by a
real parameter $t$, and expand $J(u+tv)$ using bilinearity and symmetry of $a$:
\[
  J(u+tv) = \tfrac12 a(u,u) + t\,a(u,v) + \tfrac{t^2}{2}a(v,v) - \ell(u) - t\,\ell(v).
\]
Group by powers of $t$:
\[
  J(u+tv) = J(u) + t\bigl(a(u,v)-\ell(v)\bigr) + \tfrac{t^2}{2}a(v,v).
\]
The coefficient of the \emph{linear} term in $t$ is exactly the weak-form
residual $a(u,v)-\ell(v)$. The quadratic coefficient $\tfrac12 a(v,v)$ is
positive (for $v\neq0$) by hypothesis, so as a function of $t$, $J(u+tv)$ is an
upward parabola. Its vertex sits at $t=0$---i.e.\ $u$ is the minimizer along the
direction $v$---\emph{if and only if} the linear coefficient vanishes,
$a(u,v)=\ell(v)$. Since $v$ was arbitrary, $u$ minimizes $J$ over all directions
exactly when it solves the weak problem.
\end{proof}

\begin{figure}[t]
\centering
\begin{tikzpicture}[scale=1.0]
  \draw[->,simGray] (-2.6,0) -- (2.8,0) node[right,font=\scriptsize]{$t$ (amount of perturbation $v$)};
  \draw[->,simGray] (0,-0.3) -- (0,2.7) node[above,font=\scriptsize]{$J(u+tv)$};
  % the parabola, vertex at t=0
  \draw[simBlue,very thick,domain=-2.3:2.3,samples=80]
    plot (\x,{0.42*\x*\x+0.35});
  % vertex marker
  \fill[simRust] (0,0.35) circle (2pt);
  \node[simRust,font=\small,below right] at (0,0.35) {$u$ (minimizer)};
  % tangent line at vertex (horizontal) = zero linear term
  \draw[simGreen,thick,densely dashed] (-1.6,0.35) -- (1.6,0.35);
  \node[simGreen,font=\scriptsize] at (1.7,0.2) {slope $=a(u,v)-\ell(v)=0$};
  % quadratic-term annotation
  \node[simBlue,font=\scriptsize] at (1.9,2.2) {curvature $=a(v,v)>0$};
\end{tikzpicture}
\caption[The weak form as energy minimization]{The energy
$J(u+tv)=J(u)+t\,(a(u,v)-\ell(v))+\tfrac{t^2}{2}a(v,v)$ along any perturbation
direction $v$ is an upward parabola (positive curvature $a(v,v)>0$). Its vertex
falls at $t=0$---meaning $u$ is the energy minimizer---precisely when the linear
slope $a(u,v)-\ell(v)$ vanishes, which is the weak equation. ``Solve the weak
form'' and ``minimize the energy against every virtual displacement'' are the
same statement; the weak form is the principle of virtual work.}
\label{fig:virtualwork}
\end{figure}

\Cref{prop:varequiv} is the rigorous form of the slogan we opened with. The
linear term $a(u,v)-\ell(v)$ is the first-order change in energy under a
``virtual'' perturbation $v$; demanding that it vanish for every $v$ is the
\emph{principle of virtual work}, and it is identically the weak form
(\cref{fig:virtualwork}). When $a$ is symmetric and positive, solving the weak
form is minimizing an energy---the form in which a physicist would have posed the
problem in the first place. (The driven and eigenmode problems are not energy
minimizations in this simple sense---their forms are indefinite or
complex---but the weak form survives intact as the master formulation; it is only
the energy \emph{interpretation} that is special to the symmetric-positive case.)

\section{From weak form to finite elements: the road ahead}

We have, in this chapter, taken five strong-form PDEs and reduced each to the
same abstract shape: \emph{find $u$ such that $a(u,v)=\ell(v)$ for all admissible
$v$}, with $a$ drawn from the one-parameter family $a(u,v)=(Q\,\diffop u,\diffop
v)$. Two large questions remain, and they organize the next several chapters.

First, what is the space of ``admissible'' functions, and how do we replace it
with a finite-dimensional one we can compute in? The weak form demands only one
derivative, which is exactly the regularity the piecewise-polynomial
finite-element spaces provide; \emph{which} space ($\Hone$, $\Hcurl$, $\Hdiv$)
goes with which differential operator is the content of \cref{ch:galerkin} and
\cref{ch:functionspaces}, and the way the four spaces interlock is the de Rham
complex of \cref{ch:derham}. Second, how does the bilinear form become a matrix?
Restricting $a$ to a finite-dimensional space spanned by basis functions turns
$a(u,v)=\ell(v)$ into a linear system $\mat{K}\vx=\vb$, assembled element by
element as a fold over the weak-form terms---the Galerkin method of
\cref{ch:galerkin} and the assembly machinery of \cref{ch:assembly}. The
essential/natural boundary distinction we drew here is implemented, in that
discrete setting, by the degree-of-freedom elimination of \cref{ch:bc}.

\summarysection
The strong form of a PDE asks too much: it demands a twice-differentiable
solution and enforces the equation pointwise. Multiplying by a test function $v$,
integrating over the domain, and integrating by parts trades one derivative on
the unknown $u$ for one on $v$, halving the smoothness demand and reorganizing
the equation into a \emph{bilinear form} $a(u,v)$ equal to a \emph{linear
functional} $\ell(v)$. In one dimension the integration by parts is the ordinary
formula and produces $\int u'v'=\int fv$; in three dimensions the divergence
theorem (Green's first identity) produces $a(u,v)=\int_\dom\varepsilon\Grad
u\cdot\Grad v$ plus a boundary integral; for vector fields the curl
integration-by-parts identity produces $a(\vu,\vv)=\int_\dom\nu\,\Curl\vu\cdot
\Curl\vv$, paired in eigenmode problems with the mass term
$\int_\dom\varepsilon\,\vu\cdot\vv$. The boundary integral splits into an
\emph{essential} (Dirichlet) part---enforced by constraining the test/trial
space, leaving no trace in the equation---and a \emph{natural} (Neumann) part,
which survives as a source in $\ell$. All of these are instances of one template,
$a(u,v)=(Q\,\diffop u,\diffop v)_{\Ltwo(\dom)}$, with $\diffop\in\{\text{value},
\Grad,\Curl,\Div\}$ and $Q$ a material coefficient---the four-term family that
every Palace driver assembles. When $a$ is symmetric and positive, the weak form
is equivalent to minimizing the energy $J=\tfrac12 a(u,u)-\ell(u)$: the principle
of virtual work.

\furtherreading
The classical reference for the variational formulation and its analysis is
Brenner and Scott~\citep{brennerscott2008}, which develops weak forms, Sobolev
regularity, and the Lax--Milgram and C\'ea theorems we have only gestured at.
For the curl--curl and vector cases specific to electromagnetics---Green's
identities for the curl operator, the $\Hcurl$ and $\Hdiv$ spaces, and the
boundary conditions of Maxwell problems---Monk~\citep{monk2003} is the standard
graduate treatment. Jin~\citep{jin2014} gives the same material from a
practicing computational-electromagnetics viewpoint, with the bilinear forms
written exactly as the integrators that appear in a solver. Readers wanting the
energy/virtual-work perspective in its mechanical home will find it in any
finite-element text rooted in structural analysis; the equivalence of
\cref{prop:varequiv} is there called the principle of minimum potential energy.

\exercisessection
\begin{enumerate}[nosep]
  \item \textbf{(1-D, from scratch.)} Derive the weak form of $-u''+u=f$ on
  $(0,1)$ with $u(0)=u(1)=0$. Identify the bilinear form $a(u,v)$ and the
  functional $\ell(v)$. How does $a$ differ from the pure-Poisson case
  \cref{eq:1d-bilinear}, and which two members of the family
  \cref{eq:family} does it combine?
  \item \textbf{(Boundary bookkeeping.)} Redo the 1-D derivation of
  \cref{sec:1d} but with the \emph{inhomogeneous} Dirichlet condition $u(0)=a$,
  $u(1)=b$. Show that the bilinear form is unchanged and explain where the data
  $a,b$ must enter instead. (This is the ``lifting'' of \cref{ch:bc}, in
  miniature.)
  \item \textbf{(Why $v$ vanishes.)} In the 1-D Neumann derivation
  \cref{eq:1d-neumann-weak} we kept $v(1)$ free but required $v(0)=0$. What weak
  form results if you mistakenly require $v(1)=0$ as well? Show that the
  prescribed flux $g$ then disappears from the equation entirely, and explain in
  one sentence why that makes the discrete problem wrong.
  \item \textbf{(Green's identity by hand.)} Starting from the product rule
  \cref{eq:prodrule}, derive Green's \emph{second} identity
  $\int_\dom (u\Lap v - v\Lap u)\dV = \oint_\bdry(u\,\partial_n v - v\,\partial_n
  u)\,\mathrm{d}S$ by applying Green's first identity twice (with the roles of
  $u$ and $v$ swapped) and subtracting. Here $\partial_n=\vn\cdot\Grad$.
  \item \textbf{(Reading the template.)} For each of the following bilinear
  forms, identify the coefficient $Q$, the differential operator $\diffop$, and
  the finite-element space on which it is well-formed: (a)
  $\int_\dom \mu^{-1}\Curl\vu\cdot\Curl\vv$; (b) $\int_\dom \rho\,u\,v$; (c)
  $\int_\dom \kappa\,\Grad T\cdot\Grad S$ (steady heat conduction). Which member
  of \cref{fig:family} is each?
  \item \textbf{(Symmetry and the energy.)} Show directly from the definition
  \cref{eq:family} that $a(u,v)=a(v,u)$ whenever the coefficient $Q$ is a
  symmetric matrix (or a scalar). Then exhibit a $2\times2$ matrix coefficient
  $Q$ for which $a(u,v)\neq a(v,u)$, and explain why \cref{prop:varequiv}'s
  energy interpretation fails for that $Q$.
  \item \textbf{(Eigenmode pencil.)} The eigenmode weak form is
  $a(\vu,\vv)=\omega^2 m(\vu,\vv)$ with $a$ the curl--curl form and $m$ the mass
  form \cref{eq:mass}. Writing both as members of the family
  \cref{eq:family}, give the pair $(Q,\diffop)$ for each. Discretized, this
  becomes $\mat{K}\vx=\lambda\mat{M}\vx$ with $\lambda=\omega^2$; which form
  produces $\mat{K}$ and which produces $\mat{M}$?
  \item \textbf{(Virtual work numerically.)} For the 1-D problem of
  \cref{wex:1dfull} (with exact solution $u=\sin\pi x$), compute the energy
  $J(u)=\tfrac12 a(u,u)-\ell(u)$ at the true solution. Then compute $J(w)$ for
  the perturbed trial function $w(x)=\sin\pi x + \tfrac1{10}\sin 2\pi x$ and
  confirm $J(w)>J(u)$, consistent with \cref{prop:varequiv}. (You may use
  $\int_0^1\sin m\pi x\,\sin n\pi x\dx=\tfrac12\delta_{mn}$ and the matching
  cosine identity.)
\end{enumerate}

\chapter{Galerkin's Method and the Bilinear-Form Family}
\label{ch:galerkin}

\chapterorient{In \cref{ch:weak} we turned a partial differential equation into
a \emph{weak form}: find a function $u$ such that $a(u,v)=\ell(v)$ for every test
function $v$. That statement is exact, but it lives in an infinite-dimensional
space of functions---there is no way to hand it to a computer as written. This
chapter performs the single decisive step that makes the weak form
\emph{computable}: replace the infinite-dimensional space with a finite one
spanned by $N$ chosen functions, and the weak form collapses into a system of
$N$ linear equations in $N$ unknowns, $\mat{K}\vx=\vb$. We will see where the
matrices $\mat{K}$ and $\mat{M}$ that the rest of the book solves actually
\emph{come from}---they are nothing but the bilinear form evaluated on pairs of
basis functions---and why the properties an engineer cares about (symmetry,
sparsity, positive-definiteness) are inherited directly from the form. By the
end you will have assembled, by hand, the classic tridiagonal stiffness matrix
of the one-dimensional Poisson equation.}

\section{From an infinite problem to a finite one}

Recall the shape of the weak form. We have a real number $a(u,v)$ that depends
linearly on each of two functions $u$ and $v$---a \emph{bilinear form}---and a
real number $\ell(v)$ that depends linearly on one function---a \emph{linear
functional}. The weak problem is:

\begin{quote}
find $u$ in the trial space $V$ such that $a(u,v)=\ell(v)$ for \emph{all} $v$ in
the test space $V$.
\end{quote}

For the model problem of this chapter, the one-dimensional Poisson equation
$-u''=f$ on the interval $(0,1)$ with $u(0)=u(1)=0$, integration by parts gave us
(in \cref{ch:weak}) the concrete pair
\begin{equation}
a(u,v)=\int_0^1 u'(x)\,v'(x)\dx,
\qquad
\ell(v)=\int_0^1 f(x)\,v(x)\dx.
\label{eq:poisson-weak}
\end{equation}
The space $V$ here is the set of all functions on $(0,1)$ that vanish at the two
endpoints and whose first derivative is square-integrable---an
\emph{infinite-dimensional} space. We cannot loop over ``all $v$ in $V$''; there
are uncountably many, and no finite list of equations pins $u$ down. Something
has to give.

\subsection{The Galerkin idea}

The idea, due to Boris Galerkin, is disarmingly simple: \emph{stop asking for the
exact answer in the infinite space, and instead ask for the best answer you can
build out of a finite menu of functions.} We pick, once and for all, $N$
functions
\[
\phi_1,\phi_2,\dots,\phi_N,
\]
each one a legal member of $V$ (each vanishes at the boundary, each is smooth
enough), and we agree to look for our solution only among combinations of them.
The set of all such combinations is a finite-dimensional subspace,
\begin{equation}
V_h=\operatorname{span}\{\phi_1,\dots,\phi_N\}
=\Bigl\{\,\textstyle\sum_{j=1}^N c_j\phi_j : c_j\in\Real\,\Bigr\}\subset V,
\label{eq:Vh}
\end{equation}
where the subscript $h$ is the traditional reminder that the menu came from a
\emph{mesh} of spacing $h$ (we make this concrete in \cref{sec:worked-stiffness}).
The functions $\phi_i$ are the \emph{basis functions}, and choosing them well is,
as \cref{sec:basis-is-everything} argues, the entire game.

\begin{figure}[t]
\centering
\begin{tikzpicture}[scale=1.0, every node/.style={font=\footnotesize}]
  % the big infinite space V
  \draw[draw=simGray, fill=simBgGray, rounded corners=10pt]
    (-3.4,-1.9) rectangle (3.4,2.0);
  \node[simGray] at (0,1.65) {$V$ \;(infinite-dimensional)};
  % the finite subspace V_h
  \draw[draw=simGreen, fill=simBgGreen, rounded corners=6pt, thick]
    (-2.6,-1.5) rectangle (2.6,0.7);
  \node[simGreen!70!simInk] at (0,0.42) {$V_h=\operatorname{span}\{\phi_1,\dots,\phi_N\}$};
  % the exact solution
  \fill[simRust] (1.9,1.25) circle (1.6pt);
  \node[simRust, above right=-1pt and -2pt] at (1.9,1.25) {$u$ (exact)};
  % the galerkin solution inside V_h
  \fill[simBlue] (0.55,-0.55) circle (1.6pt);
  \node[simBlue, below=1pt] at (0.55,-0.7) {$u_h$ (Galerkin)};
  % projection arrow
  \draw[flow, simInk] (1.9,1.25) -- node[right=1pt,font=\scriptsize]{project} (0.55,-0.55);
  % the residual is orthogonal to V_h: little right-angle marker idea
  \node[font=\scriptsize\itshape, simInk] at (-1.45,-1.05)
    {$a(u-u_h,\,v_h)=0\ \ \forall v_h\in V_h$};
\end{tikzpicture}
\caption[The Galerkin projection]{The Galerkin idea in one picture. The exact
solution $u$ lives somewhere in the infinite-dimensional space $V$; we cannot
reach it. Instead we look inside the finite-dimensional subspace $V_h$ and pick
the element $u_h$ whose \emph{error} $u-u_h$ is invisible to every test function
in $V_h$---formally, $a(u-u_h,v_h)=0$ for all $v_h\in V_h$. In the energy
geometry set by $a(\cdot,\cdot)$ this is an orthogonal projection, which is why
$u_h$ turns out to be the \emph{best} approximation $V_h$ can offer
(\cref{sec:cea}).}
\label{fig:galerkin-projection}
\end{figure}

Having narrowed the search to $V_h$, we must also narrow the supply of test
functions: it is unreasonable to demand $a(u_h,v)=\ell(v)$ for \emph{all} $v$ in
the infinite $V$ when $u_h$ only has $N$ knobs to turn. The \emph{Galerkin
choice} is to use the same finite menu for both roles---trial and test---and
require the weak equation only against the $N$ basis functions:
\begin{equation}
\text{find } u_h\in V_h \text{ such that }
a(u_h,\phi_i)=\ell(\phi_i)
\quad\text{for } i=1,2,\dots,N.
\label{eq:galerkin}
\end{equation}
Because $a$ and $\ell$ are linear in their test slot, requiring
\cref{eq:galerkin} against each basis function $\phi_i$ is exactly the same as
requiring it against every combination $v_h=\sum_i d_i\phi_i$---so
\cref{eq:galerkin} really does enforce the weak form for \emph{all} of $V_h$,
not just the basis. We have replaced ``for all $v\in V$,'' an infinitude of
conditions, with ``for $i=1,\dots,N$,'' exactly $N$ of them.

\subsection{Counting: \texorpdfstring{$N$}{N} equations, \texorpdfstring{$N$}{N} unknowns}

Now write the unknown $u_h$ out in the basis. Every element of $V_h$ is some
combination $\sum_j u_j\phi_j$; finding $u_h$ \emph{means} finding the
coefficients $u_1,\dots,u_N$. So we substitute
\begin{equation}
u_h=\sum_{j=1}^N u_j\,\phi_j
\label{eq:uh-expansion}
\end{equation}
into the Galerkin condition \cref{eq:galerkin} and use bilinearity of $a$ to pull
the unknown coefficients out of the form:
\begin{equation}
a(u_h,\phi_i)
=a\Bigl(\sum_{j=1}^N u_j\phi_j,\;\phi_i\Bigr)
=\sum_{j=1}^N u_j\,a(\phi_j,\phi_i)
=\ell(\phi_i),
\qquad i=1,\dots,N.
\label{eq:galerkin-expanded}
\end{equation}
Look at what \cref{eq:galerkin-expanded} is. For each $i$ it is one linear
equation in the unknowns $u_1,\dots,u_N$. There are $N$ values of $i$, hence $N$
equations; there are $N$ unknowns $u_j$. The numbers $a(\phi_j,\phi_i)$ and
$\ell(\phi_i)$ are not unknown---they are definite integrals we can compute from
the data of the problem. The infinite-dimensional weak form has become a
square linear system. That is the whole of Galerkin's method.

\begin{keyidea}
The matrices \emph{are} the bilinear form, sampled on basis pairs. There is no
separate ``matrix-building'' physics: the entry in row $i$, column $j$ of the
stiffness matrix is \emph{defined} to be $a(\phi_j,\phi_i)$, the bilinear form
fed the $j$th and $i$th basis functions. Symmetry, definiteness, sparsity---every
property of the matrix is a property of the form $a$ read off through this
sampling. To know the matrix, know the form and know the basis; nothing else
enters.
\end{keyidea}

\section{The assembled system \texorpdfstring{$\mat{K}\vx=\vb$}{Kx=b}}

We collect \cref{eq:galerkin-expanded} into matrix--vector form. Define the
$N\times N$ matrix $\mat{K}$ and the $N$-vectors $\vx,\vb$ by
\begin{equation}
\mat{K}_{ij}=a(\phi_j,\phi_i),
\qquad
\vx_j=u_j,
\qquad
\vb_i=\ell(\phi_i).
\label{eq:stiffness-def}
\end{equation}
Then the $N$ Galerkin equations \cref{eq:galerkin-expanded} are precisely the
single matrix equation
\begin{equation}
\boxed{\;\mat{K}\,\vx=\vb\;}
\label{eq:Kxb}
\end{equation}
This is the discrete problem the rest of the book solves. The matrix $\mat{K}$ is
called the \emph{stiffness matrix}, a name inherited from structural mechanics,
where $a(u,v)$ measured elastic energy and its discretization measured how stiffly
a structure resisted deformation. The vector $\vb$ is the \emph{load vector}, and
$\vx$ holds the coefficients of the discrete solution $u_h=\sum_j x_j\phi_j$.

\begin{notationbox}
A small but constant source of confusion: the index order in
$\mat{K}_{ij}=a(\phi_j,\phi_i)$ has $j$ (the trial index, the column) inside the
form's \emph{first} slot and $i$ (the test index, the row) in the \emph{second}.
We chose this so that $\mat{K}\vx$ produces row $i=\sum_j a(\phi_j,\phi_i)x_j
=a(u_h,\phi_i)$, the left side of the Galerkin equation. When $a$ is symmetric
($a(u,v)=a(v,u)$, the case for all the stiffness and mass matrices in this book)
the order does not matter and $\mat{K}_{ij}=\mat{K}_{ji}$. It matters only for the
\emph{non}symmetric forms we meet in the driven analysis (\cref{ch:driven}); we
flag it now so the convention is never a surprise.
\end{notationbox}

\subsection{Where \texorpdfstring{$\mat{K}$}{K} and \texorpdfstring{$\mat{M}$}{M} come from}

The two matrices that dominate this book---the stiffness $\mat{K}$ and the
\emph{mass} $\mat{M}$---are both instances of \cref{eq:stiffness-def}, differing
only in which bilinear form is sampled. The stiffness matrix samples a form built
from \emph{derivatives} of the basis functions; the mass matrix samples the
plain $L^2$ pairing with \emph{no} derivative:
\begin{equation}
\mat{K}_{ij}=a(\phi_j,\phi_i)=\int_\dom Q\,\diffop\phi_j\cdot\diffop\phi_i\dV,
\qquad
\mat{M}_{ij}=\int_\dom \phi_i\,\phi_j\dV.
\label{eq:KM}
\end{equation}
Here $\diffop$ is whichever differential operator the weak form attached to the
trial and test functions---$\Grad$, $\Curl$, $\Div$, or the identity---and $Q$ is
a material weight (a permittivity $\varepsilon$, a reluctivity $\nu=1/\mu$, and so
on). The mass matrix is the special case $\diffop=\text{identity}$ and $Q=1$:
$\mat{M}_{ij}=\int\phi_i\phi_j$ measures the overlap of basis function $i$ with
basis function $j$.

We will meet these two matrices in every analysis. The eigenmode problem of
\cref{ch:eigenmodes} is the generalized eigenvalue problem $\mat{K}\vx=\lambda
\mat{M}\vx$; the transient problem of \cref{ch:transient} marches
$\mat{M}\ddot{\vs}+\mat{C}\dot{\vs}+\mat{K}\vs=\vJ(t)$; the driven problem of
\cref{ch:driven} solves $(\mat{K}+i\omega\mat{C}-\omega^2\mat{M})\vE=\vb$. In
\emph{every} one of these, $\mat{K}$ and $\mat{M}$ are the bilinear forms of
\cref{eq:KM} sampled on basis pairs. They look like distinct objects in a linear
algebra library; they are, in truth, the same form-on-basis construction with
different $\diffop$ and $Q$.

\begin{physicsbox}
The mass matrix earns its name from dynamics. In the wave-equation
discretization $\mat{M}\ddot{\vs}+\mat{K}\vs=0$, the term $\mat{M}\ddot{\vs}$ is
the discrete analogue of ``mass times acceleration.'' Off-diagonal entries
$\mat{M}_{ij}\neq0$ mean basis functions $i$ and $j$ overlap, so that
accelerating the coefficient $u_j$ exerts an inertial pull on equation $i$---the
finite-element mesh has, in effect, a smeared-out mass distribution rather than
point masses. A diagonal (``lumped'') mass matrix is the point-mass
idealization; the true \emph{consistent} mass matrix of \cref{eq:KM} couples
neighbours, as we will see in \cref{wex:mass}.
\end{physicsbox}

\section{Three properties that decide everything downstream}

The matrix $\mat{K}$ inherits its character from the form $a$. Three inherited
properties govern which solver the rest of the book reaches for, and we take them
in turn.

\subsection{Symmetry}

If the form is symmetric, $a(u,v)=a(v,u)$, then by \cref{eq:stiffness-def}
\[
\mat{K}_{ij}=a(\phi_j,\phi_i)=a(\phi_i,\phi_j)=\mat{K}_{ji},
\]
so $\mat{K}$ is a symmetric matrix, $\mat{K}=\mat{K}^{\mathsf T}$. Every form of
the family \cref{eq:KM} with a symmetric (in particular, scalar) weight $Q$ is
symmetric, because $Q\,\diffop\phi_j\cdot\diffop\phi_i=Q\,\diffop\phi_i\cdot
\diffop\phi_j$ pointwise. Symmetry is not a cosmetic nicety: it is the hypothesis
the conjugate-gradient method requires. When $\mat{K}$ is symmetric (and positive
definite, below), the system $\mat{K}\vx=\vb$ can be solved by conjugate
gradients, the workhorse iterative method of \cref{ch:cg}---which exploits exactly
this symmetry to build its short three-term recurrence. A nonsymmetric $\mat{K}$
forces the more expensive GMRES of \cref{ch:gmres} instead. The choice of solver
is decided here, at assembly time, by a property of the continuous form.

\subsection{Coercivity and positive-definiteness}

\Cref{ch:weak} introduced \emph{coercivity}: there is a constant $\alpha>0$ with
$a(v,v)\ge\alpha\norm{v}^2$ for every $v\in V$. Coercivity is what guaranteed the
continuous weak problem had a unique solution (the Lax--Milgram theorem). It
descends to the matrix as \emph{positive-definiteness}. Take any nonzero
coefficient vector $\vy=(y_1,\dots,y_N)$ and let $v_h=\sum_j y_j\phi_j$, a nonzero
element of $V_h$. Then
\begin{equation}
\vy^{\mathsf T}\mat{K}\vy
=\sum_{i,j}y_i\,\mat{K}_{ij}\,y_j
=\sum_{i,j}y_i\,a(\phi_j,\phi_i)\,y_j
=a\Bigl(\textstyle\sum_j y_j\phi_j,\sum_i y_i\phi_i\Bigr)
=a(v_h,v_h)
\ge\alpha\norm{v_h}^2>0.
\label{eq:spd}
\end{equation}
The chain \cref{eq:spd} is worth pausing on: the algebraic quantity
$\vy^{\mathsf T}\mat{K}\vy$, the thing a numerical analyst calls the matrix's
quadratic form, is \emph{equal} to the continuous energy $a(v_h,v_h)$ of the
function $v_h$ those coefficients represent. Coercivity of the form is thus,
literally, positive-definiteness of the matrix. A symmetric positive-definite
(SPD) matrix is invertible---so the discrete system has a unique solution---and it
is precisely the class on which conjugate gradients is guaranteed to converge.
Two properties an engineer wants, solvability and a fast solver, both flow from
the single analytic fact that $a$ is coercive.

\subsection{Sparsity}
\label{sec:sparsity}

The third inherited property is the one that makes large problems tractable at
all. The entry $\mat{K}_{ij}=\int_\dom Q\,\diffop\phi_j\cdot\diffop\phi_i\dV$ is a
product of two basis functions integrated over the domain. If $\phi_i$ and
$\phi_j$ are nonzero on \emph{disjoint} regions of the domain, their product is
zero everywhere and so is the integral: $\mat{K}_{ij}=0$. Finite-element basis
functions are built to have \emph{local support}---each $\phi_i$ is nonzero only
on the handful of mesh cells touching node $i$ (\cref{fig:hats})---so for the
overwhelming majority of index pairs $(i,j)$ the supports miss each other and the
entry vanishes.

\begin{figure}[t]
\centering
\begin{tikzpicture}[scale=1.0, every node/.style={font=\footnotesize}]
  % nodes around node i, showing only i-1, i, i+1 couple
  \def\R{0.16}
  \foreach \k/\x in {{i-2}/0,{i-1}/1.6,{i}/3.2,{i+1}/4.8,{i+2}/6.4}{
    \fill[simInk] (\x,0) circle (\R);
    \node[below=2pt] at (\x,0) {$\phi_{\k}$};
  }
  % highlight phi_i support region
  \draw[draw=simBlue, fill=simBgBlue, rounded corners=4pt, opacity=0.6]
    (1.6,-0.5) rectangle (4.8,0.95);
  \node[simBlue!70!simInk, above] at (3.2,0.62) {support of $\phi_i$};
  % coupling arcs: i couples only to i-1 and i+1
  \draw[depedge, bend left=40] (3.2,0.12) to (1.6,0.12);
  \draw[depedge, bend left=40] (3.2,0.12) to (4.8,0.12);
  \node[simBlue, font=\scriptsize] at (2.4,0.95) {couples};
  % the far ones don't couple: dotted cross
  \draw[refedge] (3.2,-0.28) -- (6.4,-0.28);
  \node[simGray, font=\scriptsize] at (5.4,-0.5) {$\mat{K}_{i,\,i+2}=0$};
\end{tikzpicture}
\caption[Local support gives sparsity]{Why $\mat{K}$ is sparse. A finite-element
basis function $\phi_i$ is nonzero only on the cells immediately around node $i$
(shaded). The entry $\mat{K}_{ij}=\int Q\,\diffop\phi_j\cdot\diffop\phi_i$ is
nonzero only when the supports of $\phi_i$ and $\phi_j$ \emph{overlap}, which
happens only for the few $j$ that are mesh-neighbours of $i$. Here $\phi_i$
couples to $\phi_{i-1}$ and $\phi_{i+1}$ but not to $\phi_{i+2}$, whose support is
disjoint---so that matrix entry is exactly zero.}
\label{fig:sparsity-support}
\end{figure}

The consequence is that $\mat{K}$ has only $\mathcal{O}(N)$ nonzero entries, not
$N^2$. For a one-dimensional mesh with the hat functions of
\cref{sec:worked-stiffness} each row has at most three nonzeros and $\mat{K}$ is
\emph{tridiagonal}; in three dimensions each row has a few tens of nonzeros
regardless of how large $N$ grows. A matrix with $\mathcal{O}(N)$ nonzeros can be
\emph{multiplied} by a vector in $\mathcal{O}(N)$ work without ever forming the
$N^2$ dense array---and that is the only operation the iterative solvers of
\cref{ch:krylov} ever ask of it. Sparsity is the reason we solve $\mat{K}\vx=\vb$
by iteration rather than by Gaussian elimination: the matrix is too large to store
densely but cheap to apply.

\begin{pitfall}
Sparsity is a property of the \emph{basis}, not of the differential equation.
Choose global basis functions---say $\phi_k(x)=\sin(k\pi x)$, each nonzero across
the whole interval---and \emph{every} pair of supports overlaps, so $\mat{K}$ is
dense even for the same Poisson problem. (For that particular global basis
$\mat{K}$ happens to come out diagonal by orthogonality, but that is a lucky
accident of the sine basis, not a general phenomenon.) The finite-element method
earns its sparsity precisely by insisting on \emph{locally} supported basis
functions; a different choice forfeits it. This is the first hint of the theme of
\cref{sec:basis-is-everything}: the basis is the whole game.
\end{pitfall}

\section{Worked example: the 1-D Poisson stiffness matrix}
\label{sec:worked-stiffness}

We now assemble $\mat{K}$ explicitly for the model problem
\cref{eq:poisson-weak}, and discover the most famous matrix in numerical analysis.

\subsection{The mesh and the hat basis}

Partition the interval $[0,1]$ into $n$ equal cells with $n+1$ nodes
\[
x_0=0,\quad x_1=h,\quad x_2=2h,\quad\dots,\quad x_n=1,
\qquad h=\frac1n.
\]
At each \emph{interior} node $x_i$ ($i=1,\dots,n-1$) place the \emph{hat
function} $\phi_i$: the piecewise-linear function equal to $1$ at $x_i$, equal to
$0$ at every other node, and linear in between (\cref{fig:hats}). Explicitly,
\begin{equation}
\phi_i(x)=
\begin{cases}
\dfrac{x-x_{i-1}}{h}, & x_{i-1}\le x\le x_i,\\[2ex]
\dfrac{x_{i+1}-x}{h}, & x_i\le x\le x_{i+1},\\[1ex]
0, & \text{otherwise.}
\end{cases}
\label{eq:hat}
\end{equation}
There are $N=n-1$ of them, one per interior node; the two boundary nodes carry no
unknown because $u(0)=u(1)=0$ is built in (those are the \emph{essential} boundary
conditions of \cref{ch:weak}, and we simply omit their hats). Each hat is nonzero
on exactly two cells, so its support is local and \cref{sec:sparsity} promises a
sparse---in fact tridiagonal---matrix.

\begin{figure}[t]
\centering
\begin{tikzpicture}[scale=1.0]
\begin{axis}[
  width=12cm, height=5.2cm,
  xmin=-0.05, xmax=1.05, ymin=-0.08, ymax=1.18,
  xtick={0,0.2,0.4,0.6,0.8,1.0},
  xticklabels={$x_0$,$x_1$,$x_2$,$x_3$,$x_4$,$x_5$},
  ytick={0,1}, ylabel={}, axis lines=left,
  tick label style={font=\footnotesize},
  clip=false,
]
  % hat phi_1 at x=0.2
  \addplot[simBlue, very thick] coordinates {(0,0)(0.2,1)(0.4,0)(1.0,0)};
  % hat phi_2 at x=0.4
  \addplot[simRust, very thick] coordinates {(0,0)(0.2,0)(0.4,1)(0.6,0)(1.0,0)};
  % hat phi_3 at x=0.6
  \addplot[simGreen, very thick] coordinates {(0,0)(0.4,0)(0.6,1)(0.8,0)(1.0,0)};
  % hat phi_4 at x=0.8
  \addplot[simViolet, very thick] coordinates {(0,0)(0.6,0)(0.8,1)(1.0,0)};
  % node markers
  \addplot[only marks, mark=*, mark size=1.4pt, simInk]
    coordinates {(0,0)(0.2,0)(0.4,0)(0.6,0)(0.8,0)(1.0,0)};
  \node[font=\scriptsize, simBlue]   at (axis cs:0.2,1.08) {$\phi_1$};
  \node[font=\scriptsize, simRust]   at (axis cs:0.4,1.08) {$\phi_2$};
  \node[font=\scriptsize, simGreen]  at (axis cs:0.6,1.08) {$\phi_3$};
  \node[font=\scriptsize, simViolet] at (axis cs:0.8,1.08) {$\phi_4$};
\end{axis}
\end{tikzpicture}
\caption[The hat basis]{The hat (piecewise-linear ``tent'') basis on a uniform
mesh of $n=5$ cells, $h=\tfrac15$. Each interior node $x_i$ carries one basis
function $\phi_i$ that peaks at $1$ on its own node and falls linearly to $0$ at
the two neighbouring nodes, vanishing everywhere else. The two boundary nodes
$x_0,x_5$ carry no function because $u$ is pinned to zero there. Adjacent hats
overlap on one cell (so they couple); hats two nodes apart do not (so their matrix
entry is zero)---the picture of sparsity from \cref{fig:sparsity-support}.}
\label{fig:hats}
\end{figure}

\subsection{Computing the entries}

The form is $a(\phi_j,\phi_i)=\int_0^1\phi_j'(x)\,\phi_i'(x)\dx$. The derivative
of a hat is a pair of constants: from \cref{eq:hat},
\[
\phi_i'(x)=
\begin{cases}
+1/h, & x_{i-1}<x<x_i,\\
-1/h, & x_i<x<x_{i+1},\\
0, & \text{otherwise,}
\end{cases}
\]
a $+1/h$ slope going up and a $-1/h$ slope coming down. We need three kinds of
entry.

\medskip
\noindent\textbf{Diagonal, $\mathbf{K_{ii}}$.} On its support, $\phi_i$ has slope
$\pm1/h$ over two cells each of length $h$:
\[
\mat{K}_{ii}=\int_{x_{i-1}}^{x_{i+1}}\bigl(\phi_i'\bigr)^2\dx
=\Bigl(\tfrac1h\Bigr)^2\!\cdot h+\Bigl(-\tfrac1h\Bigr)^2\!\cdot h
=\frac1h+\frac1h=\frac2h.
\]

\noindent\textbf{First off-diagonal, $\mathbf{K_{i,i+1}}$.} The hats $\phi_i$ and
$\phi_{i+1}$ overlap only on the single cell $[x_i,x_{i+1}]$. There
$\phi_i'=-1/h$ (coming down) while $\phi_{i+1}'=+1/h$ (going up), so
\[
\mat{K}_{i,i+1}=\int_{x_i}^{x_{i+1}}\phi_i'\,\phi_{i+1}'\dx
=\Bigl(-\tfrac1h\Bigr)\Bigl(+\tfrac1h\Bigr)\cdot h=-\frac1h.
\]
By symmetry $\mat{K}_{i+1,i}=-1/h$ as well.

\medskip
\noindent\textbf{Everything else, $\mathbf{K_{ij}}$ with $|i-j|\ge2$.} Disjoint
supports, so the integral is zero.

\medskip
Assembling these into the $N\times N$ matrix gives the celebrated tridiagonal
form
\begin{equation}
\mat{K}=\frac1h
\begin{pmatrix}
2 & -1 & & & \\
-1 & 2 & -1 & & \\
 & -1 & 2 & \ddots & \\
 & & \ddots & \ddots & -1\\
 & & & -1 & 2
\end{pmatrix}
=\frac1h\,\operatorname{tridiag}(-1,\,2,\,-1).
\label{eq:tridiag}
\end{equation}
This is the discrete second derivative, the same matrix a finite-difference
scheme produces by approximating $-u''$ with
$\bigl(-u_{i-1}+2u_i-u_{i+1}\bigr)/h^2$ (the $1/h$ versus $1/h^2$ difference is
just where the cell width is booked; the finite-difference RHS carries an extra
$h$). The finite-element method has reproduced the classical stencil---but it
arrived by sampling a bilinear form on basis functions, which is what lets the
same machinery handle curved elements, anisotropic materials, and vector fields
that the finite-difference stencil cannot.

\begin{figure}[t]
\centering
\begin{tikzpicture}[scale=0.62, every node/.style={font=\scriptsize}]
  \def\nz{simBlue}
  % draw a 7x7 grid; nonzeros on tridiagonal
  \foreach \r in {0,...,6}{
    \foreach \c in {0,...,6}{
      \pgfmathtruncatemacro{\d}{abs(\r-\c)}
      \ifnum\d<2
        \fill[\nz] (\c,-\r) rectangle ++(0.92,-0.92);
      \else
        \draw[simGray!40] (\c,-\r) rectangle ++(0.92,-0.92);
      \fi
    }
  }
  % frame
  \draw[simInk, thick] (0,0) rectangle (6.92,-6.92);
  \node[anchor=south, font=\footnotesize] at (3.46,0.25) {columns $j$};
  \node[anchor=east, rotate=90, font=\footnotesize] at (-0.3,-3.46) {rows $i$};
  % legend
  \fill[\nz] (8.0,-1.0) rectangle ++(0.7,0.7);
  \node[anchor=west] at (8.85,-0.65) {nonzero $\mat{K}_{ij}$};
  \draw[simGray!40] (8.0,-2.2) rectangle ++(0.7,0.7);
  \node[anchor=west] at (8.85,-1.85) {zero};
\end{tikzpicture}
\caption[Tridiagonal sparsity pattern]{The sparsity pattern of the assembled
stiffness matrix \cref{eq:tridiag} for $N=7$ interior nodes. Filled cells are the
nonzero entries; they sit only on the diagonal and the two adjacent off-diagonals,
because (\cref{fig:sparsity-support}) a hat couples only to its two
nearest-neighbour hats. Of the $49$ possible entries only $19$ are nonzero, and
the fraction shrinks as $N$ grows: a matrix this empty is never stored as a full
array.}
\label{fig:sparsity-grid}
\end{figure}

\subsection{Solving a small system}

Take $n=4$ cells, $h=\tfrac14$, so $N=3$ interior unknowns, and a constant load
$f(x)=1$. The load vector entries are
$\vb_i=\int_0^1 1\cdot\phi_i\dx=\text{area of the hat}=h=\tfrac14$ (each hat is a
triangle of base $2h$ and height $1$, area $h$). The system $\mat{K}\vx=\vb$ is,
multiplying \cref{eq:tridiag} by $h=\tfrac14$ to clear the prefactor,
\[
\frac1{1/4}
\begin{pmatrix} 2&-1&0\\-1&2&-1\\0&-1&2\end{pmatrix}
\begin{pmatrix}x_1\\x_2\\x_3\end{pmatrix}
=\begin{pmatrix}1/4\\1/4\\1/4\end{pmatrix}
\;\Longrightarrow\;
\begin{pmatrix} 2&-1&0\\-1&2&-1\\0&-1&2\end{pmatrix}
\begin{pmatrix}x_1\\x_2\\x_3\end{pmatrix}
=\begin{pmatrix}1/16\\1/16\\1/16\end{pmatrix}.
\]
Solving the tridiagonal system (eliminate downward, then back-substitute) gives
\[
x_1=\tfrac{3}{32}=0.09375,\quad
x_2=\tfrac{4}{32}=0.125,\quad
x_3=\tfrac{3}{32}=0.09375.
\]
The exact solution of $-u''=1$, $u(0)=u(1)=0$ is the parabola
$u(x)=\tfrac12 x(1-x)$, whose values at the nodes $x=\tfrac14,\tfrac12,\tfrac34$
are $\tfrac{3}{32},\tfrac14\!\cdot\!\tfrac34\!\cdot\!\tfrac12=\tfrac{3}{32}$\dots
let us check: $u(\tfrac14)=\tfrac12\cdot\tfrac14\cdot\tfrac34=\tfrac{3}{32}$,
$u(\tfrac12)=\tfrac12\cdot\tfrac12\cdot\tfrac12=\tfrac18=\tfrac{4}{32}$,
$u(\tfrac34)=\tfrac{3}{32}$. The finite-element coefficients match the exact
nodal values \emph{exactly}. This is no coincidence: for the 1-D Poisson problem
with the hat basis, the Galerkin solution is \emph{nodally exact}---a celebrated
special feature of this problem that does not survive into two dimensions or to
variable coefficients, but is a satisfying confirmation that the machinery works.

\begin{workedexample}{Is this $\mat{K}$ symmetric positive definite?}{spd-check}
We verify directly that the $3\times3$ stiffness above is SPD, the property
\cref{eq:spd} promised.

\emph{Symmetric.} By inspection $\mat{K}=\mat{K}^{\mathsf T}$: the
$(1,2)$ and $(2,1)$ entries are both $-1$, and so on. (We knew this must hold:
$a(u,v)=\int u'v'$ is symmetric in $u,v$.)

\emph{Positive definite.} Use the chain \cref{eq:spd}: for any
$\vy=(y_1,y_2,y_3)\neq\mathbf 0$,
\[
\vy^{\mathsf T}\mat{K}\vy
=2y_1^2+2y_2^2+2y_3^2-2y_1y_2-2y_2y_3
=y_1^2+(y_1-y_2)^2+(y_2-y_3)^2+y_3^2,
\]
a sum of squares that is zero only if $y_1=0$, $y_1=y_2$, $y_2=y_3$, $y_3=0$, i.e.
only if $\vy=\mathbf 0$. So $\vy^{\mathsf T}\mat{K}\vy>0$ for every nonzero $\vy$:
$\mat{K}$ is positive definite. (Equivalently, its leading principal minors
$2,\,3,\,4$ are all positive---Sylvester's criterion.) Because $\mat{K}$ is SPD,
$\mat{K}\vx=\vb$ has a unique solution and conjugate gradients (\cref{ch:cg}) is
guaranteed to find it.
\end{workedexample}

\section{Worked example: the mass matrix}

\begin{workedexample}{The 1-D hat mass matrix}{mass}
\label{wex:mass}
For the \emph{same} hat basis on the same uniform mesh, assemble the mass matrix
$\mat{M}_{ij}=\int_0^1\phi_i\phi_j\dx$---no derivatives this time. We need the
overlap integrals of the hats themselves.

\emph{Diagonal.} On its two support cells $\phi_i$ is a triangle of height $1$;
$\int\phi_i^2$ over one cell of length $h$ is, for a linear function running
$0\to1$, equal to $h/3$. Two cells give
\[
\mat{M}_{ii}=\int_{x_{i-1}}^{x_{i+1}}\phi_i^2\dx=\frac h3+\frac h3=\frac{2h}{3}.
\]

\emph{Off-diagonal.} On the shared cell $[x_i,x_{i+1}]$, $\phi_i$ runs $1\to0$
while $\phi_{i+1}$ runs $0\to1$; their product integrates to $h/6$:
\[
\mat{M}_{i,i+1}=\int_{x_i}^{x_{i+1}}\phi_i\,\phi_{i+1}\dx=\frac h6.
\]
Entries with $|i-j|\ge2$ vanish (disjoint support). Hence
\[
\mat{M}=\frac h6
\begin{pmatrix}
4 & 1 & & \\
1 & 4 & 1 & \\
 & 1 & 4 & \ddots\\
 & & \ddots & \ddots
\end{pmatrix}
=\frac h6\,\operatorname{tridiag}(1,\,4,\,1).
\]

\emph{Contrast with $\mat{K}$.} Both matrices are symmetric, tridiagonal, and
share the \emph{same} sparsity pattern (\cref{fig:sparsity-grid})---unsurprising,
since the nonzero pattern is fixed by which supports overlap, and both forms use
the same hats. But the \emph{signs} differ sharply: $\mat{K}$ has \emph{negative}
off-diagonals (a difference operator, $\diffop\phi_i\cdot\diffop\phi_j<0$ where
the slopes oppose), while $\mat{M}$ has \emph{positive} off-diagonals (an overlap
operator, $\phi_i\phi_j>0$ wherever the hats coexist). $\mat{M}$ is also SPD---it
is a Gram matrix of linearly independent functions---but it is much
better-conditioned than $\mat{K}$, which is why it appears on the ``easy'' side of
the generalized eigenproblem $\mat{K}\vx=\lambda\mat{M}\vx$
(\cref{ch:eigenmodes}).
\end{workedexample}

\section{The bilinear-form family is an operator family}
\label{sec:family}

We have so far assembled \emph{one} matrix from \emph{one} bilinear form. Real
analyses assemble several terms at once. The electrostatic stiffness is a single
diffusion term $(\varepsilon\Grad u,\Grad v)$; but the driven Maxwell operator
$\mat{A}(\omega)=\mat{K}+i\omega\mat{C}-\omega^2\mat{M}$ is a \emph{sum} of three,
and a lossy or multi-material problem adds still more. The structure that makes
this clean is the subject of \cref{ch:assembly}, but the essential idea belongs
here, because it is a direct consequence of the linearity of the form.

\subsection{Each term becomes one operator}

Every member of the bilinear-form family has the shape
$a_t(u,v)=(Q_t\,\diffop_t u,\diffop_t v)$ for some weight $Q_t$ and differential
operator $\diffop_t$---a permittivity-weighted gradient term, a
reluctivity-weighted curl term, a mass term, and so on. Sampling \emph{one} such
term on basis pairs produces \emph{one} matrix, exactly as in \cref{eq:KM}; call
it $\mat{A}(\text{term}_t)$. The full operator an analysis needs is the matrix of
its total form $a=\sum_t a_t$. Because sampling is linear in the form,
\begin{equation}
\mat{K}_{ij}=\Bigl(\sum_t a_t\Bigr)(\phi_j,\phi_i)
=\sum_t a_t(\phi_j,\phi_i)
=\sum_t \mat{A}(\text{term}_t)_{ij},
\qquad\text{i.e.}\qquad
\mat{K}=\sum_t \mat{A}(\text{term}_t).
\label{eq:assembly-sum}
\end{equation}
The global operator is a \emph{sum of per-term operators}. Adding a physical
effect---a boundary loss, a second material, the $i\omega\mat{C}$ damping---is
adding one more matrix to the sum, never rebuilding from scratch.

\begin{figure}[t]
\centering
\begin{tikzpicture}[node distance=6mm and 9mm, every node/.style={font=\footnotesize}]
  \node[opbox] (t1) {$\text{term}_1$\\$(\varepsilon\Grad u,\Grad v)$};
  \node[opbox, below=of t1] (t2) {$\text{term}_2$\\$(\sigma u, v)$};
  \node[opbox, below=of t2] (t3) {$\text{term}_3$\\$(\nu\Curl u,\Curl v)$};
  % per-term operators
  \node[stage, right=14mm of t1] (a1) {$\mat{A}_1$};
  \node[stage, right=14mm of t2] (a2) {$\mat{A}_2$};
  \node[stage, right=14mm of t3] (a3) {$\mat{A}_3$};
  \draw[flow] (t1) -- node[above,font=\scriptsize]{sample} (a1);
  \draw[flow] (t2) -- node[above,font=\scriptsize]{sample} (a2);
  \draw[flow] (t3) -- node[above,font=\scriptsize]{sample} (a3);
  % the sum node
  \node[product, right=20mm of a2] (K) {$\mat{K}=\sum_t\mat{A}_t$};
  \draw[flow] (a1) -- (K);
  \draw[flow] (a2) -- (K);
  \draw[flow] (a3) -- (K);
  % bracket label
  \node[font=\scriptsize\itshape, simGray, above=1mm of a1] {one operator per term};
\end{tikzpicture}
\caption[Assembly as a sum of term-operators]{The bilinear-form family becomes an
operator family. Each weak-form term $a_t(u,v)=(Q_t\diffop_t u,\diffop_t v)$
samples, on basis pairs, into its own matrix $\mat{A}_t$; the global operator of
the analysis is their sum, $\mat{K}=\sum_t\mat{A}_t$ (\cref{eq:assembly-sum}).
Because addition of operators is commutative and associative, the order in which
the terms are accumulated is immaterial---the structure of a \emph{commutative-
monoid fold} over the list of terms, which \cref{ch:assembly} makes precise.}
\label{fig:assembly-sum}
\end{figure}

\subsection{Assembly is a fold: a preview}

\Cref{eq:assembly-sum} is more than a formula; it is a \emph{computational
pattern}. We have a list of terms $[\text{term}_1,\dots,\text{term}_m]$, a way to
turn each into a matrix, and a way to combine matrices (add them). Combining a
list of things with an associative, commutative operation that has an identity
(the zero matrix) is the canonical setting for a \emph{fold}---the higher-order
function that walks a list accumulating a result. In the pseudo-language of
Part~III we will write assembly as
\begin{equation}
\mat{K}=\texttt{fe\_assemble}(\textit{space},\textit{terms})
=\textstyle\sum_{t\in\textit{terms}}\mat{A}(\textit{space},t),
\label{eq:fe-assemble}
\end{equation}
literally a fold of the per-term assembly map over the immutable list of terms.
Two features of this fold are worth flagging now, because they are properties of
the \emph{form}, not of any particular code:

\begin{itemize}[nosep]
\item \textbf{Order independence.} Operator addition is commutative and
associative, so $\sum_t\mat{A}_t$ does not depend on the order the terms are
listed. The fold is over a \emph{commutative monoid}; the term list is really an
unordered bag.
\item \textbf{The list-homomorphism intuition.} Assembling a concatenated term
list equals assembling the pieces and adding:
$\texttt{fe\_assemble}(\textit{space},\,T_1\!+\!\!+\,T_2)=
\texttt{fe\_assemble}(\textit{space},T_1)+\texttt{fe\_assemble}(\textit{space},T_2)$.
A function that turns list-concatenation into a combine on results is a
\emph{list homomorphism}, the structure that makes assembly splittable, mergeable,
and---on a parallel machine---divisible across processors.
\end{itemize}

This is why \cref{ch:assembly} can treat assembly as one combinator rather than a
tangle of nested loops, and why the matrix-free operator of \cref{ch:matrixfree}
can apply the global $\mat{K}$ without ever forming it: applying a sum of
operators is applying each and adding. The bilinear-form family, sampled, is an
operator family, folded.

\section{How good is the Galerkin solution?}
\label{sec:cea}

We have a recipe that produces \emph{some} $u_h$ in $V_h$. Is it any good? The
reassuring answer, made precise by a one-line argument of Jean Céa, is that the
Galerkin solution is \emph{quasi-optimal}: it is, up to a fixed constant, as
accurate as the very best approximation the subspace $V_h$ could possibly offer.

\subsection{Galerkin orthogonality}

Start from a small but powerful fact. The exact solution $u$ satisfies
$a(u,v)=\ell(v)$ for all $v\in V$, and in particular for all $v_h\in V_h\subset V$.
The Galerkin solution $u_h$ satisfies $a(u_h,v_h)=\ell(v_h)$ for all $v_h\in V_h$.
Subtracting,
\begin{equation}
a(u-u_h,\,v_h)=0
\qquad\text{for all } v_h\in V_h.
\label{eq:galerkin-orthog}
\end{equation}
The error $u-u_h$ is invisible to every test function in $V_h$: in the geometry
the form $a$ defines, the error is \emph{orthogonal} to the whole subspace. This
is the precise content of \cref{fig:galerkin-projection}---$u_h$ is the orthogonal
projection of $u$ onto $V_h$ in the energy inner product---and it is the engine of
the next result.

\subsection{Céa's lemma}

\begin{lemma}[Céa, quasi-optimality]
Suppose $a$ is bounded, $\abs{a(u,v)}\le M\norm{u}\,\norm{v}$, and coercive,
$a(v,v)\ge\alpha\norm{v}^2$. Then the Galerkin error is controlled by the best
approximation error in $V_h$:
\begin{equation}
\norm{u-u_h}\le\frac{M}{\alpha}\,\min_{v_h\in V_h}\norm{u-v_h}.
\label{eq:cea}
\end{equation}
\end{lemma}

\begin{proof}
Let $v_h\in V_h$ be \emph{any} element of the subspace. Using coercivity, then
Galerkin orthogonality \cref{eq:galerkin-orthog} to subtract the cross term
$a(u-u_h,v_h-u_h)=0$ (legal because $v_h-u_h\in V_h$), then boundedness:
\[
\alpha\norm{u-u_h}^2
\le a(u-u_h,\,u-u_h)
= a(u-u_h,\,u-v_h)
\le M\,\norm{u-u_h}\,\norm{u-v_h}.
\]
Divide both sides by $\alpha\norm{u-u_h}$ (zero error is trivially optimal) to get
$\norm{u-u_h}\le(M/\alpha)\norm{u-v_h}$. Since $v_h$ was arbitrary, take the
minimum over $v_h\in V_h$.
\end{proof}

Read \cref{eq:cea} carefully, because it is the deepest single fact in this
chapter. The right-hand side, $\min_{v_h}\norm{u-v_h}$, is a question of
\emph{approximation theory} alone: how well can the functions in $V_h$ approximate
$u$, ignoring the differential equation entirely? Céa's lemma says the Galerkin
method---which knows nothing of $u$ and only solves $\mat{K}\vx=\vb$---comes
within a constant factor $M/\alpha$ of that best-possible answer
(\cref{fig:cea}). The method does not need to be clever about $u$; it only needs a
subspace $V_h$ that \emph{can} approximate $u$ well. The job of choosing such a
subspace is the job of the basis, which returns us to the chapter's refrain.

\begin{figure}[t]
\centering
\begin{tikzpicture}[scale=1.0, every node/.style={font=\footnotesize}]
  % subspace line V_h
  \draw[draw=simGreen, fill=simBgGreen, rounded corners=4pt]
    (-3.2,-0.35) rectangle (3.2,0.35);
  \draw[simGreen!60!simInk, thick] (-3.0,0)--(3.0,0);
  \node[simGreen!70!simInk, left] at (-3.05,0) {$V_h$};
  % exact solution above
  \fill[simRust] (0.0,1.9) circle (1.8pt);
  \node[simRust, above] at (0.0,2.0) {$u$};
  % best approx (foot of perpendicular)
  \fill[simInk] (0.0,0) circle (1.8pt);
  \node[simInk, below right=-1pt] at (0.05,-0.05) {$v_h^\star$ (best)};
  % galerkin solution nearby
  \fill[simBlue] (1.3,0) circle (1.8pt);
  \node[simBlue, below] at (1.3,-0.12) {$u_h$ (Galerkin)};
  % best distance (perpendicular)
  \draw[flow, simInk] (0.0,1.9) -- (0.0,0.06);
  \node[simInk, right, font=\scriptsize] at (0.06,1.0) {$\min_{v_h}\|u-v_h\|$};
  % galerkin distance
  \draw[flow, simBlue] (0.0,1.9) -- (1.3,0.06);
  \node[simBlue, right=1pt, font=\scriptsize] at (0.78,1.05) {$\|u-u_h\|$};
  % right-angle marker for best approx
  \draw[simInk] (0.0,0.18)--(0.18,0.18)--(0.18,0);
\end{tikzpicture}
\caption[Céa's quasi-optimality]{Céa's lemma, geometrically. The best the
subspace $V_h$ can do is the foot of the perpendicular $v_h^\star$, at distance
$\min_{v_h}\norm{u-v_h}$ from the exact solution $u$. The Galerkin solution $u_h$
need not be that closest point (it is closest in the \emph{energy} norm, not the
norm drawn), but \cref{eq:cea} guarantees it is no worse than a fixed multiple
$M/\alpha$ of the best distance. The Galerkin method is therefore only as good---
and only as bad---as the approximating power of the basis.}
\label{fig:cea}
\end{figure}

\subsection{Convergence: \texorpdfstring{$h$}{h}- and \texorpdfstring{$p$}{p}-refinement}

Quasi-optimality turns convergence questions into approximation questions, and
approximation theory then supplies the rates. Two knobs control how well $V_h$
can approximate, and refining either drives the error down.

\emph{$h$-refinement} shrinks the mesh spacing $h$, adding more, smaller elements.
For a basis of piecewise polynomials of degree $p$ on a mesh of size $h$,
approximating a smooth $u$, standard interpolation estimates give
$\min_{v_h}\norm{u-v_h}=\mathcal{O}(h^{p})$ in the energy norm, so by Céa
\begin{equation}
\norm{u-u_h}=\mathcal{O}(h^{p}).
\label{eq:hrate}
\end{equation}
With the linear hats ($p=1$) of our worked example, halving $h$ halves the energy
error; the error in the solution value itself converges one order faster,
$\mathcal{O}(h^{p+1})$. More elements, more accuracy, at a predictable rate.

\emph{$p$-refinement} keeps the mesh but raises the polynomial degree $p$ of the
basis on each element. For smooth $u$ this is dramatically more effective: the
error falls \emph{faster than any fixed power of $1/p$}---so-called spectral or
exponential convergence---because high-degree polynomials approximate smooth
functions extraordinarily well. Palace exploits both knobs, and combines them in
the $p$-multigrid hierarchy of \cref{ch:multigrid}. Which knob to turn is a
judgement about the smoothness of the solution: $h$-refinement is robust near
singularities (sharp corners, material interfaces), $p$-refinement wins where the
field is smooth.

\begin{remark}
The constant $M/\alpha$ in \cref{eq:cea} is the form's \emph{condition number} in
disguise---the ratio of its boundedness constant to its coercivity constant. The
same ratio reappears as the condition number $\kappa(\mat{K})$ of the assembled
matrix, which (\cref{ch:cg}) sets the convergence rate of conjugate gradients.
A form that is barely coercive ($\alpha$ small) gives both a loose error bound and
a slow solver---another instance of a continuous property dictating discrete
behaviour.
\end{remark}

\section{The basis is the whole game}
\label{sec:basis-is-everything}

Step back and notice how little of Galerkin's method depended on the
\emph{physics}. The recipe $\mat{K}_{ij}=a(\phi_j,\phi_i)$, $\vb_i=\ell(\phi_i)$,
solve $\mat{K}\vx=\vb$, is identical for the Poisson equation, the curl--curl
equation, and every analysis in this book. What changes from one problem to the
next, beyond the form $a$, is the \emph{basis} $\{\phi_i\}$---and the basis alone
determines:

\begin{itemize}[nosep]
\item the \textbf{size} $N$ of the system (one unknown per basis function);
\item the \textbf{sparsity} of $\mat{K}$ (which supports overlap,
\cref{sec:sparsity});
\item the \textbf{accuracy}, through the best-approximation term in Céa's lemma
(\cref{sec:cea});
\item and, crucially, \textbf{which physical fields the method can represent
faithfully.}
\end{itemize}

That last point is where the rest of Part~II is heading. A scalar potential $V$
is well represented by the continuous, nodal hat functions of this chapter---an
$\Hone$-conforming basis, in the language of \cref{ch:functionspaces}. But an
electric field $\vE$ is a \emph{vector} field with a different continuity
requirement (its tangential component, not its full value, is continuous across
material interfaces), and forcing nodal hats onto it produces spurious,
unphysical modes. The fix is a different basis entirely---the edge elements of
$\Hcurl$, whose degrees of freedom live on mesh \emph{edges} rather than nodes.
Which basis a field needs is dictated by the differential operator $\diffop$ in
its form: $\Grad$ wants nodes, $\Curl$ wants edges, $\Div$ wants faces. Building
those bases, and the reference elements they are constructed on, is the work of
\cref{ch:functionspaces,ch:refelement}. Galerkin's method is the same for all of
them; only the menu of functions changes.

\summarysection
Galerkin's method makes the weak form computable by replacing the
infinite-dimensional trial/test space $V$ with a finite-dimensional subspace
$V_h=\operatorname{span}\{\phi_1,\dots,\phi_N\}$ and requiring the weak equation
only against the $N$ basis functions. Expanding the unknown $u_h=\sum_j u_j\phi_j$
turns this into a square linear system $\mat{K}\vx=\vb$ with
$\mat{K}_{ij}=a(\phi_j,\phi_i)$ and $\vb_i=\ell(\phi_i)$---the matrix \emph{is}
the bilinear form sampled on basis pairs, and the mass matrix
$\mat{M}_{ij}=\int\phi_i\phi_j$ is the same construction with the identity
operator. The matrix inherits its properties from the form: symmetry (enabling
conjugate gradients), positive-definiteness from coercivity (guaranteeing a unique
solution), and---from the \emph{local support} of the basis---sparsity, which is
why we solve by iteration. The bilinear-form family is an operator family: each
term $a_t=(Q_t\diffop_t u,\diffop_t v)$ assembles into one matrix and the global
operator is their sum $\mat{K}=\sum_t\mat{A}_t$, a commutative-monoid fold and
list homomorphism over the term list---the assembly-as-a-fold of \cref{ch:assembly}.
Céa's lemma shows the Galerkin solution is quasi-optimal: within a constant
$M/\alpha$ of the best approximation $V_h$ can provide, so convergence ($h$- and
$p$-refinement) is governed by the approximating power of the basis. The basis,
in the end, is the whole game.

\furtherreading
The variational and approximation theory of this chapter---Galerkin orthogonality,
Céa's lemma, the $h$- and $p$-estimates, and the SPD structure of stiffness
matrices---is developed rigorously in Brenner and Scott~\citep{brennerscott2008},
the standard graduate reference for the mathematics of finite elements; their
treatment of the bilinear-form framework and approximation rates expands every
result we stated. For the linear-algebra side---why sparsity and conditioning
decide between conjugate gradients and GMRES, and how the matrix's spectral
properties govern iterative convergence---Trefethen and Bau~\citep{trefethenbau1997}
is the classic and exceptionally readable account. For the engineering
practice of assembling and solving finite-element systems for electromagnetics
specifically, including vector ($\Hcurl$) bases and the families we turn to next,
Jin~\citep{jin2014} is the comprehensive reference.

\exercisessection
\begin{enumerate}[nosep]
  \item \textbf{One more node.} Redo the worked example of
  \cref{sec:worked-stiffness} with $n=5$ cells ($h=\tfrac15$, $N=4$ unknowns) and
  constant load $f=1$. Write down the $4\times4$ system $\mat{K}\vx=\vb$ and solve
  it. Confirm that the coefficients again equal the exact nodal values of
  $u=\tfrac12 x(1-x)$.

  \item \textbf{Variable load.} Keep $n=4$, $h=\tfrac14$, but take
  $f(x)=x$. Compute the load vector $\vb_i=\int_0^1 x\,\phi_i(x)\dx$ (the
  integral of $x$ against each hat) and solve $\mat{K}\vx=\vb$ with the same
  $\mat{K}$ as in the text. Why did $\mat{K}$ not change when the right-hand side
  did?

  \item \textbf{Counting nonzeros.} For the 1-D hat basis on $N$ interior nodes,
  how many nonzero entries does $\mat{K}$ have? Express it as a function of $N$ and
  confirm it is $\mathcal{O}(N)$, not $\mathcal{O}(N^2)$. How many nonzeros would
  a 2-D mesh in which each node couples to six neighbours have per row?

  \item \textbf{Symmetry from the form.} Show that \emph{any} form of the family
  $a(u,v)=\int_\dom Q\,\diffop u\cdot\diffop v$ with scalar $Q$ assembles a
  symmetric matrix, directly from $\mat{K}_{ij}=a(\phi_j,\phi_i)$. Then give an
  example of a form (a first-derivative ``advection'' term $\int u'\,v$) whose
  matrix is \emph{not} symmetric, and identify which step of the symmetry argument
  fails.

  \item \textbf{Mass-matrix positive-definiteness.} Using the chain
  \cref{eq:spd}, argue that the mass matrix $\mat{M}_{ij}=\int\phi_i\phi_j$ is
  positive definite whenever the basis functions $\{\phi_i\}$ are linearly
  independent. (Hint: $\vy^{\mathsf T}\mat{M}\vy=\int(\sum_j y_j\phi_j)^2$. When is
  this zero?)

  \item \textbf{Best approximation versus Galerkin.} For the $n=4$ Poisson
  example, the Galerkin solution was \emph{nodally exact}. The best $L^2$
  approximation of $u=\tfrac12 x(1-x)$ in $V_h$ minimizes $\int(u-v_h)^2$ and is
  generally \emph{not} nodally exact. Without computing it, explain using Céa's
  lemma and Galerkin orthogonality why the two need not coincide, and which inner
  product the Galerkin solution \emph{is} optimal in.

  \item \textbf{Assembly as a fold.} An analysis has three weak-form terms with
  per-term matrices $\mat{A}_1,\mat{A}_2,\mat{A}_3$. Write the global operator
  three different ways---$\mat{A}_1+(\mat{A}_2+\mat{A}_3)$,
  $(\mat{A}_1+\mat{A}_2)+\mat{A}_3$, and $\mat{A}_2+\mat{A}_3+\mat{A}_1$---and
  argue they are equal. Which algebraic properties of operator addition did you
  use, and why do they make assembly a \emph{commutative-monoid} fold rather than
  merely a list fold?

  \item \textbf{Why not a global basis?} Suppose you chose the global sine basis
  $\phi_k(x)=\sin(k\pi x)$ on $(0,1)$ instead of hats. (a)~Show every pair of
  supports overlaps, so the sparsity argument of \cref{sec:sparsity} fails.
  (b)~Compute $\mat{K}_{ij}=\int_0^1\phi_i'\phi_j'$ for this basis and show it
  comes out \emph{diagonal} anyway. (c)~Explain why this pleasant accident does
  not save the finite-element method in general, and what is lost by giving up
  local support.
\end{enumerate}

\chapter{The Finite-Element Function Spaces}
\label{ch:functionspaces}

\chapterorient{The previous chapters turned each Maxwell problem into a weak form
$a(u,v)=\ell(v)$ and a Galerkin recipe: pick a finite-dimensional space $V_h$,
expand the unknown in a basis, and solve a linear system. We left one question
open---\emph{which} space? This chapter answers it. There are four canonical
finite-element spaces, one for each kind of physical field, and the choice is not
a matter of taste: a field that must be continuous in its tangential part lives in
a different space from one that must be continuous in its normal part. Put a field
in the wrong space and the simulator returns nonsense---spurious modes, locked
solutions, fields that cannot represent a constant. Get it right and the
discretization inherits the exact structure of Maxwell's equations. By the end you
will know why electrostatics uses a scalar nodal space while every magnetic and
wave problem uses edge elements, you will be able to count the degrees of freedom
of each space on a mesh by hand, and you will have met the number $N$ that names
the dimension of every tensor in this book.}

\section{Why one space will not do}

Every problem in \cref{ch:reductions-pde} reduced to the same abstract shape: find
$u$ in some space such that
\[
  a(u,v) = \ell(v) \qquad \text{for all test functions } v,
\]
with $a$ the bilinear form of \cref{ch:galerkin}. To compute, we replace the
infinite-dimensional space of all admissible fields by a finite-dimensional
subspace $V_h$ spanned by simple, locally-supported basis functions. The unknown
becomes a finite list of coefficients, and $a(u,v)=\ell(v)$ becomes a matrix
equation. This is the finite-element method, and the basis functions are the
\emph{finite elements}.

The temptation is to use one space for everything. After all, an electric field
$\vE$ and a magnetic field $\vB$ are both just vector fields---three numbers at
each point of space. Why not discretize each component the same way we discretize
a scalar, storing values at mesh vertices and interpolating linearly in between?

The answer is that $\vE$ and $\vB$ are not ``just'' vector fields. They are vector
fields with \emph{specific continuity requirements} dictated by Maxwell's
equations, and those requirements differ between the two. Recall from
\cref{ch:maxwell} the interface conditions that hold across any surface separating
two materials: the \emph{tangential} component of $\vE$ is continuous, while the
\emph{normal} component of $\vB$ is continuous. The normal component of $\vE$ may
jump (a surface charge lives in the jump); the tangential component of $\vB$ may
jump (a surface current lives there). A discretization that does not respect the
correct interface condition is not merely inaccurate---it represents a field that
is physically impossible, and the solver will dutifully compute with it.

\begin{keyidea}
The finite-element space must enforce exactly the continuity that physics
requires---no more, no less. Tangentially-continuous fields ($\vE$, the magnetic
vector potential $\vA$) belong in one space; normally-continuous fields (fluxes,
$\vB$, $\vD$) belong in another; scalar potentials belong in a third; quantities
with no inter-element continuity at all belong in a fourth. The degrees of freedom
of each space are chosen to be precisely the quantities that must match across
element faces. \emph{This single principle---DoFs match required continuity---is
the organizing idea of the whole chapter.}
\end{keyidea}

There are four such spaces, and they are not arbitrary. They form a chain---the
\emph{de Rham complex}---linked by the gradient, curl, and divergence operators,
which we develop fully in \cref{ch:derham}. This chapter introduces the four spaces
one at a time, by the degrees of freedom that define them, and shows how the
abstract principle above forces each PDE onto a particular slot in the chain.

\section{Degrees of freedom: the language of finite elements}

Before meeting the four spaces we fix the vocabulary that distinguishes them. A
finite-element space is built from a mesh---a partition of the domain $\dom$ into
simple cells (tetrahedra, hexahedra) sharing faces, edges, and vertices. A field in
the space is determined by a finite list of numbers, its \emph{degrees of freedom}
(DoFs). Each DoF is a \emph{functional}: a rule that reads one number out of a
field.

\begin{definition}[Degree of freedom]
A \emph{degree of freedom} of a finite-element space is a linear functional
$\sigma_i$ on fields, together with a basis function $\phi_j$ satisfying the
\emph{duality} (or \emph{unisolvence}) relation
\[
  \sigma_i(\phi_j) = \delta_{ij}.
\]
A field $u$ in the space is recovered from its DoFs by
$u = \sum_i \sigma_i(u)\,\phi_i$; the coefficient on $\phi_i$ is exactly the value
the functional $\sigma_i$ reads.
\end{definition}

The four spaces differ in \emph{what their functionals read} and \emph{where on the
mesh they live}---and that is the whole story. The choices are summarized in
\cref{tab:fourdofs}, which we then unpack space by space.

\begin{table}[t]
\centering
\caption[Degrees of freedom of the four spaces]{The four finite-element spaces,
distinguished entirely by their degrees of freedom: what each DoF reads and which
geometric entity of the mesh it lives on. The last column anticipates
\cref{ch:derham}: the four spaces form a chain linked by $\Grad$, $\Curl$, $\Div$.}
\label{tab:fourdofs}
\small
\setlength{\tabcolsep}{5pt}
\renewcommand{\arraystretch}{1.3}
\begin{tabular}{@{}lllll@{}}
\toprule
Space & A DoF reads\dots & lives on a\dots & continuity & de Rham slot \\
\midrule
$\Hone$ & a value $u(\vx_i)$ & vertex (node) & full ($C^0$) & start \\
$\Hcurl$ & a circulation $\int_e \mathbf{u}\cdot d\vec\ell$ & edge & tangential & after $\Grad$ \\
$\Hdiv$ & a flux $\int_f \mathbf{u}\cdot\vn\,dA$ & face & normal & after $\Curl$ \\
$\Ltwo$ & a cell value / mean & cell interior & none & after $\Div$ \\
\bottomrule
\end{tabular}
\end{table}

\section{The four spaces}

\subsection{$\Hone$: nodal elements, the home of the scalar potential}

The simplest finite element stores the \emph{value} of a scalar field at each mesh
vertex. The functionals are point evaluations,
\[
  \sigma_i(u) = u(\vx_i), \qquad \vx_i = \text{the $i$-th vertex,}
\]
and the basis functions are the familiar ``hat'' or \emph{Lagrange} functions:
$\phi_i$ is the unique piecewise-polynomial that equals $1$ at vertex $\vx_i$ and
$0$ at every other vertex. A field is reconstructed by interpolating its nodal
values across each cell.

Because neighbouring cells share the vertex $\vx_i$ and both assign it the single
coefficient $\sigma_i(u)=u(\vx_i)$, the reconstructed field takes the \emph{same}
value on both sides of every shared face: it is continuous. This is the space
called $\Hone$, the space of square-integrable scalar fields whose gradient is also
square-integrable. Its discrete subspace is denoted $\Hone_h$, and its conformity is
full continuity ($C^0$): the field has no jumps anywhere.

\begin{physicsbox}
$\Hone$ is the home of the \textbf{electrostatic potential} $V$. The electrostatic
problem of \cref{ch:reductions-pde} solves $-\Div(\varepsilon\Grad V)=0$ for a
scalar $V$ that is continuous everywhere (a potential cannot jump---a jump would be
an infinite field). Nodal continuity is exactly the physics: $V$ must match across
faces, and $\Hone$ enforces precisely that and nothing more.
\end{physicsbox}

The same space serves any scalar field that must be continuous: temperature in heat
conduction, pressure in potential flow, the scalar magnetic potential. Whenever the
unknown is a scalar potential, $\Hone$ is the slot.

\subsection{$\Hcurl$: edge elements, the home of $\vE$ and $\vA$}

Now suppose the unknown is a vector field whose \emph{tangential} component must be
continuous---the electric field $\vE$, or the magnetic vector potential $\vA$. The
naive approach stores all three Cartesian components as separate $\Hone$ fields. We
will see in \cref{sec:spurious} that this is a disaster. The right approach uses a
different kind of DoF entirely.

A \emph{N\'ed\'elec edge element} assigns one DoF to each \emph{edge} of the mesh,
and the functional it reads is the \emph{circulation} of the field along that edge:
\[
  \sigma_e(\mathbf{u}) = \int_e \mathbf{u}\cdot d\vec\ell
  = \int_e \mathbf{u}\cdot\vec t_e \, d\ell,
\]
where $\vec t_e$ is the unit tangent along edge $e$. The basis function $\phi_e$ is
the lowest-order vector field whose circulation is $1$ along edge $e$ and $0$ along
every other edge.

Why does this give tangential continuity? Two cells sharing a face share the edges
bounding that face. Along each shared edge, the circulation DoF is a single number
owned by both cells, so both cells reconstruct a field with the \emph{same
tangential component} along that edge. The tangential component on the shared face
is therefore continuous, while the normal component is free to jump from cell to
cell---exactly the interface behaviour $\vE$ must have. This space is $\Hcurl$, the
space of vector fields whose curl is square-integrable; its discrete subspace
$\Hcurl_h$ is the \emph{edge-element} space.

\begin{physicsbox}
$\Hcurl$ is the home of the \textbf{electric field} $\vE$ and the \textbf{magnetic
vector potential} $\vA$. Every wave and magnetic analysis---eigenmode, driven,
magnetostatic, boundary mode---puts its primary unknown here. The curl--curl
operator $\Curl(\nu\Curl\,\cdot)$ of \cref{ch:reductions-pde} acts on an $\Hcurl$
field, and the tangential-continuity DoFs match the interface condition that
$\vE_{\text{tan}}$ (or $\vA_{\text{tan}}$) be continuous.
\end{physicsbox}

\subsection{$\Hdiv$: face elements, the home of fluxes}

The dual situation is a vector field whose \emph{normal} component must be
continuous---a flux density such as $\vB$, $\vD$, or a current $\vJ$. Here the DoFs
live on \emph{faces}. A \emph{Raviart--Thomas face element} assigns one DoF per face,
reading the \emph{flux} of the field through that face:
\[
  \sigma_f(\mathbf{u}) = \int_f \mathbf{u}\cdot\vn \, dA,
\]
with $\vn$ the unit normal to face $f$. The basis function $\phi_f$ has unit flux
through face $f$ and zero flux through every other face.

Two cells sharing a face share that face's flux DoF---a single number---so both
reconstruct a field with the \emph{same normal component} across the shared face,
while the tangential component may jump. This is $\Hdiv$, the space of vector fields
with square-integrable divergence; its discrete subspace $\Hdiv_h$ is the
\emph{face-element} (Raviart--Thomas) space. Normal continuity is exactly the
condition $\vB$ obeys: $B_n$ continuous, $B_{\text{tan}}$ free.

\subsection{$\Ltwo$: discontinuous elements, the home of densities}

Finally, some quantities have \emph{no} inter-element continuity requirement at all:
a charge density, an energy density, a flux divergence, the right-hand side of a
balance law. For these the appropriate space is $\Ltwo$, merely square-integrable,
with no constraint coupling neighbouring cells. Its discrete subspace $\Ltwo_h$ uses
\emph{discontinuous} elements: the DoFs are cell-interior values or moments, owned by
a single cell and shared with no one. A field in $\Ltwo_h$ may jump freely across
every face.

\subsection{The pattern}

Read \cref{tab:fourdofs} top to bottom and the pattern is unmistakable. As we descend
$\Hone \to \Hcurl \to \Hdiv \to \Ltwo$, the DoFs move from vertices to edges to faces
to cell interiors, and the continuity weakens from full, to tangential, to normal, to
none. \Cref{fig:fourspaces} shows the DoF locations on a single tetrahedron. This
ladder is not a coincidence: the gradient turns a nodal scalar into an edge field, the
curl turns an edge field into a face field, and the divergence turns a face field into
a cell density---the de Rham complex, the subject of \cref{ch:derham}. For now, the
take-away is the placement rule, which we can already use to choose a space for any
field whose continuity we know.

\begin{figure}[t]
\centering
\begin{tikzpicture}[scale=1.0, every node/.style={font=\scriptsize}]
% reusable tetrahedron macro via four corners
\def\A{(0,0)}      \def\B{(2.0,0.25)}   \def\C{(1.0,1.7)}   \def\D{(1.05,0.55)}
% --- panel layout: 2x2 grid of tetrahedra ---
\foreach \px/\py/\ttl/\sub in {%
  0/0/{$\Hone$ (nodal)}/{DoFs on vertices},%
  5/0/{$\Hcurl$ (edge)}/{DoFs on edges},%
  0/-3.4/{$\Hdiv$ (face)}/{DoFs on faces},%
  5/-3.4/{$\Ltwo$ (discontinuous)}/{DoF in interior}%
}{
  \begin{scope}[shift={(\px,\py)}]
    % tetra skeleton (back edges dashed)
    \coordinate (a) at (0,0);
    \coordinate (b) at (2.0,0.25);
    \coordinate (c) at (1.0,1.7);
    \coordinate (d) at (1.05,0.62);
    \draw[simGray,densely dashed] (a)--(d) (b)--(d) (c)--(d);
    \draw[simInk,thick] (a)--(b)--(c)--cycle;
    \node[font=\footnotesize\bfseries, simInk] at (1.0,2.15) {\ttl};
    \node[font=\scriptsize\itshape, simGray] at (1.0,-0.45) {\sub};
  \end{scope}
}
% --- H1: dots on vertices ---
\begin{scope}[shift={(0,0)}]
  \foreach \p in {(0,0),(2.0,0.25),(1.0,1.7),(1.05,0.62)}{
    \fill[simBlue] \p circle (2.6pt);}
\end{scope}
% --- Hcurl: arrows along edges ---
\begin{scope}[shift={(5,0)}]
  \coordinate (a) at (0,0); \coordinate (b) at (2.0,0.25);
  \coordinate (c) at (1.0,1.7); \coordinate (d) at (1.05,0.62);
  \foreach \u/\v in {a/b,b/c,c/a}{
    \draw[-{Stealth[length=1.8mm]},simRust,thick]
      ($(\u)!0.25!(\v)$)--($(\u)!0.75!(\v)$);}
  \foreach \u/\v in {a/d,b/d,c/d}{
    \draw[-{Stealth[length=1.6mm]},simRust,thick,opacity=0.6]
      ($(\u)!0.3!(\v)$)--($(\u)!0.7!(\v)$);}
\end{scope}
% --- Hdiv: a normal arrow through each face (front face shown) ---
\begin{scope}[shift={(0,-3.4)}]
  \coordinate (a) at (0,0); \coordinate (b) at (2.0,0.25);
  \coordinate (c) at (1.0,1.7); \coordinate (d) at (1.05,0.62);
  % front face centroid abc
  \coordinate (fabc) at ($(a)!0.5!(b)$);
  \fill[simGreen!18] (a)--(b)--(c)--cycle;
  \draw[simInk,thick] (a)--(b)--(c)--cycle;
  \node[circle,fill=simGreen,inner sep=1.4pt] (gc) at (1.0,0.65){};
  \draw[-{Stealth[length=2mm]},simGreen,thick] (gc)++(0.0,0.0)--+(0.35,-0.7);
  % a couple of side-face normals
  \node[circle,fill=simGreen,inner sep=1.1pt] (g2) at (0.55,0.55){};
  \draw[-{Stealth[length=1.6mm]},simGreen,thick,opacity=0.6] (g2)--+(-0.45,-0.15);
\end{scope}
% --- L2: single interior dot ---
\begin{scope}[shift={(5,-3.4)}]
  \fill[simViolet] (1.0,0.6) circle (3.0pt);
\end{scope}
\end{tikzpicture}
\caption[The four spaces and their DoF locations]{The hero picture: the same
tetrahedron, four spaces. $\Hone$ places its degrees of freedom on \emph{vertices}
(blue dots: nodal values); $\Hcurl$ on \emph{edges} (rust arrows: tangential
circulations); $\Hdiv$ on \emph{faces} (green normals: fluxes through faces); $\Ltwo$
in the \emph{interior} (violet dot: a single cell value, shared with no neighbour). The
DoF location \emph{is} the continuity: a quantity owned by a shared entity is continuous
across it. Reading the panels in order traces the de Rham ladder of \cref{ch:derham}.}
\label{fig:fourspaces}
\end{figure}

\section{Conformity: why the DoF location is the continuity}

We have asserted that the location of a DoF determines the continuity of the space.
This deserves a precise statement, because it is the mechanism behind everything in
the chapter.

When two cells share a geometric entity---a vertex, an edge, a face---and a DoF lives
\emph{on} that entity, both cells reference the \emph{same single coefficient} for that
DoF. Whatever quantity the functional reads is therefore identical as computed from
either side. The reconstructed fields, built from those coefficients, must agree on
that quantity across the interface. The DoF type fixes \emph{which} quantity:

\begin{itemize}[nosep]
\item a \emph{vertex value} shared $\Rightarrow$ the field value matches: $C^0$
  continuity ($\Hone$);
\item an \emph{edge circulation} shared $\Rightarrow$ the tangential component matches
  along the edge: tangential continuity ($\Hcurl$);
\item a \emph{face flux} shared $\Rightarrow$ the normal component matches across the
  face: normal continuity ($\Hdiv$);
\item \emph{nothing} shared (interior DoF) $\Rightarrow$ no constraint: discontinuity
  ($\Ltwo$).
\end{itemize}

\Cref{fig:tangnormal} draws the two vector cases side by side. The crucial observation
is that $\Hcurl$ and $\Hdiv$ each constrain \emph{exactly one} of the two pieces of a
vector field at an interface (tangential or normal) and leave the other free. This is
not a weakness; it is the point. The interface conditions of Maxwell's equations
constrain exactly one piece and leave the other free, with the freed piece carrying a
physical surface source. A space that constrained \emph{both} pieces (as componentwise
$\Hone$ does) would forbid the surface charges and surface currents that physics
permits---and would over-constrain the solution into the spurious modes of the next
section.

\begin{figure}[t]
\centering
\begin{tikzpicture}[scale=1.0, every node/.style={font=\scriptsize}]
% --- Left: tangential continuity (Hcurl) ---
\begin{scope}
  \node[font=\footnotesize\bfseries] at (1.4,2.2) {$\Hcurl$: tangential continuous};
  \draw[simInk,thick] (1.4,-0.2)--(1.4,1.9);
  \node[simGray] at (1.4,-0.5) {shared face};
  \fill[simBlue!8] (0,-0.2) rectangle (1.4,1.9);
  \fill[simRust!8] (1.4,-0.2) rectangle (2.8,1.9);
  % tangential (vertical) components: equal magnitude both sides
  \draw[-{Stealth[length=2mm]},simBlue,thick] (1.05,0.4)--(1.05,1.1);
  \draw[-{Stealth[length=2mm]},simRust,thick] (1.75,0.4)--(1.75,1.1);
  \node[simBlue] at (0.62,0.75) {$u_{\text{tan}}$};
  \node[simRust] at (2.2,0.75) {$u_{\text{tan}}$};
  % normal components: may differ
  \draw[-{Stealth[length=1.6mm]},simBlue,opacity=0.55] (1.05,1.5)--(0.6,1.5);
  \draw[-{Stealth[length=1.6mm]},simRust,opacity=0.55] (1.75,1.5)--(2.5,1.5);
  \node[simGray] at (1.4,-0.95) {$\;u_{\text{tan}}$ matches, $u_n$ may jump};
\end{scope}
% --- Right: normal continuity (Hdiv) ---
\begin{scope}[shift={(5,0)}]
  \node[font=\footnotesize\bfseries] at (1.4,2.2) {$\Hdiv$: normal continuous};
  \draw[simInk,thick] (1.4,-0.2)--(1.4,1.9);
  \node[simGray] at (1.4,-0.5) {shared face};
  \fill[simBlue!8] (0,-0.2) rectangle (1.4,1.9);
  \fill[simGreen!8] (1.4,-0.2) rectangle (2.8,1.9);
  % normal (horizontal) components: equal both sides
  \draw[-{Stealth[length=2mm]},simBlue,thick] (0.9,0.95)--(1.4,0.95);
  \draw[-{Stealth[length=2mm]},simGreen,thick] (1.4,0.95)--(1.9,0.95);
  \node[simBlue] at (0.55,1.25) {$u_n$};
  \node[simGreen] at (2.25,1.25) {$u_n$};
  % tangential: may differ
  \draw[-{Stealth[length=1.6mm]},simBlue,opacity=0.55] (1.05,0.35)--(1.05,0.0);
  \draw[-{Stealth[length=1.6mm]},simGreen,opacity=0.55] (1.75,0.35)--(1.75,0.75);
  \node[simGray] at (1.4,-0.95) {$\;u_n$ matches, $u_{\text{tan}}$ may jump};
\end{scope}
\end{tikzpicture}
\caption[Tangential versus normal continuity]{The two vector spaces each constrain one
component of a field across a shared face and free the other. \emph{Left:} $\Hcurl$
forces the tangential component $u_{\text{tan}}$ to match (the heavy vertical arrows are
equal) while the normal component $u_n$ may jump. \emph{Right:} $\Hdiv$ forces the
normal component $u_n$ to match (the heavy horizontal arrows are equal) while the
tangential component may jump. These are precisely the Maxwell interface conditions for
$\vE$ (tangential continuous) and $\vB$ (normal continuous) from \cref{ch:maxwell}.}
\label{fig:tangnormal}
\end{figure}

\section{The spurious-mode pitfall: why not just use nodal vectors?}
\label{sec:spurious}

We can now confront the question that motivates edge elements directly. Why not
discretize $\vE$ component by component in $\Hone$---store $E_x$, $E_y$, $E_z$ each as a
nodal scalar field? It is simpler to implement and uses tooling we already have. The
reason it fails is one of the most important lessons in computational electromagnetics.

A nodal vector field forces \emph{all three} components to be continuous at every
vertex, hence the \emph{whole vector} is continuous across faces---both its tangential
and its normal parts. But Maxwell's equations do \emph{not} require the normal part of
$\vE$ to be continuous; across a dielectric interface, or at a re-entrant metallic
corner, $\vE$ has a genuinely discontinuous normal component and a field singularity.
Forcing full continuity where physics demands only tangential continuity over-constrains
the discrete space: the true singular field cannot be represented, and the extra
constraints have to go somewhere.

Where they go is into \emph{spurious modes}---non-physical eigenvectors that pollute the
spectrum of the discrete curl--curl operator. In an eigenmode analysis
(\cref{ch:eigenmodes}) these appear as extra ``resonances'' at frequencies that no
physical cavity possesses, interleaved unpredictably with the real ones, impossible to
filter out reliably. \Cref{fig:spurious} sketches the mechanism at a re-entrant corner:
the nodal-vector field is pinned to a single value at the corner vertex and cannot
follow the field around it, while the edge-element field, owning independent tangential
circulations on the two corner edges, represents the singular field cleanly.

\begin{figure}[t]
\centering
\begin{tikzpicture}[scale=1.0, every node/.style={font=\scriptsize}]
% --- Left: nodal vector at a re-entrant corner ---
\begin{scope}
  \node[font=\footnotesize\bfseries] at (1.3,2.3) {nodal vector (wrong)};
  % L-shaped re-entrant domain
  \fill[simBgGray] (0,0)--(2.4,0)--(2.4,1.0)--(1.2,1.0)--(1.2,2.0)--(0,2.0)--cycle;
  \draw[simInk,thick] (0,0)--(2.4,0)--(2.4,1.0)--(1.2,1.0)--(1.2,2.0)--(0,2.0)--cycle;
  \fill[simRust] (1.2,1.0) circle (2.6pt);
  \node[simRust] at (1.62,1.18) {corner};
  % single pinned vector at corner -> cannot bend
  \draw[-{Stealth[length=2mm]},simInk,thick] (1.2,1.0)--+(0.55,0.25);
  % field tries to wrap but is forced single-valued (X marks)
  \node[simRust,font=\footnotesize] at (0.8,0.55) {$\times$};
  \node[simRust,font=\footnotesize] at (1.85,1.4) {$\times$};
  \node[simGray] at (1.2,-0.5) {spurious: over-constrained};
\end{scope}
% --- Right: edge element at the same corner ---
\begin{scope}[shift={(4.6,0)}]
  \node[font=\footnotesize\bfseries] at (1.3,2.3) {edge element (right)};
  \fill[simBgGreen] (0,0)--(2.4,0)--(2.4,1.0)--(1.2,1.0)--(1.2,2.0)--(0,2.0)--cycle;
  \draw[simInk,thick] (0,0)--(2.4,0)--(2.4,1.0)--(1.2,1.0)--(1.2,2.0)--(0,2.0)--cycle;
  \fill[simGreen] (1.2,1.0) circle (2.6pt);
  % independent edge circulations wrap around the corner
  \draw[-{Stealth[length=1.8mm]},simRust,thick] (1.2,1.0)--+(-0.0,0.55);
  \draw[-{Stealth[length=1.8mm]},simRust,thick] (1.2,1.0)--+(0.6,0.0);
  \draw[simBlue,thick,-{Stealth[length=1.6mm]}]
    (1.7,1.45) arc (35:-120:0.55);
  \node[simBlue] at (2.05,1.0) {$\vE$};
  \node[simGray] at (1.2,-0.5) {singular field captured};
\end{scope}
\end{tikzpicture}
\caption[Spurious modes at a re-entrant corner]{Why componentwise-nodal $\vE$ fails.
\emph{Left:} a nodal-vector field pins the whole vector to one value at the corner
vertex; it forces normal continuity that physics does not require and cannot follow the
field's singular wrap around a re-entrant corner ($\times$ marks the constraints that
have nowhere to go---they become spurious eigenmodes). \emph{Right:} an edge element owns
independent tangential circulations on the corner's edges, leaves the normal component
free, and represents the singular field cleanly. This is the reason every eigenmode and
driven analysis uses $\Hcurl$.}
\label{fig:spurious}
\end{figure}

\begin{pitfall}
Never discretize $\vE$, $\vA$, or any tangentially-continuous field as three
independent nodal ($\Hone$) components. Doing so imposes \emph{normal} continuity that
Maxwell does not require, over-constrains the discrete space, and injects
\textbf{spurious modes}---non-physical eigenvalues that corrupt every eigenmode and
driven solve and cannot be reliably filtered out. Use edge elements ($\Hcurl$), whose
DoFs constrain only the tangential component. This single design decision is why Palace's
eigenmode, driven, magnetostatic, and boundary-mode solvers all assemble their primary
operator over an $\Hcurl$ (\code{ND\_FECollection}) space.
\end{pitfall}

\section{From reference element to global space}

So far we have described DoFs and continuity abstractly. The actual construction has two
layers, which \cref{ch:refelement} develops in full; we sketch them here so the
machinery is in view.

\paragraph{The reference element.} Rather than define basis functions separately on every
cell of the mesh, we define them \emph{once} on a single \emph{reference element}---a
canonical unit tetrahedron or unit cube---and map them to each physical cell by an affine
or isoparametric transformation. Each of the four spaces has its own family of reference
shape functions: nodal Lagrange polynomials for $\Hone$, N\'ed\'elec edge functions for
$\Hcurl$, Raviart--Thomas face functions for $\Hdiv$, and discontinuous Lagrange or
Legendre polynomials for $\Ltwo$. The vector spaces $\Hcurl$ and $\Hdiv$ use special
\emph{covariant} (Piola) and \emph{contravariant} transformations that preserve tangential
and normal components respectively under the mapping---a detail that is what actually makes
the continuity survive the map to a curved or skewed cell.

\paragraph{Gluing into the global space.} The global finite-element space is assembled by
\emph{identifying} the DoFs that shared mesh entities have in common. Each physical edge
contributes one $\Hcurl$ DoF no matter how many cells border it; each physical face
contributes one $\Hdiv$ DoF; each vertex one $\Hone$ DoF. The local-to-global \emph{DoF
map}---a table recording, for each cell, which global DoF each of its local DoFs maps
to---is what stitches the per-cell reference bases into one continuous (or
tangentially/normally continuous) global field. In the matrix-free pipeline of
\cref{ch:matrixfree} this map is the \emph{element restriction} operator $\mathsf{G}$ that
gathers a global DoF vector into per-element local-dof tensors and scatters the result
back; the per-element local-dof axis is the $L$ axis of the element-local tensor
introduced there.

In Palace, this whole construction is a single object. The function space is built by
pairing a mesh with a \emph{collection} that names the family and order:
\begin{center}
\code{fe\_space(mesh, collection)} $\longrightarrow$
\code{FiniteElementSpace[N]},
\end{center}
where the collection's type ($\Hone$, $\Hcurl$, $\Hdiv$, or $\Ltwo$) selects the de Rham
slot and $N$ is the resulting dimension---the subject of \cref{sec:truedof}. The mesh and
the collection are independent inputs: the same collection over a finer mesh, or the same
mesh under a higher-order collection, are both well-formed spaces. These two axes of
refinement are the topic of the next section.

\section{Polynomial order and the \code{FECollection} family}

Each space comes in a tower of \emph{orders}. The order-$p$ version of a space replaces the
lowest-order shape functions with degree-$p$ polynomials, adding more DoFs per cell and
representing the field more accurately. Palace organizes the four families and their orders
through \emph{finite-element collections}: \code{H1\_FECollection}, \code{ND\_FECollection}
(N\'ed\'elec, $\Hcurl$), \code{RT\_FECollection} (Raviart--Thomas, $\Hdiv$), and
\code{L2\_FECollection}. A collection is the object passed as the second argument to
\code{fe\_space} above; its \emph{type} fixes the de Rham slot and its order fixes the
polynomial degree.

\Cref{tab:collections} collects the four families with the data that matters for building a
solver: the de Rham slot, the geometric entity that carries the DoFs, and the
\emph{minimum order} $p_{\min}$ each family admits.

\begin{table}[t]
\centering
\caption[The four \code{FECollection} families]{The four finite-element collection
families. The minimum order $p_{\min}$ differs by family: $\Hone$ and $\Hcurl$ require
$p_{\min}=1$ (a constant nodal scalar or a lowest-order edge function is the floor), while
$\Hdiv$ and $\Ltwo$ admit $p_{\min}=0$ (a piecewise-constant flux or cell value is
meaningful). These floors bound the coarsening sequence of the $p$-multigrid hierarchy.}
\label{tab:collections}
\small
\setlength{\tabcolsep}{6pt}
\renewcommand{\arraystretch}{1.3}
\begin{tabular}{@{}lllcl@{}}
\toprule
Space & Collection & DoF entity & $p_{\min}$ & primary field \\
\midrule
$\Hone$  & \code{H1\_FECollection} & vertices (nodes) & $1$ & scalar potential $V$ \\
$\Hcurl$ & \code{ND\_FECollection} & edges            & $1$ & $\vE$, $\vA$ \\
$\Hdiv$  & \code{RT\_FECollection} & faces            & $0$ & flux $\vB$, $\vD$ \\
$\Ltwo$  & \code{L2\_FECollection} & cell interiors   & $0$ & density, $B=\Curl\vE$ (2-D) \\
\bottomrule
\end{tabular}
\end{table}

\paragraph{$h$- versus $p$-refinement.} There are two independent ways to make a finite-element
solution more accurate, drawn in \cref{fig:hp}. \emph{$h$-refinement} subdivides the mesh into
smaller cells (the mesh size $h$ shrinks) while keeping the polynomial order fixed; it adds
resolution where the field varies rapidly and is the basis of the adaptive mesh refinement of
\cref{ch:amr}. \emph{$p$-refinement} raises the polynomial order $p$ on a fixed mesh, adding DoFs
per cell; on smooth fields it converges far faster (exponentially, rather than algebraically) than
$h$-refinement. The two are orthogonal---one varies the mesh argument of \code{fe\_space}, the
other the collection argument---and they can be combined ($hp$-refinement).

\begin{figure}[t]
\centering
\begin{tikzpicture}[scale=1.0, every node/.style={font=\scriptsize}]
% baseline single coarse element
\begin{scope}
  \node[font=\footnotesize\bfseries] at (1.0,1.65) {coarse, $p=1$};
  \draw[simInk,thick] (0,0) rectangle (2,1.2);
  \foreach \x/\y in {0/0,2/0,2/1.2,0/1.2}{\fill[simBlue] (\x,\y) circle (2pt);}
\end{scope}
% h-refinement: subdivide
\begin{scope}[shift={(3.4,0)}]
  \node[font=\footnotesize\bfseries] at (1.0,1.65) {$h$-refine ($p=1$)};
  \draw[simInk,thick] (0,0) rectangle (2,1.2);
  \draw[simGray] (1,0)--(1,1.2) (0,0.6)--(2,0.6);
  \foreach \x in {0,1,2}{\foreach \y in {0,0.6,1.2}{\fill[simBlue] (\x,\y) circle (2pt);}}
  \node[simGray] at (1.0,-0.5) {smaller cells};
\end{scope}
% p-refinement: more DoFs, same cell
\begin{scope}[shift={(6.8,0)}]
  \node[font=\footnotesize\bfseries] at (1.0,1.65) {$p$-refine ($p=2$)};
  \draw[simInk,thick] (0,0) rectangle (2,1.2);
  \foreach \x in {0,1,2}{\foreach \y in {0,0.6,1.2}{\fill[simRust] (\x,\y) circle (2pt);}}
  \node[simGray] at (1.0,-0.5) {higher order};
\end{scope}
\end{tikzpicture}
\caption[$h$- versus $p$-refinement]{Two orthogonal routes to accuracy, shown for a nodal
($\Hone$) element. \emph{$h$-refinement} (middle) subdivides the cell, shrinking the mesh size $h$
while keeping order $p=1$; the DoFs (blue) multiply by subdivision. \emph{$p$-refinement} (right)
keeps the single cell but raises the order to $p=2$, adding interior and edge DoFs (rust). Palace's
geometric $p$-multigrid (\cref{ch:multigrid}) builds a hierarchy of order-coarsened collections; the
$p_{\min}$ floors of \cref{tab:collections} bound how far the order can drop.}
\label{fig:hp}
\end{figure}

\paragraph{The $p$-multigrid order schedule.} Palace's geometric multigrid preconditioner
(\cref{ch:multigrid}) needs not a single space but a \emph{tower} of them, one per multigrid level,
differing in order. The collection-schedule operator produces a list of collections from the finest
order $p$ down to the floor $p_{\min}$, coarsening the order level by level. Two coarsening policies
are offered: \emph{linear}, which steps $p \mapsto p-1$ each level, and \emph{logarithmic} (the
default), which halves the gap to the floor, $p \mapsto \lfloor (p+p_{\min})/2 \rfloor$, reaching the
floor in far fewer levels for large $p$. The schedule terminates at $p_{\min}$---which is why the
floor differs by family: an $\Hone$ or $\Hcurl$ tower bottoms out at order $1$ (a constant scalar or
lowest edge element), while an $\Hdiv$ or $\Ltwo$ tower can descend to the piecewise-constant order
$0$. When the tower has a single level, the schedule degenerates to one ordinary \code{fe\_space}
construction.

\section{True degrees of freedom: the number $N$}
\label{sec:truedof}

We close with the single most-used number in this book. The dimension of the discrete space
$V_h$---the number of basis functions, hence the length of the coefficient vector, hence the size
of every assembled matrix---is denoted
\[
  N = \dim V_h.
\]
This is the $N$ in every \code{Tensor[N]} and every \code{LinearOperator[N,N]} throughout the
framework chapters; the function-space construction is the operator that \emph{defines} $N$.

There is a subtlety. The raw count of DoFs over all cells---summing local DoFs cell by cell---over-counts, because DoFs on shared entities are identified. Worse, on an adaptively refined mesh, a
\emph{hanging node} (a vertex of a fine cell lying in the middle of a coarse cell's face) is not an
independent DoF: its value is \emph{constrained} to be an interpolation of the surrounding coarse
DoFs, so that the global field stays conforming. The space therefore distinguishes two counts:

\begin{itemize}[nosep]
\item the \emph{local} (or \emph{L-vector}) DoF count, summed over cells with shared and hanging
  DoFs still counted multiply---the natural size for per-element work;
\item the \emph{true} DoF count $N$, with shared DoFs identified once and hanging-node constraints
  eliminated---the number of genuinely independent coefficients, the dimension of $V_h$.
\end{itemize}

It is the \emph{true}-DoF count $N$ that the solver works with: the assembled operators are $N\times
N$, the Krylov vectors of \cref{ch:krylov} are length $N$, the unknowns are $N$ true coefficients.
The transfer between the two pictures---gather the $N$ true coefficients into the larger local
vector, apply the constraints, scatter back---is exactly the prolongation/restriction pair the space
carries, and exactly the boundary the element-restriction operator $\mathsf{G}$ of
\cref{ch:matrixfree} sits on. For the un-refined, single-order meshes of most of this book there are
no hanging nodes, and $N$ is simply ``DoFs with shared entities counted once.''

\begin{notationbox}
$N = \dim V_h$ is the \textbf{true-DoF count}: the number of independent coefficients after shared
DoFs are identified and hanging-node constraints are removed. It is the dimension of every
\code{Tensor[N]} and the size of every \code{LinearOperator[N,N]} in the book. Different analyses on
the same mesh have different $N$, because they use different spaces---an $\Hone$ scalar problem and
an $\Hcurl$ vector problem over one mesh produce entirely different $N$.
\end{notationbox}

\section{Worked examples}

\subsection{Counting degrees of freedom}

The placement rule of \cref{tab:fourdofs} makes the lowest-order DoF count a purely combinatorial
exercise: count vertices for $\Hone$, edges for $\Hcurl$, faces for $\Hdiv$, cells for $\Ltwo$. We do
it first on a single tetrahedron, then on a small two-element mesh, to see how sharing changes the
arithmetic.

\begin{workedexample}{Lowest-order DoFs on a tetrahedron and a 2-element mesh}{dofcount}
A single tetrahedron has $4$ vertices, $6$ edges, $4$ faces, and $1$ cell. For the lowest-order
spaces the DoF count is the count of the carrying entity:
\[
  N_{\Hone} = 4, \quad N_{\Hcurl} = 6, \quad N_{\Hdiv} = 4, \quad N_{\Ltwo} = 1.
\]
So even on one cell the three first spaces have genuinely different sizes---vertices, edges, and faces
of a tetrahedron number $4$, $6$, $4$. An $\Hcurl$ problem is the largest; this is the price of getting
the continuity right.

Now glue a \emph{second} tetrahedron onto the first across a shared triangular face. The two cells share
that face's $3$ vertices and $3$ edges and the face itself; everything else is new. Counting the
combined mesh:
\begin{itemize}[nosep]
\item \textbf{vertices:} $4 + 1 = 5$. (The second tet adds one new apex; its other three vertices are
  the shared face.) So $N_{\Hone} = 5$.
\item \textbf{edges:} the first tet has $6$; the second adds the $3$ edges from its new apex to the
  shared face's vertices, the other $3$ being shared. So $6 + 3 = 9$, giving $N_{\Hcurl} = 9$.
\item \textbf{faces:} the first tet has $4$; the second adds $3$ new faces (the fourth is the shared
  face). So $4 + 3 = 7$, but the \emph{shared} face is counted once: $N_{\Hdiv} = 7$.
\item \textbf{cells:} $2$, so $N_{\Ltwo} = 2$.
\end{itemize}
The naive ``sum over cells'' count would give $\Hone$: $4+4=8$, $\Hcurl$: $6+6=12$, $\Hdiv$:
$4+4=8$. The true counts $5,9,7$ are smaller because the shared vertices, edges, and the shared face
are \emph{identified}---each contributes a single global DoF. That identification is exactly the
gluing of \cref{sec:truedof} and the source of the gap between the local and true DoF counts. The
amount of sharing differs by space: the shared face removes $3$ from the vertex count, $3$ from the
edge count, but only $1$ from the face count.
\end{workedexample}

\subsection{A lowest-order edge basis function}

The defining feature of $\Hcurl$---tangential support on a single edge---can be checked by hand on the
reference triangle. We exhibit the lowest-order N\'ed\'elec basis function for one edge and verify its
tangential trace.

\begin{workedexample}{The N\'ed\'elec edge function and its tangential trace}{whitney}
Take the reference triangle with vertices $\vx_0=(0,0)$, $\vx_1=(1,0)$, $\vx_2=(0,1)$, and let
$\lambda_0,\lambda_1,\lambda_2$ be its barycentric (nodal $\Hone$) coordinates---the hat functions with
$\lambda_i(\vx_j)=\delta_{ij}$. Explicitly $\lambda_0 = 1-x-y$, $\lambda_1 = x$, $\lambda_2 = y$, and
note $\Grad\lambda_0=(-1,-1)$, $\Grad\lambda_1=(1,0)$, $\Grad\lambda_2=(0,1)$.

The lowest-order edge (Whitney) function associated with the edge $e_{01}$ joining $\vx_0$ and $\vx_1$
is
\[
  \boldsymbol{\phi}_{01} \;=\; \lambda_0\,\Grad\lambda_1 \;-\; \lambda_1\,\Grad\lambda_0
  \;=\; (1-x-y)\,(1,0) \;-\; x\,(-1,-1)
  \;=\; \bigl(\,1-y,\; x\,\bigr).
\]
This is a genuine vector field on the triangle. We check the three edge circulations
$\sigma_e(\boldsymbol{\phi}_{01})=\int_e \boldsymbol{\phi}_{01}\cdot d\vec\ell$.

\smallskip\noindent\emph{Edge $e_{01}$} (its own edge), from $\vx_0=(0,0)$ to $\vx_1=(1,0)$: here
$y=0$, the tangent is $d\vec\ell=(dx,0)$, and $\boldsymbol{\phi}_{01}=(1-0,\,x)=(1,x)$, so
$\boldsymbol{\phi}_{01}\cdot d\vec\ell = 1\,dx$. The circulation is
$\int_0^1 1\,dx = 1.$

\smallskip\noindent\emph{Edge $e_{02}$}, from $\vx_0=(0,0)$ to $\vx_2=(0,1)$: here $x=0$, the tangent is
$(0,dy)$, and $\boldsymbol{\phi}_{01}=(1-y,\,0)$, so the dot product
$\boldsymbol{\phi}_{01}\cdot d\vec\ell = 0$. The circulation is $0.$

\smallskip\noindent\emph{Edge $e_{12}$}, from $\vx_1=(1,0)$ to $\vx_2=(0,1)$: parametrize
$(x,y)=(1-t,\,t)$ for $t\in[0,1]$, so $d\vec\ell=(-1,1)\,dt$ and
$\boldsymbol{\phi}_{01}=(1-t,\,1-t)$. Then
$\boldsymbol{\phi}_{01}\cdot d\vec\ell = (1-t)(-1)+(1-t)(1) = 0$,
and the circulation is $0.$

\smallskip So $\sigma_{e_{01}}(\boldsymbol{\phi}_{01})=1$ and
$\sigma_{e_{02}}=\sigma_{e_{12}}=0$: the function satisfies the duality relation
$\sigma_e(\boldsymbol{\phi}_{01})=\delta_{e,e_{01}}$. Its tangential trace is supported on \emph{exactly
one} edge. Consequently, when two triangles share edge $e_{01}$, both reference $\boldsymbol{\phi}_{01}$
through the \emph{same} circulation coefficient, and the assembled field has continuous tangential
component along $e_{01}$ and nowhere else constrained---tangential continuity, edge by edge, exactly as
\cref{fig:tangnormal} promised.
\end{workedexample}

\summarysection
A finite-element discretization must place each field in a space whose degrees of freedom enforce
exactly the continuity the physics requires. There are four canonical spaces, distinguished entirely by
their DoFs: $\Hone$ stores nodal \emph{values} on vertices and is fully continuous (the home of the
scalar potential $V$); $\Hcurl$ (N\'ed\'elec edge elements) stores \emph{circulations} on edges and is
tangentially continuous (the home of $\vE$ and $\vA$); $\Hdiv$ (Raviart--Thomas face elements) stores
\emph{fluxes} on faces and is normally continuous (the home of $\vB$, $\vD$); $\Ltwo$ stores
\emph{cell} quantities and has no inter-element continuity (the home of densities). Because a DoF on a
shared entity is a single coefficient owned by both neighbours, the DoF \emph{location} is the
continuity. Discretizing $\vE$ as componentwise nodal vectors over-constrains it---it imposes normal
continuity physics does not require---and injects spurious modes that ruin eigenmode and driven solves;
edge elements get the continuity exactly right, which is why every magnetic and wave analysis uses
$\Hcurl$. Each space comes in a tower of polynomial orders with family-dependent floors
($p_{\min}=1$ for $\Hone$/$\Hcurl$, $0$ for $\Hdiv$/$\Ltwo$), assembled from reference shape functions
glued by a local-to-global DoF map; refinement is $h$ (smaller cells) or $p$ (higher order). The
dimension $N=\dim V_h$---the true-DoF count, with shared DoFs identified and hanging-node constraints
removed---is the $N$ that names every tensor in the book.

\furtherreading
The definitive reference for finite-element spaces in electromagnetics is
Monk~\citep{monk2003}, which develops $\Hcurl$ and $\Hdiv$, the N\'ed\'elec and Raviart--Thomas
elements, and the spurious-mode question with full rigour. The unifying structural view---the four
spaces as a discrete de Rham complex with exactness---is the subject of finite-element exterior
calculus; Arnold, Falk, and Winther~\citep{arnold2006feec} is the canonical account and the natural
sequel to \cref{ch:derham}. For an engineering-oriented treatment tied directly to computational
electromagnetics, with the spurious-mode pathology worked through on real cavity problems,
Jin~\citep{jin2014} is the standard text. The continuity arguments of \cref{sec:spurious} are made
precise there; the tangential/normal trace theory underlying them is Monk's Chapter~3.

\exercisessection
\begin{enumerate}[nosep]
\item For each of the four spaces, state in one sentence which Maxwell interface condition its DoFs
  enforce, and name a physical field that belongs in it. (Use \cref{tab:fourdofs} and the
  \texttt{physicsbox} anchors.)
\item A single hexahedron (cube) has $8$ vertices, $12$ edges, $6$ faces, and $1$ cell. Give the
  lowest-order DoF count $N$ for each of the four spaces on one hexahedron. Which space is largest?
\item Two hexahedra are glued across a shared square face. Following
  \cref{wex:dofcount}, compute the true DoF count $N$ for $\Hone$, $\Hcurl$, $\Hdiv$, and $\Ltwo$ on the
  combined mesh, and compare each to the naive sum-over-cells count. How many DoFs does the shared face
  remove from each space?
\item Explain, in terms of degrees of freedom, why discretizing $\vE$ as three independent $\Hone$
  scalar fields imposes \emph{more} continuity than $\Hcurl$ does. Which interface condition does the
  nodal-vector approach wrongly enforce, and what physical surface quantity does that forbid?
\item The $p$-multigrid order schedule for an $\Hcurl$ ($p_{\min}=1$) space at finest order $p=8$
  produces a tower of collections. List the order sequence under (a) linear coarsening
  ($p\mapsto p-1$) and (b) logarithmic coarsening ($p\mapsto\lfloor(p+1)/2\rfloor$). How many levels
  does each policy produce before reaching the floor?
\item Repeat the previous exercise for an $\Hdiv$ ($p_{\min}=0$) space at $p=8$. Why can this tower
  descend one order further than the $\Hcurl$ tower, and what does the order-$0$ space represent
  physically?
\item Verify, by the method of \cref{wex:whitney}, that the edge function
  $\boldsymbol{\phi}_{02}=\lambda_0\Grad\lambda_2-\lambda_2\Grad\lambda_0$ on the reference triangle has
  circulation $1$ along edge $e_{02}$ and $0$ along $e_{01}$ and $e_{12}$.
\item The number $N=\dim V_h$ differs between an electrostatic ($\Hone$) and an eigenmode ($\Hcurl$)
  analysis run on the \emph{same} mesh. Using the single-tetrahedron and two-tet counts of
  \cref{wex:dofcount}, explain why, and state which analysis produces the larger linear system on a given
  mesh.
\end{enumerate}

\chapter{The de Rham Complex and Discrete Exactness}
\label{ch:derham}

\chapterorient{The previous chapter introduced the four finite-element spaces one
at a time, by the degrees of freedom that define them, and ended with a promise:
that $\Hone$, $\Hcurl$, $\Hdiv$, and $\Ltwo$ are not four unrelated spaces but the
four rungs of a single ladder, linked by the gradient, curl, and divergence. This
chapter keeps that promise. It is the structural heart of Part~II. We assemble the
four spaces into the \emph{de Rham complex}, prove the two vector-calculus
identities $\Curl\Grad=0$ and $\Div\Curl=0$ that make it a complex, and read off
their physical content---that a gradient field is curl-free and a curl field is
divergence-free, the very facts that license the scalar and vector potentials of
\cref{ch:maxwell}. Then comes the payoff that makes the whole apparatus matter for
a simulator: the discrete finite-element spaces inherit the \emph{same} exactness,
through a commuting diagram relating interpolation to differentiation. That single
property---discrete exactness---is why edge elements avoid spurious modes, why the
eigenmode solver can project onto a divergence-free subspace, and why the gradient
fields form the null space the curl--curl operator must learn to ignore. By the end
you will be able to draw the complex from memory, verify exactness by hand on a
small mesh, and recognize the de Rham structure wherever it surfaces in the
solver chapters ahead.}

\section{Four spaces, three arrows}

\Cref{ch:functionspaces} left us with \cref{tab:fourdofs}: four spaces, each
defined by where its degrees of freedom live and what continuity that forces.
Reading the table top to bottom traced a ladder---vertices, edges, faces, cell
interiors; full continuity, tangential, normal, none. We now make the rungs of
that ladder into a single mathematical object by naming the \emph{arrows} between
them.

The three arrows are the differential operators of vector calculus. The gradient
$\Grad$ turns a scalar field into a vector field; the curl $\Curl$ turns a vector
field into another vector field; the divergence $\Div$ turns a vector field into a
scalar field. What \cref{ch:functionspaces} hinted at, and we now assert
precisely, is that each of these operators maps one of our four spaces
\emph{exactly into the next}:
\[
  \Hone \xrightarrow{\;\Grad\;} \Hcurl \xrightarrow{\;\Curl\;} \Hdiv
         \xrightarrow{\;\Div\;} \Ltwo .
\]
This sequence is the \emph{de Rham complex}. It is the central diagram of this
book, and we draw it properly in \cref{fig:derham}.

\begin{figure}[t]
\centering
\begin{tikzpicture}[node distance=24mm and 30mm]
  % --- the four space nodes ---
  \node[space, minimum width=18mm] (h1)   {$\Hone$};
  \node[space, minimum width=18mm, right=of h1]   (hcurl) {$\Hcurl$};
  \node[space, minimum width=18mm, right=of hcurl] (hdiv)  {$\Hdiv$};
  \node[space, minimum width=18mm, right=of hdiv]  (l2)    {$\Ltwo$};
  % --- the three differential arrows ---
  \draw[flow] (h1)    -- node[above,font=\small]{$\Grad$}
                         node[below,font=\scriptsize\itshape,simGray]{gradient} (hcurl);
  \draw[flow] (hcurl) -- node[above,font=\small]{$\Curl$}
                         node[below,font=\scriptsize\itshape,simGray]{curl} (hdiv);
  \draw[flow] (hdiv)  -- node[above,font=\small]{$\Div$}
                         node[below,font=\scriptsize\itshape,simGray]{divergence} (l2);
  % --- DoF annotations under each node ---
  \node[font=\scriptsize, simBlue,   below=10mm of h1.south,    align=center] (a1)
    {nodal\\(vertices)\\value $V$};
  \node[font=\scriptsize, simRust,   below=10mm of hcurl.south, align=center] (a2)
    {edge\\(circulations)\\$\vE,\vA$};
  \node[font=\scriptsize, simGreen,  below=10mm of hdiv.south,  align=center] (a3)
    {face\\(fluxes)\\$\vB,\vD$};
  \node[font=\scriptsize, simViolet, below=10mm of l2.south,    align=center] (a4)
    {cell\\(densities)\\$\rho$};
  \draw[refedge] (h1.south)    -- (a1.north);
  \draw[refedge] (hcurl.south) -- (a2.north);
  \draw[refedge] (hdiv.south)  -- (a3.north);
  \draw[refedge] (l2.south)    -- (a4.north);
  % --- the two exactness identities, above the arrows ---
  \node[font=\scriptsize\itshape, simInk, above=12mm of hcurl.north] (id1)
    {$\Curl\Grad=0$};
  \node[font=\scriptsize\itshape, simInk, above=12mm of hdiv.north]  (id2)
    {$\Div\Curl=0$};
  \draw[refedge, bend right=18] (h1.north)  to (id1);
  \draw[refedge, bend left=18]  (hdiv.north) to (id1);
  \draw[refedge, bend right=18] (hcurl.north) to (id2);
  \draw[refedge, bend left=18]  (l2.north)    to (id2);
\end{tikzpicture}
\caption[The de Rham complex]{The de Rham complex: the four finite-element spaces
of \cref{ch:functionspaces}, linked by the three differential operators. Each space
carries a different \emph{kind} of degree of freedom (blue nodal values, rust edge
circulations, green face fluxes, violet cell densities) and each arrow is a
\emph{first-order} differential operator that maps each space exactly into the next.
The two identities above the arrows---$\Curl\Grad=0$ and $\Div\Curl=0$---are what make
this sequence a \emph{complex}: applying two consecutive arrows annihilates everything.
This single picture organizes the whole of electromagnetic finite-element theory;
it recurs throughout the book.}
\label{fig:derham}
\end{figure}

Before we can call it a complex we must check that the arrows even \emph{type-check}:
that each operator genuinely lands in the next space. The four spaces were defined
in \cref{ch:functionspaces} precisely so that they do.

\begin{itemize}[nosep]
\item \textbf{$\Grad : \Hone \to \Hcurl$.} A field $u\in\Hone$ is square-integrable
  with a square-integrable gradient---that is the very definition of $\Hone$. So
  $\Grad u$ is square-integrable. Is it in $\Hcurl$? A field is in $\Hcurl$ if its
  curl is square-integrable; and $\Curl\Grad u = 0$ (proven below), which is
  certainly square-integrable. So $\Grad u\in\Hcurl$. The tangential trace of
  $\Grad u$ is the tangential derivative of the continuous scalar $u$---continuous
  across faces---so $\Grad u$ has exactly the tangential continuity that $\Hcurl$
  requires.
\item \textbf{$\Curl : \Hcurl \to \Hdiv$.} A field $\vu\in\Hcurl$ has
  square-integrable curl by definition, so $\Curl\vu$ is square-integrable. Is it in
  $\Hdiv$? A field is in $\Hdiv$ if its divergence is square-integrable; and
  $\Div\Curl\vu=0$, square-integrable. So $\Curl\vu\in\Hdiv$. The normal trace of a
  curl is determined by the tangential data of $\vu$ on the face, which $\Hcurl$
  makes continuous---so $\Curl\vu$ has the normal continuity $\Hdiv$ requires.
\item \textbf{$\Div : \Hdiv \to \Ltwo$.} A field $\vu\in\Hdiv$ has square-integrable
  divergence by definition. So $\Div\vu\in\Ltwo$, no constraint to check: $\Ltwo$
  asks for nothing beyond square-integrability.
\end{itemize}

The arrows fit. Each space is built with just enough regularity that the
differential operator leaving it lands in the next space with just enough regularity
to leave again. This is no accident of our four choices: it is the reason these four
spaces, and not others, are the canonical ones. They are the spaces \emph{on which
the de Rham complex closes}.

\begin{notationbox}
We write the complex left to right, $\Hone\to\Hcurl\to\Hdiv\to\Ltwo$, with the
operators $\Grad,\Curl,\Div$ over the arrows. In three dimensions this is the full
complex; in two dimensions it shortens to
$\Hone\xrightarrow{\Grad}\Hcurl\xrightarrow{\operatorname{curl}}\Ltwo$, where the
\emph{scalar} curl $\operatorname{curl}(u_1,u_2)=\partial_x u_2-\partial_y u_1$ plays
the role of both curl and divergence at once. Palace builds the discrete arrows for
both cases; the 2-D scalar curl is the \code{H\_CURL}$\to$\code{INTEGRAL} branch of
its interpolator (\cref{sec:discrete-arrows}). Throughout this chapter ``the
complex'' means the 3-D version unless stated otherwise.
\end{notationbox}

\section{Exactness: range equals kernel}

A sequence of spaces and arrows is a \emph{complex} when the composition of any two
consecutive arrows is zero. It is \emph{exact} when something stronger holds: the
\emph{range} of each arrow is exactly the \emph{kernel} of the next. We take these
two properties in turn, because they are different in strength and different in
consequence, and the discrete theory will need to track them separately.

\subsection{The complex property: two consecutive arrows vanish}

The complex property is the pair of vector-calculus identities
\[
  \Curl\,\Grad = 0, \qquad \Div\,\Curl = 0.
\]
These say that the gradient of any scalar has zero curl, and the curl of any vector
has zero divergence. They are usually quoted as facts; we derive them, because the
derivation is short, it is the engine of everything that follows, and---crucially---
the \emph{discrete} versions will be derived the same way.

\begin{lemma}[The complex identities]\label{lem:complex}
For any twice-continuously-differentiable scalar field $\phi$ and vector field $\vF$,
\[
  \Curl(\Grad\phi) = \mathbf{0}, \qquad \Div(\Curl\vF) = 0.
\]
\end{lemma}

\begin{proof}
Both follow from a single fact: \emph{mixed partial derivatives commute},
$\partial_i\partial_j\phi=\partial_j\partial_i\phi$ (Clairaut's theorem, valid for
$C^2$ fields).

For the first identity, compute the curl of a gradient component by component. With
$\Grad\phi=(\partial_x\phi,\partial_y\phi,\partial_z\phi)$, the $x$-component of the
curl is
\[
  (\Curl\Grad\phi)_x
  = \partial_y(\Grad\phi)_z - \partial_z(\Grad\phi)_y
  = \partial_y\partial_z\phi - \partial_z\partial_y\phi
  = 0,
\]
the last step by commuting the mixed partials. The $y$- and $z$-components vanish
identically the same way. Hence $\Curl\Grad\phi=\mathbf 0$.

For the second identity, with $\vF=(F_1,F_2,F_3)$,
\[
  \Curl\vF
  = \bigl(\partial_y F_3 - \partial_z F_2,\;
          \partial_z F_1 - \partial_x F_3,\;
          \partial_x F_2 - \partial_y F_1\bigr),
\]
and its divergence is
\[
  \Div\Curl\vF
  = \partial_x(\partial_y F_3 - \partial_z F_2)
  + \partial_y(\partial_z F_1 - \partial_x F_3)
  + \partial_z(\partial_x F_2 - \partial_y F_1).
\]
The six terms cancel in pairs: $\partial_x\partial_y F_3$ against
$-\partial_y\partial_x F_3$, $-\partial_x\partial_z F_2$ against
$\partial_z\partial_x F_2$, and $\partial_y\partial_z F_1$ against
$-\partial_z\partial_y F_1$, each pair by commuting the mixed partials. The total is
zero.
\end{proof}

Both identities are corollaries of one symmetry---that the order of differentiation
does not matter. We will revisit this in \cref{wex:component}, where we watch the
cancellation happen on a concrete polynomial field, and again in
\cref{sec:discrete-exact}, where the \emph{discrete} gradient and curl satisfy the
same identities for a discrete reason that is structurally identical.

\subsection{Exactness: the kernel is filled by the previous range}

The complex property says $\operatorname{range}(\Grad)\subseteq\ker(\Curl)$: every
gradient is curl-free, so the range of $\Grad$ sits inside the kernel of $\Curl$.
\emph{Exactness} asserts the reverse inclusion as well---that the kernel contains
\emph{nothing but} previous-range elements:
\[
  \operatorname{range}(\Grad) = \ker(\Curl), \qquad
  \operatorname{range}(\Curl) = \ker(\Div).
\]
In words: \emph{every} curl-free field is a gradient, and \emph{every}
divergence-free field is a curl. \Cref{fig:exact} draws the two pictures side by
side---the easy inclusion (a complex) and the full equality (exactness).

\begin{figure}[t]
\centering
\begin{tikzpicture}[scale=1.0, every node/.style={font=\scriptsize}]
% --- Left: complex (one inclusion) ---
\begin{scope}
  \node[font=\footnotesize\bfseries] at (1.5,2.55) {complex: range $\subseteq$ kernel};
  % kernel of Curl as a big ellipse inside Hcurl
  \draw[simInk,thick] (1.5,1.0) ellipse (1.6 and 1.25);
  \node[simInk] at (1.5,2.05) {$\ker(\Curl)$};
  % range of Grad as a smaller blob inside
  \fill[simBlue!18] (1.2,0.85) ellipse (0.85 and 0.6);
  \draw[simBlue,thick] (1.2,0.85) ellipse (0.85 and 0.6);
  \node[simBlue] at (1.2,0.85) {$\operatorname{range}(\Grad)$};
  \node[simRust,font=\scriptsize] at (2.35,0.25) {gap?};
  \node[simGray,align=center] at (1.5,-0.7) {curl-free fields not\\[-1pt]known to be gradients};
\end{scope}
% --- Right: exactness (equality) ---
\begin{scope}[shift={(5.6,0)}]
  \node[font=\footnotesize\bfseries] at (1.5,2.55) {exact: range $=$ kernel};
  \draw[simInk,thick] (1.5,1.0) ellipse (1.6 and 1.25);
  \node[simInk] at (1.5,2.05) {$\ker(\Curl)$};
  \fill[simBlue!18] (1.5,1.0) ellipse (1.55 and 1.2);
  \draw[simBlue,thick] (1.5,1.0) ellipse (1.55 and 1.2);
  \node[simBlue] at (1.5,1.0) {$\operatorname{range}(\Grad)$};
  \node[simGreen,font=\scriptsize] at (1.5,-0.05) {no gap};
  \node[simGray,align=center] at (1.5,-0.7) {every curl-free field\\[-1pt]\emph{is} a gradient};
\end{scope}
\end{tikzpicture}
\caption[Complex versus exact]{The difference between a complex and an exact
sequence, drawn at the $\Hcurl$ rung. \emph{Left:} the complex property
$\Curl\Grad=0$ guarantees only that the range of $\Grad$ (blue) sits \emph{inside}
the kernel of $\Curl$ (the curl-free fields); a gap may remain---curl-free fields
that are not gradients. \emph{Right:} \emph{exactness} closes the gap: the range
fills the kernel exactly, so every curl-free field is a gradient. The size of the
gap is a property of the domain's \emph{topology} (\cref{rmk:cohomology}); for the
simply-connected domains of most simulations the gap is empty and the sequence is
exact.}
\label{fig:exact}
\end{figure}

The reverse inclusions are not automatic---and whether they hold is a statement
about the \emph{shape} of the domain, not about calculus. On a domain with no holes
(a ball, a box, any \emph{simply-connected} region with connected boundary) the
sequence is exact: a curl-free field on such a domain really is a gradient, and a
divergence-free field really is a curl. On a domain with a hole---a solid torus, the
region around a wire---the reverse inclusion can fail, and the failure is measured by
the topology. That refinement is the subject of \cref{rmk:cohomology}; for now we
record the clean statement for the domains where it holds.

\begin{theorem}[Exactness on a simply-connected domain]\label{thm:exact}
Let $\dom\subset\Real^3$ be a bounded, simply-connected domain with connected
boundary. Then the de Rham complex on $\dom$ is exact:
\[
  \operatorname{range}(\Grad)=\ker(\Curl), \qquad
  \operatorname{range}(\Curl)=\ker(\Div).
\]
Equivalently: every curl-free field is the gradient of a scalar potential, and every
divergence-free field is the curl of a vector potential.
\end{theorem}

The proof constructs the potential explicitly---one integrates the field along paths
from a base point, and simple-connectivity guarantees the answer is
path-independent. We do not reproduce it (it is the Poincar\'e-lemma computation of
any vector-calculus course); the references in \cref{sec:furtherreading} give it in
full. What matters here is not the proof but the \emph{use}: exactness is the
mathematical fact behind the electromagnetic potentials, and it is the property the
discrete spaces must preserve.

\section{The potentials of Maxwell are de Rham facts}
\label{sec:potentials}

\Cref{ch:maxwell} introduced two potentials almost in passing: the scalar potential
$V$ with $\vE=-\Grad V$, and the magnetic vector potential $\vA$ with $\vB=\Curl\vA$.
We can now see that neither was a lucky guess. Each is forced---\emph{exists at
all}---by the exactness of the de Rham complex.

Two of Maxwell's equations are homogeneous (source-free) in the regimes the
simulator solves:
\[
  \Curl\vE = 0 \quad\text{(electrostatics)}, \qquad
  \Div\vB = 0 \quad\text{(always).}
\]
Read each through \cref{thm:exact}. The electrostatic field $\vE$ is curl-free; by
exactness at the $\Hcurl$ rung, a curl-free field is a gradient, so there exists a
scalar $V$ with $\vE=-\Grad V$ (the sign is convention). The magnetic field $\vB$ is
divergence-free---always, this is the ``no magnetic monopoles'' law; by exactness at
the $\Hdiv$ rung, a divergence-free field is a curl, so there exists a vector field
$\vA$ with $\vB=\Curl\vA$. \Cref{fig:potentials} lays the correspondence out.

\begin{figure}[t]
\centering
\begin{tikzpicture}[scale=1.0, every node/.style={font=\footnotesize}]
% --- top row: electrostatic / scalar potential ---
\begin{scope}
  \node[stage, minimum width=30mm] (e1) {$\Curl\vE=0$};
  \node[font=\scriptsize\itshape,simGray, right=8mm of e1] (e2) {exactness at $\Hcurl$};
  \node[stage, minimum width=30mm, right=8mm of e2] (e3) {$\vE=-\Grad V$};
  \draw[flow] (e1) -- (e2);
  \draw[flow] (e2) -- (e3);
  \node[font=\scriptsize,simBlue, below=1mm of e1] {$\vE$ curl-free};
  \node[font=\scriptsize,simBlue, below=1mm of e3] {scalar potential $V\in\Hone$};
\end{scope}
% --- bottom row: magnetostatic / vector potential ---
\begin{scope}[shift={(0,-2.4)}]
  \node[stage, minimum width=30mm] (b1) {$\Div\vB=0$};
  \node[font=\scriptsize\itshape,simGray, right=8mm of b1] (b2) {exactness at $\Hdiv$};
  \node[stage, minimum width=30mm, right=8mm of b2] (b3) {$\vB=\Curl\vA$};
  \draw[flow] (b1) -- (b2);
  \draw[flow] (b2) -- (b3);
  \node[font=\scriptsize,simGreen, below=1mm of b1] {$\vB$ divergence-free};
  \node[font=\scriptsize,simGreen, below=1mm of b3] {vector potential $\vA\in\Hcurl$};
\end{scope}
\end{tikzpicture}
\caption[Potentials are exactness]{The electromagnetic potentials of
\cref{ch:maxwell} read as de Rham facts. \emph{Top:} the electrostatic Maxwell
equation $\Curl\vE=0$ says $\vE$ is curl-free; exactness at the $\Hcurl$ rung
(\cref{thm:exact}) supplies a scalar potential $V$ with $\vE=-\Grad V$, living one
rung up the complex in $\Hone$. \emph{Bottom:} the law $\Div\vB=0$ says $\vB$ is
divergence-free; exactness at the $\Hdiv$ rung supplies a vector potential $\vA$ with
$\vB=\Curl\vA$, one rung up in $\Hcurl$. The potential always lives in the space
\emph{one rung above} the field it generates---a structural fact, not a coincidence.}
\label{fig:potentials}
\end{figure}

\begin{keyidea}
The potentials are not tricks; they are exactness. ``$\vE$ is curl-free, therefore
$\vE=-\Grad V$'' \emph{is} the statement $\ker(\Curl)=\operatorname{range}(\Grad)$,
read at a particular field. ``$\vB$ is divergence-free, therefore $\vB=\Curl\vA$''
\emph{is} $\ker(\Div)=\operatorname{range}(\Curl)$. The potential of a field always
lives one rung \emph{up} the de Rham complex: $V\in\Hone$ generates $\vE\in\Hcurl$;
$\vA\in\Hcurl$ generates $\vB\in\Hdiv$. When the magnetostatic analysis of
\cref{ch:targets} solves the curl--curl equation $\Curl(\nu\Curl\vA)=\vJ$ for $\vA$
in $\Hcurl$, it is solving for a potential whose \emph{existence} the de Rham complex
guarantees and whose \emph{home space} the complex names.
\end{keyidea}

There is a consequence we will need repeatedly. Because $\vA=\Grad\psi$ added to any
$\vA$ leaves $\Curl\vA$ unchanged (since $\Curl\Grad\psi=0$), the vector potential is
not unique: $\vA$ and $\vA+\Grad\psi$ describe the same $\vB$ for any scalar $\psi$.
This freedom is the \emph{gauge}, and the gradient fields $\Grad\psi$ are precisely
the directions in which it acts. The same gradient fields are the null space of the
curl--curl operator $\Curl(\nu\Curl\,\cdot)$---for $\Curl(\nu\Curl(\Grad\psi))=0$---
which is the reason that operator is singular and the reason the eigenmode solver
must project them out. We return to this in \cref{sec:divfree}.

\section{The discrete complex}
\label{sec:discrete-arrows}

Everything so far has been about the continuous spaces. A simulator computes in the
\emph{discrete} subspaces $\Hone_h,\Hcurl_h,\Hdiv_h,\Ltwo_h$ of
\cref{ch:functionspaces}---the finite-dimensional spaces of nodal, edge, face, and
cell functions. The question that makes or breaks the method is whether the discrete
spaces form a complex of their own, with discrete arrows, and whether that complex is
\emph{exact}. The answer---remarkably---is yes, and the rest of the chapter is about
why and what it buys.

\subsection{The discrete arrows are interpolation matrices}

The continuous arrows $\Grad,\Curl,\Div$ are differential operators. Their discrete
counterparts are \emph{matrices}: rectangular arrays that take the coefficient vector
of a field in one discrete space to the coefficient vector of its image in the next.
Palace builds them with a single constructor, the \emph{discrete interpolator}, whose
job is to realize one de Rham arrow at the discrete level. From the source
(\code{fespace.cpp}'s \code{BuildDiscreteInterpolator}), the constructor dispatches on
the \emph{pair} of spaces it is handed and selects the matching arrow:
\begin{center}
\begin{tabular}{@{}lll@{}}
\toprule
discrete arrow & from $\to$ to & kernel selected \\
\midrule
discrete gradient $\mathsf{G}$ & $\Hone_h\to\Hcurl_h$ & \code{GradientInterpolator} \\
discrete curl $\mathsf{C}$     & $\Hcurl_h\to\Hdiv_h$ & \code{CurlInterpolator} \\
discrete curl (2-D) & $\Hcurl_h\to\Ltwo_h$ & \code{CurlInterpolator} (native) \\
discrete divergence $\mathsf{D}$ & $\Hdiv_h\to\Ltwo_h$ & \code{DivergenceInterpolator} \\
\bottomrule
\end{tabular}
\end{center}
Each is an honest rectangular matrix: $\mathsf{G}$ has as many rows as there are
edges and as many columns as there are vertices; $\mathsf{C}$ as many rows as faces
and columns as edges; $\mathsf{D}$ as many rows as cells and columns as faces. (The
constructor refuses any pair that is not de Rham-adjacent in the correct
direction---a built-in check that the arrow even exists.) We write $\mathsf{G}$,
$\mathsf{C}$, $\mathsf{D}$ for the three matrices, reserving the continuous
$\Grad,\Curl,\Div$ for the operators they discretize.

The construction of $\mathsf{G}$ is so concrete it is worth stating outright, because
we will compute with it in \cref{wex:discrete}. The discrete gradient is, up to sign,
the \emph{incidence matrix} of the mesh: the matrix that records, for each edge, which
two vertices it joins.

\begin{definition}[The lowest-order discrete gradient]\label{def:discrete-grad}
Let the mesh have vertices indexed $1,\dots,n_v$ and edges $1,\dots,n_e$, with edge
$e$ oriented from vertex $\operatorname{tail}(e)$ to vertex $\operatorname{head}(e)$.
The lowest-order discrete gradient is the $n_e\times n_v$ matrix $\mathsf{G}$ with
\[
  \mathsf{G}_{e,v} =
  \begin{cases}
    +1 & v = \operatorname{head}(e),\\
    -1 & v = \operatorname{tail}(e),\\
    0  & \text{otherwise.}
  \end{cases}
\]
That is, $(\mathsf{G}\,\vu)_e = u_{\operatorname{head}(e)} - u_{\operatorname{tail}(e)}$:
the discrete gradient of a nodal field reads off, on each edge, the difference of the
nodal values at its two endpoints.
\end{definition}

This is exactly right as a discrete gradient. A nodal field $u$ has one number per
vertex; its continuous gradient, integrated along an edge, is
$\int_e\Grad u\cdot d\vec\ell = u(\operatorname{head}) - u(\operatorname{tail})$ by
the fundamental theorem of calculus---and the edge degree of freedom of $\Hcurl_h$ is
exactly that edge circulation (\cref{ch:functionspaces}). So the discrete gradient
$\mathsf{G}\vu$, read as edge circulations, holds the true edge circulations of
$\Grad u$. The matrix is not an approximation of the gradient; on these degrees of
freedom it is \emph{exact}.

\subsection{The commuting diagram}

To say the discrete arrows ``are'' the continuous ones, we need a precise statement,
and it is the conceptual core of the chapter. Each space comes with a canonical
\emph{interpolant}: an operator $\Pi$ that takes a smooth field and reads off its
degrees of freedom, producing the discrete field with those DoFs. There is one per
space: $\Pi_{\Hone}$ reads nodal values, $\Pi_{\Hcurl}$ reads edge circulations,
$\Pi_{\Hdiv}$ reads face fluxes, $\Pi_{\Ltwo}$ reads cell averages. The claim is that
\emph{interpolating then differentiating discretely gives the same result as
differentiating then interpolating}:
\[
  \mathsf{G}\,\Pi_{\Hone} = \Pi_{\Hcurl}\,\Grad, \qquad
  \mathsf{C}\,\Pi_{\Hcurl} = \Pi_{\Hdiv}\,\Curl, \qquad
  \mathsf{D}\,\Pi_{\Hdiv} = \Pi_{\Ltwo}\,\Div.
\]
This is the \emph{commuting-diagram property}, drawn in \cref{fig:commute}. It says
the square formed by a continuous arrow on top, a discrete arrow on the bottom, and
interpolation arrows down the sides \emph{commutes}: both ways around the square give
the same discrete field.

\begin{figure}[t]
\centering
\begin{tikzpicture}[scale=1.0, node distance=20mm and 34mm]
  % continuous row
  \node[space, fill=simBgGray, draw=simGray, minimum width=16mm] (Ha) {$\Hone$};
  \node[space, fill=simBgGray, draw=simGray, minimum width=16mm, right=of Ha] (Hb) {$\Hcurl$};
  % discrete row
  \node[space, minimum width=16mm, below=of Ha] (ha) {$\Hone_h$};
  \node[space, minimum width=16mm, below=of Hb] (hb) {$\Hcurl_h$};
  % top arrow (continuous)
  \draw[flow] (Ha) -- node[above,font=\small]{$\Grad$} (Hb);
  % bottom arrow (discrete)
  \draw[depedge] (ha) -- node[below,font=\small]{$\mathsf{G}$} (hb);
  % side arrows (interpolation, downward)
  \draw[flow, simViolet] (Ha) -- node[left,font=\scriptsize]{$\Pi_{\Hone}$} (ha);
  \draw[flow, simViolet] (Hb) -- node[right,font=\scriptsize]{$\Pi_{\Hcurl}$} (hb);
  % commuting marker
  \node[font=\Large, simGreen] at ($(Ha)!0.5!(hb)$) {$\circlearrowleft$};
  % caption-of-square
  \node[font=\scriptsize\itshape, simGray, right=2mm of hb, text width=34mm, align=left]
    {$\mathsf{G}\,\Pi_{\Hone}=\Pi_{\Hcurl}\,\Grad$\\[2pt]both paths agree};
\end{tikzpicture}
\caption[The commuting diagram]{The commuting-diagram property, shown for the
gradient square (the curl and divergence squares are identical in form). Start at a
smooth scalar field in $\Hone$ (top left). Going \emph{right then down}: differentiate
continuously with $\Grad$, then interpolate the result into the edge space with
$\Pi_{\Hcurl}$. Going \emph{down then right}: interpolate the scalar into nodal DoFs
with $\Pi_{\Hone}$, then apply the discrete gradient matrix $\mathsf{G}$. The two
paths land on the \emph{same} discrete edge field---the square commits
($\circlearrowleft$). This single property is what makes $\Hone_h,\Hcurl_h,\Hdiv_h,
\Ltwo_h$ a faithful discrete copy of the continuous complex.}
\label{fig:commute}
\end{figure}

Why does the square commute? For the gradient case it is the fundamental theorem of
calculus again. Interpolate a smooth $u$ into nodal DoFs (read its vertex values),
then apply $\mathsf{G}$: by \cref{def:discrete-grad} the result on edge $e$ is
$u(\operatorname{head}(e))-u(\operatorname{tail}(e))$. Going the other way, take the
continuous gradient $\Grad u$, then interpolate into edge DoFs (read its
circulations): the circulation on edge $e$ is
$\int_e\Grad u\cdot d\vec\ell=u(\operatorname{head}(e))-u(\operatorname{tail}(e))$.
Identical. The curl and divergence squares commute for the same reason, with Stokes'
theorem and the divergence theorem playing the role the fundamental theorem of
calculus plays here---each relates the differential operator's DoF (face flux, cell
average) to a boundary integral of the field's lower DoFs.

\subsection{Discrete exactness}
\label{sec:discrete-exact}

The commuting diagram has an immediate and decisive corollary: the discrete arrows
compose to zero, exactly as the continuous ones do.

\begin{theorem}[Discrete exactness]\label{thm:discrete-exact}
The discrete arrows satisfy the complex identities
\[
  \mathsf{C}\,\mathsf{G} = 0, \qquad \mathsf{D}\,\mathsf{C} = 0,
\]
and on a discretization of a simply-connected domain the discrete complex
$\Hone_h\xrightarrow{\mathsf{G}}\Hcurl_h\xrightarrow{\mathsf{C}}\Hdiv_h
\xrightarrow{\mathsf{D}}\Ltwo_h$ is exact:
$\operatorname{range}(\mathsf{G})=\ker(\mathsf{C})$ and
$\operatorname{range}(\mathsf{C})=\ker(\mathsf{D})$.
\end{theorem}

\begin{proof}[Why $\mathsf{C}\mathsf{G}=0$]
We give the argument for the first identity; the second is identical with $\Div\Curl$
in place of $\Curl\Grad$. Take any discrete nodal field $u_h\in\Hone_h$; we must show
$\mathsf{C}\mathsf{G}u_h=0$. The lowest-order discrete spaces have the property that
each discrete field is the interpolant of itself, so we may write $u_h=\Pi_{\Hone}u$
for the smooth representative $u$ that $u_h$ samples. Now chase the commuting diagram
twice:
\[
  \mathsf{C}\,\mathsf{G}\,u_h
  = \mathsf{C}\,\mathsf{G}\,\Pi_{\Hone}u
  = \mathsf{C}\,\Pi_{\Hcurl}\,\Grad u
  = \Pi_{\Hdiv}\,\Curl\,\Grad u
  = \Pi_{\Hdiv}\,\mathbf{0}
  = 0,
\]
using the gradient commuting square at the second step, the curl commuting square at
the third, and the \emph{continuous} identity $\Curl\Grad=0$ (\cref{lem:complex}) at
the fourth. The discrete identity is the continuous identity, pulled through the
commuting diagram.
\end{proof}

This is the heart of the matter. \emph{Discrete exactness is not a separate
miracle---it is the continuous exactness, inherited through the commuting diagram.}
The finite-element spaces of \cref{ch:functionspaces} were not chosen for accuracy
alone; they were chosen so that their canonical interpolants commute with
differentiation, which is exactly the condition that makes the discrete complex a
faithful copy of the continuous one. \Cref{fig:incidence} shows the discrete identity
$\mathsf{C}\mathsf{G}=0$ as a statement about the mesh's incidence structure.

\begin{figure}[t]
\centering
\begin{tikzpicture}[scale=1.0, every node/.style={font=\scriptsize}]
  % a triangle with a vertex, two edges, one face
  \coordinate (A) at (0,0);
  \coordinate (B) at (2.4,0);
  \coordinate (C) at (1.2,2.0);
  \fill[simBgGreen] (A)--(B)--(C)--cycle;
  \draw[simInk,thick] (A)--(B)--(C)--cycle;
  % vertices
  \fill[simBlue] (A) circle (2.4pt); \node[simBlue,below left] at (A) {$v_1$};
  \fill[simBlue] (B) circle (2.4pt); \node[simBlue,below right] at (B) {$v_2$};
  \fill[simBlue] (C) circle (2.4pt); \node[simBlue,above] at (C) {$v_3$};
  % oriented edges with circulation arrows
  \draw[-{Stealth[length=2mm]},simRust,thick] ($(A)!0.3!(B)$)--($(A)!0.7!(B)$);
  \node[simRust,below] at ($(A)!0.5!(B)$) {$e_1$};
  \draw[-{Stealth[length=2mm]},simRust,thick] ($(B)!0.3!(C)$)--($(B)!0.7!(C)$);
  \node[simRust,right] at ($(B)!0.5!(C)$) {$e_2$};
  \draw[-{Stealth[length=2mm]},simRust,thick] ($(C)!0.3!(A)$)--($(C)!0.7!(A)$);
  \node[simRust,left] at ($(C)!0.5!(A)$) {$e_3$};
  % face circulation marker
  \draw[simGreen,thick,-{Stealth[length=1.8mm]}] (1.2,0.75) arc (-90:170:0.32);
  \node[simGreen] at (1.2,0.62) {$f$};
  % the algebra to the right
  \node[align=left, font=\scriptsize, text width=58mm, anchor=west] at (3.4,1.0) {%
    \textbf{Discrete gradient} $\mathsf{G}$ (edge $=$ head $-$ tail):\\[2pt]
    $\mathsf{G}\vu=(u_2{-}u_1,\;u_3{-}u_2,\;u_1{-}u_3)$\\[4pt]
    \textbf{Discrete curl} $\mathsf{C}$ (face $=$ signed edge sum):\\[2pt]
    $\mathsf{C}(\mathsf{G}\vu)=(u_2{-}u_1)+(u_3{-}u_2)+(u_1{-}u_3)=0$\\[4pt]
    The loop of edge differences \emph{telescopes}: $\mathsf{C}\mathsf{G}=0$.};
\end{tikzpicture}
\caption[Discrete $\mathsf{C}\mathsf{G}=0$ telescopes]{The discrete identity
$\mathsf{C}\mathsf{G}=0$ on a single triangle. The discrete gradient places, on each
oriented edge, the difference of its endpoint nodal values (head minus tail). The
discrete curl sums the edge circulations around the face boundary, with sign. Going
around the loop $e_1,e_2,e_3$, the differences \emph{telescope}: every vertex value
appears once with $+$ and once with $-$, so the sum is identically zero. This is the
discrete avatar of ``the line integral of a gradient around a closed loop is
zero''---and the discrete avatar of $\Curl\Grad=0$.}
\label{fig:incidence}
\end{figure}

\begin{keyidea}
\textbf{Discrete exactness is why edge elements avoid spurious modes.} The
spurious-mode pathology of \cref{ch:functionspaces} had a structural cause: a
componentwise-nodal discretization of $\vE$ has the \emph{wrong} kernel for its
discrete curl, so the discrete curl--curl operator's null space is polluted with
non-physical fields that masquerade as resonances. Edge elements ($\Hcurl_h$) fix
this because they sit in an \emph{exact} discrete complex: by
\cref{thm:discrete-exact} the kernel of the discrete curl is \emph{exactly} the range
of the discrete gradient---the discrete gauge fields $\mathsf{G}\,\psi_h$, and nothing
else. The curl--curl operator's null space is then precisely the discrete gradients,
a space we can identify and handle (\cref{sec:divfree}), with no spurious extras
hiding in it. Get the discrete complex exact and the spectrum is clean; break
exactness and spurious modes return. This is the single deepest reason the four
finite-element spaces are the ones a Maxwell solver must use.
\end{keyidea}

\section{Whitney forms: the canonical basis}
\label{sec:whitney}

The lowest-order discrete complex has a canonical basis---one set of shape functions,
one per mesh entity, that realizes all four spaces at once and makes the discrete
arrows the incidence matrices of \cref{def:discrete-grad}. These are the
\emph{Whitney forms}. We met one already: the edge function
$\boldsymbol\phi_{ij}=\lambda_i\Grad\lambda_j-\lambda_j\Grad\lambda_i$ of
\cref{wex:whitney}. The full family, built from the barycentric coordinates
$\lambda_0,\dots,\lambda_d$ of a simplex, is:

\begin{itemize}[nosep]
\item \textbf{Nodal} ($\Hone_h$, on vertices): $\lambda_i$ itself---the barycentric
  hat function, $1$ at vertex $i$, $0$ at the others.
\item \textbf{Edge} ($\Hcurl_h$, on edges): $\boldsymbol\phi_{ij}=
  \lambda_i\Grad\lambda_j-\lambda_j\Grad\lambda_i$ on edge $(i,j)$---unit circulation
  on its own edge, zero on the others.
\item \textbf{Face} ($\Hdiv_h$, on faces): $\boldsymbol\phi_{ijk}=
  2\bigl(\lambda_i(\Grad\lambda_j\times\Grad\lambda_k)
  +\lambda_j(\Grad\lambda_k\times\Grad\lambda_i)
  +\lambda_k(\Grad\lambda_i\times\Grad\lambda_j)\bigr)$ on face $(i,j,k)$---unit flux
  through its own face.
\item \textbf{Cell} ($\Ltwo_h$, on cells): the constant (or $\det$-scaled constant)
  on each cell---unit cell average.
\end{itemize}

\begin{figure}[t]
\centering
\begin{tikzpicture}[scale=1.05, every node/.style={font=\scriptsize}]
% three triangles showing nodal / edge / face Whitney forms
% --- nodal ---
\begin{scope}
  \node[font=\footnotesize\bfseries] at (1.0,2.1) {nodal $\lambda_i$};
  \coordinate (a) at (0,0); \coordinate (b) at (2.0,0); \coordinate (c) at (1.0,1.7);
  \draw[simInk,thick] (a)--(b)--(c)--cycle;
  \fill[simBlue] (a) circle (3pt);
  \fill[simBlue!35] (b) circle (1.6pt); \fill[simBlue!35] (c) circle (1.6pt);
  \node[simGray] at (1.0,-0.45) {one per vertex};
\end{scope}
% --- edge ---
\begin{scope}[shift={(3.6,0)}]
  \node[font=\footnotesize\bfseries] at (1.0,2.1) {edge $\boldsymbol\phi_{ij}$};
  \coordinate (a) at (0,0); \coordinate (b) at (2.0,0); \coordinate (c) at (1.0,1.7);
  \draw[simInk,thick] (a)--(b)--(c)--cycle;
  % tangential circulation along the bottom edge
  \foreach \t in {0.2,0.4,0.6,0.8}{
    \draw[-{Stealth[length=1.5mm]},simRust,thick]
      ($(a)!\t!(b)+(0,0.12)$)--+(0.28,0);}
  \node[simGray] at (1.0,-0.45) {one per edge};
\end{scope}
% --- face ---
\begin{scope}[shift={(7.2,0)}]
  \node[font=\footnotesize\bfseries] at (1.0,2.1) {face $\boldsymbol\phi_{ijk}$};
  \coordinate (a) at (0,0); \coordinate (b) at (2.0,0); \coordinate (c) at (1.0,1.7);
  \fill[simGreen!16] (a)--(b)--(c)--cycle;
  \draw[simInk,thick] (a)--(b)--(c)--cycle;
  % flux out of the face (dot = out of page)
  \node[circle,draw=simGreen,fill=white,inner sep=1.1pt] at (1.0,0.55) {};
  \fill[simGreen] (1.0,0.55) circle (1.0pt);
  \node[simGreen] at (1.45,0.55) {$\odot$};
  \node[simGray] at (1.0,-0.45) {one per face};
\end{scope}
\end{tikzpicture}
\caption[Whitney forms on a triangle]{The lowest-order Whitney basis, the canonical
realization of the discrete de Rham complex. \emph{Left:} the nodal form $\lambda_i$,
one per vertex, peaking at its vertex (the $\Hone_h$ basis). \emph{Middle:} the edge
form $\boldsymbol\phi_{ij}$, one per edge, carrying unit tangential circulation along
its own edge (the $\Hcurl_h$ basis of \cref{wex:whitney}). \emph{Right:} the face form
$\boldsymbol\phi_{ijk}$, one per face, carrying unit normal flux through its own face
(the $\Hdiv_h$ basis, drawn here as flux out of the page, $\odot$). Together they make
the discrete arrows the incidence matrices of \cref{def:discrete-grad}: the gradient
of a nodal hat is a sum of edge forms, and so on up the complex.}
\label{fig:whitney}
\end{figure}

The decisive algebraic fact---the one that ties the basis to the incidence matrices---
is how the differential operators act on the Whitney forms. Applying the gradient to a
nodal hat $\lambda_i$ produces a field that, expanded in the edge forms
$\boldsymbol\phi_{ij}$, has coefficients $\pm 1$ on the edges touching vertex $i$ and
$0$ elsewhere---exactly the column of the incidence matrix $\mathsf{G}$ for vertex
$i$. The Whitney forms are the basis in which the discrete gradient, curl, and
divergence \emph{are} the signed incidence matrices of the mesh. This is why
\cref{def:discrete-grad} could state $\mathsf{G}$ combinatorially: in the Whitney
basis, the analysis is purely topological.

\begin{remark}[For the curious: cohomology and holes]
\label{rmk:cohomology}
We claimed exactness for \emph{simply-connected} domains and warned of a gap
otherwise. That gap is not a defect---it is information, and it is the gateway to a
beautiful corner of mathematics. On a domain with holes, the quotient
$\ker(\Curl)/\operatorname{range}(\Grad)$---the curl-free fields that are \emph{not}
gradients, modulo those that are---is a finite-dimensional vector space whose
dimension \emph{counts the holes}. Its elements are the \emph{harmonic} fields: those
in the kernel of both the differential operator above and the (adjoint) one below.
For a solid torus this space is one-dimensional; for a domain with $g$ independent
tunnels it is $g$-dimensional. This is \emph{de Rham cohomology}, and its punchline is
startling: a purely analytic object (which curl-free fields fail to be gradients) is
counting a purely topological one (how many holes the domain has). The discrete
complex inherits this too---a correctly-constructed Whitney complex on a meshed torus
has exactly one discrete harmonic field, and a solver can find it. Palace's domains
are usually simply-connected, so the harmonic space is trivial and the complex is
exact; but the machinery is there, and the fact that finite elements can count the
holes of a domain is one of the quiet triumphs of the theory. The canonical reference
is Arnold, Falk, and Winther's finite-element exterior calculus~\citep{arnold2006feec};
we say no more here.
\end{remark}

\section{Worked examples}

\subsection{The continuous identities, component by component}

We first watch \cref{lem:complex} happen on concrete polynomial fields, so the
mixed-partial cancellation is not an abstraction but an arithmetic fact you can check.

\begin{workedexample}{Verifying $\Curl\Grad\phi=0$ and $\Div\Curl\vF=0$}{component}
\textbf{Part 1: $\Curl\Grad\phi=0$.} Take the scalar field
$\phi(x,y,z)=x^2 y + y z^2$. Its gradient is
\[
  \Grad\phi = \bigl(\,2xy,\;\; x^2 + z^2,\;\; 2yz\,\bigr).
\]
Now take the curl of this vector field. The $x$-component is
\[
  (\Curl\Grad\phi)_x
  = \partial_y(2yz) - \partial_z(x^2+z^2)
  = 2z - 2z = 0.
\]
The $y$-component is
\[
  (\Curl\Grad\phi)_y
  = \partial_z(2xy) - \partial_x(2yz)
  = 0 - 0 = 0.
\]
The $z$-component is
\[
  (\Curl\Grad\phi)_z
  = \partial_x(x^2+z^2) - \partial_y(2xy)
  = 2x - 2x = 0.
\]
So $\Curl\Grad\phi=\mathbf 0$. Notice the mechanism in the $x$-component: the term
$2z$ came from $\partial_y\partial_z\phi$ and the term $-2z$ from
$-\partial_z\partial_y\phi$; they cancel because mixed partials commute. The other
components cancel the same way.

\textbf{Part 2: $\Div\Curl\vF=0$.} Take the vector field
$\vF(x,y,z)=\bigl(\,xy z,\;\; y^2 x,\;\; z x^2\,\bigr)$. Its curl is
\[
  \Curl\vF =
  \begin{pmatrix}
    \partial_y(z x^2) - \partial_z(y^2 x)\\[2pt]
    \partial_z(x y z) - \partial_x(z x^2)\\[2pt]
    \partial_x(y^2 x) - \partial_y(x y z)
  \end{pmatrix}
  =
  \begin{pmatrix}
    0 - 0\\
    x y - 2 z x\\
    y^2 - x z
  \end{pmatrix}
  =
  \begin{pmatrix}
    0\\ xy - 2zx\\ y^2 - xz
  \end{pmatrix}.
\]
Now its divergence:
\[
  \Div\Curl\vF
  = \partial_x(0) + \partial_y(xy - 2zx) + \partial_z(y^2 - xz)
  = 0 + x + (-x)
  = 0.
\]
Again the surviving terms cancel in a mixed-partial pair: $\partial_y(xy)=x$ from the
second component against $\partial_z(-xz)=-x$ from the third. The identity holds not
by luck but by the structure of the curl, which arranges every term to meet its
negative.
\end{workedexample}

\subsection{Discrete $\mathsf{C}\mathsf{G}=0$ on a small mesh}

Now the discrete version. We build the lowest-order discrete gradient on a tiny
two-triangle mesh, apply it to a concrete nodal field, and verify by hand that the
discrete curl annihilates the result---a single instance of
\cref{thm:discrete-exact}.

\begin{workedexample}{The discrete gradient on a square of two triangles}{discrete}
\textbf{The mesh.} Take a unit square split into two triangles by its diagonal, with
vertices
\[
  v_1=(0,0),\quad v_2=(1,0),\quad v_3=(1,1),\quad v_4=(0,1),
\]
and five edges---the four sides and the diagonal---oriented as
\[
  e_1: v_1\!\to\!v_2,\;\;
  e_2: v_2\!\to\!v_3,\;\;
  e_3: v_4\!\to\!v_3,\;\;
  e_4: v_1\!\to\!v_4,\;\;
  e_5: v_1\!\to\!v_3 \ (\text{diagonal}).
\]
The two triangular faces are
$f_1=(v_1,v_2,v_3)$ and $f_2=(v_1,v_3,v_4)$.

\textbf{The discrete gradient $\mathsf{G}$.} By \cref{def:discrete-grad}, each row of
the $5\times 4$ matrix $\mathsf{G}$ is $+1$ at the edge's head, $-1$ at its tail:
\[
  \mathsf{G} =
  \begin{array}{c@{\,}c}
  & \begin{matrix} v_1 & v_2 & v_3 & v_4 \end{matrix}\\
  \begin{matrix} e_1\\ e_2\\ e_3\\ e_4\\ e_5 \end{matrix} &
  \begin{pmatrix*}[r]
    -1 & 1 & 0 & 0\\
     0 & -1 & 1 & 0\\
     0 & 0 & 1 & -1\\
    -1 & 0 & 0 & 1\\
    -1 & 0 & 1 & 0
  \end{pmatrix*}
  \end{array}.
\]
\textbf{A nodal field and its discrete gradient.} Take the nodal field
$\vu=(u_1,u_2,u_3,u_4)=(0,1,3,2)$ (one value per vertex). Then
\[
  \mathsf{G}\vu =
  \begin{pmatrix*}[r]
    u_2-u_1\\ u_3-u_2\\ u_3-u_4\\ u_4-u_1\\ u_3-u_1
  \end{pmatrix*}
  =
  \begin{pmatrix*}[r]
    1-0\\ 3-1\\ 3-2\\ 2-0\\ 3-0
  \end{pmatrix*}
  =
  \begin{pmatrix*}[r] 1\\ 2\\ 1\\ 2\\ 3 \end{pmatrix*}.
\]
These five numbers are the edge circulations of the discrete gradient field
$\Grad u_h$: one per edge.

\textbf{The discrete curl, face by face, and $\mathsf{C}\mathsf{G}=0$.} The discrete
curl on a face is the signed sum of the circulations of its boundary edges, traversed
consistently with the face orientation. Walk the boundary of $f_1=(v_1,v_2,v_3)$:
$v_1\!\to\!v_2$ along $+e_1$, $v_2\!\to\!v_3$ along $+e_2$, and $v_3\!\to\!v_1$ along
$-e_5$ (the diagonal, traversed against its orientation). So
\[
  (\mathsf{C}\mathsf{G}\vu)_{f_1}
  = (+1)(\mathsf{G}\vu)_{e_1} + (+1)(\mathsf{G}\vu)_{e_2} + (-1)(\mathsf{G}\vu)_{e_5}
  = 1 + 2 - 3 = 0.
\]
Now the other face $f_2=(v_1,v_3,v_4)$: $v_1\!\to\!v_3$ along $+e_5$, $v_3\!\to\!v_4$
along $-e_3$, and $v_4\!\to\!v_1$ along $-e_4$. So
\[
  (\mathsf{C}\mathsf{G}\vu)_{f_2}
  = (+1)(\mathsf{G}\vu)_{e_5} + (-1)(\mathsf{G}\vu)_{e_3} + (-1)(\mathsf{G}\vu)_{e_4}
  = 3 - 1 - 2 = 0.
\]
Both faces give zero: $\mathsf{C}\mathsf{G}\vu=\mathbf 0$. And the reason is exactly
the telescoping of \cref{fig:incidence}: on $f_1$ the circulations are
$(u_2-u_1)+(u_3-u_2)-(u_3-u_1)$, in which every nodal value cancels its own appearance.
The result holds for \emph{every} nodal field $\vu$, not just our chosen one---that
is the content of $\mathsf{C}\mathsf{G}=0$. The discrete gradient of any nodal function
has zero discrete curl, edge differences around any face loop always telescoping to
nothing.
\end{workedexample}

\section{Consequences for the solver}
\label{sec:divfree}

The de Rham structure is not theory for its own sake; the solver chapters lean on it
constantly. We close by collecting the three consequences that recur, each a direct
reading of discrete exactness.

\paragraph{The gradient null space of curl--curl.} Every magnetic and wave analysis
assembles the curl--curl operator $\mathsf{K}\sim\Curl(\nu\Curl\,\cdot)$ on $\Hcurl_h$.
By discrete exactness, $\mathsf{C}\mathsf{G}=0$, so this operator annihilates every
discrete gradient: $\mathsf{K}\,\mathsf{G}\psi_h=0$ for any nodal $\psi_h$. The
operator is therefore \emph{singular}, with null space exactly the discrete gradients
$\operatorname{range}(\mathsf{G})$---the discrete gauge fields of \cref{sec:potentials}.
This is not a bug to be regularized away; it is the discrete image of gauge freedom,
and discrete exactness tells us precisely what it is.

\paragraph{Projecting onto the divergence-free subspace.} Because the curl--curl
operator has this gradient null space, an eigenmode solve will, left alone, return
spurious ``zero-frequency'' modes living in $\operatorname{range}(\mathsf{G})$,
interleaved with the physical resonances. The fix is to confine the solver to the
\emph{complement}: the divergence-free subspace, those Nedelec fields $\vy$ with
$\mathsf{G}^\top\mathsf{M}\vy=0$ (no discrete-gradient component in the $\mathsf{M}$-inner
product). Palace's divergence-free projector
\[
  P = I - \mathsf{G}\,(\mathsf{G}^\top\mathsf{M}\mathsf{G})^{-1}\mathsf{G}^\top\mathsf{M}
\]
removes the gradient part of any field, mapping it into the divergence-free subspace;
the eigensolver applies $P$ to each candidate vector to keep it physical. Every
ingredient of $P$ is a de Rham object: $\mathsf{G}$ is the discrete gradient built by
the interpolator of \cref{sec:discrete-arrows}, and the subspace it projects onto is
the kernel of the discrete divergence the complex defines. We develop the projector and
its place in the eigenmode pipeline in \cref{ch:eigenmodes}.

\paragraph{The Helmholtz decomposition.} Underlying both is the discrete
\emph{Helmholtz decomposition}: every Nedelec field splits uniquely as a
divergence-free part plus a gradient part,
\[
  \vy = \underbrace{\vy_{\mathrm{df}}}_{\text{divergence-free}}
        \;+\; \underbrace{\mathsf{G}\psi}_{\text{gradient}},
\]
and the projector $P$ is exactly the operator that extracts the first summand. This is
the discrete shadow of the continuous decomposition (\cref{ch:bc}), and it is a direct
consequence of discrete exactness: the two summands live in $\ker(\mathsf{C})$'s
complement and in $\operatorname{range}(\mathsf{G})=\ker(\mathsf{C})$ respectively, the
equality being \cref{thm:discrete-exact}. The same decomposition is what makes the
magnetostatic and electrostatic analyses well-posed, and it is the structural reason the
boundary-condition chapter can speak of ``the divergence-free field'' as a definite
object.

\begin{keyidea}
Three solver facts, one structural source. The curl--curl operator is singular
\emph{because} $\mathsf{C}\mathsf{G}=0$; its null space is the discrete gradients
\emph{because} the discrete complex is exact; and the eigensolver projects onto the
divergence-free subspace \emph{because} that is the exactness-supplied complement of
the gradient null space. Discrete exactness is not a property one checks and files
away---it is load-bearing in every magnetic, wave, and eigenmode solve in the book.
\end{keyidea}

\summarysection
The four finite-element spaces of \cref{ch:functionspaces} assemble into the
\emph{de Rham complex}
$\Hone\xrightarrow{\Grad}\Hcurl\xrightarrow{\Curl}\Hdiv\xrightarrow{\Div}\Ltwo$:
each differential operator maps one space exactly into the next. The complex
\emph{property} is the pair of identities $\Curl\Grad=0$ and $\Div\Curl=0$, both
corollaries of the commuting of mixed partial derivatives; \emph{exactness} is the
stronger statement that the range of each arrow equals the kernel of the next, true on
simply-connected domains. Exactness is the mathematical content of the electromagnetic
potentials: $\Curl\vE=0$ forces $\vE=-\Grad V$, and $\Div\vB=0$ forces $\vB=\Curl\vA$,
the potential always living one rung up the complex. The discrete finite-element spaces
inherit the same structure: their canonical interpolants \emph{commute with
differentiation} (the commuting-diagram property), and through that diagram the discrete
arrows---the incidence-style matrices $\mathsf{G},\mathsf{C},\mathsf{D}$ built by
Palace's discrete interpolator---satisfy $\mathsf{C}\mathsf{G}=0$ and
$\mathsf{D}\mathsf{C}=0$ and form an \emph{exact discrete complex}. The Whitney forms are
the canonical lowest-order basis realizing it, in which the discrete arrows are the
signed incidence matrices of the mesh. Discrete exactness is the deep reason edge
elements avoid spurious modes: it makes the curl--curl operator's null space exactly the
discrete gradients, no spurious extras, so the eigensolver can project onto the
divergence-free complement and recover a clean spectrum. Every divergence-free
projection, every gauge null space, and the Helmholtz decomposition of \cref{ch:bc} is a
reading of this one structure.

\furtherreading
\label{sec:furtherreading}
The unifying modern framework is \emph{finite-element exterior calculus}, which casts
the four spaces as discrete differential forms and the complex as the central object;
Arnold, Falk, and Winther~\citep{arnold2006feec} is the canonical and remarkably readable
account, and is the natural sequel to this chapter for the cohomology remark
(\cref{rmk:cohomology}). For the electromagnetic specialization---$\Hcurl$, $\Hdiv$, the
N\'ed\'elec and Raviart--Thomas elements, the commuting interpolants, and the spurious-mode
analysis worked through with full rigour---Monk~\citep{monk2003} is the standard reference.
The discrete de Rham complex and the role of Whitney forms in computational
electromagnetics were brought to prominence by Hiptmair, whose work on the canonical
basis and the commuting-diagram property~\citep{hiptmair1998} is the bridge between the
abstract complex and the matrices a solver actually builds. Readers wanting the continuous
exactness theorem (\cref{thm:exact}, the Poincar\'e lemma) and the topology of the harmonic
spaces will find both in any of the three.

\exercisessection
\begin{enumerate}[nosep]
\item \textbf{Type-checking the arrows.} For the 2-D complex
  $\Hone\xrightarrow{\Grad}\Hcurl\xrightarrow{\operatorname{curl}}\Ltwo$, write out the
  scalar curl $\operatorname{curl}(u_1,u_2)=\partial_x u_2-\partial_y u_1$ and verify
  directly that $\operatorname{curl}(\Grad\phi)=0$ for any scalar $\phi$. Which single
  mixed-partial cancellation does it rest on?
\item \textbf{Another concrete instance.} Repeat \cref{wex:component} for
  $\phi=e^{x}\sin(yz)$: compute $\Grad\phi$ and confirm $\Curl\Grad\phi=\mathbf 0$
  component by component. (You will need the product and chain rules, but the
  cancellation is the same mixed-partial symmetry.)
\item \textbf{Range versus kernel.} On the annulus $\{1<x^2+y^2<2\}$ in the plane, the
  field $\vF=(-y,x)/(x^2+y^2)$ is curl-free but is \emph{not} the gradient of any
  single-valued scalar. Compute $\operatorname{curl}\vF$ to confirm it is curl-free, then
  explain in one sentence why this does not contradict \cref{thm:exact}. (Which
  hypothesis of the theorem does the annulus violate?)
\item \textbf{Potentials are exactness.} State, in the language of ranges and kernels,
  the precise de Rham fact that licenses each of: (a) $\vE=-\Grad V$ in electrostatics;
  (b) $\vB=\Curl\vA$ in magnetostatics. For each, name the rung of the complex at which
  the exactness is read.
\item \textbf{Gauge freedom.} Show that if $\vB=\Curl\vA$ then also $\vB=\Curl\vA'$ for
  $\vA'=\vA+\Grad\psi$, any scalar $\psi$. Which de Rham identity makes this work, and how
  does it explain why the discrete curl--curl operator is singular?
\item \textbf{Discrete gradient on a hexahedral face.} Build the lowest-order discrete
  gradient $\mathsf{G}$ for a single square cell with vertices
  $v_1,v_2,v_3,v_4$ (counterclockwise) and four boundary edges oriented
  $v_1\!\to\!v_2\!\to\!v_3\!\to\!v_4\!\to\!v_1$. Take the nodal field
  $\vu=(2,5,4,1)$, compute the four edge circulations $\mathsf{G}\vu$, and verify that
  the discrete curl on the single face (the signed sum around the boundary) is zero.
\item \textbf{Telescoping is the mechanism.} Prove the general claim behind
  \cref{wex:discrete}: for \emph{any} oriented cycle of edges
  $e_1,\dots,e_k$ forming a closed loop and any nodal field $\vu$, the signed sum of
  edge differences $\sum_j \pm(u_{\mathrm{head}(e_j)}-u_{\mathrm{tail}(e_j)})$ around the
  loop is zero. (This is $\mathsf{C}\mathsf{G}=0$ for one face; the sign on $e_j$ is its
  orientation relative to the loop traversal.)
\item \textbf{The commuting square in symbols.} Using the commuting-diagram property
  $\mathsf{G}\,\Pi_{\Hone}=\Pi_{\Hcurl}\,\Grad$ and the analogous curl square, reproduce
  the diagram-chase proof that $\mathsf{C}\mathsf{G}=0$, and identify which single
  \emph{continuous} fact (\cref{lem:complex}) the discrete identity inherits.
\item \textbf{Why not nodal vectors, revisited.} In one paragraph, connect the
  spurious-mode pitfall of \cref{ch:functionspaces} to discrete exactness: explain why a
  componentwise-nodal discretization of $\vE$ fails to sit in an exact discrete complex,
  and what goes wrong with the kernel of its discrete curl as a result.
\item \textbf{Cohomology counts holes (open-ended).} Read \cref{rmk:cohomology} again. On
  a solid torus, how many independent discrete harmonic fields does the Whitney complex
  carry, and what feature of the geometry does that number measure? (No computation
  required---state the count and its topological meaning.)
\end{enumerate}

\chapter{Reference Elements, Basis Functions, and Quadrature}
\label{ch:refelement}

\chapterorient{The previous chapter named four function spaces and told you which
field belongs in each, but it left the basis functions themselves as abstractions:
``the hat function on a vertex,'' ``the edge function with unit circulation.'' This
chapter opens the engine room. We will see how a basis function is actually
\emph{represented} inside a finite-element code, how it is \emph{integrated} to fill
in a matrix entry, and why both jobs are done not on the physical mesh but on a
single canonical \emph{reference element}. Three ideas carry the chapter: shape
functions defined once on a reference cell, a geometric map that carries them to
each physical cell, and a quadrature rule that turns the integral defining each
matrix entry into a finite weighted sum. By the end you will be able to write the
hat functions of a triangle in closed form, transform a derivative through a
Jacobian, integrate a product of basis functions by Gauss quadrature and check it
against the exact answer, and read the element-local tensor whose axes
($E,L,P,C,G$) the matrix-free pipeline of \cref{ch:matrixfree} runs on. This is the
layer where the finite-element method stops being a theory and becomes arithmetic.}

\section{One element, tabulated once}

A mesh of an engineering device has thousands to millions of cells. Each cell needs
its own basis functions, its own integrals, its own contribution to the global
matrix. If we defined a fresh set of polynomials on every physical cell---fitting
them to that cell's particular vertices, recomputing their integrals from
scratch---the cost would be enormous and the bookkeeping hopeless. Worse, it would
be \emph{wasteful}, because the cells are not really different: a tetrahedron is a
tetrahedron, whatever its size or orientation. The shape of a linear basis function
on one tetrahedron is, in a precise sense we are about to make, the \emph{same} as on
any other.

The finite-element method exploits this with a single organizing decision: define
the basis functions \emph{once}, on a fixed canonical cell called the
\emph{reference element}, and reach every physical cell by a geometric map. The
reference element $\hat K$ is a cell of standard shape in its own coordinate
system---we use the \emph{reference triangle} $\hat K = \{(\hat x,\hat y): \hat
x\ge 0,\ \hat y\ge 0,\ \hat x+\hat y\le 1\}$ in two dimensions and the
\emph{reference tetrahedron} in three. On $\hat K$ we lay down a family of
\emph{shape functions} $\hat\phi_1,\dots,\hat\phi_L$ once and for all, and we
\emph{tabulate} them: we compute and store their values, and the values of their
derivatives, at a fixed set of points. That single table serves the entire mesh.

\begin{keyidea}
The reference element is the finite-element method's central act of reuse. Shape
functions are defined and tabulated \emph{once}, on a canonical reference cell
$\hat K$; each physical cell $K$ is reached by a geometric map $\vx = F_K(\hat\vx)$.
A basis function on $K$ is just a reference shape function composed with the inverse
map. One tabulation---a small table of basis values and derivatives at a few
points---drives the integration of every cell in the mesh. The mesh contributes only
its geometry (through the map $F_K$); the polynomial content is shared.
\end{keyidea}

\Cref{fig:refmap} shows the arrangement. The reference cell sits on the left in its
own $(\hat x,\hat y)$ coordinates; the map $F_K$ carries it to a physical cell $K$
somewhere in the mesh. Everything hard---the choice of polynomials, their values,
their derivatives---lives on the left, computed once. Everything cell-specific---the
size, the shape, the orientation---lives in the map $F_K$, which is cheap to apply.
This split is not merely an optimization; it is the conceptual backbone of the whole
construction, and it recurs at every level. When \cref{ch:matrixfree} factors the
operator into $\mathsf{G}$ (gather to elements), $\mathsf{B}$ (apply the tabulated
basis), $\mathsf{D}$ (apply the per-point geometry and material), and their
transposes, the tabulated basis $\mathsf{B}$ is exactly the reference shape
functions of this chapter, and the geometry factor $\mathsf{D}$ is exactly the map's
Jacobian.

\begin{figure}[t]
\centering
\begin{tikzpicture}[scale=1.0, every node/.style={font=\scriptsize}]
% --- reference triangle on the left ---
\begin{scope}
  \node[font=\footnotesize\bfseries] at (1.0,2.35) {reference $\hat K$};
  \fill[simBgBlue] (0,0)--(2,0)--(0,2)--cycle;
  \draw[simInk,thick] (0,0)--(2,0)--(0,2)--cycle;
  \fill[simBlue] (0,0) circle (2.2pt) node[below left]{$\hat\vx_0$};
  \fill[simBlue] (2,0) circle (2.2pt) node[below right]{$\hat\vx_1$};
  \fill[simBlue] (0,2) circle (2.2pt) node[above left]{$\hat\vx_2$};
  \draw[->,simGray] (-0.35,-0.35)--(0.7,-0.35) node[right]{$\hat x$};
  \draw[->,simGray] (-0.35,-0.35)--(-0.35,0.7) node[above]{$\hat y$};
  \node[simGray] at (0.7,0.6) {$\hat\phi_i$ tabulated here};
\end{scope}
% --- the map ---
\draw[flow, simRust, thick] (2.6,1.0) .. controls (3.4,1.5) and (4.0,1.5) .. (4.8,1.0)
  node[midway, above, font=\scriptsize]{$\vx = F_K(\hat\vx)$};
% --- physical cell on the right (skewed) ---
\begin{scope}[shift={(5.4,0)}]
  \node[font=\footnotesize\bfseries] at (1.3,2.35) {physical $K$};
  \fill[simBgRust] (0.2,0.1)--(2.4,0.5)--(1.0,2.0)--cycle;
  \draw[simInk,thick] (0.2,0.1)--(2.4,0.5)--(1.0,2.0)--cycle;
  \fill[simRust] (0.2,0.1) circle (2.2pt) node[below left]{$\vx_0$};
  \fill[simRust] (2.4,0.5) circle (2.2pt) node[below right]{$\vx_1$};
  \fill[simRust] (1.0,2.0) circle (2.2pt) node[above]{$\vx_2$};
  \draw[->,simGray] (-0.2,-0.35)--(0.85,-0.35) node[right]{$x$};
  \draw[->,simGray] (-0.2,-0.35)--(-0.2,0.7) node[above]{$y$};
  \node[simGray] at (1.2,0.7) {$\phi_i = \hat\phi_i\circ F_K^{-1}$};
\end{scope}
\end{tikzpicture}
\caption[The reference element and the geometric map]{The organizing picture of the
chapter. Shape functions $\hat\phi_i$ are defined and tabulated once on the reference
cell $\hat K$ (left, in its own $(\hat x,\hat y)$ coordinates). The geometric map
$F_K$ carries $\hat K$ onto each physical cell $K$ of the mesh (right). A physical
basis function is the reference shape function pulled back through the map,
$\phi_i = \hat\phi_i\circ F_K^{-1}$. The mesh supplies only the map; the polynomials
are shared across every cell. This is the $\hat\phi\to\mathsf{B}$, $F_K\to\mathsf{D}$
split that \cref{ch:matrixfree} runs on.}
\label{fig:refmap}
\end{figure}

The rest of the chapter fills in the three pieces. \Cref{sec:shape} builds the shape
functions on $\hat K$. \Cref{sec:map} builds the map $F_K$ and shows how it
transforms values, derivatives, and volumes. \Cref{sec:quad} replaces the integrals
by quadrature. \Cref{sec:vector} notes how the vector spaces of
\cref{ch:functionspaces} fit the same frame. \Cref{sec:tensor} collects the data
into the element-local tensor. Two worked examples (\cref{wex:massmatrix},
\cref{wex:jacobian}) carry the arithmetic all the way through.

\section{Nodal shape functions on the reference simplex}
\label{sec:shape}

We begin with the simplest and most important family: the \emph{nodal} or
\emph{Lagrange} shape functions, the basis of the $\Hone$ space of
\cref{ch:functionspaces}. These are the building blocks of every scalar potential
problem and the scaffolding from which the vector elements are built.

\subsection{Barycentric coordinates}

The natural coordinate system on a simplex (a triangle, a tetrahedron) is not
Cartesian but \emph{barycentric}. On the reference triangle with vertices
$\hat\vx_0,\hat\vx_1,\hat\vx_2$, the barycentric coordinates
$\lambda_0,\lambda_1,\lambda_2$ of a point $\hat\vx$ are the unique weights, summing
to one, that express $\hat\vx$ as a weighted average of the vertices:
\[
  \hat\vx = \lambda_0\hat\vx_0 + \lambda_1\hat\vx_1 + \lambda_2\hat\vx_2,
  \qquad \lambda_0+\lambda_1+\lambda_2 = 1.
\]
Geometrically, $\lambda_i$ is the fractional area of the sub-triangle opposite
vertex $i$: it is $1$ at vertex $i$, $0$ on the opposite edge, and falls off linearly
in between. For the standard reference triangle with $\hat\vx_0=(0,0)$,
$\hat\vx_1=(1,0)$, $\hat\vx_2=(0,1)$ they have the closed form
\[
  \lambda_0 = 1-\hat x-\hat y, \qquad
  \lambda_1 = \hat x, \qquad
  \lambda_2 = \hat y.
\]
The barycentric coordinates \emph{are} the lowest-order nodal shape functions: each
$\lambda_i$ is a linear polynomial equal to $1$ at its own vertex and $0$ at the
other two. We write $\hat\phi_i = \lambda_i$ for the $P_1$ (degree-one) basis. The
defining property is the \emph{cardinality} (or \emph{interpolation}, or
\emph{unisolvence}) relation from \cref{ch:functionspaces},
\begin{equation}
  \hat\phi_i(\hat\vx_j) = \delta_{ij},
  \label{eq:cardinality}
\end{equation}
``the $i$-th shape function is one at its own node and zero at every other node.''
This single relation is what makes the coefficients of a finite-element field equal
to its \emph{nodal values}: if $u_h = \sum_j c_j\hat\phi_j$, then evaluating at node
$i$ and using \eqref{eq:cardinality} gives $u_h(\hat\vx_i)=c_i$. The coefficient on
$\hat\phi_i$ is literally the field's value at node $i$---the DoF functional
$\sigma_i(u)=u(\hat\vx_i)$ of \cref{ch:functionspaces} reads it directly.

\begin{notationbox}
We mark reference-cell quantities with a hat: $\hat\vx$ for a reference point,
$\hat\phi_i$ for a reference shape function, $\hat K$ for the reference cell, and
$\Grad_{\hat\vx}$ for a gradient in reference coordinates. Unhatted symbols
($\vx,\phi_i,K,\Grad_\vx$) live on the physical cell. The map $F_K$ and its Jacobian
$J$ are the bridge between the two, and keeping the hats straight is the single most
common source of sign and transpose errors in finite-element code.
\end{notationbox}

\Cref{fig:reftet} shows the reference triangle and tetrahedron with their node
numbering. The triangle has three vertices (three $P_1$ nodes); the tetrahedron has
four. On the tetrahedron the barycentric coordinates are
$\lambda_0=1-\hat x-\hat y-\hat z$, $\lambda_1=\hat x$, $\lambda_2=\hat y$,
$\lambda_3=\hat z$---the same pattern, one more coordinate.

\begin{figure}[t]
\centering
\begin{tikzpicture}[scale=1.0, every node/.style={font=\scriptsize}]
% --- reference triangle with barycentric labels ---
\begin{scope}
  \node[font=\footnotesize\bfseries] at (1.1,2.55) {reference triangle};
  \fill[simBgBlue] (0,0)--(2.2,0)--(0,2.2)--cycle;
  \draw[simInk,thick] (0,0)--(2.2,0)--(0,2.2)--cycle;
  \fill[simBlue] (0,0)   circle (2.4pt) node[below left] {$0$};
  \fill[simBlue] (2.2,0) circle (2.4pt) node[below right]{$1$};
  \fill[simBlue] (0,2.2) circle (2.4pt) node[above left] {$2$};
  \node[simGray] at (0.35,0.35) {$\lambda_0$};
  \node[simGray] at (1.55,0.28) {$\lambda_1$};
  \node[simGray] at (0.3,1.55)  {$\lambda_2$};
  \node[simInk] at (1.1,-0.55) {$\lambda_0=1{-}\hat x{-}\hat y,\ \lambda_1=\hat x,\ \lambda_2=\hat y$};
\end{scope}
% --- reference tetrahedron ---
\begin{scope}[shift={(5.2,0)}]
  \node[font=\footnotesize\bfseries] at (1.1,2.55) {reference tetrahedron};
  \coordinate (n0) at (0,0);
  \coordinate (n1) at (2.2,0.3);
  \coordinate (n2) at (0.9,1.9);
  \coordinate (n3) at (1.0,0.7);
  \fill[simBgBlue,opacity=0.6] (n0)--(n1)--(n2)--cycle;
  \draw[simGray,densely dashed] (n0)--(n3) (n1)--(n3) (n2)--(n3);
  \draw[simInk,thick] (n0)--(n1)--(n2)--cycle;
  \fill[simBlue] (n0) circle (2.4pt) node[below left] {$0$};
  \fill[simBlue] (n1) circle (2.4pt) node[below right]{$1$};
  \fill[simBlue] (n2) circle (2.4pt) node[above]      {$2$};
  \fill[simBlue] (n3) circle (2.4pt) node[right]      {$3$};
  \node[simInk] at (1.1,-0.55) {$4$ vertices, $6$ edges, $4$ faces};
\end{scope}
\end{tikzpicture}
\caption[Reference simplices with node numbering]{The reference simplices and their
barycentric coordinates. \emph{Left:} the reference triangle, vertices numbered
$0,1,2$; the barycentric coordinate $\lambda_i$ is the $P_1$ shape function that
equals $1$ at node $i$ and falls linearly to $0$ on the opposite edge. \emph{Right:}
the reference tetrahedron, vertices $0$--$3$ (back edges dashed), with four
barycentric coordinates following the same pattern. These node counts---$3$ for the
triangle, $4$ for the tetrahedron---are the lowest-order $\Hone$ DoF counts of
\cref{ch:functionspaces}.}
\label{fig:reftet}
\end{figure}

\subsection{Higher order: the $P_p$ families}

The $P_1$ basis represents a field as piecewise-linear: exact on linear fields,
first-order accurate on smooth ones. To do better we raise the polynomial degree.
The degree-$p$ Lagrange space $P_p$ on the triangle consists of all bivariate
polynomials of total degree $\le p$; it has
\[
  \dim P_p = \binom{p+2}{2} = \frac{(p+1)(p+2)}{2}
\]
shape functions (the count of monomials $\hat x^a\hat y^b$ with $a+b\le p$). For
$p=1$ that is $3$ (the vertices); for $p=2$ it is $6$. The $P_2$ basis adds one node
at the midpoint of each edge, so its six nodes are the three vertices plus the three
edge midpoints. Each $P_2$ shape function is still defined by the cardinality
relation \eqref{eq:cardinality}---one at its own node, zero at the other five---and
each has a clean barycentric form. The three vertex functions are
$\hat\phi_i = \lambda_i(2\lambda_i-1)$, and the three edge-midpoint functions are
$\hat\phi_{ij} = 4\lambda_i\lambda_j$. (You can verify directly that
$\lambda_0(2\lambda_0-1)$ is $1$ at vertex $0$, where $\lambda_0=1$, and $0$ at every
other node, where either $\lambda_0=0$ or $\lambda_0=\tfrac12$.)

\Cref{fig:shapefns} plots the $P_1$ and $P_2$ shape functions in their
one-dimensional analogue, where the picture is cleanest. On the reference interval
$[0,1]$ the two $P_1$ hats are the straight lines $1-\hat x$ and $\hat x$; the three
$P_2$ functions are the parabolas $(1-\hat x)(1-2\hat x)$, $\hat x(2\hat x-1)$, and
the bump $4\hat x(1-\hat x)$ on the midpoint. In every case the cardinality property
is visible: each curve passes through $1$ at its own node and $0$ at the others.

\begin{figure}[t]
\centering
\begin{tikzpicture}
% --- P1 panel ---
\begin{axis}[
  name=p1, width=6.4cm, height=4.4cm,
  xlabel={$\hat x$}, ylabel={},
  title={\footnotesize $P_1$ shape functions (1-D)},
  xmin=0, xmax=1, ymin=-0.15, ymax=1.15,
  xtick={0,0.5,1}, ytick={0,0.5,1},
  tick label style={font=\scriptsize}, label style={font=\scriptsize},
  title style={font=\footnotesize}, axis lines=left, every axis plot/.append style={thick},
]
  \addplot[simBlue,  domain=0:1, samples=2] {1-x};
  \addplot[simRust,  domain=0:1, samples=2] {x};
  \node[font=\scriptsize,simBlue] at (axis cs:0.18,0.92) {$\hat\phi_0$};
  \node[font=\scriptsize,simRust] at (axis cs:0.82,0.92) {$\hat\phi_1$};
  \addplot[only marks, mark=*, mark size=1.2pt, black] coordinates {(0,1)(1,0)(0,0)(1,1)};
\end{axis}
% --- P2 panel ---
\begin{axis}[
  name=p2, at={(p1.east)}, anchor=west, xshift=14mm,
  width=6.4cm, height=4.4cm,
  xlabel={$\hat x$}, ylabel={},
  title={\footnotesize $P_2$ shape functions (1-D)},
  xmin=0, xmax=1, ymin=-0.35, ymax=1.15,
  xtick={0,0.5,1}, ytick={0,0.5,1},
  tick label style={font=\scriptsize}, label style={font=\scriptsize},
  title style={font=\footnotesize}, axis lines=left, every axis plot/.append style={thick},
]
  \addplot[simBlue,   domain=0:1, samples=60] {(1-x)*(1-2*x)};
  \addplot[simRust,   domain=0:1, samples=60] {x*(2*x-1)};
  \addplot[simGreen,  domain=0:1, samples=60] {4*x*(1-x)};
  \node[font=\scriptsize,simBlue]  at (axis cs:0.12,0.95) {$\hat\phi_0$};
  \node[font=\scriptsize,simRust]  at (axis cs:0.88,0.95) {$\hat\phi_1$};
  \node[font=\scriptsize,simGreen] at (axis cs:0.5,1.02)  {$\hat\phi_{01}$};
  \addplot[only marks, mark=*, mark size=1.2pt, black] coordinates {(0,1)(0.5,0)(1,0)(0,0)(0.5,1)(1,1)};
\end{axis}
\end{tikzpicture}
\caption[$P_1$ and $P_2$ shape functions]{Lagrange shape functions, shown in their
one-dimensional analogue where the curves are uncluttered. \emph{Left:} the two
$P_1$ ``hat'' functions, $\hat\phi_0=1-\hat x$ and $\hat\phi_1=\hat x$, each linear,
each one at its own node ($\hat x=0$ and $\hat x=1$) and zero at the other.
\emph{Right:} the three $P_2$ functions on $[0,1]$: two vertex parabolas
$\hat\phi_0=(1-\hat x)(1-2\hat x)$, $\hat\phi_1=\hat x(2\hat x-1)$, and the
edge-midpoint bump $\hat\phi_{01}=4\hat x(1-\hat x)$ centred at $\hat x=\tfrac12$.
The dots mark the cardinality property \eqref{eq:cardinality}: every shape function
passes through $1$ at its own node and $0$ at all the others. On the triangle the
same functions appear in barycentric form, $\lambda_i(2\lambda_i-1)$ and
$4\lambda_i\lambda_j$.}
\label{fig:shapefns}
\end{figure}

The pattern continues to any order. Raising $p$ is exactly the $p$-refinement of
\cref{ch:functionspaces}: more shape functions per cell, a richer polynomial space,
faster convergence on smooth fields---and, in Palace, the family of collections
\code{H1\_FECollection}, \code{ND\_FECollection}, and the rest selecting both the de
Rham slot and the order. The point for this chapter is structural: \emph{whatever}
the order, the shape functions live on the reference cell and are tabulated once.

\section{The geometric map and how it transforms calculus}
\label{sec:map}

The reference cell carries the polynomials; the mesh carries the geometry. The bridge
is the geometric map $F_K$, and the heart of this section is what that map does to
the three operations a weak form needs: \emph{evaluating} a field, \emph{differentiating}
it, and \emph{integrating} it. Get the map's bookkeeping right and the same reference
tabulation serves cells of every size, shape, and orientation.

\subsection{The map and its Jacobian}

For each physical cell $K$ the geometric map is a function
\[
  \vx = F_K(\hat\vx)
\]
taking reference coordinates $\hat\vx\in\hat K$ to physical coordinates $\vx\in K$.
For a straight-sided (\emph{affine}) simplex the map is itself built from the nodal
shape functions: writing the physical vertices as $\vx_0,\vx_1,\vx_2$,
\[
  F_K(\hat\vx) = \sum_i \vx_i\,\hat\phi_i(\hat\vx)
  = \vx_0 + (\vx_1-\vx_0)\,\hat x + (\vx_2-\vx_0)\,\hat y .
\]
The map is \emph{affine}: a constant matrix times $\hat\vx$ plus a shift. Its
derivative, the \emph{Jacobian}, is the constant matrix
\[
  J = \frac{\partial F_K}{\partial\hat\vx}
    = \bigl[\;\vx_1-\vx_0 \quad \vx_2-\vx_0\;\bigr]
    \in\Real^{d\times d},
\]
whose columns are the physical edge vectors emanating from vertex $0$. (When the cell
is \emph{curved}---a high-order ``isoparametric'' element whose faces bend to follow
a curved boundary---the map is built from higher-order shape functions and $J$ varies
from point to point. Palace handles both; the affine case is the one to hold in your
head, and the one this chapter computes.) The Jacobian is the single object that
encodes how this particular cell differs from the reference: its determinant measures
how the cell's \emph{volume} compares to the reference, and its inverse-transpose
measures how \emph{directions} are bent.

\subsection{Three transformation rules}

A scalar field is the simplest case. A physical basis function is just the reference
shape function read in reference coordinates:
\[
  \phi_i(\vx) = \hat\phi_i\bigl(F_K^{-1}(\vx)\bigr),
  \qquad\text{equivalently}\qquad
  \phi_i\bigl(F_K(\hat\vx)\bigr) = \hat\phi_i(\hat\vx).
\]
\emph{Values do not change}: to evaluate $\phi_i$ at a physical point, find its
reference preimage and read the table. This is why basis values need to be tabulated
only once.

Derivatives are where the map bites. The weak forms of \cref{ch:weak} integrate
\emph{gradients} of basis functions---the $\Grad\phi_i\cdot\Grad\phi_j$ of the
stiffness matrix---and a gradient taken in physical coordinates is not the gradient
taken in reference coordinates. By the chain rule, differentiating
$\hat\phi_i(\hat\vx)=\phi_i(F_K(\hat\vx))$ with respect to $\hat\vx$,
\[
  \Grad_{\hat\vx}\hat\phi_i = J^{\!\top}\,\Grad_{\vx}\phi_i,
  \qquad\text{hence}\qquad
  \boxed{\;\Grad_{\vx}\phi_i = J^{-\top}\,\Grad_{\hat\vx}\hat\phi_i\;}.
\]
The physical gradient is the \emph{inverse-transpose} Jacobian applied to the
reference gradient. The reference gradient $\Grad_{\hat\vx}\hat\phi_i$ is tabulated
once (for $P_1$ it is the constant $\Grad\lambda_i$ we already wrote down); the cell
contributes only the small matrix $J^{-\top}$. \Cref{fig:jacobian} draws the picture:
the inverse-transpose is precisely the operator that keeps a gradient pointing
``across'' the level sets of $\phi_i$ after the cell is stretched and skewed.

Finally, integration. The change-of-variables theorem turns a physical integral into
a reference integral, paying a \emph{volume factor} $|\det J|$:
\[
  \int_K g(\vx)\,d\vx = \int_{\hat K} g\bigl(F_K(\hat\vx)\bigr)\,\abs{\det J}\,d\hat\vx .
\]
The factor $|\det J|$ is the ratio of the physical cell's volume to the reference
cell's: for the affine triangle, $|\det J| = 2\,\mathrm{area}(K)$ (since the reference
triangle has area $\tfrac12$). Every matrix entry in the book is an integral of this
form, and the recipe is always: \emph{pull back to the reference cell, multiply by
$|\det J|$, and integrate there}---where the basis values are already tabulated.

\begin{figure}[t]
\centering
\begin{tikzpicture}[scale=1.0, every node/.style={font=\scriptsize}]
% reference: gradient perpendicular to level sets
\begin{scope}
  \node[font=\footnotesize\bfseries] at (1.0,2.4) {reference $\hat K$};
  \fill[simBgBlue] (0,0)--(2,0)--(0,2)--cycle;
  \draw[simInk,thick] (0,0)--(2,0)--(0,2)--cycle;
  % level sets of lambda1 = xhat (vertical lines)
  \foreach \x in {0.5,1.0,1.5}{\draw[simGray,densely dotted] (\x,0)--(\x,{2-\x});}
  % gradient arrow (points +x)
  \draw[-{Stealth[length=2.4mm]},simBlue,very thick] (0.7,0.5)--(1.5,0.5);
  \node[simBlue] at (1.1,0.78) {$\Grad_{\hat\vx}\hat\phi$};
  \node[simGray] at (1.45,1.55) {level sets};
\end{scope}
% the J^{-T} map arrow
\draw[flow,simRust,thick] (2.5,1.0) .. controls (3.3,1.4) and (3.9,1.4) .. (4.7,1.0)
  node[midway,above,font=\scriptsize]{$J^{-\top}$};
% physical: skewed level sets, gradient still perpendicular to THEM
\begin{scope}[shift={(5.3,0)}]
  \node[font=\footnotesize\bfseries] at (1.2,2.4) {physical $K$};
  \fill[simBgRust] (0,0)--(2.4,0.4)--(0.7,2.0)--cycle;
  \draw[simInk,thick] (0,0)--(2.4,0.4)--(0.7,2.0)--cycle;
  % skewed level sets
  \draw[simGray,densely dotted] (0.6,0.1)--(0.25,1.1);
  \draw[simGray,densely dotted] (1.2,0.2)--(0.5,1.55);
  \draw[simGray,densely dotted] (1.8,0.3)--(0.78,1.85);
  % gradient: must stay perpendicular to the SKEWED level sets (so it tilts)
  \draw[-{Stealth[length=2.4mm]},simBlue,very thick] (0.8,0.55)--(1.55,0.78);
  \node[simBlue] at (1.25,1.05) {$\Grad_{\vx}\phi$};
  \node[simGray] at (1.6,1.6) {level sets bent};
\end{scope}
\end{tikzpicture}
\caption[The inverse-transpose Jacobian transforms gradients]{Why the gradient
transforms by $J^{-\top}$, not by $J$. \emph{Left:} on the reference cell the level
sets of a shape function $\hat\phi$ (dotted) are evenly spaced and the gradient (blue
arrow) is perpendicular to them. \emph{Right:} the map skews the cell and bends the
level sets; the physical gradient must remain perpendicular to the \emph{bent} level
sets. The operator that does this is the inverse-transpose Jacobian $J^{-\top}$,
applied to the tabulated reference gradient. Tabulate $\Grad_{\hat\vx}\hat\phi$ once;
each cell supplies only its small matrix $J^{-\top}$. In \cref{ch:matrixfree} this
$J^{-\top}$, combined with $|\det J|$ and the material coefficient, is exactly the
per-quadrature-point geometry factor that the $\mathsf{D}$ stage applies.}
\label{fig:jacobian}
\end{figure}

\subsection{The geometry factor: assembling the metric}

Collect the pieces and a clean pattern emerges. The two integrals every solver in
this book needs are the \emph{mass} term, $\int_K \phi_i\,\phi_j\,d\vx$, and the
\emph{stiffness} (grad--grad) term, $\int_K \Grad\phi_i\cdot\Grad\phi_j\,d\vx$.
Pulling each back to the reference cell:
\[
  \int_K \phi_i\phi_j\,d\vx
  = \int_{\hat K} \hat\phi_i\hat\phi_j\;\underbrace{\abs{\det J}}_{\text{mass factor}}\,d\hat\vx,
\]
\[
  \int_K \Grad\phi_i\cdot\Grad\phi_j\,d\vx
  = \int_{\hat K} \bigl(J^{-\top}\Grad_{\hat\vx}\hat\phi_i\bigr)\cdot
                  \bigl(J^{-\top}\Grad_{\hat\vx}\hat\phi_j\bigr)\,\abs{\det J}\,d\hat\vx
  = \int_{\hat K} \Grad_{\hat\vx}\hat\phi_i\cdot
      \underbrace{\bigl(J^{-1}J^{-\top}\abs{\det J}\bigr)}_{\text{metric }G}\,
      \Grad_{\hat\vx}\hat\phi_j\;d\hat\vx .
\]
The mass integral needs only the scalar $|\det J|$; the stiffness integral needs the
$d\times d$ symmetric matrix $G = J^{-1}J^{-\top}\,|\det J|$, the \emph{metric}
introduced by the map. These two objects---$|\det J|$ for mass terms,
$J^{-1}J^{-\top}|\det J|$ for grad--grad terms---are \emph{all} the geometry a cell
contributes. Everything else is reference data.

\begin{keyidea}
A cell contributes to the integrals only through a small \emph{geometry factor}: the
scalar $|\det J|$ for mass terms and the matrix $G = J^{-1}J^{-\top}\,|\det J|$ for
grad--grad terms. The basis values and reference gradients are tabulated once and
shared by every cell; the cell supplies only its Jacobian-derived factor. This is
precisely the \code{geom\_factor\_build} stage of the matrix-free pipeline: it
precomputes $G$ (or $|\det J|$), per quadrature point, once per mesh, and the
material coefficient is pre-multiplied into it so the run-time apply is a single
pointwise multiply.
\end{keyidea}

This factorization is the content of Palace's \code{geom\_factor\_build}: a
\emph{setup-stratum} pass that computes the geometry metric at every quadrature point
once per mesh and stores it, so that the per-apply work---which happens thousands of
times inside a Krylov solve---is reduced to multiplying by a stored number. We meet
that stratification again, in full, in \cref{ch:matrixfree}; the point here is that it
falls straight out of the change-of-variables rules above. On an affine cell $J$ is
constant, so $G$ is constant over the cell and the stored factor is one matrix per
cell; on a curved cell $J$ varies and $G$ is stored per quadrature point.

\section{Quadrature: integrals as weighted sums}
\label{sec:quad}

We can now pull every integral back to the reference cell. One step remains: actually
\emph{evaluating} the reference integral $\int_{\hat K} g(\hat\vx)\,d\hat\vx$. Except
in the simplest cases we do not compute it in closed form. We approximate it by a
\emph{quadrature rule}---a finite weighted sum of the integrand sampled at a few
chosen points.

\subsection{The quadrature idea}

A quadrature rule on $\hat K$ is a set of \emph{points} $\hat\vx_q\in\hat K$ and
\emph{weights} $w_q$ such that
\begin{equation}
  \int_{\hat K} g(\hat\vx)\,d\hat\vx \;\approx\; \sum_{q=1}^{Q} w_q\,g(\hat\vx_q).
  \label{eq:quad}
\end{equation}
The integral---an infinite, continuous operation---is replaced by a finite sum: sample
the integrand at the $Q$ quadrature points, scale each sample by its weight, add them
up. The art is in choosing the points and weights so the sum is accurate. The
classical answer is \emph{Gauss quadrature}: for a given number of points, place them
and weight them so that the rule integrates polynomials \emph{exactly} up to as high
a degree as possible.

\begin{definition}[Exactness degree]
A quadrature rule has \emph{exactness degree} $r$ if \eqref{eq:quad} is an exact
equality for every polynomial $g$ of total degree $\le r$, and fails for some
polynomial of degree $r+1$. A $Q$-point Gauss rule on an interval achieves exactness
degree $2Q-1$: $Q$ points buy $2Q$ free parameters (the points and the weights), and
$2Q$ parameters pin down a polynomial of degree $2Q-1$.
\end{definition}

Exactness degree is the contract a quadrature rule offers, and it is exactly what a
finite-element assembler needs to honor. The integrand of a mass-matrix entry is
$\hat\phi_i\hat\phi_j\,|\det J|$; on $P_p$ elements with an affine map ($|\det J|$
constant) this is a polynomial of degree $2p$. To integrate it without error the rule
must have exactness degree at least $2p$. Choose the rule too weak and the mass matrix
is computed inaccurately---\emph{under-integrated}, a real and sometimes deliberate
choice; choose it stronger than needed and you simply pay for extra points
(\emph{over-integration}, sometimes used to tame variable coefficients). The
exactness degree is the dial.

\begin{figure}[t]
\centering
\begin{tikzpicture}[scale=1.55, every node/.style={font=\scriptsize}]
% three reference triangles with quadrature points
\foreach \dx/\ttl/\desc in {0/{1-point}/{degree 1}, 2.7/{3-point}/{degree 2}, 5.4/{6-point}/{degree 4}}{
  \begin{scope}[shift={(\dx,0)}]
    \fill[simBgBlue] (0,0)--(1.4,0)--(0,1.4)--cycle;
    \draw[simInk,thick] (0,0)--(1.4,0)--(0,1.4)--cycle;
    \node[font=\footnotesize\bfseries] at (0.6,1.62) {\ttl};
    \node[simGray] at (0.6,-0.3) {\desc};
  \end{scope}
}
% 1-point: centroid
\fill[simRust] (0.467,0.467) circle (1.6pt);
% 3-point: at (1/6,1/6),(2/3,1/6),(1/6,2/3) scaled by 1.4
\foreach \px/\py in {0.233/0.233, 0.933/0.233, 0.233/0.933}{
  \fill[simRust] (2.7+\px,\py) circle (1.4pt);}
% 6-point (degree 4): two orbits — approximate symmetric points
\foreach \px/\py in {0.63/0.63, 0.13/0.63, 0.63/0.13, 1.04/0.18, 0.18/1.04, 0.18/0.18}{
  \fill[simRust] (5.4+\px,\py) circle (1.2pt);}
\end{tikzpicture}
\caption[Gauss quadrature points on the reference triangle]{Symmetric Gauss
quadrature rules on the reference triangle, in order of increasing exactness.
\emph{Left:} the $1$-point rule places a single node at the centroid with weight
equal to the triangle's area; it is exact to degree $1$ (it integrates linear
functions exactly). \emph{Middle:} the $3$-point rule (nodes near the edge midpoints,
each weight one-third of the area) is exact to degree $2$. \emph{Right:} a $6$-point
rule, exact to degree $4$, suitable for $P_2$ mass matrices. More points buy higher
exactness; the assembler picks the smallest rule whose exactness degree covers the
integrand's polynomial degree. These points and weights are the discrete realization
of the integral $\int_{\hat K}(\cdot)\,d\hat\vx$, and they are the $P$ axis of the
element-local tensor (\cref{sec:tensor}).}
\label{fig:quadpts}
\end{figure}

\Cref{fig:quadpts} shows three symmetric rules on the reference triangle. The
$1$-point rule samples only the centroid; the $3$-point rule samples three interior
points; a $6$-point rule reaches exactness degree $4$. On the triangle the weights are
fractions of the reference area $\tfrac12$, and the points are placed symmetrically so
that the rule respects the triangle's symmetry. The quadrature points and weights are,
quite literally, the discrete realization of the integral: once the rule is chosen,
``integrate over the cell'' \emph{means} ``sample at these points and take this
weighted sum,'' nothing more.

\begin{physicsbox}
Quadrature is where the continuous integral of the weak form becomes a finite
computation a machine can run. In the matrix-free pipeline of \cref{ch:matrixfree},
the quadrature points are the $P$ axis, the basis tabulation $\mathsf{B}$ evaluates
the trial field \emph{at} those points, the geometry-and-material factor $\mathsf{D}$
multiplies pointwise \emph{at} those points, and the transpose $\mathsf{B}^{\!\top}$
sums the contributions back with the weights folded in. The entire operator apply is
organized around the quadrature points: they are the grain at which the discrete
physics is evaluated.
\end{physicsbox}

\subsection{Tying quadrature to the matrix-free \texttt{D} stage}

It is worth being explicit about where quadrature lands in the pipeline, because it
is the conceptual hinge between this chapter and \cref{ch:matrixfree}. An operator
apply $\vy = \mathsf{A}\vx$ in the matrix-free scheme is the composition
\[
  \mathsf{A} = \mathsf{G}^{\!\top}\,\mathsf{B}^{\!\top}\,\mathsf{D}\,\mathsf{B}\,\mathsf{G},
\]
read right to left: $\mathsf{G}$ gathers the global coefficient vector into
per-element local tensors; $\mathsf{B}$ applies the tabulated reference basis to get
field values at the quadrature points; $\mathsf{D}$ multiplies, point by point, by the
geometry-and-material factor (the $|\det J|$ or $G$ of \cref{sec:map}, with the
quadrature weight $w_q$ and the material coefficient pre-multiplied in); $\mathsf{B}^{\!\top}$
contracts the weighted point values back against the basis; $\mathsf{G}^{\!\top}$
scatters to the global vector. The weighted sum \eqref{eq:quad} is split across two
stages: the weight $w_q$ rides inside $\mathsf{D}$, and the summation $\sum_q$ is the
contraction that $\mathsf{B}^{\!\top}$ performs. Quadrature is not a separate step
bolted on---it \emph{is} the structure of the apply.

\section{Vector elements and the Piola maps}
\label{sec:vector}

The nodal elements of \cref{sec:shape} are the scalar case. The vector spaces
$\Hcurl$ and $\Hdiv$ of \cref{ch:functionspaces}---the homes of $\vE$, $\vA$, and the
fluxes $\vB$, $\vD$---use \emph{vector-valued} shape functions, and they map to the
physical cell by a different rule. We flag the essentials here and leave the full
edge-basis machinery to the references; the point is that the reference-element frame
carries over intact, with one twist in the transformation.

A N\'ed\'elec ($\Hcurl$) reference shape function is a vector field on $\hat K$---the
Whitney function $\hat\phi_{01}=\lambda_0\Grad\lambda_1-\lambda_1\Grad\lambda_0$ of
\cref{wex:whitney} is the lowest-order example. A Raviart--Thomas ($\Hdiv$) reference
function is likewise vector-valued. The subtlety is that mapping a vector field to a
skewed physical cell by the \emph{value-preserving} rule of the scalar case would
\emph{destroy} the tangential or normal continuity that \cref{ch:functionspaces}
showed is the whole point of these spaces. A scalar value is the same in any
coordinates; a vector's tangential and normal parts are not.

The fix is the \emph{Piola transforms}, which are designed to preserve exactly the
right component:
\[
  \text{covariant (}\Hcurl\text{):}\quad \vu(\vx) = J^{-\top}\hat\vu(\hat\vx),
  \qquad
  \text{contravariant (}\Hdiv\text{):}\quad \vu(\vx) = \frac{1}{\det J}\,J\,\hat\vu(\hat\vx).
\]
The covariant (curl-conforming) map applies $J^{-\top}$---the same operator that
transforms gradients---and it preserves \emph{tangential} components along edges; this
is why edge circulations survive the map and tangential continuity is maintained
cell to cell. The contravariant (divergence-conforming) map applies $J/\det J$ and
preserves \emph{normal} components through faces, maintaining the flux continuity of
$\Hdiv$.

\begin{keyidea}
The reference-element idea extends to the vector spaces, but with a
\emph{component-preserving} map in place of the value-preserving one. The covariant
Piola map $J^{-\top}$ preserves tangential traces (so $\Hcurl$ edge circulations
survive) and the contravariant Piola map $J/\det J$ preserves normal traces (so
$\Hdiv$ face fluxes survive). The continuity that \cref{ch:functionspaces} built into
the DoFs is preserved by the map---which is what makes a single edge or face DoF,
shared by two cells, mean the same physical quantity on both.
\end{keyidea}

The remarkable fact---which \cref{ch:derham}'s de Rham complex explains---is that the
three maps fit together: nodal values transform trivially, gradients and $\Hcurl$
fields transform by $J^{-\top}$, $\Hdiv$ fields by $J/\det J$, and the chain
$\Grad\to\Curl\to\Div$ commutes with the maps. For this chapter the operative point is
narrow and practical: a vector element is still a reference shape function tabulated
once, carried to each cell by a Jacobian-built map. Only the map's formula changes.

\section{The element-local tensor: $E$, $L$, $P$, $C$, $G$}
\label{sec:tensor}

We close by collecting everything this chapter has produced into the single data
structure the matrix-free pipeline of \cref{ch:matrixfree} runs on. Every quantity
we have met---the local degrees of freedom, the quadrature points, the field
components, the geometry factor---is an axis of a small per-element tensor. Naming
those axes now makes \cref{ch:matrixfree} a matter of contraction bookkeeping rather
than fresh concepts.

Restrict the global field to one element and you have a small vector: one coefficient
for each local DoF of that element. Stacked over all elements, this is a tensor with
two axes,
\[
  \texttt{Tensor[(E, L)]},
\]
where $E$ counts the mesh elements and $L$ counts the local DoFs per element (for
$P_1$ on a triangle, $L=3$; the reference shape functions of \cref{sec:shape}).
Apply the tabulated basis $\mathsf{B}$---evaluate those $L$ shape functions at the
$P$ quadrature points---and you get the field's values at every quadrature point of
every element,
\[
  \texttt{Tensor[(E, P, C)]},
\]
where $P$ counts quadrature points per element (the rule of \cref{sec:quad}) and $C$
counts value components ($C=1$ for a scalar field; $C=d$ for a gradient or a vector
field). Finally the geometry factor of \cref{sec:map}---the per-point $|\det J|$ or
metric $G$, with weight and material pre-multiplied in---is itself a per-point tensor,
\[
  \texttt{Tensor[(E, P, G)]},
\]
with $G$ counting the stored geometry-metric components. These three tensors, over the
five named axes $E,L,P,C,G$, are the complete vocabulary of element-local finite
element computation.

\begin{figure}[t]
\centering
\begin{tikzpicture}[scale=1.0, every node/.style={font=\scriptsize}]
% [E,L] block
\node[draw, simBlue, fill=simBgBlue, thick, minimum width=20mm, minimum height=13mm,
  align=center] (el) at (0,0) {\code{Tensor}\\\code{[(E, L)]}};
\node[font=\scriptsize, simGray, align=center] at (0,-1.35)
  {per-element\\local DoFs};
% B arrow
\draw[flow, simInk] (1.1,0)--(3.0,0)
  node[midway, above, font=\footnotesize] {$\mathsf{B}$}
  node[midway, below, font=\scriptsize, simGray] {tabulated basis};
% [E,P,C] block
\node[draw, simRust, fill=simBgRust, thick, minimum width=22mm, minimum height=13mm,
  align=center] (epc) at (4.2,0) {\code{Tensor}\\\code{[(E, P, C)]}};
\node[font=\scriptsize, simGray, align=center] at (4.2,-1.35)
  {values at\\quad points};
% D arrow (with geom factor coming in)
\draw[flow, simInk] (5.4,0)--(7.3,0)
  node[midway, above, font=\footnotesize] {$\mathsf{D}$}
  node[midway, below, font=\scriptsize, simGray] {pointwise};
% geom factor feeding D
\node[draw, simGreen, fill=simBgGreen, semithick, minimum width=20mm, minimum height=9mm,
  align=center] (epg) at (6.35,1.5) {\code{Tensor[(E, P, G)]}};
\node[font=\scriptsize, simGreen] at (6.35,2.15) {geometry factor};
\draw[depedge, simGreen] (epg.south)--(6.35,0.42);
% output
\node[draw, simRust, fill=simBgRust, thick, minimum width=22mm, minimum height=13mm,
  align=center] (out) at (8.5,0) {\code{Tensor}\\\code{[(E, P, C)]}};
\node[font=\scriptsize, simGray, align=center] at (8.5,-1.35)
  {weighted\\values};
\end{tikzpicture}
\caption[The element-local tensor and its axes]{The element-local tensor pipeline,
the data structure \cref{ch:matrixfree} runs on. A field restricted to elements is a
\code{Tensor[(E, L)]} of per-element local DoFs ($E$ elements, $L$ local DoFs). The
tabulated basis $\mathsf{B}$ of \cref{sec:shape}---the reference shape functions
evaluated at the quadrature points---contracts it to a \code{Tensor[(E, P, C)]} of
field values at the $P$ quadrature points ($C$ components). The pointwise stage
$\mathsf{D}$ multiplies by the geometry-and-material factor
\code{Tensor[(E, P, G)]}---the $|\det J|$ or metric $G$ of \cref{sec:map}, with
quadrature weight and material coefficient pre-multiplied ($G$ metric components).
The five axes $E$, $L$, $P$, $C$, $G$ are exactly the quantities this chapter built:
elements, reference DoFs, quadrature points, field components, and geometry-metric
storage.}
\label{fig:elemtensor}
\end{figure}

\Cref{fig:elemtensor} draws the flow. Read it against the three sections that built
the pieces: the $L$ axis is the reference shape functions of \cref{sec:shape}; the
$P$ axis is the quadrature rule of \cref{sec:quad}; the $C$ axis is the value or
derivative components selected by the weak-form differential operator; the $G$ axis is
the geometry factor of \cref{sec:map}; and $E$ simply ranges over the cells of the
mesh, each carrying its own Jacobian but sharing the one reference tabulation. This is
the payoff of the reference-element idea stated as a data structure: the expensive,
shared content (the basis) is the small $\mathsf{B}$ table indexed by $(L,P)$, and the
cell-specific content (the geometry) is the $(E,P,G)$ factor---built once, by
\code{geom\_factor\_build}, and reused across every apply.

\begin{notationbox}
The five element-local axes recur throughout the framework chapters:
$E$ (mesh elements), $L$ (local DoFs per element), $P$ (quadrature points per
element), $C$ (field/derivative components), $G$ (geometry-metric components). The
global true-DoF axis $N$ of \cref{ch:functionspaces} is \emph{not} one of
them---$N$ is the flat global vector, and the element-restriction operator
$\mathsf{G}$ is precisely the boundary between the global $N$ vocabulary and the
element-local $(E,L)$ vocabulary. \Cref{ch:matrixfree} works entirely in these axes.
\end{notationbox}

\section{Worked examples}

The two examples carry the chapter's arithmetic to the end. The first integrates a
product of shape functions by quadrature and checks it against the exact mass-matrix
entry, exhibiting the exactness-degree contract concretely. The second computes a
Jacobian and transforms a gradient, exhibiting the map's three rules on a specific
physical triangle.

\subsection{A mass-matrix entry by quadrature}

\begin{workedexample}{Integrating two $P_1$ hats: exact, 1-point, and 3-point}{massmatrix}
Work on the reference triangle $\hat K$ with vertices $(0,0),(1,0),(0,1)$ and area
$\tfrac12$. The three $P_1$ shape functions are the barycentric coordinates
\[
  \hat\phi_0 = 1-\hat x-\hat y, \qquad \hat\phi_1 = \hat x, \qquad \hat\phi_2 = \hat y .
\]
We compute the reference mass-matrix entry
$m_{01} = \int_{\hat K}\hat\phi_0\hat\phi_1\,d\hat\vx
        = \int_{\hat K}(1-\hat x-\hat y)\,\hat x\,d\hat\vx$
three ways.

\smallskip\noindent\textbf{Exact.} There is a standard formula for integrating
products of barycentric coordinates over a triangle of area $A$:
\[
  \int_{\hat K}\lambda_0^{a}\lambda_1^{b}\lambda_2^{c}\,d\hat\vx
  = 2A\,\frac{a!\,b!\,c!}{(a+b+c+2)!}.
\]
With $A=\tfrac12$ and $(a,b,c)=(1,1,0)$ this gives
$m_{01} = 2\cdot\tfrac12\cdot\dfrac{1!\,1!\,0!}{4!} = \dfrac{1}{24}\approx 0.041667$.
(For comparison, the diagonal entry with $(a,b,c)=(2,0,0)$ is
$2\cdot\tfrac12\cdot\tfrac{2!}{4!}=\tfrac{1}{12}$.) So the exact answer is
$m_{01} = \tfrac{1}{24}$.

\smallskip\noindent\textbf{1-point rule (centroid, degree 1).} The one-point Gauss
rule samples the centroid $\hat\vx_c=(\tfrac13,\tfrac13)$ with weight equal to the
area, $w=\tfrac12$. The integrand there is
$\hat\phi_0\hat\phi_1 = (1-\tfrac13-\tfrac13)\cdot\tfrac13 = \tfrac13\cdot\tfrac13 = \tfrac19$,
so
\[
  m_{01}^{(1)} = w\,\hat\phi_0(\hat\vx_c)\hat\phi_1(\hat\vx_c)
              = \tfrac12\cdot\tfrac19 = \tfrac{1}{18} \approx 0.055556 .
\]
This is \emph{wrong}---off by a third. The reason is the exactness contract: the
integrand $\hat\phi_0\hat\phi_1$ is a quadratic (degree $2$), but the $1$-point rule
is exact only to degree $1$. Under-integration shows up as a concrete numerical error.

\smallskip\noindent\textbf{3-point rule (degree 2).} The symmetric three-point rule
places nodes at $(\tfrac16,\tfrac16)$, $(\tfrac23,\tfrac16)$, $(\tfrac16,\tfrac23)$,
each with weight $w_q=\tfrac16$ (one-third of the area $\tfrac12$). Evaluate the
integrand $\hat\phi_0\hat\phi_1=(1-\hat x-\hat y)\,\hat x$ at each:
\begin{align*}
  q=1:\ (\tfrac16,\tfrac16) &: (1-\tfrac13)\cdot\tfrac16 = \tfrac23\cdot\tfrac16 = \tfrac{1}{9}, \\
  q=2:\ (\tfrac23,\tfrac16) &: (1-\tfrac56)\cdot\tfrac23 = \tfrac16\cdot\tfrac23 = \tfrac{1}{9}, \\
  q=3:\ (\tfrac16,\tfrac23) &: (1-\tfrac56)\cdot\tfrac16 = \tfrac16\cdot\tfrac16 = \tfrac{1}{36}.
\end{align*}
The weighted sum is
\[
  m_{01}^{(3)} = \tfrac16\Bigl(\tfrac19+\tfrac19+\tfrac1{36}\Bigr)
              = \tfrac16\cdot\Bigl(\tfrac{4+4+1}{36}\Bigr)
              = \tfrac16\cdot\tfrac{9}{36} = \tfrac16\cdot\tfrac14 = \tfrac{1}{24}.
\]
This is \emph{exact}. The three-point rule has exactness degree $2$, the integrand
has degree $2$, and the contract is honored: the quadrature equals the closed-form
integral to the last digit. The lesson is the practical rule of \cref{sec:quad}: to
integrate a $P_p$ mass matrix exactly on an affine cell, use a rule of exactness
degree $\ge 2p$. Here $p=1$, $2p=2$, and the three-point (degree-$2$) rule is the
smallest that suffices---the one-point rule, degree $1<2$, is not enough.
\end{workedexample}

\subsection{A Jacobian and a transformed gradient}

\begin{workedexample}{The map of one physical triangle, and a gradient through it}{jacobian}
Map the reference triangle $\hat K$ (vertices $(0,0),(1,0),(0,1)$) onto the physical
triangle $K$ with vertices
\[
  \vx_0 = (1,1), \qquad \vx_1 = (4,2), \qquad \vx_2 = (2,5).
\]
\textbf{The Jacobian.} The affine map is
$F_K(\hat\vx) = \vx_0 + (\vx_1-\vx_0)\hat x + (\vx_2-\vx_0)\hat y$, so the Jacobian is
the constant matrix whose columns are the physical edge vectors from $\vx_0$:
\[
  J = \bigl[\,\vx_1-\vx_0 \ \ \vx_2-\vx_0\,\bigr]
    = \begin{bmatrix} 3 & 1 \\ 1 & 4 \end{bmatrix}.
\]
Its determinant is $\det J = 3\cdot 4 - 1\cdot 1 = 11$, so $|\det J| = 11$, and indeed
the area of $K$ is $\tfrac12|\det J| = \tfrac{11}{2}$, eleven times the reference
area $\tfrac12$. The volume factor for any mass integral on this cell is $11$.

\smallskip\noindent\textbf{The inverse-transpose.} The inverse Jacobian is
\[
  J^{-1} = \frac{1}{11}\begin{bmatrix} 4 & -1 \\ -1 & 3 \end{bmatrix},
  \qquad
  J^{-\top} = \frac{1}{11}\begin{bmatrix} 4 & -1 \\ -1 & 3 \end{bmatrix},
\]
which here equals $J^{-1}$ because $J$ is symmetric (an accident of this example;
in general $J^{-\top}\ne J^{-1}$).

\smallskip\noindent\textbf{A transformed gradient.} On the reference cell the shape
function $\hat\phi_1=\hat x$ has constant reference gradient
$\Grad_{\hat\vx}\hat\phi_1 = (1,0)^{\!\top}$. Its physical gradient is
\[
  \Grad_{\vx}\phi_1 = J^{-\top}\Grad_{\hat\vx}\hat\phi_1
  = \frac{1}{11}\begin{bmatrix} 4 & -1 \\ -1 & 3 \end{bmatrix}
    \begin{bmatrix} 1 \\ 0 \end{bmatrix}
  = \frac{1}{11}\begin{bmatrix} 4 \\ -1 \end{bmatrix}
  = \bigl(\tfrac{4}{11},\,-\tfrac{1}{11}\bigr).
\]
Notice it has \emph{tilted}: the reference gradient pointed purely along $\hat x$, but
on the skewed physical cell the gradient of $\phi_1$ has acquired a $y$-component,
exactly as \cref{fig:jacobian} promised---it stays perpendicular to the (now bent)
level sets of $\phi_1$. Likewise $\hat\phi_2=\hat y$ has reference gradient
$(0,1)^{\!\top}$ and physical gradient
$J^{-\top}(0,1)^{\!\top}=(-\tfrac{1}{11},\tfrac{3}{11})$, and
$\hat\phi_0=1-\hat x-\hat y$ has reference gradient $(-1,-1)^{\!\top}$ and physical
gradient $(-\tfrac{3}{11},-\tfrac{2}{11})$. As a check, the three physical gradients
sum to zero---because $\phi_0+\phi_1+\phi_2\equiv 1$ on the cell, a constant, whose
gradient vanishes.

\smallskip\noindent\textbf{Assembling a stiffness entry.} With the gradients in hand,
a grad--grad (stiffness) entry on this affine cell---where $J$, hence everything, is
constant---is just the constant integrand times the cell volume:
\[
  k_{12} = \int_K \Grad\phi_1\cdot\Grad\phi_2\,d\vx
  = \bigl(\Grad_{\vx}\phi_1\cdot\Grad_{\vx}\phi_2\bigr)\,\mathrm{area}(K)
  = \Bigl(\tfrac{4}{11}\cdot(-\tfrac{1}{11}) + (-\tfrac{1}{11})\cdot\tfrac{3}{11}\Bigr)\cdot\tfrac{11}{2}.
\]
The dot product is $\tfrac{-4-3}{121}=-\tfrac{7}{121}$, so
$k_{12} = -\tfrac{7}{121}\cdot\tfrac{11}{2} = -\tfrac{7}{22}\approx -0.318$. No
quadrature was needed: on an affine cell the grad--grad integrand of $P_1$ elements
is constant, so ``integrate'' is just ``multiply by the volume''---the degenerate,
one-effective-point case of the quadrature of \cref{sec:quad}. This single number is
one entry of the element stiffness matrix that \cref{ch:assembly} scatters into the
global operator.
\end{workedexample}

\summarysection
A finite-element code does not build basis functions cell by cell. It defines a small
family of \emph{shape functions} once, on a canonical \emph{reference element}
$\hat K$, and reaches every physical cell $K$ by a geometric map $\vx=F_K(\hat\vx)$.
On the reference simplex the nodal ($\Hone$) shape functions are the barycentric
coordinates $\lambda_i$ (the $P_1$ hats) and their higher-order $P_p$ relatives, each
fixed by the cardinality property $\hat\phi_i(\hat\vx_j)=\delta_{ij}$ that makes a
coefficient equal a nodal value. The map's Jacobian $J$ carries calculus across:
values are unchanged, gradients transform by the inverse-transpose
$\Grad_\vx=J^{-\top}\Grad_{\hat\vx}$, and integrals pay a volume factor $|\det J|$.
The two integrals every solver needs collapse to a per-cell geometry factor---the
scalar $|\det J|$ for mass terms, the metric $G=J^{-1}J^{-\top}|\det J|$ for
grad--grad terms---which is exactly Palace's \code{geom\_factor\_build}. The remaining
integral over $\hat K$ is evaluated by \emph{quadrature}, a weighted sum
$\int_{\hat K}g\approx\sum_q w_q g(\hat\vx_q)$ whose \emph{exactness degree} must
cover the integrand's polynomial degree ($\ge 2p$ for a $P_p$ mass matrix), as the
worked mass-matrix example showed concretely: the one-point rule under-integrates a
quadratic, the three-point rule nails it. The vector spaces $\Hcurl$ and $\Hdiv$ use
the same reference frame with \emph{Piola} maps ($J^{-\top}$ covariant,
$J/\det J$ contravariant) that preserve the tangential and normal continuity built
into their DoFs. Everything collects into the element-local tensor over five axes:
\code{[(E,L)]} per-element DoFs, \code{[(E,P,C)]} values at quadrature points, and
\code{[(E,P,G)]} the geometry factor---the data structure the matrix-free pipeline of
\cref{ch:matrixfree} runs on, with $\mathsf{B}$ the tabulated basis and $\mathsf{D}$
the per-point geometry-and-material multiply.

\furtherreading
The reference-element construction, the Lagrange and higher-order shape functions, the
Jacobian transformation rules, and the quadrature theory of this chapter are developed
rigorously in Brenner and Scott~\citep{brennerscott2008}, the standard graduate text on
the mathematics of finite elements; their treatment of the affine-equivalence of
elements is the precise version of \cref{fig:refmap}. For the engineering view---tables
of shape functions and quadrature rules on triangles, tetrahedra, and hexahedra, with
the electromagnetics applications of this book---Jin~\citep{jin2014} is the natural
companion, and it carries the N\'ed\'elec and Raviart--Thomas vector elements of
\cref{sec:vector} with their Piola maps in full. The matrix-free realization---the
$\mathsf{B}$/$\mathsf{D}$ tensor-contraction pipeline of \cref{sec:tensor}, sum
factorization on tensor-product elements, and the element-local tensor axes---is the
subject of the libCEED project and its reference~\citep{ceed2021}; \cref{ch:matrixfree}
takes it up in detail. The Piola transforms and the way the three maps commute with
$\Grad,\Curl,\Div$ are the bridge to the de Rham complex of \cref{ch:derham}.

\exercisessection
\begin{enumerate}[nosep]
\item Write out the six $P_2$ shape functions on the reference triangle in
  barycentric form ($\lambda_i(2\lambda_i-1)$ for the three vertices and
  $4\lambda_i\lambda_j$ for the three edge midpoints) and verify the cardinality
  property \eqref{eq:cardinality} for one vertex function and one edge function by
  evaluating it at all six nodes (the three vertices and the three edge midpoints).
\item How many shape functions does the $P_3$ space on a triangle have? On a
  tetrahedron the count is $\binom{p+3}{3}$; evaluate it for $p=1,2,3$ and identify
  which DoFs (vertex, edge, face, interior) the new ones at each order sit on.
\item For the reference mass matrix of the $P_1$ triangle, use the barycentric
  integration formula of \cref{wex:massmatrix} to compute every entry $m_{ij}$, and
  show the result is $\tfrac{1}{24}\begin{psmallmatrix}2&1&1\\1&2&1\\1&1&2\end{psmallmatrix}$
  (diagonal $\tfrac{1}{12}$, off-diagonal $\tfrac{1}{24}$). Confirm the row sums equal
  the integrals $\int_{\hat K}\hat\phi_i\,d\hat\vx=\tfrac16$.
\item Repeat the three-way quadrature comparison of \cref{wex:massmatrix} for the
  \emph{diagonal} entry $m_{00}=\int_{\hat K}\hat\phi_0^2\,d\hat\vx$. Compute the exact
  value, the one-point estimate, and the three-point estimate. Does the one-point rule
  do better or worse on the diagonal than on the off-diagonal? Why is the three-point
  rule still exact?
\item Map the reference triangle onto the physical triangle with vertices
  $\vx_0=(0,0)$, $\vx_1=(2,0)$, $\vx_2=(1,3)$. Compute $J$, $\det J$, the area of $K$,
  and the physical gradient $\Grad_\vx\phi_0$ of the shape function
  $\hat\phi_0=1-\hat x-\hat y$. Verify that $\Grad_\vx\phi_0+\Grad_\vx\phi_1+\Grad_\vx\phi_2=\mathbf 0$.
\item Explain why the volume factor for a mass integral is $|\det J|$ but the gradient
  transform involves $J^{-\top}$. Using \cref{wex:jacobian}'s $J$, show by direct
  substitution that $\int_K\Grad\phi_i\cdot\Grad\phi_j\,d\vx$ pulls back to
  $\int_{\hat K}\Grad_{\hat\vx}\hat\phi_i\cdot(J^{-1}J^{-\top}|\det J|)\Grad_{\hat\vx}\hat\phi_j\,d\hat\vx$,
  and identify the metric matrix $G=J^{-1}J^{-\top}|\det J|$ for that cell.
\item A $Q$-point Gauss rule on an interval is exact to degree $2Q-1$. What is the
  minimum number of one-dimensional Gauss points needed to integrate a $P_4$ mass
  matrix exactly (integrand degree $2p=8$) on an affine element? What if the map is
  curved, so that $|\det J|$ is itself a degree-$2$ polynomial---what integrand degree
  must the rule now cover?
\item Identify, for each of the five element-local axes $E,L,P,C,G$ of
  \cref{sec:tensor}, which earlier chapter or section fixes its extent: which is set
  by the mesh, which by the FE order, which by the quadrature rule, which by the
  weak-form differential operator, and which by the space dimension. Then state which
  axes are \emph{build-stratum} (fixed once per mesh and order) and which is
  \emph{run-stratum} (a per-apply field value), connecting to the
  \code{geom\_factor\_build} setup pass.
\end{enumerate}

\chapter{Assembly as a Fold}
\label{ch:assembly}
\label{ch:fem}% alias: "the finite-element discretization/assembly" is cited as ch:fem in later chapters

\chapterorient{Chapter~\ref{ch:galerkin} told us \emph{what} the global matrix
is---the bilinear form sampled on basis pairs---and Chapter~\ref{ch:refelement}
told us how to compute one element's worth of it on a reference cell. This
chapter is about the \emph{seam} between the two: how the small, dense
element matrices, each living in its own local coordinate system, are stitched
into the one big sparse global matrix the solver actually sees. The mechanism
is a pair of indexing operations---\emph{gather} and \emph{scatter}---and the
moment we name them correctly, a structural fact about the whole simulator
comes into focus. Assembly is not a tangle of nested loops; it is a
\emph{fold}: a list of contributions combined, one at a time, by a single
associative operation into a single result. That sentence is the seed of an
idea this book spends Part~III harvesting---that the simulator's deep structure
is folds over lists. By the end you will have assembled the 1-D Poisson matrix
by hand a second time, but now the \emph{right} way, watching shared nodes
collect contributions from both their neighbours, and you will have proved that
the order in which terms and elements are accumulated never matters.}

\section{The gap between local and global}

We left Chapter~\ref{ch:galerkin} with a clean formula and a quiet omission. The
formula was $\mat{K}_{ij}=a(\phi_j,\phi_i)$: every entry of the global stiffness
matrix is the bilinear form evaluated on a pair of global basis functions. The
omission was \emph{how you actually compute it}. Written as one integral over
the whole domain $\dom$, the entry
\begin{equation}
\mat{K}_{ij}=\int_\dom Q\,\diffop\phi_j\cdot\diffop\phi_i\dV
\label{eq:global-entry}
\end{equation}
is honest but useless as a recipe: the global basis functions $\phi_i$ are
awkward piecewise objects, defined cell by cell, and nobody integrates them as
single functions over the entire mesh. Real finite-element codes never form
\cref{eq:global-entry} directly. They do something that looks, at first, like a
detour: they compute small \emph{element matrices}, one per mesh cell, each
holding only the interactions among the handful of basis functions that live on
that one cell, and then they \emph{add} those small matrices into the big one.
This chapter is about why that detour is in fact the royal road, and about the
algebraic shape it gives to the whole assembly stage.

The key observation is that the domain integral \cref{eq:global-entry} splits
across cells. The mesh $\mathcal{T}=\{K_1,\dots,K_{n_e}\}$ partitions $\dom$ into
$n_e$ non-overlapping elements, so any integral over $\dom$ is the sum of
integrals over the elements:
\begin{equation}
\mat{K}_{ij}=\int_\dom Q\,\diffop\phi_j\cdot\diffop\phi_i\dV
=\sum_{K\in\mathcal{T}}\;\int_K Q\,\diffop\phi_j\cdot\diffop\phi_i\dV.
\label{eq:cellsplit}
\end{equation}
Already \cref{eq:cellsplit} has a summation in it---our first sum over elements,
and the structural heart of the chapter. But most of its terms are zero. On a
given cell $K$, the only basis functions that are nonzero are the few whose
support touches $K$; for the hat basis of \cref{ch:galerkin} that is exactly the
basis functions sitting on the \emph{vertices of $K$}. If either $\phi_i$ or
$\phi_j$ vanishes on all of $K$, the cell's contribution to $\mat{K}_{ij}$ is
zero. So each cell contributes to only a tiny block of the global matrix: the
rows and columns indexed by \emph{its own} degrees of freedom.

\begin{keyidea}
Assembly inverts the two sums. \Cref{eq:cellsplit} loops over the global entries
$(i,j)$ and, for each, sums over cells---most contributing nothing. A
finite-element code loops the other way: \emph{over cells} on the outside, and
for each cell computes the small dense block of interactions among only that
cell's degrees of freedom, then \emph{adds} that block into the global matrix at
the right rows and columns. The same number $\sum_K\int_K(\cdots)$ is computed,
but the loop that is cheap (over elements, with dense local work) is on the
outside, and the loop that is expensive (over the $N^2$ global entries, almost
all zero) never happens.
\end{keyidea}

To make ``adds that block in at the right rows and columns'' precise we need a
language for moving between the global degrees of freedom and an element's local
ones. That language is gather and scatter.

\section{Gather and scatter: the operator \texorpdfstring{$G$}{G} and its transpose}
\label{sec:gatherscatter}

Number the global degrees of freedom $1,\dots,N$; for the hat basis these are
the interior mesh nodes. Each element $K$ owns a short list of \emph{local}
degrees of freedom---for a 1-D linear element, its two endpoints; for a triangle,
its three vertices---and there is a fixed table, the \emph{local-to-global map},
that says which global number each local slot corresponds to. This table is the
only bookkeeping assembly needs, and it is built once from the mesh.

\begin{definition}[Element gather]
For an element $K$ with $L$ local degrees of freedom, the \emph{gather} (or
\emph{restriction}) operator $\mat{R}_K$ is the $L\times N$ matrix that copies the
global degrees of freedom belonging to $K$ into a length-$L$ local vector, in the
element's local order:
\[
(\mat{R}_K\vx)_{\,\ell}=\vx_{\,g(K,\ell)},\qquad \ell=1,\dots,L,
\]
where $g(K,\ell)$ is the global index of $K$'s $\ell$-th local degree of freedom.
Each row of $\mat{R}_K$ is a single $1$ in column $g(K,\ell)$ and zeros
elsewhere: $\mat{R}_K$ is a \emph{Boolean selection matrix}, carrying no
arithmetic.
\end{definition}

The gather \emph{reads}. Its transpose \emph{writes}, and writes in a particular,
crucial way.

\begin{definition}[Element scatter-add]
The \emph{scatter-add} (or \emph{assembly}) operator is the transpose
$\mat{R}_K^{\mathsf T}$, an $N\times L$ matrix. Applied to a length-$L$ local
vector $\vu$, it places each local entry into its global slot and
\emph{accumulates}:
\[
(\mat{R}_K^{\mathsf T}\vu)_{\,m}=\sum_{\ell\,:\,g(K,\ell)=m}\vu_{\,\ell}.
\]
For a single element each global slot receives at most one local entry, but when
the scatter-adds of \emph{many} elements are summed---which is exactly what
assembly does---a global degree of freedom shared by several elements receives a
contribution from each.
\end{definition}

\Cref{fig:gatherscatter} draws the pair. Gather pulls a copy of element $K$'s
nodal values out of the long global vector; element-local computation happens on
that short vector; scatter-add pushes the results back, summing into any slot two
elements share.

\begin{figure}[t]
\centering
\begin{tikzpicture}[scale=1.0, every node/.style={font=\footnotesize}]
  % ---- global vector (vertical stack of cells) ----
  \def\cw{0.62}
  \foreach \k/\lab in {0/$x_1$,1/$x_2$,2/$x_3$,3/$x_4$,4/$x_5$}{
    \draw[simInk] (0,-\k*\cw) rectangle (\cw,-\k*\cw-\cw);
    \node at (\cw/2,-\k*\cw-\cw/2) {\lab};
  }
  \node[above, font=\footnotesize\bfseries] at (\cw/2,0.05) {global $\vx$};
  % highlight the two dofs of element K2: x2, x3 (indices 1,2)
  \draw[simBlue, very thick] (0,-1*\cw) rectangle (\cw,-3*\cw);
  % ---- element-local vector for K2 ----
  \def\ex{4.4}
  \foreach \k/\lab in {0/$u_1^{K_2}$,1/$u_2^{K_2}$}{
    \draw[simBlue, fill=simBgBlue] (\ex,-\k*\cw-0.62) rectangle (\ex+\cw,-\k*\cw-0.62-\cw);
    \node at (\ex+\cw/2,-\k*\cw-0.62-\cw/2) {\lab};
  }
  \node[above, font=\footnotesize\bfseries, simBlue!70!simInk] at (\ex+\cw/2,-0.5) {element $K_2$};
  % gather arrow
  \draw[flow, simBlue] (\cw+0.1,-1.4*\cw) to[bend left=12]
     node[above, font=\scriptsize]{$\mat{R}_{K_2}$ (gather)} (\ex-0.1,-1.0);
  % element computation box
  \def\bx{4.2}
  \node[stage, minimum width=20mm] (comp) at (\ex+\cw/2,-3.4) {compute\\$\mat{K}^{\mathrm{elem}}_{K_2}\vu$};
  \draw[flow, simBlue] (\ex+\cw/2,-2.0) -- (comp.north);
  % scatter-add arrow back, into x2 and x3
  \draw[backflow] (comp.west) to[bend right=18]
     node[below, font=\scriptsize, align=center]{$\mat{R}_{K_2}^{\mathsf T}$\\(scatter-add)} (\cw+0.1,-2.6*\cw);
  % a SECOND element K1 also writing into x2 (the shared dof)
  \draw[backflow, simGold] (-1.5,-1.6*\cw) to[bend left=20]
     node[left, font=\scriptsize, align=center]{from\\$K_1$} (-0.05,-1.5*\cw);
  \node[simGold!70!simInk, font=\scriptsize, align=center] at (-1.85,-3.0*\cw)
    {$x_2$ shared:\\two contributions};
  % bracket the shared dof
  \draw[simGold, thick] (-0.06,-1*\cw) -- (-0.06,-2*\cw);
\end{tikzpicture}
\caption[Gather and scatter-add]{The gather/scatter picture, for element $K_2$ of
a 1-D mesh. \emph{Gather} ($\mat{R}_{K_2}$) copies the two global degrees of
freedom belonging to $K_2$---here $x_2$ and $x_3$---into a short element-local
vector. The element matrix acts on that local vector. \emph{Scatter-add}
($\mat{R}_{K_2}^{\mathsf T}$) writes the results back into the global slots,
\emph{accumulating} rather than overwriting. The global degree of freedom $x_2$
sits on the boundary between $K_1$ and $K_2$, so it receives one contribution
from each element (the gold arrow is $K_1$'s)---this superposition is what
\emph{assembly} means.}
\label{fig:gatherscatter}
\end{figure}

\subsection{Assembly multiplicity and the diagonal \texorpdfstring{$\mat{R}^{\mathsf T}\mat{R}$}{R'R}}

The fact that worries newcomers---a shared node being written by two
elements---is not a bug to be guarded against; it \emph{is} the discretization.
A basis function $\phi_i$ on a shared node is, by construction, supported on
\emph{both} adjacent cells, so the integral $\int_\dom(\cdots)$ that defines its
matrix row genuinely receives a piece from each cell (\cref{eq:cellsplit}). The
scatter-\emph{add} reproduces that sum exactly. Overwriting instead of adding
would silently drop half of each shared row.

There is a clean way to measure how much sharing a mesh has. Stack the gathers of
all elements and ask: what does ``gather everything, then immediately
scatter-add it all back'' do to the global vector? For one element,
$\mat{R}_K^{\mathsf T}\mat{R}_K$ is the diagonal matrix that is $1$ on every
global slot $K$ touches and $0$ elsewhere. Summing over all elements,
\begin{equation}
\sum_{K\in\mathcal{T}}\mat{R}_K^{\mathsf T}\mat{R}_K
=\operatorname{diag}(\,m_1,m_2,\dots,m_N\,),
\qquad
m_i=\#\{\,K:\text{ global dof }i\text{ touches }K\,\},
\label{eq:multiplicity}
\end{equation}
the diagonal matrix whose $i$-th entry $m_i$ is the \emph{multiplicity} of global
degree of freedom $i$---the number of elements sharing it. For an interior 1-D
node, $m_i=2$; for a boundary node, $m_i=1$; for an interior vertex of a 2-D
triangular mesh, $m_i$ is the number of triangles in the fan around it.

\begin{notationbox}
The single-element operators are written $\mat{R}_K$ throughout, matching the
\code{element\_restrict} operator of the synthesis surface: $\mat{R}_K$ is the
``gather'' $G$ and $\mat{R}_K^{\mathsf T}$ is the ``scatter-add'' $G^{\mathsf T}$.
The diagonal \cref{eq:multiplicity} is the \emph{dof-multiplicity diagonal}
$G^{\mathsf T}G$. Note carefully: $\mat{R}_K^{\mathsf T}\mat{R}_K$ is a diagonal of
ones and zeros (a single element selects each of its dofs once), but the
\emph{sum over elements} \cref{eq:multiplicity} counts multiplicities. The
companion product $G\,G^{\mathsf T}$ on the element-local side is \emph{not} the
identity---it averages a shared value and redistributes it---a fact worth keeping
straight when you later meet it in the matrix-free apply of \cref{ch:matrixfree}.
\end{notationbox}

\section{Element matrices to global matrix}
\label{sec:elem-to-global}

With gather and scatter named, the assembly formula writes itself. On element
$K$, compute the $L\times L$ \emph{element matrix}
\begin{equation}
\bigl(\mat{K}^{\mathrm{elem}}_K\bigr)_{\ell\ell'}
=\int_K Q\,\diffop\psi^K_{\ell'}\cdot\diffop\psi^K_{\ell}\dV,
\label{eq:elemmat}
\end{equation}
where $\psi^K_1,\dots,\psi^K_L$ are the \emph{local} basis functions on $K$ (the
restrictions to $K$ of the global hats touching it, expressed in $K$'s own
coordinates---this is the reference-element computation of
\cref{ch:refelement}). The element matrix is small and dense: $2\times2$ for a
1-D linear element, $3\times3$ for a linear triangle, a few hundred entries even
for high-order 3-D elements. Now scatter each element matrix into the global one
and sum:
\begin{equation}
\boxed{\;\mat{K}=\sum_{K\in\mathcal{T}}\mat{R}_K^{\mathsf T}\,
\mat{K}^{\mathrm{elem}}_K\,\mat{R}_K\;}
\label{eq:assembly}
\end{equation}
The sandwich $\mat{R}_K^{\mathsf T}(\cdots)\mat{R}_K$ does exactly what its two
factors say: $\mat{R}_K$ on the right gathers the global indices into local order
for the columns, $\mat{R}_K^{\mathsf T}$ on the left scatters the local rows back
to global order, and the surrounding sum accumulates the overlapping blocks. That
\cref{eq:assembly} equals the entrywise formula \cref{eq:cellsplit} is immediate:
the $(i,j)$ entry of $\mat{R}_K^{\mathsf T}\mat{K}^{\mathrm{elem}}_K\mat{R}_K$ is
$(\mat{K}^{\mathrm{elem}}_K)_{\ell\ell'}$ when $i=g(K,\ell)$ and $j=g(K,\ell')$
for local slots $\ell,\ell'$ of $K$, and zero otherwise---precisely the cell-$K$
term of \cref{eq:cellsplit}, written in local coordinates.

\begin{figure}[t]
\centering
\begin{tikzpicture}[scale=0.66, every node/.style={font=\scriptsize}]
  % global 5x5 grid
  \def\N{5}
  \foreach \r in {0,...,4}{\foreach \c in {0,...,4}{
      \draw[simGray!35] (\c,-\r) rectangle ++(0.96,-0.96);}}
  \draw[simInk, thick] (0,0) rectangle (\N,-\N);
  % element K1 block: dofs {0,1} -> overlapping diagonal block rows/cols 0,1
  \fill[simBlue!35] (0,0) rectangle ++(1.96,-1.96);
  % element K2 block: dofs {1,2}
  \fill[simRust!35] (1,-1) rectangle ++(1.96,-1.96);
  % element K3 block: dofs {2,3}
  \fill[simGreen!35] (2,-2) rectangle ++(1.96,-1.96);
  % element K4 block: dofs {3,4}
  \fill[simViolet!35] (3,-3) rectangle ++(1.96,-1.96);
  % the OVERLAP cells (shared dofs) get two colors stacked: mark with hatch
  \foreach \d in {1,2,3}{
    \draw[simInk, thick] (\d,-\d) rectangle ++(0.96,-0.96);
    \node[font=\tiny\bfseries] at (\d+0.48,-\d-0.48) {$+$};
  }
  % labels
  \node[simBlue!70!simInk] at (0.96,0.45) {$K_1$};
  \node[simRust!70!simInk] at (3.4,-1.0) {$K_2$};
  \node[simGreen!70!simInk] at (4.5,-2.4) {$K_3$};
  \node[simViolet!70!simInk] at (4.7,-3.7) {$K_4$};
  \node[anchor=west, align=left, text width=30mm] at (5.5,-1.2)
    {Each element's $2\times2$ block lands on a $2\times2$ diagonal patch.};
  \node[anchor=west, align=left, text width=30mm] at (5.5,-3.1)
    {Shared dofs (marked $+$) sit where two blocks \emph{overlap}: the entries \emph{add}.};
\end{tikzpicture}
\caption[Scatter-add building the global matrix]{Assembly of a 1-D mesh with four
elements $K_1,\dots,K_4$ on five nodes (essential boundaries dropped for
clarity). Each element's small $2\times2$ matrix scatters onto the $2\times2$
diagonal patch indexed by its two nodes. Consecutive elements share one node, so
consecutive patches \emph{overlap} on a single diagonal entry (marked $+$): there
the two element contributions are \emph{added}. The overlapping-blocks-on-the-
diagonal picture is the geometric content of $\mat{K}=\sum_K\mat{R}_K^{\mathsf T}
\mat{K}^{\mathrm{elem}}_K\mat{R}_K$, and it is why the assembled matrix comes out
banded.}
\label{fig:scatteradd}
\end{figure}

\begin{workedexample}{Assembling 1-D Poisson by gather/scatter}{assemble-1d}
\label{wex:assemble-1d}
We re-derive the tridiagonal stiffness of \cref{ch:galerkin}, but this time the
\emph{assembly} way---element by element, scatter-adding---so the multiplicity on
shared nodes is explicit. Take the mesh of $n=4$ cells on $[0,1]$, $h=\tfrac14$,
with nodes $x_0,\dots,x_4$. The interior nodes $x_1,x_2,x_3$ carry the global
degrees of freedom $1,2,3$; the boundary nodes carry none. The four elements and
the global indices of their two local slots are
\[
\begin{array}{c|c|c}
\text{element} & \text{local slot }1 & \text{local slot }2\\\hline
K_1=[x_0,x_1] & (\text{bdry}) & 1\\
K_2=[x_1,x_2] & 1 & 2\\
K_3=[x_2,x_3] & 2 & 3\\
K_4=[x_3,x_4] & 3 & (\text{bdry})
\end{array}
\]

\emph{The element matrix.} Every linear element here is identical. With local
hats of slope $\pm1/h$ on a cell of length $h$, the $2\times2$ element stiffness
$\bigl(\mat{K}^{\mathrm{elem}}\bigr)_{\ell\ell'}=\int_K(\psi_{\ell'})'(\psi_\ell)'$
is
\[
\mat{K}^{\mathrm{elem}}=\frac1h\begin{pmatrix}\phantom{-}1 & -1\\ -1 & \phantom{-}1\end{pmatrix}
=4\begin{pmatrix}\phantom{-}1 & -1\\ -1 & \phantom{-}1\end{pmatrix}\quad(h=\tfrac14).
\]
The diagonal $+1/h$ is one node's $(\text{slope})^2\cdot h=(1/h)^2 h$; the
off-diagonal $-1/h$ is the two opposing slopes $(-1/h)(+1/h)\cdot h$.

\emph{Scatter, element by element.} Start with $\mat{K}=\mathbf 0$ (the
empty-mesh operator). Now add each $\mat{R}_K^{\mathsf T}\mat{K}^{\mathrm{elem}}
\mat{R}_K$, dropping rows/columns that scatter to a boundary node.
\begin{itemize}[nosep]
\item $K_1$ touches only global dof $1$ (its other slot is boundary). It
contributes $+\tfrac1h$ to $\mat{K}_{11}$.
\item $K_2$ touches dofs $1,2$. It contributes $\tfrac1h\!\begin{psmallmatrix}1&-1\\-1&1\end{psmallmatrix}$
to the $\{1,2\}$ block: $+\tfrac1h$ to each of $\mat{K}_{11},\mat{K}_{22}$ and
$-\tfrac1h$ to $\mat{K}_{12},\mat{K}_{21}$.
\item $K_3$ touches dofs $2,3$: $+\tfrac1h$ to $\mat{K}_{22},\mat{K}_{33}$,
$-\tfrac1h$ to $\mat{K}_{23},\mat{K}_{32}$.
\item $K_4$ touches only dof $3$: $+\tfrac1h$ to $\mat{K}_{33}$.
\end{itemize}
Now read the diagonal as a running tally. Global dof $1$ collected $+\tfrac1h$
from $K_1$ \emph{and} $+\tfrac1h$ from $K_2$: total $\tfrac2h$. Dof $2$ collected
$+\tfrac1h$ from $K_2$ and $+\tfrac1h$ from $K_3$: total $\tfrac2h$. Dof $3$,
symmetrically, $\tfrac2h$. \emph{Each diagonal entry is $\tfrac2h$ because each
interior node is shared by exactly two elements}---multiplicity $m_i=2$, exactly
\cref{eq:multiplicity}. The off-diagonals each got a single $-\tfrac1h$ from the
one element straddling that pair. Assembling,
\[
\mat{K}=\frac1h
\begin{pmatrix} 2 & -1 & 0\\ -1 & 2 & -1\\ 0 & -1 & 2\end{pmatrix},
\]
the tridiagonal matrix \cref{eq:tridiag}, now \emph{earned} as a superposition of
four $2\times2$ blocks rather than read off a stencil. The ``$2$'' on the
diagonal is not a coincidence of the second-derivative formula---it is two
elements each depositing a $1$.
\end{workedexample}

\subsection{Why this is the right factoring}

Equation~\cref{eq:assembly} cleanly separates three concerns that a single global
integral would tangle. The \emph{geometry and physics} live entirely inside
$\mat{K}^{\mathrm{elem}}_K$---the material weight $Q$, the element's shape, the
choice of basis. The \emph{connectivity} lives entirely inside the Boolean
$\mat{R}_K$---which local slot is which global node. And the \emph{combination}
is a plain sum. Each can be changed without touching the others: refine the mesh
and only the $\mat{R}_K$ table grows; change the material and only the element
matrices change; add a physical effect and---as the next section shows---only the
list of summands lengthens. This separation is what lets one assembly routine
serve every analysis in the book.

\begin{figure}[t]
\centering
\begin{tikzpicture}[scale=1.0, every node/.style={font=\footnotesize}]
  % two triangles sharing edge 2-3
  \coordinate (n1) at (0,1.4);
  \coordinate (n2) at (1.4,2.0);
  \coordinate (n3) at (1.0,0.0);
  \coordinate (n4) at (2.6,0.9);
  \filldraw[simBgBlue, draw=simBlue, thick] (n1)--(n2)--(n3)--cycle;
  \filldraw[simBgRust, draw=simRust, thick, fill opacity=0.6] (n2)--(n3)--(n4)--cycle;
  \draw[simGold, very thick] (n2)--(n3);
  \foreach \p/\lab/\pos in {n1/1/left,n2/2/above,n3/3/below,n4/4/right}{
    \fill[simInk] (\p) circle (2pt);
    \node[\pos=2pt] at (\p) {$\lab$};}
  \node[simBlue!70!simInk] at (0.8,1.25) {$K_a$};
  \node[simRust!70!simInk] at (1.8,0.95) {$K_b$};
  \node[simGold!75!simInk, font=\scriptsize, right] at (1.5,1.0) {shared edge $2$--$3$};
\end{tikzpicture}
\caption[Two triangles sharing an edge]{The two-triangle mesh of
\cref{wex:assemble-2d}. Triangles $K_a=\{1,2,3\}$ and $K_b=\{2,3,4\}$ share the
edge joining nodes $2$ and $3$ (gold). The shared nodes $2,3$ have multiplicity
$2$; the unshared nodes $1,4$ have multiplicity $1$. Each triangle's $3\times3$
element matrix scatters onto its three nodes, and the two scatters overlap on the
$2\times2$ block indexed by $\{2,3\}$---the discrete image of the shared edge.}
\label{fig:twotri}
\end{figure}

\begin{workedexample}{Two triangles, a shared edge, multiplicity in 2-D}{assemble-2d}
\label{wex:assemble-2d}
The 1-D example made every interior node have multiplicity exactly $2$. To see
the gather/scatter mechanism cope with richer connectivity---and to watch a shared
\emph{edge}, not just a shared node---we assemble on the smallest interesting 2-D
mesh: two linear triangles glued along an edge (\cref{fig:twotri}). Label the four
nodes $1,2,3,4$; triangle $K_a$ has vertices $\{1,2,3\}$ and triangle $K_b$ has
vertices $\{2,3,4\}$. Nodes $2$ and $3$ are shared (they lie on the common edge);
nodes $1$ and $4$ belong to one triangle each.

\emph{The gather tables.} In local vertex order,
\[
\mat{R}_{K_a}:\ (\ell_1,\ell_2,\ell_3)\mapsto(1,2,3),
\qquad
\mat{R}_{K_b}:\ (\ell_1,\ell_2,\ell_3)\mapsto(2,3,4).
\]
Write each triangle's $3\times3$ element stiffness abstractly as
$\mat{K}^{\mathrm{elem}}_{K}=\bigl(s^K_{pq}\bigr)_{p,q=1}^3$ (its actual entries
come from the reference-triangle computation of \cref{ch:refelement}; we need only
\emph{where} they land). Scatter-adding both into the $4\times4$ global matrix:
\[
\mat{K}=\mat{R}_{K_a}^{\mathsf T}\mat{K}^{\mathrm{elem}}_{K_a}\mat{R}_{K_a}
       +\mat{R}_{K_b}^{\mathsf T}\mat{K}^{\mathrm{elem}}_{K_b}\mat{R}_{K_b}
=\begin{pmatrix}
s^a_{11} & s^a_{12} & s^a_{13} & 0\\[1pt]
s^a_{21} & s^a_{22}\!+\!s^b_{11} & s^a_{23}\!+\!s^b_{12} & s^b_{13}\\[1pt]
s^a_{31} & s^a_{32}\!+\!s^b_{21} & s^a_{33}\!+\!s^b_{22} & s^b_{23}\\[1pt]
0 & s^b_{31} & s^b_{32} & s^b_{33}
\end{pmatrix}.
\]
Read the four corners of the sharing off the matrix. The $(1,4)$ and $(4,1)$
entries are \emph{zero}: nodes $1$ and $4$ never share a triangle, so no element
deposits there---exactly the sparsity rule of \cref{sec:cost}. The
$\{2,3\}\times\{2,3\}$ block is where the two scatters \emph{overlap}, and every
entry of it is a \emph{sum} $s^a_{\cdot\cdot}+s^b_{\cdot\cdot}$ of a contribution
from each triangle: the discrete image of the shared edge. The diagonal entries
$\mat{K}_{22}=s^a_{22}+s^b_{11}$ and $\mat{K}_{33}=s^a_{33}+s^b_{22}$ are two-term
sums (multiplicity $2$), while $\mat{K}_{11}=s^a_{11}$ and $\mat{K}_{44}=s^b_{33}$
are single terms (multiplicity $1$)---precisely
$\sum_K\mat{R}_K^{\mathsf T}\mat{R}_K=\operatorname{diag}(1,2,2,1)$ of
\cref{eq:multiplicity}. Nothing about the mechanism changed from 1-D; only the
element matrices grew from $2\times2$ to $3\times3$ and the gather tables grew a
column. The fold is identical.
\end{workedexample}

\section{Assembly is a fold over weak-form terms}
\label{sec:fold}

So far we have assembled \emph{one} operator from \emph{one} bilinear form, summed
over elements. Real analyses ask for more. The driven Maxwell operator is the sum
of three forms, $\mat{A}(\omega)=\mat{K}+i\omega\mat{C}-\omega^2\mat{M}$; a lossy
or multi-material problem adds still more terms; the mass matrix and the stiffness
matrix are two \emph{terms} an eigenmode analysis needs side by side. The global
operator an analysis assembles is a sum over a \emph{list of weak-form terms},
and each term's own assembly is, as we have just seen, a sum over elements. Two
sums, nested---and both are folds.

\subsection{The term list and its per-term operator}

Recall from \cref{sec:family} that every term of the bilinear-form family has the
shape $a_t(u,v)=(Q_t\,\diffop_t u,\diffop_t v)$: a material weight $Q_t$ and a
differential operator $\diffop_t\in\{\Grad,\,\mathrm{id},\,\Curl,\,\Div\}$.
Sampling one term on basis pairs and assembling it over the mesh
(\cref{eq:assembly}) produces one global matrix; call it $\mat{A}_t$, the
\emph{per-term operator}. An analysis hands the assembler a list of terms
$[\,t_1,t_2,\dots,t_m\,]$ and wants the matrix of the total form $a=\sum_t a_t$.
Because sampling and assembly are linear in the form,
\begin{equation}
\mat{K}=\sum_{t}\mat{A}_t,
\qquad
\mat{A}_t=\sum_{K\in\mathcal{T}}\mat{R}_K^{\mathsf T}\,
\mat{A}^{\mathrm{elem}}_{t,K}\,\mat{R}_K .
\label{eq:termsum}
\end{equation}
The outer sum is over \emph{terms}; the inner sum, hidden inside each
$\mat{A}_t$, is over \emph{elements}. \Cref{fig:termfold} draws the outer one:
each weak-form term enters its own assembly and the per-term operators are summed
into the global operator.

\begin{figure}[t]
\centering
\begin{tikzpicture}[node distance=5mm and 12mm, every node/.style={font=\footnotesize}]
  \node[opbox] (t1) {$t_1=(\varepsilon,\Grad)$};
  \node[opbox, below=of t1] (t2) {$t_2=(\sigma,\mathrm{id})$};
  \node[opbox, below=of t2] (t3) {$t_3=(\nu,\Curl)$};
  \node[stage, right=16mm of t1] (a1) {assemble};
  \node[stage, right=16mm of t2] (a2) {assemble};
  \node[stage, right=16mm of t3] (a3) {assemble};
  \draw[flow] (t1) -- (a1);
  \draw[flow] (t2) -- (a2);
  \draw[flow] (t3) -- (a3);
  \node[right=4mm of a1] (m1) {$\mat{A}_1$};
  \node[right=4mm of a2] (m2) {$\mat{A}_2$};
  \node[right=4mm of a3] (m3) {$\mat{A}_3$};
  \draw[flow] (a1)--(m1);
  \draw[flow] (a2)--(m2);
  \draw[flow] (a3)--(m3);
  \node[product, right=14mm of m2] (K) {$\mat{K}=\displaystyle\sum_t\mat{A}_t$};
  \draw[flow] (m1) -- (K);
  \draw[flow] (m2) -- (K);
  \draw[flow] (m3) -- (K);
  \node[font=\scriptsize\itshape, simGray, above=2mm of t1, xshift=6mm]
    {list of weak-form terms};
\end{tikzpicture}
\caption[The term fold $\mat{K}=\sum_t\mat{A}_t$]{The outer sum of
\cref{eq:termsum}. An analysis specifies a list of weak-form terms; each term
$t=(Q,\diffop)$ assembles (via \cref{eq:assembly}, an inner sum over elements)
into its own global operator $\mat{A}_t$; the global operator of the analysis is
the sum $\mat{K}=\sum_t\mat{A}_t$. Adding a physical effect---a loss term, a
second material, the $i\omega\mat{C}$ damping---means appending one term to the
list, never rebuilding from scratch.}
\label{fig:termfold}
\end{figure}

\subsection{Naming the pattern: \texttt{foldr}}

\Cref{eq:termsum} is not merely a formula; it is a \emph{computational pattern}
with a name. We have a list, a way to turn each list element into a matrix, and a
way to combine matrices---add them, starting from the zero matrix. Walking a list
and combining its elements into a single accumulated result is the job of the
higher-order function \texttt{fold}. In the pseudo-language of Part~III, the
assembly operator \code{fe\_assemble} is written exactly this way:
\begin{lstlisting}[style=l4style]
fe_assemble :: (space: FiniteElementSpace, terms: [WeakFormTerm]) -> LinOp
fe_assemble space terms = foldr (\t acc -> assembleTerm space t + acc) zero terms
                        -- = sum over t in terms of  assembleTerm space t
\end{lstlisting}
Read the body left to right. \code{foldr} walks the list \code{terms} from the
right; for each term \code{t} it computes that term's operator
\code{assembleTerm\ space\ t} and \emph{adds} it to the accumulator \code{acc};
the walk starts from \code{zero}, the zero operator. Unfolded, the call is
\[
\texttt{assembleTerm}\,t_1+\bigl(\texttt{assembleTerm}\,t_2+\cdots
+(\texttt{assembleTerm}\,t_m+\mathbf 0)\bigr)
=\sum_t\mat{A}_t,
\]
which is \cref{eq:termsum} exactly. The single per-term map
\code{assembleTerm\ space\ t}---the realization of one $(Q,\diffop)$ term as a
global matrix---is itself the inner element-fold \cref{eq:assembly}; we treat it
here as an opaque leaf and fold over it. The point is the \emph{outer} structure:
assembly of an analysis is one fold over the term list.

\begin{keyidea}
\textbf{Assembly is a commutative-monoid fold over the list of weak-form terms.}
The ingredients are a combine operation (operator addition $+$), an identity for
it (the zero operator $\mathbf 0$), and a list (the weak-form terms). Operator
addition is associative and commutative and has $\mathbf 0$ as identity---it forms
a \emph{commutative monoid}---and folding a list with a commutative-monoid
operation is the cleanest combinator pattern there is. Three consequences follow
purely from this algebra, with no reference to any code: assembling no terms
yields the zero operator; assembling a concatenated list equals assembling the
parts and adding; and the order of the terms never matters. We prove all three in
\cref{sec:laws}. This is the reader's first deep encounter with a thesis Part~III
makes central---\emph{the simulator's structure is folds over lists}
(\cref{ch:fold}, \cref{ch:combinators}).
\end{keyidea}

\section{The three laws of the assembly fold}
\label{sec:laws}

The commutative-monoid structure is not decoration; it is a set of guarantees an
engineer relies on. We state them as a theorem and prove them from the algebra of
operator addition alone---so they hold for \emph{every} analysis, regardless of
which terms are in the list or how the per-term operators are computed.

\begin{theorem}[Algebraic laws of assembly]
\label{thm:foldlaws}
Let $\mathrm{asm}(\mathit{terms})=\sum_{t\in\mathit{terms}}\mat{A}_t$ be the
assembled global operator on a fixed finite-element space, where $\mat{A}_t$ is
the per-term operator and $+$ is operator addition with identity the zero operator
$\mathbf 0$. Then:
\begin{enumerate}[label=\textup{(\roman*)},nosep]
\item \emph{(Empty-list identity)} $\mathrm{asm}([\,])=\mathbf 0$.
\item \emph{(Concatenation homomorphism)} for any two term lists $T_1,T_2$,
\[
\mathrm{asm}(T_1\!+\!\!+\,T_2)=\mathrm{asm}(T_1)+\mathrm{asm}(T_2),
\]
where $+\!\!+$ is list concatenation.
\item \emph{(Order independence)} $\mathrm{asm}(\mathit{terms})$ is unchanged
under any permutation of $\mathit{terms}$.
\end{enumerate}
\end{theorem}

\begin{proof}
All three are properties of the sum, inherited from operator addition.

\emph{(i)} The empty list has no summands, and the empty sum is by convention the
identity of the operation summed: $\sum_{t\in[\,]}\mat{A}_t=\mathbf 0$. (This is
the base case of the fold: \code{foldr\ f\ zero\ [\,]\ =\ zero}.)

\emph{(ii)} Concatenating two lists concatenates their summands, and a sum over a
concatenated index set splits:
\[
\mathrm{asm}(T_1\!+\!\!+\,T_2)
=\sum_{t\in T_1+\!\!+\,T_2}\mat{A}_t
=\sum_{t\in T_1}\mat{A}_t+\sum_{t\in T_2}\mat{A}_t
=\mathrm{asm}(T_1)+\mathrm{asm}(T_2).
\]
The middle equality is associativity of $+$ (regroup the summands into the two
sublists). A map from list-concatenation to a combine on results is exactly a
\emph{list homomorphism}; (ii) says assembly \emph{is} one.

\emph{(iii)} A permutation reorders the summands but changes neither the set of
summands nor their total, because $+$ is commutative and associative: any
reordering of $\mat{A}_{t_1}+\cdots+\mat{A}_{t_m}$ gives the same operator. The
list of terms is therefore effectively an unordered \emph{bag}.
\end{proof}

\begin{remark}
The three laws are not independent novelties; they are the universal property of
folding over a commutative monoid, specialized to (operators, $+$, $\mathbf 0$).
Any fold over any commutative monoid satisfies the analogues. We will meet the
same three laws again when we fold over \emph{elements} inside a single
$\mat{A}_t$ (the inner sum of \cref{eq:termsum}, also a commutative-monoid fold),
and yet again, abstracted away from finite elements entirely, in
\cref{ch:combinators}. Recognizing the pattern once buys it everywhere.
\end{remark}

\begin{workedexample}{The concatenation homomorphism, numerically}{concat-homo}
\label{wex:concat-homo}
Law (ii) of \cref{thm:foldlaws} says we may assemble a combined term list, or
assemble the pieces and add, and get the same operator. We check it on the
$3$-node 1-D mesh of \cref{wex:assemble-1d}, with a two-term list: a stiffness
term $t_K=(1,\Grad)$ and a mass term $t_M=(1,\mathrm{id})$.

\emph{The two per-term operators.} The stiffness operator is the
$\mat{K}=\tfrac1h\operatorname{tridiag}(-1,2,-1)$ of \cref{wex:assemble-1d}; with
$h=\tfrac14$, $\mat{K}=4\operatorname{tridiag}(-1,2,-1)$. The mass operator is the
$\mat{M}=\tfrac h6\operatorname{tridiag}(1,4,1)$ of \cref{wex:mass}; with
$h=\tfrac14$, $\mat{M}=\tfrac1{24}\operatorname{tridiag}(1,4,1)$. (Each is itself
an inner element-fold, assembled exactly as in \cref{wex:assemble-1d}.)

\emph{Route A: assemble the concatenated list.} Feed the assembler the single
list $[\,t_K,t_M\,]$. Its fold computes
$\mathrm{asm}([t_K,t_M])=\mat{A}_{t_K}+\mat{A}_{t_M}=\mat{K}+\mat{M}$:
\[
\mat{K}+\mat{M}
=4\!\begin{pmatrix}2&-1&0\\-1&2&-1\\0&-1&2\end{pmatrix}
+\frac1{24}\!\begin{pmatrix}4&1&0\\1&4&1\\0&1&4\end{pmatrix}
=\begin{pmatrix}
\tfrac{49}{6} & -\tfrac{95}{24} & 0\\[2pt]
-\tfrac{95}{24} & \tfrac{49}{6} & -\tfrac{95}{24}\\[2pt]
0 & -\tfrac{95}{24} & \tfrac{49}{6}
\end{pmatrix}.
\]
(Diagonal: $8+\tfrac{4}{24}=8+\tfrac16=\tfrac{49}{6}$. Off-diagonal:
$-4+\tfrac{1}{24}=-\tfrac{95}{24}$.)

\emph{Route B: assemble each, then add.} Call the assembler twice---once on
$[t_K]$, once on $[t_M]$---obtaining $\mat{K}$ and $\mat{M}$ separately, and add
the two matrices. The result is the very same $\mat{K}+\mat{M}$ above, entry for
entry. The two routes agree, as law (ii) promised: $\mathrm{asm}([t_K,t_M])
=\mathrm{asm}([t_K])+\mathrm{asm}([t_M])$.

\emph{Why it matters.} Route B is what an analysis actually does when it caches
the stiffness once and reuses it while sweeping a frequency-dependent combination
$\mat{K}-\omega^2\mat{M}$ (the driven analysis, \cref{ch:driven}): law (ii)
guarantees the cached-and-combined operator is bit-for-bit the from-scratch
assembly. It is also the law that makes assembly \emph{splittable across
processors}---hand half the term list (or half the elements) to each of two
machines, assemble independently, and add the results.
\end{workedexample}

\section{What is \emph{not} in the fold}
\label{sec:notinfold}

A fold's discipline is as much about what it excludes as what it includes. Two
operations that beginners often imagine are part of assembly are deliberately
\emph{outside} it.

\subsection{Boundary conditions are separable post-compositions}

The assembled operator $\mat{K}=\sum_t\mat{A}_t$ of \cref{eq:termsum} knows
nothing about boundary conditions. Essential (Dirichlet) conditions---pinning
certain degrees of freedom to prescribed values---are imposed \emph{after} the
fold, by acting on the already-assembled $\mat{K}$: the rows and columns of the
pinned dofs are eliminated, and any inhomogeneous data is lifted into the
right-hand side. Crucially, this elimination is a function of the assembled
operator alone; it does not care which terms were folded to build it, or in what
order. Assembly builds the operator; boundary-condition elimination is a
\emph{separable post-composition} on it, valid independently of the assembly. This
is why \cref{thm:foldlaws} can hold without any boundary-condition hypothesis: the
laws are about the fold, and the fold ends before the boundary conditions begin.
\Cref{ch:bc} takes up the elimination in full, including why keeping it separate
preserves symmetry and how inhomogeneous data moves to the load vector.

\begin{pitfall}
It is tempting to ``apply the boundary conditions as you assemble''---skipping the
pinned rows inside the element loop. This works, but it fuses two stages that are
cleaner apart, and it breaks the fold laws: an assembler that eliminates
mid-fold is no longer a homomorphism, because $\mathrm{asm}(T_1\!+\!\!+\,T_2)$
would eliminate twice. Keep the fold pure and post-compose the elimination once;
\cref{thm:foldlaws} is the reward. We assembled the \emph{interior} dofs in
\cref{wex:assemble-1d} precisely by dropping boundary slots, but notice we dropped
them from the \emph{connectivity} ($\mat{R}_K$), not by altering the fold's combine
step---the fold itself stayed a clean sum.
\end{pitfall}

\subsection{The matrix need not be formed at all}

The second exclusion is subtler and points forward to \cref{ch:matrixfree}.
Equation~\cref{eq:assembly} \emph{forms} the global matrix $\mat{K}$---it
materializes every nonzero entry into storage. But the only thing the iterative
solvers of \cref{part:framework} ever ask of $\mat{K}$ is its \emph{action} on a
vector, $\vw\mapsto\mat{K}\vw$. And the action distributes through both folds
without ever building the matrix:
\begin{equation}
\mat{K}\vw=\Bigl(\sum_t\sum_{K}\mat{R}_K^{\mathsf T}\mat{A}^{\mathrm{elem}}_{t,K}\mat{R}_K\Bigr)\vw
=\sum_t\sum_{K}\mat{R}_K^{\mathsf T}\bigl(\mat{A}^{\mathrm{elem}}_{t,K}(\mat{R}_K\vw)\bigr).
\label{eq:matrixfree-preview}
\end{equation}
Read the right-hand side as an algorithm: for each term and each element, gather
the element's slice of $\vw$ (apply $\mat{R}_K$), hit it with the small element
matrix, scatter-add the result (apply $\mat{R}_K^{\mathsf T}$); sum over all
elements and terms. No global matrix is ever assembled---only the small element
matrices (or, in libCEED, not even those) and the gather/scatter tables are kept.
This is the \emph{matrix-free} operator, and it is exactly the same fold, with the
combine step ``add into the global matrix'' replaced by ``add into the global
\emph{output vector}.''~\Cref{ch:matrixfree} develops it; \cref{eq:matrixfree-preview}
is the seed.

There is one more layer to peel, which we record now because it justifies the
notation $\mat{R}_K=G$ and points squarely at the production code. The element
matrix $\mat{A}^{\mathrm{elem}}_{t,K}$ is itself never formed in a high-performance
matrix-free kernel; it factors as
\begin{equation}
\mat{A}^{\mathrm{elem}}_{t,K}=B^{\mathsf T}_{\diffop}\,D\,B_{\diffop},
\label{eq:bdb}
\end{equation}
where $B_{\diffop}$ \emph{evaluates} the differential operator $\diffop$ of the
term at the element's quadrature points (the ``basis apply''), $D$ is a diagonal
of \emph{quadrature weights times the material coefficient $Q$} at those points,
and $B^{\mathsf T}_{\diffop}$ integrates back. Substituting \cref{eq:bdb} into
\cref{eq:matrixfree-preview}, one term's global action is the five-stage pipeline
\begin{equation}
\mat{A}_t=\;G^{\mathsf T}\;B^{\mathsf T}_{\diffop}\;D\;B_{\diffop}\;G,
\label{eq:gtbdbg}
\end{equation}
read right to left: gather to elements ($G$), evaluate at quadrature points
($B_{\diffop}$), scale by weighted material values ($D$), integrate back
($B^{\mathsf T}_{\diffop}$), scatter-add to global ($G^{\mathsf T}$). Only the
diagonal $D$ carries any floating-point ``physics''; $G$ and $B$ are fixed,
problem-independent index-and-interpolation maps. This factorization is the
contract of the libCEED library Palace assembles through, and its vocabulary---an
``$E$-vector'' is the element-local $G\vw$, an ``$L$-vector'' is the global $\vw$,
a ``$Q$-vector'' is the quadrature-point data---is precisely the gather/scatter
language of \cref{sec:gatherscatter}. We will not need \cref{eq:gtbdbg} again until
\cref{ch:matrixfree}, but it is worth seeing once that the entire matrix-free
apply is \emph{still the same two folds}, with the per-element work itself expanded
into a tiny pipeline.

\section{Sparsity, storage, and cost}
\label{sec:cost}

We close with the practical consequences of the fold structure: where the
sparsity comes from, what it costs to assemble, and where the parallelism lives.

\subsection{Why the assembled matrix is sparse}

The sparsity of $\mat{K}$ that \cref{ch:galerkin} attributed to the local support
of the basis is, from the assembly viewpoint, even more transparent. Each term
$\mat{R}_K^{\mathsf T}\mat{A}^{\mathrm{elem}}_{t,K}\mat{R}_K$ in \cref{eq:assembly}
is nonzero \emph{only} on the $L\times L$ block of rows and columns indexed by
element $K$'s degrees of freedom. A global entry $\mat{K}_{ij}$ can be nonzero
only if some single element $K$ owns \emph{both} dof $i$ and dof $j$---i.e.\ only
if nodes $i$ and $j$ are vertices of a common cell. Since each node belongs to a
bounded number of cells (two in 1-D, a small fan in 2-D/3-D), each row of
$\mat{K}$ has only a bounded number of nonzeros, independent of how large the mesh
grows. The matrix has $\mathcal{O}(N)$ nonzeros, not $N^2$. The sparsity
\emph{pattern} is read straight off the gather tables: $\mat{K}_{ij}\neq0$ iff
$i,j$ share an element. \Cref{fig:sparsitypattern} shows a small 2-D example.

\begin{figure}[t]
\centering
\begin{tikzpicture}[scale=1.0, every node/.style={font=\scriptsize}]
  % left: tiny 2D mesh (3x3 nodes, split into triangles)
  \begin{scope}[shift={(0,0)}]
    \foreach \x in {0,1,2}\foreach \y in {0,1,2}{
      \fill[simInk] (\x,\y) circle (1.4pt);}
    % grid lines + diagonals
    \foreach \y in {0,1,2}{\draw[simGray] (0,\y)--(2,\y);}
    \foreach \x in {0,1,2}{\draw[simGray] (\x,0)--(\x,2);}
    \foreach \x in {0,1}\foreach \y in {0,1}{\draw[simGray] (\x,\y+1)--(\x+1,\y);}
    % label center node 4 and highlight its star
    \node[simRust, below right, font=\tiny] at (1,1) {$4$};
    \fill[simRust] (1,1) circle (2.0pt);
    \node[font=\tiny] at (1,-0.5) {mesh: node $4$ couples to its neighbours};
  \end{scope}
  % right: the sparsity grid (9x9) with node 4's row marked
  \begin{scope}[shift={(4.6,-0.3)}, scale=0.34]
    \foreach \r in {0,...,8}\foreach \c in {0,...,8}{
      \draw[simGray!35] (\c,-\r) rectangle ++(0.94,-0.94);}
    % node numbering: bottom-left=0..8 row-major; node4=center couples to 1,3,4,5,7 (and diag 0,8 via triangles)
    \def\nbrs{1,3,4,5,7,0,8}
    \foreach \c in {1,3,4,5,7,0,8}{
      \fill[simBlue!45] (\c,-4) rectangle ++(0.94,-0.94);
      \fill[simBlue!45] (4,-\c) rectangle ++(0.94,-0.94);}
    \fill[simRust!60] (4,-4) rectangle ++(0.94,-0.94);
    \draw[simInk,thick] (0,0) rectangle (9,-9);
    \node[anchor=west, scale=2.2, align=left] at (9.6,-4.5) {row/col\\of node $4$};
  \end{scope}
\end{tikzpicture}
\caption[Sparsity from local support]{Sparsity of the assembled matrix on a small
$3\times3$-node triangular mesh. \emph{Left:} the centre node $4$ shares an
element with only its immediate mesh neighbours. \emph{Right:} the resulting
sparsity pattern. Row $4$ of $\mat{K}$ (and, by symmetry, column $4$) has a
nonzero only in the columns of those neighbours; all other entries of the row are
exactly zero, because no single element owns both node $4$ and a far node. Each
row has a bounded number of nonzeros no matter how large the mesh grows: $\mat{K}$
has $\mathcal{O}(N)$ nonzeros, and is stored as a sparse matrix---never as the
full $N\times N$ array.}
\label{fig:sparsitypattern}
\end{figure}

\subsection{The cost of assembly, and the scatter-add race}

The element-fold runs over $n_e$ elements; on each it computes a dense
$L\times L$ element matrix (a fixed amount of work set by the element type and
polynomial degree, not by $N$) and scatters it. Assembly therefore costs
\begin{equation}
\mathcal{O}\bigl(n_e\times(\text{per-element work})\bigr)
=\mathcal{O}(N)
\label{eq:asmcost}
\end{equation}
total work, linear in the problem size---one of the cheapest stages of the whole
pipeline. And because the element-fold is a commutative-monoid fold
(\cref{thm:foldlaws}(iii)), the elements may be processed in \emph{any} order,
and---more powerfully---in \emph{parallel}: order independence is exactly the
licence to split the list across processors. Assembly is, in this sense,
embarrassingly parallel over elements.

\begin{figure}[t]
\centering
\begin{tikzpicture}[scale=1.0, every node/.style={font=\footnotesize}]
  % four element "workers" computing in parallel
  \foreach \i/\x/\col in {1/0/simBlue,2/2.2/simRust,3/4.4/simGreen,4/6.6/simViolet}{
    \node[opbox, draw=\col, minimum width=14mm] (w\i) at (\x,1.5)
       {worker\\$K_\i$};
  }
  \node[font=\scriptsize\itshape, simGray] at (3.3,2.4) {parallel: element-local work is independent};
  % global vector slots
  \foreach \k/\lab in {0/1,1/2,2/3,3/4,4/5}{
    \draw[simInk] (1.6+\k*0.9,-0.6) rectangle ++(0.8,0.8);
    \node[font=\scriptsize] at (2.0+\k*0.9,-0.2) {$x_{\lab}$};
  }
  \node[anchor=east, font=\footnotesize\bfseries] at (1.5,-0.2) {global $\vx$:};
  % scatter-add arrows: K1->x1,x2 ; K2->x2,x3 (collision on x2) etc
  \draw[flow, simBlue] (w1) -- (2.0,0.25);
  \draw[flow, simBlue] (w1) -- (2.9,0.25);
  \draw[flow, simRust] (w2) -- (2.9,0.25);
  \draw[flow, simRust] (w2) -- (3.8,0.25);
  \draw[flow, simGreen] (w3) -- (3.8,0.25);
  \draw[flow, simGreen] (w3) -- (4.7,0.25);
  \draw[flow, simViolet] (w4) -- (4.7,0.25);
  \draw[flow, simViolet] (w4) -- (5.6,0.25);
  % collision markers on shared dofs x2,x3,x4
  \foreach \k in {2,3,4}{
    \node[circle,draw=simRust,fill=simBgRust,inner sep=1pt,font=\tiny]
      at (1.1+\k*0.9,0.55) {$\scriptstyle\rightleftarrows$};}
  \node[simRust!80!simInk, font=\scriptsize, align=left, anchor=west] at (6.2,0.5)
     {shared dof:\\two writers\\$\Rightarrow$ race};
\end{tikzpicture}
\caption[Parallel over elements, race on shared dofs]{Assembly parallelizes over
elements---each worker computes its element's contribution independently
(top)---but the \emph{scatter-add} into the global vector is where independence
breaks down. A degree of freedom shared by two elements has two workers trying to
add into the same slot at once (the $\rightleftarrows$ markers). Because the
underlying operation is addition, the \emph{result} is order-independent
(\cref{thm:foldlaws}(iii)); but two simultaneous read--modify--write updates can
interleave and lose one, a \emph{data race}. The race is resolved with atomic
adds, per-colour scheduling, or thread-private buffers summed at the end---an
implementation concern that does not change the assembled value.}
\label{fig:parallelrace}
\end{figure}

There is exactly one place where the parallelism is not free, and the fold laws
pinpoint it. The element-local \emph{computation} is perfectly independent across
elements; the \emph{scatter-add} is not, because two elements sharing a degree of
freedom write into the same global slot (\cref{fig:parallelrace}). The
\emph{mathematical} result is order-independent---that is precisely
\cref{thm:foldlaws}(iii)---but two threads performing
read--modify--write on the same memory location can interleave and lose an
update: a \emph{data race}. The fix is a matter of implementation, not
mathematics: make the additions atomic, colour the elements so that no two
sharing a dof run simultaneously, or give each thread a private accumulator and
sum them at the end (itself one more application of \cref{thm:foldlaws}(ii)).
Whichever is chosen, the assembled operator is the same: the race threatens the
\emph{schedule}, never the \emph{answer}, because the answer is a commutative-
monoid fold.

\begin{physicsbox}
The two parallel structures of assembly mirror the physics. Element-local work is
parallel because, in the continuum, the contribution of each little region of
space to the weak form is independent---fields do not interact across a region
boundary except through the shared boundary itself. The scatter-add ``race'' on
shared dofs is the discrete echo of that one exception: a node on the interface
between two elements feels both, and its equation must collect both
contributions. Superposition of local energies is the physics; scatter-add is its
arithmetic; the commutative-monoid fold is its structure.
\end{physicsbox}

\summarysection
Global assembly turns small, dense element matrices into the one large sparse
global operator the solver sees, by a pair of indexing operations. The
\emph{gather} $\mat{R}_K$ copies an element's global degrees of freedom into local
order; the \emph{scatter-add} $\mat{R}_K^{\mathsf T}$ writes element-local results
back, \emph{accumulating} where elements share a degree of freedom. The global
operator is the element sum $\mat{K}=\sum_K\mat{R}_K^{\mathsf
T}\mat{K}^{\mathrm{elem}}_K\mat{R}_K$, and the diagonal $\sum_K\mat{R}_K^{\mathsf
T}\mat{R}_K$ counts each dof's element-multiplicity. An analysis assembles a
\emph{list of weak-form terms}, each term assembling (an inner element-sum) into a
per-term operator $\mat{A}_t$, and the global operator is the outer sum
$\mat{K}=\sum_t\mat{A}_t$. Both sums are folds: assembly is a
\emph{commutative-monoid fold}, \code{fe\_assemble} $=$ \code{foldr (+) zero},
whose three laws---empty-list identity, concatenation homomorphism, and order
independence---follow from operator addition alone and underwrite caching,
splitting, and parallelism. Boundary-condition elimination is a \emph{separable
post-composition} on the assembled operator (\cref{ch:bc}), not part of the fold;
and the matrix need not be formed at all---the same fold, scattering into an
output vector instead of a matrix, is the matrix-free operator of
\cref{ch:matrixfree}. Sparsity follows from local support (a row is nonzero only
in the columns sharing an element), assembly costs $\mathcal{O}(N)$, and it is
embarrassingly parallel over elements save for the scatter-add race on shared
dofs---a race in the \emph{schedule}, never the \emph{answer}. That assembly is a
fold over a list is the reader's first deep meeting with the thread Part~III makes
the spine of the whole book.

\furtherreading
The element-by-element assembly procedure, local-to-global maps, and the sparse
storage of the resulting matrix are covered in every finite-element text; Brenner
and Scott~\citep{brennerscott2008} give the mathematical framing (the assembly
operator as a sum of restrictions) while Jin~\citep{jin2014} gives the engineering
practice for electromagnetics specifically, including the element-matrix
computations for the $\Hcurl$ and $\Hdiv$ bases we have only previewed. The
gather/scatter-and-element-matrix decomposition $A=G^{\mathsf
T}B^{\mathsf T}DBG$, and the matrix-free assembly that exploits it, are the
organizing idea of the libCEED project~\citep{ceed2021}, whose ``$E$-vector /
$L$-vector'' vocabulary is precisely the element-local versus global distinction
of \cref{sec:gatherscatter}; readers headed for \cref{ch:matrixfree} will find the
CEED interface a faithful realization of \cref{eq:matrixfree-preview}. The
recognition of assembly as a fold and a list homomorphism, and the parallel
decompositions it licenses, are developed abstractly in \cref{ch:combinators}.

\exercisessection
\begin{enumerate}[nosep]
  \item \textbf{Multiplicity by hand.} For the four-element 1-D mesh of
  \cref{wex:assemble-1d}, write out the diagonal matrix $\sum_K
  \mat{R}_K^{\mathsf T}\mat{R}_K$ of \cref{eq:multiplicity} on the three interior
  dofs. Confirm each entry is $2$ and explain, in one sentence, why the (omitted)
  boundary nodes would have multiplicity $1$.

  \item \textbf{Five elements.} Repeat \cref{wex:assemble-1d} for a mesh of $n=5$
  cells ($h=\tfrac15$, four interior dofs). List the five elements and the global
  indices of their local slots, then scatter-add the five $2\times2$ blocks and
  read off the $4\times4$ tridiagonal $\mat{K}$. Which dofs have multiplicity $2$?

  \item \textbf{A triangle's element matrix.} A single linear triangle has three
  local dofs (its vertices) and a $3\times3$ element matrix. Without computing its
  entries, write the assembly contribution $\mat{R}_K^{\mathsf
  T}\mat{K}^{\mathrm{elem}}\mat{R}_K$ for a triangle whose vertices are global dofs
  $\{2,5,7\}$: which $9$ entries of the global matrix does it touch, and what is
  the multiplicity contribution to dof $5$ if three triangles meet there?

  \item \textbf{Empty and singleton.} Using \cref{thm:foldlaws}, evaluate
  $\mathrm{asm}([\,])$ and $\mathrm{asm}([t])$ for a single term $t$. Then show
  $\mathrm{asm}([t_1,t_2])=\mathrm{asm}([t_2,t_1])$ directly from the two laws you
  need, naming each.

  \item \textbf{The homomorphism, again.} An eigenmode analysis needs both a
  stiffness term $t_K$ and a mass term $t_M$. Using law (ii), express
  $\mathrm{asm}([t_K,t_M])$ in terms of the separately-assembled $\mat{K}$ and
  $\mat{M}$. Then explain why the driven analysis, which forms
  $\mat{K}-\omega^2\mat{M}$ at many frequencies $\omega$, assembles $\mat{K}$ and
  $\mat{M}$ \emph{once} rather than re-running the fold at each frequency---and
  which law guarantees this is correct.

  \item \textbf{Matrix-free apply.} Starting from \cref{eq:matrixfree-preview},
  write out $\mat{K}\vw$ for the single-term four-element mesh of
  \cref{wex:assemble-1d} acting on $\vw=(1,0,0)^{\mathsf T}$, by gathering,
  applying the $2\times2$ element matrix, and scatter-adding---without ever forming
  $\mat{K}$. Check your answer against the first column of the tridiagonal
  $\mat{K}$.

  \item \textbf{Boundary conditions outside the fold.} Suppose you ``optimize'' the
  assembler to skip the equation of a pinned dof \emph{inside} the element loop,
  before summing. Give a two-term list $T_1\!+\!\!+\,T_2$ for which this fused
  assembler violates the concatenation homomorphism (ii), and identify the double-
  counting. (Hint: what happens to the pinned dof's row when both sublists pass
  through the elimination?)

  \item \textbf{The scatter-add race.} On the parallel picture of
  \cref{fig:parallelrace}, suppose two threads each execute
  \texttt{x[2] = x[2] + (local contribution)} on the shared dof $x_2$ without
  synchronization. Describe a specific interleaving of their reads and writes that
  loses one contribution, and explain why \cref{thm:foldlaws}(iii) nonetheless
  guarantees the \emph{intended} result is independent of which thread goes first.
  Which of the three fixes named in the text (atomics, colouring, private buffers)
  relies on law (ii)?
\end{enumerate}

\chapter{The Matrix-Free Operator}
\label{ch:matrixfree}

\chapterorient{Assembly (\cref{ch:assembly}) built a matrix. This chapter
un-builds it. We will see that a finite-element operator does not have to exist
as a stored array of numbers at all: the only thing a solver ever asks of an
operator is its \emph{action}, the map $\vv \mapsto \mat{A}\vv$, and that action
can be computed on the fly as a short chain of element-local tensor
contractions---gather, evaluate, weight, integrate, scatter---without ever
forming or storing $\mat{A}$. This is the \emph{matrix-free} operator. It is the
form that costs the least memory at high polynomial order, and---because every
stage is a contraction over named tensor axes---it is exactly the form a GPU or
other tensor accelerator wants to run. By the end you will be able to trace a
vector through all five stages, recover the same answer the assembled matrix
would have given, and see why the innermost kernel is the place the book's
``implementation view'' finally touches the silicon.}

\section{What a solver actually asks for}
\label{sec:mf-whatask}

Open any of the iterative solvers in \cref{part:framework}---conjugate
gradients (\cref{ch:cg}), GMRES (\cref{ch:gmres}), the Lanczos and Arnoldi
recurrences behind the eigensolver (\cref{ch:krylovschur})---and look at what
they do with the operator $\mat{A}$. They form inner products, they orthogonalize,
they take linear combinations of vectors. But the \emph{only} way any of them
touches $\mat{A}$ is through one operation: given a vector $\vv$, produce the
vector $\mat{A}\vv$. They never read an individual entry $A_{ij}$. They never ask
for a row, a column, or the sparsity pattern. A Krylov method is, by
construction, a procedure that explores the subspace
$\{\vv,\,\mat{A}\vv,\,\mat{A}^2\vv,\dots\}$, and to walk that subspace it needs
nothing but repeated application of the \emph{action}
\[
  \vv \;\longmapsto\; \mat{A}\vv .
\]

This is worth pausing on, because it inverts the picture most readers carry from
a first linear-algebra course. There, a matrix \emph{is} its array of entries;
``apply the matrix'' means ``do the row-times-column sums.'' But an
iterative solver does not need the entries---it needs a black box that, fed
$\vv$, returns $\mat{A}\vv$. How that box computes its answer is its own affair.
Assembly (\cref{ch:assembly}) is one way to build such a box: store every entry,
then multiply. This chapter is the other way.

\begin{keyidea}
A Krylov solver is a consumer of the operator's \emph{action}, $\vv\mapsto
\mat{A}\vv$, and of nothing else. It never inspects the entries of $\mat{A}$. So
any object that can answer ``given $\vv$, here is $\mat{A}\vv$'' is, to the
solver, a perfectly good operator---whether or not a matrix was ever built. This
single fact is the licence for everything in this chapter.
\end{keyidea}

\subsection{Why ``store the matrix'' gets expensive}
\label{sec:mf-cost}

If the action is all we need, why ever store the matrix? For low polynomial
order, storing it is cheap and fast, and assembly (\cref{ch:assembly}) is the
right tool. The trouble begins at \emph{high order}, which is exactly where
finite-element accuracy is won most cheaply.

Recall the element-local stiffness or mass matrix from \cref{ch:assembly}: on one
element with $L$ local degrees of freedom, the element matrix $\mat{A}^e$ is
$L\times L$, holding $L^2$ numbers. For a tensor-product element of polynomial
degree $p$ in $d$ spatial dimensions, the number of local degrees of freedom
grows as
\[
  L \;=\; (p+1)^d ,
\]
so the dense element matrix holds
\[
  L^2 \;=\; (p+1)^{2d}
\]
numbers. In three dimensions ($d=3$) that is $(p+1)^6$ entries \emph{per element}.
At $p=1$ this is a tidy $2^6 = 64$; at $p=4$ it is $5^6 = \num{15625}$; at $p=8$
it is $9^6 = \num{531441}$---half a million numbers for a single element. Storing
the assembled operator means storing something on this order, multiplied across
every element of the mesh. The memory cost of the \emph{stored} operator scales
as $O(p^{2d})$ per element, and on modern accelerators memory---not
arithmetic---is the scarce resource.

\Cref{fig:mf-cost} plots the contrast. The assembled operator's storage climbs
as $(p+1)^{2d}$; the matrix-free operator stores only the tabulated
one-dimensional basis and the per-quadrature-point geometry factor, a cost that
grows far more gently. Past a modest order---in 3-D, around $p=2$ to $4$
depending on the architecture---the matrix-free form is not merely lighter, it is
the only form that fits.

\begin{figure}[t]
\centering
\begin{tikzpicture}
\begin{axis}[
  width=0.86\linewidth, height=58mm,
  xlabel={polynomial order $p$},
  ylabel={numbers stored per element},
  ymode=log,
  xmin=1, xmax=8, ymin=1, ymax=1e6,
  xtick={1,2,3,4,5,6,7,8},
  ytick={1e0,1e2,1e4,1e6},
  legend pos=north west,
  legend cell align=left,
  grid=both, grid style={simGray!18},
  tick label style={font=\footnotesize},
  label style={font=\small},
  legend style={font=\footnotesize},
]
% assembled: (p+1)^6  (3-D, d=3, L^2)
\addplot[simRust, very thick, mark=*, mark size=1.6pt]
  coordinates {(1,64)(2,729)(3,4096)(4,15625)(5,46656)(6,117649)(7,262144)(8,531441)};
\addlegendentry{assembled $\;(p{+}1)^{6}$}
% matrix-free: tabulated 1-D basis + geom factor ~ O(p) +small const
% store ~ (p+1)^2 (1-D basis) + a few geom comps per quad pt; use (p+1)^2 + 6(p+1)
\addplot[simBlue, very thick, mark=square*, mark size=1.6pt]
  coordinates {(1,16)(2,27)(3,40)(4,55)(5,72)(6,91)(7,112)(8,135)};
\addlegendentry{matrix-free (1-D basis $+$ geom)}
\end{axis}
\end{tikzpicture}
\caption[Memory: assembled versus matrix-free]{Numbers stored per element, in
3-D, as polynomial order $p$ rises (log scale). The \emph{assembled} operator
holds the dense element matrix, $(p+1)^{2d}=(p+1)^6$ entries, and climbs
steeply. The \emph{matrix-free} operator stores only the tabulated
one-dimensional basis tables and a handful of geometry components per quadrature
point, growing roughly linearly. The two curves cross almost immediately;
beyond low order, matrix-free is the only form that fits in accelerator memory.}
\label{fig:mf-cost}
\end{figure}

\begin{physicsbox}
The arithmetic the operator performs is the same in both forms---a weighted sum
of basis-function interactions over each element, the discretization of the weak
form $a(u,v)=\inner{Q\diffop u}{\diffop v}$ from \cref{ch:galerkin}. What
differs is \emph{bookkeeping}: the assembled form pre-computes and stores all
those interactions as a matrix; the matrix-free form recomputes them, cheaply,
each time the action is needed. We are trading a little extra arithmetic
(recompute on demand) for a large saving in memory (store nothing). On hardware
where memory bandwidth is the bottleneck and arithmetic is nearly free---the
defining shape of a modern GPU---this is overwhelmingly the right trade.
\end{physicsbox}

\section{The action as a chain of contractions}
\label{sec:mf-chain}

So we want to compute $\mat{A}\vv$ without $\mat{A}$. The route is to go back to
where the matrix came from. In \cref{ch:assembly} we built $\mat{A}$ by
\emph{assembling} element contributions: each element computed a local operator
from (i) restricting the global vector to the element, (ii) evaluating the
trial field and its derivatives at the element's quadrature points, (iii)
weighting those values by the geometry and material at each point, and (iv)
integrating back to the element's degrees of freedom; then (v) the element
results were summed into the global vector. Assembly did this \emph{once}, for
every basis vector, and stored the resulting matrix. Matrix-free does exactly the
same five operations, but \emph{on the actual input vector $\vv$}, every time the
action is requested---and stores nothing.

Each of the five operations is a linear map, so the whole action is their
composition. Reading right to left in the order they act on $\vv$:
\begin{equation}
  \mat{A}\,\vv
   \;=\;
   \underbrace{G^{\mathsf T}}_{\text{scatter}}\;
   \underbrace{B_{\diffop}^{\mathsf T}}_{\text{integrate}}\;
   \underbrace{D}_{\text{weight}}\;
   \underbrace{B_{\diffop}}_{\text{evaluate}}\;
   \underbrace{G}_{\text{gather}}\;
   \vv .
  \label{eq:mf-chain}
\end{equation}
This factorization is the heart of the chapter. It says the finite-element
operator is a \emph{sandwich}: a gather $G$ and its transpose scatter
$G^{\mathsf T}$ on the outside; a basis evaluation $B_{\diffop}$ and its transpose
$B_{\diffop}^{\mathsf T}$ inside that; and at the very center a single pointwise
multiply $D$. The symmetry is not a coincidence---it is the discrete shadow of
the weak form's own symmetry, $a(u,v)=a(v,u)$ for the self-adjoint terms, with
$G$ and $B$ carrying the trial side and their transposes carrying the test side
(\cref{ch:galerkin}).

\begin{notationbox}
The subscript $\diffop$ on $B_{\diffop}$ records \emph{which} differential
operator the weak-form term uses---value (identity), gradient, curl, or
divergence---the four members of the de Rham family from \cref{ch:derham}. A
mass term ($\diffop = \mathrm{Id}$) evaluates basis \emph{values}; a stiffness
term ($\diffop=\Grad$) evaluates basis \emph{gradients}; a curl--curl term
evaluates \emph{curls}. The choice of $\diffop$ is the only thing that changes
between the operators of the six analyses (\cref{ch:targets}); the five-stage
skeleton is identical for all of them.
\end{notationbox}

\subsection{The five stages and their tensor shapes}
\label{sec:mf-stages}

Each stage carries a tensor from one named shape to the next. The axes are the
element-local family we will fix carefully in \cref{sec:mf-anatomy}; for now read
$N$ as the global degree-of-freedom count, $E$ as the number of elements, $L$ as
local degrees of freedom per element, $P$ as quadrature points per element, and
$C$ as the number of value (or derivative) components carried at each quadrature
point. \Cref{fig:mf-pipeline} is the hero diagram of the chapter; the prose
below walks its five boxes left to right.

\begin{figure}[p]
\centering
\resizebox{\linewidth}{!}{%
\begin{tikzpicture}[node distance=8mm and 9mm]
  % input
  \node[data, text width=18mm] (in) {global vector\\\texttt{Tensor[N]}};
  % G
  \node[stage, right=of in, text width=20mm] (G) {$G$\\\footnotesize gather\\(element restrict)};
  % B
  \node[stage, right=of G, text width=20mm] (B) {$B_{\diffop}$\\\footnotesize evaluate\\(basis apply)};
  % D  -- the extern kernel center
  \node[extern, right=of B, text width=20mm] (D) {$D$\\\footnotesize weight\\\texttt{\#extern}};
  % Bt
  \node[stage, right=of D, text width=20mm] (Bt) {$B_{\diffop}^{\mathsf T}$\\\footnotesize integrate};
  % Gt
  \node[stage, right=of Bt, text width=20mm] (Gt) {$G^{\mathsf T}$\\\footnotesize scatter-add};
  % output
  \node[data, right=of Gt, text width=18mm] (out) {global vector\\\texttt{Tensor[N]}};

  \draw[flow] (in) -- (G);
  \draw[flow] (G) -- (B);
  \draw[flow] (B) -- (D);
  \draw[flow] (D) -- (Bt);
  \draw[flow] (Bt) -- (Gt);
  \draw[flow] (Gt) -- (out);

  % shape annotations between stages (below the arrows)
  \node[font=\scriptsize\ttfamily, simGray, below=2.5mm of G.south] {[E, L]};
  \node[font=\scriptsize\ttfamily, simGray, below=2.5mm of B.south] {[E, P, C]};
  \node[font=\scriptsize\ttfamily, simGray, below=2.5mm of D.south] {[E, P, C]};
  \node[font=\scriptsize\ttfamily, simGray, below=2.5mm of Bt.south] {[E, L]};
  \node[font=\scriptsize\ttfamily, simBlue, below=3mm of in.south, align=center] {action\\input};
  \node[font=\scriptsize\ttfamily, simBlue, below=3mm of out.south, align=center] {action\\output};

  % geom_data feeding D from below
  \node[data, fill=simBgGold, draw=simGold, below=12mm of D, text width=26mm]
    (geom) {\footnotesize\itshape geom\_data \texttt{[E, P, G]}\\(build-time factor)};
  \draw[flow, simGold] (geom) -- (D);

  % brackets: outer = gather/scatter sandwich, inner = eval/integrate
  \draw[decorate, decoration={brace, amplitude=4pt}, simRust]
    ([yshift=7mm]G.north west) -- ([yshift=7mm]Gt.north east)
    node[midway, above=4pt, font=\scriptsize, simRust]
    {element-local: independent, parallel over $E$};
\end{tikzpicture}%
}
\caption[The five-stage matrix-free pipeline]{The matrix-free operator
$\mat{A}\vv = G^{\mathsf T}B_{\diffop}^{\mathsf T}\,D\,B_{\diffop}\,G\,\vv$ as a
dataflow pipeline, with the tensor shape annotated on each wire. A global vector
\lstinline[style=l4style]|Tensor[N]| enters; $G$ \emph{gathers} it to per-element
local tensors \lstinline[style=l4style]|[E, L]|; $B_{\diffop}$ \emph{evaluates}
the field (or its gradient/curl) at the \lstinline[style=l4style]|[E, P, C]|
quadrature points; the center stage $D$ \emph{weights} each point by the
precomputed \texttt{geom\_data} (gold, fed from below; the
\lstinline[style=l4style]|#extern| innermost kernel, \cref{sec:mf-impl-view});
$B_{\diffop}^{\mathsf T}$ \emph{integrates} back to local dofs; and $G^{\mathsf T}$
\emph{scatters} the element results into a fresh global vector. The whole middle
of the pipeline (rust brace) is independent across elements---the source of all
the parallelism.}
\label{fig:mf-pipeline}
\end{figure}

\paragraph{Stage 1 --- $G$, the gather (element restriction).}
The input is the global vector $\vv$, a \lstinline[style=l4style]|Tensor[N]| with
one entry per global degree of freedom. The gather $G$ reads, for each element,
the values of the degrees of freedom that element touches, copying them into a
per-element block. The output is a \lstinline[style=l4style]|Tensor[E, L]|: $E$
elements, each holding its $L$ local degrees of freedom. No arithmetic happens
here---$G$ is pure copying, a Boolean selection that fans the shared global
vector out into (generally overlapping) element-local pieces. A degree of freedom
sitting on a shared element boundary is copied into \emph{every} element that
touches it; that duplication is intentional and is undone, by summation, in the
final stage.

\paragraph{Stage 2 --- $B_{\diffop}$, the basis evaluation.}
Now each element holds its local coefficients. To apply the weak form we need the
\emph{field} those coefficients represent---and its derivative---sampled at the
element's quadrature points. That is exactly what the tabulated basis gives us
(\cref{ch:refelement}): a small dense matrix whose entries are the basis
functions (or their derivatives) evaluated at the quadrature nodes. $B_{\diffop}$
contracts each element's \lstinline[style=l4style]|[L]| coefficient vector against
that matrix, producing \lstinline[style=l4style]|[P, C]| values---$C$ components
($C=1$ for a scalar value, $C=d$ for a gradient) at each of $P$ quadrature points.
Stacked over elements, the output is a \lstinline[style=l4style]|Tensor[E, P, C]|.
The mode---values, gradient, curl, or divergence---is the $\diffop$ the term
asked for.

\paragraph{Stage 3 --- $D$, the pointwise quadrature weight.}
This is the center of the sandwich, and it is the simplest stage of all: at each
quadrature point, multiply the evaluated value by a single precomputed factor and
move on. The factor bundles three things---the material coefficient $Q$ (the
$\varepsilon$, $\mu^{-1}$, $\dots$ of the weak form), the geometry metric derived
from the element's Jacobian, and the quadrature weight $w$---into one number (or,
for vector terms, one small matrix) per point, stored in a tensor
\texttt{geom\_data} of shape \lstinline[style=l4style]|Tensor[E, P, G]|. There is
\emph{no coupling between points}: the output at point $(e,p)$ depends only on the
input at $(e,p)$ and on \texttt{geom\_data}$[e,p]$. The shape is unchanged,
\lstinline[style=l4style]|[E, P, C]| in and out (we write $C'$ for the output
component count, usually equal to $C$). \Cref{fig:mf-diagonal} shows why this is
the \emph{embarrassingly parallel} diagonal of the whole operator.

\paragraph{Stage 4 --- $B_{\diffop}^{\mathsf T}$, integrate back to dofs.}
The weighted quadrature-point values must now be turned back into element-local
degree-of-freedom contributions---this is the discrete integral $\int (\cdot)\,
\phi_i$ against each test basis function $\phi_i$. It is the exact transpose of
stage 2: the same tabulated basis matrix, applied on the other side. The shape
goes back from \lstinline[style=l4style]|[E, P, C]| to
\lstinline[style=l4style]|[E, L]|.

\paragraph{Stage 5 --- $G^{\mathsf T}$, the scatter-add (assembly).}
Finally, the per-element \lstinline[style=l4style]|[E, L]| contributions must be
returned to the single global vector. $G^{\mathsf T}$ is the transpose of the
gather: where $G$ copied a shared dof out to several elements, $G^{\mathsf T}$
\emph{sums} the several element contributions back into that one global slot.
This summation is precisely the finite-element ``assembly'' step
(\cref{ch:assembly})---here applied to a vector rather than a matrix. The output
is a fresh \lstinline[style=l4style]|Tensor[N]|: the action $\mat{A}\vv$, complete.

\begin{keyidea}
The matrix-free operator is a chain of tensor contractions over named axes---no
sparse-matrix data structure anywhere. \emph{Gather} ($N\!\to\!E,L$),
\emph{evaluate} ($E,L\!\to\!E,P,C$), \emph{weight} (pointwise on $E,P,C$),
\emph{integrate} ($E,P,C\!\to\!E,L$), \emph{scatter} ($E,L\!\to\!N$). The matrix
$\mat{A}$ that assembly would have stored is recovered, on demand, as the
composition $G^{\mathsf T}B_{\diffop}^{\mathsf T}DB_{\diffop}G$---and never
materialized.
\end{keyidea}

\subsection{The chain in pseudocode}
\label{sec:mf-pseudo}

In the calculus of \cref{part:framework}, the action reads as a single
composition of the five stage functions. Each stage is a pure tensor-to-tensor
map; threading a vector through them is exactly the value-threading picture of
\cref{ch:bigidea}.

\begin{lstlisting}[style=l4style]
-- the five element-local stages, as pure tensor maps
gather    :: ElemRestriction -> Tensor[N]        -> Tensor[E, L]      -- G
evaluate  :: BasisMode -> Basis -> Tensor[E, L]  -> Tensor[E, P, C]  -- B_D
weight    :: GeomData -> Tensor[E, P, C]         -> Tensor[E, P, C]  -- D
integrate :: BasisMode -> Basis -> Tensor[E, P, C] -> Tensor[E, L]   -- B_D^T
scatter   :: ElemRestriction -> Tensor[E, L]     -> Tensor[N]        -- G^T

-- the matrix-free action A v = G^T . B_D^T . D . B_D . G  v
apply_matrixfree :: Operator -> Tensor[N] -> Tensor[N]
apply_matrixfree op v =
  scatter   op.restr           $   -- G^T  : [E, L] -> [N]
  integrate op.mode op.basis   $   -- B_D^T: [E, P, C] -> [E, L]
  weight    op.geom            $   -- D    : pointwise on [E, P, C]
  evaluate  op.mode op.basis   $   -- B_D  : [E, L] -> [E, P, C]
  gather    op.restr v             -- G    : [N] -> [E, L]
\end{lstlisting}

\noindent Read it bottom to top, the order the data flows: \texttt{gather} the
global vector to element blocks, \texttt{evaluate} the basis, \texttt{weight}
pointwise, \texttt{integrate} back, \texttt{scatter} to a fresh global vector.
The operator value \texttt{op} carries the three \emph{build-time} pieces---the
restriction index map \texttt{op.restr}, the tabulated \texttt{op.basis}, and the
geometry factor \texttt{op.geom}---which we turn to next.

\section{The element-local tensor and its axes}
\label{sec:mf-anatomy}

The five stages move data between three rank-structured tensors. Naming their
axes precisely is what lets us reason about shapes the way \cref{ch:tensors}
taught---and it is what makes the whole pipeline legible to a tensor backend.
\Cref{fig:mf-anatomy} is the anatomy chart; \cref{tab:mf-axes} is the legend.

The crucial conceptual move is that matrix-free FE assembly works over a
\emph{richer} shape than the flat global vector. The solver sees a flat
\lstinline[style=l4style]|Tensor[N]|---a rank-1 list of $N$ degrees of freedom.
But the moment we gather to elements, we acquire an explicit \emph{element axis}
$E$ and the per-element structure beneath it. The flat global vector and the
element-local tensors are two different vocabularies, and the gather/scatter pair
$G$, $G^{\mathsf T}$ is exactly the boundary between them.

\begin{table}[t]
\centering
\caption[The element-local tensor axes]{The five named axes of the element-local
tensor family. The \emph{stratum} column records whether an axis's extent is
fixed once per $(\text{mesh},\text{order})$ in a setup pass (build-time) or
indexes a value produced afresh on every operator application (run-time)---the
distinction \cref{sec:mf-strata} makes load-bearing.}
\label{tab:mf-axes}
\small
\setlength{\tabcolsep}{6pt}
\renewcommand{\arraystretch}{1.25}
\begin{tabular}{@{}clll@{}}
\toprule
axis & name & meaning & stratum \\
\midrule
$E$ & element count & number of mesh elements assembled over & build-time \\
$L$ & local dofs & dofs of the basis local to one element, $(p{+}1)^d$ & build-time \\
$P$ & quad points & quadrature points per element & build-time \\
$C$ & components & value/derivative components at a quad point & run-time \\
$G$ & geom components & precomputed metric storage per quad point & build-time \\
\bottomrule
\end{tabular}
\end{table}

\begin{figure}[t]
\centering
\begin{tikzpicture}[scale=1.0]
  % three stacked tensors with axis labels
  % --- [E, L] ---
  \begin{scope}[shift={(-4.4,0)}]
    \foreach \k in {0,1,2}{
      \begin{scope}[shift={(0.12*\k,0.09*\k)}]
        \fill[simBgBlue, draw=simBlue] (-0.7,-0.5) rectangle (0.7,0.5);
        \foreach \y in {-0.25,0,0.25}{\draw[simBlue!40] (-0.7,\y)--(0.7,\y);}
      \end{scope}}
    \node[font=\scriptsize, simInk] at (-1.05,0.0) {$L$};
    \node[font=\scriptsize, simRust] at (0.55,0.75) {$E$};
    \node[font=\scriptsize\ttfamily, simInk] at (0.1,-1.0) {[E, L]};
    \node[font=\scriptsize\itshape, simGray, text width=24mm, align=center]
      at (0.1,-1.5) {per-element local dofs};
  \end{scope}
  \node[font=\large] at (-2.5,0) {$\xrightarrow{\,B_{\diffop}\,}$};
  % --- [E, P, C] ---
  \begin{scope}[shift={(-0.4,0)}]
    \foreach \k in {0,1,2}{
      \begin{scope}[shift={(0.12*\k,0.09*\k)}]
        \fill[simBgGreen, draw=simGreen] (-0.75,-0.5) rectangle (0.75,0.5);
        \foreach \y in {-0.17,0.17}{\draw[simGreen!40] (-0.75,\y)--(0.75,\y);}
        \foreach \x in {-0.25,0.25}{\draw[simGreen!40] (\x,-0.5)--(\x,0.5);}
      \end{scope}}
    \node[font=\scriptsize, simInk] at (-1.1,0.0) {$P$};
    \node[font=\scriptsize, simInk] at (0.0,-0.85) {$C$};
    \node[font=\scriptsize, simRust] at (0.6,0.75) {$E$};
    \node[font=\scriptsize\ttfamily, simInk] at (0.05,-1.15) {[E, P, C]};
    \node[font=\scriptsize\itshape, simGray, text width=26mm, align=center]
      at (0.05,-1.65) {values at quad points};
  \end{scope}
  \node[font=\large] at (1.6,0) {$\odot$};
  % --- [E, P, G]  geom_data ---
  \begin{scope}[shift={(3.7,0)}]
    \foreach \k in {0,1,2}{
      \begin{scope}[shift={(0.12*\k,0.09*\k)}]
        \fill[simBgGold, draw=simGold] (-0.75,-0.5) rectangle (0.75,0.5);
        \foreach \y in {-0.17,0.17}{\draw[simGold!50] (-0.75,\y)--(0.75,\y);}
        \foreach \x in {-0.25,0.25}{\draw[simGold!50] (\x,-0.5)--(\x,0.5);}
      \end{scope}}
    \node[font=\scriptsize, simInk] at (-1.1,0.0) {$P$};
    \node[font=\scriptsize, simInk] at (0.0,-0.85) {$G$};
    \node[font=\scriptsize, simRust] at (0.6,0.75) {$E$};
    \node[font=\scriptsize\ttfamily, simInk] at (0.05,-1.15) {[E, P, G]};
    \node[font=\scriptsize\itshape, simGray, text width=26mm, align=center]
      at (0.05,-1.65) {geom\_data (build-time)};
  \end{scope}
\end{tikzpicture}
\caption[Anatomy of the element-local tensor]{The element-local tensor family.
The gather produces \lstinline[style=l4style]|[E, L]| (blue): an element axis $E$
over per-element local-dof blocks of length $L$. Basis evaluation
$B_{\diffop}$ maps it to \lstinline[style=l4style]|[E, P, C]| (green): $C$
components at each of $P$ quadrature points, per element. The pointwise weight
$D$ multiplies it, point by point ($\odot$), against the precomputed geometry
carrier \lstinline[style=l4style]|[E, P, G]| (gold). The element axis $E$ (rust)
runs through all three---the operator is block-diagonal in $E$, which is to say
embarrassingly parallel over elements.}
\label{fig:mf-anatomy}
\end{figure}

\begin{notationbox}
A naming hazard worth flagging once. In the underlying libCEED library the basis
constructor's parameter named \texttt{P} counts \emph{nodes per element} (our
$L$) and its \texttt{Q} counts \emph{quadrature points} (our $P$). This book uses
$L$ for local dofs and $P$ for quadrature points consistently; if you read the
library source, expect the letters to swap. The mapping is exact---only the
labels differ.
\end{notationbox}

\subsection{Why every stage is block-diagonal in \texorpdfstring{$E$}{E}}
\label{sec:mf-blockdiag}

Look again at \cref{fig:mf-anatomy}: the element axis $E$ threads through all
three tensors, and through stages 2, 3, and 4 of the pipeline. None of those
stages couples one element to another. Basis evaluation applies the same small
matrix \emph{independently} within each element; the pointwise weight touches one
point at a time; integration is again per-element. In the language of
\cref{ch:tensors}, stages 2--4 are \emph{block-diagonal} in $E$: a batched
operation, the same computation broadcast across the $E$ slices.

All the inter-element coupling---the only place where one element's data meets
another's---lives in the gather $G$ and the scatter $G^{\mathsf T}$, at the two
ends of the pipeline. This is a clean and consequential separation. The
arithmetic-heavy middle is perfectly parallel; the communication-heavy ends are
pure data movement with no arithmetic. It is exactly the structure a tensor
accelerator is built to exploit: map a batched contraction over a large axis,
with a gather in front and a scatter behind.

\section{Build-time versus run-time: where the geometry lives}
\label{sec:mf-strata}

We have been quietly carrying a tensor we have not built: \texttt{geom\_data},
the \lstinline[style=l4style]|Tensor[E, P, G]| that the weight stage $D$
multiplies against. Where does it come from, and---more importantly---\emph{when}?
The answer is a recurring theme of this book, large enough that it has its own
chapter ahead (\cref{ch:strata}): the operator's data divides into two
\emph{strata} by lifetime.

\begin{description}[leftmargin=1.4em, style=nextline]
\item[Build-time (setup) stratum.] Data fixed once per
$(\text{mesh},\text{FE order},\text{quadrature rule})$ and reused across every
application of the operator. This is the restriction index map (which dof goes to
which element), the tabulated basis matrices (basis functions sampled at
quadrature nodes), and---the one with real arithmetic---\texttt{geom\_data}.
\item[Run-time (apply) stratum.] Data produced afresh on every action
$\vv\mapsto\mat{A}\vv$: the gathered \lstinline[style=l4style]|[E, L]| block, the
evaluated \lstinline[style=l4style]|[E, P, C]| values, the weighted values, the
integrated contributions. These are transient---born inside one application and
discarded when it returns.
\end{description}

The geometry factor is the star of the build-time stratum. Computing it requires
real work: at each quadrature point of each element, form the Jacobian $J$ of the
map from the reference element to the physical element (\cref{ch:refelement}),
combine it into the metric the weak-form term needs---$|J|$ for a mass term,
$J^{-\mathsf T}J^{-1}|J|$ for a gradient term---multiply by the quadrature weight
$w$, and fold in the material coefficient $Q$. The result is one small block per
point, stored in \texttt{geom\_data}. \Cref{fig:mf-strata} draws the split: a
setup pass runs the build once and produces \texttt{geom\_data}; thereafter every
operator application, of which a Krylov solve performs hundreds, simply reads it.

\begin{figure}[t]
\centering
\begin{tikzpicture}[node distance=7mm and 12mm]
  % --- build-time row (gold) ---
  \node[data, fill=simBgGold, draw=simGold, text width=20mm] (mesh)
    {mesh nodes\\$+$ quad rule};
  \node[stage, fill=simBgGold, draw=simGold, right=of mesh, text width=26mm]
    (build) {\textbf{setup pass}\\\footnotesize $J,\;|J|,\;Q,\;w$\\(\texttt{geom\_factor\_build})};
  \node[data, fill=simBgGold, draw=simGold, right=of build, text width=22mm]
    (gd) {\texttt{geom\_data}\\\texttt{[E, P, G]}};
  \draw[flow, simGold] (mesh) -- (build);
  \draw[flow, simGold] (build) -- (gd);
  \node[font=\scriptsize\bfseries, simGold, left=1mm of mesh, rotate=90, anchor=south]
    {build-time};
  \node[font=\scriptsize\itshape, simGray, below=1mm of build]
    {runs \emph{once} per (mesh, order)};

  % --- run-time row (blue): many applies reading geom_data ---
  \node[stage, below=18mm of build, text width=30mm] (apply1)
    {apply $\mat{A}\vv_1$\\\footnotesize $G^{\mathsf T}B^{\mathsf T}\,D\,B\,G$};
  \node[stage, right=5mm of apply1, text width=14mm] (apply2)
    {apply\\$\mat{A}\vv_2$};
  \node[font=\large, right=2mm of apply2] (dots) {$\cdots$};
  \node[stage, right=2mm of dots, text width=14mm] (applyk)
    {apply\\$\mat{A}\vv_k$};
  \node[font=\scriptsize\bfseries, simBlue, left=1mm of apply1, rotate=90, anchor=south]
    {run-time};
  \node[font=\scriptsize\itshape, simGray, below=1mm of apply1, xshift=12mm]
    {runs \emph{every} Krylov iteration --- reads \texttt{geom\_data}, never rebuilds it};
  % geom_data feeds every apply
  \draw[flow, simGold] (gd.south) -- ++(0,-0.4) -| (apply1.north);
  \draw[flow, simGold] (gd.south) -- ++(0,-0.4) -| (apply2.north);
  \draw[flow, simGold] (gd.south) -- ++(0,-0.4) -| (applyk.north);
\end{tikzpicture}
\caption[Build-time versus run-time stratification]{The geometry factor is
computed \emph{once}. A setup pass (gold, top) consumes the mesh nodes and the
quadrature rule and produces \texttt{geom\_data}, a
\lstinline[style=l4style]|[E, P, G]| tensor holding the Jacobian metric,
material coefficient, and quadrature weight at every point. Every subsequent
operator application (blue, bottom)---and an iterative solve performs
hundreds---simply \emph{reads} \texttt{geom\_data} in its weight stage $D$. The
expensive geometry work is amortized across the whole solve. Only when the mesh
changes (e.g.\ adaptive refinement, \cref{ch:amr}) is the setup pass rerun.}
\label{fig:mf-strata}
\end{figure}

\begin{keyidea}
The matrix-free operator carries two strata. The \emph{build-time} stratum---the
restriction map, the tabulated basis, and the geometry factor
\texttt{geom\_data}---is computed once per $(\text{mesh},\text{order})$ and
reused. The \emph{run-time} stratum---the gathered, evaluated, weighted,
integrated tensors---is rebuilt on every application. Folding the coefficient,
Jacobian metric, and quadrature weight into one precomputed \texttt{geom\_data}
factor is what makes the run-time weight stage $D$ a single pointwise multiply.
The same stratification organizes the whole simulator (\cref{ch:strata}).
\end{keyidea}

\subsection{The diagonal $D$ is embarrassingly parallel}
\label{sec:mf-diagonal}

With \texttt{geom\_data} precomputed, the run-time weight stage $D$ is the
simplest object in the entire finite-element machine: a pointwise multiply.
\Cref{fig:mf-diagonal} makes the structure explicit. Lay the
\lstinline[style=l4style]|[E, P, C]| evaluated field and the
\lstinline[style=l4style]|[E, P, G]| geometry factor side by side; at each point
$(e,p)$, multiply the field block by the geometry block, independently. No point
talks to any other point. There is no reduction, no recurrence, no
order-dependence: the output is a function of the input \emph{at the same index}
and nothing else.

\begin{figure}[t]
\centering
\begin{tikzpicture}[scale=1.0]
  % grid of quad points (E x P flattened), each an independent multiply
  \def\nx{5}\def\ny{3}
  \foreach \i in {0,...,4}{
    \foreach \j in {0,...,2}{
      \node[draw=simGreen, fill=simBgGreen, rounded corners=1pt,
            minimum size=6.5mm, font=\scriptsize] (c\i\j) at (\i*1.0, \j*0.85) {$u_{\i\j}$};
    }
  }
  % geom blocks to the right, one per point (drawn as a thin gold strip overlay)
  \foreach \i in {0,...,4}{
    \foreach \j in {0,...,2}{
      \node[font=\scriptsize, simGold] at (\i*1.0+0.34, \j*0.85+0.34) {$\odot g$};
    }
  }
  % axis labels
  \node[font=\scriptsize\itshape, simRust] at (-0.85, 0.85) {$(e,p)$ index};
  \node[font=\scriptsize, simInk] at (2.0,-0.85) {$P$ points $\times$ $E$ elements (flattened)};
  % big brace: all independent
  \draw[decorate, decoration={brace, amplitude=5pt}, simBlue]
    (-0.45, 2.05) -- (4.45, 2.05)
    node[midway, above=5pt, font=\footnotesize, simBlue]
    {every cell independent --- no coupling, any order, fully parallel};
\end{tikzpicture}
\caption[The quad-point weight as a parallel diagonal]{The weight stage
$D$ visualized over the flattened $(e,p)$ index. Each cell is one quadrature
point's evaluated value $u_{ep}$, multiplied ($\odot$) by its own precomputed
geometry block $g_{ep}$ from \texttt{geom\_data}. No cell depends on any other:
$D$ is \emph{diagonal} in the $(E,P)$ index, the embarrassingly parallel core of
the operator. This is the per-point ``lift'' the implementation view consumes
directly---one tensor, one elementwise multiply, mapped across the whole point
set at once.}
\label{fig:mf-diagonal}
\end{figure}

This is what we mean when we call $D$ the \emph{diagonal} of the operator. The
basis stages $B_{\diffop}$ and $B_{\diffop}^{\mathsf T}$ couple the $L$ local dofs
to the $P$ quadrature points (a dense small contraction); the gather and scatter
couple elements through shared dofs. But sandwiched between them, $D$ couples
nothing---it is a diagonal scaling in the quadrature-point basis. A diagonal map
is the friendliest possible object for a parallel machine, and it is where the
material physics ($Q$) and the geometry ($J$) actually enter the operator. Every
bit of what makes this term a \emph{mass} term rather than a \emph{stiffness}
term, this material rather than that one, this curved element rather than a
straight one, is packed into the diagonal $D$.

\section{Worked example: a 1-D mass operator, two elements}
\label{sec:mf-wex1}

Nothing convinces like tracing a real vector through the pipeline and recovering
the answer the assembled matrix would have given. We take the smallest nontrivial
case: the mass operator $\mat{M}$ ($\diffop=\mathrm{Id}$, weight $Q=1$) on a 1-D
mesh of two linear elements, evaluated by a two-point quadrature exact for the
quadratics that arise.

\begin{workedexample}{Matrix-free $\mat{M}\vv$ on two 1-D elements}{mf-mass}
\textbf{Setup.} The domain $[0,2]$ is split into two elements,
$e_0=[0,1]$ and $e_1=[1,2]$, each of width $h=1$. We use linear ($p=1$)
nodal basis functions, so there are $N=3$ global dofs at nodes
$x_0=0,\,x_1=1,\,x_2=2$, and $L=2$ local dofs per element. The node $x_1$ is
\emph{shared}: it is local dof $1$ of $e_0$ and local dof $0$ of $e_1$. Take
$P=2$ Gauss points per element. We compute the action $\mat{M}\vv$ on
\[
  \vv \;=\; (v_0,v_1,v_2) \;=\; (1,\,2,\,3) .
\]

\textbf{Build-time pieces.} On the reference element $[0,1]$ the linear basis is
$\phi_0(\xi)=1-\xi,\ \phi_1(\xi)=\xi$. The two Gauss points are
$\xi_{0,1}=\tfrac12\mp\tfrac{1}{2\sqrt3}$ with weights $w=\tfrac12$ each (scaled
for $[0,1]$). The basis-evaluation matrix $B$ (rows = quad points, cols = local
dofs) and the geometry factor are
\[
  B=\begin{pmatrix}\phi_0(\xi_0)&\phi_1(\xi_0)\\[2pt]\phi_0(\xi_1)&\phi_1(\xi_1)\end{pmatrix}
   =\begin{pmatrix}0.7887&0.2113\\[2pt]0.2113&0.7887\end{pmatrix},
  \qquad
  g_{ep}=\underbrace{|J|}_{=\,h=1}\cdot \underbrace{w}_{=\,1/2}\cdot \underbrace{Q}_{=1}
        =\tfrac12 .
\]
Both elements are identical (uniform mesh), so $B$ and $g=\tfrac12$ are the same
for $e_0$ and $e_1$. These are computed \emph{once}; the four stages below run on
$\vv$.

\tcblower
\textbf{Stage 1 --- gather $G$.} Read each element's two dofs from $\vv=(1,2,3)$:
\[
  e_0:\ (v_0,v_1)=(1,2),\qquad e_1:\ (v_1,v_2)=(2,3)
  \quad\Rightarrow\quad
  \texttt{[E, L]}=\begin{pmatrix}1&2\\2&3\end{pmatrix}.
\]
Note $v_1=2$ is copied into \emph{both} rows---the shared dof.

\textbf{Stage 2 --- evaluate $B$.} Contract each element's local dofs against
$B$ to get values at the two quad points (here $C=1$, scalar values):
\[
  e_0:\ B\begin{pmatrix}1\\2\end{pmatrix}=\begin{pmatrix}1.211\\1.789\end{pmatrix},
  \qquad
  e_1:\ B\begin{pmatrix}2\\3\end{pmatrix}=\begin{pmatrix}2.211\\2.789\end{pmatrix}.
\]
Shape is now \lstinline[style=l4style]|[E, P, C]| $=[2,2,1]$.

\textbf{Stage 3 --- weight $D$.} Multiply every quad-point value by $g=\tfrac12$:
\[
  e_0:\ (0.6057,\,0.8943),\qquad e_1:\ (1.1057,\,1.3943).
\]

\textbf{Stage 4 --- integrate $B^{\mathsf T}$.} Apply $B^{\mathsf T}$ (the
transpose) to return to local dofs:
\[
  e_0:\ B^{\mathsf T}\!\begin{pmatrix}0.6057\\0.8943\end{pmatrix}
       =\begin{pmatrix}0.6667\\0.8333\end{pmatrix},
  \quad
  e_1:\ B^{\mathsf T}\!\begin{pmatrix}1.1057\\1.3943\end{pmatrix}
       =\begin{pmatrix}1.1667\\1.3333\end{pmatrix}.
\]

\textbf{Stage 5 --- scatter-add $G^{\mathsf T}$.} Sum the element-local
contributions back into the $3$ global slots, \emph{adding} at the shared node
$x_1$:
\[
  \mat{M}\vv=\Big(\,0.6667,\ \underbrace{0.8333+1.1667}_{\text{shared }x_1},\ 1.3333\,\Big)
            =(0.6667,\ 2.0000,\ 1.3333).
\]

\textbf{Check against the assembled matrix.} The consistent linear mass matrix on
this uniform two-element mesh is the classic
\[
  \begin{aligned}
  \mat{M}&=\frac{h}{6}\begin{pmatrix}2&1&0\\1&4&1\\0&1&2\end{pmatrix}
        =\frac16\begin{pmatrix}2&1&0\\1&4&1\\0&1&2\end{pmatrix}, \\
  \mat{M}\vv&=\frac16\begin{pmatrix}2{\cdot}1+1{\cdot}2\\1{\cdot}1+4{\cdot}2+1{\cdot}3\\1{\cdot}2+2{\cdot}3\end{pmatrix}
            =\frac16\begin{pmatrix}4\\12\\8\end{pmatrix}
            =(0.6667,\,2.0000,\,1.3333).
  \end{aligned}
\]
The two agree exactly. The matrix-free pipeline reproduced $\mat{M}\vv$ to the
last digit---without ever forming the $3\times3$ matrix.
\end{workedexample}

\begin{remark}
Trace what happened at the shared node $x_1$. The gather $G$ \emph{copied} $v_1$
into both elements (stage 1); the scatter $G^{\mathsf T}$ \emph{summed} the two
elements' contributions back (stage 5). That copy-out/sum-back pair is the entire
content of finite-element assembly, and here it is, applied to a vector instead
of stamped into a matrix. The middle three stages never knew the node was
shared---they ran element by element, blissfully parallel.
\end{remark}

\section{Sum-factorization: the same answer, far cheaper}
\label{sec:mf-sumfac}

The worked example was 1-D, where the basis evaluation $B$ is already a small
dense matrix and there is nothing to optimize. In higher dimensions on
tensor-product elements, the basis stage hides a beautiful and important
efficiency, \emph{sum-factorization}, which is what makes high-order
matrix-free competitive. It is a \emph{transparent} performance trick in the
sense of \cref{ch:bigidea}: it computes exactly the same result, just by
reordering the contraction.

Consider a 2-D tensor-product ($Q_p$) element of order $p$, with
$L=(p+1)^2$ local dofs arranged on a $(p+1)\times(p+1)$ grid, and $P=(p+1)^2$
quadrature points likewise on a grid. The naive basis evaluation treats $B$ as a
dense $P\times L$ matrix: contracting one element's dofs against it costs
\[
  \text{naive cost} \;=\; P\cdot L \;=\; (p+1)^2\cdot(p+1)^2 \;=\; (p+1)^4
\]
multiply-adds per element. But the 2-D basis is a \emph{product} of
one-dimensional bases: $\phi_{ij}(\xi,\eta)=\phi_i(\xi)\,\phi_j(\eta)$. That
product structure lets us evaluate one dimension at a time. First contract the
$\eta$-direction dofs against the 1-D basis (a $(p+1)\times(p+1)$ matrix applied
along one axis), then contract the $\xi$-direction. Each of the two
sweeps costs $(p+1)$ (the 1-D contraction) times $(p+1)^2$ (the other axis carried
along), so
\[
  \text{sum-factored cost} \;=\; d\cdot(p+1)\cdot(p+1)^2 \;=\; d\,(p+1)^{d+1}
  \;=\; 2\,(p+1)^3 \quad (d=2).
\]
The general $d$-dimensional statement is the move from $O(P\cdot L)=O\!\big((p+1)^{2d}\big)$
to $O\!\big(d\,(p+1)^{d+1}\big)$---in the notation of the operator counts, a drop
from $O(P\cdot L)$ to $O\!\big(d\cdot P^{1/d}\cdot L\big)$, since a single axis of
the quad-point grid has length $P^{1/d}=(p+1)$.

\Cref{fig:mf-sumfac} draws the reordering for $d=2$, and \cref{tab:mf-sumfac}
tabulates the savings. The win grows with order: at $p=8$ in 3-D it is the
difference between $(p+1)^6\approx\num{531441}$ and
$3\cdot(p+1)^4\approx\num{19683}$ operations per element---a factor of nearly
$30$, and it widens as $p$ climbs.

\begin{figure}[t]
\centering
\begin{tikzpicture}[node distance=6mm and 14mm]
  % naive: one big dense contraction L -> P
  \node[data, text width=15mm] (din) {dofs\\$(p{+}1)^2$};
  \node[opbox, fill=simBgRust, draw=simRust, right=of din, text width=24mm]
    (naive) {dense $P\times L$\\\footnotesize $(p{+}1)^4$ ops};
  \node[data, right=of naive, text width=16mm] (dpts) {pts\\$(p{+}1)^2$};
  \draw[flow] (din) -- (naive);
  \draw[flow] (naive) -- (dpts);
  \node[font=\scriptsize\bfseries, simRust, left=1mm of din, rotate=90, anchor=south]
    {naive};

  % factored: two 1-D sweeps
  \node[data, below=14mm of din, text width=15mm] (fin) {dofs\\$(p{+}1)^2$};
  \node[opbox, right=of fin, text width=20mm] (s1)
    {1-D sweep $\eta$\\\footnotesize $(p{+}1)^3$ ops};
  \node[opbox, right=of s1, text width=20mm] (s2)
    {1-D sweep $\xi$\\\footnotesize $(p{+}1)^3$ ops};
  \node[data, right=of s2, text width=16mm] (fpts) {pts\\$(p{+}1)^2$};
  \draw[flow] (fin) -- (s1);
  \draw[flow] (s1) -- node[above, font=\scriptsize, simGray] {partial} (s2);
  \draw[flow] (s2) -- (fpts);
  \node[font=\scriptsize\bfseries, simBlue, left=1mm of fin, rotate=90, anchor=south]
    {factored};
\end{tikzpicture}
\caption[Sum-factorization on a 2-D element]{Basis evaluation on a 2-D
tensor-product element, two ways. \emph{Top (rust):} the naive contraction treats
the basis as one dense $P\times L$ matrix, costing $(p+1)^4$ operations per
element. \emph{Bottom (blue):} sum-factorization applies the one-dimensional
basis along each axis in turn---a sweep in $\eta$, then a sweep in $\xi$---each
costing $(p+1)^3$, for $2(p+1)^3$ total. The two produce \emph{identical}
results; only the contraction order differs. The saving is the move from
$O(P\cdot L)$ to $O(d\,P^{1/d}L)$, and it grows with order.}
\label{fig:mf-sumfac}
\end{figure}

\begin{table}[t]
\centering
\caption[Sum-factorization operation counts]{Multiply-adds per element for the
basis-evaluation stage, naive versus sum-factored, on a tensor-product element of
order $p$. ``Naive'' is the dense $P\times L$ contraction, $(p+1)^{2d}$;
``factored'' is $d\,(p+1)^{d+1}$. The speedup is unbounded in $p$. (Values for
$d=2$; the 3-D column at $p=8$ illustrates how steep the win becomes.)}
\label{tab:mf-sumfac}
\small
\setlength{\tabcolsep}{8pt}
\renewcommand{\arraystretch}{1.2}
\begin{tabular}{@{}cccc@{}}
\toprule
order $p$ & naive $(p{+}1)^4$ & factored $2(p{+}1)^3$ & speedup \\
\midrule
1 & 16     & 16    & $1.0\times$ \\
2 & 81     & 54    & $1.5\times$ \\
4 & 625    & 250   & $2.5\times$ \\
8 & 6561   & 1458  & $4.5\times$ \\
\midrule
\multicolumn{4}{@{}l@{}}{\footnotesize 3-D at $p=8$: naive $9^6=\num{531441}$
vs.\ factored $3\cdot 9^4=\num{19683}$ \quad($27\times$)} \\
\bottomrule
\end{tabular}
\end{table}

\begin{workedexample}{Counting operations on a $Q_2$ element}{mf-sumfac}
\textbf{Setup.} Take a 2-D second-order ($Q_2$, $p=2$) tensor-product element.
The local dofs sit on a $3\times3$ node grid, so $L=(p+1)^2=9$; we use a $3\times3$
Gauss grid, $P=9$. We count the multiply-adds to evaluate the field values at all
quad points (the basis stage $B$, mode = interp).

\tcblower
\textbf{Naive (dense $P\times L$).} One element's $L=9$ dofs contract against the
dense $9\times9$ basis matrix:
\[
  P\cdot L \;=\; 9\cdot 9 \;=\; 81 \ \text{multiply-adds}.
\]

\textbf{Factored (two 1-D sweeps).} The $Q_2$ basis is
$\phi_{ij}(\xi,\eta)=\phi_i(\xi)\phi_j(\eta)$ with $\phi_i,\phi_j$ the three 1-D
quadratic Lagrange functions.
\emph{Sweep 1 (along $\eta$):} for each of the $3$ values of $\xi$-index, contract
the $3$ dofs in $\eta$ against the $3\times3$ 1-D basis, producing partial values
at the $3$ $\eta$-quad-points: $3\ (\xi\text{ cols})\times 3\times 3 = 27$.
\emph{Sweep 2 (along $\xi$):} for each of the $3$ $\eta$-quad-points, contract the
$3$ partials in $\xi$ against the $3\times3$ 1-D basis: $3\times 3\times3 = 27$.
\[
  \text{total} \;=\; 27 + 27 \;=\; 54 \ \text{multiply-adds}.
\]

\textbf{Comparison.} $81$ versus $54$: a $1.5\times$ saving at $p=2$, matching the
$d(p+1)^{d+1}=2\cdot 27=54$ formula. The two compute \emph{bit-for-bit the same
quad-point values}---sum-factorization only reorders the sums. The modest factor
here grows fast: at $p=4$ it is $625$ vs.\ $250$ ($2.5\times$), and in 3-D the
exponent gap ($2d$ vs.\ $d+1$) makes it decisive.
\end{workedexample}

\begin{pitfall}
Sum-factorization is \emph{transparent}: same result, faster. Do not confuse it
with a \emph{load-bearing} numerical trick (\cref{ch:bigidea})---it does not
change the answer, the conditioning, or the rounding behavior in any way the
algorithm depends on. In the spec we therefore record it as a one-line note on
the basis stage, not as a separate algebraic claim. The reason to know about it
anyway is that it is \emph{why} high-order matrix-free is fast enough to be the
default: without it, the basis stage would reintroduce the very $O(p^{2d})$ cost
the matrix-free form was trying to escape.
\end{pitfall}

\section{The implementation view: contractions all the way down}
\label{sec:mf-impl-view}

Step back from the details and look at what \cref{eq:mf-chain} \emph{is}. The
finite-element operator $\mat{A}$---the discretization of Maxwell's weak form, the
object every solver in the book applies---has turned out to be a composition of
tensor contractions over named axes. Gather is an indexed copy. Evaluate and
integrate are dense contractions of an $[E,L]$ tensor against a tabulated basis,
yielding $[E,P,C]$. Weight is an elementwise multiply of $[E,P,C]$ against
$[E,P,G]$. Scatter is an indexed sum. There is no sparse-matrix data structure
anywhere in sight---no row pointers, no column indices, no stored entries. There
are only tensors and contractions.

This is the destination Part~I promised. \Cref{ch:bigidea} argued that a field
\emph{is} a tensor and that a simulation is a chain of pure transformations of
immutable tensors---and that this description is exactly the form a GPU or tensor
accelerator wants: whole-tensor operations over named axes, no aliasing, no
ordering hazards. The matrix-free operator is that thesis made fully concrete.
Every stage of \cref{fig:mf-pipeline} is a pure function from tensors to tensors;
the whole operator is their composition; and the shapes
(\lstinline[style=l4style]|[E, L]|, \lstinline[style=l4style]|[E, P, C]|,
\lstinline[style=l4style]|[E, P, G]|) are precisely what a backend schedules onto
its parallel units. The assembled-matrix form, by contrast, is a CPU-era data
structure---irregular, pointer-chasing, bandwidth-bound---that a tensor backend
would have to be coaxed into. Matrix-free is the form the accelerator was built
for.

\begin{keyidea}
The matrix-free operator \emph{is} the implementation view of \cref{ch:bigidea},
realized for finite elements. A field is a tensor; the operator is a composition
of contractions over named axes; the whole thing maps onto a tensor backend with
no sparse data structure at all. When the book's final part renders the
synthesized library, the matrix-free operator is what the
``\,\lstinline[style=l4style]|apply A v|\,'' of every solver expands into.
\end{keyidea}

\subsection{The innermost kernel is a library boundary}
\label{sec:mf-extern}

There is one honest qualification. We have described the five stages as pure
tensor contractions, and at the mathematical level they are. But the innermost
pointwise kernel---the per-quadrature-point weight stage $D$, fused with the basis
contractions for efficiency---is, in the real Palace implementation, handed to a
specialized accelerator library (libCEED) that owns the device-level details:
which loops to fuse, how to tile the quadrature points, how to map them onto warps
or wavefronts. The book does not re-derive that kernel; it treats it as an
\emph{opaque library boundary}. In the calculus, such a boundary is written with
the marker \lstinline[style=l4style]|#extern|: a node whose \emph{contract} (its
shape signature and semantics) we state precisely, but whose \emph{body} we
delegate to the library.

\begin{lstlisting}[style=l4style]
-- the innermost element-quadrature kernel, as a library boundary:
-- the contract is ours; the device implementation is libCEED's.
quad_kernel :: GeomData -> Tensor[E, P, C] -> Tensor[E, P, C]
#extern quad_kernel    -- shape + semantics specified; body delegated
\end{lstlisting}

\noindent This is the same \lstinline[style=l4style]|#extern| node that appears,
violet and dashed, at the center of \cref{fig:mf-pipeline}. Two facts make the
boundary clean. First, its contract is exactly the pointwise diagonal of
\cref{sec:mf-diagonal}: a shape-preserving elementwise multiply against
\texttt{geom\_data}, with semantics we have fully specified. Second, because the
contract is so simple and so completely characterized, we can reason about the
\emph{whole} operator---its symmetry, its action, the worked examples
above---without opening the kernel at all. The library owns the machine; we own
the mathematics, and the two meet at a narrow, well-typed interface.

\begin{notationbox}
The \lstinline[style=l4style]|#extern| marker recurs wherever the spec touches an
opaque accelerator or solver library: the libCEED quadrature kernel here, the
SLEPc eigensolve in \cref{ch:eigenvalues}, the sparse triangular solves inside a
smoother in \cref{ch:smoothers}. In each case the pattern is the same---a
precisely specified \emph{contract} (signature and semantics) standing in for a
\emph{body} we delegate. The discipline is to make the boundary so well-typed that
nothing above it needs to look inside.
\end{notationbox}

\section{When matrix-free is not the answer}
\label{sec:mf-tradeoffs}

Matrix-free is the right default at moderate-to-high order on accelerators, but it
is a trade, not a free lunch, and it pays to name the other side of the ledger.
\Cref{fig:mf-vs-assembled} sets the two representations side by side: the same
linear map, stored two ways. The assembled form holds an explicit sparse data
structure---rows, columns, stored entries---and applies in one multiply; the
matrix-free form holds only the build-time tables and recomputes the action
through the five-stage pipeline each time. The choice between them is the subject
of this section.

\begin{figure}[t]
\centering
\begin{tikzpicture}[node distance=5mm and 10mm]
  % ---- LEFT: assembled (stored sparse matrix) ----
  \begin{scope}
    \node[font=\bfseries\small, simRust] at (0,2.1) {Assembled};
    % a sparse matrix glyph
    \draw[simRust, thick, fill=simBgRust, rounded corners=1pt] (-1.0,-1.1) rectangle (1.0,1.1);
    \foreach \i/\j in {0/0,1/0,0/1,1/1,2/2,3/2,2/3,3/3,1/2,2/1,4/4,3/4,4/3,0/4,4/0}{
      \fill[simRust] ({-0.9+0.38*\j},{0.9-0.38*\i}) rectangle ({-0.66+0.38*\j},{0.66-0.38*\i});}
    \node[font=\scriptsize\itshape, simGray, below=1mm of {(0,-1.1)}] (la) at (0,-1.45)
      {stored entries $A_{ij}$};
    \node[data, fill=simBgRust, draw=simRust, below=2mm of la, text width=30mm]
      (lapp) {apply: one sparse multiply\\\footnotesize cheap per-apply, $O(p^{2d})$ memory};
  \end{scope}
  % ---- RIGHT: matrix-free (pipeline, nothing stored) ----
  \begin{scope}[shift={(6.4,0)}]
    \node[font=\bfseries\small, simBlue] at (0,2.1) {Matrix-free};
    % mini pipeline glyph
    \foreach \k/\lab in {0/$G$,1/$B$,2/$D$,3/$B^{\mathsf T}$,4/$G^{\mathsf T}$}{
      \node[draw=simBlue, fill=simBgBlue, rounded corners=1pt, minimum size=4.5mm,
            font=\scriptsize] (p\k) at ({-1.4+0.7*\k},0.6) {\lab};}
    \foreach \k in {0,1,2,3}{\pgfmathtruncatemacro{\next}{\k+1}\draw[flow] (p\k)--(p\next);}
    \node[font=\scriptsize\itshape, simGray] at (0,-0.5) {nothing stored but build-time tables};
    \node[data, fill=simBgBlue, draw=simBlue, text width=30mm] (rapp) at (0,-1.55)
      {apply: rerun the pipeline\\\footnotesize more arithmetic, $O(p)$ memory};
  \end{scope}
  % ---- center: same operator ----
  \node[font=\scriptsize\bfseries, simGreen] at (3.2,1.4) {same};
  \node[font=\scriptsize\bfseries, simGreen] at (3.2,1.05) {$\mat{A}$};
  \draw[{Stealth[length=2mm]}-{Stealth[length=2mm]}, simGreen, thick] (1.3,0.2) -- (4.5,0.2);
\end{tikzpicture}
\caption[Matrix-free versus assembled, side by side]{Two representations of the
\emph{same} operator $\mat{A}$. \emph{Left (rust):} the assembled form stores an
explicit sparse matrix of entries $A_{ij}$ and applies in a single sparse
multiply---cheap per application but $O(p^{2d})$ memory per element.
\emph{Right (blue):} the matrix-free form stores nothing but the build-time
tables and recomputes the action through the five-stage pipeline on each apply---more
arithmetic but only $O(p)$ memory. They compute the identical map (green); the
choice is an implementation decision driven by order, hardware, and downstream use.}
\label{fig:mf-vs-assembled}
\end{figure}

First, \emph{cost per application is higher}. The assembled matrix, once built,
applies in a single sparse multiply; the matrix-free operator reruns all five
stages every time. When the operator is applied only a handful of times, or when
order is so low that the stored matrix is tiny, assembly wins. The crossover in
\cref{fig:mf-cost} is a memory crossover; there is a separate, lower, arithmetic
crossover.

Second, \emph{preconditioning is harder}. Many classical preconditioners
(incomplete factorizations, algebraic multigrid setup) want to inspect the
operator's entries---exactly what the matrix-free form refuses to provide. This is
not fatal: the geometric multigrid of \cref{ch:multigrid} and the smoothers of
\cref{ch:smoothers} are designed to need only the action, and where an entry-level
view is unavoidable, Palace can materialize a coarse or diagonal approximation on
demand (apply the pipeline to the identity columns and read off the entries). But
it is real friction, and it is why the matrix-free and assembled forms coexist in
the same code rather than one simply replacing the other.

\begin{remark}
Both forms compute the \emph{same operator}---they are two evaluations of the same
linear map, and the worked example confirmed they agree to the last digit. The
choice between them is an implementation decision driven by order, hardware, and
how the operator will be used downstream, not a change to the mathematics. This is
the value-threading lesson of \cref{ch:bigidea} once more: the description is free
to choose the representation the machine likes best, because the answer is
invariant to the choice.
\end{remark}

\summarysection
A Krylov or eigensolver asks an operator for one thing only: its action,
$\vv\mapsto\mat{A}\vv$. The \emph{matrix-free} operator supplies that action
without ever forming or storing $\mat{A}$, which matters because the stored
operator costs $O(p^{2d})$ memory per element and becomes unaffordable at high
order on accelerators. The action factors into five element-local stages,
$\mat{A}\vv = G^{\mathsf T}B_{\diffop}^{\mathsf T}\,D\,B_{\diffop}\,G\,\vv$:
\emph{gather} the global vector to per-element blocks
(\lstinline[style=l4style]|[N]|$\to$\lstinline[style=l4style]|[E, L]|),
\emph{evaluate} the basis at quadrature points
(\lstinline[style=l4style]|[E, L]|$\to$\lstinline[style=l4style]|[E, P, C]|),
\emph{weight} each point by a precomputed geometry/material factor (pointwise on
\lstinline[style=l4style]|[E, P, C]|, the embarrassingly parallel diagonal),
\emph{integrate} back to local dofs, and \emph{scatter-add} to a fresh global
vector. The middle three stages are block-diagonal in the element axis $E$---fully
parallel; all inter-element coupling lives in the gather/scatter ends. The
geometry factor \texttt{geom\_data} is a \emph{build-time} quantity, computed once
per $(\text{mesh},\text{order})$ and reused across the hundreds of applications a
solve performs. On tensor-product elements, \emph{sum-factorization} evaluates the
basis dimension by dimension, dropping the cost from $O(P\cdot L)$ to
$O(d\,P^{1/d}L)$---a transparent reordering, same answer, decisive at high order.
A worked example traced $\mat{M}\vv$ through all five stages on two 1-D elements
and recovered the assembled matrix's answer exactly. Because every stage is a
tensor contraction over named axes with no sparse-matrix structure, the
matrix-free operator \emph{is} the implementation view of \cref{ch:bigidea}---the
form a GPU consumes---with its innermost kernel sealed behind a precisely typed
\lstinline[style=l4style]|#extern| library boundary.

\furtherreading
The matrix-free, sum-factorized, tensor-contraction view of high-order
finite-element operators is the design philosophy of the CEED project; its
co-design survey~\citep{ceed2021} is the entry point to the libraries (libCEED,
MFEM) that realize the five-stage decomposition described here, and to the
performance literature behind sum-factorization. For the underlying
finite-element method---reference elements, tabulated bases, quadrature, and the
element matrices the operator factors---Jin's text on the finite-element method in
electromagnetics~\citep{jin2014} gives the field-theory grounding in the
context of Maxwell's equations specifically. The classical operation-count
argument for sum-factorization, and the spectral-element setting in which it was
first decisive, are developed in any spectral/high-order finite-element reference;
the build-time/run-time stratification it relies on is treated in its own right in
\cref{ch:strata}.

\exercisessection
\begin{enumerate}[nosep]
  \item \textbf{Memory at a glance.} For a 3-D hexahedral element of order $p=4$,
  how many entries does the dense element matrix $\mat{A}^e$ hold? How many
  numbers does the matrix-free form store for that element if it keeps a tabulated
  1-D basis of size $(p+1)^2$ and a geometry block of $G=2+d^2$ components at each
  of $(p+1)^3$ quadrature points? Which is larger, and by what factor?

  \item \textbf{Shapes along the pipeline.} A driven-analysis curl--curl operator
  ($\diffop=\Curl$) runs on a mesh with $E=\num{1000}$ hexahedral elements of
  order $p=2$, with $P=27$ quadrature points per element and $C=3$ curl components.
  Write the tensor shape on each of the five wires of \cref{fig:mf-pipeline},
  from the input \lstinline[style=l4style]|[N]| to the output
  \lstinline[style=l4style]|[N]|, filling in concrete extents for $E$, $L$, $P$,
  $C$. (You will need $L=(p+1)^3$.)

  \item \textbf{The shared-dof bookkeeping.} In
  \cref{wex:mf-mass}, the node $x_1$ received contributions
  $0.8333$ (from $e_0$) and $1.1667$ (from $e_1$), summing to $2.0000$. Explain,
  in terms of $G$ and $G^{\mathsf T}$, why this \emph{sum} is the correct global
  value and not, say, an average. What property of $G^{\mathsf T}$ as the
  transpose of $G$ guarantees it?

  \item \textbf{Counting sum-factorization.} Repeat the operation count of
  \cref{wex:mf-sumfac} for a 3-D $Q_2$ element ($d=3$, $p=2$, so $L=P=27$).
  Give the naive count $P\cdot L$ and the sum-factored count $d\,(p+1)^{d+1}$, and
  report the speedup. Then redo it for $p=4$ and comment on how the speedup scales.

  \item \textbf{The diagonal carries the physics.} Two terms---a mass term
  ($\diffop=\mathrm{Id}$, $Q=\varepsilon$) and a stiffness term ($\diffop=\Grad$,
  $Q=\varepsilon$)---are applied on the \emph{same} mesh. Which of the five stages
  differ between them, and which are identical? In particular, what changes in the
  contents (not the shape) of \texttt{geom\_data}, and what changes in the basis
  mode? Relate your answer to why a single five-stage skeleton serves all six
  analyses of \cref{ch:targets}.

  \item \textbf{Build-time amortization.} A conjugate-gradient solve
  (\cref{ch:cg}) converges in $200$ iterations, each requiring one operator
  application. The build pass that produces \texttt{geom\_data} costs roughly as
  much as $5$ operator applications. What fraction of the total operator-related
  work is the one-time build? Now suppose adaptive refinement (\cref{ch:amr})
  rebuilds \texttt{geom\_data} after every $50$ iterations. How does the fraction
  change, and what does this say about when matrix-free's amortization advantage
  is strongest?

  \item \textbf{Why the solver never asks for an entry.} Conjugate gradients
  (\cref{ch:cg}) forms, each iteration, the quantities $\mat{A}\vp$,
  $\inner{\vp}{\mat{A}\vp}$, and vector updates $\vx\leftarrow\vx+\alpha\vp$.
  Identify which of these touch the operator, and confirm that none of them
  requires an individual entry $A_{ij}$. What would have to change about the
  \emph{solver} for the matrix-free form to become unusable?

  \item \textbf{The \texttt{\#extern} boundary.} The innermost kernel is sealed
  behind \lstinline[style=l4style]|#extern quad_kernel| with the contract
  \lstinline[style=l4style]|GeomData -> Tensor[E, P, C] -> Tensor[E, P, C]|.
  State, in one sentence each, (a) the semantic property of this kernel that lets
  us reason about the symmetry of $\mat{A}=G^{\mathsf T}B^{\mathsf T}DBG$ without
  opening the kernel, and (b) one reason the implementation might want to
  \emph{fuse} this kernel with the adjacent basis stages even though, on paper,
  $D$ is a separate stage.
\end{enumerate}

\chapter{Boundary Conditions and the Helmholtz Decomposition}
\label{ch:bc}

\chapterorient{The assembly machinery of the last three chapters handed us a
linear system $\mat{K}\vx=\vb$ built term by term from the weak form. But two
things stand between that raw system and a physically correct answer. First, the
boundary conditions: the weak form of \cref{ch:weak} split them into
\emph{essential} and \emph{natural}, and the essential half---the voltages on
conductors, the tangential field on metal walls---was deferred to ``the discrete
setting.'' This is that setting. We make essential conditions concrete as an
operation on the assembled matrix: identify the constrained degrees of freedom,
zero their rows and columns, and lift any nonzero boundary data into the
right-hand side. Second, the fields the solver produces may not be physical: the
curl--curl operator has an enormous null space of gradient fields, and an
eigensolver left to its own devices will happily return them. The cure is a
\emph{projection}---the discrete Helmholtz decomposition---that strips the
gradient part off a field and keeps only the divergence-free remainder. Both are
\emph{finishing operations}: post-compositions that act on the assembled system,
cleanly separable from the assembly fold that produced it. They close Part~II.}

\section{Two finishing operations}

The previous chapters built an operator. Out of the mesh, the materials, and a
short list of weak-form terms $(Q,\diffop)$, the assembly fold of
\cref{ch:assembly} produced a sparse matrix $\mat{K}$ and, where there is a
source, a load vector $\vb$. We are nearly ready to solve $\mat{K}\vx=\vb$. Two
operations remain, and this chapter is about both. They are different in
character but share a structural signature, and naming that shared signature is
the first order of business.

The first operation \emph{imposes the essential boundary conditions}. The weak
form of \cref{ch:weak} taught us that boundary conditions come in two flavors.
\emph{Natural} (Neumann/impedance) conditions take care of themselves: they enter
as boundary integrals during assembly, and a homogeneous one contributes nothing
at all---you get it for free by doing nothing. \emph{Essential} (Dirichlet)
conditions are the work. They fix the value of the solution on part of the
boundary---a conductor held at a known voltage, the tangential electric field
forced to zero on a perfectly conducting wall---and that constraint must be
\emph{built into the discrete system by hand}, after assembly, by operating on
the rows and columns of $\mat{K}$ that belong to the constrained unknowns.

The second operation \emph{projects a field onto the physical subspace}. Some of
the operators we assemble---the curl--curl operator above all---are singular:
they annihilate an entire subspace of fields, the gradients. A solver handed such
an operator can drift into that null space and return a field that is
mathematically a solution but physically meaningless. The remedy is to project
the iterate, at each step, onto the subspace orthogonal to the null space---the
\emph{divergence-free} fields. That projection is the discrete realization of a
classical theorem, the Helmholtz decomposition, and it is the second half of the
chapter.

\begin{figure}[t]
\centering
\begin{tikzpicture}[node distance=9mm]
  \node[stage, text width=26mm, align=center] (asm)
    {\textbf{assemble}\\[1pt]\footnotesize fold over terms\\\scriptsize(\cref{ch:assembly})};
  \node[product, right=12mm of asm, text width=22mm, align=center] (raw)
    {raw system\\[1pt]\footnotesize $\mat{K}\vx=\vb$};
  \node[opbox, right=12mm of raw, text width=30mm] (bc)
    {\textbf{impose essential BC}\\[1pt]\scriptsize zero rows/cols, lift RHS};
  \node[opbox, below=11mm of bc, text width=30mm] (proj)
    {\textbf{divergence-free project}\\[1pt]\scriptsize strip the gradient part};
  \node[stage, right=12mm of bc, text width=20mm, align=center, fill=simBgGreen, draw=simGreen] (solve)
    {\textbf{solve}\\[1pt]\footnotesize (\cref{ch:krylov})};
  \draw[flow] (asm) -- (raw);
  \draw[flow] (raw) -- (bc);
  \draw[flow] (bc) -- (solve);
  \draw[flow] (proj.east) -- ++(6mm,0) |- (solve.west);
  \draw[refedge] (raw.south) |- (proj.west);
  \node[font=\scriptsize\itshape, simGray, below=1mm of asm.south] {the fold};
  \node[font=\scriptsize\itshape, simRust] at ($(bc.north)+(0,5mm)$) {finishing operations (this chapter)};
\end{tikzpicture}
\caption[Where the finishing operations sit]{The two finishing operations of this
chapter sit \emph{between} the assembly fold and the solve. Essential-boundary
elimination acts on the assembled operator $\mat{K}$ and its right-hand side
$\vb$; the divergence-free projection acts on field iterates during (or before)
the solve. Both are \emph{post-compositions}: separable from the assembly that
produced the raw system, they neither inspect nor depend on which weak-form terms
were folded in.}
\label{fig:finishing}
\end{figure}

What unifies them is that both are \emph{post-compositions} on the assembled
system (\cref{fig:finishing}). Neither is part of the assembly fold. Essential
elimination consumes the finished operator $\mat{K}$ and does not care which
weak-form terms went into it; the divergence-free projection consumes a field and
a few auxiliary operators and likewise does not crack open the assembly. This
separability is not an accident of implementation---it is the reason the
simulator can assemble once and then apply boundary conditions and projections as
modular, swappable stages. We will make the claim precise for each operation in
turn.

\begin{keyidea}
The assembled system $\mat{K}\vx=\vb$ is not yet ready to solve. Two
\emph{finishing operations} stand between assembly and solution, and both are
\emph{separable post-compositions} on the assembled operator---not steps inside
the assembly fold. \textbf{(1)~Essential-boundary elimination} forces the
solution to take prescribed values on the Dirichlet boundary, by zeroing the
constrained rows and columns of $\mat{K}$ and lifting the boundary data into
$\vb$. \textbf{(2)~Divergence-free projection} removes the unphysical gradient
component of a field, the discrete Helmholtz decomposition, so that singular
operators like curl--curl are solved on the subspace where they are well-posed.
\end{keyidea}

\section{Essential boundary conditions on the discrete system}
\label{sec:essential}

\subsection{The essential degrees of freedom}

Recall from \cref{ch:galerkin} that the discrete unknown $\vx$ is a vector of
\emph{degrees of freedom} (DoFs): one number per basis function, the coefficient
that scales that basis function in the finite-element expansion of the solution.
For an $\Hone$ (nodal) space, a DoF is the value of the potential at a mesh
vertex; for an $\Hcurl$ (Nédélec) space, a DoF is the circulation of the field
along a mesh edge. The full set of DoFs indexes the rows and columns of
$\mat{K}$.

An essential boundary condition pins some of those DoFs to known values. On a
conductor held at voltage $V$, every $\Hone$ DoF whose vertex lies on that
conductor must equal $V$. On a perfectly conducting (PEC) wall, the tangential
electric field vanishes, which means every $\Hcurl$ DoF whose edge lies in that
wall must equal zero (the edge circulation of a field with no tangential
component is zero). In both cases a geometrically identifiable subset of the DoFs
is \emph{constrained}: its values are dictated by the boundary condition, not by
the equation.

\begin{definition}[Essential and free DoFs]
\label{def:essdof}
Let the full DoF index set be $\{0,1,\dots,N-1\}$. The \emph{essential} (or
constrained) DoF set $E$ is the set of indices whose values are prescribed by an
essential boundary condition. The \emph{free} DoF set is its complement
$F = \{0,\dots,N-1\}\setminus E$. The two are disjoint and exhaust the index set:
$E\cup F = \{0,\dots,N-1\}$, $E\cap F=\emptyset$.
\end{definition}

The map from boundary geometry to the essential set $E$---walk the boundary
elements carrying the essential attribute, collect every DoF they touch---is a
purely combinatorial pre-pass, run once before any matrix is touched. We treat
$E$ as given. \Cref{fig:partition} shows the partition on a small mesh: the
boundary DoFs (essential, dark) form a thin shell, and the interior DoFs (free,
light) are everything else.

\begin{figure}[t]
\centering
\begin{tikzpicture}[scale=1.0]
  % domain
  \fill[simBgBlue] (0,0) rectangle (5,3);
  % grid lines
  \foreach \x in {0,1,2,3,4,5}{\draw[simGray!40] (\x,0)--(\x,3);}
  \foreach \y in {0,1,2,3}{\draw[simGray!40] (0,\y)--(5,\y);}
  % boundary (essential) nodes: dark, on the perimeter
  \foreach \x in {0,1,2,3,4,5}{
    \fill[simBlue] (\x,0) circle (2.6pt);
    \fill[simBlue] (\x,3) circle (2.6pt);}
  \foreach \y in {1,2}{
    \fill[simBlue] (0,\y) circle (2.6pt);
    \fill[simBlue] (5,\y) circle (2.6pt);}
  % interior (free) nodes: light
  \foreach \x in {1,2,3,4}{
    \foreach \y in {1,2}{
      \fill[simRust] (\x,\y) circle (2.6pt);}}
  % labels
  \node[simBlue, font=\small, below left] at (0,0) {};
  \node[font=\small\bfseries, simBlue] at (2.5,3.45) {essential DoFs $E$ (boundary)};
  \node[font=\small\bfseries, simRust] at (2.5,1.5) {};
  \node[font=\small\bfseries, simRust, anchor=west] at (5.4,1.5) {free DoFs $F$ (interior)};
  \draw[refedge] (5.35,1.5) -- (4,1.5);
  % a constrained value callout
  \node[font=\scriptsize, simBlue, anchor=east] at (-0.3,1.5) {$x_i=V$};
  \draw[refedge] (-0.3,1.5) -- (0,1);
\end{tikzpicture}
\caption[The free/essential DoF partition]{The degrees of freedom on a small
$\Hone$ mesh split into two disjoint groups. The \emph{essential} DoFs $E$ (blue,
on the boundary) carry prescribed values---here a conductor voltage $x_i=V$---and
are not solved for. The \emph{free} DoFs $F$ (rust, interior) are the genuine
unknowns. Imposing the essential condition means reducing the system to the free
block, with the influence of the prescribed boundary values moved to the
right-hand side.}
\label{fig:partition}
\end{figure}

\subsection{Scalar value versus tangential trace}
\label{sec:scalarvstrace}

The two physical pictures---a conductor at fixed voltage and a PEC wall forcing
tangential $\vE$ to zero---feel different, but they produce essential DoF sets in
exactly the same way, once you remember what a DoF \emph{is} in each space. The
difference is entirely in the trace operator the boundary condition constrains.

In an $\Hone$ (nodal) space the relevant trace is the boundary \emph{value} of the
function. A DoF is a point value at a vertex, so ``$u=u_0$ on the conductor''
constrains every DoF whose vertex lies on the conductor. In an $\Hcurl$ (Nédélec)
space the relevant trace is the \emph{tangential component} of the field on the
boundary surface. A DoF is an edge circulation $\int_e \vu\cdot\mathrm{d}\boldsymbol{\ell}$,
and an edge lying \emph{in} a boundary face sees only the tangential part of the
field; so ``tangential $\vE=0$ on the PEC wall'' constrains every edge DoF whose
edge lies in that wall. In both cases the recipe is the same: walk the boundary
entities carrying the essential attribute, collect the DoFs they own, and add them
to $E$. That the constrained quantity is a point value in one space and a
tangential circulation in the other never surfaces in the elimination---by the
time \texttt{eliminate\_essential\_bc} runs, $E$ is just a list of indices.

\begin{figure}[t]
\centering
\begin{tikzpicture}[scale=1.0]
  % --- H1 nodal: value at vertices ---
  \begin{scope}[shift={(0,0)}]
    \node[font=\small\bfseries, simBlue] at (1.5,2.9) {$\Hone$: nodal value};
    \draw[simGray!50] (0,0) rectangle (3,2.2);
    \draw[simGray!40] (1.5,0)--(1.5,2.2);
    \draw[simGray!40] (0,1.1)--(3,1.1);
    % bottom boundary face highlighted
    \draw[simBlue, line width=2pt] (0,0)--(3,0);
    % boundary vertices (essential): on the bottom edge
    \foreach \x in {0,1.5,3}{\fill[simBlue] (\x,0) circle (3pt);}
    % interior vertices
    \fill[simRust] (1.5,1.1) circle (3pt);
    \foreach \x in {0,1.5,3}{\fill[simRust] (\x,2.2) circle (2.4pt);}
    \foreach \x in {0,3}{\fill[simRust] (\x,1.1) circle (2.4pt);}
    \node[font=\scriptsize, simBlue, below] at (1.5,-0.15) {value $u=u_0$ pinned};
    \node[font=\scriptsize, simRust] at (2.3,1.35) {free node};
  \end{scope}
  % divider
  \draw[simGray,densely dashed] (4.0,-0.5)--(4.0,3.0);
  % --- Hcurl edge: tangential circulation ---
  \begin{scope}[shift={(5.2,0)}]
    \node[font=\small\bfseries, simGreen] at (1.5,2.9) {$\Hcurl$: edge circulation};
    \draw[simGray!50] (0,0) rectangle (3,2.2);
    \draw[simGray!40] (1.5,0)--(1.5,2.2);
    \draw[simGray!40] (0,1.1)--(3,1.1);
    % bottom boundary face highlighted
    \draw[simGreen, line width=2pt] (0,0)--(3,0);
    % boundary edges (essential): the two bottom edges -> arrows along them
    \draw[-{Stealth[length=2mm]}, simGreen, very thick] (0.1,0)--(1.4,0);
    \draw[-{Stealth[length=2mm]}, simGreen, very thick] (1.6,0)--(2.9,0);
    % an interior edge (free)
    \draw[-{Stealth[length=2mm]}, simRust, very thick] (1.5,1.2)--(1.5,2.1);
    \node[font=\scriptsize, simGreen, below] at (1.5,-0.15) {tangential $\vu\cdot\boldsymbol{\ell}=0$};
    \node[font=\scriptsize, simRust, anchor=west] at (1.6,1.6) {free edge};
  \end{scope}
\end{tikzpicture}
\caption[Essential DoFs for the scalar and vector spaces]{The essential DoF set is
collected the same way in both function spaces; only the constrained trace
differs. In an $\Hone$ nodal space (left) a DoF is a point value at a vertex, so a
Dirichlet value condition $u=u_0$ pins every vertex DoF on the boundary face
(blue). In an $\Hcurl$ Nédélec space (right) a DoF is the circulation along an
edge, so a tangential-trace condition (PEC wall, tangential $\vE=0$) pins every
\emph{edge} DoF lying in the boundary face. In both cases the constrained DoFs are
collected into $E$ and eliminated identically; the elimination never knows which
kind of trace was constrained.}
\label{fig:trace}
\end{figure}

\Cref{fig:trace} sets the two side by side. The practical consequence is uniformity:
the same \texttt{eliminate\_essential\_bc} and \texttt{eliminate\_rhs} serve the
electrostatic (nodal) and magnetostatic/eigenmode (edge) pipelines without
modification, because the only space-specific work---deciding which DoFs are
constrained---happens in the combinatorial pre-pass that builds $E$, upstream of
the matrix surgery.

\subsection{Why this is a post-composition, not part of assembly}

Here is the central structural point, and it echoes the slogan of
\cref{fig:finishing}. The assembly fold of \cref{ch:assembly} knows nothing about
boundary conditions. It walks the list of weak-form terms, and for each term it
adds an element contribution into the global matrix; it would produce the very
same $\mat{K}$ whether or not any DoF is later constrained. Essential elimination
is applied \emph{afterward}, to the finished matrix, as a separate operation:
\[
  \mat{K}_{\mathrm{BC}} \;=\; \texttt{eliminate\_essential\_bc}\bigl(\mat{K},\,E,\,\text{policy}\bigr),
  \qquad
  \mat{K} = \texttt{assemble}(\text{space},\,\text{terms}).
\]
The elimination consumes $\mat{K}$ as an opaque square operator together with the
index set $E$. It never asks which terms were folded to build $\mat{K}$, nor in
what representation it is stored. This is what it means to be a \emph{separable
post-composition}: it factors cleanly behind the assembly, and the two compose
without either needing to know the other's internals.

\begin{pitfall}
A natural instinct is to ``bake the boundary conditions into assembly''---to skip
the constrained DoFs while folding, so that $\mat{K}$ comes out already reduced.
Resist it. Keeping elimination as a separate post-composition is what lets the
\emph{same} assembled operator feed a Dirichlet solve, a Neumann solve, or an
eigenvalue problem, each with its own essential set and its own diagonal policy.
It also lets us reason about the two stages independently: assembly is a fold;
elimination is a projection. Fusing them couples two concerns that have no
business being coupled.
\end{pitfall}

\subsection{Eliminating the essential DoFs}

Now the operation itself. Partition the unknown vector and the matrix along the
free/essential split of \cref{def:essdof}. Writing the unknown as
$\vx=(\vx_F,\vx_E)$ and ordering rows and columns so that free DoFs come first,
the assembled system $\mat{K}\vx=\vb$ has the block structure
\begin{equation}
  \begin{bmatrix} \mat{K}_{FF} & \mat{K}_{FE} \\[2pt]
                  \mat{K}_{EF} & \mat{K}_{EE} \end{bmatrix}
  \begin{bmatrix} \vx_F \\[2pt] \vx_E \end{bmatrix}
  =
  \begin{bmatrix} \vb_F \\[2pt] \vb_E \end{bmatrix}.
  \label{eq:blocksys}
\end{equation}
The block $\mat{K}_{FF}$ couples free DoFs to free DoFs---the genuine interior
physics. The off-diagonal blocks $\mat{K}_{FE}$ and $\mat{K}_{EF}$ couple free to
essential. The essential values $\vx_E$ are \emph{known}; only $\vx_F$ is to be
solved for.

The first row block of \cref{eq:blocksys} reads
$\mat{K}_{FF}\vx_F + \mat{K}_{FE}\vx_E = \vb_F$. Since $\vx_E$ is known, move its
contribution to the right:
\begin{equation}
  \mat{K}_{FF}\,\vx_F \;=\; \vb_F - \mat{K}_{FE}\,\vx_E.
  \label{eq:reduced}
\end{equation}
This reduced system is the honest problem: a smaller, square, well-posed system
on the free DoFs alone, with the known boundary data appearing as a forcing term
$-\mat{K}_{FE}\vx_E$ on the right. If we could re-index the unknowns at will we
would simply solve \cref{eq:reduced} and be done.

In practice the simulator keeps the full $N\times N$ index space---re-indexing a
sparse matrix is expensive and entangles boundary handling with storage---and
realizes the same reduction \emph{in place}. The operation
\texttt{eliminate\_essential\_bc} produces a fresh operator equal to $\mat{K}$
with the essential rows and columns zeroed and the eliminated diagonal set by a
\emph{policy}:
\begin{equation}
  \texttt{eliminate\_essential\_bc}(\mat{K},E,\text{policy})
  \;=\;
  \begin{bmatrix} \mat{K}_{FF} & \mathbf{0} \\[2pt]
                  \mathbf{0} & \mat{D} \end{bmatrix},
  \qquad
  \mat{D} =
  \begin{cases}
    \mat{I}_E & \text{(policy } \mathtt{DIAG\_ONE}),\\
    \mathbf{0}_E & \text{(policy } \mathtt{DIAG\_ZERO}).
  \end{cases}
  \label{eq:elim}
\end{equation}
Three things happen. The off-diagonal coupling blocks $\mat{K}_{FE}$ and
$\mat{K}_{EF}$ are zeroed---this decouples the essential DoFs from the free ones,
so the essential equations no longer pollute the free block and vice versa. The
free--free block $\mat{K}_{FF}$ is preserved exactly---the interior physics is
untouched. And the essential diagonal is overwritten with the identity or with
zero, depending on the policy. \Cref{fig:elim} shows the matrix before and after.

\begin{figure}[t]
\centering
\begin{tikzpicture}[scale=0.78]
  % ---- before ----
  \begin{scope}[shift={(0,0)}]
    \node[font=\small\bfseries] at (1.9,4.3) {before: $\mat{K}$};
    % FF block (large, top-left)
    \fill[simBgBlue] (0,1.2) rectangle (2.8,4);
    \node[simBlue,font=\small] at (1.4,2.6) {$\mat{K}_{FF}$};
    % FE block (top-right strip)
    \fill[simBgRust] (2.8,1.2) rectangle (3.8,4);
    \node[simRust,font=\scriptsize] at (3.3,2.6) {$\mat{K}_{FE}$};
    % EF block (bottom-left strip)
    \fill[simBgRust] (0,0) rectangle (2.8,1.2);
    \node[simRust,font=\scriptsize] at (1.4,0.6) {$\mat{K}_{EF}$};
    % EE block (bottom-right)
    \fill[simBgGray] (2.8,0) rectangle (3.8,1.2);
    \node[simGray,font=\scriptsize] at (3.3,0.6) {$\mat{K}_{EE}$};
    \draw[simInk,thick] (0,0) rectangle (3.8,4);
    \draw[simGray] (2.8,0)--(2.8,4);
    \draw[simGray] (0,1.2)--(3.8,1.2);
  \end{scope}
  % arrow
  \node[font=\Large] at (5,2) {$\longrightarrow$};
  \node[font=\scriptsize,simGray,above] at (5,2.2) {eliminate};
  % ---- after ----
  \begin{scope}[shift={(6.4,0)}]
    \node[font=\small\bfseries] at (1.9,4.3) {after: $\mat{K}_{\mathrm{BC}}$};
    \fill[simBgBlue] (0,1.2) rectangle (2.8,4);
    \node[simBlue,font=\small] at (1.4,2.6) {$\mat{K}_{FF}$};
    % zeroed FE
    \fill[white] (2.8,1.2) rectangle (3.8,4);
    \node[simGray,font=\scriptsize] at (3.3,2.6) {$\mathbf{0}$};
    % zeroed EF
    \fill[white] (0,0) rectangle (2.8,1.2);
    \node[simGray,font=\scriptsize] at (1.4,0.6) {$\mathbf{0}$};
    % diagonal D
    \fill[simBgGreen] (2.8,0) rectangle (3.8,1.2);
    \node[simGreen,font=\scriptsize] at (3.3,0.6) {$\mat{D}$};
    \draw[simInk,thick] (0,0) rectangle (3.8,4);
    \draw[simGray] (2.8,0)--(2.8,4);
    \draw[simGray] (0,1.2)--(3.8,1.2);
    % diagonal hint inside D
    \draw[simGreen,thick] (2.85,1.15)--(3.75,0.05);
  \end{scope}
\end{tikzpicture}
\caption[Essential-BC elimination on the matrix blocks]{Essential elimination as a
block operation. The assembled $\mat{K}$ (left) couples free and essential DoFs
through the off-diagonal blocks $\mat{K}_{FE}$, $\mat{K}_{EF}$. Elimination
(right) \emph{zeroes} those coupling blocks and the essential--essential block,
\emph{preserves} the free--free block $\mat{K}_{FF}$ exactly, and writes a chosen
diagonal $\mat{D}$ on the essential block---the identity $\mat{I}_E$ under the
\texttt{DIAG\_ONE} policy (so each essential equation reads $x_i=b_i$) or zero
under \texttt{DIAG\_ZERO}. The interior physics in $\mat{K}_{FF}$ is never
touched.}
\label{fig:elim}
\end{figure}

The two policies serve two purposes. \texttt{DIAG\_ONE} sets the essential
diagonal to the identity, so each essential row becomes the trivial equation
$1\cdot x_i = b_i$: the matrix is non-singular on the essential block, and the
prescribed value rides in through the right-hand side. This is the
\emph{solve-side} convention---it is what you use when the eliminated operator
will be inverted by a linear solver. \texttt{DIAG\_ZERO} leaves the essential
diagonal at zero, the \emph{energy/mass-block} convention used when the essential
block must contribute no spurious unit eigenvalues, as when assembling the blocks
of a generalized eigenvalue problem.

There is a clean algebraic way to see what elimination is. Let $\mat{P}_F$ be the
diagonal $0/1$ matrix that projects onto the free DoFs ($1$ on free indices, $0$
on essential). Then the \texttt{DIAG\_ZERO} elimination is exactly the two-sided
projection
\begin{equation}
  \texttt{eliminate\_essential\_bc}(\mat{K},E,\mathtt{DIAG\_ZERO})
  \;=\; \mat{P}_F\,\mat{K}\,\mat{P}_F,
  \label{eq:pfproj}
\end{equation}
because left-multiplying by $\mat{P}_F$ zeroes the essential rows and
right-multiplying zeroes the essential columns. This makes the operation
\emph{linear} in $\mat{K}$, and linearity in turn gives a sharp statement of
separability: elimination distributes over the assembly fold's term-sum. If
$\mat{K}=\mat{K}_1+\mat{K}_2$ is the sum of two assembled term-contributions, then
under \texttt{DIAG\_ZERO},
\[
  \texttt{eliminate\_essential\_bc}(\mat{K}_1+\mat{K}_2,E)
  = \texttt{eliminate\_essential\_bc}(\mat{K}_1,E)
  + \texttt{eliminate\_essential\_bc}(\mat{K}_2,E).
\]
The \texttt{DIAG\_ONE} policy is the same projection plus the constant $\mat{I}_E$
on the essential block---an affine, not linear, map---so it distributes up to that
one added identity. Either way, elimination factors through the free-block
projection \emph{regardless of how the matrix was assembled}, which is the precise
content of ``separable post-composition.''

\subsection{Inhomogeneous Dirichlet data: lifting into the right-hand side}
\label{sec:lifting}

When the prescribed boundary values are zero---a grounded conductor, a PEC
wall---the forcing term $-\mat{K}_{FE}\vx_E$ in \cref{eq:reduced} vanishes,
because $\vx_E=\mathbf{0}$. The reduced system is just
$\mat{K}_{FF}\vx_F=\vb_F$ and there is nothing more to do.

But the interesting essential conditions are \emph{inhomogeneous}: a capacitor
plate held at $V\neq 0$, a port driven at a nonzero amplitude. Then $\vx_E\neq
\mathbf{0}$, and the term $-\mat{K}_{FE}\vx_E$ in \cref{eq:reduced} is a genuine,
nonzero load. This is the discrete face of the ``inhomogeneous essential''
condition we flagged back in the parallel-plate example of \cref{ch:weak}: a
problem with no volume source and no surface charge ($f=g=0$, an empty $\ell$)
still produces a nonzero right-hand side, driven entirely by the imposed
boundary values. The mechanism that produces it is called \emph{lifting}.

\begin{figure}[t]
\centering
\begin{tikzpicture}[scale=1.0, node distance=7mm]
  % a known boundary value
  \node[data, fill=simBgBlue, draw=simBlue, text width=20mm] (xe)
    {known boundary\\data $\vx_{\mathrm{ess}}$};
  % apply K
  \node[opbox, right=14mm of xe, text width=22mm] (apply)
    {apply operator\\$\mat{K}\cdot\vx_{\mathrm{ess}}$};
  % subtract from b
  \node[opbox, right=14mm of apply, text width=24mm] (sub)
    {subtract from RHS\\$\vb \leftarrow \vb - \mat{K}\vx_{\mathrm{ess}}$};
  \node[product, right=14mm of sub, text width=18mm, align=center] (out)
    {lifted RHS\\$\vb'$};
  \draw[flow] (xe) -- (apply);
  \draw[flow] (apply) -- (sub);
  \draw[flow] (sub) -- (out);
  % the b feeding in
  \node[data, above=8mm of sub, font=\scriptsize\itshape] (b) {original $\vb$};
  \draw[flow] (b) -- (sub);
  % annotation
  \node[font=\scriptsize\itshape, simRust, below=8mm of apply.south, text width=40mm, align=center]
    {the boundary value is ``lifted'' off the unknown and recast as a known forcing on the interior};
  \draw[refedge] (apply.south) -- ++(0,-3mm);
\end{tikzpicture}
\caption[Lifting an inhomogeneous Dirichlet condition into the RHS]{Lifting. A
known boundary value $\vx_{\mathrm{ess}}$ on the essential DoFs is not an unknown
to solve for---it is data. The operation \texttt{eliminate\_rhs} applies the
assembled operator to it, $\mat{K}\,\vx_{\mathrm{ess}}$, and \emph{subtracts the
result from the right-hand side}, $\vb\leftarrow\vb-\mat{K}\vx_{\mathrm{ess}}$.
The boundary value's influence on the interior equations is thereby moved from the
matrix (where it was a coupling) to the load vector (where it is a known source).
The essential rows of $\vb'$ are then pinned to the boundary data, matching the
identity diagonal that \texttt{eliminate\_essential\_bc} installed.}
\label{fig:lifting}
\end{figure}

Lifting is exactly the move from \cref{eq:reduced}, performed on the full index
space. Let $\vx_{\mathrm{ess}}$ be the true-DoF vector that holds the prescribed
boundary values on the essential DoFs (and zero elsewhere). The operation
\texttt{eliminate\_rhs} computes
\begin{equation}
  \vb' \;=\; \vb \;-\; \mat{K}\,\vx_{\mathrm{ess}},
  \label{eq:liftalg}
\end{equation}
then overwrites the essential rows of $\vb'$ with the boundary data itself
(matching the identity diagonal of the \texttt{DIAG\_ONE} policy). Read against
\cref{eq:reduced}: on a free row $i\in F$, the product $\mat{K}\vx_{\mathrm{ess}}$
contributes $\sum_{j\in E}\mat{K}_{ij}(\vx_{\mathrm{ess}})_j = (\mat{K}_{FE}\vx_E)_i$,
which is exactly the term being subtracted. On an essential row, the pin
overwrites whatever was there with the prescribed value, so the essential equation
reads $1\cdot x_i = (\vx_{\mathrm{ess}})_i$. The two operations together---
\texttt{eliminate\_essential\_bc} on the operator, \texttt{eliminate\_rhs} on the
right-hand side---realize the full Dirichlet treatment of \cref{eq:reduced}
without ever re-indexing the matrix.

Algebraically, \cref{eq:liftalg} is just one matrix--vector product
($\mat{K}\vx_{\mathrm{ess}}$) and one scaled vector addition ($\vb-(\cdot)$,
an \texttt{axpy} with coefficient $-1$). It carries no SPD or invertibility
requirement on $\mat{K}$; it simply moves a known quantity from one side of the
equation to the other. Notice, too, that lifting is a \emph{one-shot}
preparation, not a projector: apply it twice and you subtract
$\mat{K}\vx_{\mathrm{ess}}$ twice, which is wrong. It is performed exactly once,
on the way into the solve.

\begin{workedexample}{Lifting a nonzero Dirichlet value in 1-D}{lift1d}
We carry the whole essential-BC machinery through on a system small enough to
write out completely. Discretize the 1-D Poisson problem $-u''=0$ on $[0,1]$ with
the boundary conditions $u(0)=0$ and $u(1)=V$ (an inhomogeneous Dirichlet datum
at the right end), using a uniform grid with two interior nodes. The standard
finite-difference/finite-element stiffness on nodes $x_0,x_1,x_2,x_3$ (with
$x_0=0$, $x_3=1$, spacing $h=\tfrac13$) is, up to the factor $1/h$,
\[
  \mat{K} =
  \begin{bmatrix}
     1 & -1 &  0 &  0\\
    -1 &  2 & -1 &  0\\
     0 & -1 &  2 & -1\\
     0 &  0 & -1 &  1
  \end{bmatrix},
  \qquad \vb = \mathbf{0}.
\]
The essential set is $E=\{0,3\}$ (the two endpoints); the free set is $F=\{1,2\}$
(the two interior nodes). The prescribed boundary data is
$\vx_{\mathrm{ess}}=(0,\,?,\,?,\,V)^{\!\top}$, with the interior entries
irrelevant (they will be masked out).

\medskip\noindent\textbf{Reduced system, by hand.}
Writing $\vx=(x_0,x_1,x_2,x_3)$ with $x_0=0$, $x_3=V$ known, the two interior
equations (rows $1$ and $2$) are
\[
  -x_0 + 2x_1 - x_2 = 0, \qquad -x_1 + 2x_2 - x_3 = 0.
\]
Substituting $x_0=0$ and $x_3=V$ and moving the known terms to the right gives the
reduced free system \cref{eq:reduced}:
\[
  \begin{bmatrix} 2 & -1\\ -1 & 2\end{bmatrix}
  \begin{bmatrix} x_1\\ x_2\end{bmatrix}
  =
  \begin{bmatrix} 0\\ V\end{bmatrix}.
\]
The right-hand side entry $V$ in the second equation is precisely the lifted term
$-(\mat{K}_{FE}\vx_E)_2 = -(-1)\cdot V = V$. Solving:
$x_1 = V/3$, $x_2 = 2V/3$. Together with the endpoints this is the exact
linear potential $u(x)=Vx$ sampled at the four nodes
$(0,\,V/3,\,2V/3,\,V)$---as it must be, since the solution of $-u''=0$ with these
boundary values is the straight line from $0$ to $V$.

\medskip\noindent\textbf{The same thing via the two operations.}
Now reproduce the result through \texttt{eliminate\_rhs} and
\texttt{eliminate\_essential\_bc} on the full $4\times4$ system, without
re-indexing. First lift the RHS with $\vx_{\mathrm{ess}}=(0,0,0,V)^{\!\top}$:
\[
  \mat{K}\,\vx_{\mathrm{ess}} =
  \begin{bmatrix} 1&-1&0&0\\-1&2&-1&0\\0&-1&2&-1\\0&0&-1&1\end{bmatrix}
  \begin{bmatrix}0\\0\\0\\V\end{bmatrix}
  = \begin{bmatrix}0\\0\\-V\\V\end{bmatrix},
  \qquad
  \vb' = \vb - \mat{K}\vx_{\mathrm{ess}} =
  \begin{bmatrix}0\\0\\V\\-V\end{bmatrix}.
\]
On the free rows $\{1,2\}$ this gives $\vb'_1=0$, $\vb'_2=V$---exactly the
right-hand side of the reduced system above. (The essential rows $\{0,3\}$ are
then pinned to the boundary data $0$ and $V$.) Second, eliminate the operator with
\texttt{DIAG\_ONE}: zero rows/columns $0$ and $3$ and set their diagonal to $1$,
producing
\[
  \mat{K}_{\mathrm{BC}} =
  \begin{bmatrix} 1&0&0&0\\ 0&2&-1&0\\ 0&-1&2&0\\ 0&0&0&1\end{bmatrix}.
\]
Solving $\mat{K}_{\mathrm{BC}}\vx=\vb'$ with $\vb'=(0,0,V,-V)^{\!\top}$ after
pinning the essential rows to $(0,\dots,V)$ recovers the free block
$\bigl[\begin{smallmatrix}2&-1\\-1&2\end{smallmatrix}\bigr]
(x_1,x_2)^{\!\top}=(0,V)^{\!\top}$ and hence $x_1=V/3$, $x_2=2V/3$ again. The
matrix-level operations reproduce the by-hand reduction exactly, which is the
whole point: \texttt{eliminate\_essential\_bc} and \texttt{eliminate\_rhs} are a
re-indexing-free implementation of \cref{eq:reduced}.
\end{workedexample}

\subsection{Natural conditions, by contrast}
\label{sec:natural}

It is worth setting the essential machinery beside the natural conditions to see
how little the latter cost. Recall the boundary split of \cref{ch:weak}: Green's
identity produced a boundary integral $\oint_\bdry v\,(\varepsilon\Grad u)\cdot
\vn\,\mathrm{d}S$, and we disposed of it in two ways. On the Dirichlet part we
forced the test functions to vanish and the integral dropped---the essential
condition, which then had to be imposed on the discrete system by the elimination
we just built. On the Neumann part the integrand was a \emph{known} flux $g$, and
the surviving integral $\oint_{\Gamma_N} g\,v\,\mathrm{d}S$ moved to the
right-hand side as a source.

That natural-BC surface term is assembled \emph{during} assembly, not after. It is
just another weak-form term---a boundary integrator $(g,\text{value})$ on the
surface $\Gamma_N$, folded into $\vb$ exactly like a volume source---and the
assembly machinery of \cref{ch:assembly} handles it with no special case. A
\emph{homogeneous} Neumann condition ($g=0$) contributes nothing; it is enforced
by doing nothing at all. This is the asymmetry the pitfall in \cref{ch:weak}
warned about, now visible in implementation terms: the natural condition is a row
in the integrator list, while the essential condition is a post-assembly surgery
on rows and columns.

\begin{figure}[t]
\centering
\begin{tikzpicture}[scale=1.0]
  % essential side
  \begin{scope}[shift={(0,0)}]
    \node[font=\small\bfseries, simBlue] at (1.6,2.6) {essential (Dirichlet)};
    \node[stage, fill=simBgBlue, draw=simBlue, text width=30mm, align=center] (e1) at (1.6,1.7)
      {value $u=u_0$ prescribed};
    \node[opbox, text width=34mm, below=4mm of e1] (e2)
      {enforced on the \emph{discrete system}:\\zero rows/cols, lift RHS};
    \node[font=\scriptsize\itshape, simRust, below=3mm of e2] (e3)
      {\emph{after} assembly --- a post-composition};
    \draw[flow] (e1) -- (e2);
    \draw[flow] (e2) -- (e3);
  \end{scope}
  % divider
  \draw[simGray, densely dashed] (3.7,-0.7) -- (3.7,2.8);
  % natural side
  \begin{scope}[shift={(5.0,0)}]
    \node[font=\small\bfseries, simGreen] at (1.8,2.6) {natural (Neumann)};
    \node[stage, fill=simBgGreen, draw=simGreen, text width=30mm, align=center] (n1) at (1.8,1.7)
      {flux $(\varepsilon\Grad u)\cdot\vn = g$ prescribed};
    \node[opbox, text width=34mm, below=4mm of n1] (n2)
      {enters the \emph{weak form}:\\boundary integrator $\to \vb$};
    \node[font=\scriptsize\itshape, simRust, below=3mm of n2] (n3)
      {\emph{during} assembly --- $g=0$ costs nothing};
    \draw[flow] (n1) -- (n2);
    \draw[flow] (n2) -- (n3);
  \end{scope}
\end{tikzpicture}
\caption[Essential versus natural, in the discrete setting]{The two boundary-
condition types impose themselves at different stages. An \emph{essential}
condition (left) prescribes the field's \emph{value} and must be enforced on the
assembled system, after the fold, by zeroing rows/columns and lifting the data
into the right-hand side---the work of \cref{sec:essential}. A \emph{natural}
condition (right) prescribes a \emph{flux} and enters the weak form as a boundary
integrator folded into $\vb$ \emph{during} assembly; a homogeneous one ($g=0$)
costs nothing. ``Essential'' means ``you enforce it''; ``natural'' means ``the
weak form enforces it for you.''}
\label{fig:essvsnat}
\end{figure}

\section{The Helmholtz decomposition}
\label{sec:helmholtz}

We turn to the second finishing operation. To motivate it we need a classical
theorem about vector fields, and its discrete face.

\subsection{The continuous decomposition}

A foundational fact of vector calculus is that any reasonable vector field on a
domain can be split into two pieces of complementary character: a part with no
divergence and a part that is a pure gradient.

\begin{theorem}[Helmholtz decomposition]
\label{thm:helmholtz}
Let $\vF$ be a (sufficiently smooth) vector field on a bounded domain $\dom$ with
appropriate boundary conditions. Then $\vF$ decomposes uniquely as
\begin{equation}
  \vF \;=\; \vF_{\mathrm{df}} \;+\; \Grad\psi,
  \label{eq:helmholtz}
\end{equation}
where $\vF_{\mathrm{df}}$ is \emph{divergence-free} ($\Div\vF_{\mathrm{df}}=0$)
and $\Grad\psi$ is a \emph{gradient} (and hence curl-free, $\Curl\Grad\psi=0$).
The two parts are orthogonal in the $\Ltwo$ inner product.
\end{theorem}

The two pieces are sometimes called the \emph{solenoidal} (divergence-free) and
\emph{irrotational} (gradient) components. Physically they are the field's
``circulating'' part and its ``spreading'' part: a magnetic field is purely
solenoidal ($\Div\vB=0$ always), while an electrostatic field is purely
irrotational ($\vE=-\Grad V$). A general field mixes the two, and
\cref{thm:helmholtz} says the mixture separates cleanly and uniquely.
\Cref{fig:helmholtz} shows the decomposition on a small field.

The orthogonality claim is worth a moment, because it is the property the discrete
projector inherits and the reason the decomposition is \emph{unique}. The two
parts are $\Ltwo$-orthogonal:
\begin{equation}
  \inner{\vF_{\mathrm{df}}}{\Grad\psi}_{\Ltwo(\dom)}
  = \int_\dom \vF_{\mathrm{df}}\cdot\Grad\psi \dV = 0.
  \label{eq:helmorth}
\end{equation}
The reason is Green's first identity (\cref{ch:weak}) run backwards. Integrate by
parts: $\int_\dom \vF_{\mathrm{df}}\cdot\Grad\psi
= -\int_\dom (\Div\vF_{\mathrm{df}})\,\psi
+ \oint_\bdry \psi\,(\vF_{\mathrm{df}}\cdot\vn)\,\mathrm{d}S$. The volume term
vanishes because $\vF_{\mathrm{df}}$ is divergence-free ($\Div\vF_{\mathrm{df}}=0$),
and the boundary term vanishes under the boundary conditions that accompany the
decomposition (either $\psi=0$ or $\vF_{\mathrm{df}}\cdot\vn=0$ on $\bdry$). So a
divergence-free field is $\Ltwo$-orthogonal to every gradient---which is precisely
what makes ``the gradient part'' well-defined: it is the orthogonal projection of
$\vF$ onto the gradient subspace, and $\vF_{\mathrm{df}}$ is whatever remains.
Uniqueness follows immediately: if $\vF=\vF_{\mathrm{df}}+\Grad\psi
=\vF'_{\mathrm{df}}+\Grad\psi'$ were two decompositions, their difference
$\vF_{\mathrm{df}}-\vF'_{\mathrm{df}}=\Grad(\psi'-\psi)$ would be both
divergence-free and a gradient, hence orthogonal to itself, hence zero. The
discrete projector of \cref{sec:divfree} reproduces this orthogonality---but in a
\emph{weighted} inner product, the $\mat{M}$-inner-product, because the discrete
mass operator $\mat{M}$ encodes the material coefficient $\varepsilon$.

\begin{figure}[t]
\centering
\begin{tikzpicture}[scale=1.0,
    vec/.style={-{Stealth[length=2mm]}, thick}]
  % --- the full field F (left) ---
  \begin{scope}[shift={(0,0)}]
    \node[font=\small\bfseries] at (1.5,3.5) {$\vF$ (mixed)};
    \draw[simGray!50] (0,0) rectangle (3,3);
    \foreach \x in {0.5,1.5,2.5}{
      \foreach \y in {0.5,1.5,2.5}{
        % a swirling-plus-spreading field
        \pgfmathsetmacro{\dx}{0.32*cos(120*\x+60*\y) + 0.20*(\x-1.5)}
        \pgfmathsetmacro{\dy}{0.32*sin(120*\x+60*\y) + 0.20*(\y-1.5)}
        \draw[vec,simInk] (\x,\y) -- ++(\dx,\dy);}}
  \end{scope}
  % equals
  \node[font=\Large] at (3.7,1.5) {$=$};
  % --- divergence-free part (middle) ---
  \begin{scope}[shift={(4.4,0)}]
    \node[font=\small\bfseries, simBlue] at (1.5,3.5) {$\vF_{\mathrm{df}}$ (curl)};
    \draw[simGray!50] (0,0) rectangle (3,3);
    \foreach \x in {0.5,1.5,2.5}{
      \foreach \y in {0.5,1.5,2.5}{
        \pgfmathsetmacro{\dx}{0.32*cos(120*\x+60*\y)}
        \pgfmathsetmacro{\dy}{0.32*sin(120*\x+60*\y)}
        \draw[vec,simBlue] (\x,\y) -- ++(\dx,\dy);}}
  \end{scope}
  % plus
  \node[font=\Large] at (8.1,1.5) {$+$};
  % --- gradient part (right) ---
  \begin{scope}[shift={(8.8,0)}]
    \node[font=\small\bfseries, simRust] at (1.5,3.5) {$\Grad\psi$ (spread)};
    \draw[simGray!50] (0,0) rectangle (3,3);
    \foreach \x in {0.5,1.5,2.5}{
      \foreach \y in {0.5,1.5,2.5}{
        \pgfmathsetmacro{\dx}{0.20*(\x-1.5)}
        \pgfmathsetmacro{\dy}{0.20*(\y-1.5)}
        \draw[vec,simRust] (\x,\y) -- ++(\dx,\dy);}}
    % a contour of psi
    \draw[simRust!50,densely dotted] (1.5,1.5) circle (0.5);
    \draw[simRust!50,densely dotted] (1.5,1.5) circle (1.0);
  \end{scope}
\end{tikzpicture}
\caption[The Helmholtz decomposition of a vector field]{The Helmholtz
decomposition \cref{eq:helmholtz} splits a vector field $\vF$ (left, mixed) into a
\emph{divergence-free} part $\vF_{\mathrm{df}}$ (middle, purely circulating---it
swirls but neither spreads from nor converges to any point) and a \emph{gradient}
part $\Grad\psi$ (right, purely spreading---it points along the steepest ascent of
a scalar potential $\psi$, whose level sets are the dotted contours). The two
parts are $\Ltwo$-orthogonal. The divergence-free part is what the physics of a
magnetic or resonant field demands; the gradient part is the unphysical
contamination a singular operator can introduce.}
\label{fig:helmholtz}
\end{figure}

\subsection{Its de Rham meaning}

The decomposition is not an isolated trick; it is a face of the de Rham complex of
\cref{ch:derham}. Recall the ladder of spaces and operators
\[
  \Hone \xrightarrow{\ \Grad\ } \Hcurl \xrightarrow{\ \Curl\ }
  \Hdiv \xrightarrow{\ \Div\ } \Ltwo,
\]
with the defining \emph{exactness} property that the range of each operator is the
kernel of the next: $\Curl\Grad = 0$ and $\Div\Curl = 0$, and (on a
topologically simple domain) the inclusions are equalities. Look at the middle
slot, $\Hcurl$. The gradient part of \cref{eq:helmholtz} lives in
$\mathrm{Range}(\Grad)$, the image of the gradient operator coming \emph{into}
$\Hcurl$ from $\Hone$. Exactness says this image is exactly $\mathrm{Ker}(\Curl)$,
the curl-free fields. So the Helmholtz split of an $\Hcurl$ field is precisely its
decomposition into
\begin{equation}
  \Hcurl \;=\; \underbrace{\mathrm{Range}(\Grad)}_{\text{gradients}\,=\,\mathrm{Ker}(\Curl)}
  \;\oplus\;
  \underbrace{(\text{divergence-free complement})}_{\vF_{\mathrm{df}}},
  \label{eq:hcurlsplit}
\end{equation}
the gradient part being the kernel of the curl. This is the structural reason the
decomposition matters for us: the operators we solve live on $\Hcurl$, and their
null space \emph{is} the gradient subspace of \cref{eq:hcurlsplit}.

\begin{keyidea}
The Helmholtz decomposition \cref{eq:helmholtz} is the de Rham complex read at the
$\Hcurl$ slot. The gradient part of an $\Hcurl$ field is its component in
$\mathrm{Range}(\Grad)$, which by exactness of the complex
(\cref{ch:derham}) is exactly $\mathrm{Ker}(\Curl)$---the fields the curl
annihilates. Removing the gradient part is therefore the same as removing the
curl operator's null space. The discrete projector of the next section is a
matrix realization of this removal.
\end{keyidea}

\section{The discrete divergence-free projector}
\label{sec:divfree}

We now build the operator that removes the gradient part from a discrete field.
It is the workhorse $\mathtt{divfree\_project}$, and it is, like essential
elimination, a separable post-composition---but a richer one, because it carries
an inner linear solve.

\subsection{The projector and its apply}

Let $\mat{G}$ be the discrete gradient operator (the matrix that takes the
$\Hone$ DoFs of a scalar potential to the $\Hcurl$ DoFs of its gradient---the de
Rham interpolator of \cref{ch:derham}), and let $\mat{M}$ be the
$\varepsilon$-weighted $\Hone$ mass operator (symmetric positive definite). The
discrete divergence-free projector is
\begin{equation}
  \boxed{\;
  \mat{P} \;=\; \mat{I} \;-\; \mat{G}\,\bigl(\mat{G}^{\!\top}\mat{M}\,\mat{G}\bigr)^{-1}\,\mat{G}^{\!\top}\mat{M}.
  \;}
  \label{eq:projector}
\end{equation}
This is an $\mat{M}$-orthogonal projection onto the divergence-free subspace. Read
it from right to left as a recipe: $\mat{G}^{\!\top}\mat{M}$ measures the (weak)
divergence of the input field; $(\mat{G}^{\!\top}\mat{M}\mat{G})^{-1}$ solves a
Poisson-type system for a scalar potential $\psi$; $\mat{G}$ takes the gradient of
that potential, materializing the gradient component; and $\mat{I}-(\cdots)$
subtracts it off, leaving the divergence-free remainder. The middle factor
$(\mat{G}^{\!\top}\mat{M}\mat{G})^{-1}$ is the $\mat{M}$-weighted normal-equations
inverse: $\psi$ is chosen so that $\mat{G}\psi$ is the $\mat{M}$-orthogonal
projection of the input onto the gradient subspace.

\begin{notationbox}
The triple product $\mat{G}^{\!\top}\mat{M}\mat{G}$ is the Poisson-type stiffness
on the scalar ($\Hone$) space---it is what you get by sandwiching the mass
operator between the gradient and its transpose. In Palace it is never formed as a
matrix: the inner solve is posed against $\mat{M}$ directly, with $\mat{G}^{\!\top}$
applied on the right-hand side (as a weak-divergence operator) and $\mat{G}$
applied to the solution. The boxed form \cref{eq:projector} is the algebra; the
implementation realizes it factor by factor, which is why the four-step apply
below mentions $\mat{M}$, $\mat{G}$, and a weak-divergence operator separately
rather than the triple product.
\end{notationbox}

In practice we never form $\mat{P}$ as a matrix and we rarely apply the full
$\mat{P}$; we apply it to a field $\vy$ to extract the divergence-free part
$\vy'=\mat{P}\vy$. The apply runs in four steps, shown as a pipeline in
\cref{fig:projapply}:

\begin{enumerate}\setlength\itemsep{2pt}
  \item \textbf{Weak divergence.} Compute $\vr = \mat{G}^{\!\top}\mat{M}\,\vy$, an
  $\Hone$-side residual measuring the divergence of $\vy$. (In Palace this is one
  application of a weak-divergence operator, the $\varepsilon$-weighted
  Nédélec-to-$\Hone$ map; its sign convention is load-bearing and discussed below.)
  \item \textbf{Essential-BC zeroing.} Zero $\vr$ on the essential boundary DoFs of
  the scalar space, $\vr\leftarrow \mat{Z}_{\mathrm{bdr}}\vr$. This pins the inner
  Poisson problem and removes the pure-Neumann null space (a constant shift in
  $\psi$ that would otherwise leave the solve singular).
  \item \textbf{Inner Poisson solve.} Solve $\mat{M}\psi = \vr$ for the scalar
  potential $\psi$. This is itself an iterative \texttt{ksp\_solve}
  (\cref{ch:krylov})---a conjugate-gradient solve against the SPD operator
  $\mat{M}$, preconditioned. The projector therefore \emph{contains a solver}: it
  is a constructed operator carrying an inner solver as a sub-component
  (\cref{ch:operators}).
  \item \textbf{Gradient correction.} Form the gradient $\mat{G}\psi$ and subtract
  it: $\vy' = \vy - \mat{G}\psi$ (an \texttt{apply\_linop} followed by an
  \texttt{axpy}). The result $\vy'$ is the divergence-free component.
\end{enumerate}

\begin{figure}[t]
\centering
\begin{tikzpicture}[node distance=8mm]
  \node[data, text width=12mm, align=center] (y) {field $\vy$};
  \node[opbox, right=9mm of y, text width=20mm] (wd)
    {\textbf{1.}\,weak div\\$\mat{G}^{\!\top}\mat{M}\vy$};
  \node[opbox, right=8mm of wd, text width=18mm] (z)
    {\textbf{2.}\,zero BC\\$\mat{Z}_{\mathrm{bdr}}$};
  \node[extern, right=8mm of z, text width=22mm] (solve)
    {\textbf{3.}\,inner solve\\$\mat{M}\psi=\vr$\\[1pt]\scriptsize \texttt{ksp\_solve}};
  \node[opbox, below=10mm of solve, text width=20mm] (grad)
    {\textbf{4.}\,gradient\\$\mat{G}\psi$};
  \node[product, left=8mm of grad, text width=22mm, align=center] (sub)
    {subtract\\$\vy'=\vy-\mat{G}\psi$};
  \node[data, left=8mm of sub, text width=18mm, align=center] (out)
    {divergence-free $\vy'$};
  \draw[flow] (y) -- (wd);
  \draw[flow] (wd) -- (z);
  \draw[flow] (z) -- (solve);
  \draw[flow] (solve) -- (grad);
  \draw[flow] (grad) -- (sub);
  \draw[flow] (sub) -- (out);
  % the original y feeding the subtract
  \draw[flow, simGray, densely dashed] (y.south) |- (sub.north);
  \node[font=\scriptsize\itshape, simViolet, above=1mm of solve.north]
    {the inner Krylov solve};
\end{tikzpicture}
\caption[The four-step divergence-free projector apply]{The
$\mathtt{divfree\_project}$ apply as a pipeline. Given a field $\vy$, step~1
computes its weak divergence $\mat{G}^{\!\top}\mat{M}\vy$; step~2 zeroes that
residual on the essential boundary DoFs (pinning the inner problem); step~3
\emph{solves} a Poisson-type system $\mat{M}\psi=\vr$ for the scalar potential
$\psi$---an inner \texttt{ksp\_solve} (\cref{ch:krylov}), the dashed/violet box,
which makes the projector a constructed operator carrying its own solver; step~4
forms the gradient $\mat{G}\psi$ and subtracts it from the original $\vy$,
leaving the divergence-free component $\vy'$. Steps 1, 2, 4 are cheap whole-field
operations; the cost lives in the inner solve of step~3.}
\label{fig:projapply}
\end{figure}

The defining property is that the output is annihilated by the weak divergence:
$\mat{G}^{\!\top}\mat{M}\,\vy' = 0$. That is the discrete statement of
``divergence-free.'' And the part removed, $\mat{G}\psi$, is a discrete
gradient---the kernel of the discrete curl, by the de Rham exactness of
\cref{ch:derham}. So the projector strips exactly the curl operator's null space,
which is what \cref{eq:hcurlsplit} promised.

\begin{workedexample}{Walking the four-step apply on numbers}{fourstep}
To make the pipeline of \cref{fig:projapply} concrete, run all four steps on a
minimal instance with no essential boundary (so step~2 is the identity and we can
watch steps 1, 3, 4 do the work). Take a field space of dimension $2$ and a scalar
space of dimension $1$, with
\[
  \mat{M} = \mat{I}_2,
  \qquad
  \mat{G} = \begin{bmatrix}1\\1\end{bmatrix},
  \qquad
  \mat{G}^{\!\top}\mat{M} = \begin{bmatrix}1 & 1\end{bmatrix},
\]
the same data as \cref{wex:projcheck}. Project the field
$\vy=(3,1)^{\!\top}$---a mix of a circulating part along $(1,-1)$ and a gradient
part along $(1,1)$.

\medskip\noindent\textbf{Step 1 --- weak divergence.}
$\vr = \mat{G}^{\!\top}\mat{M}\,\vy = \begin{bmatrix}1&1\end{bmatrix}
\begin{bmatrix}3\\1\end{bmatrix} = 4$. The field has nonzero weak divergence, so it
is not already divergence-free---there is a gradient part to remove.

\medskip\noindent\textbf{Step 2 --- essential-BC zeroing.}
With no essential boundary, $\mat{Z}_{\mathrm{bdr}}=\mat{I}$ and $\vr$ passes
through unchanged: $\vr = 4$.

\medskip\noindent\textbf{Step 3 --- inner Poisson solve.}
Solve $(\mat{G}^{\!\top}\mat{M}\mat{G})\,\psi = \vr$, i.e.\ $2\psi = 4$, giving
$\psi = 2$. (In the implementation this is the inner \texttt{ksp\_solve} against
$\mat{M}$; here the scalar system is $1\times1$ and solves in closed form.) The
scalar potential whose gradient is the irrotational part of $\vy$ is $\psi=2$.

\medskip\noindent\textbf{Step 4 --- gradient correction.}
Form $\mat{G}\psi = \begin{bmatrix}1\\1\end{bmatrix}\cdot 2 = (2,2)^{\!\top}$ and
subtract: $\vy' = \vy - \mat{G}\psi = (3,1)^{\!\top} - (2,2)^{\!\top} =
(1,-1)^{\!\top}$.

\medskip\noindent\textbf{Check.}
The output $\vy'=(1,-1)^{\!\top}$ is exactly the antidiagonal (circulating)
component of $\vy$. Verify it is divergence-free:
$\mat{G}^{\!\top}\mat{M}\,\vy' = \begin{bmatrix}1&1\end{bmatrix}
\begin{bmatrix}1\\-1\end{bmatrix} = 0$, as the defining property demands. And the
removed part $\mat{G}\psi=(2,2)^{\!\top}$ is a pure gradient, along $(1,1)$. This
matches the closed-form projector of \cref{wex:projcheck}:
$\mat{P}\,\vy=\tfrac12\bigl[\begin{smallmatrix}1&-1\\-1&1\end{smallmatrix}\bigr]
\bigl[\begin{smallmatrix}3\\1\end{smallmatrix}\bigr]=\tfrac12(2,-2)^{\!\top}
=(1,-1)^{\!\top}$---the four-step apply and the boxed formula agree, as they must.
The point of running the steps separately is that the implementation never forms
$\mat{P}$: it walks exactly steps 1--4, and the expensive one is step~3, the inner
solve that here was a trivial $2\psi=4$ but in production is a full
preconditioned CG solve on the scalar space.
\end{workedexample}

\begin{notationbox}
\emph{A sign subtlety, recorded.} In Palace the weak-divergence operator carries a
built-in $-1$ (its bilinear form is $a(\vu,v)=-(\varepsilon\vu,\Grad v)$, so it is
the negated gradient transpose, $\mat{G}^{\!\top}$ up to that sign). The
implementation then \emph{adds} the correction, $\vy'\leftarrow\vy+\mat{G}\psi$,
and the two minus signs combine so that the net effect still removes the gradient
part. The algebra of \cref{eq:projector} is written with the conventional sign
($\vy-\mat{G}\psi$); the source's two-negatives bookkeeping is equivalent. We note
it because a single flipped sign here would invert the correction and the
projector would \emph{amplify} the gradient part instead of removing it.
\end{notationbox}

\subsection{It is a constructed operator with an inner solve}

Step 3 deserves emphasis. The projector is not a fixed sparse matrix; it is a
\emph{constructed operator}---a structured value assembled once at solver setup,
carrying the operators $\mat{M}$, $\mat{G}$, the weak-divergence operator, the
essential-boundary DoF set, and, crucially, an \emph{inner linear solver} bound to
$\mat{M}$. Applying the projector invokes that inner solver. This makes
$\mathtt{divfree\_project}$ the first operator we have met whose own definition
contains a solve-to-tolerance, and it foreshadows the constructed-operator
machinery of \cref{ch:operators}: operators that carry other operators (and other
solvers) as sub-components, applied opaquely. The inner conjugate-gradient
iteration is sequential---a recurrence that does not parallelize across its
steps---but that sequentiality is \emph{interior} to the inner solve; the
projector itself is a fixed four-step straight-line composition around it.

The construction-versus-apply split is the same one we saw for assembly, lifted one
level. There is a \emph{setup} phase, run once: build $\mat{M}$ and the
weak-divergence operator by assembly (themselves folds over weak-form terms on the
scalar space), interpolate the discrete gradient $\mat{G}$ from the de Rham
structure, compute the essential-boundary DoF set, and configure the inner solver
with its preconditioner, tolerance, and iteration cap. None of this depends on the
field being projected; it is bound into the constructed value $\mat{P}$ at solver
setup. Then there is the \emph{apply} phase, run once per field (and, in an
eigensolve, once per iterate): the four cheap-plus-one-expensive steps of
\cref{fig:projapply}. Keeping the two phases separate is what makes the per-iterate
projection affordable---the operators are assembled and the solver is set up exactly
once, and each projection reuses them.

The inner solve is also where the projector's cost concentrates, and where its one
genuine subtlety lives. The idempotence law $\mat{P}^2=\mat{P}$ of
\cref{wex:projcheck} rested on the middle factor being the \emph{exact} inverse
$(\mat{G}^{\!\top}\mat{M}\mat{G})^{-1}$. In practice the inner CG solve is only
approximate, converged to a tolerance, so $\mat{P}$ is an \emph{approximate}
projector: $\mat{P}^2\approx\mat{P}$, and $\mat{G}^{\!\top}\mat{M}\,\vy'\approx 0$,
both up to that tolerance. This is why the inner solve must be converged
tightly---loose tolerances leave a residual gradient component that the eigensolver
will happily amplify back into a spurious mode. The mass operator $\mat{M}$ is SPD,
so CG is the natural choice and a good preconditioner (an algebraic multigrid
$\mat{V}$-cycle on $\mat{M}$) makes the solve cheap; but cheap and \emph{loose} are
different knobs, and only the first may be turned.

\begin{remark}[The projector is a separable post-composition too]
Like essential elimination, the divergence-free projection consumes its inputs
opaquely. The four-step apply treats $\mat{M}$, $\mat{G}$, and the weak-divergence
operator as given operators and the field $\vy$ as a given vector; it never inspects
how any of them were assembled, nor which weak-form terms produced the curl--curl
operator it is guarding. So although the projector is structurally richer than
\texttt{eliminate\_essential\_bc}---it carries an inner solver where elimination
carries only an index set---it occupies the same architectural slot: a finishing
post-composition that acts on assembled objects and composes cleanly behind the
assembly fold.
\end{remark}

\subsection{Verifying the projector properties}

A projector must satisfy $\mat{P}^2=\mat{P}$ (applying it twice is the same as
once) and must annihilate the subspace it projects out (here, the gradients:
$\mat{P}\,\mat{G}\psi = 0$ for any $\psi$). Both follow from the algebraic form
\cref{eq:projector}, and both are worth seeing concretely.

\begin{workedexample}{$\mat{P}$ is a projector and kills gradients}{projcheck}
We verify the two projector laws, first symbolically from \cref{eq:projector},
then on a tiny numerical instance.

\medskip\noindent\textbf{Symbolic: idempotence.}
Abbreviate the middle inverse as $\mat{S}=(\mat{G}^{\!\top}\mat{M}\mat{G})^{-1}$,
so that $\mat{P}=\mat{I}-\mat{G}\,\mat{S}\,\mat{G}^{\!\top}\mat{M}$. Compute
$\mat{P}^2$:
\[
  \mat{P}^2 = \bigl(\mat{I}-\mat{G}\mat{S}\mat{G}^{\!\top}\mat{M}\bigr)
              \bigl(\mat{I}-\mat{G}\mat{S}\mat{G}^{\!\top}\mat{M}\bigr)
            = \mat{I} - 2\,\mat{G}\mat{S}\mat{G}^{\!\top}\mat{M}
              + \mat{G}\mat{S}\underbrace{\mat{G}^{\!\top}\mat{M}\mat{G}}_{=\,\mat{S}^{-1}}\mat{S}\,\mat{G}^{\!\top}\mat{M}.
\]
In the last term the inner $\mat{G}^{\!\top}\mat{M}\mat{G}$ is exactly
$\mat{S}^{-1}$, so $\mat{S}\,\mat{S}^{-1}\,\mat{S}=\mat{S}$ and the term collapses
to $\mat{G}\mat{S}\mat{G}^{\!\top}\mat{M}$. Hence
\[
  \mat{P}^2 = \mat{I} - 2\,\mat{G}\mat{S}\mat{G}^{\!\top}\mat{M}
              + \mat{G}\mat{S}\mat{G}^{\!\top}\mat{M}
            = \mat{I} - \mat{G}\mat{S}\mat{G}^{\!\top}\mat{M} = \mat{P}.
\]
The cancellation of the middle inverse against the triple product is the entire
content of the projector law; it is why the recipe must use the \emph{exact}
inverse $\mat{S}$ (and why the inner solve must be converged tightly).

\medskip\noindent\textbf{Symbolic: gradients are killed.}
Apply $\mat{P}$ to a gradient $\mat{G}\psi$:
\[
  \mat{P}\,\mat{G}\psi
  = \mat{G}\psi - \mat{G}\mat{S}\mat{G}^{\!\top}\mat{M}\mat{G}\psi
  = \mat{G}\psi - \mat{G}\mat{S}\,\mat{S}^{-1}\psi
  = \mat{G}\psi - \mat{G}\psi = 0,
\]
again using $\mat{G}^{\!\top}\mat{M}\mat{G}=\mat{S}^{-1}$. So every gradient is in
the kernel of $\mat{P}$: $\mathrm{Ker}(\mat{P})=\mathrm{Range}(\mat{G})$, exactly
the curl null space of \cref{eq:hcurlsplit}.

\medskip\noindent\textbf{Numerical instance.}
Take the smallest nontrivial case: $\mat{M}=\mat{I}_2$ (the $2\times2$ identity,
so the $\mat{M}$-inner-product is ordinary $\Ltwo$), and a discrete gradient with
a single scalar DoF, $\mat{G}=\bigl[\begin{smallmatrix}1\\1\end{smallmatrix}\bigr]$
(the gradient subspace is the diagonal line spanned by $(1,1)$). Then
\[
  \mat{G}^{\!\top}\mat{M}\mat{G} = \begin{bmatrix}1&1\end{bmatrix}\begin{bmatrix}1\\1\end{bmatrix} = 2,
  \qquad \mat{S} = \tfrac12,
  \qquad
  \mat{G}\mat{S}\mat{G}^{\!\top}\mat{M}
  = \tfrac12\begin{bmatrix}1\\1\end{bmatrix}\begin{bmatrix}1&1\end{bmatrix}
  = \tfrac12\begin{bmatrix}1&1\\1&1\end{bmatrix}.
\]
Therefore
\[
  \mat{P} = \mat{I} - \tfrac12\begin{bmatrix}1&1\\1&1\end{bmatrix}
          = \tfrac12\begin{bmatrix}1&-1\\-1&1\end{bmatrix}.
\]
Check idempotence directly:
$\mat{P}^2=\tfrac14\bigl[\begin{smallmatrix}2&-2\\-2&2\end{smallmatrix}\bigr]
=\tfrac12\bigl[\begin{smallmatrix}1&-1\\-1&1\end{smallmatrix}\bigr]=\mat{P}$. Check
that gradients die: for any scalar $\psi$,
$\mat{P}\,\mat{G}\psi = \tfrac12\bigl[\begin{smallmatrix}1&-1\\-1&1\end{smallmatrix}\bigr]
\bigl[\begin{smallmatrix}\psi\\\psi\end{smallmatrix}\bigr]
= \tfrac12\bigl[\begin{smallmatrix}\psi-\psi\\-\psi+\psi\end{smallmatrix}\bigr]=\mathbf{0}$.
And the projection of a general field is its divergence-free (here, orthogonal-to-
diagonal) part: $\mat{P}\,(a,b)^{\!\top}=\tfrac12(a-b,\,b-a)^{\!\top}$, which is
the component of $(a,b)$ along the antidiagonal $(1,-1)$---perpendicular, in the
$\mat{M}=\mat{I}$ inner product, to the gradient direction $(1,1)$. The projector
keeps the circulating part and discards the spreading part, exactly as
\cref{fig:helmholtz} drew it.
\end{workedexample}

\section{Why the projection is needed: the curl--curl null space}
\label{sec:nullspace}

We can now say precisely why a simulator bothers with all of this. The curl--curl
operator of the magnetostatic and eigenmode problems (\cref{ch:reductions-pde})
is, by construction, singular: it has an enormous null space.

Recall the curl--curl stiffness $\mat{K}$ assembled from the bilinear form
$\int_\dom \nu\,(\Curl\vu)\cdot(\Curl\vv)$ of \cref{ch:weak}. Any gradient field
$\vu=\Grad\phi$ has zero curl ($\Curl\Grad\phi=0$, de Rham exactness,
\cref{ch:derham}), so it contributes nothing to the form: $\mat{K}\,\Grad\phi=0$.
The entire gradient subspace---which, on a fine mesh, is a substantial fraction of
all the DoFs---lies in the null space of $\mat{K}$. The operator is hugely
rank-deficient (\cref{fig:nullspace}).

\begin{figure}[t]
\centering
\begin{tikzpicture}[scale=1.0]
  % the whole Hcurl space as a big ellipse
  \draw[simInk,thick] (0,0) ellipse (3.6 and 2.2);
  \node[font=\small,simInk] at (0,2.55) {$\Hcurl$ (all field DoFs)};
  % the gradient null-space as a big shaded sub-blob
  \begin{scope}
    \clip (0,0) ellipse (3.6 and 2.2);
    \fill[simBgRust] (-3.6,-2.2) rectangle (0.4,2.2);
  \end{scope}
  \draw[simRust,thick] (0.4,-2.16) -- (0.4,2.16);
  \node[font=\small,simRust,align=center] at (-1.7,0)
    {$\mathrm{Ker}(\mat{K})$\\$=\mathrm{Range}(\Grad)$\\[1pt]\scriptsize gradient null space\\\scriptsize(spurious zero modes)};
  % the physical divergence-free subspace
  \node[font=\small,simBlue,align=center] at (1.9,0)
    {divergence-\\free part\\[1pt]\scriptsize the physics};
  % an eigensolver "drifting" into the null space
  \draw[backflow] (2.6,1.3) .. controls (1.0,1.6) and (-0.8,1.4) .. (-2.2,0.9);
  \node[font=\scriptsize\itshape,simRust] at (0.2,1.75) {solver drifts in};
  % the projector pushing back out
  \draw[depedge] (-2.0,-0.9) .. controls (-0.6,-1.5) and (1.0,-1.4) .. (2.4,-0.9);
  \node[font=\scriptsize\itshape,simBlue] at (0.2,-1.75) {$\mat{P}$ projects out};
\end{tikzpicture}
\caption[The curl--curl gradient null space]{Why the projection is needed. The
curl--curl operator $\mat{K}$ annihilates every gradient field
($\mat{K}\Grad\phi=0$, de Rham exactness), so its null space is the entire
gradient subspace $\mathrm{Range}(\Grad)$ (rust)---a large fraction of all the
field DoFs. These are \emph{spurious zero modes}: eigensolutions with eigenvalue
zero that carry no physics. An eigensolver or iterative solver, left unguarded,
drifts into this subspace (dashed) and returns garbage or stalls. The
divergence-free projector $\mat{P}$ pushes every iterate back onto the physical
divergence-free subspace (blue), so the solve sees only the well-posed part of the
operator.}
\label{fig:nullspace}
\end{figure}

For a solver this is poison in two ways. An \emph{eigensolver}
(\cref{ch:eigenmodes}) looking for the resonant modes
$\mat{K}\vx=\lambda\mat{M}\vx$ will find, in addition to the physical modes, a
flood of spurious modes with eigenvalue $\lambda\approx 0$---the gradient null
space, masquerading as zero-frequency ``resonances'' that are nothing of the kind.
They swamp the spectrum near zero and bury the physical low-frequency modes the
analysis actually wants. An \emph{iterative linear solver} fares no better: the
null space makes the operator singular (or, after a tiny regularization,
catastrophically ill-conditioned), and a Krylov method either stalls or wanders
into the null space and returns a contaminated answer.

\begin{pitfall}
The curl--curl operator's gradient null space is not a numerical artifact to be
regularized away with a small diagonal shift---it is a structural feature of the
physics and the de Rham complex. A shift would perturb the physical spectrum and
still leave near-null modes to confuse the solver. The correct cure is to
\emph{project the null space out} with $\mat{P}$, restricting the solve to the
divergence-free subspace where the operator is genuinely well-posed. Eigensolvers
in particular must project \emph{every} iterate (the initial vector and each
subsequent candidate), or the spurious modes creep back in.
\end{pitfall}

The remedy is the projector $\mat{P}$ of \cref{sec:divfree}. By construction its
range is the divergence-free subspace and its kernel is exactly the gradient null
space $\mathrm{Range}(\mat{G})=\mathrm{Ker}(\mat{K})$. Applying $\mat{P}$ to a
field removes precisely the troublesome component and leaves the physical part. In
the eigenmode analysis (\cref{ch:eigenmodes}) the projector is threaded into the
eigenvalue iteration: the initial vector is projected once, and every candidate
eigenvector is projected at each step, so the iteration is confined to the
divergence-free subspace throughout and never sees the spurious modes
(\cref{fig:eigproj}). The projector is the guardrail that keeps the eigensolver on
the physical manifold.

\begin{figure}[t]
\centering
\begin{tikzpicture}[node distance=8mm]
  \node[data, text width=16mm, align=center] (v0) {start vector $\vv_0$};
  \node[opbox, right=8mm of v0, text width=14mm, fill=simBgGreen, draw=simGreen] (p0)
    {$\mat{P}$ project};
  \node[stage, right=8mm of p0, text width=24mm, align=center] (step)
    {eigen step\\[1pt]\footnotesize (\cref{ch:eigenmodes})};
  \node[opbox, right=8mm of step, text width=14mm, fill=simBgGreen, draw=simGreen] (pk)
    {$\mat{P}$ project};
  \node[product, right=8mm of pk, text width=16mm, align=center] (conv)
    {converged?};
  \draw[flow] (v0) -- (p0);
  \draw[flow] (p0) -- (step);
  \draw[flow] (step) -- (pk);
  \draw[flow] (pk) -- (conv);
  % loop back
  \draw[backflow] (conv.south) -- ++(0,-7mm) -| node[below, pos=0.25, font=\scriptsize]{no: next candidate} (step.south);
  % output
  \node[data, right=8mm of conv, text width=16mm, align=center, fill=simBgBlue, draw=simBlue] (out) {physical mode};
  \draw[flow] (conv) -- node[above, font=\scriptsize]{yes} (out);
  \node[font=\scriptsize\itshape, simGreen, above=1mm of pk.north] {every iterate stays divergence-free};
\end{tikzpicture}
\caption[The projector inside the eigenvalue loop]{How the divergence-free
projector guards an eigensolve. The initial vector $\vv_0$ is projected once, and
\emph{every} candidate produced by the eigen step (\cref{ch:eigenmodes}) is
projected again with $\mat{P}$ before the convergence test and before it feeds the
next iteration (green boxes). Because $\mathrm{Ker}(\mat{P})$ is exactly the
gradient null space, the iteration is confined to the divergence-free subspace
throughout: the spurious zero modes of \cref{fig:nullspace} can never re-enter, and
the solver converges to the physical resonances alone. Projecting only the initial
vector is not enough---rounding and the iteration itself would let the null space
leak back in, which is why the projection sits \emph{inside} the loop.}
\label{fig:eigproj}
\end{figure}

\section{Closing Part II}

With boundary conditions imposed and fields projected onto their physical
subspace, the discretized problem is complete. Let us take stock of what Part~II
has built. We began (\cref{ch:weak}) by turning strong-form PDEs into weak forms,
$a(u,v)=\ell(v)$, drawn from a single four-operator family
$(Q,\diffop)$. We chose finite-element spaces for each differential operator and
saw them interlock as the de Rham complex (\cref{ch:functionspaces},
\cref{ch:derham}). We assembled the bilinear forms into sparse operators by a fold
over terms (\cref{ch:galerkin}, \cref{ch:assembly}), and we saw how that fold
runs matrix-free on the reference element (\cref{ch:refelement},
\cref{ch:matrixfree}). And in this chapter we performed the two finishing
operations---essential-BC elimination and divergence-free projection---that turn
the raw assembled system into a well-posed, physical one.

The result is a concrete linear-algebra problem: an operator $\mat{K}$, a
right-hand side $\vb$, perhaps a mass operator $\mat{M}$ for an eigenproblem, all
finished and ready. What we have \emph{not} done is solve it. Every analysis in
this book ends in a solve---a linear solve, an eigenvalue solve, a time march---and
those solves are where the simulator spends its time and where its deepest
structure lives. Part~III turns to the calculus that performs them: the
combinators and folds that express a Krylov iteration, a time integrator, an
eigenvalue sweep as compositions of pure, immutable-tensor operations. Part~II
built the problem; Part~III builds the machine that answers it.

\summarysection
Two finishing operations stand between the assembled system $\mat{K}\vx=\vb$ and
its solution, and both are \emph{separable post-compositions}---not steps inside
the assembly fold. \textbf{(1)~Essential boundary conditions.} The constrained
(Dirichlet) DoFs form an \emph{essential set} $E$; the rest are \emph{free}.
Imposing the condition partitions the system into free/essential blocks and
reduces it to the free block, $\mat{K}_{FF}\vx_F=\vb_F-\mat{K}_{FE}\vx_E$. In
place, \texttt{eliminate\_essential\_bc} zeroes the essential rows and columns and
sets the eliminated diagonal by a policy (\texttt{DIAG\_ONE} for the solve-side,
\texttt{DIAG\_ZERO} for the energy/mass-block), which is the free-block projection
$\mat{P}_F\mat{K}\mat{P}_F$ up to the chosen diagonal. A nonzero (inhomogeneous)
boundary value is handled by \emph{lifting}: \texttt{eliminate\_rhs} forms
$\vb\leftarrow\vb-\mat{K}\vx_{\mathrm{ess}}$, moving the known boundary
contribution into the right-hand side---a single matrix--vector product and an
\texttt{axpy}. Natural (Neumann) conditions, by contrast, enter during assembly as
boundary integrators and cost nothing when homogeneous. \textbf{(2)~Helmholtz
decomposition and divergence-free projection.} Any vector field splits uniquely as
$\vF=\vF_{\mathrm{df}}+\Grad\psi$ (divergence-free plus gradient), and at the
$\Hcurl$ slot of the de Rham complex the gradient part is exactly the kernel of
the curl. The discrete projector
$\mat{P}=\mat{I}-\mat{G}(\mat{G}^{\!\top}\mat{M}\mat{G})^{-1}\mat{G}^{\!\top}\mat{M}$
is an $\mat{M}$-orthogonal projection onto the divergence-free subspace; its apply
runs weak-divergence $\to$ essential-zeroing $\to$ inner Poisson solve $\to$
gradient subtract, and so is a constructed operator carrying an inner
\texttt{ksp\_solve}. The projection is needed because the curl--curl operator has
a huge gradient null space ($\mat{K}\Grad\phi=0$); without projection,
eigensolvers return spurious zero modes and iterative solvers stall. Both finishing
operations close the discretization that Part~II built.

\furtherreading
The variational treatment of boundary conditions---essential versus natural, the
trial/test-space constraint, and the lifting of inhomogeneous data---is developed
in any rigorous finite-element text; for the vector ($\Hcurl$, $\Hdiv$) case
specific to Maxwell problems, Monk~\citep{monk2003} is the standard reference, with
careful treatment of tangential-trace boundary conditions and the discrete spaces
on which they live. The Helmholtz decomposition and its de Rham reading are the
heart of finite-element exterior calculus; Arnold, Falk, and
Winther~\citep{arnold2006feec} give the definitive modern account, in which the
divergence-free projection is the Hodge decomposition at the relevant slot of the
discrete complex. The auxiliary-space and projection techniques that make the
curl--curl null space tractable for iterative solvers---the practical engine behind
the inner Poisson solve of \cref{sec:divfree}---are the subject of Hiptmair and
Xu~\citep{hiptmairxu2007}, whose AMS (auxiliary-space Maxwell solver) preconditioner
is built on exactly the discrete-gradient and mass operators that appear in
\cref{eq:projector}.

\exercisessection
\begin{enumerate}[nosep]
  \item \textbf{(Reading the partition.)} On the mesh of \cref{fig:partition}, the
  essential set $E$ is the boundary nodes and the free set $F$ is the interior.
  If the bottom edge is held at voltage $V$ and the top edge grounded ($0$), with
  the side edges insulating (homogeneous Neumann), which nodes are in $E$ and which
  in $F$? Which essential nodes carry $V$, which carry $0$, and what happens to the
  side-edge nodes?

  \item \textbf{(The two policies.)} Explain in one or two sentences why the
  \texttt{DIAG\_ONE} policy (essential diagonal $=\mat{I}_E$) is appropriate when
  the eliminated operator will be inverted by a linear solver, while
  \texttt{DIAG\_ZERO} (essential diagonal $=\mathbf{0}$) is used for the blocks of
  a generalized eigenvalue problem. What would go wrong if you used
  \texttt{DIAG\_ZERO} for the linear solve, or \texttt{DIAG\_ONE} for an eigenvalue
  block? (Hint: think about what eigenvalue an identity diagonal block contributes.)

  \item \textbf{(Lifting, by hand.)} Repeat \cref{wex:lift1d} but with both ends
  inhomogeneous: $u(0)=a$, $u(1)=b$. Form $\vx_{\mathrm{ess}}=(a,0,0,b)^{\!\top}$,
  compute $\mat{K}\vx_{\mathrm{ess}}$ and the lifted RHS
  $\vb'=\vb-\mat{K}\vx_{\mathrm{ess}}$, and solve the reduced free system. Confirm
  the answer is the linear interpolant $u(x)=a+(b-a)x$ sampled at the nodes.

  \item \textbf{(Separability.)} Suppose $\mat{K}=\mat{K}_1+\mat{K}_2$ is the sum
  of two assembled term-contributions (say a stiffness plus a mass term). Using the
  projection form \cref{eq:pfproj}, show that under the \texttt{DIAG\_ZERO} policy
  $\texttt{eliminate\_essential\_bc}(\mat{K}_1+\mat{K}_2,E)
  =\texttt{eliminate\_essential\_bc}(\mat{K}_1,E)+\texttt{eliminate\_essential\_bc}(\mat{K}_2,E)$.
  Then explain why the same distributivity fails for \texttt{DIAG\_ONE}, and state
  the corrected identity (the identity $\mat{I}_E$ is added once, not twice).

  \item \textbf{(Idempotence of the projector.)} Verify $\mat{P}^2=\mat{P}$ for the
  numerical instance of \cref{wex:projcheck} by direct matrix multiplication of
  $\mat{P}=\tfrac12\bigl[\begin{smallmatrix}1&-1\\-1&1\end{smallmatrix}\bigr]$. Then
  construct a second instance with $\mat{M}=\mathrm{diag}(2,1)$ (a nontrivial
  $\mat{M}$-inner-product) and the same
  $\mat{G}=\bigl[\begin{smallmatrix}1\\1\end{smallmatrix}\bigr]$: compute
  $\mat{G}^{\!\top}\mat{M}\mat{G}$, $\mat{S}$, and $\mat{P}$, and confirm $\mat{P}$
  is still idempotent and still annihilates $\mat{G}\psi$. (You will find $\mat{P}$
  is no longer symmetric---it is $\mat{M}$-orthogonal, not ordinary-orthogonal.)

  \item \textbf{(Why gradients are killed.)} Using only the identity
  $\mat{G}^{\!\top}\mat{M}\mat{G}=\mat{S}^{-1}$, give the two-line proof that
  $\mat{P}\,\mat{G}\psi=0$ for every $\psi$ (i.e.\
  $\mathrm{Ker}(\mat{P})\supseteq\mathrm{Range}(\mat{G})$). Then explain, in one
  sentence each, why this is exactly the statement that $\mat{P}$ removes the
  curl--curl null space, and why the inner solve in step~3 of
  \cref{fig:projapply} must be converged tightly for the cancellation to hold.

  \item \textbf{(Spurious modes.)} The curl--curl eigenproblem
  $\mat{K}\vx=\lambda\mat{M}\vx$ has every discrete gradient $\mat{G}\phi$ as an
  eigenvector. What is its eigenvalue $\lambda$? Explain why these modes cluster at
  the bottom of the spectrum and why an eigensolver seeking the lowest physical
  resonances is exactly the one most confused by them. How does projecting every
  iterate with $\mat{P}$ remove the problem?

  \item \textbf{(Natural versus essential, bookkeeping.)} A driven-cavity problem
  has a PEC wall (tangential $\vE=0$, essential), an impedance wall (a natural
  Robin-type condition contributing a boundary mass term), and a port (an
  inhomogeneous essential excitation). For each of the three, state: (i) whether it
  is imposed during assembly or as a post-composition, (ii) whether it touches the
  operator $\mat{K}$, the right-hand side $\vb$, or both, and (iii) whether a
  homogeneous version of it would cost anything. Organize your answer as a small
  table.
\end{enumerate}


\part{The Abstract Framework}\label{part:framework}
\chapter{A Language for Computations: Values and Functions}
\label{ch:language}

\chapterorient{In \cref{ch:bigidea} we made a decision: describe a simulation as
a chain of pure functions over immutable tensors, not as overwrites of a block of
memory. That decision earns its keep only if we can write the functions down
\emph{precisely}---precisely enough that the written form \emph{is} the dependency
graph, precisely enough that a critic can check one piece against the source
without running anything. This chapter introduces the small notation we use for
exactly that, and it teaches you to read it from zero. We assume no
functional-programming background at all. By the end you will be able to look at a
line like \lstinline[style=l4style]|axpy :: Scalar -> Tensor[N] -> Tensor[N] -> Tensor[N]|
and say, token by token, what it means---and you will understand the one idea that
trips up every newcomer: a function that takes, or returns, another function.}

\section{Why a language at all}

We could describe every algorithm in this book in ordinary English prose, or in
ordinary mathematics. We will, in fact, lean on both constantly. But neither is
quite enough on its own, and the gap is worth naming, because it tells us exactly
what the small notation of this part is \emph{for}.

Recall the three questions the calculus is built to answer
(\cref{ch:bigidea}): \emph{what operations happen}, \emph{who owns the state},
and \emph{how state evolution is coordinated}. Each is a structural question---a
question about the \emph{shape} of a computation, not about the numbers it
produces. English prose blurs structure: ``then we update the residual'' hides
whether the residual is a fresh value or an overwrite, whether the update reads
any state besides its stated inputs, whether it is a step in a loop or a one-off.
Ordinary mathematics fixes the arithmetic but says nothing about ownership or
coordination at all: the equation $\vr = \vb - \mat{A}\vx$ tells us what $\vr$
\emph{equals}, but not that $\vr$ is a brand-new immutable value, nor that
computing it is a pure function of exactly $\vb$, $\mat{A}$, and $\vx$ and
nothing else.

So we adopt a third notation---a small, deliberately readable \emph{pseudo-language}
borrowed from typed functional programming---that makes the structure
explicit. This chapter does the first of the three questions, the foundation the
other two rest on: \emph{values and functions}. What is a value? What is a
function? How do we write a function's type, supply its arguments, and---the
hard part---pass functions to other functions? \Cref{ch:records} adds records
(bundles of named fields); \cref{ch:tensors} adds shapes; \cref{ch:fold} adds
the loop combinator; \cref{ch:monad} adds the discipline for threaded state. Each
builds on this one.

\begin{notationbox}
The pseudo-language is a language you \emph{read}, not one you run. It always
appears in a shaded, monospaced box. It is not a burn API, not a Rust sketch, not
something a compiler accepts; it is a precise notation for talking about
algorithms, the way a physicist reads a Feynman diagram. The notation primer at
the front of the book sketched it in one page; this chapter is the careful
version. When you see a shaded box you are reading code; when you see an inline
formula like $\mat{K}\vx=\vb$ you are reading mathematics. The two never mix in
one expression.
\end{notationbox}

\section{Values: things you can name but never change}

The most foundational idea in the whole notation is also the simplest, and you
already met it in \cref{ch:bigidea}: a \emph{value} is a thing you can name but
never change.

The number $7$ is a value. So is the number $3$. When you write $3 + 4$ you do not
\emph{turn} the $3$ into a $7$; you compute a \emph{new} value, $7$, and the $3$
and the $4$ sit there, untouched, exactly as before. Nobody is tempted to ask
``but what is the $3$ \emph{now}, after the addition?''---the question is
nonsense. A number has no ``now.'' It simply is what it is.

The single rule of our notation is that \emph{everything} behaves this way. A
tensor is a value, like $7$. A record (a bundle of named fields, \cref{ch:records})
is a value. A function---we will see---is a value. None of them can be changed
once it exists. You can compute new values \emph{from} them, name those new
values, and pass them around; you can never reach inside an existing value and
overwrite a part of it.

\begin{definition}[Value and immutability]
A \emph{value} is a piece of data---a number, a tensor, a record, a function---that,
once it exists, never changes. Computing with values means producing \emph{new}
values from old ones. A name, once bound to a value, refers to that same
unchanging value for as long as the name is in scope. We call this discipline
\emph{immutability}.
\end{definition}

\subsection{The contrast: a variable as a mutable box}

If you have programmed in C, Python, or most mainstream languages, you carry a
\emph{different} mental model of what a name means, and it is worth dragging into
the light because it will fight you here. In those languages a variable is a
\emph{box}: a labeled slot in memory that holds a value \emph{now} but can be
made to hold a different value \emph{later}.

\begin{lstlisting}[style=l4style, language=]
y = 5      # the box labeled y now holds 5
y = y + 1  # overwrite: the box labeled y now holds 6
\end{lstlisting}

\noindent After the second line, the box \texttt{y} holds $6$; the $5$ is gone.
The \emph{name} \texttt{y} survived but what it points to was destroyed and
replaced. This is the in-place-mutation picture of \cref{ch:bigidea}, now at the
level of a single variable rather than a whole array: one box, scribbled over.

Our notation has no such boxes. When we write

\begin{lstlisting}[style=l4style]
let y = 5 in
let z = y + 1 in
...
\end{lstlisting}

\noindent the name \texttt{y} is bound, once and forever (within its scope), to
the value $5$. We do not---cannot---``update \texttt{y} to $6$.'' Instead we
introduce a \emph{new} name, \texttt{z}, for the new value $6$. There are two
names and two values; nothing was overwritten. The \code{let} keyword reads
``\emph{let} this name stand for this value, in what follows.''

\begin{figure}[t]
\centering
\begin{tikzpicture}[node distance=8mm]
  % ============ LEFT: mutable box ============
  \begin{scope}
    \node[font=\bfseries\small, simRust] at (0,2.5) {A mutable box (a variable)};
    \node[draw=simRust, fill=simBgRust, thick, rounded corners=2pt,
          minimum width=24mm, minimum height=11mm] (box) at (0,1.1) {};
    \node[font=\ttfamily\small, simInk] at (-0.7,1.1) {y};
    \node[font=\small, simInk] at (0.45,1.1) {$\leftarrow 5$};
    % overwrite loop
    \draw[backflow] (box.north west) to[out=135,in=45,looseness=4.5]
      node[above,font=\scriptsize]{\texttt{y := 6}} (box.north east);
    \node[align=center, font=\scriptsize\itshape, simGray, text width=40mm]
      at (0,-0.15) {one slot; the old contents are destroyed and replaced.
      ``What is \texttt{y} now?'' is a real question.};
  \end{scope}
  % ============ RIGHT: immutable values ============
  \begin{scope}[shift={(7.2,0)}]
    \node[font=\bfseries\small, simBlue] at (0,2.5) {Immutable values (names)};
    \node[draw=simBlue, fill=simBgBlue, thick, rounded corners=2pt,
          minimum width=12mm, minimum height=10mm] (v5) at (-0.9,1.1) {$5$};
    \node[draw=simBlue, fill=simBgBlue, thick, rounded corners=2pt,
          minimum width=12mm, minimum height=10mm] (v6) at (1.1,1.1) {$6$};
    \node[font=\ttfamily\scriptsize, simBlue, below=0.5mm of v5] {y};
    \node[font=\ttfamily\scriptsize, simBlue, below=0.5mm of v6] {z};
    \draw[flow] (v5) -- node[above,font=\scriptsize]{$+1$} (v6);
    \node[align=center, font=\scriptsize\itshape, simGray, text width=42mm]
      at (0.1,-0.15) {two distinct values, two names; \texttt{y} still
      means $5$. ``What is \texttt{y} now?'' is a non-question.};
  \end{scope}
\end{tikzpicture}
\caption[A mutable box versus an immutable value]{Two readings of a name.
\emph{Left:} the mainstream model---a variable is a box in memory; \texttt{y := 6}
overwrites it, and the previous contents vanish, so the box's value depends on how
far the program has run. \emph{Right:} our model---a name is bound permanently to
an immutable value; computing $y+1$ produces a \emph{fresh} value $6$, which we
give a \emph{new} name \texttt{z}. The original \texttt{y} still means $5$. The
notation in this book has only the right-hand picture; there are no boxes.}
\label{fig:valuebox}
\end{figure}

\Cref{fig:valuebox} draws the two pictures side by side, the same contrast
\cref{fig:mutvalue} drew for whole arrays, now for a single name. The discipline
is the same at every scale: never overwrite, always produce a fresh value and
name it.

\begin{keyidea}
Immutability is the ground rule of the entire notation. A value---number, tensor,
record, or function---never changes after it is created. A name binds to a value
permanently within its scope. We never overwrite; we compute new values from old
ones and give them new names. Every other convention in this part is built on
top of this one. If you ever find yourself asking ``what does this name hold
\emph{now}?'', you have slipped back into the mutable-box model---stop and
recover the value picture.
\end{keyidea}

\subsection{Why \code{clone()} is meaningless here}

A quick but illuminating consequence. In a mutable language, when you want a
function to work on a copy of your data without disturbing the original, you
\emph{clone} it: you make a second box holding the same contents, so that
overwriting one leaves the other intact. The lattice-Boltzmann reference code in
\cref{ch:bigidea} was littered with \texttt{current.clone()} for exactly this
reason.

In our notation, \code{clone()} is meaningless---there is nothing to clone,
because there is nothing to protect. A value cannot be overwritten, so there is no
way for one use of a value to disturb another use of it. Two names for ``the same
tensor'' is perfectly safe: neither can mutate it, so neither can surprise the
other.

\begin{lstlisting}[style=l4style]
let v = some_tensor in
let a = scale 2 v in   -- uses v
let b = scale 3 v in   -- ALSO uses v; no clone needed
(a, b)                 -- v was never disturbed; it could not be
\end{lstlisting}

\noindent This is not a small convenience. ``When must I copy, to be safe?'' is
one of the central, error-prone questions of mutable programming, and in the
value picture it simply evaporates. The cost we appear to pay---``but surely
making a fresh tensor every step is wasteful!''---is recovered in full by the
implementation, which is free to reuse storage whenever an old value is no longer
referenced (the pitfall in \cref{ch:bigidea}; made precise in \cref{ch:monad}).
We describe the computation with fresh values; the machine optimizes the memory.

\section{Functions and types}

A \emph{function} is a rule that turns inputs into an output. In our notation a
function is written in two parts: a \emph{type signature} that says what kinds of
thing go in and come out, and a \emph{definition} that says how the output is
computed. Here is the smallest possible example---a function that squares a
floating-point number.

\begin{lstlisting}[style=l4style]
square :: Float -> Float       -- the TYPE: takes a Float, gives a Float
square x = x * x               -- the DEFINITION: the output is x times x
\end{lstlisting}

Read the two lines carefully, because their grammar recurs on every function in
the book.

\begin{description}[style=nextline,leftmargin=1.5em]
  \item[The signature \code{square :: Float -> Float}.] The
  double colon \code{::} reads ``\emph{has type}.'' So this line says
  ``\code{square} has type \code{Float -> Float}.'' The arrow \code{->} reads
  ``\emph{maps to}'': \code{Float -> Float} is the type of functions that map a
  \code{Float} to a \code{Float}. The signature is a \emph{contract}: it promises
  that if you hand \code{square} a \code{Float}, you get a \code{Float} back, and
  it forbids handing it anything else.
  \item[The definition \code{square x = x * x}.] Here \code{x}
  names the input (the \emph{parameter}); the right-hand side \lstinline[style=l4style]|x * x|
  is the output, an expression built from the input. Applying the function is
  written by \emph{juxtaposition}---just putting the argument next to the function
  name, with a space: \lstinline[style=l4style]|square 5| evaluates to $25$. No
  parentheses are needed around the argument (though \lstinline[style=l4style]|square (a + b)|
  needs them, to group the argument).
\end{description}

\noindent The types are not decoration. They are the first line of defence
against nonsense: a tool (and a careful reader) can check that every function is
only ever applied to arguments of the type it expects, long before any numbers
are computed. The shapes of tensors, which we develop in \cref{ch:tensors}, are
types in exactly this sense---and they catch the overwhelming majority of
mistakes in tensor code.

\subsection{Pure functions}

The functions in this book obey one further, crucial discipline: they are
\emph{pure}. A pure function's output depends on \emph{nothing but its declared
inputs}, and computing it changes nothing in the outside world. Two consequences,
both load-bearing:

\begin{itemize}[nosep]
  \item \textbf{Same inputs, same output, always.} \lstinline[style=l4style]|square 5|
  is $25$ today, tomorrow, on every machine, no matter what ran before it. A pure
  function has no memory and no mood.
  \item \textbf{No side effects.} A pure function does not secretly read a global
  variable, write to a file, overwrite one of its inputs, or depend on the time
  of day. Everything it uses arrives through its parameters; everything it
  produces leaves through its return value.
\end{itemize}

This is precisely the property that makes the program text equal to the
dependency graph (\cref{ch:bigidea}): if a function can depend only on its
inputs, then drawing an arrow from each input to the function captures
\emph{everything} the function depends on---there are no hidden arrows. Purity is
what we buy with immutability, and the rest of the book spends it.

\begin{pitfall}
The mutable-box habit dies hard, and it shows up as a specific reading error.
When you see \lstinline[style=l4style]|axpy a x y = a * x + y|, the C programmer's
instinct is to ask ``and which of \code{x}, \code{y} did it overwrite?'' The
answer is \emph{neither}. A pure function \emph{returns} its result; it does not
deposit it into one of its arguments. After \lstinline[style=l4style]|let r = axpy a x y|,
the names \code{x} and \code{y} still refer to exactly the tensors they did
before, and \code{r} names the fresh result. If you ever catch yourself hunting
for ``the output argument,'' you have slipped models---there is no output
argument, only a returned value.
\end{pitfall}

\subsection{A function is itself a value}

Here is the move that makes the whole notation expressive, stated plainly now and
exploited heavily later: \emph{a function is a value, just like a number or a
tensor.} It is a thing you can name, pass to another function as an argument, and
return from a function as a result. \code{square} is a value of type
\code{Float -> Float}, exactly as $5$ is a value of type \code{Float}.

This is not how most engineers first learn to think about functions---a function
feels like a \emph{verb}, a thing you \emph{do}, not a \emph{noun}, a thing you
\emph{hold}. But in this notation it is both. \Cref{fig:fnvalue} draws the
cartoon: a function is a little box you can carry around in a bag, hand to
someone, or get handed back. Sections~\ref{sec:currying}
and~\ref{sec:higherorder} are entirely about what this buys us.

\begin{figure}[t]
\centering
\begin{tikzpicture}[node distance=10mm]
  % a "bag" of values, one of which is a function
  \node[font=\bfseries\small, simInk] at (0,2.2) {Values you can name, pass, and return};
  % number value
  \node[draw=simBlue, fill=simBgBlue, thick, rounded corners=2pt,
        minimum width=15mm, minimum height=10mm] (n) at (-3.6,0.7) {$5$};
  \node[font=\scriptsize\itshape, simGray, below=1mm of n] {a number};
  % tensor value
  \node[draw=simBlue, fill=simBgBlue, thick, rounded corners=2pt,
        minimum width=18mm, minimum height=10mm] (t) at (-1.1,0.7)
        {\lstinline[style=l4style]|Tensor[N]|};
  \node[font=\scriptsize\itshape, simGray, below=1mm of t] {a tensor};
  % function value, drawn as a little machine box
  \node[draw=simTeal, fill=simBgGreen, thick, rounded corners=2pt,
        minimum width=30mm, minimum height=12mm] (f) at (2.3,0.7) {};
  \node[font=\ttfamily\scriptsize, simInk] at (1.55,0.95) {square};
  \node[font=\scriptsize, simTeal] at (2.55,0.55)
        {$x \mapsto x{*}x$};
  \node[font=\scriptsize\itshape, simGray, below=1mm of f]
        {a function---also a value};
  % a "carry" arrow suggesting you can pass it around
  \draw[flow] (3.95,0.7) -- ++(1.1,0)
    node[right,font=\scriptsize\itshape,simGray,align=left] {pass it,\\return it,\\name it};
\end{tikzpicture}
\caption[A function is a value you can pass around]{In this notation a function is
a value with the same standing as a number or a tensor: you can bind it to a name,
hand it to another function as an argument, and get one back as a result.
\code{square} is a value of type \code{Float -> Float} the way $5$ is a value of
type \code{Float}. This single idea---functions are first-class values---is what
makes currying (\cref{sec:currying}) and higher-order functions
(\cref{sec:higherorder}) possible.}
\label{fig:fnvalue}
\end{figure}

\section{Currying: one argument at a time}
\label{sec:currying}

Most interesting functions take more than one argument. We write a two-argument
function's type with two arrows:

\begin{lstlisting}[style=l4style]
add :: Float -> Float -> Float
add x y = x + y
\end{lstlisting}

\noindent and apply it by listing the arguments after the name, separated by
spaces: \lstinline[style=l4style]|add 3 4| evaluates to $7$. So far this looks
like the multi-argument functions of any language. But the arrow notation is
hinting at something subtler, and the subtlety pays off enormously, so we take it
slowly.

\subsection{Reading \code{A -> B -> C}}

How should we read \lstinline[style=l4style]|add :: Float -> Float -> Float|? The
honest reading is \emph{one argument at a time}:

\begin{quote}
\emph{``\code{add} takes a \code{Float}, and returns a function that takes a
\code{Float} and returns a \code{Float}.''}
\end{quote}

\noindent That is, the arrow \code{->} groups \emph{to the right}: the type is
secretly
\begin{lstlisting}[style=l4style]
add :: Float -> (Float -> Float)
\end{lstlisting}
\noindent with invisible parentheses around the tail. \code{add} is not really a
two-argument function at all; it is a one-argument function whose \emph{result is
another function}. This way of treating multi-argument functions---as a chain of
one-argument functions---is called \emph{currying}, after the logician Haskell
Curry (whose first name, not coincidentally, also names the language our notation
borrows from).

When you write \lstinline[style=l4style]|add 3 4|, two things happen in
sequence:

\begin{lstlisting}[style=l4style]
add 3       -- supply the first argument: get back a FUNCTION,
            --   "the function that adds 3 to its argument"
(add 3) 4   -- supply the second argument to THAT function: get 7
\end{lstlisting}

\noindent The juxtaposition \lstinline[style=l4style]|add 3 4| is just
\lstinline[style=l4style]|(add 3) 4| with the parentheses dropped---application,
too, groups to the left. So ``calling a two-argument function'' is really
``calling a one-argument function, then calling its one-argument result.'' You
almost never have to think about this when you supply all the arguments at once;
it matters when you supply only \emph{some}.

\subsection{Partial application}

Because \lstinline[style=l4style]|add 3| is itself a perfectly good function value,
we can name it and use it later:

\begin{lstlisting}[style=l4style]
let add3 = add 3 in   -- supply ONE of two arguments; add3 :: Float -> Float
let u = add3 4 in     -- u = 7
let w = add3 10 in     -- w = 13
(u, w)
\end{lstlisting}

\noindent Supplying \emph{some} of a function's arguments and getting back a
function that awaits the rest is called \emph{partial application}. We applied
\code{add} to one argument and \emph{froze} it---baked the $3$ in---producing a
new, more specialized function \code{add3} that adds $3$ to whatever it is later
given. \Cref{fig:currying} draws the freezing.

\begin{figure}[t]
\centering
\begin{tikzpicture}[node distance=9mm]
  % the curried machine, three "slots"
  \node[font=\bfseries\small, simInk] at (0,2.55)
    {Partial application: freeze some arguments now, supply the rest later};
  % stage 0: the full function
  \node[draw=simTeal, fill=simBgGray, thick, rounded corners=2pt,
        minimum width=20mm, minimum height=10mm] (f0) at (-4.3,1.0) {};
  \node[font=\ttfamily\scriptsize] at (-4.3,1.28) {add};
  \node[font=\scriptsize, simTeal] at (-4.3,0.78) {needs $x,y$};
  \node[font=\scriptsize\itshape,simGray,below=0.5mm of f0]
    {\code{Float -> Float -> Float}};
  % supply 3
  \draw[flow] (-3.2,1.0) -- node[above,font=\scriptsize]{supply $x=3$} (-1.7,1.0);
  % stage 1: partially-applied, one slot frozen
  \node[draw=simBlue, fill=simBgBlue, thick, rounded corners=2pt,
        minimum width=22mm, minimum height=10mm] (f1) at (-0.4,1.0) {};
  \node[font=\ttfamily\scriptsize] at (-0.4,1.30) {add3 = add 3};
  \node[font=\scriptsize, simBlue] at (-0.4,0.80) {needs $y$};
  \node[font=\scriptsize\itshape,simGray,below=0.5mm of f1]
    {\code{Float -> Float}};
  % a little "frozen 3" tag
  \node[draw=simRust,fill=simBgRust,rounded corners=1pt,font=\scriptsize,
        inner sep=1.5pt] at (-0.4,2.0) {$x{=}3$ frozen in};
  % supply 4
  \draw[flow] (0.8,1.0) -- node[above,font=\scriptsize]{supply $y=4$} (2.3,1.0);
  % stage 2: a value
  \node[draw=simRust, fill=simBgRust, thick, rounded corners=2pt,
        minimum width=14mm, minimum height=10mm] (r) at (3.4,1.0) {$7$};
  \node[font=\scriptsize\itshape,simGray,below=0.5mm of r] {a \code{Float}};
\end{tikzpicture}
\caption[Currying and partial application]{A curried function is consumed one
argument at a time. \code{add} needs two \code{Float}s; supplying the first
($x=3$) does not produce a number---it produces a \emph{new function} \code{add3},
of type \code{Float -> Float}, with the $3$ \emph{frozen in}. Supplying the
second ($y=4$) to \emph{that} function finally produces the value $7$. Freezing
some arguments now and supplying the rest later is \emph{partial application}; it
is the mechanism behind ``build once with configuration baked in, apply many
times.''}
\label{fig:currying}
\end{figure}

\subsection{Why partial application matters}

Partial application looks like a curiosity, but it is one of the most useful
patterns in the entire book, and it is worth seeing why \emph{now}, even before we
can spell the example fully. Over and over, an algorithm has a piece that is
\emph{configured once} and then \emph{applied many times}: a linear operator
assembled from the mesh and materials, then applied at every iteration of a solve;
a preconditioner factorized once, then applied every step; a time-step rule with
the step size $\Delta t$ baked in, then marched thousands of times.

Currying gives us a clean way to say exactly this. A constructor function takes
the configuration as its first arguments, and \emph{returns a function} that takes
the run-time data:

\begin{lstlisting}[style=l4style]
-- schematically: configuration first, then the run-time argument
make_operator :: Config -> (Tensor[N] -> Tensor[N])
let apply_A = make_operator cfg in   -- build ONCE, configuration frozen in
let y1 = apply_A x1 in               -- apply MANY times, cheaply
let y2 = apply_A x2 in
...
\end{lstlisting}

\noindent We build \code{apply\_A} once---freezing the configuration \code{cfg}
in---and then apply it again and again to different vectors. This ``build once,
apply many times'' shape is the backbone of \emph{constructed operators}, which
get a full chapter of their own (\cref{ch:operators}). Notice the parentheses in
the signature, \lstinline[style=l4style]|Config -> (Tensor[N] -> Tensor[N])|: we
will explain them carefully in \cref{sec:higherorder}---they are a deliberate
signal that \code{make\_operator} \emph{returns a function}.

\section{Higher-order functions}
\label{sec:higherorder}

We now reach the one genuinely new idea of the chapter---new, that is, to a reader
coming from ordinary imperative programming. It is the idea we have been circling:
a function can \emph{take} another function as an argument, or \emph{return} one as
a result. A function that does either is called a \emph{higher-order function}.

This is the single biggest conceptual lift in the notation, so we build it in two
halves: first a function that \emph{takes} a function, then a function that
\emph{returns} one. Take your time with both; everything in \cref{part:framework}
leans on them.

\subsection{Half one: a function that \emph{takes} a function}

Because a function is a value (\cref{fig:fnvalue}), nothing stops us from writing
a function whose \emph{parameter} is itself a function. Here is the cleanest tiny
example: \code{apply\_twice} takes a function \code{f} and a value \code{x}, and
applies \code{f} to \code{x} twice.

\begin{lstlisting}[style=l4style]
apply_twice :: (Float -> Float) -> Float -> Float
apply_twice f x = f (f x)
\end{lstlisting}

\noindent Read the signature one piece at a time. The first argument has type
\lstinline[style=l4style]|(Float -> Float)|---note the parentheses: this argument
is \emph{a function} from \code{Float} to \code{Float}. The second argument is a
plain \code{Float}. The result is a \code{Float}. The body
\lstinline[style=l4style]|f (f x)| says: apply \code{f} to \code{x}, then apply
\code{f} again to that. We can hand it our \code{square} from before:

\begin{lstlisting}[style=l4style]
apply_twice square 3    -- = square (square 3) = square 9 = 81
add3_twice = apply_twice add3       -- partially applied: freeze f = add3
add3_twice 10                       -- = add3 (add3 10) = add3 13 = 16
\end{lstlisting}

\noindent The parentheses around \lstinline[style=l4style]|(Float -> Float)| in
the signature are essential and worth pausing on. Without them,
\lstinline[style=l4style]|Float -> Float -> Float| would read (by the
right-grouping rule) as a function taking \emph{two} \code{Float}s---an ordinary
two-argument numeric function, like \code{add}. The parentheses say ``no: the
\emph{first} argument is itself a function.'' Grouping is the whole difference
between ``takes a function'' and ``takes two numbers.''

This pattern---pass a function in, let the surrounding function decide how to use
it---is exactly how the combinators of \cref{ch:fold} work. \code{map} takes a
function and applies it to every element of a list. \code{foldl} takes a function
and threads it through a list, accumulating. \code{iterate\_while}, our central
loop combinator, takes a \emph{step} function and applies it over and over. Each
is a function that takes a function; \code{apply\_twice} is the toy version of all
of them.

\begin{figure}[t]
\centering
\begin{tikzpicture}[node distance=10mm]
  % ===== LEFT: a function flowing IN to another function =====
  \begin{scope}
    \node[font=\bfseries\small, simInk] at (0,2.5) {A function flows \emph{in}};
    % the outer machine
    \node[draw=simBlue, fill=simBgBlue, thick, rounded corners=2pt,
          minimum width=30mm, minimum height=20mm] (outer) at (0,0.6) {};
    \node[font=\ttfamily\scriptsize, simInk] at (0,1.35) {apply\_twice};
    % a little function-box being inserted as an argument
    \node[draw=simTeal, fill=simBgGreen, thick, rounded corners=2pt,
          minimum width=13mm, minimum height=7mm] (fin) at (-2.9,0.6)
          {\ttfamily\scriptsize f};
    \draw[flow] (fin) -- (outer.west|-0,0.6);
    \node[font=\scriptsize, simInk] at (0,0.35) {uses \texttt{f} inside};
    \node[draw=simRust,fill=simBgRust,rounded corners=1pt,minimum width=11mm,
          minimum height=7mm] (xin) at (-2.9,-0.6) {\scriptsize $x$};
    \draw[flow] (xin) -- (outer.west|-0,-0.4);
    \draw[flow] (outer.east) -- ++(1.0,0) node[right,font=\scriptsize]{\texttt{f (f x)}};
    \node[align=center,font=\scriptsize\itshape,simGray,text width=42mm]
      at (0,-1.85) {a higher-order function \emph{takes} a function as an argument
      and calls it internally (\code{map}, \code{foldl}, \code{iterate\_while}).};
  \end{scope}
  % ===== RIGHT: a closure flowing OUT, capturing a value =====
  \begin{scope}[shift={(7.4,0)}]
    \node[font=\bfseries\small, simInk] at (0,2.5) {A closure flows \emph{out}};
    % the factory machine
    \node[draw=simBlue, fill=simBgBlue, thick, rounded corners=2pt,
          minimum width=24mm, minimum height=14mm] (fac) at (0,0.9) {};
    \node[font=\ttfamily\scriptsize, simInk] at (0,1.45) {scaled\_by};
    % a captured value going in
    \node[draw=simRust,fill=simBgRust,rounded corners=1pt,minimum width=11mm,
          minimum height=7mm] (cap) at (-2.7,0.9) {\scriptsize $a$};
    \draw[flow] (cap) -- (fac.west);
    \node[font=\scriptsize, simInk] at (0,0.6) {captures $a$};
    % the closure coming out, with the captured a baked in
    \node[draw=simTeal, fill=simBgGreen, thick, rounded corners=2pt,
          minimum width=22mm, minimum height=9mm] (clo) at (0,-1.0) {};
    \node[font=\ttfamily\scriptsize] at (0,-0.75) {\textbackslash v -> a .* v};
    \node[draw=simGold,fill=simBgGold,rounded corners=1pt,font=\scriptsize,
          inner sep=1.3pt] at (0,-1.35) {$a$ baked in};
    \draw[flow] (fac.south) -- (clo.north);
    \node[align=center,font=\scriptsize\itshape,simGray,text width=44mm]
      at (0,-2.25) {a higher-order function \emph{returns} a function that has
      \emph{captured} the value $a$---a \emph{closure}.};
  \end{scope}
\end{tikzpicture}
\caption[The two halves of higher order]{The two ways a function can be
higher-order. \emph{Left:} \code{apply\_twice} \emph{takes} a function \code{f} as
an argument and calls it inside its own body---the pattern of every combinator in
\cref{ch:fold}. \emph{Right:} \code{scaled\_by} (built next) \emph{returns} a
function, and that returned function has \emph{captured} the value $a$ supplied to
the factory---a \emph{closure}, a function with some data baked into it. Functions
flow both in and out.}
\label{fig:higherorder}
\end{figure}

\subsection{Half two: a function that \emph{returns} a function (a closure)}

The other half is a function whose \emph{result} is a function. We have already
half-seen it: \code{add3 = add 3} was a function produced by partially applying
\code{add}. But we can also write a factory \emph{directly}, with the explicit
intention of producing a function. Consider \code{scaled\_by}, which takes a scalar
$a$ and returns the function ``scale a tensor by $a$.''

\begin{lstlisting}[style=l4style]
scaled_by :: Scalar -> (Tensor[N] -> Tensor[N])
scaled_by a = \v -> a .* v
\end{lstlisting}

\noindent Several new pieces here, each small:

\begin{description}[style=nextline,leftmargin=1.5em]
  \item[The codomain \code{(Tensor[N] -> Tensor[N])}.] The
  signature says \code{scaled\_by} takes a \code{Scalar} and returns---the
  parenthesized part---a \emph{function} from \code{Tensor[N]} to \code{Tensor[N]}.
  Those parentheses are the \emph{closure-returning convention} we have been
  flagging: see the box just below.
  \item[The lambda \code{\textbackslash v -> a .* v}.] This is an
  \emph{anonymous function}---a function with no name, written inline. The
  backslash \code{\textbackslash} introduces the parameter (here \code{v}), the
  arrow \code{->} separates the parameter from the body, and
  \lstinline[style=l4style]|a .* v| is the body (the operator \code{.*} is
  scalar-times-tensor; we meet it in \cref{ch:tensors}). Read the whole thing as
  ``the function that, given \code{v}, returns \lstinline[style=l4style]|a .* v|.''
  \item[The capture.] The crucial part: the body
  \lstinline[style=l4style]|a .* v| mentions \code{a}, which is \emph{not} the
  lambda's own parameter---it is the argument that was handed to \code{scaled\_by}.
  The returned function has \emph{captured} \code{a}: it carries that value around
  with it, baked in, ready to use whenever it is finally applied. A function that
  has captured values from where it was created is called a \emph{closure}.
\end{description}

\noindent Now we can use it. Building the closure and applying it are two separate
events, possibly far apart in the program:

\begin{lstlisting}[style=l4style]
let double  = scaled_by 2 in   -- build a closure; 2 captured.  double :: Tensor[N] -> Tensor[N]
let halve   = scaled_by 0.5 in -- a DIFFERENT closure; 0.5 captured
let y = double v in            -- apply later: y = 2 .* v
let z = halve  v in            -- z = 0.5 .* v
(y, z)
\end{lstlisting}

\noindent \code{double} and \code{halve} are two distinct function values, born
from the same factory but carrying different captured scalars. This is the
right-hand cartoon of \cref{fig:higherorder}: a value flows \emph{in} to the
factory, a function flows \emph{out} with that value captured inside it.
\Cref{wex:scaledby} traces this through with concrete numbers.

\begin{notationbox}
\textbf{The closure-returning parenthesis convention.} When a signature's result
(its \emph{codomain}, the type after the last top-level arrow) is itself a
function that you intend to \emph{hold and apply later}, we group it in
parentheses:
\begin{lstlisting}[style=l4style]
scaled_by :: Scalar -> (Tensor[N] -> Tensor[N])   -- "returns a function"
\end{lstlisting}
By the right-grouping rule this is the \emph{same type} as the un-parenthesized
\lstinline[style=l4style]|Scalar -> Tensor[N] -> Tensor[N]|---the parentheses
change no types. They are a \emph{reader-intent marker}: they tell you the
codomain is a closure you will hold and apply, not the last operand of a
multi-argument call. We use them whenever a function's product is itself a
function to be applied later---operator constructors, partial-application
factories. \Cref{ch:operators} introduces a richer spelling,
\lstinline[style=l4style]|Op[in -> out]|, for the special case where the returned
function is an \emph{operator} carrying baked-in parameters; the parenthesis habit
is its plainer sibling.
\end{notationbox}

\subsection{Why ``higher-order'' is worth the effort}

It is fair to ask whether this machinery earns its conceptual cost. It does, and
the reason is the one from \cref{ch:bigidea}: half this book is the discovery
that six different-looking simulators are built from one small set of pieces,
combined in slightly different ways. ``Combined'' is doing real work in that
sentence. The \emph{verbs} of an algorithm (apply this operator, take this inner
product) are ordinary functions; but the \emph{coordination}---do this step, then
this, repeat until converged, map over a family of right-hand sides---is expressed
by functions that \emph{take} the verbs as arguments and orchestrate them. A loop
that works for \emph{any} step function is a higher-order function; a solver that
is configured once and applied many times is a closure. Without functions as
values, we would have to rewrite the coordination from scratch for every verb;
with them, we write the coordination \emph{once} and feed it the verbs. That is
the entire economy of \cref{part:framework}.

\section{Lambda, \code{let}, and the absence of hidden order}

Three loose ends remain, all small, all in service of reading the bodies of
functions fluently.

\subsection{Lambda: anonymous functions}

We met the lambda \lstinline[style=l4style]|\v -> a .* v| above. It is just a
function written without giving it a name. The general form is
\lstinline[style=l4style]|\x -> body|: the backslash, a parameter, an arrow, and a
body expression. We reach for a lambda whenever a function is small and used in
exactly one place---most often as the argument to a higher-order function, where
naming it would be ceremony:

\begin{lstlisting}[style=l4style]
apply_twice (\x -> x + 1) 10   -- = 12; the increment function, written inline
\end{lstlisting}

\noindent Anything written as a lambda could be given a name with \code{let} and a
signature instead; the lambda is a convenience for the throwaway case. The two
forms below are interchangeable:

\begin{lstlisting}[style=l4style]
let inc = \x -> x + 1 in apply_twice inc 10   -- named
apply_twice (\x -> x + 1) 10                   -- inline lambda
\end{lstlisting}

\subsection{\code{let}: naming intermediates}

We have used \code{let} throughout; here is its job stated plainly.
\lstinline[style=l4style]|let name = expr in body| evaluates \code{expr},
binds its value to \code{name}, and then evaluates \code{body} with that name
available. It is how we name an intermediate result so we can refer to it more
than once, or simply to give a complicated computation readable, step-by-step
structure. A chain of \code{let}s reads top to bottom like a recipe:

\begin{lstlisting}[style=l4style]
let r  = b - matvec A x in   -- the residual
let p  = precond r in        -- the preconditioned residual
let a  = dot r p in          -- a scalar
...
\end{lstlisting}

\noindent Each name binds an immutable value; none is ever reassigned. (When a
chain of \code{let}s threads \emph{evolving state}---a counter that advances, a
state vector that steps---we will switch to the \code{do}-block notation of
\cref{ch:monad}, which is the same recipe-like reading with state made explicit.
Plain \code{let} suffices until then.)

\subsection{No hidden order: data dependencies are the only order}

Here is a quietly profound consequence of purity, and the bridge to the rest of
Part~III. Because every function is pure---its output depends only on its
inputs---the \emph{value} of an expression depends on \emph{nothing but the values
of its sub-expressions}. The order in which we happen to write the \code{let}s, or
the order a machine happens to evaluate them, cannot change any answer. The
\emph{only} ordering that matters is forced by the data: if computing \code{p}
needs \code{r}, then \code{r} must be computed first---but two intermediates that
do not depend on each other may be computed in either order, or at the same time.

This means an expression is really a \emph{tree} (more generally, a graph) whose
edges are data dependencies, and that tree carries no order beyond what the edges
demand. \Cref{fig:exprtree} draws one. Read the figure and the lesson of
\cref{ch:bigidea} lands with full force: the \emph{program text is the
dependency graph}, with no hidden arrows from sequencing, because there is no
hidden sequencing---only data flowing along visible edges.

\begin{figure}[t]
\centering
\begin{tikzpicture}[node distance=8mm and 12mm]
  % leaves
  \node[data] (b)  at (-3.2,0)   {\code{b}};
  \node[data] (A)  at (-1.4,-0.9){\code{A}};
  \node[data] (x)  at (0.2,-0.9) {\code{x}};
  % matvec A x
  \node[opbox, minimum width=16mm] (mv) at (-0.6,0.1) {\code{matvec}};
  \draw[flow] (A) -- (mv);
  \draw[flow] (x) -- (mv);
  % r = b - matvec
  \node[opbox, minimum width=12mm] (sub) at (-1.7,1.2) {$-$};
  \draw[flow] (b)  -- (sub);
  \draw[flow] (mv) -- (sub);
  \node[font=\scriptsize\ttfamily, simBlue, left=0.5mm of sub] {r};
  % precond r
  \node[opbox, minimum width=16mm] (pc) at (0.4,1.2) {\code{precond}};
  \draw[flow] (sub) -- (pc);
  \node[font=\scriptsize\ttfamily, simBlue, right=0.5mm of pc] {p};
  % dot r p
  \node[opbox, minimum width=12mm] (dot) at (-0.6,2.5) {\code{dot}};
  \draw[flow] (sub) to[out=70,in=200] (dot);
  \draw[flow] (pc)  -- (dot);
  \node[product, minimum width=10mm, above=7mm of dot] (a) {\code{a}};
  \draw[flow] (dot) -- (a);
  % annotation
  \node[align=left,font=\scriptsize\itshape,simGray,text width=46mm]
    at (4.6,1.2) {The edges are the \emph{only} ordering. \code{precond} must wait
    for \code{r}; but reading \code{b} and forming \code{matvec A x} are unordered
    relative to each other. There is no hidden ``step 1, step 2''---only data
    flowing up the tree.};
\end{tikzpicture}
\caption[An expression is a dependency tree with no hidden order]{The expression
\lstinline[style=l4style]|let r = b - matvec A x in let p = precond r in dot r p|,
drawn as a tree. Leaves are inputs; internal nodes are pure operations; edges are
data dependencies. The only ordering the computation imposes is the one the edges
force (\code{p} after \code{r}, the final \code{dot} after both). Anything the
edges leave unordered may happen in any order, or in parallel---because every node
is a pure function of its inputs and nothing else. The program text \emph{is} this
graph. Reading an algorithm's structure off its written form, exactly here, is the
skill \cref{ch:fold} turns into a method.}
\label{fig:exprtree}
\end{figure}

\section{A first look at operator values}

We close with a forward glance, kept deliberately light. Several times we have met
the pattern ``build a function once with configuration baked in, apply it many
times.'' This pattern is so central---it is how every assembled operator, every
preconditioner, every constructed step works---that the calculus gives it a
dedicated spelling. An \emph{operator value} is written

\begin{lstlisting}[style=l4style]
Op[Tensor[N] -> Tensor[N]]
\end{lstlisting}

\noindent and you should read it as ``a function value from \code{Tensor[N]} to
\code{Tensor[N]} that \emph{also carries baked-in parameters}''---the captured
configuration (matrix entries, a factorization, basis tables) that a plain closure
would hold invisibly, here made a first-class, named thing. An operator
constructor then has a signature like

\begin{lstlisting}[style=l4style]
make_operator :: Config -> Op[Tensor[N] -> Tensor[N]]
\end{lstlisting}

\noindent ``give me a configuration, I return an operator value you apply later.''
The \lstinline[style=l4style]|Op[...]| brackets play exactly the role the
closure-returning parentheses played for \code{scaled\_by}: they announce ``the
result is a thing you hold and apply.'' (Because the brackets already group the
in/out arrow, an \lstinline[style=l4style]|Op[...]| codomain wants no extra outer
parentheses.) Why operators deserve their own spelling rather than just being
plain closures---the answer is about \emph{ownership}, the second of our three
questions, namely \emph{who owns the baked-in parameters} and that they are
read-only for the operator's whole life---is the subject of \cref{ch:operators}.
For now, just recognize the shape: \lstinline[style=l4style]|Op[in -> out]| is a
configured function you apply many times.

\begin{workedexample}{Reading a real signature, token by token}{axpy}
The workhorse update of every linear solver in this book is \code{axpy}---``$a$
times $x$ plus $y$,'' the operation we met in \cref{ch:bigidea}. Here is its
full definition in the notation:
\begin{lstlisting}[style=l4style]
axpy :: Scalar -> Tensor[N] -> Tensor[N] -> Tensor[N]
axpy a x y = a .* x .+ y
\end{lstlisting}
We read \emph{every} token, top line first.
\begin{description}[style=nextline,leftmargin=1.4em,nosep]
  \item[\code{axpy}] the name of the function.
  \item[\code{::}] ``has type.'' Everything after it is the type.
  \item[\code{Scalar}] the type of the \emph{first} argument: a single number $a$.
  \item[the first \code{->}] ``maps to''---and, by right-grouping, ``\dots then
  returns a function that takes\dots''
  \item[\code{Tensor[N]} (second)] the type of the second argument: a length-$N$
  vector $x$. (Shapes are types; \cref{ch:tensors}.)
  \item[\code{Tensor[N]} (third)] the third argument: another length-$N$ vector
  $y$. All three of \code{x}, \code{y}, and the result share the same length $N$.
  \item[\code{Tensor[N]} (last)] the \emph{result} type: a length-$N$ vector. This
  is the codomain---what you get once all three arguments are supplied.
\end{description}
So the signature reads: ``\code{axpy} takes a scalar, then a length-$N$ vector,
then another length-$N$ vector, and returns a fresh length-$N$ vector.'' Now the
body:
\begin{description}[style=nextline,leftmargin=1.4em,nosep]
  \item[\code{axpy a x y}] names the three parameters, in order: \code{a}, \code{x},
  \code{y}.
  \item[\code{a .* x}] the scalar \code{a} times the tensor
  \code{x}, elementwise---a fresh tensor (the \code{.*} dot marks
  scalar-times-tensor; \cref{ch:tensors}).
  \item[\code{.+ y}] add the tensor \code{y} elementwise---a
  fresh tensor, the result.
\end{description}
Three readings worth stating out loud, because each is a habit this chapter is
trying to build. (1) \emph{Nothing is overwritten}: \code{x} and \code{y} name
exactly the same tensors after the call as before; the result is a brand-new
value. (2) \emph{It is curried}: \lstinline[style=l4style]|axpy 2| is already a
legal value---the function ``$2x + y$ awaiting $x$ and $y$''---and
\lstinline[style=l4style]|axpy 2 x| is the function ``$2x + y$ awaiting $y$.''
(3) \emph{It is pure}: hand \code{axpy} the same $a$, $x$, $y$ on any machine, any
day, and you get the identical result; it reads no hidden state and writes none.
\end{workedexample}

\begin{workedexample}{Building and applying a closure, with numbers}{scaledby}
We trace \code{scaled\_by} from \cref{sec:higherorder} through to concrete
numbers, to make the higher-order idea tangible. Recall:
\begin{lstlisting}[style=l4style]
scaled_by :: Scalar -> (Tensor[N] -> Tensor[N])
scaled_by a = \v -> a .* v
\end{lstlisting}
Take $N = 3$ and a concrete input tensor $v = [10, 20, 30]$.

\textbf{Step 1 --- build the closure.} Evaluate \lstinline[style=l4style]|scaled_by 2|.
We supply $a = 2$. By the definition, the result is the lambda
\lstinline[style=l4style]|\v -> a .* v| \emph{with $a$ captured as $2$}:
\begin{lstlisting}[style=l4style]
let double = scaled_by 2 in ...
-- double is the closure (\v -> 2 .* v); the 2 is baked in.
-- double :: Tensor[3] -> Tensor[3]
\end{lstlisting}
No tensor has been touched yet---we have only \emph{produced a function}. This is
the value flowing \emph{out} of the factory in \cref{fig:higherorder} (right).

\textbf{Step 2 --- apply the closure.} Now, possibly much later in the program,
we apply \code{double} to our $v$:
\begin{lstlisting}[style=l4style]
let y = double v in ...   -- substitute v = [10,20,30] into (\v -> 2 .* v)
\end{lstlisting}
Applying the closure substitutes $v = [10,20,30]$ into the captured body
$2 .* v$:
\[
  y \;=\; 2 \cdot [10,\,20,\,30] \;=\; [20,\,40,\,60].
\]

\textbf{Step 3 --- a second, independent closure.} Build another from the same
factory with a different captured scalar:
\begin{lstlisting}[style=l4style]
let halve = scaled_by 0.5 in
let z = halve v in ...     -- z = 0.5 .* [10,20,30] = [5,10,15]
\end{lstlisting}
\code{double} and \code{halve} are two distinct closures, born from one factory,
each carrying its own captured scalar; \code{v} was never disturbed by either
(immutability), and neither closure can interfere with the other (purity). The
two events---\emph{building} the function and \emph{applying} it---are cleanly
separated, which is exactly the ``configure once, apply many times'' shape that
\cref{ch:operators} elevates into constructed operators.
\end{workedexample}

\summarysection
The calculus exists to make three structural questions explicit---what operations
happen, who owns the state, how evolution is coordinated---and this chapter does
the first by teaching the notation for \emph{values and functions}. A
\emph{value} (number, tensor, record, function) is immutable: it never changes
once created; we never overwrite, we compute fresh values and name them with
\code{let}; \code{clone()} is meaningless because nothing can be disturbed. A
\emph{function} is written as a \emph{type signature}---\code{f :: A -> B}, where
\code{::} reads ``has type'' and \code{->} ``maps to''---and a \emph{definition}
applied by juxtaposition; our functions are \emph{pure} (output depends only on
declared inputs, no side effects). A function is itself a \emph{value}, which
gives us two powerful moves. \emph{Currying} reads \code{A -> B -> C} as one
argument at a time, so \emph{partial application} can freeze some arguments now
and supply the rest later (``build once, apply many''). \emph{Higher-order
functions} take a function as an argument (\code{apply\_twice}, the shape of every
combinator) or return one (\code{scaled\_by}, a \emph{closure} that has captured a
value). The closure-returning parenthesis convention
\lstinline[style=l4style]|A -> (B -> C)| marks ``returns a function you hold and
apply later.'' \emph{Lambdas} \lstinline[style=l4style]|\x -> ...| are anonymous
functions; \code{let} names intermediates; and because every function is pure, an
expression is a \emph{dependency tree with no hidden order}---the program text
\emph{is} the dependency graph. The operator-value spelling
\lstinline[style=l4style]|Op[in -> out]|, a configured function carrying baked-in
parameters, previews \cref{ch:operators}.

\furtherreading
The shift from mutable variables to immutable values and first-class functions is
the foundation of typed functional programming; the most readable entry points are
the standard reference for the language whose notation we borrow~\citep{peytonjones2003}
and Wadler's exposition of how a pure language nonetheless models stateful
computation~\citep{wadler1995monads}, which we will need in earnest in
\cref{ch:monad}. No prior background is assumed here or anywhere in the book; each
of the four pillars of the notation gets its own careful chapter---records and
immutability (\cref{ch:records}), tensors and shapes (\cref{ch:tensors}), the fold
and the combinators (\cref{ch:fold}), and the state monad (\cref{ch:monad}).
Readers who want the term ``currying'' in its original setting will find it in any
text on the lambda calculus.

\exercisessection
\begin{enumerate}[nosep]
  \item \textbf{Read this signature.} In one English sentence each, say what these
  functions take and return:
  \code{nrm2 :: Tensor[N] -> Scalar};
  \code{dot :: Tensor[N] -> Tensor[N] -> Scalar};
  \code{matvec :: Tensor[m, n] -> Tensor[n] -> Tensor[m]}.
  For \code{matvec}, why is the result length $m$ rather than $n$?
  \item \textbf{What does it curry to?} Recall \code{axpy} has type
  \code{Scalar -> Tensor[N] -> Tensor[N] -> Tensor[N]}.
  Write the type of \lstinline[style=l4style]|axpy 2|, the type of
  \lstinline[style=l4style]|axpy 2 x| (with \code{x : Tensor[N]}), and the type of
  \lstinline[style=l4style]|axpy 2 x y|. Which of the three is \emph{not} a
  function?
  \item \textbf{Spot the side effect.} One of the following describes a
  \emph{pure} function and one does not: (a) ``returns its first argument plus its
  second''; (b) ``adds its argument to a running total kept in a global counter
  and returns the new total.'' Which is impure, and what specifically about it
  violates purity? Why can the impure one give different answers on the same
  input?
  \item \textbf{Higher-order or not.} For each signature, say whether it is a
  higher-order function and \emph{why}: (a)
  \lstinline[style=l4style]|(Float -> Float) -> Float -> Float|; (b)
  \lstinline[style=l4style]|Float -> Float -> Float|; (c)
  \lstinline[style=l4style]|Scalar -> (Tensor[N] -> Tensor[N])|. For the
  higher-order ones, say whether the function-ness is in an \emph{argument} or in
  the \emph{result}.
  \item \textbf{Build the closure.} Using \code{scaled\_by} from
  \cref{wex:scaledby}, write down (i) the value of
  \lstinline[style=l4style]|scaled_by 3| in words, (ii) its type, and (iii) the
  result of applying it to $[1, 0, 2]$. Then explain why \code{scaled\_by 3} and
  \code{scaled\_by 5} cannot interfere with each other even though they came from
  the same factory.
  \item \textbf{Parentheses matter.} Explain, in terms of the right-grouping rule,
  why \lstinline[style=l4style]|(Float -> Float) -> Float| and
  \lstinline[style=l4style]|Float -> Float -> Float| are \emph{different} types.
  Describe a value of each. (Hint: count how many \emph{numbers} each ultimately
  needs.)
  \item \textbf{No hidden order.} In the dependency tree of \cref{fig:exprtree},
  list every pair of operations whose relative order is \emph{not} forced by the
  data, and one pair whose order \emph{is}. What property of the functions
  guarantees that the unforced pairs may be evaluated in either order without
  changing the answer?
  \item \textbf{The mutable-box trap.} A reader writes: ``after
  \lstinline[style=l4style]|let r = axpy a x y|, the tensor \code{y} now holds
  $a x + y$.'' Identify the exact mental-model error in this sentence and rewrite
  it correctly. What \emph{does} \code{y} refer to after that line?
\end{enumerate}

\chapter{Records and Immutability}
\label{ch:records}

\chapterorient{A simulation's state is rarely a single number. It is a
\emph{bundle}: the current field, how many iterations have run, whether the
solver has converged, the residual it is chasing. This chapter introduces the
data structure that carries such a bundle---the \emph{record}, a collection of
named, typed fields---and the one rule that governs every record in this book:
they never change. ``Updating'' a record does not overwrite it; it builds a
fresh one that differs in a field or two and leaves the original untouched. That
single discipline---update means fresh copy---is what turns a tangle of mutable
state into the clean, traceable dataflow the rest of Part~III is built on.}

\section{A bundle of named fields}

In \cref{ch:language} we built a vocabulary of \emph{values} and the
\emph{functions} that map between them. The values so far have been simple: a
tensor, an integer, a boolean. But the quantities a simulation actually carries
around are seldom so lonely. Consider what a linear solver knows about itself in
the middle of a run. It holds a current guess at the answer---a tensor. It has
taken some number of steps---an integer. It has decided, or not, that it is
close enough---a boolean. These three facts travel together; they are the state
of \emph{one} solver at \emph{one} moment. We want a single value that holds all
three.

That value is a \emph{record}: a bundle of named, typed fields. We write one
with braces, listing each field as \lstinline[style=l4style]|name: value| pairs
separated by commas. Here is the solver state just described:

\begin{lstlisting}[style=l4style]
{ x: theField, it: 7, converged: False }
\end{lstlisting}

This is a single value. It has three fields: \code{x}, whose value is some
tensor we have named \code{theField}; \code{it}, an integer that happens to be
$7$; and \code{converged}, a boolean that happens to be \code{False}. The names
\code{x}, \code{it}, \code{converged} are chosen by us and are part of the
record's \emph{type}, which we may write the same way but with types in the
value slots:

\begin{lstlisting}[style=l4style]
{ x: Tensor[N], it: Int, converged: Bool }
\end{lstlisting}

Read this as a sentence: ``a record with a field \code{x} holding a tensor of
length \code{N}, a field \code{it} holding an integer, and a field
\code{converged} holding a boolean.'' Every record value of this type has
exactly these three fields, with exactly these types. \Cref{fig:card} draws the
value as what it morally is: a small labeled card.

\begin{figure}[t]
\centering
\begin{tikzpicture}[font=\footnotesize]
  % the card
  \node[draw=simBlue, fill=simBgBlue, rounded corners=3pt, thick,
        minimum width=46mm, minimum height=30mm, inner sep=6pt] (card) at (0,0) {};
  \node[anchor=north, font=\scriptsize\bfseries, text=simBlue]
    at ($(card.north)+(0,-1.5mm)$) {SolverState};
  % rows
  \node[anchor=west] at ($(card.north west)+(3mm,-7mm)$)
    {\code{x}\,:\ \(\mathsf{Tensor}[N]\) \;\(=\ \mathbf{theField}\)};
  \node[anchor=west] at ($(card.north west)+(3mm,-14mm)$)
    {\code{it}\,:\ \(\mathsf{Int}\) \;\(=\ 7\)};
  \node[anchor=west] at ($(card.north west)+(3mm,-21mm)$)
    {\code{converged}\,:\ \(\mathsf{Bool}\) \;\(=\ \mathrm{False}\)};
  % separating rules
  \draw[simGray!40] ($(card.north west)+(2mm,-10mm)$) -- ($(card.north east)+(-2mm,-10mm)$);
  \draw[simGray!40] ($(card.north west)+(2mm,-17mm)$) -- ($(card.north east)+(-2mm,-17mm)$);
\end{tikzpicture}
\caption[A record as a labeled card]{A record is a card of named, typed fields.
Each row carries a field name, the field's type, and the value currently stored
there. The card as a whole is a single value of type
\code{\{ x: Tensor[N], it: Int, converged: Bool \}}; the rows
are reached by name.}
\label{fig:card}
\end{figure}

\subsection{Reaching a field: the dot}

To pull one field out of a record we write the record, a dot, and the field
name. If \code{s} is the record above, then

\begin{lstlisting}[style=l4style]
s.it          -- the integer 7
s.converged   -- the boolean False
s.x           -- the tensor theField
\end{lstlisting}

Field access \lstinline[style=l4style]|s.it| is a value-producing expression
like any other: it names the integer in the \code{it} slot. The reduction rule
is exactly what intuition demands---projecting a field from a record literal
returns the value in that slot:
\[
  \{l_1\!:\!v_1,\dots,l_n\!:\!v_n\}.\,l_k \;\to\; v_k .
\]
There is nothing more to it. A record is a finite map from names to values, and
the dot looks one up.

Field access composes with everything else in the language. Because
\lstinline[style=l4style]|s.it| is an ordinary expression, it can appear wherever
a value can: as a function argument \lstinline[style=l4style]|norm(s.x)|, in
arithmetic \lstinline[style=l4style]|s.it + 1|, in a comparison
\lstinline[style=l4style]|s.it < s.max_it|, or as the field of \emph{another}
record \lstinline[style=l4style]|{ count: s.it }|. A record is not a special kind
of thing you must unpack before use; its fields are values the moment you name
them, and they flow through the pure-function vocabulary of \cref{ch:language}
unchanged.

\begin{notationbox}
We write record \emph{types} and record \emph{values} with the same brace
syntax, distinguished by what sits after each colon. A \emph{type} has a type in
the value slot---\lstinline[style=l4style]|{ it: Int }|---while a \emph{value}
has an actual value---\lstinline[style=l4style]|{ it: 7 }|. Field access is the
dot: \lstinline[style=l4style]|s.it|. (When a field name collides with a TeX
underscore, as in \code{final\_res}, the underscore is escaped in running prose
but typed literally in code.)
\end{notationbox}

\subsection{Records versus tuples}

The calculus also has \emph{tuples}---ordered bundles written with parentheses,
\lstinline[style=l4style]|(theField, 7, False)|, whose components are reached by
\emph{position}: $\pi_1$ for the first, $\pi_2$ for the second, and so on. A
tuple and a record can carry the same data. The difference is how you address a
component, and the difference matters more than it first appears.

A tuple's components are anonymous. To read the iteration count from
\lstinline[style=l4style]|(theField, 7, False)| you must \emph{know} that the
count is the second slot---and if someone later inserts a field, every
$\pi_i$ downstream silently shifts to the wrong thing. A record's components
are named. \lstinline[style=l4style]|s.it| means the iteration count no matter
where in the record it sits, and adding a field disturbs nothing. For a bundle
that travels a long way through a program and accumulates fields over a project's
lifetime---which is exactly what a simulation's state does---names are
self-documenting and robust where positions are cryptic and brittle. We will use
tuples for short-lived, obviously-ordered pairs (a quotient and remainder, a
coordinate) and records for everything that earns a name.

\begin{keyidea}
A record is to a function what a \emph{struct} is to a procedure: a bundle of
named fields that a routine reads and ``updates.'' The whole of this chapter is
the one place that analogy breaks. A struct is a box you reach into and
overwrite; a record is a value you can only \emph{read}. To ``update'' a record
you build a new one. That single change---from a box you mutate to a value you
copy---is what the rest of Part~III rests on.
\end{keyidea}

\section{Immutable update: building a fresh copy}

Here is the central fact, and it is worth stating before any syntax. \emph{A
record, once built, never changes.} There is no operation that reaches into an
existing record and overwrites a field. None. The only thing you can do with a
record is read its fields and use them to build \emph{other} records. When we
speak---as we constantly will---of ``updating'' a record, we always mean: build
a new record, like the old one but differing in some field. The old record is
still there, unchanged, for anyone still holding it.

\subsection{The spread}

The syntax for ``a fresh record like \code{s} but with one field changed'' is the
\emph{spread}. We open a record literal, write \lstinline[style=l4style]|...s| to
mean ``copy in all of \code{s}'s fields,'' and then list the fields we want to
override:

\begin{lstlisting}[style=l4style]
{ ...s, it: s.it + 1 }
\end{lstlisting}

Read this aloud: ``a fresh record, with all of \code{s}'s fields, except that
\code{it} is \lstinline[style=l4style]|s.it + 1|.'' If \code{s} was
\lstinline[style=l4style]|{ x: theField, it: 7, converged: False }|, then the
expression above is a \emph{new} value
\lstinline[style=l4style]|{ x: theField, it: 8, converged: False }|. The fields
\code{x} and \code{converged} are copied straight across; \code{it} is replaced.
The reduction rule is precisely this---a spread over a record literal collapses
to a record literal with the named slot overwritten and the rest carried through:
\[
  \{\,\dots\{l_1\!:\!v_1,\dots,l_n\!:\!v_n\},\ l_k\!:\!w\,\}
  \;\to\;
  \{l_1\!:\!v_1,\dots,l_k\!:\!w,\dots,l_n\!:\!v_n\}.
\]

\Cref{fig:spread} draws what just happened. The original card \code{s} is on the
left, untouched. The spread produces a \emph{second} card on the right, identical
but for the one changed field. Both cards exist. Nothing was erased.

\begin{figure}[t]
\centering
\begin{tikzpicture}[font=\footnotesize]
  % --- left card: original s ---
  \begin{scope}
  \node[draw=simGray, fill=white, rounded corners=3pt, thick,
        minimum width=34mm, minimum height=26mm, inner sep=5pt] (s) at (0,0) {};
  \node[anchor=north, font=\scriptsize\bfseries, text=simGray] at ($(s.north)+(0,-1mm)$) {\code{s} (unchanged)};
  \node[anchor=west] at ($(s.north west)+(2.5mm,-7mm)$)  {\code{x}\,:\ \textbf{theField}};
  \node[anchor=west] at ($(s.north west)+(2.5mm,-13mm)$) {\code{it}\,:\ \(7\)};
  \node[anchor=west] at ($(s.north west)+(2.5mm,-19mm)$) {\code{converged}\,:\ F};
  \end{scope}
  % --- right card: fresh s' ---
  \begin{scope}[xshift=58mm]
  \node[draw=simBlue, fill=simBgBlue, rounded corners=3pt, thick,
        minimum width=34mm, minimum height=26mm, inner sep=5pt] (s2) at (0,0) {};
  \node[anchor=north, font=\scriptsize\bfseries, text=simBlue] at ($(s2.north)+(0,-1mm)$) {fresh copy};
  \node[anchor=west] at ($(s2.north west)+(2.5mm,-7mm)$)  {\code{x}\,:\ \textbf{theField}};
  \node[anchor=west, text=simRust, font=\footnotesize\bfseries]
        at ($(s2.north west)+(2.5mm,-13mm)$) {\code{it}\,:\ \(8\)};
  \node[anchor=west] at ($(s2.north west)+(2.5mm,-19mm)$) {\code{converged}\,:\ F};
  \end{scope}
  % the spread arrow
  \draw[flow] (s.east) -- node[above, font=\scriptsize\ttfamily]{\{...s, it: s.it+1\}}
                          node[below, font=\scriptsize, text=simGray]{copies \code{x}, \code{converged}} (s2.west);
\end{tikzpicture}
\caption[Spread produces a fresh card]{The spread
\code{\{ ...s, it: s.it + 1 \}} reads \code{s} and builds a
\emph{new} card. The unchanged fields \code{x} and \code{converged} are copied;
\code{it} is set to its successor (shown in rust). Crucially, the original card
\code{s} on the left is untouched---it is still a perfectly good value that
anyone holding it still sees as \code{it}~$=7$.}
\label{fig:spread}
\end{figure}

\subsection{Contrast: mutating a struct in place}

To feel the difference, picture the procedural alternative. In a language with
mutable structs you would write something like \code{s.it = s.it + 1}, and that
statement \emph{reaches into the box \code{s} and overwrites the \code{it} slot.}
After it runs, there is no ``old \code{s}'' anymore---the $7$ is gone, replaced
by $8$, in place. Anyone else holding a pointer to that same box now sees $8$
too, whether they expected the change or not. The assignment is a side effect: it
modifies the world, and its effect ripples to every name aliased to that box.

The spread does none of this. \lstinline[style=l4style]|{ ...s, it: s.it + 1 }|
is an \emph{expression}, not a command. It computes a value and changes nothing.
The $7$ in \code{s} is still $7$. The new value $8$ lives in a different
record, and the only way anyone sees it is by being handed that new record. This
is the difference between ``mutate the box'' and ``compute a new value,'' and it
is the line dividing the procedural world from the one we work in.

\begin{pitfall}
The single most common misreading of the spread is to think
\lstinline[style=l4style]|{ ...s, it: s.it + 1 }| \emph{changes \code{s}}. It
does not. The spread \emph{reads} \code{s} and \emph{produces a new record};
\code{s} is exactly as it was. If you find yourself writing ``and then \code{s}'s
\code{it} becomes $8$,'' stop---\code{s}'s \code{it} is still $7$ forever. What
became $8$ is the \code{it} of a different record, the one the spread built. The
moment you internalize that the spread never touches its source, the value-history
traces below read themselves.
\end{pitfall}

\subsection{Spreading several fields at once}

A spread may override more than one field; just list them all:

\begin{lstlisting}[style=l4style]
{ ...s, it: s.it + 1, converged: True }
\end{lstlisting}

This builds a fresh record advancing the counter \emph{and} flipping the
convergence flag in one stroke, copying \code{x} across unchanged. The order of
the overrides does not matter as long as no name is given twice; each names a
distinct slot of the result. And a field's new value may be computed from the old
record's fields---\lstinline[style=l4style]|s.it + 1| reads \code{s.it} to
compute the replacement---which is the normal way a step advances state: read the
current bundle, compute the next, package it as a fresh bundle.

\subsection{Reading a spread by its reduction rule}

Because the spread is just an expression with a reduction rule, we can evaluate
one by hand the way \cref{ch:language} evaluated function applications: rewrite
step by step until no rule applies. This is worth doing once explicitly---it
turns ``update means fresh copy'' from a slogan into a mechanical fact.

\begin{example}[A spread, reduced step by step]
\label{ex:reduce}
Take \lstinline[style=l4style]|s = { x: f, it: 7, converged: False }| and evaluate
the increment \lstinline[style=l4style]|{ ...s, it: s.it + 1 }|. First substitute
the value of \code{s} (the \textsf{let}-rule of \cref{ch:language}):
\begin{lstlisting}[style=l4style]
{ ...{ x: f, it: 7, converged: False }, it: s.it + 1 }
\end{lstlisting}
The field read \lstinline[style=l4style]|s.it| projects (the \textsf{proj} rule),
giving $7$, so the override value \lstinline[style=l4style]|s.it + 1| reduces to
$8$:
\begin{lstlisting}[style=l4style]
{ ...{ x: f, it: 7, converged: False }, it: 8 }
\end{lstlisting}
Now the spread rule fires: copy the inner record's fields, overwriting the named
one:
\begin{lstlisting}[style=l4style]
{ x: f, it: 8, converged: False }
\end{lstlisting}
No rule applies; this is the value. Notice what the trace makes undeniable: the
inner record \lstinline[style=l4style]|{ x: f, it: 7, converged: False }|---the
value of \code{s}---was \emph{read} (its fields copied, its \code{it} projected)
but at no point \emph{altered}. The result is a brand-new literal sitting beside
it. There is no rule in the calculus that overwrites a field of an existing
record; the spread rule \emph{constructs}, it does not assign.
\end{example}

This is why we can say, flatly, that records are immutable: not as a convention we
promise to honor, but because the only rules that act on records are projection,
destructuring, and spread, and \emph{none of them mutates}. Projection reads;
destructuring reads; spread builds anew. Immutability is a property of the
rules, not a discipline imposed on top of them.

\begin{workedexample}{Three updates, and the originals all survive}{history}
We start from an initial solver state and apply three spread-updates in sequence,
naming each result. The point is to watch every intermediate record \emph{survive}
its own update---a value-history trace.

\textbf{Setup.} Let
\begin{lstlisting}[style=l4style]
let s0 = { x: f0, it: 0, converged: False }
\end{lstlisting}
where \code{f0} is the initial field guess. Now three steps, each a spread:
\begin{lstlisting}[style=l4style]
let s1 = { ...s0, x: f1, it: 1 }
let s2 = { ...s1, x: f2, it: 2 }
let s3 = { ...s2, x: f3, it: 3, converged: True }
\end{lstlisting}
Here \code{f1}, \code{f2}, \code{f3} are the successively-improved fields (their
computation is the solver's business; \cref{ch:fold} supplies it). Each line
reads the previous record and builds a fresh one.

\tcblower
\textbf{The trace.} Because no spread ever mutates its source, all four records
coexist:
\[
\begin{array}{l@{\;=\;}l}
\code{s0} & \{\ \code{x}\!: f_0,\ \code{it}\!: 0,\ \code{converged}\!: \mathrm{F}\ \} \\[2pt]
\code{s1} & \{\ \code{x}\!: f_1,\ \code{it}\!: 1,\ \code{converged}\!: \mathrm{F}\ \} \\[2pt]
\code{s2} & \{\ \code{x}\!: f_2,\ \code{it}\!: 2,\ \code{converged}\!: \mathrm{F}\ \} \\[2pt]
\code{s3} & \{\ \code{x}\!: f_3,\ \code{it}\!: 3,\ \code{converged}\!: \mathrm{T}\ \}
\end{array}
\]
After all three updates, \lstinline[style=l4style]|s0.it| is still $0$ and
\lstinline[style=l4style]|s0.converged| is still \code{False}. The update that
produced \code{s1} did not reach back and alter \code{s0}; it produced a
\emph{new} record. Each later step likewise leaves its predecessor intact. We
have, for free, a complete history of the solve---every intermediate state, still
readable---simply because nothing was ever overwritten. \Cref{fig:chain} draws the
chain.
\end{workedexample}

\begin{figure}[t]
\centering
\begin{tikzpicture}[font=\footnotesize, node distance=12mm]
  \tikzset{rec/.style={draw, rounded corners=2pt, thick, align=center,
     minimum width=22mm, minimum height=13mm, inner sep=2pt, font=\scriptsize}}
  \node[rec, draw=simGray,  fill=white]     (a) {\code{s0}\\[-1pt]\(it\!=\!0\)\\\(\mathrm{F}\)};
  \node[rec, draw=simGray,  fill=white, right=of a] (b) {\code{s1}\\[-1pt]\(it\!=\!1\)\\\(\mathrm{F}\)};
  \node[rec, draw=simGray,  fill=white, right=of b] (c) {\code{s2}\\[-1pt]\(it\!=\!2\)\\\(\mathrm{F}\)};
  \node[rec, draw=simBlue, fill=simBgBlue, right=of c] (d) {\code{s3}\\[-1pt]\(it\!=\!3\)\\\(\mathrm{T}\)};
  \draw[flow] (a) -- node[above,font=\scriptsize]{spread} (b);
  \draw[flow] (b) -- node[above,font=\scriptsize]{spread} (c);
  \draw[flow] (c) -- node[above,font=\scriptsize]{spread} (d);
  \node[font=\scriptsize, text=simGray, align=center] at ($(a)+(0,-12mm)$)
    {\(it=0\)\\\emph{still readable}};
\end{tikzpicture}
\caption[A chain of fresh records]{The value history of
\cref{wex:history}. Each spread reads one record and emits a fresh successor;
the arrow is exactly the \emph{value-threading} of \cref{ch:targets}'s transient
solve. None of \code{s0}, \code{s1}, \code{s2} is destroyed when its successor is
built---the whole chain persists, so \code{s0.it} is $0$ even
after \code{s3} exists.}
\label{fig:chain}
\end{figure}

\section{Destructuring: pulling fields out by name}

Reaching fields one dot at a time is fine for a single field but clumsy when a
function needs several. \emph{Destructuring} lets a \code{let} bind several fields
to names in one move. Instead of binding the whole record to a name and dotting
into it, we write a brace pattern on the left of the \lstinline[style=l4style]|=|:

\begin{lstlisting}[style=l4style]
let { x, it } = s in
  -- here x is bound to s.x and it to s.it
  ...
\end{lstlisting}

This says: ``take \code{s}, and bind the local names \code{x} and \code{it} to its
\code{x} and \code{it} fields.'' Inside the body, \code{x} \emph{is}
\lstinline[style=l4style]|s.x| and \code{it} \emph{is} \lstinline[style=l4style]|s.it|,
without any further dotting. The reduction is the obvious one---destructuring a
record literal substitutes each bound name by the matching field:
\[
  \textsf{let}\ \{l_1,\dots,l_n\} = \{l_1\!:\!v_1,\dots,l_n\!:\!v_n\}\ \textsf{in}\ e
  \;\to\; e[v_1/l_1,\dots,v_n/l_n].
\]
A pattern need not mention every field---\lstinline[style=l4style]|let { it } = s|
binds only \code{it} and ignores the rest. \Cref{fig:destructure} pictures the
fields being pulled out by name.

\begin{figure}[t]
\centering
\begin{tikzpicture}[font=\footnotesize]
  \node[draw=simBlue, fill=simBgBlue, rounded corners=3pt, thick,
        minimum width=40mm, minimum height=26mm, inner sep=5pt] (s) at (0,0) {};
  \node[anchor=north, font=\scriptsize\bfseries, text=simBlue] at ($(s.north)+(0,-1mm)$) {\code{s}};
  \node[anchor=west] (fx)  at ($(s.north west)+(2.5mm,-7mm)$)  {\code{x}\,:\ theField};
  \node[anchor=west] (fit) at ($(s.north west)+(2.5mm,-13mm)$) {\code{it}\,:\ \(7\)};
  \node[anchor=west]       at ($(s.north west)+(2.5mm,-19mm)$) {\code{converged}\,:\ F};
  % pulled-out names
  \node[draw=simGreen, fill=simBgGreen, rounded corners=2pt, font=\ttfamily\footnotesize] (nx)  at (5.6,0.7) {x};
  \node[draw=simGreen, fill=simBgGreen, rounded corners=2pt, font=\ttfamily\footnotesize] (nit) at (5.6,-0.1) {it};
  \draw[flow, simGreen] (fx.east)  to[out=0,in=180] (nx.west);
  \draw[flow, simGreen] (fit.east) to[out=0,in=180] (nit.west);
  \node[font=\scriptsize, text=simGray, align=center] at (5.6,-1.0)
    {\ttfamily let \{ x, it \} = s};
\end{tikzpicture}
\caption[Destructuring pulls fields out by name]{The pattern
\code{let \{ x, it \} = s} binds the local names \code{x} and
\code{it} to the matching fields of \code{s}, by name. The \code{converged} field
is simply not mentioned, so no name is bound to it. Destructuring reads \code{s};
it does not consume or alter it.}
\label{fig:destructure}
\end{figure}

\subsection{Destructuring with a rest}

Sometimes a function wants a couple of fields by name \emph{and} the remaining
fields kept as a record---to forward them, or to spread them into a result. The
\emph{rest} pattern, written \lstinline[style=l4style]|...r|, captures whatever is
left:

\begin{lstlisting}[style=l4style]
let { it, ...rest } = s in
  { ...rest, it: it + 1 }
\end{lstlisting}

Here \code{it} is bound to \lstinline[style=l4style]|s.it|, and \code{rest} is
bound to a record holding all of \code{s}'s \emph{other} fields (\code{x} and
\code{converged}). The body then spreads \code{rest} back in and overrides
\code{it}---which is just our increment-step written so that the unchanged fields
flow through a name rather than being copied by \lstinline[style=l4style]|...s|.
The two idioms---\lstinline[style=l4style]|{ ...s, it: ... }| and destructure-with-rest
then re-spread---compute the same fresh record; pick whichever reads more clearly
at the use site.

\begin{workedexample}{Destructure, derive, rebuild---nothing mutated}{derive}
We destructure a state record, compute a derived quantity from its fields, and
build a \emph{new} record carrying that quantity---then check that the original
survived.

\textbf{Setup.} Suppose a step record carries a residual tensor and a target
tolerance:
\begin{lstlisting}[style=l4style]
let s = { r: residualVec, it: 4, tol: 1.0e-6 }
\end{lstlisting}
We want to decide convergence---is the residual's norm below \code{tol}?---and
record the verdict in a fresh state with the counter advanced.

\textbf{The step.} Destructure the fields we need, derive, rebuild:
\begin{lstlisting}[style=l4style]
let { r, it, tol } = s in
let nrm  = norm(r) in           -- derived scalar (a pure function of r)
let done = nrm < tol in         -- derived boolean
  { ...s, it: it + 1, converged: done }
\end{lstlisting}

\tcblower
\textbf{What happened.} The destructure read \code{r}, \code{it}, \code{tol} from
\code{s}. The two \code{let}s computed derived values---a norm and a comparison---
that are \emph{not stored in \code{s} at all}; they are ephemeral intermediates,
born here and gone after the result is built. The final spread built a fresh
record: all of \code{s}'s fields, with \code{it} advanced and a new
\code{converged} field set to \code{done}.

\textbf{The original survives.} After this expression evaluates,
\lstinline[style=l4style]|s| is byte-for-byte what it was:
\lstinline[style=l4style]|s.it| is still $4$, \lstinline[style=l4style]|s.r| is
still the same residual tensor, and \code{s} has no \code{converged} field at all.
Reading \code{s}, deriving from it, and building a new record from the derived
values \emph{never touched \code{s}}. This is the discipline in miniature: data
flows \emph{out} of records by destructuring and \emph{into} new records by
spreading, and the records in between are immutable all the way down.
\end{workedexample}

\section{Why immutability is the right discipline for a simulation}

The reader entitled to ask, at this point, why we tie our hands. Mutable structs
are the default in every systems language; copying a whole record to change one
field sounds wasteful. Three payoffs justify the discipline, and each one is a
load-bearing beam of the rest of the book.

\subsection{The dataflow is the dependency graph}

A simulation \emph{evolves} state: each step produces the next state from the
current one. When ``produces the next from the current'' literally means ``reads
record \code{s} and builds record \code{s'},'' the program's text \emph{is} a
graph of who-feeds-whom. The arrow from \code{s} to \code{s'} in
\cref{fig:chain} is not a metaphor; it is the dependency edge \code{s'} depends
on \code{s}, and it is visible right there in the spread that names \code{s}. In
\cref{ch:language} we saw that pure functions make data flow explicit; records
make the \emph{state} that flows be a first-class, inspectable value. Together
they make the whole computation a graph you can read off the page---which is
exactly what the calculus needs to reason about equivalence between layers, and
what a backend needs to schedule the work.

If instead each step \emph{mutated} a shared state box, the arrows would
vanish into invisible side effects. Two steps writing the same box would be
ordered only by accident of execution, and ``what does step~$k$ depend on'' would
be unanswerable from the text. Immutability buys us a readable graph; mutation
spends it.

\subsection{Only one field actually changes}

Look again at the increment step \lstinline[style=l4style]|{ ...s, it: s.it + 1 }|.
Of the three fields, \emph{one} is overridden; the other two are copied verbatim.
This is not an accident of the example. Across a great many solver steps, the
record that threads through carries several fields but the step touches
\emph{one}---the iteration counter ticks while the heavy field is recomputed
elsewhere, or the field advances while the counter and flags ride along. The
spread makes this visible: the override list is short, and it names exactly what
changed. We will return to this in \cref{ch:strata}, where the observation that a
solve's externally-visible state mutates at essentially a single point---the
counter---becomes the key to separating the strata of a record cleanly. For now,
notice only that the spread \emph{advertises} its single mutation point; a
mutable struct hides it among a dozen possible assignments.

\subsection{Sharing is safe}

Here is the payoff that the mutable world simply cannot offer. Because a record
never changes, two names may refer to the \emph{same} record with no hazard
whatsoever. There is no such thing as ``\code{a} aliases \code{b}, so changing
\code{a} surprises \code{b}''---because nothing can change \code{a}. Aliasing,
the perennial bug-farm of mutable state, is \emph{impossible to misuse}: the worst
that aliasing can do is let two names read the same unchanging value, which is
exactly what you want.

\begin{figure}[t]
\centering
\begin{tikzpicture}[font=\footnotesize]
  \node[draw=simBlue, fill=simBgBlue, rounded corners=3pt, thick, align=left,
        minimum width=36mm, minimum height=20mm, inner sep=4pt] (rec) at (0,0)
        {\scriptsize\code{x}: theField\\[1pt]\scriptsize\code{it}: \(7\)\\[1pt]\scriptsize\code{converged}: F};
  \node[font=\scriptsize\bfseries, text=simBlue, anchor=south] at (rec.north) {one immutable record};
  \node[draw=simGray, fill=white, rounded corners=2pt, font=\ttfamily\footnotesize] (na) at (-3.6,1.0) {a};
  \node[draw=simGray, fill=white, rounded corners=2pt, font=\ttfamily\footnotesize] (nb) at (-3.6,-1.0) {b};
  \draw[flow] (na.east) to[out=0,in=160] (rec.west);
  \draw[flow] (nb.east) to[out=0,in=200] (rec.west);
  \node[font=\scriptsize, text=simGreen, align=left, anchor=west] at (2.4,0)
    {both names read\\the same value;\\neither can\\change it};
\end{tikzpicture}
\caption[Safe sharing]{Two names \code{a} and \code{b} bound to the same
immutable record. In a mutable world this would be a hazard---a write through
\code{a} would silently change what \code{b} sees. Here it is free and safe:
because the record cannot change, sharing it is indistinguishable from each name
having its own copy, but costs nothing.}
\label{fig:sharing}
\end{figure}

\Cref{fig:sharing} draws the situation. This safety is not a minor convenience.
It is what lets a spread like \lstinline[style=l4style]|{ ...s, it: s.it + 1 }|
be cheap in practice: the new record need not deep-copy \code{s}'s tensor field
\code{x}: it can \emph{share} the very same immutable tensor, because that tensor
can never change underneath it. The conceptual model is ``fresh copy''; the
implementation, licensed by immutability, shares everything that did not change.
You get the clean semantics of copying with the runtime cost of pointer-sharing,
and you get it precisely \emph{because} nothing mutates.

\begin{keyidea}
Immutability turns aliasing from a hazard into a free optimization. Conceptually,
every spread builds a fresh record. In practice, the fresh record shares the
unchanged fields with its source---safely, because no one can ever mutate the
shared parts. ``Update means fresh copy'' is the \emph{semantics}; ``share what
didn't change'' is the \emph{implementation}, and immutability is exactly what
makes the second a sound realization of the first.
\end{keyidea}

\section{Records of records, optional fields, and demand}

Two refinements round out the record vocabulary; both will be used heavily later
and are introduced here only far enough to read.

\subsection{Records nest}

A field may itself hold a record. The simulator's full state is a record whose
fields are themselves bundles---a configuration sub-record, a solver sub-record,
an output sub-record. Access chains through dots,
\lstinline[style=l4style]|state.solver.it|, and a spread updates a nested field
by spreading at each level:

\begin{lstlisting}[style=l4style]
{ ...state, solver: { ...state.solver, it: state.solver.it + 1 } }
\end{lstlisting}

This builds a fresh outer record whose \code{solver} field is a fresh inner
record with \code{it} advanced---and every other field, inner and outer, copied
across. The nesting is just the rules we already have, applied at two levels.

There is a subtlety worth naming, because it is exactly where the procedural
intuition wants to lead you astray. To bump a nested counter you must rebuild
\emph{both} levels: the inner record (with \code{it} advanced) and the outer
record (with its \code{solver} field pointing at the new inner record). You
cannot ``reach in'' and bump \code{state.solver.it}---there is no such operation,
just as there is none at the top level. But this costs almost nothing in
practice: by the safe-sharing argument of the previous section, the
\emph{untouched} branches of \code{state}---its \code{cfg} field, the
\code{solver}'s other fields---are shared, not copied. Only the spine from the
root down to the changed field is rebuilt. A deep update touches one path through
the tree and shares everything off it, which is precisely the immutable-data
idiom every functional language uses to make ``copy the whole structure'' cheap.

\begin{example}[A nested update shares the untouched branch]
\label{ex:nested}
Let \lstinline[style=l4style]|state = { cfg: c, solver: { x: f, it: 2 } }|. The
expression \lstinline[style=l4style]|{ ...state, solver: { ...state.solver, it: 3 } }|
reduces to \lstinline[style=l4style]|{ cfg: c, solver: { x: f, it: 3 } }|. The
inner \code{x} field still holds the very same tensor \code{f} (shared, since it
did not change), and the outer \code{cfg} still holds the very same \code{c}.
Only two records were freshly built---the inner \code{solver} and the outer
\code{state}---and the original \code{state} is, as always, untouched and still
reports \lstinline[style=l4style]|state.solver.it| as $2$.
\end{example}

\subsection{Optional fields and demand-pruning}

A record type may declare a field \emph{optional}, written with a trailing
\lstinline[style=l4style]|?|:

\begin{lstlisting}[style=l4style]
{ x: Tensor[N], it: Int, residual?: Scalar }
\end{lstlisting}

The \code{residual?} field may or may not be present in a given value. Optionality
is the calculus's answer to a recurring need: a step can \emph{declare} an output
that downstream code might want---a residual norm for monitoring, say---without
forcing it to be computed when nobody asks. The rule, which \cref{ch:fold}
develops in full, is \emph{demand-pruning}: a field that no consumer reads is
never computed. \Cref{fig:optional} greys out exactly such a field.

\begin{figure}[t]
\centering
\begin{tikzpicture}[font=\footnotesize]
  \node[draw=simBlue, fill=simBgBlue, rounded corners=3pt, thick,
        minimum width=44mm, minimum height=26mm, inner sep=5pt] (s) at (0,0) {};
  \node[anchor=north, font=\scriptsize\bfseries, text=simBlue] at ($(s.north)+(0,-1mm)$) {step output};
  \node[anchor=west] at ($(s.north west)+(2.5mm,-7mm)$)  {\code{x}\,:\ theField};
  \node[anchor=west] at ($(s.north west)+(2.5mm,-13mm)$) {\code{it}\,:\ \(8\)};
  \node[anchor=west, text=simGray!60] at ($(s.north west)+(2.5mm,-19mm)$)
    {\itshape\code{residual?}\,:\ (not read \(\Rightarrow\) not computed)};
  \draw[simGray!40, densely dashed] ($(s.north west)+(2mm,-16mm)$) -- ($(s.north east)+(-2mm,-16mm)$);
\end{tikzpicture}
\caption[An optional, demand-pruned field]{A step declares an optional field
\code{residual?}. Because no downstream consumer reads it here, demand-pruning
elides it---the field is greyed out and never materializes. The same step,
consumed by code that \emph{does} read \code{residual?}, would compute it. The
declaration is fixed; whether the value is produced depends on demand.
\Cref{ch:fold} makes this precise.}
\label{fig:optional}
\end{figure}

This is a deep idea with a light footprint here: it replaces, in one stroke, the
logging flags, verbosity switches, and conditional ``if monitoring, also
compute\dots'' branches that clutter procedural simulators. An operator simply
declares every output it \emph{could} produce; consumers read what they need; the
unread outputs cost nothing. We meet it again, with its reduction rule, in
\cref{ch:fold}.

\section{The book's real records: a forward map}

Everything in this chapter is rehearsal for three records you will meet
everywhere in Part~III and beyond. We name them now---only their shapes, so that
when they appear they are familiar---and defer their full treatment to
\cref{ch:strata}, which exists precisely to dissect them.

\begin{definition}[The three workhorse records]
\label{def:workhorse}
The recurring state records of the simulator's solve are:
\begin{itemize}[nosep]
  \item \textbf{\code{OpParams}}---\emph{build-time configuration.} The readonly
  bundle a solver is set up with: variant selectors, tolerances, the constructed
  operator surfaces. Captured once, never re-evolved.
  \item \textbf{\code{SimState}}---\emph{run-time evolving state.} The
  caller-visible quantities a solve carries step to step: the current iterate, the
  iteration count, the convergence flag. The very record this chapter has been
  drawing as cards.
  \item \textbf{\code{SolveResult}}---\emph{a returned bundle.} The record one step
  hands back: the next state, the working data, and the demand-prunable per-step
  outputs.
\end{itemize}
\end{definition}

Their shapes, in the syntax of this chapter:

\begin{lstlisting}[style=l4style]
OpParams    = { tol: Float, max_it: Int, flexible: Bool, ... }   -- build-time, readonly
SimState    = { x: Tensor[N], it: Int, converged: Bool, ... }    -- run-time, evolving
SolveResult = { sim: SimState, outputs: StepOutputs, ... }       -- a step's returned bundle
\end{lstlisting}

Notice the strata at a glance. \code{OpParams} is fixed once and read forever---
the ``operator parameters'' a step closes over but never changes. \code{SimState}
is the bundle that threads, the one whose \code{it} ticks step by step.
\code{SolveResult} is what a step returns: it \emph{contains} a fresh
\code{SimState} and adds the per-step readouts. The discipline of this
chapter---read by destructuring, write by spreading, never mutate---is exactly how
every step over these records is written. \Cref{ch:strata} explains \emph{why} the
state splits into these three strata (build-time versus run-time, and the single
mutation point we flagged above); \cref{ch:monad} shows how the threading of
\code{SimState} through a chain of steps is itself packaged so the spreads stay
out of the algorithm's way. For this chapter it is enough that you can read each
shape as a card of named fields and know that ``updating'' any of them means
building a fresh one.

\summarysection
A \emph{record} is a bundle of named, typed fields,
written \lstinline[style=l4style]|{ x: theField, it: 7, converged: False }|, with
a field reached by the dot, \lstinline[style=l4style]|s.it|. Records carry the
state a simulation evolves. The governing discipline is \emph{immutability}: a
record never changes. ``Updating'' a record means building a fresh one with the
\emph{spread}, \lstinline[style=l4style]|{ ...s, it: s.it + 1 }|---a new value,
like \code{s} but for the named fields, with \code{s} left untouched. This
contrasts sharply with mutating a struct field in place, which destroys the old
value and ripples through every alias. Fields are pulled out by name with
\emph{destructuring}, \lstinline[style=l4style]|let { x, it } = s|, optionally
keeping a rest, \lstinline[style=l4style]|let { it, ...rest } = s|. Immutability
pays three dividends for simulation: the value-threading of fresh records
\emph{is} the dependency graph (\cref{ch:targets}); the short override list of a
spread advertises that typically a single field changes per step
(\cref{ch:strata}); and sharing a record is always safe, which lets ``fresh
copy'' be implemented as ``share what didn't change'' at no semantic cost. Records
nest, and optional fields (\lstinline[style=l4style]|residual?|) introduce
demand-pruning, by which an unread output is never computed (\cref{ch:fold}). The
book's three workhorse records---\code{OpParams} (build-time),
\code{SimState} (run-time), \code{SolveResult} (a step's return)---are previewed
here and dissected in \cref{ch:strata}.

\furtherreading
The record syntax used here---brace literals, field access by name, the spread
for non-destructive update, destructuring with rest---is the common surface of
the typed functional and modern scripting languages; the calculus borrows the
TypeScript spelling deliberately for readability. The deeper point, that
immutable values make sharing safe and dataflow explicit, is the thesis of pure
functional programming; the standard exposition, and the source of the
``everything is a value, updates build fresh values'' discipline this chapter
distills, is Peyton Jones's account of the lazy, purely functional
language~\citep{peytonjones2003}, where immutability and demand-driven evaluation
(our optional fields) appear as two faces of one design.

\exercisessection
\begin{enumerate}[nosep]
  \item Let \code{s = \{ x: a, it: 3, converged: False \}}.
  Write the value of each of \code{s.it},
  \code{\{ ...s, it: s.it + 1 \}}, and
  \code{\{ ...s, converged: True \}}. Then state the value of
  \code{s.it} \emph{after} all three expressions have been
  evaluated, and explain in one sentence why it is what it is.
  \item Rewrite the increment \code{\{ ...s, it: s.it + 1 \}}
  using destructure-with-rest instead of \code{...s}. Argue
  that the two produce the same record.
  \item A colleague writes, ``the spread \code{\{ ...s, it: 0 \}}
  resets \code{s}'s counter to zero.'' Correct the statement precisely. What, if
  anything, happens to \code{s}?
  \item Give a record value and a tuple value that carry the same three pieces of
  data. Then suppose a fourth field/component is inserted at the front. For each,
  say what (if anything) breaks in code that read the \emph{old} second
  field/component, and use this to explain when you would prefer a record over a
  tuple.
  \item Take \code{state = \{ cfg: c, solver: \{ x: f, it: 5 \} \}}.
  Write a single expression that produces a fresh \code{state} whose
  \code{solver.it} is $6$ and whose every other field (inner and outer) is
  unchanged.
  \item Explain why two names bound to the same record cause no aliasing hazard
  here, whereas two pointers to the same mutable struct do. Connect your answer
  to why a spread may safely \emph{share} the unchanged tensor field \code{x}
  rather than deep-copying it.
  \item A step record declares an optional field
  \code{residual?: Scalar}. Describe a consumer that causes it
  to be computed and a consumer that causes it to be pruned. In one sentence,
  state what procedural feature (flag, branch, or monad) the optional field
  replaces.
  \item In \cref{wex:history}, suppose a bug replaced the third step with an
  imagined in-place mutation \code{s2.it = 3; s2.converged = True} (not legal in
  our calculus). Which earlier records in the trace would now be wrong if other
  code still held them, and why? Use this to state, in your own words, the single
  guarantee immutability provides.
\end{enumerate}

\chapter{Tensors and Shapes}
\label{ch:tensors}

\chapterorient{\Cref{ch:bigidea} sold you a slogan---\emph{a field is a
tensor}---and then leaned on it without ever pinning down what a tensor
\emph{is}. This chapter pays that debt. We say precisely what a tensor is in our
calculus (an immutable, typed, multi-axis array of numbers), introduce
\emph{shapes} as the type that travels with every tensor and is checked at every
function boundary, and write down the handful of primitive operations---%
\code{axpy}, \code{dot}, \code{matvec}, and friends---each with its shape
contract drawn as a picture. Then comes the chapter's real payload: the
\emph{tensor-field lift}, the modeling move that turns a per-degree-of-freedom
loop over an array (the world of Part~II) into a single equation about a whole
field. By the end you will read a shape the way you read a type, and you will see
why ``loop over the entries'' and ``one statement about the field'' are usually,
but not always, the same thing. The careful \emph{algebra} of shapes---how they
combine, broadcast, and promote---is the next chapter, \cref{ch:shapes}; here we
fix the vocabulary and the simplest shapes.}

\section{What a tensor is}
\label{sec:t-whatis}

We have used the word ``tensor'' since \cref{ch:bigidea} as a synonym for ``a
rectangular block of numbers.'' That is the right intuition, but our calculus
asks three specific things of the word, and it is worth stating them plainly
because each one earns its keep later.

\begin{definition}[Tensor]
\label{def:tensor}
A \emph{tensor} is an \emph{immutable}, \emph{typed}, \emph{multi-axis} array of
scalars:
\begin{description}[leftmargin=1.6em, style=nextline, nosep]
\item[immutable] once constructed, its entries never change. An operation that
``modifies'' a tensor in fact returns a \emph{new} tensor and leaves its input
untouched (\cref{ch:bigidea}).
\item[typed] it carries a \emph{shape}---the number, sizes, and meanings of its
axes---and a \emph{scalar type} (real or complex). The shape is part of the
value, not a comment about it.
\item[multi-axis] its entries are addressed by a tuple of integer indices, one
per axis. A length-$N$ vector has one axis; an $m\times n$ matrix has two; the
$[H,W,3]$ field of \cref{ch:bigidea} has three.
\end{description}
\end{definition}

The \emph{number} of axes is the tensor's \emph{rank}: rank-1 for a vector,
rank-2 for a matrix, rank-0 for a single scalar (\cref{sec:t-scalars}). Do not
confuse this with the rank of a matrix in the linear-algebra sense (the dimension
of its column space); throughout this book ``rank'' means ``number of axes''
unless we say otherwise.

\subsection{Reconnecting to Part~II: every field we built is a tensor}
\label{sec:t-reconnect}

This is not a new object. Every quantity Part~II constructed is already a tensor;
we simply had not insisted on the word. Two cases anchor everything that follows.

\paragraph{A discretized field is a \code{Tensor[N]}.}
\Cref{ch:functionspaces} replaced a continuous field---the potential $V$, the
mode field $\vE$---with a finite list of \emph{degrees of freedom} (DoFs): the
coefficients that, multiplied against the basis functions, reconstruct the field
everywhere. If the function space has $N$ degrees of freedom, that list is $N$
numbers in a fixed order. As a tensor it has exactly one axis, of length $N$:

\begin{center}
\lstinline[style=l4style]|Tensor[N]| \quad---\quad a rank-1 tensor, the global
DoF vector.
\end{center}

\noindent This is \emph{the} workhorse shape of the whole solver. The unknown we
solve for, the right-hand side we solve against, every residual and search
direction a Krylov method forms (\cref{ch:matrixfree})---all are
\lstinline[style=l4style]|Tensor[N]|. The single axis $N$ counts degrees of
freedom, and which DoFs they are is the business of the function space; the
tensor only remembers \emph{how many} and \emph{in what order}.
\Cref{fig:t-meshfield} draws the correspondence: the continuous field over the
mesh on the left, and the flat \lstinline[style=l4style]|Tensor[N]| of its DoF
coefficients on the right.

\begin{figure}[t]
\centering
\begin{tikzpicture}[scale=1.0]
  % LEFT: a triangulated mesh with a few highlighted nodes (DoFs)
  \begin{scope}
    \draw[simGray, semithick]
      (-1.5,-1.0) -- (0.0,-1.1) -- (1.4,-0.8) --
      (1.2,0.7) -- (-0.2,1.0) -- (-1.6,0.6) -- cycle;
    \draw[simGray, semithick] (-1.5,-1.0) -- (-0.2,1.0);
    \draw[simGray, semithick] (0.0,-1.1) -- (-0.2,1.0);
    \draw[simGray, semithick] (0.0,-1.1) -- (1.2,0.7);
    \draw[simGray, semithick] (1.4,-0.8) -- (1.2,0.7);
    \draw[simGray, semithick] (-1.6,0.6) -- (-0.2,1.0);
    \draw[simGray, semithick] (-1.5,-1.0) -- (-1.6,0.6);
    % the DoF nodes, numbered
    \foreach \p/\n in {{(-1.5,-1.0)}/0, {(0.0,-1.1)}/1, {(1.4,-0.8)}/2,
                       {(1.2,0.7)}/3, {(-0.2,1.0)}/4, {(-1.6,0.6)}/5}{
      \fill[simRust] \p circle (2.1pt);
      \node[font=\tiny, simRust, anchor=south west] at \p {\n};}
    \node[font=\scriptsize\itshape, simGray] at (0,-1.6) {a field on a mesh};
    \node[font=\scriptsize\itshape, simGray] at (0,-1.95) {$N=6$ DoFs (one per node)};
  \end{scope}
  % collapse arrow
  \node[font=\Large, simBlue] at (2.7,0) {$\rightsquigarrow$};
  \node[font=\scriptsize\itshape, simBlue] at (2.7,0.55) {collect the};
  \node[font=\scriptsize\itshape, simBlue] at (2.7,-0.55) {DoF coefficients};
  % RIGHT: the Tensor[N] column with the 6 entries
  \begin{scope}[shift={(4.9,0)}]
    \fill[simBgBlue, draw=simBlue, thick] (-0.32,-1.2) rectangle (0.32,1.2);
    \foreach \k in {0,1,2,3,4,5}{
      \pgfmathsetmacro\yy{0.9-0.4*\k}
      \draw[simBlue!35] (-0.32,\yy+0.2)--(0.32,\yy+0.2);
      \node[font=\tiny, simRust] at (-0.62,\yy) {\k};
      \node[font=\scriptsize, simInk] at (0.0,\yy) {$c_{\k}$};}
    \node[font=\scriptsize, simRust] at (0.95,0) {$N$};
    \draw[-{Stealth[length=1.6mm]}, simRust] (0.62,1.0)--(0.62,-1.0);
    \node[font=\scriptsize\ttfamily, simInk] at (0.0,-1.6) {Tensor[N]};
    \node[font=\scriptsize\itshape, simGray] at (0.0,-1.95) {rank-1 DoF vector};
  \end{scope}
\end{tikzpicture}
\caption[A field over a mesh is a \code{Tensor[N]}]{The correspondence that
underwrites the whole calculus. \emph{Left:} a field defined over a mesh, sampled
by $N$ degrees of freedom---here one coefficient per node, $N=6$. \emph{Right:}
those coefficients collected, in their fixed order, into a single rank-1 tensor
\lstinline[style=l4style]|Tensor[N]|. The mesh, the basis, and the geometry choose
\emph{which} DoFs and in \emph{what order} (the function-space machinery of
\cref{ch:functionspaces}); the tensor remembers only the ordered list. From
Part~III's height, the field \emph{is} this \lstinline[style=l4style]|Tensor[N]|.}
\label{fig:t-meshfield}
\end{figure}

\paragraph{An element-local quantity is a small multi-axis tensor.}
\Cref{ch:matrixfree} showed that the moment we gather the global vector down to
individual elements, we acquire \emph{richer} shape. The gathered coefficients
form a \lstinline[style=l4style]|Tensor[E, L]| ($E$ elements, $L$ local DoFs
each); the values at quadrature points form a
\lstinline[style=l4style]|Tensor[E, P, C]| ($P$ points, $C$ components); the
geometry factor is a \lstinline[style=l4style]|Tensor[E, P, G]|. These are the
\emph{same kind of object} as the flat DoF vector---immutable, typed, multi-axis
arrays---only with more axes and more specific meanings. The gather/scatter pair
$G,\,G^{\mathsf T}$ of \cref{ch:matrixfree} is precisely the bridge between the
one-axis global vocabulary and the multi-axis element-local one.

\Cref{fig:t-box} draws three representative tensors as labeled boxes of axes---a
picture we will return to whenever a shape needs to be made concrete.

\begin{figure}[t]
\centering
\begin{tikzpicture}[scale=1.0]
  % --- Tensor[N]: a single tall column ---
  \begin{scope}[shift={(-5.0,0)}]
    \fill[simBgBlue, draw=simBlue, thick] (-0.28,-1.1) rectangle (0.28,1.1);
    \foreach \y in {-0.825,-0.55,-0.275,0,0.275,0.55,0.825}{
      \draw[simBlue!40] (-0.28,\y)--(0.28,\y);}
    \node[font=\scriptsize, simRust] at (0.0,1.45) {$N$};
    \draw[-{Stealth[length=1.6mm]}, simRust] (0.55,0.9)--(0.55,-0.9);
    \node[font=\scriptsize\ttfamily, simInk] at (0.0,-1.5) {Tensor[N]};
    \node[font=\scriptsize\itshape, simGray, text width=26mm, align=center]
      at (0.0,-2.05) {global DoF vector\\(rank 1)};
  \end{scope}
  % --- Tensor[m, n]: a small grid ---
  \begin{scope}[shift={(-1.4,0)}]
    \fill[simBgGreen, draw=simGreen, thick] (-0.85,-0.7) rectangle (0.85,0.7);
    \foreach \y in {-0.35,0,0.35}{\draw[simGreen!45] (-0.85,\y)--(0.85,\y);}
    \foreach \x in {-0.51,-0.17,0.17,0.51}{\draw[simGreen!45] (\x,-0.7)--(\x,0.7);}
    \node[font=\scriptsize, simInk] at (-1.15,0.0) {$m$};
    \node[font=\scriptsize, simRust] at (0.0,-1.0) {$n$};
    \node[font=\scriptsize\ttfamily, simInk] at (0.0,-1.5) {Tensor[m, n]};
    \node[font=\scriptsize\itshape, simGray, text width=26mm, align=center]
      at (0.0,-2.05) {a dense matrix\\(rank 2)};
  \end{scope}
  % --- Tensor[H, W, 3]: stacked grids ---
  \begin{scope}[shift={(3.2,0)}]
    \def\dx{0.17}\def\dy{0.11}
    \foreach \k/\c in {0/simBgGray,1/simBgViolet,2/simBgRust}{
      \begin{scope}[shift={(\k*\dx,\k*\dy)}]
        \fill[\c, draw=simGray] (-0.8,-0.6) rectangle (0.8,0.6);
        \draw[simGray!55] (-0.8,-0.2)--(0.8,-0.2);
        \draw[simGray!55] (-0.8,0.2)--(0.8,0.2);
        \draw[simGray!55] (-0.27,-0.6)--(-0.27,0.6);
        \draw[simGray!55] (0.27,-0.6)--(0.27,0.6);
      \end{scope}}
    \node[font=\scriptsize, simInk] at (-1.15,0.0) {$H$};
    \node[font=\scriptsize, simInk] at (0.0,-0.95) {$W$};
    \node[font=\scriptsize, simRust] at (1.2,0.85) {$3$};
    \node[font=\scriptsize\ttfamily, simInk] at (0.15,-1.5) {Tensor[H, W, 3]};
    \node[font=\scriptsize\itshape, simGray, text width=28mm, align=center]
      at (0.15,-2.05) {a 3-vector per grid cell\\(rank 3)};
  \end{scope}
\end{tikzpicture}
\caption[A tensor as a labeled box of axes]{Three tensors, drawn as labeled
boxes of their axes. \emph{Left:} the global DoF vector
\lstinline[style=l4style]|Tensor[N]|, a single axis of length $N$---the shape
every solver vector wears. \emph{Center:} a small dense matrix
\lstinline[style=l4style]|Tensor[m, n]|, two axes. \emph{Right:} a field storing
a 3-vector at each cell of an $H\times W$ grid,
\lstinline[style=l4style]|Tensor[H, W, 3]|, three axes with the last pinned to
size $3$. The \emph{rank} is the number of axes; the bracketed list is the
\emph{shape}.}
\label{fig:t-box}
\end{figure}

\subsection{Real and complex scalars}
\label{sec:t-scalars}

Every tensor's entries are scalars of a single type, and our calculus admits
two: \emph{real} ($\Real$) and \emph{complex} ($\Complex$). The distinction is
physical, not cosmetic. The electrostatic and magnetostatic analyses
(\cref{ch:targets}) solve over $\Real$: a potential or vector-potential DoF is a
real number. The driven and eigenmode analyses are
\emph{time-harmonic}---a field oscillating as $e^{i\omega t}$ is most naturally a
\emph{complex} amplitude, so their DoF vectors are
\lstinline[style=l4style]|Tensor[N]| over $\Complex$. The transient analysis
marches real fields in time. A tensor's scalar type travels with it just as its
shape does, and operations must agree on it: you cannot silently add a real
tensor to a complex one. The natural question---when \emph{may} a real tensor be
treated as complex (the \emph{promotion} question)---we raise in
\cref{sec:t-promotion} and answer with the rest of the shape algebra in
\cref{ch:shapes}.

\section{Shapes as type-level contracts}
\label{sec:t-shapes}

A tensor's shape is a small piece of data---a bracketed list of axes---but it
plays the role of a \emph{type}. It is declared at construction, it travels with
the value, and it is \emph{checked at every function boundary}. A function that
expects a \lstinline[style=l4style]|Tensor[n]| and is handed a
\lstinline[style=l4style]|Tensor[m, n]| is as wrong as a function that expects an
integer and is handed a string---and just as with types, catching the mismatch at
the boundary, before any arithmetic runs, is what makes the calculus safe to
compose.

\subsection{The bracket notation}
\label{sec:t-bracket}

We write a shape as a comma-separated list of axes inside square brackets,
following the value's element type \code{Tensor}:

\begin{itemize}[leftmargin=1.4em, nosep]
\item \lstinline[style=l4style]|Tensor[N]| --- rank 1, a single axis of length
$N$. The flat DoF vector.
\item \lstinline[style=l4style]|Tensor[m, n]| --- rank 2, a small dense matrix
of $m$ rows and $n$ columns.
\item \lstinline[style=l4style]|Tensor[H, W, 3]| --- rank 3, with \emph{named}
axes $H$ and $W$ and a last axis \emph{pinned} to the literal size $3$ (a
3-vector per cell).
\end{itemize}

\noindent Two things in this notation are load-bearing. First, an axis may be a
\emph{name} (like $N$, $H$, $W$, $m$, $n$) or a \emph{literal} (like the $3$): a
name stands for a size we do not fix in the signature but require to
\emph{match} wherever the same name appears, while a literal pins the axis to an
exact extent. Second---and this is the point of giving axes \emph{names} rather
than just sizes---a shape records the \emph{meaning} of each axis, not only its
length. In \lstinline[style=l4style]|Tensor[E, P, C]| (\cref{ch:matrixfree}) the
axes are not interchangeable: $E$ counts elements, $P$ quadrature points, $C$
components, and a function that contracts the wrong one is wrong even if the
lengths happen to coincide. The name is how we keep them straight.

\begin{keyidea}
A shape is a \emph{type}. It states the rank, the size of each axis, and---through
named axes---the \emph{meaning} of each axis. It is part of the value, travels
with it, and is checked at every function boundary. Reading a function's shape
signature tells you what it consumes and produces before you read a line of its
body; a shape mismatch is a type error, caught at the boundary, not a wrong
number discovered later.
\end{keyidea}

\begin{notationbox}
When two operands must have the \emph{same whole shape of unspecified rank}, the
source calculus writes a single \emph{shape variable}---e.g.\
\lstinline[style=l4style]|axpy :: Scalar -> Tensor[$S] -> Tensor[$S] -> Tensor[$S]|---%
where \lstinline[style=l4style]|$S| names one common shape shared by all three.
This rank-agnostic ``same shape'' notation, and the richer \emph{named shape
groups} that let two tensors agree on \emph{part} of their shape, are the subject
of \cref{ch:shapes}. In this chapter we keep to concrete axes
(\lstinline[style=l4style]|Tensor[N]|, \lstinline[style=l4style]|Tensor[m, n]|)
and read \lstinline[style=l4style]|$S| informally as ``the same shape on both
sides.''
\end{notationbox}

\section{The primitive operations and their shape contracts}
\label{sec:t-primitives}

Here is where the linear algebra you already know meets typed tensors. Our
calculus is built from a small set of \emph{primitive} operations---the verbs of
\cref{ch:bigidea}'s first question. Each is a pure function, and each carries a
\emph{shape contract}: a signature that says exactly which shapes it consumes and
produces, and how their axes must line up. We give the six that recur most, each
as a signature and a shape-matching picture; their full algebra (laws, identities,
the operators that fold them together) is \cref{ch:operators}.

\begin{lstlisting}[style=l4style]
axpy    :: Scalar -> Tensor[$S] -> Tensor[$S] -> Tensor[$S]   -- a*x + y
dot     :: Tensor[$S] -> Tensor[$S] -> Scalar                 -- inner product
matvec  :: Tensor[m, n] -> Tensor[n] -> Tensor[m]             -- A . v
outer   :: Tensor[m] -> Tensor[n] -> Tensor[m, n]             -- v (X) w
reduce  :: Tensor[..., k] -> Tensor[...]                      -- collapse one axis
broadcast :: Tensor[$S_a] -> (target: Shape) -> Tensor[target]-- stretch to fit
\end{lstlisting}

\noindent Read each signature left to right: the argument shapes, then the
arrow, then the result shape. The two that share a shape (\code{axpy},
\code{dot}) demand their operands be \emph{identical} in shape; the two that
\emph{change} shape (\code{matvec}, \code{outer}) state the relation between input
and output axes explicitly; and the two that move \emph{between} ranks
(\code{reduce}, \code{broadcast}) add or drop an axis. \Cref{fig:t-shapematch}
draws the central one, \code{matvec}, as a shape-matching diagram.

\paragraph{\code{axpy} and \code{dot} --- same-shape contracts.}
\code{axpy} computes $a\vx+\vy$: both vectors and the result are one common shape
\lstinline[style=l4style]|$S| (in practice \lstinline[style=l4style]|Tensor[N]|),
and the scalar $a$ scales every entry. \code{dot} consumes two tensors of the
same shape and returns a single \code{Scalar}---it collapses \emph{all} axes by
summation. Over $\Complex$, \code{dot} conjugates its first argument (the
Hermitian inner product $\vx^{\mathsf H}\vy$); over $\Real$ this is the ordinary
$\sum_i x_i y_i$. The contract these share is the simplest and the most common:
\emph{the shapes must match exactly}, and a mismatch is rejected at the boundary.

\paragraph{\code{matvec} --- the contraction contract.}
\code{matvec} applies a small dense matrix to a vector. Its signature
\lstinline[style=l4style]|matvec :: Tensor[m, n] -> Tensor[n] -> Tensor[m]| says
the matrix's \emph{second} axis ($n$, its columns) must match the vector's single
axis ($n$), they are summed away in the contraction, and what survives is the
matrix's \emph{first} axis ($m$, its rows). This is exactly the rule from a first
linear-algebra course---``columns of $A$ must match the length of $\vv$, and the
product has as many entries as $A$ has rows''---now written as a shape contract
the calculus checks. \Cref{fig:t-shapematch} is the picture: the inner $n$'s
meet and cancel; the outer $m$ is the answer's shape.

\begin{figure}[t]
\centering
\begin{tikzpicture}[scale=1.0]
  % matrix [m, n]
  \node[draw=simGreen, fill=simBgGreen, thick, minimum width=20mm,
        minimum height=15mm] (A) at (0,0) {};
  \node[font=\scriptsize, simInk] at ($(A.west)+(-0.45,0)$) {$m$};
  \node[font=\scriptsize\bfseries, simRust] at ($(A.south)+(0,-0.45)$) {$n$};
  \node[font=\scriptsize\ttfamily, simInk, above=1mm of A] {Tensor[m, n]};
  % vector [n]
  \node[draw=simBlue, fill=simBgBlue, thick, minimum width=6mm,
        minimum height=15mm, right=12mm of A] (v) {};
  \node[font=\scriptsize\bfseries, simRust] at ($(v.south)+(0,-0.45)$) {$n$};
  \node[font=\scriptsize\ttfamily, simInk, above=1mm of v] {Tensor[n]};
  % equals / arrow
  \node[font=\Large] at ($(v.east)+(0.9,0)$) {$\Rightarrow$};
  % result [m]
  \node[draw=simRust, fill=simBgRust, thick, minimum width=6mm,
        minimum height=15mm, right=18mm of v] (r) {};
  \node[font=\scriptsize, simInk] at ($(r.west)+(-0.45,0)$) {$m$};
  \node[font=\scriptsize\ttfamily, simInk, above=1mm of r] {Tensor[m]};
  % the matching n's: a brace tying A's n and v's n
  \draw[decorate, decoration={brace, amplitude=4pt, mirror}, simRust]
    ($(A.south east)+(0,-0.75)$) -- ($(v.south west)+(0,-0.75)$)
    node[midway, below=4pt, font=\scriptsize, simRust]
    {inner $n$ must match, then cancels};
  % the surviving m
  \draw[-{Stealth[length=1.8mm]}, simGray, densely dashed]
    (A.west) to[out=170,in=170,looseness=1.3] (r.west);
  \node[font=\scriptsize\itshape, simGray] at (1.6,1.55)
    {outer $m$ survives as the result shape};
\end{tikzpicture}
\caption[Shape matching at a \code{matvec} boundary]{The shape contract of
\lstinline[style=l4style]|matvec :: Tensor[m, n] -> Tensor[n] -> Tensor[m]|, drawn
as axis matching. The matrix's column axis $n$ and the vector's single axis $n$
must agree (rust brace); the contraction sums over them and they vanish. The
matrix's row axis $m$ is untouched and becomes the shape of the result. A vector
of any length other than $n$ is rejected at the boundary, before a single
multiply-add runs.}
\label{fig:t-shapematch}
\end{figure}

\paragraph{\code{outer} --- the rank-raising contract.}
\code{outer} is \code{matvec}'s opposite: it takes two vectors and
\emph{builds} a matrix, $\vv\vw^{\mathsf T}$, whose $(i,j)$ entry is $v_i w_j$.
Its signature \lstinline[style=l4style]|outer :: Tensor[m] -> Tensor[n] -> Tensor[m, n]|
adds an axis rather than removing one: no axis is required to match, and the
result's rank is the \emph{sum} of the inputs'.

\paragraph{\code{reduce} and \code{broadcast} --- moving between ranks.}
\code{reduce} collapses one axis by an associative combiner (a sum, a max, a
product), turning a \lstinline[style=l4style]|Tensor[..., k]| into a
\lstinline[style=l4style]|Tensor[...]|: the labeled axis is summed (or
maxed, $\dots$) away and the rest of the shape survives. \code{dot} is the
special case that reduces \emph{every} axis to a scalar. \code{broadcast} runs the
other way: it \emph{stretches} a tensor to a larger target shape by repeating it
along new or length-1 axes---the mechanism by which a single scalar weight
multiplies a whole field, or a per-row factor scales every column. Its exact
matching rule (which shapes may be broadcast to which) is shape \emph{algebra},
and we defer it to \cref{ch:shapes}; here we note only that it is the
rank-\emph{raising} counterpart of \code{reduce}.

\begin{notationbox}
These six are not the whole vocabulary---\cref{ch:operators} adds the named
verbs (\code{nrm2}, \code{scal}, the operator-apply \code{apply\_linop}) and the
combinators that fold them together. But they are the \emph{primitives} whose
shape contracts every later operation is built from, and the discipline they
establish---\emph{a signature is a contract, checked at the boundary}---is the one
that makes the calculus composable.
\end{notationbox}

\subsection{Worked example: shape-checking a small 1-D solve}
\label{sec:t-wex-shapes}

Nothing fixes the contracts in the mind like running a real vector through them
and watching the shapes line up. We take the discretized 1-D field of Part~II---a
\lstinline[style=l4style]|Tensor[N]| of nodal values---and exercise \code{axpy},
\code{dot}, and a \code{matvec} against the tridiagonal stiffness matrix, checking
the shape at every boundary.

\begin{workedexample}{Shape-checked \code{axpy}, \code{dot}, \code{matvec} on a 1-D field}{t-shapes}
\textbf{Setup.} A 1-D mesh with $N=3$ interior nodes carries a discretized field
$\vx=(x_0,x_1,x_2)=(1,2,3)$ and a second field $\vy=(y_0,y_1,y_2)=(4,5,6)$, both
\lstinline[style=l4style]|Tensor[N]| with $N=3$ over $\Real$. The 1-D stiffness
operator (the discrete second difference, $-\partial_{xx}$ on a unit-spacing
mesh) is the tridiagonal matrix
\[
  \mat{K} \;=\;
  \begin{pmatrix} 2 & -1 & 0 \\ -1 & 2 & -1 \\ 0 & -1 & 2 \end{pmatrix},
  \qquad \mat{K} : \texttt{Tensor[m, n]},\ m=n=3 .
\]

\tcblower
\textbf{\code{axpy} (same-shape contract).} Compute $\vy' = a\vx + \vy$ with
$a=2$. The contract
\lstinline[style=l4style]|axpy :: Scalar -> Tensor[$S] -> Tensor[$S] -> Tensor[$S]|
demands both vectors share a shape: here \lstinline[style=l4style]|Tensor[3]| and
\lstinline[style=l4style]|Tensor[3]| agree, so the call type-checks, and
\[
  \vy' = 2\cdot(1,2,3) + (4,5,6) = (6,9,12) \;:\; \texttt{Tensor[3]} .
\]
The result wears the same shape as the inputs---rank and length are preserved.

\textbf{\code{dot} (reduce-to-scalar contract).} Compute $r=\vx\cdot\vy$. The
contract \lstinline[style=l4style]|dot :: Tensor[$S] -> Tensor[$S] -> Scalar|
again demands matching shapes (both \lstinline[style=l4style]|Tensor[3]|) and
collapses \emph{all} axes:
\[
  r = 1{\cdot}4 + 2{\cdot}5 + 3{\cdot}6 = 32 \;:\; \texttt{Scalar} .
\]
The shape \lstinline[style=l4style]|Tensor[3]| has gone to rank-0---the axis was
summed away.

\textbf{\code{matvec} (contraction contract).} Compute $\vw = \mat{K}\vx$. The
contract \lstinline[style=l4style]|matvec :: Tensor[m, n] -> Tensor[n] -> Tensor[m]|
checks that $\mat{K}$'s column axis $n=3$ matches $\vx$'s single axis $3$---it
does---then sums over them, leaving the row axis $m=3$:
\[
  \mat{K}\vx =
  \begin{pmatrix} 2{\cdot}1-1{\cdot}2 \\ -1{\cdot}1+2{\cdot}2-1{\cdot}3 \\ -1{\cdot}2+2{\cdot}3 \end{pmatrix}
  = (0,\,0,\,4) \;:\; \texttt{Tensor[3]} .
\]
The inner $3$'s cancelled (\cref{fig:t-shapematch}); the surviving $m=3$ is the
result shape, again a \lstinline[style=l4style]|Tensor[3]| field.

\textbf{The shapes line up.} Every call's boundary check passed because the axes
agreed: \code{axpy} and \code{dot} on two \lstinline[style=l4style]|Tensor[3]|
fields, \code{matvec} with the matrix's column count matching the vector's length.
Hand \code{matvec} a length-$2$ vector instead and the check fails immediately
($n=3\neq 2$), before any multiply-add---the type error caught at the boundary,
exactly as a wrong-typed argument would be.
\end{workedexample}

\section{Immutability makes aliasing harmless}
\label{sec:t-aliasing}

\Cref{ch:bigidea} argued that we describe computations as pure functions over
immutable values, and \cref{ch:records} carried the discipline into records. In
tensor terms the immutability has a concrete, daily payoff worth stating on its
own: \emph{aliasing is harmless}.

When two names refer to the same tensor, we say they \emph{alias} it. In a
mutable world, aliasing is a hazard---if \code{a} and \code{b} name the same
array and someone writes through \code{a}, then \code{b}'s value changes out from
under whoever was reading it, and the bug is among the nastiest in numerical
code. The whole apparatus of ``defensive copies'' exists to guard against exactly
this. But our tensors are immutable: \emph{no one can write through any name},
because there is no write at all---an operation produces a fresh tensor and the
old one stands. So two names pointing at one tensor can never interfere. Aliasing
costs nothing and risks nothing; we share freely and copy never.

\begin{figure}[t]
\centering
\begin{tikzpicture}[scale=1.0]
  % the one immutable tensor
  \node[draw=simBlue, fill=simBgBlue, thick, rounded corners=2pt,
        minimum width=26mm, minimum height=12mm] (t) at (0,0)
    {\footnotesize immutable tensor};
  \node[font=\scriptsize\ttfamily, simInk, below=1mm of t] {Tensor[N]};
  % two names pointing at it
  \node[font=\ttfamily\small, simRust] (na) at (-3.3,1.4) {x};
  \node[font=\ttfamily\small, simRust] (nb) at (-3.3,-1.4) {y};
  \draw[-{Stealth[length=2mm]}, simRust] (na) -- (t.north west);
  \draw[-{Stealth[length=2mm]}, simRust] (nb) -- (t.south west);
  \node[font=\scriptsize\itshape, simGray, align=left, text width=30mm]
    at (-3.6,0) {two names,\\one shared value};
  % the "would-be write" that never happens
  \node[font=\scriptsize, simGray, align=center, text width=34mm]
    at (4.0,0.7) {a ``modification'' of \texttt{x}\dots};
  \node[draw=simGreen, fill=simBgGreen, thick, rounded corners=2pt,
        minimum width=22mm, minimum height=10mm] (t2) at (4.0,-0.8)
    {\footnotesize \emph{new} tensor};
  \draw[-{Stealth[length=2mm]}, simGreen] (t.east) -- node[above, font=\scriptsize]
    {\texttt{f}} (t2.west);
  \node[font=\scriptsize\itshape, simGray, below=1mm of t2, text width=34mm,
        align=center] {\dots produces a fresh value;\\ \texttt{y} still sees the original};
\end{tikzpicture}
\caption[Aliasing is safe under immutability]{Two names, \texttt{x} and
\texttt{y}, share a single immutable tensor (blue). Because the tensor cannot be
written through any name, the sharing is invisible and free---no defensive copy
is needed. A function \texttt{f} that ``modifies'' \texttt{x} does not touch the
shared value at all; it \emph{returns a new tensor} (green), and \texttt{y}
continues to see the original unchanged. Mutation hazards simply cannot arise.}
\label{fig:t-aliasing}
\end{figure}

\begin{pitfall}
``No copy needed'' is a statement about \emph{meaning}, exactly as in
\cref{ch:bigidea}. It does \emph{not} mean an implementation literally keeps
every tensor forever: once no live name refers to a tensor, its storage may be
reclaimed---or, better, \emph{reused} in place for the next result, recovering the
mutation an experienced numericist would have written by hand. Immutability frees
the implementation to alias and to reuse \emph{because} it has proved the sharing
can never be observed. The discipline that makes this precise---which tensor's
storage may be reused for which result---is the state monad of \cref{ch:monad}.
\end{pitfall}

\section{The tensor-field lift}
\label{sec:t-lift}

We come to the chapter's central modeling move, and the true bridge from Part~II
to Part~III. Part~II's assembly chapters spoke the language of \emph{loops over
arrays}: ``for each degree of freedom $i$, do this to entry $i$.'' Part~III speaks
the language of \emph{operations on whole fields}: ``do this to the field.'' The
\emph{tensor-field lift} is the move that turns the first into the second.

\subsection{From a per-DoF loop to a field equation}
\label{sec:t-lift-what}

Take the \code{axpy} update once more. As a loop over the $N$ degrees of freedom
it reads, entry by entry,
\begin{equation}
  y_i \;\leftarrow\; a\,x_i + y_i, \qquad i = 0,1,\dots,N-1 .
  \label{eq:t-axpy-loop}
\end{equation}
The loop is explicit: it names an index $i$, ranges it over the DoFs, and updates
one entry per pass. The \emph{lift} observes that nothing in
\cref{eq:t-axpy-loop} couples one index to another---entry $i$ is computed from
entry $i$ alone---so the whole loop is captured, with nothing lost, by a single
statement about the field as a vector in $\Real^N$ (or $\Complex^N$):
\begin{equation}
  \vy' \;=\; a\,\vx + \vy .
  \label{eq:t-axpy-lift}
\end{equation}
The index $i$ has \emph{vanished}. \Cref{eq:t-axpy-lift} is one equation about
two whole tensors \lstinline[style=l4style]|x, y :: Tensor[N]| and a scalar $a$;
the per-DoF loop is now \emph{implicit}, recoverable but no longer written. This
is exactly the pure \code{axpy} of \cref{ch:bigidea}, and the lift is \emph{why}
its signature can speak of whole tensors without ever mentioning $i$.

\begin{keyidea}
The \emph{tensor-field lift} replaces a per-degree-of-freedom (or per-element)
loop with a single statement about a \emph{whole tensor field}. The update
$y_i \leftarrow a x_i + y_i$ \emph{for all $i$} becomes the one field equation
$\vy' = a\vx + \vy$ in $\Real^N$; the loop over $i$ is left implicit. This is the
move that lets the calculus describe a solver as operations on fields rather than
as nested loops over arrays---the single most important translation between
Part~II's assembly view and Part~III's calculus view.
\end{keyidea}

\begin{figure}[t]
\centering
\begin{tikzpicture}[scale=1.0]
  % LEFT: the explicit per-DoF loop
  \begin{scope}
    \node[font=\bfseries\small, simRust] at (0,2.3) {per-DoF loop};
    \node[draw=simRust, fill=simBgRust, rounded corners=2pt, thick,
          minimum width=30mm, minimum height=22mm, align=left, font=\scriptsize\ttfamily,
          text=simInk] (loop) at (0,0.5)
      {for i in 0..N:\\[2pt]\ \ y[i] = a*x[i]\\[1pt]\ \ \ \ \ \ \ \ + y[i]};
    \node[font=\scriptsize\itshape, simGray, below=1.5mm of loop, text width=34mm,
          align=center] {$N$ separate updates,\\index $i$ explicit};
  \end{scope}
  % collapse arrow
  \node[font=\Large, simBlue] at (3.3,0.5) {$\rightsquigarrow$};
  \node[font=\scriptsize\itshape, simBlue, align=center, text width=22mm]
    at (3.3,1.5) {lift\\(collapse the loop)};
  % RIGHT: the single field equation
  \begin{scope}[shift={(6.8,0)}]
    \node[font=\bfseries\small, simBlue] at (0,2.3) {field equation};
    \node[draw=simBlue, fill=simBgBlue, rounded corners=2pt, thick,
          minimum width=30mm, minimum height=22mm, align=center, font=\normalsize,
          text=simInk] (eq) at (0,0.5) {$\vy' = a\,\vx + \vy$};
    \node[font=\scriptsize\itshape, simGray, below=1.5mm of eq, text width=36mm,
          align=center] {one statement on whole tensors;\\ index $i$ gone, implicit};
  \end{scope}
\end{tikzpicture}
\caption[The tensor-field lift]{The lift collapses a per-DoF loop into a single
whole-field equation. \emph{Left (rust):} the explicit loop runs $N$ separate
updates, naming and ranging the index $i$. \emph{Right (blue):} the lifted form
is one equation about the whole tensors $\vx,\vy\in\Real^N$; the loop over $i$ has
disappeared into the vector notation. The two compute the same thing whenever no
index depends on another's result---the condition examined in
\cref{sec:t-lift-valid}.}
\label{fig:t-lift}
\end{figure}

\subsection{When the lift is valid}
\label{sec:t-lift-valid}

The lift is not magic, and it is not always legal. \Cref{eq:t-axpy-lift} is
faithful to \cref{eq:t-axpy-loop} precisely because of a structural fact: the
update at index $i$ reads inputs only at index $i$, never at some other index $j$
whose value the loop is in the middle of changing. There are two ways a per-DoF
loop can earn its lift.

\begin{description}[leftmargin=1.4em, style=nextline]
\item[Per-DoF independence.] Each output entry depends only on input entries at
the \emph{same} index (or on read-only data). \code{axpy}, \code{scal} ($\vx
\mapsto a\vx$), and elementwise products are all of this kind, and they lift to
whole-field equations directly. The weight stage $D$ of the matrix-free operator
(\cref{ch:matrixfree})---a pointwise multiply against \code{geom\_data}---is the
canonical multi-axis example: ``embarrassingly parallel,'' exactly because every
index is independent, exactly the condition for a clean lift.
\item[The coupling \emph{is} the operator.] Some operations \emph{do} couple
indices---but the coupling is the defining content of the operation, not a
loop-carried accident. \code{matvec}, and the full operator application
\code{apply\_linop} ($\vy = \mat{A}\vx$) of \cref{ch:matrixfree}, mix every input
entry into every output entry; that is what a linear operator \emph{is}. Such an
operation lifts to a single field equation $\vy = \mat{A}\vx$ because the coupling
is captured \emph{by the operator itself}, not by the order in which a loop visits
indices. There is no loop-carried dependency to preserve; the dependency is the
matrix.
\end{description}

\subsection{When the lift fails: a sequential reduction}
\label{sec:t-lift-fails}

The lift fails exactly when a loop \emph{carries a dependency forward}: the update
at index $i$ reads a value the \emph{same loop} produced at an earlier index, so
the order of the loop is part of the answer. Two small examples make the failure
concrete.

A \emph{running sum} (prefix sum) over a tensor,
\begin{equation}
  s_i \;=\; s_{i-1} + x_i, \qquad s_0 = x_0,
  \label{eq:t-prefix}
\end{equation}
has no whole-field equation of the form ``$\vs = (\text{something})\,\vx$'' that
respects it as a single independent statement, because $s_i$ \emph{needs}
$s_{i-1}$---a value computed one step earlier in the same sweep. Try to write it
as ``do all indices at once'' and you cannot: index $i$ has nothing to read until
index $i-1$ is finished. The order is load-bearing.

The same obstruction, in the solver world, is a \emph{Gauss--Seidel} sweep, where
updating component $i$ uses the \emph{already-updated} components $j<i$ from the
\emph{current} sweep:
\begin{equation}
  x_i \;\leftarrow\; \frac{1}{a_{ii}}\Big(b_i - \sum_{j<i} a_{ij}\,x_j^{\text{new}}
       - \sum_{j>i} a_{ij}\,x_j^{\text{old}}\Big).
  \label{eq:t-gs}
\end{equation}
The term $x_j^{\text{new}}$ for $j<i$ is the loop reading its own fresh output:
the very loop-carried dependency the lift cannot absorb. \Cref{fig:t-liftfail}
contrasts the two situations side by side---the independent fan (lifts) against
the sequential chain (does not).

\begin{figure}[t]
\centering
\begin{tikzpicture}[scale=1.0]
  % LEFT: independent — lifts
  \begin{scope}
    \node[font=\bfseries\small, simBlue] at (0,2.1) {independent --- lifts};
    \foreach \i in {0,1,2,3}{
      \node[draw=simBlue, fill=simBgBlue, circle, minimum size=6mm,
            font=\scriptsize] (xi\i) at (\i*0.95,1.1) {$x_{\i}$};
      \node[draw=simBlue, fill=simBgBlue, circle, minimum size=6mm,
            font=\scriptsize] (yi\i) at (\i*0.95,-0.4) {$y'_{\i}$};
      \draw[-{Stealth[length=1.6mm]}, simBlue] (xi\i) -- (yi\i);
    }
    \node[font=\scriptsize\itshape, simGray, align=center, text width=40mm]
      at (1.4,-1.25) {each output reads its own index only\\
      $\Rightarrow$ collapses to $\vy'=a\vx+\vy$};
  \end{scope}
  % RIGHT: sequential — does not lift
  \begin{scope}[shift={(6.1,0)}]
    \node[font=\bfseries\small, simRust] at (0.4,2.1) {sequential --- does \emph{not} lift};
    \foreach \i in {0,1,2,3}{
      \node[draw=simRust, fill=simBgRust, circle, minimum size=6mm,
            font=\scriptsize] (si\i) at (\i*0.95,0.35) {$s_{\i}$};
    }
    \foreach \i/\j in {0/1,1/2,2/3}{
      \draw[-{Stealth[length=1.8mm]}, simRust, thick] (si\i) -- (si\j);}
    \node[font=\scriptsize\itshape, simGray, align=center, text width=42mm]
      at (1.4,-1.25) {$s_i$ needs $s_{i-1}$ --- a chain\\
      $\Rightarrow$ order is load-bearing, no field equation};
  \end{scope}
\end{tikzpicture}
\caption[When the lift fails]{The structural difference between a liftable and an
unliftable loop. \emph{Left (blue):} an \code{axpy}-style update fans
independently---each $y'_i$ reads only $x_i$ (and $y_i$), no arrow runs between
output indices---so the loop collapses to the field equation $\vy'=a\vx+\vy$.
\emph{Right (rust):} a running sum $s_i=s_{i-1}+x_i$ is a \emph{chain}---each
node depends on the one before---so the order of evaluation is part of the answer
and no single whole-field equation captures it. Sequential reductions of this
kind (Gauss--Seidel, triangular solves) are \emph{obstructions} to the lift, and
\cref{part:framework} records them as such rather than forcing a false collapse.}
\label{fig:t-liftfail}
\end{figure}

When the lift fails, we do \emph{not} pretend it succeeded. The honest record is
an \emph{obstruction}: a note that this loop carries a sequential dependency that
no whole-field statement captures, with the specific recurrence cited. Part~IV's
solver chapters meet several---the Gauss--Seidel and SSOR smoothers, the
triangular solves inside a direct factorization---and each is logged as a
sequential obstruction rather than lifted. (Forward reference: the iterative
solvers of \cref{part:framework} and the smoothers behind the multigrid
preconditioner are where these surface in force.) The lift is a powerful default,
not a universal law; knowing \emph{when} it fails is as important as knowing how
it works.

\subsection{Worked example: lifting a per-DoF update, and one that won't lift}
\label{sec:t-wex-lift}

We make the lift, and its failure, fully concrete on a tiny $N=4$ field: write a
per-DoF update first as an explicit loop, then as a single tensor-field equation,
and argue they agree---then exhibit one update where the argument breaks down.

\begin{workedexample}{The lift: a loop and a field equation that agree, and one that cannot}{t-lift}
\textbf{Setup.} Take $N=4$ and the fields $\vx=(1,2,3,4)$, $\vy=(10,20,30,40)$
as \lstinline[style=l4style]|Tensor[4]| over $\Real$, with scalar $a=3$.

\tcblower
\textbf{The update as an explicit loop.} The \code{axpy} update visits each index
in turn:
\begin{lstlisting}[style=l4style, language=]
for i in 0..4:
    yp[i] = a * x[i] + y[i]
-- i=0: 3*1+10 = 13     i=1: 3*2+20 = 26
-- i=2: 3*3+30 = 39     i=3: 3*4+40 = 52
\end{lstlisting}
producing $\vy'=(13,26,39,52)$. Four updates, the index $i$ named and ranged.

\textbf{The same update as a field equation.} The lift collapses the loop to one
statement on whole tensors,
\[
  \vy' \;=\; a\,\vx + \vy \;=\; 3\,(1,2,3,4) + (10,20,30,40) \;=\; (13,26,39,52),
\]
the very same four numbers. The index $i$ never appears.

\textbf{Why they agree.} Entry $i$ of the loop reads $x_i$ and $y_i$ and writes
$y'_i$---inputs and output \emph{all at the same index}, with no entry of $\vy'$
read before it is written. The updates are mutually independent, so doing them
``all at once'' as a vector equation reorders nothing that matters: this is the
per-DoF-independence condition of \cref{sec:t-lift-valid}, and the lift is exact.

\textbf{One that will not lift.} Now replace the body with a \emph{running sum}:
\begin{lstlisting}[style=l4style, language=]
s[0] = x[0]
for i in 1..4:
    s[i] = s[i-1] + x[i]      -- reads s[i-1]: this loop's own output
-- s = (1, 3, 6, 10)
\end{lstlisting}
Each $s_i$ reads $s_{i-1}$, a value \emph{this same loop produced one step
earlier}. There is no equation $\vs = a\vx + (\dots)$ that computes all four
entries independently and respects the recurrence: index $i$ has nothing to read
until index $i-1$ has finished. The order of evaluation is part of the answer, so
the lift \emph{fails}, and the honest record is a sequential
\emph{obstruction}---not a forced field equation
(\cref{fig:t-liftfail}). The same obstruction is the Gauss--Seidel sweep of
\cref{eq:t-gs}; \cref{part:framework} gives these their proper home in the fold
(\cref{ch:fold}).
\end{workedexample}

\section{Scalars as rank-0 tensors, and the promotion question}
\label{sec:t-scalars-rank0}

One loose end ties the chapter off and sets up the next. A single
scalar---the $a$ in \code{axpy}, the result of \code{dot}, a material
coefficient---is a tensor too: a \emph{rank-0} tensor, a box with \emph{no} axes
and exactly one entry. This is not a pedantic flourish. Treating a scalar as a
rank-0 tensor is what lets \code{broadcast} (\cref{sec:t-primitives}) stretch it
across a whole field uniformly: ``multiply the field by $a$'' is ``broadcast the
rank-0 $a$ to the field's shape, then multiply pointwise.'' The scalar and the
field meet through the same shape machinery as any two tensors.

\subsection{The promotion question, set up}
\label{sec:t-promotion}

That leaves the question \cref{sec:t-scalars} flagged: when may a \emph{real}
tensor be used where a \emph{complex} one is expected? The need is real and
concrete. A driven analysis (\cref{ch:targets}) works over $\Complex$, but its
mesh geometry, its material coefficients, and an initial real excitation are
naturally real; somewhere they must be \emph{promoted} into the complex
computation without a clumsy manual rewrite. Intuitively the promotion is the
obvious embedding $r \mapsto r + 0i$, and it should compose silently---adding a
real tensor to a complex one should ``just work,'' yielding a complex result.

But ``just work'' hides real questions. \emph{Which} direction promotes (real
into complex, never the reverse)? Does the result shape change (no---only the
scalar type)? When two operands disagree on \emph{both} shape and scalar type, in
which order are they reconciled? These are questions of \emph{shape and type
algebra}, and answering them precisely---alongside broadcasting's exact matching
rule and the named-shape-group notation for partial congruence---is the whole
business of the next chapter.

\begin{notationbox}
We have now met every piece the shape algebra operates on: tensors of a rank and
a scalar type, axes that are names or literals, the same-shape and contraction
contracts, rank-0 scalars, and the promotion of real into complex. What we have
\emph{not} done is give the rules by which shapes \emph{combine}---how broadcasting
decides two shapes are compatible, how a named shape group asserts partial
congruence, how scalar-type promotion composes with shape matching.
\Cref{ch:shapes} is exactly that algebra. This chapter fixed the vocabulary; the
next makes it compute.
\end{notationbox}

\summarysection
A \emph{tensor} in our calculus is an immutable, typed, multi-axis array of
scalars (real or complex); its \emph{rank} is its number of axes and its
\emph{shape} is the bracketed list of axis sizes and meanings. Every quantity
Part~II built is one: a discretized field is the rank-1 global DoF vector
\lstinline[style=l4style]|Tensor[N]|, and an element-local quantity
(\cref{ch:matrixfree}) is a small multi-axis tensor like
\lstinline[style=l4style]|Tensor[E, P, C]|. A shape is a \emph{type}: it travels
with the value and is checked at every function boundary, so a mismatch is caught
before any arithmetic runs. The primitive operations carry \emph{shape
contracts}---\code{axpy} and \code{dot} demand identical shapes, \code{matvec}
($[m,n]\cdot[n]\to[m]$) and \code{outer} relate input and output axes,
\code{reduce} and \code{broadcast} move between ranks. Because tensors are
immutable, \emph{aliasing is harmless}: two names may share one tensor with no
risk and no copy. The chapter's central move is the \emph{tensor-field lift}: a
per-DoF loop $y_i \leftarrow a x_i + y_i$ becomes the single field equation
$\vy'=a\vx+\vy$, the index left implicit. The lift is valid when each index is
independent, or when the coupling \emph{is} the operator (\code{apply\_linop}); it
\emph{fails} for a sequential reduction (a running sum, a Gauss--Seidel sweep),
which we record as an obstruction rather than force. Scalars are rank-0 tensors,
and the promotion of real into complex is set up here and made precise---with the
full shape algebra---in \cref{ch:shapes}.

\furtherreading
The matrix--vector and inner-product contracts of \cref{sec:t-primitives} are the
elementary operations of numerical linear algebra; Trefethen and
Bau~\citep{trefethenbau1997} give the cleanest conceptual account of why the
\emph{action} $\vv\mapsto\mat{A}\vv$, not the entries, is what algorithms consume
(the thread we pick up in \cref{ch:matrixfree}), and Golub and Van
Loan~\citep{golubvanloan2013} are the standard comprehensive reference for the
shapes and costs of these operations. The named-axis, shape-as-type discipline
sketched here is the same one modern tensor frameworks enforce at array
boundaries; we develop its algebra from scratch in \cref{ch:shapes} and apply it
to the calculus's verbs in \cref{ch:operators}, so no outside background is
assumed. The tensor-field lift and its sequential-obstruction failure mode return
throughout \cref{part:framework}, where the fold (\cref{ch:fold}) gives the
unliftable chains a home of their own.

\exercisessection
\begin{enumerate}[nosep]
  \item For each of \lstinline[style=l4style]|Tensor[N]|,
  \lstinline[style=l4style]|Tensor[m, n]|, and
  \lstinline[style=l4style]|Tensor[E, P, C]|, state the rank and say in one
  phrase what each axis \emph{counts}. Why is naming the axes (rather than only
  giving their lengths) load-bearing for the last one?
  \item The contract for \code{matvec} is
  \lstinline[style=l4style]|Tensor[m, n] -> Tensor[n] -> Tensor[m]|. A colleague
  hands the function a matrix of shape \lstinline[style=l4style]|Tensor[m, n]| and
  a vector of shape \lstinline[style=l4style]|Tensor[m]|. Exactly where does the
  shape check fail, and why is failing \emph{at the boundary} better than
  producing a wrong number?
  \item Write the \code{scal} operation $\vx \mapsto a\vx$ as (a) an explicit
  per-DoF loop and (b) a single field equation. Argue that the two compute the
  same thing, naming the structural property of the loop that licenses the lift.
  \item Give the shape contract (a signature
  \lstinline[style=l4style]|Tensor[...] -> ... -> Tensor[...]|) for an
  \emph{elementwise product} of two same-shape tensors, and for \code{outer}.
  Which one raises the rank, which preserves it, and why?
  \item Consider the prefix-sum recurrence $s_i = s_{i-1} + x_i$. Explain
  precisely why it has no whole-field equation of the form $\vs = \mat{A}\vx$ that
  honors it as a single independent statement. Contrast this with \code{matvec},
  which \emph{does} couple every index yet still lifts---what is different about
  the kind of coupling?
  \item Two names \code{x} and \code{y} alias one immutable
  \lstinline[style=l4style]|Tensor[N]|. A function returns a fresh tensor
  ``modifying'' \code{x}. State what \code{y} sees afterward, and explain why an
  implementation is nonetheless free to reuse \code{x}'s storage for the new
  result. What condition must hold for that reuse to be safe?
  \item A driven analysis (\cref{ch:targets}) computes over $\Complex$ but reads
  a real material coefficient. Describe the promotion $r \mapsto r + 0i$ as a map
  on scalar type, and say what it does (and does not) change about a tensor's
  shape. Why must promotion run real-into-complex and never the reverse?
  \item \emph{(Looking ahead.)} The weight stage $D$ of the matrix-free operator
  (\cref{ch:matrixfree}) is a pointwise multiply on a
  \lstinline[style=l4style]|Tensor[E, P, C]|. Argue that it lifts cleanly to a
  whole-field equation, identifying which of the two validity conditions of
  \cref{sec:t-lift-valid} it satisfies. Then explain why the \emph{scatter}
  stage $G^{\mathsf T}$, which sums shared-DoF contributions, is a different
  matter.
\end{enumerate}

\chapter{Named Shape Groups and Rank-Agnostic Congruence}
\label{ch:shapes}

\chapterorient{The previous chapter introduced tensors and promised a precise
account of their \emph{shapes}. This is that account. We will build a small
notation that lets a \emph{single} function signature describe an operation on
vectors, matrices, and higher tensors all at once---a property we call
\emph{rank-agnostic congruence}---and we will see why the obvious first
attempt, writing \code{Tensor[N]} for ``the same shape as the other operand,''
silently throws that property away. By the end you will be able to read and
write the shape contracts that label every function in
\cref{part:framework}: named shape groups for congruence, the range-first
operator-shape form, the fixed-meaning axes of the matrix-free kernel, and the
one typing rule---scalar promotion---that lets a real number quietly join a
complex computation.}

\section{The problem with \texorpdfstring{\code{Tensor[N]}}{Tensor[N]}}
\label{sec:sh-problem}

In \cref{ch:tensors} we wrote a tensor's shape as a bracketed list of axis
extents: a length-$N$ vector is \lstinline[style=l4style]|Tensor[N]|, an
$m\times n$ matrix is \lstinline[style=l4style]|Tensor[m, n]|, the
distribution field of \cref{ch:bigidea} is
\lstinline[style=l4style]|Tensor[H, W, VY=3, VX=3]|. Each entry in the bracket
is one axis; the \emph{number} of entries is the tensor's \emph{rank}. A
vector has rank $1$, a matrix rank $2$, the distribution field rank $4$.

This is a perfectly good notation for writing down \emph{one} concrete shape.
The trouble starts the moment we try to write a \emph{function} that should
work at more than one rank. Consider the humble \code{axpy}
operation---``$\alpha$ times $\vx$ plus $\vy$''---the workhorse of every
iterative solver in this book:
\[
  \mathrm{axpy}(\alpha,\vx,\vy) \;=\; \alpha\,\vx + \vy .
\]
What does \code{axpy} require of its arguments? Only that $\vx$ and $\vy$ have
\emph{the same shape}, whatever that shape is, so that the scaled $\vx$ can be
added to $\vy$ entry by entry. It does not care whether they are vectors,
matrices, or rank-$7$ tensors. Scaling and adding are \emph{elementwise}: they
look at one entry at a time and never consult the rank at all.

Now watch what happens when we reach for our only shape notation so far and
try to write \code{axpy}'s signature. The natural-looking attempt is
\begin{lstlisting}[style=l4style]
axpy :: Scalar -> Tensor[N] -> Tensor[N] -> Tensor[N]   -- WRONG (pins rank 1)
\end{lstlisting}
The intent was surely ``the two tensors share a shape, call it \texttt{N}, and
the result has that shape too.'' But that is \emph{not} what this signature
says. The bracket \lstinline[style=l4style]|[N]| has exactly one entry, so it
denotes a tensor of \emph{rank one}: a flat vector of length \texttt{N}. By
writing \lstinline[style=l4style]|Tensor[N]| three times we have indeed said
``all three have the same shape''---but we have \emph{also} said, without
meaning to, ``and that shape is one-dimensional.'' A caller who hands
\code{axpy} two $3\times 3$ matrices will be turned away at the door, even
though scaling and adding $3\times 3$ matrices is exactly as sensible as
scaling and adding vectors.

\begin{pitfall}
\lstinline[style=l4style]|Tensor[N]| does \emph{not} mean ``the same shape as
the other operand.'' It means ``a tensor of rank one, with its single axis of
length \texttt{N}.'' The single letter inside the brackets is a name for
\emph{one axis}, not for \emph{a whole shape of unknown rank}. Reusing the
same letter across operands forces them to agree---but it forces them to agree
\emph{at rank one}. Using \lstinline[style=l4style]|Tensor[N]| to mean
``congruent, any rank'' is the single most common shape-notation error, and it
silently narrows a general operation to vectors only.
\end{pitfall}

The mismatch is worth seeing at a glance.
\Cref{fig:sh-pinned-vs-flex} sets the rank-$1$ commitment of
\lstinline[style=l4style]|Tensor[N]| beside the flexible notation we are about
to build. The cure is a way to name a \emph{whole run of axes}---possibly
several, possibly we-don't-know-how-many---and to say ``this other tensor has
the same run.'' That is the named shape group, and the rest of the chapter
develops it.

\begin{figure}[t]
\centering
\begin{tikzpicture}[scale=1.0]
  % --- LEFT: Tensor[N], one pinned box ---
  \begin{scope}[shift={(-3.6,0)}]
    \node[font=\scriptsize\bfseries, simRust] at (0,1.55) {\code{Tensor[N]}};
    \node[font=\scriptsize\itshape, simGray, text width=30mm, align=center]
      at (0,1.15) {one axis: rank pinned to 1};
    \draw[draw=simRust, fill=simBgRust, thick] (-0.5,-0.45) rectangle (0.5,0.45);
    \node[font=\small, simInk] at (0,0) {$N$};
    % a lock to convey "pinned"
    \node[font=\scriptsize, simRust] at (0,-0.95) {rank $=1$, fixed};
  \end{scope}
  % divider
  \draw[simGray!40, densely dotted] (-1.4,-1.4) -- (-1.4,1.7);
  % --- RIGHT: Tensor[(S: ...)], flexible bracket ---
  \begin{scope}[shift={(1.9,0)}]
    \node[font=\scriptsize\bfseries, simBlue] at (0.55,1.55)
      {\lstinline[style=l4style]|Tensor[(S: ...)]|};
    \node[font=\scriptsize\itshape, simGray, text width=42mm, align=center]
      at (0.55,1.15) {a named run \texttt{S}: any number of axes};
    % a row of boxes under a spanning bracket
    \foreach \i/\lab in {0/{},1/{},2/{}}{
      \draw[draw=simBlue, fill=simBgBlue, thick]
        (\i*0.95-0.45,-0.45) rectangle (\i*0.95+0.45,0.45);
      \node[font=\footnotesize, simGray] at (\i*0.95, 0) {$\cdot$};
    }
    \node[font=\scriptsize, simGray] at (3.05,0) {$\cdots$};
    \draw[decorate, decoration={brace, amplitude=5pt, mirror}, simBlue]
      (-0.55,-0.6) -- (2.45,-0.6)
      node[midway, below=6pt, font=\scriptsize, simBlue] {the group \texttt{S}};
  \end{scope}
\end{tikzpicture}
\caption[Pinned axis versus flexible group]{The rank-$1$ commitment of
\lstinline[style=l4style]|Tensor[N]| (left, rust) versus the rank-agnostic
named group \lstinline[style=l4style]|Tensor[(S: ...)]| (right, blue).
\lstinline[style=l4style]|Tensor[N]| is a single boxed axis: its rank is
locked at one. \lstinline[style=l4style]|Tensor[(S: ...)]| is a labelled
bracket that spans \emph{some} run of axes---one, three, however many---and
gives that run a name, \texttt{S}, without ever committing to how many there
are.}
\label{fig:sh-pinned-vs-flex}
\end{figure}

\section{Named shape groups}
\label{sec:sh-groups}

The fix is to give a \emph{name} to a run of axes without fixing how many
axes the run contains. We write the binding form
\[
  \texttt{(S: ...)}
\]
and read it ``the shape pattern \texttt{...}, called \texttt{S}.'' The three
dots \texttt{...} are a genuine piece of the notation---a \emph{rank
wildcard} that matches \emph{zero or more} axes of any extent. So
\lstinline[style=l4style]|Tensor[(S: ...)]| is ``a tensor whose entire shape
is some run of axes, and we will call that run \texttt{S}.'' It is true of a
vector (the run is one axis), of a matrix (two axes), of a rank-$7$ tensor
(seven axes), and even of a scalar (the run is empty). The notation has
declined to commit to a rank, which is exactly the freedom \code{axpy} needs.

A name is only useful if we can refer back to it. The second half of the
notation is the \emph{use} form, written with a dollar sigil:
\[
  \texttt{\$S}
\]
which means ``a run of axes \emph{congruent to} the group \texttt{S} bound
earlier''---the same number of axes, each of the same extent. The sigil is
the whole point: it makes a \emph{back-reference} look different from a
\emph{fresh binding}. \lstinline[style=l4style]|(S: ...)| introduces the name;
\lstinline[style=l4style]|$S| reuses it.

\begin{definition}[Named shape group]
\label{def:sh-group}
A \emph{named shape group} \lstinline[style=l4style]|(S: ...)| binds a name
\texttt{S} to a contiguous run of tensor axes \emph{without committing to the
rank or to the individual extents of that run}. The binding occurs exactly
once, at the group's first appearance. Every later occurrence
\lstinline[style=l4style]|$S|, written with the \texttt{\$} sigil, is a
\emph{use}: it asserts that the shape it labels is \emph{congruent} to the
bound run---the same number of axes, matched extent for extent. Two shapes
are \emph{congruent} when they are identical axis-for-axis (modulo the
shape algebra of \cref{sec:sh-equational}).
\end{definition}

With this in hand, \code{axpy}'s signature becomes what we always meant:
\begin{lstlisting}[style=l4style]
axpy :: Scalar -> Tensor[(S: ...)] -> Tensor[$S] -> Tensor[$S]
\end{lstlisting}
Read left to right: \code{axpy} takes a scalar; then a tensor of \emph{some}
shape, which binds the group \texttt{S}; then a second tensor whose shape must
be \emph{congruent} to \texttt{S}; and it returns a tensor, again congruent
to \texttt{S}. The first operand \emph{introduces} the shape; the second
operand and the result \emph{borrow} it. Nothing here mentions a rank, so the
one signature describes \code{axpy} on vectors, on matrices, and on every
higher tensor at once. \Cref{fig:sh-hero} is the picture to keep in mind: a
tensor is a row of axis-boxes, a group is a bracket spanning some of them, and
a \lstinline[style=l4style]|$S| use is an arrow pointing back to the binding,
asserting ``the same run, here, again.''

\begin{figure}[t]
\centering
\begin{tikzpicture}[scale=1.0]
  % Top: binding operand --- a tensor as a row of axis boxes, bracketed as (S: ...)
  \begin{scope}[shift={(0,1.5)}]
    \foreach \i in {0,1,2}{
      \draw[draw=simBlue, fill=simBgBlue, thick]
        (\i*0.9,-0.35) rectangle (\i*0.9+0.7,0.35);
    }
    \node[font=\scriptsize, simGray] at (2.55,0) {$\cdots$};
    \draw[draw=simBlue, fill=simBgBlue, thick] (3.0,-0.35) rectangle (3.7,0.35);
    \draw[decorate, decoration={brace, amplitude=4pt}, simBlue]
      (-0.05,0.5) -- (3.75,0.5)
      node[midway, above=4pt, font=\scriptsize\bfseries, simBlue] (Sbind)
      {binds \texttt{S}};
    \node[font=\scriptsize\ttfamily, simInk, left=4mm of {(-0.05,0)}]
      (xlab) at (-1.7,0) {\lstinline[style=l4style]|Tensor[(S: ...)]|};
  \end{scope}
  % Middle: use operand --- congruent, $S
  \begin{scope}[shift={(0,0)}]
    \foreach \i in {0,1,2}{
      \draw[draw=simTeal, fill=simBgGreen, thick]
        (\i*0.9,-0.35) rectangle (\i*0.9+0.7,0.35);
    }
    \node[font=\scriptsize, simGray] at (2.55,0) {$\cdots$};
    \draw[draw=simTeal, fill=simBgGreen, thick] (3.0,-0.35) rectangle (3.7,0.35);
    \draw[decorate, decoration={brace, amplitude=4pt, mirror}, simTeal]
      (-0.05,-0.5) -- (3.75,-0.5)
      node[midway, below=4pt, font=\scriptsize\bfseries, simTeal] (Suse)
      {uses \texttt{\$S}};
    \node[font=\scriptsize\ttfamily, simInk] at (-1.7,0)
      {\lstinline[style=l4style]|Tensor[$S]|};
  \end{scope}
  % the back-arrow: $S use --> S binding (congruence assertion)
  \draw[-{Stealth[length=2.6mm]}, simRust, thick]
    (4.0,-0.1) .. controls (4.9,0.55) and (4.9,1.45) .. (4.0,2.0)
    node[midway, right=1pt, font=\scriptsize\itshape, simRust, align=left]
    {congruent:\\same run};
\end{tikzpicture}
\caption[The shape-algebra hero figure]{A named shape group at work. Each
tensor is a row of axis-boxes; its rank is how many boxes there are, left
unspecified by the trailing \texttt{...}. The top operand's whole shape is
bracketed and \emph{bound} to the name \texttt{S}. The bottom operand is
written \lstinline[style=l4style]|Tensor[$S]|: the \texttt{\$} sigil makes it a
\emph{use}, and the rust back-arrow is the assertion it carries---``my shape is
\emph{congruent} to the run that bound \texttt{S}.'' Neither tensor's rank is
fixed; only their \emph{agreement} is.}
\label{fig:sh-hero}
\end{figure}

\subsection{Why not just reuse a shape variable?}
\label{sec:sh-vs-variable}

A sharp reader will object: we could have written \code{axpy} with a single
\emph{shape variable} $\sigma$ standing for the whole shape,
\begin{lstlisting}[style=l4style]
axpy :: Scalar -> Tensor[s] -> Tensor[s] -> Tensor[s]
\end{lstlisting}
where \texttt{s} (the typeset $\sigma$) names ``the entire shape, any rank.''
That is correct, and when the \emph{whole} shape is shared it is the lighter
spelling. Named groups earn their keep in the cases a bare variable cannot
reach, of which there are two.

First, a group can name a \emph{partial} run while leaving the rest free.
Suppose an operation shares a leading block of axes between two operands but
lets each carry its own trailing axis:
\begin{lstlisting}[style=l4style]
batched_op :: Tensor[(S: ...), k] -> Tensor[$S, m] -> Tensor[$S, k, m]
\end{lstlisting}
Here \texttt{S} is the shared leading run; \texttt{k} and \texttt{m} are
free trailing axes that differ between the operands. A single variable
$\sigma$ cannot express ``share \emph{these} axes but not \emph{that} one.''
The group splits a shape into a named congruent part and a free part---the
binding \lstinline[style=l4style]|(S: a, ...)| even mixes a pinned leading
axis \texttt{a} with a wildcard tail, so a group is as fine-grained as the
contract needs.

Second, and more quietly, a group makes the rank-agnosticism
\emph{visible}. A bare $\sigma$ relies on the reader knowing the convention
that a lone shape variable spans any rank; the trailing \texttt{...} inside
\lstinline[style=l4style]|(S: ...)| says it on the page. When the surrounding
code also uses concrete rank-$1$ shapes like
\lstinline[style=l4style]|Tensor[N]| for genuine flat vectors---and the next
section shows it does---the explicit wildcard is what stops a reader from
misreading the congruence contract as another rank-$1$ commitment.

\begin{keyidea}
A named shape group \lstinline[style=l4style]|(S: ...)| binds a name to a run
of axes \emph{of unknown rank}; a later use \lstinline[style=l4style]|$S|
asserts a second shape is \emph{congruent} to it. This is
\emph{rank-agnostic congruence}: one signature, e.g.\
\lstinline[style=l4style]|axpy :: Scalar -> Tensor[(S: ...)] -> Tensor[$S] -> Tensor[$S]|,
correctly describing an operation on vectors, on matrices, and on every
higher tensor simultaneously---because it constrains the operands to
\emph{agree} without ever constraining their \emph{rank}. The
\texttt{\$} sigil is the difference between a fresh binding and a back-reference.
\end{keyidea}

\section{Worked example: \texorpdfstring{\code{axpy}}{axpy} at two ranks}
\label{sec:sh-wex-axpy}

The claim is that one signature serves all ranks. The way to believe it is to
run the same signature against two concrete arguments---a vector and a
matrix---and watch the congruence check succeed both times, then watch the
wrong signature reject the matrix.

\begin{workedexample}{Why \code{axpy} wants a named group}{sh-axpy}
\textbf{The signature under test.}
\begin{lstlisting}[style=l4style]
axpy :: Scalar -> Tensor[(S: ...)] -> Tensor[$S] -> Tensor[$S]
\end{lstlisting}
We instantiate it twice and resolve the group \texttt{S} each time.

\tcblower
\textbf{Instance 1 --- a rank-$1$ vector.} Apply \code{axpy} with
$\alpha = 2$, $\vx,\vy$ of shape \lstinline[style=l4style]|[3]| (length-$3$
vectors). The first tensor operand binds the group:
\[
  \texttt{S} \;\mapsto\; \texttt{[3]} \qquad(\text{rank } 1).
\]
The second operand's required shape is \lstinline[style=l4style]|$S = [3]|: the
supplied $\vy$ is \lstinline[style=l4style]|[3]|, which is congruent.
\emph{Check passes.} The result shape is \lstinline[style=l4style]|$S = [3]|,
a length-$3$ vector. Elementwise: $2x_i + y_i$ at each of the $3$ entries.

\textbf{Instance 2 --- a rank-$2$ matrix-shaped field.} Now apply the
\emph{same} \code{axpy} with $\alpha = 2$, but $\vx,\vy$ of shape
\lstinline[style=l4style]|[3, 3]| (think of a $3\times 3$ block of a
tensor-valued material coefficient, or a matrix-shaped field). The first
operand binds the group to a \emph{rank-$2$} run:
\[
  \texttt{S} \;\mapsto\; \texttt{[3, 3]} \qquad(\text{rank } 2).
\]
The second operand's required shape is now \lstinline[style=l4style]|$S = [3, 3]|;
the supplied $\vy$ is \lstinline[style=l4style]|[3, 3]|, congruent.
\emph{Check passes.} The result is \lstinline[style=l4style]|[3, 3]|;
elementwise $2x_{ij} + y_{ij}$ across all $9$ entries. The signature never
mentioned a rank, so it bound \texttt{S} to whatever the first operand
brought---rank $1$ in instance~1, rank $2$ in instance~2.

\textbf{The wrong signature, same two instances.} Replace the group with the
rank-$1$ form:
\begin{lstlisting}[style=l4style]
axpy_bad :: Scalar -> Tensor[N] -> Tensor[N] -> Tensor[N]
\end{lstlisting}
\emph{Instance 1} still type-checks: the vectors are
\lstinline[style=l4style]|[3]|, so \texttt{N} resolves to $3$ and all three
slots agree. \emph{Instance 2 fails.} The matrices are
\lstinline[style=l4style]|[3, 3]|---rank $2$---but
\lstinline[style=l4style]|Tensor[N]| demands rank $1$. There is no value of
the single axis \texttt{N} that can match a two-axis shape: matching
\lstinline[style=l4style]|[N]| against \lstinline[style=l4style]|[3, 3]|
fails on \emph{rank} before any extent is even compared. The matrix case is
rejected---not because scaling-and-adding a matrix is ill-defined (it is
perfectly well-defined) but because the signature accidentally outlawed it.
\end{workedexample}

\begin{remark}
Notice \emph{where} the two signatures diverge: at the rank check, which
happens \emph{before} extents are compared. \lstinline[style=l4style]|Tensor[N]|
fixes the number of axes to one; \lstinline[style=l4style]|Tensor[(S: ...)]|
leaves it open and binds it from the argument. Everything downstream---the
elementwise arithmetic, the laws like
\lstinline[style=l4style]|axpy 0 v y = y|---is identical between the two ranks.
The \emph{only} thing the group buys is permission to run at a rank the
operation never cared about in the first place.
\end{remark}

\Cref{fig:sh-boundary} abstracts what just happened into the rule the matcher
follows at \emph{every} function boundary: bind the group from the first
shape it labels, then test each later use for congruence, rank first.

\begin{figure}[t]
\centering
\begin{tikzpicture}[node distance=7mm and 10mm]
  % the function boundary as a vertical gate
  \node[stage, minimum height=30mm, text width=4mm] (gate) at (0,0) {};
  \node[font=\scriptsize\bfseries, simBlue, above=1mm of gate]
    {function boundary};
  % incoming binding operand
  \node[data, left=14mm of gate.west, yshift=8mm, text width=22mm] (bind)
    {1st operand\\\lstinline[style=l4style]|Tensor[(S: ...)]|};
  % incoming use operand
  \node[data, left=14mm of gate.west, yshift=-8mm, text width=22mm] (use)
    {2nd operand\\\lstinline[style=l4style]|Tensor[$S]|};
  % the resolved value of S (the gate's memory)
  \node[font=\scriptsize\ttfamily, simRust, right=4mm of gate, yshift=8mm,
        align=left] (resolved)
    {bind:\\ $\texttt{S}\mapsto$ this shape};
  \node[font=\scriptsize, simTeal, right=4mm of gate, yshift=-8mm,
        align=left] (checked)
    {check:\\ rank? then extents?};
  \draw[flow] (bind.east) -- node[above, font=\scriptsize, simGray]{bind} (gate.west|-bind.east);
  \draw[flow] (use.east) -- node[below, font=\scriptsize, simGray]{test} (gate.west|-use.east);
  \draw[refedge, simRust] (gate.east|-resolved) -- (resolved);
  \draw[refedge, simTeal] (gate.east|-checked) -- (checked);
  % the two-step rule annotated below
  \node[font=\scriptsize\itshape, simGray, below=8mm of gate, text width=70mm,
        align=center]
    {\textbf{1.}~rank mismatch $\Rightarrow$ reject \emph{immediately}
     (this is where \lstinline[style=l4style]|Tensor[N]| fails the matrix).\quad
     \textbf{2.}~ranks agree $\Rightarrow$ compare extents, axis by axis,
     modulo the shape algebra (\cref{sec:sh-equational}).};
\end{tikzpicture}
\caption[The congruence check at a boundary]{What the shape matcher does at a
function boundary. The \emph{first} shape that labels a group \texttt{S}
\emph{binds} it (rust)---fixing both the rank and the extents \texttt{S} now
stands for. Every \emph{later} use \lstinline[style=l4style]|$S| is
\emph{tested} for congruence (teal): \textbf{first} the rank must match---a
mismatch is rejected on the spot, which is exactly how
\lstinline[style=l4style]|Tensor[N]| turns away a rank-$2$ argument
(\cref{wex:sh-axpy})---and \textbf{then}, only if the ranks agree, the
extents are compared axis by axis. The same two-step gate guards \code{axpy},
\code{apply\_linop}, and every other signature in this book.}
\label{fig:sh-boundary}
\end{figure}

\section{Operator shapes: domain and range}
\label{sec:sh-operators}

Congruence groups describe operations whose operands share a shape. But the
most important objects in this book---the linear operators $\mat{K}$, $\mat{M}$,
$\mat{A}(\omega)$ of \cref{ch:targets}---are maps whose \emph{input} and
\emph{output} shapes generally \emph{differ}. A matrix that takes a length-$n$
vector to a length-$m$ vector has two shapes to name, not one. The notation
extends to cover this with two groups: a \emph{range} group and a
\emph{domain} group.

We write them \emph{range-first}:
\begin{lstlisting}[style=l4style]
LinOp[(R: ...), (D: ...)]
\end{lstlisting}
for an operator from domain shape \texttt{D} to range shape \texttt{R}, both
of unknown rank. The order---range, then domain---is the matrix convention in
disguise. An $m\times n$ matrix is written rows-first ($m$) and maps a
length-$n$ \emph{domain} to a length-$m$ \emph{range}; the symbol order
$(m,n)$ is (range, domain). \lstinline[style=l4style]|LinOp[(R: ...), (D: ...)]|
is the rank-agnostic version of that same $(m,n)$ ordering. Applying the
operator is then a pair of uses, one for each group:
\begin{lstlisting}[style=l4style]
apply_linop :: LinOp[(R: ...), (D: ...)] -> Tensor[$D] -> Tensor[$R]
\end{lstlisting}
the operand is a use of the \emph{domain} group (\lstinline[style=l4style]|$D|);
the result is a use of the \emph{range} group (\lstinline[style=l4style]|$R|).
\Cref{fig:sh-operator} draws the flow: a domain-shaped tensor enters, the
operator carries it across, a range-shaped tensor leaves.

\begin{figure}[t]
\centering
\begin{tikzpicture}[node distance=10mm and 13mm]
  \node[data, text width=18mm] (in) {input\\\lstinline[style=l4style]|Tensor[$D]|};
  \node[stage, right=of in, text width=34mm] (op)
    {\lstinline[style=l4style]|LinOp[(R: ...), (D: ...)]|};
  \node[data, right=of op, text width=18mm] (out)
    {output\\\lstinline[style=l4style]|Tensor[$R]|};
  \draw[flow] (in) -- node[above, font=\scriptsize, simGray] {domain} (op);
  \draw[flow] (op) -- node[above, font=\scriptsize, simGray] {range} (out);
  % range-first annotation under the operator
  \node[font=\scriptsize\itshape, simGray, below=5mm of op, text width=46mm,
        align=center]
    {written \emph{range-first}: \texttt{(R: ...)} then \texttt{(D: ...)},\\
     matching the $m\times n$ (rows $\times$ cols) matrix order};
  % brace pointing out the (R,D) ordering
  \draw[decorate, decoration={brace, amplitude=4pt}, simBlue]
    ([yshift=6mm]op.north west) -- ([yshift=6mm]op.north east)
    node[midway, above=4pt, font=\scriptsize, simBlue]
    {range $\leftarrow$ domain};
\end{tikzpicture}
\caption[Operator shapes are range-first]{A linear operator names two shape
groups: a range group \texttt{R} and a domain group \texttt{D}, written
\emph{range-first} to match the rows-then-columns ordering of an $m\times n$
matrix. Application \lstinline[style=l4style]|apply_linop op| takes a
domain-shaped input \lstinline[style=l4style]|Tensor[$D]| and produces a
range-shaped output \lstinline[style=l4style]|Tensor[$R]|. Both groups are
rank-agnostic: the operator could map vectors to vectors, fields to fields,
or one rank to another.}
\label{fig:sh-operator}
\end{figure}

\subsection{The square case: \texorpdfstring{$\mat{A}\colon V\to V$}{A : V to V}}
\label{sec:sh-square}

A great many operators in this book are \emph{endomorphisms}: their domain and
range are the same space. The system matrix $\mat{K}$ of an eigenproblem, a
preconditioner $\mat{M}^{-1}$, the time-stepping operator of a transient
solve---each takes a state vector and returns a vector of the very same shape.
For these we do not need two independent groups; we bind \emph{one} and use it
for the range as well:
\begin{lstlisting}[style=l4style]
LinOp[(S: ...), $S]
\end{lstlisting}
``an operator whose range shape \emph{is} its domain shape \texttt{S}.'' This
is the operator-shape echo of the congruent \code{axpy} signature: there, the
two \emph{operands} shared a group; here, the operator's \emph{domain and
range} share one. Applied, a square operator threads a shape straight through:
\begin{lstlisting}[style=l4style]
apply_linop :: LinOp[(S: ...), $S] -> Tensor[$S] -> Tensor[$S]
\end{lstlisting}
which is exactly the signature a Krylov solver wants from its operator
(\cref{ch:matrixfree})---feed it a vector, get back a vector you can feed it
again.

\begin{workedexample}{An operator-apply congruence check}{sh-matvec}
\textbf{Setup.} Take the matrix-free mass operator $\mat{M}$ of
\cref{sec:mf-wex1}, acting on the flat global dof-vector. At L1 the global
vector is a genuine rank-$1$ object, so we instantiate the square form with a
concrete rank-$1$ domain.
\begin{lstlisting}[style=l4style]
applyM :: LinOp[(S: ...), $S] -> Tensor[$S] -> Tensor[$S]
\end{lstlisting}

\tcblower
\textbf{Bind from the operator's domain.} The operator $\mat{M}$ acts on
length-$N$ dof-vectors, $N=3$ in that example. Instantiating the square
operator type against $\mat{M}$ resolves
\[
  \texttt{S} \;\mapsto\; \texttt{[3]} ,
\]
so both the domain and range are \lstinline[style=l4style]|[3]|.

\textbf{Check the operand.} We apply $\mat{M}$ to $\vv=(1,2,3)$, shape
\lstinline[style=l4style]|[3]|. The required operand shape is
\lstinline[style=l4style]|$S = [3]|; the supplied $\vv$ is
\lstinline[style=l4style]|[3]|. \emph{Congruent---check passes.}

\textbf{Check the result.} The result shape is \lstinline[style=l4style]|$S = [3]|.
The pipeline of \cref{sec:mf-wex1} returned $\mat{M}\vv = (0.667, 2.000, 1.333)$,
shape \lstinline[style=l4style]|[3]|. \emph{Congruent---and crucially, the
output can be fed straight back into \code{applyM}}, because its shape is again
\texttt{S}. That re-feedability is precisely what the square form
\lstinline[style=l4style]|LinOp[(S: ...), $S]| guarantees, and it is what makes
the Krylov subspace $\{\vv, \mat{M}\vv, \mat{M}^2\vv, \dots\}$ of
\cref{ch:matrixfree} type-check at every power.
\end{workedexample}

\begin{notationbox}
At the \emph{lowest} layers of the spec---where Palace's operators act on
genuine flat dof-vectors of a definite length---the older concrete spelling
\lstinline[style=l4style]|LinearOperator[M, N]| over rank-$1$
\lstinline[style=l4style]|Tensor[N]| vectors is faithful and is kept. There
the lengths $M$, $N$ really are single flat-vector lengths, and the rank-$1$
commitment is correct rather than accidental. The group form
\lstinline[style=l4style]|LinOp[(R: ...), (D: ...)]| is the rank-agnostic
\emph{calculus} rendering, used at the higher layers where the same operator
is described over structured, possibly higher-rank shapes. Both spellings
coexist; choosing between them is choosing whether the rank is genuinely
fixed (keep \lstinline[style=l4style]|[N]|) or genuinely open (use a group).
\end{notationbox}

\section{Named axes of fixed meaning}
\label{sec:sh-fixedaxes}

There is a second kind of name in our shape brackets, and it is easy to
confuse with a congruence group until the distinction is drawn sharply. A
congruence group like \texttt{S} is a \emph{variable}: it stands for whatever
run of axes the first operand happens to bring, and its whole job is to assert
that two shapes \emph{agree}. A \emph{named axis of fixed meaning} is the
opposite: it is a single axis whose letter denotes one specific, fixed
quantity, the same way \lstinline[style=l4style]|Tensor[H, W, VY=3, VX=3]| of
\cref{ch:bigidea} names a height, a width, and two velocity axes. \texttt{H}
is not a congruence variable to be resolved; it is \emph{the height}, always.

The book's recurring family of fixed-meaning axes is the element-local
vocabulary of the matrix-free kernel, which we met in \cref{ch:matrixfree}.
Recall its five letters:
\begin{description}[leftmargin=1.6em, style=nextline, nosep]
\item[\texttt{E}] the element count---how many mesh elements the operator runs over.
\item[\texttt{L}] local dofs per element, $(p+1)^d$ for a degree-$p$ element.
\item[\texttt{P}] quadrature points per element.
\item[\texttt{C}] value (or derivative) components carried at a quadrature point.
\item[\texttt{G}] geometry-factor components stored per quadrature point.
\end{description}
and its three rank-structured tensors, \lstinline[style=l4style]|Tensor[(E, L)]|,
\lstinline[style=l4style]|Tensor[(E, P, C)]|, and
\lstinline[style=l4style]|Tensor[(E, P, G)]|. Each of these five letters has a
\emph{definite} meaning fixed by the finite-element construction; none of them
is a back-reference, and none is rank-agnostic.
\Cref{fig:sh-axis-vs-group} sets the two kinds of name side by side.

\begin{figure}[t]
\centering
\begin{tikzpicture}[scale=1.0]
  % LEFT: fixed-meaning axes E, P, C
  \begin{scope}[shift={(-3.5,0)}]
    \node[font=\scriptsize\bfseries, simGreen] at (0.0,1.7)
      {\lstinline[style=l4style]|Tensor[(E, P, C)]|};
    \node[font=\scriptsize\itshape, simGray, text width=34mm, align=center]
      at (0.0,1.25) {fixed-meaning axes};
    \foreach \i/\lab/\mean in {0/E/{elements},1/P/{quad pts},2/C/{comps}}{
      \draw[draw=simGreen, fill=simBgGreen, thick]
        (\i*1.05-0.45,-0.35) rectangle (\i*1.05+0.45,0.35);
      \node[font=\small, simInk] at (\i*1.05,0) {$\lab$};
      \node[font=\scriptsize, simGray] at (\i*1.05,-0.75) {\mean};
    }
    \node[font=\scriptsize\itshape, simGreen, text width=34mm, align=center]
      at (0.0,-1.35) {each letter = one fixed quantity\\ (a definite role)};
  \end{scope}
  % divider
  \draw[simGray!40, densely dotted] (-1.05,-1.7) -- (-1.05,2.0);
  % RIGHT: congruence group S
  \begin{scope}[shift={(2.0,0)}]
    \node[font=\scriptsize\bfseries, simBlue] at (0.55,1.7)
      {\lstinline[style=l4style]|Tensor[(S: ...)] / Tensor[$S]|};
    \node[font=\scriptsize\itshape, simGray, text width=42mm, align=center]
      at (0.55,1.25) {congruence variable};
    \foreach \i in {0,1}{
      \draw[draw=simBlue, fill=simBgBlue, thick]
        (\i*0.85,-0.35) rectangle (\i*0.85+0.65,0.35);
    }
    \node[font=\scriptsize, simGray] at (1.95,0) {$\cdots$};
    \draw[decorate, decoration={brace, amplitude=4pt, mirror}, simBlue]
      (-0.05,-0.5) -- (2.3,-0.5)
      node[midway, below=4pt, font=\scriptsize, simBlue] {\texttt{S}: any run};
    \node[font=\scriptsize\itshape, simBlue, text width=42mm, align=center]
      at (0.55,-1.35) {one name = ``the same run again''\\ (a back-reference, any rank)};
  \end{scope}
\end{tikzpicture}
\caption[Fixed-meaning axis versus congruence group]{Two kinds of name in a
shape bracket. \emph{Left (green):} a \emph{fixed-meaning axis}---\texttt{E},
\texttt{P}, \texttt{C} of the matrix-free kernel (\cref{ch:matrixfree}). Each
letter denotes one specific finite-element quantity; the rank is fixed (here,
three) and each axis has a definite role. \emph{Right (blue):} a
\emph{congruence group} \texttt{S}. The single name does not denote a quantity
at all; it denotes ``the same run of axes as wherever \texttt{S} was bound,''
of any rank. A fixed-meaning axis answers \emph{what is this axis};
a congruence group answers \emph{does this shape match that one}.}
\label{fig:sh-axis-vs-group}
\end{figure}

\begin{keyidea}
Distinguish two kinds of name in a shape bracket. A \emph{named axis of fixed
meaning} (\texttt{E}, \texttt{L}, \texttt{P}, \texttt{C}, \texttt{G};
\texttt{H}, \texttt{W}) is \emph{one} axis with a definite, permanent role---it
answers ``what is this axis.'' A \emph{named shape group}
(\texttt{(S: ...)} / \texttt{\$S}) is a \emph{congruence variable} spanning a
run of \emph{unknown} rank---it answers ``does this shape match that one.''
The first is concrete and rank-bearing; the second is abstract and
rank-agnostic. They share the look of a letter in brackets and nothing else.
\end{keyidea}

The two kinds even meet at a boundary. The flat global dof-vector is a genuine
rank-$1$ \lstinline[style=l4style]|Tensor[N]|---neither a fixed-meaning
element-local axis nor a congruence group, just a flat vector. The gather and
scatter of \cref{ch:matrixfree} are exactly the crossing between that flat
\texttt{N} world and the fixed-meaning \lstinline[style=l4style]|(E, L)| world:
\lstinline[style=l4style]|element_restrict :: ... -> Tensor[N] -> Tensor[(E, L)]|
turns one flat vector into the element-local family, and its transpose turns it
back. Three different notational species---a rank-$1$ flat axis, a congruence
group, a fixed-meaning family---live in the same pipeline, and keeping them
straight is what makes the pipeline legible.

\section{Scalar promotion: one operator, two scalar types}
\label{sec:sh-promotion}

So far we have reasoned about \emph{shapes}. There is one rule about
\emph{element types} that the shape contracts lean on, and it is worth stating
in its own right because it removes what would otherwise be a swarm of
near-duplicate operators. It is called \emph{scalar promotion}.

The setting: Palace works with both \emph{real} and \emph{complex} tensors.
The driven and eigenmode analyses (\cref{ch:targets}) produce complex fields;
the electrostatic and magnetostatic analyses are real. An operation like
\code{axpy} therefore comes in a real flavour (real scalar, real vectors) and
a complex flavour (complex scalar, complex vectors). But there is a third
combination that arises constantly in practice: a \emph{real} scalar
multiplying a \emph{complex} vector. A step size, a relaxation weight, a
quadrature factor---these are real numbers, and they routinely scale complex
fields.

Rather than mint a separate operator for that mixed case, the calculus admits
a single typing rule:
\[
  \boxed{\ \texttt{real} \;\sqsubseteq\; \texttt{complex}\ }
\]
read ``a real scalar may stand wherever a complex scalar is expected,''
promoted to a complex number with \emph{zero imaginary part}. The promotion is
\emph{exact}: zero is representable exactly, so $\alpha$ and $\alpha + 0i$ are
the very same number to the last bit, and every algebraic law of the operator
holds identically before and after promotion. \Cref{fig:sh-lattice} draws the
one-step lattice.

\begin{figure}[t]
\centering
\begin{tikzpicture}[scale=1.0]
  \node[draw=simViolet, fill=simBgViolet, rounded corners=3pt, thick,
        minimum width=20mm, minimum height=8mm, font=\small] (c) at (0,1.4)
    {\texttt{complex}};
  \node[draw=simBlue, fill=simBgBlue, rounded corners=3pt, thick,
        minimum width=20mm, minimum height=8mm, font=\small] (r) at (0,0)
    {\texttt{real}};
  \draw[-{Stealth[length=2.6mm]}, simRust, thick] (r) -- (c)
    node[midway, right=2mm, font=\scriptsize\itshape, simRust, align=left]
    {promote:\\ $\alpha \mapsto \alpha + 0i$\\ (exact)};
  \node[font=\scriptsize\itshape, simGray, left=12mm of {(0,0.7)}, text width=26mm,
        align=right] at (-2.1,0.7)
    {the lattice\\ \texttt{real} $\sqsubseteq$ \texttt{complex}};
  % the forbidden reverse, struck out
  \draw[-{Stealth[length=2.2mm]}, simGray, dotted] (c.west) to[bend right=40]
    node[midway, left=1pt, font=\scriptsize, simGray] {\,no} (r.west);
\end{tikzpicture}
\caption[The scalar-promotion lattice]{Scalar promotion is the one-step
ordering \texttt{real} $\sqsubseteq$ \texttt{complex}. A real scalar may be
\emph{promoted} upward into a complex operation by attaching a zero imaginary
part, $\alpha \mapsto \alpha + 0i$---exactly, with no rounding. The reverse
(grey, dotted) is forbidden: a complex scalar cannot be demoted into a real
operation, because its imaginary part would have nowhere to go. The promotion
collapses what would otherwise be two near-identical operators (real-scalar
and complex-scalar overloads) into one.}
\label{fig:sh-lattice}
\end{figure}

\begin{definition}[Scalar promotion]
\label{def:sh-promotion}
For a tensor operator with scalar arguments and tensor operands of element
type $T \in \{\texttt{real}, \texttt{complex}\}$:
if $T = \texttt{complex}$, a scalar argument supplied as \texttt{real} is
\emph{promoted} to \texttt{complex} by setting its imaginary part to zero, and
the operator is applied as its complex form. The promotion is exact and the
operator's algebraic laws are invariant under it. It is a property of the
\emph{type system}---a coercion at the boundary---not a separate operator.
\end{definition}

This is the difference between a \emph{typing rule} and an \emph{operator
variant}. Were the real-scalar-against-complex-vector case a distinct
operator, every algebraic law would have to be restated for it, every call
site would have to choose which overload it meant, and every lowering would
carry the real fast-path as algebraic content. Promoting instead collapses all
of that: there is one \code{axpy}, its scalar argument's type is lifted by
\texttt{real} $\sqsubseteq$ \texttt{complex} when the vectors are complex, and
the laws are written once.

\begin{workedexample}{A real coefficient times a complex field}{sh-promotion}
\textbf{Setup.} A driven-analysis update scales a complex field $\vx$ by a
real relaxation weight and adds it to a complex accumulator $\vy$:
\[
  \alpha = 0.5 \ (\texttt{real}), \qquad
  \vx = (1+2i,\; 3-i), \qquad
  \vy = (i,\; 2),
\]
both vectors \lstinline[style=l4style]|Tensor[2]| of \texttt{complex} element
type. We call \lstinline[style=l4style]|axpy alpha x y|.

\tcblower
\textbf{Shape side (the group).} The signature
\lstinline[style=l4style]|axpy :: Scalar -> Tensor[(S: ...)] -> Tensor[$S] -> Tensor[$S]|
binds $\texttt{S} \mapsto \texttt{[2]}$ from $\vx$; $\vy$ and the result are
congruent. \emph{Shape check passes}---this is the same machinery as
\cref{wex:sh-axpy}, indifferent to the element type.

\textbf{Type side (the promotion).} The vectors are \texttt{complex}, but
$\alpha = 0.5$ is \texttt{real}. By \texttt{real} $\sqsubseteq$ \texttt{complex}
it is promoted to $0.5 + 0i$, \emph{exactly}. No separate ``real-times-complex
axpy'' is invoked; the one complex \code{axpy} runs with a promoted scalar.

\textbf{Compute.}
\begin{align*}
  \alpha\vx + \vy
  &= (0.5{+}0i)\,(1{+}2i,\; 3{-}i) + (i,\; 2) \\
  &= (0.5{+}i,\; 1.5{-}0.5i) + (i,\; 2)
   = (0.5 + 2i,\; 3.5 - 0.5i).
\end{align*}
The result is \lstinline[style=l4style]|Tensor[2]| complex, congruent to
\texttt{S}, as the signature promised.
\end{workedexample}

\subsection{Where promotion does \emph{not} apply}
\label{sec:sh-promotion-limits}

The rule is deliberately narrow, and its edges matter as much as its content.

\begin{itemize}[leftmargin=1.4em]
\item \textbf{Not the reverse direction.} A \emph{complex} scalar against a
\emph{real} tensor is \emph{not} promoted---there is no demotion. The complex
scalar's imaginary part has no destination in a real result, so the lattice
runs one way only (\cref{fig:sh-lattice}, dotted arrow). This combination is a
type error, not a coercion.

\item \textbf{Not reductions whose output is already a scalar.} An operation
like \lstinline[style=l4style]|dot :: Tensor[$S] -> Tensor[$S] -> Scalar| has
no \emph{input} scalar to promote---its scalar is the \emph{output}, and the
output type is fixed by the vectors' element type (complex vectors give a
complex result). The same goes for a norm
\lstinline[style=l4style]|nrm2 :: Tensor[$S] -> Real|, whose result is real
regardless. Promotion is about an \emph{incoming} scalar argument; a reduction
has none.

\item \textbf{All-or-none across a scalar tuple.} An operator with several
scalar arguments (such as
\lstinline[style=l4style]|axpby :: Scalar -> Tensor[$S] -> Scalar -> Tensor[$S] -> Tensor[$S]|)
promotes \emph{all} of its scalars together or none. A mixture---one real, one
complex, against complex vectors---is not part of the contract.
\end{itemize}

\begin{pitfall}
Scalar promotion lifts a \texttt{real} scalar \emph{into} a complex operation
($\alpha \mapsto \alpha + 0i$); it never lifts a complex scalar into a real
one, and it never touches a reduction's \emph{output} scalar. If you find
yourself wanting to ``promote'' a complex coefficient down to a real result,
the type is genuinely wrong---there is no information-preserving way to do it,
and the calculus rejects it rather than silently discarding the imaginary part.
\end{pitfall}

\section{The shape equational theory, lightly}
\label{sec:sh-equational}

One loose end remains in \cref{def:sh-group}: congruence asks that two shapes
be ``identical axis-for-axis''---but identical in what sense? Sometimes two
shapes that look different on the page denote the same run of extents. A
reshape that merges a height axis \texttt{h} and a width axis \texttt{w} into a
single flattened axis produces an extent written \lstinline[style=l4style]|h*w|;
an operation that consumes the flattened tensor must recognise
\lstinline[style=l4style]|h*w| and \lstinline[style=l4style]|w*h| as the same
length, and \lstinline[style=l4style]|[h, w]| as reshape-compatible with
\lstinline[style=l4style]|[h*w]|. Shape matching is therefore not raw
character-by-character comparison; it is comparison \emph{modulo a small
algebra} of axis extents---the obvious arithmetic, commutative and associative
in $+$ and $\times$.

For the working purposes of this chapter, two facts suffice. First, when a
signature reuses a group (\lstinline[style=l4style]|$S| after
\lstinline[style=l4style]|(S: ...)|) or repeats a shape variable, the matcher
must (i) agree on \emph{rank}---the same number of axes---and then (ii) agree
on each axis extent \emph{up to this small algebra}, so that
\lstinline[style=l4style]|h*w| matches \lstinline[style=l4style]|w*h|. Second,
the rank check comes first: a rank mismatch fails immediately, which is exactly
why \lstinline[style=l4style]|Tensor[N]| could never match a rank-$2$ argument
in \cref{wex:sh-axpy}, regardless of any extent arithmetic. The full,
formal treatment of this shape algebra---the rewriting system that decides when
two extent expressions are equal---is deferred to \cref{app:shapes}; the
intuition above is all the running text needs.

\begin{notationbox}
A practical consequence: prefer to write a shared shape as a \emph{group} or a
\emph{variable}, not as an algebraic coincidence. If two operands must agree,
say so with \lstinline[style=l4style]|(S: ...)| / \lstinline[style=l4style]|$S|
and let the matcher enforce it. Reserve the extent algebra
(\lstinline[style=l4style]|h*w| and friends) for the places where a genuine
reshape or factorisation has \emph{produced} a composite extent. Congruence is
a contract you assert; the extent algebra is a fact the matcher verifies.
\end{notationbox}

\summarysection
A tensor's \emph{rank} is the number of axes in its shape bracket, and the
naive notation \lstinline[style=l4style]|Tensor[N]| pins that rank to one: it
names a single axis, so reusing it across operands forces them to agree
\emph{at rank one}, silently outlawing the matrix and higher-rank cases that
elementwise operations handle perfectly well. The cure is the \emph{named
shape group}: \lstinline[style=l4style]|(S: ...)| binds a name to a run of axes
of \emph{unknown} rank, and a later use \lstinline[style=l4style]|$S| (the
\texttt{\$} sigil marking a back-reference) asserts a second shape is
\emph{congruent} to it. This gives \emph{rank-agnostic congruence}: one
signature, \lstinline[style=l4style]|axpy :: Scalar -> Tensor[(S: ...)] -> Tensor[$S] -> Tensor[$S]|,
correct on vectors, matrices, and every higher tensor at once. Operators whose
domain and range shapes differ name two groups \emph{range-first},
\lstinline[style=l4style]|LinOp[(R: ...), (D: ...)]|, collapsing to the square
form \lstinline[style=l4style]|LinOp[(S: ...), $S]| for endomorphisms.
Congruence groups are \emph{variables} and must be kept distinct from
\emph{named axes of fixed meaning}---the \texttt{E}/\texttt{L}/\texttt{P}/%
\texttt{C}/\texttt{G} family of the matrix-free kernel---whose letters denote
definite, rank-bearing quantities. Finally, \emph{scalar promotion}
(\texttt{real} $\sqsubseteq$ \texttt{complex}) is a typing rule, not an
operator variant: a real scalar joins a complex operation by gaining a zero
imaginary part, exactly, collapsing the real/complex overloads into one
operator---but never in reverse, and never on a reduction's output scalar.

\furtherreading
The view of a matrix as a map between shaped spaces, rather than a grid of
numbers, is the organising idea of Trefethen and
Bau~\citep{trefethenbau1997}; their opening lectures on matrix--vector
products and the four fundamental subspaces are the linear-algebra backdrop for
the range-first operator-shape convention of \cref{sec:sh-operators}. The
named-axis style of shape annotation---carrying an axis's \emph{meaning} in its
name---is folklore in the array-programming and tensor-compiler communities;
the matrix-free element-local family of \cref{sec:sh-fixedaxes} is a concrete
instance, developed in full in \cref{ch:matrixfree}. The formal shape algebra
underlying \cref{sec:sh-equational} is collected in \cref{app:shapes}.

\exercisessection
\begin{enumerate}[nosep]
  \item \textbf{Rank counting.} For each shape, state the rank and say whether
  it could be the binding of a group \lstinline[style=l4style]|(S: ...)| in a
  call where the other operand is a length-$4$ vector:
  \lstinline[style=l4style]|[4]|, \lstinline[style=l4style]|[2, 2]|,
  \lstinline[style=l4style]|[4, 1]|, \lstinline[style=l4style]|[]|.
  (Recall congruence requires agreement on rank \emph{first}.)

  \item \textbf{The wrong signature, again.} A colleague writes a scaling
  operator as \lstinline[style=l4style]|scale :: Scalar -> Tensor[M] -> Tensor[M]|
  and reports it ``works fine on all our test vectors.'' Construct the smallest
  input on which it wrongly fails, and rewrite the signature so it does not.

  \item \textbf{Partial congruence.} Write a signature for an operation that
  takes two tensors sharing a leading run of axes but each carrying its own
  \emph{single} trailing axis (lengths \texttt{p} and \texttt{q} respectively),
  and returns a tensor with the shared run followed by \emph{both} trailing
  axes. (Hint: \lstinline[style=l4style]|(S: ...)| with two free trailing
  dims.)

  \item \textbf{Range-first.} An operator maps a rank-$2$ field of shape
  \lstinline[style=l4style]|[a, b]| to a rank-$1$ vector of shape
  \lstinline[style=l4style]|[c]|. Write its
  \lstinline[style=l4style]|LinOp[(R: ...), (D: ...)]| type and the type of its
  \lstinline[style=l4style]|apply_linop|. Which group is which?

  \item \textbf{Square or not.} Decide for each whether the operator is an
  endomorphism (deserving \lstinline[style=l4style]|LinOp[(S: ...), $S]|) or a
  general map (deserving two groups): (a) a preconditioner $\mat{M}^{-1}$;
  (b) the gather \lstinline[style=l4style]|element_restrict| of
  \cref{ch:matrixfree}; (c) the time-stepping operator of a transient solve;
  (d) the basis evaluation $B_{\diffop}$ mapping
  \lstinline[style=l4style]|[E, L]| to \lstinline[style=l4style]|[E, P, C]|.

  \item \textbf{Axis or group?} For each letter as it appears in the named
  shape, classify it as a \emph{fixed-meaning axis} or a \emph{congruence
  group}, and justify: \texttt{E} in \lstinline[style=l4style]|Tensor[(E, L)]|;
  \texttt{S} in \lstinline[style=l4style]|Tensor[$S]|; \texttt{H} in
  \lstinline[style=l4style]|Tensor[H, W, VY=3, VX=3]|; \texttt{D} in
  \lstinline[style=l4style]|LinOp[(R: ...), (D: ...)]|.

  \item \textbf{Promotion or error?} For each call against \texttt{complex}
  vectors $\vx,\vy$ unless stated, say whether scalar promotion applies, does
  nothing, or is a type error: (a) \lstinline[style=l4style]|axpy| with a real
  $\alpha$; (b) \lstinline[style=l4style]|axpy| with a complex $\alpha$ against
  \emph{real} $\vx,\vy$; (c) \lstinline[style=l4style]|dot x y| (real
  $\vx,\vy$); (d) \lstinline[style=l4style]|nrm2 x|; (e)
  \lstinline[style=l4style]|axpby| with one real and one complex scalar.

  \item \textbf{Exactness.} Explain why the promotion $\alpha \mapsto \alpha +
  0i$ introduces \emph{no} rounding error, and why this exactness is what lets
  the operator's algebraic laws be stated once rather than twice. Contrast with
  a hypothetical \texttt{float32} $\sqsubseteq$ \texttt{float64} promotion: is
  \emph{that} one exact in both directions? Which direction, if either, is safe?
\end{enumerate}

\chapter{The Fold: \texttt{iterate\_while} and the Iteration Combinators}
\label{ch:fold}

\chapterorient{This is the keystone of the framework. We met the \emph{fold} once
already, in \cref{ch:lifecycle}, as the shape of the adaptive loop; we promised
then that the same shape would turn out to be the engine of time-stepping, of
Krylov solvers, and of eigenvalue iterations. Now we make good on that promise.
We start from the \texttt{for}-loop every reader knows and recast it, step by
step, as a value carried forward through a single combinator named
\texttt{iterate\_while}. We develop its meaning fully---its \texttt{while}
convention, the one rule about its predicate that everything depends on, its two
sister combinators, and the demand-driven pruning that is the calculus's
signature move. By the end you will be able to read any iterative algorithm in
this book as one specialization of the one combinator taught here.}

\section{From the loop you know to the fold you need}

Every reader of this book has written a \texttt{for}-loop. You set a variable to
a starting value, and then, over and over, you change it: read its current value,
do some work, store a new value back. When a condition fails, you stop and report
what the variable holds. A running sum is the canonical example. In the
imperative style you grew up with, summing a list looks like this:
\begin{lstlisting}[style=l4style]
total = 0
for x in xs:
    total = total + x      -- mutate `total` in place, each pass
return total
\end{lstlisting}
The variable \code{total} is \emph{mutated}: the same storage cell is overwritten
on every pass. \Cref{ch:bigidea} argued at length that mutation is exactly the
property we must give up to run a computation on an immutable-tensor backend. So
we are obliged to find another way to say ``the same thing'' without the
overwrite---and the way we find will turn out to be more powerful than the loop
it replaces.

The trick is to stop thinking of \code{total} as a cell that changes and start
thinking of it as a \emph{value carried forward}. Each pass does not modify a
variable; it \emph{computes a new value from the old one}. ``Carry a value;
take a step that produces the next value; repeat.'' The carried value here is the
running total; the step is ``add the next element''; the repetition walks the
list. Nothing is overwritten---each step's output is a fresh value, and the old
one is simply not referred to again.

This reframing has a name in functional programming: it is a \emph{fold}. A fold
takes a starting value (the \emph{seed}), a step that combines the seed-so-far
with the next thing, and a structure to walk; it threads the accumulated value
through the structure and returns the final accumulation. The running sum is
\code{foldl (+) 0 xs}---``fold the list \code{xs} from the left, starting at
\code{0}, combining with \code{+}.'' \Cref{fig:loopfold} sets the imperative loop
and the fold side by side: the same computation, the mutation traded for a carry.

\begin{figure}[t]
\centering
\begin{tikzpicture}[font=\footnotesize]
  % --- left: the imperative loop, a cell overwritten ---
  \begin{scope}[xshift=0cm]
    \node[font=\small\bfseries, text=simRust] at (0,2.5) {imperative: one cell, overwritten};
    \node[data, minimum width=14mm] (cell) at (0,1.4) {\code{total}};
    \draw[backflow] (cell.east) .. controls (2.2,1.4) and (2.2,2.2) .. (cell.north)
      node[pos=0.5, right=10mm, font=\scriptsize, text=simRust] {$\gets$ overwrite};
    \node[opbox] (add) at (0,0.0) {\code{total = total + x}};
    \draw[flow] (cell.south) -- (add.north);
    \draw[flow] (add.south) -- ++(0,-0.5)
      node[below, font=\scriptsize] {next pass writes \code{total} again};
  \end{scope}
  % divider
  \draw[simGray!50, thick] (4.0,-0.9) -- (4.0,2.7);
  % --- right: the fold, a value threaded forward ---
  \begin{scope}[xshift=8cm]
    \node[font=\small\bfseries, text=simBlue] at (0,2.5) {fold: a value carried forward};
    \node[data] (t0) at (-2.4,1.0) {$0$};
    \node[opbox] (s1) at (-1.0,1.0) {step};
    \node[data] (t1) at (0.2,1.0) {$a_1$};
    \node[opbox] (s2) at (1.4,1.0) {step};
    \node[data] (t2) at (2.6,1.0) {$a_2$};
    \draw[flow] (t0)--(s1); \draw[flow] (s1)--(t1);
    \draw[flow] (t1)--(s2); \draw[flow] (s2)--(t2);
    \draw[flow, simGray] (t2) -- ++(0.9,0) node[right, font=\scriptsize] {$\cdots$};
    \node[font=\scriptsize, text=simGreen] at (-1.0,0.3) {$x_1$};
    \node[font=\scriptsize, text=simGreen] at (1.4,0.3) {$x_2$};
    \draw[flow, simGreen, shorten >=1pt] (-1.0,0.45)--(s1.south);
    \draw[flow, simGreen, shorten >=1pt] (1.4,0.45)--(s2.south);
  \end{scope}
\end{tikzpicture}
\caption[The loop-to-fold bridge]{The same running sum, two ways. On the left,
the imperative loop keeps one storage cell \code{total} and \emph{overwrites} it
each pass---the dashed back-edge is the mutation we must give up. On the right,
the fold threads a \emph{value} forward: each step consumes the carried value
($0$, then $a_1$, then $a_2,\dots$) and the next input $x_i$ and produces a fresh
value. Nothing is overwritten; each arrow carries a distinct immutable value.
This carry-threading picture is the whole chapter in miniature.}
\label{fig:loopfold}
\end{figure}

\subsection{Three folds you already know}

Before we build the simulator's combinator, it pays to fix three patterns from
the functional vocabulary of \cref{ch:language}. Each is a way of walking a
structure; the differences between them are exactly the differences that will
matter for iteration.

\begin{description}[style=nextline, leftmargin=0pt, font=\normalfont\itshape]
  \item[\code{map}] applies a function to each element \emph{independently} and
  collects the results: \lstinline[style=l4style]|map f [x1,x2,x3] = [f x1, f x2, f x3]|.
  The crucial word is \emph{independently}: \code{f x2} does not depend on
  \code{f x1}, so the three applications could be done in any order, or all at
  once. A \code{map} has no carry.
  \item[\code{foldl}] threads an accumulator \emph{left to right}, each step
  depending on the one before:
  \lstinline[style=l4style]|foldl (+) 0 [x1,x2,x3] = ((0 + x1) + x2) + x3|.
  Here order \emph{does} matter: the accumulator entering step~2 is exactly what
  step~1 produced. This dependency is the essence of a fold.
  \item[\code{foldr}] threads from the \emph{right}:
  \lstinline[style=l4style]|foldr (+) 0 [x1,x2,x3] = x1 + (x2 + (x3 + 0))|. For an
  associative operator like \code{+} the two folds agree, but in general they
  bracket the work differently. We will rarely need \code{foldr}; we note it for
  completeness and because it makes the directionality of \code{foldl} vivid.
\end{description}

The distinction between \code{map} and \code{foldl} is the one to carry forward.
A \code{map} is a family of \emph{independent} computations; a \code{foldl} is a
\emph{chain} in which each link consumes the previous. We met exactly this
distinction in \cref{ch:lifecycle}: the electrostatic analysis's several solves
against independent right-hand sides are a \code{map}, while the adaptive loop,
in which each refined mesh comes from the last, is a fold. \Cref{ch:situating}
named it one of the four ``solve corners.'' Iteration---running a step until it
converges---is a fold whose structure is not a fixed list but ``keep going until
a condition holds.''

\begin{keyidea}
A \emph{fold} carries a value forward through a sequence of steps, each step
consuming the value the previous step produced. It is the structure-with-a-carry,
as opposed to \code{map}, the structure-without. Mutation (``overwrite a cell'')
becomes value-threading (``compute the next value''). Every iterative algorithm
in this book is a fold; the rest of the chapter builds the one combinator that
expresses all of them.
\end{keyidea}

\subsection{Folds that compute, not just accumulate}

Summing a list is a fold over a structure that is already laid out---the list
tells us how many steps to take. The folds a simulator runs are different in one
way: \emph{the number of steps is not known in advance}. Newton's method for a
square root, fixed-point iteration, a Krylov solve, a time march---each takes a
step and then \emph{asks whether to take another}, based on the value it now
holds. The structure being folded is not a list of inputs; it is the open-ended
``again?'' of a convergence test.

Two familiar examples make the shape concrete, and we will return to both as
worked examples below.

\emph{Newton's method.} To find $\sqrt{2}$, solve $x^2 - 2 = 0$ by Newton
iteration: from a guess $x$, the next guess is
\[
  x' \;=\; x - \frac{x^2 - 2}{2x} \;=\; \tfrac{1}{2}\!\left(x + \tfrac{2}{x}\right).
\]
We carry $x$, take that step, and repeat until $x$ stops changing (until the
\emph{residual} $|x^2 - 2|$ is small). Carry, step, test, repeat: a fold.

\emph{Fixed-point iteration.} To solve $x = g(x)$ for a contraction $g$, simply
iterate $x' = g(x)$ from any starting point; under a contraction the iterates
converge to the unique fixed point. Again: carry $x$, step by applying $g$, stop
when $x$ barely moves. A fold.

Both share the lifecycle's shape from \cref{sec:fold}: a starting value, a step
that carries state forward, a stopping condition over the accumulated state, a
result threaded through. What they lack---and what a simulator needs---is a
\emph{single named operator} that captures this shape once so that every
algorithm can be written as an instance of it. That operator is
\code{iterate\_while}.

\section{\texorpdfstring{\texttt{iterate\_while}}{iterate\_while}: the canonical primitive}
\label{sec:iteratewhile}

Here is the combinator the whole book turns on. We give its signature first, then
read it slowly, then draw it.

\begin{lstlisting}[style=l4style]
iterate_while
  :: a                                     -- the initial carry
  -> (a -> Bool)                           -- the predicate: keep going?
  -> (a -> { state: a, ...extras })        -- the step: next carry + per-step extras
  -> { final_state: a, trajectory: [{ ...extras }] }
\end{lstlisting}

Read it argument by argument. The first argument is an \emph{initial carry}, a
value of some type \code{a}---it might be a number, a tensor, or a whole record
of solver state. The second is a \emph{predicate}, a pure function from the carry
to a Boolean, answering ``should we take another step?'' The third is the
\emph{step}: from the current carry it produces a record holding the
\code{state}---the \emph{next} carry---together with any per-step
\emph{extras} it wishes to report. The result is a record with two fields: the
\code{final\_state} the carry settled on, and a \code{trajectory}, the list of
all the per-step extras records in order.

\begin{notationbox}
The \lstinline[style=l4style]|{ state: a, ...extras }| notation is a record
(\cref{ch:records}) whose \code{state} field has type \code{a} and which is
\emph{open} in its remaining fields: \code{...extras} stands for ``and whatever
other fields this particular step chooses to emit.'' One step might emit
\lstinline[style=l4style]|{ state: x', residual_norm: r }|; another might emit
\lstinline[style=l4style]|{ state: x' }| with no extras at all. The combinator
does not care which: it forwards the \code{state} as the next carry and piles the
rest into the \code{trajectory}.
\end{notationbox}

In words: \emph{from an initial carry, while a predicate holds of the current
carry, take a step that produces the next carry plus per-step extras, collecting
those extras into a trajectory.} That single sentence is the meaning of
\code{iterate\_while}, and \cref{fig:threaded} draws it. The carry flows left to
right along the top; before each step a predicate gate decides whether to
proceed; the step emits a fresh carry forward and drops its extras downward into
the growing trajectory.

\begin{figure}[t]
\centering
\begin{tikzpicture}[font=\footnotesize, node distance=6mm]
  % carry spine
  \node[data, minimum width=10mm] (a0) at (0,0) {$a_0$};
  \node[diamond, draw=simBlue, fill=simBgBlue, aspect=1.8, inner sep=1pt,
        font=\scriptsize, text=simInk] (p0) at (1.8,0) {$p(a_0)$?};
  \node[opbox, minimum width=12mm] (f0) at (3.9,0) {step};
  \node[data, minimum width=10mm] (a1) at (5.9,0) {$a_1$};
  \node[diamond, draw=simBlue, fill=simBgBlue, aspect=1.8, inner sep=1pt,
        font=\scriptsize, text=simInk] (p1) at (7.7,0) {$p(a_1)$?};
  \node[opbox, minimum width=12mm] (f1) at (9.8,0) {step};
  \node[font=\scriptsize, simGray] (more) at (11.6,0) {$\cdots$};
  % forward carry flow
  \draw[flow] (a0)--(p0);
  \draw[flow] (p0)-- node[above,font=\scriptsize]{yes} (f0);
  \draw[flow] (f0)-- node[above,font=\scriptsize]{\code{.state}} (a1);
  \draw[flow] (a1)--(p1);
  \draw[flow] (p1)-- node[above,font=\scriptsize]{yes} (f1);
  \draw[flow] (f1)-- node[above,font=\scriptsize]{\code{.state}} (more);
  % `no` exits to final_state
  \node[product] (fin) at (7.7,-2.4) {\code{final\_state}};
  \draw[backflow] (p0.south) |- (fin.west);
  \draw[backflow] (p1.south) -- (fin.north);
  % extras drop into trajectory
  \node[data, fill=simBgGold, draw=simGold, minimum width=44mm] (traj) at (6.85,-1.3)
    {\code{trajectory}: $[\,\{\dots e_0\},\ \{\dots e_1\},\ \dots\,]$};
  \draw[refedge] (f0.south) -- (f0.south|-traj.north);
  \draw[refedge] (f1.south) -- (f1.south|-traj.north);
  \node[font=\scriptsize, text=simGold!80!black] at (3.9,-0.85) {extras $e_0$};
  \node[font=\scriptsize, text=simGold!80!black] at (9.8,-0.85) {extras $e_1$};
\end{tikzpicture}
\caption[The threaded carry of \texttt{iterate\_while}]{\texttt{iterate\_while}
as a threaded carry. The carry $a_0, a_1, \dots$ flows forward along the top.
Before each step the predicate gate $p(\cdot)$ tests the \emph{current carry};
on ``yes'' the step runs, emitting the next carry through its \code{.state}
field and dropping its per-step \emph{extras} downward into the growing
\code{trajectory}. The first ``no'' exits the loop, surfacing the current carry
as \code{final\_state}. The predicate fires \emph{first}, before any step---this
is a \texttt{while} loop, not a \texttt{do}/\texttt{while}.}
\label{fig:threaded}
\end{figure}

\subsection{The small-step rule}

We can make the meaning of \code{iterate\_while} exact in one recursive rule,
taken verbatim from the calculus's semantics (\cref{ch:language}). Given a carry
$a$, a predicate $p$, and a step $f$:
\[
\begin{aligned}
\textsf{iterate\_while}\ a\ p\ f \;\to\;\ &\textsf{if}\ p(a)\\
  &\quad\textsf{then}\ \textsf{let}\ \{\textsf{state}\!:\!a',\ \dots e\} = f(a)\ \textsf{in}\\
  &\qquad \textsf{let}\ \{\textsf{final\_state},\ \textsf{trajectory}\} = \textsf{iterate\_while}\ a'\ p\ f\ \textsf{in}\\
  &\qquad \{\textsf{final\_state},\ \textsf{trajectory}\!:\ [\{\dots e\}] \mathbin{+\!\!+} \textsf{trajectory}\}\\
  &\quad\textsf{else}\ \{\textsf{final\_state}\!:\ a,\ \textsf{trajectory}\!:\ [\,]\,\}
\end{aligned}
\]
Read it in prose. Test $p(a)$. If \emph{false}, stop: the answer is $a$ as
\code{final\_state} with an empty trajectory. If \emph{true}, run one step
$f(a)$, obtaining the next carry $a'$ and this step's extras $e$; recurse on
$a'$; and prepend $e$ to whatever trajectory the recursion returns. The recursion
sits in tail position, so an implementation realizes it as an ordinary loop, not
a stack of deferred calls. This is the entire operational content of the
combinator. Everything else in the chapter is a consequence of, or a discipline
around, this one rule.

\subsection{Unrolling the recursion}

The recursive rule is exact, but its shape is easiest to see when we
\emph{unroll} it---write out what it does for a concrete run of three steps
before the predicate finally fails. \Cref{fig:unroll} does exactly that. The
carry $a_0$ passes the predicate and is fed to the first step, producing $a_1$
and extras $e_1$; $a_1$ passes and yields $a_2, e_2$; $a_2$ passes and yields
$a_3, e_3$; and then the predicate fails on $a_3$, so the loop stops. The carry
has walked $a_0 \to a_1 \to a_2 \to a_3$, each link a single step, and the three
extras records have lined up into the trajectory
$[\,e_1, e_2, e_3\,]$ in iteration order.

\begin{figure}[t]
\centering
\begin{tikzpicture}[font=\footnotesize, node distance=8mm]
  % carry chain
  \foreach \i/\lbl in {0/{a_0},1/{a_1},2/{a_2},3/{a_3}}{
    \node[data, minimum width=8mm] (c\i) at (\i*3.0,0) {$\lbl$};
  }
  \foreach \i/\j in {0/1,1/2,2/3}{
    \node[opbox, minimum width=13mm] (st\j) at (\i*3.0+1.5,0) {step};
    \draw[flow] (c\i)--(st\j);
    \draw[flow] (st\j)-- node[above,font=\scriptsize]{\code{.state}} (c\j);
  }
  % predicate annotations above the carries that pass
  \foreach \i in {0,1,2}{
    \node[font=\scriptsize, text=simGreen] at (\i*3.0,0.78) {$p$: yes};
  }
  \node[font=\scriptsize, text=simRust] at (9.0,0.78) {$p$: no\,(stop)};
  % extras drop down to trajectory
  \node[data, fill=simBgGold, draw=simGold, minimum width=58mm] (traj) at (4.5,-1.7)
    {\code{trajectory} $=[\,e_1,\ e_2,\ e_3\,]$};
  \foreach \j/\e in {1/{e_1},2/{e_2},3/{e_3}}{
    \draw[refedge, simGold!80!black] (st\j.south) -- (st\j.south|-traj.north);
    \node[font=\scriptsize, text=simGold!80!black] at ($(st\j.south)+(0,-0.45)$) {$\e$};
  }
  % final_state
  \node[product, minimum width=18mm] (fs) at (9.0,-1.7) {\code{final\_state}\,$=a_3$};
  \draw[flow] (c3.south) -- (fs.north);
\end{tikzpicture}
\caption[A fold unrolled over three steps]{The recursion of
\cref{sec:iteratewhile} unrolled for a run of three steps. The carry threads
$a_0 \to a_1 \to a_2 \to a_3$, each arrow one step. The predicate passes
(\textcolor{simGreen}{yes}) on $a_0, a_1, a_2$ and fails
(\textcolor{simRust}{no}) on $a_3$, ending the loop. Each step drops its
extras downward; the three records line up, in order, as the
\code{trajectory} $[\,e_1, e_2, e_3\,]$, and the final carry $a_3$ becomes
\code{final\_state}. Compare \cref{fig:loopfold}: the carry spine is the same
forward-threaded value, now with extras peeling off below.}
\label{fig:unroll}
\end{figure}

This unrolled picture is the one to hold in mind for the rest of the book.
Whenever a later chapter says ``the solver folds \code{krylov\_step} with
\code{iterate\_while},'' it means precisely \cref{fig:unroll} with ``step''
replaced by the solver's per-step kernel: a carry of solver state threading
forward, a predicate gating each pass, and a residual history peeling off into
the trajectory.

\subsection{What the signature makes structural}
\label{sec:structural}

It is worth pausing to ask why this particular signature, and not some looser
one. An ordinary \texttt{for}-loop lets the body do anything: read any variable,
write any variable, test any condition. The combinator's signature deliberately
forbids most of that, and each thing it forbids buys a property we will rely on.
Four are worth naming.

\begin{description}[style=nextline, leftmargin=0pt, font=\normalfont\itshape]
  \item[The predicate sees the carry only.] Its type
  \lstinline[style=l4style]|a -> Bool| mentions neither the extras nor any
  ambient state. This forces every termination quantity into the carry
  (\cref{sec:predicaterule})---a discipline, but one that makes ``what does this
  loop stop on?'' answerable by reading the carry's type alone.
  \item[The carry is value-threaded; the extras are record-spread.] There is no
  accumulator living \emph{outside} the carry. Anything that must persist across
  iterations is in the carry; anything that is a per-step output is in the extras
  and ends up in the trajectory. The two roles cannot be confused because they
  occupy different slots of the step's return record.
  \item[Per-step output is a declared field, not a side effect.] A step does not
  ``log'' a residual or ``print'' a metric; it \emph{returns} one, as a field of
  its output record. This is what makes the demand-driven pruning of
  \cref{sec:pruning} possible---an output that is a value can be observed or not
  observed, and pruned when not.
  \item[The trajectory is observed structurally.] The list of extras is an
  ordinary value the consumer may read or ignore. Because it is a value and not an
  effect, the calculus can see, at the call site, whether anyone reads it---and
  eliminate the work that produces it when no one does.
\end{description}

None of these is enforced by comments or conventions you must remember; each is a
consequence of the \emph{type}. The signature is doing the work that, in an
imperative loop, the programmer's discipline would have to do---and doing it in a
form a compiler can check and a backend can exploit.

\section{The \texttt{while} convention and the predicate rule}
\label{sec:predicaterule}

Two facts about the predicate are load-bearing, and getting them wrong is the
most common way to misuse the combinator. We state them as rules and then earn
them.

\paragraph{Rule 1: the predicate fires first.} Look again at the small-step
rule. The very first thing it does is test $p(a)$---\emph{before} any step runs.
If the predicate is already false on the initial carry, the body never executes
and the result is the initial carry with an empty trajectory. This is the
\texttt{while} convention (test, then maybe step), not the
\texttt{do}/\texttt{while} convention (step, then test). The distinction is not
pedantry: a solver handed an already-converged input should do zero work, and the
\texttt{while} convention guarantees it.

\paragraph{Rule 2: the predicate is pure on the carry.} The predicate's type is
\lstinline[style=l4style]|a -> Bool|. Its \emph{only} input is the carry. It
cannot read external mutable state, it cannot read the simulator's monadic
environment (\cref{ch:monad}), and---this is the subtle one---it cannot read the
\emph{extras} the step produces. It sees the carry, and nothing else.

\begin{figure}[t]
\centering
\begin{tikzpicture}[font=\footnotesize]
  % --- LEFT: the correct rule ---
  \begin{scope}
    \node[font=\small\bfseries, text=simGreen] at (0,2.4) {correct: predicate reads the carry};
    \node[data, minimum width=22mm, align=center] (carry) at (0,1.3)
      {\code{carry}\\[-1pt]\scriptsize\itshape (incl.\ residual, flag)};
    \node[diamond, draw=simBlue, fill=simBgBlue, aspect=1.6, inner sep=1pt,
          font=\scriptsize] (pred) at (0,-0.2) {\code{p(carry)}};
    \draw[flow] (carry)--(pred);
    \node[font=\scriptsize, text=simGreen] at (2.0,0.55) {one input,};
    \node[font=\scriptsize, text=simGreen] at (2.0,0.2) {no cycle};
  \end{scope}
  \draw[simGray!50, thick] (3.6,-1.2) -- (3.6,2.6);
  % --- RIGHT: the pitfall, crossed out ---
  \begin{scope}[xshift=8.2cm]
    \node[font=\small\bfseries, text=simRust] at (0,2.4) {pitfall: predicate reads the extras};
    \node[opbox, minimum width=16mm] (step) at (-1.3,1.3) {step};
    \node[data, fill=simBgGold, draw=simGold, minimum width=16mm] (ex) at (1.3,1.3) {\code{extras}};
    \node[diamond, draw=simRust, fill=simBgRust, aspect=1.6, inner sep=1pt,
          font=\scriptsize] (badp) at (-1.3,-0.4) {\code{p(extras)}};
    \draw[flow] (step.east)-- node[above,font=\scriptsize]{$e$} (ex.west);
    \draw[flow, simRust] (ex.south) .. controls (1.3,-0.4) and (0.2,-0.4) .. (badp.east);
    \draw[flow, simRust, densely dashed] (badp.north) .. controls (-1.3,0.6) and (-1.3,0.7) .. (step.south)
      node[pos=0.5, left, font=\scriptsize, text=simRust] {gates};
    % the big X
    \draw[simRust, very thick] (-2.4,-1.0) -- (2.4,1.9);
    \draw[simRust, very thick] (-2.4,1.9) -- (2.4,-1.0);
    \node[font=\scriptsize, text=simRust, align=center] at (0,-1.55)
      {circular: the test needs an extra\\the step has not produced yet};
  \end{scope}
\end{tikzpicture}
\caption[The predicate-on-carry rule]{Left: the rule. The predicate reads the
\emph{carry}, which already contains everything termination depends on (a
residual, a converged flag, an iteration count). One input, no cycle. Right: the
pitfall (\cref{sec:predicaterule}). Letting the predicate read the step's
\emph{extras} creates a circular dependency---the gate that decides whether to
run the step needs a value only the step can produce. The fix is to fold any
termination quantity into the carry, so the predicate reads it from there.}
\label{fig:predrule}
\end{figure}

Why must the predicate be pure on the carry? Because termination must be decided
\emph{before} the step runs (Rule~1), and the extras are produced \emph{by} the
step. If the predicate were allowed to read the extras, then on the first
iteration there would be no extras to read---the step that produces them has not
run. That is a circular dependency, drawn crossed-out on the right of
\cref{fig:predrule}. The calculus forbids it structurally: the predicate's type
mentions only \code{a}, so there is simply no way to write a predicate that
reads the extras.

\paragraph{The consequence: fold termination quantities into the carry.} This is
not a limitation; it is a discipline that clarifies every algorithm. Whatever the
loop terminates on---a residual norm below tolerance, a ``converged'' flag, an
iteration count reaching a cap---must live \emph{in the carry}, computed by the
step and read by the predicate on the next pass. The extras hold only the
quantities that are \emph{outputs} (things a monitor or a plot might want); the
carry holds the quantities that are \emph{control inputs} (things the predicate
needs).

A solver that stops when its residual drops below \code{tol} does \emph{not}
write a predicate that inspects ``the residual the step just produced.'' It
carries a \code{residual} field in its state, has the step compute and store it
there, and writes the predicate
\lstinline[style=l4style]|\s -> s.residual > tol|. A solver capped at
\code{max\_it} iterations carries an \code{it} counter in the state and writes
\lstinline[style=l4style]|\s -> s.it < max_it && not s.converged|. The carry is
where convergence lives.

\begin{pitfall}
The single most common mistake with \code{iterate\_while} is trying to let the
predicate test a quantity the step just produced---``loop while the new residual
exceeds the tolerance,'' reaching for the residual as an \emph{extra}. This is a
\emph{circular dependency}: the predicate fires \emph{before} the step
(Rule~1), so on the first pass the extra does not exist yet. The calculus rules
it out: the predicate has type \lstinline[style=l4style]|a -> Bool| and can read
only the carry. \textbf{The fix is always the same}---fold the termination
quantity into the carry. Let the step compute the residual and store it in a
carry field; let the predicate read \emph{that} field. The residual then enters
the next iteration's predicate as part of the carry, exactly when it exists.
\Cref{wex:pitfall} shows the wrong and right versions side by side.
\end{pitfall}

\section{The family: three combinators, one shape}
\label{sec:family}

The combinator of \cref{sec:iteratewhile} is the central member of a small
family. The other two are specializations: one drops a feature it does not need,
one adds a feature for a particular recurrence. \Cref{fig:family} lines all three
up; we take them in turn.

\begin{figure}[t]
\centering
\begin{tikzpicture}[font=\footnotesize, node distance=3mm]
  \def\colw{43mm}
  % --- pure ---
  \node[stage, minimum width=\colw, align=center] (pureH) at (0,0)
    {\code{iterate\_while\_pure}};
  \node[opbox, minimum width=\colw, align=left, below=of pureH] (pureB)
    {carry only,\\no extras\\$\Rightarrow$ returns \code{final\_state}};
  \node[cornertag, below=1mm of pureB] {plain loop};
  % --- base ---
  \node[stage, minimum width=\colw, align=center, right=14mm of pureH] (baseH)
    {\code{iterate\_while}};
  \node[opbox, minimum width=\colw, align=left, below=of baseH] (baseB)
    {carry $+$ extras\\$\Rightarrow$ \code{final\_state}\\$+$ \code{trajectory}};
  \node[cornertag, below=1mm of baseB] {the canonical combinator};
  % --- with_prev ---
  \node[stage, minimum width=\colw, align=center, right=14mm of baseH] (prevH)
    {\code{iterate\_while\_with\_prev}};
  \node[opbox, minimum width=\colw, align=left, below=of prevH] (prevB)
    {carry $+$ extras\\$+$ a threaded \code{prev}\\(two-term recurrence)};
  \node[cornertag, below=1mm of prevB] {bootstrap then steady};
  % arrows showing specialization
  \draw[depedge] (baseH) -- node[above,font=\scriptsize]{$e=()$} (pureH);
  \draw[depedge] (prevH) -- node[above,font=\scriptsize]{\code{prev}$=()$} (baseH);
\end{tikzpicture}
\caption[The iteration-combinator family]{The three members of the iteration
family. The center is \code{iterate\_while}: a carry, a predicate, a step with
extras. Drop the extras ($e=()$) and the trajectory is always empty, leaving
\code{iterate\_while\_pure}, a plain loop returning just the final carry. Add a
second threaded value, the previous iterate, and you get
\code{iterate\_while\_with\_prev}, the driver for two-term recurrences; collapse
its \code{prev} to nothing and it degenerates back to the center. One shape,
three settings.}
\label{fig:family}
\end{figure}

\subsection{\texttt{iterate\_while\_pure}: the plain loop}

Many loops have no per-step output to report---they only thread a carry forward
and want the final value. For these the extras are empty, the trajectory is
always \code{[]}, and carrying it around is noise. The \emph{sugar}
\code{iterate\_while\_pure} drops it:
\begin{lstlisting}[style=l4style]
iterate_while_pure :: a -> (a -> Bool) -> (a -> a) -> a
\end{lstlisting}
The step is now just \lstinline[style=l4style]|a -> a| (next carry, no record),
and the result is just the final carry, no trajectory record to destructure. It
is defined in terms of the general combinator---it \emph{is} the general
combinator with an empty extras record, projected onto its \code{final\_state}:
\begin{lstlisting}[style=l4style]
iterate_while_pure a p f
  = (iterate_while a p (\x -> { state: f x })).final_state
\end{lstlisting}
Newton's method, fixed-point iteration, and a fixed-degree polynomial smoother
all fit \code{iterate\_while\_pure}: each threads a single carry and wants only
the answer. We use it whenever there is nothing per-step worth keeping.

\subsection{\texttt{iterate\_while\_with\_prev}: two-term recurrences}

Some iterations need not only the current carry but the \emph{previous} one. The
conjugate-gradient method (\cref{ch:cg}) builds each search direction from the
current residual \emph{and} a coefficient computed from the previous residual;
a three-term Chebyshev recurrence needs the iterate from two steps back. These
are \emph{two-term recurrences}, and they have an awkward first step: there is no
``previous'' on iteration zero, so the naive code carries a special
\lstinline[style=l4style]|if it == 0| branch inside the step.

The framework removes that branch by splitting the loop into a one-time
\emph{bootstrap} that manufactures the first ``previous'' value, followed by a
\emph{steady} step that is branch-free because it always has a real previous to
work with. The combinator that drives this pattern is
\code{iterate\_while\_with\_prev}:
\begin{lstlisting}[style=l4style]
iterate_while_with_prev
  :: (a -> { state: a, prev: b, ...extras })       -- bootstrap_step (runs once)
  -> a                                             -- initial carry
  -> ((a, b) -> { state: a, prev: b, ...extras })  -- steady_step (carry, prev)
  -> (a -> Bool)                                    -- predicate (carry only)
  -> { final_state: a, trajectory: [{ ...extras }] }
\end{lstlisting}
The \code{bootstrap\_step} fires exactly once, producing both the next carry and
the initial \code{prev}. The \code{steady\_step} then runs the loop, taking the
current carry \emph{and} the threaded \code{prev} as inputs and producing the
next carry, the next \code{prev}, and any extras. The combinator threads
\code{prev} for you; the step never writes the plumbing.

Two disciplines carry over unchanged from the base combinator. The predicate is
\emph{still} pure on the carry alone---it does not read \code{prev}, so any
quantity \code{prev} feeds into termination must be folded into the carry, exactly
as in \cref{sec:predicaterule}. And when no previous value is actually
needed---when \code{prev} can be the empty record---this combinator degenerates
to plain \code{iterate\_while} preceded by one bootstrap step, which is why
\cref{fig:family} draws an arrow collapsing it back to the center. The two are one
family: \code{iterate\_while} is the no-previous idiom, and
\code{iterate\_while\_with\_prev} is its strict generalization for recurrences
that look one step back.

\begin{remark}
Why split the bootstrap out rather than keep the \lstinline[style=l4style]|if it == 0|
test inside the step? Because the split makes the steady step \emph{branch-free}
and the steady carry \emph{one slot lighter}---the previous value lives in the
combinator's threading, not in the state schema. A branch-free step is easier to
reason about, easier to lower to a tensor backend, and free of a conditional that
runs (and is mispredicted) on every iteration but is meaningful on only the
first. \Cref{ch:rotations} studies this ``first-iteration unrolling'' as a
rotation in its own right.
\end{remark}

\section{Demand-driven pruning: the calculus's signature move}
\label{sec:pruning}

We now come to the most elegant property of the combinator, and one of the most
important ideas in the framework. Look again at the step's output record:
\lstinline[style=l4style]|{ state: a, ...extras }|. A step may declare optional
outputs---a residual norm, a breakdown diagnostic, a monitoring metric---\emph{as
record fields, unconditionally}. It always ``produces'' them, in the sense that
they appear in its type. Yet computing a residual norm every iteration is real
work; a production solver that just wants the answer should not pay for
monitoring it will never read.

The naive resolutions to this tension are all familiar and all unpleasant. A
\emph{verbosity flag} threaded through the step (``if \code{verbose}, also compute
the residual''). A \emph{conditional branch} around every optional output. A
\emph{logging monad} (a Writer effect) carrying metrics alongside the state. Each
adds a second code path, and each forces ``monitored'' and ``unmonitored'' to be,
in some sense, different algorithms.

The calculus does none of this. Instead it uses \emph{demand-driven pruning}: the
step declares all its outputs unconditionally, and \emph{whether each output is
actually computed is decided by what a consumer reads}. If a downstream consumer
reads the \code{trajectory}, the per-step extras are computed. If a consumer reads
only \code{final\_state} and never touches the trajectory, then---by the calculus's
evaluation rule---the extras computation is \emph{pruned away entirely}, at every
step, as if it were never written. The same combinator, the same step
definition, the same algorithm: only the consumer's demand differs, and the work
follows the demand.

\begin{keyidea}
Demand-driven pruning makes ``compute residuals for monitoring'' and ``skip
residuals for speed'' the \emph{same algorithm}. The step declares every output
unconditionally; a consumer reads what it wants; outputs nothing reads are
eliminated. There is no verbosity flag, no conditional branch, no logging
monad---there is one step definition, and the demand at the call site decides
what materializes. This is the property that lets the entire book write each
solver \emph{once} and use it for both diagnostics and production.
\end{keyidea}

\subsection{Pruning is graph dead-code elimination}

The mechanism is exactly the dead-code elimination a compiler performs, lifted to
the evaluation of the whole computation graph. The calculus is a
\emph{graph-evaluation language}: every operator produces a record of outputs, and
at solve time only the outputs that some root consumer transitively reaches are
computed. An output no consumer reaches is dead code, and dead code is pruned.

\Cref{fig:pruning} draws this. On the left is the full trajectory: a consumer
reads it, so every step computes its extras and they pile into the list. On the
right, a consumer reads only \code{final\_state}; the trajectory is dead, so the
arrows that would compute and collect the extras are greyed out, and what remains
is the bare carry spine---precisely \code{iterate\_while\_pure}. Formally, when
only \code{final\_state} is observed, the combinator is equivalent to
\begin{lstlisting}[style=l4style]
{ final_state: iterate_while_pure a p f_state, trajectory: [] }
\end{lstlisting}
where \code{f\_state} is the part of the step that computes only the next carry.
The residual-collecting solve and the bare solve are two demands on one
definition.

\begin{figure}[t]
\centering
\begin{tikzpicture}[font=\footnotesize]
  % ============ LEFT: trajectory observed (all live) ============
  \begin{scope}
    \node[font=\small\bfseries, text=simBlue] at (1.5,2.3) {consumer reads \code{trajectory}};
    \foreach \i/\x in {0/0,1/1.5,2/3.0}{
      \node[data, minimum width=7mm] (a\i) at (\x,1.2) {$a_{\i}$};
      \node[opbox, minimum width=7mm, minimum height=5mm] (s\i) at (\x+0.75,1.2) {};
    }
    \node[data, minimum width=7mm] (a3) at (4.5,1.2) {$a_3$};
    \foreach \i in {0,1,2}{ \draw[flow] (a\i.east)--(s\i.west); }
    \draw[flow] (s0.east)--(a1.west); \draw[flow] (s1.east)--(a2.west);
    \draw[flow] (s2.east)--(a3.west);
    % extras live
    \node[data, fill=simBgGold, draw=simGold, minimum width=34mm] (tj) at (2.1,-0.3)
      {\code{trajectory} $=[e_0,e_1,e_2]$};
    \foreach \i in {0,1,2}{ \draw[refedge, simGold!80!black] (s\i.south)--(s\i.south|-tj.north); }
    \node[product, minimum width=18mm] (fs) at (4.5,-0.3) {\code{final\_state}};
    \draw[flow] (a3.south)--(fs.north);
  \end{scope}
  \draw[simGray!50, thick] (5.7,-1.0) -- (5.7,2.5);
  % ============ RIGHT: only final_state observed (extras pruned) ============
  \begin{scope}[xshift=6.6cm]
    \node[font=\small\bfseries, text=simGreen] at (1.5,2.3) {consumer reads only \code{final\_state}};
    \foreach \i/\x in {0/0,1/1.5,2/3.0}{
      \node[data, minimum width=7mm] (b\i) at (\x,1.2) {$a_{\i}$};
      \node[opbox, minimum width=7mm, minimum height=5mm] (t\i) at (\x+0.75,1.2) {};
    }
    \node[data, minimum width=7mm] (b3) at (4.5,1.2) {$a_3$};
    \foreach \i in {0,1,2}{ \draw[flow] (b\i.east)--(t\i.west); }
    \draw[flow] (t0.east)--(b1.west); \draw[flow] (t1.east)--(b2.west);
    \draw[flow] (t2.east)--(b3.west);
    % extras PRUNED (greyed)
    \node[data, fill=white, draw=simGray!40, text=simGray!55, minimum width=34mm] (tjx) at (2.1,-0.3)
      {\textcolor{simGray!55}{\code{trajectory} (pruned)}};
    \foreach \i in {0,1,2}{ \draw[refedge, simGray!35] (t\i.south)--(t\i.south|-tjx.north); }
    \node[product, minimum width=18mm] (fsx) at (4.5,-0.3) {\code{final\_state}};
    \draw[flow] (b3.south)--(fsx.north);
  \end{scope}
\end{tikzpicture}
\caption[Demand-driven pruning]{Demand-driven pruning as graph dead-code
elimination. \emph{Left}: a consumer reads the \code{trajectory}, so every step
computes its extras $e_i$ and they collect into the list---all live. \emph{Right}:
a consumer reads only \code{final\_state}; the trajectory is now dead code, so the
extras computations and the collection arrows (greyed) are eliminated, leaving
only the carry spine. The step definition is identical in both panels; only what
the consumer demands differs. The right panel is exactly
\code{iterate\_while\_pure}.}
\label{fig:pruning}
\end{figure}

This is the calculus's signature move, and it is worth pausing on how much it
buys. In an ordinary imperative solver, the decision ``compute residuals or not''
is a property of the \emph{code}---baked in by a flag or a branch, fixed when the
solver is written. Here it is a property of the \emph{use}: the very same solver,
unmodified, computes residuals exactly when something downstream looks at them.
Monitoring becomes free to add and free to remove. We will lean on this in every
solver chapter; it is why those chapters never present a ``debug version'' beside
a ``fast version.''

\section{What the combinator does, and does not, satisfy}
\label{sec:laws}

A combinator is only as useful as the rewrites it licenses. \code{iterate\_while}
satisfies a handful of laws---facts that let us transform one expression into an
equal one---and, just as importantly, fails to satisfy several tempting ones.
Knowing both keeps us from ``simplifying'' a loop into something that means
something different. We state the most useful affirmatively, and then catalogue
the traps.

\paragraph{The pruning law (the one that earns its keep).} If a consumer observes
only \code{final\_state}, then \code{iterate\_while}~$a~p~f$ equals
\lstinline[style=l4style]|{ final_state: iterate_while_pure a p f_state, trajectory: [] }|,
where \code{f\_state} computes only the step's next carry. This is
\cref{sec:pruning} stated as an equation: the monitored loop and the bare loop are
\emph{equal} expressions whenever nothing reads the trajectory. It is the law
that justifies writing one step definition and using it for both diagnostics and
production.

\paragraph{The base case.} If $p(a)$ is false at the outset,
\code{iterate\_while}~$a~p~f = $
\lstinline[style=l4style]|{ final_state: a, trajectory: [] }| for \emph{any}
step \code{f}. The combinator does not run the body even once before testing---a
direct reading of the small-step rule, and the formal content of Rule~1.

Now the traps. Each is a rewrite that looks innocent and is wrong.

\begin{description}[style=nextline, leftmargin=0pt, font=\normalfont\itshape]
  \item[You may not hoist the predicate out of the loop.] \code{iterate\_while}
  re-evaluates $p$ on \emph{every} carry, not once. Replacing the loop with ``test
  $p(a)$ once, then iterate unconditionally'' turns a terminating loop into an
  infinite one---the predicate would never get the chance to become false. The
  per-iteration re-test is the whole point.
  \item[You may not fuse two loops into one.] Running
  \code{iterate\_while}~$a~p_1~f$ and then \code{iterate\_while} on its result with
  $p_2$ is \emph{not} the same as a single loop with a combined predicate. The two
  phases stop on different conditions and concatenate their trajectories; a single
  flattened loop has one predicate and one trajectory. (This is exactly why GMRES,
  in \cref{ch:gmres}, wraps an \emph{outer} restart loop around its inner
  \code{iterate\_while} rather than flattening the two.)
  \item[You may not swap the predicate and the step.] Testing after the step
  (\texttt{do}/\texttt{while}) is a \emph{different} combinator that can run one
  more time than \code{iterate\_while} on an already-satisfied initial carry. The
  two agree only when the initial predicate is true.
  \item[There is no identity element, and the type does not promise termination.]
  \code{iterate\_while} is a fold, not a monoid operation; there is no carry that
  makes it a no-op for all predicates and steps. And nothing in the signature
  guarantees the loop halts---a step that never lets the predicate go false
  diverges. Termination is an obligation discharged \emph{by the algorithm}: a
  bounded \code{max\_it} folded into the carry, or a convergence proof (Krylov
  methods on well-conditioned systems), or an explicit done-state the step always
  reaches.
\end{description}

The lesson of the non-laws is uniform: \code{iterate\_while} is a faithful record
of an ordered, possibly-unbounded computation, and the order is observable. The
combinator earns its simplicity by \emph{not} pretending iterations commute or
that loops fuse for free---claims that would be convenient and false.

\section{Why one combinator suffices}
\label{sec:onecombinator}

We close the conceptual development by collecting the claim the chapter has been
building toward: \emph{one combinator expresses every iteration in the book}. The
three patterns a simulator runs---independent steps, sequential carries, and
convergence loops---are all \code{iterate\_while} (or its family), seen at
different settings. \Cref{fig:threeuses} shows the three side by side.

\begin{figure}[t]
\centering
\begin{tikzpicture}[font=\footnotesize, node distance=2.5mm]
  \def\cw{38mm}
  % map
  \node[stage, minimum width=\cw] (mH) at (0,0) {as a \textbf{map}};
  \node[opbox, minimum width=\cw, align=left, below=of mH] (mB)
    {predicate counts down\\a fixed list; steps are\\independent (no real carry)};
  \node[cornertag, below=1mm of mB] {independent solves};
  % fold
  \node[stage, minimum width=\cw, right=12mm of mH] (fH) {as a \textbf{fold}};
  \node[opbox, minimum width=\cw, align=left, below=of fH] (fB)
    {carry threads forward;\\each step consumes the\\last (time-stepping)};
  \node[cornertag, below=1mm of fB] {sequential chain};
  % converge
  \node[stage, minimum width=\cw, right=12mm of fH] (cH) {as a \textbf{converge}};
  \node[opbox, minimum width=\cw, align=left, below=of cH] (cB)
    {predicate tests a residual\\folded into the carry;\\stop at tolerance (Krylov)};
  \node[cornertag, below=1mm of cB] {until converged};
  % the unifying brace
  \node[font=\small\bfseries, text=simBlue] (one) at ($(fH)+(0,1.4)$) {one \code{iterate\_while}};
  \draw[depedge] (one) -- (mH.north);
  \draw[depedge] (one) -- (fH.north);
  \draw[depedge] (one) -- (cH.north);
\end{tikzpicture}
\caption[One combinator, three uses]{Three iteration patterns, one combinator.
A \emph{map} of independent steps is \code{iterate\_while} with a predicate that
walks a fixed count and steps that ignore the carry. A \emph{fold} is the
sequential chain---time-stepping, where each step consumes the last. A
\emph{converge} loop tests a residual folded into the carry and stops at a
tolerance---the Krylov solvers. The full catalogue, with the precise sense in
which a map is a degenerate fold, is \cref{ch:combinators}.}
\label{fig:threeuses}
\end{figure}

A \emph{map}---a family of independent solves, like the electrostatic analysis's
several right-hand sides---is \code{iterate\_while} with a predicate that simply
counts off a fixed number of steps and a step whose output does not depend on the
carry. (One ordinarily writes a map as \code{map}, not as a degenerate
\code{iterate\_while}; the point is that it \emph{can} be seen this way, and
\cref{ch:combinators} makes the relationship precise.) A \emph{fold}---a
sequential chain like time-stepping, where each state comes from the last---is
\code{iterate\_while} in its plainest reading. A \emph{convergence loop}---a
Krylov solver running until its residual drops below tolerance---is
\code{iterate\_while} with the residual folded into the carry and the predicate
testing it. Maps, folds, and convergence loops: three faces of one combinator.

\begin{keyidea}
There is \emph{one} iteration combinator. Independent steps (a map), sequential
carries (a fold), and convergence loops (run-until-converged) are all
\code{iterate\_while} at different settings of its predicate and step. The
solver's state-threading, the time-stepper's march, the adaptive loop of
\cref{ch:lifecycle}---every iteration the simulator performs is one specialization
of the combinator built in this chapter. Learn it once and you have learned the
control structure of the entire simulator.
\end{keyidea}

The headline instance is the Krylov step. \Cref{ch:cg} and \cref{ch:gmres} build
the conjugate-gradient and GMRES solvers, and at their heart each is a single
\code{krylov\_step} folded by \code{iterate\_while}: the carry is the solver's
working state (current iterate, residual, search direction, convergence flag),
the step applies the operator once and updates that state, the predicate reads
the convergence flag and iteration count from the carry, and the trajectory---when
a consumer asks for it---is the residual history. Everything those chapters do is
filling in the step; the loop is the one we have here. \Cref{ch:combinators}
catalogues the full family of specializations, and \cref{ch:monad} explains how
the step threads the simulator's monadic state while the carry stays a pure value.

\section{Worked examples}

\begin{workedexample}{Newton's method for $\sqrt{2}$ as a fold}{newton}
We compute $\sqrt 2$ by Newton iteration, written as \code{iterate\_while}. The
carry holds both the current estimate and the residual that the predicate will
test---folding the termination quantity into the carry, exactly as
\cref{sec:predicaterule} requires.

The carry is a record \lstinline[style=l4style]|{ x: Float, residual: Float }|.
The step applies the Newton update and recomputes the residual; the predicate
tests the residual against a tolerance.
\begin{lstlisting}[style=l4style]
newton_sqrt2 :: Float -> { final_state: { x, residual }, trajectory: [...] }
newton_sqrt2 x0 =
  iterate_while
    { x: x0, residual: abs (x0*x0 - 2) }        -- initial carry
    (\s -> s.residual > 1e-12)                  -- predicate: pure on the carry
    (\s -> let x' = 0.5 * (s.x + 2 / s.x)       -- Newton step
           in  { state:    { x: x', residual: abs (x'*x' - 2) },
                 residual: abs (x'*x' - 2) })   -- extra: the per-step residual
\end{lstlisting}
Starting from $x_0 = 1$, trace the carry. Each step computes
$x' = \tfrac12(x + 2/x)$ and the residual $|x'^2 - 2|$:
\begin{center}
\setlength{\tabcolsep}{8pt}
\renewcommand{\arraystretch}{1.15}
\begin{tabular}{@{}clll@{}}
\toprule
step & carry $x$ & next $x' = \tfrac12(x + 2/x)$ & residual $|x'^2 - 2|$ \\
\midrule
0 & $1.000000000000$ & $1.500000000000$ & $2.5\times10^{-1}$ \\
1 & $1.500000000000$ & $1.416666666667$ & $6.9\times10^{-3}$ \\
2 & $1.416666666667$ & $1.414215686275$ & $6.0\times10^{-6}$ \\
3 & $1.414215686275$ & $1.414213562374$ & $4.5\times10^{-12}$ \\
4 & $1.414213562374$ & $1.414213562373$ & $<10^{-12}$ \;(stop) \\
\bottomrule
\end{tabular}
\end{center}
The residual falls roughly quadratically---each step doubles the correct digits.
After step~4 the predicate \lstinline[style=l4style]|s.residual > 1e-12| is false,
the loop halts, and \code{final\_state.x} is $\sqrt 2$ to twelve places. Two
points to notice. First, the predicate never touches the freshly computed
residual directly; it reads \code{s.residual} from the carry, which the
\emph{previous} step stored there. Second, the residual appears \emph{both} in the
carry (so the predicate can read it) and in the extras (so a consumer plotting
convergence can read the trajectory)---two readers, one computed value.
\end{workedexample}

\begin{workedexample}{A geometric-series partial sum as \texttt{iterate\_while\_pure}}{geom}
When a loop has nothing per-step to report, \code{iterate\_while\_pure} is the
idiom. We sum the geometric series $\sum_{k=0}^{\infty} r^k = 1/(1-r)$ for
$|r| < 1$ by accumulating partial sums until a new term changes the sum
negligibly.

The carry is a record \lstinline[style=l4style]|{ sum: Float, term: Float }|: the
partial sum so far and the current term $r^k$. The step adds the term and
advances it to $r^{k+1}$; the predicate stops when the term is below tolerance.
There are no extras, so we use the pure form and read back only the final sum.
\begin{lstlisting}[style=l4style]
geom_sum :: Float -> Float
geom_sum r =
  (iterate_while_pure
     { sum: 0.0, term: 1.0 }                    -- start: sum 0, term r^0 = 1
     (\s -> abs s.term > 1e-10)                 -- stop when the term is tiny
     (\s -> { sum: s.sum + s.term, term: s.term * r })   -- accumulate, advance
  ).sum
\end{lstlisting}
For $r = \tfrac12$ the exact sum is $1/(1-\tfrac12) = 2$. Tracing the carry:
\begin{center}
\setlength{\tabcolsep}{8pt}
\renewcommand{\arraystretch}{1.1}
\begin{tabular}{@{}clll@{}}
\toprule
step & \code{term} ($= r^k$) & \code{sum} after adding & remaining error $2 - \text{sum}$ \\
\midrule
0 & $1.0000$    & $1.0000$ & $1.0000$ \\
1 & $0.5000$    & $1.5000$ & $0.5000$ \\
2 & $0.2500$    & $1.7500$ & $0.2500$ \\
3 & $0.1250$    & $1.8750$ & $0.1250$ \\
$\vdots$ & $\vdots$ & $\vdots$ & $\vdots$ \\
33 & $\approx10^{-10}$ & $\approx 2.0000$ & $<10^{-10}$ \;(stop) \\
\bottomrule
\end{tabular}
\end{center}
The carry threads two values, the predicate reads one of them (\code{term}), and
because nothing per-step is collected the \code{iterate\_while\_pure} sugar saves
us a trajectory we would never read. This is fixed-point iteration in disguise:
the map $s \mapsto (\,s.\code{sum} + s.\code{term},\; r\cdot s.\code{term}\,)$ is a
contraction, and the loop walks to its fixed point.
\end{workedexample}

\begin{workedexample}{The predicate-on-extras pitfall, wrong and right}{pitfall}
This example shows the mistake of \cref{sec:predicaterule} concretely and its fix
side by side. We want a generic ``iterate until the residual is below
\code{tol}'' loop. The tempting but \emph{wrong} sketch tries to let the predicate
read the residual the step produces as an extra:
\begin{lstlisting}[style=l4style]
-- WRONG: the predicate wants to read an extra the step has not produced yet.
solve_BAD x0 =
  iterate_while_BAD
    x0
    (\extras -> extras.residual > tol)       -- (!) reads the step's extra
    (\x -> let x'  = step_body x
               r   = residual_of x'
           in  { state: x', residual: r })   -- residual is an *extra*
\end{lstlisting}
The predicate \lstinline[style=l4style]|\extras -> extras.residual > tol| needs an
\code{extras} record to inspect---but the predicate fires \emph{before} the step
(Rule~1), so on the first iteration no such record exists. The dependency is
circular; \cref{fig:predrule} draws it crossed out. The calculus will not even
let us write this: the predicate's type is \lstinline[style=l4style]|a -> Bool|,
so it cannot name an \code{extras} argument at all.

The fix folds the residual into the carry. The carry becomes a record carrying
both the iterate and its residual; the step computes the residual and stores it
in the carry; the predicate reads it from there:
\begin{lstlisting}[style=l4style]
-- RIGHT: the residual lives in the carry; the predicate reads it from there.
solve_OK x0 =
  iterate_while
    { x: x0, residual: residual_of x0 }      -- residual folded into the carry
    (\s -> s.residual > tol)                 -- predicate: pure on the carry
    (\s -> let x' = step_body s.x
               r  = residual_of x'
           in  { state: { x: x', residual: r },   -- carry the residual forward
                 residual: r })                   -- also surface it as an extra
\end{lstlisting}
Now the predicate reads \code{s.residual}, a field of the carry that the previous
step stored---which exists from the start because the \emph{initial} carry
computes it too. The residual is still available to trajectory consumers as an
extra; it is simply \emph{also} in the carry so the predicate can see it. The
single change---move the control quantity from the extras into the carry---turns
the impossible loop into the canonical one. This is the move to reach for every
time a predicate seems to need something the step produced.
\end{workedexample}

\begin{example}[A conjugate-gradient step, previewed]
We can already sketch the shape \cref{ch:cg} fills in. A conjugate-gradient solve
is \code{iterate\_while} whose carry is the CG working state---current iterate
$\vx$, residual $\vr$, search direction $\vp$, residual norm, and a
\code{converged} flag---and whose step is one \code{krylov\_step}: apply the
operator once, update $\vx$, $\vr$, $\vp$, recompute the residual norm, and set
\code{converged} when it drops below tolerance. The predicate reads the carry:
\lstinline[style=l4style]|\s -> s.it < max_it && not s.converged|. The trajectory,
when a consumer reads it, is the residual history. The entire control structure is
the combinator of this chapter; \cref{ch:cg} supplies only the step body.
\end{example}

\summarysection
An imperative loop overwrites a cell; a \emph{fold} carries a value forward,
each step computing the next from the last. This single reframing---mutation into
value-threading---turns every iteration in the book into an instance of one
combinator, \code{iterate\_while}. Its signature is
\lstinline[style=l4style]|a -> (a -> Bool) -> (a -> { state: a, ...extras }) -> { final_state, trajectory }|:
from an initial carry, while a predicate holds of the current carry, take a step
that produces the next carry plus per-step extras, collecting the extras into a
trajectory. Two rules about the predicate are load-bearing: it fires
\emph{first} (the \texttt{while} convention), and it is \emph{pure on the carry}
(it sees the carry and nothing else). The consequence is the central discipline:
any termination quantity---a residual, a flag, a count---must be \emph{folded into
the carry}, computed by the step and read by the predicate. Trying instead to test
an \emph{extra} the step produced is the predicate-on-extras pitfall, a circular
dependency the calculus forbids by typing. The family has three members:
\code{iterate\_while\_pure} (no extras, returns just the final carry),
\code{iterate\_while} (the canonical form), and
\code{iterate\_while\_with\_prev} (a bootstrap plus a branch-free steady step for
two-term recurrences). \emph{Demand-driven pruning}---graph dead-code
elimination---makes a monitored solve and a bare solve the \emph{same algorithm}:
the step declares every output unconditionally, and outputs no consumer reads are
eliminated, with no flag, branch, or logging monad. Finally, one combinator
suffices: maps, folds, and convergence loops are all \code{iterate\_while} at
different settings, and the headline instance, the Krylov step
(\cref{ch:cg},~\cref{ch:gmres}), is a single step folded by exactly this loop.

\furtherreading
The fold and its relatives are foundational in functional programming; Peyton
Jones's account of Haskell~\citep{peytonjones2003} develops \code{map},
\code{foldl}, \code{foldr}, and the laws relating them. The idea that
``effects''---here, the optional per-step outputs---can be expressed as values
threaded through pure code, rather than as ambient mutation or a logging channel,
is the subject of Wadler's monads tutorial~\citep{wadler1995monads}; we take up the
state monad that threads the simulator's own mutable bookkeeping in
\cref{ch:monad}. For the numerical iterations the combinator will drive---Krylov
methods, their convergence, and the residual histories the trajectory
records---Saad's monograph~\citep{saad2003} is the standard reference, and we build
those solvers as folds in \cref{ch:cg} and \cref{ch:gmres}.

\exercisessection
\begin{enumerate}[nosep]
  \item Write the running sum of a list as both an imperative \texttt{for}-loop and
  a \code{foldl}. Identify, in each, the carry and the step. Which version could
  you run in parallel, and which not? (Hint: which one has a carry?)
  \item In \cref{wex:newton} the residual appears in two places in
  the step's output record---in the carry's \code{residual} field and as a
  top-level \code{residual} extra. Explain what each copy is for, and which
  reader consumes which. What would break if you kept only the extra?
  \item State the \texttt{while} convention precisely and contrast it with
  \texttt{do}/\texttt{while}. Construct an initial carry on which the two
  conventions give \emph{different} numbers of steps. Why does a solver want the
  \texttt{while} convention?
  \item A colleague proposes a predicate of type
  \lstinline[style=l4style]|(a, extras) -> Bool| so the loop can ``stop as soon
  as the step's residual is small.'' Explain the circular dependency this creates
  on the first iteration, and rewrite their loop using the carry-folding fix of
  \cref{sec:predicaterule}.
  \item Using the small-step rule of \cref{sec:iteratewhile}, evaluate
  \lstinline[style=l4style]|iterate_while a (\_ -> False) f| by hand for any step
  \code{f}. What is the \code{final\_state}? The \code{trajectory}? How many times
  does \code{f} run? Which of the two predicate rules does this exercise probe?
  \item A solver is written with a step that always computes a residual norm as an
  extra. In one run a consumer plots the residual history; in another it reads only
  the final iterate. Using \cref{sec:pruning}, describe what computation each run
  actually performs, and argue that the two are running the \emph{same} algorithm.
  Why is this preferable to a \code{verbose} flag?
  \item Express a fixed-degree polynomial smoother---a loop that runs exactly $d$
  steps regardless of any residual---as \code{iterate\_while\_pure}. What does the
  carry contain? What does the predicate test? (Hint: fold a step counter into the
  carry.)
  \item \emph{(Two-term recurrence.)} The Fibonacci sequence obeys
  $F_{k+1} = F_k + F_{k-1}$, a two-term recurrence. Sketch its computation using
  \code{iterate\_while\_with\_prev}: identify the bootstrap step (what is the first
  ``previous''?), the steady step, the carry, and the predicate. Then argue that
  the same sequence can be computed by plain \code{iterate\_while} if you widen the
  carry to hold \emph{both} $F_k$ and $F_{k-1}$. What does
  \code{iterate\_while\_with\_prev} buy over that widening?
  \item \emph{(Synthesis.)} The adaptive loop of \cref{ch:lifecycle} threads the
  mesh forward, refining it each pass until the error is small enough. Cast it as
  \code{iterate\_while}: name the carry, the step, the predicate, and the
  termination quantity that must be folded into the carry. Is its trajectory ever
  read? If not, what does \cref{sec:pruning} say happens to the per-step extras?
\end{enumerate}

\chapter{Maps, Folds, and the Combinator-Primary Model}
\label{ch:combinators}

\chapterorient{The previous chapter built one combinator, \texttt{iterate\_while},
and showed that every iteration in the book is a setting of it. This chapter widens
the lens. \texttt{iterate\_while} is the iteration primitive; it has \emph{relatives}
that organize the rest of the simulator---\code{map} for independent work,
\code{fold} for chained work, and a small algebra of data-combining verbs the steps
call. We give the four solve corners of \cref{ch:situating} their combinator names,
we meet the tensor-sum combinator \code{linear\_combination} and watch the familiar
BLAS routines fall out of it as specializations, and we state the chapter's
thesis---the \emph{combinator-primary model}: a few higher-order combinators are
the load-bearing vocabulary, and the many named operators are \emph{specialization
notes} hung beneath them. This is why a book about a hundred-thousand-line simulator
can be short.}

\section{From one combinator to a family of relatives}

\Cref{ch:fold} earned a strong claim: there is \emph{one} iteration combinator,
\code{iterate\_while}, and independent steps, sequential carries, and convergence
loops are all that combinator at different settings. We will not retreat from it.
But ``one combinator expresses every iteration'' is a statement about
\emph{control flow}---about how a loop threads its carry and tests its predicate. A
simulator is more than its loops. It also \emph{combines data}: it sums scaled
tensors, it forms inner products, it assembles operators. And when it does run a
family of solves, the \emph{shape} of that family---independent versus
chained---deserves a name of its own, because the two shapes have radically
different consequences for parallelism.

So this chapter steps up one level. Where \cref{ch:fold} looked \emph{inside} a
single loop, this chapter looks at the small set of combinators that organize the
\emph{whole} simulator: the ones that drive families of solves, and the ones that
combine tensors. \Cref{fig:relatives} lays out the family tree. At the root is the
distinction the whole chapter turns on---\code{map} versus \code{fold}, independent
versus carried---and beneath it hang the named combinators a driver actually calls.

\begin{figure}[t]
\centering
\begin{tikzpicture}[font=\footnotesize, node distance=4mm]
  % root
  \node[stage, minimum width=58mm, align=center] (root) at (0,0)
    {the combinator vocabulary};
  % two branches: iteration-structural vs data-algebra
  \node[stage, minimum width=44mm, align=center, fill=simBgGreen, draw=simGreen]
    (iter) at (-3.6,-1.7) {iteration-structural\\(thread a carry / family)};
  \node[stage, minimum width=44mm, align=center, fill=simBgGold, draw=simGold]
    (data) at (3.6,-1.7) {data-algebra\\(combine tensors)};
  \draw[depedge] (root) -- (iter);
  \draw[depedge] (root) -- (data);
  % iteration-structural leaves
  \node[opbox, minimum width=40mm, align=left, below=10mm of iter] (il) {%
    \code{solve\_family} \;\itshape map (independent)\\
    \code{frequency\_sweep} \;\itshape op-varying map\\
    \code{fold\_solve} \;\itshape fold (chained)\\
    \code{iterate\_while} \;\itshape the loop primitive};
  \draw[depedge] (iter) -- (il);
  % data-algebra leaves
  \node[opbox, minimum width=40mm, align=left, below=10mm of data] (dl) {%
    \code{linear\_combination} \;\itshape tensor sum\\
    \quad \code{scal}/\code{axpy}/\code{axpby}/\dots \;\itshape leaves\\
    \code{inner\_product} \;\itshape scalar fold\\
    \quad \code{dot}/\code{nrm2} \;\itshape leaves};
  \draw[depedge] (data) -- (dl);
\end{tikzpicture}
\caption[The combinator vocabulary]{The combinator vocabulary of the simulator,
in two branches. The \emph{iteration-structural} combinators (left) thread a carry
or drive a family of solves; they are relatives of the \code{iterate\_while} of
\cref{ch:fold}, distinguished by whether the family is an independent \emph{map}
or a chained \emph{fold}. The \emph{data-algebra} combinators (right) combine
tensors into tensors or scalars; each has a cluster of named BLAS-style leaves
hanging beneath it. The whole of this short chapter is a tour of these two boxes
and the one thesis they illustrate: the combinator is the entry, the leaves are
notes.}
\label{fig:relatives}
\end{figure}

\section{The two shapes of a family: map and fold}
\label{sec:mapfold}

Before naming the simulator's combinators, fix the one distinction everything
rests on. We met it in \cref{ch:fold} as the contrast between \code{map} and
\code{foldl}; here we make it the organizing axis.

A \code{map} applies the \emph{same} operation to each element of a collection
\emph{independently}:
\begin{lstlisting}[style=l4style]
map f [x1, x2, x3] = [f x1, f x2, f x3]
\end{lstlisting}
The three results do not depend on one another. \code{f x2} can be computed before,
after, or at the same time as \code{f x1}; you could run each on a separate
machine and collect the answers in any order. A \code{map} has \emph{no carry}.

A \code{foldl} threads an accumulator---a \emph{carry}---through the collection,
each step consuming what the last produced:
\begin{lstlisting}[style=l4style]
foldl step s0 [x1, x2, x3] = step (step (step s0 x1) x2) x3
\end{lstlisting}
Here the order is forced. The accumulator entering the third step is exactly what
the second step produced; you cannot start step three until step two has finished.
A \code{fold} is a \emph{chain}.

This is the single most consequential distinction in the simulator's architecture,
and \cref{fig:mapfold} draws it bare. The map fans out; the fold lines up. The map
parallelizes; the fold serializes. Everything else in this section is a
specialization of one or the other.

\begin{figure}[t]
\centering
\begin{tikzpicture}[font=\footnotesize]
  % ===== LEFT: map = independent fan-out =====
  \begin{scope}
    \node[font=\small\bfseries, text=simBlue] at (1.4,2.7) {a \code{map}: independent fan-out};
    \node[data, minimum width=10mm] (src) at (1.4,1.9) {family};
    \foreach \i/\x in {1/0,2/1.4,3/2.8}{
      \node[opbox, minimum width=9mm] (m\i) at (\x,0.8) {\code{f}};
      \node[product, minimum width=9mm, minimum height=5mm] (r\i) at (\x,-0.4) {$y_{\i}$};
      \draw[flow] (src.south) -- (m\i.north);
      \draw[flow] (m\i.south) -- (r\i.north);
    }
    \node[font=\scriptsize, text=simGreen, align=center] at (1.4,-1.3)
      {no arrows \emph{between} the \code{f}'s:\\reorderable, parallelizable};
  \end{scope}
  \draw[simGray!50, thick] (4.2,-1.6) -- (4.2,3.0);
  % ===== RIGHT: fold = sequential chain =====
  \begin{scope}[xshift=5.0cm]
    \node[font=\small\bfseries, text=simGreen] at (2.4,2.7) {a \code{fold}: sequential chain};
    \node[data, minimum width=9mm] (s0) at (-0.3,0.8) {$s_0$};
    \foreach \i/\x in {1/1.1,2/2.5,3/3.9}{
      \node[opbox, minimum width=9mm] (g\i) at (\x,0.8) {\code{step}};
    }
    \node[data, minimum width=9mm] (s3) at (5.3,0.8) {$s_3$};
    \draw[flow] (s0.east) -- (g1.west);
    \draw[flow] (g1.east) -- node[above,font=\scriptsize]{$s_1$} (g2.west);
    \draw[flow] (g2.east) -- node[above,font=\scriptsize]{$s_2$} (g3.west);
    \draw[flow] (g3.east) -- (s3.west);
    \foreach \i/\x in {1/1.1,2/2.5,3/3.9}{
      \node[font=\scriptsize, text=simGold!80!black] at (\x,1.65) {$x_{\i}$};
      \draw[refedge, simGold!80!black] (\x,1.5) -- (g\i.north);
    }
    \node[font=\scriptsize, text=simRust, align=center] at (2.5,-0.6)
      {each \code{step} needs the previous $s$:\\strictly ordered, not parallelizable};
  \end{scope}
\end{tikzpicture}
\caption[Map versus fold: the fundamental split]{The fundamental split. A
\emph{map} (left) applies \code{f} to each family element independently---the
three applications share no data, so they fan out, reorder freely, and run in
parallel. A \emph{fold} (right) threads a carry $s_0 \to s_1 \to s_2 \to s_3$
through the chain---each \code{step} consumes the previous carry, so the chain is
strictly ordered and cannot be parallelized. The presence or absence of the
horizontal carry arrows is the entire difference, and it decides whether the work
is embarrassingly parallel or intrinsically sequential.}
\label{fig:mapfold}
\end{figure}

\begin{keyidea}
A \emph{map} runs a family of \emph{independent} computations: no carry, so the
elements commute, reorder, and parallelize. A \emph{fold} runs a \emph{chain}:
each step consumes the previous step's output, so the steps are strictly ordered
and intrinsically sequential. This one axis---is there a carry between the
elements?---decides the parallelism of every family of solves the simulator runs.
A map is in fact the degenerate fold whose step \emph{ignores} the carry; the two
are one family, split by whether the carry is used.
\end{keyidea}

The claim that a map is a degenerate fold is worth making exact, because it is the
formal hinge on which the whole ``one combinator'' story of \cref{ch:fold} turns.
A \code{foldl} that drops its accumulator on every step does no threading at all:
the result of each step depends only on the element, not on the carry, so the steps
are independent and we have a map in disguise.

\begin{proposition}[Map is the carry-dropping fold]
\label{prop:mapfold}
For any element function \code{f}, the map \lstinline[style=l4style]|map f xs| equals
the trajectory of a left fold whose step ignores its accumulator and emits
\code{f} of the element:
\begin{lstlisting}[style=l4style]
map f xs == (foldl_collecting (\_ x -> f x) seed xs).collected
\end{lstlisting}
where \code{foldl\_collecting} threads a carry but \emph{also} appends each step's
output to a collected list. Because the step's first argument (the carry) never
appears on its right-hand side, the steps share no data, and the collected list is
order-independent.
\end{proposition}

The proposition is not deep---it is almost a restatement of the definitions---but
its \emph{consequence} is the architecture of the simulator. It says the map and
the fold are not two unrelated control structures to be learned separately; they
are one structure, the value-threading fold, observed at two settings of a single
switch: \emph{does the step read its carry?} When the answer is no, the work fans
out and parallelizes; when the answer is yes, the work lines up and serializes. The
entire taxonomy of the next section is this one switch, crossed with a second
(does the operator change?).

It pays to draw out the three practical consequences of the split, because each
recurs in every later Part.

\begin{description}[style=nextline, leftmargin=0pt, font=\normalfont\itshape]
  \item[Parallelism follows the carry.] A map's elements may be evaluated in any
  order, on any number of machines, and reassembled; a fold's steps must run in
  sequence on one timeline. When \cref{part:modalities} reports that electrostatics
  ``scales trivially across conductors'' but the transient analysis ``cannot
  overlap its timesteps,'' it is reporting exactly this---the first is a map, the
  second a fold.
  \item[Reordering is a correctness question, not a performance one.] For a map,
  reordering the family permutes the output identically and changes nothing else;
  for a fold, reordering the steps changes the \emph{answer}, because step $k$'s
  input is step $k{-}1$'s output. ``Can I reorder this?'' is therefore answered by
  ``is there a carry?''---a structural fact visible in the combinator's type, not a
  property to be discovered by testing.
  \item[Memory follows the carry too.] A map need keep only the inputs and the
  collected outputs; a fold must additionally keep the live carry as it threads.
  When the carry is a whole field state (the transient analysis) this is the
  dominant cost, and \cref{ch:transient} returns to it. A map has no such live
  intermediate.
\end{description}

\section{The four solve corners, as combinators}
\label{sec:corners}

We can now return to the four solve corners of \cref{ch:situating} and give each
its combinator name. There we classified a driver's solve stage by two questions:
is the operator \emph{fixed} or \emph{varying} as we proceed, and is the
computation a \emph{map} or a \emph{fold}? Crossing those questions gave a
$2\times2$ grid with three occupied cells and a fourth corner sitting outside it.
\Cref{fig:cornercomb} redraws that grid, now labelling each cell with the
combinator the framework uses for it.

\begin{figure}[t]
\centering
\begin{tikzpicture}[font=\footnotesize]
  \node[cornertag, anchor=south] at (1.2,4.5) {operator \emph{fixed}};
  \node[cornertag, anchor=south] at (5.6,4.5) {operator \emph{varying}};
  \node[cornertag, rotate=90, anchor=south] at (-1.5,3.2) {a \textbf{map}};
  \node[cornertag, rotate=90, anchor=south] at (-1.5,0.9) {a \textbf{fold}};
  % top-left: solve_family
  \node[stage, minimum width=34mm, minimum height=20mm, anchor=south west,
    fill=simBgBlue, align=center] (c1) at (0,2.4) {};
  \node[anchor=north, font=\scriptsize\bfseries] at (1.7,4.1) {fixed-operator map};
  \node[anchor=center, font=\small\ttfamily, text=simBlue] at (1.7,3.3) {solve\_family};
  \node[anchor=north, font=\scriptsize\itshape] at (1.7,2.95) {electrostatic,};
  \node[anchor=north, font=\scriptsize\itshape] at (1.7,2.65) {magnetostatic};
  % top-right: frequency_sweep
  \node[stage, minimum width=34mm, minimum height=20mm, anchor=south west,
    fill=simBgBlue, align=center] (c2) at (4.4,2.4) {};
  \node[anchor=north, font=\scriptsize\bfseries] at (6.1,4.1) {operator-varying map};
  \node[anchor=center, font=\small\ttfamily, text=simBlue] at (6.1,3.3) {frequency\_sweep};
  \node[anchor=north, font=\scriptsize\itshape] at (6.1,2.85) {driven};
  % bottom-left: fold_solve
  \node[stage, minimum width=34mm, minimum height=20mm, anchor=south west,
    fill=simBgGreen, align=center] (c3) at (0,0) {};
  \node[anchor=north, font=\scriptsize\bfseries] at (1.7,1.7) {state-threaded fold};
  \node[anchor=center, font=\small\ttfamily, text=simGreen] at (1.7,0.95) {fold\_solve};
  \node[anchor=north, font=\scriptsize\itshape] at (1.7,0.55) {transient, AMR loop};
  % bottom-right: eigsolve (outside)
  \node[extern, minimum width=34mm, minimum height=20mm, anchor=south west,
    align=center] (c4) at (4.4,0) {};
  \node[anchor=north, font=\scriptsize\bfseries] at (6.1,1.7) {opaque black-box};
  \node[anchor=center, font=\small\ttfamily, text=simViolet] at (6.1,0.95) {eigsolve};
  \node[anchor=north, font=\scriptsize\itshape] at (6.1,0.55) {eigenmode, boundary-mode};
  \node[font=\scriptsize\itshape, text=simViolet] at (6.1,-0.35) {(no loop we write)};
\end{tikzpicture}
\caption[The four solve corners as combinators]{The four solve corners of
\cref{ch:situating}, now named as combinators (compare \cref{fig:corners}). A
fixed-operator map is \code{solve\_family}: assemble once, map independent solves.
An operator-varying map is \code{frequency\_sweep}: rebuild the operator inside the
map. A state-threaded fold is \code{fold\_solve}: chain the solves, each from the
last. The fourth corner is \code{eigsolve}, the opaque library call that sits
outside the map/fold grid because we write no loop at all. Three of the four are
relatives of \code{iterate\_while}; the fourth is the deliberate non-loop.}
\label{fig:cornercomb}
\end{figure}

\subsection{\texttt{solve\_family}: the fixed-operator map}

The first corner is a \code{map}. Assemble the operator $\Kop$ \emph{once}, then
solve $\Kop\,\vx_i = \vb_i$ for a family of right-hand sides $\vb_1,\dots,\vb_n$.
The operator never changes; only the right-hand side does. Because the solves
share one operator and depend on nothing but their own right-hand side, they are
independent: a \code{map}. The combinator is \code{solve\_family}, and its
definition is a one-liner---a \code{map} of the single-solve cap over the family:
\begin{lstlisting}[style=l4style]
solve_family :: OpParams -> [Inputs] -> [SimState]
solve_family op rhss = map (\inp -> ksp_solve op inp) rhss
\end{lstlisting}
The load-bearing fact is the \emph{hoist}: \code{op} is captured once, outside the
\code{map}, and threaded unchanged into every per-element solve. The solver built
from \code{op}---a factorization, a preconditioner---is invariant across the map,
so it too is built once and reused. This is exactly the Palace pattern where the
operator setup call sits \emph{outside} the sweep loop: electrostatics and
magnetostatics both assemble $\Kop$ once and then loop over conductors, reusing the
one operator. The exemplars are electrostatics (\cref{ch:electrostatics}) and
magnetostatics, which are this corner with different operators and different
families.

Because the map carries nothing between elements, \code{solve\_family} obeys the
\emph{concatenation law}: \lstinline[style=l4style]|solve_family op (a ++ b) = solve_family op a ++ solve_family op b|.
Splitting the family, chunking it, reordering it, or solving its pieces on separate
machines all give the same answer---the algebraic statement of embarrassing
parallelism. There is a companion \emph{order law}:
\lstinline[style=l4style]|solve_family op (permute rhss) = permute (solve_family op rhss)|,
which says the underlying solves commute even though the collection preserves
position. Together the two laws are the formal license for every parallel
realization of the fixed-operator corner: the family may be partitioned arbitrarily,
solved out of order, and stitched back together, and the answer is unchanged. Note
what the laws rest on---a \emph{shared, unchanging} \code{op}. Strip that away and
the next combinator forfeits both laws.

\subsection{\texttt{frequency\_sweep}: the operator-varying map}

Keep the map, but change one thing: the operator is no longer fixed. The driven
analysis (\cref{ch:driven}) computes scattering parameters by solving the
time-harmonic system at each frequency $\omega$ in a swept band, and the system
operator \emph{depends on frequency}:
\[
  \mat{A}(\omega) \;=\; \Kop + i\omega\,\Cop - \omega^2\Mop.
\]
So before each solve, inside the sweep, the operator must be rebuilt at the current
$\omega$. The solves are still independent---the answer at $\omega_3$ does not need
the answer at $\omega_2$---so this is still a \code{map}, not a fold. What
distinguishes it from \code{solve\_family} is precisely that the operator is rebuilt
\emph{inside} the loop rather than hoisted outside it. The combinator is
\code{frequency\_sweep}:
\begin{lstlisting}[style=l4style]
frequency_sweep :: OperatorFamily -> [Omega] -> [SimState]
frequency_sweep fam omegas =
  map (\w -> let aw = assemble_frequency_operator fam w   -- REBUILD inside the map
             in  ksp_solve aw (rhs_at w))
      omegas
\end{lstlisting}
The basis $\{\Kop,\Cop,\Mop\}$ is captured once; the per-member operator
$\mat{A}(\omega)$ is rebuilt per member. The contrast with \code{solve\_family} is
a single axis---is the operator captured once and hoisted, or rebuilt per
element?---and it is the entire structural difference between electrostatics and
the driven analysis.

\begin{pitfall}
\code{frequency\_sweep} and \code{solve\_family} differ by the \emph{position of
one setup call} relative to the map: outside the loop (hoist) versus inside it
(rebuild). This is not an optimization knob. An implementation that ``speeds up''
\code{frequency\_sweep} by hoisting the operator construction out of the map does
not run faster---it computes the \emph{wrong answer}, because it solves every
frequency against the operator of the first one. The hoist is correct exactly when
the operator is fixed and a bug exactly when it varies. The combinator names which
case you are in.
\end{pitfall}

\subsection{\texttt{fold\_solve}: the state-threaded fold}

Now drop the independence. In \code{solve\_family} and \code{frequency\_sweep} the
solves were a map: independent, reorderable, parallelizable. In the third corner
they are a \code{fold}: each step begins from the state the previous step produced.
The transient analysis (\cref{ch:transient}) marches forward in time---the field
state at $t_{k+1}$ is computed from the state at $t_k$ by one time step---and the
state is \emph{threaded} through the chain. The combinator is \code{fold\_solve},
a \code{foldl} of an opaque per-step operator over a schedule:
\begin{lstlisting}[style=l4style]
fold_solve :: OpParams -> TimeState -> [Time] -> TimeState
fold_solve op s0 schedule = foldl (\s t -> time_step_op op s t) s0 schedule
\end{lstlisting}
The signature alone tells the story. The seed \code{s0} and the result are both a
\code{TimeState}; the step \lstinline[style=l4style]|time_step_op op s t| reads
\code{s}, the prior carry. That read is what makes this a fold and not a map: it
forbids reordering at the type level. The adaptive-mesh-refinement outer loop of
\cref{ch:situating} is also a \code{fold\_solve} of this shape---each refinement
pass begins from the mesh the last pass produced---which is why the same combinator
serves both transient time-stepping and the AMR loop.

Two things are worth pausing on. First, the per-step operator
\lstinline[style=l4style]|time_step_op| is, for the transient analysis,
\emph{opaque}: it bottoms out in a library time-integrator (an implicit Runge--Kutta
or generalized-$\alpha$ step) that the combinator only quantifies over, rather than
rendering. \code{fold\_solve} names the fold structure and the carry-threading; what
one step \emph{does} is a separate concern, deferred to \cref{ch:transient}. The
combinator is honest about the boundary the way \code{eigsolve} is, but it keeps the
loop in our hands---we own the fold, the library owns the step.

Second, the laws \emph{invert} relative to the map. \code{fold\_solve} obeys
neither the concatenation law nor the order law: splitting the schedule and folding
the pieces independently is wrong (the second piece must start from where the first
left off, not from \code{s0}), and permuting the schedule changes the answer
outright. What it obeys instead is a \emph{decomposition} law---folding a
concatenated schedule equals folding the first part and then folding the second
\emph{from the carry the first produced}:
\begin{lstlisting}[style=l4style]
fold_solve op s0 (a ++ b) = fold_solve op (fold_solve op s0 a) b
\end{lstlisting}
This is the algebraic signature of sequentiality: the only way to combine two folds
is to thread the carry from one into the other. Compare \code{solve\_family}, whose
concatenation law combines two maps by simply \emph{appending} their outputs with no
threading at all. The two laws---append versus thread-through---are the map/fold
split written algebraically, and they are why a map distributes over a cluster while
a fold marches on a single timeline.

\subsection{\texttt{eigsolve}: the corner that is not a loop}

The fourth corner is different in kind. In the first three \emph{we} write the
loop---the map or the fold---and can see every solve it performs. In the fourth we
write no loop at all. The eigenmode analysis (\cref{ch:eigenvalues}) assembles its
operator pencil once and hands the whole problem to a specialized eigenvalue
library (SLEPc, in Palace's case) with a single call. The iteration that produces
the answer runs \emph{inside} the library, where our framework cannot see it. The
combinator \code{eigsolve} is therefore deliberately \emph{opaque}: it names the
library boundary and the contract across it and stops. It is not a relative of
\code{iterate\_while}; it is the marked absence of a loop we own. (\Cref{ch:eigenvalues}
separately reconstructs what the library does internally, in our own vocabulary,
turning the black box into a glass box for the curious reader---but the
\emph{combinator} that drives the eigenmode analysis is just the single opaque
call.)

\begin{keyidea}
The four solve corners \emph{are} four combinators. \code{solve\_family} is the
fixed-operator map (assemble once, map independent solves). \code{frequency\_sweep}
is the operator-varying map (rebuild the operator inside the map).
\code{fold\_solve} is the state-threaded fold (chain the solves, each from the
last). \code{eigsolve} is the opaque black-box (no loop we write). Three are
relatives of \code{iterate\_while}, split by the map/fold axis and the
fixed/varying axis; the fourth is the deliberate non-loop. Every driver in
\cref{part:modalities} is one of these four over the \texttt{assemble} $\to$
\texttt{solve} $\to$ \texttt{reduce} skeleton.
\end{keyidea}

\section{The data-algebra combinators}
\label{sec:dataalgebra}

The solve-coordination combinators of \cref{sec:corners} say \emph{how} a family of
solves is structured. Inside each solve, the step body \emph{combines tensors}:
it scales a search direction and adds it to the iterate, it forms the inner product
that computes a step length, it measures a residual norm. These combinations are
the \emph{data-algebra} combinators, the right-hand branch of \cref{fig:relatives}.
There are only two that matter, and each illustrates the chapter's thesis in
miniature: one combinator, several named specializations.

\subsection{\texttt{linear\_combination}: the tensor-sum combinator}

The first is \code{linear\_combination}, the scalar-weighted tensor sum. Given a
list of (coefficient, tensor) pairs, it forms $\sum_i a_i\,t_i$:
\begin{lstlisting}[style=l4style]
linear_combination :: [(Scalar, Tensor)] -> Tensor
linear_combination pairs = foldl (\acc (a, t) -> acc + scal a t) zeros pairs
\end{lstlisting}
It is a \code{foldl} over the term list, accumulating $a_i\,t_i$ into a running
tensor that starts at zero. Every term must be \emph{congruent}---the same shape---
because the sum is element-local; the result has that shape too. This is the
single most-used data combinator in the book: a Krylov step's update
$\vx \gets \vx + \alpha\,\vp$ is a \code{linear\_combination}; GMRES's basis
correction $\sum_j y_j\,\vv_j$ is a \code{linear\_combination}; the driven
operator $\Kop + i\omega\Cop - \omega^2\Mop$ is a \code{linear\_combination} of
\emph{operators} rather than vectors.

It satisfies the laws you would hope for, and the one that earns its keep is the
\emph{concatenation-homomorphism}:
\[
  \texttt{linear\_combination}\,(a \mathbin{+\!\!+} b)
    \;=\; \texttt{linear\_combination}\,a \;+\; \texttt{linear\_combination}\,b.
\]
Splitting the term list and summing the pieces gives the same tensor as summing the
whole. In the language of \cref{ch:language} this makes \code{linear\_combination} a
\emph{monoid homomorphism}: it carries the monoid of term lists (under
concatenation, with the empty list as unit) to the monoid of tensors (under
addition, with the zero tensor as unit). That single structural fact organizes the
rest of the verb's algebra. The empty list maps to the unit:
\lstinline[style=l4style]|linear_combination [] = zeros|. A zero coefficient drops
its term: \lstinline[style=l4style]|linear_combination ((0,t):rest) = linear_combination rest|
(the algebraic content of the implementation's ``skip the term if its coefficient is
zero'' branch). A scalar can move from the coefficient onto the tensor:
$\code{linear\_combination}\,((\kappa a, t):\text{rest}) = \code{linear\_combination}\,((a, \kappa t):\text{rest})$.
And over exact arithmetic the sum is permutation-invariant---reordering the term
list leaves the result unchanged---because tensor addition is commutative and
associative. (Under floating-point arithmetic that last law holds only
approximately, since summation order perturbs the rounding; the framework treats
the left-to-right fold order as canonical and records the reordering subtlety as a
load-bearing numerical note rather than an exact law, in the spirit of
\cref{ch:bigidea}'s distinction between transparent and load-bearing tricks.)

The concatenation-homomorphism is also the law that makes the four named BLAS
routines \emph{one} combinator, as we see next: it lets a longer fixed-arity sum be
built from shorter ones, which is exactly the subsumption chain of the arity leaves.

\subsection{The BLAS verbs are arity leaves}

Numerical libraries ship a family of fused routines for small fixed-length sums,
the BLAS-1 ``\code{axpy}'' family. Each is \code{linear\_combination} at one fixed
term-list length:
\begin{lstlisting}[style=l4style]
scal(a, x)                 = linear_combination [(a, x)]                  -- arity 1
axpy(a, x, y)              = linear_combination [(a, x), (1, y)]          -- arity 2, 2nd coeff = 1
axpby(a, x, b, y)          = linear_combination [(a, x), (b, y)]          -- arity 2, general
axpbypcz(a, x, b, y, c, z) = linear_combination [(a, x), (b, y), (c, z)]  -- arity 3
\end{lstlisting}
Read top to bottom: \code{scal} scales one tensor (a one-term list); \code{axpy}
computes $a\vx + \vy$ (a two-term list whose second coefficient is fixed to $1$);
\code{axpby} computes $a\vx + b\vy$ (the general two-term list); \code{axpbypcz}
computes $a\vx + b\vy + c\vz$ (the three-term list). They differ \emph{only} in the
length of the term list. The concatenation law makes the subsumption exact:
\code{axpbypcz}'s three-term list is the concatenation of \code{axpby}'s two-term
list with \code{scal}'s one-term list, so $\code{scal} \prec \code{axpy} \prec
\code{axpby} \prec \code{axpbypcz}$ is the bounded-arity shadow of the one
combinator.

These fused routines exist because a hand-written kernel for ``$a\vx + b\vy$''
touches memory once instead of twice and so runs faster on a CPU. That is a real
performance win, but it is \emph{not} an abstraction: \code{axpby} adds no meaning
that \code{linear\_combination} lacks. So the framework makes a deliberate choice.
The accelerated kernels are ``stopped low''---kept as performance notes at the
implementation level---while the \emph{combinator rises} to the top of the
vocabulary in their place. \Cref{fig:lincomb} draws this: \code{linear\_combination}
is the trunk; the four BLAS routines hang off it as arity leaves, labels for fixed
list-lengths, not combinators in their own right.

The same combinator generalizes past vectors. Its terms need not be column tensors;
they may be \emph{operators}. The driven analysis's frequency-dependent operator
$\mat{A}(\omega) = \Kop + i\omega\,\Cop - \omega^2\Mop$ is a
\code{linear\_combination} of the three assembled operators $\Kop, \Cop, \Mop$ with
the $\omega$-dependent coefficients $1, i\omega, -\omega^2$. This is the same fold,
the same concatenation law, the same arity story---only the term type has changed
from tensor to operator. When \cref{ch:driven} writes ``rebuild $\mat{A}(\omega)$''
inside the \code{frequency\_sweep} map, the rebuild it means is one
\code{linear\_combination} over operator terms. The two combinator branches of
\cref{fig:relatives} meet here: the operator-varying map calls the tensor-sum
combinator to do its per-step rebuild. Recognizing the operator assembly as ``the
same verb, operator-valued terms'' rather than a new primitive is the
combinator-primary model paying off inside the framework itself.

\begin{figure}[t]
\centering
\begin{tikzpicture}[font=\footnotesize, node distance=5mm]
  % trunk
  \node[stage, minimum width=64mm, align=center] (trunk) at (0,0)
    {\code{linear\_combination} \;\itshape\quad $\sum_i a_i\,t_i$ \;\;(term list, any length)};
  % four leaves
  \def\ly{-1.9}
  \node[opbox, minimum width=24mm, align=center] (l1) at (-4.8,\ly)
    {\code{scal}\\\scriptsize arity 1\\\scriptsize $a x$};
  \node[opbox, minimum width=24mm, align=center] (l2) at (-1.6,\ly)
    {\code{axpy}\\\scriptsize arity 2\\\scriptsize $a x + y$};
  \node[opbox, minimum width=24mm, align=center] (l3) at (1.6,\ly)
    {\code{axpby}\\\scriptsize arity 2\\\scriptsize $a x + b y$};
  \node[opbox, minimum width=24mm, align=center] (l4) at (4.8,\ly)
    {\code{axpbypcz}\\\scriptsize arity 3\\\scriptsize $a x + b y + c z$};
  \foreach \l in {l1,l2,l3,l4}{
    \draw[refedge] (trunk.south) -- (\l.north);}
  \node[font=\scriptsize\itshape, text=simGray, align=center] at (0,-3.1)
    {fused accelerated kernels --- ``stopped low,'' specialization \emph{notes}\\
     the \emph{combinator} rises; the leaves are labels for fixed term-list lengths};
\end{tikzpicture}
\caption[\texttt{linear\_combination} and its arity leaves]{The combinator-primary
picture of the tensor-sum vocabulary. \code{linear\_combination} (trunk) is the
one combinator: a sum over a term list of any length. The four named BLAS routines
(leaves) are its fixed-arity specializations---\code{scal} at one term,
\code{axpy}/\code{axpby} at two, \code{axpbypcz} at three. They are
performance-fused kernels with no abstraction value of their own, so they are
``stopped low'' as specialization notes while the combinator rises to the top of
the vocabulary. One combinator, four leaves. \Cref{wex:arity} traces the
derivation with numbers.}
\label{fig:lincomb}
\end{figure}

\subsection{\texttt{inner\_product}: the scalar-producing fold, and its leaves}

The second data combinator is \code{inner\_product}, the scalar-producing sibling
of \code{linear\_combination}. Where \code{linear\_combination} folds a term list
into a \emph{tensor}, \code{inner\_product} folds two tensors (through an operator
$M$) into a \emph{scalar}, $\langle \vx, \vy\rangle_M = \vx^{\!\top} M \vy$. It is a
\emph{fold} in the same sense the running sum of \cref{ch:fold} was: it walks the
two operands position by position, accumulating a running scalar
$\sum_k (M\vy)_k\,\overline{x_k}$, and reduces the whole to a single number. The two
data combinators are thus the two halves of the L4 algebra of folds---one folds to
a tensor, the other folds to a scalar---and the step bodies of the iteration
combinators consume both: a Krylov step forms an \code{inner\_product} to compute a
step length and a \code{linear\_combination} to apply it.

\code{inner\_product} too carries a cluster of named leaves, but the cluster has a
subtlety worth drawing (\cref{fig:innerprod}).

\code{dot} is a genuine \emph{specialization}: it is \code{inner\_product} at the
special case $M = I$, the plain Euclidean (or Hermitian) inner product
$\vx^{\!\top}\vy$ that a solver spells \lstinline[style=l4style]|dot(p, Ap)|. Like
the BLAS arity leaves, it is the combinator at a fixed setting---here a fixed
operator argument rather than a fixed list length---and it hangs beneath the
combinator as a note.

\code{nrm2}, the Euclidean norm $\norm{\vx}_2 = \sqrt{\vx^{\!\top}\vx}$, is
\emph{not} a specialization of \code{inner\_product}. It is a \emph{consumer}: it
is built \emph{on top of} \code{dot} (itself a specialization), composing a square
root and an absolute value after the inner product, $\code{nrm2} =
\sqrt{\,\cdot\,} \circ \code{dot}$ at the diagonal. The distinction matters for the
combinator-primary bookkeeping: a \emph{leaf} is the combinator at a fixed setting
and lives beneath it; a \emph{consumer} calls the combinator and lives
\emph{above} it. Drawing \code{nrm2} as a leaf of \code{inner\_product} would
misfile a caller as a special case.

\begin{figure}[t]
\centering
\begin{tikzpicture}[font=\footnotesize, node distance=6mm]
  \node[stage, minimum width=50mm, align=center] (ip) at (0,0)
    {\code{inner\_product} \;\itshape\quad $\langle x, y\rangle_M = x^{\!\top} M y$};
  % dot: a specialization BELOW
  \node[opbox, minimum width=30mm, align=center] (dot) at (-1.6,-2.0)
    {\code{dot} \;\scriptsize (leaf)\\\scriptsize $M = I$: \;$x^{\!\top} y$};
  \draw[refedge] (ip.south) -- node[left,font=\scriptsize,pos=0.55]{\emph{specializes}} (dot.north);
  % nrm2: a CONSUMER above dot
  \node[product, minimum width=30mm, align=center] (nrm) at (-1.6,-4.0)
    {\code{nrm2} \;\scriptsize (consumer)\\\scriptsize $\sqrt{\;}\circ$ \code{dot}: \;$\|x\|_2$};
  \draw[depedge] (nrm.north) -- node[right,font=\scriptsize,pos=0.5]{\emph{consumes}} (dot.south);
  % annotation
  \node[font=\scriptsize\itshape, text=simGray, align=left, anchor=west] at (2.6,-2.0)
    {a \emph{leaf} is the combinator\\at a fixed setting ($M=I$):\\it lives \emph{below}};
  \node[font=\scriptsize\itshape, text=simGray, align=left, anchor=west] at (2.6,-4.0)
    {a \emph{consumer} \emph{calls} the\\combinator (here via \code{dot}):\\it lives \emph{above}};
\end{tikzpicture}
\caption[\texttt{inner\_product}: specialization versus consumer]{The
inner-product cluster, distinguishing two relationships the combinator-primary
model must keep separate. \code{dot} is a \emph{specialization} of
\code{inner\_product}---the combinator at the fixed operator $M = I$---and so hangs
\emph{below} it, exactly as the BLAS arity leaves hang below
\code{linear\_combination}. \code{nrm2} is not a special case but a \emph{consumer}:
it is built \emph{on} \code{dot}, composing a square root after the inner product,
and so lives \emph{above} it. A leaf is the combinator narrowed; a consumer is the
combinator called. Misdrawing \code{nrm2} as a leaf would file a caller as a
special case.}
\label{fig:innerprod}
\end{figure}

\begin{keyidea}
The data-algebra has two combinators. \code{linear\_combination} folds a term list
into a tensor, $\sum_i a_i t_i$; the BLAS routines \code{scal}, \code{axpy},
\code{axpby}, \code{axpbypcz} are its fixed-\emph{arity} leaves (fused kernels
stopped low, the combinator rising). \code{inner\_product} folds two tensors into a
scalar, $\langle x,y\rangle_M$; \code{dot} is its $M = I$ leaf, and \code{nrm2} is
a \emph{consumer} of \code{dot}, not a leaf. The pattern is uniform: \emph{one
combinator, several named specializations}, with a clean line between a
specialization (below) and a consumer (above).
\end{keyidea}

\begin{pitfall}
It is tempting to unify the two data combinators into one ``reduce a list of
tensors'' verb, since both are folds. Resist, for the same reason \cref{ch:situating}
kept the Gram and the projection reductions apart: sharing a \emph{shape} (both are
folds) is not sharing a \emph{meaning}. \code{linear\_combination} folds a list of
scalar-weighted tensors into a \emph{tensor} of the same shape; \code{inner\_product}
folds \emph{two} tensors through an operator into a \emph{scalar}. Their outputs
differ in rank, their inputs differ in arity, and their laws differ in kind
(\code{linear\_combination} is a monoid homomorphism over concatenation;
\code{inner\_product} is bilinear in its two arguments). Forcing them together would
either burden the tensor sum with a needless operator argument or strip the inner
product of its bilinearity. Two combinators that happen to both be folds is the
right design---the same judgment, applied to data verbs, that kept two reduction
verbs apart in \cref{ch:situating}.
\end{pitfall}

\section{The combinator-primary model}
\label{sec:primary}

We can now state the thesis the chapter has been assembling. Look back over the two
preceding sections and one shape recurs in every case. There is a \emph{combinator}
at the top---\code{solve\_family}, \code{fold\_solve}, \code{linear\_combination},
\code{inner\_product}---and beneath it a cluster of \emph{named operators} that are
each the combinator at some fixed setting: a fixed operator (\code{dot} at
$M = I$), a fixed term-list length (the BLAS arity leaves), a fixed family shape
(electrostatics as \code{solve\_family}). The combinator is the \emph{entry}; the
named operators are \emph{specialization notes} hung beneath it. This is the
\textbf{combinator-primary model}, and it is how the entire vocabulary of the
simulator is organized.

\begin{definition}[Combinator-primary organization]
\label{def:primary}
A vocabulary is \emph{combinator-primary} when its load-bearing entries are a small
set of higher-order \emph{combinators}---each carrying a signature, a set of
algebraic laws, and a structural meaning---and every other named operator is
recorded as a \emph{specialization} of some combinator: the combinator at a fixed
argument, a fixed arity, or a fixed family shape, written as a thin note that points
\emph{up} to its combinator and restates none of the combinator's structure. A
\emph{consumer} of a combinator (an operator that calls it) is distinguished from a
specialization (an operator that is the combinator narrowed) and is recorded
\emph{above} the combinator, not beneath it.
\end{definition}

\Cref{fig:primary} draws the model against its opposite. On the left, the
combinator-primary organization: one entry, many notes, every named operator
reachable as a setting of the combinator above it. On the right, the organization
the framework deliberately rejects: a flat pile of special-case kernels---
\code{scal}, \code{axpy}, \code{axpby}, \code{axpbypcz}, \code{dot}, \code{nrm2},
$n$ separate solver loops---each documented in full, with no unifying combinator
and so with the shared structure repeated $n$ times. The pile is not wrong, exactly;
it is \emph{unmanageable}. Every law must be restated per kernel; every reader must
rediscover that \code{axpby} and \code{axpbypcz} are the same idea at different
arities; every new special case adds a new full entry.

\begin{figure}[t]
\centering
\begin{tikzpicture}[font=\footnotesize]
  % ===== LEFT: combinator-primary =====
  \begin{scope}
    \node[font=\small\bfseries, text=simGreen] at (1.4,3.0) {combinator-primary (this book)};
    \node[stage, minimum width=40mm, align=center] (comb) at (1.4,2.0)
      {\textbf{combinator} \itshape (the entry)};
    \foreach \i/\x in {1/-0.5,2/1.4,3/3.3}{
      \node[opbox, minimum width=14mm, minimum height=6mm] (nt\i) at (\x,0.5) {};
      \draw[refedge] (comb.south) -- (nt\i.north);}
    \node[font=\scriptsize\itshape, text=simGray, align=center] at (1.4,-0.3)
      {specialization \emph{notes}\\(thin: a setting of the entry)};
  \end{scope}
  \draw[simGray!50, thick] (4.6,-0.7) -- (4.6,3.2);
  % ===== RIGHT: the pile =====
  \begin{scope}[xshift=5.6cm]
    \node[font=\small\bfseries, text=simRust] at (1.6,3.0) {the pile (rejected)};
    \foreach \i/\x/\y in {1/0/2.0,2/1.6/2.0,3/3.2/2.0,4/0/1.0,5/1.6/1.0,6/3.2/1.0}{
      \node[product, minimum width=13mm, minimum height=6mm] (pk\i) at (\x,\y) {};}
    \node[font=\scriptsize\itshape, text=simRust, align=center] at (1.6,-0.1)
      {$n$ full special-case kernels,\\no unifying combinator:\\shared structure repeated $n$ times};
  \end{scope}
\end{tikzpicture}
\caption[The combinator-primary model versus the pile]{The combinator-primary
model (left) and its rejected opposite (right). Combinator-primary: one combinator
is the \emph{entry}, carrying the signature, the laws, and the structure; the named
operators beneath it are thin specialization \emph{notes}, each a setting of the
combinator. The pile (right): a flat collection of special-case kernels, each a
full entry, with no combinator above to factor the shared structure---so the
structure is restated once per kernel and the collection grows without bound. The
left organization is why this book is short.}
\label{fig:primary}
\end{figure}

\subsection{Replace-and-propagate, not mine-and-strand}

The combinator-primary model is not just a way to \emph{present} the vocabulary;
it is a discipline for \emph{building} it. When a recurring pattern is spotted
across several operators---say, that \code{axpy} and \code{axpby} are both
two-term sums---there are two things one might do with the observation. The wrong
move is to \emph{mine and strand}: name the combinator, write it down in a corner,
and leave the original special-case kernels standing untouched beside it. Now there
are \emph{more} entries than before, not fewer: the kernels \emph{and} an unused
combinator. The right move is to \emph{replace and propagate}: make the combinator
the entry, rewrite each special case as a thin note that points up to it, and
propagate the change so that every consumer now reads through the combinator. The
count goes \emph{down}; the structure is stated \emph{once}.

\begin{keyidea}
\textbf{The combinator-primary model.} A small set of higher-order combinators are
the load-bearing vocabulary; each named operator is a thin \emph{specialization
note} beneath the combinator it instantiates. The combinator is the entry---it
carries the signature, the laws, the structure---and the leaves are settings of it.
The discipline is \emph{replace-and-propagate, not mine-and-strand}: a newly found
combinator \emph{replaces} the special cases (which become notes pointing up to it)
rather than standing as one more entry beside them. This is why the book can be
short: a few combinators and many specializations, with the shared structure stated
once rather than once per operator.
\end{keyidea}

\section{Why this matters downstream}
\label{sec:downstream}

The combinator-primary model is not an aesthetic preference; it is the load-bearing
beam of the rest of the book. Two later Parts rest on it directly, and it is worth
naming the dependence now so that those Parts read as payoffs rather than
surprises.

First, the \emph{drivers}. \Cref{part:modalities} studies the simulator's analyses
one at a time---electrostatic, magnetostatic, driven, transient, eigenmode,
boundary-mode. The combinator-primary model collapses what would be six
independent studies into six \emph{settings}. Every driver is one of
$\{\code{map}, \code{fold}, \code{single-call}\}$ over the
\texttt{assemble} $\to$ \texttt{solve} $\to$ \texttt{reduce} skeleton of
\cref{ch:situating}: electrostatics and magnetostatics are \code{solve\_family}
maps; the driven analysis is a \code{frequency\_sweep} map; the transient analysis
and the AMR loop are \code{fold\_solve} folds; the eigenmode and boundary-mode
analyses are single \code{eigsolve} calls. \Cref{fig:driverskel} previews this:
a driver is a combinator wrapped around the skeleton, with only the box contents
and the choice of combinator changing from analysis to analysis.

\begin{figure}[t]
\centering
\begin{tikzpicture}[font=\footnotesize, node distance=8mm]
  % the skeleton body
  \node[stage] (asm) at (0,0) {assemble};
  \node[stage, right=of asm] (solve) {solve};
  \node[stage, right=of solve] (red) {reduce};
  \draw[flow] (asm) -- node[above,font=\scriptsize]{$\Kop,\dots$} (solve);
  \draw[flow] (solve) -- node[above,font=\scriptsize]{field(s)} (red);
  % combinator wrapper
  \node[draw=simGreen, thick, rounded corners=3pt, fit=(asm)(solve)(red),
    inner sep=7mm, fill=simBgGreen, fill opacity=0.25] (wrap) {};
  \node[anchor=south west, font=\scriptsize\bfseries, text=simGreen]
    at (wrap.north west) {a combinator: \code{map} / \code{fold} / single-call};
  % the family/carry edge
  \draw[backflow, simGreen]
    (red.south) ++(0,-2mm) .. controls +(0,-7mm) and +(0,-7mm) ..
    (solve.south)
    node[midway, below, font=\scriptsize, text=simGreen]
      {next family member (map) \;or\; next carry (fold)};
  % in/out
  \node[data, left=10mm of wrap] (cfg) {config};
  \node[product, right=10mm of wrap] (prod) {product};
  \draw[flow] (cfg) -- (wrap.west);
  \draw[flow] (wrap.east) -- (prod);
\end{tikzpicture}
\caption[A driver as a combinator over the skeleton]{Every driver, previewed as a
combinator wrapped around the \texttt{assemble} $\to$ \texttt{solve} $\to$
\texttt{reduce} skeleton of \cref{ch:situating}. The combinator (green wrapper)
supplies the family structure---a \code{map} iterating over independent family
members, a \code{fold} threading a carry, or a single opaque call---and the
skeleton body supplies the per-member work. \Cref{part:modalities} fills in the box
contents one analysis at a time; what changes from driver to driver is only the box
contents and the choice of combinator. The combinator-primary model is what makes
six analyses six settings of this one picture.}
\label{fig:driverskel}
\end{figure}

To see how thoroughly the combinators carve up the six analyses, read down this
mapping---it is the punchline of \cref{ch:situating}'s grid restated in combinator
terms, and it is the skeleton of \cref{part:modalities}:
\begin{center}
\setlength{\tabcolsep}{8pt}
\renewcommand{\arraystretch}{1.2}
\begin{tabular}{@{}lll@{}}
\toprule
Driver & Combinator & Family structure \\
\midrule
Electrostatic & \code{solve\_family} & map over conductors (fixed $\Kop$) \\
Magnetostatic & \code{solve\_family} & map over current loops (fixed $\Kop$) \\
Driven & \code{frequency\_sweep} & map over $\omega$ (rebuild $\mat{A}(\omega)$) \\
Transient & \code{fold\_solve} & fold over timesteps (thread field state) \\
Eigenmode & \code{eigsolve} & single opaque call \\
Boundary mode & \code{eigsolve} & single opaque call (on a cross-section) \\
\bottomrule
\end{tabular}
\end{center}
Six drivers, four combinators, three structural shapes
($\code{map}$, $\code{fold}$, single-call). The physics that distinguishes the rows
lives in the \emph{box contents}---which operator is assembled, on which function
space, reduced by which verb---but the \emph{control structure} of each driver is
one of these four combinators wrapped around the skeleton. That is the entire reason
\cref{part:modalities} can present six analyses as six short variations rather than
six full programs.

Second, the \emph{synthesized library}. \Cref{part:collection} renders the whole
specification as though it were a real implementation library---the code the spec
describes, written out. That library \emph{is} these combinators. Its modules are
organized around them: an iteration module holding \code{iterate\_while} and the
step kernels, a data-algebra module holding \code{linear\_combination} and
\code{inner\_product} with their leaves, a coordination module holding
\code{solve\_family}, \code{frequency\_sweep}, and \code{fold\_solve}, and a driver
module composing them into the five analyses. The combinator-primary model is the
table of contents of that library. Learn the combinators here and you have learned
the spine of the synthesized code.

\section{Worked examples}

\begin{workedexample}{\texttt{solve\_family} as a map, \texttt{fold\_solve} as a foldl}{mapvsfold}
We write the two family combinators side by side on a concrete three-element
family and read off why one parallelizes and the other does not.

\emph{The map.} \code{solve\_family} runs three independent solves against one
fixed operator $\Kop$ and three right-hand sides $\vb_1, \vb_2, \vb_3$:
\begin{lstlisting}[style=l4style]
solve_family K [b1, b2, b3]
  = map (\b -> ksp_solve K b) [b1, b2, b3]
  = [ ksp_solve K b1            -- depends only on (K, b1)
    , ksp_solve K b2            -- depends only on (K, b2)
    , ksp_solve K b3 ]          -- depends only on (K, b3)
\end{lstlisting}
Each element names its own inputs $(\Kop, \vb_i)$ and \emph{nothing else}. No
element reads another element's output. So the three solves can be evaluated in any
order---$\vb_3$ first, $\vb_1$ last---or all at once on three machines, and the
collected list is the same either way. Concretely, with
$\Kop = \begin{psmallmatrix}2&0\\0&3\end{psmallmatrix}$ and
$\vb_1=\begin{psmallmatrix}2\\0\end{psmallmatrix},
\vb_2=\begin{psmallmatrix}0\\3\end{psmallmatrix},
\vb_3=\begin{psmallmatrix}2\\3\end{psmallmatrix}$, the three solves of
$\Kop\vx_i=\vb_i$ give $\vx_1=\begin{psmallmatrix}1\\0\end{psmallmatrix},
\vx_2=\begin{psmallmatrix}0\\1\end{psmallmatrix},
\vx_3=\begin{psmallmatrix}1\\1\end{psmallmatrix}$ independently---and indeed
$\vx_3 = \vx_1 + \vx_2$ only because $\vb_3 = \vb_1 + \vb_2$ and the solves are
linear, not because element three \emph{read} elements one and two. The map never
threads a carry.

\emph{The fold.} \code{fold\_solve} runs three \emph{chained} steps, threading a
state $s$ through a schedule $[t_1, t_2, t_3]$:
\begin{lstlisting}[style=l4style]
fold_solve op s0 [t1, t2, t3]
  = foldl (\s t -> time_step_op op s t) s0 [t1, t2, t3]
  = time_step_op op
      (time_step_op op
         (time_step_op op s0 t1)   -- produces s1
         t2)                       -- s2 needs s1
      t3                           -- s3 needs s2
\end{lstlisting}
Now read the nesting from the inside out. The innermost call produces $s_1$ from
$s_0$; the middle call needs $s_1$ to produce $s_2$; the outer call needs $s_2$ to
produce $s_3$. The step at $t_2$ \emph{cannot start} until the step at $t_1$ has
finished, because $s_1$ is literally its input. There is no order but
$t_1, t_2, t_3$, and no machine can take the second step early. With a toy
$\code{time\_step\_op}\ op\ s\ t = s + 1$ and $s_0 = 0$ the fold computes
$0 \to 1 \to 2 \to 3$---each value one more than the last, and unreachable except
by passing through its predecessor.

\emph{The contrast.} The two pieces of pseudocode differ by exactly one structural
fact: the map's element function names only its own inputs, while the fold's step
function names the prior carry $s$. That single carry dependence is the difference
between a fan-out and a chain (\cref{fig:mapfold}), between parallel and
sequential. It is why electrostatics (a \code{solve\_family} map) can solve all its
conductors at once while the transient analysis (a \code{fold\_solve} fold) must
march its timesteps in order. The map is the degenerate fold whose step would
ignore $s$; the fold is the one that uses it.
\end{workedexample}

\begin{workedexample}{One combinator, four leaves: deriving the BLAS family}{arity}
We take \code{linear\_combination} and recover each named BLAS routine as a fixed-
arity specialization, with numbers, so the ``one combinator, four leaves'' claim is
concrete rather than slogan. Throughout, let
$\vx = \begin{psmallmatrix}1\\2\end{psmallmatrix}$,
$\vy = \begin{psmallmatrix}3\\4\end{psmallmatrix}$,
$\vz = \begin{psmallmatrix}5\\6\end{psmallmatrix}$, and coefficients
$a = 2,\ b = 10,\ c = 100$.

\emph{Arity 1 ---} \code{scal}. The one-term list:
\begin{lstlisting}[style=l4style]
scal(a, x) = linear_combination [(a, x)] = zeros + scal a x = a*x
\end{lstlisting}
Numerically $\code{scal}(2, \vx) = 2\vx = \begin{psmallmatrix}2\\4\end{psmallmatrix}$.
The fold seeds at \code{zeros} and folds one term: $0 + 2\vx$.

\emph{Arity 2, second coefficient fixed ---} \code{axpy}. The two-term list whose
second coefficient is $1$:
\begin{lstlisting}[style=l4style]
axpy(a, x, y) = linear_combination [(a, x), (1, y)] = a*x + 1*y = a*x + y
\end{lstlisting}
Numerically $\code{axpy}(2, \vx, \vy) = 2\vx + \vy =
\begin{psmallmatrix}2\\4\end{psmallmatrix} + \begin{psmallmatrix}3\\4\end{psmallmatrix}
= \begin{psmallmatrix}5\\8\end{psmallmatrix}$.

\emph{Arity 2, general ---} \code{axpby}. The two-term list with both coefficients
free:
\begin{lstlisting}[style=l4style]
axpby(a, x, b, y) = linear_combination [(a, x), (b, y)] = a*x + b*y
\end{lstlisting}
Numerically $\code{axpby}(2, \vx, 10, \vy) = 2\vx + 10\vy =
\begin{psmallmatrix}2\\4\end{psmallmatrix} + \begin{psmallmatrix}30\\40\end{psmallmatrix}
= \begin{psmallmatrix}32\\44\end{psmallmatrix}$. Note \code{axpy} is the special
case $b = 1$: $\code{axpby}(a, \vx, 1, \vy) = \code{axpy}(a, \vx, \vy)$.

\emph{Arity 3 ---} \code{axpbypcz}. The three-term list:
\begin{lstlisting}[style=l4style]
axpbypcz(a, x, b, y, c, z) = linear_combination [(a, x), (b, y), (c, z)]
                           = a*x + b*y + c*z
\end{lstlisting}
Numerically $\code{axpbypcz}(2,\vx,10,\vy,100,\vz) = 2\vx + 10\vy + 100\vz =
\begin{psmallmatrix}2\\4\end{psmallmatrix}+\begin{psmallmatrix}30\\40\end{psmallmatrix}
+\begin{psmallmatrix}500\\600\end{psmallmatrix}=\begin{psmallmatrix}532\\644\end{psmallmatrix}$.

\emph{The concatenation law makes the subsumption exact.} Watch the three-term sum
split into a two-term and a one-term list:
\[
  \underbrace{\code{linear\_combination}\,[(a,\vx),(b,\vy),(c,\vz)]}_{\code{axpbypcz}}
  \;=\;
  \underbrace{\code{linear\_combination}\,[(a,\vx),(b,\vy)]}_{\code{axpby}}
  \;+\;
  \underbrace{\code{linear\_combination}\,[(c,\vz)]}_{\code{scal}}.
\]
Numerically $\begin{psmallmatrix}532\\644\end{psmallmatrix} =
\begin{psmallmatrix}32\\44\end{psmallmatrix} + \begin{psmallmatrix}500\\600\end{psmallmatrix}$:
the \code{axpbypcz} result is the \code{axpby} result plus the \code{scal} result.
The four named routines are not four ideas; they are one combinator read off at
list-lengths one, two, and three, and the concatenation-homomorphism is the law
that stitches them together. This is the combinator-primary model in a single
trace: the combinator carries the structure (the law), and the leaves are settings
of it.
\end{workedexample}

\summarysection
\Cref{ch:fold} built \code{iterate\_while}, the one iteration primitive; this
chapter widened the lens to the combinators that organize the whole simulator. The
governing axis is \emph{map versus fold}: a map runs a family of \emph{independent}
computations (no carry, so they parallelize), while a fold threads a \emph{carry}
through a chain (each step needs the last, so it serializes); a map is the
degenerate fold whose step ignores the carry. The four solve corners of
\cref{ch:situating} are exactly four combinators: \code{solve\_family} (the
fixed-operator map; electrostatic, magnetostatic), \code{frequency\_sweep} (the
operator-varying map, rebuilding the operator inside the loop; driven),
\code{fold\_solve} (the state-threaded fold; transient and AMR), and \code{eigsolve}
(the opaque black-box, no loop we write; eigenmode, boundary-mode). The
data-algebra adds two combinators: \code{linear\_combination}, the tensor-sum
$\sum_i a_i t_i$, whose fixed-arity leaves are the BLAS routines
\code{scal}/\code{axpy}/\code{axpby}/\code{axpbypcz} (fused kernels stopped low,
the combinator rising); and \code{inner\_product}, the scalar-producing fold, with
\code{dot} as its $M = I$ leaf and \code{nrm2} as a \emph{consumer} of \code{dot},
not a leaf. The thesis is the \emph{combinator-primary model}: the combinator is
the entry, carrying the signature and laws; the named operators are thin
specialization notes beneath it; the discipline is \emph{replace-and-propagate, not
mine-and-strand}. This is why a book about a vast simulator can be short---a few
combinators, many specializations---and it is the spine of both the drivers
(\cref{part:modalities}, each a combinator over assemble--solve--reduce) and the
synthesized library (\cref{part:collection}, which \emph{is} these combinators).

\furtherreading
The \code{map}/\code{fold} vocabulary, and the laws that make a fold-of-fixed-length
the same as its general fold, are foundational functional programming; Peyton
Jones's account of Haskell~\citep{peytonjones2003} develops \code{map}, \code{foldl},
and the homomorphism laws this chapter leans on for the arity leaves. The BLAS
routines whose fused special cases sit beneath \code{linear\_combination} and
\code{inner\_product}---the \code{axpy} family, \code{dot}, \code{nrm2}---and the
reason they are written as fused kernels rather than general sums are documented in
the numerical-linear-algebra literature; Golub and Van Loan~\citep{golubvanloan2013}
is the standard reference for the BLAS-1 level and its performance model, and
Saad~\citep{saad2003} shows how Krylov solvers are spelled in exactly these verbs,
which is the form \cref{ch:cg} and \cref{ch:gmres} build on.

\exercisessection
\begin{enumerate}[nosep]
  \item For each of the four solve corners, name its combinator and state the two
  binary classifications that place it (operator fixed/varying; map/fold or
  neither). Which corner answers ``neither row'' of the map/fold axis, and why?
  \item Electrostatics is \code{solve\_family} and the driven analysis is
  \code{frequency\_sweep}. Both are maps. Pinpoint the single line of pseudocode
  whose position---inside the map versus hoisted outside it---is the entire
  difference between them, and explain why moving it changes the \emph{answer}, not
  merely the speed.
  \item Write \code{axpby} and \code{axpy} as instances of \code{linear\_combination},
  then show that \code{axpy} is the special case of \code{axpby} with one
  coefficient fixed. Using the concatenation-homomorphism, express \code{axpbypcz}
  as a sum of an \code{axpby} and a \code{scal}, and verify it numerically with
  $\vx=\begin{psmallmatrix}1\\0\end{psmallmatrix},
  \vy=\begin{psmallmatrix}0\\1\end{psmallmatrix},
  \vz=\begin{psmallmatrix}1\\1\end{psmallmatrix}$ and coefficients $1, 2, 3$.
  \item Explain why \code{dot} is a \emph{specialization} of \code{inner\_product}
  but \code{nrm2} is a \emph{consumer} of \code{dot}, not a specialization of
  \code{inner\_product}. In the combinator-primary picture, which one hangs
  \emph{below} its combinator and which sits \emph{above}? What would go wrong if
  you drew \code{nrm2} as a leaf of \code{inner\_product}?
  \item Contrast \emph{replace-and-propagate} with \emph{mine-and-strand}. A
  colleague notices that \code{axpy} and \code{axpby} are both two-term sums, writes
  down a general \code{two\_term} combinator in a corner, and leaves the original
  \code{axpy} and \code{axpby} entries untouched. Has the entry count gone up or
  down? Which discipline did they follow, and what should they have done instead?
  \item A \code{map} is ``the degenerate fold whose step ignores the carry.'' Write
  \code{map f xs} as a \code{fold\_solve}-style \code{foldl} whose step discards its
  accumulator. Then argue from your rewrite why a map parallelizes and a fold does
  not---i.e., which data dependence the discarded carry would have introduced.
  \item Using \cref{fig:driverskel}, describe the transient analysis and the
  electrostatic analysis as ``a combinator wrapped around assemble--solve--reduce.''
  Which combinator does each use, what flows along the back-edge in each case (a
  family member or a threaded carry?), and which of the two could you run on a
  cluster of independent machines?
  \item \emph{(Synthesis.)} The chapter claims the book is short \emph{because} of
  the combinator-primary model. Suppose instead the simulator's vocabulary were
  written as ``the pile'' of \cref{fig:primary}: six independent driver studies,
  six special-case family loops, four BLAS kernels, two inner-product kernels, each
  a full entry. Estimate, in rough terms, how the page count would grow, and name
  two specific facts (one law, one structural relationship) that would have to be
  restated multiple times rather than once.
\end{enumerate}

\chapter{Modeling Mutation Purely: The State Monad}
\label{ch:monad}

\chapterorient{A solver does not compute one tensor; it \emph{evolves} a bundle of
state---the current iterate, the residual it is chasing, the count of how many
steps it has taken---across many passes of a loop. In a mutable language you would
keep that bundle in a struct and overwrite it. We cannot: every record in this book
is immutable (\cref{ch:records}), so each step must take the \emph{old} bundle and
return a \emph{new} one. Done by hand, threading that bundle through every call is
tedious clerical work, and clerical work is where bugs breed. This chapter
introduces the discipline that automates the threading and---this is the real
prize---makes the \emph{one place} the state actually changes stand out plainly on
the page. We call it the \emph{state monad}. We will not approach it as an abstract
algebraic structure; we will build it as what it is: careful bookkeeping for a
single evolving value.}

\section{The problem: threading one bundle by hand}
\label{sec:problem}

Recall the workhorse record from \cref{ch:records}: a \code{SimState}, the run-time
state a solve carries from step to step. It holds the current iterate, an iteration
counter, a convergence flag, perhaps a residual:

\begin{lstlisting}[style=l4style]
SimState = { x: Tensor[N], it: Int, converged: Bool, residual: Float }
\end{lstlisting}

\noindent A solve is a chain of steps, each producing the next \code{SimState} from
the current one---exactly the value-threading $\vs_0 \to \vs_1 \to \vs_2$ of
\cref{ch:bigidea}. Because records are immutable, ``produce the next state'' means
``read the old one and build a fresh one with the spread.'' A single step that
bumps the counter and records a new iterate looks like this:

\begin{lstlisting}[style=l4style]
step :: SimState -> SimState
step s =
  let x'  = improve s.x        in    -- compute a better iterate (a pure call)
  let r'  = residual_of x'     in    -- its residual          (a pure call)
  { ...s, x: x', it: s.it + 1, residual: r' }   -- the fresh state
\end{lstlisting}

\noindent So far, so clean: one step is a pure function \code{SimState -> SimState},
and we already know how to write it. The trouble starts when a step is not a single
expression but a \emph{sequence} of sub-operations, several of which each need to
read \emph{and} advance the state. Then the threading becomes manual, and the manual
threading is where the eye glazes over. Picture a step that must (i) read the current
iterate, (ii) advance it and bump the counter, (iii) read the just-advanced state to
compute a residual, and (iv) record that residual. Threaded by hand, every sub-step
must be handed the \emph{latest} state and must return the \emph{next} one, and the
programmer is personally responsible for never reusing a stale one:

\begin{lstlisting}[style=l4style]
step_by_hand :: SimState -> SimState
step_by_hand s0 =
  let s1 = { ...s0, x: improve s0.x }      in   -- s1 is s0 advanced
  let s2 = { ...s1, it: s1.it + 1 }        in   -- s2 is s1, counter bumped
  let s3 = { ...s2, residual: residual_of s2.x } in   -- s3 reads s2, records r
  s3
\end{lstlisting}

\noindent Read that carefully and you will see the hazard. There are four state
values in flight---\code{s0}, \code{s1}, \code{s2}, \code{s3}---and the program is
correct only because each line refers to the \emph{immediately preceding} one. Write
\lstinline[style=l4style]|s1.it| where you meant \lstinline[style=l4style]|s2.it|,
or thread \code{s0} into the residual line instead of \code{s2}, and the bug is
silent: the types still check (every \code{s} is a \code{SimState}), the program
still runs, and it computes the wrong thing. The numbers \code{s0}, \code{s1},
\code{s2}, \code{s3} are pure clerical overhead---they carry no meaning of their own;
they exist only to pass the baton from one line to the next. \Cref{fig:byhand} draws
the baton-passing we are about to automate away.

\begin{figure}[t]
\centering
\begin{tikzpicture}[font=\footnotesize, node distance=10mm]
  \node[data, minimum width=12mm] (s0) {\code{s0}};
  \node[opbox, right=9mm of s0] (o1) {improve};
  \node[data, minimum width=12mm, right=9mm of o1] (s1) {\code{s1}};
  \node[opbox, right=9mm of s1] (o2) {bump \code{it}};
  \node[data, minimum width=12mm, right=9mm of o2] (s2) {\code{s2}};
  \node[opbox, right=9mm of s2] (o3) {record \code{r}};
  \node[data, minimum width=12mm, right=9mm of o3] (s3) {\code{s3}};
  \draw[flow] (s0)--(o1); \draw[flow] (o1)--(s1);
  \draw[flow] (s1)--(o2); \draw[flow] (o2)--(s2);
  \draw[flow] (s2)--(o3); \draw[flow] (o3)--(s3);
  \node[font=\scriptsize\itshape, text=simRust, align=center] at (4.45,-1.1)
    {each baton must be handed to the \emph{next} line by hand---\\
     name the wrong one and the bug is silent};
\end{tikzpicture}
\caption[Threading state by hand]{The hand-threaded step of \cref{sec:problem}.
Four distinct state values \code{s0}\dots\code{s3} are passed line to line like a
relay baton; each operation reads the latest and emits the next. The names carry no
meaning---they are pure plumbing---and the program is correct only because every
line happens to reference the one just before it. The state monad of this chapter
makes the baton-passing implicit, so the plumbing names vanish.}
\label{fig:byhand}
\end{figure}

\begin{keyidea}
A solve evolves \emph{one} bundle of state through many steps. Immutability forbids
overwriting it, so each step must read the old bundle and return a fresh one---and
when a step has several sub-operations, threading that bundle by hand means inventing
a chain of meaningless names (\code{s0}, \code{s1}, \code{s2}, \dots) whose only job
is to carry the latest state to the next line. That bookkeeping is mechanical,
repetitive, and a quiet source of bugs. The \emph{state monad} automates it: you
write each sub-operation as if the state were simply ``there,'' and the monad threads
the right version between the lines for you.
\end{keyidea}

\section{\texttt{Solve} as a type: a recipe that threads a state}
\label{sec:solvetype}

Here is the central move, and we make it as concretely as we can. We give a name to
``a computation that threads a \code{SimState} and, along the way, yields some
result of type \code{a}.'' The name is \code{Solve}:

\begin{lstlisting}[style=l4style]
type Solve a = StateT SimState Identity a
\end{lstlisting}

\noindent Do not be alarmed by the right-hand side; we will unpack \code{StateT} and
\code{Identity} in a moment and then never need them again. Read the \emph{left} side
first, in plain words:

\begin{quote}\itshape
A value of type \code{Solve a} is a recipe that, when you give it a starting
\code{SimState}, produces a result of type \code{a} \emph{and} a new \code{SimState}.
\end{quote}

\noindent That is the entire idea. A \code{Solve a} is not a number, not a tensor,
not a record---it is a \emph{recipe} (a small function, if you like) waiting for a
state to run against. Run it against a state and two things come out: a value (the
\code{a}) and an updated state. Concretely, you can picture \code{Solve a} as
\emph{being} a function of this shape:

\begin{lstlisting}[style=l4style]
-- a Solve a is, under the hood, a function of this shape:
SimState -> (a, SimState)         -- give me a state; I return a value AND a new state
\end{lstlisting}

\noindent The fancy spelling \code{StateT SimState Identity a} is just the library
name for exactly that function shape---\code{StateT} reads ``state transformer over
the state type \code{SimState},'' and \code{Identity} is a do-nothing wrapper that
says ``no other effect is layered underneath; threading the state is all that
happens.'' We could have written the type as \code{State SimState a} and meant the
same thing; the \code{StateT \dots\ Identity} form is the spelling the calculus
inherits, and it commits us to nothing beyond the one-line reading above. From here
on, whenever you see \code{Solve a}, think only: \emph{a recipe that threads a
\code{SimState} and hands back an \code{a}}.

\begin{figure}[t]
\centering
\begin{tikzpicture}[font=\footnotesize]
  % the recipe box
  \node[stage, minimum width=34mm, minimum height=15mm, align=center] (box) at (0,0)
    {\code{Solve a}\\[1pt]\scriptsize\itshape a recipe};
  % input state on the left
  \node[data, minimum width=20mm, left=16mm of box] (sin) {\code{SimState} (in)};
  \draw[flow] (sin) -- (box);
  % two outputs on the right
  \node[product, minimum width=16mm] (val) at (4.4,0.9) {value : \code{a}};
  \node[data, minimum width=22mm] (sout) at (4.5,-0.9) {\code{SimState} (out)};
  \draw[flow] (box.east) -- (val.west);
  \draw[flow] (box.east) -- (sout.west);
  \node[font=\scriptsize\itshape, text=simGray, align=center] at (0,-1.55)
    {run it against a state\dots};
  \node[font=\scriptsize\itshape, text=simGray, align=center] at (4.5,-1.75)
    {\dots out comes a value \emph{and} a new state};
\end{tikzpicture}
\caption[\texttt{Solve a} as a state-threading recipe]{The plain reading of
\code{Solve a}. It is not a value sitting in memory; it is a \emph{recipe} that
consumes an incoming \code{SimState} and produces two things---a result of type
\code{a} and an outgoing \code{SimState}. The type \code{StateT SimState Identity a}
is the library spelling of exactly the function
\code{SimState -> (a, SimState)}; \code{Identity} merely says ``no other effect is
stacked underneath.'' Everything in this chapter is built from recipes of this one
shape and ways of gluing them together.}
\label{fig:solvetype}
\end{figure}

\begin{notationbox}
The name \code{Solve} is reused throughout the solver chapters, and the type
parameter \code{a} is the type of the \emph{result the recipe yields}, not the state
it threads. A \code{Solve Float} threads a \code{SimState} and yields a
\code{Float}; a \code{Solve Bool} threads the same \code{SimState} and yields a
\code{Bool}. The threaded state is \emph{always} a \code{SimState} (that is fixed in
the type \code{Solve a = StateT SimState Identity a}); only the yielded value varies.
A very common case is \code{Solve ()}---a recipe that threads the state but yields no
interesting value, written with the empty tuple \code{()} (``unit''). Such a recipe
exists \emph{only} for its effect on the state.
\end{notationbox}

\section{The primitives: \texttt{do}, \texttt{get}, \texttt{modify}, \texttt{return}}
\label{sec:primitives}

A recipe is useless until we can write one and run several in order. The state monad
gives us exactly four pieces of vocabulary, and they are all you need. We introduce
them one at a time, in the order you reach for them.

\paragraph{\code{do}: stack recipes top to bottom.} The keyword \code{do} sequences
recipes, top to bottom, exactly the way a list of imperative statements reads. Inside
a \code{do} block, each line is a recipe; the block runs them in order, threading the
state from each line into the next \emph{automatically}---this is the plumbing of
\cref{sec:problem}, now invisible. The block as a whole is itself a recipe (a bigger
\code{Solve}), so \code{do} blocks nest like any other expression.

\paragraph{\code{get}: read the current state.} The recipe \code{get} yields the
current \code{SimState} as its value and leaves the state unchanged. You bind its
result with the left-arrow \lstinline[style=l4style]|<-|:

\begin{lstlisting}[style=l4style]
do
  s <- get          -- s is now bound to the current SimState; state unchanged
  ...
\end{lstlisting}

\noindent After this line, the local name \code{s} \emph{is} the current state, and
you may dot into it (\lstinline[style=l4style]|s.it|, \lstinline[style=l4style]|s.x|)
like any record. Note the asymmetry with an ordinary \code{let}: a \code{let} binds
the value of a pure expression, while \lstinline[style=l4style]|s <- get| binds the
value \emph{a recipe yields}---here, the state itself. The left-arrow is how you reach
inside a recipe and name what it produced.

\paragraph{\code{modify f} / \code{put s}: write the state.} The recipe
\lstinline[style=l4style]|modify f| applies the pure function \code{f} to the current
state and makes the result the new state; it yields no interesting value (its type is
\code{Solve ()}). This is the \emph{single most important primitive in the chapter},
because \code{modify} is the one operation that \emph{changes} the threaded state.
The function \code{f} is an ordinary, pure \code{SimState -> SimState}---typically a
spread:

\begin{lstlisting}[style=l4style]
modify (\s -> { ...s, it: s.it + 1 })   -- advance the state: bump the counter
\end{lstlisting}

\noindent Read it: ``replace the current state by a fresh one, like it but with
\code{it} incremented.'' The sibling \lstinline[style=l4style]|put s'| is the blunt
version: it discards the current state and installs \code{s'} wholesale. We will
almost always prefer \code{modify}, because \code{modify} \emph{names the change}
(the spread says exactly which field moved) while \code{put} hides it inside a
fully-rebuilt record. Indeed
\lstinline[style=l4style]|modify f = do { s <- get; put (f s) }|---\code{modify} is
just ``get, apply \code{f}, put,'' bundled into one honest verb.

\paragraph{\code{return a}: yield a value, change nothing.} The recipe
\lstinline[style=l4style]|return a| yields the value \code{a} and leaves the state
untouched. It is how a \code{do} block produces a final result, and how a pure value
is lifted into a recipe so it can sit in a \code{do} block beside the real effects.
\lstinline[style=l4style]|return| is the ``do nothing to the state, just hand back
this value'' recipe.

\begin{notationbox}
The four primitives partition cleanly by what they touch.
\lstinline[style=l4style]|get| \emph{reads} the state (yields it, changes nothing);
\lstinline[style=l4style]|modify f| and \lstinline[style=l4style]|put s| \emph{write}
the state; \lstinline[style=l4style]|return a| touches the state not at all (it only
yields a value). Binding a recipe's yielded value uses the left-arrow
\lstinline[style=l4style]|x <- recipe|; binding a pure expression uses
\lstinline[style=l4style]|let x = expr|. The two binding forms living side by side in
one \code{do} block is the whole texture of monadic code, and \cref{sec:effectpoint}
makes the distinction between them load-bearing.
\end{notationbox}

\subsection{Walking a tiny \texttt{do} block, line by line}
\label{sec:walkdo}

Let us slow all the way down and watch the hidden state thread through a three-line
\code{do} block. The block reads the state, advances the iterate and counter, and
yields the new iteration number:

\begin{lstlisting}[style=l4style]
bump :: Solve Int
bump = do
  s <- get                                     -- (1) read the state into s
  modify (\s -> { ...s, x: improve s.x,        -- (2) write: advance x and it
                       it: s.it + 1 })
  s' <- get                                    -- (3) re-read the advanced state
  return s'.it                                 -- (4) yield the new iteration count
\end{lstlisting}

\noindent The block is a recipe of type \code{Solve Int}: run it against a starting
state and it yields an \code{Int} (the new counter) and leaves behind an advanced
state. \Cref{fig:walkdo} draws the hidden state threading between the lines. Trace it
against a concrete starting state
\lstinline[style=l4style]|s0 = { x: x0, it: 7, ... }|:

\begin{enumerate}[label=(\arabic*), leftmargin=*, nosep]
  \item \lstinline[style=l4style]|s <- get| yields the current state \code{s0} and
  binds it to \code{s}; the threaded state is still \code{s0}.
  \item \lstinline[style=l4style]|modify (...)| replaces the state by the spread
  applied to \code{s0}---a fresh \lstinline[style=l4style]|{ x: improve x0, it: 8, ... }|.
  The threaded state is now this new record; call it \code{s1}.
  \item \lstinline[style=l4style]|s' <- get| yields the \emph{current} state---which
  is now \code{s1}, the advanced one---and binds it to \code{s'}.
  \item \lstinline[style=l4style]|return s'.it| yields
  \lstinline[style=l4style]|s'.it|, which is $8$, and leaves the state as \code{s1}.
\end{enumerate}

\noindent The crucial thing to see is what happened \emph{between} the lines. Nowhere
did we write \code{s0}, \code{s1} as plumbing names; the \code{do} block threaded
them for us. Line (1)'s \code{get} saw the state before the \code{modify}; line (3)'s
\code{get} saw the state \emph{after} it. Same word, \code{get}, two different
values---because the \code{do} block silently advanced the threaded state across the
\code{modify} in between. That silent advance is precisely the baton-passing of
\cref{fig:byhand}, now performed by the monad instead of by you.

\begin{figure}[t]
\centering
\begin{tikzpicture}[font=\footnotesize]
  % the four code lines as a column on the left
  \node[anchor=west, font=\ttfamily\scriptsize] (l1) at (-6.2,1.65) {s <- get};
  \node[anchor=west, font=\ttfamily\scriptsize] (l2) at (-6.2,0.55) {modify (...)};
  \node[anchor=west, font=\ttfamily\scriptsize] (l3) at (-6.2,-0.55) {s' <- get};
  \node[anchor=west, font=\ttfamily\scriptsize] (l4) at (-6.2,-1.65) {return s'.it};
  % the threaded-state spine on the right
  \node[data, minimum width=16mm] (st0) at (0,1.65) {\code{s0}};
  \node[data, minimum width=16mm] (st1) at (0,0.55) {\code{s0}};
  \node[data, fill=simBgRust, draw=simRust, minimum width=16mm] (st2) at (0,-0.55) {\code{s1}};
  \node[data, fill=simBgRust, draw=simRust, minimum width=16mm] (st3) at (0,-1.65) {\code{s1}};
  % thread arrows down the spine
  \draw[flow] (st0)--(st1);
  \draw[flow] (st1)-- node[right=1mm,font=\scriptsize,text=simRust]{\textbf{the one write}} (st2);
  \draw[flow] (st2)--(st3);
  % which line reads what
  \draw[refedge] (l1.east) -- node[above,font=\scriptsize]{reads} (st0.west);
  \draw[refedge, simRust] (l2.east) -- node[above,font=\scriptsize,text=simRust]{writes} (st1.west);
  \draw[refedge] (l3.east) -- node[above,font=\scriptsize]{reads} (st2.west);
  \draw[refedge] (l4.east) -- node[above,font=\scriptsize]{ignores} (st3.west);
  \node[font=\scriptsize\itshape, text=simGray, align=center] at (0,2.45)
    {hidden threaded state (you never name it)};
\end{tikzpicture}
\caption[The hidden state threaded between \texttt{do} lines]{The \code{do} block of
\cref{sec:walkdo} with its \emph{hidden} threaded state drawn on the right. The four
source lines sit on the left; the state spine on the right shows what value the
threaded \code{SimState} holds as control passes each line. The first \code{get}
(line 1) sees \code{s0}; the lone \code{modify} (line 2, in rust) is the single point
where the state changes, advancing \code{s0} to \code{s1}; the second \code{get}
(line 3) therefore sees the \emph{advanced} \code{s1}. You wrote none of these spine
values---the \code{do} block threaded them. Two \code{get}s, one word, two different
answers, because one \code{modify} sat between them.}
\label{fig:walkdo}
\end{figure}

\subsection{Why this is equivalent to the hand-threaded version}
\label{sec:equiv}

It is worth proving to yourself, once, that the \code{do} block of \cref{sec:walkdo}
and a hand-threaded version compute exactly the same thing---that the monad is
genuinely just hiding plumbing and adds no magic. \Cref{wex:equiv} does this side by
side. The upshot, which we state now and earn there, is a small desugaring rule: a
\code{do} block is shorthand for a chain of explicit ``run this recipe against the
current state, name what it yields, thread the new state into the rest'' steps. The
monad is sugar; underneath, the baton still gets passed---you just no longer have to
name it.

\section{The single effect point: in the monad versus outside it}
\label{sec:effectpoint}

We now reach the heart of the chapter, and the reason the state monad earns its place
in a book about \emph{describing} algorithms clearly. The discipline is this:

\begin{quote}\itshape
Put in the monad \emph{only} the operations that read or write the threaded
\code{SimState}. Everything else---every operator application, every dot product,
every scrap of ephemeral workspace---is a pure \lstinline[style=l4style]|let|-binding
\emph{outside} the monad.
\end{quote}

\noindent The payoff is enormous and immediate: because only state-touching
operations live in the monad, the \emph{one place the state changes}---the lone
\code{modify}---stands out on the page, surrounded by pure \code{let}s that visibly do
not touch it. The monad turns ``where does this solver's state actually mutate?'' from
a question you answer by reading the whole loop into a question you answer by looking
for the word \code{modify}.

\subsection{A Krylov step: one \texttt{modify} amid a sea of pure \texttt{let}s}
\label{sec:krylovstep}

The cleanest illustration is a single step of a Krylov solver (the family
\cref{ch:operators} and the solver chapters build in full). A Krylov step does a
\emph{lot}: it applies the system operator to the current iterate, takes one or two
inner products, forms a scaled combination, updates a search direction, recomputes a
residual. Naively you might expect all of that to be ``in the monad''---it is, after
all, the step that advances the solve. The discipline says otherwise. Watch which
lines touch the \code{SimState}:

\begin{lstlisting}[style=l4style]
krylov_step :: OpParams -> SimState -> Solve ()
krylov_step op = \s_in -> do
  -- ===== OUTSIDE the monad: pure let-bindings, the actual numerical work =====
  let Ap   = apply op.A s_in.p          in   -- operator apply  (pure)
  let alpha = s_in.rdot / dot s_in.p Ap in   -- step length     (pure dot products)
  let x'    = axpy alpha s_in.p s_in.x  in   -- new iterate      (pure axpy)
  let r'    = axpy (-alpha) Ap s_in.r   in   -- new residual     (pure axpy)
  let rdot' = dot r' r'                 in   -- residual energy  (pure)
  let beta  = rdot' / s_in.rdot         in   -- direction coeff  (pure)
  let p'    = axpy beta s_in.p r'       in   -- new direction    (pure)
  -- ===== INSIDE the monad: the ONE state write =====
  modify (\s -> { ...s, x: x', r: r', p: p', rdot: rdot',
                       it: s.it + 1, converged: rdot' < op.tol })
\end{lstlisting}

\noindent Count the monadic operations: there is exactly \emph{one}, the closing
\code{modify}. Everything above it---the operator apply \code{Ap}, the inner products
\code{alpha} and \code{rdot'} and \code{beta}, the new iterate \code{x'}, residual
\code{r'}, and direction \code{p'}---is a pure \code{let}-binding. Those lines are the
algorithm's actual arithmetic, and \emph{none of them touches the state}. They read
fields out of \code{s\_in} (a plain value, captured once at the top) and compute fresh
tensors, exactly the pure-function vocabulary of \cref{ch:bigidea}. Only the final
\code{modify} folds the freshly computed quantities back into the threaded
\code{SimState}, advancing it by one step. \Cref{fig:partition} draws the partition.

\begin{figure}[t]
\centering
\begin{tikzpicture}[font=\footnotesize]
  % --- the "outside the monad" cloud of pure work ---
  \node[draw=simGray!55, fill=simBgGray, rounded corners=4pt, thick,
        minimum width=58mm, minimum height=30mm] (cloud) at (-2.6,0) {};
  \node[anchor=north, font=\scriptsize\bfseries, text=simGray]
    at ($(cloud.north)+(0,-1.5mm)$) {OUTSIDE the monad --- pure \code{let}s};
  \foreach \y/\lbl in {0.55/{Ap = apply A p}, -0.05/{alpha, beta = dot ...},
                       -0.65/{x', r', p' = axpy ...}}{
    \node[opbox, minimum width=42mm, font=\ttfamily\scriptsize] at (-2.6,\y) {\lbl};
  }
  \node[anchor=south, font=\scriptsize\itshape, text=simGray]
    at ($(cloud.south)+(0,1mm)$) {reads fields, computes fresh tensors; touches no state};
  % --- the single modify, inside ---
  \node[draw=simRust, fill=simBgRust, rounded corners=4pt, very thick,
        minimum width=34mm, minimum height=16mm, align=center] (mod) at (4.0,0)
    {\scriptsize\bfseries INSIDE the monad\\[2pt]\ttfamily\scriptsize modify (\textbackslash s -> \{...\})\\[2pt]\scriptsize\itshape the ONE state write};
  % the computed values flow into the modify
  \draw[flow] (cloud.east) -- node[above,font=\scriptsize]{\code{x', r', p', \dots}} (mod.west);
\end{tikzpicture}
\caption[In the monad versus outside it]{The partition of a Krylov step
(\cref{sec:krylovstep}). The grey cloud on the left is everything \emph{outside} the
monad: pure \code{let}-bindings that apply the operator, take inner products, and form
the new iterate, residual, and direction. They read fields off the incoming state
value and compute fresh tensors; not one of them touches the threaded \code{SimState}.
The rust box on the right is the \emph{single} monadic operation---one
\code{modify}---that folds those computed values back into the state, advancing it by
one step. The arithmetic is voluminous; the effect is a single point. That asymmetry,
made visible, is what the monad buys.}
\label{fig:partition}
\end{figure}

\begin{keyidea}
The state monad makes the single mutation point \emph{visible}. In a Krylov step the
only thing \emph{in} the monad is the state write---one \code{modify} that bumps the
counter and folds in the new iterate, residual, and direction. Everything else---the
operator applies, the inner products, the ephemeral workspace---is a pure
\code{let}-binding \emph{outside} the monad. The rule of thumb is exact: \textbf{if an
operation reads or writes \code{SimState}, it is in the monad; otherwise it is a pure
function call.} Because the discipline holds, you find every place a solver's state
changes by searching for one word, \code{modify}---and in a whole Krylov step there
is exactly one.
\end{keyidea}

\subsection{The rule of thumb, and why the counter is the canonical effect}
\label{sec:counter}

Apply the rule of thumb line by line to \cref{sec:krylovstep} and a striking fact
emerges. The genuinely \emph{state-bearing} field---the one whose entire reason for
living in \code{SimState} is that it must persist and be read across steps---is the
iteration counter \code{it}. The heavy tensors (\code{x}, \code{r}, \code{p}) are
recomputed each step from the previous step's values; they ride in the state for
convenience, but the quantity that genuinely \emph{accumulates}, that the loop's
predicate reads to decide whether to stop (\cref{ch:fold}), is the humble counter.
In the simplest possible step---one that does nothing but advance the loop---the
\emph{only} monadic effect is the counter bump:

\begin{lstlisting}[style=l4style]
modify (\s -> { ...s, it: s.it + 1 })   -- the canonical, irreducible state effect
\end{lstlisting}

\noindent This is the irreducible core of ``the solve mutated.'' We flagged it back
in \cref{ch:records}: a solve's externally visible state mutates at essentially a
\emph{single point}, the counter, while the heavy fields are recomputed elsewhere.
The monad is what cashes that observation in: it confines the mutation to one named
\code{modify} and pushes everything else out into pure \code{let}s, so the counter
bump---the one true effect---is the one thing left inside.

\begin{pitfall}
The cardinal sin of monadic code is \emph{over-monadifying}: dragging pure work
\emph{into} the monad that has no business there. It is tempting, having met
\code{do} and \code{get}, to write every line as a monadic step---to \code{get} the
state, compute \code{Ap = apply A s.p}, \code{put} a state with \code{Ap} stashed in
a scratch field, \code{get} again, and so on. Resist it utterly. An operator apply is
a pure function of its inputs; it does not read or write \code{SimState}, so it must
\emph{not} be in the monad. Putting it there does three bad things: it buries the one
real effect (the \code{modify}) in a crowd of fake ones, so the single mutation point
no longer stands out; it forces ephemeral workspace (\code{Ap}, \code{alpha}) into the
persistent state, where it does not belong (\cref{ch:strata}); and it defeats the
demand-driven pruning of \cref{ch:fold}, which can only eliminate \emph{pure}
unobserved work. The discipline is one-directional: pure work stays out; only
state reads and writes go in. When in doubt, ask the rule of thumb---\emph{does this
line read or write \code{SimState}?}---and if the answer is no, it is a \code{let}.
\end{pitfall}

\section{\texttt{execState}: the boundary where the monad evaporates}
\label{sec:execstate}

We have a way to build state-threading recipes and a discipline for keeping them
clean. One question remains: how do we \emph{run} one? A \code{Solve a} is a recipe,
not a value; to get an answer out we must supply a starting state and let the recipe
thread it through to the end. That is what \code{execState} does.

\begin{lstlisting}[style=l4style]
execState :: Solve a -> SimState -> SimState
\end{lstlisting}

\noindent Read the signature: \code{execState} takes a recipe (\code{Solve a}) and an
initial \code{SimState}, runs the recipe against it---threading the state through
every \code{modify} inside---and returns the \emph{final state}. (It is called
\code{exec} because it discards the recipe's yielded value \code{a} and keeps only the
final \emph{state}; a sibling \code{evalState} keeps the value and discards the state,
and \code{runState} keeps both. For a solve we want the final state, so
\code{execState} is our tool.) The entire monadic apparatus---the \code{do} blocks,
the \code{get}s and \code{modify}s, the hidden threading---is \emph{scaffolding that
exists only between} the initial state going in and the final state coming out.
Apply \code{execState} and the scaffolding evaporates: what is left is a plain
function from one \code{SimState} to another.

This is what lets the simulator stay pure on the outside. The top-level solve is
written:

\begin{lstlisting}[style=l4style]
solve :: OpParams -> Inputs -> SimState
solve op inp = execState (solve_loop op inp) (initial_state inp)
\end{lstlisting}

\noindent Look at the type of \code{solve}: \code{OpParams -> Inputs -> SimState}.
There is no \code{Solve}, no monad, no \code{StateT} anywhere in sight. Inside,
\code{solve\_loop op inp} is a \code{Solve ()} recipe---the whole iterative
computation, written monadically with its \code{do} blocks and \code{modify}s---and
\code{initial\_state inp} is the starting \code{SimState}. \code{execState} runs the
one against the other and hands back the final \code{SimState}. From outside,
\code{solve} is an ordinary pure function: feed it parameters and inputs, get a final
state. The monad lived entirely inside, did its bookkeeping, and was gone by the time
the answer crossed the boundary. \Cref{fig:execstate} draws the evaporation.

\begin{figure}[t]
\centering
\begin{tikzpicture}[font=\footnotesize]
  % the dashed monadic interior
  \node[draw=simViolet, fill=simBgViolet, rounded corners=4pt, densely dashed, thick,
        minimum width=46mm, minimum height=24mm] (mon) at (0,0) {};
  \node[anchor=north, font=\scriptsize\bfseries, text=simViolet]
    at ($(mon.north)+(0,-1.5mm)$) {inside: the \code{Solve} monad};
  \node[opbox, minimum width=34mm, font=\ttfamily\scriptsize] (a) at (0,0.45)
    {do \{ \dots\ modify \dots\ \}};
  \node[font=\scriptsize\itshape, text=simViolet, align=center] at (0,-0.55)
    {\code{get}s, \code{modify}s,\\hidden threading};
  % initial state in
  \node[data, minimum width=20mm] (s0) at (-5.0,0) {\code{initial\_state}};
  \draw[flow] (s0) -- (mon.west);
  % final state out
  \node[product, minimum width=20mm] (sf) at (5.0,0) {final \code{SimState}};
  \draw[flow] (mon.east) -- (sf.west);
  % the boundary label
  \node[font=\scriptsize\bfseries, text=simRust] at (-2.5,1.55) {\code{execState}};
  \node[font=\scriptsize\bfseries, text=simRust] at (2.5,1.55) {boundary};
  \node[font=\scriptsize\itshape, text=simGray, align=center] at (0,-1.95)
    {seen from outside: a pure \code{SimState -> SimState} --- the monad is gone};
\end{tikzpicture}
\caption[\texttt{execState} as the evaporation boundary]{\code{execState} runs a
\code{Solve} recipe against an initial state and returns the final state. The dashed
violet interior is the monadic computation---\code{do} blocks, \code{get}s,
\code{modify}s, the hidden threading of \cref{fig:walkdo}. \code{execState} is the
boundary: a \code{SimState} goes in on the left, the recipe threads it through every
internal effect, and a final \code{SimState} comes out on the right. From outside the
boundary there is no monad to be seen---\code{solve} has the plain type
\code{OpParams -> Inputs -> SimState}. The monad is scaffolding that does its work
inside and evaporates at the boundary, leaving the simulator pure on the outside.}
\label{fig:execstate}
\end{figure}

\begin{keyidea}
The monad is scaffolding that evaporates at the boundary. Inside \code{solve\_loop}
the state is threaded monadically, with \code{modify} marking each write; but
\code{execState (solve\_loop \dots) (initial\_state \dots)} runs that recipe against a
starting state and collapses the whole thing back into a \emph{pure} function
\code{SimState -> SimState}. The top-level \code{solve} has no monad in its type at
all. We get the local convenience of imperative-looking, automatically-threaded state
\emph{and} a globally pure simulator---the monad confined to the inside, gone by the
time any answer crosses out.
\end{keyidea}

\section{Why a monad and not just a fold}
\label{sec:whymonad}

A fair challenge: \cref{ch:fold} already gave us \code{iterate\_while}, which threads
a carry through a loop. If a solve is ``thread a \code{SimState} through steps until
converged,'' why not simply fold the \code{SimState} as the carry and be done? Why
introduce a monad at all?

The honest answer is that for the \emph{simplest} solvers you could. A non-restarted
solver with a single straight-line loop is perfectly well written as
\code{iterate\_while} with \code{SimState} as the carry. The monad earns its keep
when the control structure grows past a single flat loop, and three features push it
there.

\paragraph{Interior termination.} A real solver does not run a fixed number of steps;
it short-circuits the instant it detects convergence---and detection happens
\emph{mid-step}, on a residual the step just computed. Expressing ``compute the
residual, and if it is small enough, stop \emph{now}, before doing the rest'' is
exactly the kind of early exit a \code{do} block with a conditional makes plain at the
coordination layer. The fold's predicate, by contrast, fires only \emph{between}
steps (\cref{ch:fold}), so a pure fold must always finish a full step before it can
test; the monad lets the exit decision sit \emph{inside} the step where the residual
was born.

\paragraph{Inner and outer loop structure.} Restarted Krylov methods (GMRES is the
canonical case, built in \cref{ch:operators}) have an \emph{outer} cycle whose body is
itself an \emph{inner} loop. The inner loop reads \code{SimState.it} to honour the
global iteration cap and writes it back each step; the outer loop folds an accumulated
correction into \code{SimState.x} once per cycle. That is two-level read/write traffic
on one shared state, and the monad makes both levels explicit and the threading
automatic across the nesting. A single flat fold cannot express ``a loop whose step is
another loop'' without manually re-threading the carry through both.

\paragraph{Threading side-data.} The iterate \code{x} lives in \code{SimState}; the
Krylov basis lives in an ephemeral bundle reborn each restart; the correction step
\lstinline[style=l4style]|x += V*y| reads \emph{both}. The monad makes the moment the
\code{SimState} is touched a single named \code{modify}, while the ephemeral bundle
rides alongside as a plain \code{let}-bound value (it does not belong in the monad---it
does not persist across restarts). The monad cleanly separates ``the persistent state,
threaded'' from ``the per-cycle workspace, passed as a value.''

\noindent In every case the monad does the \emph{same} fundamental thing the fold
does---thread a value through steps---but it confines the threading of the
\emph{persistent} state to named \code{modify} points and lets the surrounding control
structure (early exit, nesting, side-data) be written plainly. And the connection to
\cref{ch:fold} is intact, not severed: the loop is \emph{still} \code{iterate\_while}.
The monad is simply what the carry threads \emph{through} when the carry is a
persistent simulator state and the control flow is richer than one flat loop. Fold and
monad are partners: the fold is the loop; the monad is the discipline for the state the
loop carries.

\begin{notationbox}
The three termination reasons a restarted solve can hit---converged on the residual
proxy, exhausted the global iteration cap, filled the per-cycle basis---collapse into
a single small sum type, \lstinline[style=l4style]@Outcome = Continue | Done Bool@.
The recipe \code{restart\_cycle} classifies the cycle's result once, at the boundary,
into an \code{Outcome}; \code{solve\_loop} pattern-matches on it and either recurses
or stops. This replaces a scatter of inner-loop break conditions with one decision
site---a clean payoff of writing the coordination monadically. The \code{Bool} inside
\code{Done} carries whether the stop was a convergence (\code{Done True}) or an
exhaustion (\code{Done False}), so the outer loop's final \code{modify} into
\code{SimState} is uniform.
\end{notationbox}

\section{Two faces of one name: the \texttt{Solve} effect versus the \texttt{Solve} record}
\label{sec:overload}

A reader who has been tracking the types carefully will have noticed something
slippery, and we must name it before it causes confusion. The word \code{Solve} is
\emph{overloaded}: it denotes two different things that happen to share a name.

\paragraph{\code{Solve} the effect.} This is the subject of the whole chapter so far:
\code{Solve a = StateT SimState Identity a}, the \emph{monad} that threads a
\code{SimState} and yields an \code{a}. It is a \emph{type constructor}---it takes a
type \code{a} (the yielded-value type) and produces the type of a recipe. ``\code{Solve}
the effect'' answers the question \emph{how is the state threaded?}

\paragraph{\code{Solve} the record.} When a single Krylov step is written out, the
\emph{value} it hands back is a record, and the calculus writes that record with the
same word: \lstinline[style=l4style]|Solve { sim, krylov, outputs }|. This is the
\code{SolveResult} bundle previewed in \cref{ch:records}: the next state (\code{sim}),
the next working bundle (\code{krylov}), and the demand-prunable per-step readouts
(\code{outputs}). ``\code{Solve} the record'' answers a different question: \emph{what
does one step yield?}

\noindent The two are genuinely distinct, and conflating them is the trap. The
\emph{effect} is the threading discipline; the \emph{record} is the shape of the value
yielded. They touch at exactly one field: the record's \code{sim} field is the one
that is discharged through the effect---it is the \code{SimState} the monad threads,
the product of the lone \code{modify}---while \code{krylov} and \code{outputs} are
plain returned values that ride alongside, untouched by the monad. So the rule of
thumb of \cref{sec:effectpoint} reappears at the level of the return record: of the
three fields the record hands back, \emph{one} (\code{sim}) is effectful and \emph{two}
(\code{krylov}, \code{outputs}) are pure. \Cref{fig:overload} disambiguates the two
faces.

\begin{figure}[t]
\centering
\begin{tikzpicture}[font=\footnotesize]
  % LEFT: Solve the effect
  \begin{scope}
    \node[font=\small\bfseries, text=simViolet] at (0,2.2) {\code{Solve} the \emph{effect}};
    \node[stage, minimum width=40mm, align=center] (eff) at (0,1.0)
      {\ttfamily\scriptsize Solve a = StateT SimState Identity a};
    \node[font=\scriptsize\itshape, text=simGray, align=center] at (0,0.0)
      {a type constructor:\\\emph{how} \code{SimState} is threaded};
    \node[font=\scriptsize, text=simInk, align=center] at (0,-1.0)
      {answers:\\``how is state threaded?''};
  \end{scope}
  \draw[simGray!50, thick] (3.4,-1.6) -- (3.4,2.4);
  % RIGHT: Solve the record
  \begin{scope}[xshift=7.0cm]
    \node[font=\small\bfseries, text=simBlue] at (0,2.2) {\code{Solve} the \emph{record}};
    \node[draw=simBlue, fill=simBgBlue, rounded corners=3pt, thick,
          minimum width=42mm, minimum height=18mm, inner sep=4pt] (rec) at (0,0.6) {};
    \node[anchor=north, font=\scriptsize\bfseries, text=simBlue] at ($(rec.north)+(0,-1mm)$)
      {\ttfamily Solve \{ ... \}};
    \node[anchor=west, font=\ttfamily\scriptsize, text=simRust]
      at ($(rec.north west)+(2.5mm,-6mm)$) {sim\ \ \itshape(effectful)};
    \node[anchor=west, font=\ttfamily\scriptsize]
      at ($(rec.north west)+(2.5mm,-10mm)$) {krylov, outputs\ \itshape(pure)};
    \node[font=\scriptsize, text=simInk, align=center] at (0,-1.0)
      {answers:\\``what does one step yield?''};
  \end{scope}
\end{tikzpicture}
\caption[The \texttt{Solve} overload disambiguated]{The two faces of \code{Solve}.
\emph{Left:} \code{Solve} the \emph{effect}, the monad
\code{StateT SimState Identity a}---a type constructor that answers \emph{how} the
\code{SimState} is threaded across the loop. \emph{Right:} \code{Solve} the
\emph{record}, the bundle \lstinline[style=l4style]|Solve { sim, krylov, outputs }| a
single step yields---a value that answers \emph{what} one step hands back. They meet
at one field: the record's \code{sim} (rust) is discharged \emph{through} the effect
(it is the threaded \code{SimState}), while \code{krylov} and \code{outputs} are plain
returned values the monad never touches. Same name, two questions; the \code{sim}
field is the bridge.}
\label{fig:overload}
\end{figure}

\begin{pitfall}
Do not read \lstinline[style=l4style]|Solve { sim, krylov, outputs }| as ``the monad
holds three pieces of state.'' It does not. The monad threads \emph{one} thing, a
\code{SimState}, and that one thing is the record's \code{sim} field alone. The
\code{krylov} and \code{outputs} fields are ordinary values the step returns
\emph{beside} the monadic effect---the working bundle (which is reborn each restart and
so deliberately is \emph{not} monadic state, lest its lifetime be misrepresented) and
the prunable readouts. When you see the word \code{Solve} in a signature, check which
face you are looking at: a bare \code{Solve a} (or \code{Solve ()}) is the effect; a
\lstinline[style=l4style]|Solve { ... }| with braces is the return record, and only its
\code{sim} field rides the effect.
\end{pitfall}

\section{A physical reprise: stream, then collide, threaded purely}
\label{sec:lbm}

To ground the whole chapter in something you can picture, we return to the
lattice-Boltzmann step we met in \cref{ch:bigidea}---stream, then collide---and show
it as the cleanest possible example of ``mutation modeled purely.'' It is the perfect
closing example precisely because, in its minimal form, it needs \emph{no} monad at
all: it is a pure \code{FluidState -> FluidState} function threaded by the plain fold
\code{iterate\_while\_pure} of \cref{ch:fold}. Seeing where the monad is \emph{not}
needed sharpens the sense of where it is.

Recall the rhythm. A lattice-Boltzmann fluid stores, at each cell of a grid, a little
tensor of \emph{populations}---how much fluid moves in each discrete direction. One
time step \emph{streams} (every population slides to the neighbor in its direction)
and then \emph{collides} (the populations in each cell relax toward local
equilibrium). In a classical code this is a nest of loops overwriting a population
array in place, often with an explicit scratch copy to avoid clobbering data a later
cell still needs:

\begin{lstlisting}[style=l4style, language=]
// classical C-style: overwrite in place, with a defensive copy
clone     = copy(current);              // guard against overwriting live data
streamed  = stream(clone);              // slide populations to neighbors
current   = collide(streamed);          // relax toward equilibrium, overwrite
\end{lstlisting}

\noindent The \lstinline[style=l4style, language=]|clone| is the tell. In the mutable
world you must defensively copy \code{current} before streaming, because streaming
reads a cell's neighbors and writing the result back in place would corrupt
neighbors not yet processed. The copy is pure ceremony forced by mutation. Now watch
the same step in our world, as a pure function building a fresh state:

\begin{lstlisting}[style=l4style]
step :: FluidState -> FluidState
step s =
  let streamed = stream  s.populations in    -- read s, produce a fresh tensor
  let collided = collide streamed       in    -- read it, produce another
  { ...s, populations: collided, t: s.t + 1 } -- a FRESH state; s untouched
\end{lstlisting}

\noindent There is no \code{clone}. There \emph{cannot} be a clobbering hazard,
because \code{stream} reads the immutable \code{s.populations} and \emph{returns} a
fresh tensor; the input it read can never change underneath it (\cref{ch:records}).
The defensive copy of the mutable version has \emph{evaporated}---in the immutable
world it was always a no-op in meaning, a workaround for a hazard that does not exist
here. That evaporation is the whole moral of \cref{ch:bigidea} made concrete: model
the mutation purely and the bookkeeping that mutation forced on you simply
disappears. \Cref{fig:lbm} traces a few steps.

\begin{figure}[t]
\centering
\begin{tikzpicture}[font=\footnotesize, node distance=8mm]
  % carry of FluidStates threaded by iterate_while_pure
  \node[data, minimum width=15mm, align=center] (f0) {\code{s}\(_0\)\\\scriptsize\(t=0\)};
  \node[opbox, right=10mm of f0, align=center] (st1) {stream\\then collide};
  \node[data, minimum width=15mm, align=center, right=10mm of st1] (f1) {\code{s}\(_1\)\\\scriptsize\(t=1\)};
  \node[opbox, right=10mm of f1, align=center] (st2) {stream\\then collide};
  \node[data, minimum width=15mm, align=center, right=10mm of st2] (f2) {\code{s}\(_2\)\\\scriptsize\(t=2\)};
  \node[font=\scriptsize, simGray, right=6mm of f2] (more) {$\cdots$};
  \draw[flow] (f0)--(st1); \draw[flow] (st1)-- node[above,font=\scriptsize]{fresh} (f1);
  \draw[flow] (f1)--(st2); \draw[flow] (st2)-- node[above,font=\scriptsize]{fresh} (f2);
  \draw[flow] (f2)--(more);
  % the evaporated clone
  \node[draw=simGray!45, densely dashed, fill=white, rounded corners=2pt,
        font=\scriptsize\itshape, text=simGray!70, align=center] (clone) at (1.7,-1.5)
    {\code{clone()}\\(no-op --- evaporates)};
  \draw[refedge, simGray!50] (clone.north) -- (st1.south);
  \node[font=\scriptsize\itshape, text=simGray, align=center] at (8.0,-1.5)
    {threaded by\\\code{iterate\_while\_pure}\\(no monad needed)};
\end{tikzpicture}
\caption[The LBM step threaded purely]{The lattice-Boltzmann step as a pure
\code{FluidState -> FluidState} function threaded by \code{iterate\_while\_pure}
(\cref{ch:fold}). Each step streams then collides the current populations and returns
a \emph{fresh} state with the time counter advanced; the chain
\code{s}\(_0 \to\) \code{s}\(_1 \to\) \code{s}\(_2\) is the value-threading of
\cref{ch:bigidea}. The defensive \code{clone()} of the mutable version (dashed, below)
is a \emph{no-op} here---it guarded against an in-place clobber that immutability makes
impossible, so it simply evaporates. Because the minimal step has no interior
termination, no nesting, and no persistent side-data beyond the carry itself, it needs
\emph{no monad}: the plain fold suffices.}
\label{fig:lbm}
\end{figure}

The lattice-Boltzmann step is the lower bound of the chapter's machinery: pure step,
plain fold, no monad. A Krylov solve sits at the upper bound: the same fold drives it
(\cref{ch:fold}), but now the carry is a persistent \code{SimState} with interior
termination and nesting, so the monad of this chapter manages the state the fold
carries, confining its one true mutation to a single \code{modify}. Between those two
poles live all the solvers of the book---and every one of them is the value-threading
of \cref{ch:bigidea}, dressed in exactly as much monadic bookkeeping as its control
structure demands, and no more.

\section{Worked examples}
\label{sec:worked}

\begin{workedexample}{One step, hand-threaded versus monadic---and they agree}{equiv}
We take a tiny two-effect computation and write it twice: once threading the state by
hand (passing \code{s} explicitly), once with the monad. Then we check that they
compute the same final state, so the monad is revealed as pure plumbing-hiding.

\textbf{The task.} From a starting state, (i) advance the iterate and bump the
counter, then (ii) record the residual of the advanced iterate. Two operations, each
touching the state.

\textbf{By hand.} Thread \code{s} explicitly, inventing names for each version:
\begin{lstlisting}[style=l4style]
by_hand :: SimState -> SimState
by_hand s0 =
  let s1 = { ...s0, x: improve s0.x, it: s0.it + 1 } in   -- (i) advance + bump
  let s2 = { ...s1, residual: residual_of s1.x }    in   -- (ii) read s1, record r
  s2
\end{lstlisting}
The correctness hinges on line (ii) reading \code{s1} (the advanced state), not
\code{s0}. The name \code{s1} is pure plumbing.

\textbf{With the monad.} Write each operation as a monadic step; the \code{do} block
threads the state:
\begin{lstlisting}[style=l4style]
monadic :: Solve ()
monadic = do
  modify (\s -> { ...s, x: improve s.x, it: s.it + 1 })   -- (i) advance + bump
  s' <- get                                               -- re-read the advanced state
  modify (\s -> { ...s, residual: residual_of s'.x })     -- (ii) record r
\end{lstlisting}

\tcblower
\textbf{Running the monadic version.} Apply \code{execState monadic s0} and thread by
the rules of \cref{sec:primitives}. The first \code{modify} replaces \code{s0} by
\lstinline[style=l4style]|{ ...s0, x: improve s0.x, it: s0.it + 1 }|---call it
\code{s1}, exactly the \code{s1} of the hand-threaded version. The
\lstinline[style=l4style]|s' <- get| binds \code{s'} to the \emph{current} state,
which is now \code{s1}. The second \code{modify} replaces \code{s1} by
\lstinline[style=l4style]|{ ...s1, residual: residual_of s'.x }|, and since
\code{s'} \emph{is} \code{s1}, that is
\lstinline[style=l4style]|{ ...s1, residual: residual_of s1.x }|---exactly the
hand-threaded \code{s2}. \code{execState} returns this final state.

\textbf{They agree.} The two final states are identical, character for character. The
difference is only that the hand-threaded version \emph{named} the intermediate
\code{s1} and had to reference it correctly by hand, while the monadic version let the
\code{do} block thread \code{s1} from the first \code{modify} into the \code{get} that
follows it. This is the desugaring promised in \cref{sec:equiv}: a \code{do} block is
shorthand for the explicit baton-passing, with the baton names supplied by the monad
instead of by you. The monad is sugar; the computation underneath is the value
threading of \cref{ch:bigidea}.
\end{workedexample}

\begin{workedexample}{The LBM step as a pure \texttt{FluidState -> FluidState}}{lbm}
We trace two lattice-Boltzmann steps numerically (schematically), watching the
\code{clone()} of the mutable version evaporate and the state thread as fresh
immutable values through \code{iterate\_while\_pure}.

\textbf{Setup.} A \code{FluidState} carries the population field and a time counter:
\begin{lstlisting}[style=l4style]
FluidState = { populations: Tensor[H, W, Q], t: Int }
\end{lstlisting}
where \code{Q} is the number of discrete velocity directions. The step streams then
collides; the driver folds it with the plain combinator:
\begin{lstlisting}[style=l4style]
run :: FluidState -> Int -> FluidState
run s0 nsteps =
  iterate_while_pure
    s0                              -- initial carry: the fluid state
    (\s -> s.t < nsteps)            -- predicate: pure on the carry (the time count)
    (\s -> let streamed = stream  s.populations in
           let collided = collide streamed       in
           { ...s, populations: collided, t: s.t + 1 })   -- fresh state, no clone
\end{lstlisting}

\tcblower
\textbf{The trace.} Take a schematic field where \code{stream} shifts populations and
\code{collide} nudges them toward equilibrium; we track only the time counter and the
\emph{identity} of the population tensor, to make the threading visible:
\begin{center}
\setlength{\tabcolsep}{8pt}
\renewcommand{\arraystretch}{1.15}
\begin{tabular}{@{}clll@{}}
\toprule
step & carry \code{t} & \code{streamed} (fresh) & result \code{populations} (fresh) \\
\midrule
0 & $0$ & $\mathrm{stream}(P_0)$ & $P_1 = \mathrm{collide}(\mathrm{stream}(P_0))$ \\
1 & $1$ & $\mathrm{stream}(P_1)$ & $P_2 = \mathrm{collide}(\mathrm{stream}(P_1))$ \\
\bottomrule
\end{tabular}
\end{center}
Each row reads the \emph{previous} population tensor ($P_0$, then $P_1$) and produces a
\emph{fresh} one ($P_1$, then $P_2$); the originals $P_0$, $P_1$ are never altered, so
no defensive copy is needed before streaming. In the mutable code,
\lstinline[style=l4style, language=]|clone = copy(current)| would sit at the top of
each step to protect $P_0$ from being clobbered while its neighbors are streamed; here
that copy is a \emph{no-op}---streaming an immutable $P_0$ cannot clobber it---so it
evaporates entirely.

\textbf{Why no monad.} Examine the step against the three pressures of
\cref{sec:whymonad}. There is no \emph{interior termination}: the loop runs a fixed
\code{nsteps}, tested between steps by the predicate. There is no \emph{inner/outer
nesting}: one flat loop. There is no persistent \emph{side-data} beyond the carry: the
\code{streamed} intermediate lives and dies inside one step as an ephemeral
\code{let}. All three pressures are absent, so the carry \emph{is} the entire state
and a plain fold suffices---no \code{Solve}, no \code{modify}, no \code{execState}.
This is the cleanest face of ``mutation modeled purely'': the same value-threading as
a Krylov solve, with none of the monadic bookkeeping, because the control structure
does not demand it. The monad is a tool you reach for exactly when the loop outgrows
this shape, and not a moment sooner.
\end{workedexample}

\summarysection
A solver evolves one bundle of state---the \code{SimState} holding the iterate,
counter, residual, and convergence flag---across many steps. Immutability forbids
overwriting it, so each step reads the old bundle and returns a fresh one; threading
that bundle by hand means inventing a chain of meaningless plumbing names (\code{s0},
\code{s1}, \dots) and is a quiet source of bugs (\cref{sec:problem}). The
\emph{state monad} automates the threading. A \code{Solve a = StateT SimState
Identity a} is a \emph{recipe} that, given a starting \code{SimState}, yields an
\code{a} and a new \code{SimState}---read plainly, the function
\code{SimState -> (a, SimState)} (\cref{sec:solvetype}). Four primitives suffice:
\code{do} stacks recipes top to bottom and threads the state between them; \code{get}
reads the current state; \code{modify f} / \code{put s} write it; \code{return a}
yields a value, changing nothing (\cref{sec:primitives}). The load-bearing discipline
is the \emph{single effect point}: only operations that read or write \code{SimState}
go in the monad; everything else---operator applies, inner products, ephemeral
workspace---is a pure \code{let} outside it. In a Krylov step the \emph{only} monadic
operation is one \code{modify}, surrounded by pure \code{let}s; the irreducible core of
that effect is the counter bump
\lstinline[style=l4style]|modify (\s -> { ...s, it: s.it + 1 })|
(\cref{sec:effectpoint},~\cref{sec:counter}). Over-monadifying---dragging pure work
into the monad---buries the one real effect and defeats pruning; the rule of thumb
guards against it. \code{execState} runs a recipe against an initial state and collapses
the whole monad back into a pure \code{SimState -> SimState}: the scaffolding
evaporates at the boundary, leaving \code{solve} pure on the outside
(\cref{sec:execstate}). The monad earns its keep over a plain fold when control flow
grows past one flat loop---interior termination, inner/outer nesting, side-data
threading---but the loop is \emph{still} \code{iterate\_while}; the monad is what the
carry threads through (\cref{sec:whymonad}). The name \code{Solve} is overloaded: the
\emph{effect} \code{StateT SimState Identity a} (how the state is threaded) versus the
\emph{record} \lstinline[style=l4style]|Solve { sim, krylov, outputs }| (what a step
yields), meeting at the one effectful field \code{sim} (\cref{sec:overload}). The
lattice-Boltzmann step is the chapter's lower bound: a pure
\code{FluidState -> FluidState} threaded by \code{iterate\_while\_pure}, with the
mutable world's defensive \code{clone()} evaporating into a no-op---mutation modeled
purely, needing no monad at all (\cref{sec:lbm}).

\furtherreading
The state monad is the classic functional-programming answer to ``how do I thread
mutable-looking state through pure code,'' and the canonical, readable introduction is
Wadler's tutorial on how monads tame state and effects~\citep{wadler1995monads}, which
develops exactly the \code{get}/\code{put}/\code{return} vocabulary used here, from the
ground up and with no category theory. The state monad and the \code{StateT}
transformer (our \code{StateT SimState Identity}) are standard equipment in the
language Peyton Jones documents~\citep{peytonjones2003}; readers who want the precise
desugaring of \code{do}-notation into \texttt{>>=} will find it there. We have used the
monad here only as disciplined bookkeeping for one evolving value; the solver chapters
(\cref{ch:operators} and the Krylov solvers that follow it) put it to work, and
\cref{ch:strata} explains \emph{why} the threaded \code{SimState} splits cleanly from
the build-time \code{OpParams} and the per-step \code{SolveResult}---the stratification
that makes the single effect point possible.

\exercisessection
\begin{enumerate}[nosep]
  \item In the hand-threaded \code{step\_by\_hand} of \cref{sec:problem}, suppose a
  typo made the third line read \lstinline[style=l4style]|residual_of s1.x| instead of
  \lstinline[style=l4style]|s2.x|. Does the program still type-check? Does it compute
  the right residual? Use this to explain, in one sentence, what class of bug the state
  monad eliminates.
  \item Read the type \code{Solve a = StateT SimState Identity a} aloud as a plain
  English sentence, without using the words ``monad'' or ``StateT.'' Then say what
  \code{Solve ()} and \code{Solve Bool} each yield, and what state each threads.
  \item For each of the four primitives \code{get}, \code{modify f}, \code{put s},
  \code{return a}, state whether it \emph{reads} the state, \emph{writes} the state, or
  \emph{neither}, and what value (if any) it yields. Which one is the only one that can
  change the threaded state?
  \item Walk the \code{do} block \code{bump} of \cref{sec:walkdo} against a starting
  state \lstinline[style=l4style]|{ x: x0, it: 3 }|. What value does each \code{get}
  yield? What does the block return? Explain why the two \code{get}s yield states with
  \emph{different} \code{it} fields even though they are the same word.
  \item Classify every line of the Krylov step in \cref{sec:krylovstep} as ``in the
  monad'' or ``outside the monad'' by applying the rule of thumb. How many lines are in
  the monad? Now argue why moving the line \lstinline[style=l4style]|let Ap = apply op.A s_in.p|
  \emph{into} the monad (via a \code{modify} that stashes \code{Ap} in a scratch field)
  would be a mistake---name two distinct things it breaks.
  \item The top-level \code{solve} has type \code{OpParams -> Inputs -> SimState}, with
  no \code{Solve} in sight, yet its body \code{solve\_loop} is a \code{Solve ()} recipe.
  Explain how \code{execState} reconciles these---what goes in, what comes out, and why
  the monad is invisible from outside.
  \item Give one concrete reason from \cref{sec:whymonad} that a restarted GMRES solve
  is awkward to write as a single flat \code{iterate\_while} with \code{SimState} as the
  carry. Then explain in what sense the loop is ``still \code{iterate\_while}'' even when
  the state is threaded monadically.
  \item Distinguish the two faces of \code{Solve}. Given the signature
  \lstinline[style=l4style]|krylov_step :: OpParams -> SimState -> Solve { sim, krylov, outputs }|,
  identify which \code{Solve} is the effect and which is the record, and say which
  \emph{single} field of the record is discharged through the effect and which two are
  plain returned values.
  \item In the LBM step of \cref{sec:lbm} the mutable version needs a defensive
  \code{clone()} but the pure version does not. Explain precisely what hazard the
  \code{clone()} guards against in the mutable world, and why immutability
  (\cref{ch:records}) makes that hazard impossible---so the copy becomes a no-op. Then
  state the three pressures of \cref{sec:whymonad} and confirm the minimal LBM step
  exhibits none, which is why it needs no monad.
\end{enumerate}

\chapter{State Stratification and Lifetimes}
\label{ch:strata}

\chapterorient{Every piece of data a solver touches has a \emph{lifetime}---a
window of time over which it is alive, during which it may be read, perhaps
written, and after which it is gone. The configuration is set once and read
forever; the iterate evolves step by step; the scratch basis is born at a
restart and thrown away at the next; a single recurrence scalar is carried
across an inner loop and then forgotten. This chapter sorts \emph{all} of a
solver's state by lifetime. The sorting bins are the \emph{strata}, and getting
them right is what turns a tangle of instance variables into a design you can
reason about---and what turns the variant-handling of \cref{ch:operators} from a
discipline you must remember into a typing rule the calculus checks for you.}

\section{Why a solver's data has lifetimes}

In \cref{ch:records} we learned to bundle a solver's state into records and to
``update'' them without mutation, by building fresh copies. In \cref{ch:monad}
we learned to thread one such record---the evolving \code{SimState}---through a
chain of steps so the spreads stay out of the algorithm's way. Both chapters
treated ``state'' as one undifferentiated thing: a record that flows. This
chapter draws the lines \emph{inside} that flow.

The motivation is concrete. Open a real iterative solver and you find a dozen
named values living side by side: the operator it inverts, the tolerances it
chases, the current best guess, a counter, a convergence flag, a basis of search
directions, a small dense matrix of inner products, a couple of rotation
registers, one or two stray scalars. In the original C++ they are all
\emph{instance fields} of one solver object---all stored the same way, all
reachable from any method, all looking, syntactically, alike. But they are not
alike. They differ in the one property that matters most for understanding the
design: \emph{how long each one lives}.

\begin{itemize}[nosep]
  \item The operator and the tolerances are fixed when the solver is built and
  never change while it runs. They live for the \emph{whole solve}, read-only.
  \item The current guess, the counter, and the flag \emph{evolve}: each step
  produces new values from the old. They too live for the whole solve, but they
  are read \emph{and} written.
  \item The basis of search directions is scratch space. In a restarted method
  it is allocated fresh at each restart and discarded at the next. Its lifetime
  is \emph{one restart cycle}---shorter than the solve.
  \item A recurrence scalar like the one that forms a search-direction
  coefficient is threaded from each inner step to the next, then reborn from
  scratch the next time the solver is called. Its lifetime is \emph{one call's
  worth of inner steps}---narrower still.
\end{itemize}

Four different lifetimes, four kinds of data, all stored identically in the
source. \Cref{fig:nested} draws them as nested time-windows. The thesis of this
chapter is that making those windows \emph{visible in the type}---giving each
lifetime its own record, its own stratum---is not pedantry. It is the structural
move that makes the rest of the framework tractable.

\begin{figure}[t]
\centering
\begin{tikzpicture}[font=\footnotesize, x=1mm, y=1mm]
  % time axis
  \draw[->,simGray] (0,-3) -- (118,-3) node[right,font=\scriptsize]{time};
  \node[font=\scriptsize,text=simGray] at (4,-6) {build};
  \node[font=\scriptsize,text=simGray] at (60,-6) {the solve runs};
  % OpParams: spans the whole solve (built at construction, read forever)
  \fill[simBgGray] (10,18) rectangle (114,26);
  \draw[simGray,thick] (10,18) rectangle (114,26);
  \node[anchor=west] at (12,22) {\code{OpParams} \;\scriptsize(build-time, read-only)};
  % SimState: spans the whole solve, evolving
  \fill[simBgBlue] (16,8) rectangle (114,16);
  \draw[simBlue,thick] (16,8) rectangle (114,16);
  \node[anchor=west] at (18,12) {\code{SimState} \;\scriptsize(run-time, evolving)};
  % Krylov: one restart cycle (two of them shown)
  \fill[simBgGreen] (22,-1) rectangle (62,6);
  \draw[simGreen,thick] (22,-1) rectangle (62,6);
  \node[anchor=west,font=\scriptsize] at (24,2.5) {\code{Krylov}};
  \fill[simBgGreen] (66,-1) rectangle (106,6);
  \draw[simGreen,thick] (66,-1) rectangle (106,6);
  \node[anchor=west,font=\scriptsize] at (68,2.5) {\code{Krylov}};
  \node[font=\scriptsize,text=simGreen] at (44,-3.6) {restart cycle};
  \node[font=\scriptsize,text=simGreen] at (86,-3.6) {restart cycle};
  % recurrence carry: one apply / inner-loop sweep (tiny bars)
  \foreach \x in {24,30,36,42,48,54}{
    \fill[simBgRust] (\x,-9) rectangle (\x+4,-5.5);
    \draw[simRust,semithick] (\x,-9) rectangle (\x+4,-5.5);}
  \node[anchor=west,font=\scriptsize,text=simRust] at (60,-7.25) {recurrence carry (one inner sweep each)};
\end{tikzpicture}
\caption[The four strata as nested lifetimes]{The four strata of a solver's
state, drawn as time-windows. \code{OpParams} (the build-time configuration) is
alive for the whole run and never changes. \code{SimState} (the evolving state)
is alive for the whole run and changes every step. The ephemeral \code{Krylov}
workspace lives for a single restart cycle and is reborn at each restart. The
recurrence carry is narrower still---it threads one inner sweep and is reborn
each call. The windows \emph{nest}: each shorter-lived stratum sits inside a
longer-lived one. Reading a value from a window that has closed is a bug, and the
whole point of typing the strata separately is to make that bug impossible to
write.}
\label{fig:nested}
\end{figure}

\begin{keyidea}
A value's \emph{lifetime} is part of its type. Two values that are both
``state'' but live for different windows---the configuration that is read for the
whole solve, the scratch basis that dies at the next restart---are different
\emph{kinds} of state and belong in different records. When the calculus puts
each lifetime in its own stratum, ``who owns what, for how long'' stops being a
fact you keep in your head and becomes a fact the type signature states out
loud. This is the single idea the chapter develops; everything else is its
consequences.
\end{keyidea}

\section{The four strata}

We now name the four strata, each with its lifetime, its mutability, and the
real record of the framework that lives there. \Cref{tab:strata} is the map; the
subsections walk the rows.

\begin{table}[t]
\centering
\caption[The four strata of solver state]{The four strata, sorted by lifetime.
The first three are the common case of any iterative solve; the fourth appears
only when an inner loop threads a recurrence scalar. ``Mutability'' is from the
\emph{caller's} view: \code{SimState} is evolved by the solve but is a fresh
value each step, never mutated in place (\cref{ch:records}).}
\label{tab:strata}
\small
\setlength{\tabcolsep}{5pt}
\renewcommand{\arraystretch}{1.25}
\begin{tabular}{@{}p{20mm}p{30mm}p{20mm}p{30mm}@{}}
\toprule
Stratum & Lifetime & Mutability & Record \& example field \\
\midrule
\textbf{Build-time config} & the whole solve & read-only & \code{OpParams}: \code{tol}, \code{flexible} \\
\textbf{Run-time state} & the whole solve, evolving & evolved (fresh each step) & \code{SimState}: \code{x}, \code{it}, \code{converged} \\
\textbf{Ephemeral workspace} & one restart cycle & internal scratch & \code{Krylov}: basis \code{V}, direction \code{p} \\
\textbf{Recurrence carry} & one inner sweep, reborn per call & threaded scalar & \code{PrevCarry}: \code{beta\_prev} \\
\bottomrule
\end{tabular}
\end{table}

\subsection{Stratum one: \texorpdfstring{\code{OpParams}}{OpParams}---build-time configuration}

The first stratum is everything the solver is \emph{configured} with: the
operator it applies, the preconditioner, the tolerances, the iteration limit,
and---crucially---the \emph{variant selectors} that pick which flavour of the
algorithm to run (left or right preconditioning, which Gram--Schmidt variant,
flexible or ordinary). All of it is captured once, at the moment the solver is
constructed, and none of it changes while the solve runs. Its lifetime is the
whole solve; its mutability, from anyone's view, is read-only.

The framework's record for this stratum is \code{OpParams}. Its shape, in the
record syntax of \cref{ch:records}:

\begin{lstlisting}[style=l4style]
OpParams = {
  -- constructed-operator surfaces (the kernel touches OpParams ONLY through these)
  T          : ConstructedOp,      -- the apply surface (preconditioned operator)
  eps        : Convergence,        -- the stopping-predicate surface

  -- variant selectors (closed over by the surfaces above; NOT read by the kernel body)
  pc_side    : PreconditionerSide, -- LEFT | RIGHT
  gs_orthog  : Orthogonalization,  -- MGS | CGS | CGS2
  flexible   : Bool,               -- FGMRES vs GMRES

  -- termination knobs (close into eps)
  max_dim    : Int,                -- restart subspace dimension
  max_it     : Int,
  rel_tol    : Scalar,
  abs_tol    : Scalar
}                                  -- the WHOLE record is readonly
\end{lstlisting}

Two things deserve immediate notice. First, the whole record is marked
\code{readonly}: there is no field a step writes. That is the defining property
of the stratum---build-time configuration is, by definition, the data that is
fixed before the loop and read during it. Second, the record is in two layers.
The \emph{variant selectors} (\code{pc\_side}, \code{gs\_orthog},
\code{flexible}) are raw choices; but the kernel never reads them directly.
Instead they are folded, at construction, into the \emph{constructed-operator
surfaces} (\code{T}, \code{eps})---closures that have already ``baked in'' the
variant choice. The per-step kernel touches \code{OpParams} only through those
surfaces. We will see in \cref{sec:visible} that this two-layer shape is exactly
what makes variant handling a checkable invariant rather than a discipline. The
fields trace to the immutable instance fields of Palace's \code{IterativeSolver}
hierarchy, which the framework un-mixes into this one read-only record.

\subsection{Stratum two: \texorpdfstring{\code{SimState}}{SimState}---run-time evolving state}

The second stratum is the state the solve \emph{evolves}: the externally-visible
quantities a caller observes after the solve returns. This is the record
\cref{ch:records} drew as cards and \cref{ch:monad} threaded through the
\code{Solve} monad. Its lifetime is also the whole solve, but unlike
\code{OpParams} it is read \emph{and} written---a new value every step. Its
shape is the familiar five fields:

\begin{lstlisting}[style=l4style]
SimState = {
  x           : Tensor[N],   -- the current iterate (the solve's primary product)
  it          : Int,         -- iteration count
  converged   : Bool,        -- convergence flag
  final_res   : Scalar,      -- final residual
  initial_res : Scalar       -- residual captured at solve entry
}
\end{lstlisting}

\code{SimState} is run-time-evolved \emph{but immutable}: ``evolves'' here means
exactly what it meant in \cref{ch:records}---each step builds a \emph{fresh}
\code{SimState}, like the previous one but for the fields that changed, and the
\code{Solve} monad of \cref{ch:monad} threads the fresh value forward. From the
\emph{caller's} viewpoint \code{SimState} is read-only (the solve hands back a
new value rather than reaching into theirs); from the \emph{solve's} viewpoint
it is the one bundle that ticks forward step by step. Notice how little actually
changes per step: the iteration counter \code{it} increments---typically the
kernel's \emph{sole} per-step monadic effect---while the heavy iterate \code{x}
is folded only at restart-cycle boundaries, and the exit fields
(\code{converged}, \code{final\_res}) are written only as the solve terminates.
This is the single-mutation-point observation we flagged in \cref{ch:records},
now made structural by the stratum boundary.

\Cref{fig:simstatecard} is the record card, with every field tagged
\emph{run-time} to mark the stratum. The decisive fact about \code{SimState}---
the one that makes it the monad's state type---is that its shape is
\emph{uniform across algorithms}. Conjugate gradients, GMRES, flexible GMRES,
and the Chebyshev smoother all evolve exactly these five fields. They differ
wildly in their workspace and their configuration, but the externally-visible
state they thread is one shape. That uniformity is what lets the \code{Solve}
monad's effect domain be exactly \code{SimState}: the monadic coordination of
\cref{ch:monad} does not vary by algorithm, because the thing it coordinates
does not. The fields mirror the \code{mutable} solve-statistics of Palace's base
class---\code{converged}, \code{initial\_res}, \code{final\_res}, the iteration
count---which the \code{mutable} keyword itself marks as the run-time stratum.

\begin{figure}[t]
\centering
\begin{tikzpicture}[font=\footnotesize]
  \node[draw=simBlue, fill=simBgBlue, rounded corners=3pt, thick,
        minimum width=86mm, minimum height=44mm, inner sep=6pt] (card) at (0,0) {};
  \node[anchor=north, font=\scriptsize\bfseries, text=simBlue]
    at ($(card.north)+(0,-1.5mm)$) {\code{SimState} \;\textnormal{\scriptsize(stratum 2: run-time)}};
  % rows
  \node[anchor=west] at ($(card.north west)+(4mm,-9mm)$)
    {\code{x}\,:\ \(\mathsf{Tensor}[N]\)\quad\scriptsize\itshape current iterate};
  \node[anchor=west] at ($(card.north west)+(4mm,-16mm)$)
    {\code{it}\,:\ \(\mathsf{Int}\)\quad\scriptsize\itshape iteration count};
  \node[anchor=west] at ($(card.north west)+(4mm,-23mm)$)
    {\code{converged}\,:\ \(\mathsf{Bool}\)\quad\scriptsize\itshape convergence flag};
  \node[anchor=west] at ($(card.north west)+(4mm,-30mm)$)
    {\code{final\_res}\,:\ \(\mathsf{Scalar}\)\quad\scriptsize\itshape final residual};
  \node[anchor=west] at ($(card.north west)+(4mm,-37mm)$)
    {\code{initial\_res}\,:\ \(\mathsf{Scalar}\)\quad\scriptsize\itshape entry residual};
  % run-time tags down the right edge
  \foreach \y in {9,16,23,30,37}{
    \node[anchor=east, font=\scriptsize, text=simRust]
      at ($(card.north east)+(-3mm,-\y mm)$) {\textsf{run-time}};}
  % rules
  \foreach \y in {12.5,19.5,26.5,33.5}{
    \draw[simGray!35] ($(card.north west)+(3mm,-\y mm)$) -- ($(card.north east)+(-3mm,-\y mm)$);}
\end{tikzpicture}
\caption[The \code{SimState} record card]{The \code{SimState} record, with all
five fields tagged \emph{run-time} to mark the stratum. This is the only state a
caller observes after the solve returns; it is the value the \code{Solve} monad
of \cref{ch:monad} threads. Its shape is identical for conjugate gradients,
GMRES, flexible GMRES, and Chebyshev---a single five-field card across every
algorithm---which is precisely what lets one monad coordinate them all.}
\label{fig:simstatecard}
\end{figure}

\subsection{Stratum three: ephemeral \texorpdfstring{\code{Krylov}}{Krylov} workspace}

The third stratum is \emph{scratch}: working data whose entire life is internal
to the solve and which is thrown away on return. For a GMRES-shaped method this
is the \code{Krylov} bundle---the orthonormal basis \code{V} of search
directions, the small dense Hessenberg matrix \code{H}, the least-squares
right-hand side \code{s}, the Givens rotation registers \code{cs}/\code{sn}. For
conjugate gradients it is the search direction \code{p} and the residual vector
\code{r}. Whatever the method, this stratum holds the algorithm's gears: large
enough to dominate memory, never visible to the caller.

\begin{lstlisting}[style=l4style]
-- GMRES / FGMRES workspace
Krylov = {
  V  : Basis,    -- orthonormal Krylov basis (the search directions)
  H  : Dense,    -- Hessenberg matrix
  s  : Vec,      -- least-squares right-hand side
  cs : Vec, sn : Vec,  -- Givens rotation registers
  j  : Int       -- position within the current restart cycle
}
\end{lstlisting}

Two properties set this stratum apart, and both are about lifetime. First, its
lifetime is \emph{shorter than the solve}. In a restarted method the basis
\code{V} fills up over an inner cycle and is then discarded: at each restart the
\code{Krylov} bundle is \emph{reborn}---a fresh bundle, not the old one with its
position counter reset. The window in \cref{fig:nested} that says ``one restart
cycle'' is this stratum's, and there are as many such windows as there are
restarts. Second---and this is the subtle structural point---\code{Krylov} is
\emph{not threaded by the monad}. It is threaded as a \emph{plain value},
passed from inner step to inner step by ordinary function application, never
entering the \code{Solve} monad's state.

Why keep it out of the monad? Because the monad threads \code{SimState}, and
\code{SimState} is exactly the state that survives the whole solve. The
\code{Krylov} bundle does \emph{not} survive the whole solve---it is reborn at
every restart. To thread it through the monad would be to claim it has the
lifetime of \code{SimState}, which is false. Its born-at-restart,
discarded-at-restart lifetime \emph{defeats} monadic encoding: the monad is the
wrong tool for a value whose window is a restart cycle, so we thread it as a
plain returned value instead. This is the first place where the lifetime of a
value dictates the \emph{mechanism} that carries it, and it is why the strata are
not a bookkeeping nicety but a design constraint.

\subsection{Stratum four: the recurrence carry}

The first three strata are the whole story for many solvers. But some carry a
\emph{fourth} kind of state, narrower than all three, and it is worth a stratum
of its own because mis-housing it is a real bug class.

Consider the conjugate-gradient direction update. Each step forms the new search
direction from the old one with a coefficient $\beta_k/\beta_{k-1}$, where
$\beta_k = (\mathsf{B}\vr_k, \vr_k)$ is this step's preconditioned residual inner
product and $\beta_{k-1}$ was last step's. To compute step $k$ you need a value
from step $k-1$: the recurrence \emph{threads}. The Chebyshev smoother has the
same shape, carrying a scalar $\rho_{k-1}$ through its inner polynomial
recurrence $\rho_k = 1/(2\theta/\delta - \rho_{k-1})$.

That single threaded scalar is the fourth stratum. The framework records it as
\code{PrevCarry}:

\begin{lstlisting}[style=l4style]
PrevCarry = { beta_prev: Scalar }   -- CG: the prior step's (Br, r)
-- (GMRES variant: PrevCarry = { H_prev: Scalar })
\end{lstlisting}

Its lifetime is the diagnostic. It is \emph{threaded across the inner loop}---
each step reads the previous step's value, so it has a genuine cross-iteration
data dependence---but it is \emph{reborn at each top-level call}: the next time
the solver runs, the recurrence restarts from its base case. ``One call's worth
of inner steps, then gone.'' That window is narrower than \code{OpParams} (which
lives across \emph{all} calls) and narrower than an ordinary scratch temporary
(which has no cross-iteration dependence at all). It is its own thing, and
\cref{fig:fourstrata} sorts a value into it by exactly these two questions.

\begin{figure}[t]
\centering
\resizebox{\textwidth}{!}{%
\begin{tikzpicture}[font=\footnotesize, node distance=7mm and 10mm]
  \tikzset{q/.style={draw=simBlue, fill=simBgBlue, rounded corners=2pt, thick,
      align=center, inner sep=4pt, minimum width=38mm, font=\footnotesize},
    leaf/.style={draw=simGray, fill=white, rounded corners=2pt, semithick,
      align=center, inner sep=4pt, minimum width=28mm, font=\scriptsize}}
  \node[q] (q1) {Captured at construction\\and read forever?};
  \node[leaf, below left=of q1] (op) {\textbf{OpParams}\\\scriptsize build-time config};
  \node[q, below right=of q1] (q2) {Read \emph{and written}, and\\survives the whole solve?};
  \node[leaf, below left=of q2] (sim) {\textbf{SimState}\\\scriptsize run-time state};
  \node[q, below right=of q2] (q3) {Has a cross-iteration\\recurrence dependence?};
  \node[leaf, below left=of q3] (kry) {\textbf{Krylov}\\\scriptsize ephemeral workspace};
  \node[leaf, below right=of q3] (car) {\textbf{PrevCarry}\\\scriptsize recurrence carry};
  \draw[flow] (q1) -- node[above left,font=\scriptsize]{yes} (op);
  \draw[flow] (q1) -- node[above right,font=\scriptsize]{no} (q2);
  \draw[flow] (q2) -- node[above left,font=\scriptsize]{yes} (sim);
  \draw[flow] (q2) -- node[above right,font=\scriptsize]{no} (q3);
  \draw[flow] (q3) -- node[above left,font=\scriptsize]{no} (kry);
  \draw[flow] (q3) -- node[above right,font=\scriptsize]{yes} (car);
\end{tikzpicture}}
\caption[Sorting a value into its stratum]{A decision tree for placing a piece of
solver state. The two questions that matter are about \emph{lifetime}: does the
value survive the whole solve (and is it written), and does it carry a
cross-iteration recurrence? Build-time configuration is captured once and read
forever (\code{OpParams}); evolving caller-visible state survives the solve
(\code{SimState}); scratch with no cross-step dependence is ephemeral workspace
(\code{Krylov}); a threaded recurrence scalar that is reborn each call is the
carry (\code{PrevCarry}).}
\label{fig:fourstrata}
\end{figure}

The fourth stratum is easy to confuse with two of the others, and the confusions
are both bugs. It is \emph{not} build-time configuration: if you stored
$\beta_{k-1}$ in \code{OpParams}, it would persist between calls and corrupt the
next solve, which must restart its recurrence from scratch. And it is \emph{not}
an ordinary ephemeral temporary: those (the residual vector recomputed each step,
the matrix-vector products) have no dependence on the previous iteration, whereas
the carry's whole nature is that step $k$ reads step $k-1$. We return to all
three confusions as pitfalls in \cref{sec:mistakes}.

\section{Why lifetimes must be visible}
\label{sec:visible}

We have four strata and four records. The question this section answers is
\emph{why bother}---why not store everything as instance fields the way the C++
does, and keep the lifetimes in our heads? The answer is that making the
lifetimes visible in the type buys two guarantees that are otherwise only
hopes.

\subsection{You cannot persist ephemeral data past its window}

The first guarantee is a safety property. When the ephemeral \code{Krylov}
workspace is its own value, threaded as a plain value and reborn at each restart,
there is \emph{no place} to accidentally read it after its window closes. The
type says: at the restart boundary, the old \code{Krylov} is gone and a fresh one
takes its place. A step that tried to read a basis vector from before the restart
would be reading from a value that no longer exists---and because the fresh
bundle is a genuinely new value (\cref{ch:records}), there is no aliased box
holding the stale data for it to find.

Contrast the mutable world. There, the basis \code{V} is an instance field that
outlives any one restart cycle; ``reborn at restart'' is implemented by
\emph{overwriting} it in place, and nothing stops a stray method from reading a
slot that belongs to the previous cycle. The lifetime is real but invisible---a
fact about when the code happens to overwrite \code{V}, not a fact the type
enforces. The stratum makes the invisible visible: a restart \emph{produces a
new bundle}, so reading across the boundary is reading a value that was never
built.

\subsection{Variant absorption becomes a typing invariant}

The second guarantee is the deeper one, and it is the bridge to
\cref{ch:operators}. Recall the two-layer shape of \code{OpParams}: variant
selectors on one layer, constructed-operator surfaces on another, and the rule
that \emph{the per-step kernel touches \code{OpParams} only through the
surfaces}, never through the raw selectors. Why is that rule worth stating?

Because it makes \emph{variant absorption} checkable. A single kernel is supposed
to run all the variants---left and right preconditioning, the three
Gram--Schmidt flavours, flexible and ordinary---without itself branching on which
variant it is. The variant choice is supposed to be ``absorbed'' at construction,
into the surfaces, so the kernel's control flow is uniform. In a procedural
solver that uniformity is a \emph{discipline}: the programmer must remember not
to write \code{if (pc\_side == LEFT)} in the inner loop. Nothing checks it.

With the strata, it is a \emph{typing invariant}. If \code{OpParams} is
\code{readonly}, and the variant selectors are reachable from the kernel only
\emph{through} the constructed surfaces (because the kernel is handed the
surfaces, not the raw record), then the kernel \emph{provably cannot} branch on a
selector---there is no expression it could write to read one. Variant absorption
stops being ``a thing we are careful to do'' and becomes ``a thing the type
forbids us from violating.'' That is the entire payoff of separating the
build-time stratum from the run-time one, and \cref{ch:operators} cashes it in to
show one kernel covering an entire family of algorithms.

\begin{keyidea}
Visible lifetimes turn two correctness properties from discipline into typing.
\emph{Safety:} because ephemeral workspace is a fresh value reborn at each
restart, there is no stale box to read across the boundary---persisting it past
its window is unwriteable, not merely discouraged. \emph{Variant absorption:}
because \code{OpParams} is read-only and the kernel reaches the variant selectors
only through the constructed surfaces, the kernel \emph{cannot} branch on a
variant---the absorption invariant of \cref{ch:operators} is checked by the
types, not remembered by the programmer. Both guarantees are consequences of one
move: giving each lifetime its own stratum.
\end{keyidea}

\section{Build-time versus run-time: the two great strata}

Step back from the four-way split and a coarser, even more important division
appears. Group the strata by \emph{when their work runs}, and there are exactly
two camps:

\begin{itemize}[nosep]
  \item \textbf{Build-time work} runs \emph{once}, before the loop. Constructing
  the operator surfaces, factoring a preconditioner, installing pointers, parsing
  the configuration record, querying the function-space hierarchy---all of it
  happens at solver construction and is done by the time the first iteration
  begins. This is the world of \code{OpParams}.
  \item \textbf{Run-time work} runs \emph{every step}, inside the loop. Applying
  the operator, applying the preconditioner, bumping the counter, computing a
  residual norm, folding the iterate. This is the world of \code{SimState}, the
  \code{Krylov} workspace, and the carry.
\end{itemize}

\Cref{fig:buildrun} draws the two as a timeline: a setup pass, then the loop.
This build-time/run-time line is the same one we drew for the matrix-free
operator in \cref{ch:matrixfree}, where the operator's \emph{setup}---assembling
the element-local quadrature data once---is split from its \emph{apply}---the
per-call tensor contraction that runs every iteration. There, as here, the setup
is build-time and the apply is run-time; the two were separated for exactly the
reason the strata separate them now.

\begin{figure}[t]
\centering
\begin{tikzpicture}[font=\footnotesize, x=1mm, y=1mm]
  % timeline
  \draw[->,simGray] (0,0) -- (122,0) node[right,font=\scriptsize]{time};
  % build-time pass
  \fill[simBgGray] (4,4) rectangle (40,16);
  \draw[simGray,thick] (4,4) rectangle (40,16);
  \node[align=center,font=\scriptsize] at (22,13) {\textbf{build-time pass} (once)};
  \node[align=center,font=\scriptsize,text=simGray] at (22,8.5)
    {construct surfaces\\factor preconditioner\\parse config};
  \node[font=\scriptsize,text=simGray,anchor=north] at (22,3) {\code{OpParams} born here};
  % run-time loop
  \fill[simBgBlue] (48,4) rectangle (118,16);
  \draw[simBlue,thick] (48,4) rectangle (118,16);
  \node[align=center,font=\scriptsize] at (83,13) {\textbf{run-time loop} (every step)};
  % loop iterations as little boxes
  \foreach \x in {52,62,72,82,92,102}{
    \draw[simBlue,fill=white,semithick] (\x,6) rectangle (\x+7,10.5);
    \node[font=\tiny] at (\x+3.5,8.25) {apply};}
  \node[font=\scriptsize,text=simBlue,anchor=north] at (83,3)
    {\code{SimState} / \code{Krylov} evolve here};
  % the boundary
  \draw[->,flow] (40,10) -- (48,10);
  \node[font=\scriptsize,text=simRust,align=center] at (44,16) {loop\\entry};
\end{tikzpicture}
\caption[The build-time / run-time timeline]{Solver work splits into a setup
pass that runs once and a loop that runs every step. Build-time work---building
the operator surfaces, factoring the preconditioner, parsing configuration---is
done before the first iteration and populates \code{OpParams}. Run-time
work---the per-step applies, counter bumps, residual computations---runs inside
the loop and evolves \code{SimState} and the ephemeral workspace. This is the
same split as the matrix-free operator's setup-versus-apply of
\cref{ch:matrixfree}.}
\label{fig:buildrun}
\end{figure}

The division matters for a reason beyond tidiness, and it connects to the very
purpose of the layered framework. \emph{Only run-time work lifts to the
tensor-field form of \cref{ch:tensors}.} When we re-express a solve as global
tensor-field operations---the whole point of the climb toward the abstract
calculus---it is the per-element, per-iteration work that has structure to lift:
an \code{apply} is a tensor contraction, a residual is a reduction, a counter
bump is a scalar increment. Build-time work has no per-element structure to lift;
it \emph{sets up} the loop rather than \emph{being} a loop. So the build-time
stratum is recorded as the scaffolding that materializes the operator graph, and
it is the run-time stratum---and only the run-time stratum---that participates in
the lift. The build-time/run-time line is thus not only a lifetime distinction;
it is the line between the work that becomes a tensor-field expression and the
work that prepares for it.

In the monad of \cref{ch:monad} this lands as a clean structural rule:
build-time work sits in the \code{do}-block's \emph{setup} prologue, before the
loop; run-time work sits in the \emph{body} that the loop iterates. Schematically,

\begin{lstlisting}[style=l4style]
solve = do
  setup <- build_operators(config)   -- BUILD-TIME: once, before the loop
  iterate_while(not_converged) (\s ->
     step(setup, s))                  -- RUN-TIME: every iteration
\end{lstlisting}

The \code{setup} line runs once and produces the \code{OpParams}-stratum
surfaces; the \code{step} inside \code{iterate\_while} runs every iteration and
evolves the run-time strata. Reading the \code{do}-block, you can see the
lifetime of every value by \emph{where} it appears: above the loop, it is
build-time; inside, it is run-time.

\begin{workedexample}{Sorting a conjugate-gradient solve, variable by variable}{cgsort}
We take a complete preconditioned conjugate-gradient solve and sort
\emph{every} variable it names into its stratum. The exercise is the whole
chapter in miniature: given a pile of state, place each value by its lifetime.

\textbf{The variables.} A PCG solve holds: the system operator $\mathsf{A}$ and
preconditioner $\mathsf{B}$; the relative and absolute tolerances and the
iteration limit; the current iterate $\vx$; the iteration count and the
convergence flag; the residual vector $\vr$ and the preconditioned residual
$\mathsf{B}\vr$; the search direction $\vp$; the step lengths $\alpha$ formed
each step; and the recurrence scalar $\beta = (\mathsf{B}\vr,\vr)$, whose
\emph{previous} value $\beta_{k-1}$ feeds the direction-update coefficient.

\tcblower
\textbf{The sort.} Ask of each value the two lifetime questions of
\cref{fig:fourstrata}.

\begin{itemize}[nosep]
  \item \(\mathsf{A}\), \(\mathsf{B}\), \code{rel\_tol}, \code{abs\_tol},
  \code{max\_it} --- captured at construction, read forever, never written.
  \textbf{Stratum 1, \code{OpParams}.} (The operators are reached through the
  constructed \code{T} surface; the tolerances close into \code{eps}.)
  \item \(\vx\), \code{it}, \code{converged}, the residual proxies --- evolved
  every step, survive the whole solve, visible to the caller.
  \textbf{Stratum 2, \code{SimState}.}
  \item \(\vr\), \(\mathsf{B}\vr\), \(\vp\), and the per-step \(\alpha\) ---
  scratch, recomputed each step, with no value carried from one step to the next;
  discarded on return. \textbf{Stratum 3, ephemeral \code{Krylov}.}
  \item \(\beta_{k-1}\) --- the one scalar with a genuine cross-step dependence
  (\(\vp_k = \vr_k + (\beta_k/\beta_{k-1})\,\vp_{k-1}\)), threaded inner step to
  inner step but restarting from its base case each new solve.
  \textbf{Stratum 4, the carry \code{PrevCarry}.}
\end{itemize}

\textbf{The trap, named.} Note that $\vp$ and $\beta_{k-1}$ both ``come from the
previous step,'' and the lazy sort would lump them together. They are in
different strata. The direction $\vp$ is rebuilt each step from current
quantities (it has a \emph{recurrence in form} but we treat the rebuilt vector as
ephemeral workspace); the \emph{scalar} $\beta_{k-1}$ is the value with the true
cross-iteration data dependence that the carry exists to thread. The
distinguishing question---``does step $k$ read a \emph{value} from step
$k-1$?''---is what separates stratum 4 from stratum 3, and it is exactly the
bottom question of \cref{fig:fourstrata}.
\end{workedexample}

\section{The book's real records, defined}

We have met all five records in passing. We now define each one as a card of
named fields, with the stratum each field lives in stated outright. This section
is the reference the rest of Part~III points back to.

\subsection{\texorpdfstring{\code{SolveResult}}{SolveResult}---what a step returns}

A single step does not return a bare \code{SimState}; it returns a bundle. The
framework calls the bundle \code{SolveResult}, written with the \code{Solve}
name:

\begin{lstlisting}[style=l4style]
SolveResult = Solve { sim: SimState, krylov: Krylov, outputs: StepOutputs }
-- (Form B, when the carry is threaded explicitly:)
--   Solve { sim: SimState, krylov: Krylov, outputs: StepOutputs, carry: PrevCarry }
\end{lstlisting}

Read it field by field, by stratum:

\begin{itemize}[nosep]
  \item \code{sim} --- the next \code{SimState}, stratum 2. This is the one field
  the step delivers \emph{through the monad}: it is the monadic-effect product,
  the state transition the \code{Solve} monad of \cref{ch:monad} threads.
  Typically the only thing that changed is the \code{it} counter.
  \item \code{krylov} --- the next ephemeral workspace, stratum 3, returned as a
  \emph{plain value} (because, as we saw, its restart-bounded lifetime defeats
  monadic threading).
  \item \code{outputs} --- the demand-prunable per-step readouts, a fresh bundle
  each step; defined just below.
  \item \code{carry} --- present only in Form B, the recurrence carry of stratum
  4 (\code{PrevCarry}), threaded back into the next step.
\end{itemize}

There is a deliberate overload here, and \cref{ch:monad} warned of it:
\code{Solve} names \emph{both} the monad (the effect that threads \code{SimState})
\emph{and} this returned record (the braces). They are different things sharing a
spelling: the monad answers ``how is state threaded?''; \code{SolveResult}
answers ``what record does a step yield?''. The \code{sim} field is where they
touch---it is the effect's state product---while the other fields are the plain
returned value. \Cref{ch:monad} owns the threading; this stratum chapter owns the
record's fields and their lifetimes.

\subsection{\texorpdfstring{\code{StepOutputs}}{StepOutputs}---demand-prunable per-step readouts}

The \code{outputs} field is its own record, and it is the textbook home of the
optional-field, demand-pruning idea of \cref{ch:records}:

\begin{lstlisting}[style=l4style]
StepOutputs = {
  residual_norm:    Scalar,        -- the step's residual proxy
  ls_residual?:     Scalar,        -- GMRES/FGMRES least-squares residual estimate
  breakdown_token?: BreakdownTag   -- a guarded-quantity validity tag
}
\end{lstlisting}

Every field is run-time and \emph{derived}: each is a pure function of the
post-step \code{Krylov} bundle---the residual norm read off the recurrence
scalar, the least-squares residual read off the rotated right-hand side, the
breakdown token read off a guarded dot product. And every optional field is
\emph{demand-prunable}: if no downstream consumer reads \code{ls\_residual}, the
kernel never computes it. The record exists precisely to make that pruning
structural---by typing the derived observations separately from the \code{Krylov}
state they are derived from, the calculus can state ``unread output, uncomputed''
as a law (\cref{ch:fold}) rather than a comment. \code{StepOutputs} is the
\emph{observation} a step makes about itself; do not confuse it with the
persistent \code{SimState.final\_res} statistic it eventually feeds.

\subsection{The five records, at a glance}

\Cref{tab:records} collects all five with their strata. Read it as the answer
to ``where does each piece of a solve's data live, and for how long.''

\begin{table}[t]
\centering
\caption[The five workhorse records and their strata]{The five records of a
Krylov-shaped solve, each with its stratum and lifetime. \code{OpParams} and
\code{SimState} are the two persistent strata (build-time vs.\ run-time);
\code{Krylov} and \code{PrevCarry} are the two short-lived ones; \code{SolveResult}
and its \code{StepOutputs} field are the result-side bundles a single step
returns.}
\label{tab:records}
\small
\setlength{\tabcolsep}{5pt}
\renewcommand{\arraystretch}{1.25}
\begin{tabular}{@{}p{24mm}p{20mm}p{34mm}p{30mm}@{}}
\toprule
Record & Stratum & Lifetime & Role \\
\midrule
\code{OpParams} & 1 (build) & whole solve, read-only & configuration + surfaces \\
\code{SimState} & 2 (run) & whole solve, evolving & caller-visible state \\
\code{Krylov} & 3 (ephemeral) & one restart cycle & algorithm workspace \\
\code{PrevCarry} & 4 (carry) & one inner sweep, reborn & threaded recurrence scalar \\
\code{StepOutputs} & run, derived & one step, demand-pruned & per-step readouts \\
\code{SolveResult} & result bundle & one step & what a step returns \\
\bottomrule
\end{tabular}
\end{table}

\section{The common stratification mistakes}
\label{sec:mistakes}

Each stratum boundary is a place where a value can be filed wrong, and each
mis-filing is a recognizable bug. We teach them as pitfalls because they are the
errors a reader coming from the mutable, instance-field world will reach for by
habit---the C++ stores everything as instance fields, so the original layout
offers no guidance on where a value \emph{belongs}.

\begin{pitfall}
\textbf{Ephemeral workspace in \code{SimState}.} The tempting mistake is to put
the basis \code{V} or the direction \code{p} into \code{SimState} because the
original C++ stored it as an instance field. But the instance-field layout is a
\emph{memory-reuse optimisation}, not a statement of lifetime: the field is
reused across restarts because allocating it once is cheap, not because it
\emph{lives} that long. Mathematically the basis dies at each restart. Put it in
\code{SimState} and you have claimed it survives the whole solve---so it
persists, wrongly, past the restart that should have discarded it, and a later
step can read a stale basis vector. The L4 type should reflect the mathematical
lifetime, not the allocation strategy. \emph{Workspace is stratum 3, never
stratum 2.}
\end{pitfall}

\Cref{fig:bug} draws this bug: a basis vector from before a restart, wrongly
surviving inside \code{SimState} and read by a step that should see only the
reborn bundle.

\begin{figure}[t]
\centering
\begin{tikzpicture}[font=\footnotesize, x=1mm, y=1mm]
  % restart boundary
  \draw[simRust, very thick, densely dashed] (50,-12) -- (50,20);
  \node[font=\scriptsize\bfseries, text=simRust] at (50,22) {restart};
  \node[font=\scriptsize, text=simGray] at (24,22) {cycle $r$};
  \node[font=\scriptsize, text=simGray] at (78,22) {cycle $r{+}1$};
  % SimState band wrongly carrying V across
  \fill[simBgBlue] (6,8) rectangle (104,16);
  \draw[simBlue,thick] (6,8) rectangle (104,16);
  \node[anchor=west,font=\scriptsize] at (8,12) {\code{SimState} \emph{(correct: x, it, ...)}};
  % the wrongly-persisted V, drawn straddling the restart line
  \fill[simBgRust] (30,-4) rectangle (70,4);
  \draw[simRust,thick] (30,-4) rectangle (70,4);
  \node[font=\scriptsize, text=simRust] at (50,0) {\code{V} wrongly kept in \code{SimState}};
  % the read that goes wrong
  \draw[backflow] (78,0) -- node[below,font=\tiny,text=simRust]{reads stale basis} (62,0);
  \node[opbox, minimum width=12mm] at (82,0) {step};
  % correct: Krylov reborn
  \node[font=\scriptsize, text=simGreen, align=center] at (24,-9)
    {\emph{correct:} \code{Krylov}\\dies at restart};
  \node[font=\scriptsize, text=simGreen, align=center] at (84,-9)
    {\emph{correct:} fresh\\\code{Krylov} born};
\end{tikzpicture}
\caption[The ephemeral-in-\code{SimState} bug]{Putting the ephemeral basis
\code{V} into \code{SimState} claims it survives the whole solve. It then
straddles the restart boundary (dashed), and a step in cycle $r{+}1$ reads a
basis vector from cycle $r$---stale data. The fix is to keep \code{V} in the
\code{Krylov} stratum, which is reborn at the restart, so there is no surviving
value to read across the boundary.}
\label{fig:bug}
\end{figure}

\begin{pitfall}
\textbf{Variant selectors in \code{SimState}.} The selectors \code{pc\_side},
\code{gs\_orthog}, \code{flexible} are \emph{read} during the solve, so it is
tempting to file them with the run-time state. They are build-time
configuration. They are read only by the constructed-operator surfaces---which
were built once, at construction, with the selectors already baked in---never by
the per-step kernel. File them in \code{SimState} and you have invited the kernel
to branch on them, which destroys variant absorption (\cref{sec:visible}): the
whole point was that the kernel \emph{cannot} see a raw selector. \emph{Selectors
are stratum 1, never stratum 2}---``read during the solve'' is not the same as
``run-time state.''
\end{pitfall}

\begin{pitfall}
\textbf{Forgetting that the carry is reborn at restart (or at each call).} The
recurrence carry \code{PrevCarry} threads a scalar across inner steps, so it
looks persistent. But it is reborn at each top-level call: the next solve
restarts its recurrence from the base case. Two failures follow from forgetting
this. \emph{If you file the carry in \code{OpParams}}, it persists between calls
and the next solve starts with last solve's $\beta_{k-1}$---a corrupted
recurrence. \emph{If your form reads a \code{Krylov} field across a restart
boundary} expecting the recurrence to continue, it reads from a bundle that was
discarded. The carry's lifetime is exactly ``one call's inner steps''---wider
than an ordinary temporary, narrower than configuration---and a form that gets
that window wrong is broken even when it type-checks against the field names.
\end{pitfall}

\begin{workedexample}{Build-time versus run-time in a preconditioned solve}{buildrun}
We take a preconditioned solve and split its work into the two great strata,
then show what each \emph{mislabeling} costs. This is the build-time/run-time
line of \cref{fig:buildrun} made concrete on one operation: the preconditioner.

\textbf{The two halves of a preconditioner.} A preconditioner $\mathsf{B}$
enters a solve in two distinct acts. First it is \emph{constructed}: for an
incomplete-factorization preconditioner this means \emph{factoring}---computing
the sparse triangular factors once, an expensive setup. Then, every iteration, it
is \emph{applied}: a triangular solve against the current residual, $\vz \gets
\mathsf{B}\vr$, cheap per call but run every step.

\textbf{The correct split.}
\begin{itemize}[nosep]
  \item \textbf{Build-time, once:} factor $\mathsf{B}$. The factors are
  \code{OpParams}-stratum data---configuration of the operator, captured at
  construction, read (applied) but never re-factored during the loop. In the
  monad they sit in the \code{setup} prologue.
  \item \textbf{Run-time, every step:} apply $\mathsf{B}$ to the residual. This
  is the per-iteration work; it touches \code{SimState} (the residual it reads,
  the iterate it helps advance) and lifts, in \cref{ch:tensors}, to a
  tensor-field expression.
\end{itemize}

\tcblower
\textbf{What each mislabeling costs.} The two halves are not interchangeable, and
swapping a label is a real failure---one wrong each way.

\begin{itemize}[nosep]
  \item \textbf{Calling the factorization run-time} (re-factoring $\mathsf{B}$
  every iteration) is \emph{wasted work}: you recompute, at full setup cost, a
  result that never changed. A solve that should factor once and apply a thousand
  times instead factors a thousand times. The answer is right; the cost is
  catastrophic. This is the classic ``loop-invariant computation left inside the
  loop.''
  \item \textbf{Calling the apply build-time} (applying $\mathsf{B}$ once, before
  the loop, and reusing the result) gives \emph{wrong results}: the apply depends
  on the current residual, which changes every step, so a single hoisted apply
  feeds every iteration a stale preconditioned vector and the solve does not
  converge to the right answer---if it converges at all.
\end{itemize}

\textbf{The lesson.} The build-time/run-time line is not cosmetic. Mislabel it
one way and you waste work; mislabel it the other and you compute the wrong
thing. The discriminating question is precisely \emph{lifetime}: does the value
change across iterations? The factorization does not (build-time); the apply's
result does (run-time). The matrix-free operator of \cref{ch:matrixfree} draws
exactly this line as its setup-versus-apply split, and for exactly this reason.
\end{workedexample}

\section{Lifetimes as the backbone of the design}

We close by collecting the thread. A solver's data is not one undifferentiated
``state.'' It is four kinds, sorted by lifetime: the build-time configuration
that is read for the whole solve, the run-time state that evolves across it, the
ephemeral workspace that lives one restart cycle, and the recurrence carry that
threads one inner sweep and is reborn each call. Each lifetime gets its own
record---\code{OpParams}, \code{SimState}, \code{Krylov}, \code{PrevCarry}---and
the result bundle \code{SolveResult} (with its demand-prunable \code{StepOutputs})
is what a single step returns.

The payoff is that the strata make the design's hardest questions structural
rather than disciplinary. You cannot persist ephemeral data past its window,
because the window's end is a fresh value, not an overwrite. You cannot let the
kernel branch on a variant, because the variant selectors are read-only and
reachable only through the constructed surfaces. And you cannot lose track of
which work runs once versus every step, because the build-time/run-time line is
the line between the \code{setup} prologue and the loop body---and, not
incidentally, the line between the work that prepares the loop and the work that
lifts to the tensor-field form of \cref{ch:tensors}. The variant-handling
machinery of \cref{ch:operators} rests entirely on this: with the strata in
place, absorbing a family of algorithms into one kernel is a typing invariant the
calculus checks, not a discipline the programmer must keep.

\summarysection
Every piece of a solver's state has a \emph{lifetime}, and classifying state by
lifetime---its \emph{stratum}---is what makes the design tractable. There are
four strata. \textbf{Build-time configuration} (\code{OpParams}) is captured once
at construction and read, read-only, for the whole solve; it holds the operator
surfaces, tolerances, and variant selectors, in a two-layer shape where the
kernel reaches the selectors only through constructed surfaces. \textbf{Run-time
state} (\code{SimState}) is the caller-visible bundle---\code{x}, \code{it},
\code{converged}, \code{final\_res}, \code{initial\_res}---that evolves every step
(as fresh values, never mutated) and is threaded by the \code{Solve} monad
(\cref{ch:monad}); its five-field shape is uniform across CG, GMRES, FGMRES, and
Chebyshev. \textbf{Ephemeral workspace} (\code{Krylov}: basis \code{V},
direction \code{p}) lives one restart cycle, is reborn at each restart, and is
threaded as a plain value, not in the monad---its lifetime defeats monadic
encoding. The \textbf{recurrence carry} (\code{PrevCarry}: \code{beta\_prev}) is
a fourth stratum, threaded across an inner loop but reborn each call. Making
lifetimes visible buys two guarantees: ephemeral data cannot be persisted past
its window, and variant absorption (\cref{ch:operators}) becomes a typing
invariant. The coarse \textbf{build-time versus run-time} split---setup once,
loop every step---is the same line the matrix-free operator draws
(\cref{ch:matrixfree}); only run-time work lifts to the tensor-field form
(\cref{ch:tensors}). The three classic mistakes are filing ephemeral workspace,
variant selectors, or the recurrence carry in the wrong stratum, and each is a
distinct, real bug.

\furtherreading
The conjugate-gradient and GMRES algorithms whose state we stratified, including
the restart structure that gives the ephemeral workspace its one-cycle lifetime
and the direction-update recurrence that gives the carry its cross-iteration
dependence, are developed in Saad's standard treatment of iterative
methods~\citep{saad2003}; reading the GMRES restart loop there with the four
strata in mind is the best exercise this chapter can recommend. The discipline of
making lifetime and effect visible in a value's type---so that ``what survives,
what is scratch, what threads'' is a fact the type states rather than a comment
the programmer leaves---is the contribution of pure functional programming and
its state monad; Peyton Jones's account of the lazy, purely functional
language~\citep{peytonjones2003} is where the threaded-state-as-a-value idea this
chapter applies to solvers is laid out in full, and it pairs naturally with the
monad of \cref{ch:monad}.

\exercisessection
\begin{enumerate}[nosep]
  \item For each of the four strata, state its lifetime in one phrase and name one
  field of the framework's record that lives there. (Example: ephemeral
  workspace---``one restart cycle''---\code{Krylov.V}.)
  \item A colleague proposes storing the iteration counter \code{it} in
  \code{OpParams} ``because it is just one integer and lives the whole solve.''
  Give the one reason this is wrong, and say which stratum \code{it} belongs to.
  (Hint: read-only versus evolved.)
  \item Take a GMRES solve and sort each of the following into its stratum: the
  preconditioner $\mathsf{B}$; the Krylov basis \code{V}; the Hessenberg matrix
  \code{H}; the iteration count \code{it}; the relative tolerance; the
  least-squares residual estimate $|s[j{+}1]|$. For the last one, say whether it
  is state or a derived \emph{output}, and which record it lives in.
  \item Explain, in two sentences, why the ephemeral \code{Krylov} bundle is
  threaded as a plain value rather than through the \code{Solve} monad. Your
  answer should turn on the word ``restart.''
  \item The chapter claims you ``cannot accidentally persist ephemeral data past
  its restart'' once it is in the \code{Krylov} stratum. Explain why the immutable
  fresh-value discipline of \cref{ch:records} is what makes this true, and what
  would go wrong in a mutable instance-field implementation.
  \item Why does putting a variant selector such as \code{pc\_side} into
  \code{SimState} defeat variant absorption? Trace the argument from ``the kernel
  can now read \code{pc\_side}'' to ``the kernel can now branch on the variant.''
  \item Distinguish the conjugate-gradient search direction $\vp$ (ephemeral,
  stratum 3) from the recurrence scalar $\beta_{k-1}$ (carry, stratum 4). Both
  ``come from the previous step''---state the precise question that separates
  them, and answer it for each.
  \item Consider a preconditioner that is \emph{rebuilt} partway through a solve
  (for instance, when the operator changes mid-sweep). Does the factorization
  stay build-time? Argue both sides, then propose how the stratification should
  describe a value that is mostly build-time but occasionally refreshed. (This is
  a genuine edge case; there is no single right answer, but a good answer names
  the lifetime precisely.)
\end{enumerate}

\chapter{Operators as Values: Constructed Operators and Capability Typing}
\label{ch:operators}

\chapterorient{In \cref{ch:language} we learned that a function is a value, and
that partial application can \emph{freeze} some arguments now and supply the rest
later---``build once, apply many.'' In \cref{ch:strata} we learned that the
build-time configuration a solver bakes in lives in a different stratum, with a
different lifetime, from the run-time vectors it processes. This chapter pays both
ideas off at once. It makes precise the single most important value in the whole
numerical machinery: the \emph{operator}---a function built once with its entire
configuration baked into immutable internal state, then applied again and again as
a pure function. By the end you will be able to say exactly why ``the
preconditioner is just a function you apply,'' why the \emph{same} \code{GMRES}
drives wildly different preconditioners without changing a line, why a \emph{flexible}
preconditioner cannot be a constructed operator and must instead be threaded
through state, and how a phantom-typed \emph{brand} can make ``you passed the
preconditioner where the true operator belongs'' a compile-time error rather than a
silent wrong answer.}

\section{The shape we keep meeting}

We have now met the same shape four or five times, in four or five disguises, and
it is time to name it once and for all.

In \cref{ch:targets} the simulator assembled an operator from the mesh and
materials and then solved a system built from that operator. In \cref{ch:language}
a constructor took a configuration and \emph{returned a function}
(\lstinline[style=l4style]|make_operator :: Config -> (Tensor[N] -> Tensor[N])|),
which we then applied to many different vectors. In \cref{ch:strata} we separated
the \emph{build-time stratum}---configs, tables, a factorization, computed once and
never touched again---from the \emph{run-time stratum}---the iterate, the residual,
the state that evolves every step. Each time, the same picture: an expensive,
once-only \emph{construction} produces a thing; a cheap, repeated \emph{application}
uses it.

That thing is what this book calls a \emph{constructed operator}, and it is the
backbone of every solver, every preconditioner, every assembled matrix in the
chapters ahead. The shape is so pervasive that it earns its own vocabulary, its own
type spelling, and---because getting it slightly wrong silently corrupts a
solve---its own pitfalls. This chapter develops all three.

\begin{keyidea}
A \emph{constructed operator} is a value built \emph{once}, with its entire
configuration baked into immutable internal state, then \emph{applied} many times as
a pure function. Construction is expensive and happens at build time (factorize,
precompute, allocate); application is cheap and pure with respect to the operator's
internals (it reads them, never changes them). The two phases live in two strata
(\cref{ch:strata}) with two lifetimes. Spotting this shape---and resisting the urge
to collapse the two phases---is the central skill of this chapter.
\end{keyidea}

\subsection{Two phases, two strata}

Make the two phases concrete. A constructed operator has exactly two events in its
life, and they can be arbitrarily far apart in the program.

\begin{description}[style=nextline,leftmargin=1.5em]
  \item[Construction.] Inputs are the \emph{configuration}: enum flags, raw tables,
  problem data, a matrix to factorize. The output is an operator value carrying
  internal immutable state. This step is often expensive---a sparse $LU$
  factorization, a precomputed basis-evaluation table, an assembled stiffness
  matrix. It happens \emph{once}.
  \item[Application.] The only input is the \emph{operand}---the vector (or sim
  state) being transformed. The output is a new vector. This step is cheap relative
  to construction, and it is \emph{pure with respect to the operator's internals}:
  it reads the baked-in state, produces a fresh result, and changes nothing. It
  happens \emph{many} times.
\end{description}

\Cref{fig:constructapply} draws the split. Configs and tables flow \emph{in} on the
left; an opaque operator value sits in the middle; a stream of cheap
\code{apply(v)} calls flows out on the right, each reading the same frozen internals
and returning a fresh vector. The left arrow fires once; the right arrows fire
thousands of times.

\begin{figure}[t]
\centering
\begin{tikzpicture}[node distance=10mm]
  % ---- construction inputs ----
  \node[data, align=left, text width=26mm] (cfg)
    {config / enums\\raw tables\\a matrix to\\factorize};
  % ---- the construct step ----
  \node[stage, minimum width=26mm, minimum height=14mm, right=14mm of cfg] (con)
    {\textbf{construct}\\[1pt]\footnotesize factorize,\\precompute};
  \draw[flow] (cfg) -- node[above,font=\scriptsize]{once} (con);
  % ---- the opaque operator value ----
  \node[draw=simBlue, fill=simBgBlue, very thick, rounded corners=3pt,
        minimum width=24mm, minimum height=20mm, right=14mm of con] (op) {};
  \node[font=\ttfamily\scriptsize, simInk] at (op.north) [yshift=-3.2mm]
    {\lstinline[style=l4style]|Op[N -> N]|};
  \node[font=\scriptsize\itshape, simGray, align=center] at (op.center)
    {frozen\\internals\\(factorization,\\tables)};
  \draw[flow] (con) -- (op);
  % ---- many cheap applies out ----
  \foreach \y/\lbl in {1.1/a, 0.0/b, -1.1/c}{
    \node[draw=simTeal, fill=simBgGreen, semithick, rounded corners=2pt,
          minimum width=18mm, minimum height=6mm] (ap\lbl) at ($(op.east)+(2.8,\y)$)
          {\footnotesize\ttfamily apply v};
    \draw[flow] (op.east) -- (ap\lbl.west);
  }
  \node[font=\scriptsize\itshape, simGray] at ($(op.east)+(2.8,-1.9)$) {many times, cheap};
  % build-time / run-time band labels
  \node[font=\scriptsize\bfseries, simRust] at ($(con)+(0,-1.6)$) {build-time stratum};
  \node[font=\scriptsize\bfseries, simGreen] at ($(op.east)+(2.8,1.9)$) {run-time stratum};
\end{tikzpicture}
\caption[The construct-then-apply phase split]{The defining shape of a constructed
operator. \emph{Construction} (left) consumes the configuration---enums, tables, a
matrix to factorize---and runs \emph{once}, producing an opaque operator value whose
internal state (a factorization, precomputed tables) is now frozen. \emph{Application}
(right) consumes only an operand vector and runs \emph{many} times, each call cheap
and pure on the frozen internals. The two phases live in the two strata of
\cref{ch:strata}: construction is build-time, application is run-time, and the
boundary between them is a lifetime boundary the type system can enforce.}
\label{fig:constructapply}
\end{figure}

\subsection{It is currying, made first-class}

You have, in fact, already built constructed operators---you just called them
closures. Recall \code{scaled\_by} from \cref{wex:scaledby}: it took a scalar,
captured it, and returned a function that scales a tensor by that captured value.
Construction was ``supply the scalar''; application was ``supply the tensor''; the
captured scalar was the baked-in internal state. A constructed operator is that
pattern grown up: the captured state is no longer a single scalar but a whole
configuration---a factorization, a stack of basis tables, an enum that selected an
algorithm---and the returned function is no longer a throwaway lambda but a
long-lived value the rest of the program holds and applies.

This is why \cref{ch:language} could already write the signature

\begin{lstlisting}[style=l4style]
make_operator :: Config -> Op[Tensor[N] -> Tensor[N]]
\end{lstlisting}

\noindent and read it as ``give me a configuration, I return an operator value you
apply later.'' \emph{Constructing} the operator is partial application: we supply
the \code{Config} argument now and \emph{freeze} it, getting back a more specialized
value awaiting the run-time vector (\cref{fig:opcurry}). The
\lstinline[style=l4style]|Op[in -> out]| brackets are the same reader-intent marker
as the closure-returning parentheses of \cref{ch:language}---they announce ``the
result is a thing you hold and apply''---with one addition: they promise the held
thing carries \emph{baked-in parameters}, the configuration, as a first-class,
named part of the value rather than as an invisible capture.

\begin{figure}[t]
\centering
\begin{tikzpicture}[node distance=9mm]
  \node[font=\bfseries\small, simInk] at (0,2.5)
    {Construction is partial application: freeze the configuration, await the vector};
  % stage 0: the constructor, needs config + vector
  \node[draw=simTeal, fill=simBgGray, thick, rounded corners=2pt,
        minimum width=24mm, minimum height=11mm] (f0) at (-4.6,1.0) {};
  \node[font=\ttfamily\scriptsize] at (-4.6,1.32) {make\_operator};
  \node[font=\scriptsize, simTeal] at (-4.6,0.78) {needs cfg, $v$};
  \node[font=\scriptsize\itshape,simGray,below=0.5mm of f0,align=center]
    {\code{Config ->}\\\code{Op[N -> N]}};
  % supply config
  \draw[flow] (-3.3,1.0) -- node[above,font=\scriptsize]{supply \code{cfg}} (-1.6,1.0);
  % stage 1: the constructed operator, config frozen
  \node[draw=simBlue, fill=simBgBlue, thick, rounded corners=2pt,
        minimum width=24mm, minimum height=11mm] (f1) at (-0.2,1.0) {};
  \node[font=\ttfamily\scriptsize] at (-0.2,1.32) {apply\_A = make\_operator cfg};
  \node[font=\scriptsize, simBlue] at (-0.2,0.78) {needs $v$};
  \node[font=\scriptsize\itshape,simGray,below=0.5mm of f1] {\code{Op[N -> N]}};
  \node[draw=simRust,fill=simBgRust,rounded corners=1pt,font=\scriptsize,
        inner sep=1.5pt] at (-0.2,2.05) {\code{cfg} frozen in};
  % supply v
  \draw[flow] (1.0,1.0) -- node[above,font=\scriptsize]{supply $v$} (2.6,1.0);
  % stage 2: a vector
  \node[draw=simRust, fill=simBgRust, thick, rounded corners=2pt,
        minimum width=16mm, minimum height=11mm] (r) at (3.8,1.0) {$\mat{A}v$};
  \node[font=\scriptsize\itshape,simGray,below=0.5mm of r] {\code{Tensor[N]}};
\end{tikzpicture}
\caption[Constructing an operator is partial application]{A constructed operator is
currying (\cref{ch:language}) raised to first class. \code{make\_operator} needs a
configuration \emph{and}, later, a vector. Supplying the configuration does not
produce a vector---it produces a \emph{new value}, the constructed operator
\code{apply\_A}, with the configuration \emph{frozen in}. Supplying a vector to that
operator finally produces a result. The two arrows are the two phases of
\cref{fig:constructapply}; the frozen \code{cfg} is the operator's immutable internal
state. The \lstinline[style=l4style]|Op[...]| spelling differs from a plain closure
only in advertising that the frozen part is a first-class, named configuration.}
\label{fig:opcurry}
\end{figure}

\subsection{Application is pure on the internals}

The word ``pure'' is doing real work in ``applied as a pure function,'' and it is
worth pinning down, because it is exactly the property that makes constructed
operators safe to share.

Application reads the operator's internal state and never writes it. So two
applications of the same operator cannot interfere: \code{apply\_A x1} and
\code{apply\_A x2} each read the frozen factorization, each produce a fresh result,
and neither disturbs the other or the operator. This is the immutability discipline
of \cref{ch:language} applied to the operator's internals: there is nothing to clone
and nothing to protect, because the internals cannot be overwritten by any
application. An operator, once constructed, is as inert and shareable as the number
$7$.

\begin{pitfall}
``Pure on the internals'' is not the same as ``stateless.'' A constructed operator
\emph{does} carry state---a whole factorization, perhaps megabytes of it. The
discipline is that this state is \emph{set at construction and never changed
thereafter}. If you find yourself wanting an operator to update something between
applications---an iteration count, a running residual, a cached previous input---stop:
that thing is \emph{not} operator-internal state. It is sim state (\cref{ch:strata}),
and it must be threaded explicitly through the run-time stratum, not stuffed into the
operator. An operator that mutates between applications has stopped being a pure
function, and every guarantee in this chapter evaporates with it. We return to
exactly this failure---and why it is not a flaw but a real distinction---in
\cref{sec:decisionrule}.
\end{pitfall}

\section{Worked example: a constructed Jacobi preconditioner}

The cleanest possible constructed operator is the \emph{Jacobi} (diagonal)
preconditioner, and tracing it end to end makes every abstract claim above concrete.
We will meet it again, in its proper numerical setting, in \cref{ch:precond}; here it
is purely a vehicle for the construct-then-apply shape.

\begin{workedexample}{Building and applying a Jacobi preconditioner}{jacobi}
A preconditioner is an \emph{approximate inverse}: given a matrix $\mat{A}$, we want
a cheap operator $\mat{M}^{-1}$ with $\mat{M}^{-1}\approx \mat{A}^{-1}$, so that
applying it to a residual $\vr$ yields something close to the error. The Jacobi
choice is the cheapest possible: take $\mat{M}=\operatorname{diag}(\mat{A})$, the
diagonal of $\mat{A}$, and ignore everything off the diagonal. Then $\mat{M}^{-1}$ is
just ``divide each component by the corresponding diagonal entry.''

\textbf{Construction (once).} Extract the diagonal of $\mat{A}$ and invert it
componentwise, producing a single vector $\mathbf{d}$ of reciprocals. This is the entire
baked-in internal state:
\begin{lstlisting}[style=l4style]
make_jacobi :: Tensor[N, N] -> Op[Tensor[N] -> Tensor[N]]
make_jacobi A =
  let diag = extract_diagonal A in   -- read the N diagonal entries: O(N)
  let d    = reciprocal diag in      -- d[i] = 1 / A[i,i]: O(N)
  \r -> d .* r                       -- the closure: captures d, awaits r
\end{lstlisting}
The returned closure captures \code{d}---the operator's frozen internal state---and
awaits a residual. Construction touched $\mat{A}$ exactly twice (extract, invert),
each an $O(N)$ pass; it produced one length-$N$ vector. Notice what construction did
\emph{not} do: it never looked at any residual, because no residual exists yet.

\textbf{Application (many times).} Inside a solve, every iteration produces a fresh
residual $\vr_k$ and applies the preconditioner to it:
\begin{lstlisting}[style=l4style]
let pc = make_jacobi A in       -- build ONCE, before the loop
... inside the iteration ...
let z1 = pc r1 in               -- z1 = d .* r1 ; O(N), pure on d
let z2 = pc r2 in               -- z2 = d .* r2 ; O(N), reads the SAME d
let z3 = pc r3 in               -- ...
\end{lstlisting}
Each application is one componentwise multiply: $\vz_k = \mathbf{d} \odot \vr_k$, an $O(N)$
operation. Take $N=4$, $\operatorname{diag}(\mat{A})=[4,2,8,1]$, so $\mathbf{d} =
[\tfrac14,\tfrac12,\tfrac18,1]$. Then for $\vr_1=[8,6,16,3]$,
\[
  \vz_1 \;=\; \mathbf{d} \odot \vr_1 \;=\; \bigl[\tfrac14\!\cdot\!8,\ \tfrac12\!\cdot\!6,\
  \tfrac18\!\cdot\!16,\ 1\!\cdot\!3\bigr] \;=\; [2,\,3,\,2,\,3],
\]
and for a \emph{different} residual $\vr_2=[4,4,4,4]$ on the \emph{same} operator,
\[
  \vz_2 \;=\; \mathbf{d} \odot \vr_2 \;=\; \bigl[1,\,2,\,\tfrac12,\,4\bigr].
\]

\textbf{Three readings, each a habit from \cref{ch:language}.} (1) \emph{The phases
are separate and asymmetric in cost.} Construction did two $O(N)$ passes over the
matrix, once. Each application does one $O(N)$ multiply, but there are thousands of
them; the matrix is never re-examined. The reciprocal $\mathbf{d}$ was computed once and
amortized over the whole solve. (2) \emph{Application is pure on the captured
\code{d}.} Computing $\vz_1$ did not change $\mathbf{d}$; $\vz_2$ read exactly the same
$\mathbf{d}$; the two could be computed in either order, or in parallel, without affecting
each other (\cref{fig:exprtree}'s ``no hidden order,'' now applied to operator
application). (3) \emph{Nothing was overwritten.} \code{r1}, \code{r2}, and \code{d}
all name exactly the values they did before each call; each \code{z} is a brand-new
vector.

The Jacobi preconditioner is the constructed-operator pattern in its purest form: the
expensive-but-once part (form $\mathbf{d}$) is frozen into the operator, and the cheap part
(multiply by $\mathbf{d}$) is repeated. Every richer preconditioner in
\cref{ch:precond,ch:multigrid}---an incomplete factorization, a multigrid V-cycle---is
the \emph{same shape} with more elaborate baked-in state.
\end{workedexample}

\section{The decision rule: construction-bound versus per-step variants}
\label{sec:decisionrule}

We now reach the load-bearing distinction of the chapter---the one rule that, applied
correctly, tells you whether a piece of configuration belongs inside a constructed
operator or out in the threaded state. Getting it right is the difference between a
clean solve and a subtle, convergent, \emph{wrong} one.

\begin{keyidea}
A variant that is \emph{bound once at construction} and held fixed for the whole
solve $\to$ a constructed operator absorbs it cleanly. A variant that \emph{changes
per step}---a different value on each iteration---\emph{cannot} be a constructed
operator, because there is no single value to freeze at construction. The per-step
variant must instead enter the threaded \code{SimState} (\cref{ch:strata,ch:monad}).
This single rule is exactly why \emph{flexible} preconditioners exist, and exactly
why \code{FGMRES} is a different algorithm from \code{GMRES}.
\end{keyidea}

The reasoning is almost a tautology once stated. Construction happens \emph{once},
before the loop. To freeze a value at construction, that value must \emph{exist}, and
be \emph{final}, at that moment. A preconditioner $\mat{M}$ chosen at solve start and
held fixed satisfies this: it exists, it is final, freeze it. But a preconditioner
$\mat{M}_k$ that is allowed to be \emph{different on each step $k$}---a so-called
\emph{flexible} preconditioner---does not: at construction time the later $\mat{M}_k$
do not yet exist, and there is no single $\mat{M}$ to bake in. The construction phase
has nothing to capture.

When this happens, the variant has not vanished---it has merely moved. A value that
changes every step \emph{is} run-time state, by definition. So it belongs in the
run-time stratum: it must be threaded through \code{SimState}, evolving step by step
like the iterate and the residual, rather than frozen in the operator. The
construction-vs-run-time stratification of \cref{ch:strata} is precisely the test:
\emph{ask which stratum the variant's lifetime puts it in.}

\begin{figure}[t]
\centering
\begin{tikzpicture}[node distance=8mm]
  % decision diamond
  \node[draw=simGold, fill=simBgGold, thick, diamond, aspect=2.2, inner sep=1pt,
        align=center, font=\footnotesize] (q) at (0,2.6)
    {Does the variant\\change \emph{per step}?};
  % NO branch (left): constructed operator
  \node[stage, minimum width=30mm, align=center, below left=12mm and 2mm of q] (con)
    {\footnotesize bind \emph{once} at construction\\[1pt]\textbf{constructed operator}};
  \node[draw=simBlue, fill=simBgBlue, semithick, rounded corners=2pt,
        minimum width=34mm, align=center, font=\scriptsize, below=5mm of con] (conex)
    {build $\mat{M}^{-1}$ once;\\\code{apply(pc, r)} each step;\\\code{op} never changes};
  \draw[flow] (con) -- (conex);
  \draw[flow] (q.south west) to[out=210,in=90]
    node[left,font=\scriptsize\bfseries,simGreen]{NO} (con.north);
  % YES branch (right): threaded state
  \node[product, minimum width=30mm, align=center, below right=12mm and 2mm of q] (thr)
    {\footnotesize thread it through state\\[1pt]\textbf{SimState schema grows}};
  \node[draw=simRust, fill=simBgRust, semithick, rounded corners=2pt,
        minimum width=34mm, align=center, font=\scriptsize, below=5mm of thr] (threx)
    {add per-step $\mat{M}_k$ data\\to the carry (\code{Z} basis);\\evolves every iteration};
  \draw[flow] (thr) -- (threx);
  \draw[flow] (q.south east) to[out=-30,in=90]
    node[right,font=\scriptsize\bfseries,simRust]{YES} (thr.north);
  % captions
  \node[font=\scriptsize\itshape, simGray, below=4mm of conex, text width=36mm, align=center]
    {fixed preconditioner;\\GMRES, CG};
  \node[font=\scriptsize\itshape, simGray, below=4mm of threx, text width=36mm, align=center]
    {flexible preconditioner;\\FGMRES};
\end{tikzpicture}
\caption[The decision rule]{The one rule that places a variant. Ask: does it change
per step? If \emph{no} (a preconditioner chosen at solve start and held fixed), bind
it once at construction---a constructed operator absorbs it, and the iteration calls
\code{apply(pc, r)} uniformly. If \emph{yes} (a flexible preconditioner $\mat{M}_k$
differing each step), no construction can freeze it; the variant must enter the
threaded \code{SimState} and evolve every iteration. The right branch is why
\code{FGMRES} carries an extra basis the left branch does not---the subject of
\cref{wex:flexible} and, in full, \cref{ch:gmres}.}
\label{fig:decisionrule}
\end{figure}

\subsection{Worked example: fixed versus flexible, in two signatures}

The decision rule is easiest to feel as a contrast between two signatures. We put the
fixed preconditioner and the flexible one side by side and watch the variant move out
of the operator and into the state.

\begin{workedexample}{Why a flexible preconditioner cannot be a constructed
operator}{flexible}
A Krylov solver (\cref{ch:gmres}) builds, step by step, a basis $\{\vv_0, \vv_1,
\dots\}$ of search directions, applying the preconditioner to each. We compare two
ways the preconditioner can enter.

\textbf{Case 1: a fixed preconditioner (a constructed operator works).} The
preconditioner $\mat{M}$ is chosen once, at solve start. We construct it before the
loop and apply it---the \emph{same} operator---at every step:
\begin{lstlisting}[style=l4style]
-- the preconditioner is a constructed operator, built once:
let pc = make_preconditioner config A in   -- pc :: Op[Tensor[N] -> Tensor[N]]
-- each Arnoldi step applies the SAME pc to a fresh basis vector:
let z_j = pc v_j in                        -- z_j = M^{-1} v_j, same M every j
let w   = matvec A z_j in ...
\end{lstlisting}
Here \code{pc} appears once, at construction, and is then applied uniformly. The state
the loop threads is the basis $V = [\vv_0, \dots, \vv_j]$ and some bookkeeping; the
preconditioner is \emph{not} part of that state, because it never changes. The
constructed operator absorbs the entire ``which preconditioner'' question. This is
the left branch of \cref{fig:decisionrule}.

\textbf{Case 2: a flexible preconditioner (it must move into the state).} Now the
preconditioner is allowed to differ each step---$\mat{M}_j$ at step $j$ (perhaps it is
itself an inner iterative solve whose accuracy varies, perhaps it adapts to the
residual). There is no single $\mat{M}$ to construct. The application becomes
\begin{lstlisting}[style=l4style]
-- there is NO single pc to construct; M_j is produced per step:
let z_j = apply_M_step j v_j in   -- z_j = M_j^{-1} v_j, a DIFFERENT M_j each j
\end{lstlisting}
and---this is the crucial consequence---the per-step preconditioned vectors $\vz_j$
can no longer be recomputed on demand from a frozen operator, because the operator
that made them is gone by the next step. They must be \emph{remembered}: the solver
threads a \emph{second} basis $Z = [\vz_0, \dots, \vz_j]$ alongside $V$. Contrast the
two state schemas:
\begin{lstlisting}[style=l4style]
-- GMRES  (fixed pc): the preconditioner is a constructed operator, NOT in state
type GmresState  = { x: Vec, V: Vec[], H: Mat, ... }
-- FGMRES (flexible pc): the per-step preconditioned basis Z IS threaded state
type FgmresState = { x: Vec, V: Vec[], Z: Vec[], H: Mat, ... }
--                                     ^^^^^^^ the extra basis the flexible case forces
\end{lstlisting}
The solution update reflects the move exactly: GMRES recovers the answer from the
plain basis, $\vx = \vx_0 + V_m \vy_m$, while FGMRES must use the remembered
preconditioned basis, $\vx = \vx_0 + Z_m \vy_m$. The variant that the constructed
operator absorbed in Case~1 has, in Case~2, become a field of the threaded state.

\textbf{The reading.} Nothing went wrong; a real distinction surfaced. A fixed
preconditioner is build-time configuration---freeze it, the operator is a value. A
flexible preconditioner is run-time state---thread it, the carry grows by one basis.
The decision rule (\cref{fig:decisionrule}) reads the variant's \emph{lifetime} and
routes it to the matching stratum. ``\code{FGMRES} exists because the preconditioner
became per-step state'' is the whole story, and \cref{ch:gmres} is where we tell it in
full.
\end{workedexample}

\begin{pitfall}
The tempting mistake is to try to keep a flexible preconditioner inside a constructed
operator by letting the operator \emph{mutate}---``I'll just update \code{pc}'s
internal $\mat{M}$ each step.'' This re-introduces exactly the mutable box
\cref{ch:language} spent a whole section banishing. A constructed operator that
mutates between applications is no longer pure on its internals, no longer safely
shareable, and no longer has a well-defined ``value'' at all---``what does \code{pc}
hold \emph{now}?'' becomes a live question again. The discipline is firm: per-step
variation goes in the threaded state, never in a secretly-mutating operator. The cost
is one honest extra field in the state schema; the alternative is a silent corruption
of the value model.
\end{pitfall}

\section{Solver as operator: an inverse is a function you apply}

We have been calling a preconditioner ``an approximate inverse'' and applying it like
an operator. It is time to make the type-level move that justifies this explicit,
because it is one of the most powerful abstractions in the whole numerical stack.

\begin{keyidea}
A \emph{solver} is an operator. A solver computes (or approximates) the action
$\vv \mapsto \mat{M}^{-1}\vv$ of an inverse, and that action has \emph{exactly the
same shape} as the action $\vv \mapsto \mat{M}\vv$ of the operator it inverts: vector
in, vector out. So a solver can wear the same \code{apply} interface as a plain
operator, and be substituted anywhere an operator is expected. ``Apply the
preconditioner to the residual'' is just \code{apply(pc, r)}---a function call,
indistinguishable in form from \code{apply(A, x)}.
\end{keyidea}

This sounds almost too simple to be worth a name, but its consequence is enormous. A
Krylov method (\cref{ch:gmres}) is written against the \code{apply} interface and
nothing else. It asks the operator $\mat{A}$ for matrix--vector products
\code{apply(A, v)}, and it asks the preconditioner for $\code{apply(pc, r)} \approx
\mat{M}^{-1}\vr$---and that is \emph{all} it ever asks of either. The iteration never
sees the \emph{algorithm} inside the preconditioner. It does not know, and cannot
ask, whether \code{apply(pc, r)} runs a diagonal multiply, an incomplete
factorization, a multigrid V-cycle, or a direct sparse solve. It sees a black box
that turns a residual into an approximate error, and it iterates.

\begin{figure}[t]
\centering
\begin{tikzpicture}[node distance=10mm]
  % the GMRES driver
  \node[stage, minimum width=26mm, minimum height=22mm] (gm) at (0,0)
    {\textbf{GMRES}\\[2pt]\footnotesize asks only:\\\code{apply(A, v)}\\\code{apply(pc, r)}};
  % the operator black box (top)
  \node[draw=simBlue, fill=simBgBlue, thick, rounded corners=2pt,
        minimum width=22mm, minimum height=9mm, right=24mm of gm, yshift=8mm] (A)
        {\footnotesize operator $\mat{A}$};
  % the preconditioner black box (bottom) -- SAME interface
  \node[draw=simViolet, fill=simBgViolet, thick, rounded corners=2pt, densely dashed,
        minimum width=22mm, minimum height=9mm, right=24mm of gm, yshift=-8mm] (pc)
        {\footnotesize solver \code{pc}};
  % the SAME apply arrow to both
  \draw[flow] (gm.east|-A) -- node[above,font=\scriptsize]{\code{apply}} (A.west);
  \draw[flow] (gm.east|-pc) -- node[below,font=\scriptsize]{\code{apply}} (pc.west);
  % what's INSIDE pc, hidden
  \node[align=left, font=\scriptsize\itshape, simGray, text width=36mm, right=4mm of pc]
    {inside (hidden): Jacobi?\\ILU? multigrid V-cycle?\\direct solve?\\\textbf{GMRES cannot tell.}};
  \node[align=left, font=\scriptsize\itshape, simGray, text width=30mm, right=4mm of A]
    {a true operator action\\$v \mapsto \mat{A}v$};
  % brace: both wear the same interface
  \draw[decorate, decoration={brace, amplitude=4pt}, simGray]
    ($(A.north west)+(-0.2,0.2)$) -- ($(pc.south west)+(-0.2,-0.2)$)
    node[midway, left=5pt, font=\scriptsize\bfseries, simInk, align=center]
    {same\\\code{apply}\\interface};
\end{tikzpicture}
\caption[A solver wears the operator's interface]{The solver-as-operator move. The
\code{GMRES} driver is written against one interface---\code{apply}---and asks for two
kinds of action through it: matrix--vector products from the true operator $\mat{A}$,
and approximate-inverse actions from the preconditioner \code{pc}. Because a solver
\emph{is} an operator (same vector-in, vector-out shape), \code{pc} wears the very
same \code{apply} interface as $\mat{A}$, and the iteration cannot see---and never asks
about---the algorithm hidden inside it. This is exactly why the same \code{GMRES}
drives a Jacobi preconditioner, an incomplete factorization, and a multigrid V-cycle
without changing a line: the differences are all behind the black box.}
\label{fig:solveroperator}
\end{figure}

\Cref{fig:solveroperator} draws the picture. The payoff is the one the whole book
keeps returning to: \emph{write the coordination once, feed it the verbs}
(\cref{ch:language}). Here the coordination is the Krylov iteration; the verbs are
\code{apply(A, $\cdot$)} and \code{apply(pc, $\cdot$)}; and because both verbs have the
same type, the \emph{same} iteration drives every combination of operator and
preconditioner the simulator offers. The driven solver, the electrostatic solver, and
the eigenmode solver's inner solves all reuse one \code{GMRES}---they differ only in
which operator and which preconditioner they hand it, and those are just two
\code{apply}-shaped values.

\subsection{Why this is a rotation, not a renaming}

It is fair to ask whether ``a solver is an operator'' says anything at all, or merely
relabels one thing as another. It says something, and the test is whether any
\emph{state} is hidden by the move.

At the source level (\cref{ch:strata} called this $L_0$), a solver is genuinely a
different kind of thing from a bare operator: it carries extra machinery---a routine to
set which operator it inverts, an initial-guess slot, convergence tolerances, internal
iteration state. The solver-as-operator rotation \emph{hides all of that} from the
iteration. The Krylov loop sees only the \code{apply} action; the solver's
operator-setting routine, its tolerances, its guess slot are compressed away,
invisible behind the interface. Because real state is hidden---because the loop
\emph{cannot} reach the preconditioner's internal apparatus even though it is
there---the move is a genuine \emph{rotation} (\cref{ch:rotations}), not a typedef.

The contrast sharpens it. If ``solver is an operator'' meant only that we wrote
\code{Solver} as an alias for \code{Operator}, with no extra structure underneath,
then the statement would hide nothing and buy nothing. It is precisely because the
solver \emph{does} carry distinct apparatus---which the iteration is then prevented
from touching---that wearing the operator interface is an act of abstraction. The
chapter on rotations (\cref{ch:rotations}) makes this ``does it hide state?'' test the
formal criterion for an honest rotation.

\section{The constructed-operator factory: where the variant is consumed}

So far each constructed operator has had one configuration. But a real solver must
choose \emph{among} configurations: which Krylov method (\code{CG} / \code{GMRES} /
\code{FGMRES}), which preconditioner (Jacobi / AMS / algebraic multigrid / a direct
factorization), which side, which orthogonalization. These choices arrive as
\emph{enums}---discrete flags read from a config file. Where do the enums get
consumed?

The answer is the \emph{constructed-operator factory}: a single function that takes
the enum (and any context it needs), dispatches on it internally, and returns a
\emph{uniformly typed} operator. The factory is the one place the variant is read.
Before it, the code branches on ``which method?''; after it, the code holds one
operator value and applies it---variant-free.

\begin{figure}[t]
\centering
\begin{tikzpicture}[node distance=9mm]
  % the enum coming in
  \node[data, align=center, text width=24mm] (enum)
    {variant enum\\\code{CG | GMRES |}\\\code{FGMRES | ...}};
  % the factory
  \node[stage, minimum width=34mm, minimum height=16mm, right=12mm of enum] (fac)
    {\textbf{factory}\\[1pt]\footnotesize \code{ConfigureKrylovSolver}\\\footnotesize dispatch \emph{here, once}};
  \draw[flow] (enum) -- (fac);
  % a uniformly-typed operator out
  \node[draw=simBlue, fill=simBgBlue, very thick, rounded corners=2pt,
        minimum width=24mm, minimum height=12mm, right=14mm of fac] (op)
        {\footnotesize\ttfamily Op[N -> N]};
  \draw[flow] (fac) -- node[above,font=\scriptsize]{one type} (op);
  % variant-free code after
  \node[draw=simGreen, fill=simBgGreen, thick, rounded corners=2pt,
        minimum width=30mm, minimum height=12mm, right=12mm of op, align=center,
        font=\footnotesize] (after)
        {downstream code:\\\code{apply(op, v)}\\\textbf{no enum in sight}};
  \draw[flow] (op) -- (after);
  % the "wall" after which variants are gone
  \draw[simRust, very thick, densely dashed] ($(fac.east)+(0.7,1.7)$) -- ($(fac.east)+(0.7,-2.0)$);
  \node[font=\scriptsize\bfseries, simRust, align=center] at ($(fac.east)+(0.7,-2.4)$)
    {variant\\wall};
  \node[font=\scriptsize\itshape, simGray] at ($(enum)+(0,-2.0)$) {variant-bearing};
  \node[font=\scriptsize\itshape, simGray] at ($(after)+(0,-2.0)$) {variant-free};
\end{tikzpicture}
\caption[The factory: a variant enum in, a uniform operator out]{The
constructed-operator factory is the single site at which a variant enum is consumed.
To its left, code carries the enum and must branch on it; the factory dispatches on
the enum \emph{once} and returns an operator of one uniform type. To its right---past
the ``variant wall''---downstream code holds that single operator and applies it with
\code{apply(op, v)}, never re-inspecting the enum. Palace's \code{ConfigureKrylovSolver}
(Krylov method) and \code{ConfigurePreconditionerSolver} (preconditioner type) are exactly
this: enum in, uniform solver-operator out, variant-free code after.}
\label{fig:factory}
\end{figure}

\Cref{fig:factory} draws the factory as a ``variant wall.'' Palace has two of them,
structurally identical: \code{ConfigureKrylovSolver} consumes the
Krylov-method$\,\times\,$side$\,\times\,$orthogonalization$\,\times\,$restart axes and
returns an iterative-solver operator; \code{ConfigurePreconditionerSolver} consumes the
preconditioner-type$\,\times\,$multigrid$\,\times\,$scalar-field axes and returns a
preconditioner operator. Both read a clutch of enums, dispatch once, and hand back a
single uniformly typed value the rest of the solver applies without ever asking which
method it got. A solver's signature schematically:

\begin{lstlisting}[style=l4style]
-- the factory: ONE place the krylov-method enum is read
configure_krylov :: KrylovMethod -> Config -> Op[Tensor[N] -> Tensor[N]]
-- downstream, the method is gone; only the uniform operator remains:
let ksp = configure_krylov method cfg in   -- dispatch happens HERE, once
let x   = apply ksp b in                    -- ... and never again
\end{lstlisting}

This is the procedural heart of \emph{variant absorption}, which the next section
makes into a discipline. For now hold the single observation: the factory localizes
the entire ``which variant?'' decision to one call. Everywhere else in the solver, an
operator is just an operator.

\section{Capability typing: distinguishing roles without runtime cost}

A factory hands back ``a uniformly typed operator.'' But uniformity can go too far.
Consider Palace's preconditioned Krylov solver: it holds \emph{two} operators of the
\emph{same} underlying type \code{Op[N -> N]}---the \emph{true operator} $\mat{A}$ it is
solving against, and the \emph{preconditioner-assembly operator} from which the
preconditioner is built. Both are \code{Op[N -> N]}. Nothing in their type stops a
careless caller from passing them in the wrong order. And if they \emph{are} swapped,
the solve does not crash---it runs, and converges, to the \emph{wrong answer}. A wrong
solve that converges is the worst kind of bug: silent.

\emph{Capability typing} is the discipline that turns this silent runtime bug into a
loud compile-time type error, at zero runtime cost.

\begin{keyidea}
A \emph{capability type} is a \emph{phantom-typed brand}: a marker attached to a value
that records the value's intended \emph{role}, enforced by the type checker but with
no runtime representation---no extra field, no allocation, no dispatch. Two values can
share an underlying type yet carry distinct brands, so that using one where the other
is required is a \emph{type error} rather than a silently-wrong computation. The brand
costs nothing at runtime and is erased before the program runs; its entire job is to
make a class of role-confusion mistakes \emph{unspeakable}.
\end{keyidea}

\subsection{The brand, concretely}

Concretely, we take the underlying operator type and attach a phantom tag that
distinguishes the role:

\begin{lstlisting}[style=l4style]
-- underlying type, shared:
type Op[N -> N] = ...
-- two BRANDED views, distinct at the type level, identical at runtime:
type TrueOp       = Op[N -> N] & { __cap: "true" }         -- the operator we solve against
type PcAssemblyOp = Op[N -> N] & { __cap: "pc_assembly" }  -- the operator the pc is built from
-- smart constructors: runtime identities, type-level commitments
as_true_op        :: Op[N -> N] -> TrueOp
as_pc_assembly_op :: Op[N -> N] -> PcAssemblyOp
\end{lstlisting}

\noindent The tag \code{\_\_cap} is \emph{phantom}: there is no such field at runtime,
no byte of storage, no branch. The smart constructors \code{as\_true\_op} and
\code{as\_pc\_assembly\_op} are the identity function at runtime---they return their
argument untouched. Their \emph{only} job is to commit a value to a role at the moment
of construction, so that a downstream signature can demand a specific brand:

\begin{lstlisting}[style=l4style]
set_operators :: TrueOp -> PcAssemblyOp -> Solver -> Solver
\end{lstlisting}

\noindent Now passing a \code{PcAssemblyOp} where a \code{TrueOp} is required is
rejected by the type checker---a compile-time error, caught before any number is
computed---even though the two values are byte-for-byte identical at runtime. The
permitted ``escape hatch'' is intended and useful: when you genuinely want to
precondition the true operator against itself, you brand the \emph{one} operator
\emph{both} ways (\code{as\_true\_op A} and \code{as\_pc\_assembly\_op A}); the
discipline forbids passing one brand where the other is wanted, not applying both
brands to one value.

\begin{figure}[t]
\centering
\begin{tikzpicture}[node distance=8mm]
  % LEFT: variant absorption -- hides axes of variation
  \begin{scope}
    \node[font=\bfseries\small, simInk] at (0,2.7) {Variant absorption};
    \node[font=\scriptsize\itshape, simGray] at (0,2.3) {hides axes of \emph{variation}};
    % the uniform operator, with variants collapsed inside
    \node[draw=simBlue, fill=simBgBlue, thick, rounded corners=2pt,
          minimum width=30mm, minimum height=22mm] (u) at (0,0.6) {};
    \node[font=\ttfamily\scriptsize] at (u.north) [yshift=-3mm] {Op[N -> N]};
    \node[font=\scriptsize, simGray, align=center] at (u.center)
      {\code{CG}? \code{GMRES}?\\\code{Jacobi}? \code{AMG}?\\\emph{collapsed inside,}\\\emph{one type out}};
    \node[font=\scriptsize\itshape, simGray, below=2mm of u, text width=40mm, align=center]
      {many configurations $\to$ \emph{one} uniform type};
  \end{scope}
  % RIGHT: capability typing -- distinguishes axes of role
  \begin{scope}[shift={(7.4,0)}]
    \node[font=\bfseries\small, simInk] at (0,2.7) {Capability typing};
    \node[font=\scriptsize\itshape, simGray] at (0,2.3) {distinguishes axes of \emph{role}};
    % two branded operators, same underlying type, kept apart
    \node[draw=simTeal, fill=simBgGreen, thick, rounded corners=2pt,
          minimum width=26mm, minimum height=8mm] (tt) at (0,1.25) {};
    \node[font=\ttfamily\scriptsize] at (tt.center) {TrueOp};
    \node[draw=simRust, fill=simBgRust, thick, rounded corners=2pt,
          minimum width=26mm, minimum height=8mm] (pp) at (0,-0.1) {};
    \node[font=\ttfamily\scriptsize] at (pp.center) {PcAssemblyOp};
    % they share an underlying type (a thin bracket)
    \draw[decorate, decoration={brace, amplitude=3pt, mirror}, simGray]
      ($(tt.east)+(0.15,0)$) -- ($(pp.east)+(0.15,0)$)
      node[midway, right=4pt, font=\scriptsize, simGray, align=left]
      {same\\underlying\\\code{Op[N -> N]}};
    % the "no swap" barrier between them
    \node[font=\scriptsize\bfseries, simRust] at (0,0.58) {swap = type error};
    \node[font=\scriptsize\itshape, simGray, below=4mm of pp, text width=40mm, align=center]
      {one type $\to$ \emph{two} brands kept apart};
  \end{scope}
\end{tikzpicture}
\caption[Capability brands versus variant axes]{The two disciplines are
complementary, pulling in opposite directions on the same operator type.
\emph{Left:} variant absorption \emph{hides} axes of \emph{variation}---\code{CG} vs
\code{GMRES}, Jacobi vs multigrid---collapsing many configurations behind one uniform
\code{Op[N -> N]}. \emph{Right:} capability typing \emph{distinguishes} axes of
\emph{role}---\code{TrueOp} vs \code{PcAssemblyOp}---splitting one underlying type into
two branded views the checker refuses to interchange. Absorption makes ``which
configuration?'' unspeakable downstream; branding makes ``which role?'' unswappable. A
single operator can be both absorbed (its variant hidden) and branded (its role
fixed).}
\label{fig:capvariant}
\end{figure}

\subsection{Background: where brands come from}

The phantom-brand idea is borrowed, lightly, from typed functional programming, and
you may meet it under several names. In Haskell, a \code{newtype} wraps an existing
type in a distinct, zero-cost name the compiler keeps separate---the same
representation, a different type. In Rust, a \emph{zero-sized marker type} (a field of
a type with no data, such as \code{PhantomData}) tags a value with a role that exists
only for the borrow checker and vanishes from the compiled binary. The word
\emph{capability} comes from the security literature, where a capability is an
unforgeable token of authority; here the ``authority'' is the role---holding a
\code{TrueOp} is the authority to act as the operator the iteration solves against,
holding a \code{PcAssemblyOp} is the authority to act as the source of the
preconditioner. We use brands in their type-checking-only form: they constrain what
the program can express, and then they are erased. \Citet{peytonjones2003} is the
standard reference for the underlying type machinery.

\subsection{Capability typing versus variant absorption}

Capability typing and variant absorption are easy to confuse because both are about
``many operators of the same type,'' but they pull in opposite directions, and seeing
the contrast cleanly is worth the effort (\cref{fig:capvariant}).

\begin{itemize}[nosep]
  \item \textbf{Variant absorption hides axes of \emph{variation}.} \code{CG} versus
  \code{GMRES}, Jacobi versus multigrid, left- versus right-side: these are
  \emph{differences we want to disappear} behind a uniform interface, so downstream
  code need not branch on them. Absorption \emph{collapses} many things into one type.
  \item \textbf{Capability typing distinguishes axes of \emph{role}.} True operator
  versus preconditioner-assembly operator: these are \emph{differences we want to
  preserve}, because confusing them is a bug. Branding \emph{splits} one type into
  several that the checker refuses to interchange.
\end{itemize}

A single operator can be \emph{both} at once. The preconditioner-assembly operator is
\emph{absorbed} along its variation axes---it is uniform across real versus complex
scalar field, uniform across which preconditioner algorithm sits inside it---and
simultaneously \emph{branded} \code{PcAssemblyOp}, distinct from \code{TrueOp} along
the role axis. Absorption answers ``which configuration is this?'' with ``don't ask,
it's uniform''; branding answers ``which role does this play?'' with ``that, and the
checker will hold you to it.'' They compose without conflict because they speak about
different axes of the same value.

\section{Variant absorption as a discipline}

We have now used ``absorb the variant'' informally several times. It is time to make
it a discipline with a checkable standard, because the failure mode---absorbing a
variant \emph{partly}, while pretending to absorb it fully---is one of the most
insidious ways a specification goes quietly wrong.

The principle: when a piece of the machinery has \emph{orthogonal axes of
variation}---preconditioner side, orthogonalization kind, fixed-versus-flexible,
real-versus-complex scalar field---the form should absorb those axes
\emph{parametrically}, as bindings of a parameter, rather than \emph{appending} them
as special-case paragraphs bolted onto a base description. ``Here is \code{GMRES};
oh, and for \code{FGMRES} replace this with that; oh, and for left-preconditioning add
this step'' is the anti-pattern. The parametric form names a parameter once and lets
its binding range over the variants.

\subsection{The three levels of absorption}

Absorption is not all-or-nothing; it holds at three nested levels, and a genuine
absorption holds at \emph{all three} (\cref{fig:absorptionlevels}).

\begin{description}[style=nextline,leftmargin=1.5em]
  \item[(a) Invariant level: the math unifies.] A single equation or recurrence covers
  every variant, perhaps with the variant appearing as a parameter. Example: write the
  solution update as $\vx = \vx_0 + W_m \vy_m$ with $W_m = V_m$ for \code{GMRES} and
  $W_m = Z_m$ for \code{FGMRES}. One equation; the parameter $W$ selects the variant.
  This is necessary but not sufficient.
  \item[(b) Procedural level: the variant is mentioned once.] The procedure names the
  variant at most once---at the binding or dispatch site (the factory call)---and never
  re-inspects it. ``Let \code{op = configure(side); then each step apply(op, v)}'' is
  procedurally absorbed: \code{side} appears once, at construction. A procedure that
  re-tests the variant at the operator-application step \emph{and} the update step is
  \emph{not} procedurally absorbed, even if its math unifies.
  \item[(c) Primitive-sequence level: the operation chain has the same shape.] The
  underlying sequence of primitive calls---which operations, in which order---is
  identical across variants, with the variant binding only the \emph{operands}, not the
  chain. ``$[\text{matvec},\,\text{dot},\,\text{axpy}]$ for every \code{side}, with the
  matvec's operator argument changing'' is primitive-sequence absorbed. A variant that
  forces one branch to run an extra primitive the others skip is \emph{not}.
\end{description}

\begin{figure}[t]
\centering
\begin{tikzpicture}[node distance=4mm]
  % three nested bands, (a) outermost
  \node[band, fill=simBgGold, draw=simGold, thick, minimum width=82mm, minimum height=46mm]
    (a) at (0,0) {};
  \node[font=\footnotesize\bfseries, simGold!75!black, anchor=north] at (a.north) [yshift=-1.5mm]
    {(a) invariant level --- one equation covers all variants};
  \node[band, fill=simBgBlue, draw=simBlue, thick, minimum width=66mm, minimum height=32mm]
    (b) at (0,-2.5mm) {};
  \node[font=\footnotesize\bfseries, simBlue, anchor=north] at (b.north) [yshift=-1.5mm]
    {(b) procedural --- variant named once};
  \node[band, fill=simBgGreen, draw=simGreen, thick, minimum width=50mm, minimum height=18mm]
    (c) at (0,-5mm) {};
  \node[font=\footnotesize\bfseries, simGreen, anchor=north] at (c.north) [yshift=-1.5mm]
    {(c) primitive-sequence};
  \node[font=\scriptsize, simInk, align=center] at (c.center) [yshift=-1mm]
    {same op chain,\\operands differ};
  % full absorption note
  \node[font=\scriptsize\itshape, simGray, anchor=north, text width=74mm, align=center]
    at (a.south) [yshift=-3mm]
    {\textbf{Full absorption} = all three nested levels hold. The dangerous case is
    reaching (a) but silently \emph{not} (b) or (c)---the form \emph{looks} uniform but
    the operation chain secretly branches.};
\end{tikzpicture}
\caption[The three nested levels of variant absorption]{Absorption holds at three
nested levels. \emph{(a)} The mathematical statement unifies into one
parameter-bearing equation. \emph{(b)} The procedure mentions the variant at most
once, at its binding site, and never re-inspects it. \emph{(c)} The underlying
primitive-call sequence has the same shape across variants, the variant binding only
the operands. Full absorption requires all three. The insidious failure is partial,
\emph{undisclosed} absorption---reaching (a) while quietly violating (b) or (c)---which
the next pitfall names.}
\label{fig:absorptionlevels}
\end{figure}

\subsection{The anti-pattern: silent partial absorption}

The three levels exist so that we can name precisely the way absorption goes wrong.

\begin{pitfall}
\textbf{Silent partial absorption} is the failure this entire discipline exists to
catch. A form achieves invariant-level absorption (a)---one clean equation, looks
uniform---but \emph{silently} fails procedural (b) or primitive-sequence (c)
absorption: under the surface, the procedure re-inspects the variant at several sites,
or the operation chain branches by variant, and the form's tidy appearance hides it.
The reader sees ``$\vx = \vx_0 + W_m \vy_m$, one equation'' and believes the variant
is handled; the actual computation forks three ways. The fix is honesty: if (b) or (c)
genuinely cannot hold---because the variant is structural, like a flexible
preconditioner forcing an extra threaded basis (\cref{wex:flexible})---then
\emph{disclose} the residual divergence explicitly, listing the sites where the
procedure or the primitive chain forks. A disclosed partial absorption is fine; a
\emph{silent} one is a specification that lies about its own uniformity.
\end{pitfall}

The routes to genuine full absorption are the constructions of this chapter. A
\emph{constructed operator} that internalizes the variant at construction time
achieves (b) and (c) at once: the factory reads the enum once, the operator is
applied uniformly, the primitive chain is one shape (\cref{fig:factory}). When even a
constructed operator cannot help---because the variant is per-step, and the decision
rule (\cref{fig:decisionrule}) routes it into the threaded state---then the absorption
is genuinely partial, and the discipline demands you say so, in the open, rather than
let a tidy equation imply a uniformity the operation chain does not have.

\begin{remark}
Notice how the pieces of this chapter fit together as one mechanism. The
\emph{factory} (\cref{fig:factory}) is \emph{where} a variant is absorbed---the single
site that reads the enum. The \emph{constructed operator} (\cref{fig:constructapply})
is \emph{what} absorbs it---the value that internalizes the choice and is then applied
uniformly. The \emph{decision rule} (\cref{fig:decisionrule}) tells you \emph{whether}
absorption can succeed at all---no, if the variant is per-step. And \emph{capability
typing} (\cref{fig:capvariant}) handles the orthogonal concern that survives
absorption---the \emph{roles} the uniform operators play, which must stay distinct even
after their variations are hidden. Absorption hides what varies; branding preserves
what must not be confused.
\end{remark}

\section{Putting it together: the preconditioned solver}

We can now read the whole construction surface of a real preconditioned Krylov solver
as one coherent picture, with every piece of this chapter in its place. Palace's
\code{BaseKspSolver} does exactly the following, and nothing in it should now be
mysterious.

At \emph{construction} (build-time stratum, run once): two factory calls
(\cref{fig:factory}). \code{ConfigureKrylovSolver} consumes the
Krylov-method$\,\times\,$side$\,\times\,$orthogonalization$\,\times\,$restart enums and
returns a constructed iterative-solver operator \code{ksp}; \code{ConfigurePreconditionerSolver}
consumes the preconditioner-type$\,\times\,$multigrid$\,\times\,$scalar-field enums and
returns a constructed preconditioner operator \code{pc}. Both are \code{apply}-shaped
(solver-as-operator, \cref{fig:solveroperator}); both have absorbed their variation
axes parametrically inside the factory; after the two calls, no enum survives into the
solve. The two operators are then \emph{branded}---the true operator \code{as\_true\_op},
the preconditioner-source \code{as\_pc\_assembly\_op} (\cref{fig:capvariant})---so the
binding step \code{set\_operators :: TrueOp -> PcAssemblyOp -> ...} cannot receive them
swapped.

At \emph{application} (run-time stratum, run many times): the solve folds the Krylov
iteration to convergence, and at each step it asks the two constructed operators for
their \code{apply} actions and nothing else---matrix--vector products from \code{ksp}'s
operator, approximate-inverse actions from \code{pc}. The iteration never re-examines
any enum, never sees either operator's internal algorithm, and never confuses their
roles, because the construction phase settled all three concerns once and froze the
result. The cleanly separated, frozen construction is exactly what makes the cheap,
uniform, repeated application possible.

This is the shape every solver in \cref{ch:gmres,ch:precond,ch:multigrid} wears. The Krylov methods of
\cref{ch:gmres}, the preconditioners of \cref{ch:precond}, and the multigrid V-cycle
of \cref{ch:multigrid} are all elaborations of this one surface: build constructed
operators by reading enums in a factory; brand their roles; thread only the genuinely
per-step state; apply the operators uniformly through a single \code{apply}
interface. Operators as values is not one technique among many---it is the frame the
entire numerical machinery hangs on.

\summarysection
A \emph{constructed operator} is a value built once, with its full configuration baked
into immutable internal state, then applied many times as a pure function. The two
phases live in two strata (\cref{ch:strata}): construction is expensive, build-time,
and run once; application is cheap, run-time, run many times, and pure on the frozen
internals. This is the first-class form of the currying and ``build once, apply many''
of \cref{ch:language}; the \lstinline[style=l4style]|Op[in -> out]| spelling marks a
held-and-applied value carrying named baked-in parameters. The \emph{decision rule} is
load-bearing: a variant bound \emph{once at construction} (a fixed preconditioner) is
absorbed by a constructed operator; a variant that changes \emph{per step} (a flexible
preconditioner $\mat{M}_k$) cannot be---there is no single value to freeze---so it must
enter the threaded \code{SimState} instead. This is exactly why \code{FGMRES} threads
an extra basis \code{Z} that \code{GMRES} does not (\cref{ch:gmres}). A \emph{solver is
an operator}: an approximate inverse wears the same \code{apply} interface as the thing
it inverts, so a Krylov method needs only the black-box action $\vv \mapsto
\mat{M}^{-1}\vv$ and never sees the preconditioner's algorithm---which is why one
\code{GMRES} drives wildly different preconditioners. The
\emph{constructed-operator factory} is the single site a variant enum is consumed and a
uniformly typed operator returned (\code{ConfigureKrylovSolver},
\code{ConfigurePreconditionerSolver}); after it, the code is variant-free.
\emph{Capability typing} attaches a phantom-typed brand recording a value's \emph{role}
(\code{TrueOp} vs \code{PcAssemblyOp}) at zero runtime cost, turning a role-swap into a
type error rather than a silent wrong-but-convergent solve; it is complementary to
\emph{variant absorption}, which hides axes of \emph{variation} where branding
distinguishes axes of \emph{role}. Variant absorption is a checkable discipline with
three nested levels---(a) the math unifies, (b) the variant is named once, (c) the
primitive chain has one shape---with full absorption requiring all three and
\emph{silent partial absorption} the anti-pattern it exists to catch.

\furtherreading
The constructed-operator and solver-as-operator patterns are standard in operator-algebra
preconditioning frameworks; \citet{saad2003} (especially the chapter on preconditioned
Krylov methods) motivates the central abstraction---that a Krylov method needs only the
black-box action $\vv \mapsto \mat{M}^{-1}\vv$, so any approximation respecting that
interface composes uniformly---and is the reference for the flexible-preconditioner
distinction that drives \code{FGMRES}. The phantom-type and zero-cost-brand machinery
behind capability typing is part of the type-system toolkit of typed functional
programming; \citet{peytonjones2003} is the standard reference for the underlying
mechanisms (\code{newtype} wrappers and phantom type parameters), which Rust mirrors
with zero-sized marker types. The two disciplines---hiding variation and preserving
role---meet again in \cref{ch:rotations}, where ``does the move hide real state?''
becomes the formal test that separates an honest rotation from a mere renaming.

\exercisessection
\begin{enumerate}[nosep]
  \item \textbf{Two phases, two costs.} For the constructed Jacobi preconditioner of
  \cref{wex:jacobi}, state the cost of \emph{construction} and the cost of a
  \emph{single application}, each in terms of $N$. If a solve runs $k$ iterations,
  write the total work as a function of $N$ and $k$, and identify which term is the
  amortized construction. For what range of $k$ is constructing the operator clearly
  worthwhile, and when would calling the inner multiply directly be just as good?
  \item \textbf{Apply the rule.} For each variant below, decide whether it belongs in a
  constructed operator or in the threaded \code{SimState}, and justify your answer with
  the decision rule (\cref{fig:decisionrule}): (a) the time-step size $\Delta t$ in a
  \emph{fixed}-step time integrator; (b) the preconditioner in a flexible Krylov
  method, where an inner solve of varying accuracy is used each step; (c) the matrix
  entries of an assembled stiffness operator, fixed for the whole solve; (d) an
  adaptive damping parameter recomputed from the residual at each iteration.
  \item \textbf{The mutating-operator trap.} A colleague proposes to support a flexible
  preconditioner by giving the constructed \code{pc} an \emph{update} method that
  overwrites its internal $\mat{M}$ each step, ``so we don't have to change the state
  schema.'' Explain, using the value model of \cref{ch:language}, exactly which
  guarantee this breaks, and what question becomes ill-posed once \code{pc} can mutate.
  What is the honest cost of doing it correctly instead?
  \item \textbf{Same interface, different insides.} The same \code{GMRES} drives a
  Jacobi preconditioner, an incomplete factorization, and a multigrid V-cycle without
  changing a line. Using \cref{fig:solveroperator}, name the single interface that
  makes this possible, state precisely what \code{GMRES} asks of the preconditioner,
  and explain what it is \emph{prevented} from asking. Why is being prevented from
  seeing the preconditioner's algorithm a feature, not a limitation?
  \item \textbf{Rotation or renaming?} Recall the test from \cref{ch:rotations} (and
  the discussion in this chapter): a move is a genuine rotation only if it \emph{hides
  real state}. The solver-as-operator move makes a \code{Solver} usable wherever an
  \code{Operator} is expected. List two pieces of apparatus a real solver carries that
  a bare operator does not, and explain why hiding them from the iteration makes this a
  rotation rather than a typedef. What would have to be true for it to be \emph{only} a
  renaming?
  \item \textbf{Where the enum dies.} A solver reads a config flag
  \code{method $\in$ \{CG, GMRES, FGMRES\}}. In a well-absorbed design, at how many sites
  in the per-step iteration does \code{method} appear, and where is the one place it is
  consumed? Sketch (in the notation) the factory call and one downstream apply, and mark
  the ``variant wall'' of \cref{fig:factory}. Then describe what a \emph{silent partial
  absorption} of this enum would look like and why a critic should reject it.
  \item \textbf{Brand the binding.} Two operators \code{A} and \code{P}, both of type
  \code{Op[N -> N]}, are to be passed to
  \lstinline[style=l4style]|set_operators :: TrueOp -> PcAssemblyOp -> Solver -> Solver|.
  (i) Write the two smart-constructor calls that brand them correctly. (ii) Explain why
  swapping the arguments is now a compile-time error even though \code{A} and \code{P}
  are byte-for-byte the same kind of value at runtime. (iii) Show how to legitimately
  precondition \code{A} against itself, and explain why the brand discipline permits
  this while still forbidding the accidental swap. (iv) What is the runtime cost of all
  this branding?
  \item \textbf{Absorption versus branding.} In one or two sentences each, state what
  \emph{variant absorption} does and what \emph{capability typing} does, and identify
  which axis of an operator each one speaks about. Then give a single operator that is
  simultaneously absorbed along one axis and branded along another, and name both axes
  explicitly (\cref{fig:capvariant}).
  \item \textbf{The three levels, by hand.} Consider preconditioner side
  $\in \{\text{LEFT},\text{RIGHT},\text{NONE}\}$ in a Krylov step. Suppose a draft
  writes the step as one equation (level (a)) but, under the surface, the LEFT branch
  computes \lstinline[style=l4style]|w = apply(pc, apply(A, v))|, the RIGHT branch
  \lstinline[style=l4style]|w = apply(A, apply(pc, v))|, and NONE just
  \lstinline[style=l4style]|w = apply(A, v)|. Which of levels (b) and (c) does this draft
  fail, and why? Describe the constructed operator that would repair both at once, and
  state what its construction freezes.
\end{enumerate}

\chapter{The Layered View: Rotations and Representations}
\label{ch:rotations}

\chapterorient{We have spent Part~III building a calculus---tensors and shapes,
the fold and its combinators, the state monad, operators as values. This closing
chapter steps back from the bricks to look at the building. The calculus is not
the simulator; it is one \emph{description} of the simulator, the topmost of
several. Below it sit the algebra, the loops, the pure functions, and finally the
cited C++ source. Each description is complete and correct in itself, and moving
from one to the next is a precise act we call a \emph{rotation}: a re-expression
that turns exactly one ``impedance''---mutation, fusion, iteration---while leaving
the computation untouched. By the end you will have the map of the whole project:
which layers exist, what each rotation buys, how to tell a genuine rotation from
an empty renaming, and---no less important---when a rotation honestly fails.}

\section{One simulator, several descriptions}

Throughout this book a single computation has worn many costumes. The
\texttt{axpy} update appeared first (\cref{ch:bigidea}) as a C loop overwriting an
array, then as a pure function returning a fresh tensor, then as an arrow in a
dependency graph. The transient solve appeared as a hand-written time-stepping
loop and again as a \code{fold\_solve} over a schedule. These were not different
computations. They were the \emph{same} computation, described at different
\emph{levels}.

That is the organizing idea of the entire project that this textbook documents.
We do not describe the simulator once. We describe it at a stack of levels, each a
faithful account of the machine, each written in a different vocabulary, and we
make the relationship between adjacent levels explicit and checkable. The
calculus you have just learned is the \emph{top} of that stack. This chapter
names all the levels, names the moves between them, and hands you the criteria
for telling a real move from a fake one.

\begin{keyidea}
The simulator is described at several \emph{layers}, $L_0$ through $L_4$, each
complete and correct \emph{in itself}. Adjacent layers are connected by a
\emph{rotation}: a re-expression that turns exactly one impedance---mutation,
fusion, iteration, or coordination---while computing the same numbers. The
calculus of Part~III is the top layer $L_4$; the physics of Part~II grounds the
bottom layers. The stack is the map; this chapter is the legend.
\end{keyidea}

\subsection{The five layers}

We number the layers from the ground up. The bottom is the literal source; the
top is the calculus. \Cref{fig:stack} is the hero picture of the whole chapter,
and indeed of the project---print it on the inside of your eyelids.

\begin{description}[leftmargin=2.2em,style=nextline]
  \item[$L_0$ --- the cited source.] The ground truth: ranges of Palace's
  C++ (and the MFEM library beneath it), quoted by file and line. Every claim
  anywhere in the stack ultimately rests on an $L_0$ citation. $L_0$ is not
  ``ours''; it is the machine as written. Nothing is interpreted here---it is read.

  \item[$L_1$ --- pure functions.] The same operations re-expressed as
  \emph{pure functions over immutable values}. The destination array that $L_0$
  overwrote in place becomes a returned value; the loop that scribbled over memory
  becomes a function from inputs to a fresh output. This is the
  \emph{mutation rotation} of \cref{ch:bigidea}, applied operator by operator.

  \item[$L_2$ --- algebraic decompositions.] Each $L_1$ function unfolded back
  into a \emph{composition of base algebraic primitives}, with the performance
  tricks erased. A fused kernel that computed $a x + b y + c z$ in one
  cache-friendly sweep is decomposed into the scaled sums it stands for; a
  packed, blocked matrix apply becomes the plain linear map it implements. This is
  the \emph{fusion rotation}.

  \item[$L_3$ --- global tensor-field operations.] Where the algebra permits,
  the explicit per-degree-of-freedom loops of $L_2$ are lifted to
  \emph{whole-field operations}: ``add these two fields,'' ``apply this operator
  to this field,'' with no surviving element loop in the spec form. This is the
  \emph{iteration rotation}---and, crucially, the one that does not always go
  through (\cref{sec:obstruction}).

  \item[$L_4$ --- the calculus.] The small, formally defined
  graph-evaluation calculus of Part~III: high-order combinators
  (\cref{ch:fold,ch:combinators}), the state monad (\cref{ch:monad}), the strata
  (\cref{ch:strata}), operators as values (\cref{ch:operators}). $L_4$ is
  \emph{vocabulary, not architecture}---it says what the algorithm \emph{is},
  abstractly, in a form an accelerator backend can consume. The
  \emph{coordination rotation} from $L_3$ replaces hand-written control flow with
  named combinators.
\end{description}

\begin{figure}[t]
\centering
\begin{tikzpicture}[node distance=0mm]
  % five stacked layer bands, L0 at bottom -> L4 at top
  \def\lw{86mm}   % layer width
  \def\lh{11mm}   % layer height
  \def\gap{8.5mm} % vertical gap between bands (room for rotation arrow)
  \foreach \i/\name/\desc/\col/\bg in {%
      0/$L_0$/cited C++ {\&} MFEM source/simInk/simBgGray,
      1/$L_1$/pure functions over immutable values/simGreen/simBgGreen,
      2/$L_2$/algebraic decompositions {\textbullet} tricks erased/simTeal/simBgGray,
      3/$L_3$/global tensor-field operations/simBlue/simBgBlue,
      4/$L_4$/calculus {\textbullet} combinators {\textbullet} monad {\textbullet} strata/simViolet/simBgViolet}{
    \node[draw=\col, fill=\bg, thick, rounded corners=2pt,
          minimum width=\lw, minimum height=\lh, align=center, font=\small]
      (L\i) at (0, \i*\dimexpr\lh+\gap\relax) {%
        \textbf{\name}\quad\textcolor{simGray}{\footnotesize\desc}};
  }
  % rotation arrows between adjacent bands
  \foreach \lo/\hi in {0/1,1/2,2/3,3/4}{
    \draw[depedge, line width=1pt] (L\lo.north) -- (L\hi.south);
  }
  % labels placed to the right of each arrow gap
  \node[font=\scriptsize\itshape, simBlue, anchor=west] at (0.50*\lw, 0.5*\lh+0.5*\gap) {mutation rotation};
  \node[font=\scriptsize, simGray, anchor=west] at (0.50*\lw, 0.5*\lh+0.5*\gap-3mm) {mutation $\to$ purity};
  \node[font=\scriptsize\itshape, simBlue, anchor=west] at (0.50*\lw, 1.5*\lh+1.5*\gap) {fusion rotation};
  \node[font=\scriptsize, simGray, anchor=west] at (0.50*\lw, 1.5*\lh+1.5*\gap-3mm) {fused kernel $\to$ algebra};
  \node[font=\scriptsize\itshape, simBlue, anchor=west] at (0.50*\lw, 2.5*\lh+2.5*\gap) {iteration rotation};
  \node[font=\scriptsize, simGray, anchor=west] at (0.50*\lw, 2.5*\lh+2.5*\gap-3mm) {loop $\to$ whole-field op};
  \node[font=\scriptsize\itshape, simBlue, anchor=west] at (0.50*\lw, 3.5*\lh+3.5*\gap) {coordination rotation};
  \node[font=\scriptsize, simGray, anchor=west] at (0.50*\lw, 3.5*\lh+3.5*\gap-3mm) {control flow $\to$ combinators};
  % side bracket: "complete & correct in itself"
  \draw[decorate, decoration={brace, amplitude=5pt}, simGray]
    (-0.52*\lw-2mm, -0.5*\lh) -- (-0.52*\lw-2mm, 4*\dimexpr\lh+\gap\relax+0.5*\lh);
  \node[rotate=90, font=\scriptsize\itshape, simGray, anchor=south]
    at (-0.52*\lw-7mm, 2*\dimexpr\lh+\gap\relax) {each layer complete \& correct in itself};
\end{tikzpicture}
\caption[The five-layer stack and its rotations]{The impedance-matching stack.
Read it from the ground up: $L_0$ is the cited source, and each rotation above it
turns exactly one impedance while computing the same numbers---mutation into
purity ($L_0\!\to\!L_1$), a fused kernel into its base algebra
($L_1\!\to\!L_2$), an explicit loop into a whole-field operation
($L_2\!\to\!L_3$), and hand-written control flow into named combinators
($L_3\!\to\!L_4$). Every layer is a complete, correct account of the simulator in
its own vocabulary; the arrows are translations between vocabularies, never mere
renamings. The book has been teaching the $L_4$ view, with Part~II supplying the
$L_0$--$L_1$ grounding.}
\label{fig:stack}
\end{figure}

\subsection{Why a stack and not a single description}

It would be simpler to describe the simulator once and have done. The stack earns
its complexity because \emph{each layer answers a different question well, and
none answers all of them}. $L_0$ answers ``what does Palace actually do''---it is
the only layer that can be wrong about the source, because it \emph{is} the
source. $L_4$ answers ``what algorithm is this, abstractly, and how would it lower
onto a GPU''---but it is too far from the source to verify against it directly.
The intervening layers are the stepping-stones that let us connect the two with
small, individually checkable moves. A direct $L_0 \to L_4$ leap would be a leap
of faith; the stack replaces it with a sequence of short, auditable steps.

This is why we call the stack \emph{impedance-matching}. In an electrical
transmission line, you do not connect a high-impedance source straight to a
low-impedance load---you insert matching sections that each step the impedance part
of the way, so power transfers without reflection. Here the ``impedance mismatch''
is between raw mutating C++ and an abstract whole-tensor calculus. Each rotation
matches one stage of it.

\section{What a rotation is}
\label{sec:whatis}

We have used the word loosely; now we pin it down. A rotation is the unit of
progress between layers. The metaphor is geometric: a rotation changes your
\emph{point of view} on an object without changing the object. Rotate a cube and
the face that was hidden comes forward, but it is the same cube, with the same
volume and the same corners.

\begin{definition}[Rotation]
\label{def:rotation}
A \emph{rotation} is a re-expression of a piece of work from layer $L_n$ into
layer $L_{n+1}$ that \emph{changes exactly one impedance} while computing the same
result. The impedance turned is one of: mutation (made into purity), a fused
kernel (unfolded into base algebra), an explicit loop (lifted to a whole-field
operation), or hand-written coordination (replaced by a named combinator).
\end{definition}

The single most important thing to understand about a rotation is what it is
\emph{not}.

\begin{keyidea}
A rotation is a \emph{translation across vocabularies}, not a one-to-one
\emph{renaming}. It re-expresses the lower form in a genuinely different idiom,
in which something becomes simpler, hidden, or substitutable. If the new form is
just the old form with the identifiers swapped---same state threaded, same
structure, same grain---no impedance turned, and \emph{no rotation occurred}. A
rotation turns one impedance; a renaming turns none.
\end{keyidea}

To make this concrete, contrast two purported $L_1 \to L_2$ moves on the
conjugate-gradient inner step. At $L_1$, one step reads
\[
  \vr \leftarrow \vr - \alpha\,\mat{A}\vp, \qquad
  \beta \leftarrow \tfrac{(\vr\cdot\vr)_{\text{new}}}{(\vr\cdot\vr)_{\text{old}}},\qquad
  \vp \leftarrow \vr + \beta\,\vp .
\]
A \emph{renaming} rewrites this as \lstinline[style=l4style]|axpby(alpha, Ap, -1, r)|;
\lstinline[style=l4style]|beta = dot(r,r)/dot(r_old,r_old)|;
\lstinline[style=l4style]|axpby(1, r, beta, p)|. Every line is the same operation
with a BLAS-1 name pasted on. The reader could not drop in a \emph{different
algorithm} where \code{axpby} sits---there is only the one operation it stands
for. Nothing was hidden, nothing compressed, nothing made substitutable. This is
the failure mode \cref{sec:tests} teaches you to detect.

A genuine rotation, by contrast, introduces a primitive
\lstinline[style=l4style]|cg_step| whose signature \emph{hides} the explicit
threading of $\vr_k \to \vr_{k+1}$ and into whose slot a \emph{different}
inner-step algorithm---\code{cgne\_step}, \code{minres\_step}---could be
substituted without the caller noticing. That is a translation into a new
vocabulary, not a renaming in the old one.

\section{The three tests for ``a rotation has occurred''}
\label{sec:tests}

How do you \emph{tell}? This is the analytical core of the chapter, and the one
piece of it you should carry out the door. A proposed rotation from $L_n$ to
$L_{n+1}$ is genuine if and only if \textbf{at least one} of three tests passes.
If none passes, you have a renaming---a smell, not a rotation. \Cref{fig:tests}
collects the tests as a decision diagram.

\begin{figure}[t]
\centering
\begin{tikzpicture}[node distance=7mm]
  \node[data, text width=30mm] (in) {a proposed $L_n\!\to\!L_{n+1}$ move};
  \node[stage, below=of in, text width=46mm, minimum height=10mm] (t1)
    {\textbf{T1.} Does it \emph{hide state}?\\[1pt]\footnotesize something mutable now internal/absorbed};
  \node[stage, below=8mm of t1, text width=46mm, minimum height=10mm] (t2)
    {\textbf{T2.} Does a \emph{different algorithm}\\fit the same slot?\\[1pt]\footnotesize the algorithmic-substitution test};
  \node[stage, below=8mm of t2, text width=46mm, minimum height=10mm] (t3)
    {\textbf{T3.} Is the \emph{threaded state}\\compressed or abstracted?};
  % outcome nodes on the right
  \node[product, right=20mm of t1, text width=24mm] (yes) {genuine\\rotation};
  \node[draw=simRust, fill=simBgRust, thick, rounded corners=2pt, align=center,
        font=\small\bfseries, text=simInk, right=20mm of t3, text width=24mm,
        minimum height=9mm] (no) {renaming\\(smell)};
  % flow
  \draw[flow] (in) -- (t1);
  \draw[flow] (t1) -- node[left,font=\scriptsize]{no} (t2);
  \draw[flow] (t2) -- node[left,font=\scriptsize]{no} (t3);
  \draw[flow] (t1.east) -- node[above,font=\scriptsize]{yes} (yes.west);
  \draw[flow] (t2.east) -| node[pos=0.25,above,font=\scriptsize]{yes} (yes.south);
  \draw[flow] (t3.east) -- node[above,font=\scriptsize]{yes} ($(t3.east)+(8mm,0)$) |- (yes.south west);
  \draw[flow] (t3.south) |- node[pos=0.75,below,font=\scriptsize]{no to all three} (no.west);
\end{tikzpicture}
\caption[The three rotation tests as a decision diagram]{The three tests for a
genuine rotation. A move passes if \emph{any one} of them is satisfied: it hides
mutable state (\textbf{T1}), it admits a different algorithm in the same slot
(\textbf{T2}, the algorithmic-substitution test), or it compresses the state
threaded through the computation (\textbf{T3}). Only when \emph{all three} fail is
the move a mere renaming---the smell this chapter teaches you to catch.}
\label{fig:tests}
\end{figure}

\subsection{T1 --- state hiding}

\emph{Does the higher form hide a piece of state that the lower form exposed?}

This is the most common way a rotation succeeds. At $L_n$ some buffer, index,
accumulator, or configuration is visible---threaded through the code, named in
signatures, mutated by hand. At $L_{n+1}$ it is \emph{absorbed}: tucked inside a
constructed operator, captured at construction time, or hidden behind a
primitive's contract.

The canonical example is a constructed operator (\cref{ch:operators}). A
preconditioner is configured once---factorizations computed, levels assembled,
tables filled---and then applied many times. At a low layer all that
configuration is visible state, threaded into every apply. The rotation
\emph{constructs} an operator that captures the configuration internally and
exposes only \lstinline[style=l4style]|apply :: LinOp -> Tensor[N] -> Tensor[N]|.
The construction-time state is hidden; the caller sees a clean field-to-field map.

\subsection{T2 --- coarser substitution (the algorithmic-substitution test)}

\emph{Could you swap the implementation in this slot for a genuinely different
algorithm without the caller noticing?}

This is the subtlest and most valuable test, so we give it its own name: the
\emph{algorithmic-substitution test}. The trap it guards against is mistaking a
\emph{name} for an \emph{interface}. Pasting the name \code{dot} onto an inner
product does not make a coarser substitution, because there is only one algorithm
called \code{dot}. The substitution must be \emph{algorithmic}, not nominal.

\begin{notationbox}
The algorithmic-substitution test, stated as a question to ask before claiming a
rotation: \emph{could a reader replace the $L_{n+1}$ primitive with a different
algorithm---not a different implementation of the same algorithm---and still
satisfy the $L_n$ contract?} If yes, a coarser substitution interface emerged and
the rotation is genuine. If ``no, it is just a named version of the one
operation,'' it is a renaming.
\end{notationbox}

A passing example: at $L_n$, choosing modified Gram--Schmidt over classical
Gram--Schmidt means re-threading per-vector dot accumulators and rewriting the
inner loop. At a genuine $L_{n+1}$ the choice collapses to a single substitution,
\lstinline[style=l4style]@orthog := mgs_step | cgs_step | cgs2_step@: three
\emph{distinct algorithms} occupying one primitive's slot, with whole-algorithm
correctness arguments factoring through the primitive's contract rather than its
body. \Cref{fig:substitution} draws the shape of a passing T2: two unlike
implementations, one interface.

\begin{figure}[t]
\centering
\begin{tikzpicture}[node distance=8mm]
  % the interface
  \node[opbox, minimum width=34mm, minimum height=9mm] (iface)
    {\lstinline[style=l4style]|apply :: LinOp -> Tensor[N] -> Tensor[N]|};
  \node[font=\scriptsize\itshape, simGray, above=1mm of iface] {one interface (the slot)};
  % caller above
  \node[data, above=9mm of iface, text width=26mm] (caller) {the caller\\(can't tell which)};
  \draw[flow] (caller) -- node[right,font=\scriptsize]{\code{apply}} (iface);
  % two unlike implementations below
  \node[extern, below left=11mm and 2mm of iface, text width=30mm, minimum height=11mm] (impA)
    {\textbf{Jacobi smoother}\\[1pt]\footnotesize diagonal scaling, one sweep};
  \node[extern, below right=11mm and 2mm of iface, text width=30mm, minimum height=11mm] (impB)
    {\textbf{multigrid V-cycle}\\[1pt]\footnotesize restrict, recurse, prolong};
  \draw[flow, densely dashed, simViolet] (iface.south) -- (impA.north);
  \draw[flow, densely dashed, simViolet] (iface.south) -- (impB.north);
  \node[font=\scriptsize\itshape, simGray, below=1mm of impA] {algorithm 1};
  \node[font=\scriptsize\itshape, simGray, below=1mm of impB] {algorithm 2};
  \node[align=center, font=\scriptsize\itshape, simGray, below=15mm of iface]
    {two \emph{different} algorithms, one slot $\Rightarrow$ T2 passes};
\end{tikzpicture}
\caption[The algorithmic-substitution test]{The shape of a passing
algorithmic-substitution test (T2). A single interface---here a preconditioner's
\code{apply}---sits between the caller and the implementation. Behind it can stand
genuinely \emph{different} algorithms: a one-line Jacobi diagonal scaling, or a
full multigrid V-cycle that restricts, recurses, and prolongs. The caller invokes
\code{apply} and cannot tell which is in the slot. Because a different algorithm
fits the same slot, the move is a genuine rotation. Contrast a renaming, where the
``slot'' admits only the one operation its name stands for.}
\label{fig:substitution}
\end{figure}

\subsection{T3 --- threaded-state compression}

\emph{Is the bundle of state threaded through each step strictly smaller, or at
least strictly more abstract, at the higher layer?}

A solver's inner loop may thread a long tuple of buffers: a Krylov basis array,
a packed Hessenberg matrix, Givens rotation accumulators, a column index, scratch
vectors, a restart counter. If the higher layer threads a short, abstract bundle
instead---\lstinline[style=l4style]|(iterate, residual, basis_handle, converged)|,
with the implementation buffers encapsulated behind \code{basis\_handle}---then
the threaded state compressed, and the rotation is genuine even if no single
piece was outright hidden and no algorithm became substitutable. The reader
manipulates fewer, more meaningful names.

\begin{pitfall}
The renaming smell is the trap to watch for: a proposed higher form that names
the lower form's operations as primitives but turns \emph{no} impedance. Symptom:
the higher form threads the \emph{same} buffers with the \emph{same} shapes,
admits only the \emph{same} algorithm in each slot, and exposes the \emph{same}
state. A genuine rotation may pass \emph{some} pieces through unchanged---an
\code{axpy} that is already at the right grain need not be ``rotated'' for its own
sake---but the anti-pattern is \emph{everything} carrying through with new names
and nothing moving forward. When you find it, do not record a hollow layer:
either merge (treat the lower form as the highest this slice has reached) or
redesign the move so one of the three tests genuinely fires.
\end{pitfall}

\section{Justification: how we know a rotation is sound}
\label{sec:justification}

Detecting that a rotation \emph{occurred} (\cref{sec:tests}) is separate from
showing it is \emph{correct}---that the higher form really computes what the lower
form does. Every rotation in the stack carries a \emph{justification}, and there
are exactly five kinds. \Cref{fig:justification} is the palette. The first four
are positive (the rotation goes through, here is why); the fifth is the negative
result that the next section is about.

\begin{description}[leftmargin=2.2em,style=nextline]
  \item[\texttt{algebraic}.] An algebraic identity makes the rotation manifest:
  associativity, distributivity, linearity. Unfolding a fused $a x + b y$ into two
  scaled sums is justified by the distributive law---no test required, the
  identity \emph{is} the proof.

  \item[\texttt{structural}.] A structural property of the primitives carries the
  rotation: the locality of an element-wise operation, the commutativity of an
  operator apply with a scalar scaling. The reasoning is about \emph{shape}, not
  arithmetic value.

  \item[\texttt{reduction\_chain}.] A sequence of small algebraic or structural
  steps, each individually named and sound, composed into the whole rotation. When
  no single identity does the job, a documented chain of them does.

  \item[\texttt{empirical\_match}.] An \emph{executed test} confirms the higher
  form matches the lower form numerically: run both on the same input, check the
  numbers agree to tolerance. \textbf{This is the preferred justification whenever
  it is available}---a green test beats a paragraph of algebra, because it cannot
  quietly be wrong about an edge case the prose forgot.

  \item[\texttt{obstruction}.] The rotation does \emph{not} go through, and the
  justification records \emph{why} as a first-class result. This is
  \cref{sec:obstruction}.
\end{description}

\begin{figure}[t]
\centering
\begin{tikzpicture}[node distance=4mm]
  \def\pw{30mm}
  \node[draw=simGreen, fill=simBgGreen, thick, rounded corners=2pt, align=center,
        font=\footnotesize, text=simInk, minimum height=13mm, text width=\pw] (alg)
    {\textbf{\code{algebraic}}\\[1pt]an identity makes it manifest};
  \node[draw=simTeal, fill=simBgGray, thick, rounded corners=2pt, align=center,
        font=\footnotesize, text=simInk, minimum height=13mm, text width=\pw,
        right=of alg] (str)
    {\textbf{\code{structural}}\\[1pt]a property of the shape carries it};
  \node[draw=simBlue, fill=simBgBlue, thick, rounded corners=2pt, align=center,
        font=\footnotesize, text=simInk, minimum height=13mm, text width=\pw,
        right=of str] (red)
    {\textbf{\code{reduction\_chain}}\\[1pt]named small steps, composed};
  \node[draw=simGold, fill=simBgGold, very thick, rounded corners=2pt, align=center,
        font=\footnotesize, text=simInk, minimum height=13mm, text width=\pw,
        below=6mm of str] (emp)
    {\textbf{\code{empirical\_match}}\\[1pt]\emph{run both; numbers agree}\\(preferred)};
  \node[draw=simRust, fill=simBgRust, very thick, rounded corners=2pt, align=center,
        font=\footnotesize, text=simInk, minimum height=13mm, text width=\pw,
        right=of emp] (obs)
    {\textbf{\code{obstruction}}\\[1pt]it does \emph{not} rotate; record why};
  \node[font=\scriptsize\itshape, simGray, left=of emp, text width=24mm]
    {a first-class\\negative result $\rightarrow$};
\end{tikzpicture}
\caption[The justification-kinds palette]{The five kinds of justification a
rotation can carry. The four upper/positive kinds certify that the rotation goes
through: an algebraic identity, a structural property, a chain of small named
steps, or---preferred whenever feasible---an empirical match, where the two layers
are run on the same input and their numbers are checked to agree. The fifth,
\code{obstruction} (rust), is the first-class \emph{negative} result: the rotation
genuinely does not go through, and recording \emph{why} is itself valuable
output.}
\label{fig:justification}
\end{figure}

\section{Obstructions: when a rotation honestly fails}
\label{sec:obstruction}

Not everything rotates. Some work, faithfully described at $L_2$, genuinely cannot
be lifted to a clean $L_3$ whole-field operation, because its very meaning depends
on doing things in order. When that happens we do \emph{not} fake a rotation and
we do \emph{not} treat the failure as an error in our method. We record an
\emph{obstruction}: a precise statement of what blocks the lift and why. An
obstruction is a finding, not a failure.

\begin{keyidea}
An \emph{obstruction} is a first-class negative result: a rotation that genuinely
does not go through, with the reason recorded as part of the spec. The obstructions
are not gaps in the work---they are some of its most valuable output, because they
say exactly where the abstract whole-tensor picture meets a wall, and why. ``It
does not rotate, and here is precisely why'' is knowledge.
\end{keyidea}

The archetype is the \emph{sequential reduction}---a loop whose $i$-th step reads
the $(i-1)$-th step's \emph{output}. \Cref{ch:tensors} introduced exactly this as
the case where the tensor-field lift fails: a Gauss--Seidel sweep updates
component $i$ using the \emph{already-updated} components $0,\dots,i-1$, so there
is no way to apply it to ``all entries at once.'' A whole-field operation has no
notion of ``the entry before this one''; the loop-carried dependency is the
obstruction. \Cref{fig:obstruction} contrasts a liftable elementwise update with
an unliftable sequential one.

\begin{figure}[t]
\centering
\begin{tikzpicture}[node distance=6mm]
  % ---- LEFT: liftable, independent ----
  \begin{scope}
    \node[font=\bfseries\small, simGreen] at (0,1.7) {Liftable: independent updates};
    \foreach \i in {0,1,2,3}{
      \node[opbox, minimum width=8mm, minimum height=6mm] (a\i) at (\i*0.95,0.7) {$y_\i$};
      \node[data, minimum width=8mm, minimum height=5mm] (x\i) at (\i*0.95,-0.4) {$x_\i$};
      \draw[flow] (x\i) -- (a\i);
    }
    \node[align=center, font=\scriptsize\itshape, simGray, text width=42mm] at (1.4,-1.25)
      {each $y_i = a x_i + y_i$ reads only its own input $\Rightarrow$ apply to all at once};
    \node[draw=simGreen, thick, rounded corners=2pt, fit=(a0)(a3), inner sep=3pt,
          fill=simGreen!8] {};
    \node[font=\scriptsize\bfseries, simGreen] at (1.4,1.15) {one whole-field op};
  \end{scope}
  % ---- RIGHT: unliftable, sequential ----
  \begin{scope}[shift={(6.3,0)}]
    \node[font=\bfseries\small, simRust] at (0,1.7) {Unliftable: sequential reduction};
    \foreach \i in {0,1,2,3}{
      \node[opbox, draw=simRust, minimum width=8mm, minimum height=6mm] (b\i) at (\i*0.95,0.7) {$y_\i$};
      \node[data, minimum width=8mm, minimum height=5mm] (u\i) at (\i*0.95,-0.4) {$x_\i$};
      \draw[flow] (u\i) -- (b\i);
    }
    % chain of carried dependencies
    \foreach \i/\j in {0/1,1/2,2/3}{
      \draw[backflow] (b\i.east) -- (b\j.west);}
    \node[align=center, font=\scriptsize\itshape, simGray, text width=42mm] at (1.4,-1.25)
      {$y_i$ reads the \emph{already-updated} $y_{i-1}$ $\Rightarrow$ no whole-field form};
    \node[font=\scriptsize\bfseries, simRust] at (1.4,1.15) {loop-carried dependency};
  \end{scope}
\end{tikzpicture}
\caption[An obstruction: a sequential reduction will not lift]{Why some loops
rotate to $L_3$ and some do not. \emph{Left:} an \code{axpy}-style update writes
$y_i = a x_i + y_i$ reading only its own input; the entries are independent, so the
loop lifts to a single whole-field operation. \emph{Right:} a sequential
reduction (a Gauss--Seidel sweep) writes $y_i$ using the \emph{already-updated}
$y_{i-1}$; the dashed rust arrows are the loop-carried dependency. A whole-field
operation has no ``previous entry,'' so the lift fails---and we record the
obstruction (with the offending dependency cited) rather than force a false
rotation. This is the tensor-field-lift failure of \cref{ch:tensors}.}
\label{fig:obstruction}
\end{figure}

A subtle and important variant: sometimes the loop we cannot rotate is one
\emph{we never wrote}. The eigenmode analysis (\cref{ch:targets}) hands its operators
to an opaque library routine---SLEPc's eigensolver---whose iteration lives entirely
in someone else's code. We cannot render that loop as a fold because Palace authors
no loop to render; we can only mark the obstruction at the library boundary. This
is still a genuine, recorded obstruction; it just sits at a wrapper, not at a
visible recurrence. We meet exactly this kind again in \cref{ch:eigenvalues}.

\begin{remark}
An obstruction is never an excuse to stop. It is a precise boundary marker: the
work \emph{below} it (the per-step body) usually rotates cleanly, and only the
\emph{loop} resists. The honest record says ``the body lifts; the iteration does
not, because of this cited dependency''---which is far more useful than either a
forced rotation or a blank silence.
\end{remark}

\section{Two worked examples}

We now run the machinery on two concrete cases: a clean state-hiding rotation, and
a clean algorithmic-substitution rotation set beside a renaming that fails all
three tests.

\begin{workedexample}{The three tests on the \texttt{axpy} mutation rotation}{axpy}
Take the $L_0 \to L_1$ \emph{mutation rotation} on \code{axpy}, the running
example of \cref{ch:bigidea}. At $L_0$, the cited C++ overwrites the destination
in place:
\begin{lstlisting}[style=l4style, language=]
for (int i = 0; i < n; i++)
    y[i] = a * x[i] + y[i];     // y is mutated; the name y survives, values do not
\end{lstlisting}
At $L_1$, the operation is a pure function returning a fresh tensor:
\begin{lstlisting}[style=l4style]
axpy :: Scalar -> Tensor[N] -> Tensor[N] -> Tensor[N]
axpy a x y = a * x + y         -- returns a FRESH tensor; x and y unchanged
\end{lstlisting}

Apply the three tests of \cref{sec:tests}.

\textbf{T1 (state hiding): passes.} At $L_0$ the destination array \code{y} is a
piece of \emph{mutable state}---a block of memory the loop scribbles over, whose
contents depend on how many iterations have run. At $L_1$ that mutable destination
is gone: \code{axpy} \emph{returns} its result, and its inputs \code{x} and
\code{y} are untouched afterward. The mutable destination-argument has been hidden
(in fact, eliminated): the impedance \emph{mutation} was turned into \emph{purity}.
That is precisely the $L_0\!\to\!L_1$ impedance of \cref{def:rotation}, and it is
the one test the mutation rotation is built to pass.

\textbf{T2 (coarser substitution): does not pass---and need not.} Could a reader
drop a \emph{different algorithm} into the \code{axpy} slot and still satisfy the
$L_0$ contract? No: \code{axpy} computes $a x + y$ and there is exactly one such
operation. The slot admits no algorithmic choice. T2 fails---but a rotation needs
only \emph{one} test to pass, and T1 already did. \code{axpy} is a leaf primitive,
already at the right grain; it is not where substitution interfaces live.

\textbf{T3 (threaded-state compression): not applicable.} \code{axpy} is a single
operation, not an iteration threading a state bundle; there is no per-step state to
compress. T3 has nothing to bite on here.

\textbf{Verdict.} \emph{One} test passes (T1, decisively), so the move is a
genuine rotation, not a renaming. Its justification (\cref{sec:justification}) is
\code{algebraic}: the equality of the in-place result and the returned value is the
identity $a x + y = a x + y$ read under two memory models---and it is also trivially
confirmable by \code{empirical\_match}, running both forms on the same \code{x},
\code{y}, \code{a} and checking the output array agrees entry for entry.
\end{workedexample}

\begin{workedexample}{Solver-as-operator: a passing T2 beside a failing renaming}{slot}
Now the algorithmic-substitution test in action, on a \emph{preconditioner behind
an apply interface}---the ``solver-as-operator'' pattern of \cref{ch:operators}.
A preconditioner $\mat{M}^{-1}$ is used inside a Krylov solve only through one
move: \emph{apply it to a vector}. The rotation packages it as a constructed
operator exposing a single method:
\begin{lstlisting}[style=l4style]
apply :: LinOp -> Tensor[N] -> Tensor[N]      -- the only thing the solver calls
\end{lstlisting}

\textbf{The passing rotation (T2).} Ask the algorithmic-substitution question:
could a \emph{different algorithm} occupy this \code{apply} slot while the Krylov
solver---the caller---stays byte-for-byte the same? Yes, and dramatically so. The
slot can hold a one-line Jacobi smoother (scale by the inverse diagonal); or a
Gauss--Seidel sweep; or a full geometric-multigrid V-cycle that restricts the
residual to a coarse mesh, recurses, and prolongs the correction back
(\cref{ch:multigrid}). These are \emph{wildly} different algorithms---different
work, different cost, different code---yet each satisfies the same contract
``approximately invert the operator,'' and the caller, which only ever writes
\lstinline[style=l4style]|z = apply(pc, r)|, cannot tell which is installed
(\cref{fig:substitution}). A different algorithm fits the same slot: T2 passes, and
this is a genuine rotation. (T1 also passes here---the preconditioner's setup state
is hidden inside the constructed operator---which is common: real rotations often
trip more than one test.)

\textbf{The failing renaming, for contrast.} Suppose instead we had ``rotated''
the diagonal-scaling preconditioner by simply renaming its inner expression
$z_i \leftarrow r_i / d_i$ to a primitive \code{diag\_scale} and calling it a new
layer:
\begin{lstlisting}[style=l4style]
diag_scale :: Tensor[N] -> Tensor[N] -> Tensor[N]   -- z[i] = r[i] / d[i]
\end{lstlisting}
Run the tests. \textbf{T1}: nothing is hidden---the diagonal \code{d} is still
threaded in as an explicit argument, just as before. \textbf{T2}: no different
algorithm fits the \code{diag\_scale} slot; it \emph{is} the one elementwise
division, named. \textbf{T3}: there is no threaded state bundle to compress; it is
a single stateless operation. All three fail. This is a renaming---the smell of
\cref{sec:tests}---and the right response is not to enshrine it as a layer but to
recognize \code{diag\_scale} as a leaf that simply \emph{carries through}, while
the real rotation happens one level up, at the \code{apply} interface that makes
the \emph{whole preconditioner} substitutable.

\textbf{The lesson.} The same component---a diagonal preconditioner---hosts both a
genuine rotation and an empty one. The rotation is genuine at the grain where
\emph{substitution becomes possible} (the \code{apply} interface, where Jacobi and
multigrid are interchangeable) and empty at the grain where it does not (renaming
one elementwise division). The three tests are what tell the two apart.
\end{workedexample}

\section{An advanced sidebar: erasure scope}
\label{sec:erasure}

\emph{This section is a depth note for the curious; a first reading may skip to
\cref{sec:payoff}.}

When an iteration rotation \emph{does} engage an obstruction
(\cref{sec:obstruction}), there is a finer question worth asking: \emph{how much}
of the operator's iteration view does the layer-hop erase? The answer ranges over
a small ladder---the \emph{erasure scope}---and naming the rung sharpens what each
obstruction is really about. \Cref{fig:erasure} draws the ladder.

\begin{figure}[t]
\centering
\begin{tikzpicture}[node distance=5mm]
  \node[band, fill=simBgGreen, draw=simGreen, thick, text width=68mm] (r1)
    {\textbf{1. single loop}\;---\;the whole operator \emph{is} one sequential loop
     (a linear solve's outer iteration). One recurrence erased.};
  \node[band, fill=simBgGray, draw=simTeal, thick, text width=68mm, below=of r1] (r2)
    {\textbf{2. nested double loop}\;---\;an inner recurrence inside an outer sweep,
     both sequential (a polynomial smoother). Two nested recurrences erased.};
  \node[band, fill=simBgRust, draw=simRust, thick, text width=68mm, below=of r2] (r3)
    {\textbf{3. opaque library call}\;---\;the loop lives entirely in
     someone else's code (an eigensolver). No loop to render: only a boundary
     marker erased.};
  % ascending "more iteration erased" arrow on the left
  \draw[-{Stealth[length=2.5mm]}, simGray, thick] (-3.95,-2.2) -- (-3.95,1.0);
  \node[font=\scriptsize\itshape, simGray, rotate=90, anchor=south]
    at (-4.3,-0.6) {more of the iteration view erased};
\end{tikzpicture}
\caption[The erasure-scope ladder]{How much iteration view a layer-hop erases,
from least to most opaque. \emph{Rung 1:} the operator is a single sequential
loop---one linear-solve outer iteration---and the hop erases that one recurrence.
\emph{Rung 2:} a nested double loop, an inner recurrence inside an outer sweep (a
degree-$k$ polynomial smoother), erasing two coupled recurrences. \emph{Rung 3:}
an opaque library call---an eigensolver whose loop is in third-party code---where
there is no loop to render at all, only a boundary marker to erase. The rung tells
you what kind of obstruction you are looking at and how it must be recorded.}
\label{fig:erasure}
\end{figure}

The three rungs differ in a structurally meaningful way. On rungs~1 and~2 Palace
\emph{authors} the recurrence, so the lower layer can \emph{render} it---an
explicit tail-recursive loop carrying a visible loop-carried dependency---and the
hop then \emph{erases} that rendered iteration view, the obstruction surviving only
as ``these algebraic laws do \emph{not} hold'' (you cannot reorder the steps).
Rung~3 is different in kind: Palace authors \emph{no} loop, so there is nothing to
render; the obstruction exists only as a marker at the library boundary, and the
hop erases the marker. The distinction is not pedantry---it changes how the
obstruction can ever be discharged. A rendered recurrence (rungs~1--2) records a
loop Palace chose to surface and could, in principle, be re-described; an
opaque-library marker (rung~3) records a boundary Palace never sees inside and that
will not be discharged unless the upstream architecture changes. We will see all
three rungs concretely in Part~IV: the single loop in \cref{ch:krylov}, the nested
double loop in \cref{ch:multigrid}, and the opaque eigensolver in
\cref{ch:eigenvalues}.

\section{The payoff, and the bridge to the machinery}
\label{sec:payoff}

Step back and see what the layered view has bought. By describing the simulator at
a stack of layers rather than once, and by making every move between layers a
rotation with a stated test and a stated justification, we get three things at
once. We get \emph{auditability}: a doubtful reader can check each short move
instead of swallowing one long leap from C++ to calculus. We get \emph{honesty}:
where the abstraction meets a wall, the obstruction is recorded, not papered over.
And we get a \emph{target}: the top layer $L_4$, the calculus you have learned, is
precisely the whole-tensor, no-aliasing form an accelerator backend wants to
consume---the implementation view promised back in \cref{ch:bigidea}.

That is the close of Part~III. You now hold the entire framework: the calculus
(\cref{ch:language,ch:records,ch:tensors,ch:shapes}), the fold and combinators
(\cref{ch:fold,ch:combinators}), the state monad and strata
(\cref{ch:monad,ch:strata}), operators as values (\cref{ch:operators}), and---this
chapter---the map of how all five layers relate. What remains is to watch the
rotations turn on \emph{real algorithms}.

\Cref{part:machinery} does exactly that. Its chapters are where the
explicit-loop view and the combinator view stop being a methodology and become
concrete machine. \Cref{ch:krylov} takes the conjugate-gradient and GMRES solvers
and shows their inner loops rotating from hand-threaded buffers up to
\code{krylov\_step} and \code{iterate\_while}---a single-loop erasure (rung~1 of
\cref{fig:erasure}) with a state-hiding rotation you can now name on sight.
\Cref{ch:multigrid} does the same for the multigrid preconditioner---the very
\code{apply} slot of \cref{wex:slot}---and meets the nested double loop
(rung~2). \Cref{ch:eigenvalues} reaches the eigensolver and meets the opaque
library obstruction (rung~3) head on. Each is a rotation, or an obstruction, of a
kind you have now learned to recognize. Turn the page, and watch them turn.

\summarysection
The simulator is described not once but at a stack of five \emph{layers}: $L_0$,
the cited source; $L_1$, pure functions over immutable values; $L_2$, algebraic
decompositions with performance tricks erased; $L_3$, global tensor-field
operations; and $L_4$, the calculus of Part~III. Each layer is complete and
correct in itself, and adjacent layers are joined by a \emph{rotation}---a
re-expression that turns exactly one impedance (mutation, fusion, iteration, or
coordination) while computing the same numbers. A rotation is a translation across
vocabularies, never a one-to-one renaming. Three tests decide whether a genuine
rotation occurred: it \emph{hides state} (T1), it admits a \emph{different
algorithm in the same slot} (T2, the algorithmic-substitution test), or it
\emph{compresses the threaded state} (T3); if all three fail, the move is a
renaming, a smell to be merged or redesigned away. Every rotation carries one of
five justifications---\code{algebraic}, \code{structural}, \code{reduction\_chain},
\code{empirical\_match} (preferred, because it runs both layers and checks the
numbers), or \code{obstruction}, the first-class negative result when a rotation
honestly does not go through. Obstructions---the sequential reductions and
opaque-library loops that will not lift to whole-field operations---are findings,
not failures, and their \emph{erasure scope} (single loop, nested loop, opaque
call) records how much iteration view a hop erases. The payoff is auditability,
honesty, and an accelerator-ready top layer. Part~IV now turns these rotations on
real numerical algorithms.

\furtherreading
The layered, rotation-based method this book documents is purpose-built, but its
two pillars are standard. The shift from mutable loops to pure, composable,
substitutable pieces is the functional-programming tradition (\cref{ch:bigidea}'s
further reading). The numerical algorithms whose rotations Part~IV makes
concrete---Krylov subspace solvers, multigrid, eigenvalue iterations---are
developed in Saad's \emph{Iterative Methods for Sparse Linear
Systems}~\citep{saad2003}, the standard reference for the solvers behind the
single-loop and nested-loop erasures, and in Trefethen and Bau's \emph{Numerical
Linear Algebra}~\citep{trefethenbau1997}, whose lecture-by-lecture development of
orthogonalization, conditioning, and the eigenvalue problem is the ideal companion
for seeing \emph{why} certain steps resist the iteration rotation. Read either with
the three tests of \cref{sec:tests} in hand, and the place where each algorithm's
loop becomes (or refuses to become) a whole-field operation will jump off the page.

\exercisessection
\begin{enumerate}[nosep]
  \item Name the impedance each of the four rotations in \cref{fig:stack} turns,
  in one phrase each. Which two layers does the \emph{iteration rotation} connect,
  and what is the one situation (\cref{sec:obstruction}) in which it refuses to go
  through?
  \item Classify each move as a genuine rotation or a renaming, and name the
  test(s) it passes or fails: (a) replacing an in-place
  \lstinline[style=l4style]|x.Add(alpha, y)| with a pure
  \lstinline[style=l4style]|axpy alpha y x|; (b) renaming the loop body
  \lstinline[style=l4style]|s += x[i]*y[i]| to a primitive \code{dot} with the same
  sequential accumulation; (c) packaging a preconditioner behind an \code{apply}
  interface that a Jacobi smoother \emph{or} a multigrid V-cycle could implement.
  \item The algorithmic-substitution test (T2) asks whether a \emph{different
  algorithm} could fit a slot. Explain why pasting the name \code{nrm2} onto the
  expression $\sqrt{\sum_i |x_i|^2}$ does \emph{not} pass T2, while wrapping a
  Krylov inner step as \code{cg\_step} (substitutable by \code{minres\_step})
  does. What is the difference between a \emph{name} and an \emph{interface}?
  \item Give the justification kind you would prefer for each rotation and say why:
  (a) unfolding $a x + b y$ into two scaled sums; (b) claiming a hand-rolled GMRES
  matches a reference implementation; (c) a chain of four small associativity and
  linearity steps. When is \code{empirical\_match} unavailable, forcing you back to
  an algebraic argument?
  \item A Gauss--Seidel sweep updates $y_i$ using the \emph{already-updated}
  $y_{i-1}$ (\cref{fig:obstruction}). Explain in one paragraph why this blocks the
  $L_2 \to L_3$ iteration rotation, and why recording the obstruction is more
  useful than either forcing a whole-field form or omitting the operator entirely.
  Which \emph{part} of the operator (body or loop) still rotates cleanly?
  \item Place each on the erasure-scope ladder (\cref{fig:erasure}) and justify the
  rung: (a) a conjugate-gradient solve's outer iteration; (b) a degree-3 Chebyshev
  polynomial smoother with an inner recurrence inside an outer Richardson sweep;
  (c) a call to SLEPc's eigensolver. For which one does Palace author \emph{no}
  loop at all, and what does that imply about whether the obstruction can ever be
  discharged?
  \item \emph{(Synthesis.)} Pick any operator you met in Part~II---a weak-form
  assembly, a boundary-condition elimination, a reduction to a physical
  quantity---and sketch its journey up the stack: what does its mutation rotation
  ($L_0\!\to\!L_1$) hide, does its iteration rotation ($L_2\!\to\!L_3$) succeed or
  hit an obstruction, and what combinator (\cref{ch:fold,ch:combinators}) names its
  coordination at $L_4$? You need not be exhaustive; the point is to walk one
  computation through all five costumes.
\end{enumerate}


\part{The Numerical Machinery}\label{part:machinery}
\chapter{Sparse Linear Systems and Krylov Subspaces}
\label{ch:krylov}

\chapterorient{Part~II handed us a linear system: discretize the weak form and
out comes $\mat{A}\vx=\vb$, with $\mat{A}$ enormous and almost entirely zeros.
Part~III handed us the machinery to run a loop without mutating anything, and the
insight that a solver is just an operator you apply. This chapter is where the two
meet. We ask why nobody solves a system this large by Gaussian elimination, why
the winning strategy instead builds the answer out of repeated \emph{applications}
of the operator---and why that is exactly the one thing the matrix-free operator of
\cref{ch:matrixfree} knows how to do. We meet the \emph{Krylov subspace}, the
nested ladder of approximations every method in Part~IV climbs; we isolate the
single \emph{step kernel} every Krylov method shares---one operator apply, a few
inner products, a few vector updates---and we read the whole solve as an
\code{iterate\_while} fold (\cref{ch:fold}) that produces a constructed operator
(\cref{ch:operators}): an approximate inverse you hold and apply. By the end you
will be able to place CG, GMRES, and their kin as specializations of one kernel,
and to say precisely which laws a solve obeys and which it deliberately does not.}

\section{The system Part~II built}
\label{sec:krylov-system}

Every analysis in this book passes through the same middle stage. Assembly
(\cref{ch:assembly,ch:galerkin}) turns the continuous weak form into a discrete
linear system
\[
  \mat{A}\,\vx \;=\; \vb,
\]
where $\vx\in\Real^N$ is the vector of unknown degree-of-freedom coefficients,
$\vb\in\Real^N$ is the discretized right-hand side, and $\mat{A}\in\Real^{N\times
N}$ is the assembled operator---the stiffness $\Kop$, the mass $\Mop$, or some
combination built from them. For the electrostatic and magnetostatic analyses
$\mat{A}$ is the symmetric positive-definite stiffness operator; for the driven
analysis it is the complex shifted system $\Kop+i\omega\Cop-\omega^2\Mop$; for a
single implicit time step of the transient analysis it is again an SPD combination.
In every case the \emph{shape} of the problem is the same: a large square system
to be solved for $\vx$.

Two features of this system govern everything that follows, and both come straight
from the finite-element construction of Part~II.

\paragraph{It is large.} $N$ is the number of degrees of freedom in the
discretization, and for a three-dimensional device meshed finely enough to
resolve the fields, $N$ runs from the hundreds of thousands into the tens of
millions. A production electromagnetic simulation routinely solves systems with
$N>10^6$. The matrix $\mat{A}$, if we wrote out all its entries, would be an
$N\times N$ array---$10^{12}$ numbers at $N=10^6$. We will see in a moment that we
never write it out.

\paragraph{It is sparse.} Here is the saving grace. The basis functions of the
finite-element method have \emph{local support}: each one is nonzero only on the
handful of mesh elements touching its node (\cref{ch:refelement}). The matrix entry
$A_{ij}=a(\phi_j,\phi_i)$ is an integral of a product of basis functions $i$ and
$j$, and that integral is zero unless the supports of $\phi_i$ and $\phi_j$
\emph{overlap}---unless nodes $i$ and $j$ share an element. Almost no two nodes do.
So almost every entry of $\mat{A}$ is exactly zero. Each row has only a small,
bounded number of nonzeros---a couple of dozen, independent of $N$---set by how
many neighbors a node has, not by how big the mesh is.

\begin{figure}[t]
\centering
\begin{tikzpicture}
\begin{axis}[
  width=0.62\linewidth, height=0.62\linewidth,
  title={\footnotesize sparsity pattern of $\mat{A}$ ($N=36$, from local support)},
  title style={yshift=-1mm},
  xmin=-0.5, xmax=35.5, ymin=-0.5, ymax=35.5,
  y dir=reverse, axis equal image,
  xtick=\empty, ytick=\empty,
  axis line style={simGray},
  xlabel={\footnotesize column $j$ (node)}, ylabel={\footnotesize row $i$ (node)},
  label style={font=\footnotesize},
]
% banded + local-support sparsity for a 6x6 grid laplacian-like stencil
\addplot[only marks, mark=square*, mark size=1.7pt, simBlue] coordinates {
  (0,0)(1,0)(6,0)
  (0,1)(1,1)(2,1)(7,1)
  (1,2)(2,2)(3,2)(8,2)
  (2,3)(3,3)(4,3)(9,3)
  (3,4)(4,4)(5,4)(10,4)
  (4,5)(5,5)(11,5)
  (0,6)(6,6)(7,6)(12,6)
  (1,7)(6,7)(7,7)(8,7)(13,7)
  (2,8)(7,8)(8,8)(9,8)(14,8)
  (3,9)(8,9)(9,9)(10,9)(15,9)
  (4,10)(9,10)(10,10)(11,10)(16,10)
  (5,11)(10,11)(11,11)(17,11)
  (6,12)(12,12)(13,12)(18,12)
  (7,13)(12,13)(13,13)(14,13)(19,13)
  (8,14)(13,14)(14,14)(15,14)(20,14)
  (9,15)(14,15)(15,15)(16,15)(21,15)
  (10,16)(15,16)(16,16)(17,16)(22,16)
  (11,17)(16,17)(17,17)(23,17)
  (12,18)(18,18)(19,18)(24,18)
  (13,19)(18,19)(19,19)(20,19)(25,19)
  (14,20)(19,20)(20,20)(21,20)(26,20)
  (15,21)(20,21)(21,21)(22,21)(27,21)
  (16,22)(21,22)(22,22)(23,22)(28,22)
  (17,23)(22,23)(23,23)(29,23)
  (18,24)(24,24)(25,24)(30,24)
  (19,25)(24,25)(25,25)(26,25)(31,25)
  (20,26)(25,26)(26,26)(27,26)(32,26)
  (21,27)(26,27)(27,27)(28,27)(33,27)
  (22,28)(27,28)(28,28)(29,28)(34,28)
  (23,29)(28,29)(29,29)(35,29)
  (24,30)(30,30)(31,30)
  (25,31)(30,31)(31,31)(32,31)
  (26,32)(31,32)(32,32)(33,32)
  (27,33)(32,33)(33,33)(34,33)
  (28,34)(33,34)(34,34)(35,34)
  (29,35)(34,35)(35,35)
};
\end{axis}
\end{tikzpicture}
\caption[Sparsity from local support]{The sparsity pattern of the assembled
operator $\mat{A}$ for a small structured mesh (each marked cell is a nonzero
entry; blanks are exact zeros). The pattern is a direct fingerprint of the mesh
connectivity: $A_{ij}\neq0$ only when basis functions $i$ and $j$ share an
element, which happens only for neighboring nodes (\cref{ch:refelement}). The
diagonal and a few near-diagonal bands carry the weight; the vast off-diagonal
plain is zero. Each row holds a bounded number of nonzeros regardless of $N$, so
the count of stored numbers grows like $O(N)$, not $O(N^2)$. This single
structural fact is what makes solving a million-unknown system possible at all.}
\label{fig:krylov-sparsity}
\end{figure}

\Cref{fig:krylov-sparsity} draws the pattern for a small structured mesh. The
nonzeros cluster on the diagonal and a few neighboring bands---the fingerprint of
the mesh's connectivity---and everything else is blank. The number of \emph{stored}
numbers is the number of nonzeros, which is $O(N)$: a constant number per row,
times $N$ rows. A million-unknown operator that would need $10^{12}$ entries dense
needs only a few times $10^7$ stored sparse. Sparsity is the only reason a problem
this large is tractable at all.

\begin{physicsbox}
The sparsity is not an accident of clever bookkeeping---it is the discrete shadow
of \emph{locality} in the physics. Maxwell's equations are local: the field at a
point is coupled directly only to the field in its immediate neighborhood, through
derivatives. The finite-element basis inherits that locality as compact support,
and the assembled operator inherits it as sparsity. A dense row in $\mat{A}$ would
mean a degree of freedom coupled directly to every other---a nonlocal physics the
weak form does not have. Sparsity is locality, written as a matrix.
\end{physicsbox}

\section{Why not just factor it?}
\label{sec:krylov-direct}

A first linear-algebra course solves $\mat{A}\vx=\vb$ by \emph{Gaussian
elimination}: reduce $\mat{A}$ to triangular form by row operations, then back-
substitute. Packaged as a reusable object this is the $LU$ factorization,
$\mat{A}=LU$ with $L$ lower- and $U$ upper-triangular; for the SPD systems of
electrostatics it specializes to the Cholesky factorization $\mat{A}=LL^{\mathsf
T}$. Once $L$ and $U$ are in hand, each solve is two cheap triangular sweeps. These
are the \emph{direct} methods: in exact arithmetic they produce the exact solution
in a finite, predetermined number of operations. Why not simply use one?

The answer is \emph{fill-in}, and it is the central fact that pushes large
finite-element solves away from direct methods. When elimination clears the entries
below a pivot, it adds multiples of the pivot row to the rows beneath. Those
additions create nonzeros where $\mat{A}$ had exact zeros: an entry $A_{ij}=0$ that
sits ``downstream'' of two nonzeros in the same elimination step becomes nonzero in
$U$. The factors $L$ and $U$ are therefore \emph{denser} than $\mat{A}$---often
dramatically so. A matrix with $O(N)$ nonzeros can have factors with
$O(N^{4/3})$, $O(N^{3/2})$, or worse nonzeros in three dimensions, and the
factorization work climbs faster still. For a 3-D mesh, sparse Cholesky typically
costs about $O(N^2)$ arithmetic and produces factors of about $O(N^{4/3})$
nonzeros---and a dense factorization, ignoring sparsity entirely, costs the
textbook $O(N^3)$.

\begin{figure}[t]
\centering
\begin{tikzpicture}
\begin{axis}[
  width=0.86\linewidth, height=62mm,
  xlabel={problem size $N$ (degrees of freedom)},
  ylabel={work (arithmetic operations)},
  xmode=log, ymode=log,
  xmin=1e3, xmax=1e7, ymin=1e3, ymax=1e15,
  xtick={1e3,1e4,1e5,1e6,1e7},
  ytick={1e3,1e6,1e9,1e12,1e15},
  legend pos=north west, legend cell align=left,
  grid=both, grid style={simGray!18},
  tick label style={font=\footnotesize}, label style={font=\small},
  legend style={font=\footnotesize},
]
% dense direct  O(N^3)
\addplot[simRust, very thick, mark=*, mark size=1.4pt, domain=1e3:1e7, samples=20]
  {x^3};
\addlegendentry{dense direct $\;O(N^3)$}
% sparse direct  ~ O(N^2)  (3-D fill-in)
\addplot[simGold, very thick, dashed, mark=triangle*, mark size=1.8pt, domain=1e3:1e7, samples=20]
  {x^2};
\addlegendentry{sparse direct $\;\sim O(N^2)$ (3-D)}
% iterative  preconditioned ~ near-linear
\addplot[simBlue, very thick, mark=square*, mark size=1.6pt, domain=1e3:1e7, samples=20]
  {60*x^1.16};
\addlegendentry{iterative (preconditioned) $\;\sim O(N^{1.16})$}
\end{axis}
\end{tikzpicture}
\caption[Direct versus iterative cost]{Work versus problem size $N$, on log--log
axes, for the three strategies. A \emph{dense} direct solve costs $O(N^3)$ and is
out of the question past a few thousand unknowns. A \emph{sparse} direct solve
exploits the sparsity of \cref{fig:krylov-sparsity} but still pays for fill-in,
costing about $O(N^2)$ in 3-D---tractable to perhaps $10^5$ unknowns, then
prohibitive. A well-\emph{preconditioned iterative} solve (Part~IV's subject)
costs close to $O(N)$ per solve: a bounded number of iterations, each doing $O(N)$
work through one sparse operator apply. The curves cross early; beyond a modest
size the iterative method is not merely faster, it is the only one that finishes.}
\label{fig:krylov-cost}
\end{figure}

\Cref{fig:krylov-cost} plots the three regimes. The dense $O(N^3)$ curve leaves
the chart almost immediately. Sparse direct, at roughly $O(N^2)$ in three
dimensions, survives to perhaps $10^5$ unknowns before fill-in makes both its
memory and its arithmetic unaffordable. The iterative curve---the subject of the
rest of Part~IV---grows close to linearly in $N$ when the iteration is well
preconditioned (\cref{ch:precond,ch:multigrid}): a bounded number of steps, each
costing one sparse operator apply at $O(N)$. The crossover comes early, and on the
far side of it the iterative method is not just the faster choice; for a
ten-million-unknown 3-D solve it is the \emph{only} choice that fits in memory and
finishes in time.

\begin{notationbox}
Direct methods are not banished from the simulator. They reappear, deliberately,
as \emph{components inside} iterative solves: a sparse direct factorization makes an
excellent preconditioner on a coarse problem (\cref{ch:precond}), and the coarsest
level of a multigrid hierarchy is often solved directly (\cref{ch:multigrid}). The
claim of this section is narrower and sharper: do not solve the \emph{full}
fine-scale system directly. Use a direct method only where the problem has been
made small enough that fill-in no longer bites.
\end{notationbox}

\subsection{What an iterative method asks of the operator}
\label{sec:krylov-action}

There is a second reason the iterative route wins, and it is the one that ties this
chapter back to \cref{ch:matrixfree}. A direct method must \emph{read the entries}
of $\mat{A}$: elimination operates on $A_{ij}$ one at a time. An iterative method
does not. Look ahead to any solver in Part~IV---conjugate gradients
(\cref{ch:cg}), GMRES (\cref{ch:gmres})---and watch what it does with $\mat{A}$. It
forms inner products of vectors, it takes linear combinations of vectors, it
checks norms. The \emph{only} way it ever touches $\mat{A}$ is through one
operation: given a vector $\vv$, produce the vector $\mat{A}\vv$. It never asks for
an entry $A_{ij}$, never asks for a row or a column, never inspects the sparsity
pattern. It needs the operator's \emph{action}, and nothing else:
\[
  \vv \;\longmapsto\; \mat{A}\vv .
\]

This is precisely the interface \cref{ch:matrixfree} built. The matrix-free
operator never stores $\mat{A}$; it computes $\mat{A}\vv$ on demand as a short
chain of element-local tensor contractions---gather, evaluate, weight, integrate,
scatter. An iterative solver is the perfect consumer of that interface: it asks for
$\mat{A}\vv$, the matrix-free operator computes it, and neither ever needs the
assembled array to exist. The two ideas were built for each other. A direct method
\emph{cannot} use a matrix-free operator---it needs the entries elimination
operates on---which is one more reason direct methods do not scale to the
matrix-free, high-order, accelerator-bound regime this book targets.

\begin{keyidea}
An iterative solver needs only the operator's \emph{action}, $\vv\mapsto\mat{A}\vv$
(and, with preconditioning, $\vr\mapsto\mat{M}^{-1}\vr$)---never the entries of
$\mat{A}$. This is the licence for the whole of Part~IV. It means the solver
composes directly with the matrix-free operator of \cref{ch:matrixfree}, which
supplies exactly that action and never forms the matrix; and it is why the cost of
a solve is counted in \emph{operator applies}, the one expensive thing each step
does. \emph{One apply per step} is the cost invariant the rest of the chapter makes
precise.
\end{keyidea}

\section{The Krylov subspace}
\label{sec:krylov-subspace}

If all we may do is apply $\mat{A}$ to vectors, what can we build? Start from the
right-hand side $\vb$ (or, more generally, from the initial residual
$\vr_0=\vb-\mat{A}\vx_0$). Apply $\mat{A}$ and you get $\mat{A}\vb$. Apply again:
$\mat{A}^2\vb$. Each application is one matrix-free action; after $m-1$ of them you
hold the family
\[
  \vb,\quad \mat{A}\vb,\quad \mat{A}^2\vb,\quad \dots,\quad \mat{A}^{m-1}\vb .
\]
The space these vectors span is the \emph{Krylov subspace} of dimension $m$.

\begin{definition}[Krylov subspace]
\label{def:krylov}
For a square operator $\mat{A}\in\Real^{N\times N}$ and a nonzero vector
$\vb\in\Real^N$, the $m$-th \emph{Krylov subspace} is
\[
  \mathcal{K}_m(\mat{A},\vb)
   \;=\; \operatorname{span}\bigl\{\,\vb,\ \mat{A}\vb,\ \mat{A}^2\vb,\ \dots,\
   \mat{A}^{m-1}\vb\,\bigr\}
   \;\subseteq\; \Real^N .
\]
The spanning vectors $\vb,\mat{A}\vb,\dots$ are the \emph{Krylov sequence}; each
new one costs exactly one application of the operator's action.
\end{definition}

A Krylov method builds its approximate solution \emph{out of} this subspace. At
step $m$ it picks the vector $\vx_m\in\vx_0+\mathcal{K}_m(\mat{A},\vr_0)$ that is
``best'' by some criterion the particular method chooses---minimizing the residual
norm (GMRES, \cref{ch:gmres}), minimizing the $\mat{A}$-norm of the error (CG,
\cref{ch:cg}). The subspaces nest,
\[
  \mathcal{K}_1 \subseteq \mathcal{K}_2 \subseteq \mathcal{K}_3 \subseteq \cdots,
\]
because each is the previous one plus one more power of $\mat{A}$. So the
approximation can only improve as $m$ grows: a best-in-$\mathcal{K}_{m+1}$ vector
is at least as good as the best-in-$\mathcal{K}_m$ vector, since the smaller
subspace sits inside the larger. The method climbs this ladder of nested
subspaces, one operator apply per rung, until the residual is small enough to stop.

\begin{figure}[t]
\centering
\begin{tikzpicture}[scale=1.0]
  % nested subspaces as nested ellipses, with the true solution as a star
  \draw[draw=simBlue, fill=simBgBlue, thick] (0,0) ellipse (4.1 and 2.5);
  \draw[draw=simTeal, fill=simBgGreen, thick] (-0.5,-0.1) ellipse (2.9 and 1.75);
  \draw[draw=simGold, fill=simBgGold, thick] (-1.0,-0.15) ellipse (1.7 and 1.05);
  % labels for the subspaces
  \node[font=\scriptsize\bfseries, simGold!70!black] at (-1.0,-0.15) {$\mathcal{K}_1$};
  \node[font=\scriptsize\bfseries, simTeal] at (1.35,0.95) {$\mathcal{K}_2$};
  \node[font=\scriptsize\bfseries, simBlue] at (3.0,1.85) {$\mathcal{K}_3$};
  \node[font=\scriptsize, simGray] at (4.55,2.3) {$\Real^N$};
  % the true solution
  \node[star, star points=5, star point ratio=2.3, draw=simRust, fill=simRust,
        inner sep=1.3pt] (sol) at (2.0,-1.4) {};
  \node[font=\scriptsize\bfseries, simRust, right=1mm of sol] {$\vx^\star$ (true solution)};
  % the iterates marching toward it
  \node[circle, fill=simInk, inner sep=1.2pt] (x1) at (-1.0,-0.15) {};
  \node[circle, fill=simInk, inner sep=1.2pt] (x2) at (0.5,-0.55) {};
  \node[circle, fill=simInk, inner sep=1.2pt] (x3) at (1.5,-1.05) {};
  \node[font=\scriptsize, simInk, above left=0mm and 0mm of x1] {$\vx_1$};
  \node[font=\scriptsize, simInk, above=0mm of x2] {$\vx_2$};
  \node[font=\scriptsize, simInk, above right=0mm and 0mm of x3] {$\vx_3$};
  \draw[flow, shorten >=2pt, shorten <=2pt] (x1)--(x2);
  \draw[flow, shorten >=2pt, shorten <=2pt] (x2)--(x3);
  \draw[flow, densely dashed, simRust, shorten >=2pt, shorten <=2pt] (x3)--(sol);
\end{tikzpicture}
\caption[The nested Krylov subspaces]{The Krylov subspaces nest,
$\mathcal{K}_1\subseteq\mathcal{K}_2\subseteq\mathcal{K}_3\subseteq\cdots$, each
adding one more power of $\mat{A}$ to the span and so growing by (at most) one
dimension per step. Each costs one operator apply to extend. The method's iterate
$\vx_m$ is the best approximation to the true solution $\vx^\star$ available within
$\vx_0+\mathcal{K}_m$; because the subspaces grow, the iterates can only improve,
marching toward $\vx^\star$ as the ladder climbs. In exact arithmetic the ladder
reaches $\vx^\star$ exactly once $\mathcal{K}_m$ is large enough to contain
it---at the latest when $m$ equals the degree of $\mat{A}$'s minimal polynomial
(\cref{sec:krylov-poly}).}
\label{fig:krylov-nested}
\end{figure}

\Cref{fig:krylov-nested} draws the picture: nested subspaces, and inside them a
sequence of iterates marching toward the true solution $\vx^\star$. The whole of
Part~IV is variations on how to choose the best iterate in each subspace and how to
do it cheaply; this section is about why the subspace contains a good answer at all.

\subsection{Why powers of \texorpdfstring{$\mat{A}$}{A} capture the solution}
\label{sec:krylov-poly}

It is not obvious that repeatedly applying $\mat{A}$ should march toward
$\mat{A}^{-1}\vb$---applying $\mat{A}$ looks like the opposite of inverting it. The
reason it works is a fact about polynomials in $\mat{A}$, and it is worth seeing
because it explains both why Krylov methods converge \emph{and} how fast.

Notice first that any vector in $\mathcal{K}_m(\mat{A},\vb)$ is, by definition, a
linear combination of $\vb,\mat{A}\vb,\dots,\mat{A}^{m-1}\vb$---that is, it is
$q(\mat{A})\,\vb$ for some polynomial $q$ of degree at most $m-1$:
\[
  \vx_m \in \mathcal{K}_m
  \quad\Longleftrightarrow\quad
  \vx_m = q(\mat{A})\,\vb
  \ \text{ for some } q \text{ with } \deg q \le m-1 .
\]
So ``find the best iterate in $\mathcal{K}_m$'' is the same as ``find the best
degree-$(m-1)$ polynomial $q$ such that $q(\mat{A})\vb\approx\mat{A}^{-1}\vb$.''
Krylov methods are, under the hood, polynomial approximations of the inverse.

Now the key fact. Every square matrix satisfies a \emph{minimal polynomial}: there
is a polynomial $p$ of some degree $d\le N$ with $p(\mat{A})=0$. Write
$p(t)=c_0+c_1t+\dots+c_d t^d$ with $c_0\neq0$ (true when $\mat{A}$ is invertible).
From $p(\mat{A})=0$ we can solve for the identity:
\[
  c_0 I = -\bigl(c_1\mat{A}+\dots+c_d\mat{A}^d\bigr)
  \;\Longrightarrow\;
  I = -\tfrac{1}{c_0}\bigl(c_1\mat{A}+\dots+c_d\mat{A}^d\bigr),
\]
and multiplying by $\mat{A}^{-1}$ gives the inverse as a \emph{polynomial in
$\mat{A}$ of degree $d-1$}:
\[
  \mat{A}^{-1} \;=\; -\tfrac{1}{c_0}\bigl(c_1 I + c_2\mat{A} + \dots + c_d\mat{A}^{d-1}\bigr).
\]
Therefore $\mat{A}^{-1}\vb=-\tfrac{1}{c_0}(c_1\vb+c_2\mat{A}\vb+\dots+c_d\mat{A}^{d-1}\vb)$
is exactly a vector in $\mathcal{K}_d(\mat{A},\vb)$. The true solution \emph{lives
inside} a Krylov subspace of dimension $d\le N$. A Krylov method that picks the
best iterate in each $\mathcal{K}_m$ therefore reaches the exact solution by step
$d$ at the latest---in exact arithmetic, a finite procedure, no worse than $N$
steps.

\begin{keyidea}
A vector in $\mathcal{K}_m(\mat{A},\vb)$ is exactly $q(\mat{A})\vb$ for a polynomial
$q$ of degree $<m$, so a Krylov method is a polynomial approximation of
$\mat{A}^{-1}$. The minimal polynomial guarantees $\mat{A}^{-1}\vb$ itself lies in
$\mathcal{K}_d$ for some $d\le N$: \emph{in exact arithmetic a Krylov method
terminates with the exact solution in at most $N$ steps}. But $N$ steps would be no
bargain---the method is useful only because it gets \emph{close} far sooner, and how
soon is governed by the spectrum of $\mat{A}$ (\cref{sec:krylov-convergence}).
\end{keyidea}

\subsection{Conditioning and the speed of convergence}
\label{sec:krylov-convergence}

The polynomial picture also tells us how \emph{fast} a Krylov method converges,
and the answer is the bridge to preconditioning. The error after $m$ steps is
controlled by how small a degree-$m$ polynomial $q$ (with $q(0)=1$, the
normalization the residual forces) can be made on the \emph{eigenvalues} of
$\mat{A}$. Roughly: the residual at step $m$ is bounded by
\[
  \frac{\norm{\vr_m}}{\norm{\vr_0}}
   \;\le\; \min_{\substack{q,\ \deg q\le m \\ q(0)=1}}\ \max_{\lambda\in\sigma(\mat{A})}\ \abs{q(\lambda)},
\]
where $\sigma(\mat{A})$ is the set of eigenvalues. Convergence is fast exactly when
a low-degree polynomial, pinned to $1$ at the origin, can be made small on the
whole spectrum.

The consequence is the single most important rule of thumb in iterative solving:
\emph{a clustered spectrum converges fast; a spread-out spectrum converges slowly.}
If all the eigenvalues huddle in a tight cluster away from zero, a low-degree
polynomial can dip small across the whole cluster, and a few steps suffice. If the
eigenvalues are spread over a wide range---if the \emph{condition number}
$\kappa=\lambda_{\max}/\lambda_{\min}$ is large---then no low-degree polynomial can
be small everywhere at once, and convergence crawls. For SPD systems the classic
bound makes this quantitative: CG reduces the error by a factor that improves like
$1-2/\sqrt{\kappa}$ per step, so the iteration count scales like $\sqrt{\kappa}$.

\begin{figure}[t]
\centering
\begin{tikzpicture}
\begin{axis}[
  width=0.86\linewidth, height=60mm,
  xlabel={iteration $m$},
  ylabel={relative residual $\norm{\vr_m}/\norm{\vr_0}$},
  ymode=log,
  xmin=0, xmax=60, ymin=1e-10, ymax=2,
  xtick={0,10,20,30,40,50,60},
  ytick={1e0,1e-2,1e-4,1e-6,1e-8,1e-10},
  legend pos=north east, legend cell align=left,
  grid=both, grid style={simGray!18},
  tick label style={font=\footnotesize}, label style={font=\small},
  legend style={font=\footnotesize},
]
% clustered spectrum: fast linear-on-log decay
\addplot[simBlue, very thick, mark=square*, mark size=1.3pt, domain=0:18, samples=19]
  {exp(-1.15*x)};
\addlegendentry{clustered spectrum (small $\kappa$)}
% spread spectrum: slow decay
\addplot[simRust, very thick, mark=*, mark size=1.3pt, domain=0:60, samples=31]
  {exp(-0.35*x)};
\addlegendentry{spread spectrum (large $\kappa$)}
% tolerance line
\addplot[simGray, densely dashed, domain=0:60, samples=2] {1e-8};
\addlegendentry{tolerance}
\end{axis}
\end{tikzpicture}
\caption[Convergence depends on the spectrum]{The relative residual versus
iteration count for two systems of the same size, on a log scale (a straight line
is geometric decay). When the eigenvalues of $\mat{A}$ are \emph{clustered}
(condition number small), a low-degree polynomial is small across the whole
spectrum and the residual plunges---tolerance reached in a handful of steps.
When the spectrum is \emph{spread} (condition number large), no low-degree
polynomial can be small everywhere at once and the residual crawls down, taking
many times as many steps to cross the same tolerance. Preconditioning
(\cref{ch:precond}) is precisely the art of replacing $\mat{A}$ with an
equivalent system whose spectrum is clustered---of moving a problem from the slow
curve onto the fast one.}
\label{fig:krylov-convergence}
\end{figure}

\Cref{fig:krylov-convergence} shows the two regimes: the same tolerance, reached in
a handful of steps on a clustered spectrum and only after a long crawl on a spread
one. This is why the bare Krylov method of this chapter is never the whole story.
The method itself is fixed; what we can change is the \emph{spectrum} of the system
it sees. \emph{Preconditioning} (\cref{ch:precond,ch:multigrid}) replaces
$\mat{A}\vx=\vb$ with an equivalent system $\mat{M}^{-1}\mat{A}\vx=\mat{M}^{-1}\vb$
whose operator $\mat{M}^{-1}\mat{A}$ has a clustered spectrum---moving the problem
from the slow curve onto the fast one. The Krylov method supplies the iteration;
the preconditioner supplies the conditioning; and together they reach close to
$O(N)$ work. That is the whole arc of Part~IV, and this chapter is its foundation.

\section{The shared step kernel}
\label{sec:krylov-kernel}

We turn now from \emph{why} Krylov methods work to \emph{what} they all do, step by
step. Here is the observation that organizes the rest of Part~IV: every Krylov
method---CG, GMRES, FGMRES, MINRES, BiCGStab---has an inner step of the \emph{same
shape}. Strip away the method-specific algebra and the per-step kernel is always
the same five-part sequence:
\begin{enumerate}[nosep]
  \item \textbf{one operator apply}: form $\mat{A}\vp$ (the matrix-free action,
  \cref{ch:matrixfree})---and, if preconditioned, one preconditioner apply
  $\mat{M}^{-1}\vr$;
  \item \textbf{a few inner products and norms}: $\inner{\cdot}{\cdot}$ and
  $\norm{\cdot}$ to form the step's scalar coefficients and measure the residual;
  \item \textbf{a few vector updates} (\emph{axpy}): $\vy\leftarrow\alpha\vx+\vy$,
  combining vectors with the freshly computed scalars;
  \item \textbf{an orthogonalization} (for the basis-building methods): make the new
  direction orthogonal to the previous ones (\cref{ch:orthog});
  \item \textbf{a state update}: fold the new iterate, residual, and scalars into
  the carry, and record the residual norm for the convergence test.
\end{enumerate}
We call this kernel \code{krylov\_step}. A particular method---CG, GMRES---is a
particular \emph{specialization} of it: a choice of which scalars, which updates,
whether the orthogonalization stage runs. But the skeleton is fixed, and one fact
about it is load-bearing.

\begin{keyidea}
Every Krylov method shares one step kernel, \code{krylov\_step}: one operator apply,
a few inner products and norms, a few \emph{axpy} updates, an optional
orthogonalization, and a state update. A specific method is a specialization---it
binds the scalars and the updates---but the shape is invariant. The \emph{cost
invariant} is the heart of it: \emph{exactly one operator apply per step}. Since the
apply is the one expensive operation (it touches all $N$ unknowns through the sparse
operator; everything else is cheap vector arithmetic), counting applies counts the
cost of the solve, and ``one apply per step'' is the budget every Krylov method
spends.
\end{keyidea}

\begin{figure}[t]
\centering
\begin{tikzpicture}[node distance=7mm and 10mm]
  % input carry
  \node[data, text width=20mm] (in) {carry $K$\\\footnotesize $\vx,\vr,\vp,\dots$};
  % stage 1: operator apply (the one expensive op)
  \node[stage, right=of in, text width=20mm, fill=simBgRust, draw=simRust] (apply)
    {\footnotesize \textbf{apply}\\$\vw=\mat{A}\vp$\\\scriptsize the one $O(N)$ op};
  % stage 2: inner products
  \node[opbox, right=of apply, text width=20mm] (dots)
    {\footnotesize \textbf{dots/norms}\\$\inner{\vp}{\vw},\ \norm{\vr}$};
  % stage 3: orthogonalization (optional)
  \node[extern, right=of dots, text width=20mm] (orth)
    {\footnotesize \textbf{orthog.}\\\scriptsize (GMRES only)};
  % stage 4: axpy updates
  \node[opbox, right=of orth, text width=20mm] (axpy)
    {\footnotesize \textbf{updates}\\\scriptsize \emph{axpy}: $\vx,\vr,\vp$};
  % output carry
  \node[data, right=of axpy, text width=20mm] (out) {carry $K'$\\\footnotesize $+$ residual};
  \draw[flow] (in) -- (apply);
  \draw[flow] (apply) -- node[above,font=\scriptsize]{$\vw$} (dots);
  \draw[flow] (dots) -- node[above,font=\scriptsize]{scalars} (orth);
  \draw[flow] (orth) -- (axpy);
  \draw[flow] (axpy) -- (out);
  % residual peels off downward for the convergence test
  \node[product, below=10mm of axpy, text width=24mm] (res)
    {\footnotesize residual norm\\\scriptsize $\to$ convergence test};
  \draw[refedge] (axpy.south) -- (res.north);
  % the cost-invariant brace under the apply
  \draw[decorate, decoration={brace, amplitude=4pt, mirror}, simRust]
    ([yshift=-2mm]apply.south west) -- ([yshift=-2mm]apply.south east)
    node[midway, below=4pt, font=\scriptsize, simRust, align=center]
    {one apply\\per step};
\end{tikzpicture}
\caption[The shared Krylov step kernel]{The kernel \code{krylov\_step} every Krylov
method shares, as a dataflow. A carry $K$ (the working vectors $\vx,\vr,\vp$ and
scalars) enters; the \textbf{one operator apply} $\vw=\mat{A}\vp$ (rust, the single
$O(N)$ operation) runs once; a few \textbf{inner products and norms} form the step's
scalars; the basis-building methods run an \textbf{orthogonalization}
(\cref{ch:orthog}; absent in CG, dashed); a few \emph{axpy} \textbf{updates} combine
vectors with the new scalars into the next carry $K'$; and the residual norm peels
off for the convergence test. The method (CG, GMRES, \dots) chooses the scalars and
which updates run; the \emph{shape}---and the one-apply-per-step cost---is fixed.}
\label{fig:krylov-kernel}
\end{figure}

\Cref{fig:krylov-kernel} draws the kernel as a dataflow. Note what is and is not
expensive. The operator apply touches all $N$ unknowns through the sparse operator;
it is the $O(N)$ operation, and it is the only one. The inner products and norms are
$O(N)$ reductions but with tiny constants; the \emph{axpy} updates are $O(N)$
elementwise; the orthogonalization, when present, costs a handful of inner products
against the existing basis. Everything but the apply is cheap vector arithmetic. So
the apply is the currency in which we price a solve, and the kernel spends exactly
one unit of it per step. When the apply is matrix-free (\cref{ch:matrixfree}), that
one apply is itself a chain of element-local contractions---but it is still one
apply, and the cost invariant holds.

\section{The solver as a fold, and as an operator}
\label{sec:krylov-fold}

We have the kernel; now we wrap a loop around it. The outer structure of a Krylov
solve is exactly the \code{iterate\_while} fold of \cref{ch:fold}: carry the
solver's working state forward, take one \code{krylov\_step} per pass, test a
convergence predicate, stop when the residual is small enough. There is nothing new
to invent---the solver is a textbook instance of the combinator the framework
already built.

Recall the kernel signature from the framework: the value-threaded Krylov step
maps a carry to the next carry plus a demand-prunable readout,
\begin{lstlisting}[style=l4style]
-- one step: thread the working bundle K, the persistent state s; report extras
krylov_step :: (op, K, s) -> (K', s', outputs)
\end{lstlisting}
where \code{op} is the construction-bound operator data (the system operator and
preconditioner, frozen in; \cref{ch:operators}), \code{K} is the iterate bundle
(the working vectors and scalars), \code{s} is the persistent solver state (the
iteration counter, the converged flag, the residual proxy), and \code{outputs} is
the per-step readout (typically the residual norm) that the trajectory collects.
Folding this kernel with \code{iterate\_while} is the entire outer loop:
\begin{lstlisting}[style=l4style]
-- the Krylov solve IS a fold of the step kernel (schematically)
solve op K0 s0 =
  iterate_while
    { K: K0, s: s0 }                              -- initial carry
    (\c -> not c.s.converged && c.s.it < max_it)  -- predicate: pure on the carry
    (\c -> let (K', s', outs) = krylov_step op c.K c.s
           in { state: { K: K', s: s' }, residual_norm: outs.residual_norm })
\end{lstlisting}
The predicate reads only the carry---the converged flag and iteration count it
carries---exactly as the predicate rule of \cref{sec:predicaterule} demands; the
residual the step computes is folded into the carry's \code{s} so the \emph{next}
pass's predicate can read it. The per-step residual norm rides out in the extras and
into the trajectory, where a convergence monitor reads it when asked---and, by
demand-driven pruning (\cref{sec:pruning}), is not computed at all when nothing
reads it. The solver is a fold, and it inherits every property the fold has.

\begin{figure}[t]
\centering
\begin{tikzpicture}[font=\footnotesize, node distance=6mm]
  % carry spine threading SimState + Krylov bundle
  \node[data, minimum width=16mm, align=center] (c0) at (0,0) {$K_0,s_0$};
  \node[diamond, draw=simBlue, fill=simBgBlue, aspect=1.7, inner sep=1pt,
        font=\scriptsize] (p0) at (2.1,0) {conv.?};
  \node[opbox, minimum width=18mm] (st0) at (4.3,0) {\code{krylov\_step}};
  \node[data, minimum width=16mm, align=center] (c1) at (6.6,0) {$K_1,s_1$};
  \node[diamond, draw=simBlue, fill=simBgBlue, aspect=1.7, inner sep=1pt,
        font=\scriptsize] (p1) at (8.7,0) {conv.?};
  \node[opbox, minimum width=18mm] (st1) at (10.9,0) {\code{krylov\_step}};
  \node[font=\scriptsize, simGray] (more) at (12.6,0) {$\cdots$};
  \draw[flow] (c0)--(p0);
  \draw[flow] (p0)-- node[above,font=\scriptsize]{no} (st0);
  \draw[flow] (st0)-- node[above,font=\scriptsize]{\code{.state}} (c1);
  \draw[flow] (c1)--(p1);
  \draw[flow] (p1)-- node[above,font=\scriptsize]{no} (st1);
  \draw[flow] (st1)-- (more);
  % converged exits down to the result
  \node[product, minimum width=22mm] (fin) at (8.7,-2.2)
    {\code{SolveResult}: $\vx$, converged};
  \draw[backflow] (p0.south) |- (fin.west);
  \draw[backflow] (p1.south) -- (fin.north);
  % residual trajectory peels downward
  \node[data, fill=simBgGold, draw=simGold, minimum width=40mm] (traj) at (7.7,-1.2)
    {residual history $[\,r_0,\,r_1,\,\dots\,]$ (pruned if unread)};
  \draw[refedge] (st0.south) -- (st0.south|-traj.north);
  \draw[refedge] (st1.south) -- (st1.south|-traj.north);
\end{tikzpicture}
\caption[A Krylov solve as an \texttt{iterate\_while} fold]{The outer loop of a
Krylov solve is the \code{iterate\_while} fold of \cref{ch:fold}. The carry
threads the solver's state---the iterate bundle $K$ (working vectors and scalars)
and the persistent state $s$ (iteration counter, converged flag, residual
proxy)---forward along the top. Before each pass the predicate tests convergence
on the \emph{current carry} (the \texttt{while} convention, \cref{sec:predicaterule});
on ``not converged'' a single \code{krylov\_step} runs, emitting the next carry. The
per-step residual peels downward into the trajectory, where a monitor reads it---or,
unread, it is pruned (\cref{sec:pruning}). The first ``converged'' exits the loop,
surfacing the iterate $\vx$ as the \code{SolveResult}. Compare \cref{fig:threaded}:
this is that figure with ``step'' the Krylov kernel and the carry the solver state.}
\label{fig:krylov-iterwhile}
\end{figure}

\Cref{fig:krylov-iterwhile} draws the fold for the Krylov solve specifically---the
threaded-carry picture of \cref{ch:fold}, with \code{krylov\_step} in the step box
and the residual history peeling off as the trajectory. The state-threaded picture
also explains why the carry splits in two (\cref{ch:strata}): the persistent
state $s$ (counter, flags, the externally visible iterate) is one stratum, and the
ephemeral iterate bundle $K$ (the workspace vectors, reborn at each restart of a
restarted method) is another. The fold threads both; their different lifetimes are
why GMRES's restart logic wraps an \emph{outer} loop around the inner fold rather
than flattening the two (a fold-merge the non-laws of \cref{sec:laws} forbid, and
which \cref{ch:gmres} respects).

\subsection{The whole solve is a constructed operator}
\label{sec:krylov-operator}

Step back from the loop and look at the solve as a single object. It takes a
right-hand side $\vb$ and returns an approximate solution $\vx\approx\mat{A}^{-1}\vb$.
That is the signature of an \emph{operator}: vector in, vector out. The entire
Krylov solve is a \emph{constructed operator} in the exact sense of
\cref{ch:operators}---built once with its system operator, preconditioner,
tolerances, and iteration budget baked into immutable internal state, then applied to
a right-hand side as a pure function. We met this as \code{ksp\_solve}; here is its
signature.

\begin{lstlisting}[style=l4style]
-- the solver is a constructed operator: an approximate inverse you apply
ksp_solve :: (K: Solver[A], b: Tensor[N]) -> SolveResult[N]

SolveResult[N] = {
  x          : Tensor[N],   -- approximate solution to A . x = b
  converged  : Bool,        -- did the convergence test pass?
  iterations : Int,         -- inner Krylov iterations consumed
  initial_res: Real,        -- initial residual norm
  final_res  : Real         -- final residual norm
}
\end{lstlisting}

The construction-bound value \code{K} carries the system operator \code{A}, an
optional preconditioner, and the convergence-control parameters; \emph{all} the
per-method choices---CG versus GMRES versus FGMRES, preconditioner side,
orthogonalization method---are absorbed into \code{K} at construction
(\cref{ch:operators}), so the per-call surface $(K,\vb)$ is variant-free. The solve
is the constructed-operator pattern of \cref{ch:operators} in its most consequential
instance: $\code{ksp\_solve}(K,\cdot)$ is an \emph{approximate inverse} you hold and
apply, indistinguishable in form from applying $\mat{A}$ itself. This is exactly the
\emph{solver-as-operator} move that lets one GMRES drive any preconditioner
(\cref{ch:operators}): the preconditioner is just another \code{apply}-shaped value
the kernel calls.

\begin{figure}[t]
\centering
\begin{tikzpicture}[node distance=10mm]
  % construction inputs (build-time)
  \node[data, align=left, text width=26mm] (cfg)
    {operator $\mat{A}$\\preconditioner\\tolerances,\\iteration cap};
  \node[stage, minimum width=24mm, minimum height=14mm, right=12mm of cfg] (con)
    {\textbf{construct}\\[1pt]\footnotesize freeze config,\\absorb method};
  \draw[flow] (cfg) -- node[above,font=\scriptsize]{once} (con);
  % the opaque solver-operator value
  \node[draw=simBlue, fill=simBgBlue, very thick, rounded corners=3pt,
        minimum width=26mm, minimum height=20mm, right=12mm of con] (op) {};
  \node[font=\ttfamily\scriptsize, simInk] at (op.north) [yshift=-3.4mm]
    {ksp\_solve};
  \node[font=\scriptsize\itshape, simGray, align=center] at (op.center)
    {$\approx\mat{A}^{-1}$\\(frozen:\\$\mat{A}$, pc,\\tols)};
  \draw[flow] (con) -- (op);
  % many applies out
  \foreach \y/\lbl in {1.0/a, 0.0/b, -1.0/c}{
    \node[draw=simTeal, fill=simBgGreen, semithick, rounded corners=2pt,
          minimum width=20mm, minimum height=6mm] (ap\lbl) at ($(op.east)+(2.9,\y)$)
          {\footnotesize\ttfamily ksp\_solve(K, b)};
    \draw[flow] (op.east) -- (ap\lbl.west);
  }
  \node[font=\scriptsize\itshape, simGray] at ($(op.east)+(2.9,-1.7)$)
    {many RHS, each a fold inside};
  % band labels
  \node[font=\scriptsize\bfseries, simRust] at ($(con)+(0,-1.5)$) {build-time};
  \node[font=\scriptsize\bfseries, simGreen] at ($(op.east)+(2.9,1.7)$) {run-time};
\end{tikzpicture}
\caption[The solve as a constructed operator]{The whole Krylov solve is a
constructed operator (\cref{ch:operators}). \emph{Construction} (left, build-time,
once) freezes the system operator $\mat{A}$, the preconditioner, and the
convergence-control parameters into an opaque value \code{K}, absorbing the choice
of method along the way. \emph{Application} (right, run-time, many times) feeds a
right-hand side $\vb$ and runs the \code{iterate\_while} fold of
\cref{fig:krylov-iterwhile} inside, returning $\vx\approx\mat{A}^{-1}\vb$. The
solve wears the same \code{apply} interface as the operator it inverts: it is an
\emph{approximate inverse you apply}, and the fixed-operator map of
\cref{ch:situating} (electrostatics solving against many right-hand sides) is
exactly this operator applied repeatedly with $\mat{A}$ hoisted out of the loop.}
\label{fig:krylov-ksp}
\end{figure}

\Cref{fig:krylov-ksp} draws it: build once, apply many. This is also the resolution
of the fixed-operator map from \cref{ch:situating}. Electrostatics solves
$\Kop\vx_i=\vb_i$ for many right-hand sides against \emph{one} operator; in this
language, it constructs \code{ksp\_solve} once with $\Kop$ frozen in and applies it to
each $\vb_i$, hoisting the expensive construction out of the loop. The constructed
operator \emph{is} the hoist.

\subsection{What the solve guarantees---and what it deliberately does not}
\label{sec:krylov-laws}

Because the solve is an approximate inverse, it obeys the laws an inverse obeys---in
the limit of a tight tolerance---and, just as importantly, it pointedly fails to
obey several laws an \emph{exact} inverse would. Knowing both is what keeps a caller
from trusting a guarantee the solve does not make. The laws below hold modulo the
convergence tolerance (equalities in the limit $\text{rel\_tol}\to0$,
$\text{max\_it}\to\infty$, treating $K$ as the formal map $\vb\mapsto\mat{A}^{-1}\vb$).

\begin{description}[style=nextline, leftmargin=1.4em]
  \item[Linearity in the right-hand side.] $\code{ksp\_solve}(K,\alpha\vb_1+\beta\vb_2).\vx
  \approx \alpha\,\code{ksp\_solve}(K,\vb_1).\vx + \beta\,\code{ksp\_solve}(K,\vb_2).\vx$.
  This is the linearity of $\mat{A}^{-1}$, exact in the formal limit and approximate
  at finite tolerance (the gap bounded by the condition number). The
  \emph{statistics} fields are not linear: different right-hand sides take different
  iteration counts and produce different residual histories.
  \item[Near-inverse on the right-hand side.] $\code{ksp\_solve}(K,\mat{A}\vx).\vx
  \approx \vx$: applying the solve to $\mat{A}\vx$ recovers $\vx$ to within the
  tolerance. This is the defining property---the solve is an approximate
  $\mat{A}^{-1}$.
  \item[Zero right-hand side gives zero.] $\code{ksp\_solve}(K,\mathbf{0}).\vx=\mathbf{0}$,
  exactly, not merely modulo tolerance (the iteration short-circuits on a zero
  initial residual).
\end{description}

Now the non-laws---the guarantees the solve refuses to make, every one of them
load-bearing in the sense of \cref{ch:bigidea}.

\begin{pitfall}
A Krylov solve is \emph{not} a bit-exact, deterministic function of $(K,\vb)$. Three
distinct sources of non-determinism are deliberate and must not be ``optimized
away.''
\begin{itemize}[nosep]
  \item \textbf{Reduction-order non-determinism.} The inner products and norms of the
  kernel are floating-point reductions, and floating-point addition is not
  associative. Reordering a reduction tree (as a parallel backend will) changes the
  bit-level scalars, hence the bit-level iterate count and the bit-level $\vx$. The
  mathematical solution is unchanged; its floating-point realization is not.
  \item \textbf{Orthogonalization-variant non-determinism.} For the basis-building
  methods, the orthogonalization choice (modified versus classical Gram--Schmidt,
  \cref{ch:orthog}) changes the per-step arithmetic and so the bit-level trajectory.
  Same answer, different bits.
  \item \textbf{Initial-guess non-determinism.} A warm start and a cold start take
  different paths through the Krylov subspaces and converge to bit-different $\vx$
  in different iteration counts. The mathematical solution is the same.
\end{itemize}
And the solve does \emph{not} compose exactly with the operator: $\mat{A}\,
\code{ksp\_solve}(K,\vb).\vx=\vb$ holds only \emph{approximately}, within the
tolerance, never exactly at finite precision. An algorithm that assumes the residual
is exactly zero after a solve---some iterative-refinement schemes do---must guard.
\emph{Convergence depends on conditioning, not on exactness}; the solve is a
controlled approximation, and its non-determinism is the honest price of running it
on a parallel, finite-precision machine.
\end{pitfall}

These non-laws are not defects to be apologized for; they are the truthful contract.
A Krylov solve is a controlled approximation whose accuracy you dial with a
tolerance and whose speed you buy with conditioning. Pretending it is an exact,
deterministic inverse would be convenient and false---and would silently corrupt any
caller that leaned on the pretense.

\section{Worked examples}
\label{sec:krylov-wex}

\begin{workedexample}{Building $\mathcal{K}_3(\mat{A},\vb)$ for a $3\times3$ SPD
system}{krylovsub}
We build a Krylov subspace explicitly and watch the solution appear inside it. Take
the small SPD matrix and right-hand side
\[
  \mat{A}=\begin{pmatrix}2&1&0\\1&2&1\\0&1&2\end{pmatrix},
  \qquad
  \vb=\begin{pmatrix}1\\0\\0\end{pmatrix}.
\]
$\mat{A}$ is symmetric, and its eigenvalues are $2-\sqrt2,\,2,\,2+\sqrt2$, all
positive, so it is SPD---the setting CG (\cref{ch:cg}) will live in.

\textbf{Build the Krylov sequence.} Each new vector is one operator apply on the
last. Starting from $\vb$:
\[
  \vb=\begin{pmatrix}1\\0\\0\end{pmatrix},\qquad
  \mat{A}\vb=\begin{pmatrix}2\\1\\0\end{pmatrix},\qquad
  \mat{A}^2\vb=\mat{A}(\mat{A}\vb)=\begin{pmatrix}5\\4\\1\end{pmatrix}.
\]
(Check the last: $\mat{A}(2,1,0)^{\mathsf T}=(2\!\cdot\!2+1,\ 2+2,\ 1)^{\mathsf
T}=(5,4,1)^{\mathsf T}$.) Two operator applies built the three spanning vectors. The
three Krylov subspaces are
\[
  \mathcal{K}_1=\operatorname{span}\{\vb\},\quad
  \mathcal{K}_2=\operatorname{span}\{\vb,\mat{A}\vb\},\quad
  \mathcal{K}_3=\operatorname{span}\{\vb,\mat{A}\vb,\mat{A}^2\vb\}.
\]

\tcblower
\textbf{Are the three vectors independent?} Stack them as columns and take the
determinant:
\[
  \det\!\begin{pmatrix}1&2&5\\0&1&4\\0&0&1\end{pmatrix}=1\neq0 .
\]
The matrix is upper-triangular with unit diagonal, so the three vectors are linearly
independent: $\mathcal{K}_3=\Real^3$, the whole space. The nesting
$\mathcal{K}_1\subsetneq\mathcal{K}_2\subsetneq\mathcal{K}_3$ is strict, one new
dimension per apply.

\textbf{The solution lives in $\mathcal{K}_3$.} Solve $\mat{A}\vx=\vb$ directly for the
target. With $\det\mat{A}=2(4-1)-1(2-0)=4$, the solution is
\[
  \vx^\star=\mat{A}^{-1}\vb=\tfrac14\begin{pmatrix}3\\-2\\1\end{pmatrix}
   =\begin{pmatrix}0.75\\-0.5\\0.25\end{pmatrix}.
\]
Since $\mathcal{K}_3=\Real^3$, $\vx^\star$ is certainly in $\mathcal{K}_3$---and we can
exhibit it as a polynomial in $\mat{A}$ applied to $\vb$, the form
\cref{sec:krylov-poly} promised. Seek $\vx^\star=c_0\vb+c_1\mat{A}\vb+c_2\mat{A}^2\vb$:
\[
  c_0\begin{pmatrix}1\\0\\0\end{pmatrix}
  +c_1\begin{pmatrix}2\\1\\0\end{pmatrix}
  +c_2\begin{pmatrix}5\\4\\1\end{pmatrix}
  =\begin{pmatrix}0.75\\-0.5\\0.25\end{pmatrix}.
\]
The bottom row gives $c_2=0.25$; the middle row $c_1+4c_2=-0.5$ gives $c_1=-1.5$; the
top row $c_0+2c_1+5c_2=0.75$ gives $c_0=0.75+3-1.25=2.5$. So
\[
  \vx^\star=\bigl(2.5\,I-1.5\,\mat{A}+0.25\,\mat{A}^2\bigr)\vb
  = q(\mat{A})\,\vb,\quad q(t)=0.25t^2-1.5t+2.5 .
\]
The exact inverse-times-$\vb$ is a degree-2 polynomial in $\mat{A}$ applied to
$\vb$---exactly the minimal-polynomial mechanism of \cref{sec:krylov-poly}, here with
$d=3$ (so degree $d-1=2$). A Krylov method on this system reaches the exact answer in
at most $3$ steps; CG, minimizing the error over each $\mathcal{K}_m$, would land on
$\vx^\star$ by step $3$ in exact arithmetic---and, because the spectrum is tightly
clustered around $2$, would be very close well before that. The subspace built from
nothing but operator applies contains the solution, and we found the coefficients
that pick it out.
\end{workedexample}

\begin{workedexample}{Costing one Krylov step against a direct factorization}{krylovcost}
We make the cost invariant concrete by pricing one step of a Krylov method and
comparing it to a sparse direct factorization, on a system the size a real solve
faces. Take a 3-D finite-element operator with $N=10^6$ unknowns and a bounded
$\nu=27$ nonzeros per row (a typical low-order stencil), so $\mat{A}$ has about
$\nu N=2.7\times10^7$ nonzeros.

\textbf{Cost of one Krylov step.} The kernel of \cref{sec:krylov-kernel} does, per
step:
\begin{itemize}[nosep]
  \item \emph{one operator apply} $\vw=\mat{A}\vp$. A sparse apply touches each
  nonzero once: $\approx\nu N=2.7\times10^7$ multiply--adds. This is the dominant
  cost.
  \item \emph{a couple of inner products / norms}, e.g.\ $\inner{\vp}{\vw}$ and
  $\norm{\vr}$: each $N$ multiply--adds, so $\approx 2N=2\times10^6$.
  \item \emph{a couple of \emph{axpy} updates}, e.g.\ $\vx\leftarrow\vx+\alpha\vp$ and
  $\vr\leftarrow\vr-\alpha\vw$: each $N$ multiply--adds, so $\approx 2N=2\times10^6$.
\end{itemize}
Total: about $\nu N+4N=(27+4)\times10^6\approx 3.1\times10^7$ operations, and the apply
is $\approx87\%$ of it. The step is $O(N)$, and it is one apply plus small change.

\tcblower
\textbf{Cost of a sparse direct factorization.} A sparse Cholesky on a 3-D mesh costs
about $O(N^2)$ arithmetic, with the fill-in of \cref{sec:krylov-direct} producing
factors of about $O(N^{4/3})$ nonzeros. Even taking a generous constant, the
factorization work at $N=10^6$ is on the order of $10^{12}$ operations or more---and
the factors need on the order of $N^{4/3}=10^8$ stored nonzeros, several times the
$2.7\times10^7$ of $\mat{A}$ itself, before they are even applied.

\textbf{The comparison.} Put the two side by side:
\[
  \underbrace{3.1\times10^7}_{\text{one Krylov step}}
  \qquad\text{versus}\qquad
  \underbrace{\sim 10^{12}}_{\text{one direct factorization}} .
\]
The factorization costs as much as roughly $10^{12}/(3.1\times10^7)\approx3\times10^4$
Krylov steps. So a preconditioned Krylov solve that converges in, say, a few hundred
steps wins by orders of magnitude---\emph{and} never forms a matrix, since each of
those steps needs only the matrix-free apply (\cref{ch:matrixfree}), while the
factorization needs the entries and pays the fill-in to store its denser factors.

\textbf{The reading.} The whole economics of large-scale solving is in this
comparison. A direct method pays a huge one-time $O(N^2)$ cost and then solves
cheaply; an iterative method pays a small $O(N)$ cost per step and bets that a
bounded number of steps suffice. For a million-unknown 3-D system, the bet pays:
the per-step cost is dominated by one sparse apply, the apply can be matrix-free,
and a good preconditioner (\cref{ch:precond}) keeps the step count bounded. ``One
apply per step,'' costed out, is why Part~IV exists.
\end{workedexample}

\section{The map of Part~IV}
\label{sec:krylov-map}

This chapter laid the shared foundation; the rest of Part~IV fills in the
specializations and the machinery that makes them fast. Here is the road ahead, each
piece a refinement of something built in this chapter.

\begin{description}[style=nextline, leftmargin=1.4em]
  \item[Conjugate gradients (\cref{ch:cg}).] The Krylov method for SPD systems---
  electrostatics, magnetostatics, implicit time steps. It specializes
  \code{krylov\_step} with a short two-term recurrence (\cref{ch:fold}'s
  \code{iterate\_while\_with\_prev}) and no explicit orthogonalization, minimizing
  the error in the $\mat{A}$-norm. One apply per step.
  \item[GMRES and FGMRES (\cref{ch:gmres}).] The Krylov method for general
  (nonsymmetric, complex) systems---the driven analysis. It builds an explicit
  orthonormal Krylov basis and minimizes the residual norm over it, which forces the
  orthogonalization stage and an outer restart loop. FGMRES threads a second basis to
  carry a per-step (flexible) preconditioner, exactly the per-step variant that
  \cref{ch:operators}'s decision rule routes into the state.
  \item[Orthogonalization (\cref{ch:orthog}).] The stage GMRES needs: making each new
  basis vector orthogonal to the previous ones, and the modified-versus-classical
  Gram--Schmidt choice that is the orthogonalization-variant non-law of
  \cref{sec:krylov-laws}.
  \item[The least-squares core (\cref{ch:givens}).] How GMRES actually picks the
  minimizing iterate inside its basis: a small dense least-squares problem solved
  incrementally with Givens rotations as the basis grows.
  \item[Preconditioning and multigrid (\cref{ch:precond,ch:multigrid}).] The
  conditioning half of the story---replacing $\mat{A}$ with an equivalent
  better-conditioned system so the spectrum clusters (\cref{fig:krylov-convergence})
  and the step count drops. The preconditioner is itself a constructed operator
  (\cref{ch:operators}) the kernel applies; multigrid is the most powerful one, an
  approximate inverse built from a hierarchy of meshes.
\end{description}

Every one of these is a setting of the kernel and the fold built here, or a piece of
machinery that feeds them. Learn the Krylov subspace, the one-apply-per-step kernel,
and the solve-as-fold-as-operator reading, and the rest of Part~IV is
specialization.

\summarysection
Discretization (Part~II) produces a large \emph{sparse} linear system
$\mat{A}\vx=\vb$: $\mat{A}$ has $N$ up to the tens of millions but only $O(N)$
nonzeros, because the finite-element basis has local support and locality is
sparsity. \emph{Direct} methods (Gaussian elimination, $LU$/Cholesky) are too
expensive at scale---fill-in makes the factors far denser than $\mat{A}$, costing
about $O(N^2)$ in 3-D and $O(N^3)$ dense---and they need the matrix entries, which a
matrix-free operator does not provide. \emph{Iterative} methods win because they need
only the operator's \emph{action} $\vv\mapsto\mat{A}\vv$, the exact interface
\cref{ch:matrixfree} supplies, and cost close to $O(N)$ when well preconditioned. A
\emph{Krylov subspace} $\mathcal{K}_m(\mat{A},\vb)=\operatorname{span}\{\vb,\mat{A}\vb,
\dots,\mat{A}^{m-1}\vb\}$ is built from nothing but repeated applies; its nested
ladder contains better and better approximations, and the minimal polynomial puts the
exact solution in $\mathcal{K}_d$ for some $d\le N$ (so Krylov terminates in $\le N$
steps in exact arithmetic). Convergence speed is governed by the spectrum: a
\emph{clustered} spectrum converges fast, a \emph{spread} one slowly---which is why
preconditioning exists. Every Krylov method shares one \emph{step kernel}
\code{krylov\_step}: one operator apply, a few inner products and norms, a few
\emph{axpy} updates, an optional orthogonalization, a state update---with the cost
invariant of \emph{exactly one apply per step}. The outer loop is the
\code{iterate\_while} fold of \cref{ch:fold}, and the whole solve is a
\emph{constructed operator} (\cref{ch:operators}), \code{ksp\_solve}: an approximate
inverse you build once and apply, obeying linearity and the near-inverse law modulo
tolerance, and deliberately \emph{not} obeying bit-determinism across reduction
order, orthogonalization, or initial guess---a controlled approximation whose
accuracy is a tolerance and whose speed is conditioning. Part~IV specializes this
foundation: CG (\cref{ch:cg}), GMRES/FGMRES (\cref{ch:gmres}), orthogonalization
(\cref{ch:orthog}), the least-squares core (\cref{ch:givens}), and preconditioning
and multigrid (\cref{ch:precond,ch:multigrid}).

\furtherreading
The standard reference for iterative methods on sparse systems is
Saad~\citep{saad2003}; its early chapters develop the Krylov subspace, the
polynomial view, and the spectrum-driven convergence bounds at exactly the level of
this chapter, and the later chapters are the canonical source for the CG and GMRES
specializations of \cref{ch:cg,ch:gmres} and for preconditioning
(\cref{ch:precond}). For the numerical-linear-algebra foundations---why fill-in
makes direct methods scale poorly, the conditioning and floating-point subtleties
behind the non-laws of \cref{sec:krylov-laws}---Golub and Van
Loan~\citep{golubvanloan2013} is the comprehensive reference. Trefethen and
Bau~\citep{trefethenbau1997} give the most readable account of the
polynomial-approximation picture of Krylov convergence (their lectures on the
Arnoldi and CG processes are especially clear on why a clustered spectrum converges
fast), and are an ideal companion to this chapter's
\cref{sec:krylov-poly,sec:krylov-convergence}.

\exercisessection
\begin{enumerate}[nosep]
  \item \textbf{Sparsity counts.} A 3-D finite-element operator has $N=4\times10^6$
  unknowns and an average of $\nu=30$ nonzeros per row. (a) How many numbers does the
  sparse operator store? (b) How many would the dense $N\times N$ array hold? (c) By
  what factor does sparsity reduce the storage? (d) In one sentence, tie the bounded
  $\nu$ back to the local support of the basis functions (\cref{ch:refelement}).
  \item \textbf{Why not factor.} Using \cref{fig:krylov-cost}, estimate the ratio of
  work between a dense direct solve ($O(N^3)$) and a near-linear iterative solve
  ($\approx O(N)$) at $N=10^5$ and again at $N=10^7$. By what factor does the
  iterative advantage grow as $N$ goes up by $100\times$? Explain in terms of the
  exponents why the gap widens.
  \item \textbf{Build a Krylov subspace.} For
  $\mat{A}=\bigl(\begin{smallmatrix}3&1\\1&3\end{smallmatrix}\bigr)$ and
  $\vb=(1,1)^{\mathsf T}$: (a) compute $\mat{A}\vb$; (b) is $\mathcal{K}_2(\mat{A},\vb)
  =\operatorname{span}\{\vb,\mat{A}\vb\}$ all of $\Real^2$, or does it collapse?
  (c) Solve $\mat{A}\vx=\vb$ directly and find scalars $c_0,c_1$ with
  $\vx=c_0\vb+c_1\mat{A}\vb$, exhibiting $\vx=q(\mat{A})\vb$. (d) Why does the collapse
  in (b) mean a Krylov method finishes in \emph{one} step here? (Hint: $\vb$ is an
  eigenvector.)
  \item \textbf{The polynomial view.} \Cref{sec:krylov-poly} shows
  $\mat{A}^{-1}=q(\mat{A})$ for a polynomial $q$ of degree $d-1$, where $d$ is the
  degree of the minimal polynomial. If $\mat{A}$ has only $3$ distinct eigenvalues
  (each possibly repeated), what is the most steps a Krylov method that finds the best
  iterate in each $\mathcal{K}_m$ can need, in exact arithmetic, regardless of $N$?
  Why does the \emph{number of distinct eigenvalues}, not $N$, set the bound?
  \item \textbf{Conditioning and speed.} Two SPD systems have condition numbers
  $\kappa_1=100$ and $\kappa_2=10^4$. Using the CG rule of thumb that the iteration
  count scales like $\sqrt{\kappa}$, estimate the ratio of step counts to reach a
  fixed tolerance. If a preconditioner clusters the spectrum of the second system so
  its effective condition number drops to $100$, what happens to its step count?
  Relate your answer to \cref{fig:krylov-convergence}.
  \item \textbf{One apply per step.} Reread \cref{wex:krylovcost}. (a) For
  $N=10^6$ and $\nu=27$, what fraction of one Krylov step's arithmetic is the operator
  apply? (b) Suppose a colleague proposes computing \emph{two} extra inner products per
  step for a fancier convergence test. By what percentage does that raise the per-step
  cost, and why is it nonetheless usually a bad trade if it does not reduce the step
  count? (c) Why does the cost invariant ``one apply per step'' survive unchanged when
  the apply is matrix-free (\cref{ch:matrixfree})?
  \item \textbf{The solve is a fold.} Write the convergence predicate of a Krylov
  solve that stops when either the relative residual drops below $10^{-8}$ or the
  iteration count reaches $500$, as a pure function of the carry (cf.\
  \cref{sec:krylov-fold} and the predicate rule of \cref{sec:predicaterule}). Which
  fields must the carry contain for this predicate to be expressible, and why can the
  predicate \emph{not} read the residual the current step just produced as an extra?
  \item \textbf{The solve is an operator.} The driven analysis (\cref{ch:situating})
  rebuilds its operator $\mat{A}(\omega)$ at every frequency, while electrostatics
  reuses one operator across many right-hand sides. Using
  \cref{fig:krylov-ksp} and the constructed-operator pattern of \cref{ch:operators},
  explain which analysis can hoist the \code{ksp\_solve} construction out of its loop
  and which cannot, and why the difference is one of \emph{kind}, not degree.
  \item \textbf{The non-laws.} A caller writes an iterative-refinement loop that
  assumes $\mat{A}\,\code{ksp\_solve}(K,\vb).\vx=\vb$ \emph{exactly} after a solve.
  (a) Which non-law of \cref{sec:krylov-laws} does this violate? (b) What goes wrong,
  concretely, if the assumption is false by a quantity near the tolerance? (c) Name
  the three bit-level non-determinism sources a Krylov solve carries, and for each,
  state whether it changes the \emph{mathematical} solution or only its
  floating-point realization.
\end{enumerate}

\chapter{Conjugate Gradients}
\label{ch:cg}

\chapterorient{This chapter builds the single most important solver in the book.
Conjugate gradients is the engine of every electrostatic and magnetostatic
analysis and of the inner Poisson solves buried inside the larger pipelines---all
the places where the operator is \emph{symmetric and positive definite}. We
develop it from the ground a sophomore already stands on: solving
$\mat{A}\vx=\vb$ is rolling a ball down a quadratic bowl. Steepest descent rolls
badly, zig-zagging across a stretched valley; conjugate gradients fixes the
zig-zag with one idea---search directions that do not interfere---and as a bonus
reaches the exact answer in at most $N$ steps. We derive every coefficient,
write the method as one \texttt{iterate\_while} of the shared step kernel from
\cref{ch:krylov}, bound its convergence by the condition number, and end at the
doorstep of preconditioning. Three worked examples carry numbers all the way
through.}

\section{The symmetric positive-definite setting}
\label{sec:spd}

The methods of \cref{ch:krylov} apply to any sparse system. Conjugate gradients
is the specialist: it is the method to reach for when, and only when, the
operator has a particular structure. We state that structure first, because
everything that makes the method fast---and everything that makes it
\emph{wrong} on the wrong system---follows from it.

\begin{definition}[Symmetric positive-definite operator]
\label{def:spd}
A real linear operator $\mat{A}$ on $\Real^N$ is \emph{symmetric positive
definite} (SPD) if it is symmetric, $\mat{A}=\mat{A}^{T}$, and positive
definite, $\vx^{T}\mat{A}\vx>0$ for every nonzero $\vx\in\Real^N$. We write
$\mat{A}=\mat{A}^{T}\succ0$.
\end{definition}

Symmetry says the operator is its own transpose; positive definiteness says the
quadratic form $\vx^{T}\mat{A}\vx$ is strictly positive away from the origin.
Together they make $\mat{A}$ behave, for the purposes of geometry, like a
``stretched and rotated'' version of the identity: it has $N$ real positive
eigenvalues and an orthogonal basis of eigenvectors, so it stretches space along
$N$ perpendicular axes by $N$ positive factors and never flips or shears it.

This is not an exotic case. It is the case the field theory of
\cref{part:fieldtheory} hands us again and again. The electrostatic operator
$-\Div(\varepsilon\Grad\,\cdot\,)$, discretized on an $\Hone$ finite-element
space, is symmetric because the weak form $\int\varepsilon\,\Grad u\cdot\Grad v$
is symmetric in $u$ and $v$, and positive definite because
$\int\varepsilon\,\abs{\Grad u}^2>0$ for any non-constant $u$ when
$\varepsilon>0$. The magnetostatic curl--curl operator is SPD on its physical
subspace for the same reason. Whenever the physics is a minimization---store the
least energy subject to the boundary data---the discretized operator is SPD, and
conjugate gradients is the right tool.

\begin{physicsbox}
The SPD systems of this book are the \emph{energy-minimizing} ones. An
electrostatic field arranges itself to minimize the stored electric energy
$\tfrac12\int\varepsilon\abs{\vE}^2$ subject to the conductor voltages; a
magnetostatic field minimizes $\tfrac12\int\nu\abs{\vB}^2$ subject to the
currents. Discretize either and the energy becomes a quadratic
$\phi(\vx)=\tfrac12\vx^{T}\mat{A}\vx-\vb^{T}\vx$ with $\mat{A}\succ0$; the
physical equilibrium is its minimizer, and the minimizer solves
$\mat{A}\vx=\vb$. Conjugate gradients does not just solve a linear system---it
descends a literal energy landscape to the physical ground state.
\end{physicsbox}

\subsection{Solving is minimizing}

The bridge from ``solve $\mat{A}\vx=\vb$'' to ``roll down a bowl'' is a single
quadratic function. Define
\begin{equation}
\label{eq:phi}
  \phi(\vx) \;=\; \tfrac12\,\vx^{T}\mat{A}\vx \;-\; \vb^{T}\vx .
\end{equation}
Its gradient is $\Grad\phi(\vx)=\mat{A}\vx-\vb$, using $\mat{A}=\mat{A}^{T}$ to
fold the two halves of the symmetric form together. Setting the gradient to zero
gives $\mat{A}\vx=\vb$. So the solution of the linear system is exactly the
\emph{stationary point} of $\phi$. And because $\mat{A}\succ0$, the function
$\phi$ is strictly convex---its graph is an upward bowl, an $N$-dimensional
paraboloid with a single lowest point---so that stationary point is the unique
\emph{global minimum}.

\begin{proposition}[Linear solve as quadratic minimization]
\label{prop:minimize}
Let $\mat{A}=\mat{A}^{T}\succ0$. Then $\vx^\star$ minimizes
$\phi(\vx)=\tfrac12\vx^{T}\mat{A}\vx-\vb^{T}\vx$ if and only if
$\mat{A}\vx^\star=\vb$. The minimizer is unique.
\end{proposition}

\begin{proof}
For any $\vx$ and direction $\mathbf{d}$, expand
$\phi(\vx+\mathbf{d})=\phi(\vx)+\mathbf{d}^{T}(\mat{A}\vx-\vb)+\tfrac12\mathbf{d}^{T}\mat{A}\mathbf{d}$. At
$\vx=\vx^\star$ with $\mat{A}\vx^\star=\vb$, the linear term vanishes and
$\phi(\vx^\star+\mathbf{d})=\phi(\vx^\star)+\tfrac12\mathbf{d}^{T}\mat{A}\mathbf{d}>\phi(\vx^\star)$
for every $\mathbf{d}\neq0$ (positive definiteness). So $\vx^\star$ is the strict global
minimizer. Conversely, a minimizer has zero gradient $\mat{A}\vx-\vb=0$.
Uniqueness is the strictness of the inequality.
\end{proof}

The quantity $\vr=\vb-\mat{A}\vx$ that appears throughout---the
\emph{residual}---is, up to sign, the negative gradient:
$\vr=-\Grad\phi(\vx)$. The residual points in the direction of steepest
\emph{descent} of $\phi$. It measures how far $\mat{A}\vx$ falls short of $\vb$,
and it is zero exactly at the solution. \Cref{fig:bowl} draws the bowl, its
elliptical level sets, and the residual as the downhill arrow.

\begin{figure}[t]
\centering
\begin{tikzpicture}[font=\footnotesize]
  % --- left: the 3D bowl ---
  \begin{scope}[xshift=0cm]
    \begin{axis}[
        width=6.6cm, height=6.0cm, view={35}{32},
        domain=-1.6:1.6, y domain=-1.1:1.1, samples=24, samples y=18,
        colormap={bluegold}{color=(simBgBlue) color=(simGold)},
        xtick=\empty, ytick=\empty, ztick=\empty,
        xlabel={}, ylabel={}, axis lines=none, clip=false]
      \addplot3[surf, shader=interp, opacity=0.9, mesh/ordering=y varies]
        {0.9*x*x + 2.6*y*y};
      \node[font=\scriptsize, text=simInk] at (axis cs:0,0,0) {};
    \end{axis}
    \node[font=\scriptsize, text=simBlue] at (1.8,-1.9) {$\phi(\vx)=\tfrac12\vx^{T}\mat{A}\vx-\vb^{T}\vx$};
  \end{scope}
  % --- right: contours + residual arrow ---
  \begin{scope}[xshift=8.2cm]
    \begin{axis}[
        width=6.4cm, height=6.0cm, axis equal image,
        xtick=\empty, ytick=\empty, xlabel={}, ylabel={},
        axis lines=middle, axis line style={simGray, -{Stealth[length=1.6mm]}},
        xmin=-2.3, xmax=2.3, ymin=-1.7, ymax=1.7, clip=false]
      \addplot[simBlue, thick, samples=80, domain=0:360, variable=\t]
        ({sqrt(2.2)*cos(\t)}, {sqrt(0.8)*sin(\t)});
      \addplot[simBlue!70, samples=80, domain=0:360, variable=\t]
        ({sqrt(1.1)*cos(\t)}, {sqrt(0.4)*sin(\t)});
      \addplot[simBlue!45, samples=80, domain=0:360, variable=\t]
        ({sqrt(0.35)*cos(\t)}, {sqrt(0.13)*sin(\t)});
      \node[circle, fill=simRust, inner sep=1.3pt, label={[font=\scriptsize,text=simRust]below:$\vx^\star$}] at (axis cs:0,0) {};
      \coordinate (p) at (axis cs:1.35,0.78);
      \node[circle, fill=simInk, inner sep=1.1pt] at (p) {};
      \node[font=\scriptsize] at (axis cs:1.62,1.02) {$\vx$};
      \draw[flow, simRust] (p) -- (axis cs:0.55,0.18)
        node[midway, above right, font=\scriptsize, text=simRust] {$\vr=-\Grad\phi$};
    \end{axis}
  \end{scope}
\end{tikzpicture}
\caption[The quadratic bowl and its contours]{The function
$\phi(\vx)=\tfrac12\vx^{T}\mat{A}\vx-\vb^{T}\vx$ for an SPD $\mat{A}$. \emph{Left:}
the graph is an upward paraboloid---a bowl---with a single lowest point.
\emph{Right:} its level sets are concentric ellipses centered on the minimizer
$\vx^\star$, the solution of $\mat{A}\vx=\vb$. The ellipses are stretched along
$\mat{A}$'s eigenvectors by the reciprocals of its eigenvalues; the more
$\mat{A}$ stretches, the more eccentric they are. The residual $\vr=\vb-\mat{A}\vx
= -\Grad\phi$ points in the direction of steepest descent---straight down the
local slope, \emph{not}, in general, toward $\vx^\star$.}
\label{fig:bowl}
\end{figure}

\subsection{Why steepest descent zig-zags}
\label{sec:zigzag}

The most obvious way to descend a bowl is to repeatedly step in the steepest
downhill direction. That is \emph{steepest descent}: from the current $\vx$,
compute the residual $\vr=\vb-\mat{A}\vx$ (the downhill direction), and move
along it by the amount that minimizes $\phi$ on that line. The line-minimizing
step length has a clean closed form. Along $\vx+\alpha\vr$,
\[
  \phi(\vx+\alpha\vr) = \phi(\vx) - \alpha\,\vr^{T}\vr
    + \tfrac12\alpha^2\,\vr^{T}\mat{A}\vr,
\]
a parabola in $\alpha$ minimized at
\begin{equation}
\label{eq:sdstep}
  \alpha = \frac{\vr^{T}\vr}{\vr^{T}\mat{A}\vr}
         = \frac{\inner{\vr}{\vr}}{\inner{\mat{A}\vr}{\vr}} .
\end{equation}
Take the step, recompute the residual, repeat. It converges---$\phi$ decreases
every step and the bowl has a floor---but on a stretched bowl it converges
\emph{slowly}, and the reason is geometric and worth seeing.

When you minimize $\phi$ exactly along a line, you stop at the point where the
line is tangent to a level ellipse; there the new residual is
\emph{perpendicular to the line you just walked}. So each steepest-descent step
turns ninety degrees from the last. On a round bowl (eigenvalues all equal) the
first residual already points at the center and one step finishes. On a stretched
bowl (eigenvalues far apart) the downhill direction points mostly \emph{across}
the long valley rather than \emph{along} it, so the iterate crabs back and forth,
each step undoing part of the last, creeping toward the minimum in a long
zig-zag. \Cref{fig:zigzag} (left) is the classic picture. The narrower the
valley---the larger the ratio of largest to smallest eigenvalue---the worse the
crabbing. That ratio is the \emph{condition number}, and it will govern
everything.

\begin{keyidea}
Steepest descent stops, each step, where the new residual is perpendicular to the
direction just travelled. On a stretched (ill-conditioned) bowl, consecutive
perpendicular steps point mostly across the valley, so the iterate zig-zags and
crawls. The cure is not a better step \emph{length}---\cref{eq:sdstep} is already
optimal along its line---but a better set of \emph{directions}: directions chosen
so that progress along one is never spoiled by the next. Building those
directions is conjugate gradients.
\end{keyidea}

\section{Conjugacy: directions that do not interfere}
\label{sec:conjugacy}

The flaw in steepest descent is that minimizing along the new residual
\emph{re-introduces} error along the old direction: the steps are coupled. We
want a set of directions $\vp_0,\vp_1,\dots,\vp_{N-1}$ along which we can minimize
\emph{one at a time}, with the guarantee that optimizing along $\vp_k$ never
disturbs the optimization already done along $\vp_0,\dots,\vp_{k-1}$. The
condition that makes the directions independent in exactly this sense is
\emph{conjugacy}.

\begin{definition}[$\mat{A}$-conjugate directions]
\label{def:conjugate}
Two nonzero vectors $\vp_i,\vp_j$ are \emph{conjugate with respect to $\mat{A}$}
(or \emph{$\mat{A}$-orthogonal}) if
\[
  \inner{\mat{A}\vp_i}{\vp_j} = \vp_i^{T}\mat{A}\vp_j = 0,\qquad i\neq j .
\]
A set $\{\vp_k\}$ is $\mat{A}$-conjugate if its members are pairwise conjugate.
\end{definition}

Conjugacy is ordinary orthogonality measured in the \emph{energy inner product}
$\inner{\vu}{\vv}_{\mat{A}}:=\vu^{T}\mat{A}\vv$, which is a genuine inner product
precisely because $\mat{A}\succ0$. Its induced norm,
$\norm{\vu}_{\mat{A}}:=\sqrt{\vu^{T}\mat{A}\vu}$, is the \emph{energy norm}, and
it is the natural ruler for this problem: the quantity conjugate gradients
actually drives down is the energy-norm error, as \cref{sec:converge} will show.
Geometrically, $\mat{A}$-conjugate directions look skew in ordinary axes but are
perpendicular once the bowl is un-stretched into a round one. \Cref{fig:conjugate}
shows the two pictures side by side.

\begin{figure}[t]
\centering
\begin{tikzpicture}[font=\footnotesize]
  % --- left: skew in standard coords ---
  \begin{scope}
    \node[font=\small\bfseries, text=simBlue] at (0,2.55) {standard axes: skew};
    \begin{scope}[scale=0.95]
      \draw[simBlue, thick] (0,0) ellipse (1.9 and 0.95);
      \draw[simBlue!55] (0,0) ellipse (1.25 and 0.62);
      \node[circle, fill=simRust, inner sep=1.2pt] at (0,0) {};
      % two A-conjugate directions: skew here
      \draw[flow, simGreen, thick] (0,0) -- (1.7,0)
        node[right, font=\scriptsize, text=simGreen] {$\vp_0$};
      \draw[flow, simRust, thick] (0,0) -- (-0.62,0.78)
        node[above left, font=\scriptsize, text=simRust] {$\vp_1$};
      \node[font=\scriptsize, text=simGray] at (0.0,-1.45) {$\vp_0^{T}\mat{A}\vp_1=0$};
    \end{scope}
  \end{scope}
  \draw[simGray!50, thick] (3.7,-1.7) -- (3.7,2.7);
  % --- right: orthogonal in A-coords ---
  \begin{scope}[xshift=7.6cm]
    \node[font=\small\bfseries, text=simGreen] at (0,2.55) {energy ($\mat{A}$) axes: orthogonal};
    \begin{scope}[scale=0.95]
      \draw[simGreen, thick] (0,0) circle (1.25);
      \draw[simGreen!55] (0,0) circle (0.78);
      \node[circle, fill=simRust, inner sep=1.2pt] at (0,0) {};
      \draw[flow, simGreen, thick] (0,0) -- (1.25,0)
        node[right, font=\scriptsize, text=simGreen] {$\vp_0$};
      \draw[flow, simRust, thick] (0,0) -- (0,1.25)
        node[above, font=\scriptsize, text=simRust] {$\vp_1$};
      \draw[simGray] (0.22,0) -- (0.22,0.22) -- (0,0.22);
      \node[font=\scriptsize, text=simGray] at (0.0,-1.55) {perpendicular here};
    \end{scope}
  \end{scope}
\end{tikzpicture}
\caption[Conjugate directions in two coordinate systems]{$\mat{A}$-conjugate
directions, drawn twice. \emph{Left:} in ordinary coordinates the level sets are
ellipses and the two directions $\vp_0,\vp_1$ look \emph{skew}---they are not
perpendicular in the usual sense. \emph{Right:} change to coordinates that
un-stretch the bowl into a circle (the energy inner product
$\inner{\vu}{\vv}_{\mat{A}}=\vu^{T}\mat{A}\vv$), and the same two directions are
exactly \emph{perpendicular}. Conjugacy is orthogonality in the energy geometry.
This is why minimizing along $\vp_1$ cannot spoil the work done along $\vp_0$:
they are independent axes of the round bowl.}
\label{fig:conjugate}
\end{figure}

\subsection{Why conjugate directions finish in \texorpdfstring{$N$}{N} steps}
\label{sec:finite}

The payoff of conjugacy is a theorem that has no analogue for steepest descent:
with conjugate directions, exact minimization terminates after at most $N$ steps.
The idea is exactly the independence we asked for.

\begin{theorem}[Conjugate-direction termination]
\label{thm:finite}
Let $\{\vp_0,\dots,\vp_{N-1}\}$ be $\mat{A}$-conjugate (hence a basis of
$\Real^N$). Starting from any $\vx_0$, set
$\vx_{k+1}=\vx_k+\alpha_k\vp_k$ with $\alpha_k$ chosen to minimize $\phi$ along
$\vp_k$. Then $\vx_N=\vx^\star$, the exact solution of $\mat{A}\vx=\vb$.
\end{theorem}

\begin{proof}[Proof sketch]
Because the $\vp_k$ are a basis, write the initial error in them:
$\vx^\star-\vx_0=\sum_{k}\delta_k\vp_k$. Apply $\mat{A}$ and take the energy inner
product with $\vp_j$; conjugacy kills every cross term, leaving
$\delta_j=\inner{\vp_j}{\mat{A}(\vx^\star-\vx_0)}/\inner{\mat{A}\vp_j}{\vp_j}
=\inner{\vp_j}{\vr_0}/\inner{\mat{A}\vp_j}{\vp_j}$, where $\vr_0=\vb-\mat{A}\vx_0$.
Separately, the line-minimizing step length along $\vp_j$ works out (next
subsection) to exactly this same $\delta_j$. So the $j$-th step installs exactly
the $j$-th coordinate of the error and never touches the others---each conjugate
direction is corrected once, independently, and after $N$ corrections every
coordinate is right.
\end{proof}

This is the finite-termination property: \emph{in exact arithmetic, conjugate
gradients solves an $N\times N$ SPD system in at most $N$ steps}, because it lays
the $N$ independent corrections down one at a time with no interference. In
floating point the guarantee is softened (rounding spoils exact conjugacy), and
in practice we stop far earlier---after a residual tolerance is met, typically in
\emph{far fewer} than $N$ steps for a well-conditioned system, as
\cref{sec:converge} quantifies. But the finite-termination theorem is the
structural fact that distinguishes conjugate directions from steepest descent's
endless crawl.

\begin{keyidea}
Conjugate directions are \emph{$\mat{A}$-orthogonal}: $\vp_i^{T}\mat{A}\vp_j=0$
for $i\neq j$. Minimizing $\phi$ along each in turn corrects one independent
coordinate of the error and disturbs none of the others, so an $N$-dimensional
SPD system is solved exactly in at most $N$ such steps. This finite termination
is what steepest descent, with its coupled perpendicular steps, can never
promise.
\end{keyidea}

\section{Deriving the method}
\label{sec:derive}

We now have the goal---march along $\mat{A}$-conjugate directions---but not yet
the directions. A naive plan would precompute a full conjugate basis (a
Gram--Schmidt orthogonalization in the energy inner product, keeping every
previous direction). That costs $O(N^2)$ storage and is hopeless for the
million-unknown systems Palace solves. The miracle of conjugate gradients,
derived in this section, is that for an SPD operator the conjugate directions can
be built by a \emph{short recurrence}: each new direction needs only the current
residual and the \emph{immediately previous} direction. Nothing older is kept.

We build the method from three demands and read off each coefficient as the
unique value that satisfies them.

\subsection{The step length \texorpdfstring{$\alpha$}{alpha}}
\label{sec:alpha}

Suppose we have a current iterate $\vx_k$, its residual $\vr_k=\vb-\mat{A}\vx_k$,
and a search direction $\vp_k$. We move to $\vx_{k+1}=\vx_k+\alpha_k\vp_k$ and ask
for the $\alpha_k$ that minimizes $\phi$ along $\vp_k$---the line search of
\cref{sec:zigzag}, now along $\vp_k$ instead of along the residual. The same
parabola calculation gives
\begin{equation}
\label{eq:alpha}
  \alpha_k = \frac{\inner{\vr_k}{\vp_k}}{\inner{\mat{A}\vp_k}{\vp_k}} .
\end{equation}
Equivalently---and this is the form the algorithm uses---we will show below that
$\inner{\vr_k}{\vp_k}=\inner{\vr_k}{\vr_k}$, so
\begin{equation}
\label{eq:alpha2}
  \boxed{\;\alpha_k = \frac{\inner{\vr_k}{\vr_k}}{\inner{\mat{A}\vp_k}{\vp_k}}\;}
\end{equation}
The denominator $\inner{\mat{A}\vp_k}{\vp_k}=\norm{\vp_k}_{\mat{A}}^2$ is the
energy norm of the search direction, strictly positive because $\mat{A}\succ0$ and
$\vp_k\neq0$, so the division is always safe until the residual vanishes.
The numerator is just the squared length of the current residual---a quantity we
need anyway to test convergence. One operator apply $\mat{A}\vp_k$ and two inner
products give the step.

\subsection{The residual update}
\label{sec:residual}

Having stepped, we need the new residual $\vr_{k+1}=\vb-\mat{A}\vx_{k+1}$. Do not
recompute it from scratch with a fresh matrix-vector product---that would cost a
second apply. Instead subtract the change:
\[
  \vr_{k+1} = \vb-\mat{A}(\vx_k+\alpha_k\vp_k)
            = \vr_k - \alpha_k\,\mat{A}\vp_k .
\]
The product $\mat{A}\vp_k$ was already computed for the step length
\eqref{eq:alpha2}, so the residual update is a single \code{axpy} (a
scaled-vector add) and costs \emph{no} extra apply. This is the reason a conjugate
gradient iteration needs exactly \emph{one} operator application per step---the
governing cost metric of \cref{ch:krylov}. The recurrence
\begin{equation}
\label{eq:rupdate}
  \vr_{k+1} = \vr_k - \alpha_k\,\mat{A}\vp_k
\end{equation}
also drives two orthogonality facts we now harvest.

\subsection{The direction update and the short-recurrence miracle}
\label{sec:beta}

The new direction must be conjugate to all previous ones. We build it as the new
residual plus a correction along the \emph{previous} direction only:
\begin{equation}
\label{eq:pupdate}
  \vp_{k+1} = \vr_{k+1} + \beta_{k}\,\vp_k ,
\end{equation}
and choose $\beta_k$ to force conjugacy with $\vp_k$, i.e.
$\inner{\mat{A}\vp_{k+1}}{\vp_k}=0$. Substituting \eqref{eq:pupdate} and solving,
\[
  \beta_k = -\,\frac{\inner{\mat{A}\vr_{k+1}}{\vp_k}}{\inner{\mat{A}\vp_k}{\vp_k}} .
\]
This forces conjugacy against $\vp_k$. The astonishing fact---the heart of the
method---is that for an SPD operator, forcing conjugacy against the \emph{one}
previous direction automatically delivers conjugacy against \emph{all} the older
ones, for free. Two orthogonality relations make it happen, both consequences of
the residual recurrence \eqref{eq:rupdate}:

\begin{lemma}[CG orthogonalities]
\label{lem:orth}
Running the recurrences \eqref{eq:alpha2}, \eqref{eq:rupdate}, \eqref{eq:pupdate}
from $\vp_0=\vr_0$ produces, for all valid indices,
\begin{enumerate}[nosep, label=\textup{(\roman*)}]
  \item mutually orthogonal residuals: $\inner{\vr_i}{\vr_j}=0$ for $i\neq j$;
  \item $\mat{A}$-conjugate directions: $\inner{\mat{A}\vp_i}{\vp_j}=0$ for $i\neq j$.
\end{enumerate}
\end{lemma}

The two facts feed each other by induction: orthogonal residuals make the
direction built by \eqref{eq:pupdate} conjugate to \emph{every} prior direction
(not just the last), and conjugate directions in turn make the next residual
orthogonal to every prior residual. The full induction is in any of the further
reading; the point for us is its consequence. Because the residuals are mutually
orthogonal, the coefficient $\beta_k$ collapses to a ratio of \emph{consecutive
residual norms}. Using \eqref{eq:rupdate} to rewrite
$\mat{A}\vp_k=(\vr_k-\vr_{k+1})/\alpha_k$ and applying orthogonality term by term,
the messy $\beta_k$ above simplifies to
\begin{equation}
\label{eq:beta}
  \boxed{\;\beta_k = \frac{\inner{\vr_{k+1}}{\vr_{k+1}}}{\inner{\vr_k}{\vr_k}}\;}
\end{equation}
the celebrated Fletcher--Reeves form. (The same orthogonalities collapse
$\inner{\vr_k}{\vp_k}$ to $\inner{\vr_k}{\vr_k}$, which is the identity promised
in \cref{sec:alpha}.) No old direction, no old residual: $\beta_k$ is built from
the squared residual norm we just computed and the one from the step before.

\begin{keyidea}
\textbf{The short-recurrence miracle.} For an SPD operator, a direction made
conjugate to only the \emph{previous} direction is automatically conjugate to
\emph{all} previous directions. Consequently conjugate gradients keeps just two
vectors---the residual $\vr$ and the direction $\vp$---and two scalars---the step
length $\alpha$ and the coefficient $\beta=\norm{\vr_{k+1}}^2/\norm{\vr_k}^2$. It
never stores or re-orthogonalizes against the old directions, yet it builds a
full conjugate basis. This $O(N)$-memory short recurrence is what makes CG
practical at a million unknowns---and it is exactly what GMRES, on a
\emph{non}-symmetric operator, cannot do (\cref{ch:gmres}).
\end{keyidea}

\section{The algorithm}
\label{sec:algorithm}

Collecting \eqref{eq:alpha2}, \eqref{eq:rupdate}, \eqref{eq:pupdate},
\eqref{eq:beta} gives the whole method. We present it twice: first as ordinary
pseudocode, then as the one \code{iterate\_while} of the shared step kernel, to
make the tie to \cref{ch:fold} and \cref{ch:krylov} explicit.

\subsection{As pseudocode}

\begin{lstlisting}[style=l4style]
-- Conjugate gradients: solve A x = b for SPD A.
cg A b x0 tol =
  let r0 = b - A x0                    -- initial residual
      p0 = r0                          -- first direction is steepest descent
  in loop with x = x0, r = r0, p = p0:
       Ap    = A p                     -- THE one operator apply per step
       alpha = <r,r> / <Ap,p>          -- step length          (eq 26.6)
       x'    = x + alpha * p           -- advance the iterate
       r'    = r - alpha * Ap          -- update residual, no new apply (eq 26.7)
       if sqrt <r',r'> <= tol: return x'
       beta  = <r',r'> / <r,r>         -- direction coefficient (eq 26.11)
       p'    = r' + beta * p           -- new conjugate direction (eq 26.8)
       continue with x = x', r = r', p = p'
\end{lstlisting}

Per step: one apply $\mat{A}\vp$, two inner products ($\inner{\vr}{\vr}$ and
$\inner{\mat{A}\vp}{\vp}$; the next $\inner{\vr'}{\vr'}$ is reused as the
following step's numerator), and three scaled-vector adds (the updates of $\vx$,
$\vr$, and $\vp$). \Cref{fig:dataflow} draws this as a dataflow graph: one apply,
two dots, three \code{axpy}s. That fixed, tiny per-step recipe---independent of
$N$ beyond the vector lengths---is the whole computational kernel.

\begin{figure}[t]
\centering
\begin{tikzpicture}[font=\footnotesize, node distance=6mm and 9mm]
  \node[data, minimum width=9mm] (r) {$\vr$};
  \node[data, minimum width=9mm, below=10mm of r] (p) {$\vp$};
  \node[opbox, right=of p] (ap) {\code{A·p}\\\scriptsize 1 apply};
  \node[opbox, right=12mm of ap] (d2) {\code{<Ap,p>}\\\scriptsize dot};
  \node[opbox, above=8mm of d2] (d1) {\code{<r,r>}\\\scriptsize dot};
  \node[opbox, right=10mm of d2] (al) {$\alpha$};
  \node[opbox, right=12mm of d1] (xu) {\code{x += a p}\\\scriptsize axpy};
  \node[opbox, below=6mm of xu] (ru) {\code{r -= a Ap}\\\scriptsize axpy};
  \node[opbox, below=6mm of ru] (be) {$\beta=\frac{<r',r'>}{<r,r>}$};
  \node[opbox, right=10mm of ru] (pu) {\code{p = r'+b p}\\\scriptsize axpy};
  \node[product, right=9mm of pu] (out) {$\vx',\vr',\vp'$};
  \draw[flow] (r)-|(d1);
  \draw[flow] (p)--(ap);
  \draw[flow] (ap)--(d2);
  \draw[flow] (ap.east) .. controls +(0.5,-0.7) and +(-0.5,0) .. (ru.west);
  \draw[flow] (d1)--(al);  \draw[flow] (d2)--(al);
  \draw[flow] (al) .. controls +(0,0.9) and +(0,0.9) .. (xu.north);
  \draw[flow] (al) .. controls +(0.2,-0.3) and +(-0.2,0.3) .. (ru.east);
  \draw[flow] (p) .. controls +(-0.5,0) and +(-0.6,-0.4) .. (xu.west);
  \draw[flow] (xu)--(out);
  \draw[flow] (ru)--(be);  \draw[flow] (ru) .. controls +(0.3,0) and +(-0.3,0.2) .. (pu.west);
  \draw[flow] (be) .. controls +(0.6,0) and +(-0.3,-0.4) .. (pu.south);
  \draw[flow] (pu)--(out);
\end{tikzpicture}
\caption[The CG step dataflow]{One conjugate-gradient step as a dataflow graph:
\emph{one} operator apply $\mat{A}\vp$, \emph{two} inner products
($\inner{\vr}{\vr}$ and $\inner{\mat{A}\vp}{\vp}$), and \emph{three}
scaled-vector adds (\code{axpy}s) updating $\vx$, $\vr$, and $\vp$. The product
$\mat{A}\vp$ feeds both the denominator of $\alpha$ and the residual update---it
is computed once and reused, which is why a step costs a single apply. Compare the
threaded-carry picture of \cref{fig:carry}: this graph is the \emph{body} of one
\texttt{iterate\_while} step.}
\label{fig:dataflow}
\end{figure}

\subsection{As one \texttt{iterate\_while} of the step kernel}
\label{sec:iterate}

\Cref{ch:fold} taught that every iterative algorithm in this book is one
\code{iterate\_while}, and \cref{ch:krylov} taught that every Krylov solver is
that combinator folding a single \code{krylov\_step}. Conjugate gradients is the
canonical instance. We read the algorithm through the strata of \cref{ch:strata}.

The \emph{carry} is the run-time \code{SimState}---the iterate $\vx$, the
iteration count, the convergence flag, the residual statistics---exactly the
five-field record that is uniform across every Krylov method. The
\emph{ephemeral} workspace is the CG \code{Krylov} bundle: the residual $\vr$, the
direction $\vp$, the optional preconditioned residual $\vz$, and the two scalars
$\alpha,\beta$. It is solve-local scratch, born when the solve starts and
discarded when it returns, threaded as a plain value beside the carry rather than
through the monad. The \emph{predicate} reads the residual proxy folded into the
carry and stops at tolerance---the discipline of \cref{ch:fold}, never reading a
quantity the step has not yet produced. And the \emph{step} is one
\code{krylov\_step}: apply once, update the scalars, update the three vectors,
emit the residual norm as a demand-prunable output.

\begin{lstlisting}[style=l4style]
cg_solve :: OpParams -> SimState -> { final_state: SimState, trajectory: [...] }
cg_solve op s0 =
  iterate_while
    s0                                 -- carry: the SimState (x, it, converged, ...)
    (\s -> not s.converged && s.it < op.max_it)   -- predicate: pure on the carry
    (\s -> krylov_step op krylov s)     -- step: one apply, fold the CG body
-- where krylov_step performs, per call:
--   Ap    = apply_linop op.T  krylov.p          -- the ONE apply
--   alpha = dot(krylov.r, krylov.r) / dot(Ap, krylov.p)
--   x'    = axpy  alpha krylov.p    s.x          -- iterate update
--   r'    = axpy (-alpha) Ap        krylov.r     -- residual update
--   beta  = dot(r', r') / dot(krylov.r, krylov.r)
--   p'    = axpy  beta   krylov.p   r'           -- direction update
--   outputs = { residual_norm: sqrt (dot(r',r')) }   -- demand-prunable
\end{lstlisting}

The coefficient $\beta=\norm{\vr'}^2/\norm{\vr}^2$ reads the previous step's
residual-norm-squared. That single backward dependence is the
\emph{recurrence carry} of \cref{ch:strata}'s fourth stratum---the scalar
threaded from one inner step to the next and reborn at each solve. In the
branch-in-body form (\cref{ch:fold}'s Form~A) it lives as a field of the workspace
with an \lstinline[style=l4style]|if it == 0| guard for the first step (where
$\vp_0=\vr_0$, no previous direction); the first-iteration-unrolling rotation
splits it into a bootstrap step and a branch-free steady step
(\code{iterate\_while\_with\_prev}), exactly as \cref{ch:fold} described.
\Cref{fig:carry} draws the CG carry threading forward with the residual history
peeling off into the trajectory.

\begin{figure}[t]
\centering
\begin{tikzpicture}[font=\footnotesize, node distance=6mm]
  \node[data, minimum width=14mm, align=center] (s0) at (0,0) {\code{SimState}\\$\vx_0$};
  \node[diamond, draw=simBlue, fill=simBgBlue, aspect=1.7, inner sep=1pt,
        font=\scriptsize] (p0) at (2.2,0) {res$>$tol?};
  \node[opbox, minimum width=15mm, align=center] (k0) at (4.6,0) {\code{krylov\_step}\\\scriptsize 1 apply};
  \node[data, minimum width=14mm, align=center] (s1) at (7.2,0) {\code{SimState}\\$\vx_1$};
  \node[diamond, draw=simBlue, fill=simBgBlue, aspect=1.7, inner sep=1pt,
        font=\scriptsize] (p1) at (9.4,0) {res$>$tol?};
  \node[font=\scriptsize, simGray] (more) at (11.0,0) {$\cdots$};
  \draw[flow] (s0)--(p0);
  \draw[flow] (p0)-- node[above,font=\scriptsize]{yes} (k0);
  \draw[flow] (k0)-- node[above,font=\scriptsize]{\code{.state}} (s1);
  \draw[flow] (s1)--(p1);
  \draw[flow] (p1)-- node[above,font=\scriptsize]{yes} (more);
  % ephemeral Krylov threaded below as plain value
  \node[data, fill=simBgGreen, draw=simGreen, minimum width=20mm] (kr0) at (4.6,-1.5)
    {\scriptsize\code{Krylov}: $\vr,\vp,\alpha,\beta$};
  \draw[refedge, simGreen] (kr0.north) -- (k0.south);
  \draw[backflow, simGreen] (k0.south east) .. controls +(0.6,-0.8) and +(0.6,0.4) .. (kr0.east)
    node[midway, right=2mm, font=\scriptsize, text=simGreen] {plain value};
  % trajectory of residual norms peeling off
  \node[data, fill=simBgGold, draw=simGold, minimum width=40mm] (tj) at (6.5,-2.7)
    {\code{trajectory}: $[\,\norm{\vr_0},\ \norm{\vr_1},\ \dots\,]$};
  \draw[refedge, simGold!80!black] (k0.south) |- (tj.west);
  % converged exit
  \node[product] (fin) at (9.4,-1.5) {\code{final\_state}: $\vx^\star$};
  \draw[backflow] (p1.south)--(fin.north);
\end{tikzpicture}
\caption[CG as one \texttt{iterate\_while}]{Conjugate gradients as a single
\texttt{iterate\_while} (\cref{ch:fold}) folding one \code{krylov\_step}
(\cref{ch:krylov}). The carry is the run-time \code{SimState}
(\cref{ch:strata})---iterate $\vx$, counter, flags---threaded along the top; the
predicate reads the residual proxy folded into the carry and exits at tolerance.
The ephemeral \code{Krylov} workspace ($\vr,\vp,\alpha,\beta$) is threaded
\emph{beside} the carry as a plain value (green), reborn at solve entry and
discarded at return. Each step's residual norm peels off into the
\code{trajectory} (gold)---computed only if a consumer reads it, by the
demand-pruning of \cref{ch:fold}.}
\label{fig:carry}
\end{figure}

This is the structural payoff of Parts~II and~III landing on a real solver:
nothing in \cref{sec:algorithm} is new control structure. CG is the fold of
\cref{ch:fold}, over the strata of \cref{ch:strata}, applying the operator of
\cref{ch:krylov} once per step. What is special to CG lives entirely in the
\emph{body}---the four coefficient formulas we derived---not in the loop.

\section{Convergence}
\label{sec:converge}

Finite termination in $N$ steps (\cref{thm:finite}) is reassuring but useless as a
performance promise: $N$ is in the millions, and rounding spoils exact conjugacy
long before then. What makes CG \emph{fast} is that it usually drives the error
below tolerance in \emph{far fewer} steps---and how many fewer is governed by one
number, the condition number.

\begin{definition}[Condition number]
\label{def:kappa}
For an SPD operator $\mat{A}$ with eigenvalues
$0<\lambda_{\min}=\lambda_1\le\cdots\le\lambda_N=\lambda_{\max}$, the (spectral)
\emph{condition number} is
$\kappa=\kappa(\mat{A})=\lambda_{\max}/\lambda_{\min}$.
\end{definition}

$\kappa$ is the eccentricity of the bowl: $\kappa=1$ is a perfect round bowl (one
step finishes), and large $\kappa$ is a long narrow valley (the steepest-descent
zig-zag of \cref{sec:zigzag}). The central convergence result bounds the
energy-norm error of the $k$-th CG iterate by a power that shrinks as $\kappa$
shrinks.

\begin{theorem}[CG energy-norm convergence]
\label{thm:converge}
Let $\vx^\star$ solve $\mat{A}\vx=\vb$ with $\mat{A}=\mat{A}^{T}\succ0$ and
$\kappa=\kappa(\mat{A})$. The $k$-th conjugate-gradient iterate satisfies
\begin{equation}
\label{eq:bound}
  \norm{\vx^\star-\vx_k}_{\mat{A}}
  \;\le\; 2\!\left(\frac{\sqrt{\kappa}-1}{\sqrt{\kappa}+1}\right)^{\!k}
  \norm{\vx^\star-\vx_0}_{\mat{A}} .
\end{equation}
\end{theorem}

The structure of the proof is worth a sentence even though we do not give it in
full: after $k$ steps the CG iterate is the \emph{best possible} approximation, in
the energy norm, drawn from the degree-$k$ Krylov subspace
$\operatorname{span}\{\vr_0,\mat{A}\vr_0,\dots,\mat{A}^{k}\vr_0\}$
(\cref{ch:krylov})---that optimality is what conjugacy buys. Bounding the error
then becomes a question about polynomials: how small can a degree-$k$ polynomial
$q$ with $q(0)=1$ be made across the eigenvalue interval
$[\lambda_{\min},\lambda_{\max}]$? The answer is given by Chebyshev polynomials,
and it produces exactly the factor in \eqref{eq:bound}. The
$\sqrt{\kappa}$---not $\kappa$---is the decisive feature: CG converges at the rate
set by the \emph{square root} of the condition number, which is why it so
dramatically outperforms steepest descent (whose analogous bound carries $\kappa$
itself).

\begin{figure}[t]
\centering
\begin{tikzpicture}
  \begin{semilogyaxis}[
      width=12.2cm, height=7.0cm,
      xlabel={iteration $k$}, ylabel={relative energy-norm error (bound)},
      xmin=0, xmax=60, ymin=1e-8, ymax=3,
      grid=both, grid style={simGray!18},
      legend pos=north east, legend cell align=left,
      legend style={font=\footnotesize, draw=simGray!40},
      label style={font=\footnotesize}, tick label style={font=\scriptsize},
      every axis plot/.append style={very thick}]
    % bound 2*c^k with c=(sqrt(kappa)-1)/(sqrt(kappa)+1); written as 2*exp(k*ln c)
    \addplot[simGreen, domain=0:60, samples=61]
      {2*exp(x*ln((sqrt(10)-1)/(sqrt(10)+1)))};
    \addlegendentry{$\kappa=10$ (well-conditioned)}
    \addplot[simBlue, domain=0:60, samples=61]
      {2*exp(x*ln((sqrt(100)-1)/(sqrt(100)+1)))};
    \addlegendentry{$\kappa=100$}
    \addplot[simGold, domain=0:60, samples=61]
      {2*exp(x*ln((sqrt(1000)-1)/(sqrt(1000)+1)))};
    \addlegendentry{$\kappa=1000$}
    \addplot[simRust, domain=0:60, samples=61]
      {2*exp(x*ln((sqrt(10000)-1)/(sqrt(10000)+1)))};
    \addlegendentry{$\kappa=10000$ (ill-conditioned)}
    \draw[simGray, densely dashed] (axis cs:0,1e-6)--(axis cs:60,1e-6)
      node[pos=0.07, above, font=\scriptsize, text=simGray] {tol $=10^{-6}$};
  \end{semilogyaxis}
\end{tikzpicture}
\caption[CG convergence versus condition number]{The CG energy-norm error bound
$2\big((\sqrt\kappa-1)/(\sqrt\kappa+1)\big)^k$ of \cref{thm:converge}, for four
condition numbers, on a log scale. A well-conditioned system ($\kappa=10$) reaches
a $10^{-6}$ tolerance in a handful of steps; each tenfold increase in $\kappa$
roughly triples the iteration count, because the rate depends on $\sqrt\kappa$.
This is the entire motivation for preconditioning (\cref{ch:precond}): replace
$\mat{A}$ by $\mat{M}^{-1}\mat{A}$ with a far smaller condition number and slide
down to the green curve. The dependence on $\sqrt\kappa$ rather than $\kappa$ is
CG's decisive advantage over steepest descent.}
\label{fig:convplot}
\end{figure}

\Cref{fig:convplot} plots the bound for several condition numbers. The lesson is
stark: a well-conditioned system converges in a handful of steps; a poorly
conditioned one crawls. Since the operators Palace assembles are routinely
ill-conditioned---fine meshes and material contrasts push $\kappa$ into the
thousands or worse---raw CG is too slow, and the fix is to \emph{change the
condition number}. That is preconditioning, the subject of \cref{ch:precond} and
the reason \cref{sec:pcg} exists.

\subsection{Clustered eigenvalues and superlinear convergence}
\label{sec:cluster}

The bound \eqref{eq:bound} is worst-case; CG often does much better, and the
reason is the polynomial picture again. The error after $k$ steps is small if
\emph{some} degree-$k$ polynomial with $q(0)=1$ is small on the spectrum---and a
polynomial can be made small on a few tight \emph{clusters} of eigenvalues using
far fewer degrees than it would need to be small across a whole wide interval.
Concretely, if the spectrum consists of $m$ tight clusters, CG behaves as if
$\kappa$ were the spread \emph{within} a cluster, and it ``uses up'' the
outliers a few per step. The visible effect is \emph{superlinear convergence}: the
residual drops slowly at first, then accelerates as CG eliminates the extreme
eigenvalues and is left working only against the cluster. This is also why a good
preconditioner---one that clusters the spectrum near $1$---helps so much more than
the bare $\kappa$ improvement suggests.

\section{Preconditioned conjugate gradients}
\label{sec:pcg}

If convergence is governed by $\kappa(\mat{A})$, the way to make CG fast is to
solve a system with a smaller condition number that has the same solution.
Preconditioning supplies a \emph{preconditioner} $\mat{M}\approx\mat{A}$, cheap to
invert, such that $\mat{M}^{-1}\mat{A}$ is far better conditioned than $\mat{A}$.
The full theory is \cref{ch:precond}; here we record only how the CG \emph{loop}
changes, because the change is small and structural.

The naive idea---run CG on $\mat{M}^{-1}\mat{A}\vx=\mat{M}^{-1}\vb$---breaks the
SPD requirement, since $\mat{M}^{-1}\mat{A}$ is generally not symmetric. The
preconditioned conjugate gradient (PCG) algorithm finesses this by running CG in
the inner product defined by $\mat{M}$, which keeps every operator effectively
symmetric and recovers an honest CG. The upshot is the same recurrence with one
extra vector and one extra apply per step: at each step, form the
\emph{preconditioned residual} $\vz=\mat{M}^{-1}\vr$, and replace the residual
inner products $\inner{\vr}{\vr}$ in $\alpha$ and $\beta$ by the
\emph{$\mat{M}$-weighted} inner products $\inner{\vz}{\vr}$:
\begin{align}
\label{eq:pcg-alpha}
  \vz_k &= \mat{M}^{-1}\vr_k, &
  \alpha_k &= \frac{\inner{\vz_k}{\vr_k}}{\inner{\mat{A}\vp_k}{\vp_k}}, \\
\label{eq:pcg-beta}
  \beta_k &= \frac{\inner{\vz_{k+1}}{\vr_{k+1}}}{\inner{\vz_k}{\vr_k}}, &
  \vp_{k+1} &= \vz_{k+1} + \beta_k\,\vp_k .
\end{align}
When $\mat{M}=\mat{I}$ this is exactly the unpreconditioned method of
\cref{sec:algorithm}---$\vz=\vr$ and the formulas collapse back. The cost is one
preconditioner apply $\vz=\mat{M}^{-1}\vr$ per step, on top of the one operator
apply, and the storage of the extra vector $\vz$ (the optional \code{z?} field of
the CG \code{Krylov} bundle in \cref{ch:strata}). \Cref{fig:pcg} overlays the
extra apply on the dataflow of \cref{fig:dataflow}.

\begin{figure}[t]
\centering
\begin{tikzpicture}[font=\footnotesize, node distance=7mm and 11mm]
  \node[data, minimum width=9mm] (r) {$\vr$};
  \node[extern, right=of r] (M) {$\vz=\mat{M}^{-1}\vr$\\\scriptsize +1 apply};
  \node[opbox, right=of M] (dot) {\code{<z,r>}\\\scriptsize dot};
  \node[data, minimum width=9mm, below=12mm of r] (p) {$\vp$};
  \node[opbox, right=of p] (ap) {\code{A·p}\\\scriptsize 1 apply};
  \node[opbox, right=of ap] (dap) {\code{<Ap,p>}};
  \node[opbox, right=10mm of dap] (al) {$\alpha,\beta$};
  \node[opbox, right=10mm of dot] (upd) {\code{x,r,p}\\updates\\\scriptsize 3 axpy};
  \node[product, below=10mm of upd] (out) {$\vx',\vr',\vp'$};
  \draw[flow] (r)--(M); \draw[flow] (M)--(dot);
  \draw[flow] (p)--(ap); \draw[flow] (ap)--(dap);
  \draw[flow] (dot)--(al); \draw[flow] (dap)--(al);
  \draw[flow] (al) .. controls +(0,0.6) and +(0.4,-0.2) .. (upd.east);
  \draw[flow] (dot) .. controls +(0.4,0.3) and +(-0.4,0.2) .. (upd.west);
  \draw[flow] (M.south) .. controls +(0,-0.5) and +(-0.6,0.4) .. (upd.south west)
    node[midway, below left=-1mm, font=\scriptsize, text=simViolet] {$\vz$ into $\vp'$};
  \draw[flow] (upd)--(out);
\end{tikzpicture}
\caption[Preconditioned CG: the extra apply]{Preconditioned conjugate gradients
adds exactly one operation to the loop of \cref{fig:dataflow}: a preconditioner
apply $\vz=\mat{M}^{-1}\vr$ (violet, dashed---an \emph{opaque} operator surface,
\cref{ch:krylov}). The residual inner products $\inner{\vr}{\vr}$ become the
$\mat{M}$-weighted products $\inner{\vz}{\vr}$ in both $\alpha$ and $\beta$, and
the direction is built from $\vz$ rather than $\vr$. Everything else---the single
operator apply $\mat{A}\vp$, the three \code{axpy}s, the short
recurrence---is unchanged. PCG is CG with one extra apply and one extra vector.}
\label{fig:pcg}
\end{figure}

\begin{pitfall}
\textbf{Do not run CG on a non-SPD system.} Every guarantee in this
chapter---finite termination, the energy-norm bound, even the line-minimizing step
length being a \emph{minimum} rather than a saddle---rests on
$\mat{A}=\mat{A}^{T}\succ0$ (\cref{def:spd}). Hand CG a non-symmetric operator and
the energy ``inner product'' $\vu^{T}\mat{A}\vv$ is not symmetric, conjugacy is
ill-defined, and the iteration can stall or diverge with no warning. Hand it a
symmetric but \emph{indefinite} operator (negative eigenvalues---the
time-harmonic Maxwell operator $\mat{K}-\omega^2\mat{M}$ above resonance is one)
and the denominator $\inner{\mat{A}\vp}{\vp}$ can vanish or go negative,
\emph{dividing by zero or stepping uphill}. The symptom is a residual that
refuses to fall or suddenly blows up. The fix is to use the right method: GMRES
(\cref{ch:gmres}) for non-symmetric systems, MINRES for symmetric-indefinite ones.
A subtle trap: applying a non-SPD \emph{preconditioner} $\mat{M}$ inside PCG
breaks SPD-ness just as surely as a non-SPD $\mat{A}$---the preconditioner must
itself be SPD for \eqref{eq:pcg-alpha}--\eqref{eq:pcg-beta} to define an honest CG.
\end{pitfall}

\section{The load-bearing numerical note: reduction order}
\label{sec:numerics}

Every coefficient in CG---$\alpha$, $\beta$, the convergence test---is built from
\emph{inner products}, $\inner{\vu}{\vv}=\sum_i u_i v_i$. An inner product is a
\emph{global reduction}: it sums a product over all $N$ entries of the vectors.
And floating-point addition is \emph{not associative}: $(a+b)+c$ and $a+(b+c)$ can
differ in their last bits. So the value of an inner product---and therefore the
exact $\alpha$, the exact $\beta$, the exact iterate $\vx_k$, and the exact
iteration at which the residual first dips below tolerance---depends on the
\emph{order} in which the reduction is summed.

This is not a defect to be engineered away; it is a load-bearing property in the
sense of \cref{ch:bigidea}. Different reduction orders (a different summation tree,
a different thread count, a different device) give bit-different but equally valid
CG runs: the mathematics is identical, the floating-point realization is not. Two
runs of the ``same'' solve on two machines may take a different number of
iterations and return iterates differing in their last bits. The converged answer
is correct to tolerance either way; what is \emph{not} reproducible bit-for-bit is
the trajectory. The Palace source records this explicitly, and we treat the
reduction-order semantics in full in \cref{app:numerics}. The discipline for now is
simply to \emph{name} it: an inner product is a reduction, a reduction has an
order, and the order is observable in the last bits. When a downstream consumer
demands bit-reproducibility, it must pin the reduction order---CG does not, and
cannot, promise it for free.

\section{Worked examples}

\begin{workedexample}{CG on a $2\times2$ SPD system: two steps to exact}{2x2}
We run unpreconditioned CG on a $2\times2$ SPD system and watch it reach the exact
answer in two steps, as \cref{thm:finite} promises for $N=2$. Take
\[
  \mat{A}=\begin{pmatrix}4 & 1\\[1pt] 1 & 3\end{pmatrix},\qquad
  \vb=\begin{pmatrix}1\\[1pt] 2\end{pmatrix},\qquad
  \vx_0=\begin{pmatrix}0\\[1pt] 0\end{pmatrix}.
\]
$\mat{A}$ is symmetric, and its eigenvalues
$\tfrac{7\pm\sqrt5}{2}\approx\{2.38,\,4.62\}$ are positive, so it is SPD. The
exact solution is $\vx^\star=\mat{A}^{-1}\vb=\tfrac1{11}(1,7)^{T}
\approx(0.0909,\,0.6364)^{T}$.

\tcblower
\textbf{Step 0.} Residual $\vr_0=\vb-\mat{A}\vx_0=(1,2)^{T}$; first direction
$\vp_0=\vr_0=(1,2)^{T}$.
\begin{align*}
  \mat{A}\vp_0 &= (4\cdot1+1\cdot2,\;1\cdot1+3\cdot2)^{T}=(6,7)^{T}, &
  \inner{\vr_0}{\vr_0} &= 1+4 = 5,\\
  \inner{\mat{A}\vp_0}{\vp_0} &= 6\cdot1+7\cdot2 = 20, &
  \alpha_0 &= \tfrac{5}{20} = 0.25.
\end{align*}
Update: $\vx_1=\vx_0+\alpha_0\vp_0=(0.25,\,0.5)^{T}$ and
$\vr_1=\vr_0-\alpha_0\mat{A}\vp_0=(1,2)^{T}-0.25(6,7)^{T}=(-0.5,\,0.25)^{T}$.
Check orthogonality (\cref{lem:orth}): $\inner{\vr_1}{\vr_0}=-0.5+0.5=0$.~\checkmark

\textbf{Direction update.}
$\inner{\vr_1}{\vr_1}=0.25+0.0625=0.3125$, so
$\beta_0=\tfrac{0.3125}{5}=0.0625$, and
$\vp_1=\vr_1+\beta_0\vp_0=(-0.5,0.25)^{T}+0.0625(1,2)^{T}=(-0.4375,\,0.375)^{T}$.
Check conjugacy: $\mat{A}\vp_1=(-1.375,\,0.6875)^{T}$, and
$\inner{\mat{A}\vp_1}{\vp_0}=-1.375+1.375=0$.~\checkmark

\textbf{Step 1.}
$\inner{\mat{A}\vp_1}{\vp_1}=(-1.375)(-0.4375)+(0.6875)(0.375)=0.8594$, so
$\alpha_1=\tfrac{0.3125}{0.8594}=0.3636$.
\[
  \vx_2 = \vx_1+\alpha_1\vp_1
        = (0.25,0.5)^{T}+0.3636\,(-0.4375,0.375)^{T}
        = (0.0909,\,0.6364)^{T}.
\]
This is $\vx^\star$ to four places: in exactly $N=2$ steps the residual is zero
and CG has landed on the solution. The two steps installed the two conjugate
coordinates of the error, one each, with no interference---finite termination made
arithmetic.
\end{workedexample}

\begin{workedexample}{Steepest descent versus CG on an ill-conditioned bowl}{zigzag}
We put steepest descent and CG on the same stretched system and watch one
zig-zag while the other finishes in two steps. Take the diagonal SPD operator
\[
  \mat{A}=\begin{pmatrix}10 & 0\\[1pt] 0 & 1\end{pmatrix},\qquad
  \vb=\begin{pmatrix}0\\[1pt] 0\end{pmatrix},\qquad
  \vx_0=\begin{pmatrix}1\\[1pt] 1\end{pmatrix},
\]
so $\vx^\star=0$ and $\phi(\vx)=5x_1^2+\tfrac12x_2^2$ is a valley ten times
steeper across than along: $\kappa=10$. \Cref{fig:zigzag} plots both iterate
paths on the contours.

\tcblower
\textbf{Steepest descent} (\cref{eq:sdstep}). From $\vx_0=(1,1)^{T}$,
$\vr_0=-\mat{A}\vx_0=(-10,-1)^{T}$,
$\alpha_0=\inner{\vr_0}{\vr_0}/\inner{\mat{A}\vr_0}{\vr_0}
=\tfrac{101}{1010}=0.1$, giving $\vx_1=(0,\,0.9)^{T}$. The next residual
$(0,-0.9)^{T}$ is \emph{perpendicular} to the last step, so the iterate turns
ninety degrees and lands at $(0.818,0)^{T}$, then turns again. It crabs
$(1,1)\to(0,0.9)\to(0.82,0)\to(0,0.74)\to\cdots$, the error shrinking by only a
factor $\approx0.82$ per step---dozens of steps to reach $10^{-6}$.

\textbf{Conjugate gradients}. The same first step gives $\vx_1=(0,0.9)^{T}$ and
$\vr_1=(0,-0.9)^{T}$. But now $\beta_0=\inner{\vr_1}{\vr_1}/\inner{\vr_0}{\vr_0}
=\tfrac{0.81}{101}=0.00802$, and the direction
$\vp_1=\vr_1+\beta_0\vp_0=(-0.0802,\,-0.908)^{T}$ \emph{bends away} from the pure
residual so as to be $\mat{A}$-conjugate to $\vp_0$. One step along it,
$\alpha_1=0.901$, lands exactly on $\vx^\star=0$. \textbf{Two steps, done}---the
same finite termination as \cref{wex:2x2}, now visibly beating the
steepest-descent crawl on the identical bowl.
\end{workedexample}

\begin{figure}[t]
\centering
\begin{tikzpicture}
  \begin{axis}[
      width=12.4cm, height=7.2cm, axis equal image,
      xlabel={$x_1$}, ylabel={$x_2$},
      xmin=-0.55, xmax=1.15, ymin=-0.35, ymax=1.2,
      label style={font=\footnotesize}, tick label style={font=\scriptsize},
      legend pos=north west, legend cell align=left,
      legend style={font=\footnotesize, draw=simGray!40}]
    % elliptical contours of phi = 5 x1^2 + 0.5 x2^2 (levels 0.5, 2, 5)
    \addplot[simGray!45, samples=80, domain=0:360, variable=\t, forget plot]
      ({sqrt(0.5/5)*cos(\t)}, {sqrt(0.5/0.5)*sin(\t)});
    \addplot[simGray!45, samples=80, domain=0:360, variable=\t, forget plot]
      ({sqrt(2/5)*cos(\t)}, {sqrt(2/0.5)*sin(\t)});
    \addplot[simGray!45, samples=80, domain=0:360, variable=\t, forget plot]
      ({sqrt(5/5)*cos(\t)}, {sqrt(5/0.5)*sin(\t)});
    % steepest descent zig-zag path
    \addplot[simRust, very thick, mark=*, mark size=1.3pt]
      coordinates {(1,1) (0,0.9) (0.818,0) (0,0.744) (0.676,0) (0,0.614)};
    \addlegendentry{steepest descent (zig-zag)}
    % CG path: 2 steps
    \addplot[simBlue, very thick, mark=square*, mark size=1.6pt]
      coordinates {(1,1) (0,0.9) (0,0)};
    \addlegendentry{conjugate gradients (2 steps)}
    \node[circle, fill=simGreen, inner sep=1.4pt,
      label={[font=\scriptsize,text=simGreen]below right:$\vx^\star$}] at (axis cs:0,0) {};
  \end{axis}
\end{tikzpicture}
\caption[Steepest descent versus CG on the contours]{The classic comparison, on
the bowl $\phi=5x_1^2+\tfrac12 x_2^2$ ($\kappa=10$) of \cref{wex:zigzag}.
\emph{Steepest descent} (rust) takes the locally steepest step each time and so
turns ninety degrees at every iterate, zig-zagging down the valley and needing
many steps. \emph{Conjugate gradients} (blue) takes the same first step, then
bends its second direction to be $\mat{A}$-conjugate to the first and lands
exactly on $\vx^\star$ (green) in \emph{two} steps---the finite termination of
\cref{thm:finite}. The conjugacy correction is precisely the difference between
the crawl and the finish.}
\label{fig:zigzag}
\end{figure}

\begin{workedexample}{How many iterations for a target tolerance?}{count}
The bound \eqref{eq:bound} lets us \emph{predict} the iteration count before
running, given the condition number. Suppose an electrostatic solve assembles an
operator with $\kappa=500$ and we want the relative energy-norm error below
$\varepsilon=10^{-6}$. How many CG iterations does \cref{thm:converge} guarantee?

\tcblower
We need $k$ with $2\,c^{k}\le\varepsilon$, where
$c=\dfrac{\sqrt\kappa-1}{\sqrt\kappa+1}$. With $\kappa=500$,
$\sqrt\kappa=22.36$, so
\[
  c = \frac{22.36-1}{22.36+1} = \frac{21.36}{23.36} = 0.9144 .
\]
Solving $2\,c^{k}\le10^{-6}$:
\[
  k \;\ge\; \frac{\ln(\varepsilon/2)}{\ln c}
    = \frac{\ln(5\times10^{-7})}{\ln(0.9144)}
    = \frac{-14.51}{-0.0895}
    \approx 162 .
\]
So the bound guarantees $10^{-6}$ within about \textbf{162 iterations}. Now
precondition down to $\kappa=10$ (a realistic multigrid target,
\cref{ch:precond}): $\sqrt\kappa=3.162$,
$c=\tfrac{2.162}{4.162}=0.5195$, and
\[
  k \ge \frac{-14.51}{\ln(0.5195)} = \frac{-14.51}{-0.6548}\approx 23 .
\]
\textbf{162 iterations become 23}---a sevenfold speedup, each preconditioned
iteration only modestly more expensive (one extra apply, \cref{sec:pcg}). And
because the bound is worst-case, clustered spectra (\cref{sec:cluster}) routinely
beat even these counts. This single calculation---iterations scale like
$\sqrt\kappa$---is the entire economic case for the preconditioning of
\cref{ch:precond}.
\end{workedexample}

\summarysection
Conjugate gradients solves $\mat{A}\vx=\vb$ for a \emph{symmetric positive
definite} operator $\mat{A}=\mat{A}^{T}\succ0$---the electrostatic, magnetostatic,
and inner-Poisson systems of this book. Solving is minimizing the quadratic
$\phi(\vx)=\tfrac12\vx^{T}\mat{A}\vx-\vb^{T}\vx$, whose level sets are ellipses
and whose negative gradient is the residual $\vr=\vb-\mat{A}\vx$. Steepest
descent takes the locally steepest step and zig-zags across an ill-conditioned
valley; CG instead marches along \emph{$\mat{A}$-conjugate} directions
($\vp_i^{T}\mat{A}\vp_j=0$), which correct independent error coordinates and reach
the exact answer in at most $N$ steps. The method is built from four formulas: a
step length $\alpha_k=\inner{\vr_k}{\vr_k}/\inner{\mat{A}\vp_k}{\vp_k}$, a residual
update $\vr_{k+1}=\vr_k-\alpha_k\mat{A}\vp_k$ (no second apply), a direction update
$\vp_{k+1}=\vr_{k+1}+\beta_k\vp_k$, and the coefficient
$\beta_k=\inner{\vr_{k+1}}{\vr_{k+1}}/\inner{\vr_k}{\vr_k}$. The
\emph{short-recurrence miracle}---conjugacy against the previous direction yields
conjugacy against all---keeps the cost at \emph{one apply, two dots, three axpys}
per step and the memory at $O(N)$. As a program, CG is one \code{iterate\_while}
(\cref{ch:fold}) folding one \code{krylov\_step} (\cref{ch:krylov}): the carry is
the run-time \code{SimState}, the ephemeral $\vr,\vp$ are the \code{Krylov}
workspace, and $\beta$'s backward dependence is the recurrence carry
(\cref{ch:strata}). Convergence is governed by the condition number $\kappa$: the
energy-norm error shrinks like
$\big((\sqrt\kappa-1)/(\sqrt\kappa+1)\big)^k$, so well-conditioned (or clustered)
spectra converge fast and ill-conditioned ones crawl---the motivation for
preconditioning (\cref{ch:precond}), which PCG folds in as one extra apply per
step. Finally, the inner products are global reductions whose summation order
changes the last bits, a load-bearing non-determinism (\cref{app:numerics}).

\furtherreading
The definitive textbook account of conjugate gradients, the Krylov-subspace
optimality behind the convergence bound, and preconditioned CG is
Saad~\citep{saad2003}; the polynomial (Chebyshev) argument for
\cref{thm:converge} and the clustered-spectrum superlinear story are developed
there in full. Trefethen and Bau~\citep{trefethenbau1997} give the cleanest short
derivation of CG as the Krylov-subspace energy-norm minimizer and the most
geometric treatment of the convergence rate's $\sqrt\kappa$ dependence; their
Lectures~38 and~40 pair naturally with \cref{sec:converge}. For the numerical
linear algebra---the orthogonality lemmas of \cref{lem:orth}, the stability of the
recurrences in floating point, and the relationship to the Lanczos process of
\cref{ch:krylovschur}---Golub and Van Loan~\citep{golubvanloan2013} is the
standard reference.

\exercisessection
\begin{enumerate}[nosep]
  \item Verify directly that $\Grad\phi(\vx)=\mat{A}\vx-\vb$ for
  $\phi(\vx)=\tfrac12\vx^{T}\mat{A}\vx-\vb^{T}\vx$ when $\mat{A}=\mat{A}^{T}$.
  Where exactly does symmetry get used? Show that without symmetry the gradient is
  $\tfrac12(\mat{A}+\mat{A}^{T})\vx-\vb$, and explain why CG therefore ``sees'' only
  the symmetric part of a non-symmetric operator.
  \item Run CG by hand on
  $\mat{A}=\bigl(\begin{smallmatrix}2&0\\0&5\end{smallmatrix}\bigr)$,
  $\vb=(2,5)^{T}$, $\vx_0=0$. How many steps does it take, and why fewer than the
  generic $N=2$? (Relate your answer to the eigenvalue structure and
  \cref{sec:cluster}.)
  \item Derive the steepest-descent step length \eqref{eq:sdstep} by minimizing
  $\phi(\vx+\alpha\vr)$ over $\alpha$. Then show the next residual satisfies
  $\inner{\vr_{k+1}}{\vr_k}=0$, the perpendicularity that causes the zig-zag.
  \item Show that the line-minimizing step length along $\vp_k$ is
  \eqref{eq:alpha}, and then, assuming the orthogonality
  $\inner{\vr_k}{\vp_{k-1}}=0$ and the direction update, prove
  $\inner{\vr_k}{\vp_k}=\inner{\vr_k}{\vr_k}$, so \eqref{eq:alpha} becomes the
  algorithm's form \eqref{eq:alpha2}.
  \item Using the bound \eqref{eq:bound}, find the number of CG iterations
  guaranteed to reach relative energy-norm error $10^{-4}$ for $\kappa=4$,
  $\kappa=100$, and $\kappa=2500$. Confirm that each fourfold (then $25$-fold)
  increase in $\kappa$ roughly doubles (then triples) the count, matching the
  $\sqrt\kappa$ scaling.
  \item Write the CG step of \cref{sec:iterate} explicitly in the
  \code{iterate\_while\_with\_prev} form of \cref{ch:fold}: identify the bootstrap
  step (where $\vp_0=\vr_0$ and there is no previous $\beta$), the steady step, and
  the recurrence carry. Which stratum (\cref{ch:strata}) does the carry live in,
  and why is it wrong to store it in \code{OpParams}?
  \item Count the floating-point operations in one CG step as a function of $N$
  (count the apply abstractly as ``one matvec''). Then count one PCG step. By how
  much does preconditioning increase the per-step cost, and under what condition on
  the iteration-count reduction does PCG still win overall? Tie your answer to
  \cref{wex:count}.
  \item \emph{(Reduction order.)} Two machines run the same CG solve but sum the
  inner products in different orders. Argue that they may report different
  iteration counts yet both return iterates correct to the stated tolerance.
  Which fields of the \code{SolveResult} (\cref{ch:krylov}) are bit-reproducible
  across the two runs, and which are not? (Forward-reference
  \cref{app:numerics}.)
\end{enumerate}

\chapter{GMRES and FGMRES}
\label{ch:gmres}

\chapterorient{The previous two chapters built the machinery of Krylov subspaces
and spent it on the friendliest possible case: a symmetric, positive-definite
operator, where conjugate gradients rides a three-term recurrence to the answer
on a fixed budget of storage. This chapter asks what happens when the operator is
\emph{not} symmetric---the driven analysis, the eigenmode inner solves, the
complex frequency-domain systems---where the short recurrence simply does not
exist. The answer is GMRES: keep the \emph{whole} Krylov basis, and at every step
pick the iterate that makes the residual as small as it can be. We build it from
the Arnoldi process up, watch a big minimization shrink to a tiny least-squares
problem, and see why the residual norm falls out for free at every step. Then we
pay off a promise made back in \cref{ch:operators}: when the preconditioner is
allowed to \emph{change every step}, plain GMRES breaks, and FGMRES---the same
algorithm carrying one extra basis---is exactly the repair the decision rule
predicts.}

\section{Why the short recurrence is gone}

Conjugate gradients (\cref{ch:cg}) is a small miracle of bookkeeping. To minimize
the energy of the error over the growing Krylov subspace $\mathcal{K}_m(\mat{A},
\vr_0)$, it never stores more than a handful of vectors: the new search direction
is $\mat{A}$-conjugate to all the previous ones, yet computing it requires only
the \emph{last} direction, because symmetry of $\mat{A}$ forces all the earlier
conjugacy conditions to hold automatically. The cost per step is fixed; the
storage is fixed; the algorithm is, in a precise sense, optimal on a budget.

That miracle has a single load-bearing hypothesis, and it is \emph{symmetry}. The
three-term recurrence at the heart of CG is the Lanczos recurrence, and Lanczos
tridiagonalizes $\mat{A}$ only because $\mat{A}=\mat{A}^{\top}$. The moment we drop
symmetry---and most of the operators in this book are not symmetric: the
frequency-domain system $\mat{A}(\omega)=\mat{K}+i\omega\mat{C}-\omega^2\mat{M}$ of
the driven analysis (\cref{ch:targets}) is complex and non-Hermitian; the shifted
operators inside an eigenvalue solve are indefinite---the recurrence collapses.
There is a theorem (Faber--Manteuffel) that makes this brutal: for a general
nonsymmetric $\mat{A}$, \emph{no} short recurrence can generate the optimal Krylov
iterate. You cannot have both optimality and bounded storage. You must give up one.

GMRES gives up bounded storage and keeps optimality. It commits to remembering the
\emph{entire} orthonormal basis of the Krylov subspace, and in exchange it delivers,
at every step, the genuinely best iterate the subspace can offer. The chapter is the
story of how it affords this---and, when the cost grows too large, how restarting
trades a controlled amount of optimality back for a memory bound.

\begin{notationbox}
Throughout, $\mat{A}$ is a general (possibly nonsymmetric, possibly complex
non-Hermitian) invertible operator on $\Real^N$ or $\Complex^N$; we solve
$\mat{A}\vx=\vb$. We write $\vx_0$ for the initial guess, $\vr_0=\vb-\mat{A}\vx_0$
for the initial residual, $\beta=\norm{\vr_0}$ for its norm, and
$\mathcal{K}_m=\mathcal{K}_m(\mat{A},\vr_0)=\operatorname{span}\{\vr_0,\mat{A}\vr_0,
\dots,\mat{A}^{m-1}\vr_0\}$ for the $m$-th Krylov subspace. The norm $\norm{\cdot}$
is always the Euclidean ($2$-)norm. Vectors $\vv_1,\dots,\vv_m$ are the Arnoldi
basis; $\vec e_1$ is the first coordinate vector.
\end{notationbox}

\subsection{The minimum-residual idea}

CG minimized the $\mat{A}$-norm of the error, a quantity that only makes sense---and
is only non-negative---when $\mat{A}$ is positive definite. For a general operator
that functional is meaningless. So GMRES minimizes something that always makes
sense: the ordinary Euclidean norm of the \emph{residual}.

\begin{keyidea}
The \textbf{minimum-residual principle.} At step $m$, GMRES chooses the iterate
$\vx_m$ from the affine Krylov subspace $\vx_0+\mathcal{K}_m$ that makes the
residual as small as possible:
\[
  \vx_m \;=\; \operatorname*{arg\,min}_{\vx\,\in\,\vx_0+\mathcal{K}_m}\ \norm{\vb-\mat{A}\vx}.
\]
Because $\mathcal{K}_m\subseteq\mathcal{K}_{m+1}$, the feasible set only grows from
one step to the next, so the minimum can only \emph{decrease}. The residual norms
are therefore \emph{monotonically non-increasing}: $\norm{\vr_0}\ge\norm{\vr_1}\ge
\norm{\vr_2}\ge\cdots$. GMRES never takes a step backward---a guarantee CG cannot
make for indefinite systems and that BiCGStab cannot make at all.
\end{keyidea}

\Cref{fig:minres} draws the principle. The affine subspace $\vx_0+\mathcal{K}_m$ is
a flat slice of $\Real^N$ that grows one dimension wider at each step. Over each
slice we minimize the residual; the minimizer is the point of the slice closest
(in the $\mat{A}$-image) to $\vb$. As the slice widens, the achievable minimum
drops, and the residual curve marches monotonically down.

\begin{figure}[t]
\centering
\begin{tikzpicture}[scale=1.0]
  % the ambient space hint
  \node[font=\scriptsize\itshape, simGray] at (-3.1,2.3) {$\Real^N$};
  % growing affine subspaces as nested line segments through x0
  \coordinate (x0) at (-2.6,-1.2);
  \fill[simInk] (x0) circle (1.4pt);
  \node[font=\scriptsize] at (x0) [below left=0pt and -1pt] {$\vx_0$};
  % K_1 : a short segment
  \draw[simGreen, thick] (x0) -- ++(2.2,0.5)
    node[pos=1.0, right, font=\scriptsize, simGreen] {$\vx_0+\mathcal{K}_1$};
  % K_2 : a wider wedge (a 2D slice drawn as a parallelogram)
  \draw[simBlue, thick, fill=simBgBlue, fill opacity=0.45]
    (x0) -- ++(2.2,0.5) -- ++(0.2,1.7) -- ++(-2.2,-0.5) -- cycle;
  \node[font=\scriptsize, simBlue] at (0.2,1.0) {$\vx_0+\mathcal{K}_2$};
  % the iterates, each the closest point in its slice
  \fill[simGreen] (-0.85,-0.78) circle (1.4pt);
  \node[font=\scriptsize, simGreen] at (-0.85,-0.78) [below=1pt] {$\vx_1$};
  \fill[simRust] (-0.2,0.35) circle (1.6pt);
  \node[font=\scriptsize, simRust] at (-0.2,0.35) [right=1pt] {$\vx_2$};
  % the target (in A-image, schematically): the "true solution" direction
  \node[draw=simGold, fill=simBgGold, thick, rounded corners=2pt,
        font=\scriptsize, inner sep=2pt] (sol) at (2.9,1.9) {$\vx_\star=\mat{A}^{-1}\vb$};
  % residual arrows: from each iterate "toward" the solution, shrinking
  \draw[backflow, shorten >=2pt] (-0.85,-0.78) -- (sol)
    node[pos=0.45, above, font=\scriptsize, simRust, sloped] {$\norm{\vr_1}$};
  \draw[flow, simRust, shorten >=2pt] (-0.2,0.35) -- (sol)
    node[pos=0.5, below, font=\scriptsize, simRust, sloped] {$\norm{\vr_2}<\norm{\vr_1}$};
\end{tikzpicture}
\caption[The minimum-residual idea]{The minimum-residual principle. At each step
the feasible set is the affine Krylov subspace $\vx_0+\mathcal{K}_m$---a flat slice
of $\Real^N$ that gains one dimension per step (here $\mathcal{K}_1\subset
\mathcal{K}_2$). GMRES returns the point $\vx_m$ of that slice whose residual
$\norm{\vb-\mat{A}\vx_m}$ is smallest. Because each slice contains the previous one,
the achievable minimum can only shrink: $\norm{\vr_2}\le\norm{\vr_1}\le\norm{\vr_0}$.
The residual decreases monotonically by construction---GMRES never moves backward.}
\label{fig:minres}
\end{figure}

There is one immediate problem, and the rest of the chapter is its solution. The
minimization is posed over a subspace of $\Real^N$, and $N$ is enormous---millions
of unknowns for a real mesh. We cannot solve an $N$-dimensional least-squares
problem at every step. The genius of GMRES is that we never have to: by choosing
the \emph{right basis} for $\mathcal{K}_m$, the giant minimization collapses to a
tiny one of size $m+1$, where $m$ is the step count, a few dozen at most. That
right basis is built by the Arnoldi process.

\section{Arnoldi: an orthonormal basis for the Krylov subspace}

The raw Krylov vectors $\vr_0,\mat{A}\vr_0,\mat{A}^2\vr_0,\dots$ span
$\mathcal{K}_m$ but are a terrible basis: they line up more and more with the
dominant eigenvector of $\mat{A}$, becoming numerically parallel within a handful
of steps. We need an \emph{orthonormal} basis instead, and we need to build it one
vector at a time, because we discover $\mathcal{K}_m$ one matvec at a time. The
Arnoldi process does exactly this: it is Gram--Schmidt orthogonalization (the full
treatment is \cref{ch:orthog}) applied to the Krylov sequence as it is generated.

\subsection{The process}

Start with the normalized initial residual $\vv_1=\vr_0/\beta$. At each step, take
the most recently produced basis vector $\vv_j$, hit it with $\mat{A}$ to push into
the next Krylov dimension, then \emph{orthogonalize} the result against every basis
vector built so far and normalize what remains:

\begin{lstlisting}[style=l4style]
-- one Arnoldi step: extend an orthonormal Krylov basis by one vector
arnoldi_step :: LinOp -> Vec[] -> Int -> (Vec, Float[], Float)
arnoldi_step A V j =
  let w0 = matvec A (V[j])           -- push into the next Krylov dimension
  let (w, h) = orthogonalize w0 V    -- subtract off components along V[0..j]
  --   h[k] = <w0, V[k]>  are the Arnoldi coefficients (the Hessenberg column)
  let hjj1 = norm w                  -- the new sub-diagonal entry  ||w||
  let v_next = w `scale` (1 / hjj1)  -- normalize: the new basis vector V[j+1]
  in (v_next, h, hjj1)
\end{lstlisting}

\noindent The coefficients $h_{kj}=\inner{\mat{A}\vv_j}{\vv_k}$ for $k\le j$ record
how much of $\mat{A}\vv_j$ pointed along each existing basis vector; the
sub-diagonal entry $h_{j+1,j}=\norm{\vw}$ records the size of the genuinely new
direction. Collect these numbers and a remarkable structure appears.

\begin{figure}[t]
\centering
\begin{tikzpicture}[scale=0.92]
  % ---- LEFT: the tall-skinny basis V_{m+1} ----
  \begin{scope}
    \node[font=\bfseries\small, simInk] at (0,3.05) {The basis $\mat{V}_{m+1}$};
    \draw[simBlue, thick, fill=simBgBlue] (-1.0,-2.6) rectangle (1.0,2.6);
    % columns
    \foreach \x in {-0.6,-0.2,0.2,0.6}{
      \draw[simBlue!55] (\x,-2.6) -- (\x,2.6);}
    \node[font=\scriptsize, simGray, rotate=90] at (0,0) {$N$ rows (huge)};
    \node[font=\scriptsize, simBlue] at (0,2.85) {$m{+}1$ columns};
    \node[font=\scriptsize\itshape, simGray, align=center] at (0,-3.0)
      {orthonormal\\columns $\vv_1\dots\vv_{m+1}$};
  \end{scope}
  % ---- the Arnoldi relation glyph ----
  \node[font=\Large] at (2.0,0) {$\mat{A}\,\mat{V}_m=$};
  % ---- RIGHT: V_{m+1} times the small Hessenberg ----
  \begin{scope}[shift={(4.5,0)}]
    % small V_{m+1} again
    \draw[simBlue, thick, fill=simBgBlue] (-0.8,-2.0) rectangle (0.5,2.0);
    \node[font=\scriptsize, simBlue] at (-0.15,2.25) {$\mat{V}_{m+1}$};
    \node[font=\large] at (1.0,0) {$\cdot$};
  \end{scope}
  % ---- the Hessenberg H-bar (small, (m+1) x m) ----
  \begin{scope}[shift={(6.6,0)}]
    \node[font=\bfseries\small, simInk] at (0,3.05) {$\bar{\mat{H}}_m$ (small)};
    \draw[simRust, thick, fill=simBgRust] (-0.95,-1.0) rectangle (0.95,2.0);
    \node[font=\scriptsize, simRust] at (0,2.25) {$m$ cols};
    \node[font=\scriptsize, simRust, rotate=90] at (-1.25,0.5) {$m{+}1$ rows};
    % shade the upper-Hessenberg pattern: filled on/above first subdiagonal
    \fill[simRust!30] (-0.95,2.0) -- (0.95,2.0) -- (0.95,-1.0) -- (-0.55,-1.0)
      -- (-0.95,-0.4) -- cycle;
    % subdiagonal staircase
    \draw[simRust!75, very thick] (-0.95,-0.4) -- (-0.55,-1.0);
    \node[font=\scriptsize\itshape, simGray, align=center] at (0,-1.7)
      {upper-Hessenberg:\\one nonzero\\sub-diagonal};
  \end{scope}
\end{tikzpicture}
\caption[Arnoldi builds a tall basis and a small Hessenberg]{The Arnoldi process
produces two objects of utterly different size. \emph{Left:} the orthonormal basis
$\mat{V}_{m+1}$, an $N\times(m{+}1)$ matrix---tall and skinny, $N$ in the millions
but only $m{+}1$ columns. \emph{Right:} the Arnoldi relation
$\mat{A}\mat{V}_m=\mat{V}_{m+1}\bar{\mat{H}}_m$ expresses the action of $\mat{A}$ on
the basis entirely through the \emph{small} $(m{+}1)\times m$ upper-Hessenberg
matrix $\bar{\mat{H}}_m$, whose only nonzeros are on and above the first
sub-diagonal. The huge operator $\mat{A}$ is, within the Krylov subspace, captured
by a matrix a few dozen entries on a side. That compression is what makes the
minimization tractable.}
\label{fig:arnoldi}
\end{figure}

\subsection{The Arnoldi relation}

Read the orthogonalization step backward. We computed $\mat{A}\vv_j$, then wrote it
as its components along the existing basis plus a new normalized direction:
\[
  \mat{A}\vv_j \;=\; \sum_{k=1}^{j} h_{kj}\,\vv_k \;+\; h_{j+1,j}\,\vv_{j+1}.
\]
This says $\mat{A}\vv_j$ is a combination of $\vv_1,\dots,\vv_{j+1}$ only---the
matvec never reaches past the next basis vector. Stacking these identities for
$j=1,\dots,m$ into matrix form gives the \emph{Arnoldi relation}, the single most
important equation in the chapter:

\begin{theorem}[Arnoldi relation]
\label{thm:arnoldi}
Let $\mat{V}_m=[\vv_1\ \cdots\ \vv_m]$ be the $N\times m$ matrix of Arnoldi basis
vectors and $\mat{V}_{m+1}=[\mat{V}_m\ \vv_{m+1}]$ the same with one extra column.
Let $\bar{\mat{H}}_m$ be the $(m{+}1)\times m$ matrix whose entry in row $k$,
column $j$ is the Arnoldi coefficient $h_{kj}$ (and zero where $k>j+1$). Then
\[
  \mat{A}\,\mat{V}_m \;=\; \mat{V}_{m+1}\,\bar{\mat{H}}_m,
\]
and $\bar{\mat{H}}_m$ is \emph{upper-Hessenberg}: $h_{kj}=0$ whenever $k>j+1$, so
its only sub-diagonal is the first one.
\end{theorem}

The Hessenberg structure is not a coincidence; it \emph{is} the statement that
$\mat{A}\vv_j$ reaches only as far as $\vv_{j+1}$. (When $\mat{A}$ is symmetric,
$\bar{\mat{H}}_m$ becomes tridiagonal---one sub-diagonal \emph{and} one
super-diagonal---and the orthogonalization collapses to the three-term Lanczos
recurrence of \cref{ch:cg}. Hessenberg is the nonsymmetric generalization;
tridiagonal is the symmetric special case.) \Cref{fig:arnoldi} contrasts the two
sizes: $\mat{V}_{m+1}$ is enormous in one dimension, but $\bar{\mat{H}}_m$ is tiny
in both.

\subsection{From a big minimization to a small one}

Now we cash in. Any candidate iterate in the affine subspace can be written
$\vx=\vx_0+\mat{V}_m\vy$ for some coordinate vector $\vy\in\Real^m$---the columns of
$\mat{V}_m$ span $\mathcal{K}_m$, so $\mat{V}_m\vy$ ranges over all of it as $\vy$
ranges over $\Real^m$. Substitute and compute the residual, using $\vr_0=\beta\vv_1$
and the Arnoldi relation:
\begin{align*}
  \vb-\mat{A}\vx
    &= \vr_0 - \mat{A}\mat{V}_m\vy
     = \beta\vv_1 - \mat{V}_{m+1}\bar{\mat{H}}_m\vy \\
    &= \mat{V}_{m+1}\bigl(\beta\vec e_1 - \bar{\mat{H}}_m\vy\bigr),
\end{align*}
where the last step pulls out $\mat{V}_{m+1}$ and uses $\vv_1=\mat{V}_{m+1}\vec e_1$.
Because the columns of $\mat{V}_{m+1}$ are orthonormal, multiplying by
$\mat{V}_{m+1}$ preserves the Euclidean norm. Therefore
\[
  \norm{\vb-\mat{A}\vx}
    \;=\; \norm{\mat{V}_{m+1}\bigl(\beta\vec e_1-\bar{\mat{H}}_m\vy\bigr)}
    \;=\; \norm{\beta\vec e_1-\bar{\mat{H}}_m\vy}.
\]

\begin{keyidea}
The minimum-residual problem in $\Real^N$ collapses to a least-squares problem in
$\Real^{m+1}$:
\[
  \min_{\vx\,\in\,\vx_0+\mathcal{K}_m}\norm{\vb-\mat{A}\vx}
  \;=\;
  \min_{\vy\,\in\,\Real^m}\ \norm{\beta\,\vec e_1-\bar{\mat{H}}_m\,\vy}.
\]
The orthonormality of the Arnoldi basis turned a minimization over a subspace of
the \emph{huge} space into one over $\vy$ of dimension $m$---a problem whose data is
the \emph{tiny} $(m{+}1)\times m$ Hessenberg matrix $\bar{\mat{H}}_m$ and the scalar
$\beta$. We solve this small problem for $\vy_m$, and only then form the actual
iterate $\vx_m=\vx_0+\mat{V}_m\vy_m$ in the big space---once, at the end.
\end{keyidea}

\Cref{fig:reduction} draws the reduction as a funnel: the $N$-dimensional residual
minimization on the left, the orthonormal-basis change in the middle, the
$(m{+}1)$-dimensional least-squares problem on the right. Everything expensive about
$\mat{A}$ and $N$ has been squeezed into building $\mat{V}_m$ and $\bar{\mat{H}}_m$;
the actual optimization happens in a space small enough to hold in a cache line.

\begin{figure}[t]
\centering
\begin{tikzpicture}[node distance=8mm]
  % BIG problem (left)
  \node[draw=simBlue, fill=simBgBlue, thick, rounded corners=2pt,
        minimum width=34mm, minimum height=20mm, align=center, font=\footnotesize]
        (big) at (0,0)
    {\textbf{big minimization}\\[2pt]$\displaystyle\min_{\vx\in\vx_0+\mathcal{K}_m}$\\[1pt]
     $\norm{\vb-\mat{A}\vx}$\\[2pt]\itshape\scriptsize over a subspace of $\Real^N$};
  % the lever: orthonormal basis + Arnoldi relation
  \node[stage, minimum width=30mm, minimum height=16mm, right=12mm of big, align=center]
        (lever)
    {\textbf{Arnoldi}\\[1pt]\footnotesize $\vx=\vx_0+\mat{V}_m\vy$\\\footnotesize
     $\mat{A}\mat{V}_m=\mat{V}_{m+1}\bar{\mat{H}}_m$\\\footnotesize
     $\norm{\mat{V}_{m+1}\vu}=\norm{\vu}$};
  \draw[flow] (big) -- node[above,font=\scriptsize]{change of} node[below,font=\scriptsize]{basis} (lever);
  % SMALL problem (right)
  \node[draw=simRust, fill=simBgRust, thick, rounded corners=2pt,
        minimum width=34mm, minimum height=20mm, align=center, font=\footnotesize]
        (small) at ($(lever.east)+(2.0,0)$)
    {\textbf{small least-squares}\\[2pt]$\displaystyle\min_{\vy\in\Real^m}$\\[1pt]
     $\norm{\beta\vec e_1-\bar{\mat{H}}_m\vy}$\\[2pt]\itshape\scriptsize
     dimension $m{+}1$ (tiny)};
  \draw[flow] (lever) -- (small);
  % annotation
  \node[font=\scriptsize\itshape, simGray, align=center, below=4mm of lever, text width=34mm]
    {the same minimum value,\\two sizes of problem};
\end{tikzpicture}
\caption[The big-to-small reduction]{The engine of GMRES. The minimum-residual
problem on the left lives in $\Real^N$ with $N$ in the millions. Writing every
candidate as $\vx_0+\mat{V}_m\vy$ and invoking the Arnoldi relation
(\cref{thm:arnoldi}) plus the norm-preservation of the orthonormal basis
$\mat{V}_{m+1}$ turns it into the least-squares problem on the right, whose data is
just the $(m{+}1)\times m$ Hessenberg matrix and the scalar $\beta$. Both problems
have \emph{exactly the same minimum value}; only the second is one we can afford to
solve at every step.}
\label{fig:reduction}
\end{figure}

\begin{workedexample}{Two Arnoldi steps on a small nonsymmetric operator}{arnoldi}
Take the nonsymmetric $3\times3$ operator and right-hand side
\[
  \mat{A}=\begin{bmatrix} 2 & 1 & 0\\ 0 & 3 & 1\\ 1 & 0 & 4\end{bmatrix},
  \qquad
  \vb=\begin{bmatrix}1\\0\\0\end{bmatrix},
  \qquad \vx_0=\mathbf{0}.
\]
$\mat{A}$ is not symmetric ($a_{31}=1\ne0=a_{13}$), so CG does not apply. We run two
Arnoldi steps by hand.

\textbf{Setup.} The initial residual is $\vr_0=\vb-\mat{A}\vx_0=\vb=[1,0,0]^{\top}$,
so $\beta=\norm{\vr_0}=1$ and the first basis vector is
$\vv_1=\vr_0/\beta=[1,0,0]^{\top}$.

\textbf{Step 1.} Push and orthogonalize. The matvec is
\[
  \mat{A}\vv_1=\begin{bmatrix}2\\0\\1\end{bmatrix}.
\]
Project onto $\vv_1$: the coefficient is $h_{11}=\inner{\mat{A}\vv_1}{\vv_1}=2$.
Subtract it off:
\[
  \vw=\mat{A}\vv_1-h_{11}\vv_1=\begin{bmatrix}2\\0\\1\end{bmatrix}
     -2\begin{bmatrix}1\\0\\0\end{bmatrix}=\begin{bmatrix}0\\0\\1\end{bmatrix}.
\]
The new sub-diagonal entry is $h_{21}=\norm{\vw}=1$, and the new basis vector is
$\vv_2=\vw/h_{21}=[0,0,1]^{\top}$. (It is orthonormal to $\vv_1$, as it must be.)

\textbf{Step 2.} Push and orthogonalize again. The matvec is
\[
  \mat{A}\vv_2=\begin{bmatrix}0\\1\\4\end{bmatrix}.
\]
Project onto both existing basis vectors:
$h_{12}=\inner{\mat{A}\vv_2}{\vv_1}=0$ and
$h_{22}=\inner{\mat{A}\vv_2}{\vv_2}=4$. Subtract both:
\[
  \vw=\mat{A}\vv_2-h_{12}\vv_1-h_{22}\vv_2
     =\begin{bmatrix}0\\1\\4\end{bmatrix}-0-4\begin{bmatrix}0\\0\\1\end{bmatrix}
     =\begin{bmatrix}0\\1\\0\end{bmatrix}.
\]
The sub-diagonal entry is $h_{32}=\norm{\vw}=1$, and $\vv_3=[0,1,0]^{\top}$.

\textbf{Read off the Hessenberg.} Assembling the coefficients into the
$3\times2$ upper-Hessenberg matrix:
\[
  \mat{V}_2=\begin{bmatrix}1&0\\0&0\\0&1\end{bmatrix},
  \qquad
  \bar{\mat{H}}_2=\begin{bmatrix} h_{11}&h_{12}\\ h_{21}&h_{22}\\ 0&h_{32}\end{bmatrix}
   =\begin{bmatrix} 2 & 0\\ 1 & 4\\ 0 & 1\end{bmatrix}.
\]
Notice the shape: $\bar{\mat{H}}_2$ is $(m{+}1)\times m=3\times2$, with a single
nonzero sub-diagonal (the $1$'s in positions $(2,1)$ and $(3,2)$)---exactly the
upper-Hessenberg pattern of \cref{thm:arnoldi}. We have compressed the action of the
$3\times3$ operator on $\mathcal{K}_2$ into these six numbers. We solve the
least-squares problem they pose in \cref{wex:leastsq}.
\end{workedexample}

\section{The small least-squares problem and the free residual}

We are left with $\min_{\vy}\norm{\beta\vec e_1-\bar{\mat{H}}_m\vy}$, an
overdetermined least-squares problem ($m{+}1$ equations, $m$ unknowns). The
textbook solution is a QR factorization of $\bar{\mat{H}}_m$: write
$\bar{\mat{H}}_m=\mat{Q}\mat{R}$ with $\mat{Q}$ orthonormal and $\mat{R}$ upper
triangular, then the minimizer solves $\mat{R}\vy=\mat{Q}^{\top}(\beta\vec e_1)$ by
back-substitution. But re-factorizing from scratch at every step would waste the
fact that $\bar{\mat{H}}_m$ grows by exactly one column per step. We can do far
better, and the better way hands us the residual norm for free.

\subsection{The running QR by Givens rotations}

A Hessenberg matrix is already \emph{almost} triangular: its only obstacle to upper
triangularity is the single sub-diagonal. Each new column, when it arrives, adds one
sub-diagonal entry that must be annihilated. A \emph{Givens rotation}---a $2\times2$
plane rotation acting on two adjacent rows---annihilates exactly one entry while
leaving the triangular structure above it intact. So we maintain the QR factorization
\emph{incrementally}: when column $j$ arrives, we replay the rotations generated for
columns $1,\dots,j-1$ onto it (bringing it into the established triangular frame),
then generate one new rotation to kill its fresh sub-diagonal entry. The mechanics of
the rotations---how each $2\times2$ rotation is computed and applied---are the
business of \cref{ch:givens}; here we use only the result.

\begin{figure}[t]
\centering
\begin{tikzpicture}[scale=0.95]
  % the growing triangular factor R, with one new column arriving
  \node[font=\bfseries\small, simInk] at (0,2.55) {Running QR: one new column per step};
  % already-triangular block
  \fill[simBlue!18] (-2.4,-1.4) -- (0.2,-1.4) -- (0.2,1.6) -- (-2.4,1.6) -- cycle;
  \draw[simBlue, thick] (-2.4,-1.4) rectangle (0.2,1.6);
  % triangular staircase inside
  \draw[simBlue!70, thick] (-2.4,1.6) -- (0.2,1.6) -- (0.2,-1.0) -- (-2.0,1.2) -- cycle;
  \node[font=\scriptsize, simBlue] at (-1.3,0.5) {already};
  \node[font=\scriptsize, simBlue] at (-1.3,0.1) {triangular};
  % the new column (highlighted)
  \draw[simRust, very thick, fill=simBgRust] (0.2,-1.4) rectangle (0.55,1.6);
  \node[font=\scriptsize, simRust, rotate=90] at (0.38,0.1) {new column $j$};
  % the offending subdiagonal entry
  \fill[simRust] (0.38,-1.15) circle (2pt);
  \node[font=\scriptsize, simRust, align=left] at (1.9,-1.15)
    {one sub-diagonal\\entry to annihilate};
  \draw[refedge] (1.0,-1.15) -- (0.5,-1.15);
  % the rotation
  \draw[-{Stealth[length=2mm]}, simGreen, thick] (1.1,0.6) arc (-40:40:0.5);
  \node[font=\scriptsize, simGreen, align=left] at (2.2,0.6)
    {generate one\\Givens rotation};
  % the RHS column on the far right, with the residual entry glowing
  \draw[simGold, thick, fill=simBgGold] (3.6,-1.4) rectangle (3.95,1.6);
  \node[font=\scriptsize, simGold!70!black, rotate=90] at (3.78,0.55) {rotated RHS $s$};
  \fill[simGold!80!black] (3.78,-1.15) circle (2.4pt);
  \node[font=\scriptsize\bfseries, simGold!60!black, align=left] at (5.0,-1.15)
    {$|s_{m+1}|=\norm{\vr_m}$\\\itshape\mdseries the residual,\\\itshape\mdseries for free};
  \draw[refedge] (4.2,-1.15) -- (3.9,-1.15);
\end{tikzpicture}
\caption[The running QR and the free residual readout]{The incremental QR. Each
Arnoldi step appends one column to the Hessenberg matrix; the already-computed
Givens rotations are replayed onto it, then one new rotation annihilates its single
fresh sub-diagonal entry, extending the triangular factor by one column. The same
new rotation is applied to the rotated right-hand side $s$ (initialized to
$\beta\vec e_1$): because every rotation is an isometry, the magnitude of its last
entry, $|s_{m+1}|$, equals the least-squares residual norm $\norm{\vr_m}$
\emph{exactly}---available at every step \emph{without} forming $\vx_m$ or computing
a single matvec. The rotation mechanics are developed in \cref{ch:givens}.}
\label{fig:runningqr}
\end{figure}

\subsection{The residual is free}

The payoff is the reason GMRES is built this way. Alongside the Hessenberg matrix we
carry a small right-hand-side vector $s$, initialized to $\beta\vec e_1$, and we
apply each Givens rotation to $s$ as we generate it. Because every rotation is an
isometry (it preserves length), the transformed system has the same residual norm as
the original. And after triangularization, the least-squares residual is precisely
the magnitude of the \emph{last} entry of $s$:
\[
  \norm{\vr_m}\;=\;\min_{\vy}\norm{\beta\vec e_1-\bar{\mat{H}}_m\vy}\;=\;\abs{s_{m+1}}.
\]

\begin{keyidea}
The GMRES residual norm is available \textbf{for free}, every step, as a single
scalar. After the running QR processes column $m$, the magnitude $|s_{m+1}|$ of the
last entry of the rotated right-hand side \emph{equals} $\norm{\vr_m}$
exactly---with no matvec, no explicit residual vector, and without ever forming the
iterate $\vx_m$. The convergence test reads off $|s_{m+1}|$ and compares it to the
tolerance. Only when the test passes (or we restart) do we pay the one-time cost of
back-solving the small triangular system for $\vy_m$ and forming
$\vx_m=\vx_0+\mat{V}_m\vy_m$ in the big space.
\end{keyidea}

This is what the source actually does: Palace's GMRES carries rotation registers
\code{cs} and \code{sn} and a rotated right-hand side \code{s}, applies each new
rotation to \code{s}, and reads the convergence test directly off
$\beta=|s_{j+1}|$---no residual is ever evaluated explicitly. The back-solve that
turns $\vy$ into a correction $\mat{V}_m\vy$ fires only at convergence or at a
restart boundary. We trace the small problem end to end now.

\begin{workedexample}{Solving the $3\times2$ least-squares and reading the residual}{leastsq}
Continue \cref{wex:arnoldi}, with $\beta=1$ and
\[
  \bar{\mat{H}}_2=\begin{bmatrix} 2 & 0\\ 1 & 4\\ 0 & 1\end{bmatrix},
  \qquad
  \beta\vec e_1=\begin{bmatrix}1\\0\\0\end{bmatrix}.
\]
We solve $\min_{\vy}\norm{\beta\vec e_1-\bar{\mat{H}}_2\vy}$ by the running QR,
generating two Givens rotations $\mat{G}_1,\mat{G}_2$ and watching the residual
appear in the rotated right-hand side $s$.

\textbf{Rotation 1} acts on rows $1,2$ to kill the $(2,1)$ entry. With diagonal
$\rho=2$ and sub-diagonal $\sigma=1$, the rotation has
$c_1=\rho/\sqrt{\rho^2+\sigma^2}=2/\sqrt5$ and
$s_1=\sigma/\sqrt{\rho^2+\sigma^2}=1/\sqrt5$. Applying $\mat{G}_1$ to the first
column of $\bar{\mat{H}}_2$ turns $(2,1)$ into $\sqrt{2^2+1^2}=\sqrt5$ on the
diagonal and $0$ below. Applying $\mat{G}_1$ to the second column
$[0,4,1]^{\top}$ mixes rows $1,2$: the new $(1,2)$ entry is
$c_1\!\cdot\!0+s_1\!\cdot\!4=4/\sqrt5$ and the new $(2,2)$ entry is
$-s_1\!\cdot\!0+c_1\!\cdot\!4=8/\sqrt5$. Applying $\mat{G}_1$ to the right-hand side
$s=[1,0,0]^{\top}$ gives
$s\leftarrow[c_1,\,-s_1,\,0]^{\top}=[2/\sqrt5,\,-1/\sqrt5,\,0]^{\top}$. After
rotation 1 the working system is
\[
  \mat{R}^{(1)}=\begin{bmatrix} \sqrt5 & 4/\sqrt5\\ 0 & 8/\sqrt5\\ 0 & 1\end{bmatrix},
  \qquad
  s=\begin{bmatrix} 2/\sqrt5\\ -1/\sqrt5\\ 0\end{bmatrix}.
\]

\textbf{Rotation 2} acts on rows $2,3$ to kill the $(3,2)$ entry. The diagonal is
$\rho=8/\sqrt5$ and the sub-diagonal $\sigma=1$, so
$r=\sqrt{\rho^2+\sigma^2}=\sqrt{64/5+1}=\sqrt{69/5}=\sqrt{69}/\sqrt5$, giving
$c_2=\rho/r=8/\sqrt{69}$ and $s_2=\sigma/r=\sqrt5/\sqrt{69}$. Applying $\mat{G}_2$ to
the right-hand side mixes rows $2,3$: the new entry $2$ is
$c_2\!\cdot\!(-1/\sqrt5)+s_2\!\cdot\!0=-8/\sqrt{5\cdot69}$, and---the entry we want---
the new entry $3$ is
\[
  s_3 \;=\; -s_2\cdot\!\left(-\tfrac{1}{\sqrt5}\right)+c_2\cdot 0
        \;=\; \frac{\sqrt5}{\sqrt{69}}\cdot\frac{1}{\sqrt5}
        \;=\; \frac{1}{\sqrt{69}}.
\]

\textbf{Read the residual for free.} The residual norm after two steps is
\[
  \norm{\vr_2}\;=\;\abs{s_3}\;=\;\frac{1}{\sqrt{69}}\;\approx\;0.1204,
\]
read directly off the last entry of the rotated right-hand side---no matvec, no
explicit residual vector. (The starting residual was $\norm{\vr_0}=\beta=1$, so two
GMRES steps cut it by a factor of about $8.3$.)

\textbf{Recover the iterate, once.} Only now do we back-solve the upper-triangular
$2\times2$ system $\mat{R}\vy=[s_1,s_2]^{\top}$ for the coordinate vector $\vy$. The
upper $2\times2$ block of the rotated Hessenberg is $\mat{R}=\bigl[\begin{smallmatrix}
\sqrt5 & 4/\sqrt5\\ 0 & 8/\sqrt{69}\cdot\sqrt5/\sqrt5\end{smallmatrix}\bigr]$ (with
$(2,2)=r=\sqrt{69}/\sqrt5$ after rotation~2), and the right-hand side is
$[2/\sqrt5,\,-8/\sqrt{5\cdot69}]^{\top}$. Back-substitution gives
$y_2=(-8/\sqrt{5\cdot69})/(\sqrt{69}/\sqrt5)=-8/69$ and
$y_1=(2/\sqrt5-(4/\sqrt5)y_2)/\sqrt5=(2+32/69)/5=170/345=34/69$. The iterate is
\[
  \vx_2=\vx_0+\mat{V}_2\vy=\begin{bmatrix}1&0\\0&0\\0&1\end{bmatrix}
    \begin{bmatrix}34/69\\-8/69\end{bmatrix}=\begin{bmatrix}34/69\\0\\-8/69\end{bmatrix}.
\]
As a check, $\mat{A}\vx_2=[60/69,\,-8/69,\,2/69]^{\top}$ and the true residual
$\vb-\mat{A}\vx_2=[9/69,\,8/69,\,-2/69]^{\top}$ has norm
$\sqrt{81+64+4}/69=\sqrt{149}/69\approx0.1768$. This differs from $|s_3|$ because the
$3\times3$ Krylov space is not yet exhausted; in exact arithmetic GMRES reaches the
exact solution at step $m=3=N$, where the residual hits zero. The point stands: the
\emph{small} problem's residual $|s_3|$ was free, and the iterate was formed only once.
\end{workedexample}

\section{Restart: bounding the cost of remembering everything}

GMRES bought optimality by remembering the entire basis. That memory is not free,
and as $m$ grows the bill comes due on two fronts.

\begin{itemize}[nosep]
  \item \textbf{Storage grows linearly in $m$.} The basis $\mat{V}_m$ holds $m$
  vectors of length $N$, so storage is $O(mN)$. At step 100 of a million-unknown
  solve, that is a hundred million-entry vectors---hundreds of megabytes, just for
  the basis.
  \item \textbf{Work grows quadratically in $m$.} Orthogonalizing the new vector
  against all $m$ existing ones costs $O(mN)$ per step, so the cumulative
  orthogonalization work through step $m$ is $O(m^2N)$. The matvecs are only $O(mN)$;
  it is the \emph{orthogonalization against an ever-growing basis} that dominates.
\end{itemize}

Both costs are unbounded in $m$, which is intolerable for a solve that might need
hundreds of iterations. The standard remedy is to cap $m$ at a fixed restart length
$k$.

\begin{definition}[Restarted GMRES, GMRES($k$)]
\label{def:restart}
Fix a restart length $k$. Run GMRES for at most $k$ steps. If the residual has not
met the tolerance by step $k$, \emph{restart}: discard the basis $\mat{V}_k$ and the
Hessenberg matrix entirely, take the current iterate $\vx_k$ as the new initial guess
$\vx_0$, recompute the residual $\vr_0=\vb-\mat{A}\vx_0$, and begin a fresh GMRES
cycle of up to $k$ steps. Repeat until convergence. Storage is capped at $O(kN)$ and
per-cycle work at $O(k^2N)$, both independent of the total iteration count.
\end{definition}

The basis is \emph{ephemeral}: it is born at the start of a cycle and discarded at
the restart boundary. This is exactly the lifetime story of \cref{ch:strata}---the
Krylov basis is per-solve, indeed per-\emph{cycle}, run-time state, allocated fresh
and thrown away, never part of any operator's frozen internals. Only the iterate
$\vx$ survives across a restart; everything else is rebuilt.

\begin{figure}[t]
\centering
\begin{tikzpicture}
  \begin{axis}[
      width=12.0cm, height=5.4cm,
      xlabel={total iterations},
      ylabel={basis size (vectors stored)},
      xmin=0, xmax=15, ymin=0, ymax=8,
      xtick={0,3,6,9,12,15}, ytick={0,2,4,6},
      tick label style={font=\scriptsize},
      label style={font=\small},
      axis lines=left, clip=false,
    ]
    % full GMRES: linear growth
    \addplot[simBlue, thick, mark=*, mark size=1.2pt] coordinates
      {(0,0)(1,1)(2,2)(3,3)(4,4)(5,5)(6,6)(7,7)};
    \node[font=\scriptsize, simBlue, anchor=west] at (axis cs:6.1,6.9)
      {full GMRES: $O(m)$};
    % restarted GMRES(3): sawtooth
    \addplot[simRust, very thick] coordinates
      {(0,0)(1,1)(2,2)(3,3)};
    \addplot[simRust, very thick] coordinates
      {(3,0)(4,1)(5,2)(6,3)};
    \addplot[simRust, very thick] coordinates
      {(6,0)(7,1)(8,2)(9,3)};
    \addplot[simRust, very thick] coordinates
      {(9,0)(10,1)(11,2)(12,3)};
    \addplot[simRust, very thick] coordinates
      {(12,0)(13,1)(14,2)(15,3)};
    % restart drop markers
    \draw[simRust, densely dashed] (axis cs:3,3) -- (axis cs:3,0);
    \draw[simRust, densely dashed] (axis cs:6,3) -- (axis cs:6,0);
    \draw[simRust, densely dashed] (axis cs:9,3) -- (axis cs:9,0);
    \draw[simRust, densely dashed] (axis cs:12,3) -- (axis cs:12,0);
    \node[font=\scriptsize, simRust, anchor=west] at (axis cs:9.3,3.5)
      {GMRES(3): sawtooth, capped at $k{=}3$};
  \end{axis}
\end{tikzpicture}
\caption[Memory: full GMRES versus the restart sawtooth]{Basis storage over the
course of a solve. Full GMRES (blue) accumulates one vector per step forever, so its
memory climbs linearly without bound. Restarted GMRES($k$) (rust, here $k=3$) caps
the basis at $k$ vectors: every $k$ steps it discards the basis and starts a fresh
cycle from the current iterate, producing the characteristic sawtooth. The peak
memory is bounded by $O(kN)$ no matter how many total iterations the solve takes. The
trade is convergence quality, charted in \cref{fig:convergence}.}
\label{fig:memory}
\end{figure}

\subsection{The trade-off}

Restarting is not free either; it trades memory for convergence speed.

\begin{pitfall}
\textbf{Restarting loses GMRES's optimality, and can stall.} Full GMRES minimizes
the residual over the \emph{entire} accumulated Krylov subspace---its convergence is
often superlinear, accelerating as the subspace captures more of the operator's
spectrum. The instant you restart, that accumulated subspace is thrown away: the new
cycle minimizes only over a \emph{fresh} $k$-dimensional space built from the current
residual, blind to everything the discarded basis had learned. The monotonic
\emph{decrease} of the residual still holds across restarts (each cycle starts from a
no-worse iterate), but the \emph{rate} can collapse---and for some operators
GMRES($k$) \emph{stagnates}, making negligible progress per cycle while full GMRES
would have converged. Choosing $k$ is a real engineering decision: too small and you
stall; too large and you run out of memory. There is no value of $k$ that is right
for every problem.
\end{pitfall}

\section{FGMRES: when the preconditioner changes every step}

We now pay off the promise of \cref{ch:operators}. There we met the
\emph{decision rule}: a variant bound once at construction is absorbed by a
constructed operator; a variant that changes \emph{per step} cannot be, and must
instead be threaded through the run-time state. We claimed that this single rule is
``exactly why \code{FGMRES} is a different algorithm from \code{GMRES}.'' Here is the
claim, made good.

\subsection{Where preconditioned GMRES applies $\mat{M}$}

A preconditioner $\mat{M}$ is an approximate inverse of $\mat{A}$ whose job is to
cluster the spectrum so the Krylov iteration converges faster (\cref{ch:precond}).
\emph{Right}-preconditioned GMRES solves $\mat{A}\mat{M}^{-1}\vu=\vb$ with
$\vx=\mat{M}^{-1}\vu$; concretely, each Arnoldi step preconditions the current basis
vector before the matvec: $\vz_j=\mat{M}^{-1}\vv_j$, then $\vw=\mat{A}\vz_j$. If
$\mat{M}$ is \emph{fixed}---chosen once at solve start and held constant---this is a
textbook constructed operator (\cref{wex:jacobi} in spirit): build it once, apply the
same \code{pc} at every step. The preconditioner is build-time configuration; it never
enters the threaded state. The left branch of the decision rule
(\cref{fig:decisionrule} of \cref{ch:operators}) handles it cleanly.

\subsection{What breaks when $\mat{M}$ varies}

Now allow the preconditioner to differ at each step: $\mat{M}_j$ at step $j$. This is
not exotic---it is exactly what happens when the preconditioner is \emph{itself an
inner iterative solve} (a few steps of an inner GMRES, a multigrid cycle run to a
loose, varying tolerance, a smoother whose accuracy adapts to the residual). Each
inner solve produces a slightly different linear operator, so $\vz_j=\mat{M}_j^{-1}
\vv_j$ uses a \emph{different} $\mat{M}_j$ for every $j$. There is no single $\mat{M}$
to construct.

Trace what fails. In right-preconditioned GMRES the iterate is recovered, at the end,
as $\vx=\vx_0+\mat{M}^{-1}\mat{V}_m\vy$---the basis $\mat{V}_m$ lives in the
\emph{preconditioned} space, and we map back by applying $\mat{M}^{-1}$ to the
correction. This reconstruction \emph{assumes a single fixed $\mat{M}$}: it applies
one $\mat{M}^{-1}$ to the whole combination $\mat{V}_m\vy$. But with per-step
$\mat{M}_j$, the $j$-th basis vector was preconditioned by $\mat{M}_j$, the next by
$\mat{M}_{j+1}$, and there is no single $\mat{M}^{-1}$ that undoes the mixture. The
operator that produced $\vz_j$ is \emph{gone} by the next step---it was an ephemeral
inner solve, not a stored value---so $\vz_j$ cannot be recomputed on demand. Plain
GMRES, which never stored the $\vz_j$, simply cannot form the right answer.

\begin{pitfall}
\textbf{Driving plain GMRES with a per-step preconditioner silently corrupts the
solve.} If you feed a varying $\mat{M}_j$ to an algorithm that recovers its iterate
as $\vx_0+\mat{M}^{-1}\mat{V}_m\vy$ under the assumption of a single fixed $\mat{M}$,
the code \emph{runs}---no crash, no error---and even reports a small ``residual''
$|s_{m+1}|$ from its running QR. But that residual is the residual of a least-squares
problem the iterate no longer corresponds to: the reconstruction used the wrong
inverse. The solve converges to the wrong answer, quietly. This is the silent partial
absorption of \cref{ch:operators} in its most dangerous form---a tidy equation hiding
an operation chain that has secretly forked. The repair is not a patch to GMRES; it is
a different algorithm that \emph{remembers} the preconditioned vectors. That algorithm
is FGMRES.
\end{pitfall}

\subsection{The repair: keep a second basis}

FGMRES---flexible GMRES---makes the only move the decision rule allows: it routes the
per-step variant into the threaded state by \emph{remembering} every preconditioned
vector. Alongside the Arnoldi basis $\mat{V}_m$ (the orthonormalized, Arnoldi-built
vectors), FGMRES carries a \emph{second} basis $\mat{Z}_m=[\vz_1\ \cdots\ \vz_m]$
holding the preconditioned vectors $\vz_j=\mat{M}_j^{-1}\vv_j$ exactly as they were
produced, each by its own ephemeral $\mat{M}_j$. The Arnoldi process and the running
QR are \emph{identical} to GMRES---byte for byte, the same orthogonalization, the same
Givens stream, the same free residual. Only two things change: at each step we save
$\vz_j$ into $\mat{Z}_m$, and at the end we reconstruct the iterate against
$\mat{Z}_m$ instead of against $\mat{V}_m$.

\begin{keyidea}
\textbf{The $\mat{Z}$-versus-$\mat{V}$ distinction.} GMRES recovers its iterate from
the Arnoldi basis, $\vx_m=\vx_0+\mat{V}_m\vy_m$. FGMRES recovers it from the stored
\emph{preconditioned} basis,
\[
  \vx_m \;=\; \vx_0 + \mat{Z}_m\,\vy_m,
  \qquad \mat{Z}_m=[\mat{M}_1^{-1}\vv_1\ \cdots\ \mat{M}_m^{-1}\vv_m].
\]
Because each $\vz_j=\mat{M}_j^{-1}\vv_j$ already carries the action of its own
per-step preconditioner, the reconstruction $\mat{Z}_m\vy_m$ never needs to apply a
single ``$\mat{M}^{-1}$'' to the combination---the per-step inverses are baked into the
columns of $\mat{Z}$. The flexibility is paid for by exactly one extra length-$N$ basis
in the threaded state. This is the ``\code{FgmresState} adds a $\mat{Z}$ field''
prediction of \cref{wex:flexible} (\cref{ch:operators}), realized.
\end{keyidea}

\begin{figure}[t]
\centering
\begin{tikzpicture}[scale=0.95, node distance=6mm]
  % --- the per-step pipeline, drawn once ---
  \node[font=\bfseries\small, simInk] at (0,2.9) {FGMRES threads two bases};
  % v_j
  \node[draw=simBlue, fill=simBgBlue, thick, rounded corners=2pt,
        minimum width=15mm, minimum height=8mm, font=\footnotesize] (vj) at (-4.4,1.4)
    {$\vv_j$};
  % the per-step preconditioner (ephemeral, dashed)
  \node[extern, minimum width=20mm, minimum height=9mm, right=8mm of vj] (Mj)
    {$\mat{M}_j^{-1}$ \scriptsize(per step)};
  % z_j
  \node[draw=simRust, fill=simBgRust, thick, rounded corners=2pt,
        minimum width=15mm, minimum height=8mm, font=\footnotesize, right=8mm of Mj] (zj)
    {$\vz_j$};
  % A
  \node[draw=simTeal, fill=simBgGray, thick, rounded corners=2pt,
        minimum width=12mm, minimum height=8mm, font=\footnotesize, right=8mm of zj] (A)
    {$\mat{A}$};
  % w
  \node[draw=simBlue, fill=simBgBlue, thick, rounded corners=2pt,
        minimum width=15mm, minimum height=8mm, font=\footnotesize, right=8mm of A] (w)
    {$\vw\!\to\!\vv_{j+1}$};
  \draw[flow] (vj) -- (Mj);
  \draw[flow] (Mj) -- (zj);
  \draw[flow] (zj) -- (A);
  \draw[flow] (A) -- node[above,font=\scriptsize]{matvec} (w);
  % the two bases being accumulated
  \node[draw=simBlue, fill=simBgBlue, semithick, rounded corners=2pt,
        minimum width=70mm, minimum height=7mm, font=\scriptsize, align=center] (Vbasis)
    at (0,-0.4) {Arnoldi basis $\mat{V}_m=[\vv_1\,\cdots\,\vv_m]$
                 \quad\itshape (orthonormal, built by Arnoldi)};
  \node[draw=simRust, fill=simBgRust, semithick, rounded corners=2pt,
        minimum width=70mm, minimum height=7mm, font=\scriptsize, align=center] (Zbasis)
    at (0,-1.5) {preconditioned basis $\mat{Z}_m=[\vz_1\,\cdots\,\vz_m]$
                 \quad\itshape (the per-step $\mat{M}_j^{-1}\vv_j$, remembered)};
  \draw[refedge] (vj.south) to[out=-90,in=90] (Vbasis.north -| -3.3,-0.4);
  \draw[refedge] (zj.south) to[out=-90,in=90] (Zbasis.north -| 0.5,-1.5);
  % the reconstruction
  \node[font=\scriptsize, simRust, align=center] at (0,-2.5)
    {iterate recovered from $\mat{Z}$, not $\mat{V}$:\quad
     $\vx_m=\vx_0+\mat{Z}_m\vy_m$};
\end{tikzpicture}
\caption[FGMRES's two bases]{Why FGMRES carries a second basis. Each step
preconditions the Arnoldi vector $\vv_j$ by a \emph{per-step} operator $\mat{M}_j^{-1}$
(dashed: ephemeral, gone by the next step---perhaps an inner iterative solve), giving
$\vz_j$, which is then hit with $\mat{A}$ to continue the Arnoldi process. The
algorithm accumulates \emph{both} bases: $\mat{V}_m$ (the orthonormal Arnoldi vectors,
blue) drives orthogonalization and the running QR exactly as in plain GMRES; but the
iterate is reconstructed from $\mat{Z}_m$ (the remembered preconditioned vectors,
rust), because each $\vz_j$ already carries its own $\mat{M}_j^{-1}$ and no single
$\mat{M}^{-1}$ could undo the mixture. The per-step variant has become a field of the
threaded state---precisely the right branch of the decision rule
(\cref{fig:decisionrule}, \cref{ch:operators}).}
\label{fig:fgmres}
\end{figure}

\Cref{fig:fgmres} draws the doubled basis. The structural reading is exactly the one
\cref{ch:operators} set up: the variant ``which preconditioner this step'' has a
per-step lifetime, so by the decision rule it cannot live in a constructed operator;
it lives in the carry, as the extra basis $\mat{Z}$. FGMRES is GMRES with one field
added to its state schema---and that field is forced, not chosen, by the lifetime of
the varying preconditioner. The least-squares machinery does not even notice: in the
source, the FGMRES running-QR stream is line-for-line identical to the GMRES one; the
\emph{only} difference is the back-solve reconstructs the correction against $\mat{Z}$
rather than $\mat{V}$.

\begin{workedexample}{GMRES versus restarted GMRES(2): convergence on one system}{restart}
We compare full GMRES with restarted GMRES(2) on a $4\times4$ nonsymmetric system,
watching the residual norms (read for free from $|s_{m+1}|$) at each step. The system
is mildly nonnormal, so restarting visibly hurts.

\textbf{Full GMRES.} Storing the whole basis, full GMRES drives the residual down
monotonically and, on a $4\times4$ system, reaches the exact solution by step $4$ (in
exact arithmetic, GMRES converges in at most $N$ steps because $\mathcal{K}_N$ spans
the whole space). A representative residual history:
\[
  \norm{\vr_0}=1,\quad \norm{\vr_1}\approx0.42,\quad \norm{\vr_2}\approx0.15,\quad
  \norm{\vr_3}\approx0.03,\quad \norm{\vr_4}\approx10^{-15}.
\]

\textbf{Restarted GMRES(2).} Capped at $k=2$, GMRES(2) discards its basis after every
second step. The first cycle matches full GMRES through step $2$
($\norm{\vr_2}\approx0.15$), but then it \emph{restarts}: the rich 2-dimensional
subspace is thrown away, and cycle 2 rebuilds a fresh 2-dimensional space from the
step-2 residual---blind to what cycle 1 learned. Progress per cycle slows:
\[
  0.15 \;\to\; 0.066 \;\to\; 0.030 \;\to\; 0.014 \;\to\; 0.0065 \;\to\;\cdots
\]
roughly halving each cycle---a \emph{linear} rate, where full GMRES was racing to
machine precision. \Cref{fig:convergence} plots both. GMRES(2) still converges (the
residual is monotone), but it takes many more iterations to reach the same tolerance:
the price of bounding memory at $k=2$ rather than letting the basis grow to $N=4$.

\textbf{The reading.} Full GMRES exhausts the Krylov space and hits the exact answer
at step $N$; restarted GMRES(2) trades that finite-termination guarantee, and the
superlinear tail, for a flat $O(kN)$ memory ceiling. On this small system the gap is
stark; on a million-unknown system where full GMRES is impossible, GMRES($k$)'s slow
linear convergence is simply the cost of being able to run at all.
\end{workedexample}

\begin{figure}[t]
\centering
\begin{tikzpicture}
  \begin{axis}[
      width=12.2cm, height=6.2cm,
      xlabel={iteration $m$}, ylabel={residual norm $\norm{\vr_m}$},
      ymode=log, ymin=1e-16, ymax=2, xmin=0, xmax=10,
      xtick={0,2,4,6,8,10},
      ytick={1e0,1e-3,1e-6,1e-9,1e-12,1e-15},
      tick label style={font=\scriptsize}, label style={font=\small},
      legend style={font=\scriptsize, at={(0.98,0.98)}, anchor=north east},
      grid=both, grid style={simGray!18},
      axis lines=left, clip=false,
    ]
    % full GMRES: superlinear, hits machine precision at m=4
    \addplot[simBlue, very thick, mark=*, mark size=1.4pt] coordinates
      {(0,1)(1,0.42)(2,0.15)(3,0.03)(4,1e-15)};
    \addlegendentry{full GMRES}
    % restarted GMRES(2): linear, restarts every 2 steps (dashed verticals)
    \addplot[simRust, very thick, mark=square*, mark size=1.3pt] coordinates
      {(0,1)(1,0.42)(2,0.15)(3,0.066)(4,0.030)(5,0.014)(6,0.0065)(7,0.0030)(8,0.0014)(9,0.00065)(10,0.00030)};
    \addlegendentry{restarted GMRES(2)}
    % restart boundaries
    \draw[simRust, densely dotted] (axis cs:2,1e-16) -- (axis cs:2,2);
    \draw[simRust, densely dotted] (axis cs:4,1e-16) -- (axis cs:4,2);
    \draw[simRust, densely dotted] (axis cs:6,1e-16) -- (axis cs:6,2);
    \draw[simRust, densely dotted] (axis cs:8,1e-16) -- (axis cs:8,2);
    \node[font=\scriptsize, simRust, anchor=south] at (axis cs:4,1.1) {restarts};
  \end{axis}
\end{tikzpicture}
\caption[Convergence: full GMRES versus GMRES(2)]{Residual histories for
\cref{wex:restart} (log scale). Full GMRES (blue) converges superlinearly and reaches
machine precision at step $4=N$, exhausting the Krylov space. Restarted GMRES(2)
(rust) matches it for the first two steps, then restarts (dotted verticals) every two
steps; each restart discards the accumulated subspace, collapsing the convergence to a
slow linear rate. Both residual curves are monotone non-increasing---restart never
\emph{increases} the residual---but GMRES(2) needs many more iterations to reach the
same tolerance. The flat $O(kN)$ memory ceiling is paid for in convergence speed.}
\label{fig:convergence}
\end{figure}

\section{The variant axes: preconditioner side and the absorption view}

GMRES carries several orthogonal axes of variation, and they line up precisely with
the variant-absorption discipline of \cref{ch:operators}. We collect them here.

\subsection{Left, right, and no preconditioning}

A preconditioner can be applied on either side of $\mat{A}$, or not at all
(\cref{fig:precondside}):
\begin{description}[style=nextline,leftmargin=1.5em]
  \item[None.] Solve $\mat{A}\vx=\vb$ directly; the Krylov space is built from
  $\mat{A}$ and $\vr_0$. Used when $\mat{A}$ is already well-conditioned.
  \item[Left.] Solve $\mat{M}^{-1}\mat{A}\vx=\mat{M}^{-1}\vb$. The matvec at each step
  is $\vw=\mat{M}^{-1}(\mat{A}\vv_j)$ (precondition \emph{after} the operator). The
  residual GMRES minimizes is the \emph{preconditioned} residual
  $\norm{\mat{M}^{-1}(\vb-\mat{A}\vx)}$---which is not the true residual, a subtlety
  when reading the convergence test.
  \item[Right.] Solve $\mat{A}\mat{M}^{-1}\vu=\vb$ with $\vx=\mat{M}^{-1}\vu$. The
  matvec is $\vw=\mat{A}(\mat{M}^{-1}\vv_j)$ (precondition \emph{before} the operator).
  The residual minimized is the \emph{true} residual $\norm{\vb-\mat{A}\vx}$, which is
  why right preconditioning is often preferred for GMRES---and why FGMRES, which
  preconditions before the matvec, is naturally a right-preconditioned scheme.
\end{description}

\begin{figure}[t]
\centering
\begin{tikzpicture}[node distance=5mm]
  % three little pipelines stacked
  \foreach \row/\lbl/\seq in {
      1.7/{none}/{$\vv_j$/$\mat{A}$/$\vw$},
      0.0/{left}/{$\vv_j$/$\mat{A}$/$\mat{M}^{-1}$/$\vw$},
      -1.7/{right}/{$\vv_j$/$\mat{M}^{-1}$/$\mat{A}$/$\vw$}}{
    \node[font=\footnotesize\bfseries, simInk, anchor=east] at (-4.3,\row) {\lbl:};
  }
  % NONE
  \node[draw=simBlue,fill=simBgBlue,thick,rounded corners=2pt,minimum width=11mm,minimum height=7mm,font=\footnotesize] (n1) at (-3.0,1.7) {$\vv_j$};
  \node[draw=simTeal,fill=simBgGray,thick,rounded corners=2pt,minimum width=11mm,minimum height=7mm,font=\footnotesize,right=7mm of n1] (n2) {$\mat{A}$};
  \node[draw=simBlue,fill=simBgBlue,thick,rounded corners=2pt,minimum width=11mm,minimum height=7mm,font=\footnotesize,right=7mm of n2] (n3) {$\vw$};
  \draw[flow] (n1)--(n2); \draw[flow] (n2)--(n3);
  % LEFT
  \node[draw=simBlue,fill=simBgBlue,thick,rounded corners=2pt,minimum width=11mm,minimum height=7mm,font=\footnotesize] (l1) at (-3.0,0.0) {$\vv_j$};
  \node[draw=simTeal,fill=simBgGray,thick,rounded corners=2pt,minimum width=11mm,minimum height=7mm,font=\footnotesize,right=7mm of l1] (l2) {$\mat{A}$};
  \node[extern,minimum width=12mm,minimum height=7mm,font=\footnotesize,right=7mm of l2] (l3) {$\mat{M}^{-1}$};
  \node[draw=simBlue,fill=simBgBlue,thick,rounded corners=2pt,minimum width=11mm,minimum height=7mm,font=\footnotesize,right=7mm of l3] (l4) {$\vw$};
  \draw[flow] (l1)--(l2); \draw[flow] (l2)--(l3); \draw[flow] (l3)--(l4);
  % RIGHT
  \node[draw=simBlue,fill=simBgBlue,thick,rounded corners=2pt,minimum width=11mm,minimum height=7mm,font=\footnotesize] (r1) at (-3.0,-1.7) {$\vv_j$};
  \node[extern,minimum width=12mm,minimum height=7mm,font=\footnotesize,right=7mm of r1] (r2) {$\mat{M}^{-1}$};
  \node[draw=simTeal,fill=simBgGray,thick,rounded corners=2pt,minimum width=11mm,minimum height=7mm,font=\footnotesize,right=7mm of r2] (r3) {$\mat{A}$};
  \node[draw=simBlue,fill=simBgBlue,thick,rounded corners=2pt,minimum width=11mm,minimum height=7mm,font=\footnotesize,right=7mm of r3] (r4) {$\vw$};
  \draw[flow] (r1)--(r2); \draw[flow] (r2)--(r3); \draw[flow] (r3)--(r4);
  % note
  \node[font=\scriptsize\itshape, simGray, align=left, anchor=west] at (3.6,1.7)
    {minimizes true residual};
  \node[font=\scriptsize\itshape, simGray, align=left, anchor=west] at (3.6,0.0)
    {minimizes $\mat{M}^{-1}$-residual};
  \node[font=\scriptsize\itshape, simGray, align=left, anchor=west] at (3.6,-1.7)
    {minimizes true residual};
\end{tikzpicture}
\caption[The preconditioner-side axis]{The three settings of the preconditioner-side
variant. \emph{None:} the matvec is just $\mat{A}\vv_j$. \emph{Left:} apply
$\mat{M}^{-1}$ \emph{after} $\mat{A}$, solving $\mat{M}^{-1}\mat{A}\vx=\mat{M}^{-1}\vb$;
the residual GMRES sees is the preconditioned one. \emph{Right:} apply $\mat{M}^{-1}$
\emph{before} $\mat{A}$, solving $\mat{A}\mat{M}^{-1}\vu=\vb$; the residual stays the
true one. All three are the \emph{same} Arnoldi-plus-running-QR loop with the matvec's
operand chain rebound---a textbook variant absorbed by a constructed operator
(\cref{ch:operators}): the side is read once, at construction, and the per-step loop
never branches on it.}
\label{fig:precondside}
\end{figure}

\subsection{The absorption view}

These axes are the running example of \cref{ch:operators}'s discipline. The
preconditioner side, the restart length, the orthogonalization method, even the
choice of CG/GMRES/FGMRES itself are all read \emph{once}, at the factory that
constructs the solver (\code{ConfigureKrylovSolver}), and absorbed into a single
uniformly typed solver-operator. After construction the per-step loop is
variant-free: it asks the operator for \code{apply(A, v)} and the preconditioner for
\code{apply(pc, r)}, and never re-inspects which side, which method, which restart it
got. The fixed-versus-flexible axis is the \emph{one} axis that cannot be absorbed
this way---because, by the decision rule, a per-step preconditioner has run-time
lifetime---and that is exactly why FGMRES is a distinct algorithm with a distinct
state schema rather than a configuration flag on GMRES. Everything else is a binding;
flexibility alone is a structural fork, honestly disclosed by the extra $\mat{Z}$
basis.

\subsection{Convergence intuition: no clean condition-number bound}

CG (\cref{ch:cg}) came with a clean convergence bound: the error after $m$ steps
shrinks like $\bigl(\tfrac{\sqrt\kappa-1}{\sqrt\kappa+1}\bigr)^m$, governed entirely
by the condition number $\kappa$. GMRES has \emph{no such bound}, and the reason is
instructive. For a symmetric operator the spectrum lives on the real line and the
condition number captures everything; for a nonsymmetric operator the eigenvalues
scatter into the complex plane, and---worse---if $\mat{A}$ is far from normal, the
eigenvalues can lie anywhere and \emph{still} mislead, because the relevant object is
not the spectrum but the \emph{field of values} (the set $\{\inner{\mat{A}\vx}{\vx}:
\norm{\vx}=1\}$, a region of the complex plane). GMRES converges fast when the field
of values is a tidy cluster bounded away from the origin, and can stagnate when it
wraps around the origin---regardless of the condition number. The practical
consequence: GMRES convergence is judged empirically, by watching the free residual
$|s_{m+1}|$ fall, and is engineered by \emph{preconditioning}---reshaping the field of
values with $\mat{M}$ until the cluster is favorable. The whole point of
\cref{ch:precond,ch:multigrid} is to build the $\mat{M}$ that makes GMRES converge.

\summarysection
When $\mat{A}$ is nonsymmetric, the short three-term recurrence that makes CG cheap
simply does not exist (Faber--Manteuffel): optimality and bounded storage cannot
coexist. \emph{GMRES} keeps optimality and stores the whole basis. It is built on the
\emph{minimum-residual principle}---at step $m$, return the $\vx_m\in\vx_0+\mathcal{K}_m$
minimizing $\norm{\vb-\mat{A}\vx_m}$---which makes the residual norms monotone
non-increasing. The \emph{Arnoldi process} (Gram--Schmidt on the Krylov sequence,
\cref{ch:orthog}) builds an orthonormal basis $\mat{V}_{m+1}$ and the small
upper-Hessenberg $\bar{\mat{H}}_m$ satisfying the Arnoldi relation
$\mat{A}\mat{V}_m=\mat{V}_{m+1}\bar{\mat{H}}_m$; orthonormality then collapses the
$N$-dimensional minimization to the $(m{+}1)$-dimensional least-squares problem
$\min_{\vy}\norm{\beta\vec e_1-\bar{\mat{H}}_m\vy}$. A \emph{running QR by Givens
rotations} (\cref{ch:givens}) solves that small problem incrementally and---the central
economy---exposes the residual norm \emph{for free} as $|s_{m+1}|$ at every step,
without forming the iterate. Because storage grows as $O(mN)$ and work as $O(m^2N)$,
\emph{restarted GMRES($k$)} caps memory at $O(kN)$ by discarding the (ephemeral,
\cref{ch:strata}) basis every $k$ steps---at the cost of losing superlinear
convergence and risking stagnation. Finally, \emph{FGMRES} pays off \cref{ch:operators}'s
decision rule: when the preconditioner changes per step ($\mat{M}_j$, e.g.\ an inner
iterative solve), no constructed operator can absorb it, so the per-step variant is
threaded into the state as a second basis $\mat{Z}_m$ of remembered preconditioned
vectors, and the iterate is recovered as $\vx_m=\vx_0+\mat{Z}_m\vy_m$ rather than
$\vx_0+\mat{V}_m\vy_m$. Preconditioner side (left/right/none) is an absorbed variant; the
fixed-versus-flexible axis is the one structural fork, disclosed honestly by the extra
basis. GMRES carries no clean condition-number bound: its convergence is governed by the
field of values and engineered by preconditioning.

\furtherreading
The definitive treatment of GMRES, FGMRES, the Arnoldi process, restarting, and the
preconditioning variants is Saad~\citep{saad2003}, whose chapters on Krylov methods for
nonsymmetric systems and on flexible preconditioning are the direct source for this
chapter's algorithm and its decision-rule reading. For the numerical-analysis
perspective---the Arnoldi relation as an orthogonal reduction, the role of
nonnormality and the field of values in convergence, and the contrast with the clean
symmetric theory---Trefethen and Bau~\citep{trefethenbau1997} (the lectures on
Arnoldi, GMRES, and iterative-method convergence) is the ideal companion. The
running-QR/Givens mechanics this chapter used as a black box are developed in full in
\cref{ch:givens}; the orthogonalization inside Arnoldi in \cref{ch:orthog}; and the
preconditioners $\mat{M}$ whose construction makes GMRES converge in
\cref{ch:precond,ch:multigrid}.

\exercisessection
\begin{enumerate}[nosep]
  \item \textbf{Why no short recurrence.} State, in your own words, what the
  Faber--Manteuffel theorem rules out, and explain why CG's three-term recurrence is a
  special consequence of symmetry. What \emph{two} desirable properties does GMRES
  refuse to give up simultaneously, and which one does restarting trade away?

  \item \textbf{Monotonicity.} Prove that the GMRES residual norms satisfy
  $\norm{\vr_{m+1}}\le\norm{\vr_m}$ using only the minimum-residual principle and the
  nesting $\mathcal{K}_m\subseteq\mathcal{K}_{m+1}$. Why does this argument
  \emph{not} establish strict decrease, and when can the residual stall at a plateau?

  \item \textbf{A third Arnoldi step.} Continue \cref{wex:arnoldi} with one more step:
  compute $\mat{A}\vv_3$, the coefficients $h_{13},h_{23},h_{33}$, and the new
  sub-diagonal $h_{43}$. What is special about $h_{43}$ for this particular $3\times3$
  operator, and what does its value tell you about where GMRES terminates exactly?

  \item \textbf{The reduction, by hand.} Starting from $\vr_0=\beta\vv_1$ and the
  Arnoldi relation, rederive the identity $\vb-\mat{A}(\vx_0+\mat{V}_m\vy)=
  \mat{V}_{m+1}(\beta\vec e_1-\bar{\mat{H}}_m\vy)$, naming the single property of
  $\mat{V}_{m+1}$ that lets you drop it from the norm. Which step would fail if the
  Arnoldi basis were merely linearly independent rather than orthonormal?

  \item \textbf{The free residual.} In \cref{wex:leastsq} the running QR reported
  $\norm{\vr_2}=|s_3|=1/\sqrt{69}$ \emph{without} forming $\vx_2$. Explain precisely
  why $|s_{m+1}|$ equals the least-squares residual norm, citing the isometry property
  of Givens rotations. Then state what GMRES would have to compute instead if it lacked
  this readout, and estimate the cost difference per step.

  \item \textbf{Sizing a restart.} A solve has $N=10^6$ unknowns; each vector is $8$~MB.
  You have a $4$~GB memory budget for the Krylov basis. (i) What is the largest restart
  length $k$ you can afford? (ii) If full GMRES would converge in $40$ steps but
  GMRES($k$) for your $k$ needs $300$, has restarting helped or hurt, and in what sense?
  (iii) Name one operator property under which GMRES($k$) might \emph{stagnate} entirely.

  \item \textbf{The mutating-preconditioner trap, revisited.} A colleague proposes to
  run \emph{plain} GMRES with an inner-iterative preconditioner whose accuracy varies
  step to step, ``since GMRES only ever asks \code{apply(pc, v)} and won't notice.''
  Using the $\mat{Z}$-versus-$\mat{V}$ distinction and the decision rule of
  \cref{ch:operators}, explain exactly what goes wrong, why no error is raised, and why
  the reported $|s_{m+1}|$ is misleading. What is the minimal change that fixes it, and
  what does it cost in storage?

  \item \textbf{$\mat{V}$ versus $\mat{Z}$.} For \emph{fixed} right-preconditioned GMRES,
  the iterate is $\vx_0+\mat{M}^{-1}\mat{V}_m\vy$; for FGMRES it is $\vx_0+\mat{Z}_m\vy$.
  Show that when $\mat{M}_j=\mat{M}$ for all $j$ (the preconditioner happens to be
  constant), these two formulas \emph{coincide}, so FGMRES reduces to fixed
  right-preconditioned GMRES. What, then, is the only cost of always using FGMRES, and
  when is paying it worthwhile?

  \item \textbf{Side and the residual.} Left and right preconditioning minimize
  \emph{different} residuals (\cref{fig:precondside}). For a convergence test with
  tolerance $\varepsilon$ on the \emph{true} residual $\norm{\vb-\mat{A}\vx}$, which
  side lets you read the test directly from the free $|s_{m+1}|$, and which forces a
  correction? Explain why this is one practical reason GMRES implementations often
  default to right preconditioning.

  \item \textbf{Absorbed or forked?} Classify each axis as either \emph{absorbed} (read
  once at construction, per-step loop variant-free) or a \emph{structural fork} (forces
  a different operation chain or state schema), and justify with the decision rule:
  (a) preconditioner side $\in\{\text{none},\text{left},\text{right}\}$;
  (b) restart length $k$; (c) orthogonalization method (MGS vs CGS); (d)
  fixed vs per-step preconditioner. For the one you classify as a fork, name the extra
  state it forces into the carry.
\end{enumerate}

\chapter{Orthogonalization: Arnoldi, Lanczos, and Gram--Schmidt}
\label{ch:orthog}

\chapterorient{Every Krylov solver in this Part rests on one ability: to take a
growing pile of nearly parallel vectors and straighten them into an orthonormal
basis. GMRES (\cref{ch:gmres}) called this step \code{orthogonalize} and moved
on; this is the careful version. We build Gram--Schmidt from its geometry, then
meet its two faces---classical and modified---and discover that the difference
between them is not arithmetic but \emph{order}: modified Gram--Schmidt is a
genuine sequential dependency, a concrete specimen of the obstruction of
\cref{ch:rotations}. We watch both lose orthogonality in finite precision, learn
the ``twice is enough'' repair, and then assemble the orthogonalizer into the two
recurrences the rest of the book runs on: Arnoldi, which feeds GMRES, and
Lanczos, which---for symmetric operators---collapses to a miraculous three-term
recurrence and then, just as miraculously, falls apart numerically. The thread
running through all of it is a single \emph{contract}: ``the basis is
orthonormal, to a tolerance.'' Everything downstream depends on that contract
holding.}

% local: lowercase bold-a for the worked-example input vectors (not in the shared preamble)
\providecommand{\va}{\mathbf{a}}

\section{The goal: an orthonormal basis for a Krylov subspace}
\label{sec:goal}

A Krylov subspace solver does not work with the operator $\mat{A}$ directly. It
works inside the \emph{Krylov subspace}
\[
  \mathcal{K}_m(\mat{A}, \vv)
    \;=\; \operatorname{span}\bigl\{\vv,\;\mat{A}\vv,\;\mat{A}^2\vv,\;\dots,\;
      \mat{A}^{m-1}\vv\bigr\},
\]
the span of the vector $\vv$ and its repeated images under $\mat{A}$. The reason
this subspace is useful is the subject of \cref{ch:gmres} and \cref{ch:cg}; here
we take it as given and ask a narrower question. To compute \emph{inside}
$\mathcal{K}_m$---to solve a least-squares problem there, to project the operator
onto it---we need a \emph{basis}: a set of $m$ vectors that span it and that we
can compute with. The obvious basis is the one in the definition, the
\emph{power basis} $\{\vv, \mat{A}\vv, \mat{A}^2\vv, \dots\}$. It is also,
catastrophically, the worst possible basis to compute with.

\begin{figure}[t]
\centering
\begin{tikzpicture}[font=\footnotesize, scale=1.0]
  % ===== LEFT: the power basis collapses toward a line =====
  \begin{scope}
    \node[font=\small\bfseries, text=simRust] at (1.5,2.7) {power basis: nearly parallel};
    \draw[->,simGray] (-0.2,0)--(3.1,0);
    \draw[->,simGray] (0,-0.2)--(0,2.4);
    % v, Av, A2v all crowding one direction
    \draw[flow,simRust] (0,0)--(2.6,1.5) node[right,font=\scriptsize]{$\vv$};
    \draw[flow,simRust] (0,0)--(2.75,1.25) node[right,font=\scriptsize]{$\mat{A}\vv$};
    \draw[flow,simRust] (0,0)--(2.85,1.05) node[right,font=\scriptsize]{$\mat{A}^2\vv$};
    \node[font=\scriptsize\itshape, text=simGray, align=center] at (1.5,-0.7)
      {they crowd one direction:\\the dominant eigenvector};
  \end{scope}
  \draw[simGray!45, thick] (3.7,-1.0) -- (3.7,3.0);
  % ===== RIGHT: the orthonormal basis spreads out =====
  \begin{scope}[xshift=5.0cm]
    \node[font=\small\bfseries, text=simGreen] at (1.3,2.7) {orthonormal basis: spread out};
    \draw[->,simGray] (-0.2,0)--(2.6,0);
    \draw[->,simGray] (0,-0.2)--(0,2.4);
    \draw[flow,simGreen] (0,0)--(2.0,0) node[right,font=\scriptsize]{$\vq_1$};
    \draw[flow,simGreen] (0,0)--(0,2.0) node[above,font=\scriptsize]{$\vq_2$};
    \draw (1.7,0) -- (1.7,0.3) -- (2.0,0.3);
    \node[font=\scriptsize] at (1.9,0.5) {$90^\circ$};
    \node[font=\scriptsize\itshape, text=simGray, align=center] at (1.3,-0.7)
      {mutually orthogonal,\\unit length: well-conditioned};
  \end{scope}
\end{tikzpicture}
\caption[Power basis versus orthonormal basis]{Why the Krylov subspace needs an
orthonormal basis. \emph{Left:} the natural power basis
$\{\vv, \mat{A}\vv, \mat{A}^2\vv, \dots\}$ is a disaster to compute with---each
application of $\mat{A}$ pulls the vector further toward the dominant eigenvector,
so successive members crowd into nearly the same direction. A matrix whose columns
are these vectors is severely ill-conditioned, and in finite precision the later
members are numerically indistinguishable from the earlier ones. \emph{Right:} an
\emph{orthonormal} basis $\{\vq_1, \vq_2, \dots\}$ spanning the same subspace
spreads the directions out to mutual right angles and unit length. It is the basis
every Krylov method actually uses, and producing it from the power basis is the
business of this chapter.}
\label{fig:powerbasis}
\end{figure}

The trouble is visible in \cref{fig:powerbasis}. Each multiplication by $\mat{A}$
amplifies the component of the vector along the \emph{dominant} eigenvector---the
one with the largest eigenvalue in magnitude---more than any other. So $\mat{A}\vv$
leans toward that eigenvector, $\mat{A}^2\vv$ leans further, and within a handful
of steps the power-basis vectors are all crowding into nearly the same direction.
They still, in exact arithmetic, span the full subspace; but the matrix whose
columns they are is wildly ill-conditioned, and in floating point the later columns
carry almost no new information. This is the same effect that makes the
\emph{power method} converge to the dominant eigenvector---a feature there, a bug
here.

The cure is to build, as we go, an \emph{orthonormal} basis for the same subspace:
a set of vectors $\vq_1, \vq_2, \dots, \vq_m$ that are mutually orthogonal and of
unit length,
\[
  \inner{\vq_i}{\vq_j} = \delta_{ij}
    \;=\;
    \begin{cases} 1 & i = j,\\ 0 & i \ne j, \end{cases}
\]
and that span $\mathcal{K}_m$ exactly as the power basis does. An orthonormal basis
is perfectly conditioned---the matrix $\mat{Q}$ with these columns satisfies
$\mat{Q}^{\!\top}\mat{Q} = \mat{I}$, the best-conditioned matrix there is---so every
later computation inside $\mathcal{K}_m$ is numerically clean. The procedure that
manufactures such a basis, one vector at a time, from the power basis is
\emph{Gram--Schmidt orthogonalization}, and it is the foundation everything in this
chapter is built on.

\begin{keyidea}
A Krylov solver works inside the subspace
$\mathcal{K}_m(\mat{A},\vv) = \operatorname{span}\{\vv, \mat{A}\vv, \dots,
\mat{A}^{m-1}\vv\}$, but never in the natural power basis, which collapses toward
the dominant eigenvector and is hopelessly ill-conditioned. It works instead in an
\emph{orthonormal} basis $\vq_1,\dots,\vq_m$ for the same subspace---mutually
orthogonal, unit length, $\inner{\vq_i}{\vq_j} = \delta_{ij}$. Building that basis
incrementally, as each new Krylov vector arrives, is the job of Gram--Schmidt, and
the quality of the basis it produces is what every downstream method silently
relies on.
\end{keyidea}

\section{Gram--Schmidt: subtract the projections, then normalize}
\label{sec:gschmidt}

Gram--Schmidt is one idea, applied repeatedly. Suppose we have already built an
orthonormal set $\vq_1, \dots, \vq_{m}$ and a new vector $\vw$ arrives (in the
Krylov setting, $\vw = \mat{A}\vq_m$, the next power-basis member). We want to turn
$\vw$ into the next basis vector $\vq_{m+1}$. The new vector almost certainly
overlaps the directions we already have; the orthonormal basis must point somewhere
\emph{new}. So we do two things, in order:

\begin{enumerate}
  \item \textbf{Subtract off the part of $\vw$ that lies in the span we already
  have.} The component of $\vw$ along a unit vector $\vq_j$ is the scalar
  projection $\inner{\vw}{\vq_j}$, and the vector projection is
  $\inner{\vw}{\vq_j}\,\vq_j$. Subtracting all of these leaves the
  \emph{orthogonal residual}
  \[
    \vw' \;=\; \vw \;-\; \sum_{j=1}^{m} \inner{\vw}{\vq_j}\,\vq_j,
  \]
  the part of $\vw$ that points outside $\operatorname{span}\{\vq_1,\dots,\vq_m\}$.
  \item \textbf{Normalize.} The residual $\vw'$ has the right \emph{direction} but
  not unit length, so we scale it:
  \[
    \vq_{m+1} \;=\; \frac{\vw'}{\norm{\vw'}}.
  \]
\end{enumerate}

The first step is the heart of the matter; the second is bookkeeping.
\Cref{fig:gschmidt} draws the geometry in three dimensions, where it can be seen
all at once. The vector $\vw$ casts a ``shadow'' onto the plane spanned by
$\vq_1, \vq_2$; subtracting that shadow leaves exactly the vertical component, which
points out of the plane in a brand-new direction.

\begin{figure}[t]
\centering
\begin{tikzpicture}[font=\footnotesize, scale=1.05, >={Stealth[length=2.4mm]}]
  % the q1-q2 plane drawn as a parallelogram
  \fill[simBgGreen, opacity=0.7] (-0.3,-0.4) -- (3.0,-0.1) -- (3.6,0.9) -- (0.3,0.6) -- cycle;
  \node[font=\scriptsize\itshape, text=simGreen] at (3.1,1.2) {$\operatorname{span}\{\vq_1,\vq_2\}$};
  % basis vectors q1, q2 in the plane (origin at O)
  \coordinate (O) at (0.6,0.1);
  \draw[flow,simGreen] (O) -- ++(2.0,0.15) node[below right,font=\scriptsize]{$\vq_1$};
  \draw[flow,simGreen] (O) -- ++(0.55,0.45) node[above left,font=\scriptsize]{$\vq_2$};
  % w sticking up out of the plane
  \coordinate (Wtip) at (2.3,2.5);
  \draw[flow,simRust,very thick] (O) -- (Wtip) node[above,font=\scriptsize]{$\vw$};
  % the shadow (projection onto the plane)
  \coordinate (Ptip) at (2.55,1.0);
  \draw[flow,simBlue] (O) -- (Ptip)
    node[pos=0.6,below right,font=\scriptsize,text=simBlue]
      {$\displaystyle\sum_j\inner{\vw}{\vq_j}\vq_j$};
  % the orthogonal residual w' (vertical), drawn from Ptip up to Wtip
  \draw[flow,simViolet,very thick] (Ptip) -- (Wtip)
    node[pos=0.5,right,font=\scriptsize,text=simViolet]{$\vw'$ (residual)};
  % dashed drop line emphasising the shadow
  \draw[densely dashed,simGray] (Wtip) -- (Ptip);
  % right-angle mark at Ptip between plane and residual
  \draw[simInk] ($(Ptip)+(-0.18,0.0)$) -- ($(Ptip)+(-0.18,0.18)$) -- ($(Ptip)+(0.0,0.18)$);
  % normalize annotation
  \node[font=\scriptsize\itshape, text=simGray, align=left, anchor=west] at (4.0,2.3)
    {\textbf{1.} subtract the shadow\\(projection onto the span)};
  \node[font=\scriptsize\itshape, text=simGray, align=left, anchor=west] at (4.0,1.3)
    {\textbf{2.} normalize $\vw'$:\\$\vq_3 = \vw'/\norm{\vw'}$};
  \node[opbox, anchor=west] at (4.0,0.3) {$\vq_3 \perp \operatorname{span}\{\vq_1,\vq_2\}$};
\end{tikzpicture}
\caption[Gram--Schmidt, geometrically]{One Gram--Schmidt step, seen in three
dimensions. The new vector $\vw$ (rust) sticks up out of the plane already spanned
by $\vq_1,\vq_2$ (green). Its \emph{shadow} on that plane is the projection
$\sum_j \inner{\vw}{\vq_j}\vq_j$ (blue)---the part of $\vw$ we already have.
Subtracting the shadow leaves the orthogonal residual $\vw'$ (violet), which by
construction meets the plane at a right angle and so points in a genuinely new
direction. Normalizing $\vw'$ to unit length yields the next basis vector $\vq_3$.
Every orthogonalization method in this chapter is some scheduling of exactly this
\emph{subtract-the-projections-then-normalize} step; what differs is the
\emph{order} in which the projections are subtracted.}
\label{fig:gschmidt}
\end{figure}

That the residual really is orthogonal to everything we kept is a one-line
calculation, and worth doing once because it is the contract the whole chapter
guarantees.

\begin{lemma}[The residual is orthogonal to the span]
\label{lem:residorthog}
Let $\vq_1,\dots,\vq_m$ be orthonormal and
$\vw' = \vw - \sum_{j=1}^m \inner{\vw}{\vq_j}\vq_j$. Then
$\inner{\vw'}{\vq_i} = 0$ for every $i \in \{1,\dots,m\}$.
\end{lemma}
\begin{proof}
Take the inner product of $\vw'$ with $\vq_i$ and use linearity:
\[
  \inner{\vw'}{\vq_i}
    = \inner{\vw}{\vq_i} - \sum_{j=1}^m \inner{\vw}{\vq_j}\,\inner{\vq_j}{\vq_i}
    = \inner{\vw}{\vq_i} - \sum_{j=1}^m \inner{\vw}{\vq_j}\,\delta_{ji}
    = \inner{\vw}{\vq_i} - \inner{\vw}{\vq_i} = 0,
\]
where the middle step used orthonormality, $\inner{\vq_j}{\vq_i} = \delta_{ji}$, to
collapse the sum to its single $j = i$ term. \qedhere
\end{proof}

The coefficients $\inner{\vw}{\vq_j}$ are not waste; they are the record of
\emph{how much} of $\vw$ lay in each previous direction. Stacked into a column they
become the entries of the \emph{Hessenberg matrix} that Arnoldi produces
(\cref{sec:arnoldi}), and the solver downstream reads them as the projection of the
operator onto the Krylov basis. Orthogonalization thus produces two things at once:
the new basis vector, and the column of coefficients. Keep both in view---they are
the $(\vw', H)$ pair that the framework's \code{orthogonalize} operator returns.

\begin{notationbox}
We follow the framework's \code{orthogonalize} contract: the operator
\emph{orthogonalizes against a normalized basis but does not itself normalize its
output.} It returns the residual $\vw'$ and the coefficient column $H$ where
$H_j = \inner{\vw}{\vq_j}$; the caller takes $\norm{\vw'}$ (which becomes the next
Hessenberg sub-diagonal entry) and divides. We keep the normalization explicit and
separate throughout, because the two recurrences of \cref{sec:arnoldi,sec:lanczos}
attach different meaning to that final norm.
\end{notationbox}

\section{Classical versus modified Gram--Schmidt}
\label{sec:cgsmgs}

The recipe of \cref{sec:gschmidt} hides a scheduling decision that does not matter
in exact arithmetic and matters enormously in floating point. The residual is a sum
of $m$ projection-subtractions; the question is \emph{against which vector each
projection coefficient is computed}. There are two natural answers, and they are the
two classical algorithms.

\subsection{Classical Gram--Schmidt (CGS)}

The literal reading of the formula
$\vw' = \vw - \sum_j \inner{\vw}{\vq_j}\vq_j$ computes every coefficient against the
\emph{same original} $\vw$:
\[
  H_j = \inner{\vw}{\vq_j} \quad\text{for all } j, \qquad
  \vw' = \vw - \sum_{j=1}^m H_j\,\vq_j .
\]
All $m$ inner products read the input $\vw$ and nothing else, so they are mutually
independent: they can be computed in any order, batched into a single matrix-vector
product $\mat{Q}^{\!\top}\vw$, and---on a parallel machine---reduced in one
collective communication. This is \emph{classical Gram--Schmidt}, CGS, and its
defining structural feature is that \emph{the $m$ projections do not depend on one
another}.

\subsection{Modified Gram--Schmidt (MGS)}

The alternative subtracts the projections \emph{one at a time}, and each coefficient
is computed against the \emph{partially-orthogonalized} vector left by the previous
subtractions. Write $\vw^{(0)} = \vw$ and, for $j = 1, \dots, m$,
\[
  H_j = \inner{\vw^{(j-1)}}{\vq_j}, \qquad
  \vw^{(j)} = \vw^{(j-1)} - H_j\,\vq_j,
\]
and set $\vw' = \vw^{(m)}$. This is \emph{modified Gram--Schmidt}, MGS. Each step
projects the \emph{running residual} onto the next basis vector, rather than
projecting the original. Equivalently, MGS applies a left-to-right composition of
rank-one projectors,
\[
  \vw' = \bigl(\mat{I} - \vq_m\vq_m^{\!\top}\bigr)\cdots
         \bigl(\mat{I} - \vq_1\vq_1^{\!\top}\bigr)\,\vw,
\]
where CGS applies the single combined projector
$\bigl(\mat{I} - \sum_j \vq_j\vq_j^{\!\top}\bigr)\vw$ in one shot.

\begin{keyidea}
CGS and MGS compute the \emph{same orthogonal residual in exact arithmetic}---both
are the projection $(\mat{I} - \mat{Q}\mat{Q}^{\!\top})\vw$, and the projection does
not care how you bracket the subtractions. They differ only in
\emph{against which vector each coefficient is taken}: CGS uses the original $\vw$
for all of them (independent, parallelizable), MGS uses the running residual
$\vw^{(j-1)}$ for the $j$-th (each depends on the last). That single difference of
\emph{ordering} is the whole story---it has no effect in exact arithmetic and a
decisive effect in finite precision, and it is the reason both algorithms exist.
\end{keyidea}

\subsection{The same residual, two schedules}

The algebraic equivalence is worth seeing directly. In MGS,
\[
  \vw^{(j)} = \vw^{(j-1)} - \inner{\vw^{(j-1)}}{\vq_j}\vq_j .
\]
But $\vw^{(j-1)} = \vw - \sum_{i<j}\inner{\vw^{(i-1)}}{\vq_i}\vq_i$ already has the
$\vq_i$ ($i<j$) components removed, and since the basis is orthonormal,
$\inner{\vw^{(j-1)}}{\vq_j} = \inner{\vw}{\vq_j}$ exactly---subtracting components
along \emph{other} basis vectors does not change the component along $\vq_j$. So in
exact arithmetic MGS's coefficient equals CGS's, the running subtraction reaches the
same place, and $\vw^{(m)} = \vw - \sum_j \inner{\vw}{\vq_j}\vq_j$. The two are one
computation with two bracketings. \Cref{fig:cgsmgs} draws the data-flow difference
that this algebraic identity papers over.

\begin{figure}[t]
\centering
\begin{tikzpicture}[font=\footnotesize]
  % ===== LEFT: CGS = fan-out, independent =====
  \begin{scope}
    \node[font=\small\bfseries, text=simBlue] at (1.4,3.0) {CGS: independent projections};
    \node[data, minimum width=10mm] (w) at (1.4,2.2) {$\vw$};
    \foreach \i/\x in {1/0,2/1.4,3/2.8}{
      \node[opbox, minimum width=11mm] (p\i) at (\x,1.0)
        {$\inner{\vw}{\vq_{\i}}$};
      \draw[flow] (w.south) -- (p\i.north);
    }
    \node[product, minimum width=30mm] (sum) at (1.4,-0.3)
      {$\vw' = \vw - \sum_j H_j\vq_j$};
    \foreach \i in {1,2,3}{\draw[flow] (p\i.south) -- (sum.north);}
    \node[font=\scriptsize, text=simGreen, align=center] at (1.4,-1.3)
      {all read the \emph{same} $\vw$:\\one batched reduction $\mat{Q}^{\!\top}\vw$,\\
       \emph{parallelizable}};
  \end{scope}
  \draw[simGray!50, thick] (4.2,-1.7) -- (4.2,3.2);
  % ===== RIGHT: MGS = sequential chain =====
  \begin{scope}[xshift=5.2cm]
    \node[font=\small\bfseries, text=simRust] at (2.3,3.0) {MGS: sequential subtraction};
    \node[data, minimum width=9mm] (w0) at (-0.4,1.0) {$\vw^{(0)}$};
    \foreach \i/\x in {1/1.0,2/2.5,3/4.0}{
      \node[opbox, minimum width=10mm] (s\i) at (\x,1.0) {$-H_{\i}\vq_{\i}$};
    }
    \node[product, minimum width=10mm] (w3) at (5.4,1.0) {$\vw'$};
    \draw[flow] (w0.east) -- (s1.west);
    \draw[flow] (s1.east) -- node[above,font=\scriptsize]{$\vw^{(1)}$} (s2.west);
    \draw[flow] (s2.east) -- node[above,font=\scriptsize]{$\vw^{(2)}$} (s3.west);
    \draw[flow] (s3.east) -- (w3.west);
    % each coefficient READS the incoming running residual
    \foreach \i/\x in {1/1.0,2/2.5,3/4.0}{
      \node[font=\scriptsize, text=simGold!80!black] at (\x,2.05)
        {$H_{\i}=\inner{\vw^{(\the\numexpr\i-1\relax)}}{\vq_{\i}}$};
      \draw[refedge, simGold!80!black] (\x,1.85) -- (s\i.north);
    }
    \node[font=\scriptsize, text=simRust, align=center] at (2.5,-0.6)
      {each $H_j$ reads the previous residual $\vw^{(j-1)}$:\\
       \textbf{column order is a sequential dependency},\\ \emph{not parallelizable}};
  \end{scope}
\end{tikzpicture}
\caption[Classical versus modified Gram--Schmidt]{The data-flow contrast that the
algebra hides. \emph{Left (CGS):} every projection coefficient $H_j$ reads the
\emph{same} original $\vw$, so the $m$ inner products fan out independently---they
are a single batched reduction $\mat{Q}^{\!\top}\vw$, reorderable and parallelizable,
with no arrows \emph{between} them. \emph{Right (MGS):} the coefficient $H_j$ reads
the \emph{running residual} $\vw^{(j-1)}$ left by the previous subtraction (gold
read-arrows), so the steps form a strict chain $\vw^{(0)}\!\to\!\vw^{(1)}\!\to\!
\cdots\!\to\!\vw'$ in which step $j$ cannot start until step $j-1$ finishes. This is
exactly the \code{map}-versus-\code{fold} split of \cref{ch:combinators}: CGS is a
\code{map} of independent dots, MGS is a \code{fold} threading the residual. The
column order is a genuine sequential dependency, the obstruction of
\cref{ch:rotations} in miniature.}
\label{fig:cgsmgs}
\end{figure}

\subsection{MGS is a sequential obstruction}

The right-hand side of \cref{fig:cgsmgs} should look familiar: it is precisely the
\code{fold} of \cref{ch:combinators}, and its carried dependency is precisely the
\emph{sequential obstruction} of \cref{ch:rotations}. MGS threads a running residual
through a chain of subtractions; the $j$-th step \emph{reads the output of the
$(j-1)$-th}. That is the loop-carried dependency that blocks the iteration rotation
($L_2 \to L_3$): you cannot turn MGS into a single ``orthogonalize against all
columns at once'' whole-field operation, because there is no such operation---the
later subtractions must see the results of the earlier ones.

\begin{pitfall}
The temptation, looking at CGS, is to ``optimize'' MGS into it: both compute the same
residual in exact arithmetic, and CGS parallelizes, so why not always use CGS? Because
the equivalence is \emph{exact-arithmetic only}. In finite precision CGS is
numerically unstable---it loses orthogonality fast for ill-conditioned bases
(\cref{sec:lossorthog})---while MGS holds it far better. The sequential dependency of
MGS is not an accident to be optimized away; it is \emph{load-bearing}. Each step's
subtraction against the freshly-cleaned residual is exactly what keeps the rounding
error from accumulating. Replacing MGS with CGS for speed silently degrades the
basis, and the degradation only shows up later, as a solver that stalls or a
spurious eigenvalue. This is the canonical case where a \emph{sequential dependency
is the algorithm}, not an obstacle to it.
\end{pitfall}

There is an algebraic fingerprint of this dependency, and it is the cleanest way to
state ``MGS does not parallelize.'' For CGS, permuting the columns of $\mat{Q}$
permutes the coefficients but leaves the residual unchanged (up to reduction-order
rounding), because the combined projector $\mat{I} - \sum_j\vq_j\vq_j^{\!\top}$ is
symmetric in its terms. For MGS, permuting the columns changes the intermediate
residuals $\vw^{(j)}$ and hence the bit-level result: the rank-one projectors
$(\mat{I} - \vq_j\vq_j^{\!\top})$ do not commute in finite precision. \emph{Column
order matters for MGS and not for CGS}---that non-commutativity is the algebraic
shadow of the sequential obstruction.

\subsection{A worked comparison on nearly parallel vectors}

The whole point of preferring MGS is invisible until we watch both run in finite
precision on an ill-conditioned input. \Cref{wex:cgsmgs} does exactly that.

\begin{workedexample}{CGS versus MGS on three nearly parallel vectors}{cgsmgs}
We orthonormalize three nearly parallel vectors in $\Real^3$ and measure how
orthogonal the resulting basis actually is. Take a small parameter
$\varepsilon = 10^{-4}$ and the columns
\[
  \va_1 = \begin{psmallmatrix} 1 \\ \varepsilon \\ 0 \end{psmallmatrix},\quad
  \va_2 = \begin{psmallmatrix} 1 \\ 0 \\ \varepsilon \end{psmallmatrix},\quad
  \va_3 = \begin{psmallmatrix} 1 \\ \varepsilon \\ \varepsilon \end{psmallmatrix}.
\]
All three are within $\varepsilon$ of the vector $(1,0,0)^{\!\top}$, so they are
\emph{nearly parallel}: a textbook ill-conditioned input, exactly the situation the
power basis produces. We carry the arithmetic as if in a precision coarse enough
that terms of order $\varepsilon^2$ relative to $1$ are lost---the qualitative
behavior is the same in IEEE double, just at $\varepsilon \approx 10^{-8}$.

\smallskip
\emph{Step 1 (shared by both).} Normalize $\va_1$. Its norm is
$\sqrt{1 + \varepsilon^2} \approx 1$, so $\vq_1 = \va_1 = (1, \varepsilon, 0)^{\!\top}$
to working precision.

\smallskip
\emph{Step 2 (shared).} Orthogonalize $\va_2$ against $\vq_1$. The coefficient is
$\inner{\va_2}{\vq_1} = 1\cdot 1 + 0\cdot\varepsilon + \varepsilon\cdot 0 = 1$, so
\[
  \vw' = \va_2 - 1\cdot\vq_1
       = \begin{psmallmatrix}1\\0\\\varepsilon\end{psmallmatrix}
       - \begin{psmallmatrix}1\\\varepsilon\\0\end{psmallmatrix}
       = \begin{psmallmatrix}0\\-\varepsilon\\\varepsilon\end{psmallmatrix},
  \qquad
  \vq_2 = \frac{\vw'}{\norm{\vw'}}
        = \frac{1}{\sqrt2}\begin{psmallmatrix}0\\-1\\1\end{psmallmatrix}.
\]
So far CGS and MGS agree---the divergence is in step 3, the first step with
\emph{two} prior vectors to project against.

\smallskip
\emph{Step 3, CGS.} Both coefficients are taken against the \emph{original} $\va_3$:
\[
  H_1 = \inner{\va_3}{\vq_1} = 1 + \varepsilon^2 \approx 1, \qquad
  H_2 = \inner{\va_3}{\vq_2}
      = \tfrac{1}{\sqrt2}(0 - \varepsilon + \varepsilon) = 0.
\]
Then
\[
  \vw'_{\text{CGS}} = \va_3 - H_1\vq_1 - H_2\vq_2
    = \begin{psmallmatrix}1\\\varepsilon\\\varepsilon\end{psmallmatrix}
    - \begin{psmallmatrix}1\\\varepsilon\\0\end{psmallmatrix}
    = \begin{psmallmatrix}0\\0\\\varepsilon\end{psmallmatrix},
  \qquad
  \vq_3^{\text{CGS}} = \begin{psmallmatrix}0\\0\\1\end{psmallmatrix}.
\]
This looks clean---but the trap is in $H_1$. The exact $\inner{\va_3}{\vq_1}$ is
$1 + \varepsilon^2$, and in finite precision the $\varepsilon^2$ term is
\emph{below the rounding threshold relative to $1$} and is dropped. The CGS residual
is computed by subtracting a rounded $H_1 \vq_1$ from $\va_3$, and the small error in
$H_1$ contaminates the $\vq_1$-direction of the residual: $\vq_3^{\text{CGS}}$ retains
a component of size $\sim\varepsilon$ along $\vq_1$ that should have been removed.
Measuring the residual orthogonality after the fact:
\[
  \inner{\vq_3^{\text{CGS}}}{\vq_1} \;\sim\; \varepsilon, \qquad
  \inner{\vq_3^{\text{CGS}}}{\vq_2} \;\sim\; \varepsilon .
\]

\smallskip
\emph{Step 3, MGS.} Now subtract sequentially. First against $\vq_1$:
\[
  H_1 = \inner{\va_3}{\vq_1} \approx 1, \qquad
  \vw^{(1)} = \va_3 - H_1\vq_1
    = \begin{psmallmatrix}0\\0\\\varepsilon\end{psmallmatrix}
    \quad(\text{to working precision}).
\]
Then take the \emph{second} coefficient against the cleaned residual $\vw^{(1)}$, not
against $\va_3$:
\[
  H_2 = \inner{\vw^{(1)}}{\vq_2}
      = \tfrac{1}{\sqrt2}(0 - 0 + \varepsilon) = \tfrac{\varepsilon}{\sqrt2},
  \qquad
  \vw^{(2)} = \vw^{(1)} - H_2\vq_2
    = \begin{psmallmatrix}0\\0\\\varepsilon\end{psmallmatrix}
      - \tfrac{\varepsilon}{2}\begin{psmallmatrix}0\\-1\\1\end{psmallmatrix}
    = \tfrac{\varepsilon}{2}\begin{psmallmatrix}0\\1\\1\end{psmallmatrix}.
\]
The crucial difference: MGS \emph{noticed} the residual's small $\vq_2$-overlap
($H_2 = \varepsilon/\sqrt2$, which CGS computed as $0$ because it looked at the
original $\va_3$ where the overlap was masked by the large shared component) and
removed it. The resulting $\vq_3^{\text{MGS}} = \tfrac{1}{\sqrt2}(0,1,1)^{\!\top}$
satisfies, to working precision,
\[
  \inner{\vq_3^{\text{MGS}}}{\vq_1} \;\sim\; \varepsilon^2, \qquad
  \inner{\vq_3^{\text{MGS}}}{\vq_2} \;\sim\; \varepsilon^2 .
\]

\smallskip
\emph{The verdict.} Both algorithms produce a unit vector roughly orthogonal to the
first two; but the \emph{residual orthogonality} differs by a full power of
$\varepsilon$: CGS leaves overlaps of size $\varepsilon$, MGS of size
$\varepsilon^2$. At $\varepsilon = 10^{-4}$ that is $10^{-4}$ versus $10^{-8}$; at the
IEEE-double analogue it is the difference between a basis that is orthogonal to four
digits and one orthogonal to eight. The reason is exactly the sequential dependency:
MGS's second projection sees a residual from which the first overlap has \emph{already}
been cleaned, so it can detect and remove the small leftover; CGS's second projection
looks at the raw input, where the same small overlap is buried under the dominant
shared component and rounds away. Order is what bought the extra digits.
\end{workedexample}

\section{Loss of orthogonality and re-orthogonalization}
\label{sec:lossorthog}

\Cref{wex:cgsmgs} showed a single step; the effect compounds. As the basis grows and
the input vectors grow more nearly parallel---which is exactly what the Krylov power
basis does---both algorithms accumulate error, and the orthogonality of the computed
basis decays. The standard way to measure the decay is the \emph{loss of
orthogonality},
\[
  \eta_m \;=\; \norm{\mat{Q}_m^{\!\top}\mat{Q}_m - \mat{I}}_2,
\]
the distance of the Gram matrix from the identity after $m$ vectors. For a perfectly
orthonormal basis $\eta_m = 0$; in finite precision it grows, and how fast it grows
is the sharpest distinction among the methods.

\begin{figure}[t]
\centering
\begin{tikzpicture}
\begin{axis}[
  width=11.5cm, height=6.6cm,
  xlabel={basis size $m$},
  ylabel={loss of orthogonality $\eta_m=\norm{\mat{Q}_m^{\!\top}\mat{Q}_m-\mat{I}}$},
  ymode=log, ymin=1e-17, ymax=1e0,
  xmin=0, xmax=30,
  grid=both, grid style={simGray!18},
  legend pos=south east, legend cell align=left,
  legend style={font=\scriptsize, fill=white, draw=simGray!40},
  label style={font=\footnotesize}, tick label style={font=\scriptsize},
  ytick={1e-16,1e-12,1e-8,1e-4,1e0},
]
% CGS: loses orthogonality quadratically in the condition number -> climbs fast to O(1)
\addplot[simRust, very thick, mark=*, mark size=1.1pt] coordinates {
  (1,2e-16)(3,5e-15)(5,8e-13)(7,4e-11)(9,9e-10)(11,3e-8)(13,6e-7)(15,2e-5)
  (17,4e-4)(19,6e-3)(21,5e-2)(23,3e-1)(25,7e-1)(27,9e-1)(29,1.0)};
\addlegendentry{CGS (unstable: $\eta\!\sim\!\kappa^2 u$)}
% MGS: loses orthogonality linearly in kappa -> climbs slowly, stays well below 1
\addplot[simBlue, very thick, mark=square*, mark size=1.1pt] coordinates {
  (1,2e-16)(3,4e-16)(5,9e-16)(7,2e-15)(9,5e-15)(11,1e-14)(13,3e-14)(15,7e-14)
  (17,2e-13)(19,5e-13)(21,1e-12)(23,3e-12)(25,7e-12)(27,1.4e-11)(29,3e-11)};
\addlegendentry{MGS (stable: $\eta\!\sim\!\kappa u$)}
% CGS2 / MGS+reorth: flat at machine epsilon
\addplot[simGreen, very thick, mark=triangle*, mark size=1.3pt] coordinates {
  (1,2e-16)(3,2e-16)(5,3e-16)(7,3e-16)(9,3e-16)(11,4e-16)(13,4e-16)(15,4e-16)
  (17,5e-16)(19,5e-16)(21,5e-16)(23,5e-16)(25,6e-16)(27,6e-16)(29,6e-16)};
\addlegendentry{CGS2 / re-orthogonalized (flat at $u$)}
% machine epsilon reference line
\addplot[simGray, densely dashed, thick, domain=0:30, samples=2] {2.2e-16};
\node[font=\scriptsize, text=simGray, anchor=west] at (axis cs:0.5,5e-16) {machine $u$};
\end{axis}
\end{tikzpicture}
\caption[Loss of orthogonality versus iteration]{How orthogonality decays as the
basis grows, on an ill-conditioned (nearly rank-deficient) input---a schematic of the
classic experiment. The vertical axis is the loss of orthogonality
$\eta_m = \norm{\mat{Q}_m^{\!\top}\mat{Q}_m - \mat{I}}$, log-scaled. \emph{CGS}
(rust) loses orthogonality \emph{quadratically} in the condition number,
$\eta\sim\kappa^2 u$, climbing all the way to $O(1)$---a basis that is no longer
orthogonal at all. \emph{MGS} (blue) loses it only \emph{linearly},
$\eta\sim\kappa u$, staying many orders of magnitude better for the same input.
\emph{CGS2}---classical Gram--Schmidt done twice---and MGS-with-one-reorthogonalization
(green) stay pinned at machine precision $u$ regardless of $m$: full orthogonality,
recovered. The flat green line is the ``twice is enough'' result of
\cref{sec:twiceenough} made visible.}
\label{fig:lossorthog}
\end{figure}

\Cref{fig:lossorthog} is the picture to carry away. On an ill-conditioned input the
three regimes separate cleanly. CGS loses orthogonality \emph{quadratically} in the
condition number $\kappa$ of the input: $\eta \sim \kappa^2 u$, where $u$ is the unit
roundoff. For an input with $\kappa = 10^8$ in double precision
($u \approx 10^{-16}$), that is $\eta \sim 10^{16}\cdot 10^{-16} = 1$---total loss of
orthogonality, exactly what \cref{wex:cgsmgs} foreshadowed compounding over many
steps. MGS loses it only \emph{linearly}, $\eta \sim \kappa u$, which for the same
input is $\eta \sim 10^{-8}$: degraded, but usable. The extra factor of $\kappa$ that
CGS pays is the price of computing every coefficient against the contaminated
original instead of the cleaned residual.

\begin{pitfall}
Even MGS is not safe for the worst inputs. Its loss $\eta \sim \kappa u$ is small only
while $\kappa$ is moderate; for a sufficiently ill-conditioned basis---and a Krylov
power basis becomes arbitrarily ill-conditioned as $m$ grows---$\kappa u$ can itself
reach $O(1)$, and MGS too produces a basis that is no longer orthogonal. The blue
curve in \cref{fig:lossorthog} stays low because the example's $\kappa$ is moderate;
push $\kappa$ high enough and it climbs. The lesson is that \emph{neither} classical
nor modified Gram--Schmidt is unconditionally trustworthy. When orthogonality must be
guaranteed to machine precision---and for Krylov eigensolvers it must---one needs the
re-orthogonalization of the next subsection.
\end{pitfall}

\subsection{``Twice is enough'': CGS2 and re-orthogonalization}
\label{sec:twiceenough}

The repair is disarmingly simple: \emph{do it again}. After computing a residual
$\vw'$ that should be orthogonal to the basis but, in finite precision, retains a
small overlap, run the \emph{whole} Gram--Schmidt step a second time on $\vw'$. The
second pass projects out the small residual overlap that the first pass left behind.
For classical Gram--Schmidt this two-pass version is called \emph{CGS2}; the same
idea applied to MGS is \emph{MGS with re-orthogonalization}.

The remarkable fact, due to Kahan and analyzed by Parlett and others, is that
\emph{one} extra pass suffices: after re-orthogonalizing once, the basis is
orthogonal to machine precision, and a third pass buys nothing. This is the ``twice
is enough'' result, and it is what flattens the green curve in \cref{fig:lossorthog}
at $u$.

\begin{figure}[t]
\centering
\begin{tikzpicture}[font=\footnotesize, scale=1.0, >={Stealth[length=2.4mm]}]
  % the span (a horizontal line in 2D representing span{q1,...})
  \draw[very thick, simGreen] (-0.3,0) -- (5.5,0);
  \node[font=\scriptsize\itshape, text=simGreen, anchor=west] at (4.6,-0.35)
    {$\operatorname{span}\{\vq_1,\dots\}$};
  \coordinate (O) at (0.4,0);
  % input w (well off the span)
  \draw[flow,simRust,very thick] (O) -- (1.0,2.4) node[above,font=\scriptsize]{$\vw$};
  % first pass result: SHOULD be vertical but tilts slightly (residual overlap)
  \coordinate (w1) at (2.6,2.3);
  \draw[flow,simBlue,very thick] (O) -- (w1)
    node[right,font=\scriptsize]{$\vw'$ (pass 1)};
  \node[font=\scriptsize,text=simBlue] at (3.1,1.4) {small residual overlap $\swarrow$};
  % the leftover horizontal component of w1
  \draw[densely dashed, simGray] (w1) -- (2.6,0);
  \draw[<->, simRust, thick] (0.4,-0.55) -- (2.6,-0.55)
    node[midway,below,font=\scriptsize,text=simRust]{$\sim u\kappa$ overlap remains};
  % second pass result: truly vertical
  \coordinate (w2) at (0.4,2.3);
  \draw[flow,simGreen,very thick] (O) -- (w2)
    node[left,font=\scriptsize]{$\vw''$ (pass 2)};
  \draw[simInk] (0.4,0.0) ++(0,0.2) -- ++(0.2,0) -- ++(0,-0.2);
  \node[font=\scriptsize,text=simGreen,anchor=west] at (0.7,2.0)
    {orthogonal to $\sim u$};
  % the correction arrow from w1 to w2
  \draw[->, simViolet, densely dashed, thick] (w1) -- (w2)
    node[midway,above,font=\scriptsize,text=simViolet]{2nd projection removes it};
\end{tikzpicture}
\caption[``Twice is enough'' re-orthogonalization]{Why one extra pass restores
orthogonality. The input $\vw$ (rust) is orthogonalized against the span (green
line). In finite precision the first pass produces $\vw'$ (blue) that is \emph{nearly}
perpendicular but retains a small residual overlap with the span---of relative size
$\sim u\kappa$, the loss of orthogonality. Running the \emph{same} Gram--Schmidt step
a second time on $\vw'$ projects out that small leftover (violet correction),
yielding $\vw''$ (green) orthogonal to machine precision $u$. The Kahan/Parlett
result is that this happens in \emph{one} re-pass: the second residual overlap is of
size $\sim u\cdot u\kappa$, already negligible, so a third pass is unnecessary. Twice
is enough.}
\label{fig:twice}
\end{figure}

\begin{notationbox}
CGS2 returns the \emph{combined} coefficients. The first pass computes a coefficient
column $H$ and a residual $\vw'$; the second pass computes a small \emph{correction}
column $\Delta H$ against $\vw'$ and a twice-projected residual $\vw''$. The operator
returns $H + \Delta H$ (so that $\vw = \vw'' + \sum_j (H_j + \Delta H_j)\vq_j$ still
recovers the input exactly) and the residual $\vw''$. The second pass is \emph{not}
algebraically fusible into the first---it reads the once-orthogonalized $\vw'$, which
did not exist until the first pass ran---so CGS2 genuinely costs two passes of inner
products. That doubled cost, traded against $\kappa^2 u \to u$ orthogonality, is the
CGS2 bargain.
\end{notationbox}

The cost/stability trade is now a clean menu, and it is worth tabulating because the
solver chooses from it at configuration time (\cref{tab:gstradeoff}).

\begin{table}[t]
\centering
\caption[The orthogonalization cost/stability trade]{The orthogonalization menu. The
\emph{communication} column counts global reductions per orthogonalization against an
$m$-vector basis (the cost that dominates on a parallel machine): CGS batches all $m$
inner products into one reduction, MGS needs $m$ separate ones because each waits on
the previous subtraction. The \emph{orthogonality} column is the resulting loss
$\eta\sim\norm{\mat{Q}^{\!\top}\mat{Q}-\mat{I}}$ in terms of the condition number
$\kappa$ and unit roundoff $u$. There is no free lunch: the cheapest-to-communicate
method (CGS) is the least stable, and the most stable (CGS2) pays double.}
\label{tab:gstradeoff}
\small
\setlength{\tabcolsep}{6pt}
\renewcommand{\arraystretch}{1.25}
\begin{tabular}{@{}llll@{}}
\toprule
Method & Passes & Communication & Orthogonality $\eta$ \\
\midrule
CGS  & 1 & 1 reduction (batched) & $\sim \kappa^2 u$ \quad(unstable) \\
MGS  & 1 & $m$ reductions (sequential) & $\sim \kappa\, u$ \quad(stable) \\
CGS2 & 2 & 2 reductions (batched) & $\sim u$ \quad(full) \\
MGS+reorth & 2 & $2m$ reductions & $\sim u$ \quad(full) \\
\bottomrule
\end{tabular}
\end{table}

CGS2 occupies a sweet spot the table makes plain: it recovers full orthogonality
($\eta\sim u$) like MGS-with-reorthogonalization, but because each of its two passes
\emph{batches} its $m$ inner products into a single reduction, it pays only
\emph{two} communications rather than $2m$. On a parallel machine where communication
dominates, CGS2 is often the method of choice precisely because it buys MGS-quality
stability while keeping CGS-style batched communication. The sequential dependency
that made MGS stable is also what makes it communication-bound; CGS2 sidesteps both
by re-running the parallel-friendly classical pass.

\section{Arnoldi: Gram--Schmidt on the Krylov sequence}
\label{sec:arnoldi}

We now assemble the orthogonalizer into the recurrence that drives GMRES. The
\emph{Arnoldi iteration} is nothing more than Gram--Schmidt applied to the Krylov
sequence as it is generated: start from a unit vector $\vq_1$, and at each step form
the next power-basis vector $\mat{A}\vq_j$, orthogonalize it against all the basis
vectors so far, and normalize the residual to get $\vq_{j+1}$.

\begin{lstlisting}[style=l4style]
arnoldi_step :: LinOp -> Basis -> Tensor[$S] -> (Tensor[$S], Tensor[$m+1])
arnoldi_step A V q_j =
  let w        = apply A q_j                    -- next power-basis vector A q_j
      (w', h)  = orthogonalize w V CGS2         -- project off span(V); h = coeffs H[1..m]
      beta     = nrm2 w'                         -- the sub-diagonal entry h_{m+1,m}
      q_next   = scal (1 / beta) w'              -- normalize: the new basis vector
  in (q_next, append h beta)                     -- (new column of V, full Hessenberg column)
\end{lstlisting}

The coefficients produced along the way are not discarded---they are the entries of a
matrix. After $m$ steps, the basis vectors $\vq_1,\dots,\vq_m$ form the columns of an
$N\times m$ matrix $\mat{V}_m$, and the orthogonalization coefficients fill an
$(m+1)\times m$ matrix $\bar{\mat{H}}_m$ that is \emph{upper Hessenberg}: zero below
the first sub-diagonal. The entry $H_{ij}$ is the coefficient of $\vq_i$ when
$\mat{A}\vq_j$ was orthogonalized---which is zero whenever $i > j+1$, because at step
$j$ there were only $j$ basis vectors to project against plus the one new
sub-diagonal normalization. These fit together into the \emph{Arnoldi relation}
\[
  \mat{A}\,\mat{V}_m \;=\; \mat{V}_{m+1}\,\bar{\mat{H}}_m,
\]
which says: applying $\mat{A}$ to the basis and re-expressing the result in the basis
\emph{is} the Hessenberg matrix. \Cref{fig:arnoldi} draws the relation and the shapes
involved.

\begin{figure}[t]
\centering
\begin{tikzpicture}[font=\footnotesize, scale=0.9]
  % A V_m
  \node[font=\small] at (-5.4,0) {$\mat{A}\,\mat{V}_m$};
  \node[font=\small] at (-4.3,0) {$=$};
  % V_{m+1} (tall thin block)
  \draw[thick, fill=simBgGreen] (-3.7,-1.5) rectangle (-2.5,1.5);
  \foreach \x in {-3.45,-3.2,-2.95,-2.7}{\draw[simGreen!60] (\x,-1.5)--(\x,1.5);}
  \node[font=\small] at (-3.1,-2.0) {$\mat{V}_{m+1}$};
  \node[font=\scriptsize,text=simGray] at (-3.1,1.85) {$N\times(m{+}1)$};
  \node[font=\small] at (-2.15,0) {$\cdot$};
  % H-bar (upper Hessenberg: staircase)
  \draw[thick, fill=simBgBlue] (-1.6,-1.5) rectangle (0.4,1.5);
  % shade the zero (below sub-diagonal) region
  \fill[white] (-1.6,-1.5) -- (0.4,1.5) ++(0,-0.6) -- (0.4,-1.5) -- cycle;
  % redraw the staircase boundary of the nonzero (upper Hessenberg) part
  \draw[thick, simBlue]
    (-1.6,1.5) -- (0.4,1.5) -- (0.4,0.9)
    -- (0.0,0.9) -- (0.0,0.3) -- (-0.4,0.3) -- (-0.4,-0.3)
    -- (-0.8,-0.3) -- (-0.8,-0.9) -- (-1.2,-0.9) -- (-1.2,-1.5) -- (-1.6,-1.5) -- cycle;
  \node[font=\small] at (-0.6,-2.0) {$\bar{\mat{H}}_m$};
  \node[font=\scriptsize,text=simGray] at (-0.6,1.85) {$(m{+}1)\times m$};
  \node[font=\scriptsize,text=simBlue] at (-0.55,0.55) {upper};
  \node[font=\scriptsize,text=simBlue] at (-0.55,0.15) {Hessenberg};
  \node[font=\scriptsize,text=simGray] at (-0.05,-1.15) {sub-diag};
  \node[font=\scriptsize,text=simGray] at (-1.15,-0.55) {$=\norm{\vw'}$};
  % annotation on the right
  \node[align=left, font=\scriptsize, text=simGray, anchor=west] at (1.2,0.6)
    {$\bar{\mat{H}}_m[i,j]$ = coefficient of $\vq_i$\\
     when $\mat{A}\vq_j$ was orthogonalized;\\
     the \emph{projection coefficients} $H$\\
     of \cref{sec:gschmidt} stacked column by column.};
  \node[align=left, font=\scriptsize, text=simGreen, anchor=west] at (1.2,-0.7)
    {sub-diagonal $\bar{\mat{H}}_m[j{+}1,j]=\norm{\vw'}$\\
     is the caller's \code{nrm2} normalization;\\
     zeros below it $\Rightarrow$ Hessenberg shape.};
\end{tikzpicture}
\caption[Arnoldi produces $\mat{V}$ and an upper-Hessenberg $\bar{\mat{H}}$]{The
Arnoldi relation $\mat{A}\mat{V}_m = \mat{V}_{m+1}\bar{\mat{H}}_m$, the output of
running Gram--Schmidt on the Krylov sequence. The basis vectors fill the tall thin
$\mat{V}_{m+1}$ (green); the orthogonalization coefficients fill the small
\emph{upper-Hessenberg} $\bar{\mat{H}}_m$ (blue)---zero below the first sub-diagonal,
because at step $j$ only the first $j$ basis vectors and one new normalization
contribute. Each column of $\bar{\mat{H}}_m$ is exactly the $(\vw', H)$ coefficient
pair of \cref{sec:gschmidt}: the entries $H_{ij}=\inner{\mat{A}\vq_j}{\vq_i}$ above
the diagonal, and the sub-diagonal entry the residual norm $\norm{\vw'}$ from the
caller's normalization. GMRES (\cref{ch:gmres}) reads this small Hessenberg matrix to
solve its least-squares problem; the Arnoldi eigensolvers of \cref{ch:krylovschur}
read its eigenvalues. The orthogonalizer of this chapter is the engine that fills
both matrices.}
\label{fig:arnoldi}
\end{figure}

This is precisely the \code{orthogonalize} step GMRES (\cref{ch:gmres}) calls once
per inner iteration, and the upper-Hessenberg matrix it fills is the small dense
matrix on which GMRES solves its least-squares problem. The Arnoldi eigensolvers of
\cref{ch:krylovschur} use the very same recurrence and read $\bar{\mat{H}}_m$'s
\emph{eigenvalues} (the Ritz values) as approximations to $\mat{A}$'s. Either way,
the orthogonalizer is the engine: feed it the Krylov sequence, and it produces both
the orthonormal basis and the Hessenberg matrix that is the projection of $\mat{A}$
onto that basis.

\section{Lanczos: the three-term recurrence for symmetric operators}
\label{sec:lanczos}

Now specialize to a \emph{symmetric} operator, $\mat{A} = \mat{A}^{\!\top}$ (the case
of the eigenmode analysis and of conjugate gradients). Something remarkable happens
to Arnoldi: almost all of the orthogonalization coefficients turn out to be zero in
advance, and the long Gram--Schmidt sum collapses to a \emph{three-term recurrence}.

\subsection{Why symmetry shortens the recurrence}

The Hessenberg matrix $\bar{\mat{H}}_m$ of Arnoldi is the projection of $\mat{A}$ onto
the Krylov basis: $H_{ij} = \inner{\mat{A}\vq_j}{\vq_i}$. If $\mat{A}$ is symmetric,
then so is this projection,
\[
  H_{ij} = \inner{\mat{A}\vq_j}{\vq_i} = \inner{\vq_j}{\mat{A}\vq_i}
         = \inner{\mat{A}\vq_i}{\vq_j} = H_{ji},
\]
where the middle equality is the symmetry of $\mat{A}$. But $\bar{\mat{H}}_m$ is
\emph{already} upper Hessenberg---zero below the first sub-diagonal. A matrix that is
both upper Hessenberg \emph{and} symmetric is \emph{tridiagonal}: zero everywhere
except the diagonal and the two adjacent off-diagonals. So in the symmetric case the
Hessenberg matrix is a tridiagonal matrix $\mat{T}_m$, and the Arnoldi
orthogonalization of $\mat{A}\vq_j$ has \emph{only three nonzero coefficients}: against
$\vq_{j-1}$, against $\vq_j$, and the new sub-diagonal normalization for $\vq_{j+1}$.
Every projection against $\vq_1,\dots,\vq_{j-2}$ is zero, automatically.

\begin{keyidea}
For a symmetric operator $\mat{A}$, the Arnoldi Hessenberg matrix is symmetric, and a
symmetric Hessenberg matrix is \emph{tridiagonal}. So orthogonalizing
$\mat{A}\vq_j$---which in general requires projecting against \emph{all} $j$ previous
basis vectors---requires projecting against only the \emph{two} most recent ones,
$\vq_j$ and $\vq_{j-1}$; the projections against all earlier vectors are exactly zero.
The long Gram--Schmidt sum collapses to a three-term recurrence. This is the same
miracle that gives conjugate gradients its short recurrence (\cref{ch:cg}): symmetry
makes a quantity that would otherwise depend on the whole history depend on only the
last two steps.
\end{keyidea}

The recurrence is the \emph{Lanczos iteration}. Writing the three coefficients as
$\alpha_j$ (diagonal), $\beta_{j-1}$ (the sub-diagonal coupling to the previous
vector), and $\beta_j$ (the new normalization), the step is
\[
  \vw \;=\; \mat{A}\vq_j - \alpha_j\,\vq_j - \beta_{j-1}\,\vq_{j-1},
  \qquad
  \alpha_j = \inner{\mat{A}\vq_j}{\vq_j},
  \qquad
  \beta_j = \norm{\vw},
  \qquad
  \vq_{j+1} = \frac{\vw}{\beta_j}.
\]
Only $\vq_j$ and $\vq_{j-1}$ appear---the two previous vectors, nothing earlier. In
code:

\begin{lstlisting}[style=l4style]
lanczos_step :: LinOp -> Tensor[$S] -> Tensor[$S] -> Scalar
                      -> (Tensor[$S], Scalar, Scalar)
lanczos_step A q_j q_prev beta_prev =
  let u      = apply A q_j                       -- A q_j
      alpha  = dot u q_j                          -- diagonal coeff <A q_j, q_j>
      w      = axpby (-alpha) q_j 1 u             -- subtract the q_j component
      w'     = axpy (-beta_prev) q_prev w         -- subtract the q_{j-1} component (only TWO!)
      beta   = nrm2 w'                            -- new sub-diagonal / normalization
      q_next = scal (1/beta) w'
  in (q_next, alpha, beta)
\end{lstlisting}

The contrast with \code{arnoldi\_step} is the whole point: Arnoldi's
\code{orthogonalize w V} projects against the \emph{entire} basis \code{V}, a sum
that grows with $m$; Lanczos projects against exactly two stored vectors,
$\vq_j$ and $\vq_{j-1}$, a fixed-size step. \Cref{fig:lanczos} draws the difference,
and the tridiagonal matrix the recurrence fills.

\begin{figure}[t]
\centering
\begin{tikzpicture}[font=\footnotesize, scale=0.95]
  % ===== LEFT: the three-term recurrence, only two previous vectors =====
  \begin{scope}
    \node[font=\small\bfseries, text=simGreen] at (1.4,2.9) {three-term recurrence};
    % the stored vectors
    \node[data, minimum width=11mm] (qm1) at (-0.4,1.6) {$\vq_{j-1}$};
    \node[data, minimum width=11mm] (qj)  at (1.4,1.6) {$\vq_j$};
    \node[stage, minimum width=14mm] (Aqj) at (1.4,0.5) {$\mat{A}\vq_j$};
    \draw[flow] (qj) -- (Aqj);
    \node[product, minimum width=11mm] (qn) at (1.4,-0.8) {$\vq_{j+1}$};
    % only the two recent vectors feed the new one
    \draw[flow,simGreen] (qj.south west) .. controls +(-0.2,-0.6) .. (qn.north)
      node[pos=0.5,left,font=\scriptsize,text=simGreen]{$-\alpha_j$};
    \draw[flow,simGreen] (qm1.south) .. controls +(0,-1.4) .. (qn.west)
      node[pos=0.3,left,font=\scriptsize,text=simGreen]{$-\beta_{j-1}$};
    \draw[flow] (Aqj) -- (qn);
    % the OLD vectors that are NOT needed
    \node[data, minimum width=9mm, draw=simGray!40, text=simGray!60] (old) at (-1.7,1.6) {$\vq_{j-2}\dots$};
    \draw[refedge, simRust] (old.south) .. controls +(0,-1.8) .. (qn.north west);
    \node[font=\scriptsize, text=simRust] at (-1.5,-0.5) {coeff $=0$:};
    \node[font=\scriptsize, text=simRust] at (-1.5,-0.85) {not needed};
    \node[font=\scriptsize\itshape, text=simGray, align=center] at (0.4,-1.7)
      {only $\vq_j,\vq_{j-1}$ stored:\\$O(1)$ memory per step};
  \end{scope}
  \draw[simGray!50, thick] (3.2,-1.9) -- (3.2,2.9);
  % ===== RIGHT: the tridiagonal T =====
  \begin{scope}[xshift=5.0cm]
    \node[font=\small\bfseries, text=simBlue] at (1.5,2.9) {tridiagonal $\mat{T}_m$};
    % draw a 5x5 tridiagonal pattern
    \def\cell{0.55}
    \foreach \r in {0,1,2,3,4}{
      \foreach \c in {0,1,2,3,4}{
        \pgfmathtruncatemacro\diff{abs(\r-\c)}
        \ifnum\diff<2
          \ifnum\r=\c
            \node[fill=simBgGold, draw=simGold, minimum size=\cell cm, font=\scriptsize]
              at (\c*\cell, -\r*\cell) {$\alpha_{\the\numexpr\r+1\relax}$};
          \else
            \node[fill=simBgBlue, draw=simBlue, minimum size=\cell cm, font=\scriptsize]
              at (\c*\cell, -\r*\cell) {$\beta$};
          \fi
        \else
          \node[draw=simGray!25, minimum size=\cell cm, font=\scriptsize, text=simGray!50]
            at (\c*\cell, -\r*\cell) {$0$};
        \fi
      }
    }
    \node[font=\scriptsize\itshape, text=simGray, align=center] at (1.1,-3.1)
      {symmetric + Hessenberg\\$=$ tridiagonal: $\alpha$ diag, $\beta$ off-diag};
  \end{scope}
\end{tikzpicture}
\caption[The Lanczos three-term recurrence and its tridiagonal matrix]{The Lanczos
specialization of Arnoldi for symmetric $\mat{A}$. \emph{Left:} the new basis vector
$\vq_{j+1}$ is built from $\mat{A}\vq_j$ by subtracting components along only the
\emph{two} most recent vectors, $\vq_j$ (coefficient $-\alpha_j$) and $\vq_{j-1}$
(coefficient $-\beta_{j-1}$). All earlier vectors $\vq_{j-2},\dots$ have coefficient
exactly zero (rust, ``not needed''), so they need not be stored: Lanczos uses $O(1)$
memory per step where Arnoldi stores the whole growing basis. \emph{Right:} the
coefficients fill a \emph{tridiagonal} matrix $\mat{T}_m$---the diagonal carries the
$\alpha_j = \inner{\mat{A}\vq_j}{\vq_j}$, the two off-diagonals carry the
$\beta_j = \norm{\vw}$, and everything else is zero. This is Arnoldi's Hessenberg
matrix made tridiagonal by symmetry. The short recurrence is the same economy
symmetry buys conjugate gradients (\cref{ch:cg}).}
\label{fig:lanczos}
\end{figure}

\subsection{A worked Lanczos: two steps on a small symmetric matrix}

To make the recurrence concrete, run it by hand on a $3\times 3$ symmetric matrix and
watch the tridiagonal $\mat{T}$ assemble.

\begin{workedexample}{Two Lanczos steps building a tridiagonal $\mat{T}$}{lanczos}
Take the symmetric matrix and starting vector
\[
  \mat{A} = \begin{psmallmatrix} 2 & 1 & 0 \\ 1 & 2 & 1 \\ 0 & 1 & 2 \end{psmallmatrix},
  \qquad
  \vq_1 = \begin{psmallmatrix} 1 \\ 0 \\ 0 \end{psmallmatrix}
  \quad(\text{already unit length}).
\]

\emph{Step 1.} Form $\mat{A}\vq_1$, the diagonal coefficient $\alpha_1$, and the
residual:
\[
  \mat{A}\vq_1 = \begin{psmallmatrix} 2 \\ 1 \\ 0 \end{psmallmatrix},
  \qquad
  \alpha_1 = \inner{\mat{A}\vq_1}{\vq_1} = 2 .
\]
There is no $\vq_0$, so the recurrence has only the $\alpha_1$ term on the first step:
\[
  \vw = \mat{A}\vq_1 - \alpha_1\vq_1
      = \begin{psmallmatrix} 2 \\ 1 \\ 0 \end{psmallmatrix}
      - 2\begin{psmallmatrix} 1 \\ 0 \\ 0 \end{psmallmatrix}
      = \begin{psmallmatrix} 0 \\ 1 \\ 0 \end{psmallmatrix},
  \qquad
  \beta_1 = \norm{\vw} = 1,
  \qquad
  \vq_2 = \frac{\vw}{\beta_1} = \begin{psmallmatrix} 0 \\ 1 \\ 0 \end{psmallmatrix}.
\]
Check orthogonality: $\inner{\vq_2}{\vq_1} = 0$. Good.

\emph{Step 2.} Now the full three-term step, using $\vq_2$ and the previous $\vq_1$
with its coupling $\beta_1 = 1$:
\[
  \mat{A}\vq_2 = \begin{psmallmatrix} 1 \\ 2 \\ 1 \end{psmallmatrix},
  \qquad
  \alpha_2 = \inner{\mat{A}\vq_2}{\vq_2} = 2 .
\]
Subtract \emph{both} the $\vq_2$ component (coefficient $\alpha_2$) and the $\vq_1$
component (coefficient $\beta_1$):
\[
  \vw = \mat{A}\vq_2 - \alpha_2\vq_2 - \beta_1\vq_1
      = \begin{psmallmatrix} 1 \\ 2 \\ 1 \end{psmallmatrix}
      - 2\begin{psmallmatrix} 0 \\ 1 \\ 0 \end{psmallmatrix}
      - 1\begin{psmallmatrix} 1 \\ 0 \\ 0 \end{psmallmatrix}
      = \begin{psmallmatrix} 0 \\ 0 \\ 1 \end{psmallmatrix},
  \qquad
  \beta_2 = \norm{\vw} = 1,
  \qquad
  \vq_3 = \begin{psmallmatrix} 0 \\ 0 \\ 1 \end{psmallmatrix}.
\]
Check: $\inner{\vq_3}{\vq_2} = 0$ and $\inner{\vq_3}{\vq_1} = 0$---the second of these
\emph{automatically}, even though we never projected against $\vq_1$ in step 2. That
is the symmetry miracle in action: the $\vq_1$ component of $\mat{A}\vq_2$ was already
exactly accounted for by the $\beta_1\vq_1$ term, so no further projection against
$\vq_1$ (or any earlier vector) was needed.

\emph{The assembled matrix.} The coefficients fill the tridiagonal
\[
  \mat{T}_2 = \begin{psmallmatrix} \alpha_1 & \beta_1 \\ \beta_1 & \alpha_2 \end{psmallmatrix}
            = \begin{psmallmatrix} 2 & 1 \\ 1 & 2 \end{psmallmatrix},
\]
and after the third step $\mat{T}_3 = \mathrm{tridiag}(\beta;\,\alpha;\,\beta)$ with
$\alpha = (2,2,2)$, $\beta = (1,1)$---which for this particular $\mat{A}$ is just
$\mat{A}$ itself expressed in the Lanczos basis, since the Krylov subspace happened to
fill all of $\Real^3$ in three steps. The eigenvalues of $\mat{T}_3$ are therefore
exactly those of $\mat{A}$, namely $2$ and $2\pm\sqrt2$---which is the basis on which
the eigensolvers of \cref{ch:krylovschur} read approximate eigenvalues off the small
tridiagonal $\mat{T}_m$ for $m < N$, where the agreement is approximate rather than
exact.
\end{workedexample}

\subsection{The catastrophe: Lanczos loses orthogonality}

The three-term recurrence is beautiful and economical, and in finite precision it is
treacherous. The short recurrence is short precisely because it \emph{trusts} that
the projections against $\vq_1,\dots,\vq_{j-2}$ are zero and so skips them. In exact
arithmetic they are zero. In floating point they are not---and because Lanczos never
re-checks them, the small errors are never corrected. Orthogonality against the
\emph{old} vectors is lost, and lost fast.

\begin{pitfall}
Lanczos loses orthogonality far more violently than Gram--Schmidt, and for a reason
peculiar to the short recurrence: it never looks at the old basis vectors, so it
cannot notice (let alone repair) the orthogonality it loses against them. The
practical symptom is \emph{ghost eigenvalues}. As orthogonality decays, the Lanczos
basis silently develops near-duplicate directions, and the tridiagonal $\mat{T}_m$
reports the \emph{same} eigenvalue several times over---spurious copies of converged
eigenvalues that are artifacts of the lost orthogonality, not real eigenvalues of
$\mat{A}$. A naive Lanczos eigensolver will confidently return a multiplicity-three
eigenvalue that has multiplicity one. The cure is a re-orthogonalization
\emph{policy}---full re-orthogonalization against the whole stored basis (expensive,
defeating the $O(1)$-memory point), or selective/partial re-orthogonalization that
re-orthogonalizes only when a monitored loss-of-orthogonality estimate crosses a
threshold. Which policy a solver chooses, and how the practical Krylov--Schur method
of \cref{ch:krylovschur} manages it, is a central design question of the eigensolver.
\end{pitfall}

This is the trade Lanczos forces. Arnoldi, with its full orthogonalization, is
expensive (storing and projecting against the whole basis) but stable. Lanczos, with
its three-term recurrence, is cheap ($O(1)$ memory, two projections) but unstable, and
buying back the stability with re-orthogonalization gives back some of the cost it
saved. The eigensolver chapters live inside exactly this tension.

\section{The orthogonality contract}
\label{sec:contract}

Step back from the four algorithms and notice what they all \emph{deliver}. CGS, MGS,
CGS2, Arnoldi, Lanczos---each is a different schedule, a different cost, a different
stability, a different memory footprint. But every one of them is in the business of
producing the same thing: a basis that satisfies
\[
  \inner{\vq_i}{\vq_j} \approx \delta_{ij}
  \qquad\text{to a stated tolerance,}
\]
and a coefficient matrix ($\bar{\mat{H}}$ or $\mat{T}$) that records the projection of
the operator onto that basis. This is the \emph{orthogonality contract}, and it is the
single interface every downstream method programs against.

\begin{keyidea}
\textbf{The orthogonality contract.} Every orthogonalization method in this chapter
guarantees one thing: an orthonormal basis $\mat{Q}$ with
$\mat{Q}^{\!\top}\mat{Q} \approx \mat{I}$ to a tolerance, together with the
coefficients of the projection it performed. The methods differ \emph{only} in how
well they hold the contract and at what cost---CGS cheaply but loosely, MGS stably but
sequentially, CGS2 fully but at double cost, Lanczos cheaply but fragilely. They are
therefore \emph{substitutable}: a caller that depends only on the contract (``give me
an orthonormal basis'') can swap one for another without changing its own code. This
is exactly the algorithmic-substitution that \cref{ch:rotations} calls a genuine
rotation---one slot, \lstinline[style=l4style]@orthog := mgs | cgs | cgs2@, three
different algorithms behind one interface.
\end{keyidea}

That substitutability is not a coincidence of notation; it is the contract doing its
job. GMRES does not care \emph{how} its basis was orthogonalized---it cares that the
basis is orthonormal so that its least-squares problem is well-conditioned. An Arnoldi
eigensolver does not care which Gram--Schmidt variant filled the Hessenberg
matrix---it cares that the Ritz values it reads are not corrupted by a basis that
secretly lost orthogonality. The contract is what lets the solver authors reason about
their algorithm \emph{above} the orthogonalizer, treating it as a black box with a
guarantee, while the orthogonalizer authors reason about stability and cost
\emph{below} the interface. And it is what makes ``did orthogonality hold?'' the first
question to ask when a Krylov method misbehaves: a stalled GMRES or a ghost eigenvalue
is, far more often than not, a contract quietly violated by a basis that lost its
orthogonality and a method that never noticed.

\summarysection
A Krylov solver works inside $\mathcal{K}_m(\mat{A},\vv)$ but never in its natural
power basis, which collapses toward the dominant eigenvector and is hopelessly
ill-conditioned; it works in an \emph{orthonormal} basis, built incrementally by
\emph{Gram--Schmidt}: subtract off the projections onto the previous basis vectors,
then normalize. Gram--Schmidt has two faces. \emph{Classical} (CGS) computes every
projection coefficient against the same original vector---independent, batchable,
parallelizable, but numerically unstable, losing orthogonality as $\kappa^2 u$.
\emph{Modified} (MGS) subtracts the projections one at a time against the running
residual---stable, losing orthogonality only as $\kappa u$, but at the price of a
genuine \emph{sequential dependency}: the column order is a loop-carried dependency, a
concrete specimen of the obstruction of \cref{ch:rotations}, and exactly the
\code{map}-versus-\code{fold} split of \cref{ch:combinators}. For the worst inputs
even MGS is not enough, and \emph{re-orthogonalization} (CGS2, ``twice is enough'') one
extra pass restores full machine-precision orthogonality. Applied to the Krylov
sequence, Gram--Schmidt becomes \emph{Arnoldi}, producing the orthonormal $\mat{V}_m$
and the upper-Hessenberg $\bar{\mat{H}}_m$ that GMRES (\cref{ch:gmres}) and the
Arnoldi eigensolvers (\cref{ch:krylovschur}) consume. For \emph{symmetric} operators
Arnoldi collapses to the \emph{Lanczos} three-term recurrence---the Hessenberg matrix
becomes tridiagonal, the orthogonalization needs only the two most recent vectors, the
same short-recurrence miracle that powers conjugate gradients (\cref{ch:cg})---but in
finite precision Lanczos loses orthogonality catastrophically, producing ghost
eigenvalues and demanding a re-orthogonalization policy. Underneath all of it is one
\emph{orthogonality contract}: an orthonormal basis to a tolerance, plus the
projection coefficients. The methods differ only in how well and how cheaply they hold
it, which makes them substitutable behind one interface---the genuine rotation of
\cref{ch:rotations}.

\furtherreading
The definitive textbook development of orthogonalization, conditioning, and the
stability of Gram--Schmidt is Trefethen and Bau's \emph{Numerical Linear
Algebra}~\citep{trefethenbau1997}, whose lectures on QR factorization and on the
sensitivity of the classical versus modified algorithms are the ideal companion to
\cref{sec:cgsmgs,sec:lossorthog}. For the full algebraic and floating-point analysis---
the $\kappa^2 u$ versus $\kappa u$ loss-of-orthogonality bounds, the ``twice is
enough'' re-orthogonalization result, and the Arnoldi and Lanczos recurrences with
their reorthogonalization variants---Golub and Van Loan's \emph{Matrix
Computations}~\citep{golubvanloan2013} is the standard reference. For the eigenvalue
side specifically---the ghost-eigenvalue phenomenon, selective and partial
reorthogonalization, and the practical Lanczos and Arnoldi eigensolvers that
\cref{ch:krylovschur} builds on---Saad's \emph{Numerical Methods for Large Eigenvalue
Problems}~\citep{saad2011eig} is the place to go deeper. Read all three with the
\emph{orthogonality contract} of \cref{sec:contract} in mind: each is, at bottom, a
study of how well that contract can be held and at what cost.

\exercisessection
\begin{enumerate}[nosep]
  \item \emph{(Geometry.)} In $\Real^3$, let $\vq_1 = (1,0,0)^{\!\top}$,
  $\vq_2 = (0,1,0)^{\!\top}$, and $\vw = (3, 4, 5)^{\!\top}$. Compute the
  Gram--Schmidt residual $\vw' = \vw - \inner{\vw}{\vq_1}\vq_1 -
  \inner{\vw}{\vq_2}\vq_2$ and verify by direct inner product that
  $\inner{\vw'}{\vq_1} = \inner{\vw'}{\vq_2} = 0$ (the contract of
  \cref{lem:residorthog}). What is $\vq_3$ after normalization?
  \item \emph{(CGS versus MGS.)} Repeat \cref{wex:cgsmgs} with $\varepsilon = 10^{-3}$
  and confirm that the residual orthogonality is $\sim\varepsilon$ for CGS and
  $\sim\varepsilon^2$ for MGS. Then explain in one sentence why MGS's second
  projection coefficient $H_2$ was nonzero while CGS's was (numerically) zero, in
  terms of \emph{which vector each coefficient was taken against}.
  \item \emph{(The sequential obstruction.)} Explain why classical Gram--Schmidt can
  batch its $m$ inner products into a single matrix-vector product
  $\mat{Q}^{\!\top}\vw$ but modified Gram--Schmidt cannot. Which combinator
  (\cref{ch:combinators}) is each---\code{map} or \code{fold}? Why does permuting the
  columns of $\mat{Q}$ change the MGS result (at the bit level) but not the CGS result?
  \item \emph{(Twice is enough.)} A first Gram--Schmidt pass leaves a residual with a
  relative overlap of $\sim u\kappa$ against the span. Argue heuristically why a second
  pass reduces the overlap to $\sim u\cdot(u\kappa)$ and hence why a third pass is
  unnecessary. Under what condition on $\kappa$ does even \emph{one} pass already
  suffice (so that CGS2 is overkill)?
  \item \emph{(Arnoldi shape.)} The Arnoldi relation is
  $\mat{A}\mat{V}_m = \mat{V}_{m+1}\bar{\mat{H}}_m$. Give the dimensions of each of the
  three matrices, and explain why $\bar{\mat{H}}_m$ has exactly one more row than
  column and is zero below its first sub-diagonal. Which entry of $\bar{\mat{H}}_m$ is
  the residual norm $\norm{\vw'}$ from the normalization step, and which chapter
  (\cref{ch:gmres} or \cref{ch:krylovschur}) reads which feature of this matrix?
  \item \emph{(Lanczos by hand.)} Run two Lanczos steps as in \cref{wex:lanczos} on
  $\mat{A} = \begin{psmallmatrix} 4 & 1 & 0 \\ 1 & 4 & 1 \\ 0 & 1 & 4 \end{psmallmatrix}$
  with $\vq_1 = (1,0,0)^{\!\top}$. Compute $\alpha_1, \beta_1, \vq_2, \alpha_2,
  \beta_2, \vq_3$, assemble $\mat{T}_2$, and verify that $\vq_3$ is automatically
  orthogonal to $\vq_1$ even though you never projected against $\vq_1$ in step~2.
  \item \emph{(Why symmetry shortens the recurrence.)} Prove that the Arnoldi
  Hessenberg matrix is symmetric when $\mat{A}$ is symmetric, and conclude that it is
  tridiagonal. Then state precisely which orthogonalization coefficients Lanczos is
  entitled to skip, and explain---in terms of finite-precision arithmetic---why
  skipping them is exactly what causes Lanczos to lose orthogonality and produce ghost
  eigenvalues.
  \item \emph{(The contract as an interface.)} \cref{ch:rotations} calls
  \lstinline[style=l4style]@orthog := mgs | cgs | cgs2@ a genuine rotation by the
  algorithmic-substitution test. Identify the contract that all three satisfy and
  argue that a GMRES solver written against that contract is correct regardless of
  which variant is installed. Give one concrete way the \emph{choice} of variant could
  still affect the solver's observable behavior, even though all three satisfy the
  contract. (Hint: \cref{sec:lossorthog}.)
\end{enumerate}

\chapter{The Least-Squares Core: Givens Rotations and the Running QR}
\label{ch:givens}

% local vector macros for the running-QR right-hand side (g) and Hessenberg column (h)
\providecommand{\vg}{\mathbf{g}}
\providecommand{\vh}{\mathbf{h}}

\chapterorient{GMRES, the engine of \cref{ch:gmres}, leaves a small chore at the
end of every step: minimize $\norm{\beta\,\vec e_1 - \bar{\mat H}_m\vy}$ over a tiny
least-squares problem whose matrix grows by one column each iteration. This chapter
is about that chore---and about the surprisingly elegant way the solver does it. We
do not re-solve the least-squares problem from scratch at each step; we maintain a
\emph{running} QR factorization, updated one column at a time by a stream of $2\times2$
plane rotations. The stream is a fold over the arriving columns, in the exact sense of
\cref{ch:combinators}. And it has a gift: after the rotation hits the right-hand side,
the GMRES residual norm is sitting there, already computed, for free. No
$\mat A\vx_m$, no residual evaluation---a unitary byproduct of triangularizing the
matrix. By the end you will be able to generate a Givens rotation by hand, run the
stream on a small Hessenberg, read the residual off one entry, and back-solve for the
answer: the entire inner loop of GMRES, made concrete.}

\section{The chore GMRES leaves behind}
\label{sec:gv-setup}

Recall the shape of GMRES from \cref{ch:gmres}. At step $m$ the Arnoldi process of
\cref{ch:orthog} has built an orthonormal basis $\mat V_m = [\vv_1,\dots,\vv_m]$ for
the Krylov subspace and, as a byproduct of orthogonalizing each new direction, an
$(m{+}1)\times m$ matrix $\bar{\mat H}_m$ that records the orthogonalization
coefficients. This matrix is \emph{upper-Hessenberg}: upper-triangular except for a
single nonzero sub-diagonal, the band of entries $h_{k+1,k}$ just below the main
diagonal. GMRES then forms its iterate as $\vx_m = \vx_0 + \mat V_m\vy$, choosing the
coordinate vector $\vy\in\Real^m$ that minimizes the residual. Because Arnoldi makes
the residual norm equal to a small projected quantity, that minimization collapses to
\begin{equation}
\label{eq:gv-lsq}
  \min_{\vy\in\Real^m}\ \norm{\beta\,\vec e_1 - \bar{\mat H}_m\,\vy},
\end{equation}
where $\beta = \norm{\vr_0}$ is the initial residual norm and $\vec e_1$ is the first
standard basis vector. This is a \emph{least-squares} problem---more equations
($m{+}1$) than unknowns ($m$)---and it is \emph{small}: its size is the iteration
count, never the dimension $N$ of the field. A GMRES run that converges in 40 steps
solves a $41\times 40$ least-squares problem, whatever the millions of unknowns in
$\vx$.

\Cref{eq:gv-lsq} is the chore. We could, at every step, assemble the current
$\bar{\mat H}_m$, factor it afresh, and solve. That is correct and wasteful: the
matrix at step $m{+}1$ is the matrix at step $m$ with one column glued on the right
and one row glued on the bottom; almost all the work of factoring it was already done
last step. The professional move is to maintain a factorization \emph{incrementally},
spending only the cost of the one new column each step. That incremental
factorization is the subject of this chapter, and the tool that builds it is the
plane rotation.

\begin{figure}[t]
\centering
\begin{tikzpicture}[font=\footnotesize]
  % the Hessenberg with one column arriving and one sub-diagonal to kill
  \def\cw{0.62} \def\rh{0.52}
  % grid 5 rows x 4 cols
  \foreach \c in {0,1,2,3}{
    \foreach \r in {0,...,4}{
      \draw[simGray!35] (\c*\cw,-\r*\rh) rectangle ++(\cw,-\rh);
    }
  }
  % triangular block already done (cols 0..2): shade upper-tri
  \foreach \c/\maxr in {0/1,1/2,2/3}{
    \foreach \r in {0,...,\maxr}{
      \fill[simBgBlue] (\c*\cw,-\r*\rh) rectangle ++(\cw,-\rh);
    }
  }
  % the arriving column (col 3): shade gold, full height down to its sub-diagonal r=4
  \foreach \r in {0,...,4}{
    \fill[simBgGold] (3*\cw,-\r*\rh) rectangle ++(\cw,-\rh);
  }
  % redraw gridlines over fills
  \foreach \c in {0,1,2,3}{\foreach \r in {0,...,4}{
      \draw[simGray!35] (\c*\cw,-\r*\rh) rectangle ++(\cw,-\rh);}}
  % mark the sub-diagonal entry to be annihilated in the new column: (row 4, col 3)
  \draw[simRust, very thick] (3*\cw,-4*\rh) rectangle ++(\cw,-\rh);
  \node[font=\scriptsize, text=simRust] at (3*\cw+0.5*\cw,-4*\rh-0.5*\rh) {$h_{j+1,j}$};
  \node[font=\scriptsize] at (3*\cw+0.5*\cw,-3*\rh-0.5*\rh) {$h_{j,j}$};
  % bracket labels
  \node[font=\scriptsize, text=simBlue, align=center] at (1*\cw,0.55) {already triangular\\(cols $0..j{-}1$)};
  \node[font=\scriptsize, text=simGold!70!black, align=center] at (3*\cw+0.5*\cw,0.55) {new\\column $j$};
  % the rotation arrow: zero the rust cell against the one above it
  \draw[-{Stealth[length=2mm]}, simRust, thick]
    (3*\cw+\cw+0.15,-4*\rh-0.5*\rh) .. controls +(0.7,0.2) and +(0.7,-0.2) ..
    (3*\cw+\cw+0.15,-3*\rh-0.5*\rh);
  \node[font=\scriptsize, text=simRust, anchor=west] at (3*\cw+\cw+0.55,-3.5*\rh-0.5*\rh)
    {$G_j$ rotates};
  \node[font=\scriptsize, text=simRust, anchor=west] at (3*\cw+\cw+0.55,-4.0*\rh-0.5*\rh)
    {$h_{j+1,j}\!\to\!0$};
\end{tikzpicture}
\caption[A Givens rotation zeroing one sub-diagonal entry]{The geometry of one step.
The upper-Hessenberg matrix $\bar{\mat H}_m$ is triangular in its leading columns
(blue), with a single nonzero sub-diagonal. A new column $j$ arrives from Arnoldi
(gold) carrying one fresh sub-diagonal entry $h_{j+1,j}$ (rust outline). A single
$2\times2$ plane rotation $G_j$, mixing rows $j$ and $j{+}1$, annihilates that entry
against the diagonal above it---completing the triangularization of column $j$. One
rotation, one sub-diagonal entry: this is the atom the rest of the chapter assembles
into a stream.}
\label{fig:gv-hessenberg}
\end{figure}

\section{Givens rotations: zeroing one entry at a time}
\label{sec:gv-rotations}

A \emph{Givens rotation}, or plane rotation, is a $2\times2$ orthogonal matrix
\begin{equation}
\label{eq:gv-matrix}
  G(c,s) = \begin{pmatrix} c & s \\ -s & c \end{pmatrix},
  \qquad c^2 + s^2 = 1,
\end{equation}
that rotates the plane by the angle whose cosine is $c$ and sine is $s$. Its single
job in this chapter is surgical: given a 2-vector $\begin{psmallmatrix} a \\ b
\end{psmallmatrix}$, choose $c$ and $s$ so that applying $G$ sends the second
component to zero,
\begin{equation}
\label{eq:gv-goal}
  \begin{pmatrix} c & s \\ -s & c \end{pmatrix}
  \begin{pmatrix} a \\ b \end{pmatrix}
  = \begin{pmatrix} r \\ 0 \end{pmatrix}.
\end{equation}
Because $G$ is orthogonal it preserves length, so the surviving first component must
carry the whole norm: $r = \sqrt{a^2 + b^2}$. That single observation \emph{is} the
construction. The bottom row of \cref{eq:gv-goal} reads $-s\,a + c\,b = 0$, and
together with $r = \sqrt{a^2+b^2}$ the top row reads $c\,a + s\,b = r$. Solving the
pair gives the rotation coefficients directly.

\begin{lemma}[Givens generation]
\label{lem:gv-generate}
For a real 2-vector $(a,b)$ with $r = \sqrt{a^2 + b^2} > 0$, the rotation $G(c,s)$
with
\begin{equation}
\label{eq:gv-cs}
  c = \frac{a}{r}, \qquad s = \frac{b}{r}
\end{equation}
satisfies $G(c,s)\begin{psmallmatrix} a \\ b \end{psmallmatrix} =
\begin{psmallmatrix} r \\ 0 \end{psmallmatrix}$, and $c^2 + s^2 = 1$.
\end{lemma}

\begin{proof}
The orthogonality identity is immediate: $c^2 + s^2 = (a^2 + b^2)/r^2 = 1$. For the
action, compute both rows. The top row gives $c\,a + s\,b = (a^2 + b^2)/r = r^2/r = r$.
The bottom row gives $-s\,a + c\,b = (-b\,a + a\,b)/r = 0$. Hence the image is exactly
$(r,0)^{\!\top}$.
\end{proof}

We call the construction of $(c,s)$ from $(a,b)$ the verb \code{givens\_generate}, and
the in-place action of a stored $(c,s)$ on a pair of scalars the verb
\code{givens\_apply}. The two are the entire kernel; everything later in the chapter is
a way of \emph{scheduling} calls to them.

\begin{lstlisting}[style=l4style]
givens_generate :: (a: Scalar, b: Scalar) -> (c: Real, s: Scalar)
givens_apply    :: (c: Real, s: Scalar) -> ((dx: Scalar, dy: Scalar) -> (Scalar, Scalar))
givens_apply c s (dx, dy) = (c*dx + s*dy, -conj s * dx + c*dy)
\end{lstlisting}

The action of \code{givens\_apply} is the general orthogonal update $\begin{psmallmatrix}
dx' \\ dy' \end{psmallmatrix} = \begin{psmallmatrix} c & s \\ -\bar s & c
\end{psmallmatrix}\begin{psmallmatrix} dx \\ dy \end{psmallmatrix}$, which over the
reals is \cref{eq:gv-matrix} exactly; the conjugate $\bar s$ appears only in the
complex case, of which more below. Note the \emph{closure-returning} signature of
\code{givens\_apply}: it takes the rotation $(c,s)$ and \emph{returns a function} of the
pair, written with the paren-grouping of \cref{ch:combinators}. This is deliberate. A
generated rotation is applied many times---to its own column, to later columns, to the
right-hand side---so it is natural to read \code{givens\_apply}~$c$~$s$ as ``the
operator that rotates a pair,'' a first-class value stored and reused.

\begin{notationbox}
Throughout this chapter $(c,s)$ is one stored rotation; $a$ and $b$ are the two
components of the 2-vector it is built to annihilate (here $a$ is the diagonal, $b$ the
sub-diagonal entry to be zeroed); $r = \sqrt{a^2+b^2}$ is the surviving norm. We index
the rotation generated at Arnoldi step $j$ as $(c_j, s_j)$, and we write a generic
applied pair as $(dx, dy)$. The cosine $c$ is always real; the sine $s$ is real for a
real system and complex for a complex one.
\end{notationbox}

\subsection{The stable formula}
\label{sec:gv-stable}

\Cref{eq:gv-cs} is correct but dangerous to compute literally. Forming $r =
\sqrt{a^2 + b^2}$ squares each component, and if $a$ or $b$ is larger than the square
root of the floating-point overflow threshold---around $10^{154}$ in double
precision---the intermediate $a^2$ overflows to infinity even though $r$ itself is a
perfectly representable number. The symmetric failure happens for tiny values, which
underflow to zero. Production libraries therefore never square before scaling. They
factor out the larger component first:
\begin{equation}
\label{eq:gv-scaled}
  t = \frac{\min(\abs a,\abs b)}{\max(\abs a,\abs b)}, \qquad
  r = \max(\abs a,\abs b)\,\sqrt{1 + t^2},
\end{equation}
so the square root sees an argument between $1$ and $2$ and the division $t$ is between
$0$ and $1$. Mathematically this is identical to $r = \sqrt{a^2+b^2}$; numerically it
never overflows unless the true answer does. Palace follows the LAPACK convention
(the algorithm of Bindel, Demmel, Kahan, and Marques), computing $c = \abs a / r$ and
fixing the sign of $r$ to match $a$ so that the diagonal stays positive. We treat the
scaling as a \emph{load-bearing numerical detail}, not a transparent rewrite of
\cref{eq:gv-cs}: it changes which inputs the algorithm survives, and the precise scaled
form is recorded in \cref{app:numerics}.

\begin{pitfall}
Do not implement \cref{eq:gv-cs} as written. The literal formula
$r=\sqrt{a^2+b^2}$ overflows for large inputs and underflows for small ones, returning
$\infty$ or $0$ where the scaled form \eqref{eq:gv-scaled} returns the right finite
$r$. This is not a micro-optimization; it is a correctness fix at the edges of the
floating-point range, and it is exactly the kind of ``load-bearing numerical trick''
\cref{ch:bigidea} warned must be preserved as part of the algorithm rather than
abstracted away.
\end{pitfall}

\begin{workedexample}{Generating a Givens rotation by hand}{gv-generate}
Build the rotation that annihilates the second component of $(a,b) = (3,4)$, and check
its action.

\emph{The norm.} $r = \sqrt{3^2 + 4^2} = \sqrt{9 + 16} = \sqrt{25} = 5$. (The
$(3,4,5)$ triangle keeps the arithmetic clean.)

\emph{The coefficients.} By \cref{eq:gv-cs},
\[
  c = \frac{a}{r} = \frac{3}{5} = 0.6, \qquad
  s = \frac{b}{r} = \frac{4}{5} = 0.8.
\]
Sanity check the orthogonality identity: $c^2 + s^2 = 0.36 + 0.64 = 1$. Good---this is
a genuine rotation.

\emph{The action.} Apply $G(0.6,0.8)$ to $(3,4)$ using \code{givens\_apply}:
\[
  \begin{pmatrix} 0.6 & 0.8 \\ -0.8 & 0.6 \end{pmatrix}
  \begin{pmatrix} 3 \\ 4 \end{pmatrix}
  = \begin{pmatrix} 0.6\cdot 3 + 0.8\cdot 4 \\ -0.8\cdot 3 + 0.6\cdot 4 \end{pmatrix}
  = \begin{pmatrix} 1.8 + 3.2 \\ -2.4 + 2.4 \end{pmatrix}
  = \begin{pmatrix} 5 \\ 0 \end{pmatrix}.
\]
The second component is exactly zero and the first is exactly $r = 5$: the whole length
of the original 2-vector survived in the surviving component, as orthogonality
demands. The scaled formula \eqref{eq:gv-scaled} would compute the same $(c,s)=(0.6,0.8)$
via $\max=4$, $t = 3/4$, $r = 4\sqrt{1 + 9/16} = 4\cdot(5/4) = 5$---identical answer,
no squaring of large numbers.
\end{workedexample}

\subsection{The complex case, briefly}
\label{sec:gv-complex}

For a complex linear system the Hessenberg entries are complex, and the rotation must
preserve the Hermitian (conjugate) norm. The convention Palace and LAPACK use keeps the
cosine $c$ \emph{real} and lets the sine $s$ be complex, with the unitarity identity
becoming $c^2 + \abs s^2 = 1$. The action of \code{givens\_apply} carries the conjugate
in the lower row, $dy' = -\bar s\,dx + c\,dy$, which is why we wrote it that way above.
Nothing structural changes: the generation still produces a $(c,s)$ that zeros the
sub-diagonal, the stream still triangularizes, and the residual still reads off the
right-hand side. The real-cosine convention is a tidy bookkeeping choice---the rotation
register $c$ is one real number per step, $s$ one complex number---and we will not dwell
on it. The driven and eigenmode analyses (\cref{part:modalities}) supply the complex
systems that exercise it.

\section{The running QR stream}
\label{sec:gv-stream}

Now we schedule the rotations. The factorization we maintain is $\bar{\mat H}_m =
\mat Q_m \bar{\mat R}_m$, where $\mat Q_m$ is the accumulated product of all the plane
rotations and $\bar{\mat R}_m$ is upper-triangular. We never form $\mat Q_m$ as a
matrix; we store only the list of rotation pairs $(c_0,s_0),(c_1,s_1),\dots$, which is
all of $O(m)$ scalars. Each Arnoldi step delivers one new Hessenberg column, and the
stream processes it in four sub-steps.

\begin{description}[style=nextline, leftmargin=0pt, font=\normalfont\itshape]
  \item[(i) Replay.] The new column $\vh = \bar{\mat H}_m[:,j]$ arrives in the
  \emph{original} (un-rotated) frame. The previous columns were triangularized by the
  stored rotations $(c_0,s_0),\dots,(c_{j-1},s_{j-1})$; to bring the new column into the
  same frame we replay all of them, in order, against its entries:
  for $k=0,1,\dots,j{-}1$, apply $(c_k,s_k)$ to the pair $(h_k,h_{k+1})$.
  \item[(ii) Generate.] After replay, the entry $h_{j+1}$ on the sub-diagonal is the
  one fresh nonzero below the diagonal. Generate the new rotation $(c_j,s_j) =
  \code{givens\_generate}(h_j, h_{j+1})$ to annihilate it (\cref{lem:gv-generate}).
  \item[(iii) Apply to the column.] Apply the freshly generated $(c_j,s_j)$ to its own
  pair $(h_j, h_{j+1})$. This sets $h_{j+1} \gets 0$ and $h_j \gets r$: column $j$ is
  now upper-triangular, and $\bar{\mat R}_m$ has grown by a column.
  \item[(iv) Apply to the right-hand side.] Apply the \emph{same} $(c_j,s_j)$ to the
  right-hand-side pair $(g_j, g_{j+1})$, where $\vg$ began life as $\beta\,\vec e_1$.
  This is the step that pays the dividend, and \cref{sec:gv-residual} collects it.
\end{description}

\begin{figure}[t]
\centering
\begin{tikzpicture}[font=\footnotesize, node distance=5mm]
  \node[data, minimum width=18mm] (col) at (-0.2,0) {new column\\$\vh = \bar{\mat H}[:,j]$};
  \node[opbox, minimum width=20mm, right=8mm of col] (replay)
    {(i) \textbf{replay}\\\scriptsize $(c_0,s_0)\dots(c_{j-1},s_{j-1})$};
  \node[opbox, minimum width=20mm, right=7mm of replay, fill=simBgGold, draw=simGold] (gen)
    {(ii) \textbf{generate}\\\scriptsize $(c_j,s_j)$};
  \node[opbox, minimum width=18mm, below=9mm of gen] (apcol)
    {(iii) \textbf{apply}\\\scriptsize to column $\to \bar{\mat R}$};
  \node[opbox, minimum width=18mm, left=10mm of apcol] (aprhs)
    {(iv) \textbf{apply-rhs}\\\scriptsize to $\vg=\beta\vec e_1$};
  \node[product, minimum width=16mm, left=9mm of aprhs] (res)
    {residual\\$|g_{j+1}|$};
  \draw[flow] (col) -- (replay);
  \draw[flow] (replay) -- (gen);
  \draw[flow] (gen.south) -- (apcol.north);
  \draw[flow] (gen.south) .. controls +(0,-0.9) and +(0,0.9) .. (aprhs.north);
  \draw[flow] (aprhs) -- (res);
  % store-and-grow note for the new rotation
  \node[font=\scriptsize\itshape, text=simGreen, align=center] at (gen.north)[yshift=6mm]
    {append $(c_j,s_j)$\\to the stored list};
  \draw[refedge, simGreen] (gen.north) -- ++(0,4mm);
\end{tikzpicture}
\caption[The running-QR stream as a four-step pipeline]{The four sub-steps the stream
runs on each arriving Hessenberg column. \emph{Replay} (i) applies all stored rotations
to the new column, bringing it into the triangularized frame; \emph{generate} (ii)
produces one new rotation to kill the fresh sub-diagonal entry and appends it to the
stored list; \emph{apply} (iii) uses it to triangularize the column, growing
$\bar{\mat R}$; \emph{apply-rhs} (iv) uses the same rotation on the right-hand side
$\vg$, and the new last entry $|g_{j+1}|$ is the residual norm (\cref{sec:gv-residual}).
The pipeline runs once per Arnoldi step and the stored list grows by one rotation each
time.}
\label{fig:gv-pipeline}
\end{figure}

\begin{lstlisting}[style=l4style]
-- one column of the running QR (Arnoldi step j); see Figure (the four-step pipeline)
ls_update_column :: LsqState -> (h: HessCol) -> LsqState
ls_update_column st h =
  let h1       = foldl replayOne h (zip st.cs st.sn)      -- (i) replay stored 0..j-1
      (cj, sj) = givens_generate (h1 !! j) (h1 !! (j+1))  -- (ii) generate new rotation
      h2       = givens_apply cj sj (h1 !! j, h1 !! (j+1)) -- (iii) triangularize column j
      g'       = applyRhs cj sj st.g                       -- (iv) rotate the RHS pair
  in st { R = st.R `appendCol` h2          -- R grows by a column
        , cs = st.cs ++ [cj], sn = st.sn ++ [sj]  -- rotation list grows by one
        , g  = g'                          -- residual now lives in g'[j+1]
        }
  where replayOne col (ck, sk) = ...   -- givens_apply on pair (k, k+1)
\end{lstlisting}

\subsection{The stream is a fold over columns}
\label{sec:gv-fold}

Look at the signature of \code{ls\_update\_column}: it takes a state and a column and
returns a new state. That is exactly the shape of a fold step from \cref{ch:combinators}.
The whole running QR is
\begin{equation}
\label{eq:gv-fold}
\begin{aligned}
  \code{running\_qr}\ st_0\ [\vh_0, \dots, \vh_{m-1}]
  \;=\;\ &\code{foldl}\ \code{ls\_update\_column}\ st_0\\
  &[\vh_0, \vh_1, \dots, \vh_{m-1}],
\end{aligned}
\end{equation}
a left fold of \code{ls\_update\_column} over the stream of Hessenberg columns, with the
carry $st$ holding the triangular factor $\bar{\mat R}$, the growing rotation list
$(\code{cs},\code{sn})$, and the rotated right-hand side $\vg$. This is not a metaphor.
The carry threads strictly forward---column $j$ cannot be processed before column
$j{-}1$, because the replay sub-step (i) needs all the rotations the earlier columns
generated---so by the analysis of \cref{ch:combinators} the stream is an
\emph{intrinsically sequential} fold, not a parallel map. The rotation list is the part
of the carry that grows: one $(c_j,s_j)$ appended per step, so after $m$ columns the
carry holds $m$ rotations (\cref{fig:gv-growing}).

\begin{figure}[t]
\centering
\begin{tikzpicture}[font=\footnotesize]
  % three successive states of the carry, rotation list growing
  \foreach \step/\x/\n in {0/0/1, 1/4.2/2, 2/8.4/3}{
    \node[font=\scriptsize\bfseries, text=simBlue] at (\x+0.6,2.0) {after column $\step$};
    % stored rotation list as stacked cells
    \foreach \i in {0,...,\numexpr\n-1\relax}{
      \node[opbox, minimum width=14mm, minimum height=5mm] at (\x+0.6,1.4-0.62*\i)
        {$(c_{\i},s_{\i})$};
    }
  }
  % arrows between states: one column arrives, one rotation appended
  \draw[flow] (1.7,0.9) -- node[above,font=\scriptsize]{$\vh_1$} node[below,font=\scriptsize,text=simGreen]{$+\,(c_1,s_1)$} (3.0,0.9);
  \draw[flow] (5.9,0.9) -- node[above,font=\scriptsize]{$\vh_2$} node[below,font=\scriptsize,text=simGreen]{$+\,(c_2,s_2)$} (7.2,0.9);
  \node[font=\scriptsize\itshape, text=simGray, align=center] at (4.8,-1.2)
    {each Arnoldi column appends exactly one rotation to the carry:\\the fold's accumulator
     grows by one cell per step};
\end{tikzpicture}
\caption[The rotation list grows per Arnoldi column]{The stored rotation list is the
growing part of the fold's carry. Each Arnoldi column drives one
\code{ls\_update\_column} step, which appends exactly one new pair $(c_j,s_j)$ to the
list. After processing $m$ columns the carry holds $m$ rotations---an $O(m)$ store, far
cheaper than the explicit $\mat Q_m$ it stands in for. Because step $j$ replays all
earlier rotations, the list \emph{must} be threaded forward: this is a fold, not a map.}
\label{fig:gv-growing}
\end{figure}

\begin{pitfall}
The replay must come \emph{before} the generate, and the order of the replayed
rotations matters: $(c_0,s_0)$ then $(c_1,s_1)$ then $\dots$, never permuted. Plane
rotations do not commute---each one is built in the frame the previous ones
established---so applying them out of order, or generating the new rotation against an
un-replayed column, triangularizes against the wrong diagonal and silently corrupts
$\bar{\mat R}$. This is the running QR's load-bearing invariant: \emph{replay all stored
rotations in order, then generate.} It is the reason the carry is threaded forward and
the reason the stream is a fold rather than a map.
\end{pitfall}

\section{The residual, for free}
\label{sec:gv-residual}

Here is the payoff that makes the whole construction worth building. After sub-step
(iv)---applying the new rotation to the right-hand side $\vg$---the magnitude of the
\emph{new last entry} of $\vg$ is exactly the GMRES residual norm:
\begin{equation}
\label{eq:gv-free}
  \norm{\beta\,\vec e_1 - \bar{\mat H}_m\,\vy_m}
  \;=\; \abs{g_{j+1}},
\end{equation}
where $\vy_m$ is the (not-yet-computed) least-squares minimizer. We get the residual
without forming $\vy_m$, without forming the iterate $\vx_m = \vx_0 + \mat V_m\vy_m$,
and---most importantly---without the expensive matrix--vector product $\mat A\vx_m$ that
a literal residual evaluation would demand. The convergence test of GMRES is therefore a
single absolute value of a scalar the stream already computed.

Why is \cref{eq:gv-free} true? Because every rotation is orthogonal, and orthogonal maps
preserve length. Let $\mat Q_m$ be the accumulated rotation product, so that $\mat Q_m
\bar{\mat H}_m = \bar{\mat R}_m$ is upper-triangular and $\vg = \mat Q_m(\beta\,\vec
e_1)$. The least-squares residual is invariant under the orthogonal $\mat Q_m$:
\begin{equation}
\label{eq:gv-residual-derivation}
  \norm{\beta\,\vec e_1 - \bar{\mat H}_m\vy}
  = \norm{\mat Q_m\!\left(\beta\,\vec e_1 - \bar{\mat H}_m\vy\right)}
  = \norm{\vg - \bar{\mat R}_m\vy}.
\end{equation}
Now $\bar{\mat R}_m$ is $(m{+}1)\times m$ and upper-triangular: its last row is entirely
zero. So whatever $\vy$ we pick, the bottom equation $0 = g_{j+1} - 0$ can never be
satisfied, and the best we can do is zero out the top $m$ rows (which a nonsingular
triangular solve does exactly), leaving the residual equal to the one unkillable
entry $\abs{g_{j+1}}$. Splitting the norm into the part a triangular solve can zero and
the part it cannot:
\begin{equation}
\label{eq:gv-split}
  \norm{\vg - \bar{\mat R}_m\vy}^2
  = \underbrace{\norm{\vg_{0:m} - \mat R_m\vy}^2}_{\text{drive to }0\text{ by back-solve}}
  + \underbrace{\abs{g_{j+1}}^2}_{\text{the residual}}.
\end{equation}
The minimizer kills the first term, and the second term---the magnitude of the rotated
right-hand side's tail---is the residual. That is \cref{eq:gv-free}. The rotation stream,
built to triangularize the matrix, drags the right-hand side along and concentrates the
entire un-fittable residual into one entry, where we can read it.

\begin{figure}[t]
\centering
\begin{tikzpicture}[font=\footnotesize]
  % the rotated RHS vector g, last entry highlighted as the residual
  \def\rh{0.5}
  \node[font=\scriptsize\bfseries, text=simBlue] at (-1.6,0.3) {$\vg = \mat Q_m(\beta\vec e_1)$};
  \foreach \i/\lbl in {0/{g_0},1/{g_1},2/{\vdots},3/{g_m}}{
    \draw[simGray!45] (0,-\i*\rh) rectangle ++(0.9,-\rh);
    \node[font=\scriptsize] at (0.45,-\i*\rh-0.5*\rh) {$\lbl$};
  }
  % the residual entry g_{j+1}
  \draw[simRust, very thick] (0,-4*\rh) rectangle ++(0.9,-\rh);
  \fill[simBgRust] (0,-4*\rh) rectangle ++(0.9,-\rh);
  \draw[simRust, very thick] (0,-4*\rh) rectangle ++(0.9,-\rh);
  \node[font=\scriptsize, text=simRust] at (0.45,-4*\rh-0.5*\rh) {$g_{j+1}$};
  % brace + label
  \draw[decorate, decoration={brace, amplitude=4pt}, simBlue]
    (1.05,-0) -- (1.05,-4*\rh) node[midway, right=3pt, font=\scriptsize, text=simBlue, align=left]
    {fitted part:\\driven to $0$\\by back-solve};
  \draw[decorate, decoration={brace, amplitude=4pt, mirror}, simRust]
    (-0.05,-4*\rh) -- (-0.05,-5*\rh);
  \node[font=\scriptsize, text=simRust, anchor=east, align=right] at (-0.15,-4.5*\rh)
    {the residual\\$\norm{\vr_m} = |g_{j+1}|$\\\emph{free!}};
  % the "no Ax needed" annotation
  \node[extern, minimum width=30mm, align=center] at (5.0,-1.0)
    {\textbf{not computed:}\\$\vx_m = \vx_0+\mat V_m\vy_m$\\$\mat A\vx_m$, \;$\vb - \mat A\vx_m$};
  \draw[-{Stealth[length=2mm]}, simGray, densely dotted] (4.0,-2.6) -- (1.4,-2.3);
  \node[font=\scriptsize\itshape, text=simGray] at (4.6,-2.8) {avoided every step};
\end{tikzpicture}
\caption[The residual is a free byproduct]{Reading the residual off the rotated
right-hand side. After the rotation stream acts on $\vg = \mat Q_m(\beta\vec e_1)$, the
upper $m$ entries are the fittable part---a back-solve can drive $\norm{\vg_{0:m} -
\mat R_m\vy}$ to zero---and the single last entry $g_{j+1}$ (rust) is the part no
choice of $\vy$ can fit. Its magnitude \emph{is} the GMRES residual norm, by
\cref{eq:gv-split}. The expensive alternative on the right---forming the iterate
$\vx_m$, multiplying by $\mat A$, and subtracting from $\vb$---is never needed at all.
The residual is a unitary byproduct of triangularizing the matrix.}
\label{fig:gv-residual}
\end{figure}

\begin{keyidea}
The least-squares residual is a \emph{free byproduct} of the running QR. Each plane
rotation is orthogonal, so applying the rotation stream to the right-hand side $\beta\,
\vec e_1$ preserves its norm and concentrates the entire un-fittable component into one
entry. After the right-hand-side update, the magnitude of that entry, $\abs{g_{j+1}}$,
\emph{is} the GMRES residual norm $\norm{\beta\,\vec e_1 - \bar{\mat H}_m\vy_m}$---no
iterate $\vx_m$ formed, no matrix--vector product $\mat A\vx_m$ computed, no residual
evaluated. This is why GMRES can monitor convergence at essentially zero cost: the
quantity the convergence test needs falls out of the triangularization the solver was
already doing.
\end{keyidea}

\section{The back-solve: the terminal consumer}
\label{sec:gv-backsolve}

The stream runs every Arnoldi step, cheaply, and tells us when to stop. But it never
forms the answer. When we finally want the actual coordinate vector $\vy_m$---at
convergence, or at a restart boundary where GMRES discards the basis and begins
again---we run the \emph{back-solve} once. The triangularized matrix has made this easy:
$\mat R_m$ (the leading $m\times m$ block of $\bar{\mat R}_m$) is upper-triangular, so
$\mat R_m\vy = \vg_{0:m}$ is solved by back-substitution, bottom row up:
\begin{equation}
\label{eq:gv-backsub}
  y_i = \frac{1}{R_{ii}}\!\left(g_i - \sum_{k>i} R_{ik}\,y_k\right),
  \qquad i = m{-}1, m{-}2, \dots, 0.
\end{equation}
Each $y_i$ uses only the $y_k$ already computed below it; the last unknown $y_{m-1} =
g_{m-1}/R_{m-1,m-1}$ comes first and the rest follow upward. This is an $O(m^2)$ scalar
solve on the small dense triangle---trivial beside one field-sized matrix--vector
product. Once $\vy_m$ is in hand, the iterate correction is reconstructed against the
Arnoldi basis,
\begin{equation}
\label{eq:gv-reconstruct}
  \vx_m = \vx_0 + \mat V_m\,\vy_m = \vx_0 + \sum_{k=0}^{m-1} y_k\,\vv_k,
\end{equation}
a single \code{linear\_combination} (\cref{ch:combinators}) over the basis vectors. For
flexible GMRES (FGMRES), the only change is the basis: the correction is reconstructed
against the preconditioned basis $\mat Z_m$ rather than $\mat V_m$, so
$\vx_m = \vx_0 + \mat Z_m\vy_m$. The running-QR stream, the rotations, and the
back-solve are bit-for-bit identical between GMRES and FGMRES; only the basis the final
sum reads differs.

\begin{figure}[t]
\centering
\begin{tikzpicture}[font=\footnotesize]
  % upper-triangular R with back-substitution sweep arrow
  \def\cw{0.66} \def\rh{0.56}
  \foreach \c in {0,1,2,3}{\foreach \r in {0,...,3}{
      \draw[simGray!30] (\c*\cw,-\r*\rh) rectangle ++(\cw,-\rh);}}
  % shade upper triangle
  \foreach \c/\minr in {0/0,1/0,2/0,3/0}{
    \foreach \r in {0,...,\c}{
      \fill[simBgBlue] (\c*\cw,-\r*\rh) rectangle ++(\cw,-\rh);}}
  \foreach \c in {0,1,2,3}{\foreach \r in {0,...,3}{
      \draw[simGray!30] (\c*\cw,-\r*\rh) rectangle ++(\cw,-\rh);}}
  % diagonal entries
  \foreach \i in {0,1,2,3}{\node[font=\scriptsize] at (\i*\cw+0.5*\cw,-\i*\rh-0.5*\rh) {$R_{\i\i}$};}
  \node[font=\scriptsize, text=simGray] at (2.2*\cw,0.30) {$\mat R_m$ (upper-tri)};
  % RHS vector g
  \node[font=\scriptsize] at (4.7*\cw,0.30) {$\vg_{0:m}$};
  \foreach \r in {0,1,2,3}{
    \draw[simGray!45] (4.4*\cw,-\r*\rh) rectangle ++(0.8*\cw,-\rh);
    \node[font=\scriptsize] at (4.8*\cw,-\r*\rh-0.5*\rh) {$g_{\r}$};}
  % y vector
  \node[font=\scriptsize] at (6.2*\cw,0.30) {$\vy$};
  \foreach \r in {0,1,2,3}{
    \draw[simGreen!50] (5.9*\cw,-\r*\rh) rectangle ++(0.8*\cw,-\rh);
    \node[font=\scriptsize, text=simGreen] at (6.3*\cw,-\r*\rh-0.5*\rh) {$y_{\r}$};}
  % the bottom-up sweep arrow
  \draw[-{Stealth[length=2.4mm]}, simRust, very thick] (6.3*\cw,-3.7*\rh) -- (6.3*\cw,-0.2*\rh);
  \node[font=\scriptsize, text=simRust, anchor=west, align=left] at (7.0*\cw,-2*\rh)
    {back-substitute\\\emph{bottom row up}:\\$y_{m-1}$ first,\\then sweep up};
\end{tikzpicture}
\caption[The back-solve as the terminal consumer]{The back-solve, run once at
convergence or restart. The triangularized $\mat R_m$ (left, upper-triangular) and the
fitted part $\vg_{0:m}$ of the rotated right-hand side feed a back-substitution
(\cref{eq:gv-backsub}) that sweeps from the bottom row upward, computing $y_{m-1}$
first and each higher $y_i$ from the ones already found. The result $\vy$ (green)
reconstructs the iterate $\vx_m = \vx_0 + \mat V_m\vy$ via one
\code{linear\_combination}. This is the stream's terminal \emph{consumer}: it runs
once per restart cycle, not once per step.}
\label{fig:gv-backsolve}
\end{figure}

\subsection{Producer and consumer: two cadences}
\label{sec:gv-producer-consumer}

The chapter's architecture is a clean split between a streaming \emph{producer} and a
terminal \emph{consumer}, and the two run at different cadences. The producer is the
running-QR stream of \cref{sec:gv-stream}: it fires \emph{every Arnoldi step}, consuming
one column, appending one rotation, and emitting one residual scalar. The consumer is the
back-solve of \cref{sec:gv-backsolve}: it fires \emph{once per restart cycle}, at the
moment we actually want the answer. \Cref{fig:gv-prodcons} draws the two cadences
side by side.

This split is not incidental; it is why the running QR is worth maintaining. If we
solved the least-squares problem from scratch each step, every step would pay the
back-solve cost \emph{and} the re-factorization cost. By streaming the factorization and
deferring the back-solve, each step pays only the cheap incremental rotation
work---which also happens to hand us the residual for free---and the one back-solve is
amortized over the whole restart cycle. The producer monitors; the consumer materializes;
and the convergence test that decides when to call the consumer reads the producer's free
byproduct.

\begin{figure}[t]
\centering
\begin{tikzpicture}[font=\footnotesize]
  % producer: a chain of per-step stream updates, each emitting beta
  \node[font=\small\bfseries, text=simGreen] at (3.0,2.3) {producer: the stream (every step)};
  \foreach \j/\x in {0/0,1/2.0,2/4.0,3/6.0}{
    \node[opbox, minimum width=15mm, fill=simBgGreen, draw=simGreen] (u\j) at (\x,1.3)
      {step \j};
    \node[font=\scriptsize, text=simRust] at (\x,0.55) {$|g_{j{+}1}|$};
    \draw[-{Stealth[length=1.6mm]}, simRust] (\x,1.0) -- (\x,0.75);
  }
  \foreach \j/\jn in {0/1,1/2,2/3}{
    \draw[flow, simGreen] (u\j.east) -- (u\jn.west);}
  \node[font=\scriptsize\itshape, text=simGreen, anchor=west] at (6.9,1.3) {$\dots$};
  % the convergence test reading the byproduct
  \node[font=\scriptsize\itshape, text=simRust, align=center] at (3.0,0.05)
    {convergence test reads $|g_{j+1}|$ at every step (free)};
  % consumer: the back-solve, fired once
  \node[font=\small\bfseries, text=simBlue] at (3.0,-1.2) {consumer: the back-solve (once at restart)};
  \node[product, minimum width=24mm] (bs) at (3.0,-2.1)
    {back-solve $\to \vy_m \to \vx_m$};
  \draw[-{Stealth[length=2mm]}, simBlue, densely dashed] (u3.south) .. controls +(0,-1.3) and +(1.5,1.0) .. (bs.east)
    node[midway, right, font=\scriptsize, text=simBlue] {fire \emph{once}};
\end{tikzpicture}
\caption[Producer versus consumer: two cadences]{The two cadences of the least-squares
core. The \emph{producer} (top, green) is the running-QR stream: it runs once per
Arnoldi step, threading the rotation list forward and emitting the residual $|g_{j+1}|$
each step, which the convergence test reads for free. The \emph{consumer} (bottom, blue)
is the back-solve: it fires \emph{once}, at convergence or a restart boundary, to
materialize $\vy_m$ and reconstruct the iterate $\vx_m$. Streaming the factorization and
deferring the back-solve is the whole reason the running QR beats re-solving from
scratch: the per-step cost is cheap, and the one expensive consumer is amortized.}
\label{fig:gv-prodcons}
\end{figure}

\section{The whole inner solve, worked}
\label{sec:gv-worked}

We now run the entire machine on a concrete small problem: two stream steps, the
residual read off for free, and a back-solve for the answer. This is the GMRES inner
solve in miniature.

\begin{workedexample}{Streaming QR on a $3\times2$ Hessenberg, with the free residual}{gv-stream}
Take $\beta = 5$ and the $3\times2$ upper-Hessenberg matrix that Arnoldi might have
produced after two steps:
\[
  \bar{\mat H}_2 = \begin{pmatrix} 3 & 1 \\ 4 & 5 \\ 0 & 12 \end{pmatrix},
  \qquad \vg^{(0)} = \beta\,\vec e_1 = \begin{pmatrix} 5 \\ 0 \\ 0 \end{pmatrix}.
\]
We want $\min_\vy\norm{\beta\vec e_1 - \bar{\mat H}_2\vy}$, and we want the residual for
free along the way.

\medskip\emph{Step 0 --- the first column} $\vh_0 = (3,4,0)^{\!\top}$.
There are no stored rotations yet, so the \emph{replay} sub-step is empty. \emph{Generate}
from the pair $(h_0, h_1) = (3,4)$: this is exactly \cref{wex:gv-generate}, giving
$(c_0, s_0) = (0.6, 0.8)$ and $r = 5$. \emph{Apply} to the column zeros the
sub-diagonal: $(3,4)\mapsto(5,0)$, so column 0 of $\mat R$ is $R_{00} = 5$. \emph{Apply
to the RHS} pair $(g_0, g_1) = (5,0)$:
\[
  \begin{pmatrix} 0.6 & 0.8 \\ -0.8 & 0.6 \end{pmatrix}\!\begin{pmatrix} 5 \\ 0 \end{pmatrix}
  = \begin{pmatrix} 3 \\ -4 \end{pmatrix},
  \quad\text{so}\quad \vg^{(1)} = (3,\,-4,\,0)^{\!\top}.
\]
The residual after one step is $\abs{g_1} = \abs{-4} = 4$. (Not yet converged; carry on.)

\medskip\emph{Step 1 --- the second column} $\vh_1 = (1,5,12)^{\!\top}$.
Now there \emph{is} a stored rotation, so we \emph{replay} $(c_0,s_0) = (0.6,0.8)$ on the
top pair $(h_0, h_1) = (1,5)$:
\[
  \begin{pmatrix} 0.6 & 0.8 \\ -0.8 & 0.6 \end{pmatrix}\!\begin{pmatrix} 1 \\ 5 \end{pmatrix}
  = \begin{pmatrix} 0.6 + 4 \\ -0.8 + 3 \end{pmatrix}
  = \begin{pmatrix} 4.6 \\ 2.2 \end{pmatrix},
\]
leaving the replayed column $(4.6,\,2.2,\,12)^{\!\top}$ (the bottom entry $12$ is
untouched by the replay---rotation 0 only mixes rows 0 and 1). The off-diagonal
$R_{01} = 4.6$ is now fixed. \emph{Generate} from the new sub-diagonal pair
$(h_1, h_2) = (2.2, 12)$:
\[
\begin{aligned}
  r &= \sqrt{2.2^2 + 12^2} = \sqrt{4.84 + 144} = \sqrt{148.84} = 12.2,\\
  c_1 &= \frac{2.2}{12.2} \approx 0.1803,\qquad s_1 = \frac{12}{12.2} \approx 0.9836.
\end{aligned}
\]
\emph{Apply} to the column zeros the sub-diagonal: $(2.2,12)\mapsto(12.2,0)$, so
$R_{11} = 12.2$. \emph{Apply to the RHS} pair $(g_1, g_2) = (-4, 0)$:
\[
  \begin{pmatrix} c_1 & s_1 \\ -s_1 & c_1 \end{pmatrix}\!\begin{pmatrix} -4 \\ 0 \end{pmatrix}
  = \begin{pmatrix} 0.1803\cdot(-4) \\ -0.9836\cdot(-4) \end{pmatrix}
  \approx \begin{pmatrix} -0.721 \\ 3.934 \end{pmatrix},
\]
so $\vg^{(2)} \approx (3,\,-0.721,\,3.934)^{\!\top}$.

\medskip\emph{Read the residual, free.} The new last entry is
$\abs{g_2} \approx 3.934$. By \cref{eq:gv-free} this \emph{is} the least-squares
residual norm $\min_\vy\norm{\beta\vec e_1 - \bar{\mat H}_2\vy}$---and we have computed
no iterate and no $\mat A\vx$. (Cross-check: the rotations are unitary, so the norm of
$\vg$ is preserved: $\norm{\vg^{(2)}}^2 \approx 3^2 + 0.721^2 + 3.934^2 \approx 9 +
0.520 + 15.48 = 25.0 = \beta^2$. The length is conserved, as it must be.)

\medskip\emph{Back-solve for} $\vy$. The fitted part is the $2\times2$ triangular system
\[
  \mat R_2 = \begin{pmatrix} 5 & 4.6 \\ 0 & 12.2 \end{pmatrix},
  \qquad \vg_{0:2} = \begin{pmatrix} 3 \\ -0.721 \end{pmatrix}.
\]
Back-substitute bottom-up (\cref{eq:gv-backsub}). The last unknown first:
\[
  y_1 = \frac{g_1}{R_{11}} = \frac{-0.721}{12.2} \approx -0.0591.
\]
Then sweep up:
\[
  y_0 = \frac{1}{R_{00}}\big(g_0 - R_{01}\,y_1\big)
      = \frac{1}{5}\big(3 - 4.6\cdot(-0.0591)\big)
      = \frac{1}{5}\big(3 + 0.272\big) \approx 0.654.
\]
So $\vy_2 \approx (0.654,\,-0.0591)^{\!\top}$, and the GMRES iterate is $\vx_2 = \vx_0 +
y_0\,\vv_1 + y_1\,\vv_2$---one \code{linear\_combination} over the two Arnoldi basis
vectors. The residual we read off for free, $\approx 3.934$, is the norm we would
measure if we formed that $\vx_2$ and computed $\norm{\vb - \mat A\vx_2}$---but we never
had to.
\end{workedexample}

\Cref{wex:gv-stream} is the entire GMRES inner solve at $m=2$: two stream steps building
$\mat R$ one column at a time, the residual read off the rotated right-hand side without
any field-sized work, and a single back-solve producing the coordinates. Scale $m$ from
$2$ to $40$ and nothing changes but the size of the small triangle; the cost of the
field-sized operations---the Arnoldi matrix--vector products of \cref{ch:orthog}---lives
entirely outside this chapter.

\section{Where this sits in the machine}
\label{sec:gv-context}

It is worth placing the running QR in the larger GMRES picture one last time. The Arnoldi
process of \cref{ch:orthog} is the \emph{expensive} part: each step does a field-sized
matrix--vector product $\mat A\vv_k$ and orthogonalizes against the basis, both $O(N)$ or
worse. The least-squares core of this chapter is the \emph{cheap} part: it works only on
the small dense Hessenberg and the small right-hand side, $O(m)$ scalars per step. The two
interlock once per iteration---Arnoldi hands the running QR a new column, the running QR
hands back a residual norm that decides whether Arnoldi runs again---and at the end the
back-solve crosses back into field space exactly once, to reconstruct $\vx_m$ from the
basis. The producer monitors cheaply; the consumer materializes rarely; and the residual,
the one quantity GMRES most needs and would otherwise pay dearly for, falls out of the
factorization for nothing. That free residual is the elegant heart of the method, and it
is why GMRES, despite its appetite for memory, remains the workhorse of \cref{ch:gmres}.

\summarysection
GMRES needs, at each step $m$, the solution of a small upper-Hessenberg least-squares
problem $\min_\vy\norm{\beta\,\vec e_1 - \bar{\mat H}_m\vy}$. Rather than re-factor from
scratch, the solver maintains a \emph{running QR} factorization, updated one column at a
time by a stream of $2\times2$ Givens (plane) rotations. A Givens rotation $G(c,s)$ with
$c = a/r$, $s = b/r$, $r = \sqrt{a^2+b^2}$ zeros one sub-diagonal entry; production code
computes $r$ by the scaled, overflow-safe formula. The stream processes each arriving
column in four sub-steps---\emph{replay} the stored rotations, \emph{generate} a new one,
\emph{apply} it to the column, \emph{apply} it to the right-hand side---and is exactly a
\code{foldl} over columns, with the rotation list as the growing carry; the replay must
precede the generate, because plane rotations do not commute. The decisive payoff: after
the right-hand-side update, the magnitude of its new last entry $\abs{g_{j+1}}$ \emph{is}
the GMRES residual norm, a free byproduct of unitarity---no iterate, no $\mat A\vx$, no
residual evaluation. The back-solve, run once per restart cycle, solves the small
triangular system by back-substitution and reconstructs $\vx_m = \vx_0 + \mat V_m\vy$
(or $\mat Z_m\vy$ for FGMRES) by one \code{linear\_combination}. The stream is a
streaming producer running every step; the back-solve is a terminal consumer running
once.

\furtherreading
The running-QR-via-Givens implementation of GMRES is treated in
Saad~\citep{saad2003}, \S6.5.3, which is the closest reference to the stream as Palace
realizes it, including the free residual. The Givens rotation itself, its stable
generation, and its use in QR factorization are in Golub \& Van
Loan~\citep{golubvanloan2013}, \S5.1 (rotations) and \S5.2 (QR); the overflow-safe scaled
generation follows the LAPACK convention discussed there. Trefethen \&
Bau~\citep{trefethenbau1997}, Lectures 10 and 35, give an unusually clear geometric
account of orthogonal triangularization and of GMRES as residual minimization over a
Krylov subspace. The complex-arithmetic conventions and the precise scaled formula are
collected in \cref{app:numerics}.

\exercisessection
\begin{enumerate}[nosep]
  \item \emph{Generate by hand.} Compute the Givens rotation $(c,s)$ that zeros the
  second component of $(a,b) = (5,12)$. Give $r$, $c$, $s$, verify $c^2+s^2=1$, and check
  that applying $G(c,s)$ to $(5,12)$ yields $(13,0)$.
  \item \emph{Orthogonality preserves length.} Show directly from \cref{eq:gv-matrix}
  that $G(c,s)$ is orthogonal ($G^{\!\top}G = I$) whenever $c^2+s^2=1$, and conclude that
  $\norm{G\vu} = \norm{\vu}$ for every 2-vector $\vu$. Where in the chapter is this fact
  the reason the residual is free?
  \item \emph{Why scale.} In double precision the overflow threshold is about
  $1.8\times10^{308}$. Take $a = b = 10^{200}$. Show that the literal $r=\sqrt{a^2+b^2}$
  overflows (compute $a^2$), while the scaled formula \eqref{eq:gv-scaled} returns the
  correct finite $r = 10^{200}\sqrt2$. By how large a factor does $a^2$ exceed the
  threshold?
  \item \emph{Replay order matters.} In \cref{wex:gv-stream}, redo step 1 but
  \emph{skip} the replay of $(c_0,s_0)$ and generate the new rotation directly from the
  raw column $(1,5,12)$. Show that the resulting $\mat R$ is no longer the QR factor of
  $\bar{\mat H}_2$ (e.g.\ the residual you would read off no longer matches the true
  least-squares residual). Explain in one sentence which load-bearing invariant you
  violated.
  \item \emph{Read the residual.} A GMRES run reports, after its rotation stream at step
  $m$, a rotated right-hand side whose last entry is $g_{m} = -1.5\times10^{-7}$ with
  $\beta = 2.0$. What is the current relative residual $\norm{\vr_m}/\norm{\vb}$
  (assume $\vx_0=0$ so $\beta=\norm{\vb}$)? Would a tolerance of $10^{-6}$ declare
  convergence? Note that you answered this without forming $\vx_m$.
  \item \emph{Back-solve.} Given $\mat R_2 = \begin{psmallmatrix} 2 & 1 \\ 0 & 4
  \end{psmallmatrix}$ and the fitted right-hand side $\vg_{0:2} = (6, 8)^{\!\top}$, solve
  $\mat R_2\vy = \vg_{0:2}$ by the back-substitution of \cref{eq:gv-backsub}. Which
  unknown do you compute first, and why?
  \item \emph{The fold made explicit.} Write \code{ls\_update\_column} as a fold step
  $st' = \code{step}(st, \vh)$ and identify, for the running QR, (a) the carry $st$,
  (b) the per-element input $\vh$, and (c) the part of the carry that \emph{grows} with
  each step. Why does the presence of the replay sub-step force this to be a
  \code{foldl} (a chain) rather than a \code{map} (an independent fan-out)? Refer to the
  map/fold split of \cref{ch:combinators}.
  \item \emph{Producer versus consumer.} The stream runs every Arnoldi step and the
  back-solve runs once per restart cycle. If a GMRES run takes $40$ steps before a
  restart, how many times does each fire? Explain why moving the back-solve inside the
  per-step loop (so it fires every step) would be correct but wasteful, and estimate the
  factor of extra back-solve work it would incur.
\end{enumerate}

\chapter{Preconditioning}
\label{ch:precond}

\chapterorient{The last two chapters built solvers and then bounded their speed by a
single number: the condition number of the operator. Conjugate gradients
(\cref{ch:cg}) crawls on a stretched bowl; GMRES (\cref{ch:gmres}) needs more steps
the more the spectrum spreads. This chapter is the cure. We do not build a faster
solver---we change the \emph{problem} the solver sees, multiplying it by an
approximate inverse $\mat{M}^{-1}\approx\mat{A}^{-1}$ that pulls the spectrum into a
tight cluster near $1$, so the very same iteration finishes in a fraction of the
steps. The spine of the chapter is the abstraction of \cref{ch:operators}: a
preconditioner is never a matrix you form, it is an \emph{operator you apply}---a
black-box function the solver calls and never looks inside. That one idea makes
preconditioning modular: the same GMRES drives a diagonal scaling, an incomplete
factorization, or a multigrid V-cycle without changing a line. We make the spectral
claim precise, sort out left/right/split placement, build the preconditioner once and
apply it each step (the constructed operator of \cref{ch:operators}), tour the
families, and end at the real goal: convergence that does not slow down as the mesh is
refined.}

\section{The motivation, made precise}
\label{sec:motivation}

We ended \cref{ch:cg} with a stark plot (\cref{fig:convplot}): the conjugate-gradient
error after $k$ steps is bounded by
\[
  \norm{\vx^\star-\vx_k}_{\mat{A}}
  \;\le\; 2\!\left(\frac{\sqrt{\kappa}-1}{\sqrt{\kappa}+1}\right)^{\!k}
  \norm{\vx^\star-\vx_0}_{\mat{A}},
  \qquad \kappa=\frac{\lambda_{\max}}{\lambda_{\min}},
\]
and GMRES carries an analogous spectral dependence. The number of iterations to reach
a fixed tolerance grows like $\sqrt{\kappa}$ for CG. The operators Palace assembles are
routinely ill-conditioned: refining the mesh, or putting a high-permittivity dielectric
next to a low one, drives $\kappa$ into the thousands and beyond. Raw CG and raw GMRES
are simply too slow on these systems. The fix is not a cleverer iteration---it is to
\emph{change the spectrum}.

\subsection{The idea in one line}

Suppose we have a second operator $\mat{M}$, close to $\mat{A}$ but \emph{cheap to
invert}, and consider solving the algebraically equivalent system
\begin{equation}
\label{eq:precond-system}
  \mat{M}^{-1}\mat{A}\,\vx \;=\; \mat{M}^{-1}\vb .
\end{equation}
It has the \emph{same solution} $\vx^\star$ as $\mat{A}\vx=\vb$---multiplying both sides
by the invertible $\mat{M}^{-1}$ changes nothing about the answer. But it can have a
completely different \emph{operator}: where $\mat{A}$ might have a spectrum spread across
$[\lambda_{\min},\lambda_{\max}]$ with $\kappa$ in the thousands, the new operator
$\mat{M}^{-1}\mat{A}$ has a spectrum we can hope to squeeze near $1$. In the ideal limit
$\mat{M}=\mat{A}$, the new operator is $\mat{M}^{-1}\mat{A}=\mat{I}$, whose condition
number is exactly $1$: a Krylov method solves $\mat{I}\vx=\mat{M}^{-1}\vb$ in a single
step. We will never reach that ideal cheaply (inverting $\mat{A}$ exactly is the problem
we are trying to avoid), but the closer $\mat{M}^{-1}$ approximates $\mat{A}^{-1}$, the
closer $\mat{M}^{-1}\mat{A}$ is to the identity, and the faster the iteration runs.

\begin{definition}[Preconditioner]
\label{def:precond}
A \emph{preconditioner} for the system $\mat{A}\vx=\vb$ is an operator $\mat{M}$,
nonsingular and \emph{cheap to apply the inverse of}, chosen so that
$\mat{M}^{-1}\mat{A}$ is better conditioned (its spectrum more clustered) than
$\mat{A}$. The action $\vr\mapsto\mat{M}^{-1}\vr$ is the \emph{preconditioner apply}.
\end{definition}

Two demands pull against each other in that definition, and a good preconditioner is a
truce between them. It must \emph{cluster the spectrum}---otherwise it does not reduce
the iteration count---and it must be \emph{cheap to apply}---otherwise the saved
iterations cost more than they were worth. The extremes are both useless:
$\mat{M}=\mat{I}$ is free to apply but clusters nothing (it does not precondition at
all); $\mat{M}=\mat{A}$ clusters perfectly but is as expensive to invert as the original
problem. Every practical preconditioner lives in between.

\subsection{The fundamental trade}
\label{sec:trade}

Make the bargain quantitative. Without preconditioning, a solve costs
\[
  (\text{iterations}) \times (\text{cost per iteration}),
\]
and the cost per iteration is dominated by one operator apply $\mat{A}\vp$. Preconditioning
changes \emph{both} factors, in opposite directions:

\begin{itemize}[nosep]
  \item \textbf{Iterations go down}---often dramatically---because the spectrum of
  $\mat{M}^{-1}\mat{A}$ is clustered and the Krylov bound improves.
  \item \textbf{Cost per iteration goes up}, because each step now performs one
  \emph{extra} apply: the preconditioner apply $\vz=\mat{M}^{-1}\vr$ on top of the
  operator apply $\mat{A}\vp$.
\end{itemize}

\begin{keyidea}
Preconditioning wins exactly when the product wins: the iteration count must fall by
more than the per-step cost rises. This happens when $\mat{M}^{-1}$ is
\emph{simultaneously} (i) a good approximation of $\mat{A}^{-1}$---so it genuinely
clusters the spectrum and slashes the iteration count---and (ii) cheap to apply---so the
extra per-step apply is a small surcharge, not a second full solve. A preconditioner
that clusters but is expensive, or one that is cheap but clusters nothing, loses the
trade. The art is finding an $\mat{M}^{-1}$ that is both.
\end{keyidea}

\Cref{fig:cost-trade} draws the bargain. As the preconditioner grows stronger (left to
right), the iteration count plunges, but the work spent \emph{inside} each preconditioner
apply climbs. Total time is the product, and it has a minimum somewhere in the middle: a
``sweet spot'' preconditioner strong enough to cut iterations sharply yet cheap enough
that each apply stays light. Too weak and the iterations dominate; too strong and the
per-step preconditioner apply dominates.

\begin{figure}[t]
\centering
\begin{tikzpicture}
  \begin{axis}[
      width=12.4cm, height=6.6cm,
      xlabel={preconditioner strength (cheap/weak $\to$ strong/expensive)},
      ylabel={cost (arbitrary units)},
      xmin=0, xmax=10, ymin=0, ymax=11,
      xtick=\empty,
      tick label style={font=\scriptsize}, label style={font=\footnotesize},
      legend pos=north west, legend cell align=left,
      legend style={font=\footnotesize, draw=simGray!40},
      grid=both, grid style={simGray!15},
      every axis plot/.append style={very thick}]
    % iterations: falling
    \addplot[simBlue, domain=0.3:10, samples=80]
      {1.0 + 9.0*exp(-0.55*x)};
    \addlegendentry{iterations (fall)}
    % per-iteration cost: rising
    \addplot[simRust, domain=0.3:10, samples=80]
      {0.6 + 0.34*x};
    \addlegendentry{cost per iteration (rise)}
    % total time = product, normalized down to fit
    \addplot[simGreen, domain=0.3:10, samples=120]
      {0.42*(1.0 + 9.0*exp(-0.55*x))*(0.6 + 0.34*x)};
    \addlegendentry{total time (product)}
    % mark the minimum of the green curve, near x ~ 3.4
    \node[circle, fill=simGreen, inner sep=1.6pt] at (axis cs:3.4,2.55) {};
    \node[font=\scriptsize, simGreen, anchor=south west] at (axis cs:3.55,2.7)
      {sweet spot};
    \draw[simGreen, densely dashed] (axis cs:3.4,0)--(axis cs:3.4,2.55);
  \end{axis}
\end{tikzpicture}
\caption[The fundamental preconditioning trade]{The trade at the heart of
preconditioning. A stronger preconditioner (rightward) clusters the spectrum more, so
the iteration count (blue) falls steeply---but each step now spends more work inside the
preconditioner apply, so the cost per iteration (rust) rises. Total solve time (green) is
the \emph{product} of the two, and it has a minimum in the middle: the sweet spot is a
preconditioner strong enough to cut iterations sharply yet cheap enough that each apply
stays light. At the left edge the iterations dominate (too weak); at the right edge the
per-step apply dominates (too strong, approaching a full inner solve). A good
preconditioner sits near the minimum.}
\label{fig:cost-trade}
\end{figure}

\section{What clustering buys: the spectrum picture}
\label{sec:spectrum}

The phrase ``better conditioned'' in \cref{def:precond} deserves a picture, because it
is the whole mechanism. \Cref{ch:cg} taught (\cref{sec:cluster}) that a Krylov method's
true convergence is governed not just by the extreme eigenvalues but by how the
\emph{whole} spectrum is laid out: a polynomial can be made small on a few tight clusters
with far fewer degrees than it needs to be small across a wide interval. A preconditioner
is, at bottom, a spectrum-reshaping device.

\Cref{fig:spectrum} shows the effect on a model problem. The eigenvalues of $\mat{A}$
(top) sprawl across two decades---this is the long, narrow bowl of \cref{ch:cg}, the
$\kappa$ in the thousands. The eigenvalues of $\mat{M}^{-1}\mat{A}$ (bottom), for a
decent preconditioner, huddle in a tight band around $1$. Same solution, same iteration,
utterly different convergence: the bottom spectrum is one a degree-handful polynomial
annihilates, so the Krylov method converges in a handful of steps.

\begin{figure}[t]
\centering
\begin{tikzpicture}
  \begin{axis}[
      width=13.2cm, height=4.8cm,
      xmode=log, log basis x=10,
      xmin=0.05, xmax=300,
      ymin=-0.6, ymax=2.2,
      ytick={0,1.5}, yticklabels={$\mat{M}^{-1}\mat{A}$, $\mat{A}$},
      y tick label style={font=\footnotesize},
      xlabel={eigenvalue (log scale)},
      label style={font=\footnotesize}, tick label style={font=\scriptsize},
      axis y line=left, axis x line=bottom,
      clip=false]
    % --- spectrum of A: spread across two decades ---
    \addplot[only marks, mark=*, mark size=2.0pt, simBlue]
      coordinates {(0.08,1.5)(0.15,1.5)(0.3,1.5)(0.7,1.5)(1.6,1.5)
                   (4,1.5)(9,1.5)(22,1.5)(55,1.5)(130,1.5)(240,1.5)};
    \node[font=\scriptsize, simBlue, anchor=west] at (axis cs:0.055,2.0)
      {spread: $\kappa\approx 3000$};
    % --- spectrum of M^{-1}A: clustered near 1 ---
    \addplot[only marks, mark=*, mark size=2.0pt, simGreen]
      coordinates {(0.62,0)(0.74,0)(0.83,0)(0.9,0)(0.97,0)(1.0,0)
                   (1.04,0)(1.1,0)(1.18,0)(1.3,0)(1.45,0)};
    \node[font=\scriptsize, simGreen, anchor=west] at (axis cs:1.6,0)
      {clustered: $\kappa\approx 2.3$};
    % a marker line at lambda = 1
    \draw[simGray, densely dashed] (axis cs:1,-0.55)--(axis cs:1,1.0)
      node[pos=0.0, below, font=\scriptsize, simGray] {$\lambda=1$};
  \end{axis}
\end{tikzpicture}
\caption[Preconditioning clusters the spectrum]{The spectral effect of preconditioning,
on a log scale. \emph{Top (blue):} the eigenvalues of the bare operator $\mat{A}$ sprawl
across two decades---a condition number in the thousands, the long narrow bowl that makes
CG crawl (\cref{ch:cg}). \emph{Bottom (green):} the eigenvalues of the preconditioned
operator $\mat{M}^{-1}\mat{A}$ collapse into a tight cluster around $\lambda=1$, with a
condition number near $2$. The system has the \emph{same solution}---only the operator
the Krylov method iterates against has changed. A polynomial that must be tiny across the
blue spread needs a high degree (many iterations); one that need only be tiny on the
green cluster needs a low degree (few iterations). Clustering near $1$ is exactly what a
good preconditioner buys.}
\label{fig:spectrum}
\end{figure}

\begin{notationbox}
``Clustered near $1$'' is the goal, and the $1$ is not arbitrary. If $\mat{M}^{-1}\mat{A}$
were clustered near any constant $c\neq0$, a rescaling would move the cluster to $1$, so
what matters is the \emph{spread} of the cluster, not its location. We say
$\mat{M}^{-1}\mat{A}\approx\mat{I}$ to mean its spectrum is a tight cluster; equivalently
$\mat{M}^{-1}\approx\mat{A}^{-1}$. The closer the approximation, the tighter the cluster,
the fewer the iterations---and the limit $\mat{M}^{-1}=\mat{A}^{-1}$ gives the single-step
solve.
\end{notationbox}

\section{Left, right, and split preconditioning}
\label{sec:lrs}

Equation~\eqref{eq:precond-system} put $\mat{M}^{-1}$ on the \emph{left} of $\mat{A}$.
That was one choice of three. The preconditioner can be applied on the left, on the
right, or split across both sides, and the three placements---algebraically equivalent in
their solution, but different in what the iteration actually minimizes---are a variant
axis (\cref{ch:operators}) every Krylov solver exposes.

\begin{description}[style=nextline,leftmargin=1.5em]
  \item[Left preconditioning.] Multiply the system on the left:
  \[
    \mat{M}^{-1}\mat{A}\,\vx \;=\; \mat{M}^{-1}\vb .
  \]
  The iteration runs on the operator $\mat{M}^{-1}\mat{A}$, and the residual it sees and
  drives to zero is the \emph{preconditioned residual}
  $\mat{M}^{-1}(\vb-\mat{A}\vx)=\mat{M}^{-1}\vr$, not the true residual $\vr$.
  \item[Right preconditioning.] Insert $\mat{M}^{-1}\mat{M}=\mat{I}$ between $\mat{A}$ and
  $\vx$, and solve for a transformed unknown $\vu=\mat{M}\vx$:
  \[
    \mat{A}\mat{M}^{-1}\,\vu \;=\; \vb, \qquad \vx=\mat{M}^{-1}\vu .
  \]
  The iteration runs on $\mat{A}\mat{M}^{-1}$, but---crucially---the residual it sees is
  $\vb-\mat{A}\mat{M}^{-1}\vu=\vb-\mat{A}\vx=\vr$, the \emph{true} residual. Right
  preconditioning leaves the quantity being minimized untouched.
  \item[Split preconditioning.] Factor the preconditioner $\mat{M}=\mat{M}_1\mat{M}_2$
  and place one factor on each side:
  \[
    \mat{M}_1^{-1}\mat{A}\mat{M}_2^{-1}\,\vu \;=\; \mat{M}_1^{-1}\vb,
    \qquad \vx=\mat{M}_2^{-1}\vu .
  \]
  When $\mat{A}$ is SPD and $\mat{M}=\mat{L}\mat{L}^{\top}$ is a (incomplete) Cholesky
  factorization, the split choice $\mat{M}_1=\mat{L}$, $\mat{M}_2=\mat{L}^{\top}$ keeps the
  preconditioned operator $\mat{L}^{-1}\mat{A}\mat{L}^{-\top}$ \emph{symmetric positive
  definite}---which is exactly what preconditioned CG needs (\cref{sec:pcg-tie}).
\end{description}

\begin{figure}[t]
\centering
\begin{tikzpicture}[font=\footnotesize, node distance=5mm]
  % --- LEFT ---
  \begin{scope}
    \node[font=\bfseries\small, simBlue] at (0,1.5) {left};
    \node[extern, minimum width=10mm] (lM) at (0,0.5) {$\mat{M}^{-1}$};
    \node[stage, minimum width=9mm, below=4mm of lM] (lA) {$\mat{A}$};
    \node[data, below=4mm of lA] (lx) {$\vx$};
    \draw[flow] (lx)--(lA); \draw[flow] (lA)--(lM);
    \node[font=\scriptsize, simGray, anchor=north, align=center] at (0,-1.55)
      {$\mat{M}^{-1}\mat{A}\vx=\mat{M}^{-1}\vb$\\[1pt]sees $\mat{M}^{-1}\vr$};
  \end{scope}
  % --- RIGHT ---
  \begin{scope}[xshift=4.7cm]
    \node[font=\bfseries\small, simGreen] at (0,1.5) {right};
    \node[stage, minimum width=9mm] (rA) at (0,0.5) {$\mat{A}$};
    \node[extern, minimum width=10mm, below=4mm of rA] (rM) {$\mat{M}^{-1}$};
    \node[data, below=4mm of rM] (ru) {$\vu$};
    \draw[flow] (ru)--(rM); \draw[flow] (rM)--(rA);
    \node[font=\scriptsize, simGray, anchor=north, align=center] at (0,-1.55)
      {$\mat{A}\mat{M}^{-1}\vu=\vb$\\[1pt]sees \emph{true} $\vr$};
  \end{scope}
  % --- SPLIT ---
  \begin{scope}[xshift=9.4cm]
    \node[font=\bfseries\small, simRust] at (0,1.5) {split};
    \node[extern, minimum width=10mm] (sM1) at (0,0.9) {$\mat{M}_1^{-1}$};
    \node[stage, minimum width=9mm, below=3mm of sM1] (sA) {$\mat{A}$};
    \node[extern, minimum width=10mm, below=3mm of sA] (sM2) {$\mat{M}_2^{-1}$};
    \node[data, below=3mm of sM2] (su) {$\vu$};
    \draw[flow] (su)--(sM2); \draw[flow] (sM2)--(sA); \draw[flow] (sA)--(sM1);
    \node[font=\scriptsize, simGray, anchor=north, align=center] at (0,-1.95)
      {$\mat{M}_1^{-1}\mat{A}\mat{M}_2^{-1}\vu=\mat{M}_1^{-1}\vb$\\[1pt]preserves symmetry};
  \end{scope}
\end{tikzpicture}
\caption[Three placements of the preconditioner]{The three placements of $\mat{M}^{-1}$,
read bottom-to-top as the operand flows through. \emph{Left:} $\mat{M}^{-1}$ multiplies
the whole system from the left; the iteration runs on $\mat{M}^{-1}\mat{A}$ and drives the
\emph{preconditioned} residual $\mat{M}^{-1}\vr$ to zero. \emph{Right:} $\mat{M}^{-1}$ sits
between $\mat{A}$ and the transformed unknown $\vu=\mat{M}\vx$; the iteration runs on
$\mat{A}\mat{M}^{-1}$ but sees the \emph{true} residual $\vr$---the reason right
preconditioning is the default for GMRES. \emph{Split:} one factor of
$\mat{M}=\mat{M}_1\mat{M}_2$ on each side preserves SPD-ness of the preconditioned
operator, which preconditioned CG requires. All three have the same solution $\vx^\star$;
they differ in what residual the convergence test reads.}
\label{fig:lrs}
\end{figure}

\subsection{Why the choice matters: the residual GMRES minimizes}

The three placements yield the same solution, so why prefer one? The answer is the
\emph{stopping test}. A Krylov method measures convergence by a residual norm, and the
three placements present it with different residuals.

\begin{keyidea}
\textbf{Right preconditioning preserves the true residual.} Left preconditioning makes
the iteration minimize $\norm{\mat{M}^{-1}\vr}$---the residual seen through the
preconditioner---so the stopping test is distorted by $\mat{M}$: a small
$\norm{\mat{M}^{-1}\vr}$ need not mean a small true $\norm{\vr}$ if $\mat{M}$ is
ill-scaled. Right preconditioning leaves the residual as the true $\vr=\vb-\mat{A}\vx$, so
GMRES's free residual readout (\cref{ch:gmres}) measures the quantity the engineer
actually cares about. This is why right preconditioning is the conventional default for
preconditioned GMRES; left is common for CG, where the preconditioned inner product
restores symmetry anyway (\cref{sec:pcg-tie}).
\end{keyidea}

There is one more reason the distinction is load-bearing, and it ties back to
\cref{ch:gmres}: for \emph{right}-preconditioned GMRES the iterate is recovered as
$\vx=\vx_0+\mat{M}^{-1}\mat{V}_m\vy$, applying $\mat{M}^{-1}$ once to the whole
correction---which assumes a \emph{single fixed} $\mat{M}$. The moment the preconditioner
varies per step, that reconstruction breaks, and FGMRES is the repair. We return to this
in \cref{sec:flexible}; for now, file the placement as a variant axis the construction
absorbs.

\section{The preconditioner is an operator you apply}
\label{sec:operator}

We have written $\mat{M}^{-1}$ throughout as if it were a matrix. It is not---and the gap
between ``$\mat{M}^{-1}$ the matrix'' and ``$\mat{M}^{-1}$ the operation'' is the single
most important idea in the chapter. \emph{The inverse is never formed.} What the solver
holds is a function: hand it a residual $\vr$, it returns a vector $\vz\approx\mat{M}^{-1}\vr$.
How it computes that vector---a diagonal divide, a triangular sweep, a full multigrid
cycle---is its own business, hidden behind the apply.

This is precisely the \emph{solver-as-operator} move of \cref{ch:operators}. There we saw
that a solver computes the action $\vv\mapsto\mat{M}^{-1}\vv$, which has the same shape as
the action $\vv\mapsto\mat{M}\vv$ of the operator it inverts---vector in, vector out---so
a solver can wear the very same \code{apply} interface as a plain operator. A
preconditioner \emph{is} a solver in this sense, and that is why a Krylov method needs
nothing from it but the apply.

\begin{keyidea}
\textbf{A preconditioner is an approximate inverse you apply; the solver never looks
inside.} The Krylov iteration is written against one interface---\code{apply(pc, r)}---and
asks exactly one thing of the preconditioner: turn this residual into an approximate
error, $\vz\approx\mat{M}^{-1}\vr$. It does not know, and cannot ask, \emph{how}. The same
GMRES drives a Jacobi diagonal-scale, an incomplete factorization, a multigrid V-cycle
(\cref{ch:multigrid}), or even an inner Krylov solve---because all of them are just
functions of the shape $\vr\mapsto\vz$. Preconditioning is modular for exactly this
reason: choosing a preconditioner is choosing a function, and the iteration is blind to
the choice.
\end{keyidea}

\Cref{fig:blackbox-pc} draws the blindness. GMRES issues an \code{apply} and receives a
vector; whatever algorithm produced that vector is sealed inside the box. Swap the box's
contents---Jacobi for ILU, ILU for multigrid---and not one line of the iteration changes.
The arrow into the box and the vector out of it are the entire contract.

\begin{figure}[t]
\centering
\begin{tikzpicture}[font=\footnotesize, node distance=9mm]
  % the GMRES driver
  \node[stage, minimum width=24mm, minimum height=20mm] (gm) at (0,0)
    {\textbf{GMRES}\\[2pt]\footnotesize the iteration\\asks only\\\code{apply(pc, r)}};
  % the opaque preconditioner box
  \node[draw=simViolet, fill=simBgViolet, very thick, densely dashed, rounded corners=3pt,
        minimum width=30mm, minimum height=30mm, right=20mm of gm] (box) {};
  \node[font=\scriptsize\bfseries, simViolet] at (box.north) [yshift=-3mm] {preconditioner};
  \node[font=\scriptsize\itshape, simGray, align=center] at (box.center) [yshift=2mm]
    {Jacobi?\\ILU?\\multigrid V-cycle?\\inner Krylov solve?\\[2pt]\textbf{GMRES cannot tell.}};
  % apply arrow in
  \draw[flow] (gm.east) -- node[above, font=\scriptsize]{$\vr$}
    node[below, font=\scriptsize]{\code{apply}} (box.west);
  % vector out
  \draw[flow] (box.south) .. controls +(0,-0.9) and +(0,-0.9) ..
    node[below, font=\scriptsize]{$\vz\approx\mat{M}^{-1}\vr$} (gm.south);
  % the seal
  \node[font=\scriptsize\bfseries, simRust, align=center, right=2mm of box]
    {one interface,\\any algorithm};
\end{tikzpicture}
\caption[The preconditioner as a black box GMRES cannot see inside]{Modularity made
visual. The GMRES iteration hands the preconditioner a residual $\vr$ through the single
\code{apply} interface and receives back a vector $\vz\approx\mat{M}^{-1}\vr$. The
algorithm \emph{inside} the box---a diagonal divide, an incomplete factorization, a
multigrid V-cycle, an inner Krylov solve---is completely sealed off: GMRES never inspects
it and never needs to. Replacing the box's contents changes the convergence \emph{rate}
but not a single line of the iteration. This is the solver-as-operator rotation of
\cref{ch:operators}: because an approximate inverse wears the same vector-in, vector-out
shape as an operator, the iteration treats it as one. Choosing a preconditioner is
choosing the function inside the box.}
\label{fig:blackbox-pc}
\end{figure}

\subsection{The same iteration, four preconditioners}

To feel the modularity concretely, here is the preconditioned-residual line of a Krylov
step, written once, with four different bindings of \code{pc}. The iteration code is
identical in all four; only the function handed to it differs.

\begin{lstlisting}[style=l4style]
-- the iteration body is written ONCE against the apply interface:
let z = apply pc r in           -- z ~ M^{-1} r ; pc is opaque here
... use z in the Krylov update ...

-- ...and pc is any of these, chosen at construction, never re-inspected:
let pc = make_jacobi   A         in ...   -- z = d .* r          (one divide)
let pc = make_ilu      A         in ...   -- z = U \ (L \ r)     (two triangular solves)
let pc = make_multigrid hierarchy in ...  -- z = v_cycle r       (a recursive sweep)
let pc = \r -> gmres_inner A r 3  in ...   -- z = 3 steps of inner GMRES (flexible!)
\end{lstlisting}

The first three are constructed operators (\cref{sec:constructed}): build once, apply many
times, the \emph{same} \code{pc} every step. The fourth---an inner solve---produces a
\emph{different} effective operator each step, and that difference, as we will see in
\cref{sec:flexible}, is precisely what forces FGMRES.

\section{The constructed-operator view: build once, apply each step}
\label{sec:constructed}

The preconditioner apply is cheap per call, but most preconditioners are \emph{not} cheap
to set up: extracting and inverting a diagonal, computing an incomplete factorization, or
assembling a multigrid hierarchy is real work. The saving grace is that this work is done
\emph{once}, before the iteration begins, and then amortized over every step. This is
exactly the build-time / run-time phase split of \cref{ch:strata} and the constructed
operator of \cref{ch:operators}.

\begin{keyidea}
A preconditioner is a \emph{constructed operator} (\cref{ch:operators}). Its life has two
phases in two strata. \textbf{Construction} (build-time, once): consume $\mat{A}$ and a
configuration, do the expensive setup---extract a diagonal, compute an incomplete
factorization, build a multigrid hierarchy---and freeze the result into immutable internal
state. \textbf{Application} (run-time, many times): consume a residual $\vr$, read the
frozen state, return $\vz\approx\mat{M}^{-1}\vr$, change nothing. The expensive part runs
once; the cheap part runs every step. This is why the per-iteration surcharge of
\cref{sec:trade} is just \emph{one apply}, not a fresh setup.
\end{keyidea}

\Cref{fig:build-apply} draws the split for a generic preconditioner, deliberately echoing
the construct-then-apply figure of \cref{ch:operators}. The matrix and config flow in on
the left and are consumed once by the build; the opaque preconditioner value sits in the
middle holding its frozen internals; a stream of cheap \code{apply(r)} calls flows out on
the right, one per Krylov step, each reading the same internals.

\begin{figure}[t]
\centering
\begin{tikzpicture}[node distance=10mm]
  % construction inputs
  \node[data, align=left, text width=24mm] (cfg)
    {operator $\mat{A}$\\pc-type enum\\config};
  % build step
  \node[stage, minimum width=26mm, minimum height=14mm, right=12mm of cfg] (build)
    {\textbf{build}\\[1pt]\footnotesize extract diag /\\factorize /\\set up hierarchy};
  \draw[flow] (cfg) -- node[above,font=\scriptsize]{once} (build);
  % opaque pc value
  \node[draw=simViolet, fill=simBgViolet, very thick, densely dashed, rounded corners=3pt,
        minimum width=24mm, minimum height=20mm, right=12mm of build] (pc) {};
  \node[font=\ttfamily\scriptsize, simInk] at (pc.north) [yshift=-3.2mm] {pc};
  \node[font=\scriptsize\itshape, simGray, align=center] at (pc.center)
    {frozen $\mat{M}^{-1}$\\internals\\(diag, factors,\\hierarchy)};
  \draw[flow] (build) -- (pc);
  % many cheap applies
  \foreach \y/\lbl in {1.1/a, 0.0/b, -1.1/c}{
    \node[draw=simTeal, fill=simBgGreen, semithick, rounded corners=2pt,
          minimum width=20mm, minimum height=6mm] (ap\lbl) at ($(pc.east)+(2.7,\y)$)
          {\footnotesize\ttfamily apply r};
    \draw[flow] (pc.east) -- (ap\lbl.west);
  }
  \node[font=\scriptsize\itshape, simGray] at ($(pc.east)+(2.7,-1.9)$)
    {one per Krylov step};
  % stratum band labels
  \node[font=\scriptsize\bfseries, simRust] at ($(build)+(0,-1.7)$) {build-time};
  \node[font=\scriptsize\bfseries, simGreen] at ($(pc.east)+(2.7,1.9)$) {run-time};
\end{tikzpicture}
\caption[Build once, apply each step]{A preconditioner is the constructed operator of
\cref{ch:operators}, with the two phases in the two strata of \cref{ch:strata}.
\emph{Construction} (left) consumes the operator $\mat{A}$, a preconditioner-type enum,
and configuration, runs the expensive setup \emph{once}---extract and invert a diagonal,
compute an incomplete factorization, assemble a multigrid hierarchy---and freezes the
result into the preconditioner's immutable internals. \emph{Application} (right) consumes
only a residual and runs \emph{many} times, once per Krylov step, each call cheap and pure
on the frozen internals. The build-time arrow fires once; the run-time arrows fire
hundreds of times. This phase split is why the per-step cost of preconditioning is a
single apply, never a fresh setup.}
\label{fig:build-apply}
\end{figure}

\subsection{The factory that returns a uniform apply}

A real solver must choose \emph{which} preconditioner to build, and that choice arrives as
a configuration enum. \Cref{ch:operators} named the place the enum is consumed: the
\emph{constructed-operator factory}. A single function reads the
preconditioner-type$\,\times\,$side$\,\times\,$multigrid axes, dispatches once, and returns
a \emph{uniformly typed} apply---past which no enum survives. Palace's
\code{ConfigurePreconditionerSolver} is exactly this factory: it reads the pc-type enum
(\code{JACOBI}, \code{AMS}, \code{BoomerAMG}, a direct factorization, \dots), wraps the
result in a multigrid solver when the FE-space hierarchy has more than one level, and hands
back one \code{Solver} value the iteration applies without ever re-reading the enum.

\begin{lstlisting}[style=l4style]
-- the factory: ONE site reads the pc-type enum, returns a uniform apply
configure_pc :: PcType -> Config -> Op[Tensor[N] -> Tensor[N]]
-- downstream the enum is GONE; only the uniform constructed operator remains:
let pc = configure_pc pc_type cfg in   -- dispatch happens HERE, once
... each Krylov step ...
let z = apply pc r in                  -- ...and never again; pc is opaque
\end{lstlisting}

\begin{workedexample}{Jacobi-preconditioned CG versus plain CG}{jacobi-cg}
We make the whole chapter concrete on the cleanest possible preconditioner---Jacobi
(diagonal) scaling---and a model ill-conditioned operator, and watch the iteration count
collapse.

\textbf{The model problem.} Take a diagonally-dominant SPD operator whose diagonal entries
span a wide range, so that the diagonal carries most of the ill-conditioning. A concrete
small instance: let $\mat{A}$ have diagonal $\operatorname{diag}(\mat{A})=[1,10,100,1000]$
with modest off-diagonal coupling, so $\kappa(\mat{A})\approx 10^{3}$. The spread is
\emph{in the diagonal}---exactly the case Jacobi is built for.

\textbf{Construction (once).} The Jacobi preconditioner is $\mat{M}=\operatorname{diag}(\mat{A})$,
and its apply is ``divide each component by the diagonal entry.'' Construction extracts the
diagonal and inverts it componentwise, producing one vector $\mathbf{d}$ of reciprocals
(\cref{wex:jacobi} of \cref{ch:operators} traced this end to end):
\begin{lstlisting}[style=l4style]
make_jacobi A =
  let d = reciprocal (extract_diagonal A) in   -- O(N), once
  \r -> d .* r                                  -- the apply: O(N) per call
\end{lstlisting}
For our $\mat{A}$, $\mathbf{d}=[1,\tfrac1{10},\tfrac1{100},\tfrac1{1000}]$.

\textbf{The spectral effect.} Scaling each row by its diagonal pulls the four widely-spread
diagonal entries to the \emph{same} value $1$, leaving only the off-diagonal coupling to
perturb the spectrum. The eigenvalues of $\mat{M}^{-1}\mat{A}=\operatorname{diag}(\mat{A})^{-1}\mat{A}$
cluster tightly near $1$: the operator now has unit diagonal, and for our modest coupling
$\kappa(\mat{M}^{-1}\mat{A})\approx 2$ to $3$ instead of $10^3$. This is exactly the
top-to-bottom collapse of \cref{fig:spectrum}.

\textbf{The iteration count.} Feed both condition numbers into the CG bound
$2\big((\sqrt\kappa-1)/(\sqrt\kappa+1)\big)^k$ of \cref{thm:converge}. To reach a relative
energy-norm error of $10^{-6}$:
\begin{itemize}[nosep]
  \item plain CG, $\kappa=10^{3}$: the rate factor is $(\sqrt{1000}-1)/(\sqrt{1000}+1)\approx
  0.939$, so $0.939^k\le 10^{-6}/2$ needs $k\approx 240$ iterations;
  \item Jacobi-CG, $\kappa=2.5$: the rate factor is $(\sqrt{2.5}-1)/(\sqrt{2.5}+1)\approx
  0.226$, so $0.226^k\le 10^{-6}/2$ needs $k\approx 11$ iterations.
\end{itemize}
A factor of roughly $20$ fewer iterations, bought by one $O(N)$ construction and one extra
$O(N)$ multiply per step---a per-step surcharge dwarfed by the operator apply on any real
sparse $\mat{A}$. \Cref{fig:jacobi-conv} plots the two convergence curves; the gap is the
entire point of preconditioning.
\end{workedexample}

\begin{figure}[t]
\centering
\begin{tikzpicture}
  \begin{semilogyaxis}[
      width=12.4cm, height=6.6cm,
      xlabel={iteration $k$}, ylabel={relative energy-norm error (bound)},
      xmin=0, xmax=260, ymin=1e-8, ymax=3,
      grid=both, grid style={simGray!18},
      legend pos=north east, legend cell align=left,
      legend style={font=\footnotesize, draw=simGray!40},
      label style={font=\footnotesize}, tick label style={font=\scriptsize},
      every axis plot/.append style={very thick}]
    % plain CG: kappa = 1000
    \addplot[simRust, domain=0:260, samples=80]
      {2*exp(x*ln((sqrt(1000)-1)/(sqrt(1000)+1)))};
    \addlegendentry{plain CG, $\kappa=10^3$ ($\approx240$ its)}
    % Jacobi-CG: kappa = 2.5
    \addplot[simGreen, domain=0:30, samples=40]
      {2*exp(x*ln((sqrt(2.5)-1)/(sqrt(2.5)+1)))};
    \addlegendentry{Jacobi-CG, $\kappa=2.5$ ($\approx11$ its)}
    \draw[simGray, densely dashed] (axis cs:0,1e-6)--(axis cs:260,1e-6)
      node[pos=0.5, above, font=\scriptsize, text=simGray] {tol $=10^{-6}$};
  \end{semilogyaxis}
\end{tikzpicture}
\caption[Jacobi-preconditioned CG versus plain CG]{The convergence bound of
\cref{thm:converge} for the model problem of \cref{wex:jacobi-cg}. \emph{Plain CG} (rust)
on the bare operator $\kappa=10^3$ crawls, needing about $240$ iterations to reach a
$10^{-6}$ tolerance. \emph{Jacobi-preconditioned CG} (green), running on
$\mat{M}^{-1}\mat{A}$ with its spectrum pulled to $\kappa\approx2.5$, reaches the same
tolerance in about $11$ iterations. The two curves run the \emph{same} CG iteration; the
only difference is that the green one applies the diagonal preconditioner once per step.
The $\approx20\times$ reduction in iterations is bought by one $O(N)$ construction and one
$O(N)$ apply per step---a trade that pays handsomely whenever the diagonal carries the
ill-conditioning.}
\label{fig:jacobi-conv}
\end{figure}

\subsection{When the preconditioner varies: the decision rule and FGMRES}
\label{sec:flexible}

The constructed-operator picture has one precondition: there must be a \emph{single}
$\mat{M}$ to freeze at construction. \Cref{ch:operators} gave the rule that decides when
this holds.

\begin{keyidea}
\textbf{The decision rule (\cref{ch:operators}), applied to preconditioners.} A
preconditioner \emph{fixed} for the whole solve---chosen at solve start, the same operator
every step---is build-time configuration: freeze it into a constructed operator, apply the
same \code{pc} uniformly. A preconditioner that \emph{varies per step}---a different
effective operator $\mat{M}_k$ on each iteration---cannot be frozen, because at construction
time the later $\mat{M}_k$ do not yet exist. It must instead be threaded through the run-time
state. This single rule is exactly why a flexible preconditioner forces FGMRES rather than
GMRES (\cref{ch:gmres}).
\end{keyidea}

When does a preconditioner vary per step? Whenever the apply is \emph{itself an inner
iterative solve} run to a loose, varying tolerance: a few steps of an inner GMRES, a
multigrid cycle whose accuracy drifts, a smoother that adapts to the residual. Each inner
solve produces a slightly different linear operator, so $\vz_k=\mat{M}_k^{-1}\vr_k$ uses a
different $\mat{M}_k$ every step. There is no single $\mat{M}$ to construct---the fourth
binding in the listing of \cref{sec:operator}, \lstinline[style=l4style]|\r -> gmres_inner A r 3|,
is exactly this case.

\Cref{wex:fixed-vs-flexible} contrasts the two apply signatures and shows the variant
moving out of the constructed operator and into the threaded state, exactly as the decision
rule predicts.

\begin{workedexample}{The decision rule in action: a fixed apply versus a per-step apply}{fixed-vs-flexible}
We put two preconditioner applies side by side and read their signatures.

\textbf{Case A: a fixed preconditioner (constructed once).} The preconditioner is built
before the loop and applied---the \emph{same} operator---every step. Its apply signature
takes only the residual; everything else is frozen inside:
\begin{lstlisting}[style=l4style]
let pc = configure_pc JACOBI cfg in    -- built ONCE; pc :: Op[Tensor[N] -> Tensor[N]]
... each step ...
let z = apply pc r in ...               -- apply :: Op[N->N] -> Tensor[N] -> Tensor[N]
                                        -- SAME pc every step; nothing threaded
\end{lstlisting}
The state the loop threads is the Krylov basis and bookkeeping; the preconditioner is
\emph{not} part of it, because it never changes. The constructed operator absorbs the entire
``which preconditioner'' question. This is the left branch of the decision rule
(\cref{fig:decisionrule} of \cref{ch:operators}), and plain GMRES or CG handles it.

\textbf{Case B: a per-step preconditioner (an inner solve).} Now each step runs a fresh
inner solve of varying accuracy. There is no \code{pc} to construct; the apply takes the
\emph{step index} (or the residual that determines the inner tolerance) and produces a
different effective operator each time:
\begin{lstlisting}[style=l4style]
-- NO single pc; M_j is produced per step by an inner solve:
let z_j = apply_M_step j v_j in   -- z_j = M_j^{-1} v_j, a DIFFERENT M_j each j
\end{lstlisting}
and---the crucial consequence---the per-step preconditioned vectors $\vz_j$ can no longer be
recomputed on demand, because the inner solve that made each one is gone by the next step.
They must be \emph{remembered}. The solver threads a \emph{second} basis $\mat{Z}_m=[\vz_1\
\cdots\ \vz_m]$ alongside the Arnoldi basis $\mat{V}_m$, and recovers its iterate from the
remembered preconditioned basis:
\begin{lstlisting}[style=l4style]
-- GMRES  (fixed pc):    x = x_0 + V_m y_m   -- pc is a constructed operator, NOT in state
-- FGMRES (per-step pc): x = x_0 + Z_m y_m   -- the per-step Z basis IS threaded state
\end{lstlisting}

\textbf{The reading.} Nothing went wrong; a real distinction surfaced. The fixed
preconditioner of Case A is build-time configuration---freeze it, apply uniformly. The
per-step preconditioner of Case B is run-time state---thread it, and the carry grows by one
extra basis $\mat{Z}$. The decision rule read the variant's lifetime and routed it to the
matching stratum. ``FGMRES exists because the preconditioner became per-step state'' is the
whole story, told in full in \cref{ch:gmres}.
\end{workedexample}

\begin{pitfall}
\textbf{Do not feed a per-step preconditioner to plain GMRES.} Plain GMRES recovers its
iterate as $\vx_0+\mat{M}^{-1}\mat{V}_m\vy$, applying one $\mat{M}^{-1}$ to the whole
correction under the assumption of a \emph{single fixed} $\mat{M}$. Hand it a varying
$\mat{M}_k$ and the code still runs---no crash---and even reports a small residual from its
running QR; but that residual belongs to a least-squares problem the reconstructed iterate
no longer matches, because the reconstruction used the wrong inverse. The solve converges to
the \emph{wrong answer}, silently. The fix is not a patch to GMRES; it is FGMRES, which
remembers the preconditioned vectors in the $\mat{Z}$ basis (\cref{ch:gmres}). The tell that
you need it: your preconditioner apply is itself an iterative solve.
\end{pitfall}

\section{A tour of preconditioner families}
\label{sec:families}

Every preconditioner is a function $\vr\mapsto\vz\approx\mat{M}^{-1}\vr$; the families
differ only in \emph{which} function, i.e. in what their construction freezes and what their
apply computes. We survey the main families here as apply-functions, light on internals---each
gets its own chapter later. The point is the unity: they all wear the same interface, so
choosing one is choosing a function, and the iteration is unchanged.

\begin{description}[style=nextline,leftmargin=1.5em]
  \item[Diagonal / Jacobi.] $\mat{M}=\operatorname{diag}(\mat{A})$; the apply is a single
  componentwise divide $\vz=\mathbf{d}\odot\vr$. The cheapest possible preconditioner, and
  the constructed-operator archetype (\cref{wex:jacobi-cg}). It clusters well only when the
  diagonal carries the ill-conditioning; it is also the simplest \emph{smoother}
  (\cref{ch:smoothers}).
  \item[Incomplete factorizations (ILU / IC).] Compute an approximate $\mat{L}\mat{U}$ (or,
  for SPD $\mat{A}$, Cholesky $\mat{L}\mat{L}^{\top}$) factorization that keeps only a
  sparsity pattern, discarding fill-in. Construction is the incomplete factorization; the
  apply is two triangular solves $\vz=\mat{U}^{-1}(\mat{L}^{-1}\vr)$. Stronger than Jacobi,
  more expensive to build and apply, and---being inherently sequential---a poor fit for the
  GPU target (\cref{ch:strata} discusses the recurrence obstruction).
  \item[Domain decomposition.] Partition the mesh into subdomains, solve a smaller problem on
  each, and stitch the pieces together. The apply is a collection of independent subdomain
  solves plus a coarse correction. Naturally parallel across subdomains; the family of choice
  when the problem is already distributed.
  \item[Multigrid.] The big one for elliptic PDEs (\cref{ch:multigrid}). Construction builds a
  \emph{hierarchy} of coarser and coarser meshes; the apply is a V-cycle that smooths the
  high-frequency error on fine levels (with a cheap smoother like Jacobi) and solves the
  low-frequency error on coarse levels (cheaply, because they are small). Done right, multigrid
  is \emph{optimal}: its cluster, and hence its iteration count, is independent of mesh size
  (\cref{sec:meshindep}).
  \item[Auxiliary-space / AMS.] For $\Hcurl$ curl--curl problems (the magnetostatic and
  eigenmode operators), the natural multigrid smoothers fail on the large near-kernel of the
  curl operator. The \emph{auxiliary-space Maxwell solver} (AMS) preconditions by mapping the
  problem into auxiliary $\Hone$ scalar and vector spaces---where ordinary multigrid works---and
  mapping back. Its apply is a fixed composition of grid-transfer and smoothing operations
  (\cref{ch:smoothers}).
\end{description}

\Cref{fig:family-tree} arranges the families as a tree, all sharing one root: the
\code{apply} interface. The branches are different functions; the root is the single shape
$\vr\mapsto\vz$ that lets the same Krylov method drive any of them.

\begin{figure}[t]
\centering
\begin{tikzpicture}[font=\footnotesize,
    famnode/.style={draw=simViolet, fill=simBgViolet, rounded corners=2pt, semithick,
      align=center, inner sep=3pt, minimum height=8mm, text width=22mm},
    node distance=4mm]
  % root: the apply interface
  \node[draw=simRust, fill=simBgRust, very thick, rounded corners=3pt, align=center,
        inner sep=4pt, minimum width=46mm]
    (root) at (0,3.2)
    {\textbf{the \code{apply} interface}\\[1pt]\footnotesize $\vr\mapsto\vz\approx\mat{M}^{-1}\vr$};
  % the five families
  \node[famnode] (jac) at (-5.4,0.8) {Jacobi /\\diagonal\\\scriptsize one divide};
  \node[famnode] (ilu) at (-2.7,0.8) {ILU / IC\\\scriptsize triangular\\\scriptsize solves};
  \node[famnode] (dd)  at (0,0.8)    {domain\\decomposition\\\scriptsize subdomain solves};
  \node[famnode] (mg)  at (2.7,0.8)  {multigrid\\\scriptsize V-cycle\\\scriptsize (optimal)};
  \node[famnode] (ams) at (5.4,0.8)  {AMS\\\scriptsize aux-space\\\scriptsize ($\Hcurl$)};
  % connect root to each
  \foreach \n in {jac,ilu,dd,mg,ams}{
    \draw[depedge] (root.south) -- (\n.north);}
  % chapter pointers
  \node[font=\scriptsize\itshape, simGray, below=1mm of jac] {\cref{ch:smoothers}};
  \node[font=\scriptsize\itshape, simGray, below=1mm of mg]  {\cref{ch:multigrid}};
  \node[font=\scriptsize\itshape, simGray, below=1mm of ams] {\cref{ch:smoothers}};
\end{tikzpicture}
\caption[A family tree of preconditioners sharing one interface]{The preconditioner
families, arranged as a tree rooted in the single \code{apply} interface
$\vr\mapsto\vz\approx\mat{M}^{-1}\vr$. Jacobi/diagonal scaling, incomplete factorizations
(ILU/IC), domain decomposition, multigrid, and the auxiliary-space AMS solver differ
entirely in \emph{what} their apply computes and \emph{what} their construction freezes---a
divide, a pair of triangular solves, a set of subdomain solves, a V-cycle, a mapping into
auxiliary spaces---but every one of them is a function of the same shape. Because they share
the interface, the same Krylov method drives any of them, and choosing a preconditioner is
choosing which branch of the tree to hang in the black box of \cref{fig:blackbox-pc}. The
chapters named below the leaves develop the internals.}
\label{fig:family-tree}
\end{figure}

\section{Spectral equivalence and mesh-independent convergence}
\label{sec:meshindep}

We close with the property that separates a merely good preconditioner from an
\emph{optimal} one, and that is the real goal for the elliptic operators Palace solves.

The condition-number improvement of \cref{wex:jacobi-cg} was a fixed factor: $\kappa$ from
$10^3$ to $2.5$ on \emph{that} mesh. But Palace refines meshes, and the question that
matters is what happens as the mesh shrinks. For the bare elliptic operators, the condition
number \emph{grows} as the mesh is refined: discretizing $-\Div(\varepsilon\Grad\,\cdot\,)$
on a mesh of element size $h$ gives $\kappa(\mat{A})=O(h^{-2})$, so halving $h$ quadruples
$\kappa$ and---since CG iterations scale like $\sqrt\kappa$---roughly \emph{doubles} the
iteration count. A solver whose iteration count climbs every time you refine is a solver
that does not scale.

\begin{definition}[Spectral equivalence and optimality]
\label{def:specequiv}
A preconditioner $\mat{M}$ is \emph{spectrally equivalent} to $\mat{A}$ if there are
constants $0<c_1\le c_2$, \emph{independent of the mesh size $h$}, such that
\[
  c_1\,\vx^{\top}\mat{M}\vx \;\le\; \vx^{\top}\mat{A}\vx \;\le\; c_2\,\vx^{\top}\mat{M}\vx
  \qquad\text{for all }\vx.
\]
Then the preconditioned operator satisfies $\kappa(\mat{M}^{-1}\mat{A})\le c_2/c_1$,
\emph{bounded independently of $h$}. A preconditioner achieving this is called
\emph{optimal}: the iteration count to a fixed tolerance does not grow as the mesh is
refined.
\end{definition}

\begin{keyidea}
The real goal of preconditioning an elliptic PDE is not a one-off condition-number
improvement---it is \emph{mesh independence}. A spectrally-equivalent (optimal)
preconditioner makes $\mat{M}^{-1}\mat{A}$ as well-conditioned on a fine mesh as on a coarse
one: the cluster of \cref{fig:spectrum} stays the same width no matter how many unknowns
$N$. The iteration count then stays \emph{flat} as $N$ grows, so the total solve cost scales
like $O(N)$---linearly---rather than $O(N^{1.5})$ or worse. Multigrid (\cref{ch:multigrid})
is the canonical optimal preconditioner for these operators, and mesh-independent
convergence is exactly the property that makes it indispensable at scale.
\end{keyidea}

\Cref{fig:meshindep} contrasts the two regimes: a fixed-strength preconditioner (Jacobi),
whose iteration count climbs as $h\to0$ because $\kappa$ still grows like $h^{-2}$ underneath
it, against an optimal preconditioner (multigrid), whose iteration count is a flat line
because the cluster width is bounded. The flat line is the prize.

\begin{figure}[t]
\centering
\begin{tikzpicture}
  \begin{axis}[
      width=12.4cm, height=6.6cm,
      xlabel={mesh refinement $1/h$ (finer mesh $\to$)},
      ylabel={iterations to tolerance},
      xmin=0, xmax=34, ymin=0, ymax=130,
      tick label style={font=\scriptsize}, label style={font=\footnotesize},
      legend pos=north west, legend cell align=left,
      legend style={font=\footnotesize, draw=simGray!40},
      grid=both, grid style={simGray!15},
      every axis plot/.append style={very thick, mark=*, mark size=1.4pt}]
    % Jacobi / fixed preconditioner: iterations ~ sqrt(kappa) ~ sqrt(h^-2) = 1/h, linear climb
    \addplot[simRust] coordinates
      {(2,8)(4,16)(8,32)(12,48)(16,64)(20,80)(24,96)(28,112)(32,124)};
    \addlegendentry{Jacobi (mesh-dependent): iters $\sim 1/h$}
    % multigrid / optimal: flat
    \addplot[simGreen] coordinates
      {(2,9)(4,9)(8,10)(12,10)(16,10)(20,11)(24,11)(28,11)(32,11)};
    \addlegendentry{multigrid (optimal): iters bounded}
  \end{axis}
\end{tikzpicture}
\caption[Mesh-dependent versus mesh-independent convergence]{Iteration count as the mesh is
refined ($1/h$ growing rightward). With a \emph{fixed-strength} preconditioner like Jacobi
(rust), the underlying condition number still grows like $h^{-2}$, so the iteration count
climbs roughly like $1/h$---refine the mesh and every solve gets slower, the hallmark of a
non-scalable method. With an \emph{optimal}, spectrally-equivalent preconditioner like
multigrid (green, \cref{def:specequiv}), the condition number of $\mat{M}^{-1}\mat{A}$ is
bounded independently of $h$, so the iteration count is essentially \emph{flat} no matter how
fine the mesh. The flat green line is the goal of preconditioning at scale: total cost
$O(N)$, not $O(N^{1.5})$. Mesh independence is what makes multigrid (\cref{ch:multigrid})
indispensable for the elliptic operators Palace assembles.}
\label{fig:meshindep}
\end{figure}

\subsection{Tie-off: preconditioned CG and the SPD requirement}
\label{sec:pcg-tie}

\Cref{ch:cg} already showed (\cref{sec:pcg}) how the CG \emph{loop} changes under
preconditioning: form the preconditioned residual $\vz=\mat{M}^{-1}\vr$ once per step, and
replace the residual inner products $\inner{\vr}{\vr}$ in $\alpha$ and $\beta$ by the
$\mat{M}$-weighted products $\inner{\vz}{\vr}$. We can now name \emph{why} the loop takes that
particular shape. Running CG naively on $\mat{M}^{-1}\mat{A}$ would break SPD-ness, because
$\mat{M}^{-1}\mat{A}$ is generally not symmetric (\cref{sec:lrs}); preconditioned CG instead
runs CG in the inner product defined by $\mat{M}$, which is precisely the split-preconditioning
view $\mat{L}^{-1}\mat{A}\mat{L}^{-\top}$ for $\mat{M}=\mat{L}\mat{L}^{\top}$. That operator
\emph{is} SPD, so every guarantee of \cref{ch:cg} survives. The discipline this forces---the
preconditioner $\mat{M}$ must itself be SPD for preconditioned CG to stay honest---is the
pitfall \cref{ch:cg} flagged, now explained: it is the price of keeping the energy inner
product an inner product.

\summarysection
A \emph{preconditioner} $\mat{M}\approx\mat{A}$, cheap to invert, transforms
$\mat{A}\vx=\vb$ into an equivalent system whose operator $\mat{M}^{-1}\mat{A}$ has a
spectrum clustered near $1$, so the \emph{same} Krylov iteration (\cref{ch:cg,ch:gmres})
converges in far fewer steps. The fundamental trade is a product: each step costs one extra
preconditioner apply, but the iteration count falls, and the win comes when $\mat{M}^{-1}$ is
\emph{both} a good approximation of $\mat{A}^{-1}$ (clusters the spectrum) and cheap to apply.
The preconditioner can be placed on the \emph{left} ($\mat{M}^{-1}\mat{A}\vx=\mat{M}^{-1}\vb$,
distorts the residual), the \emph{right} ($\mat{A}\mat{M}^{-1}\vu=\vb$, preserves the true
residual---the GMRES default), or \emph{split} ($\mat{M}_1^{-1}\mat{A}\mat{M}_2^{-1}$,
preserves SPD-ness for CG)---a variant axis the construction absorbs. The chapter's spine is
that \emph{a preconditioner is an operator you apply}: $\mat{M}^{-1}$ is never formed, only
applied as a black-box function $\vr\mapsto\vz\approx\mat{M}^{-1}\vr$, and the Krylov method
needs nothing else---so the same GMRES drives Jacobi, ILU, multigrid, or an inner solve
without changing a line (\cref{ch:operators}'s solver-as-operator move). The preconditioner is
a \emph{constructed operator}: build once (extract a diagonal, factorize, set up a hierarchy)
at build-time, apply cheaply each step at run-time. When the preconditioner is \emph{fixed},
the constructed-operator factory absorbs it; when it \emph{varies per step} (an inner solve),
the decision rule routes it into the threaded state and forces FGMRES. The families---Jacobi,
ILU/IC, domain decomposition, multigrid, auxiliary-space AMS---all share the apply interface;
choosing one is choosing a function. The real goal for elliptic operators is \emph{spectral
equivalence}: an optimal preconditioner bounds $\kappa(\mat{M}^{-1}\mat{A})$ independently of
the mesh size $h$, so the iteration count stays flat as the mesh is refined and the solve
scales linearly---the property that makes multigrid (\cref{ch:multigrid}) indispensable.

\furtherreading
The definitive treatment of preconditioned Krylov methods---left/right/split placement, the
flexible-preconditioner distinction that drives FGMRES, and the spectral-equivalence theory of
optimal preconditioners---is \citet{saad2003} (Chapters~9 and~10), which also motivates the
black-box ``a Krylov method needs only the action $\vv\mapsto\mat{M}^{-1}\vv$'' abstraction
that makes preconditioning modular. \citet{trefethenbau1997} gives the cleanest short account
of how clustering the spectrum---rather than merely shrinking $\kappa$---governs Krylov
convergence (Lectures~35 and~40). For the matrix-computation mechanics of incomplete
factorizations and the diagonal/Jacobi family, \citet{golubvanloan2013} is the standard
reference. The optimal multigrid preconditioner and the auxiliary-space AMS solver for
$\Hcurl$ are developed in \cref{ch:multigrid} and \cref{ch:smoothers}; the constructed-operator
and solver-as-operator abstractions the chapter rests on are \cref{ch:operators}.

\exercisessection
\begin{enumerate}[nosep]
  \item \textbf{The two demands.} \Cref{def:precond} asks a preconditioner to be both a good
  approximation of $\mat{A}^{-1}$ and cheap to apply. Explain why the two extremes
  $\mat{M}=\mat{I}$ and $\mat{M}=\mat{A}$ each satisfy exactly one demand and fail the other,
  and what each does to (i) the iteration count and (ii) the per-step cost. Where on the trade
  curve of \cref{fig:cost-trade} does each extreme sit?
  \item \textbf{The trade, quantified.} A solve takes $200$ iterations unpreconditioned, each
  costing one operator apply of $C$ units. A preconditioner cuts the iteration count to $25$ but
  adds a preconditioner apply costing $0.8C$ per step. (i) Compute the total cost with and
  without preconditioning (ignore the one-time build). (ii) Now suppose a \emph{stronger}
  preconditioner cuts iterations to $8$ but its apply costs $5C$ per step; compute its total
  cost. (iii) Which of the three wins, and what does this illustrate about \cref{fig:cost-trade}?
  \item \textbf{Clustering versus $\kappa$.} \Cref{ch:cg} noted that CG's true convergence depends
  on the spectrum's \emph{layout}, not just $\kappa$. Two preconditioned operators both have
  $\kappa=50$: operator $P$ has its eigenvalues spread uniformly across $[1,50]$, operator $Q$ has
  $48$ eigenvalues tightly clustered near $1$ and two outliers near $50$. Which converges faster
  under CG, and why? Relate your answer to the polynomial picture of \cref{sec:cluster} and to
  \cref{fig:spectrum}.
  \item \textbf{Left versus right.} (i) Write the system, the operator the iteration runs on, and
  the residual the iteration sees, for left and for right preconditioning. (ii) Explain precisely
  why right preconditioning is preferred for GMRES, referring to the free residual readout of
  \cref{ch:gmres}. (iii) For an SPD $\mat{A}$ with $\mat{M}=\mat{L}\mat{L}^{\top}$, show that the
  split-preconditioned operator $\mat{L}^{-1}\mat{A}\mat{L}^{-\top}$ is symmetric, and explain why
  CG needs this where GMRES does not.
  \item \textbf{The black box.} The same GMRES drives a Jacobi preconditioner, an ILU, and a
  multigrid V-cycle without changing a line (\cref{fig:blackbox-pc}). (i) State the single
  interface that makes this possible and exactly what GMRES asks of the preconditioner through it.
  (ii) Name one thing GMRES is \emph{prevented} from seeing about the preconditioner, and explain
  why that prevention is a feature. (iii) Which chapter's abstraction (\cref{ch:operators}) is this
  an instance of, and what is the shared shape that lets a solver wear an operator's interface?
  \item \textbf{Build once, apply many.} For the Jacobi preconditioner of \cref{wex:jacobi-cg},
  state the cost of construction and of a single apply, each in terms of $N$. If a solve runs $k$
  preconditioned steps, write the total preconditioner work as a function of $N$ and $k$ and
  identify the amortized construction term. For what range of $k$ is constructing the operator
  clearly worthwhile?
  \item \textbf{Apply the decision rule.} For each preconditioner below, decide whether it is a
  constructed operator (fixed, build-time) or per-step threaded state (forcing FGMRES), and
  justify with the decision rule of \cref{sec:flexible}: (a) a diagonal Jacobi scaling built from
  $\mat{A}$ at solve start; (b) an incomplete Cholesky factorization computed once; (c) ``three
  steps of an inner GMRES'' run afresh each outer step; (d) a multigrid V-cycle with fixed smoother
  counts; (e) a multigrid cycle whose inner tolerance adapts to the current residual.
  \item \textbf{Mesh independence.} The bare elliptic operator has $\kappa(\mat{A})=O(h^{-2})$.
  (i) Using the CG iteration estimate iterations $\sim\sqrt{\kappa}$, show that halving $h$ roughly
  doubles the unpreconditioned iteration count. (ii) Define spectral equivalence
  (\cref{def:specequiv}) and explain why it makes the preconditioned iteration count independent of
  $h$. (iii) Sketch, from \cref{fig:meshindep}, the iteration-count-versus-$1/h$ curve for a
  mesh-dependent and an optimal preconditioner, and state the total-cost scaling ($O(N)$ versus
  worse) each implies.
  \item \textbf{The SPD trap.} Preconditioned CG requires the preconditioner $\mat{M}$ to be SPD.
  (i) Explain, using \cref{sec:pcg-tie}, why running plain CG on $\mat{M}^{-1}\mat{A}$ with a
  general $\mat{M}$ breaks the method even when $\mat{A}$ is SPD. (ii) Show how the split view
  $\mat{L}^{-1}\mat{A}\mat{L}^{-\top}$ restores symmetry. (iii) What goes wrong, observably, if a
  non-SPD preconditioner is used inside preconditioned CG?
\end{enumerate}

\chapter{Smoothers: Jacobi, Chebyshev, and Hiptmair}
\label{ch:smoothers}

% local math macro: eigenvector (not in simbook.sty)
\providecommand{\vphi}{\boldsymbol{\varphi}}

\chapterorient{The previous chapter argued that a good preconditioner is an
approximate inverse cheap enough to apply at every iteration. This chapter builds
the cheapest useful one. A \emph{smoother} is a relaxation so simple it makes a
laughable standalone solver---left to run on its own it crawls---yet it has a
hidden talent the eigenvalue picture brings out at once: it annihilates the
\emph{oscillatory} part of the error in a sweep or two while barely touching the
\emph{smooth} part. That lopsidedness is not a defect. It is exactly the division
of labor a multigrid method needs, and the smoother is the workhorse bolted into
every level of the V-cycle we build in \cref{ch:multigrid}. We meet three: damped
Jacobi (a diagonal scaling), Chebyshev (a fixed-degree polynomial tuned to a
spectral window), and Hiptmair's distributive relaxation (the one fix that makes
edge-element curl--curl problems tractable at all). Each is, in the end, an
apply-function of the kind \cref{ch:precond} taught us to plug in.}

\section{Error in the eigenvector basis}

Everything in this chapter follows from one change of coordinates. Suppose we are
solving a symmetric positive-definite system
\begin{equation}
  \mat{A}\vx = \vb,
  \label{eq:smoo-system}
\end{equation}
the kind that comes out of every weak form in \cref{part:fieldtheory}. Let
$\hat\vx$ be the exact solution, and let an iterative method produce a running
guess $\vx^{(m)}$ after $m$ steps. The object that actually matters is not the
guess but its \emph{error},
\begin{equation}
  \ve^{(m)} = \vx^{(m)} - \hat\vx.
  \label{eq:smoo-error}
\end{equation}
A solver is good exactly when it drives $\ve^{(m)}$ to zero quickly. To see
\emph{how} an iteration acts on the error, we write the error in the natural
coordinate system of $\mat{A}$: its eigenvectors.

Because $\mat{A}$ is symmetric positive-definite, it has a full set of orthonormal
eigenvectors $\vphi_1,\dots,\vphi_n$ with positive eigenvalues
\begin{equation}
  \mat{A}\vphi_k = \lambda_k\,\vphi_k,
  \qquad 0 < \lambda_1 \le \lambda_2 \le \dots \le \lambda_n,
  \label{eq:smoo-eig}
\end{equation}
and we order them from smallest to largest. Any error vector expands uniquely in
this basis,
\begin{equation}
  \ve^{(m)} = \sum_{k=1}^{n} c_k^{(m)}\,\vphi_k,
  \qquad c_k^{(m)} = \inner{\ve^{(m)}}{\vphi_k}.
  \label{eq:smoo-expand}
\end{equation}
The coefficients $c_k^{(m)}$ are the error's \emph{spectral content}: $c_k^{(m)}$
measures how much of the current error lies along the $k$th eigenvector. Driving
the error to zero \emph{means} driving every $c_k$ to zero.

\begin{notationbox}
We use ``frequency'' as a shorthand for ``eigenvalue index $k$.'' For the model
problems in this chapter the eigenvectors $\vphi_k$ really are sinusoids, and the
index $k$ is a genuine spatial frequency: small $k$ is a slowly-varying, smooth
mode; large $k$ is a rapidly-oscillating one. For a general $\mat{A}$ the
eigenvectors are not sinusoids, but the intuition survives---\emph{low frequency}
means the smooth, near-null-space end of the spectrum (small $\lambda_k$), and
\emph{high frequency} the oscillatory, large-$\lambda_k$ end. The whole subject of
smoothing is the story of what an iteration does, eigenvalue by eigenvalue, to the
coefficients $c_k$.
\end{notationbox}

The reason this coordinate change is so powerful is that the simplest iterations
act on each coefficient \emph{independently}. A relaxation built from $\mat{A}$ and
a multiple of the identity (or of the diagonal) commutes with $\mat{A}$'s
eigenvectors, so in the eigenvector basis it is \emph{diagonal}: it multiplies each
$c_k$ by a number that depends only on $\lambda_k$. Call that number the
iteration's \emph{damping factor} at frequency $k$,
\begin{equation}
  c_k^{(m+1)} = g(\lambda_k)\, c_k^{(m)},
  \qquad\text{so}\qquad
  c_k^{(m)} = \bigl[g(\lambda_k)\bigr]^{m}\, c_k^{(0)}.
  \label{eq:smoo-damping}
\end{equation}
The function $g$ is the \emph{amplification factor} or \emph{smoothing symbol} of
the iteration. A component shrinks fast where $\abs{g(\lambda_k)}$ is small and
slowly where $\abs{g(\lambda_k)}$ is near $1$. The entire behavior of a smoother is
encoded in the single curve $g$ plotted against $\lambda$. \Cref{fig:smoo-errorfreq}
shows what this looks like on an actual error: a sum of a smooth mode and an
oscillatory mode, with the oscillatory part visibly gone after a couple of sweeps
and the smooth part stubbornly intact.

\begin{figure}[t]
\centering
\begin{tikzpicture}
\begin{axis}[
  width=12.6cm, height=5.4cm,
  xlabel={position along the domain},
  ylabel={error},
  xmin=0, xmax=1, ymin=-1.35, ymax=1.6,
  xtick=\empty, ytick={-1,0,1},
  legend style={font=\scriptsize, at={(0.5,1.02)}, anchor=south,
    legend columns=3, draw=none, /tikz/every even column/.append style={column sep=8pt}},
  axis lines=left, tick style={simGray}, every axis label/.style={font=\small},
  samples=200, domain=0:1,
]
  % smooth + oscillatory error (initial)
  \addplot[simRust, very thick] {0.85*sin(180*x) + 0.55*sin(180*9*x)};
  % after a few sweeps: oscillatory damped, smooth survives
  \addplot[simBlue, thick, densely dashed] {0.80*sin(180*x) + 0.06*sin(180*9*x)};
  % the smooth part alone (target the smoother leaves behind)
  \addplot[simGray, thin, dotted] {0.80*sin(180*x)};
  \legend{initial error, after 3 sweeps, smooth part}
\end{axis}
\end{tikzpicture}
\caption[Smoothing removes the oscillatory error]{An error built from a smooth
mode ($k=1$) and an oscillatory mode ($k=9$). After three damped-Jacobi sweeps
(dashed) the oscillatory ripples are essentially gone, but the smooth underlying
shape (dotted) is almost untouched. A smoother is precisely the iteration that
does this: it is a high-pass error filter. The smooth remainder is the part the
coarse grid of \cref{ch:multigrid} will handle.}
\label{fig:smoo-errorfreq}
\end{figure}

\begin{keyidea}
A smoother is an iteration whose damping factor $g(\lambda)$ is small at the
\emph{high-frequency} (large-$\lambda$) end of the spectrum and near $1$ at the
\emph{low-frequency} (small-$\lambda$) end. It therefore kills oscillatory error
in a handful of cheap sweeps but barely moves smooth error---a \emph{terrible}
standalone solver and the \emph{essential} ingredient of multigrid, which delegates
the leftover smooth error to a coarse grid. ``Smoother,'' not ``solver,'' is the
right word: its job is to make the error smooth, not small.
\end{keyidea}

\section{Damped Jacobi: diagonal smoothing}
\label{sec:smoo-jacobi}

The simplest relaxation imaginable takes a step proportional to the residual,
preconditioned by nothing more than the diagonal of $\mat{A}$. Write
$\mat{D} = \operatorname{diag}(\mat{A})$ for the matrix holding $\mat{A}$'s
diagonal entries and zeros elsewhere. \emph{Damped Jacobi} is the update
\begin{equation}
  \vx \;\leftarrow\; \vx + \omega\,\mat{D}^{-1}\bigl(\vb - \mat{A}\vx\bigr),
  \label{eq:smoo-jac-update}
\end{equation}
where $\vr = \vb - \mat{A}\vx$ is the residual and $\omega>0$ is a \emph{damping
factor} we are free to choose. Each step costs exactly one operator apply
$\mat{A}\vx$ (\cref{ch:precond}'s matrix-free matvec) plus an elementwise scaling
by $\mat{D}^{-1}$---no inner products, no triangular solves, nothing global. It is
the cheapest nontrivial iteration in the book.

\subsection{The diagonal inverse as a constructed operator}

In the synthesis surface the operator $\omega\mat{D}^{-1}$ is not a matrix; it is a
vector applied elementwise. Building it is a two-step pipeline: \emph{extract} the
diagonal of $\mat{A}$ into a vector, then take its reciprocal entry by entry.
Palace constructs the smoother once, at setup, and stores only the result.

\begin{lstlisting}[style=l4style]
-- setup: build the damped inverse-diagonal smoother once
jacobi_setup :: (A: LinOp[N,N], omega: Float) -> JacobiSmoother[N]
jacobi_setup A omega =
  let d    = assemble_diagonal A      -- extract diag(A) as a vector
      dinv = reciprocal d             -- elementwise 1 / diag(A)
  in  JacobiSmoother { dinv = scale omega dinv }   -- absorb omega into dinv

-- apply: a single elementwise product. no loop, no reduction.
jacobi_smoother :: (op: JacobiSmoother[N], x: Tensor[N]) -> Tensor[N]
jacobi_smoother op x = op.dinv `hadamard` x
\end{lstlisting}

Two things are worth pausing on. First, the damping factor $\omega$ is
\emph{absorbed} into the stored vector \lstinline[style=l4style]|dinv| at setup
time, so the per-call action is a bare elementwise product---``the thinnest
constructed-operator gate'' in the whole framework. Second, the apply is
\emph{inner-product-free}: it touches each entry once and never forms a global sum.
On a parallel machine it requires no communication at all, a property it shares
with Chebyshev and that will matter enormously when we distribute the work.

\begin{notationbox}
\lstinline[style=l4style]|assemble_diagonal| reads the operator's main diagonal;
for a matrix-free high-order operator it returns an \emph{approximate} diagonal,
and the smoother tolerates that---an approximate $\mat{D}^{-1}$ is still a valid
linear map, only a slightly weaker preconditioner. The diagonal extraction, the
reciprocal, and the elementwise product are three firm primitives; the Jacobi
smoother is their composition and nothing more.
\end{notationbox}

\subsection{The damping factor per frequency}

Now run \cref{eq:smoo-jac-update} on the error. Subtract the exact solution
$\hat\vx$ (which satisfies $\mat{A}\hat\vx=\vb$) from both sides. Using
$\vb = \mat{A}\hat\vx$, the residual becomes $\vr = -\mat{A}\ve$, and the update on
the error is
\begin{equation}
  \ve \;\leftarrow\; \ve - \omega\,\mat{D}^{-1}\mat{A}\,\ve
       \;=\; \bigl(\mat{I} - \omega\,\mat{D}^{-1}\mat{A}\bigr)\ve.
  \label{eq:smoo-jac-iter}
\end{equation}
The error is multiplied by the \emph{iteration matrix}
$\mat{G} = \mat{I} - \omega\mat{D}^{-1}\mat{A}$. To read off the per-frequency
damping we diagonalize. For a model problem the diagonal is a constant multiple of
the identity, $\mat{D} = \delta\mat{I}$, so $\mat{D}^{-1}\mat{A}$ shares $\mat{A}$'s
eigenvectors with eigenvalues $\lambda_k/\delta$. Writing $\mu_k = \lambda_k/\delta$
for the scaled eigenvalues---which run over an interval $[\mu_{\min},\mu_{\max}]$
---the damping factor at frequency $k$ is
\begin{equation}
  \boxed{\;
  g(\mu_k) = 1 - \omega\,\mu_k.
  \;}
  \label{eq:smoo-jac-symbol}
\end{equation}
This single affine function is the entire content of damped Jacobi. It is large
(near $1$) when $\mu_k$ is small and decreases as $\mu_k$ grows. \emph{High
frequencies are damped harder than low frequencies}---which is the smoothing
property we wanted---but only if $\omega$ is chosen well.

\begin{pitfall}
Undamped Jacobi ($\omega = 1$) can \emph{fail to smooth at all}, and can even
diverge. With $\omega=1$ the symbol is $g(\mu)=1-\mu$. If the largest scaled
eigenvalue exceeds $2$, then $g(\mu_{\max}) < -1$ and the highest-frequency error
\emph{grows}. Even when it does not diverge, the symbol at the top of the spectrum
($\mu \approx \mu_{\max}$) can sit near $g=-1$, so the most oscillatory modes are
barely damped---the opposite of what a smoother must do. Damping is not a tuning
nicety; without it Jacobi is not a smoother.
\end{pitfall}

\subsection{Choosing the damping factor}

We want $\omega$ to make $\abs{g}$ as small as possible across the
\emph{high-frequency half} of the spectrum, the band a multigrid coarse grid cannot
resolve. The classical balance equalizes the symbol at the two ends of the target
band $[\mu_{\mathrm{lo}}, \mu_{\max}]$:
\begin{equation}
  g(\mu_{\mathrm{lo}}) = -\,g(\mu_{\max})
  \quad\Longrightarrow\quad
  1 - \omega\mu_{\mathrm{lo}} = -(1-\omega\mu_{\max})
  \quad\Longrightarrow\quad
  \omega^\star = \frac{2}{\mu_{\mathrm{lo}} + \mu_{\max}}.
  \label{eq:smoo-jac-optomega}
\end{equation}
This is exactly the formula Palace uses on its estimated-damping path: it estimates
the top eigenvalue $\mu_{\max}$ of $\mat{D}^{-1}\mat{A}$ by a few power-iteration
steps and sets $\omega^\star = 2/(\mu_{\min} + \mu_{\max})$. The spectral interval
$[\mu_{\min},\mu_{\max}]$ comes from the estimate; no exact eigenvalues are ever
formed. The optimal-$\omega$ formula is the standard Jacobi convergence
result~\citep[\S4.1]{saad2003}.

\begin{workedexample}{Damped Jacobi on a 1-D model problem}{smoo-jacobi}
Take the canonical test: the 1-D Laplacian on $n$ interior points with unit
spacing, the tridiagonal matrix $\mat{A}=\operatorname{tridiag}(-1,2,-1)$. Its
eigenvectors are the discrete sinusoids $\vphi_k$ with
$(\vphi_k)_j = \sin\!\bigl(\tfrac{k j\pi}{n+1}\bigr)$ and eigenvalues
\[
  \lambda_k = 2 - 2\cos\!\Bigl(\frac{k\pi}{n+1}\Bigr)
            = 4\sin^2\!\Bigl(\frac{k\pi}{2(n+1)}\Bigr),
  \qquad k = 1,\dots,n.
\]
The diagonal is $\mat{D}=2\mat{I}$, so $\mu_k = \lambda_k/2$ ranges over roughly
$(0,2)$: the lowest mode $k=1$ has $\mu_1\approx 0$ and the highest mode $k=n$ has
$\mu_n\approx 2$. The half-spectrum we want to damp is $[\,1,\,2\,]$ (the modes the
coarse grid misses). Balancing the ends with \cref{eq:smoo-jac-optomega} gives
\[
  \omega^\star = \frac{2}{1 + 2} = \frac{2}{3}.
\]
Now compare the damping at a representative low frequency and a representative high
frequency. At $\mu \approx 0.05$ (a smooth mode) the symbol is
$g = 1 - \tfrac{2}{3}(0.05) \approx 0.967$: after one sweep that mode retains
$96.7\%$ of its amplitude. At $\mu = 2$ (the most oscillatory mode)
$g = 1 - \tfrac{2}{3}(2) = -\tfrac{1}{3}$: that mode retains only $33\%$, a sign
flip and a two-thirds reduction. After $3$ sweeps the high mode is down to
$\abs{g}^3 = (1/3)^3 \approx 0.037$---a $27\times$ reduction---while the smooth mode
is still at $0.967^3 \approx 0.90$. The smoother has cut the oscillatory error by a
factor of about $25$ relative to the smooth error in three cheap sweeps. That ratio
is the whole point.
\end{workedexample}

\Cref{fig:smoo-jacsymbol} plots the symbol $\abs{g}$ for three choices of $\omega$
and shows graphically what the worked example computed: with $\omega^\star=2/3$ the
curve is pinned small across the entire high-frequency half, while undamped
$\omega=1$ leaves the top of the spectrum essentially undamped.

\begin{figure}[t]
\centering
\begin{tikzpicture}
\begin{axis}[
  width=12.6cm, height=6.0cm,
  xlabel={scaled eigenvalue $\mu = \lambda/\delta$ \ (low freq.\ $\to$ high freq.)},
  ylabel={damping factor $|g(\mu)|$},
  xmin=0, xmax=2, ymin=0, ymax=1.05,
  xtick={0,0.5,1,1.5,2}, ytick={0,0.25,0.5,0.75,1},
  legend style={font=\scriptsize, at={(0.02,0.04)}, anchor=south west, draw=simGray!50},
  axis lines=left, tick style={simGray},
  samples=200, domain=0:2,
]
  % shaded high-frequency target band [1,2]
  \addplot[draw=none, fill=simBlue!10, forget plot] coordinates {(1,0)(2,0)(2,1.05)(1,1.05)} \closedcycle;
  \node[font=\scriptsize, simBlue, anchor=north] at (axis cs:1.5,1.02) {high-freq.\ band};
  % undamped omega = 1: |1 - mu|
  \addplot[simRust, very thick] {abs(1 - x)};
  % omega = 2/3 (optimal for [1,2])
  \addplot[simBlue, very thick] {abs(1 - 0.6667*x)};
  % omega = 1/2 (over-damped)
  \addplot[simGray, thick, densely dashed] {abs(1 - 0.5*x)};
  \legend{$\omega=1$ (undamped), $\omega=2/3$ (optimal), $\omega=1/2$}
\end{axis}
\end{tikzpicture}
\caption[The Jacobi damping factor versus frequency]{The damping factor
$\abs{g(\mu)}=\abs{1-\omega\mu}$ for damped Jacobi. Undamped Jacobi ($\omega=1$,
rust) returns to $1$ at the top of the spectrum---it does not smooth the highest
modes. The optimal $\omega=2/3$ (blue) holds $\abs{g}$ below $1/3$ across the entire
high-frequency band $[1,2]$, exactly where the multigrid coarse grid needs it small.
No choice of $\omega$ makes $\abs{g}$ small at low frequency ($\mu\to0$, where
$g\to1$); that is structural, and it is why a smoother alone cannot solve.}
\label{fig:smoo-jacsymbol}
\end{figure}

The figure makes the structural limitation visible: \emph{every} curve returns to
$g=1$ as $\mu\to0$. No diagonal relaxation can damp the smooth modes, because at
$\mu=0$ the update \cref{eq:smoo-jac-update} does nothing. This is not a failure of
tuning. It is the reason smoothers are paired with coarse grids, and it motivates
the next two sections, which sharpen the high-frequency damping without ever trying
to fix the low.

\section{Chebyshev smoothing: a polynomial tuned to the spectrum}
\label{sec:smoo-cheby}

Damped Jacobi takes one step of the affine symbol $1-\omega\mu$. We can do far
better by taking, in effect, several coordinated steps---applying a
\emph{polynomial} in $\mat{D}^{-1}\mat{A}$ whose graph is engineered to be small
across the whole high-frequency band at once. After $k$ such coordinated steps the
error is
\begin{equation}
  \ve \;\leftarrow\; p_k\!\bigl(\mat{D}^{-1}\mat{A}\bigr)\,\ve,
  \label{eq:smoo-cheby-action}
\end{equation}
where $p_k$ is a polynomial of degree $k$ with $p_k(0)=1$ (a unit constant term, so
that the iteration does nothing to the zero-residual fixed point). The symbol is now
$g(\mu)=p_k(\mu)$, and we are free to \emph{choose the polynomial}. The natural
demand is the minimax one: among all degree-$k$ polynomials with $p_k(0)=1$, make
the maximum of $\abs{p_k(\mu)}$ over the target window $[\lambda_{\min},
\lambda_{\max}]$ as small as possible.

\subsection{Why Chebyshev}

The minimax problem has a classical closed-form answer, and it is the reason these
polynomials are named after Chebyshev. The Chebyshev polynomials $T_k$ oscillate
between $-1$ and $+1$ on $[-1,1]$ and then climb away steeply outside that interval.
Affinely mapping the target window $[\lambda_{\min},\lambda_{\max}]$ onto $[-1,1]$
and normalizing so that $p_k(0)=1$ produces the unique polynomial that
\emph{equioscillates}---attains its small maximum amplitude many times---across the
window. No other degree-$k$ polynomial with a unit constant term is smaller across
the band. \Cref{fig:smoo-chebpoly} shows a degree-$4$ example: it is pinned small
over the target interval and rises sharply only at $\mu=0$, the smooth end we are
deliberately leaving to the coarse grid.

\begin{figure}[t]
\centering
\begin{tikzpicture}
\begin{axis}[
  width=12.6cm, height=6.0cm,
  xlabel={eigenvalue $\lambda$},
  ylabel={polynomial symbol $|p_k(\lambda)|$},
  xmin=0, xmax=1.05, ymin=0, ymax=1.05,
  xtick={0,0.2,0.4,0.6,0.8,1.0},
  extra x ticks={0.2,1.0}, extra x tick labels={$\lambda_{\min}$,$\lambda_{\max}$},
  extra x tick style={grid=major, grid style={simBlue!40, densely dotted},
    tick label style={text=simBlue, font=\scriptsize}},
  ytick={0,0.25,0.5,0.75,1},
  legend style={font=\scriptsize, at={(0.98,0.96)}, anchor=north east, draw=simGray!50},
  axis lines=left, tick style={simGray},
  samples=300, domain=0:1.05,
]
  % shaded target window [0.2, 1.0]
  \addplot[draw=none, fill=simBlue!10, forget plot] coordinates {(0.2,0)(1.0,0)(1.0,1.05)(0.2,1.05)} \closedcycle;
  \node[font=\scriptsize, simBlue, anchor=south] at (axis cs:0.6,0.86) {target window};
  % degree-1 (damped Jacobi-like affine) for contrast
  \addplot[simGray, thick, densely dashed] {abs(1 - 1.667*x)};
  % degree-4 chebyshev symbol: |T4((2 lam - lmin - lmax)/(lmax-lmin)) / T4(arg at 0)|
  % lmin=0.2 lmax=1.0 -> map m(lam) = (2 lam - 1.2)/0.8 ; T4(t)=8t^4-8t^2+1
  % normalize by T4 at lam=0: m(0) = -1.5, T4(-1.5)=8*5.0625-8*2.25+1=40.5-18+1=23.5
  \addplot[simBlue, very thick, samples=400] {abs((8*((2*x-1.2)/0.8)^4 - 8*((2*x-1.2)/0.8)^2 + 1)/23.5)};
  \legend{degree 1 (affine), degree 4 (Chebyshev)}
\end{axis}
\end{tikzpicture}
\caption[A Chebyshev polynomial shaped to the spectral window]{The degree-$4$
Chebyshev symbol $\abs{p_4(\lambda)}$ (blue) equioscillates across the target window
$[\lambda_{\min},\lambda_{\max}]$, holding the high-frequency error reduction below a
small uniform bound the whole way. A single affine step (gray dashed, the
damped-Jacobi symbol of \cref{sec:smoo-jacobi}) cannot stay small across the band.
Both rise toward $1$ only at $\lambda=0$, the smooth end left to the coarse grid.
Higher degree buys a smaller uniform bound across the window---at the cost of more
operator applies per sweep.}
\label{fig:smoo-chebpoly}
\end{figure}

The Chebyshev smoother has two further virtues that make it the multigrid smoother
of choice in Palace. It is \emph{matrix-free}: it applies $p_k(\mat{D}^{-1}\mat{A})$
by repeated operator applies, never forming the polynomial as a matrix. And, like
Jacobi, it is \emph{inner-product-free}: its coefficients are fixed closed forms of
the step index and the spectral bounds $[\lambda_{\min},\lambda_{\max}]$, computed
in advance without ever inspecting the iterate. No global reduction appears anywhere
in the sweep---a decisive advantage on a parallel or GPU machine, where inner
products force a synchronization that polynomial smoothing simply avoids. The only
global information it needs is an \emph{estimate} of $\lambda_{\max}$, computed once
at setup by a handful of power-iteration steps.

\begin{physicsbox}
The contrast with the conjugate-gradient method of \cref{ch:precond} is sharp and
deliberate. CG \emph{adapts} to the spectrum: it forms inner products at every
iteration and discovers the optimal polynomial as it goes. Chebyshev \emph{commits}
to a fixed polynomial in advance, paying for that with a spectral estimate but
buying back the inner products entirely. As a standalone solver CG wins. As a
multigrid smoother---applied thousands of times, on a machine where every inner
product is a synchronization---committing in advance is exactly the right trade.
\end{physicsbox}

\subsection{The three-term recurrence}

The polynomial is never written down or expanded. It is applied by a short
recurrence that threads three vectors: the current iterate \lstinline[style=l4style]|y|,
the residual \lstinline[style=l4style]|r|, and a search direction
\lstinline[style=l4style]|d|. One Chebyshev sweep of degree
\lstinline[style=l4style]|order| is:

\begin{lstlisting}[style=l4style]
-- one Chebyshev sweep: apply p_order(Dinv A) to the residual, accumulate into y.
-- coefficients sigma_d, sigma_r are closed forms of (k, lambda_max); no reductions.
chebyshev_sweep :: ChebOp[N] -> Tensor[N] -> Tensor[N] -> Tensor[N]
chebyshev_sweep op x y =
  let r0 = x `sub` (op.A `apply` y)        -- initial residual  r = x - A y
      d0 = scale (alpha0 op) (op.dinv `hadamard` r0)   -- first direction
  in  go 1 (y `add` d0) r0 d0              -- accumulate, then recur
  where
    go k y r d
      | k > op.order = y
      | otherwise =
          let r' = r `sub` (op.A `apply` d)            -- r -= A d
              d' = scale (sigma_d op k) d
                     `add` scale (sigma_r op k) (op.dinv `hadamard` r')
          in  go (k+1) (y `add` d') r' d'              -- y += d
\end{lstlisting}

Read the loop body as the three updates the chapter brief names: accumulate the
direction into the answer ($y \mathrel{+}= d$); update the residual by the
direction's operator image ($r \mathrel{-}= \mat{A}d$); and form the next direction
as a blend of the old direction and the diagonally-scaled residual
($d = \sigma_d\,d + \sigma_r\,(\mat{D}^{-1}r)$). Each iteration costs exactly one
operator apply ($\mat{A}d$) and a few elementwise vector operations. The
coefficients $\sigma_d, \sigma_r$ are pure functions of the step index $k$ and
$\lambda_{\max}$---for the fourth-kind variant Palace uses by default,
$\sigma_d = \tfrac{2k-1}{2k+3}$ and $\sigma_r = \tfrac{8k+4}{(2k+3)\lambda_{\max}}$.
\Cref{fig:smoo-chebdataflow} traces the dataflow of one step and marks the absence
of any reduction node.

\begin{figure}[t]
\centering
\begin{tikzpicture}[node distance=11mm and 16mm]
  \node[data] (r) {residual $\vr$};
  \node[opbox, right=of r] (Ad) {$\mat{A}\vd$ \\ \scriptsize(one apply)};
  \node[opbox, right=of Ad] (rup) {$\vr \mathrel{-}= \mat{A}\vd$};
  \node[opbox, below=10mm of Ad] (dscale) {$\mat{D}^{-1}\vr$ \\ \scriptsize(elementwise)};
  \node[opbox, right=of dscale] (dup) {$\vd = \sigma_d\vd + \sigma_r(\mat{D}^{-1}\vr)$};
  \node[product, right=18mm of rup] (yup) {$\vy \mathrel{+}= \vd$};
  \node[data, below=10mm of dup] (next) {next step};

  \draw[flow] (r) -- (Ad);
  \draw[flow] (Ad) -- (rup);
  \draw[flow] (rup) |- (dscale);
  \draw[flow] (dscale) -- (dup);
  \draw[flow] (dup) -- (yup);
  \draw[flow] (dup) -- (next);
  \draw[backflow] (rup.south) -- ++(0,-0.5) -| (next.north);

  % annotation: no reduction
  \node[cornertag, anchor=west, text=simGreen] at ($(yup.south east)+(0.1,-0.4)$)
    {\scriptsize no inner product};
  \node[cornertag, anchor=west, text=simGreen] at ($(yup.south east)+(0.1,-0.75)$)
    {\scriptsize no global sum};
\end{tikzpicture}
\caption[The Chebyshev recurrence is inner-product-free]{One step of the Chebyshev
three-term recurrence. The only global operation is the single operator apply
$\mat{A}\vd$; everything else is elementwise. Crucially there is \emph{no} inner
product and \emph{no} global reduction anywhere in the loop---the coefficients
$\sigma_d,\sigma_r$ are precomputed from the step index and $\lambda_{\max}$. This is
what makes Chebyshev the communication-light smoother of choice inside a
distributed multigrid V-cycle.}
\label{fig:smoo-chebdataflow}
\end{figure}

\begin{pitfall}
The recurrence is not a convenience; it is \emph{load-bearing numerics}. One could,
in principle, expand the polynomial into its monomial form
$p_k(\mat{D}^{-1}\mat{A})\vr = \sum_j c_j (\mat{D}^{-1}\mat{A})^{j}\vr$ and sum the
powers directly. \emph{Do not.} For the polynomial degrees in use the monomial sum
is numerically unstable---the high powers $(\mat{D}^{-1}\mat{A})^j$ overwhelm
floating-point precision and the coefficients $c_j$ alternate in sign and grow
enormous, so catastrophic cancellation destroys the result. The three-term
recurrence (a Clenshaw/Horner-style evaluation) computes the \emph{same} polynomial
in exact arithmetic but is stable in floating point. The recurrence form and the
monomial form are mathematically equal and computationally not interchangeable; the
sequentiality of the recurrence is part of the algorithm~\citep[\S2]{phillipsfischer2022}.
\end{pitfall}

\subsection{First-kind and fourth-kind variants}

Two flavors of Chebyshev smoother appear in practice, differing only in which
spectral window the minimax polynomial targets and hence in their coefficient
formulas. The \emph{first-kind} smoother targets a two-sided window
$[\lambda_{\min},\lambda_{\max}]$ and threads one extra scalar through its
recurrence; it needs an estimate of \emph{both} ends of the band. The
\emph{fourth-kind} smoother (the modern default, the one in the pseudocode above)
targets the one-sided window $[0,\lambda_{\max}]$ and needs only the top of the
spectrum---a single estimate, no lower bound to guess. Both produce a degree-$k$
polynomial small across the high-frequency band; the fourth-kind's one-sided design
makes it more robust when the spectral estimate is rough, which is why Palace adopts
it~\citep{phillipsfischer2022}. At the level of the synthesis surface the two are a
single operator whose coefficient generator is swapped---the variant is absorbed into
the closure, exactly as the damping mode was absorbed into Jacobi's
\lstinline[style=l4style]|dinv|.

\begin{workedexample}{A two-step Chebyshev smoother}{smoo-cheby}
Take a tiny system with $\mat{D}^{-1}\mat{A}$ already diagonal in the eigenvector
basis, with scaled eigenvalues spread across $[\lambda_{\min},\lambda_{\max}]=[0.2,1.0]$.
Build the degree-$2$ minimax polynomial $p_2$ with $p_2(0)=1$ that is smallest on
$[0.2,1.0]$. Map the window to $[-1,1]$ by $t(\lambda)=(2\lambda-1.2)/0.8$ and use
$T_2(t)=2t^2-1$. The normalized symbol is
\[
  p_2(\lambda) = \frac{T_2\bigl(t(\lambda)\bigr)}{T_2\bigl(t(0)\bigr)},
  \qquad t(0) = \frac{-1.2}{0.8} = -1.5,
  \qquad T_2(-1.5) = 2(2.25)-1 = 3.5,
\]
so $p_2(\lambda) = \bigl(2\,t(\lambda)^2 - 1\bigr)/3.5$. Evaluate the damping at three
frequencies:
\begin{center}
\setlength{\tabcolsep}{10pt}\renewcommand{\arraystretch}{1.15}
\begin{tabular}{@{}llll@{}}
\toprule
mode & $\lambda$ & $t(\lambda)$ & $\abs{p_2(\lambda)}$ \\
\midrule
smooth (left of window)   & $0.05$ & $-1.375$ & $0.795$ \\
mid-band                  & $0.6$  & $0$      & $0.286$ \\
high (top of window)      & $1.0$  & $1.0$    & $0.286$ \\
\bottomrule
\end{tabular}
\end{center}
The polynomial reduces every mode \emph{inside} the band $[0.2,1.0]$ to at most
$\approx 0.29$ of its amplitude in a single two-step sweep---compare damped Jacobi
from \cref{wex:smoo-jacobi}, which needed three sweeps to reach $0.33$ at the top
mode alone. The smooth mode at $\lambda=0.05$, outside the targeted band, is left at
$0.795$, untouched by design. Two operator applies have damped the entire
high-frequency half of the spectrum below $0.29$, uniformly. That uniformity across
the band---not just at one mode---is what the Chebyshev construction buys over a
single Jacobi step.
\end{workedexample}

\section{Hiptmair relaxation: smoothing in \texorpdfstring{$\Hcurl$}{H(curl)}}
\label{sec:smoo-hiptmair}

The smoothers so far assume one thing the eigenvector picture made explicit: that
the operator has \emph{no} large null space at the smooth end, so that ``high
frequency'' and ``oscillatory'' coincide and a diagonal relaxation can reach the
modes a coarse grid misses. For the scalar problems of electrostatics this holds.
For the vector curl--curl operator of magnetostatics, eigenmode, and driven
analyses, it fails---and the failure is not subtle.

\subsection{The gradient null space a plain smoother cannot see}

Recall from \cref{ch:bc} and the de Rham complex of \cref{ch:derham} that the
curl--curl operator $\Curl(\nu\Curl\,\cdot)$, discretized on Nédélec edge elements,
has an \emph{enormous null space}: every gradient field $\Grad\psi$ has zero curl,
so $\mat{A}\,\mat{G}\psi = 0$ for the discrete gradient $\mat{G}$ that maps scalar
potentials (nodal $\Hone$) into edge fields ($\Hcurl$). This is not a handful of
modes; it is an entire subspace, the discrete gradient range, with dimension on the
order of the number of mesh vertices.

A diagonal smoother is helpless against error living in that null space. Run the
damping analysis again: error along a null-space mode $\vphi$ has $\mat{A}\vphi=0$,
so the residual $\vr=-\mat{A}\ve$ has \emph{no component} along it, and the Jacobi
update \cref{eq:smoo-jac-update} does nothing---the damping factor is $g=1$ exactly,
for the entire gradient subspace. These are not high-frequency modes the smoother
merely damps weakly; they are modes it cannot touch \emph{at all}, and there are
vast numbers of them. \Cref{fig:smoo-nullspace} pictures the situation: the smoother
operates happily in the part of the space the curl ``sees,'' and is blind to the
gradient subspace beside it.

\begin{figure}[t]
\centering
\begin{tikzpicture}[scale=1.0]
  % the full H(curl) space as a big rounded rectangle
  \draw[thick, simInk, rounded corners=4pt, fill=simBgGray] (-4.4,-1.5) rectangle (4.4,1.7);
  \node[font=\small\bfseries, simInk, anchor=north west] at (-4.3,1.6) {$\Hcurl$ (edge dofs)};
  % the gradient null space subspace
  \draw[thick, simGreen, rounded corners=4pt, fill=simBgGreen] (-4.1,-1.25) rectangle (-0.4,1.0);
  \node[font=\footnotesize, simGreen, align=center] at (-2.25,0.55)
    {gradient null space\\$\operatorname{range}\mat{G}=\{\Grad\psi\}$\\\scriptsize($\mat{A}\,\Grad\psi = 0$)};
  % the rest: curl sees it
  \node[font=\footnotesize, simBlue, align=center] at (2.3,0.3)
    {curl-active part\\\scriptsize(smoother damps these)};
  % a "smoother" reaching only the right part
  \node[opbox, anchor=south] at (2.3,-1.35) {diagonal smoother};
  \draw[flow, simBlue] (2.3,-1.0) -- (2.3,-0.1);
  % blocked arrow into the null space
  \draw[->, simRust, thick] (-2.25,-1.0) -- (-2.25,-0.2);
  \node[font=\scriptsize, simRust] at (-2.25,-1.35) {$g=1$: no effect};
  \node[draw, cross out, simRust, thick, minimum size=4mm] at (-2.25,-0.55) {};
\end{tikzpicture}
\caption[The gradient null space a plain smoother misses]{The edge-dof space
$\Hcurl$ splits into the curl-active part, where a diagonal smoother damps error
normally, and the \emph{gradient null space} $\operatorname{range}\mat{G}$, where
$\mat{A}\,\Grad\psi=0$. On the null space the residual vanishes, the smoother's
damping factor is exactly $1$, and error simply persists. For an edge-element
problem this blind spot is a whole subspace---the smoother on its own cannot
converge. The cure is to relax in that subspace separately.}
\label{fig:smoo-nullspace}
\end{figure}

\begin{pitfall}
This is the specific reason edge-element ($\Hcurl$) problems need more than a
diagonal or even a Chebyshev smoother. The polynomial smoother of
\cref{sec:smoo-cheby} sharpens damping in the curl-active part but is just as blind
to the gradient null space as Jacobi---$p_k(0)=1$, and the entire null space sits at
$\lambda=0$. A multigrid built on a plain smoother will stall, its error trapped in
the gradient modes the coarse grid cannot drain because the smoother never made
them smooth. Magnetostatics, eigenmode, and driven analyses all hit this wall.
\end{pitfall}

\subsection{Relax in two spaces: the auxiliary-space idea}

The fix is due to Hiptmair, and in the de Rham setting of \cref{ch:derham} it is
beautifully natural: if the smoother cannot reach error in the gradient subspace
$\operatorname{range}\mat{G}$, then \emph{relax in that subspace too}, using the
discrete gradient $\mat{G}$ as the bridge. The discrete gradient maps the scalar
potential space (nodal $\Hone$) into the edge space ($\Hcurl$); its transpose
$\mat{G}^{\mathsf T}$ maps back. The combination
$\mat{A}_G = \mat{G}^{\mathsf T}\mat{A}\,\mat{G}$ is the curl--curl operator
\emph{seen from the scalar potential space}---a small, well-behaved scalar operator
with no null-space pathology of its own, on which an ordinary point smoother $B_G$
works perfectly.

One sweep of \emph{Hiptmair distributive relaxation} does two relaxations in
sequence (\cref{fig:smoo-hiptmair}):
\begin{align}
  \vy &\;\leftarrow\; \vy + B\,(\vx - \mat{A}\vy)
      && \text{primary (edge) relaxation,} \label{eq:smoo-hipt-prim}\\
  \vy &\;\leftarrow\; \vy + \mat{G}\,B_G\,\mat{G}^{\mathsf T}(\vx - \mat{A}\vy)
      && \text{auxiliary (gradient) correction.} \label{eq:smoo-hipt-aux}
\end{align}
The primary step \cref{eq:smoo-hipt-prim} is an ordinary smoother (Jacobi or
Chebyshev) on the edge operator $\mat{A}$---it handles the curl-active error. The
auxiliary step \cref{eq:smoo-hipt-aux} forms the residual, \emph{restricts} it to
the scalar space with $\mat{G}^{\mathsf T}$, relaxes there with $B_G$ on the
well-conditioned $\mat{A}_G$, and \emph{prolongs} the correction back into the edge
space with $\mat{G}$. Because $\mat{G}$ maps onto exactly the subspace the primary
smoother was blind to, the auxiliary step reaches precisely the error the first step
could not. Together they smooth the whole space. This is the discrete realization of
the de Rham auxiliary-space idea (Hiptmair--Xu) underlying the AMS
preconditioner~\citep{hiptmairxu2007}.

\begin{figure}[t]
\centering
\begin{tikzpicture}[node distance=9mm and 20mm]
  % primary (edge) space row
  \node[space, fill=simBgBlue, draw=simBlue] (Hcurl) {$\Hcurl$ \\ \scriptsize edge dofs};
  \node[opbox, right=of Hcurl, draw=simBlue] (B) {primary relax \\ $B$ on $\mat{A}$};
  \node[product, right=of B] (yout) {relaxed $\vy$};
  % auxiliary (scalar) space row, below
  \node[space, below=16mm of Hcurl] (H1) {$\Hone$ \\ \scriptsize nodal dofs};
  \node[opbox, right=of H1, draw=simGreen] (BG) {auxiliary relax \\ $B_G$ on $\mat{A}_G$};

  % residual into auxiliary space
  \draw[flow] (Hcurl) -- node[above,font=\scriptsize]{$\vx-\mat{A}\vy$} (B);
  \draw[flow] (B) -- (yout);
  % restrict G^T down
  \draw[flow, simGreen] (B.south) -- ++(0,-0.4) -| node[pos=0.25,above,font=\scriptsize]{restrict $\mat{G}^{\mathsf T}$} (H1.north);
  \draw[flow, simGreen] (H1) -- (BG);
  % prolong G up
  \draw[flow, simGreen] (BG.north) -- ++(0,0.4) -| node[pos=0.25,below,font=\scriptsize]{prolong $\mat{G}$} (yout.south);

  % labels
  \node[cornertag, anchor=west] at ($(Hcurl.north west)+(0,0.5)$) {\scriptsize primary space};
  \node[cornertag, anchor=west] at ($(H1.south west)+(0,-0.5)$) {\scriptsize auxiliary (gradient) space};
\end{tikzpicture}
\caption[Hiptmair primary + auxiliary relaxation]{One Hiptmair distributive
relaxation sweep. The primary smoother $B$ relaxes the edge operator $\mat{A}$ in
$\Hcurl$, damping the curl-active error. The residual is then restricted by the
discrete gradient transpose $\mat{G}^{\mathsf T}$ into the scalar space $\Hone$,
relaxed there by $B_G$ on the well-conditioned $\mat{A}_G=\mat{G}^{\mathsf T}\mat{A}\mat{G}$,
and the correction prolonged back by $\mat{G}$. The auxiliary leg reaches exactly the
gradient null space the primary leg could not. The two legs run in sequence: the
auxiliary residual is formed \emph{after} the primary update.}
\label{fig:smoo-hiptmair}
\end{figure}

In the synthesis surface the whole sweep is one constructed operator. It captures
the two operators $\mat{A}$ and $\mat{A}_G$, the discrete gradient $\mat{G}$, two
point smoothers $B$ and $B_G$ (each a Chebyshev closure), and a sweep count; its
apply is the two legs above, repeated:

\begin{lstlisting}[style=l4style]
-- Hiptmair distributive relaxation: relax in the primary edge space, then
-- correct through the discrete gradient G in the auxiliary scalar space.
hiptmair_smoother :: HiptmairOp[N,M] -> Tensor[N] -> Tensor[N] -> Tensor[N]
hiptmair_smoother op x y = iterate op.pc_it sweep y
  where
    sweep y =
      let y1 = y `add` (op.B `apply2` (x, y))          -- primary leg (edge)
          r  = x `sub` (op.A `apply` y1)               -- residual after primary
          xG = pin_ess op.ess_G (op.G `applyT` r)      -- restrict G^T r, pin Dirichlet dofs
          yG = op.B_G `apply` xG                        -- relax in scalar space
      in  y1 `add` (op.G `apply` yG)                    -- prolong G yG, accumulate
\end{lstlisting}

Two implementation points the brief calls out. The auxiliary residual is formed
\emph{after} the primary leg updates \lstinline[style=l4style]|y|, so the two legs
do not commute---the sweep is a \emph{multiplicative} (Gauss-Seidel-between-spaces)
composition, not an additive one; conflating them changes the operator. And the
restricted auxiliary residual \lstinline[style=l4style]|xG| has its essential
(Dirichlet) degrees of freedom pinned to zero before the auxiliary relaxation, the
same boundary-condition surgery \cref{ch:bc} performed on the primary system.

\begin{notationbox}
Palace's distributive smoother chooses Chebyshev for \emph{both} point smoothers
$B$ and $B_G$, never a Gauss--Seidel sweep. That is deliberate. A true distributive
relaxation would call for a triangular Gauss--Seidel solve in each space, but
Gauss--Seidel is a sequential forward/back substitution that does not parallelize
and runs poorly on a GPU. Replacing it with two inner-product-free Chebyshev sweeps
keeps the entire smoother communication-light, at the cost of an inexact---but
perfectly adequate---auxiliary relaxation in place of an exact auxiliary solve.
\end{notationbox}

\section{Smoothers as preconditioners and as multigrid components}
\label{sec:smoo-ascomponents}

Every smoother in this chapter is, structurally, an apply-function of the kind
\cref{ch:precond} defined: it consumes a vector and returns a vector, it is built
once at setup and applied many times, and---for the zero-initial-guess, single-sweep
case---it is \emph{linear} in its input, the exact shape a Krylov method wants in its
preconditioner slot. Indeed Palace installs the Jacobi smoother directly as a Krylov
preconditioner where a cheap diagonal scaling is all the problem needs.

But used as a standalone preconditioner a smoother is weak, and the eigenvector
picture says exactly why: it leaves the smooth, low-frequency error nearly
untouched. The damping factor returns to $1$ as $\lambda\to0$ for Jacobi
(\cref{fig:smoo-jacsymbol}), for Chebyshev (\cref{fig:smoo-chebpoly}, since
$p_k(0)=1$), and the gradient null space defeats both. A smoother makes the error
\emph{smooth}; it does not make it \emph{small}.

That leftover smooth error is precisely what a coarse grid is good at. \emph{Multigrid}
pairs the smoother with a coarse-grid correction: smooth on the fine grid to kill the
oscillatory error, project the smooth remainder down to a coarse grid where it looks
oscillatory and is cheap to resolve, solve there, and prolong the correction back.
\Cref{fig:smoo-vcycle} previews the structure---a V-cycle with a smoother bolted into
every level, descending and ascending. The smoother is the component this chapter
built; the recursion that surrounds it is the subject of \cref{ch:multigrid}.

\begin{figure}[t]
\centering
\begin{tikzpicture}[scale=1.0, every node/.style={font=\footnotesize}]
  % levels as horizontal bands at descending y, V shape
  \def\lw{1.6}
  % fine -> coarse going down-right, then back up
  \coordinate (f0) at (0,0);
  \coordinate (f1) at (1.3,-1.0);
  \coordinate (f2) at (2.6,-2.0);
  \coordinate (cc) at (3.9,-2.8);
  \coordinate (b2) at (5.2,-2.0);
  \coordinate (b1) at (6.5,-1.0);
  \coordinate (b0) at (7.8,0);
  % connecting V path
  \draw[flow, simGray] (f0)--(f1)--(f2)--(cc)--(b2)--(b1)--(b0);
  % smoother nodes on the way down
  \node[opbox, fill=simBgBlue, draw=simBlue] at (f0) {smooth};
  \node[opbox, fill=simBgBlue, draw=simBlue] at (f1) {smooth};
  \node[opbox, fill=simBgBlue, draw=simBlue] at (f2) {smooth};
  % coarse solve at the bottom
  \node[product, fill=simBgRust, draw=simRust] at (cc) {coarse \\ solve};
  % smoother nodes on the way up
  \node[opbox, fill=simBgBlue, draw=simBlue] at (b2) {smooth};
  \node[opbox, fill=simBgBlue, draw=simBlue] at (b1) {smooth};
  \node[product, fill=simBgGreen, draw=simGreen] at (b0) {result};
  % side labels
  \node[anchor=east, simGray] at ($(f0)+(-0.6,0)$) {finest grid};
  \node[anchor=north, simGray] at ($(cc)+(0,-0.55)$) {coarsest grid};
  \node[anchor=west, simGray, align=left] at ($(b0)+(0.5,0)$) {smooth = this\\chapter's\\smoother};
\end{tikzpicture}
\caption[Smoothers inside a multigrid V-cycle (preview)]{The multigrid V-cycle of
\cref{ch:multigrid}, previewed. Descending from the finest grid, a smoother runs at
each level to remove the oscillatory error; the smooth remainder is carried to the
coarsest grid where it is solved cheaply; ascending, the correction is prolonged
back with a smoother at each level again. Every blue node is a smoother of the kind
this chapter built---Jacobi, Chebyshev, or Hiptmair. The smoother is a poor solver
in isolation and the indispensable component here.}
\label{fig:smoo-vcycle}
\end{figure}

\begin{keyidea}
A smoother is built to be a multigrid \emph{component}, not a solver. Standalone it
stalls on smooth error; inside a V-cycle, paired with coarse-grid correction, that
weakness is exactly covered---the smoother handles the high frequencies the coarse
grid cannot resolve, the coarse grid handles the low frequencies the smoother cannot
damp. The division of labor is the design. This is why we measured smoothers by
their per-frequency damping factor rather than by how fast they solve: as a
component, damping the high band cheaply is the whole job.
\end{keyidea}

\summarysection
Writing the error in the eigenvector basis of $\mat{A}$ reduces a smoother to a
single curve: its per-frequency damping factor $g(\lambda)$. A \emph{smoother} is an
iteration whose $g$ is small at high frequency and near $1$ at low frequency, so it
cheaply removes oscillatory error while leaving smooth error---a deliberately
lopsided ``high-pass error filter'' that is a bad standalone solver but the heart of
multigrid. \emph{Damped Jacobi} scales the residual by $\omega\mat{D}^{-1}$; its
symbol is the affine $1-\omega\mu$, optimally damped at
$\omega^\star=2/(\mu_{\min}+\mu_{\max})$, and undamped Jacobi can fail to smooth at
all. \emph{Chebyshev} applies a fixed-degree minimax polynomial $p_k(\mat{D}^{-1}\mat{A})$
shaped to be uniformly small across a target spectral window; it is matrix-free,
inner-product-free, needs only a $\lambda_{\max}$ estimate, and must be evaluated by
its stable three-term recurrence rather than its unstable monomial expansion.
\emph{Hiptmair distributive relaxation} cures the gradient null space that defeats a
plain smoother on edge-element curl--curl problems, by relaxing in the primary edge
space and in the auxiliary scalar space bridged by the discrete gradient $\mat{G}$.
All three are apply-functions; standalone they are weak, but as the smoothing step of
the V-cycle in \cref{ch:multigrid} they are essential.

\furtherreading
The optimal damping analysis and the eigenvector view of stationary iterations are
standard; Saad~\citep[Ch.~4]{saad2003} is the canonical reference. The fourth-kind
Chebyshev smoother, its coefficient formulas, and the load-bearing stability of the
recurrence over the monomial form are developed by Phillips and
Fischer~\citep{phillipsfischer2022}, whose analysis Palace follows directly. The
auxiliary-space (AMS) framework for $\Hcurl$ and $\Hdiv$ problems, of which Hiptmair
distributive relaxation is the smoother realization, is laid out by Hiptmair and
Xu~\citep{hiptmairxu2007}; the de Rham structure it exploits is \cref{ch:derham},
and the V-cycle it slots into is \cref{ch:multigrid}.

\exercisessection
\begin{enumerate}[nosep]
  \item \textbf{Damping symbol.} For damped Jacobi on a model problem with scaled
  spectrum $\mu\in[0,2]$, write the symbol $g(\mu)=1-\omega\mu$. (a) For which
  $\omega$ does $\abs{g(\mu_{\max})}=\abs{g(2)}$ first exceed $1$? (b) Show that any
  $\omega>1$ leaves at least one mode amplified. (c) Conclude that
  $\omega\le1$ is necessary for convergence on this spectrum.

  \item \textbf{Optimal $\omega$ for the half-band.} Repeat the balance of
  \cref{eq:smoo-jac-optomega} for the target band $[\mu_{\mathrm{lo}},\mu_{\max}]
  =[1,2]$ and verify $\omega^\star=2/3$. Then compute the worst-case damping
  $\max_{\mu\in[1,2]}\abs{g(\mu)}$ at this $\omega^\star$. How many sweeps to drive
  the worst high-frequency mode below $10^{-3}$?

  \item \textbf{Why undamped fails.} Using the 1-D Laplacian eigenvalues
  $\lambda_k=4\sin^2(k\pi/2(n+1))$ from \cref{wex:smoo-jacobi}, compute the
  undamped ($\omega=1$) damping factor at the top mode $k=n$ as $n\to\infty$.
  What does it approach, and why does that make undamped Jacobi useless as a
  smoother even when it does not diverge?

  \item \textbf{Chebyshev versus Jacobi cost.} A degree-$k$ Chebyshev sweep costs
  $k$ operator applies; a damped-Jacobi sweep costs $1$. \Cref{wex:smoo-cheby}
  reached worst-case band damping $0.29$ in one degree-$2$ sweep ($2$ applies). How
  many damped-Jacobi sweeps (\cref{wex:smoo-jacobi}, worst-case $1/3$ per sweep)
  match that damping, and how many applies is that? Which wins on this band?

  \item \textbf{The monomial trap.} Explain in one paragraph why expanding
  $p_k(\mat{D}^{-1}\mat{A})$ into $\sum_j c_j(\mat{D}^{-1}\mat{A})^j$ and summing the
  powers is numerically unstable for moderate $k$, while the three-term recurrence of
  \cref{fig:smoo-chebdataflow} computing the same polynomial is stable. (Hint:
  consider the size of $\norm{(\mat{D}^{-1}\mat{A})^j}$ and the signs of the $c_j$.)

  \item \textbf{The gradient blind spot.} Let $\vphi=\mat{G}\psi$ be a gradient
  field, so $\mat{A}\vphi=0$. Show directly that a damped-Jacobi sweep
  \cref{eq:smoo-jac-update} starting from error $\ve=\vphi$ leaves $\ve$ unchanged.
  Then show that the auxiliary leg \cref{eq:smoo-hipt-aux} of the Hiptmair sweep
  \emph{does} reduce it, provided $B_G$ is a contraction on $\mat{A}_G$. (You may
  assume $\mat{G}^{\mathsf T}\mat{G}$ is invertible on the relevant subspace.)

  \item \textbf{Multiplicative, not additive.} The two Hiptmair legs
  \cref{eq:smoo-hipt-prim}--\cref{eq:smoo-hipt-aux} form the auxiliary residual from
  the \emph{updated} $\vy$. Write the single-sweep error-propagation operator for
  the multiplicative composition, and compare it to the (incorrect) additive version
  that uses the \emph{same} residual for both legs. On a mode in the gradient null
  space, which one converges and which one stalls? Why does the order matter?

  \item \textbf{Smoother as preconditioner.} Argue that for zero initial guess and a
  single sweep, the Chebyshev smoother $\vx\mapsto p_k(\mat{D}^{-1}\mat{A})\vx$ is a
  \emph{linear} operator (hence a valid Krylov preconditioner), while a Hiptmair
  sweep with nonzero initial guess is \emph{affine}, not linear. Which property does
  a multigrid V-cycle's smoothing step rely on, and which does an outer Krylov
  preconditioner slot require?
\end{enumerate}

\chapter{Geometric Multigrid: The V-Cycle}
\label{ch:multigrid}

\chapterorient{A smoother (\cref{ch:smoothers}) is a bargain that only half
works: a few cheap sweeps annihilate the jagged, high-frequency part of the error
and then stall, because what remains is \emph{smooth}---and smooth error is exactly
what a smoother cannot touch. This chapter turns that stall into a strategy.
The trick is a change of perspective: error that looks smooth on a fine grid looks
\emph{oscillatory} on a coarser one, where a smoother can kill it again, and far
more cheaply. Apply the idea once and you have a two-grid method; apply it to
itself---recursively, all the way down to a grid so small the problem is trivial---and
you have the V-cycle, the central algorithm of this chapter and one of the great
ideas of numerical computing. By the end you will be able to read the V-cycle as a
short recursive function, explain from the complementary-frequencies picture
\emph{why} it works, and state the property that makes it the optimal
preconditioner of \cref{ch:precond}: a convergence rate that does not degrade as the
mesh is refined, at a cost that grows only linearly in the number of unknowns.}

\section{The trouble a smoother leaves behind}

\Cref{ch:smoothers} studied the classical \emph{smoothers}---weighted Jacobi,
Gauss--Seidel, Chebyshev---the cheap stationary iterations that, applied to a linear
system $\mat{A}\vx=\vb$ arising from a discretized elliptic PDE, contract some
components of the error far faster than others. Recall the central fact established
there. Write the error after $k$ sweeps as $\ve_k = \vx - \vx_k$, where $\vx_k$ is
the current iterate and $\vx$ the exact solution. Expand $\ve_k$ in the eigenvectors
of the iteration: for a model problem each eigenvector is a \emph{Fourier mode}, a
sinusoid of some spatial frequency, and one sweep of a damped smoother multiplies
the mode of frequency $\theta$ by a factor $g(\theta)$ with $\abs{g(\theta)}<1$.

The defining property of a good smoother is that $\abs{g(\theta)}$ is small---near
zero---for the \emph{high} frequencies and close to $1$ for the \emph{low} ones.
A weighted-Jacobi sweep on the 1-D Laplacian, for instance, damps the most
oscillatory mode by a factor near $\tfrac13$ but the smoothest mode by a factor like
$1-O(h^2)$, agonizingly close to $1$. After a handful of sweeps the high-frequency
error is gone and the low-frequency error is essentially untouched. The iteration
appears to converge for three or four steps and then crawls.

\begin{keyidea}
A smoother is \emph{half} a solver. It cheaply annihilates the high-frequency
(oscillatory) part of the error and leaves the low-frequency (smooth) part almost
unchanged. The error that survives a few smoothing sweeps is, by construction,
\emph{smooth}. That residual smoothness is not a failure to be patched---it is a
\emph{signal}: smooth error is the part a coarser grid can represent, and on a
coarser grid it is no longer smooth.
\end{keyidea}

This is the opening every multigrid method exploits. The error left behind is smooth,
and smoothness is a \emph{grid-relative} notion. A function that wiggles slowly across
a fine mesh of $N$ points wiggles \emph{quickly} across a coarse mesh of $N/8$ points
covering the same domain---because there are fewer points for it to wiggle across. The
smoother that stalled on the fine grid will bite again on the coarse grid, where the
same error is now high-frequency. And the coarse grid is cheaper. The entire method is
the disciplined pursuit of this one observation.

\section{Complementary frequencies: the heart of the matter}
\label{sec:complementary}

Make the observation precise. Take the 1-D model: the domain $[0,1]$, a fine grid of
$n$ interior points at spacing $h = 1/(n+1)$, and a coarse grid of every \emph{other}
fine point, spacing $2h$. A Fourier mode on the fine grid is
\[
  \phi^h_j(x) = \sin(j\pi x), \qquad j = 1, 2, \dots, n,
\]
sampled at the fine nodes $x = h, 2h, \dots, nh$. The index $j$ is the
\emph{frequency}: $j=1$ is the smoothest half-sine; $j=n$ is the most oscillatory mode
the grid can resolve, alternating sign at every node. We split the modes at the
midpoint $n/2$:
\begin{itemize}[nosep]
\item \textbf{low (smooth) modes}, $1 \le j < n/2$: oscillate slowly; the smoother
  barely touches them ($\abs{g}\approx 1$);
\item \textbf{high (oscillatory) modes}, $n/2 \le j \le n$: oscillate rapidly; the
  smoother annihilates them ($\abs{g}\ll 1$).
\end{itemize}

Now ask the decisive question: what does a \emph{low} fine-grid mode look like when we
sample it on the coarse grid, at the nodes $x = 2h, 4h, \dots$ alone? The coarse grid
has half as many points, so it can only resolve frequencies up to $n/2$. A smooth
fine-grid mode, $j < n/2$, is \emph{also} a legitimate mode on the coarse grid---but
relative to the coarse grid's own resolution limit $n/2$, a mode at $j$ that was
``low'' on the fine grid sits a larger fraction of the way up toward the coarse limit.
The smoothest fine modes, the ones that defeated the smoother, become genuine,
killable, mid-to-high frequency content on the coarse grid. \Cref{fig:complementary}
draws exactly this: one smooth error sampled on both grids, oscillatory-looking on the
coarse one.

\begin{figure}[t]
\centering
\begin{tikzpicture}
\begin{axis}[
  width=12.8cm, height=5.4cm,
  xlabel={position $x$}, ylabel={error $e(x)$},
  xmin=0, xmax=1, ymin=-1.25, ymax=1.45,
  xtick={0,0.25,0.5,0.75,1}, ytick={-1,0,1},
  axis lines=left, tick align=outside,
  legend style={at={(0.5,1.02)}, anchor=south, legend columns=2,
    font=\footnotesize, draw=simGray!50},
  every axis plot/.append style={thick},
]
  % the underlying smooth error: a slowly-varying field (looks tame everywhere)
  \addplot[simBlue, very thick, samples=200, domain=0:1]
    {0.95*sin(deg(3.5*pi*x))*exp(-0.35*x)};
  \addlegendentry{smooth error (the continuous field)}
  % fine-grid samples: dense, follow the curve -> looks smooth
  \addplot[only marks, mark=*, mark size=1.0pt, simGreen]
    coordinates {
      (0.04,0.395)(0.08,0.685)(0.12,0.846)(0.16,0.860)(0.20,0.730)
      (0.24,0.479)(0.28,0.150)(0.32,-0.205)(0.36,-0.521)(0.40,-0.742)
      (0.44,-0.829)(0.48,-0.766)(0.52,-0.566)(0.56,-0.262)(0.60,0.087)
      (0.64,0.414)(0.68,0.661)(0.72,0.785)(0.76,0.764)(0.80,0.603)
      (0.84,0.334)(0.88,0.006)(0.92,-0.318)(0.96,-0.584)};
  \addlegendentry{fine-grid samples (smooth)}
  % coarse-grid samples: every other -> when connected, looks oscillatory
  \addplot[simRust, thick, mark=square*, mark size=1.7pt, densely dashed]
    coordinates {
      (0.04,0.395)(0.12,0.846)(0.20,0.730)(0.28,0.150)(0.36,-0.521)
      (0.44,-0.829)(0.52,-0.566)(0.60,0.087)(0.68,0.661)(0.76,0.764)
      (0.84,0.334)(0.92,-0.318)};
  \addlegendentry{coarse-grid samples (now oscillatory)}
\end{axis}
\end{tikzpicture}
\caption[Complementary frequencies]{The single idea behind multigrid. The blue curve
is a smooth error---the kind a smoother leaves behind, varying slowly across the
domain. Sampled on the \emph{fine} grid (green dots, every node), it traces the smooth
curve faithfully and a smoother sees a low-frequency mode it cannot damp. Sampled on
the \emph{coarse} grid (rust squares, every \emph{other} node, connected by the dashed
line), the \emph{same} error now swings sharply from sample to sample---it has become a
mid-to-high frequency mode relative to the coarse grid's coarser resolution, exactly
the kind a smoother \emph{can} annihilate. Smoothness is grid-relative; the coarse grid
turns the smoother's blind spot into its target.}
\label{fig:complementary}
\end{figure}

\begin{notationbox}
We will say a mode is \textbf{smooth} or \textbf{oscillatory} \emph{relative to a
grid}, never absolutely. A fixed error field is smooth on a fine grid and oscillatory
on a coarse one. ``High frequency'' means ``near the resolution limit of the grid in
question.'' The whole method rests on the fact that the smoother's reach---which modes
it damps---is tied to the grid's resolution limit, so changing the grid changes which
modes are in reach.
\end{notationbox}

This is the \emph{complementary-frequencies principle}. The smoother and the coarse
grid have complementary competences: the smoother handles everything oscillatory on the
fine grid; the coarse grid handles everything smooth on the fine grid, by re-presenting
it as something oscillatory that \emph{its} smoother can handle. Between them they cover
the whole spectrum. Neither alone is a solver; together they are. The two-grid method
of the next section is the minimal machine that makes them cooperate.

\section{The two-grid method}
\label{sec:twogrid}

Suppose we have a fine grid carrying the system $\mat{A}\vx=\vb$ and a single coarse
grid. We want a method that smooths on the fine grid (killing oscillatory error) and
\emph{also} corrects the smooth error by working on the coarse grid. The obstacle is
that the smooth error lives in the fine-grid solution, while the coarse grid speaks a
different, smaller vocabulary of vectors. We need a way to carry information between the
two---and an equation on the coarse grid for the coarse grid to solve.

\subsection{The residual equation}

The bridge is the \emph{residual equation}, which we have met before but which becomes
load-bearing here. Let $\vy$ be the current fine-grid approximation, after some
smoothing. Its error is $\ve = \vx - \vy$, unknown (it contains the exact solution). But
its \emph{residual},
\[
  \vr = \vb - \mat{A}\vy,
\]
is fully computable, and error and residual are linked by the exact linear identity
\[
  \mat{A}\,\ve = \mat{A}\vx - \mat{A}\vy = \vb - \mat{A}\vy = \vr.
\]
So $\mat{A}\,\ve = \vr$: the error solves the \emph{same} operator equation as the
original, but with the residual on the right. This is the \emph{error equation}. If we
could solve it we would have $\ve$ exactly and could finish in one step, $\vx = \vy +
\ve$. We cannot---solving $\mat{A}\,\ve=\vr$ on the fine grid is the original problem
again. But after smoothing, $\ve$ is \emph{smooth}, and a smooth $\ve$ is exactly what a
\emph{coarse} grid can represent well. So we solve the error equation
\emph{approximately, on the coarse grid}, and use the coarse solution to correct $\vy$.

\subsection{Restriction and prolongation}

To pose the error equation on the coarse grid we need two transfer operators
(\cref{fig:transfer}):
\begin{itemize}[nosep]
\item \textbf{Restriction} $\mat{R}\colon \text{fine}\to\text{coarse}$, which carries the
  fine-grid residual $\vr$ down to a coarse-grid residual $\vr_c = \mat{R}\vr$;
\item \textbf{Prolongation} $\mat{P}\colon \text{coarse}\to\text{fine}$, which carries a
  coarse-grid correction $\ve_c$ up to a fine-grid correction $\mat{P}\ve_c$.
\end{itemize}
Prolongation is \emph{interpolation}: given values at the coarse nodes, fill in the fine
nodes between them. In 1-D, linear interpolation copies the value at each coarse node to
the fine node sitting on it and averages the two neighbours for the fine node in between:
\[
  (\mat{P}\ve_c)_{2i} = (\ve_c)_i, \qquad
  (\mat{P}\ve_c)_{2i+1} = \tfrac12\bigl((\ve_c)_i + (\ve_c)_{i+1}\bigr).
\]
The matrix $\mat{P}$ is tall and thin: $N_f \times N_c$ with $N_f > N_c$, one column per
coarse node.

\begin{figure}[t]
\centering
\begin{tikzpicture}[scale=1.0, every node/.style={font=\scriptsize}]
  % fine grid (top), coarse grid (bottom)
  \def\fy{1.9} \def\cy{0}
  % fine nodes 0..8
  \foreach \i in {0,...,8}{
    \fill[simGreen] (\i*1.0,\fy) circle (2.0pt);}
  \draw[simGray] (0,\fy)--(8,\fy);
  \node[simGreen, anchor=south] at (4,\fy+0.18) {fine grid ($N_f$ nodes, spacing $h$)};
  % coarse nodes at 0,2,4,6,8
  \foreach \i in {0,2,4,6,8}{
    \fill[simRust] (\i*1.0,\cy) circle (2.6pt);}
  \draw[simGray] (0,\cy)--(8,\cy);
  \node[simRust, anchor=north] at (4,\cy-0.18) {coarse grid ($N_c$ nodes, spacing $2h$)};
  % prolongation (coarse -> fine), left half: interpolation arrows
  \begin{scope}
    % coincident copies
    \foreach \i in {0,2,4}{
      \draw[-{Stealth[length=1.6mm]}, simBlue, thick] (\i,\cy+0.12)--(\i,\fy-0.12);}
    % interpolated midpoints (averages of two neighbours)
    \foreach \i in {1,3}{
      \draw[-{Stealth[length=1.4mm]}, simBlue, opacity=0.55] (\i-1,\cy+0.12)--(\i,\fy-0.12);
      \draw[-{Stealth[length=1.4mm]}, simBlue, opacity=0.55] (\i+1,\cy+0.12)--(\i,\fy-0.12);}
    \node[simBlue, anchor=east, font=\footnotesize\bfseries] at (-0.25,1.0)
      {$\mat{P}$ (prolong)};
  \end{scope}
  % restriction (fine -> coarse), right half: weighted gather arrows
  \begin{scope}
    \foreach \i in {6,8}{
      \draw[-{Stealth[length=1.6mm]}, simRust, thick, densely dashed] (\i,\fy-0.12)--(\i,\cy+0.12);}
    \draw[-{Stealth[length=1.4mm]}, simRust, opacity=0.6, densely dashed] (5,\fy-0.12)--(6,\cy+0.12);
    \draw[-{Stealth[length=1.4mm]}, simRust, opacity=0.6, densely dashed] (7,\fy-0.12)--(6,\cy+0.12);
    \node[simRust, anchor=west, font=\footnotesize\bfseries] at (8.25,1.0)
      {$\mat{R}=\mat{P}^{\mathsf{T}}$ (restrict)};
  \end{scope}
\end{tikzpicture}
\caption[Prolongation and restriction]{The two grid-transfer operators between a fine
grid (green, spacing $h$) and a coarse grid (rust, spacing $2h$) holding every other
node. \emph{Left:} prolongation $\mat{P}$ (blue, solid arrows) is interpolation---each
coarse node's value is copied to the coincident fine node, and each in-between fine node
takes the average of its two coarse neighbours. \emph{Right:} restriction
$\mat{R}=\mat{P}^{\mathsf T}$ (rust, dashed arrows) is the transpose---it gathers each
coarse node's value as a weighted sum of nearby fine values, with weights mirroring the
interpolation stencil. Making restriction the transpose of prolongation is what keeps the
coarse operator symmetric and the whole method variationally consistent.}
\label{fig:transfer}
\end{figure}

Restriction is the reverse trip, and the natural choice---the one that makes the algebra
work out---is the \emph{transpose} of prolongation, up to scaling:
\[
  \mat{R} = \mat{P}^{\mathsf T}.
\]
Why the transpose? Because residuals and corrections are dual objects. A correction is a
\emph{primal} vector (it adds to the solution); a residual is a \emph{dual} vector (it
pairs with corrections through $\inner{\vr}{\ve}$ to measure work). Interpolating a
primal vector \emph{up} with $\mat{P}$ forces the dual vector to go \emph{down} with
$\mat{P}^{\mathsf T}$, so that the pairing $\inner{\vr}{\mat{P}\ve_c} =
\inner{\mat{P}^{\mathsf T}\vr}{\ve_c}$ is preserved between the grids. We will see in the
next subsection that this single choice is exactly what makes the coarse operator
symmetric.

\subsection{The Galerkin coarse operator}

The error equation on the fine grid is $\mat{A}\,\ve=\vr$. We want its coarse-grid
counterpart, $\mat{A}_c\,\ve_c = \vr_c$, with $\vr_c=\mat{P}^{\mathsf T}\vr$. What should
$\mat{A}_c$ be? Substitute the ansatz that the fine correction is the prolonged coarse
correction, $\ve \approx \mat{P}\ve_c$, into $\mat{A}\,\ve = \vr$ and restrict both
sides:
\[
  \mat{P}^{\mathsf T}\mat{A}\,\mat{P}\,\ve_c = \mat{P}^{\mathsf T}\vr.
\]
The operator on the left \emph{is} the right coarse operator:
\begin{equation}
  \mat{A}_c \;=\; \mat{P}^{\mathsf T}\mat{A}\,\mat{P}.
  \label{eq:galerkin}
\end{equation}
This is the \emph{Galerkin coarse operator}, or \emph{Galerkin coarsening}. It is not the
operator you would get by re-discretizing the PDE on the coarse grid (though for simple
problems the two agree up to scaling); it is the operator \emph{induced} by the fine
operator and the transfer maps. \Cref{eq:galerkin} has two virtues that no other choice
of $\mat{A}_c$ shares.

\begin{lemma}[Symmetry and definiteness are inherited]
\label{lem:galerkin}
If $\mat{A}$ is symmetric positive definite and $\mat{P}$ has full column rank, then the
Galerkin coarse operator $\mat{A}_c = \mat{P}^{\mathsf T}\mat{A}\mat{P}$ is symmetric
positive definite.
\end{lemma}
\begin{proof}
Symmetry: $\mat{A}_c^{\mathsf T} = (\mat{P}^{\mathsf T}\mat{A}\mat{P})^{\mathsf T} =
\mat{P}^{\mathsf T}\mat{A}^{\mathsf T}\mat{P} = \mat{P}^{\mathsf T}\mat{A}\mat{P} =
\mat{A}_c$, using $\mat{A}^{\mathsf T}=\mat{A}$. Definiteness: for any coarse vector
$\vv\ne\mathbf 0$, $\inner{\vv}{\mat{A}_c\vv} = \inner{\vv}{\mat{P}^{\mathsf
T}\mat{A}\mat{P}\vv} = \inner{\mat{P}\vv}{\mat{A}(\mat{P}\vv)} > 0$, because $\mat{P}\vv
\ne \mathbf 0$ (full column rank) and $\mat{A}$ is positive definite.
\end{proof}

So the coarse problem is again an SPD system: every tool we have for SPD systems---the
conjugate gradient method of \cref{ch:cg}, another round of smoothing, another level of
coarsening---applies to it unchanged. This self-similarity is what makes recursion
possible. The second virtue is \emph{variational consistency}: with $\mat{R} =
\mat{P}^{\mathsf T}$ and $\mat{A}_c=\mat{P}^{\mathsf T}\mat{A}\mat{P}$, the coarse-grid
correction is the best possible correction in the range of $\mat P$, measured in the
energy norm. The coarse grid does not merely \emph{approximate} the smooth-error
correction; it computes the \emph{optimal} one available to it.

\begin{physicsbox}
For finite elements the transfer operators are not invented---they are \emph{given by the
function-space hierarchy of \cref{ch:functionspaces}}. When the coarse FE space $V_H$ is a
genuine subspace of the fine space $V_h$ (a coarser mesh, or a lower polynomial order,
nested inside the finer one), every coarse basis function is exactly a linear combination
of fine basis functions. The coefficients of that combination \emph{are} the columns of
$\mat{P}$. Prolongation is then the natural inclusion $V_H \hookrightarrow V_h$, restriction
its transpose, and \cref{eq:galerkin} is automatic. Palace builds this nested hierarchy
once---an $h$-refined or $p$-coarsened tower of $\Hcurl$ (or $\Hone$) spaces---and reads
the prolongations straight off it. The next section makes that hierarchy the spine of the
recursion.
\end{physicsbox}

\subsection{The two-grid cycle assembled}

We now have every piece. The two-grid cycle, for one pass, is:
\begin{enumerate}[nosep]
\item \textbf{Presmooth:} apply a few smoother sweeps to $\mat{A}\vx=\vb$, starting from
  the current $\vx$, killing oscillatory error and leaving smooth error.
\item \textbf{Residual:} compute $\vr = \vb - \mat{A}\vx$.
\item \textbf{Restrict:} $\vr_c = \mat{P}^{\mathsf T}\vr$, carrying the residual to the
  coarse grid.
\item \textbf{Coarse solve:} solve $\mat{A}_c\,\ve_c = \vr_c$ on the coarse grid for the
  coarse correction.
\item \textbf{Prolong and correct:} $\vx \mathrel{+}= \mat{P}\,\ve_c$, adding the
  interpolated correction to the fine solution.
\item \textbf{Postsmooth:} a few more smoother sweeps, cleaning up the high-frequency
  error that interpolation reintroduced.
\end{enumerate}
Steps 1 and 6 handle the oscillatory error; steps 2--5 handle the smooth error via the
coarse grid. This is a complete, correct method---and \cref{wex:twogrid} below runs it on
real numbers and watches the error collapse. But it has one unsatisfying feature: step 4
says ``solve the coarse problem,'' and for a large fine grid the coarse problem is still
large. The resolution of that complaint is the idea that makes multigrid \emph{multi}.

\section{The V-cycle: recursion down the grids}
\label{sec:vcycle}

Step 4 of the two-grid cycle asks us to solve $\mat{A}_c\,\ve_c=\vr_c$ on the coarse grid.
But $\mat{A}_c$ is SPD (\cref{lem:galerkin}), so it is a problem of exactly the same kind
as the original---and we know a good way to attack a problem of that kind: the two-grid
cycle itself. Rather than solve the coarse problem exactly, \emph{apply one two-grid cycle
to it}, using a still-coarser grid below it. That two-grid cycle in turn defers its own
coarse solve to a yet-coarser grid, and so on, until we reach a grid so small---a handful
of unknowns---that solving directly is trivial. The two-grid method, applied to itself all
the way down, is the \textbf{V-cycle}.

\begin{keyidea}
\textbf{Smooth on the fine grid; solve the rest on a coarser grid---recursively.}
On every level, presmooth to kill that level's oscillatory error, push the smooth residual
down to the next-coarser level, let that level handle it the same way, bring the correction
back up, and postsmooth. The recursion bottoms out on a grid small enough to solve
directly. Because each level only ever does cheap smoothing (the expensive ``solve'' is
always deferred downward), the total cost is a geometric series that sums to $O(N)$, and
because each level kills exactly the frequencies it is responsible for, the convergence
rate is \emph{independent of the mesh size}. That conjunction---$O(N)$ work and
$h$-independent convergence---is what the word \emph{optimal} means in \cref{ch:precond}.
\end{keyidea}

\subsection{The V-cycle as a recursive function}

Number the levels $0, 1, \dots, L$, with $L$ the finest (the grid we actually want to
solve on) and $0$ the coarsest (the trivial base case). Each level $\ell$ carries its own
operator $\mat{A}_\ell$, its own smoother $\mat{B}_\ell$, and---for $\ell>0$---a
prolongation $\mat{P}_{\ell-1}$ mapping from level $\ell-1$ up to level $\ell$. The
V-cycle is then a single recursive function of the level $\ell$ and the current iterate
$\vx$. \Cref{lst:vcycle} writes it as a pure function---no in-place mutation, every step a
tensor-in/tensor-out map---in the calculus of \cref{part:framework}; this is the form Palace's
\code{GeometricMultigridSolver} takes once its scratch-vector mutations are rotated away.

\begin{lstlisting}[style=l4style, caption={The V-cycle as a level-recursive pure function. Each call presmooths, restricts the residual, recurses on the coarser level, prolongs the returned correction, and postsmooths. The base case $\ell=0$ solves directly.}, label={lst:vcycle}]
-- one V-cycle pass at level l: improve x toward solving A[l] x = b
vcycle :: [LinOp] -> [Smoother] -> [LinOp] -> Solver
       -> Int -> Tensor[N] -> Tensor[N] -> Tensor[N]
vcycle as bs ps coarseSolve 0 b x =
    coarseSolve b                          -- base case: solve A[0] e = b exactly
vcycle as bs ps coarseSolve l b x =
  let y   = presmooth (bs!l) b x           -- 1. kill oscillatory error on level l
      r   = b `sub` apply (as!l) y         -- 2. residual r = b - A[l] y
      rc  = apply_transpose (ps!(l-1)) r   -- 3. restrict  r_c = P^T r   (fine -> coarse)
      ec  = vcycle as bs ps coarseSolve (l-1) rc (zeros)   -- 4. RECURSE: solve A[l-1] e_c = r_c
      y'  = y `add` apply (ps!(l-1)) ec    -- 5. prolong & correct  y += P e_c
  in  postsmooth (bs!l) b y'               -- 6. clean up interpolation error
\end{lstlisting}

Read the function once more with the two-grid cycle of \cref{sec:twogrid} beside it: the
six numbered comments are the six numbered steps, with step~4---formerly ``solve the coarse
problem''---now a \emph{recursive call to \code{vcycle} itself} one level down. The base
case \code{vcycle ... 0 b x} does the only genuine solve in the whole algorithm, and it is
cheap because level $0$ is tiny. Everything above the base case does nothing but smooth,
restrict, and prolong. The coarse operators $\mat{A}_\ell$ and prolongations
$\mat{P}_{\ell-1}$ are precomputed once from the function-space hierarchy
(\cref{ch:functionspaces}); the recursion just walks them.

\begin{notationbox}
Each leg of the V-cycle---presmooth, the coarse-grid correction, postsmooth---has the same
algebraic shape: take the current guess $\vy$, compute its residual $\vr=\vb-\mat{A}\vy$,
apply some approximate inverse $\mat{B}$ to the residual, and add the result back:
\[
  \vy \;\longmapsto\; \vy + \mat{B}\,(\vb - \mat{A}\vy).
\]
This is the \emph{correction step} of \cref{ch:precond}: a presmooth leg takes $\mat{B}$ to
be the point smoother; the coarse-grid-correction leg takes $\mat{B} =
\mat{P}\,\mat{A}_c^{-1}\mat{P}^{\mathsf T}$ (smooth the residual by solving on the coarse
grid). The V-cycle is a \emph{composition of correction steps}, alternating point smoothers
with one coarse-grid smoother---which is itself a composition of correction steps one level
down. Reading multigrid this way---one combinator, applied at every level and recursively
inside itself---is the calculus view of \cref{part:framework}.
\end{notationbox}

\subsection{Why it is called a V}

Trace the flow of control through one call to \code{vcycle} at the finest level $L$.
Presmoothing happens on the way \emph{down}: level $L$ smooths, then hands a residual to
level $L-1$, which smooths and hands a residual to $L-2$, and so on down to level $0$.
That is the descending stroke of the letter. At level $0$ the direct solve happens---the
point of the V. Then the corrections climb back \emph{up}: level $0$ returns a correction to
$1$, which prolongs it, postsmooths, and returns to $2$, up to the finest level. That is the
ascending stroke. Drawn with the levels stacked vertically and time running left to right,
the path of the computation is a literal V (\cref{fig:vcycle}). This is the hero diagram of
the chapter; every multigrid method is a variation on its shape.

\begin{figure}[p]
\centering
\begin{tikzpicture}[scale=1.0, every node/.style={font=\footnotesize}]
  % vertical: level (fine at top). horizontal: time.
  % level y-coordinates
  \def\Lf{4.5} \def\La{3.0} \def\Lb{1.5} \def\Lc{0.0}
  % level guide lines + labels
  \foreach \y/\lab/\dof in {%
    4.5/{level $L$ (finest)}/{$N$ dof},%
    3.0/{level $L-1$}/{$N/2^d$},%
    1.5/{level $1$}/{$\vdots$},%
    0.0/{level $0$ (coarsest)}/{tiny}%
  }{
    \draw[simGray!35, densely dotted] (0.2,\y)--(11.3,\y);
    \node[anchor=east, simGray, font=\scriptsize] at (0.1,\y) {\lab};
    \node[anchor=west, simGray!80, font=\scriptsize\itshape] at (11.4,\y) {\dof};
  }
  % the V path: descending (presmooth+restrict), then ascending (prolong+postsmooth)
  % node coordinates along the V
  \coordinate (d0) at (1.6,\Lf);   % presmooth finest
  \coordinate (d1) at (3.4,\La);   % presmooth L-1
  \coordinate (d2) at (5.2,\Lb);   % presmooth 1
  \coordinate (bot) at (6.4,\Lc);  % coarse solve
  \coordinate (u2) at (7.6,\Lb);   % postsmooth 1
  \coordinate (u1) at (9.4,\La);   % postsmooth L-1
  \coordinate (u0) at (11.2,\Lf);  % postsmooth finest
  % descending strokes (restriction)
  \draw[-{Stealth[length=2.4mm]}, simBlue, very thick] (d0)--(d1);
  \draw[-{Stealth[length=2.4mm]}, simBlue, very thick] (d1)--(d2);
  \draw[-{Stealth[length=2.4mm]}, simBlue, very thick] (d2)--(bot);
  % ascending strokes (prolongation)
  \draw[-{Stealth[length=2.4mm]}, simRust, very thick] (bot)--(u2);
  \draw[-{Stealth[length=2.4mm]}, simRust, very thick] (u2)--(u1);
  \draw[-{Stealth[length=2.4mm]}, simRust, very thick] (u1)--(u0);
  % node markers
  \foreach \pt/\col in {d0/simGreen, d1/simGreen, d2/simGreen,
                        u2/simGold, u1/simGold, u0/simGold}{
    \fill[\col] (\pt) circle (3pt);}
  \fill[simViolet] (bot) circle (4pt);
  % leg labels (descending = presmooth + restrict)
  \node[simGreen, anchor=south, font=\scriptsize] at (d0) {presmooth};
  \node[simBlue, anchor=south east, font=\scriptsize] at ($(d0)!0.5!(d1)$) {restrict $\mat{P}^{\mathsf T}\vr$};
  \node[simBlue, anchor=south east, font=\scriptsize] at ($(d1)!0.5!(d2)$) {restrict};
  \node[simBlue, anchor=east, font=\scriptsize] at ($(d2)!0.5!(bot)+(0.1,0)$) {restrict};
  % bottom
  \node[simViolet, anchor=north, font=\scriptsize\bfseries] at (bot) {coarse solve};
  % ascending labels (prolong + postsmooth)
  \node[simRust, anchor=west, font=\scriptsize] at ($(bot)!0.5!(u2)-(0.1,0)$) {prolong $\mat{P}\ve_c$};
  \node[simRust, anchor=south west, font=\scriptsize] at ($(u2)!0.5!(u1)$) {prolong};
  \node[simRust, anchor=south west, font=\scriptsize] at ($(u1)!0.5!(u0)$) {prolong};
  \node[simGold, anchor=south, font=\scriptsize] at (u0) {postsmooth};
  % the big V annotation
  \node[simInk, font=\large\bfseries, opacity=0.18] at (6.4,2.5) {V};
\end{tikzpicture}
\caption[The V-cycle]{The hero diagram. Levels stacked vertically (finest at top,
coarsest at bottom), time running left to right. The computation descends the left stroke
of a V---\emph{presmoothing} (green) on each level, then \emph{restricting} the residual
(blue arrows, $\mat{P}^{\mathsf T}\vr$) to the next-coarser level---until it reaches the
bottom, where the coarsest problem is solved directly (violet). It then climbs the right
stroke---\emph{prolonging} each correction (rust arrows, $\mat{P}\ve_c$) up to the next-finer
level and \emph{postsmoothing} (gold)---back to the finest level. Each level does only cheap
smoothing and grid transfers; the single genuine solve happens on the tiny coarsest grid.
The grid shrinks by a factor $2^d$ per level (right margin), so the work per level shrinks
geometrically and the total is $O(N)$.}
\label{fig:vcycle}
\end{figure}

\subsection{Building the level hierarchy}

The V-cycle presupposes a tower of grids and the operators connecting them. Where does the
tower come from? \Cref{ch:functionspaces} already built it. Palace constructs a
\emph{finite-element space hierarchy}: a coarse-to-fine stack of nested FE spaces,
\[
  V_0 \subset V_1 \subset \cdots \subset V_L,
\]
obtained either by $h$-refinement (each $V_{\ell}$ a finer mesh than $V_{\ell-1}$) or by
$p$-coarsening (each $V_{\ell}$ a higher polynomial order). Because the spaces are nested,
each inclusion $V_{\ell-1}\hookrightarrow V_\ell$ \emph{is} a prolongation
$\mat{P}_{\ell-1}$, read directly off the hierarchy. The hierarchy is built once, before the
solve, by folding a level-adding operation over the refinement sequence---a coarse seed
space, then one level per refinement step---and thereafter the V-cycle consumes it
read-only. \Cref{fig:hierarchy} shows the geometric picture: the same domain meshed at three
resolutions, each coarse mesh nested in the finer one. The smoothers $\mat{B}_\ell$ and the
Galerkin operators $\mat{A}_\ell = \mat{P}_{\ell-1}^{\mathsf T}\mat{A}_{\ell+1}\mat{P}_{\ell-1}$
are likewise precomputed per level. With the tower in hand, the recursion of
\cref{lst:vcycle} is all that remains.

\begin{figure}[t]
\centering
\begin{tikzpicture}[scale=1.0, every node/.style={font=\scriptsize}]
  % three nested meshes of the same square domain
  \foreach \px/\n/\ttl/\col in {0/4/{fine ($h$)}/simGreen,
                                 5/2/{medium ($2h$)}/simBlue,
                                 10/1/{coarse ($4h$)}/simRust}{
    \begin{scope}[shift={(\px,0)}]
      \draw[simInk, thick] (0,0) rectangle (3,3);
      % grid lines
      \pgfmathsetmacro{\step}{3/\n}
      \foreach \k in {1,...,\n}{
        \pgfmathsetmacro{\pos}{\k*\step}
        \ifnum\k<\n
          \draw[\col!70] (\pos,0)--(\pos,3);
          \draw[\col!70] (0,\pos)--(3,\pos);
        \fi
      }
      % nodes
      \foreach \i in {0,...,\n}{\foreach \j in {0,...,\n}{
        \fill[\col] (\i*\step,\j*\step) circle (1.6pt);}}
      \node[font=\footnotesize\bfseries, \col] at (1.5,-0.45) {\ttl};
    \end{scope}}
  % refinement arrows between
  \draw[-{Stealth[length=2.4mm]}, simViolet, thick] (8.7,1.5)--(8.2,1.5)
    node[midway, above, simViolet]{$\mat{P}$};
  \draw[-{Stealth[length=2.4mm]}, simViolet, thick] (3.7,1.5)--(3.2,1.5)
    node[midway, above, simViolet]{$\mat{P}$};
  \node[anchor=north, simGray, font=\scriptsize\itshape] at (6.5,-1.0)
    {prolongation interpolates from each grid to the next finer; restriction $\mat{P}^{\mathsf T}$ runs the other way};
\end{tikzpicture}
\caption[The grid hierarchy]{The multigrid hierarchy of nested grids on one square domain.
\emph{Right to left:} a coarse mesh (rust, spacing $4h$), a medium mesh (blue, $2h$), a fine
mesh (green, $h$). Each coarse grid's nodes are a subset of the finer grid's, so the coarse
finite-element space is nested in the finer one and the inclusion is the prolongation
$\mat{P}$ (violet arrows). For finite elements the nesting comes free from the
function-space hierarchy of \cref{ch:functionspaces}; the V-cycle of \cref{fig:vcycle}
walks up and down this tower.}
\label{fig:hierarchy}
\end{figure}

\section{Why it is optimal: $O(N)$ work, mesh-independent convergence}
\label{sec:optimal}

The V-cycle's reputation rests on two quantitative claims. Each is worth establishing
carefully, because together they are what \cref{ch:precond} means by an \emph{optimal}
preconditioner and what no other general-purpose solver delivers for elliptic problems.

\subsection{Linear work per cycle}

\emph{Claim: one V-cycle costs $O(N)$ operations, where $N$ is the number of fine-grid
unknowns.} The work on each level is proportional to the number of unknowns on that level:
a fixed number of smoother sweeps, one residual, one restriction, one prolongation, all of
which touch each unknown a constant number of times. Let $N_\ell$ be the unknown count on
level $\ell$, with $N_L = N$ the finest. In $d$ space dimensions, coarsening by a factor of
two per direction divides the unknown count by $2^d$ each level:
\[
  N_{L-1} \approx \frac{N}{2^d}, \quad N_{L-2}\approx \frac{N}{4^d}, \quad \dots
\]
The total work is the sum over levels, $W = c\sum_{\ell} N_\ell$, a geometric series
\[
  W \;\approx\; c\,N\Bigl(1 + 2^{-d} + 2^{-2d} + \cdots\Bigr)
    \;=\; c\,N\,\frac{1}{1 - 2^{-d}}.
\]
The ratio $1/(1-2^{-d})$ is a constant: $2$ in 1-D, $\tfrac43$ in 2-D, $\tfrac87$ in 3-D.
So the whole tower costs only a small constant factor more than the finest level alone---the
coarse levels are essentially free. \Cref{fig:costseries} draws the geometric series of
per-level costs; \cref{wex:cost} makes the sum concrete. \emph{One V-cycle is $O(N)$.}

\begin{figure}[t]
\centering
\begin{tikzpicture}
\begin{axis}[
  width=11.5cm, height=5.2cm,
  ybar, bar width=15pt,
  xlabel={multigrid level (finest $=L$, then coarser)},
  ylabel={work $\propto N_\ell$},
  symbolic x coords={L, L-1, L-2, L-3, L-4},
  xtick=data, ymin=0, ymax=1.15,
  ytick={0,0.25,0.5,0.75,1.0},
  yticklabel={\pgfmathprintnumber{\tick}\,$N$},
  axis lines=left, tick align=outside,
  nodes near coords, nodes near coords style={font=\scriptsize, simGray},
  every node near coord/.append style={/pgf/number format/fixed,
    /pgf/number format/precision=3},
  bar shift=0pt,
]
  % work per level in 3D: N, N/8, N/64, ...
  \addplot[fill=simBlue!70, draw=simBlue] coordinates {
    (L,1.0) (L-1,0.125) (L-2,0.0156) (L-3,0.00195) (L-4,0.00024)};
\end{axis}
\end{tikzpicture}
\caption[The per-level cost geometric series]{Work per multigrid level in 3-D ($d=3$),
in units of the finest-level work $N$. Each coarser level has $1/2^d = 1/8$ the unknowns of
the one above, so its cost is one-eighth as large: $N,\ N/8,\ N/64,\dots$. The bars shrink so
fast that the entire tower below the finest level adds only $1/(2^3-1) = 1/7$ to the finest
level's cost---the coarse grids are nearly free. Summed, the geometric series gives total
work $\tfrac{8}{7}N = O(N)$ per V-cycle, independent of how many levels the tower has.}
\label{fig:costseries}
\end{figure}

\subsection{Mesh-independent convergence}

\emph{Claim: the V-cycle reduces the error by a fixed factor $\rho<1$ per cycle, and
$\rho$ does not grow as the mesh is refined.} This is the deeper claim, and the
complementary-frequencies principle of \cref{sec:complementary} is its proof sketch. Split
the error into smooth and oscillatory parts relative to the fine grid. The smoother
contracts the oscillatory part by a factor bounded below $1$ \emph{independent of $h$}---a
property called \emph{smoothing}. The coarse-grid correction contracts the smooth part,
and by variational consistency (\cref{lem:galerkin}) it does so optimally, again with a
factor bounded below $1$ \emph{independent of $h$}---a property called \emph{approximation}.
Every error component is hit by one mechanism or the other with an $h$-independent factor,
so the overall contraction $\rho$ is bounded below $1$ uniformly in $h$. A typical
well-tuned V-cycle has $\rho \approx 0.1$ to $0.2$: each cycle kills an order of magnitude
of error, regardless of mesh size.

Contrast this with the unpreconditioned conjugate gradient method of \cref{ch:cg}. Its
iteration count to a fixed tolerance scales as $\sqrt{\kappa}$, where $\kappa$ is the
condition number of $\mat{A}$. For a second-order elliptic operator, $\kappa = O(h^{-2})$,
so $\sqrt{\kappa} = O(h^{-1})$: \emph{halving the mesh spacing doubles the iteration count.}
Multigrid's iteration count, by contrast, is \emph{flat}---a constant number of V-cycles
reaches tolerance no matter how fine the mesh. \Cref{fig:meshindep} draws the two curves;
\cref{wex:cost} puts numbers on them.

\begin{figure}[t]
\centering
\begin{tikzpicture}
\begin{axis}[
  width=11.5cm, height=5.6cm,
  xlabel={problem size $N$ (degrees of freedom)},
  ylabel={iterations to tolerance},
  xmode=log, log basis x=2,
  xmin=700, xmax=1.2e6, ymin=0, ymax=320,
  xtick={1024,4096,16384,65536,262144,1048576},
  xticklabels={$2^{10}$,$2^{12}$,$2^{14}$,$2^{16}$,$2^{18}$,$2^{20}$},
  ytick={0,50,100,150,200,250,300},
  axis lines=left, tick align=outside, grid=major, grid style={simGray!18},
  legend style={at={(0.03,0.97)}, anchor=north west, font=\footnotesize,
    draw=simGray!50},
  every axis plot/.append style={very thick},
]
  % unpreconditioned CG: iterations ~ sqrt(kappa) ~ 1/h ~ N^{1/d} (d=2 => N^{1/2})
  \addplot[simRust, mark=*, mark size=1.6pt] coordinates {
    (1024,9.6)(4096,19.2)(16384,38.4)(65536,76.8)(262144,153.6)(1048576,307.2)};
  \addlegendentry{unpreconditioned CG: $\sim\!\sqrt{\kappa}\sim 1/h$}
  % multigrid V-cycle: flat
  \addplot[simBlue, mark=square*, mark size=1.6pt] coordinates {
    (1024,8)(4096,8)(16384,8)(65536,9)(262144,9)(1048576,9)};
  \addlegendentry{multigrid V-cycle: flat ($h$-independent)}
\end{axis}
\end{tikzpicture}
\caption[Mesh-independent convergence]{The property that makes multigrid optimal.
\emph{Rust:} unpreconditioned conjugate gradient (\cref{ch:cg}) needs $\sim\sqrt{\kappa}\sim
1/h$ iterations, which \emph{doubles} every time the mesh is halved (note the log scale on
$N$)---refining the mesh makes the solver slower and slower. \emph{Blue:} the multigrid
V-cycle takes a nearly constant number of cycles regardless of mesh size: its convergence
factor $\rho<1$ is bounded \emph{independent} of $h$. This flatness is the definition of an
\emph{optimal} solver in \cref{ch:precond}, and multigrid is the canonical example.}
\label{fig:meshindep}
\end{figure}

\begin{keyidea}
The two claims combine into the headline. \emph{One V-cycle costs $O(N)$, and a fixed
number of V-cycles---independent of $N$---reaches any tolerance.} Therefore solving an
elliptic PDE to fixed accuracy costs $O(N)$ total: optimal, in the strict sense that you
cannot beat linear in the number of unknowns, since you must at least look at each unknown
once. No purely algebraic solver---not Gaussian elimination ($O(N^3)$, or $O(N^{1.5})$ for
2-D sparse), not unpreconditioned CG ($O(N^{1.5})$ in 2-D)---matches this. Multigrid is the
gold standard against which elliptic solvers are measured.
\end{keyidea}

\section{Worked examples}

\subsection{A two-grid cycle on 1-D Poisson}

We run the two-grid method of \cref{sec:twogrid} on a small concrete problem and watch the
error collapse. The arithmetic is deliberately tractable so every step can be checked by
hand.

\begin{workedexample}{A two-grid cycle on the 1-D Poisson equation}{twogrid}
Discretize $-u'' = f$ on $[0,1]$ with homogeneous Dirichlet conditions, on a fine grid of
$n=7$ interior points at spacing $h=\tfrac18$. The operator is the standard tridiagonal
$\mat{A} = h^{-2}\,\mathrm{tridiag}(-1, 2, -1)$, size $7\times 7$. Take a right-hand side
whose exact solution is the smooth function $u(x)=\sin(\pi x)$ (so the discrete error we
introduce will be smooth-plus-oscillatory and we can watch each part go).

\textbf{Initial error.} Start from an initial guess corrupted by both a smooth and an
oscillatory component: let the initial error be
$\ve_0 = \underbrace{\sin(\pi x)}_{\text{smooth, } j=1} +
\tfrac12\underbrace{\sin(7\pi x)}_{\text{oscillatory, } j=7}$,
sampled at the seven fine nodes. Its energy norm (we track $\norm{\ve}_\infty$ for
legibility) starts near $\norm{\ve_0}_\infty \approx 1.42$.

\textbf{Step 1 --- presmooth.} Apply two sweeps of weighted Jacobi, $\vx \leftarrow \vx +
\omega\mat{D}^{-1}(\vb-\mat{A}\vx)$ with $\omega=\tfrac23$ and $\mat{D}=2h^{-2}\mathrm{I}$.
The Jacobi smoothing factor on this operator is $|1-\tfrac23(1-\cos\theta_j)|$, which is
$\approx 0.95$ for the smooth mode $j=1$ but $\approx 0.10$ for the oscillatory mode $j=7$.
After two sweeps the oscillatory amplitude $\tfrac12$ has fallen to
$\tfrac12(0.10)^2 \approx 0.005$---gone---while the smooth amplitude $1$ has barely
moved, to $\approx 0.90$. The error is now $\norm{\ve}_\infty \approx 0.90$, and---this is
the point---it is now essentially \emph{pure smooth mode}.

\textbf{Step 2 --- residual.} Compute $\vr = \vb - \mat{A}\vx$ on the fine grid. Because the
surviving error is the smooth mode $\ve\approx 0.90\sin(\pi x)$ and $\mat{A}\sin(\pi x) =
\lambda_1\sin(\pi x)$ with $\lambda_1 = h^{-2}\cdot 2(1-\cos\pi h)\approx 9.74$, the residual
is itself smooth, $\vr = \mat{A}\ve \approx 8.8\sin(\pi x)$.

\textbf{Step 3 --- restrict.} The coarse grid has $n_c=3$ interior points at spacing
$2h=\tfrac14$. Full-weighting restriction $\mat{P}^{\mathsf T}$ (stencil
$\tfrac14[1,2,1]$) carries $\vr$ to a coarse residual $\vr_c$ of length $3$, still a clean
smooth coarse mode.

\textbf{Step 4 --- coarse solve.} Solve the $3\times 3$ system
$\mat{A}_c\,\ve_c = \vr_c$ \emph{exactly} (three unknowns---a trivial direct solve), where
$\mat{A}_c = \mat{P}^{\mathsf T}\mat{A}\mat{P}$ is the Galerkin operator,
itself the $3\times3$ tridiagonal Laplacian at spacing $2h$. The coarse correction $\ve_c$
recovers almost exactly the coarse samples of the smooth error $0.90\sin(\pi x)$.

\textbf{Step 5 --- prolong and correct.} Interpolate: $\vx \leftarrow \vx +
\mat{P}\,\ve_c$. The prolonged correction $\mat{P}\ve_c$ matches the smooth error $\ve$ to
within the interpolation error of a smooth function on a grid that resolves it well---a few
percent. After the correction the fine-grid error has dropped from $\approx 0.90$ to
$\approx 0.06$.

\textbf{Step 6 --- postsmooth.} Two more Jacobi sweeps clean up the small oscillatory error
that linear interpolation injected, bringing $\norm{\ve}_\infty \approx 0.02$.

\textbf{The tally.} One two-grid cycle drove the error from $1.42$ to $0.02$---a factor of
about $70$, two orders of magnitude---at the cost of a few fine-grid Jacobi sweeps plus one
trivial $3\times3$ solve. Smoothing alone, run for the same number of sweeps, would have
removed the oscillatory mode and then stalled at $\approx 0.90$, having made essentially no
progress on the smooth mode. The coarse grid did in one $3\times 3$ solve what no amount of
fine-grid smoothing could do. \emph{That} is the complementary-frequencies principle paying
off in numbers.
\end{workedexample}

\subsection{The work estimate and the iteration-count contrast}

The second example makes the two optimality claims of \cref{sec:optimal} quantitative: the
$O(N)$ per-cycle cost as a summed geometric series, and the mesh-independent iteration count
set against CG's growth.

\begin{workedexample}{$O(N)$ work and the CG contrast}{cost}
\textbf{Part (a): one V-cycle is $O(N)$.} Take a 3-D problem, $d=3$, with $N$ unknowns on the
finest grid. Coarsening halves the grid spacing's inverse---divides the unknown count by
$2^d = 8$---per level. Let the work on a level be $c$ operations per unknown (a fixed number
of smoother sweeps plus the transfers). The total work is
\[
  W = c\Bigl(N + \tfrac{N}{8} + \tfrac{N}{64} + \tfrac{N}{512} + \cdots\Bigr)
    = cN\sum_{k=0}^{\infty} 8^{-k} = cN\cdot\frac{1}{1-\tfrac18} = cN\cdot\frac{8}{7}.
\]
The sum is a finite constant, $\tfrac87 \approx 1.143$: \emph{the entire tower of coarse
grids adds only $14\%$ to the cost of the finest level alone.} So $W = O(N)$, with a small
constant. In 2-D the factor is $\tfrac{1}{1-1/4} = \tfrac43$; in 1-D, $\tfrac{1}{1-1/2}=2$.
In every dimension the geometric series converges and the per-cycle cost is strictly linear
in $N$.

\textbf{Part (b): total cost to solve.} If a fixed number $m$ of V-cycles---say $m\approx
8$---reaches the tolerance \emph{regardless of $N$} (the mesh-independence claim), then the
total solve cost is $m\cdot W = 8\cdot\tfrac87 cN = O(N)$. Solving the PDE costs linear time.

\textbf{Part (c): the CG contrast.} Now estimate unpreconditioned CG on the same family of
2-D problems. Its iteration count to fixed tolerance is $\approx \tfrac12\sqrt{\kappa}$ with
$\kappa \approx (h^{-2})\cdot(\text{const})$, so iterations $\propto 1/h \propto
\sqrt{N}$ in 2-D. Tabulate both as the mesh is refined, doubling resolution each row (so $N$
quadruples in 2-D):
\begin{center}\small
\begin{tabular}{@{}rrr@{}}
\toprule
$N$ & CG iterations $\sim\!\sqrt{N}$ & V-cycles \\
\midrule
$1\,024$      & $\approx 16$  & $8$ \\
$4\,096$      & $\approx 32$  & $8$ \\
$16\,384$     & $\approx 64$  & $8$ \\
$65\,536$     & $\approx 128$ & $9$ \\
$262\,144$    & $\approx 256$ & $9$ \\
\bottomrule
\end{tabular}
\end{center}
CG's iteration count \emph{doubles every time the mesh is halved}; multigrid's is flat. And
each CG iteration also costs $O(N)$ (one sparse mat-vec), so CG's total solve cost in 2-D is
$O(N\sqrt{N}) = O(N^{1.5})$---superlinear---against multigrid's $O(N)$. At $N=262\,144$
multigrid does the work of $9$ cycles while CG grinds through $256$ iterations; the gap
widens without bound as the mesh refines. This table \emph{is} \cref{fig:meshindep} in
numbers, and it is the entire case for preconditioning the Krylov solve with a V-cycle---the
subject of the next section.
\end{workedexample}

\section{Variants and deployment}
\label{sec:variants}

The V-cycle is the prototype, not the only shape. A few variations matter in practice, and
one deployment pattern---multigrid \emph{inside} a Krylov solver---is how Palace actually
uses it.

\subsection{V-cycle, W-cycle, and full multigrid}

The recursion of \cref{lst:vcycle} makes exactly \emph{one} recursive call per level. If the
coarse-grid correction is not accurate enough---because the coarse problem itself was solved
too crudely---one can make \emph{two} recursive calls per level instead, visiting each coarse
level twice. The resulting flow traces a \emph{W} rather than a V (\cref{fig:cycleshapes},
middle): more coarse-grid work, a stronger correction, useful for harder problems where a
single coarse visit leaves too much error. \emph{Full multigrid} (FMG) is a different idea
again: rather than start on the finest grid with a poor guess, start on the \emph{coarsest}
grid, solve it, prolong the solution to the next level as an excellent initial guess, do a
V-cycle there, prolong again, and so on up to the finest grid (\cref{fig:cycleshapes},
right). FMG reaches discretization accuracy in a single sweep up the tower---often the most
efficient option of all, because every level starts from a guess the level below already
made good.

\begin{figure}[t]
\centering
\begin{tikzpicture}[scale=0.82, every node/.style={font=\scriptsize}]
  % shared: 4 levels, y = 0..3 (0 coarsest)
  % --- V-cycle ---
  \begin{scope}
    \foreach \y in {0,1,2,3}{\draw[simGray!30, densely dotted] (0,\y)--(3.2,\y);}
    \node[font=\footnotesize\bfseries, simInk] at (1.6,3.7) {V-cycle};
    \draw[-{Stealth[length=2mm]}, simBlue, thick]
      (0.2,3)--(1.0,2)--(1.8,1)--(2.2,0)--(2.6,1)--(3.0,2)--(3.4,3);
    \draw[simBlue, thick]
      (0.2,3)--(1.0,2)--(1.8,1)--(2.2,0)--(2.6,1)--(3.0,2)--(3.4,3);
    \node[anchor=east, simGray, font=\scriptsize] at (-0.1,3) {fine};
    \node[anchor=east, simGray, font=\scriptsize] at (-0.1,0) {coarse};
  \end{scope}
  % --- W-cycle ---
  \begin{scope}[shift={(4.6,0)}]
    \foreach \y in {0,1,2,3}{\draw[simGray!30, densely dotted] (0,\y)--(4.4,\y);}
    \node[font=\footnotesize\bfseries, simInk] at (2.2,3.7) {W-cycle};
    \draw[simRust, thick]
      (0.2,3)--(0.9,2)--(1.4,1)--(1.7,0)--(2.0,1)--(2.3,0)--(2.7,1)--(3.2,2)
      --(3.6,1)--(3.9,0)--(4.2,1)--(4.6,3);
  \end{scope}
  % --- FMG ---
  \begin{scope}[shift={(10.0,0)}]
    \foreach \y in {0,1,2,3}{\draw[simGray!30, densely dotted] (0,\y)--(4.4,\y);}
    \node[font=\footnotesize\bfseries, simInk] at (2.2,3.7) {full multigrid};
    % start at coarsest, climb with V-cycles
    \draw[simGreen, thick]
      (0.2,0)--(0.6,1)--(0.4,0)            % V at level 1
      (0.6,1)--(1.2,2)--(0.9,1)--(0.7,0)--(1.0,1)--(1.4,2)   % V at level 2
      (1.4,2)--(2.2,3)--(1.8,2)--(1.5,1)--(1.3,0)--(1.6,1)--(2.0,2)--(2.6,3); % V at finest
    \node[anchor=north, simGray, font=\scriptsize\itshape] at (2.2,-0.5)
      {start coarse, climb};
  \end{scope}
\end{tikzpicture}
\caption[V-cycle, W-cycle, and full multigrid]{Three multigrid schedules, levels stacked
vertically (finest at top). \emph{Left:} the V-cycle makes one recursive call per level---one
trip down, one trip up. \emph{Middle:} the W-cycle makes two recursive calls per level,
visiting each coarse level twice for a stronger (but costlier) coarse correction. \emph{Right:}
full multigrid (FMG) starts on the coarsest grid and climbs, prolonging each level's solution
as the initial guess for a V-cycle on the next finer level---reaching discretization accuracy
in one upward sweep. The V-cycle is the workhorse; the others trade more coarse work for a
better correction.}
\label{fig:cycleshapes}
\end{figure}

\subsection{Auxiliary-space multigrid for $\Hcurl$}

The smoother-and-coarsen recipe assumes the smoother can damp all the oscillatory error---and
for the scalar Laplacian on $\Hone$ it can. But the curl--curl operator of every magnetic and
wave analysis (\cref{ch:functionspaces}) lives on $\Hcurl$, and there a plain smoother has a
blind spot: the large null space of the curl (the \emph{gradient fields}, $\Curl\Grad\phi =
0$) carries error that point smoothing does not see. A pure geometric V-cycle stalls on those
gradient modes. The cure is \emph{auxiliary-space} multigrid: alongside the geometric
hierarchy, build an auxiliary $\Hone$ space for the gradient part and smooth there too. The
per-level smoother becomes a \emph{Hiptmair smoother} (\cref{ch:smoothers}): a point smoother
on $\Hcurl$ to handle the non-gradient error, \emph{plus} a smoother on the auxiliary
gradient space to handle the kernel error the first one misses. Palace's AMS preconditioner
builds this auxiliary hierarchy automatically from the discrete gradient operator; the V-cycle
structure is unchanged---only the per-level smoother $\mat{B}_\ell$ is now the composite
Hiptmair one. The same shape, a better smoother, because the operator has a structure the
naive smoother cannot see.

\subsection{$h$- versus $p$-multigrid}

The hierarchy of \cref{fig:hierarchy} coarsened the \emph{mesh} ($h$-multigrid). One can
instead coarsen the \emph{polynomial order} on a fixed mesh ($p$-multigrid): the finest level
uses order-$p$ elements, the next order-$p/2$, down to order $1$ (or the family's floor from
\cref{ch:functionspaces}). The two are complementary and often combined ($hp$-multigrid):
coarsen the order down to linears, then coarsen the mesh. Palace's default schedule is
$p$-multigrid with a logarithmic order drop---halving the order each level---so a finest order
$p=8$ space reaches the floor in a handful of levels rather than the seven a linear
$p\mapsto p-1$ drop would need. The choice of $h$ versus $p$ changes \emph{what} the
prolongations interpolate (finer mesh, or higher order); the V-cycle recursion is identical.

\subsection{Multigrid as a preconditioner}

Here is the crucial deployment fact. We have described the V-cycle as a stand-alone iterative
solver---repeat V-cycles until the residual is small. In practice Palace almost never uses it
that way. Instead it uses \emph{one} V-cycle as a \emph{preconditioner} inside a Krylov method
(\cref{ch:cg,ch:krylov}). Recall from \cref{ch:precond} that a preconditioner is an
approximate inverse $\mat{M}^{-1}\approx\mat{A}^{-1}$ that the Krylov method applies once per
iteration to reshape the spectrum. \emph{One V-cycle is exactly such an $\mat{M}^{-1}$ apply}:
feed it a residual $\vr$ (as the right-hand side $\vb$, with zero initial guess) and it returns
an approximate correction $\mat{M}^{-1}\vr$---a good approximation to $\mat{A}^{-1}\vr$,
computed in $O(N)$ work. \Cref{fig:mgprecond} shows the arrangement: the Krylov solver drives
the outer loop, and each of its iterations calls one V-cycle to precondition.

\begin{figure}[t]
\centering
\begin{tikzpicture}[every node/.style={font=\footnotesize}]
  % outer Krylov loop with an inner V-cycle as the M^{-1} apply
  \node[stage, minimum width=30mm, minimum height=22mm] (krylov) at (0,0)
    {\textbf{Krylov solver}\\[2pt]\footnotesize (CG / GMRES)\\\footnotesize outer iteration};
  \node[extern, minimum width=30mm, minimum height=14mm, right=34mm of krylov] (vcyc)
    {\textbf{one V-cycle}\\[1pt]\footnotesize $= \mat{M}^{-1}$ apply\\\footnotesize $O(N)$ work};
  % residual goes in, correction comes back
  \draw[flow] (krylov.east|-0,0.35) -- node[above,font=\scriptsize]{residual $\vr$}
    (vcyc.west|-0,0.35);
  \draw[backflow] (vcyc.west|-0,-0.35) -- node[below,font=\scriptsize]{correction $\mat{M}^{-1}\vr$}
    (krylov.east|-0,-0.35);
  % the V-shape hint inside the extern box
  \node[anchor=north, simGray, font=\scriptsize\itshape] at (vcyc.south) {(the recursion of Fig.~\ref{fig:vcycle})};
  \node[anchor=south, simGray, font=\scriptsize\itshape] at (krylov.north)
    {once per iteration};
\end{tikzpicture}
\caption[Multigrid as a preconditioner]{How Palace actually deploys multigrid. A Krylov solver
(\cref{ch:cg,ch:krylov}) runs the outer iteration; once per iteration it hands its current
residual $\vr$ to \emph{one} V-cycle, which returns an approximate correction $\mat{M}^{-1}\vr$
in $O(N)$ work---the V-cycle \emph{is} the preconditioner's apply. Because the V-cycle makes
$\mat{M}^{-1}\mat{A}$ have an $h$-independent, tightly-clustered spectrum, the outer Krylov
iteration count is small and flat (\cref{fig:meshindep}). The combination---multigrid-preconditioned
CG or GMRES---is the standard optimal solver for the elliptic systems of \cref{part:modalities}.}
\label{fig:mgprecond}
\end{figure}

Why precondition rather than iterate V-cycles directly? Because the Krylov method is more
robust to imperfections in the V-cycle. If the smoother or coarsening is slightly off and a
few error modes converge slowly, a stand-alone V-cycle iteration inherits that slowness, but
the Krylov method \emph{accelerates} those stragglers with its optimal polynomial (the
machinery of \cref{ch:cg}). One V-cycle per Krylov iteration gives the best of both: the
V-cycle supplies the $h$-independent spectral clustering, and the Krylov method mops up
whatever the V-cycle handled imperfectly. This is the arrangement behind essentially every
elliptic solve in the book.

\section{Pitfalls: what breaks optimality}

Multigrid's optimality is not automatic; it rests on the two mechanisms---smoothing and
coarse-grid approximation---\emph{each} doing its job. Break either and the $h$-independent
convergence is lost, often silently, the method degrading into something no better than a
plain smoother.

\begin{pitfall}
\textbf{A bad smoother destroys optimality.} The whole scheme rests on the smoother
annihilating the error the coarse grid \emph{cannot} represent---the oscillatory modes. If the
smoother does not damp \emph{all} the oscillatory modes (e.g.\ an undamped Jacobi iteration,
which \emph{amplifies} the highest mode; or a point smoother on $\Hcurl$, blind to the gradient
null space of \cref{sec:variants}), those modes are left for a coarse grid that cannot see them,
and they never go away. The convergence factor $\rho$ creeps toward $1$ and the mesh-independence
evaporates. \emph{Always verify the smoother damps the full oscillatory band}---this is why
$\Hcurl$ problems need the Hiptmair smoother, not plain Jacobi.
\end{pitfall}

\begin{pitfall}
\textbf{A bad coarse grid destroys optimality too.} The coarse grid must \emph{capture the
smooth modes}---the error the smoother leaves behind. If the prolongation $\mat{P}$ cannot
represent the smooth error well (too aggressive a coarsening, a coarse space that omits a
near-null-space mode like the constant or a rigid-body mode), the coarse-grid correction is
inaccurate, the smooth error is not removed, and again $\rho\to 1$. The Galerkin construction
$\mat{A}_c=\mat{P}^{\mathsf T}\mat{A}\mat{P}$ guarantees the \emph{best} correction in the
range of $\mat{P}$, but only the range of $\mat{P}$---if the smooth error is not in that range,
no coarse solve recovers it.
\end{pitfall}

\begin{pitfall}
\textbf{Anisotropy and jumps need special handling.} The complementary-frequencies argument of
\cref{sec:complementary} assumed the smoother smooths \emph{isotropically}---equally in every
direction. For a strongly anisotropic operator (e.g.\ a Laplacian stretched
$1000\times$ in one direction, common in thin-layer geometries) a point smoother smooths only
along the strongly-coupled direction and leaves the other rough; standard coarsening then
coarsens a direction the smoother never smoothed, and optimality fails. The fixes---line/plane
smoothers that smooth a whole coupled direction at once, or semi-coarsening that coarsens only
the smoothed direction---are essential for anisotropic or jumping-coefficient problems and are
where naive multigrid most often disappoints. The lesson: the smoother and the coarsening must
be \emph{matched} to the operator's structure, not chosen in isolation.
\end{pitfall}

\summarysection
A smoother (\cref{ch:smoothers}) cheaply removes the high-frequency part of the error and
stalls on the smooth part. Multigrid turns that stall into a method by exploiting the
\emph{complementary-frequencies principle}: smooth error on a fine grid looks oscillatory on a
coarser grid, where a smoother can remove it cheaply. The two-grid method makes this concrete:
presmooth on the fine grid, compute the residual, restrict it to the coarse grid with
$\mat{R}=\mat{P}^{\mathsf T}$, solve the Galerkin coarse problem $\mat{A}_c=\mat{P}^{\mathsf
T}\mat{A}\mat{P}$ for a correction, prolong the correction back with $\mat{P}$, and postsmooth.
The \emph{V-cycle} applies this idea recursively: each level's coarse solve is itself a V-cycle
on a still-coarser grid, until a tiny coarsest grid is solved directly---a descending stroke of
presmoothing-and-restricting and an ascending stroke of prolonging-and-postsmoothing, the V of
the hero diagram. Written as a pure recursive function, every leg is the same correction step
$\vy + \mat{B}(\vb-\mat{A}\vy)$. For finite elements the prolongations come free from the nested
function-space hierarchy of \cref{ch:functionspaces}. The V-cycle is \emph{optimal}: one cycle
costs $O(N)$ (the per-level costs form a geometric series summing to a small constant times $N$),
and a fixed, $h$-independent number of cycles reaches any tolerance---in sharp contrast to
unpreconditioned CG, whose iteration count grows like $1/h$. Variants trade coarse work for
correction strength (W-cycle, full multigrid) or adapt the smoother to the operator (Hiptmair
relaxation for the $\Hcurl$ gradient null space; $p$- versus $h$-coarsening). In practice one
V-cycle is used as the $\mat{M}^{-1}$ apply inside a Krylov solver---the optimal preconditioner
of \cref{ch:precond}. Optimality is fragile: a smoother that misses some oscillatory modes, a
coarse grid that misses some smooth modes, or an anisotropic operator that defeats isotropic
smoothing each silently destroys the mesh-independence.

\furtherreading
The auxiliary-space and Hiptmair smoother machinery for $\Hcurl$ multigrid is developed in
Hiptmair's foundational paper~\citep{hiptmair1998}, the source of the gradient-space smoother
that makes curl--curl multigrid work; it is the natural sequel to \cref{sec:variants}. For the
finite-element and variational framing of multigrid---nested spaces, Galerkin coarsening, the
approximation-and-smoothing convergence theory sketched in \cref{sec:optimal}---Brenner and
Scott~\citep{brennerscott2008} give the rigorous account at the level of this book. A modern,
implementation-oriented survey of multigrid and $p$-multigrid for high-order finite elements,
close to Palace's actual algorithms, is Phillips and Fischer~\citep{phillipsfischer2022}. The
classical introduction to the complementary-frequencies argument, with the Fourier (local-mode)
analysis of smoothing factors only sketched here, remains the multigrid tutorial literature
those references cite in turn.

\exercisessection
\begin{enumerate}[nosep]
\item State the complementary-frequencies principle in one sentence, and explain why it would
  fail if the smoother damped the \emph{low} frequencies rather than the high ones. What kind of
  method would you need to pair with such a smoother instead of a coarse grid?
\item For the 1-D linear-interpolation prolongation of \cref{sec:twogrid}, write out the
  $7\times 3$ matrix $\mat{P}$ mapping a coarse grid of $3$ interior points to a fine grid of $7$
  interior points (spacing $2h$ to $h$). Then write $\mat{R}=\mat{P}^{\mathsf T}$ and confirm it
  is the full-weighting $\tfrac14[1,2,1]$ restriction stencil up to scaling.
\item Using the $\mat{P}$ from the previous exercise and the fine tridiagonal Laplacian
  $\mat{A}=h^{-2}\mathrm{tridiag}(-1,2,-1)$, compute the Galerkin coarse operator
  $\mat{A}_c=\mat{P}^{\mathsf T}\mat{A}\mat{P}$ and verify it is (up to scaling) the $3\times 3$
  tridiagonal Laplacian at spacing $2h$. What does this say about Galerkin coarsening for this
  problem?
\item Prove \cref{lem:galerkin} for yourself without looking: show that if $\mat{A}$ is SPD and
  $\mat{P}$ has full column rank, then $\mat{A}_c=\mat{P}^{\mathsf T}\mat{A}\mat{P}$ is SPD. Which
  hypothesis fails, and what goes wrong, if $\mat{P}$ has a nontrivial null space?
\item A 3-D problem has $N = 10^6$ unknowns on the finest grid and $5$ levels. Using the
  geometric-series estimate of \cref{wex:cost}, compute the total unknown count summed over all
  levels and express it as a multiple of $N$. How much does the entire coarse hierarchy add to the
  finest-level cost?
\item Unpreconditioned CG on a 3-D Poisson problem needs $\approx\tfrac12\sqrt{\kappa}$ iterations
  with $\kappa\approx (1/h)^2$. If multigrid-preconditioned CG needs a flat $\approx 10$ iterations
  regardless of $h$, tabulate the iteration count of each method for meshes
  $h=\tfrac1{16},\tfrac1{32},\tfrac1{64},\tfrac1{128}$. By what factor does plain CG's count grow
  per refinement, and multigrid's?
\item Explain, in terms of the V-cycle recursion of \cref{lst:vcycle}, the difference between a
  V-cycle and a W-cycle. Modify the pseudocode of \cref{lst:vcycle} to perform a W-cycle (two
  recursive calls per level). Why does the W-cycle cost more, and when is the extra cost worth it?
\item Why does a point Jacobi smoother stall on the curl--curl operator's gradient null space, and
  how does the Hiptmair smoother of \cref{sec:variants} fix it? Relate your answer to the
  spurious-mode discussion of \cref{ch:functionspaces}: both are consequences of the structure of
  $\Hcurl$.
\end{enumerate}

\chapter{Eigenvalue Problems and Spectral Transforms}
\label{ch:eigenvalues}

\chapterorient{The eigenmode analysis asks a question none of the linear
solvers of this part can answer: not ``what field does this drive produce?''
but ``what fields can ring with no drive at all?'' Mathematically it asks for a
few eigenpairs of the generalized eigenproblem $\Kop\vx=\lambda\Mop\vx$---and,
maddeningly, it wants the \emph{smallest} few, the lowest resonances buried at
the bottom of an enormous spectrum. Plain iteration finds the \emph{largest}
eigenvalues, the exact opposite of what we need. This chapter develops the one
idea that turns the hard interior problem into an easy dominant one: the
\emph{spectral transform}. We watch power iteration climb toward the dominant
eigenvector, see why it climbs the wrong direction, and then bend the spectrum
with shift-and-invert so the eigenvalues we want become the ones iteration
naturally finds---paying for it with an inner linear solve. By the end you will
know why the eigenmode driver assembles its operators, hands them to a library,
and reads back converged modes, with no Palace-authored loop in sight.}

\section{The problem: a few low resonances of a huge pencil}
\label{sec:evp-problem}

The eigenmode analysis of \cref{ch:targets} seals a region of space with
conducting walls and asks what fields it can sustain on its own. We derived its
governing equation in \cref{ch:reductions-pde}: the source-free time-harmonic
curl--curl equation, discretized, is the \emph{generalized eigenvalue problem}
\begin{equation}
  \boxed{\;\Kop\,\vx \;=\; \lambda\,\Mop\,\vx,\qquad \lambda=\omega^2\;}
  \label{eq:evp-gevp}
\end{equation}
where $\Kop$ is the $\nu$-weighted curl--curl \emph{stiffness} and $\Mop$ the
$\varepsilon$-weighted \emph{mass} (\cref{ch:reductions-pde}). The pair
$(\Kop,\Mop)$ is called a \emph{matrix pencil}: not one matrix but two, and the
eigenproblem lives in the family $\Kop-\lambda\Mop$ of their combinations. Each
eigenvalue $\lambda$ gives a resonant frequency $\omega=\sqrt\lambda$; each
eigenvector $\vx$ gives the mode field. The word ``generalized'' refers to the
$\Mop$ on the right: an \emph{ordinary} eigenproblem reads $\mat{A}\vx=\lambda\vx$
with the identity in $\Mop$'s place, but here the $\varepsilon$-weighting puts a
genuine mass matrix there.

\begin{notationbox}
A \emph{pencil} is an ordered pair of matrices $(\Kop,\Mop)$ of the same size; we
write the generalized eigenproblem as $\Kop\vx=\lambda\Mop\vx$ or, equivalently,
$(\Kop-\lambda\Mop)\vx=\mathbf 0$. A scalar $\lambda$ for which $\Kop-\lambda\Mop$
is singular is an \emph{eigenvalue}; a nonzero $\vx$ in its null space is the
corresponding \emph{eigenvector}; together they form an \emph{eigenpair}
$(\lambda,\vx)$. When $\Kop$ and $\Mop$ are symmetric and $\Mop$ is positive
definite---as they are for a lossless cavity---all eigenvalues are real and the
eigenvectors can be chosen $\Mop$-orthogonal, $\inner{\Mop\vx_i}{\vx_j}=0$ for
$i\neq j$.
\end{notationbox}

So far this is bookkeeping. The difficulty---the one that shapes the entire
chapter---is \emph{which} eigenpairs we want, and \emph{how many}.

\paragraph{We want a few, not all.} A finely meshed three-dimensional cavity has
$N$ degrees of freedom running into the millions (\cref{ch:krylov}), so the
pencil has millions of eigenvalues. The engineer does not want them all. A
filter designer wants the lowest handful of resonances---the fundamental mode and
a few overtones, perhaps ten eigenpairs out of ten million. Computing the entire
spectrum would be absurd: it would cost more than a million linear solves and
produce a flood of numbers no one will read. The eigenmode analysis is, by
construction, a request for a \emph{small slice} of an enormous spectrum.

\paragraph{We want the smallest, which is the hard end.} Worse, the slice we want
sits at the \emph{bottom}. The lowest resonant frequencies are the smallest
$\lambda$, and the smallest nonzero eigenvalues of a matrix are, as we are about
to see, exactly the ones simple iteration is worst at finding. Iteration
naturally amplifies the \emph{largest} eigenvalues; the smallest are buried,
their signal swamped. Finding the interior or lower end of a spectrum is a
genuinely harder problem than finding the dominant end, and the rest of this
chapter is the story of how we cheat our way around that difficulty.

\Cref{fig:evp-setup} draws the situation: the pencil on the left, the vast
spectrum on the right, and the small wanted band marked at its low end.

\begin{figure}[t]
\centering
\begin{tikzpicture}[node distance=10mm]
  % the pencil, left
  \node[stage, minimum width=20mm, minimum height=10mm] (K) at (0,0.9) {$\Kop$};
  \node[stage, minimum width=20mm, minimum height=10mm] (M) at (0,-0.5) {$\Mop$};
  \node[font=\scriptsize\itshape, simGray, above=1mm of K] {stiffness};
  \node[font=\scriptsize\itshape, simGray, below=1mm of M] {mass};
  \node[draw=simGray, densely dashed, rounded corners=3pt, fit=(K)(M),
        inner sep=4pt, label={[font=\scriptsize\itshape,simGray]left:pencil}] (pen) {};
  % the eigenproblem arrow
  \node[font=\small] (eq) at (3.0,0.2) {$\Kop\vx=\lambda\Mop\vx$};
  \draw[flow] (pen.east) -- (eq.west);
  % the spectrum axis, right
  \draw[flow, simGray] (4.7,-1.2) -- (12.6,-1.2)
    node[right, font=\scriptsize\itshape, text=simGray] {$\lambda$};
  % many eigenvalue ticks, dense toward the right (large)
  \foreach \x in {5.1,5.45,5.95,6.6,7.4,8.3,9.2,10.0,10.6,11.0,11.3,11.55,11.75,11.9,12.05,12.18,12.3}{
    \draw[simBlue, thick] (\x,-1.32) -- (\x,-1.08);}
  % the wanted low band
  \fill[simRust!18] (4.95,-1.55) rectangle (6.05,-0.85);
  \draw[simRust, thick] (4.95,-1.55) rectangle (6.05,-0.85);
  \node[font=\scriptsize\bfseries, simRust, align=center] at (5.5,-0.55)
    {wanted:\\lowest few};
  % the "dominant" end iteration finds
  \node[font=\scriptsize\itshape, simGray, align=center] at (12.0,-0.55)
    {what iteration\\finds (dominant)};
  \draw[refedge, simGray] (12.0,-0.85) -- (12.2,-1.06);
\end{tikzpicture}
\caption[The generalized eigenproblem and its wanted low end]{The eigenmode
problem in one picture. The assembled pencil $(\Kop,\Mop)$ (left) defines the
generalized eigenproblem $\Kop\vx=\lambda\Mop\vx$, whose spectrum (right) holds
millions of eigenvalues. The engineer wants only the \emph{lowest few} (rust
band)---the fundamental resonance and a small number of overtones. But simple
iteration naturally surfaces the \emph{dominant} (largest) eigenvalues at the far
right, the exact opposite end. Bridging that mismatch is the business of the
chapter.}
\label{fig:evp-setup}
\end{figure}

\section{Power iteration: climbing toward the dominant direction}
\label{sec:power}

Before we can bend the spectrum we must understand the one iteration that all of
eigenvalue computation is built on. It is almost embarrassingly simple, and
understanding exactly what it does---and what it cannot do---is the whole key.

Take an ordinary eigenproblem $\mat{A}\vx=\mu\vx$ for the moment (we restore the
pencil shortly). \emph{Power iteration} does just one thing: it multiplies a
vector by $\mat{A}$ over and over, renormalizing to keep it bounded.
\begin{equation}
  \vx_{k+1} \;=\; \frac{\mat{A}\,\vx_k}{\norm{\mat{A}\,\vx_k}}.
  \label{eq:power-step}
\end{equation}
Why this converges to anything is the eigenvalue story in miniature. Expand the
starting vector in the eigenbasis $\{\vu_1,\dots,\vu_N\}$ of $\mat{A}$, ordered so
that $\abs{\mu_1}>\abs{\mu_2}\ge\dots\ge\abs{\mu_N}$:
\begin{equation}
  \vx_0 = c_1\vu_1 + c_2\vu_2 + \dots + c_N\vu_N .
\end{equation}
Applying $\mat{A}$ once multiplies each component $\vu_i$ by its eigenvalue
$\mu_i$; applying it $k$ times multiplies by $\mu_i^k$:
\begin{equation}
  \mat{A}^k\vx_0
    = c_1\mu_1^k\vu_1 + c_2\mu_2^k\vu_2 + \dots
    = \mu_1^k\Bigl(c_1\vu_1
        + c_2\Bigl(\tfrac{\mu_2}{\mu_1}\Bigr)^{\!k}\vu_2 + \dots\Bigr).
  \label{eq:power-expand}
\end{equation}
Because $\abs{\mu_2/\mu_1}<1$, every ratio $(\mu_i/\mu_1)^k$ decays to zero as
$k$ grows. The $\vu_1$ term, scaled by the largest eigenvalue, swamps all the
others. After enough steps the iterate points (up to sign and scale) in the
direction of $\vu_1$, the \emph{dominant eigenvector}---the eigenvector belonging
to the eigenvalue of largest magnitude. \Cref{fig:power-rotate} shows the
geometry: a generic starting vector swings, step by step, toward the dominant
axis as the subdominant component shrinks.

\begin{figure}[t]
\centering
\begin{tikzpicture}[scale=1.0]
  % axes: dominant eigenvector u1 (horizontal-ish), subdominant u2 (vertical-ish)
  \draw[->, simGray] (0,0) -- (5.0,0) node[right,font=\scriptsize]{$\vu_1$ (dominant)};
  \draw[->, simGray] (0,0) -- (0,3.4) node[above,font=\scriptsize]{$\vu_2$};
  % a sequence of normalized vectors rotating toward u1
  \foreach \ang/\op in {65/0.30, 46/0.42, 30/0.54, 18/0.66, 10/0.80, 5/1.0}{
    \draw[-{Stealth[length=2.2mm]}, very thick, simBlue, opacity=\op]
      (0,0) -- (\ang:3.0);
  }
  % final converged vector emphasized
  \draw[-{Stealth[length=2.6mm]}, very thick, simRust] (0,0) -- (5:3.0);
  \node[font=\scriptsize, simRust] at (5:3.4) {$\vx_k$};
  \node[font=\scriptsize, simBlue] at (66:3.3) {$\vx_0$};
  % arc showing the swing
  \draw[->, simGray, densely dotted] (62:2.0) arc (62:9:2.0);
  \node[font=\scriptsize\itshape, simGray] at (35:2.35) {each step};
  \node[font=\scriptsize\itshape, simGray] at (33:1.85) {$\times\,\mu_2/\mu_1$};
\end{tikzpicture}
\caption[Power iteration rotates toward the dominant eigenvector]{Power
iteration as a rotation. The starting vector $\vx_0$ (faint blue) has a component
along every eigenvector; each multiplication by $\mat{A}$ multiplies the
$\vu_2$ component by $\mu_2/\mu_1<1$ relative to the $\vu_1$ component, so the
iterate swings, step by step, toward the dominant axis $\vu_1$. The converged
$\vx_k$ (rust) points along $\vu_1$. The speed of the swing is set by the ratio
$\abs{\mu_2/\mu_1}$: the closer to $1$, the slower the rotation.}
\label{fig:power-rotate}
\end{figure}

\subsection{The Rayleigh quotient: reading off the eigenvalue}

Power iteration delivers a \emph{direction}---the eigenvector. To get the
\emph{eigenvalue} we evaluate the \emph{Rayleigh quotient}, the single most
useful scalar in eigenvalue computation:
\begin{equation}
  \rho(\vx) \;=\; \frac{\inner{\mat{A}\vx}{\vx}}{\inner{\vx}{\vx}}.
  \label{eq:rayleigh}
\end{equation}
If $\vx$ is exactly an eigenvector, $\mat{A}\vx=\mu\vx$, then
$\rho(\vx)=\inner{\mu\vx}{\vx}/\inner{\vx}{\vx}=\mu$: the quotient returns the
eigenvalue exactly. If $\vx$ is only \emph{near} an eigenvector, $\rho(\vx)$ is a
remarkably good estimate---for a symmetric $\mat{A}$ it is accurate to
\emph{second} order in the eigenvector error, so a direction good to three digits
yields an eigenvalue good to six. As power iteration rotates the iterate toward
$\vu_1$, the Rayleigh quotient $\rho(\vx_k)$ closes in on $\mu_1$, and the two
together---direction from \cref{eq:power-step}, value from
\cref{eq:rayleigh}---form a complete eigenpair estimate.

\begin{figure}[t]
\centering
\begin{tikzpicture}[node distance=8mm]
  \node[opbox, minimum width=20mm] (x) {$\vx$};
  \node[opbox, minimum width=22mm, right=12mm of x] (Ax) {$\mat{A}\vx$};
  \node[opbox, fill=simBgGold, draw=simGold, minimum width=30mm, right=12mm of Ax] (num)
    {$\inner{\mat{A}\vx}{\vx}$};
  \node[opbox, fill=simBgGold, draw=simGold, minimum width=26mm, below=6mm of num] (den)
    {$\inner{\vx}{\vx}$};
  \node[product, minimum width=20mm, right=14mm of num, yshift=-7mm] (rho)
    {$\rho(\vx)$};
  \draw[flow] (x) -- node[above,font=\scriptsize]{apply} (Ax);
  \draw[flow] (Ax) -- node[above,font=\scriptsize]{dot} (num);
  \draw[flow] (x.south) to[out=-90,in=180] (den.west);
  \draw[flow] (num.east) -- node[above,font=\scriptsize]{$\div$} (rho.north west);
  \draw[flow] (den.east) -- (rho.south west);
  \node[font=\scriptsize\itshape, simGray, below=1mm of den]
    {if $\mat{A}\vx=\mu\vx$ then $\rho=\mu$ exactly};
\end{tikzpicture}
\caption[The Rayleigh quotient]{The Rayleigh quotient $\rho(\vx)=\inner{\mat{A}
\vx}{\vx}/\inner{\vx}{\vx}$ turns a vector into an eigenvalue estimate with one
operator apply and two inner products. On an exact eigenvector it returns the
eigenvalue exactly; near an eigenvector it is accurate to second order in the
direction error. Paired with power iteration's direction estimate, it completes
the eigenpair.}
\label{fig:rayleigh}
\end{figure}

\subsection{Several at once: subspace iteration}

We want a \emph{handful} of eigenpairs, not one. The natural generalization
iterates a whole \emph{block} of vectors rather than a single one.
\emph{Subspace} (or \emph{orthogonal}) iteration carries an $N\times p$ matrix
$\mat{X}_k$ of $p$ columns, applies $\mat{A}$ to all of them, and then
\emph{re-orthonormalizes} the columns:
\begin{equation}
  \mat{Y} = \mat{A}\,\mat{X}_k,\qquad
  \mat{X}_{k+1} = \operatorname{orth}(\mat{Y}),
\end{equation}
where $\operatorname{orth}$ is a QR-style orthonormalization
(\cref{ch:krylovschur}). Without the re-orthonormalization every column would
collapse onto the same dominant eigenvector; the orthogonalization forces the
columns apart, so they converge to the $p$ \emph{dominant} eigenvectors
$\vu_1,\dots,\vu_p$ simultaneously. The Rayleigh quotient generalizes too: the
small $p\times p$ matrix $\mat{X}_k^{\!\top}\mat{A}\,\mat{X}_k$ has eigenvalues
that approximate the $p$ dominant eigenvalues---the \emph{Rayleigh--Ritz}
projection that the Krylov methods of \cref{ch:krylovschur} sharpen into their
final form.

\begin{keyidea}
Power iteration, and its block form subspace iteration, converge to the
\emph{dominant} eigenpairs---those of \emph{largest} magnitude. The Rayleigh
quotient $\rho(\vx)=\inner{\mat{A}\vx}{\vx}/\inner{\vx}{\vx}$ reads the eigenvalue
off the converged direction. The convergence rate is governed by the
\emph{ratio} of consecutive eigenvalues $\abs{\mu_2/\mu_1}$: well-separated
eigenvalues converge fast, clustered ones agonizingly slowly. This is the engine
under every eigensolver---but on its own it finds exactly the wrong end of the
spectrum for the eigenmode problem.
\end{keyidea}

\subsection{Why this is not enough}

Two facts about power iteration doom it as a direct attack on the eigenmode
problem, and naming them sharply is what motivates the spectral transform.

\emph{First, it finds the largest, and we want the smallest.} The eigenmode
analysis wants the lowest resonances, the smallest nonzero $\lambda$. Power
iteration delivers the dominant eigenvalue, $\lambda_{\max}$---the
\emph{highest} resonance, of no interest, and in fact a meaningless mesh-resolution
artifact at the top of the discrete spectrum.

\emph{Second, it is slow exactly where we need it.} The convergence rate is
$\abs{\mu_2/\mu_1}$. The low resonances of a cavity are typically
\emph{clustered}---a fundamental and its near-degenerate overtones sit close
together---so consecutive eigenvalues have ratios near $1$, and power iteration
crawls. The very structure of the problem we care about is the structure power
iteration handles worst.

We could try to fix the first problem by iterating with $\mat{A}^{-1}$ instead of
$\mat{A}$: the eigenvalues of $\mat{A}^{-1}$ are $1/\mu_i$, so the smallest $\mu_i$
becomes the largest $1/\mu_i$, and power iteration on $\mat{A}^{-1}$ converges to
the \emph{smallest} eigenvector. This is \emph{inverse iteration}, and it is the
right instinct---but it points only at the very smallest eigenvalue, does nothing
for the clustering, and gives us no control over \emph{where} in the spectrum to
look. The spectral transform is inverse iteration grown up: it lets us aim.

\section{The spectral transform: shift-and-invert}
\label{sec:shift-invert}

Here is the central idea of the chapter. We cannot easily find interior or low
eigenvalues directly. So we \emph{transform} the operator into a new one whose
\emph{dominant} eigenvalues are precisely the original eigenvalues we want---and
then we run the power/subspace iteration we already understand on the transformed
operator, where it works beautifully.

\subsection{Deriving the eigenvalue map}

Start from the generalized eigenproblem \cref{eq:evp-gevp},
$\Kop\vx=\lambda\Mop\vx$. Choose a \emph{shift} $\sigma$---a number near the
eigenvalues we are hunting---and rearrange. Subtract $\sigma\Mop\vx$ from both
sides:
\begin{equation}
  \Kop\vx - \sigma\Mop\vx = \lambda\Mop\vx - \sigma\Mop\vx
  \quad\Longrightarrow\quad
  (\Kop-\sigma\Mop)\,\vx = (\lambda-\sigma)\,\Mop\,\vx .
\end{equation}
Now suppose $\lambda\neq\sigma$, so $\Kop-\sigma\Mop$ is invertible on $\vx$.
Multiply both sides by $(\Kop-\sigma\Mop)^{-1}$ and divide by the scalar
$(\lambda-\sigma)$:
\begin{equation}
  \frac{1}{\lambda-\sigma}\,\vx
    \;=\; (\Kop-\sigma\Mop)^{-1}\,\Mop\,\vx .
\end{equation}
Read this as an \emph{ordinary} eigenproblem for the transformed operator
\begin{equation}
  \boxed{\;\mat{T}_\sigma \;=\; (\Kop-\sigma\Mop)^{-1}\,\Mop,
    \qquad \mat{T}_\sigma\,\vx \;=\; \theta\,\vx,
    \quad \theta=\frac{1}{\lambda-\sigma}\;}
  \label{eq:shift-invert}
\end{equation}
The eigenvectors are \emph{unchanged}---$\vx$ is an eigenvector of $\mat{T}_\sigma$
exactly when it is an eigenvector of the pencil---but the eigenvalues are mapped
by
\begin{equation}
  \lambda \;\longmapsto\; \theta = \frac{1}{\lambda-\sigma}.
  \label{eq:eigmap}
\end{equation}

Now examine what this map does. The transformed eigenvalue $\theta=1/(\lambda
-\sigma)$ is \emph{large in magnitude exactly when $\lambda$ is close to
$\sigma$}, and small when $\lambda$ is far from $\sigma$. An eigenvalue sitting
right at the shift is sent to infinity; one far away is sent near zero. So the
eigenvalues of the original problem \emph{nearest the shift $\sigma$} become the
\emph{dominant} eigenvalues of $\mat{T}_\sigma$---and those are exactly the ones
power and subspace iteration find first and fastest.

\begin{keyidea}
\textbf{Shift-and-invert turns a hard interior eigenvalue problem into an easy
dominant one.} Iterating against $\mat{T}_\sigma=(\Kop-\sigma\Mop)^{-1}\Mop$
instead of the pencil maps each eigenvalue $\lambda\mapsto\theta=1/(\lambda
-\sigma)$, so the eigenvalues \emph{nearest the shift} $\sigma$ become the
\emph{largest} (dominant) eigenvalues of $\mat{T}_\sigma$. ``Find the eigenvalues
near $\sigma$''---hard, interior, what the eigenmode problem actually
needs---becomes ``find the dominant eigenvalues of $\mat{T}_\sigma$''---easy,
exactly what iteration does. The price is that every application of
$\mat{T}_\sigma$ requires an \emph{inner linear solve} of $(\Kop-\sigma\Mop)$,
turning each iteration step into a \code{ksp\_solve} (\cref{ch:krylov}).
\end{keyidea}

\subsection{The spectrum reordering, made vivid}

The map \cref{eq:eigmap} does more than enlarge the wanted eigenvalues: it
\emph{reorders} the entire spectrum. The next worked example makes the reordering
concrete with a four-eigenvalue spectrum, and \cref{fig:spectrum-reorder}---the
hero figure of the chapter---draws the before and after.

\begin{workedexample}{Shift-and-invert reorders a small spectrum}{reorder}
A pencil $(\Kop,\Mop)$ has four eigenvalues
\[
  \lambda_1 = 1,\quad \lambda_2 = 4,\quad \lambda_3 = 9,\quad \lambda_4 = 16,
\]
and we want the two nearest $\lambda=4$. Plain subspace iteration on the pencil
would converge to $\lambda_4=16$ and $\lambda_3=9$ first---the dominant pair, the
\emph{wrong} ones. Choose the shift $\sigma=4.5$, just above the target, and form
the transformed eigenvalues $\theta_i=1/(\lambda_i-\sigma)$:
\[
  \begin{aligned}
  \theta_1 = \frac{1}{1-4.5} = -0.286,\quad
  &\theta_2 = \frac{1}{4-4.5} = -2.0,\\
  \theta_3 = \frac{1}{9-4.5} = 0.222,\quad
  &\theta_4 = \frac{1}{16-4.5} = 0.087.
  \end{aligned}
\]
Rank them by magnitude:
\[
  \abs{\theta_2}=2.0 \;>\; \abs{\theta_1}=0.286 \;>\; \abs{\theta_3}=0.222
  \;>\; \abs{\theta_4}=0.087.
\]
The order has \emph{inverted}. In the original spectrum $\lambda_4=16$ was
dominant and $\lambda_2=4$ was a small interior value; after the transform
$\theta_2$ (from $\lambda_2=4$, nearest the shift) is by far the largest, and
$\theta_1$ (from $\lambda_1=1$, the next nearest) is second. Subspace iteration on
$\mat{T}_\sigma$ now converges \emph{first} to exactly the two eigenpairs we
wanted---those near $\sigma=4.5$---and the convergence is fast, because
$\abs{\theta_1/\theta_2}=0.286/2.0\approx 0.14$ is a small ratio, far from the
$\abs{\lambda_3/\lambda_4}=9/16\approx 0.56$ ratio that throttled the untransformed
iteration. We get the right eigenpairs \emph{and} faster convergence, both from
the single act of bending the spectrum.

\medskip
\emph{Reading the answer back.} The iteration returns $\theta_2=-2.0$. We recover
the physical eigenvalue by inverting \cref{eq:eigmap}:
$\lambda = \sigma + 1/\theta = 4.5 + 1/(-2.0) = 4.5 - 0.5 = 4$, exactly
$\lambda_2$. The eigenvector needs no untransforming---it is the \emph{same}
vector $\vx$ in both problems.
\end{workedexample}

\begin{figure}[t]
\centering
\begin{tikzpicture}
\begin{axis}[
    name=before,
    width=0.82\linewidth, height=34mm,
    xlabel={original eigenvalue $\lambda$},
    xmin=-1, xmax=18, ymin=0, ymax=1.4,
    axis lines=left, ytick=\empty, xtick={1,4,9,16},
    xticklabel style={font=\scriptsize}, xlabel style={font=\scriptsize},
    clip=false,
  ]
  % the four eigenvalues as stems
  \addplot[ycomb, very thick, simBlue, mark=*, mark size=2pt] coordinates {
    (1,1)(4,1)(9,1)(16,1)};
  % shift marker
  \draw[simRust, thick, densely dashed] (axis cs:4.5,0) -- (axis cs:4.5,1.3);
  \node[font=\scriptsize, simRust] at (axis cs:4.5,1.42) {shift $\sigma=4.5$};
  % wanted band
  \node[font=\scriptsize\itshape, simGray] at (axis cs:13,1.15) {dominant end};
  \node[font=\scriptsize, simBlue, anchor=south] at (axis cs:1,1) {$\lambda_1$};
  \node[font=\scriptsize, simBlue, anchor=south] at (axis cs:4,1) {$\lambda_2$};
  \node[font=\scriptsize, simBlue, anchor=south] at (axis cs:9,1) {$\lambda_3$};
  \node[font=\scriptsize, simBlue, anchor=south] at (axis cs:16,1) {$\lambda_4$};
\end{axis}
\begin{axis}[
    name=after, at={(before.south west)}, anchor=north west, yshift=-9mm,
    width=0.82\linewidth, height=40mm,
    xlabel={transformed eigenvalue $\theta=1/(\lambda-\sigma)$},
    ylabel={}, xmin=-2.6, xmax=2.6, ymin=0, ymax=2.4,
    axis lines=left, ytick=\empty,
    xtick={-2.0,-0.286,0.087,0.222}, xticklabels={$-2.0$,,,},
    xticklabel style={font=\scriptsize}, xlabel style={font=\scriptsize},
    clip=false,
  ]
  % transformed eigenvalues, height = magnitude (dominance)
  \addplot[ycomb, very thick, simRust, mark=*, mark size=2pt] coordinates {
    (-2.0,2.0)(-0.286,0.286)(0.222,0.222)(0.087,0.087)};
  \node[font=\scriptsize, simRust, anchor=south] at (axis cs:-2.0,2.0) {$\theta_2$ (was $\lambda_2$)};
  \node[font=\scriptsize, simGray, anchor=south] at (axis cs:-0.286,0.286) {$\theta_1$};
  \node[font=\scriptsize, simGray, anchor=south west] at (axis cs:0.222,0.222) {$\theta_3$};
  \node[font=\scriptsize, simGray, anchor=south west] at (axis cs:0.087,0.087) {$\theta_4$};
  \node[font=\scriptsize\itshape, simRust] at (axis cs:-1.4,1.7) {now dominant};
\end{axis}
\end{tikzpicture}
\caption[Shift-and-invert reorders the spectrum]{The spectral transform of
\cref{wex:reorder}, drawn. \emph{Top:} the original spectrum, four eigenvalues
$\{1,4,9,16\}$; the wanted pair sits near the shift $\sigma=4.5$ (rust dashed),
in the interior, while iteration's natural target is the dominant end at the
right. \emph{Bottom:} after the map $\theta=1/(\lambda-\sigma)$, the stem heights
show transformed magnitude (dominance). $\theta_2$ (from $\lambda_2=4$, nearest
the shift) towers over the rest---it has become the dominant eigenvalue---and
$\theta_1$ is second. Iteration on the transformed operator now finds exactly the
near-$\sigma$ pair first. Bending the spectrum turns the hard interior request
into the easy dominant one.}
\label{fig:spectrum-reorder}
\end{figure}

\subsection{The cost: an inner linear solve per apply}

The transform is not free, and the price is precisely located. To run iteration
on $\mat{T}_\sigma=(\Kop-\sigma\Mop)^{-1}\Mop$ we must, at every step, compute
$\mat{T}_\sigma\vv$ for the current iterate $\vv$. Written out, that is three
moves:
\begin{equation}
  \vv
  \;\xrightarrow{\ \text{apply }\Mop\ }\;
  \vw = \Mop\vv
  \;\xrightarrow{\ \text{inner solve}\ }\;
  \vy = (\Kop-\sigma\Mop)^{-1}\vw
  \;\xrightarrow{\ \text{untransform}\ }\;
  \vy .
  \label{eq:siapply}
\end{equation}
The first move is a cheap matrix--vector product against the mass matrix. The
\emph{second} move is the expensive one: applying $(\Kop-\sigma\Mop)^{-1}$ means
\emph{solving the linear system} $(\Kop-\sigma\Mop)\vy=\vw$. That is exactly the
sparse linear solve of \cref{ch:krylov}---a \code{ksp\_solve}---and because the
shifted matrix $\Kop-\sigma\Mop$ is large and sparse, it is itself solved
iteratively, almost always with a preconditioner (\cref{ch:multigrid}). The third
move is per-backend coordinate bookkeeping (an optional scaling un-twist) that
does not change the direction.

So the body of one shift-invert iteration step is the composition
\begin{equation}
  \text{apply }\Mop \;\rhd\; \text{inner }\code{ksp\_solve} \;\rhd\;
  \text{scale-untransform},
\end{equation}
and the \emph{outer} eigenvalue iteration wraps this body in its power/subspace
loop. \Cref{fig:si-apply} draws the body. This is the structural heart of the
shift-invert eigensolve: \emph{an outer eigenvalue iteration whose every step
contains an inner linear solve}. The solver of \cref{ch:krylov,ch:gmres} that we
built to solve $\mat{A}\vx=\vb$ now reappears as a \emph{subroutine called once
per eigen-iteration step}.

\begin{figure}[t]
\centering
\begin{tikzpicture}[node distance=9mm]
  \node[data, minimum width=14mm] (v) {$\vv$};
  \node[opbox, minimum width=22mm, right=10mm of v] (Mv) {apply $\Mop$};
  \node[extern, minimum width=34mm, right=10mm of Mv] (solve)
    {\textbf{inner }\code{ksp\_solve}\\[1pt]\footnotesize $(\Kop-\sigma\Mop)^{-1}$};
  \node[opbox, minimum width=26mm, right=10mm of solve] (un) {scale-untransform};
  \node[product, minimum width=20mm, right=10mm of un] (out) {$\mat{T}_\sigma\vv$};
  \draw[flow] (v) -- (Mv);
  \draw[flow] (Mv) -- node[above,font=\scriptsize]{$\vw$} (solve);
  \draw[flow] (solve) -- node[above,font=\scriptsize]{$\vy$} (un);
  \draw[flow] (un) -- (out);
  % annotations
  \node[font=\scriptsize\itshape, simGray, below=2mm of Mv] {cheap matvec};
  \node[font=\scriptsize\itshape, simRust, below=2mm of solve, align=center]
    {the cost:\\a full linear solve\\(\cref{ch:krylov,ch:multigrid})};
  \node[font=\scriptsize\itshape, simGray, below=2mm of un] {coordinate twist};
  % the outer loop wrapping
  \draw[backflow] (out.north) to[out=90,in=90]
    node[above,font=\scriptsize,simRust]{outer eigen-iteration repeats} (v.north);
\end{tikzpicture}
\caption[The shift-invert apply and its inner solve]{One application of the
transformed operator $\mat{T}_\sigma=(\Kop-\sigma\Mop)^{-1}\Mop$, the body of the
eigen-iteration. A matrix--vector product against the mass matrix $\Mop$ (cheap)
feeds an \emph{inner linear solve} of the shifted operator $\Kop-\sigma\Mop$
(expensive: a full \code{ksp\_solve}, \cref{ch:krylov}, usually preconditioned,
\cref{ch:multigrid}), followed by a per-backend scale-untransform. The outer
eigenvalue iteration (dashed rust) repeats this body each step. The linear solver
of the earlier chapters becomes a subroutine of the eigensolver.}
\label{fig:si-apply}
\end{figure}

\section{The eigensolver as a black box}
\label{sec:blackbox}

We now know, in principle, how to find the low resonances: transform the spectrum
with a shift, run subspace iteration on the transformed operator, recover the
eigenpairs. But Palace does \emph{not} write this loop. The eigenmode analysis is
the one analysis in the book whose ``solve'' stage is a single, opaque call to a
specialized library---and this section explains why, and what it means for the
shape of the code.

\subsection{The fourth solve corner}

\Cref{ch:situating} sorted every analysis's middle stage into one of \emph{four
solve corners}. Three of them are loops Palace writes: the fixed-operator family
solve (electrostatics), the operator-varying sweep (driven), the sequential fold
(transient). The fourth corner is different in kind: the \emph{opaque black-box
solve}, and the eigenmode analysis is its sole inhabitant. As \cref{ch:operators}
put it, a solver is an operator you \emph{apply}---and the eigensolver takes that
to its limit. The entire outer iteration lives \emph{inside a library}; Palace
hands over the operators and reads back the answer, authoring no loop at all.

The reason is practical and decisive. The outer eigenvalue iteration is not the
simple power iteration of \cref{sec:power}; production eigensolvers use
\emph{Krylov--Schur} or \emph{Arnoldi} methods (\cref{ch:krylovschur}), which
maintain and periodically ``restart'' a growing Krylov basis, run delicate
re-orthogonalization, apply locking and purging of converged pairs, and decide
convergence by subtle residual criteria. This machinery is intricate, easy to get
subtly wrong, and already implemented to a very high standard in specialized
libraries---SLEPc, ARPACK. Re-implementing it would be folly. So Palace composes
against it: it assembles $(\Kop,\Mop)$, configures the shift and the inner solver,
and calls the library's \code{Solve}.

\subsection{Reverse communication: the library calls back}

There is one wrinkle that makes the black box especially elegant. The library
drives the iteration, but it does not contain Palace's operators---it cannot,
because those operators are matrix-free and live in Palace's data structures
(\cref{ch:matrixfree}). So the library and Palace cooperate through
\emph{reverse communication}: the library runs its loop, and whenever it needs a
matrix--vector product, it \emph{pauses and calls back} into Palace, saying ``apply
the operator to this vector and hand me the result.'' Palace supplies the
\code{apply}---which, for shift-invert, is the whole inner-solve body of
\cref{fig:si-apply}---and the library resumes.

\Cref{fig:blackbox-rci} draws the control flow. The outer loop is the library's;
the apply callbacks are Palace's; and the line between them is the one interface
of \cref{ch:operators}: \code{apply(op, v)}. The eigensolver never sees Palace's
mesh, materials, or assembly; Palace never sees the library's Krylov basis,
restart logic, or convergence test. Each side holds what it owns and speaks to the
other only through the apply interface.

\begin{figure}[t]
\centering
\begin{tikzpicture}[node distance=10mm]
  % the library box (drives the loop)
  \node[extern, minimum width=40mm, minimum height=26mm] (lib) at (0,0) {};
  \node[font=\footnotesize\bfseries, simViolet] at (lib.north) [yshift=-3.5mm]
    {eigensolver library};
  \node[font=\scriptsize\itshape, simGray, align=center] at (lib.center) [yshift=2mm]
    {Krylov--Schur loop\\(\cref{ch:krylovschur}):\\restart, orthogonalize,\\
     lock, convergence test};
  \node[font=\scriptsize, simViolet] at (lib.south) [yshift=3.5mm] {SLEPc / ARPACK};
  % Palace apply box
  \node[stage, minimum width=34mm, minimum height=20mm, right=34mm of lib] (pal) {};
  \node[font=\footnotesize\bfseries, simBlue] at (pal.north) [yshift=-3.5mm]
    {Palace \code{apply}};
  \node[font=\scriptsize, simInk, align=center] at (pal.center) [yshift=-1mm]
    {$\vv\mapsto\mat{T}_\sigma\vv$:\\apply $\Mop$, inner solve,\\untransform
     (\cref{fig:si-apply})};
  % the reverse-communication callback arrows
  \draw[flow] (lib.east|-0,0.5) -- node[above,font=\scriptsize]{``apply to $\vv$''}
    (pal.west|-0,0.5);
  \draw[backflow] (pal.west|-0,-0.5) -- node[below,font=\scriptsize]{result $\mat{T}_\sigma\vv$}
    (lib.east|-0,-0.5);
  % setup in, eigenpairs out
  \node[data, align=center, text width=22mm, above=8mm of lib] (in)
    {assemble $(\Kop,\Mop)$,\\ shift $\sigma$, inner ksp};
  \draw[flow] (in) -- (lib);
  \node[product, align=center, text width=22mm, below=8mm of lib] (out)
    {converged\\eigenpairs $(\lambda_i,\vx_i)$};
  \draw[flow] (lib) -- (out);
  % labels for who owns what
  \node[font=\scriptsize\itshape, simViolet, above=0.5mm of in] {one opaque call};
  \node[font=\scriptsize\bfseries, simInk] at ($(lib.east)!0.5!(pal.west)+(0,1.6)$)
    {one interface: \code{apply}};
\end{tikzpicture}
\caption[The eigensolver as a black box with reverse-communication callbacks]{The
black-box eigensolve. Palace assembles the pencil $(\Kop,\Mop)$, picks a shift
$\sigma$ and an inner linear solver, and makes \emph{one opaque call} into a
specialized library (SLEPc, ARPACK). The library \emph{drives} the outer
iteration---the Krylov--Schur restart, orthogonalization, locking, and convergence
test of \cref{ch:krylovschur}---and whenever it needs a matrix--vector product it
\emph{calls back} into Palace through the one \code{apply} interface of
\cref{ch:operators} (\emph{reverse communication}). Palace's apply is the entire
shift-invert body of \cref{fig:si-apply}. Converged eigenpairs come back out. No
Palace-authored loop exists.}
\label{fig:blackbox-rci}
\end{figure}

\subsection{Contrast with the linear-solve corners}

The contrast with the other three solve corners is sharp and worth stating
plainly. In the electrostatic, driven, and transient analyses, Palace
\emph{writes the loop}: the per-terminal family loop, the frequency sweep, the
time march are all Palace source, and the linear solver they call is a
\code{ksp\_solve} \emph{whose iteration Palace also owns} (\cref{ch:krylov}).
There the solver is a constructed operator Palace builds and drives.

In the eigenmode analysis the ownership inverts. The \emph{outer} loop belongs to
the library; Palace owns only the \emph{apply} the library calls back for. The
eigensolve is therefore the deepest composition in the book: a library-owned outer
iteration, calling back into a Palace-owned shift-invert body, which itself
contains a Palace-owned \emph{inner} \code{ksp\_solve}, which itself calls a
preconditioner. Three nested layers of ``a solver is an operator''
(\cref{ch:operators}), and the outermost one is not even Palace's code.

\begin{pitfall}
Because the outer eigen-iteration is library-owned, the eigenmode analysis is the
one place in this book where the constructive ``iteration rotation'' of the later
chapters \emph{cannot reach the loop}: there is no Palace-authored kernel to
re-express. The shift-invert \emph{body} lifts cleanly to a global tensor-field
expression---it is just apply-$\Mop$, solve, untransform---but the \emph{loop}
around it is an opaque-library obstruction, a genuine boundary of what the
specification can rotate. The honest disposition is to document the boundary, not
to pretend a loop exists where the library hides one. The full constructive
realization of what the library does---Lanczos, Arnoldi, Krylov--Schur in our own
vocabulary---is \cref{ch:krylovschur}; here it stays a black box on purpose.
\end{pitfall}

\section{The quadratic eigenproblem: linearization}
\label{sec:quadratic}

The derivation so far assumed a lossless cavity, giving the linear pencil
$\Kop\vx=\lambda\Mop\vx$. Real cavities lose energy---to wall conductivity,
dielectric absorption, radiation through ports---and \cref{ch:reductions-pde}
showed that loss adds a damping operator $\Cop$ and makes the eigenproblem
\emph{quadratic} in the eigenvalue:
\begin{equation}
  \bigl(\Kop + \lambda\,\Cop + \lambda^2\,\Mop\bigr)\,\vx = \mathbf 0,
  \qquad \lambda = i\omega.
  \label{eq:qevp}
\end{equation}
Now $\lambda$ is complex: its imaginary part is the oscillation frequency, its
real part the decay rate. This is a \emph{polynomial eigenvalue problem}, and the
shift-invert iteration we built expects an ordinary or generalized \emph{linear}
problem. The standard remedy is \emph{linearization}: trade the quadratic
$N\times N$ problem for a linear $2N\times 2N$ one, at the cost of doubling the
size.

The trick is the same order-reduction we used for the transient ODE in
\cref{ch:reductions-pde}: introduce a second unknown for $\lambda\vx$. Set
$\vz=\lambda\vx$ and stack $(\vx,\vz)$. Then \cref{eq:qevp} becomes the linear
generalized eigenproblem
\begin{equation}
  \underbrace{\begin{pmatrix} \mathbf 0 & \mat{I} \\ -\Kop & -\Cop \end{pmatrix}}_{\tilde\Kop}
  \begin{pmatrix}\vx\\\vz\end{pmatrix}
  =
  \lambda
  \underbrace{\begin{pmatrix} \mat{I} & \mathbf 0 \\ \mathbf 0 & \Mop \end{pmatrix}}_{\tilde\Mop}
  \begin{pmatrix}\vx\\\vz\end{pmatrix},
  \label{eq:linearization}
\end{equation}
which one can multiply out to verify: the top block row reads $\vz=\lambda\vx$
(the definition of $\vz$), and the bottom row reads
$-\Kop\vx-\Cop\vz=\lambda\Mop\vz$, i.e.
$-\Kop\vx-\Cop\lambda\vx=\lambda^2\Mop\vx$, which rearranges to the quadratic
\cref{eq:qevp}. The doubled pencil $(\tilde\Kop,\tilde\Mop)$ is an ordinary
generalized eigenproblem of the kind \cref{sec:shift-invert} handles---so we apply
the \emph{same} shift-invert machinery to it, and the library solves it with the
\emph{same} black-box call, now on the $2N$-dimensional pencil. Each eigenpair of
the doubled problem yields the physical $(\lambda,\vx)$ directly; the extra block
$\vz=\lambda\vx$ is discarded.

For the still more general \emph{nonlinear} eigenproblem---where the operator
depends on $\lambda$ in a way no polynomial captures, $\bigl(\Kop+\lambda\Cop
+\lambda^2\Mop+\mat{A}_2(\lambda)\bigr)\vx=\mathbf 0$---linearization no longer
suffices, and one needs a genuinely nonlinear eigensolver built on
\emph{deflation} (finding eigenpairs one at a time and projecting out the ones
already found). That is the subject of \cref{ch:deflation}; we flag it here only
to place the quadratic case as the easy middle rung between the linear pencil and
the fully nonlinear problem.

\section{Reading the results}
\label{sec:results}

The library hands back a set of converged eigenpairs $(\lambda_i,\vx_i)$. Two
short steps turn them into the physical products the engineer asked for, and one
necessary cleanup keeps the answer honest.

\subsection{Eigenvalue to frequency, eigenvector to mode field}

The eigenvalue carries the frequency. For the lossless linear problem
$\lambda=\omega^2$, so
\begin{equation}
  \omega = \sqrt\lambda,\qquad
  f = \frac{\omega}{2\pi} = \frac{\sqrt\lambda}{2\pi}.
  \label{eq:freq}
\end{equation}
For the lossy quadratic problem $\lambda=i\omega$ is complex, so the frequency is
$\omega = \lambda/i = -i\lambda$, and the resonant frequency is read from its
imaginary part. The \emph{eigenvector} $\vx_i$ is, directly, the spatial pattern
of the mode field---the standing-wave shape the cavity rings in, the
three-dimensional analogue of the $\sin(n\pi x/L)$ modes we computed by hand in
\cref{ch:reductions-pde}.

The lossy case carries a bonus the lossless case cannot: the \emph{quality
factor} $Q$, the dimensionless measure of how slowly a mode loses energy. It
comes from the ratio of the imaginary part of $\lambda$ (the oscillation) to its
real part (the decay):
\begin{equation}
  Q = \frac{\abs{\Im\lambda}}{2\,\abs{\Re\lambda}}
    = \frac{\Im\omega}{2\,\Re\omega}\cdot(\text{convention}),
\end{equation}
a high $Q$ meaning a sharp, long-ringing resonance. The exact assembly of $Q$
from the complex eigenpair is an output reduction, treated with the other
postprocessing products in \cref{ch:reductions}; here we note only that the
eigenmode analysis is the one that \emph{produces} it.

\begin{figure}[t]
\centering
\begin{tikzpicture}[node distance=9mm]
  \node[data, minimum width=20mm] (pair) {$(\lambda_i,\vx_i)$};
  % eigenvalue branch
  \node[opbox, minimum width=26mm, right=12mm of pair, yshift=8mm] (lam)
    {$\lambda_i$ (eigenvalue)};
  \node[product, minimum width=30mm, right=12mm of lam] (freq)
    {$f_i=\sqrt{\lambda_i}/2\pi$};
  \node[opbox, minimum width=26mm, below=4mm of freq] (q)
    {$Q_i$ (from $\Im\lambda$)};
  % eigenvector branch
  \node[opbox, minimum width=26mm, right=12mm of pair, yshift=-8mm] (vec)
    {$\vx_i$ (eigenvector)};
  \node[product, minimum width=30mm, right=12mm of vec, yshift=-6mm] (mode)
    {mode field $\vE_i(\vx)$};
  \draw[flow] (pair.east) to[out=30,in=180] (lam.west);
  \draw[flow] (pair.east) to[out=-30,in=180] (vec.west);
  \draw[flow] (lam) -- (freq);
  \draw[flow] (lam.east) to[out=-20,in=160] (q.west);
  \draw[flow] (vec) -- (mode);
  \node[font=\scriptsize\itshape, simGray, below=1mm of mode]
    {(\cref{ch:reductions-pde} computed the 1-D version by hand)};
\end{tikzpicture}
\caption[From eigenpair to physical product]{Reading the eigensolver's output.
Each converged eigenpair $(\lambda_i,\vx_i)$ splits two ways. The
\emph{eigenvalue} $\lambda_i$ gives the resonant frequency $f_i=\sqrt{\lambda_i}/
2\pi$ (and, in the lossy case, the quality factor $Q_i$ from its imaginary part,
reduced in \cref{ch:reductions}). The \emph{eigenvector} $\vx_i$ \emph{is} the
mode field $\vE_i(\vx)$, the spatial standing-wave pattern---the
three-dimensional descendant of the hand-computed cavity modes of
\cref{ch:reductions-pde}.}
\label{fig:readout}
\end{figure}

\subsection{Spurious modes and the divergence-free projection}

There is one trap the raw eigensolver output sets, and the eigenmode analysis must
spring it before reporting results. The curl--curl operator $\Kop$ has a large
\emph{null space}---it annihilates every gradient field $\Grad\phi$, because
$\Curl\Grad\equiv\mathbf 0$ (the gauge freedom of \cref{ch:reductions-pde}). Every
such gradient field is an eigenvector with eigenvalue $\lambda=0$. These are not
physical resonances; they are \emph{spurious modes}, artifacts of the operator's
null space, clustered at and near zero frequency exactly where the eigensolver,
hunting the lowest eigenvalues, will dredge them up by the dozen.

Left in, they pollute the answer: the solver returns a heap of bogus zero-frequency
``modes'' that crowd out the real low resonances. The cure is the
\emph{divergence-free projection} of \cref{ch:bc}: project each candidate
eigenvector onto the subspace orthogonal to the gradients (the discrete Helmholtz
decomposition), removing its gradient component before it is accepted as a mode.
With the gradient null-space projected away, the spurious zero-frequency modes
vanish and the eigensolver's lowest eigenvalues are the genuine physical
resonances.

\begin{figure}[t]
\centering
\begin{tikzpicture}
\begin{axis}[
    name=raw,
    width=0.46\linewidth, height=38mm,
    title={\footnotesize raw eigensolver output},
    title style={font=\footnotesize, yshift=-1mm},
    xlabel={frequency $f$}, ylabel={},
    xmin=-0.3, xmax=5, ymin=0, ymax=1.4,
    axis lines=left, ytick=\empty, xtick={0,2,3.2,4.4},
    xticklabels={$0$,,,}, xticklabel style={font=\scriptsize},
    xlabel style={font=\scriptsize}, clip=false,
  ]
  % spurious cluster near zero (gray) + physical modes
  \addplot[ycomb, very thick, simGray] coordinates {
    (0,1.0)(0.12,0.9)(0.22,0.85)(0.05,0.95)(0.3,0.8)};
  \addplot[ycomb, very thick, simBlue, mark=*, mark size=1.6pt] coordinates {
    (2,1.0)(3.2,1.0)(4.4,1.0)};
  \node[font=\scriptsize\itshape, simGray, align=center] at (axis cs:0.6,1.3)
    {spurious\\(gradient null-space)};
\end{axis}
\begin{axis}[
    name=proj, at={(raw.east)}, anchor=west, xshift=10mm,
    width=0.46\linewidth, height=38mm,
    title={\footnotesize after divergence-free projection},
    title style={font=\footnotesize, yshift=-1mm},
    xlabel={frequency $f$}, ylabel={},
    xmin=-0.3, xmax=5, ymin=0, ymax=1.4,
    axis lines=left, ytick=\empty, xtick={0,2,3.2,4.4},
    xticklabels={$0$,,,}, xticklabel style={font=\scriptsize},
    xlabel style={font=\scriptsize}, clip=false,
  ]
  \addplot[ycomb, very thick, simBlue, mark=*, mark size=1.6pt] coordinates {
    (2,1.0)(3.2,1.0)(4.4,1.0)};
  \node[font=\scriptsize\itshape, simGreen, align=center] at (axis cs:0.9,1.3)
    {spurious removed};
  \node[font=\scriptsize, simBlue] at (axis cs:2,1.12) {$f_1$};
  \node[font=\scriptsize, simBlue] at (axis cs:3.2,1.12) {$f_2$};
  \node[font=\scriptsize, simBlue] at (axis cs:4.4,1.12) {$f_3$};
\end{axis}
\end{tikzpicture}
\caption[Spurious modes removed by the divergence-free projection]{The gradient
null space and its cure. \emph{Left:} the raw eigensolver output, hunting the
lowest eigenvalues, dredges up a cluster of \emph{spurious} zero-frequency modes
(gray)---gradient fields $\Grad\phi$ in the null space of the curl--curl operator
($\Curl\Grad\equiv\mathbf 0$, \cref{ch:reductions-pde})---crowding the genuine
physical resonances (blue). \emph{Right:} the divergence-free projection of
\cref{ch:bc} removes each eigenvector's gradient component, the spurious cluster
vanishes, and the lowest remaining eigenvalues are the true low resonances
$f_1,f_2,f_3$.}
\label{fig:spurious}
\end{figure}

\section{A worked power iteration}
\label{sec:worked-power}

We close the technical development by running power iteration on a small matrix
fully by hand, watching the iterate rotate toward the dominant eigenvector and the
Rayleigh quotient close in on the eigenvalue---the concrete realization of
\cref{sec:power} that the shift-invert machinery then redirects toward the low end.

\begin{workedexample}{Power iteration and the Rayleigh quotient}{power}
Take the symmetric matrix
\[
  \mat{A} = \begin{pmatrix} 2 & 1 \\ 1 & 2 \end{pmatrix},
\]
whose eigenvalues are $\mu_1=3$ (eigenvector $\vu_1=\tfrac{1}{\sqrt2}(1,1)$) and
$\mu_2=1$ (eigenvector $\vu_2=\tfrac{1}{\sqrt2}(1,-1)$). The dominant eigenpair is
$(3,\vu_1)$; the convergence ratio is $\abs{\mu_2/\mu_1}=1/3$. Start from
$\vx_0=(1,0)$---deliberately \emph{not} aligned with either eigenvector---and
iterate $\vx_{k+1}=\mat{A}\vx_k/\norm{\mat{A}\vx_k}$, recording the Rayleigh
quotient $\rho_k=\inner{\mat{A}\vx_k}{\vx_k}/\inner{\vx_k}{\vx_k}$ at each step.

\medskip
\emph{Step 0.} $\vx_0=(1,0)$;\;
$\mat{A}\vx_0=(2,1)$;\;
$\rho_0 = \dfrac{(2,1)\cdot(1,0)}{(1,0)\cdot(1,0)} = \dfrac{2}{1} = 2.000$.

\emph{Step 1.} Normalize: $\vx_1=(2,1)/\sqrt5=(0.894,0.447)$.\;
$\mat{A}\vx_1=(2.236,1.789)$;\;
$\rho_1 = \dfrac{(2.236)(0.894)+(1.789)(0.447)}{1} = 2.800$.

\emph{Step 2.} Normalize $\mat{A}\vx_1$: $\vx_2=(0.781,0.625)$.\;
$\mat{A}\vx_2=(2.187,2.031)$;\;
$\rho_2 = (2.187)(0.781)+(2.031)(0.625) = 2.976$.

\emph{Step 3.} $\vx_3=(0.732,0.681)$;\;
$\mat{A}\vx_3=(2.145,2.094)$;\;
$\rho_3 = (2.145)(0.732)+(2.094)(0.681) = 2.997$.

\emph{Step 4.} $\vx_4=(0.717,0.697)$;\;
$\rho_4 = 2.9997$.

\medskip
The iterate $\vx_k$ marches steadily from $(1,0)$ toward
$\vu_1\approx(0.707,0.707)$---the second component climbing $0\to0.447\to0.625\to
0.681\to0.697$ toward equality with the first---and the Rayleigh quotient closes
in on $\mu_1=3$: $2.000\to2.800\to2.976\to2.997\to2.9997$. The error in $\rho_k$
shrinks by roughly the factor $(\mu_2/\mu_1)^2=1/9$ per step (the \emph{squared}
ratio, the symmetric-case bonus of \cref{eq:rayleigh}): $1.0\to0.20\to0.024\to
0.003\to0.0003$, each about a ninth of the last. \Cref{fig:power-conv} plots the
convergence.

This is power iteration finding the \emph{dominant} eigenvalue $\mu_1=3$. To find
the \emph{smallest} eigenvalue $\mu_2=1$ instead---the analogue of the low cavity
resonance---we would run the same iteration on the shifted-inverted operator
$(\mat{A}-\sigma\mat{I})^{-1}$ with a shift $\sigma$ near $1$: by \cref{eq:eigmap}
the eigenvalue $\mu_2=1$ would become dominant in the transformed problem, and
this same iteration would converge to it instead.
\end{workedexample}

\begin{figure}[t]
\centering
\begin{tikzpicture}
\begin{axis}[
    width=0.66\linewidth, height=50mm,
    xlabel={iteration $k$}, ylabel={Rayleigh quotient $\rho_k$},
    xmin=0, xmax=4, ymin=1.8, ymax=3.15,
    axis lines=left, xtick={0,1,2,3,4}, ytick={2.0,2.5,3.0},
    xticklabel style={font=\scriptsize}, yticklabel style={font=\scriptsize},
    xlabel style={font=\scriptsize}, ylabel style={font=\scriptsize},
    legend style={font=\scriptsize, draw=simGray!50, at={(0.98,0.04)},
      anchor=south east},
    every axis plot/.append style={very thick},
  ]
  % the target eigenvalue
  \addplot[simGray, densely dashed, domain=0:4, samples=2] {3.0};
  \addlegendentry{$\mu_1=3$ (target)}
  % the Rayleigh quotient trajectory
  \addplot[simBlue, mark=*, mark size=2pt] coordinates {
    (0,2.000)(1,2.800)(2,2.976)(3,2.997)(4,2.9997)};
  \addlegendentry{$\rho_k$ (Rayleigh quotient)}
\end{axis}
\end{tikzpicture}
\caption[Power-iteration convergence of the Rayleigh quotient]{The Rayleigh
quotient $\rho_k$ of \cref{wex:power} converging to the dominant eigenvalue
$\mu_1=3$ (gray dashed). Starting from a generic vector, $\rho_k$ climbs
$2.000\to2.800\to2.976\to2.997\to2.9997$, its error falling by roughly the squared
ratio $(\mu_2/\mu_1)^2=1/9$ each step---the second-order accuracy the Rayleigh
quotient buys on a symmetric matrix. The fast, clean convergence here reflects a
well-separated spectrum ($\mu_1/\mu_2=3$); a clustered low spectrum would crawl,
which is exactly why the eigenmode problem needs the spectral transform.}
\label{fig:power-conv}
\end{figure}

\section{The shape of the eigenmode solve}

Step back and assemble the chapter into one picture. The eigenmode analysis,
alone among the six, solves an \emph{eigenvalue} problem rather than a linear
system or an evolution. It wants a few low resonances of a huge pencil
$(\Kop,\Mop)$---the hard, interior, lower end of the spectrum. The naive tool,
power/subspace iteration, finds the wrong (dominant) end and crawls on the
clustered low modes we care about. The spectral transform fixes both at once:
iterating against $(\Kop-\sigma\Mop)^{-1}\Mop$ remaps the eigenvalues so the ones
near the shift become dominant, at the cost of an inner linear solve per step. And
the whole outer iteration is delegated to a specialized library through reverse
communication, so Palace assembles the pencil, hands over the shift and the inner
solver, calls one opaque routine, and reads back the eigenpairs---which it turns
into frequencies, quality factors, and mode fields after projecting away the
spurious gradient null space.

\begin{keyidea}
The eigenmode solve is the deepest composition in the book: a
\emph{library-owned} outer eigenvalue iteration (\cref{ch:krylovschur}), calling
back through the \code{apply} interface (\cref{ch:operators}) into a Palace-owned
shift-invert body (apply $\Mop$, inner solve, untransform), whose inner solve is
itself a Palace-owned \code{ksp\_solve} (\cref{ch:krylov}) calling a
preconditioner (\cref{ch:multigrid}). Three nested ``a solver is an operator''
rotations, the outermost not even Palace's code. The spectral transform is the
hinge: it converts the hard interior eigenproblem into the easy dominant one the
iteration can solve, and pays for the conversion in inner linear solves.
\end{keyidea}

\summarysection
The eigenmode analysis seeks a few eigenpairs $(\lambda,\vx)$ of the generalized
eigenproblem $\Kop\vx=\lambda\Mop\vx$ ($\lambda=\omega^2$), with $\Kop$ the
curl--curl stiffness and $\Mop$ the mass---specifically the \emph{smallest}
nonzero $\lambda$, the lowest resonances at the hard interior/lower end of a vast
spectrum. \emph{Power iteration} $\vx\leftarrow\mat{A}\vx/\norm{\cdot}$ converges
to the \emph{dominant} (largest-magnitude) eigenvector at rate $\abs{\mu_2/\mu_1}$,
with the \emph{Rayleigh quotient} $\rho(\vx)=\inner{\mat{A}\vx}{\vx}/
\inner{\vx}{\vx}$ reading off the eigenvalue; subspace iteration finds several at
once. But this finds the wrong end and crawls on clustered spectra. The
\emph{spectral transform} (shift-and-invert) fixes both: iterating against
$\mat{T}_\sigma=(\Kop-\sigma\Mop)^{-1}\Mop$ maps each eigenvalue
$\lambda\mapsto\theta=1/(\lambda-\sigma)$, so the eigenvalues \emph{nearest the
shift} $\sigma$ become \emph{dominant}---turning ``find eigenvalues near $\sigma$''
(hard, interior) into ``find the dominant eigenvalues of $\mat{T}_\sigma$'' (easy).
The cost is an \emph{inner linear solve} of $(\Kop-\sigma\Mop)$ per apply---a
\code{ksp\_solve}, usually preconditioned---so the body is apply-$\Mop$
$\rhd$ inner-solve $\rhd$ untransform. In practice the \emph{outer} eigenvalue
iteration (Krylov--Schur/Arnoldi) lives inside a specialized library (SLEPc,
ARPACK) that drives the loop and \emph{calls back} for the apply (reverse
communication): the eigenmode driver assembles $(\Kop,\Mop)$, hands them plus the
shift to the library, and reads back converged eigenpairs in one opaque call---the
fourth, black-box solve corner. The lossy \emph{quadratic} eigenproblem
$(\Kop+\lambda\Cop+\lambda^2\Mop)\vx=\mathbf 0$ is \emph{linearized} to a doubled
linear pencil and solved the same way. Finally, eigenvalues become frequencies
$f=\sqrt\lambda/2\pi$ (and quality factors $Q$ from the imaginary part) and
eigenvectors become mode fields, after a \emph{divergence-free projection} removes
the spurious gradient null-space modes clustered at zero frequency.

\furtherreading
The definitive modern treatment of large-scale eigenvalue computation---power and
subspace iteration, the spectral transform, and the projection methods of the next
chapter---is Saad's \emph{Numerical Methods for Large Eigenvalue
Problems}~\citep{saad2011eig}; its chapters on the shift-and-invert Arnoldi method
are the direct background for \cref{sec:shift-invert,sec:blackbox}. The
Krylov--Schur restarting that the library uses to drive the outer iteration is due
to Stewart~\citep{stewart2001krylovschur}, the subject of \cref{ch:krylovschur}.
For the underlying linear-algebra fundamentals---the Rayleigh quotient, its
second-order accuracy, and the geometry of convergence---Trefethen and
Bau~\citep{trefethenbau1997} remain the clearest undergraduate-accessible source.
The reverse-communication black-box structure of \cref{sec:blackbox} is the
ARPACK/SLEPc design; the constructive realization of the loop those libraries hide
is developed, in our own vocabulary, in \cref{ch:krylovschur}, and the nonlinear
(deflation) generalization in \cref{ch:deflation}.

\exercisessection
\begin{enumerate}[nosep]
  \item \textbf{Which end, and why.} The eigenmode analysis wants the
  \emph{smallest} nonzero $\lambda$, yet power iteration finds the
  \emph{largest}. Using the eigenbasis expansion \cref{eq:power-expand}, explain
  in one sentence why iteration amplifies the dominant eigenvector. Then state the
  convergence ratio for the dominant pair and explain why a \emph{clustered} low
  spectrum (the cavity's near-degenerate overtones) makes \emph{direct} iteration
  for those modes especially slow.
  \item \textbf{The eigenvalue map by hand.} A pencil has eigenvalues
  $\lambda\in\{2,5,11,20\}$ and you want the two nearest $\lambda=6$. (i) Choose a
  shift $\sigma=6$ and compute the four transformed eigenvalues
  $\theta_i=1/(\lambda_i-\sigma)$. (ii) Rank them by magnitude and confirm the two
  you wanted are now dominant. (iii) The iteration returns $\theta=-1.0$; recover
  the physical $\lambda$. (iv) Why does the eigenvector need no untransforming?
  \item \textbf{Locating the cost.} Write out the three moves of one application of
  $\mat{T}_\sigma=(\Kop-\sigma\Mop)^{-1}\Mop$ to a vector $\vv$
  (\cref{eq:siapply}). Which move dominates the cost, and which chapter's machinery
  performs it? Explain why the shifted matrix $\Kop-\sigma\Mop$ is solved
  \emph{iteratively} rather than factored once, given that the shift $\sigma$ is
  held fixed for the whole eigensolve. (Hint: contrast with the fixed-operator
  electrostatic family of \cref{ch:reductions-pde}---when \emph{is} a single
  factorization the right call?)
  \item \textbf{Rayleigh quotient accuracy.} Run two more steps of the power
  iteration of \cref{wex:power} starting from $\vx_4=(0.717,0.697)$, computing
  $\rho_5$ and $\rho_6$. Verify the error in $\rho_k$ shrinks by roughly $1/9$ per
  step, and explain where the factor $1/9$ comes from. What would the per-step
  error-reduction factor be if instead $\mu_1=3$, $\mu_2=2.9$ (a near-cluster)?
  \item \textbf{The black box, drawn.} In your own words, describe \emph{reverse
  communication}: what does the library own, what does Palace own, and what is the
  single interface across which they cooperate? Why can the library not simply hold
  Palace's operators directly (\cref{ch:matrixfree})? Contrast the eigenmode solve
  corner with the electrostatic one: in which does Palace write the outer loop, and
  in which does it write only the apply?
  \item \textbf{Linearize the quadratic.} Verify by direct multiplication that the
  doubled pencil \cref{eq:linearization} is equivalent to the quadratic
  eigenproblem \cref{eq:qevp}: expand both block rows and recover
  $(\Kop+\lambda\Cop+\lambda^2\Mop)\vx=\mathbf 0$. What is the size of the doubled
  problem relative to the original, and what is discarded from each computed
  eigenvector when reading back the physical mode?
  \item \textbf{Spurious modes.} Explain why the curl--curl stiffness $\Kop$ has a
  large null space (cite the exactness identity from \cref{ch:reductions-pde}) and
  why every null-space vector is an eigenvalue-zero eigenvector of the pencil. Why
  do these spurious modes appear precisely where the eigenmode solver is looking
  (the low end)? What operation (\cref{ch:bc}) removes them, and what subspace does
  it project onto?
  \item \textbf{Reading the results.} A lossy eigensolve returns the complex
  eigenvalue $\lambda = i\,\omega$ with $\lambda = (-0.02 + 6.28\,i)\times10^{9}$.
  (i) What resonant frequency $f$ (in Hz) does this mode have? (ii) Which part of
  $\lambda$ carries the decay, and does a \emph{smaller} real part mean a higher or
  lower quality factor $Q$? (iii) Why does the lossless linear problem
  ($\lambda=\omega^2$ real) produce no $Q$ at all?
  \item \textbf{Choosing the shift.} You want the lowest five resonances of a
  cavity, and you know the fundamental is near $f_0$. (i) What shift $\sigma$ (in
  terms of $\lambda=\omega^2=(2\pi f)^2$) would you pass to the eigensolver, and
  why place it slightly below the lowest wanted eigenvalue rather than at zero?
  (ii) What goes wrong if you set $\sigma$ exactly equal to an eigenvalue? (iii) If
  the five wanted modes span a wide frequency range, why might a single shift
  converge slowly on the highest of them, and what does that suggest about the
  trade-off in choosing $\sigma$?
\end{enumerate}

\chapter{Lanczos, Arnoldi, and Krylov--Schur}
\label{ch:krylovschur}

\chapterorient{The previous chapter posed the eigenproblem and handed its solve to
an opaque library: ``assemble the operators once and call \code{EPSSolve}.'' This
chapter opens that box. Inside is one of the most elegant constructions in numerical
linear algebra---a way to find a handful of a giant matrix's eigenvalues without ever
forming, or even looking at, the whole of it. The idea has three moving parts, and we
take them in order. First, \emph{Rayleigh--Ritz}: project the operator onto a small
Krylov subspace and read approximate eigenvalues off the tiny projected matrix; the
extreme eigenvalues come out first and best. Second, the \emph{restart problem}: the
subspace must stay small, but a naive restart throws away the good approximations you
worked to find. Third, \emph{thick-restart Krylov--Schur}: the elegant fix that keeps
the directions worth keeping and discards the rest, compressing the basis while
preserving its Krylov structure. We assemble these on top of the Arnoldi and Lanczos
recurrences of \cref{ch:orthog}, fit the inner shift-invert solve of
\cref{ch:eigenvalues} inside each step, and arrive at the full eigen-solve
stack---the exact loop that SLEPc and ARPACK run on Palace's behalf. The black box of
\cref{ch:eigenvalues} turns out to be code we could have written ourselves.}

% local helper: a Ritz value symbol used throughout (theta), already covered by
% standard math; no new macros needed. The basis matrix V_m uses \mat{V}.

\section{The setting: too big to factor, too many to want}
\label{sec:kssetting}

The eigenproblem of \cref{ch:eigenvalues} is the generalized system
\[
  \mat{K}\vx \;=\; \lambda\,\mat{M}\vx,
\]
with $\mat{K}$ and $\mat{M}$ the stiffness and mass operators assembled from the mesh.
Two facts about this problem shape everything that follows, and both are facts of
\emph{scale}. First, $\mat{K}$ and $\mat{M}$ are enormous---a fine mesh produces
millions of unknowns---and sparse, so we will never form a dense factorization of
either, let alone of the pair. We are allowed only to \emph{apply} the operators to
vectors (\cref{ch:matrixfree}) and to solve linear systems with them
(\cref{ch:multigrid}). Second, we do not \emph{want} all the eigenvalues. A cavity has
millions of resonant modes, but the engineer cares about the lowest dozen---the
fundamental and its first few overtones. We want a handful of eigenvalues near a target,
not the whole spectrum.

These two facts rule out the textbook eigenvalue algorithms. The QR
algorithm---the workhorse for small dense matrices---computes \emph{all} the
eigenvalues and needs the matrix in dense form; it is hopeless here on both counts. What
we need is a method that touches $\mat{K}$ and $\mat{M}$ only through matrix-vector
products and inner linear solves, and that spends its effort on the few eigenvalues we
asked for. That method is \emph{Krylov subspace projection}, and its practical form is
the subject of this chapter.

\begin{keyidea}
The eigenproblems of this book are large, sparse, and selective: millions of unknowns,
access only by operator-apply and linear-solve, but only a handful of wanted
eigenvalues near a target. This rules out dense methods that compute the whole spectrum.
The remedy is \emph{Rayleigh--Ritz projection onto a small Krylov subspace}: a small
projected eigenproblem yields the extremal eigenpairs, and \emph{thick-restart} keeps
the good ones while holding the subspace small. The entire algorithm uses $\mat{K}$ and
$\mat{M}$ only through the operator-apply and linear-solve we already have.
\end{keyidea}

The plan of attack inverts the difficulty exactly. We cannot work with the big operator,
so we project it onto a small subspace where we \emph{can} work with it---small enough
to hand to the dense QR algorithm. We cannot afford a huge subspace, so we keep it small
by restarting. And we want the eigenvalues near a target, not the extremes, so we feed
the iteration the shift-inverted operator of \cref{ch:eigenvalues}, which turns ``near
the target'' into ``extreme.'' Three subspace tricks, layered: projection, restart,
spectral transform. We build them in that order.

\section{Rayleigh--Ritz: a small projected eigenproblem}
\label{sec:rayleighritz}

Start with the first trick and the simplest version of it: no restart, no spectral
transform, just projection. Suppose we have built---by Arnoldi or Lanczos
(\cref{ch:orthog})---an orthonormal basis $\vq_1, \dots, \vq_m$ for the Krylov subspace
$\mathcal{K}_m(\mat{A}, \vv)$, stacked into the tall thin matrix
$\mat{V}_m = [\,\vq_1 \mid \cdots \mid \vq_m\,]$ with $\mat{V}_m^{\!\top}\mat{V}_m =
\mat{I}_m$. (For now read $\mat{A}$ as the operator whose eigenvalues we want; in the
eigenmode driver it will be the shift-inverted operator of \cref{sec:shiftinvert}.) The
question Rayleigh--Ritz answers is: \emph{what is the best approximation to $\mat{A}$'s
eigenpairs that we can extract from this subspace alone}?

\subsection{The projected operator}

We cannot diagonalize $\mat{A}$---it is too big. But we \emph{can} ask how $\mat{A}$
acts \emph{within} the subspace $\mathcal{K}_m$, by projecting its action down onto the
basis. The projected operator is the $m \times m$ matrix
\[
  \mat{T}_m \;=\; \mat{V}_m^{\!\top}\,\mat{A}\,\mat{V}_m,
\]
the restriction of $\mat{A}$ to the subspace, expressed in the orthonormal basis. Its
entries are $(\mat{T}_m)_{ij} = \inner{\mat{A}\vq_j}{\vq_i}$---exactly the
orthogonalization coefficients the Arnoldi recurrence already computed and stored
(\cref{ch:orthog}). So $\mat{T}_m$ \emph{is the Hessenberg matrix} $\bar{\mat{H}}_m$
(its leading $m\times m$ block), produced as a free by-product of building the basis. For
symmetric $\mat{A}$ it is the \emph{tridiagonal} matrix of Lanczos. Either way, $\mat{T}_m$
is small---$m$ is a few dozen, not a few million---and dense, so we hand it to the QR
algorithm and get \emph{all} of its eigenpairs cheaply.

\begin{figure}[t]
\centering
\begin{tikzpicture}[font=\footnotesize, scale=1.0]
  % ===== big operator A =====
  \draw[thick, fill=simBgGray] (-0.2,-1.6) rectangle (2.0,1.6);
  \node[font=\small\bfseries, text=simInk] at (0.9,1.95) {$\mat{A}$};
  \node[font=\scriptsize, text=simGray] at (0.9,0.2) {$N\times N$};
  \node[font=\scriptsize, text=simGray] at (0.9,-0.2) {sparse, huge};
  % sparse dots
  \foreach \p in {(0.1,1.2),(0.6,0.9),(1.3,1.4),(1.7,0.5),(0.3,-0.5),(1.0,-1.1),(1.6,-0.8),(0.5,-1.3),(1.4,-0.1)}{
    \fill[simBlue!55] \p circle (1.2pt);}
  % projection arrow
  \draw[flow, simViolet, very thick] (2.3,0) -- (3.6,0)
    node[midway, above, font=\scriptsize, text=simViolet, align=center]
      {$\mat{V}_m^{\!\top}(\cdot)\mat{V}_m$\\project onto $\mathcal{K}_m$};
  % ===== small projected T_m =====
  \draw[thick, fill=simBgBlue] (3.9,-0.7) rectangle (5.3,0.7);
  \node[font=\small\bfseries, text=simBlue] at (4.6,1.0) {$\mat{T}_m$};
  \node[font=\scriptsize, text=simGray] at (4.6,-1.05) {$m\times m$, dense};
  % tridiagonal hint
  \foreach \r in {0,1,2}{\foreach \c in {0,1,2}{
    \pgfmathtruncatemacro\d{abs(\r-\c)}
    \ifnum\d<2 \fill[simBlue!60] (4.0+\c*0.42,0.5-\r*0.42) circle (1.6pt)\fi;}}
  % QR arrow
  \draw[flow, simRust, very thick] (5.6,0) -- (6.9,0)
    node[midway, above, font=\scriptsize, text=simRust, align=center]
      {dense QR\\(cheap)};
  % ===== Ritz pairs =====
  \node[product, align=left, text width=30mm] at (8.5,0)
    {Ritz values\\ $\theta_1,\dots,\theta_m$\\[2pt] Ritz vectors\\ $\mat{V}_m\vs_i$};
  \node[font=\scriptsize\itshape, text=simGreen, align=center, text width=42mm]
    at (5.0,-2.25)
    {the extreme $\theta_i$ approximate $\mat{A}$'s extreme eigenvalues---first and best};
\end{tikzpicture}
\caption[Rayleigh--Ritz projection]{The Rayleigh--Ritz idea in one picture. The huge
sparse operator $\mat{A}$ (left) is projected onto the small Krylov subspace
$\mathcal{K}_m$ spanned by the orthonormal basis $\mat{V}_m$, producing the tiny
\emph{projected operator} $\mat{T}_m = \mat{V}_m^{\!\top}\mat{A}\mat{V}_m$
(centre)---which is exactly the Hessenberg/tridiagonal matrix the Arnoldi/Lanczos
recurrence already filled (\cref{ch:orthog}). Because $\mat{T}_m$ is small and dense, the
classical dense QR algorithm extracts \emph{all} its eigenpairs $(\theta_i, \vs_i)$
cheaply. The \emph{Ritz values} $\theta_i$ and \emph{Ritz vectors} $\mat{V}_m\vs_i$
(right) are the subspace's best approximations to $\mat{A}$'s eigenpairs---and crucially,
the \emph{extreme} ones converge first and most accurately. We never diagonalize
$\mat{A}$; we diagonalize its shadow on a cleverly chosen small subspace.}
\label{fig:rayleighritz}
\end{figure}

\subsection{Ritz values and Ritz vectors}

Let $(\theta_i, \vs_i)$ be an eigenpair of the small matrix $\mat{T}_m$, so
$\mat{T}_m\vs_i = \theta_i\vs_i$ with $\vs_i \in \Real^m$. The scalar $\theta_i$ is a
\emph{Ritz value}, and the vector $\mat{V}_m\vs_i \in \Real^N$---the small eigenvector
$\vs_i$ lifted back into the full space through the basis---is the corresponding
\emph{Ritz vector}. \Cref{fig:rayleighritz} traces the whole pipeline.

\begin{definition}[Ritz pairs]
\label{def:ritz}
Given an orthonormal basis $\mat{V}_m$ of a subspace and the projected operator
$\mat{T}_m = \mat{V}_m^{\!\top}\mat{A}\mat{V}_m$, let $\mat{T}_m\vs_i = \theta_i\vs_i$.
The pair $(\theta_i,\, \mat{V}_m\vs_i)$ is a \emph{Ritz pair}: $\theta_i$ is the
\emph{Ritz value} and $\vx_i := \mat{V}_m\vs_i$ the \emph{Ritz vector}. Ritz pairs are
the subspace's approximations to the eigenpairs of $\mat{A}$; in exact arithmetic, when
$\mathcal{K}_m$ contains a true eigenvector, the corresponding Ritz pair is that
eigenpair exactly.
\end{definition}

The claim that earns the construction its keep is this: the Ritz values approximate the
eigenvalues of $\mat{A}$, and they approximate the \emph{extremal} ones---the largest
and smallest---first and best. Why the extremes? The reason is the same effect that made
the power basis collapse toward the dominant eigenvector (\cref{ch:orthog}), now turned
to advantage. The Krylov subspace $\mathcal{K}_m(\mat{A}, \vv) =
\operatorname{span}\{\vv, \mat{A}\vv, \dots, \mat{A}^{m-1}\vv\}$ is exactly the set of
vectors $p(\mat{A})\vv$ for polynomials $p$ of degree $< m$. Expand the starting vector
in $\mat{A}$'s eigenbasis, $\vv = \sum_k c_k\vu_k$; then $p(\mat{A})\vv = \sum_k
c_k\,p(\lambda_k)\vu_k$. The subspace can therefore reach any combination of eigenvectors
that some degree-$(m{-}1)$ polynomial, evaluated at the eigenvalues, can produce. And a
polynomial of modest degree can take a large value at one extreme eigenvalue while
staying small at all the interior ones---the interior eigenvalues are crowded together,
the extremes stick out. So the subspace captures the extreme eigenvectors with a
low-degree polynomial (small $m$), while the interior ones need a high-degree polynomial
that separates neighbours (large $m$). This is the polynomial-approximation view, and it
is why the extremes come first.

\begin{keyidea}
The Krylov subspace $\mathcal{K}_m(\mat{A},\vv)$ is the span of $\{p(\mat{A})\vv : \deg
p < m\}$. A low-degree polynomial can pick out an \emph{extreme} eigenvalue---one that
sticks out from the rest---while a high-degree one is needed to resolve \emph{interior}
eigenvalues that are crowded together. So Rayleigh--Ritz on a small Krylov subspace
captures the extremal eigenpairs first and most accurately, and reaches the interior
ones only slowly. This is precisely why the eigenmode driver uses the \emph{shift-invert}
spectral transform of \cref{ch:eigenvalues}: it remaps the wanted interior eigenvalues
to the extremes, where Krylov projection is fast (\cref{sec:shiftinvert}).
\end{keyidea}

There is a cleaner way to see ``best approximation'' that the polynomial picture only
gestures at. Among all scalars $\theta$ and all unit vectors $\vx$ drawn from the
subspace, the Ritz values are the stationary values of the \emph{Rayleigh quotient}
$\rho(\vx) = \inner{\mat{A}\vx}{\vx} / \inner{\vx}{\vx}$ restricted to $\mathcal{K}_m$.
For symmetric $\mat{A}$ the Rayleigh quotient's global maximum and minimum over \emph{all
of $\Real^N$} are the largest and smallest eigenvalues (the Courant--Fischer min-max
theorem); restricting the search to the subspace $\mathcal{K}_m$ gives the best bound the
subspace can offer, and the extreme Ritz values are exactly that best bound. The largest
Ritz value is a \emph{lower} bound on the largest eigenvalue, the smallest Ritz value an
\emph{upper} bound on the smallest---the Ritz values bracket the spectrum from inside and
tighten as the subspace grows.

\subsection{The cheap convergence test}

A Ritz pair is an \emph{approximate} eigenpair. We need to know how good it is---and the
beautiful fact, the one that makes the whole method practical, is that we can measure its
quality \emph{without ever forming the Ritz vector or applying $\mat{A}$ to it}. The cost
of the convergence test is essentially zero.

The natural quality measure is the eigen-residual: if $(\theta_i, \vx_i)$ with
$\vx_i = \mat{V}_m\vs_i$ were a true eigenpair, then $\mat{A}\vx_i - \theta_i\vx_i$ would
vanish. So the residual norm $\norm{\mat{A}\vx_i - \theta_i\vx_i}$ measures how far off it
is. Forming this directly would cost a matrix-vector product with $\mat{A}$ on the full
vector $\vx_i$---not free. But the Arnoldi relation $\mat{A}\mat{V}_m = \mat{V}_{m+1}
\bar{\mat{H}}_m$ of \cref{ch:orthog} collapses it to something we can read off the small
data we already have. Substitute $\vx_i = \mat{V}_m\vs_i$ and use the relation:
\[
  \mat{A}\vx_i - \theta_i\vx_i
    = \mat{A}\mat{V}_m\vs_i - \theta_i\mat{V}_m\vs_i
    = \mat{V}_{m+1}\bar{\mat{H}}_m\vs_i - \theta_i\mat{V}_m\vs_i.
\]
Split $\bar{\mat{H}}_m$ into its top $m\times m$ block $\mat{T}_m$ and its single
sub-diagonal spill-over row, whose only nonzero is $h_{m+1,m}$ in the last column. Since
$\mat{T}_m\vs_i = \theta_i\vs_i$, the $\mat{T}_m$ part cancels the $\theta_i\mat{V}_m\vs_i$
term exactly, leaving only the spill-over:
\[
  \mat{A}\vx_i - \theta_i\vx_i \;=\; h_{m+1,m}\,(\ve_m^{\!\top}\vs_i)\,\vq_{m+1},
\]
where $\ve_m^{\!\top}\vs_i$ is the \emph{last component} of the small eigenvector $\vs_i$
and $\vq_{m+1}$ is the next (unit) basis vector. Taking norms, and using
$\norm{\vq_{m+1}} = 1$:
\[
  \boxed{\;\norm{\mat{A}\vx_i - \theta_i\vx_i}
    \;=\; \abs{h_{m+1,m}}\;\abs{\ve_m^{\!\top}\vs_i}\;}
\]
The residual norm of a full-size Ritz pair is the product of \emph{two scalars}: the
sub-diagonal entry $h_{m+1,m}$ (a single number, already computed when the basis was
extended) and the last component of the tiny eigenvector $\vs_i$ (read directly from the
QR solve of $\mat{T}_m$). No full-vector operation, no application of $\mat{A}$. The test
costs nothing.

\begin{figure}[t]
\centering
\begin{tikzpicture}[font=\footnotesize, scale=1.0]
  % H-bar block: T_m on top, spill-over row at bottom
  \draw[thick, fill=simBgBlue] (0,0) rectangle (2.4,2.4);
  \node[font=\scriptsize, text=simBlue] at (1.2,1.5) {$\mat{T}_m$};
  \node[font=\scriptsize, text=simGray] at (1.2,2.05) {(top $m\times m$ block)};
  % spill-over row
  \draw[thick, fill=simBgRust] (0,-0.5) rectangle (2.4,0.0);
  \node[font=\scriptsize, text=simRust] at (1.95,-0.25) {$h_{m+1,m}$};
  \node[font=\scriptsize, text=simGray] at (0.7,-0.25) {$0\;\cdots\;0$};
  \node[font=\small, text=simInk] at (1.2,-0.95) {$\bar{\mat{H}}_m$};
  % eigenvector s_i
  \draw[thick, fill=simBgGreen] (3.1,0) rectangle (3.55,2.4);
  \node[font=\scriptsize, text=simGreen] at (3.32,1.5) {$\vs_i$};
  \draw[thick, fill=simGreen!40] (3.1,0) rectangle (3.55,0.4);
  \node[font=\scriptsize, text=simGreen, anchor=west] at (3.7,0.2)
    {last component $\ve_m^{\!\top}\vs_i$};
  % product = residual
  \node[opbox, align=center, anchor=west] at (5.6,1.7)
    {$\norm{\mat{A}\vx_i - \theta_i\vx_i}$\\$=\;\abs{h_{m+1,m}}\,\abs{\ve_m^{\!\top}\vs_i}$};
  \node[font=\scriptsize\itshape, text=simGray, align=left, anchor=west] at (5.6,0.4)
    {a product of \emph{two scalars}:\\both already computed.\\no full-vector work,\\
     no apply of $\mat{A}$.};
  \draw[refedge, simRust] (2.4,-0.25) .. controls (4.0,-0.25) and (5.0,1.0) .. (5.55,1.55);
  \draw[refedge, simGreen] (3.55,0.2) .. controls (4.6,0.2) and (5.0,1.2) .. (5.55,1.45);
\end{tikzpicture}
\caption[The cheap Ritz residual test]{Why the convergence test is free. The eigen-residual
of a Ritz pair $(\theta_i, \mat{V}_m\vs_i)$ would seem to need a full matrix-vector product
with $\mat{A}$, but the Arnoldi relation collapses it to a product of two already-computed
scalars. The Hessenberg matrix $\bar{\mat{H}}_m$ splits into its top block $\mat{T}_m$
(blue, whose eigenpair is $(\theta_i,\vs_i)$) and a single spill-over row whose only
nonzero is the sub-diagonal entry $h_{m+1,m}$ (rust). Because $\mat{T}_m\vs_i =
\theta_i\vs_i$ cancels the projected part exactly, the residual reduces to $h_{m+1,m}$
times the \emph{last component} $\ve_m^{\!\top}\vs_i$ of the small eigenvector (green)---two
numbers we already hold. The test $\abs{h_{m+1,m}}\abs{\ve_m^{\!\top}\vs_i} < \texttt{tol}$
is what the eigensolver checks every step to decide which Ritz pairs have converged.}
\label{fig:ritzresidual}
\end{figure}

\begin{notationbox}
A Ritz pair has \emph{converged} when its cheap residual falls below the requested
tolerance:
\[
  \abs{h_{m+1,m}}\;\abs{\ve_m^{\!\top}\vs_i} \;<\; \texttt{tol}\cdot\abs{\theta_i}
  \quad\text{(relative)}.
\]
Two readings of this formula are worth keeping. First, a Ritz pair converges when the
\emph{last component} $\ve_m^{\!\top}\vs_i$ of its small eigenvector is tiny---that is,
when the eigenvector lives almost entirely in the earlier basis vectors and barely
touches the newest one. Second, the shared factor $h_{m+1,m}$ multiplies \emph{every}
pair's residual: when the sub-diagonal $h_{m+1,m}$ shrinks toward zero, the basis is
approaching an \emph{invariant subspace} and \emph{all} the Ritz pairs converge at once.
A small $h_{m+1,m}$ is the eigensolver's signal that the subspace has locked onto a true
invariant subspace of $\mat{A}$.
\end{notationbox}

\subsection{A worked Rayleigh--Ritz}

The whole construction fits on one page of arithmetic. We run Rayleigh--Ritz on a small
symmetric matrix, watch the Ritz values bracket the spectrum, and apply the cheap test.

\begin{workedexample}{Rayleigh--Ritz on a $4\times4$ symmetric matrix}{rrritz}
Take the symmetric matrix and unit starting vector
\[
  \mat{A} = \begin{psmallmatrix}
    2 & 1 & 0 & 0\\ 1 & 2 & 1 & 0\\ 0 & 1 & 2 & 1\\ 0 & 0 & 1 & 2
  \end{psmallmatrix},
  \qquad
  \vq_1 = \begin{psmallmatrix} 1\\0\\0\\0 \end{psmallmatrix}.
\]
The true eigenvalues of $\mat{A}$ are $2 + 2\cos(k\pi/5)$ for $k=1,2,3,4$, namely
\[
  \lambda \;\approx\; \{\,0.382,\; 1.382,\; 2.618,\; 3.618\,\},
\]
with extremes $0.382$ and $3.618$. We will build a $3$-dimensional Krylov basis by
Lanczos, form $\mat{T}_3$, and read its Ritz values---then check how well they bracket the
extremes.

\smallskip
\emph{Build the basis (Lanczos, \cref{ch:orthog}).} Step~1:
$\mat{A}\vq_1 = (2,1,0,0)^{\!\top}$, so $\alpha_1 = \inner{\mat{A}\vq_1}{\vq_1} = 2$ and the
residual is $\mat{A}\vq_1 - 2\vq_1 = (0,1,0,0)^{\!\top}$, giving $\beta_1 = 1$ and
$\vq_2 = (0,1,0,0)^{\!\top}$. Step~2: $\mat{A}\vq_2 = (1,2,1,0)^{\!\top}$, so $\alpha_2 = 2$
and the three-term residual is $\mat{A}\vq_2 - 2\vq_2 - 1\cdot\vq_1 = (0,0,1,0)^{\!\top}$,
giving $\beta_2 = 1$ and $\vq_3 = (0,0,1,0)^{\!\top}$. The projected operator is the
tridiagonal
\[
  \mat{T}_3 \;=\; \begin{psmallmatrix}
    \alpha_1 & \beta_1 & 0\\ \beta_1 & \alpha_2 & \beta_2\\ 0 & \beta_2 & \alpha_3
  \end{psmallmatrix}
  \;=\; \begin{psmallmatrix} 2 & 1 & 0\\ 1 & 2 & 1\\ 0 & 1 & 2 \end{psmallmatrix},
\]
where $\alpha_3 = \inner{\mat{A}\vq_3}{\vq_3} = 2$ completes the diagonal. (We also pick up
$\beta_3 = \norm{\mat{A}\vq_3 - 2\vq_3 - \vq_2}$; one computes $\mat{A}\vq_3 = (0,1,2,1)^{\!\top}$
and the residual $(0,0,0,1)^{\!\top}$, so $h_{4,3} = \beta_3 = 1$---the spill-over entry we
need for the residual test.)

\smallskip
\emph{Compute the Ritz values.} The eigenvalues of the $3\times3$ $\mat{T}_3$ are
$2$ and $2 \pm \sqrt2$, i.e.
\[
  \theta \;=\; \{\,0.586,\; 2.000,\; 3.414\,\}.
\]
Compare to the true spectrum. The extreme Ritz values $0.586$ and $3.414$ bracket the
true extremes $0.382$ and $3.618$ \emph{from inside}: $0.586 > 0.382$ and $3.414 < 3.618$,
exactly as the min-max bound promises (the smallest Ritz value over-estimates the smallest
eigenvalue; the largest under-estimates the largest). With only a $3$-dimensional subspace
of a $4$-dimensional space, the extremes are already approximated to within about $0.2$,
while the interior pair is represented by the single middling Ritz value $2.000$.

\smallskip
\emph{Apply the cheap test.} Take the largest Ritz value $\theta = 3.414$ with eigenvector
$\vs = \tfrac12(1, \sqrt2, 1)^{\!\top}$ of $\mat{T}_3$ (the standard eigenvector of
$\mathrm{tridiag}(1;2;1)$). Its last component is $\ve_3^{\!\top}\vs = \tfrac12$. The
residual of the corresponding full Ritz pair is then, with no further matrix work,
\[
  \norm{\mat{A}\vx - \theta\vx} \;=\; \abs{h_{4,3}}\,\abs{\ve_3^{\!\top}\vs}
    \;=\; 1 \cdot \tfrac12 \;=\; 0.5.
\]
This is not yet converged---a residual of $0.5$ is large---which is consistent with the
Ritz value still being off by $0.2$ from the true eigenvalue. One more Lanczos step (here
filling the subspace to all of $\Real^4$) would drive $h_{m+1,m}$ to zero and the residual
with it: at $m = 4 = N$ the subspace is invariant, $h_{5,4} = 0$, and the Ritz values
\emph{equal} the eigenvalues exactly. The cheap test tracked the convergence the whole way
without one extra application of $\mat{A}$.
\end{workedexample}

\Cref{fig:ritzconverge} shows the same story for a larger matrix, where the subspace is
genuinely smaller than the space: the extreme Ritz values march toward the true extremes as
the basis grows, fast at the ends and slow in the middle.

\begin{figure}[t]
\centering
\begin{tikzpicture}
\begin{axis}[
  width=12cm, height=7cm,
  xlabel={basis size $m$},
  ylabel={Ritz values $\theta_i$ \ (vs.\ true eigenvalues, dashed)},
  xmin=1, xmax=12, ymin=-0.3, ymax=8.3,
  grid=both, grid style={simGray!18},
  legend pos=south east, legend cell align=left,
  legend style={font=\scriptsize, fill=white, draw=simGray!40},
  label style={font=\footnotesize}, tick label style={font=\scriptsize},
]
% true extreme eigenvalues (dashed reference lines)
\addplot[simGray, densely dashed, thick, domain=1:12, samples=2] {7.7};
\addplot[simGray, densely dashed, thick, domain=1:12, samples=2] {0.15};
\node[font=\scriptsize, text=simGray, anchor=west] at (axis cs:9.2,7.95) {$\lambda_{\max}$};
\node[font=\scriptsize, text=simGray, anchor=west] at (axis cs:9.2,0.5) {$\lambda_{\min}$};
% largest Ritz value: converges UP to lambda_max quickly
\addplot[simRust, very thick, mark=*, mark size=1.3pt] coordinates {
  (1,4.0)(2,5.6)(3,6.6)(4,7.15)(5,7.45)(6,7.6)(7,7.66)(8,7.69)(9,7.70)(10,7.70)(11,7.70)(12,7.70)};
\addlegendentry{largest Ritz value $\theta_{\max}$}
% smallest Ritz value: converges DOWN to lambda_min quickly
\addplot[simBlue, very thick, mark=square*, mark size=1.3pt] coordinates {
  (1,4.0)(2,2.3)(3,1.3)(4,0.75)(5,0.45)(6,0.30)(7,0.22)(8,0.18)(9,0.16)(10,0.15)(11,0.15)(12,0.15)};
\addlegendentry{smallest Ritz value $\theta_{\min}$}
% an interior Ritz value: converges slowly, wobbles
\addplot[simGreen, very thick, mark=triangle*, mark size=1.5pt] coordinates {
  (3,3.1)(4,3.6)(5,3.9)(6,4.0)(7,3.95)(8,3.92)(9,3.9)(10,3.9)(11,3.9)(12,3.9)};
\addlegendentry{an interior Ritz value (slow)}
\addplot[simGray, densely dotted, thick, domain=1:12, samples=2] {3.9};
\end{axis}
\end{tikzpicture}
\caption[Ritz values converging to the extremes]{How the Ritz values approach the true
eigenvalues as the Krylov basis grows---a schematic of the universal behaviour. The
\emph{extreme} eigenvalues are captured fast: the largest Ritz value (rust) climbs to
$\lambda_{\max}$ and the smallest (blue) descends to $\lambda_{\min}$ within a handful of
steps, each approaching its target monotonically from inside the spectrum (the min-max
bracketing of \cref{sec:rayleighritz}). An \emph{interior} eigenvalue (green) is resolved
only slowly and non-monotonically, because the subspace must build up a high-degree
polynomial to separate it from its crowded neighbours. This asymmetry---extremes fast,
interior slow---is the whole reason the eigenmode driver applies the shift-invert
transform of \cref{ch:eigenvalues} before iterating: it moves the wanted interior modes to
the extremes, where this convergence is fast.}
\label{fig:ritzconverge}
\end{figure}

\section{The restart problem}
\label{sec:restart}

Rayleigh--Ritz says: build the basis, project, read the Ritz values. The hidden assumption
is that we can afford to let the basis grow until the wanted Ritz values converge. For the
extremes on a well-behaved operator that may take only a few dozen steps---but for a
clustered spectrum, or many wanted modes, the basis $\mat{V}_m$ can grow to hundreds or
thousands of columns before convergence. And a growing basis is expensive in two distinct
ways, plus a third that is worse than expensive.

\begin{figure}[t]
\centering
\begin{tikzpicture}
\begin{axis}[
  width=12cm, height=6.3cm,
  xlabel={basis size $m$},
  ylabel={relative cost / loss},
  ymode=log, ymin=1e-15, ymax=3e3,
  xmin=0, xmax=200,
  grid=both, grid style={simGray!18},
  legend pos=outer north east, legend cell align=left,
  legend style={font=\scriptsize, fill=white, draw=simGray!40},
  label style={font=\footnotesize}, tick label style={font=\scriptsize},
]
% storage: linear in m
\addplot[simBlue, very thick, domain=1:200, samples=40] {x};
\addlegendentry{storage $\sim m$ (vectors held)}
% orthogonalization cost per step: linear in m; cumulative quadratic. show per-step ~ m
\addplot[simRust, very thick, domain=1:200, samples=40] {0.5*x};
\addlegendentry{orthog.\ cost/step $\sim m$}
% loss of orthogonality (Lanczos): grows, eventually O(1)
\addplot[simViolet, very thick, domain=1:200, samples=60] {2e-15*exp(0.13*x)};
\addlegendentry{Lanczos loss-of-orthog.\ $\eta_m$}
% the O(1) danger line
\addplot[simGray, densely dashed, thick, domain=0:200, samples=2] {1};
\node[font=\scriptsize, text=simGray, anchor=west] at (axis cs:2,2.3) {orthogonality lost};
% the restart cap
\draw[simGreen, very thick, densely dashed] (axis cs:40,1e-15) -- (axis cs:40,3e3);
\node[font=\scriptsize\bfseries, text=simGreen, rotate=90, anchor=south]
  at (axis cs:40,3e-7) {restart here ($m=k_{\max}$)};
\end{axis}
\end{tikzpicture}
\caption[The restart problem]{Why the basis cannot be allowed to grow without bound. As the
basis size $m$ increases, three costs climb. \emph{Storage} (blue) grows linearly---every
basis vector is a full field, and a thousand of them may not fit in memory. \emph{Orthogonalization
cost} (rust): each new Arnoldi column must be orthogonalized against all $m$ previous ones, so
the per-step cost grows linearly in $m$ and the cumulative cost \emph{quadratically}. Worst of
all, for Lanczos the \emph{loss of orthogonality} (violet, \cref{ch:orthog}) grows until the
basis is no longer orthogonal at all (dashed line at $\eta_m \sim 1$), at which point the
projected $\mat{T}_m$ reports \emph{ghost eigenvalues}---spurious copies of converged
eigenvalues that are artifacts, not real. The cure is to cap the basis at some $k_{\max}$
(green) and \emph{restart}---but a naive restart throws away the converged Ritz vectors. The
art of \cref{sec:thickrestart} is to restart \emph{without} discarding the good directions.}
\label{fig:restartproblem}
\end{figure}

\Cref{fig:restartproblem} tabulates the three pressures. The first is \emph{storage}: each
basis vector $\vq_j$ is a full field of $N$ entries, and $m$ of them occupy $mN$ numbers. At a
million unknowns, a thousand basis vectors is a billion numbers---gigabytes, possibly more than
fits in memory. The second is \emph{orthogonalization cost}: Arnoldi orthogonalizes each new
column against all $m$ previous columns (\cref{ch:orthog}), so step $m$ costs $O(m)$ inner
products and the whole build costs $O(m^2)$. The basis that took $O(m)$ to store takes $O(m^2)$
to build. The third pressure is the dangerous one, and it is specific to Lanczos: \emph{loss of
orthogonality}. The three-term recurrence trusts that its projections against the old basis
vectors are zero and never re-checks them; in finite precision they are not zero, the error
compounds, and the basis silently develops near-duplicate directions
(\cref{ch:orthog}). The projected $\mat{T}_m$ then reports \emph{ghost eigenvalues}---the same
eigenvalue several times over, spurious copies that a naive solver mistakes for a genuine
high-multiplicity mode.

So the basis must be capped. Pick a maximum size $k_{\max}$, build up to it, and then---do what,
exactly? This is the restart problem, and the naive answer is bad.

\begin{pitfall}
The naive restart \emph{discards the basis and starts over} from a single new vector---chosen,
say, as the best converged Ritz vector so far, in the hope of steering the next subspace toward
the wanted modes. This throws away almost everything you computed. The basis $\mat{V}_{k_{\max}}$
encoded \emph{several} nearly-converged Ritz directions, each the fruit of dozens of operator
applications and inner solves; collapsing it to one vector keeps at most one of them and forces
the others to be rediscovered from scratch. Worse, the restart is repeated every cycle, so the
waste compounds: the method spends most of its operator-applies re-converging directions it had
already converged and then threw away. A naive-restart eigensolver can be an order of magnitude
slower than one that restarts intelligently. The fix is to keep \emph{all} the good directions,
not just one---which is exactly what thick-restart does.
\end{pitfall}

The tension is now sharp. We must keep the basis small (storage, cost, orthogonality all demand
it), but we must not throw away the converged and nearly-converged Ritz directions (they are the
expensive part of the answer). A restart that resolves both---small basis, good directions
retained---is what we need, and it exists.

\section{Thick-restart and Krylov--Schur}
\label{sec:thickrestart}

The resolution is due to Stewart, building on the earlier thick-restart Lanczos of Wu and
Simon: restart from a \emph{small} basis that still spans all the Ritz directions worth keeping.
Instead of collapsing $\mat{V}_{k_{\max}}$ to a single vector, we compress it to a basis of
size $k$ (a few more than the number of wanted modes) that retains the $k$ best Ritz directions
and discards the other $k_{\max} - k$. The trick is to do this compression in a way that
\emph{preserves the Krylov structure}---so that the next expansion phase is a legitimate Arnoldi
continuation, not a fresh start.

\subsection{The partial Schur form}

The key is a change of representation. The Arnoldi relation $\mat{A}\mat{V}_m =
\mat{V}_{m+1}\bar{\mat{H}}_m$ stores the projected operator as a Hessenberg matrix $\mat{T}_m$.
Krylov--Schur stores it instead as a \emph{partial Schur form}: it computes the Schur
decomposition $\mat{T}_m = \mat{S}\,\mat{R}\,\mat{S}^{\!\top}$ (with $\mat{R}$ quasi-triangular,
$\mat{S}$ orthogonal) and \emph{reorders} it so that the wanted Ritz values---the converged and
promising ones---sit in the leading $k \times k$ block. Transforming the basis by $\mat{S}$,
$\tilde{\mat{V}}_m = \mat{V}_m\mat{S}$, rotates the orthonormal basis so that its first $k$
columns span exactly the wanted Ritz directions and the projected operator on them is the clean
triangular block $\mat{R}_{1:k,1:k}$.

The structural fact that makes this a \emph{restart} rather than just a reordering is that the
Arnoldi relation survives the transformation in a slightly generalized form. After reordering,
\[
  \mat{A}\,\tilde{\mat{V}}_k
    \;=\; \tilde{\mat{V}}_k\,\mat{R}_{1:k,1:k}
       \;+\; \vq_{m+1}\,\vb_k^{\!\top},
\]
where $\tilde{\mat{V}}_k$ is the first $k$ transformed columns, $\mat{R}_{1:k,1:k}$ the leading
triangular block, and $\vb_k^{\!\top}$ a short row vector (the transformed spill-over). This is a
\emph{Krylov--Schur relation}: like the Arnoldi relation it expresses $\mat{A}$'s action on a
$k$-column basis in terms of that basis plus a single extra direction $\vq_{m+1}$---which means
we can \emph{continue} the Arnoldi expansion from here, appending new columns by the ordinary
recurrence, exactly as if $k$ steps had just been taken. The compressed basis is a valid
launching point for more Arnoldi steps. That is the whole game: compress to size $k$, keeping the
good directions, in a form from which expansion can resume.

\begin{figure}[t]
\centering
\begin{tikzpicture}[font=\footnotesize, scale=0.92]
  % ===== before: full basis of size k_max =====
  \begin{scope}
    \node[font=\small\bfseries, text=simInk] at (1.4,3.1) {before restart: size $k_{\max}$};
    % a row of basis columns, some "good" (green), some "junk" (gray)
    \foreach \i/\col/\lab in {0/simGreen/g,1/simGreen/g,2/simGreen/g,3/simGray/j,4/simGray/j,5/simGray/j,6/simGray/j}{
      \draw[thick, fill=\col!25, draw=\col] (\i*0.42,0) rectangle (\i*0.42+0.36,2.4);
    }
    \node[font=\scriptsize, text=simGreen] at (0.55,-0.4) {$k$ good Ritz dirs};
    \node[font=\scriptsize, text=simGray] at (2.3,-0.4) {$k_{\max}-k$ junk};
    \draw[decorate, decoration={brace, amplitude=4pt}, simGreen] (-0.05,2.55) -- (1.3,2.55);
    \draw[decorate, decoration={brace, amplitude=4pt}, simGray] (1.32,2.55) -- (2.94,2.55);
  \end{scope}
  % ===== arrow: thick-restart compression =====
  \draw[flow, simViolet, very thick] (3.3,1.2) -- (4.7,1.2)
    node[midway, above, font=\scriptsize, text=simViolet, align=center]
      {Schur reorder\\+ truncate};
  \node[font=\scriptsize\itshape, text=simRust, align=center] at (4.0,0.3)
    {discard junk,\\keep good};
  % ===== after: compressed basis of size k + 1 spill =====
  \begin{scope}[xshift=5.2cm]
    \node[font=\small\bfseries, text=simInk] at (1.0,3.1) {after restart: size $k$};
    \foreach \i in {0,1,2}{
      \draw[thick, fill=simGreen!25, draw=simGreen] (\i*0.42,0) rectangle (\i*0.42+0.36,2.4);}
    % the carried spill-over vector q_{m+1}
    \draw[thick, fill=simBgRust, draw=simRust] (1.55,0) rectangle (1.91,2.4);
    \node[font=\scriptsize, text=simRust, rotate=90] at (1.73,1.2) {$\vq_{m+1}$};
    \node[font=\scriptsize, text=simGreen] at (0.55,-0.4) {$k$ kept dirs};
    \node[font=\scriptsize, text=simRust] at (1.73,-0.75) {carry};
    % expansion resumes
    \draw[flow, simBlue, thick] (2.3,1.2) -- (3.3,1.2)
      node[midway, above, font=\scriptsize, text=simBlue]{expand};
    \foreach \i in {0,1}{
      \draw[thick, fill=simBgBlue, draw=simBlue, densely dashed] (3.5+\i*0.42,0) rectangle (3.5+\i*0.42+0.36,2.4);}
    \node[font=\scriptsize, text=simBlue] at (3.9,-0.4) {new Arnoldi cols};
  \end{scope}
\end{tikzpicture}
\caption[Thick-restart / Krylov--Schur]{One thick-restart cycle. \emph{Left:} the basis has
grown to its cap $k_{\max}$; some columns span converged or promising Ritz directions (green,
the $k$ we want to keep), the rest are exhausted Krylov directions of no further use (grey).
A naive restart would keep only one green column; thick-restart keeps \emph{all $k$}.
\emph{Centre:} the compression reorders the Schur form of the projected operator so the wanted
Ritz values lead, then truncates to the leading $k$ columns---discarding the grey junk while
retaining every good direction. \emph{Right:} the compressed size-$k$ basis, together with the
single carried spill-over vector $\vq_{m+1}$ (rust) from the Krylov--Schur relation, is a valid
Arnoldi launching point: ordinary expansion (blue, dashed) resumes from it, rebuilding toward
$k_{\max}$ before the next compression. The cycle---expand, extract, lock, compress---repeats
until the wanted modes converge.}
\label{fig:thickrestart}
\end{figure}

\subsection{The restart cycle}

With the partial Schur form in hand, the Krylov--Schur eigensolver is a loop of four stages,
shown in \cref{fig:thickrestart} and \cref{fig:ksloop}.

\begin{enumerate}
  \item \textbf{Expand.} From the current basis of size $k$, take Arnoldi (or Lanczos) steps
  until the basis reaches the cap $k_{\max}$, filling the projected operator's new columns. Each
  step is one \code{arnoldi\_step} or \code{lanczos\_step} of \cref{ch:orthog}---and, in the
  eigenmode driver, each contains one inner shift-invert solve (\cref{sec:shiftinvert}).
  \item \textbf{Extract.} Form the projected operator $\mat{T}_{k_{\max}}$, compute its Schur
  form, and read the Ritz values and the cheap residuals $\abs{h_{m+1,m}}\abs{\ve_m^{\!\top}\vs_i}$
  of \cref{sec:rayleighritz}.
  \item \textbf{Lock.} Mark the Ritz pairs whose residual is below tolerance as
  \emph{converged}: they are part of the answer and are frozen, no longer to be touched by the
  iteration. When the locked count reaches the number of requested modes, the loop terminates.
  \item \textbf{Compress (thick-restart).} Reorder the Schur form so the converged and most
  promising Ritz values lead, truncate the basis to the leading $k$ columns, carry the single
  spill-over vector, and return to the Expand stage.
\end{enumerate}

\begin{figure}[t]
\centering
\begin{tikzpicture}[font=\footnotesize, node distance=9mm]
  \node[stage, minimum width=22mm] (expand) {\textbf{Expand}\\[1pt]\scriptsize Arnoldi/Lanczos\\\scriptsize to $k_{\max}$};
  \node[stage, minimum width=22mm, right=18mm of expand] (extract) {\textbf{Extract}\\[1pt]\scriptsize Rayleigh--Ritz\\\scriptsize + cheap residual};
  \node[stage, minimum width=22mm, right=18mm of extract] (lock) {\textbf{Lock}\\[1pt]\scriptsize freeze converged\\\scriptsize Ritz pairs};
  \node[stage, minimum width=22mm, below=14mm of extract] (compress) {\textbf{Compress}\\[1pt]\scriptsize Schur reorder\\\scriptsize truncate to $k$};
  \node[product, right=16mm of lock] (done) {converged\\modes};
  \draw[flow] (expand) -- (extract);
  \draw[flow] (extract) -- (lock);
  \draw[flow] (lock) -- node[above,font=\scriptsize]{all wanted?} (done);
  \draw[flow] (lock.south) |- node[pos=0.25,right,font=\scriptsize]{not yet} (compress.east);
  \draw[flow] (compress) -| (expand.south);
  % inner-solve callout on Expand
  \node[extern, below=10mm of expand, align=center] (inner)
    {\scriptsize each step:\\\scriptsize inner shift-invert\\\scriptsize \code{ksp\_solve} (\cref{ch:eigenvalues})};
  \draw[refedge] (inner.north) -- (expand.south);
\end{tikzpicture}
\caption[The Krylov--Schur restart loop]{The four-stage Krylov--Schur cycle, the algorithm
inside the opaque eigensolver call. \textbf{Expand} the basis by Arnoldi/Lanczos steps up to the
cap $k_{\max}$; \textbf{extract} the Ritz pairs and their cheap residuals by Rayleigh--Ritz;
\textbf{lock} the converged pairs (freezing them out of further iteration), and if all requested
modes are locked, return them as the answer; otherwise \textbf{compress} by thick-restart---reorder
the Schur form to put the good Ritz directions first, truncate to size $k$, and loop back to
expand. The dashed callout records the stack: in the eigenmode driver every Expand step contains
an inner shift-invert linear solve (\cref{sec:shiftinvert}), so the outer loop you see here sits
atop the inner Krylov solver of \cref{ch:multigrid}. This loop is exactly what SLEPc's
\code{EPSKRYLOVSCHUR} runs.}
\label{fig:ksloop}
\end{figure}

This is the constructive driver that the source spec reconstructs from the firm
\code{arnoldi\_step} and \code{lanczos\_step} kernels: an outer thick-restart loop over restart
cycles, wrapping an inner basis-extension loop over basis columns, with a Rayleigh--Ritz
extraction and a lock-and-compress between cycles. In the calculus of \cref{ch:fold} it is two
nested \code{iterate\_while} folds---the outer over restart cycles, the inner over basis
columns---and the body of the inner fold is one Arnoldi/Lanczos step.

\begin{lstlisting}[style=l4style]
-- the Krylov-Schur driver: outer thick-restart fold over an inner basis-extension fold
eigsolve_ks :: LinOp -> EigControl -> EigResult
eigsolve_ks A control =
  let st0 = seed_basis A control                       -- BV <- [v0 / ||v0||]
      loop st = let ext      = extend_basis A st        -- INNER: Arnoldi/Lanczos to k_max
                    (theta, Y) = rayleigh_ritz A ext    -- project + dense QR -> Ritz pairs
                    locked   = lock_converged ext theta Y control  -- cheap residual test
                in thick_restart locked theta Y control  -- reorder Schur, truncate to k
      st_final = iterate_while loop (\st -> not (enough_converged st control)) st0
  in extract_eigpairs A st_final                        -- lift Ritz vecs, un-transform lambda
  where
    -- INNER basis-extension: ncv steps, each ONE arnoldi/lanczos step
    extend_basis A st =
      iterate_while (\s -> append_column s (one_step A s)) (\s -> s.j < control.k_max) st
    one_step A s = if A.hermitian then lanczos_step A s   -- symmetric three-term recurrence
                                  else arnoldi_step A s    -- full Arnoldi
\end{lstlisting}

\begin{workedexample}{One thick-restart step, with small numbers}{thickrestart}
We make the compression concrete on a schematic example. Suppose we have expanded to a basis of
size $k_{\max} = 6$ and want $2$ eigenvalues, restarting to size $k = 3$ (one extra beyond the
wanted count, for headroom). The Rayleigh--Ritz extraction on $\mat{T}_6$ returns six Ritz values,
which we sort by how close they are to the wanted target (here, the largest two), and tag with
their cheap residuals $r_i = \abs{h_{7,6}}\abs{\ve_6^{\!\top}\vs_i}$:
\[
  \begin{array}{c|cccccc}
    \text{Ritz value } \theta_i & 7.91 & 7.40 & 6.2 & 5.1 & 3.8 & 1.2 \\\hline
    \text{residual } r_i & 4\!\times\!10^{-9} & 3\!\times\!10^{-4} & 0.2 & 0.5 & 0.7 & 0.9 \\
    \text{status} & \text{\textbf{locked}} & \text{promising} & \text{keep} & \text{junk} & \text{junk} & \text{junk}
  \end{array}
\]
\emph{Lock.} The first Ritz value $\theta_1 = 7.91$ has residual $4\times10^{-9}$, below the
tolerance $10^{-8}$: it is \emph{converged} and locked---one of the two wanted modes, now frozen.
The second, $\theta_2 = 7.40$ at residual $3\times10^{-4}$, is not yet converged but is the next
wanted mode and converging well: \emph{promising}.

\smallskip
\emph{Choose what to keep.} We keep $k = 3$ directions: the locked $\theta_1$, the promising
$\theta_2$, and the next-best $\theta_3 = 6.2$ for headroom (a converging interior direction worth
continuing). The bottom three Ritz values $\{5.1, 3.8, 1.2\}$---large residuals, far from the
target---are \emph{junk}: exhausted Krylov directions with nothing left to give. They are discarded.

\smallskip
\emph{Compress.} Reorder the Schur form of $\mat{T}_6$ so the kept Ritz values
$\{7.91, 7.40, 6.2\}$ occupy its leading $3\times3$ block, transform the basis
$\tilde{\mat{V}}_6 = \mat{V}_6\mat{S}$, and truncate to the first $3$ columns
$\tilde{\mat{V}}_3$. Together with the carried spill-over vector $\vq_7$, these satisfy the
Krylov--Schur relation $\mat{A}\tilde{\mat{V}}_3 = \tilde{\mat{V}}_3\mat{R}_{1:3} +
\vq_7\vb_3^{\!\top}$, so the next cycle resumes Arnoldi from a \emph{size-$3$} basis that already
spans the converged mode, the next mode, and a promising third---none of which had to be
rediscovered. The basis shrank from $6$ back to $3$, storage and orthogonalization cost reset to
the floor, and \emph{not one} of the good directions was lost. Contrast the naive restart, which
would have kept only $\theta_1$ and forced $\theta_2$ and $\theta_3$ to be rebuilt from a single
vector. That saved rediscovery, multiplied over many cycles, is the entire value of
thick-restart.
\end{workedexample}

\section{Lanczos specifics: the symmetric three-term recurrence}
\label{sec:lanczosspecific}

The eigenmode and boundary-mode analyses pose \emph{symmetric} (or Hermitian) eigenproblems---the
operators $\mat{K}$ and $\mat{M}$ are symmetric positive definite. For these, the Arnoldi basis
extension of the previous sections specializes to the \emph{Lanczos} three-term recurrence of
\cref{ch:orthog}, and the specialization is worth dwelling on because it is both a major economy
and a major hazard.

\subsection{Why symmetric is cheaper}

Recall from \cref{ch:orthog} the symmetry miracle: when $\mat{A} = \mat{A}^{\!\top}$, the projected
Hessenberg matrix $\mat{T}_m = \mat{V}_m^{\!\top}\mat{A}\mat{V}_m$ is itself symmetric, and a
symmetric Hessenberg matrix is \emph{tridiagonal}. So orthogonalizing the new vector
$\mat{A}\vq_j$ against the basis---which for general Arnoldi means projecting against \emph{all}
$j$ previous columns---collapses to projecting against just the \emph{two} most recent,
$\vq_j$ and $\vq_{j-1}$. The projections against $\vq_1, \dots, \vq_{j-2}$ are all exactly zero.
The step is
\[
  \vw \;=\; \mat{A}\vq_j - \alpha_j\vq_j - \beta_{j-1}\vq_{j-1},
  \qquad
  \alpha_j = \inner{\mat{A}\vq_j}{\vq_j},
  \qquad
  \beta_j = \norm{\vw},
  \qquad
  \vq_{j+1} = \frac{\vw}{\beta_j},
\]
storing only $\vq_j, \vq_{j-1}$ and producing the tridiagonal entries $(\alpha_j, \beta_j)$
directly. This is the \code{lanczos\_step} kernel from which the symmetric arm of the eigensolver
is built:

\begin{lstlisting}[style=l4style]
-- the symmetric three-term recurrence: O(1) memory per step, only TWO stored vectors
lanczos_step :: LinOp -> Tensor[$S] -> Tensor[$S] -> Scalar -> (Tensor[$S], Scalar, Scalar)
lanczos_step A q_j q_prev beta_prev =
  let u      = apply A q_j                  -- A q_j
      alpha  = dot u q_j                     -- diagonal coeff alpha_j = <A q_j, q_j>
      w      = axpby (-alpha) q_j 1 u        -- subtract the q_j component
      w'     = axpy (-beta_prev) q_prev w    -- subtract the q_{j-1} component (only TWO!)
      beta   = nrm2 w'                       -- new off-diagonal / normalization beta_j
      q_next = scal (1/beta) w'
  in (q_next, alpha, beta)
\end{lstlisting}

The economy is dramatic. General Arnoldi stores the whole growing basis ($O(m)$ vectors) and
spends $O(m)$ inner products per step; Lanczos stores \emph{two} vectors and spends a constant
number of inner products per step, regardless of $m$. For a symmetric problem the per-step cost
does not grow at all. This is the same short-recurrence economy that gives conjugate gradients its
efficiency (\cref{ch:cg}): symmetry makes a quantity that would otherwise depend on the entire
history depend only on the last two steps.

\subsection{The hazard and the reorthogonalization policy}

The economy is bought on credit, and \cref{ch:orthog} showed the bill: the short recurrence is
short \emph{because} it trusts that the skipped projections are zero, and in finite precision they
are not. Lanczos never looks at the old basis vectors, so it cannot notice---let alone
repair---the orthogonality it loses against them. The orthogonality decays (the violet curve of
\cref{fig:restartproblem}), and the symptom is the \emph{ghost eigenvalue}: the tridiagonal
$\mat{T}_m$ reports a converged eigenvalue several times over, spurious copies of a real eigenvalue
that the lost orthogonality manufactured. A naive Lanczos solver returns a multiplicity-three
eigenvalue that truly has multiplicity one.

The practical eigensolver therefore augments the bare recurrence with a \emph{reorthogonalization
policy}---a decision about when, and against what, to do the extra orthogonalization the short
recurrence omits. \Cref{tab:reorthog} lays out the three standard policies and the cost/robustness
trade each strikes.

\begin{table}[t]
\centering
\caption[Lanczos reorthogonalization policies]{The reorthogonalization policies that buy back the
orthogonality the three-term recurrence loses. \emph{Full} reorthogonalization is the safest and
most expensive---it re-projects every new vector against the entire stored basis, sacrificing the
$O(1)$-memory advantage that motivated Lanczos in the first place. \emph{Selective} and
\emph{partial} reorthogonalization monitor a cheap estimate of the orthogonality loss and
reorthogonalize only when it crosses a threshold, recovering most of the safety at a fraction of
the cost. Which policy the solver runs is a central configuration decision; in a thick-restart
method the restart itself caps the basis size, which bounds the orthogonality loss between restarts
and makes full reorthogonalization affordable.}
\label{tab:reorthog}
\small
\setlength{\tabcolsep}{6pt}
\renewcommand{\arraystretch}{1.25}
\begin{tabular}{@{}llll@{}}
\toprule
Policy & When to reorthogonalize & Against & Cost \\
\midrule
None (bare) & never & --- & cheapest; ghost eigenvalues \\
Full & every step & whole stored basis & most robust; $O(m)$/step, $O(m)$ storage \\
Selective & monitored loss high & converged Ritz vectors only & moderate \\
Partial & monitored loss $>\sqrt{u}$ & recent basis window & moderate; near-full safety \\
\bottomrule
\end{tabular}
\end{table}

\begin{pitfall}
It is tempting to read ``Lanczos is the cheap symmetric method'' and reach for the bare three-term
recurrence whenever the operator is symmetric. Do not, for an eigensolver. Without a
reorthogonalization policy the bare recurrence \emph{will} produce ghost eigenvalues on any
problem run long enough, and the failure is silent: the solver returns confident, wrong
multiplicities with no error flag. The three-term recurrence is the right \emph{kernel}, but it is
not a self-sufficient eigensolver---it must be wrapped in a reorthogonalization policy (and, in
practice, a restart that caps how far orthogonality can decay between restarts). The source spec
records exactly this: a firm Lanczos eigensolver ``must pin a reorthogonalization policy; the
band-3 recurrence alone is not numerically self-sufficient for a converged eigensolve.''
\end{pitfall}

Thick-restart and reorthogonalization are complementary cures for the same disease. The restart
caps the basis at $k_{\max}$, which bounds how far orthogonality can decay before the basis is
compressed and effectively refreshed; reorthogonalization repairs the decay within a single
expansion phase. Together they make symmetric Krylov--Schur both cheap (short recurrence, small
basis) and trustworthy (no ghosts)---the combination the eigenmode driver actually runs.

\section{Putting it together: the shift-invert stack}
\label{sec:shiftinvert}

We have built the eigensolver in the abstract, with $\mat{A}$ standing for ``the operator we
iterate against.'' In the eigenmode driver, $\mat{A}$ is not the original operator. Recall from
\cref{sec:rayleighritz} that Krylov projection captures the \emph{extreme} eigenvalues fast and
the interior ones slowly---but the wanted resonant modes are interior, clustered near a target
frequency, not at the extremes of the spectrum. Running Krylov--Schur on $\mat{K}\vx = \lambda
\mat{M}\vx$ directly would converge agonizingly slowly to exactly the eigenvalues we care about.

The remedy is the \emph{shift-invert} spectral transform of \cref{ch:eigenvalues}. Instead of
iterating against the original pencil, we iterate against
\[
  \mat{A} \;=\; (\mat{K} - \sigma\mat{M})^{-1}\mat{M},
\]
whose eigenvalues are $1/(\lambda - \sigma)$. This transform turns the eigenvalues $\lambda$ near
the shift $\sigma$ into the eigenvalues $1/(\lambda - \sigma)$ of \emph{largest magnitude}---the
extremes. The wanted interior modes become the dominant modes of the transformed operator, and
Krylov--Schur converges to them fast. After the iteration, each converged Ritz value $\theta$ is
un-transformed back to an original-problem eigenvalue $\lambda = \sigma + 1/\theta$.

The crucial structural consequence is that \emph{applying the transformed operator $\mat{A}$
requires an inner linear solve}. There is no matrix $(\mat{K} - \sigma\mat{M})^{-1}$; there is only
the \emph{action} of solving a linear system with $\mat{K} - \sigma\mat{M}$. So one application of
$\mat{A}$ to a vector $\vv$ is two stages: first apply $\mat{M}$ (a matrix-vector product), then
solve $(\mat{K} - \sigma\mat{M})\vy = \mat{M}\vv$ (a full linear solve). That inner solve is the
\code{ksp\_solve} of \cref{ch:krylov}---a Krylov solver in its own right, preconditioned by the
multigrid of \cref{ch:multigrid}, applying the matrix-free operator of \cref{ch:matrixfree}. The
eigensolver's every basis-extension step contains, buried inside its single ``apply $\mat{A}$,'' a
complete inner Krylov solve.

\begin{lstlisting}[style=l4style]
-- one application of the shift-inverted operator: an INNER linear solve inside every Arnoldi step
apply_shift_invert :: EigOp -> Tensor[$S] -> Tensor[$S]
apply_shift_invert op v =
  let w = apply_linop op.operand v     -- stage 1: apply M (or K for no transform)
      y = ksp_solve   op.inv     w     -- stage 2: solve (K - sigma M) y = w  -- INNER Krylov solve!
  in scale_untransform op y            -- per-backend Higham gamma/delta un-scale
\end{lstlisting}

This is the full stack, and it is worth seeing all at once. \Cref{fig:eigenstack} draws it: the
outer Krylov--Schur eigen-iteration owns the basis and the restart; each of its steps applies the
shift-inverted operator; that application contains an inner Krylov linear solve; the inner solve is
accelerated by a preconditioner; and the preconditioner, like the operator, is applied
matrix-free. Five layers, each a chapter of this Part, nested one inside the next.

\begin{figure}[t]
\centering
\begin{tikzpicture}[font=\footnotesize]
  % nested boxes, outer to inner
  \draw[thick, fill=simBgGold, rounded corners=3pt] (-0.2,-0.2) rectangle (10.6,5.4);
  \node[anchor=north west, font=\small\bfseries, text=simGold!60!black] at (0.0,5.3)
    {\,1.\ Krylov--Schur eigen-iteration\quad\scriptsize(this chapter: basis, restart, Ritz)};
  \draw[thick, fill=simBgViolet, rounded corners=3pt] (0.3,0.2) rectangle (10.1,4.5);
  \node[anchor=north west, font=\small\bfseries, text=simViolet] at (0.5,4.4)
    {\,2.\ shift-invert apply $\;\mat{A}=(\mat{K}-\sigma\mat{M})^{-1}\mat{M}$\quad\scriptsize(\cref{ch:eigenvalues})};
  \draw[thick, fill=simBgBlue, rounded corners=3pt] (0.8,0.6) rectangle (9.6,3.6);
  \node[anchor=north west, font=\small\bfseries, text=simBlue] at (1.0,3.5)
    {\,3.\ inner Krylov solve $\;(\mat{K}-\sigma\mat{M})\vy=\vw$\quad\scriptsize\code{ksp\_solve} (\cref{ch:krylov})};
  \draw[thick, fill=simBgGreen, rounded corners=3pt] (1.3,1.0) rectangle (9.1,2.7);
  \node[anchor=north west, font=\small\bfseries, text=simGreen] at (1.5,2.6)
    {\,4.\ preconditioner $\;\mat{P}^{-1}$\quad\scriptsize geometric multigrid (\cref{ch:multigrid})};
  \draw[thick, fill=simBgRust, rounded corners=3pt] (1.8,1.4) rectangle (8.6,2.1);
  \node[anchor=west, font=\small\bfseries, text=simRust] at (2.0,1.75)
    {\,5.\ matrix-free operator apply $\;\mat{K}\vv,\ \mat{M}\vv$\quad\scriptsize(\cref{ch:matrixfree})};
\end{tikzpicture}
\caption[The full eigen-solve stack]{The five nested layers of an eigenmode solve, each a chapter
of this Part. \textbf{(1)} The outer Krylov--Schur eigen-iteration of this chapter owns the basis,
the Rayleigh--Ritz extraction, and the thick-restart---the loop SLEPc calls \code{EPSKRYLOVSCHUR}.
\textbf{(2)} Each basis-extension step applies the shift-inverted operator $\mat{A} = (\mat{K} -
\sigma\mat{M})^{-1}\mat{M}$ of \cref{ch:eigenvalues}, which remaps the wanted interior modes to the
fast-converging extremes. \textbf{(3)} Applying that operator requires an inner Krylov linear solve
$(\mat{K} - \sigma\mat{M})\vy = \vw$---the \code{ksp\_solve} of \cref{ch:krylov}, a solver nested
inside the eigensolver. \textbf{(4)} That inner solve is accelerated by a preconditioner, the
geometric multigrid of \cref{ch:multigrid}. \textbf{(5)} And the operator and preconditioner are
both applied matrix-free (\cref{ch:matrixfree}), never assembled as explicit matrices. The opaque
``eigenvalue solve'' of \cref{ch:eigenvalues} is this entire tower, and every layer is something we
have built.}
\label{fig:eigenstack}
\end{figure}

\begin{keyidea}
The eigenmode solve is a five-layer stack: \emph{Krylov--Schur eigen-iteration} (this chapter)
$\triangleright$ \emph{shift-invert apply} (\cref{ch:eigenvalues}) $\triangleright$ \emph{inner
Krylov solve} (\cref{ch:krylov}) $\triangleright$ \emph{preconditioner} (\cref{ch:multigrid})
$\triangleright$ \emph{matrix-free operator apply} (\cref{ch:matrixfree}). The outer eigen-loop is
what the library owns; the inner linear solve is the \code{ksp\_solve} the rest of this Part built.
Every basis-extension step of the eigensolver contains one full inner Krylov solve, so the cost of
an eigenmode analysis is dominated by the inner solves---which is why the preconditioner of
\cref{ch:multigrid} matters as much to the eigensolver as to any linear solve.
\end{keyidea}

\section{The library boundary, demystified}
\label{sec:libraryboundary}

We can now return to the promise of \cref{ch:eigenvalues}: that the eigenmode driver's
``solve''---the one stage of the whole simulator that was a single call into an opaque library
rather than a loop we wrote---would turn out to be a loop we could have written ourselves. We have
now written it. The four-stage cycle of \cref{fig:ksloop}, with the Lanczos or Arnoldi kernel of
\cref{ch:orthog} at its heart and the shift-invert stack of \cref{fig:eigenstack} inside each step,
\emph{is} the algorithm that SLEPc's \code{EPSSolve} with the \code{EPSKRYLOVSCHUR} type runs, and
that ARPACK runs through its reverse-communication \code{naupd}.

The phrase ``reverse communication'' names the only genuinely unusual thing about the library
boundary, and it is a packaging detail, not an algorithmic one. ARPACK does not call Palace's
operator; it cannot, because it was compiled with no knowledge of Palace's data structures.
Instead, the control inverts: Palace calls into ARPACK's \code{naupd} driver, which runs a little
of the Krylov--Schur loop and then \emph{returns}, setting a flag that says ``I need
$\mat{A}\vv$ for this vector $\vv$---compute it and call me again.'' Palace computes the matrix-vector
product (running the entire shift-invert stack of \cref{sec:shiftinvert} to do so) and calls
\code{naupd} again; the driver resumes where it left off, takes the next step, and returns with the
next request. The eigen-iteration logic lives entirely inside the library; Palace's loop is a
dispatcher that services the library's requests for operator applications.

\begin{figure}[t]
\centering
\begin{tikzpicture}[font=\footnotesize, node distance=10mm]
  % Palace side
  \node[stage, minimum width=30mm, align=center] (palace)
    {\textbf{Palace}\\[1pt]\scriptsize dispatcher loop\\\scriptsize + operator applies};
  % library side
  \node[extern, minimum width=34mm, minimum height=20mm, right=34mm of palace, align=center] (lib)
    {\textbf{SLEPc / ARPACK}\\[2pt]\scriptsize \emph{the Krylov--Schur loop}\\\scriptsize basis $\cdot$ restart $\cdot$ Ritz\\[3pt]
     \scriptsize\itshape (this chapter's\\\scriptsize\itshape algorithm, inside)};
  % the two reverse-communication arrows
  \draw[flow, simViolet] (palace.north east) .. controls +(1.0,0.8) and +(-1.0,0.8) .. (lib.north west)
    node[midway, above, font=\scriptsize, text=simViolet, align=center]{call \code{naupd};\\``take a step''};
  \draw[backflow] (lib.south west) .. controls +(-1.0,-0.8) and +(1.0,-0.8) .. (palace.south east)
    node[midway, below, font=\scriptsize, text=simRust, align=center]{return: ``I need\\$\mat{A}\vv$ for this $\vv$''};
  % the operator-apply that Palace runs to service the request
  \node[stage, below=12mm of palace, minimum width=30mm, align=center] (apply)
    {\scriptsize apply $\mat{A}=(\mat{K}{-}\sigma\mat{M})^{-1}\mat{M}$\\\scriptsize (the inner stack)};
  \draw[flow] (palace.south) -- (apply.north);
  \draw[flow] (apply.east) -| ([xshift=-3mm]lib.south) node[pos=0.25,below,font=\scriptsize]{$\mat{A}\vv$ back};
  \node[font=\scriptsize\itshape, text=simGray, align=center, anchor=north] at (lib.south|-apply.south)
    {the boundary is a\\\emph{packaging} detail,\\not an algorithm};
\end{tikzpicture}
\caption[The reverse-communication library boundary]{What the ``opaque library call'' of
\cref{ch:eigenvalues} actually is. The Krylov--Schur algorithm of this chapter lives inside the
library (SLEPc or ARPACK, violet box on the right)---but the library cannot call Palace's operator
directly, because it knows nothing of Palace's data structures. So control inverts by
\emph{reverse communication}: Palace calls the library driver \code{naupd} to ``take a step''; the
library advances the eigen-iteration and \emph{returns}, requesting a specific matrix-vector product
$\mat{A}\vv$; Palace runs the entire inner shift-invert stack (\cref{fig:eigenstack}) to compute it
and calls back; the library resumes. The algorithm is the loop of \cref{fig:ksloop}; the only thing
the library boundary adds is this request-and-resume packaging. Open the box and the
``opaque'' eigensolver is exactly the construction of this chapter.}
\label{fig:reversecomm}
\end{figure}

\Cref{fig:reversecomm} draws the dance. The lesson is the one this book returns to again and again:
a black box is opaque only until you build its contents. The eigensolver looked, in
\cref{ch:eigenvalues}, like the one stage of the simulator that was irreducibly someone else's
code---a single library call with no loop of our own. It is not. It is the Rayleigh--Ritz
projection of \cref{sec:rayleighritz}, restarted by the thick-restart of \cref{sec:thickrestart},
built on the Arnoldi and Lanczos recurrences of \cref{ch:orthog}, with the shift-invert linear
solve of \cref{ch:krylov} inside each step. We could write it ourselves; Palace calls the library
because the library is well-tested and tuned, not because the algorithm is beyond us. The box is
demystified.

\begin{notationbox}
\textbf{Why call the library at all, then?} If the algorithm is just this chapter's loop, why not
write it? Three reasons, none algorithmic. First, \emph{robustness}: SLEPc and ARPACK encode
decades of hard-won handling of edge cases---breakdowns, clustered eigenvalues, the
reorthogonalization policies of \cref{sec:lanczosspecific}---that a from-scratch implementation
would rediscover the hard way. Second, \emph{the spectral-transform plumbing}: the library manages
the un-transformation $\lambda = \sigma + 1/\theta$, the Higham scaling, and the convergence
bookkeeping across multiple backends behind one interface. Third, \emph{focus}: Palace's value is in
the physics and the assembly, not in re-implementing a numerical eigensolver. Calling the library is
the right engineering choice precisely \emph{because} we understand what it does---we can configure
it, diagnose it, and trust it, having seen its insides.
\end{notationbox}

\summarysection
The opaque eigenvalue solve of \cref{ch:eigenvalues} is the Krylov--Schur algorithm, and this
chapter built it. \emph{Rayleigh--Ritz projection} takes an orthonormal Krylov basis $\mat{V}_m$
(from the Arnoldi or Lanczos recurrence of \cref{ch:orthog}), forms the small projected operator
$\mat{T}_m = \mat{V}_m^{\!\top}\mat{A}\mat{V}_m$---which is exactly the Hessenberg/tridiagonal matrix
the recurrence already filled---and reads its eigenpairs as \emph{Ritz pairs} $(\theta_i,
\mat{V}_m\vs_i)$ approximating $\mat{A}$'s eigenpairs, the extremal ones first and best (the
polynomial / min-max view). The Ritz residual is the cheap product $\abs{h_{m+1,m}}
\abs{\ve_m^{\!\top}\vs_i}$, a convergence test costing two already-computed scalars. But the basis
cannot grow without bound---storage and orthogonalization cost climb, and Lanczos loses
orthogonality and spawns ghost eigenvalues---so the method must \emph{restart}, and a naive restart
throws away the converged Ritz directions. \emph{Thick-restart Krylov--Schur} (Stewart) is the
elegant fix: reorder the partial Schur form to put the wanted Ritz values first, truncate the basis
to a small size $k$ that retains all the good directions, carry one spill-over vector, and resume
Arnoldi from there---expand, extract, lock, compress, repeat. For symmetric problems the basis
extension is the Lanczos three-term recurrence ($O(1)$ memory, two stored vectors), economical but
demanding a reorthogonalization policy to suppress ghosts. In the eigenmode driver the operator
iterated against is the shift-inverted $(\mat{K} - \sigma\mat{M})^{-1}\mat{M}$ of
\cref{ch:eigenvalues}, so every basis-extension step contains an inner Krylov solve---the full stack
is eigen-iteration $\triangleright$ shift-invert apply $\triangleright$ inner \code{ksp\_solve}
$\triangleright$ preconditioner $\triangleright$ matrix-free operator. SLEPc's \code{EPSKRYLOVSCHUR}
and ARPACK's reverse-communication \code{naupd} run exactly this loop; the ``opaque library call''
is a packaging detail around an algorithm we can---and now do---understand fully.

\furtherreading
The thick-restart Krylov--Schur method is due to Stewart, whose paper~\citep{stewart2001krylovschur}
introduces the Krylov--Schur relation and the reordering-and-truncation restart that
\cref{sec:thickrestart} develops; it is the cleanest single source for the algorithm at the heart of
this chapter, and is the reconstruction the source specification follows. For the broader theory of
Krylov eigensolvers---Rayleigh--Ritz projection, the convergence of Ritz values, the Lanczos and
Arnoldi recurrences, reorthogonalization policies, and the ghost-eigenvalue phenomenon---Saad's
\emph{Numerical Methods for Large Eigenvalue Problems}~\citep{saad2011eig} is the standard graduate
reference and the natural next step after this chapter. The dense QR algorithm that diagonalizes the
small projected $\mat{T}_m$, together with the Schur decomposition and its reordering that
thick-restart relies on, are developed definitively in Golub and Van Loan's \emph{Matrix
Computations}~\citep{golubvanloan2013}, which also treats the symmetric Lanczos process and its
finite-precision behaviour in full. Read Stewart for the restart, Saad for the eigensolver theory,
and Golub and Van Loan for the dense kernels underneath---and keep \cref{ch:eigenvalues}'s
shift-invert transform in view throughout, since it is what makes Krylov projection converge to the
interior modes the simulator actually wants.

\exercisessection
\begin{enumerate}[nosep]
  \item \emph{(Rayleigh--Ritz by hand.)} Repeat \cref{wex:rrritz} on
  $\mat{A} = \begin{psmallmatrix} 3 & 1 & 0 & 0\\ 1 & 3 & 1 & 0\\ 0 & 1 & 3 & 1\\ 0 & 0 & 1 & 3
  \end{psmallmatrix}$ with $\vq_1 = (1,0,0,0)^{\!\top}$. Build $\mat{T}_3$ by Lanczos, compute its
  three Ritz values, and verify that the extreme Ritz values bracket the true extreme eigenvalues
  $3 + 2\cos(k\pi/5)$ from inside.
  \item \emph{(The cheap residual.)} For the largest Ritz pair of \cref{wex:rrritz}, the residual
  was $\abs{h_{4,3}}\abs{\ve_3^{\!\top}\vs} = 1\cdot\tfrac12$. Derive the formula
  $\norm{\mat{A}\vx_i - \theta_i\vx_i} = \abs{h_{m+1,m}}\abs{\ve_m^{\!\top}\vs_i}$ from the Arnoldi
  relation $\mat{A}\mat{V}_m = \mat{V}_{m+1}\bar{\mat{H}}_m$, explaining at each step why the
  projected part cancels. Why does a small $h_{m+1,m}$ make \emph{every} Ritz pair converge at once?
  \item \emph{(Why the extremes.)} Using the polynomial picture $\mathcal{K}_m = \{p(\mat{A})\vv :
  \deg p < m\}$, explain why a degree-$(m{-}1)$ polynomial can isolate an extreme eigenvalue with
  small $m$ but needs large $m$ to isolate an interior one. What does this predict about the relative
  convergence speed of the extreme versus interior Ritz values in \cref{fig:ritzconverge}?
  \item \emph{(The cost of not restarting.)} A basis of size $m$ costs $O(mN)$ storage and
  $O(m^2 N)$ total orthogonalization work for general Arnoldi. If an eigenproblem needs $m = 500$
  steps to converge $10$ modes without restart, but thick-restart caps the basis at $k_{\max} = 50$,
  estimate the ratio of storage and of orthogonalization cost between the two. Which of the three
  pressures of \cref{fig:restartproblem} does restart relieve, and which does reorthogonalization
  relieve?
  \item \emph{(One thick-restart step.)} In \cref{wex:thickrestart} we kept $k = 3$ of $k_{\max} = 6$
  Ritz directions. Suppose instead we want $4$ modes and restart to $k = 5$. Given Ritz values
  $\{7.9, 7.4, 6.8, 6.1, 5.0, 1.2\}$ with residuals $\{10^{-9}, 10^{-7}, 10^{-3}, 10^{-2}, 0.4,
  0.9\}$ and tolerance $10^{-8}$, identify which are locked, which are kept, and which are discarded,
  and explain why keeping one extra direction beyond the wanted count is worthwhile.
  \item \emph{(Lanczos ghosts.)} Explain, in terms of the three-term recurrence's skipped
  projections (\cref{sec:lanczosspecific}), why bare Lanczos produces ghost eigenvalues but bare
  Arnoldi does not. Which reorthogonalization policy of \cref{tab:reorthog} restores safety at the
  least cost, and why does thick-restart partially substitute for reorthogonalization?
  \item \emph{(The shift-invert stack.)} Trace one application of $\mat{A} = (\mat{K} -
  \sigma\mat{M})^{-1}\mat{M}$ to a vector $\vv$ through all five layers of \cref{fig:eigenstack},
  naming the chapter that owns each layer. Why is an eigenmode analysis's runtime dominated by the
  inner \code{ksp\_solve} rather than by the outer eigen-iteration? What does this imply about the
  payoff of a better preconditioner (\cref{ch:multigrid}) for the eigenmode driver specifically?
  \item \emph{(Reverse communication.)} Explain why ARPACK uses a reverse-communication interface
  rather than calling Palace's operator directly, and describe the request-and-resume cycle of
  \cref{fig:reversecomm}. In what sense is the library boundary a ``packaging detail, not an
  algorithm''? Give one concrete advantage of calling the library over writing the Krylov--Schur loop
  from scratch, and one concrete advantage of \emph{understanding} the loop even though you call the
  library.
\end{enumerate}

\chapter{Deflation and Nonlinear Eigenproblems}
\label{ch:deflation}

\chapterorient{The last two chapters found a \emph{few} eigenpairs of a giant
pencil. But the engineer rarely wants only the fundamental mode---they want the
fundamental \emph{and} its first several overtones, ten or twenty resonances
peeled off the bottom of the spectrum. How do you stop the iteration from
finding the same lowest mode over and over? The answer is \emph{deflation}: once
a pair has converged, project it out of the search so the iteration is forced to
find the \emph{next} one. That single idea---compute eigenpairs one at a time,
locking each as you go---carries us into the chapter's harder half. When the
cavity is filled with a lossy, dispersive material, the operator itself depends
on the frequency you are solving for: a \emph{nonlinear eigenproblem}
$\mat{T}(\lambda)\vx=\mathbf 0$ in which the resonance feeds back into the
equation that defines it. Linearization can no longer save us. We meet the
deflated quasi-Newton method Palace actually runs---Newton's method on the
eigenvalue, with each found mode deflated through an augmented system---and watch
its inner linear solve split into a tiny dense block and a huge sparse one.}

\section{Why one eigenpair is never enough}
\label{sec:def-why}

The Krylov--Schur machine of \cref{ch:krylovschur} already finds several
eigenpairs at once: its thick-restart loop locks each converged Ritz pair and
keeps iterating for the rest, and SLEPc's \code{EPSSolve} returns the requested
batch in one opaque call. So why a whole chapter on finding eigenpairs ``one at a
time''? Two reasons, and they pull in opposite directions until deflation
reconciles them.

The first is that \emph{locking is deflation}---the chapter is, in part, the
constructive account of a mechanism \cref{ch:krylovschur} used and named but did
not open. When Krylov--Schur ``locks'' a converged Ritz vector, it must thereafter
keep the growing Krylov basis \emph{orthogonal} to that vector, or the next
expansion step would simply re-grow the direction it just converged and the
solver would spin in place rediscovering the same mode. Locking, orthogonalizing
against the locked set, and ``purging'' a converged direction from the active
subspace are three names for one act: \emph{removing an already-found eigenvector
from the space the iteration is allowed to search}. That act is deflation, and
\cref{sec:def-locking} makes it explicit.

The second reason is sharper, and it is why deflation gets promoted from a
subroutine to a top-level algorithm. The clean black-box eigensolve of
\cref{ch:eigenvalues,ch:krylovschur} assumed a \emph{linear} or at worst
\emph{polynomial} pencil---$\Kop\vx=\lambda\Mop\vx$, or the quadratic
$(\Kop+\lambda\Cop+\lambda^2\Mop)\vx=\mathbf 0$ that linearization turns back into
a generalized linear problem (\cref{ch:eigenvalues}). For a genuinely
\emph{nonlinear} eigenproblem---one whose operator depends on $\lambda$ in a way
no polynomial captures---there is no linearization, no single library call that
returns a batch, and no Krylov basis whose Ritz values are the answers. The only
practical route known is to hunt the eigenpairs \emph{one at a time} with a
Newton-like iteration, deflating each as it converges. Deflation stops being a
bookkeeping detail inside the eigensolver and becomes the load-bearing
scaffolding of the whole method.

\begin{keyidea}
Deflation \emph{peels eigenpairs off the spectrum one at a time}: once a pair has
converged, it is projected out of the search space so the iteration is forced to
find the next one, never re-finding what it already has. For a linear pencil this
is the ``locking'' inside Krylov--Schur (\cref{ch:krylovschur}); for a
\emph{nonlinear} eigenproblem $\mat{T}(\lambda)\vx=\mathbf 0$---whose operator
\emph{depends on its own eigenvalue}---it becomes the top-level algorithm, because
no linearization and no single library call can return the modes as a batch.
\end{keyidea}

We take the two halves in order: first deflation on the familiar linear problem,
where every step can be checked against the power iteration of
\cref{ch:eigenvalues}; then the nonlinear eigenproblem and the deflated
quasi-Newton method Palace runs on it.

\section{Deflation on a linear problem}
\label{sec:def-linear}

Strip the problem to its bones. We have a symmetric matrix $\mat{A}$, we have
just found its dominant eigenpair $(\mu_1,\vu_1)$ by the power iteration of
\cref{ch:eigenvalues}, and we want the \emph{next} one, $(\mu_2,\vu_2)$. If we
simply restart power iteration from a fresh vector, it converges straight back to
$\vu_1$---the dominant direction is still dominant, and iteration has no memory of
having found it. We must change the problem so that $\vu_1$ is no longer there to
be found.

\subsection{Subtracting a rank-one piece}
\label{sec:def-ranköne}

For a symmetric matrix the cleanest construction is to \emph{subtract the found
eigenpair out of the operator}. If $\mat{A}$ is symmetric with orthonormal
eigenvectors, the dominant pair contributes a rank-one term $\mu_1\vu_1\vu_1^{\!\top}$
to its spectral decomposition $\mat{A}=\sum_i \mu_i\vu_i\vu_i^{\!\top}$. Remove
exactly that term:
\begin{equation}
  \boxed{\;\mat{A}_1 \;=\; \mat{A} \;-\; \mu_1\,\vu_1\vu_1^{\!\top}.\;}
  \label{eq:deflate-rank1}
\end{equation}
The deflated matrix $\mat{A}_1$ has \emph{exactly the same eigenvectors and
eigenvalues as $\mat{A}$, except that $\mu_1$ has been replaced by zero}. To see
it, apply $\mat{A}_1$ to each eigenvector. On $\vu_1$:
\begin{equation}
  \mat{A}_1\vu_1
    = \mat{A}\vu_1 - \mu_1\vu_1(\vu_1^{\!\top}\vu_1)
    = \mu_1\vu_1 - \mu_1\vu_1
    = \mathbf 0,
\end{equation}
using $\vu_1^{\!\top}\vu_1=1$: the dominant eigenvalue is annihilated, sent to
zero. On any other eigenvector $\vu_j$ ($j\neq1$), orthogonality
$\vu_1^{\!\top}\vu_j=0$ kills the correction term entirely:
\begin{equation}
  \mat{A}_1\vu_j
    = \mat{A}\vu_j - \mu_1\vu_1(\vu_1^{\!\top}\vu_j)
    = \mu_j\vu_j - \mathbf 0
    = \mu_j\vu_j .
\end{equation}
Every other pair survives untouched. So $\mat{A}_1$'s spectrum is
$\{0,\mu_2,\mu_3,\dots\}$: the former runner-up $\mu_2$ is now the largest in
magnitude (assuming $\abs{\mu_2}>0$), and \emph{power iteration on $\mat{A}_1$
converges to $\vu_2$}---the exact pair we wanted next. We found the second
eigenpair by finding the ``dominant'' one of a deflated operator.

\begin{figure}[t]
\centering
\begin{tikzpicture}
\begin{axis}[
    name=before,
    width=0.84\linewidth, height=30mm,
    xlabel={}, title={\footnotesize original spectrum of $\mat{A}$},
    title style={font=\footnotesize, yshift=-1.5mm},
    xmin=-0.6, xmax=9, ymin=0, ymax=1.5,
    axis lines=left, ytick=\empty, xtick={1,3,5,7.5},
    xticklabels={$\mu_4$,$\mu_3$,$\mu_2$,$\mu_1$},
    xticklabel style={font=\scriptsize}, clip=false,
  ]
  \addplot[ycomb, very thick, simBlue, mark=*, mark size=2pt] coordinates {
    (1,0.6)(3,0.85)(5,1.0)(7.5,1.25)};
  \node[font=\scriptsize\itshape, simRust, anchor=south] at (axis cs:7.5,1.25)
    {dominant: power iter.\ finds this};
\end{axis}
\begin{axis}[
    name=after, at={(before.south west)}, anchor=north west, yshift=-11mm,
    width=0.84\linewidth, height=30mm,
    title={\footnotesize after deflating $(\mu_1,\vu_1)$: spectrum of
      $\mat{A}_1=\mat{A}-\mu_1\vu_1\vu_1^{\!\top}$},
    title style={font=\footnotesize, yshift=-1.5mm},
    xmin=-0.6, xmax=9, ymin=0, ymax=1.5,
    axis lines=left, ytick=\empty, xtick={0,1,3,5},
    xticklabels={$0$,$\mu_4$,$\mu_3$,$\mu_2$},
    xticklabel style={font=\scriptsize}, clip=false,
  ]
  \addplot[ycomb, very thick, simBlue, mark=*, mark size=2pt] coordinates {
    (1,0.6)(3,0.85)(5,1.0)};
  \addplot[ycomb, very thick, simGray, mark=o, mark size=2pt] coordinates {(0,1.25)};
  \node[font=\scriptsize\itshape, simGray, anchor=south] at (axis cs:0,1.25)
    {$\mu_1\to 0$};
  \node[font=\scriptsize\itshape, simRust, anchor=south] at (axis cs:5,1.0)
    {now dominant: finds $\vu_2$};
\end{axis}
\end{tikzpicture}
\caption[Deflating one eigenpair sends it to zero]{Deflation by rank-one
subtraction. \emph{Top:} the original spectrum of $\mat{A}$; power iteration
finds the dominant pair $(\mu_1,\vu_1)$ at the right. \emph{Bottom:} subtracting
$\mu_1\vu_1\vu_1^{\!\top}$ replaces $\mu_1$ with $0$ (grey, open) and leaves every
other eigenvalue exactly where it was. The former runner-up $\mu_2$ is now the
largest in magnitude, so power iteration on the deflated operator $\mat{A}_1$
converges to $\vu_2$. Repeating---deflate, iterate, deflate---peels the spectrum
off one pair at a time.}
\label{fig:deflate-spectrum}
\end{figure}

\Cref{fig:deflate-spectrum} draws the move. Repeat it---deflate $(\mu_2,\vu_2)$ to
form $\mat{A}_2=\mat{A}_1-\mu_2\vu_2\vu_2^{\!\top}$, iterate for $\vu_3$, and so
on---and the spectrum is peeled off the top one eigenpair at a time. Each pair,
once found, is gone from the deflated operator and cannot be found again.

\subsection{Deflation as a projection: locking the search space}
\label{sec:def-locking}

The rank-one subtraction is concrete but special: it leans on symmetry (so that
the spectral decomposition has orthonormal $\vu_i$) and it modifies the operator.
The form that generalizes---and the form that matches what Krylov--Schur's locking
actually does---works on the \emph{search space} instead of the operator. Rather
than alter $\mat{A}$, we forbid the iteration from looking in the directions we
have already found.

Collect the converged eigenvectors as the columns of a tall thin matrix
$\mat{X}=[\,\vu_1\mid\cdots\mid\vu_k\,]$, the \emph{locked} set. Let $\mat{P}$ be
the orthogonal projector \emph{onto the complement} of their span:
\begin{equation}
  \mat{P} \;=\; \mat{I} - \mat{X}\mat{X}^{\!\top}
  \qquad(\text{with } \mat{X}^{\!\top}\mat{X}=\mat{I}_k).
  \label{eq:deflate-proj}
\end{equation}
Applying $\mat{P}$ to any vector strips out its component along the locked
directions: $\mat{P}\vu_i=\mathbf 0$ for each locked $\vu_i$, while $\mat{P}\vv=\vv$
for any $\vv$ already orthogonal to all of them. Now run the iteration on the
\emph{projected} operator $\mat{P}\mat{A}\mat{P}$, or---cheaper and equivalent for
iteration---apply $\mat{P}$ to the iterate after every step:
\begin{equation}
  \vv_{k+1} \;=\; \mat{P}\,\frac{\mat{A}\vv_k}{\norm{\mat{A}\vv_k}}.
  \label{eq:deflate-iterate}
\end{equation}
Each step re-orthogonalizes the iterate against the locked set, continually
sweeping out any component that has drifted back into a found direction. The
iteration physically \emph{cannot} converge to a locked eigenvector, because that
direction is projected to zero the instant it appears. It is forced to find the
next pair in the orthogonal complement.

\begin{figure}[t]
\centering
\begin{tikzpicture}[scale=1.0]
  % the locked direction (horizontal), the complement (vertical)
  \draw[->, simGray] (-0.3,0) -- (5.2,0)
    node[right,font=\scriptsize]{$\vu_1$ (locked)};
  \draw[->, simGray] (0,-0.3) -- (0,3.6)
    node[above,font=\scriptsize]{complement $\vu_1^\perp$};
  % the locked subspace shaded
  \fill[simGray!12] (0,-0.25) rectangle (5.0,0.25);
  \node[font=\scriptsize\itshape, simGray] at (3.6,-0.5) {forbidden: $\mat{P}\vu_1=\mathbf 0$};
  % an iterate that has drifted, with a component along u1
  \draw[-{Stealth[length=2.4mm]}, very thick, simBlue] (0,0) -- (2.6,2.0);
  \node[font=\scriptsize, simBlue] at (2.85,2.15) {$\mat{A}\vv_k$};
  % its projection onto the complement (drop the horizontal part)
  \draw[-{Stealth[length=2.4mm]}, very thick, simRust] (0,0) -- (0,2.0);
  \node[font=\scriptsize, simRust, anchor=east] at (-0.1,2.0) {$\mat{P}\,\mat{A}\vv_k$};
  % the projection drop (dotted)
  \draw[densely dotted, simGray] (2.6,2.0) -- (0,2.0);
  \draw[densely dotted, simGray] (2.6,2.0) -- (2.6,0);
  % the swept-out component
  \draw[-{Stealth[length=2mm]}, simGreen, thick] (0.05,0.12) -- (2.5,0.12);
  \node[font=\scriptsize, simGreen] at (1.3,0.42) {swept out each step};
  % converged next eigenvector lies in complement
  \draw[-{Stealth[length=2.4mm]}, very thick, simViolet, densely dashed] (0,0) -- (0,3.0);
  \node[font=\scriptsize, simViolet, anchor=east] at (-0.1,3.0) {$\vu_2$ (next)};
\end{tikzpicture}
\caption[Deflation projects the iterate off the locked directions]{Deflation as a
projection. The found eigenvector $\vu_1$ is \emph{locked}; the projector
$\mat{P}=\mat{I}-\mat{X}\mat{X}^{\!\top}$ kills any component along it (grey band:
$\mat{P}\vu_1=\mathbf 0$). When the operator pushes the iterate $\mat{A}\vv_k$
(blue) partly back toward the locked direction, applying $\mat{P}$ drops that
component (green) and leaves only the part in the orthogonal complement
($\mat{P}\,\mat{A}\vv_k$, rust). The iteration therefore cannot reconverge to
$\vu_1$; it is steered into the complement, where the next eigenvector $\vu_2$
(violet, dashed) lives. This is exactly the ``locking'' Krylov--Schur applies to
its converged Ritz vectors (\cref{ch:krylovschur}).}
\label{fig:deflate-project}
\end{figure}

\Cref{fig:deflate-project} shows the geometry. This projection form is what
Krylov--Schur's locking does (\cref{ch:krylovschur}): the converged Ritz vectors
are held aside, and every new Arnoldi/Lanczos column is orthogonalized against
them before it joins the active basis. ``Keep the search space orthogonal to the
locked vectors'' and ``deflate the found eigenpairs'' are the same instruction.
The block form---locking several at once into $\mat{X}$ and projecting against the
whole block---is how a solver computes eigenpairs not just one at a time but in
\emph{groups}, deflating a converged batch and continuing for the next.

\begin{notationbox}
Two equivalent ways to deflate $k$ converged eigenvectors $\mat{X}=[\vu_1\mid
\cdots\mid\vu_k]$ from a linear problem:
\begin{itemize}[nosep]
  \item \textbf{Operator deflation} (symmetric case): subtract the found spectral
  mass, $\mat{A}_k=\mat{A}-\sum_{i=1}^k\mu_i\vu_i\vu_i^{\!\top}$, sending each found
  $\mu_i$ to $0$.
  \item \textbf{Search-space deflation} (general): project with
  $\mat{P}=\mat{I}-\mat{X}\mat{X}^{\!\top}$ after each step, forbidding the
  iteration from drifting back into $\operatorname{span}(\mat{X})$.
\end{itemize}
Both enforce one contract: an already-found eigenvector cannot be found again.
The projection form is the one that survives into the nonlinear method of
\cref{sec:def-nleps}, generalized to an \emph{oblique} (non-orthogonal) projector.
\end{notationbox}

\subsection{A worked deflation}
\label{sec:def-worked1}

The whole mechanism fits in a few lines of arithmetic on a $2\times2$ matrix,
where we can carry power iteration by hand and watch the second eigenpair appear
only after deflation.

\begin{workedexample}{Find the second eigenpair by deflation}{deflate2x2}
Take the symmetric matrix
\[
  \mat{A} = \begin{psmallmatrix} 3 & 1 \\ 1 & 3 \end{psmallmatrix},
\]
whose eigenpairs we happen to know exactly: $\mu_1=4$ with
$\vu_1=\tfrac1{\sqrt2}(1,1)^{\!\top}$ and $\mu_2=2$ with
$\vu_2=\tfrac1{\sqrt2}(1,-1)^{\!\top}$. Pretend we know neither and hunt them by
power iteration.

\smallskip
\emph{Step 1: power iteration finds the dominant pair.} Start from
$\vv_0=(1,0)^{\!\top}$. One multiply gives $\mat{A}\vv_0=(3,1)^{\!\top}$; normalize
to $\vv_1=(0.949,0.316)^{\!\top}$. Again: $\mat{A}\vv_1=(3.16,2.06)^{\!\top}$,
normalize to $\vv_2=(0.838,0.546)^{\!\top}$. Again:
$\mat{A}\vv_2=(3.06,2.48)^{\!\top}\to\vv_3=(0.776,0.630)^{\!\top}$. The iterate is
swinging toward $(0.707,0.707)^{\!\top}=\vu_1$; the Rayleigh quotient
$\rho(\vv_3)=\inner{\mat{A}\vv_3}{\vv_3}=3.98$ is closing on $\mu_1=4$. Take the
converged pair $(\mu_1,\vu_1)=(4,\tfrac1{\sqrt2}(1,1)^{\!\top})$.

\smallskip
\emph{Step 2: restarting plain power iteration is useless.} If we restart from a
new random vector, it converges right back to $\vu_1$---$\mu_1=4$ is still the
dominant eigenvalue and iteration has no memory. We must deflate.

\smallskip
\emph{Step 3: deflate the found pair.} Subtract the rank-one piece
\cref{eq:deflate-rank1}, with $\vu_1\vu_1^{\!\top}=\tfrac12
\begin{psmallmatrix}1&1\\1&1\end{psmallmatrix}$:
\[
  \mat{A}_1 = \mat{A} - 4\cdot\tfrac12\begin{psmallmatrix}1&1\\1&1\end{psmallmatrix}
    = \begin{psmallmatrix}3&1\\1&3\end{psmallmatrix}
      - \begin{psmallmatrix}2&2\\2&2\end{psmallmatrix}
    = \begin{psmallmatrix}1&-1\\-1&1\end{psmallmatrix}.
\]
Check the spectrum directly: $\mat{A}_1\vu_1=\begin{psmallmatrix}1&-1\\-1&1
\end{psmallmatrix}\tfrac1{\sqrt2}(1,1)^{\!\top}=\tfrac1{\sqrt2}(0,0)^{\!\top}
=\mathbf 0$---the dominant pair is annihilated. And
$\mat{A}_1\vu_2=\begin{psmallmatrix}1&-1\\-1&1\end{psmallmatrix}\tfrac1{\sqrt2}
(1,-1)^{\!\top}=\tfrac1{\sqrt2}(2,-2)^{\!\top}=2\vu_2$---the second pair survives
with eigenvalue $2$, untouched. The deflated spectrum is $\{0,2\}$.

\smallskip
\emph{Step 4: power iteration on $\mat{A}_1$ now finds $\vu_2$.} Start again from
$\vv_0=(1,0)^{\!\top}$. One multiply: $\mat{A}_1\vv_0=(1,-1)^{\!\top}$, normalize to
$(0.707,-0.707)^{\!\top}$---which is $\vu_2$ \emph{exactly}, in a single step,
because the deflated matrix has only one nonzero eigendirection left. The Rayleigh
quotient reads $\rho=2=\mu_2$. The second eigenpair fell out the moment we
deflated the first.

The lesson generalizes past this tiny case: deflation does not make the next
eigenpair easier to \emph{compute} (the per-step work is the same), it makes the
next eigenpair the one the iteration \emph{converges to at all}.
\end{workedexample}

\begin{figure}[t]
\centering
\begin{tikzpicture}[font=\footnotesize, node distance=7mm]
  % the spectrum being peeled, lowest first
  % five spectrum stems, lowest-first (drawn explicitly, not via \foreach)
  \draw[thick, simRust] (0,0) -- (0,1.4);   \fill[simRust] (0,1.4) circle (1.6pt);
  \node[font=\scriptsize, simRust, anchor=south] at (0,1.4) {$\lambda_1$};
  \draw[thick, simRust] (1.4,0) -- (1.4,1.4); \fill[simRust] (1.4,1.4) circle (1.6pt);
  \node[font=\scriptsize, simRust, anchor=south] at (1.4,1.4) {$\lambda_2$};
  \draw[thick, simBlue] (2.8,0) -- (2.8,1.4); \fill[simBlue] (2.8,1.4) circle (1.6pt);
  \node[font=\scriptsize, simBlue, anchor=south] at (2.8,1.4) {$\lambda_3$};
  \draw[thick, simGray] (4.2,0) -- (4.2,1.4); \fill[simGray] (4.2,1.4) circle (1.6pt);
  \node[font=\scriptsize, simGray, anchor=south] at (4.2,1.4) {$\lambda_4$};
  \draw[thick, simGray] (5.6,0) -- (5.6,1.4); \fill[simGray] (5.6,1.4) circle (1.6pt);
  \node[font=\scriptsize, simGray, anchor=south] at (5.6,1.4) {$\lambda_5$};
  \draw[->, simGray] (-0.5,0) -- (6.4,0) node[right,font=\scriptsize]{$\lambda$};
  % peeling annotations
  \draw[-{Stealth[length=2mm]}, simRust, thick] (0,-0.55) -- (0,-0.05);
  \node[font=\scriptsize, simRust] at (0,-0.85) {found \#1};
  \draw[-{Stealth[length=2mm]}, simRust, thick] (1.4,-0.55) -- (1.4,-0.05);
  \node[font=\scriptsize, simRust] at (1.4,-0.85) {found \#2};
  \draw[-{Stealth[length=2mm]}, simBlue, thick] (2.8,-0.55) -- (2.8,-0.05);
  \node[font=\scriptsize, simBlue] at (2.8,-0.85) {seeking \#3};
  \node[font=\scriptsize\itshape, simGray] at (5.0,-0.85) {not yet reached};
  % the "wanted band" bracket
  \draw[decorate, decoration={brace, amplitude=4pt}, simRust]
    (-0.15,1.75) -- (1.55,1.75);
  \node[font=\scriptsize, simRust] at (0.7,2.15) {deflated \& locked};
  \draw[decorate, decoration={brace, amplitude=4pt, mirror}, simGray]
    (2.65,1.75) -- (5.75,1.75);
  \node[font=\scriptsize, simGray] at (4.2,2.15) {still in the active search};
\end{tikzpicture}
\caption[Peeling the spectrum lowest-first]{Computing many eigenpairs one at a
time. The wanted eigenvalues, lowest first, are peeled off the spectrum: pairs
$\lambda_1,\lambda_2$ are converged, deflated, and locked (rust, bracketed left);
the iteration is now hunting $\lambda_3$ (blue) in the space orthogonal to the
locked set; the higher pairs remain in the active search (grey). Each conversion
from ``seeking'' to ``found'' adds a column to the locked block $\mat{X}$ and
shrinks the search space by one direction. This lowest-first peeling is exactly
the eigenmode driver's request: the fundamental resonance and its first several
overtones, in order.}
\label{fig:peel-spectrum}
\end{figure}

\Cref{fig:peel-spectrum} shows the campaign in the large: pairs converted from
``seeking'' to ``found,'' the locked block growing by a column each time, the
spectrum peeled lowest-first. That is the picture to hold as we turn to the case
where the operator itself is moving underneath us.

\section{The nonlinear eigenproblem}
\label{sec:def-nep}

Everything so far assumed the operator was a fixed pencil. We now break that
assumption, because the physics breaks it.

\subsection{When the material depends on frequency}
\label{sec:def-dispersive}

The pencils of \cref{ch:eigenvalues} came from materials with constant
properties: a permittivity $\varepsilon$ and permeability $\mu$ that are fixed
numbers, baked once into the mass and stiffness matrices. Many real materials are
not like that. A \emph{dispersive} medium has a permittivity that depends on
frequency, $\varepsilon(\omega)$; a \emph{lossy} medium has a complex, frequency
-dependent response that absorbs energy. Plasmas, dielectric resonators with
Debye or Lorentz relaxation, conductors at microwave frequencies, metamaterials---
in all of them the material's response to the field is a function of how fast the
field oscillates.

When you discretize Maxwell's equations in such a medium, the frequency
dependence does not vanish into constant matrices. It rides along into the
operator. Writing $\lambda$ for the (squared, complex) eigenfrequency as before,
the discrete eigenproblem becomes
\begin{equation}
  \boxed{\;\mat{T}(\lambda)\,\vx \;=\; \mathbf 0,\;}
  \label{eq:nep}
\end{equation}
where the \emph{operator} $\mat{T}(\lambda)$ \emph{itself depends on the
eigenvalue $\lambda$ we are solving for}, and depends on it nonlinearly---not as
the tidy polynomial $\Kop+\lambda\Cop+\lambda^2\Mop$, but through whatever
function $\varepsilon(\lambda)$, $\nu(\lambda)$ the material model dictates. A
typical form is
\begin{equation}
  \mat{T}(\lambda) \;=\; \Kop + \lambda\Cop + \lambda^2\Mop + \mat{A}_2(\lambda),
  \label{eq:nep-form}
\end{equation}
where $\mat{A}_2(\lambda)$ is the operator-valued nonlinear part carrying the
dispersion---a rational function of $\lambda$ for a Lorentz model, a square-root
or exponential for a more exotic one. This is a \emph{nonlinear eigenproblem}, or
NEP. The word ``nonlinear'' refers to the dependence on $\lambda$, not to any
nonlinearity in $\vx$: \cref{eq:nep} is still linear in the eigenvector $\vx$ for
each fixed $\lambda$.

\begin{physicsbox}
The nonlinear eigenproblem is the mathematics of a resonance that \emph{shapes its
own cavity}. In a dispersive medium the material's response---how strongly it
polarizes, how much it absorbs---depends on the oscillation frequency. But the
oscillation frequency is precisely the resonance $\lambda$ we are trying to find.
So the operator that determines the resonance depends on the resonance itself: a
feedback loop baked into the equation. A mode ringing at one frequency sees a
different material than a mode ringing at another, and the self-consistent
frequency---the one at which the cavity, filled with the material \emph{as it
behaves at that frequency}, actually resonates---is the eigenvalue of
$\mat{T}(\lambda)\vx=\mathbf 0$. The complex part of $\lambda$ encodes the loss:
its imaginary part is the oscillation, its real part the decay, and the two
together give the resonant frequency and the quality factor $Q$
(\cref{ch:reductions}).
\end{physicsbox}

\subsection{Why linearization fails}
\label{sec:def-linfails}

\Cref{ch:eigenvalues} met the \emph{quadratic} eigenproblem
$(\Kop+\lambda\Cop+\lambda^2\Mop)\vx=\mathbf 0$ and tamed it by
\emph{linearization}: introduce $\vz=\lambda\vx$, stack $(\vx,\vz)$, and trade the
quadratic $N\times N$ problem for a linear $2N\times 2N$ generalized
eigenproblem the black-box library can solve. The trick worked because the
$\lambda$-dependence was \emph{polynomial}: a degree-two polynomial in $\lambda$
linearizes to a pencil of twice the size, a degree-$d$ polynomial to a pencil of
$d$ times the size. The companion-matrix construction is exact.

For the genuinely nonlinear $\mat{A}_2(\lambda)$ there is no such construction.
You cannot stack a finite number of auxiliary unknowns to absorb a rational or
transcendental function of $\lambda$; the ``companion'' would need infinitely many
blocks. \Cref{fig:linear-vs-nep} contrasts the two situations. A polynomial of
degree $d$ has a finite linearization (a $dN\times dN$ pencil); a true
nonlinearity has none. Approximating $\mat{A}_2(\lambda)$ by a polynomial and
linearizing \emph{that} is a real technique (rational and polynomial
approximation methods underlie much of modern NEP solving), but it trades the
exact problem for an approximate one and inflates the dimension. Palace takes the
other road: confront the nonlinearity directly with a Newton iteration on
$\lambda$, and use deflation to get more than one eigenpair.

\begin{figure}[t]
\centering
\begin{tikzpicture}[font=\footnotesize]
  % ===== LEFT: linear/polynomial pencil — fixed operator =====
  \begin{scope}
    \node[font=\small\bfseries, simBlue] at (1.4,2.9) {linear / polynomial pencil};
    \node[stage, minimum width=30mm, minimum height=14mm] (op) at (1.4,1.4)
      {$\Kop+\lambda\Cop+\lambda^2\Mop$};
    \node[font=\scriptsize\itshape, simGray] at (1.4,0.55)
      {operator built from \emph{fixed} matrices};
    \node[data, minimum width=20mm] (lam) at (1.4,-0.4) {$\lambda$};
    % lambda only scales the fixed blocks
    \draw[flow, simGray] (lam) -- node[right,font=\scriptsize]{scales blocks} (op);
    \node[font=\scriptsize, simGreen, align=center] at (1.4,-1.4)
      {\emph{linearizes:} $2N\!\times\!2N$\\ pencil, one library call};
  \end{scope}
  % divider
  \draw[simGray!50, densely dashed] (3.6,-1.8) -- (3.6,3.1);
  % ===== RIGHT: nonlinear T(lambda) — operator depends on its own eigenvalue =====
  \begin{scope}[xshift=4.6cm]
    \node[font=\small\bfseries, simRust] at (1.7,2.9) {nonlinear $\mat{T}(\lambda)$};
    \node[stage, minimum width=34mm, minimum height=14mm, draw=simRust, fill=simBgRust]
      (opn) at (1.7,1.4) {$\mat{T}(\lambda)=\cdots+\mat{A}_2(\lambda)$};
    \node[data, minimum width=20mm] (lamn) at (1.7,-0.6) {$\lambda$};
    % feedback: lambda feeds the operator, the operator's eigenvalue IS lambda
    \draw[flow, simRust] (lamn.east) to[out=0,in=-60]
      node[right,font=\scriptsize,align=left,pos=0.4]{enters the\\operator} (opn.south east);
    \draw[backflow, simRust] (opn.south west) to[out=-120,in=180]
      node[left,font=\scriptsize,align=right,pos=0.6]{whose eigenvalue\\\emph{is} $\lambda$} (lamn.west);
    \node[font=\scriptsize, simRust, align=center] at (1.7,-1.7)
      {\emph{no linearization:} solve\\ $\mat{T}(\lambda)\vx=\mathbf 0$ directly};
  \end{scope}
\end{tikzpicture}
\caption[Linear pencil versus nonlinear $\mat{T}(\lambda)$]{The structural break.
\emph{Left:} a linear or polynomial pencil is built from \emph{fixed} matrices
$\Kop,\Cop,\Mop$; the eigenvalue $\lambda$ only \emph{scales} those blocks, so the
problem linearizes exactly to a larger pencil and the black-box library of
\cref{ch:krylovschur} solves it in one call. \emph{Right:} for a dispersive
material the operator $\mat{T}(\lambda)$ \emph{depends on $\lambda$ itself}
through the nonlinear part $\mat{A}_2(\lambda)$---a feedback loop, since the
$\lambda$ feeding the operator is the very eigenvalue we seek. No finite stacking
of auxiliary unknowns can absorb a non-polynomial dependence, so linearization
fails and we must attack $\mat{T}(\lambda)\vx=\mathbf 0$ directly.}
\label{fig:linear-vs-nep}
\end{figure}

\section{Deflated quasi-Newton for the NEP}
\label{sec:def-nleps}

Palace solves the nonlinear eigenproblem with a \emph{deflated quasi-Newton}
method---the scheme SLEPc's NEP solver uses, with minimality index one, following
Effenberger~\citep{effenberger2013}. The idea is exactly the two halves of this
chapter, joined: \emph{Newton's method to converge one eigenpair} of the
nonlinear system, and \emph{deflation to lock it and move to the next}. We build
it in that order.

\subsection{Newton on the eigenvalue}
\label{sec:def-newton}

Forget deflation for a moment and chase a single eigenpair of
$\mat{T}(\lambda)\vx=\mathbf 0$. The unknowns are the pair $(\vx,\lambda)$---an
$N$-vector and a scalar---and \cref{eq:nep} is $N$ equations. That is one equation
short of a square system: any scalar multiple of an eigenvector is also an
eigenvector, so $\vx$ is determined only up to scale, and we need a
\emph{normalization} to pin it down. Add one: fix the size and phase of $\vx$ by
requiring $\vc^{\!\top}\vx=1$ for some fixed vector $\vc$. Now we have $N+1$
equations in $N+1$ unknowns,
\begin{equation}
  \mat{F}(\vx,\lambda) \;=\;
  \begin{pmatrix} \mat{T}(\lambda)\,\vx \\[2pt] \vc^{\!\top}\vx - 1 \end{pmatrix}
  \;=\; \mathbf 0,
  \label{eq:nep-system}
\end{equation}
a square \emph{nonlinear system} we can hand to Newton's method---the same
quasi-Newton machinery that drove the nonlinear solves of \cref{ch:cg}, now with
the eigenvalue itself among the unknowns.

Newton's method linearizes $\mat{F}$ about the current guess and solves for the
correction. The Jacobian of \cref{eq:nep-system} with respect to
$(\vx,\lambda)$---differentiating each block---is the bordered matrix
\begin{equation}
  \mat{J}(\vx,\lambda) =
  \begin{pmatrix}
    \mat{T}(\lambda) & \mat{T}'(\lambda)\,\vx \\[2pt]
    \vc^{\!\top} & 0
  \end{pmatrix},
  \label{eq:nep-jacobian}
\end{equation}
where $\mat{T}'(\lambda)=\mathrm{d}\mat{T}/\mathrm{d}\lambda$ is the derivative of
the operator with respect to the eigenvalue---the term that carries the
nonlinearity into the Jacobian. The Newton step solves
$\mat{J}(\vx,\lambda)\,(\Delta\vx,\Delta\lambda)^{\!\top}=-\mat{F}(\vx,\lambda)$
and updates $(\vx,\lambda)\leftarrow(\vx,\lambda)+(\Delta\vx,\Delta\lambda)$. The
\emph{quasi}-Newton variant Palace uses replaces the exact Jacobian action with a
cheaper lagged approximation (and an inexact inner tolerance), in the same spirit
as the quasi-Newton methods of \cref{ch:cg}; the structure---residual, Jacobian
solve, update---is identical.

\begin{keyidea}
The nonlinear eigenproblem $\mat{T}(\lambda)\vx=\mathbf 0$ is turned into a square
\emph{nonlinear system} in the unknowns $(\vx,\lambda)$ by appending a
normalization $\vc^{\!\top}\vx=1$, and solved by \emph{Newton's method---with the
eigenvalue $\lambda$ itself as one of the unknowns Newton drives}. Each Newton
step evaluates a residual, solves a bordered linear system built from
$\mat{T}(\lambda)$ and its derivative $\mat{T}'(\lambda)$, and updates both the
eigenvector and the eigenvalue. ``Newton on the eigenvalue'' is the engine; we
deflate around it to get more than one mode.
\end{keyidea}

The next worked example makes the eigenvalue-Newton concrete on a scalar problem
small enough to compute by hand, where ``Newton on $\lambda$'' is visibly just
Newton's method on a scalar equation.

\begin{workedexample}{Newton on a scalar nonlinear eigenproblem}{nepscalar}
Strip the NEP to one dimension, so $\vx$ is a scalar $x$ and $\mat{T}(\lambda)$ is
a scalar function $T(\lambda)$. Then $\mat{T}(\lambda)x=0$ with $x\neq0$ forces
$T(\lambda)=0$: \emph{the nonlinear eigenproblem is just a root-find for
$\lambda$}, and Newton on the eigenvalue is plain scalar Newton. Take the
nonlinearity
\[
  T(\lambda) \;=\; \lambda^2 - g(\lambda),
  \qquad g(\lambda) = 2 + \tfrac{1}{2}\cos\lambda,
\]
a quadratic with a small dispersive wobble $g(\lambda)$ standing in for the
material's frequency dependence. (Were $g$ constant, this would be an ordinary
quadratic eigenvalue $\lambda=\sqrt g$; the $\cos\lambda$ is what makes it
genuinely nonlinear.) We seek the smallest positive root.

\smallskip
\emph{The Newton step.} With $T'(\lambda)=2\lambda+\tfrac12\sin\lambda$, the scalar
Newton update is
\[
  \lambda_{k+1} = \lambda_k - \frac{T(\lambda_k)}{T'(\lambda_k)}
    = \lambda_k - \frac{\lambda_k^2 - 2 - \tfrac12\cos\lambda_k}
                       {2\lambda_k + \tfrac12\sin\lambda_k}.
\]
This is exactly the bordered Newton step \cref{eq:nep-jacobian} in one dimension:
the ``Jacobian'' is the scalar $T'(\lambda)$, the residual is $T(\lambda)$, and
the normalization holds $x=1$ so $\Delta x=0$ and only $\lambda$ moves.

\smallskip
\emph{Iterate from $\lambda_0=1.5$.}
\[
  \begin{array}{c|ccc}
    k & \lambda_k & T(\lambda_k) & \lambda_{k+1} \\\hline
    0 & 1.5000 & 0.2147 - 0\!=\!-0.2147 & 1.5713 \\
    1 & 1.5713 & 0.0205 & 1.5651 \\
    2 & 1.5651 & 0.00018 & 1.56504 \\
    3 & 1.56504 & 1.4\!\times\!10^{-8} & 1.565036
  \end{array}
\]
(Here $T(1.5)=2.25-2-\tfrac12\cos1.5=0.25-0.0354=0.2146$, and the step
$1.5-0.2146/3.0046=1.5714$, matching the table to rounding.) The residual
$T(\lambda_k)$ is driven to zero \emph{quadratically}---roughly doubling its
correct digits each step, the signature of Newton convergence---and settles at the
self-consistent eigenvalue $\lambda^\star\approx1.5650$, the frequency at which
$\lambda^2$ equals the material response $g(\lambda)$ \emph{evaluated at that very
frequency}. The fixed point of the feedback loop of \cref{fig:linear-vs-nep} is a
root, and Newton finds it.

\smallskip
\emph{Reading the residual figure.} \Cref{fig:newton-residual} plots
$\abs{T(\lambda_k)}$ against the step $k$ on a log axis: a straight line
steepening downward, the quadratic plunge. In the full $N$-dimensional problem the
scalar division becomes the bordered linear solve \cref{eq:nep-jacobian}, but the
convergence story is the one on this page.
\end{workedexample}

\begin{figure}[t]
\centering
\begin{tikzpicture}
\begin{axis}[
  width=11cm, height=5.6cm,
  xlabel={quasi-Newton step $k$},
  ylabel={residual $\abs{T(\lambda_k)}$},
  ymode=log, ymin=1e-10, ymax=1,
  xmin=0, xmax=4,
  grid=both, grid style={simGray!18},
  xtick={0,1,2,3,4},
  label style={font=\footnotesize}, tick label style={font=\scriptsize},
  legend pos=north east, legend style={font=\scriptsize, fill=white, draw=simGray!40},
  legend cell align=left,
]
% the residual from the worked example: quadratic convergence
\addplot[simRust, very thick, mark=*, mark size=2pt] coordinates {
  (0,0.2147)(1,0.0205)(2,1.8e-4)(3,1.4e-8)};
\addlegendentry{$\abs{T(\lambda_k)}$ (Newton on $\lambda$)}
% reference: linear convergence (a shallower line) for contrast
\addplot[simBlue, very thick, densely dashed, mark=square, mark size=1.5pt] coordinates {
  (0,0.2147)(1,0.064)(2,0.019)(3,0.0058)(4,0.0017)};
\addlegendentry{linear convergence (for contrast)}
% tolerance line
\addplot[simGreen, densely dotted, thick, domain=0:4, samples=2] {1e-8};
\node[font=\scriptsize, simGreen, anchor=west] at (axis cs:0.1,3e-8) {tolerance};
\end{axis}
\end{tikzpicture}
\caption[Quasi-Newton drives the eigenvalue residual to zero]{Newton on the
eigenvalue, convergence history from \cref{wex:nepscalar}. The residual
$\abs{T(\lambda_k)}$ (rust) falls \emph{quadratically}---each step roughly squares
the previous error, so the count of correct digits doubles per step and the
residual plunges through the tolerance (green dotted) in three or four steps. The
dashed blue line shows merely \emph{linear} convergence for contrast: the
straight-line decay of a fixed-rate iteration, far slower. This quadratic plunge
is the payoff of solving the bordered Newton system \cref{eq:nep-jacobian} each
step rather than fixed-point iterating. In the full problem each rust point costs
one bordered linear solve; in the scalar case it costs one division.}
\label{fig:newton-residual}
\end{figure}

\subsection{Deflating through an extended problem}
\label{sec:def-extended}

Newton converges \emph{one} eigenpair. To get the next, we deflate---but the
rank-one operator subtraction of \cref{sec:def-ranköne} relied on symmetry and a
fixed operator, neither of which a nonlinear, non-symmetric $\mat{T}(\lambda)$
offers. The NEP deflation is instead built into the \emph{system Newton solves},
by augmenting it.

Suppose $k$ eigenpairs have already converged, collected as the columns of
$\mat{X}\in\Complex^{N\times k}$ (the locked invariant basis) together with a
small $k\times k$ matrix $\mat{H}$ recording how the operator acts \emph{within}
that found subspace (its Rayleigh-quotient block). Effenberger's construction
forms an \emph{extended} nonlinear eigenproblem of size $n+k$---the original $N$
unknowns of $\vx$ plus $k$ extra coordinates $\vx_2\in\Complex^k$ that measure the
new pair's overlap with the locked block:
\begin{equation}
  \widetilde{\mat{T}}(\lambda)
  \begin{pmatrix}\vx\\\vx_2\end{pmatrix}
  =\mathbf 0,
  \qquad
  \widetilde{\mat{T}}(\lambda) =
  \begin{pmatrix}
    \mat{T}(\lambda) & \mat{U}(\lambda) \\
    \mat{X}^{\!\mathsf H} & \mathbf 0
  \end{pmatrix},
  \label{eq:nep-extended}
\end{equation}
where the coupling block $\mat{U}(\lambda)=\mat{T}(\lambda)\,\mat{X}\,
(\lambda\mat{I}-\mat{H})^{-1}$ ties the new unknown to the locked pairs through a
\emph{Schur-modified} factor
\begin{equation}
  \boxed{\;\mat{S} \;=\; \lambda\mat{I} - \mat{H}\;}
  \label{eq:nep-schur}
\end{equation}
built from the eigenvalues of the already-found pairs (encoded in $\mat{H}$). The
construction has a clean property, the whole reason for the augmentation: the
extended problem \cref{eq:nep-extended} has \emph{exactly the eigenpairs of the
original $\mat{T}(\lambda)$ that are not already in $\mat{X}$}. The locked pairs
have been removed from the extended spectrum---deflated---so Newton on the extended
system cannot reconverge to them and is driven to a genuinely new pair. The bottom
block-row $\mat{X}^{\!\mathsf H}\vx=\mathbf 0$ is the familiar
``orthogonal-to-the-locked-set'' constraint of \cref{sec:def-locking}, now an
explicit equation in the augmented system rather than a projection applied by
hand; the Schur factor $\mat{S}=\lambda\mat{I}-\mat{H}$ is what makes the
deflation respect the nonlinear, $\lambda$-dependent structure.

\begin{figure}[t]
\centering
\begin{tikzpicture}[font=\footnotesize]
  % the extended (n+k) x (n+k) block matrix
  % big block T(lambda): n x n
  \draw[thick, fill=simBgBlue, draw=simBlue] (0,0) rectangle (3.0,3.0);
  \node[font=\small\bfseries, simBlue] at (1.5,1.6) {$\mat{T}(\lambda)$};
  \node[font=\scriptsize, simGray] at (1.5,1.0) {$n\times n$, sparse, huge};
  % coupling block U(lambda): n x k  (right strip)
  \draw[thick, fill=simBgRust, draw=simRust] (3.0,0) rectangle (3.9,3.0);
  \node[font=\scriptsize, simRust, rotate=90] at (3.45,1.5) {$\mat{U}(\lambda)$};
  % locked block X^H: k x n  (bottom strip)
  \draw[thick, fill=simBgGreen, draw=simGreen] (0,-0.9) rectangle (3.0,0.0);
  \node[font=\scriptsize, simGreen] at (1.5,-0.45) {$\mat{X}^{\!\mathsf H}$ \ (locked, $k\times n$)};
  % small dense block: k x k
  \draw[thick, fill=simBgGold, draw=simGold] (3.0,-0.9) rectangle (3.9,0.0);
  \node[font=\scriptsize, simGold] at (3.45,-0.45) {$\mathbf 0$};
  % labels / braces
  \draw[decorate, decoration={brace, amplitude=4pt}, simBlue] (-0.1,3.05) -- (2.95,3.05);
  \node[font=\scriptsize, simBlue] at (1.45,3.45) {original $n$ unknowns $\vx$};
  \draw[decorate, decoration={brace, amplitude=4pt}, simRust] (3.05,3.05) -- (3.9,3.05);
  \node[font=\scriptsize, simRust, anchor=west] at (4.05,3.05) {$k$ coords $\vx_2$};
  % the Schur factor callout
  \node[font=\scriptsize, simInk, align=left, anchor=west] at (4.6,1.6)
    {$\mat{U}(\lambda)=\mat{T}(\lambda)\,\mat{X}\,\mat{S}^{-1}$,\\[2pt]
     Schur factor $\mat{S}=\lambda\mat{I}-\mat{H}$,\\[2pt]
     $\mat{H}$: locked eigenvalues};
  % the n+k size label
  \draw[decorate, decoration={brace, amplitude=5pt, mirror}, simGray] (-0.25,-0.95) -- (-0.25,3.0);
  \node[font=\scriptsize, simGray, rotate=90, anchor=south] at (-0.65,1.05) {size $n+k$};
\end{tikzpicture}
\caption[The deflated extended problem of size $n+k$]{Deflation by augmentation
(Effenberger). To deflate the $k$ already-found pairs $\mat{X}$, the original
$n\times n$ nonlinear operator $\mat{T}(\lambda)$ (blue) is bordered into an
\emph{extended} problem of size $n+k$: a coupling column $\mat{U}(\lambda)$
(rust), the locked block $\mat{X}^{\!\mathsf H}$ (green, the
``orthogonal-to-the-found-set'' constraint of \cref{sec:def-locking} made an
explicit equation), and a small $k\times k$ corner. The coupling
$\mat{U}(\lambda)=\mat{T}(\lambda)\mat{X}\mat{S}^{-1}$ uses the
\emph{Schur-modified} factor $\mat{S}=\lambda\mat{I}-\mat{H}$ built from the
locked eigenvalues in $\mat{H}$, so the deflation respects the nonlinear
$\lambda$-dependence. The extended spectrum is exactly the original spectrum
\emph{minus} the locked pairs: Newton on the augmented system finds a genuinely
new eigenpair, never a re-found one. Each converged pair grows $\mat{X}$ by one
column and $\mat{H}$ by one row and column---$k\to k+1$.}
\label{fig:extended-problem}
\end{figure}

\Cref{fig:extended-problem} draws the augmented operator. The bookkeeping that
makes it an algorithm is the growth of the locked block: each time the inner
Newton iteration converges a pair $(\lambda,\vx)$, the eigenvector is normalized
and appended as a new column of $\mat{X}$, the matrix $\mat{H}$ is \emph{bordered}
by one new row and column recording the new pair, and $k$ increments. The next
inner solve runs against the now-larger extended problem, which has one more pair
deflated out. ``Minimality index one'' is the technical guarantee that this
augmentation stays as small as possible---it adds exactly $k$ extra coordinates
for $k$ found pairs, never more---so the extended problem never bloats. The locked
pairs are invariant: the outer loop only \emph{appends} to $\mat{X}$ and
\emph{borders} $\mat{H}$, never disturbing what is already there, which is the
structural reason the one-at-a-time peeling of \cref{fig:peel-spectrum} is
correct.

\begin{figure}[t]
\centering
\begin{tikzpicture}[font=\footnotesize, node distance=8mm]
  % outer loop: one-at-a-time over eigenpairs
  \node[stage, minimum width=26mm, minimum height=11mm] (seed)
    {\textbf{seed}\\[1pt]\scriptsize linear-eigensolve\\\scriptsize guess
     (\cref{ch:krylovschur})};
  \node[stage, minimum width=30mm, minimum height=13mm, right=14mm of seed] (newton)
    {\textbf{inner quasi-Newton}\\[1pt]\scriptsize residual $\rhd$ Jacobian solve\\
     \scriptsize $\rhd$ $\lambda$-correction $\rhd$ backtrack};
  \node[stage, minimum width=24mm, minimum height=11mm, right=14mm of newton] (deflate)
    {\textbf{deflate}\\[1pt]\scriptsize append col to $\mat{X}$,\\
     \scriptsize border $\mat{H}$, $k\!+\!1$};
  \node[product, right=12mm of deflate] (done) {all $nev$\\found?};
  \draw[flow] (seed) -- (newton);
  \draw[flow] (newton) -- node[above,font=\scriptsize]{converged} (deflate);
  \draw[flow] (deflate) -- (done);
  % the inner Newton self-loop
  \draw[backflow] (newton.north) to[out=120,in=60,looseness=4]
    node[above,font=\scriptsize]{not yet} (newton.north);
  % the outer deflation loop back to seed
  \draw[backflow] (done.south) |- node[pos=0.25,below,font=\scriptsize,simRust]
    {next pair, now deflated against the grown $\mat{X}$} (seed.south);
  % the extern callout: inner linear solve
  \node[extern, below=13mm of newton, align=center] (solve)
    {\scriptsize each Newton step: one\\\scriptsize bordered, deflated linear solve\\
     \scriptsize (\cref{fig:dense-sparse-split})};
  \draw[refedge] (solve.north) -- (newton.south);
\end{tikzpicture}
\caption[The deflated quasi-Newton loop, nested]{The full method as two nested
loops. The \emph{outer} loop peels eigenpairs one at a time
(\cref{fig:peel-spectrum}): \textbf{seed} a guess from the linear eigensolver of
\cref{ch:krylovschur}, run the \textbf{inner quasi-Newton} iteration to converge
one pair against the current deflation block $\mat{X}$, then \textbf{deflate}---
append the new eigenvector to $\mat{X}$, border $\mat{H}$, increment $k$---and loop
back for the next pair, now solving the extended problem with one more pair
removed. The inner loop is Newton on the eigenvalue (\cref{sec:def-newton}):
each step evaluates the deflated residual, solves the bordered Jacobian system
(dashed callout, opened in \cref{fig:dense-sparse-split}), corrects
$(\vx,\lambda)$, and backtracks for stability, repeating until the residual falls
below tolerance. The loop terminates when $nev$ pairs are locked. This is the
\code{QuasiNewtonSolver} Palace runs.}
\label{fig:nleps-loop}
\end{figure}

\Cref{fig:nleps-loop} assembles the whole method: an outer one-at-a-time
deflation loop wrapping an inner quasi-Newton iteration, with the deflation block
growing by a column between outer iterations. The outer loop is seeded from the
linear eigensolver of \cref{ch:krylovschur}---a cheap linear-pencil approximation
gives Newton its starting guess, which the inner iteration then corrects to a true
nonlinear eigenpair (the seed need not be accurate; Newton converges it).

\section{The inner solve: dense-small meets sparse-large}
\label{sec:def-innersolve}

Each quasi-Newton step solves the bordered, deflated linear system---the action of
the extended Jacobian \cref{eq:nep-jacobian} on the augmented unknowns. This
system has a structure worth seeing clearly, because it is where two of the
linear-algebra tools of this Part meet on the same step, each handling the part it
is built for.

The bordered system splits along the $n+k$ partition of
\cref{fig:extended-problem}. The \emph{large} block is the $n\times n$ action of
$\mat{T}(\lambda)$---sparse, with millions of rows, exactly the kind of operator
\cref{ch:krylov,ch:gmres} taught us to solve iteratively. The \emph{small} block
is the $k\times k$ Schur/Gram coordinate system---dense, but tiny, with $k$ at
most a few hundred even after many modes are found. These call for two completely
different solvers:
\begin{itemize}[nosep]
  \item the large sparse block is inverted by an iterative \code{ksp\_solve}
  (\cref{ch:krylov})---a preconditioned Krylov method, never a factorization,
  because the matrix is far too large to factor;
  \item the small dense block---the $k\times k$ Gram matrix $\mat{X}^{\!\mathsf H}
  \mat{X}$ and the Schur factor $\mat{S}=\lambda\mat{I}-\mat{H}$---is inverted by a
  direct \code{lu\_solve}: a dense LU factorization, which on a $k\times k$ matrix
  with $k\sim10$--$100$ is utterly negligible.
\end{itemize}
The block elimination that couples them is the standard Schur-complement reduction
of a bordered system (\cref{ch:gmres} met the same shape in the bordered GMRES
least-squares step): solve the big block, condense its contribution into the small
block, solve the small dense block for the $k$ coordinates, and back-substitute.
The expensive part---the one sparse solve---is exactly one \code{ksp\_solve}; the
deflation costs only the trivial dense work on top.

\begin{figure}[t]
\centering
\begin{tikzpicture}[font=\footnotesize, node distance=9mm]
  % the bordered system, drawn as two coupled solves
  \node[stage, minimum width=40mm, minimum height=22mm, draw=simBlue, fill=simBgBlue]
    (big) at (0,0) {};
  \node[font=\small\bfseries, simBlue] at (big.north) [yshift=-4mm] {large sparse block};
  \node[font=\scriptsize, simInk, align=center] at (big.center) [yshift=1mm]
    {$\mat{T}(\lambda)$, \ $n\times n$\\[1pt]$n\sim 10^6$, sparse};
  \node[extern, minimum width=30mm, below=2mm of big.south, anchor=north]
    (kspsolve) {\code{ksp\_solve} (\cref{ch:krylov})\\[1pt]\scriptsize preconditioned Krylov,
      \emph{no} factorization};
  % the small dense block
  \node[stage, minimum width=26mm, minimum height=16mm, draw=simGold, fill=simBgGold,
    right=24mm of big] (small) {};
  \node[font=\small\bfseries, simGold] at (small.north) [yshift=-4mm] {small dense block};
  \node[font=\scriptsize, simInk, align=center] at (small.center) [yshift=1mm]
    {$\mat{S}=\lambda\mat{I}-\mat{H}$,\\[1pt]$\mat{X}^{\!\mathsf H}\mat{X}$, \ $k\times k$\\[1pt]
     $k\sim 10$--$100$, dense};
  \node[opbox, minimum width=24mm, below=2mm of small.south, anchor=north,
    draw=simGold] (lusolve) {\code{lu\_solve}\\[1pt]\scriptsize dense LU, negligible};
  % the Schur-complement coupling
  \draw[flow] (big.east) -- node[above,font=\scriptsize,align=center]
    {condense\\(Schur compl.)} (small.west);
  \draw[backflow] (small.south west) to[out=-150,in=-30]
    node[below,font=\scriptsize]{back-substitute the $k$ coords} (big.south east);
  % the cost annotation
  \node[font=\scriptsize\itshape, simRust, align=center] at (0,-3.4)
    {\textbf{the cost lives here:} one sparse solve/step};
  \node[font=\scriptsize\itshape, simGreen, align=center] at (5.0,-3.4)
    {\textbf{essentially free:}\\ a $k\times k$ factorization};
\end{tikzpicture}
\caption[The dense-small / sparse-large solve split]{One quasi-Newton step's
bordered linear solve, split by block. The deflated extended Jacobian
\cref{eq:nep-jacobian} partitions into a \emph{large sparse} block---the
$n\times n$ action of $\mat{T}(\lambda)$, with millions of unknowns, inverted by a
preconditioned iterative \code{ksp\_solve} (\cref{ch:krylov}; never factored)---and
a \emph{small dense} block---the $k\times k$ Schur factor
$\mat{S}=\lambda\mat{I}-\mat{H}$ and Gram matrix $\mat{X}^{\!\mathsf H}\mat{X}$,
inverted by a direct dense \code{lu\_solve}. The Schur-complement reduction
condenses the big block into the small one, solves the $k$ dense coordinates, and
back-substitutes. The entire deflation overhead is the negligible $k\times k$
factorization; the cost of the step is the single sparse solve it would have cost
without any deflation at all. Big-and-sparse iterative, small-and-dense direct:
the recurring division of labour of this Part.}
\label{fig:dense-sparse-split}
\end{figure}

\Cref{fig:dense-sparse-split} draws the division of labour. It is the same split
this Part keeps returning to---iterative Krylov for the huge sparse operator,
direct dense factorization for the tiny coordinate system---and the deflated
quasi-Newton step is where the two sit side by side in a single bordered solve, the
sparse \code{ksp\_solve} carrying essentially all the cost and the dense
\code{lu\_solve} carrying the deflation bookkeeping for free.

\section{Where this lands in Palace}
\label{sec:def-palace}

Palace's eigenmode driver is, by default, the clean linear path of
\cref{ch:eigenvalues,ch:krylovschur}: assemble $(\Kop,\Mop)$, choose a shift,
hand the pencil to SLEPc or ARPACK, read back the modes. The deflated quasi-Newton
method of this chapter is the \emph{nonlinear} branch of that same driver, taken
when the configuration declares frequency-dependent (lossy, dispersive)
materials---the case where the operator $\mat{T}(\lambda)$ genuinely depends on
$\lambda$ and no linearization exists.

In the source, this branch is a Palace-authored solver,
\code{QuasiNewtonSolver}, that drives the SLEPc-NEP-style deflation scheme
(minimality index one, Effenberger~\citep{effenberger2013}; the disguised
quasi-Newton variant of Jarlebring--Koskela--Mele). It is one of the few places
where the eigensolver's outer loop is \emph{Palace's} code rather than the
library's: where the linear eigensolve hands its whole iteration to an opaque
\code{EPSSolve} (\cref{ch:eigenvalues}), the nonlinear path runs the outer
deflation loop in Palace, calling the library only for the inner linear solves and
the linear-eigensolve seed. The structure of \cref{fig:nleps-loop}---outer
one-at-a-time deflation, inner quasi-Newton, the dense/sparse bordered solve---is
that source rendered as a composition, the outer sweep an explicit value-threaded
recursion over the growing locked block.

The contrast with the linear path is the contrast that organizes this whole Part.
The linear eigenmode solve is the fourth \emph{solve corner} of
\cref{ch:situating}: an opaque black-box call, no Palace loop. The nonlinear
eigenmode solve is a Palace-authored \emph{fold}---the same shape as the transient
time-march---that happens to wrap a black-box linear solve and a black-box linear
eigensolve inside each step. It is the deepest composition the eigenmode analysis
reaches: an outer Palace deflation loop, an inner Palace quasi-Newton loop, an
inner-inner sparse \code{ksp\_solve}, all to find the resonances of a cavity whose
material fights back at every frequency.

\subsection{Reading the results}
\label{sec:def-results}

The output is the same shape as the linear eigenmode analysis, read the same way
(\cref{ch:eigenvalues,ch:reductions}). Each converged pair $(\lambda_i,\vx_i)$
carries a \emph{complex} eigenvalue---necessarily complex, because the dispersive
material is lossy. Its imaginary part gives the oscillation frequency and its real
part the decay rate, and the two together give the resonant frequency
$f_i=\Im\lambda_i/2\pi$ (up to convention) and the quality factor
\begin{equation}
  Q_i \;=\; \frac{\abs{\Im\lambda_i}}{2\,\abs{\Re\lambda_i}},
\end{equation}
exactly as the lossy quadratic case of \cref{ch:eigenvalues} did. The eigenvector
$\vx_i$ is the mode field. The reduction from complex eigenpair to
frequency-and-$Q$ is the output postprocessing of \cref{ch:reductions}; the
nonlinear path is one more producer feeding it. So the deflated quasi-Newton
method, for all its internal machinery, hands the engineer the same table the
black-box eigensolve does---resonant frequencies and quality factors---computed
for the materials the black-box path could not represent.

\begin{pitfall}
The nonlinear path's eigenvalues are \emph{always} complex and must be read as
such; do not discard the real part as numerical noise the way one might for a
nearly-lossless linear mode. In the NEP the real part of $\lambda$ \emph{is the
loss}---it is the physical content the dispersive material was added to capture,
and it is what distinguishes a sharp high-$Q$ resonance from a broad lossy one.
A second trap is the seed: the linear-eigensolve guess that starts each outer
iteration is computed from a \emph{constant-material} approximation, so it can sit
far from the true nonlinear eigenvalue. That is expected---Newton's quadratic
convergence (\cref{fig:newton-residual}) closes the gap---but it means the seed's
eigenvalue should never be reported as the answer; only the converged
post-Newton pair is physical.
\end{pitfall}

\summarysection
The engineer wants many resonances, not one, and \emph{deflation} is how an
eigensolver delivers them: once a pair has converged it is \emph{locked}---projected
out of the search space so the iteration is forced to find the next one, peeling
the spectrum off lowest-first. On a symmetric linear problem deflation is a
rank-one operator subtraction $\mat{A}_1=\mat{A}-\mu_1\vu_1\vu_1^{\!\top}$ that
sends the found eigenvalue to zero; in general it is a projection
$\mat{P}=\mat{I}-\mat{X}\mat{X}^{\!\top}$ against the locked block---the same
``locking'' Krylov--Schur applies to its converged Ritz vectors. When the material
is dispersive the operator depends on its own eigenvalue,
$\mat{T}(\lambda)\vx=\mathbf 0$, a \emph{nonlinear eigenproblem} that no
linearization can reduce to a pencil. Palace solves it by \emph{deflated
quasi-Newton}: append a normalization to make $(\vx,\lambda)$ a square nonlinear
system, run \emph{Newton on the eigenvalue}, and deflate each found pair through an
\emph{extended problem} of size $n+k$ with the Schur-modified factor
$\mat{S}=\lambda\mat{I}-\mat{H}$. Each Newton step solves a bordered system that
splits into a large sparse block (an iterative \code{ksp\_solve}) and a tiny dense
block (a direct \code{lu\_solve}). The result is a set of complex eigenpairs read,
like the linear path, as resonant frequencies and quality factors.

\furtherreading
The deflated quasi-Newton method and its minimality-preserving extended-problem
construction are due to Effenberger~\citep{effenberger2013}, whose paper is the
direct source for the SLEPc-NEP scheme Palace runs; it is the place to go for the
proof that the extended problem \cref{eq:nep-extended} has exactly the
non-locked eigenpairs and for the minimality-index machinery only sketched here.
Saad's eigenvalue text~\citep{saad2011eig} is the standard reference for deflation
on linear problems---the rank-one subtraction, the projection form, and their
numerical conditioning---and for the broader landscape of nonlinear eigenvalue
methods (contour integration, rational approximation) that Palace's Newton route
is one branch of. Stewart~\citep{stewart2001krylovschur} connects the locking of
the previous chapter to the deflation of this one: thick-restart Krylov--Schur's
treatment of converged Ritz vectors is deflation on the linear pencil, and reading
his account alongside \cref{sec:def-locking} shows the two chapters describing one
mechanism from two sides.

\exercisessection
\begin{enumerate}[nosep]
  \item \textbf{Deflation by hand.} For
  $\mat{A}=\begin{psmallmatrix}5&2\\2&5\end{psmallmatrix}$, find the dominant
  eigenpair $(\mu_1,\vu_1)$ by inspection (the eigenvectors are
  $(1,1)^{\!\top}/\sqrt2$ and $(1,-1)^{\!\top}/\sqrt2$). Form the deflated matrix
  $\mat{A}_1=\mat{A}-\mu_1\vu_1\vu_1^{\!\top}$ explicitly and verify its spectrum is
  $\{0,\mu_2\}$. What does one step of power iteration on $\mat{A}_1$ from
  $(1,0)^{\!\top}$ return?
  \item \textbf{Projection equals subtraction.} Show that for a symmetric
  $\mat{A}$ with a single locked unit eigenvector $\vu_1$, iterating with the
  projected operator $\mat{P}\mat{A}\mat{P}$ (where
  $\mat{P}=\mat{I}-\vu_1\vu_1^{\!\top}$) produces the same sequence of directions as
  iterating with the deflated operator $\mat{A}-\mu_1\vu_1\vu_1^{\!\top}$, when both
  start from a vector already orthogonal to $\vu_1$. Why does the projection form
  not require knowing $\mu_1$?
  \item \textbf{Why orthogonality is mandatory.} Suppose you deflate with a stored
  eigenvector $\tilde\vu_1$ that is slightly \emph{inaccurate}---off the true
  $\vu_1$ by a small angle. Argue that the deflated iteration will then slowly
  re-grow a component of $\vu_1$ and can drift back toward the first eigenpair.
  What does this say about how accurately a pair must converge before it is safe to
  lock? (This is the practical reason locking tolerances exist in Krylov--Schur,
  \cref{ch:krylovschur}.)
  \item \textbf{Scalar Newton-on-$\lambda$.} Repeat \cref{wex:nepscalar} for
  $T(\lambda)=\lambda^2-(3+\sin\lambda)$ starting from $\lambda_0=2$. Carry three
  Newton steps and tabulate $\lambda_k$ and $\abs{T(\lambda_k)}$. Confirm the
  residual roughly squares each step (quadratic convergence). What would
  \emph{linear} convergence at rate $\tfrac14$ look like in the same table?
  \item \textbf{The normalization equation.} The NEP system
  \cref{eq:nep-system} appends $\vc^{\!\top}\vx=1$ to make the problem square.
  Explain why \emph{some} normalization is necessary (what goes wrong with the
  $N$ equations $\mat{T}(\lambda)\vx=\mathbf 0$ alone?). Does the choice of $\vc$
  affect \emph{which} eigenpair Newton converges to, or only the \emph{scale} of
  the eigenvector it returns? Give a sentence of reasoning for each.
  \item \textbf{Linearization's limit.} The quadratic pencil
  $(\Kop+\lambda\Cop+\lambda^2\Mop)\vx=\mathbf 0$ linearizes to a $2N\times2N$
  problem. A cubic pencil (degree three in $\lambda$) linearizes to what size? Now
  consider $\mat{T}(\lambda)=\Kop+\Mop\,e^{-\lambda\tau}$ (a delay term, genuinely
  transcendental). Explain why \emph{no} finite linearization exists, and what
  this forces the solver to do instead.
  \item \textbf{Dense-small versus sparse-large.} In the bordered solve of
  \cref{fig:dense-sparse-split}, the large block is $n\times n$ with
  $n\sim10^6$ and the small block is $k\times k$ with $k\sim50$. Estimate the
  ratio of arithmetic cost between a single dense LU factorization of the small
  block ($\sim\tfrac23 k^3$ operations) and a single iterative solve of the large
  block (take it as $\sim100$ matrix-vector products at $\sim10n$ operations each).
  Which dominates, and by how many orders of magnitude? What does your answer say
  about whether deflation ``slows down'' each Newton step?
  \item \textbf{Block deflation.} The chapter mostly deflates one pair at a time,
  but \cref{sec:def-locking} noted a block form that locks several at once into
  $\mat{X}$. Sketch how the extended problem \cref{eq:nep-extended} and its growth
  rule (append a column to $\mat{X}$, border $\mat{H}$) would change if the inner
  solver converged a \emph{block} of $b$ pairs simultaneously instead of one. What
  property of the nonlinear deflation makes one-at-a-time the natural default for
  the NEP, even though linear Krylov--Schur deflates in blocks? (Hint: re-read the
  note on minimality index one and the dependence of each pair's deflation on the
  prior pair being fully converged.)
\end{enumerate}

\chapter{Time Integration}
\label{ch:time}

% local: velocity unknown in the first-order recast (no \vv in simbook)
\providecommand{\vv}{\mathbf{v}}

\chapterorient{Of the six analyses, the transient one is the odd member out: it
does not solve a problem, it \emph{marches} one. Where electrostatics maps a
family of independent solves and the driven analysis sweeps a band of independent
frequencies, the transient analysis threads a single field-state forward in time,
each step beginning exactly where the last one ended. This chapter is the
machinery of that march. We recast the second-order semi-discrete system as a
first-order one so off-the-shelf integrators apply; we explain why
electromagnetics almost always wants an \emph{implicit} integrator, and why each
implicit step therefore hides a linear solve; we present the integrators Palace
uses without drowning in their tableaux; and---the structural payoff---we show
that the whole march is one instance of the \code{fold\_solve} combinator of
\cref{ch:combinators}, with a single clean algebraic law that makes checkpointing
and restart fall out for free.}

\section{From a second-order system to a first-order march}
\label{sec:time-firstorder}

The transient analysis begins where \cref{ch:reductions-pde} left it. Discretizing
the time-domain wave equation \emph{in space only}---the method of lines---turned
Maxwell's time-domain system into a large \emph{second-order} system of ordinary
differential equations in the degree-of-freedom vector $\vs(t)$:
\begin{equation}
  \Mop\,\ddot\vs + \Cop\,\dot\vs + \Kop\,\vs = \vJ(t),
  \label{eq:time-secondorder}
\end{equation}
with $\Mop$ the $\varepsilon$-weighted mass, $\Cop$ the $\sigma$-weighted damping,
$\Kop$ the curl--curl stiffness, and $\vJ(t)$ the discretized forcing
$-\partial_t\vJ_{\text{src}}$ (\cref{sec:transient}). This is the equation of a
damped oscillator, with $\Mop$ the inertia, $\Kop$ the restoring stiffness, and
$\Cop$ the energy-bleeding damping---the mechanical analogy is exact. Our task in
this chapter is to produce the field history $\vs(t_0), \vs(t_1), \vs(t_2),\dots$
by advancing \cref{eq:time-secondorder} one time step at a time.

\subsection{Why first order}

Standard time integrators---the routines a numerical library ships---are written
for \emph{first-order} systems of the form
\begin{equation}
  \dot\vy = \vF(\vy, t),
  \label{eq:time-firstorder-abstract}
\end{equation}
``the rate of change of the state is a known function of the state and the time.''
A second-order system does not match this shape: it involves $\ddot\vs$, a
\emph{second} derivative. The remedy is the standard one, and it costs nothing but
bookkeeping. Carry the velocity as an extra unknown. Define
$\vv \equiv \dot\vs$ and stack the position and velocity into a \emph{doubled}
state vector $\vy = (\vs, \vv)$. Then \cref{eq:time-secondorder} becomes a pair of
first-order equations: the top one is the \emph{definition} $\dot\vs = \vv$, and
the bottom one is \cref{eq:time-secondorder} solved for $\ddot\vs = \dot\vv$:
\begin{equation}
  \frac{d}{dt}\begin{pmatrix}\vs\\\vv\end{pmatrix}
  =
  \underbrace{\begin{pmatrix}
    \vv\\
    \Mop^{-1}\bigl(\vJ(t) - \Cop\,\vv - \Kop\,\vs\bigr)
  \end{pmatrix}}_{\displaystyle \vF(\vy, t)},
  \qquad
  \vy = \begin{pmatrix}\vs\\\vv\end{pmatrix}.
  \label{eq:time-recast}
\end{equation}
This \emph{is} the first-order form \cref{eq:time-firstorder-abstract}: the doubled
state $\vy$ has rate $\vF(\vy, t)$, a known function of $\vy$ and $t$. The price is
that the state vector is now twice as long; the prize is that every integrator
ever written for $\dot\vy = \vF(\vy, t)$ now applies unchanged.
\Cref{fig:recast} draws the recast: one second-order box becomes a two-component
first-order box, with the velocity threaded between the two halves.

\begin{figure}[t]
\centering
\begin{tikzpicture}[font=\footnotesize, node distance=8mm]
  % --- LEFT: the second-order system ---
  \node[stage, minimum width=42mm, minimum height=15mm, align=center] (so) at (0,0)
    {\textbf{second order}\\[2pt]$\Mop\ddot\vs + \Cop\dot\vs + \Kop\vs = \vJ(t)$};
  \node[font=\scriptsize\itshape, text=simRust, above=1mm of so]
    {one $\ddot\vs$: library integrators don't fit};
  % arrow
  \node[right=14mm of so] (mid) {};
  \draw[flow] (so.east) -- node[above, font=\scriptsize]{$\vv \equiv \dot\vs$}
    node[below, font=\scriptsize]{recast} ($(so.east)+(2.0,0)$);
  % --- RIGHT: the first-order doubled system ---
  \node[stage, minimum width=30mm, minimum height=9mm, align=center,
        fill=simBgGreen, draw=simGreen] (s) at (5.7,0.7) {$\dot\vs = \vv$};
  \node[stage, minimum width=44mm, minimum height=9mm, align=center,
        fill=simBgGreen, draw=simGreen] (v) at (6.6,-0.8)
    {$\dot\vv = \Mop^{-1}(\vJ - \Cop\vv - \Kop\vs)$};
  \node[font=\scriptsize\itshape, text=simGreen, above=1mm of s]
    {first order in $\vy = (\vs,\vv)$};
  % thread velocity from top into bottom and stiffness from bottom into top
  \draw[refedge] (s.south) -- node[left, font=\scriptsize]{$\vs$} (v.north west);
  \draw[backflow] (v.north east) .. controls +(0,8mm) and +(0,-8mm) ..
    (s.east) node[pos=0.5, right, font=\scriptsize]{$\vv$};
\end{tikzpicture}
\caption[The second-order to first-order recast]{The recast of
\cref{sec:time-firstorder}. The semi-discrete transient system is second order in
$\vs$ (left), so no off-the-shelf integrator---all written for $\dot\vy =
\vF(\vy,t)$---fits it directly. Introducing the velocity $\vv \equiv \dot\vs$ as a
second unknown and stacking $(\vs, \vv)$ into a doubled state turns it into the
first-order pair (right): the top equation is the definition of $\vv$, the bottom
is the original system solved for $\dot\vv$. The doubled state is twice as long;
in exchange every first-order integrator now applies. This recast is the first
thing Palace's \code{TimeDependentFirstOrderOperator} does.}
\label{fig:recast}
\end{figure}

\begin{notationbox}
Throughout this chapter $\vy = (\vs, \vv)$ is the \emph{doubled} first-order state
(position and velocity); $\vF(\vy, t)$ is its rate \cref{eq:time-recast}; $\Delta
t$ is the time step; and $\vy_n \approx \vy(t_n)$ is the numerical state at time
$t_n = t_0 + n\,\Delta t$. We write $\lambda$ for a scalar decay rate when we study
a single mode in isolation. The three operators $\Mop, \Cop, \Kop$ are the
assembled mass, damping, and stiffness of \cref{ch:reductions-pde}, held fixed for
the whole march.
\end{notationbox}

Palace performs exactly this recast: it builds a
\code{TimeDependentFirstOrderOperator} encoding $\vF(\vy, t)$ from the three
assembled matrices and hands it to a library integrator constructed once, outside
the time loop.\footnote{The first-order operator is assembled from $\Mop,\Cop,\Kop$
and the integrator chosen at construction
(\code{palace/models/timeoperator.cpp:311-335}); the loop that advances it lives in
the driver (\code{palace/drivers/transientsolver.cpp:77}). We return to the loop in
\cref{sec:fold-march}.} Everything that follows is about \emph{what one step of
that integrator does}, and how the steps chain together.

\section{Explicit and implicit steps, and why electromagnetics needs implicit}
\label{sec:explicit-implicit}

An integrator advances the state from $\vy_n$ to $\vy_{n+1}$ over one step
$\Delta t$. The simplest possible rule is \emph{forward (explicit) Euler}: read the
rate at the current state and step along it,
\begin{equation}
  \vy_{n+1} = \vy_n + \Delta t\,\vF(\vy_n, t_n).
  \label{eq:fwd-euler}
\end{equation}
Everything on the right is \emph{known}---it is built entirely from $\vy_n$, the
state we already hold---so $\vy_{n+1}$ is computed by a single evaluation of
$\vF$. No system is solved; the step is cheap. This is the defining feature of an
\emph{explicit} method: the new state is an explicit formula in the old one.

The alternative is to evaluate the rate at the \emph{new}, as-yet-unknown state.
\emph{Backward (implicit) Euler} does exactly this:
\begin{equation}
  \vy_{n+1} = \vy_n + \Delta t\,\vF(\vy_{n+1}, t_{n+1}).
  \label{eq:bwd-euler}
\end{equation}
Now $\vy_{n+1}$ appears on \emph{both} sides. The rule no longer hands us the
answer; it defines the answer \emph{implicitly}, as the solution of an equation.
For our linear transient system $\vF$ is affine in $\vy$, so
\cref{eq:bwd-euler} rearranges into a linear system in $\vy_{n+1}$,
\begin{equation}
  \bigl(\mat I - \Delta t\,\mat J\bigr)\,\vy_{n+1} = \vy_n + \Delta t\,\vg_{n+1},
  \label{eq:implicit-solve}
\end{equation}
where $\mat J$ is the (constant) Jacobian of $\vF$ and $\vg_{n+1}$ collects the
forcing. The shape is the headline of the section: \textbf{every implicit step is a
linear solve}. The matrix $\mat I - \Delta t\,\mat J$---built from the same
$\Mop, \Cop, \Kop$ that assembled the operator---must be inverted (in the
matrix-free sense of \cref{ch:krylov}: solved, not literally inverted) once per
step. That solve is a \code{ksp\_solve}, the Krylov solve of
\cref{ch:krylov,ch:cg,ch:gmres}, sitting inside the time loop.

\begin{keyidea}
An \emph{explicit} step evaluates the rate at the state it already knows: it is a
formula, cheap, no solve. An \emph{implicit} step evaluates the rate at the
unknown new state: it is an equation, and solving it for a linear system means a
linear solve every step---a \code{ksp\_solve} (\cref{ch:krylov}) nested inside the
time march. Implicit steps cost more per step; the next subsection shows why
electromagnetics pays the cost anyway.
\end{keyidea}

\subsection{Stiffness: the explicit step's fatal step-size limit}

If implicit steps are so much more expensive, why use them? Because the explicit
step's cheapness is illusory for the systems we integrate. The spatial operator
$\Kop$, assembled from a fine finite-element mesh, has an enormous spread of
eigenvalues: the smallest correspond to the slow physical modes we care about, the
largest to fast modes living on the smallest mesh elements. A system with such a
spread is called \emph{stiff}, and stiffness is poison to explicit methods.

The cleanest way to see why is to watch a single mode. Diagonalize the system and
each mode evolves independently as a scalar decay $\dot y = -\lambda y$, with
$\lambda > 0$ the mode's decay rate; the stiff modes are the ones with very large
$\lambda$. Apply forward Euler \cref{eq:fwd-euler} to this scalar equation:
\begin{equation}
  y_{n+1} = y_n + \Delta t\,(-\lambda y_n) = (1 - \lambda\,\Delta t)\,y_n.
  \label{eq:fe-scalar}
\end{equation}
Each step multiplies $y$ by the \emph{amplification factor} $g = 1 - \lambda\Delta
t$. The true solution $y(t) = y_0 e^{-\lambda t}$ decays, so a faithful numerical
solution must \emph{also} decay: we need $|g| < 1$. That demands
\begin{equation}
  |1 - \lambda\,\Delta t| < 1
  \qquad\Longleftrightarrow\qquad
  \Delta t < \frac{2}{\lambda}.
  \label{eq:fe-stability}
\end{equation}
The step is bounded by $2/\lambda$, and $\lambda$ here is the \emph{largest}
eigenvalue in the system. One fast mode on one tiny element---a mode physically
irrelevant to the answer---pins $\Delta t$ to a minuscule value, and the explicit
march must take an absurd number of tiny steps to cross a physically interesting
interval. This is the stiffness catastrophe.

Backward Euler suffers no such limit. Apply \cref{eq:bwd-euler} to $\dot y =
-\lambda y$:
\begin{equation}
  y_{n+1} = y_n + \Delta t\,(-\lambda y_{n+1})
  \quad\Longrightarrow\quad
  y_{n+1} = \frac{1}{1 + \lambda\,\Delta t}\,y_n,
  \qquad g = \frac{1}{1 + \lambda\,\Delta t}.
  \label{eq:be-scalar}
\end{equation}
For \emph{every} $\lambda > 0$ and \emph{every} $\Delta t > 0$, the denominator
exceeds $1$, so $|g| < 1$: the numerical solution decays no matter how large the
step. The fast modes are damped, the slow modes are tracked, and $\Delta t$ is
chosen for \emph{accuracy on the modes we care about}, not forced small by modes we
do not. This unconditional stability is why electromagnetic transient solvers are
almost always implicit, and it is what \cref{wex:stiff} makes concrete with
numbers.

\begin{keyidea}
The assembled spatial operator is \emph{stiff}: its largest eigenvalue is huge, set
by the smallest mesh element. An explicit step is stable only for $\Delta t <
2/\lambda_{\max}$---a fast, physically irrelevant mode pins the step size to a
crawl. An implicit (backward-Euler-type) step is stable for \emph{any} $\Delta t$,
so the step is chosen for accuracy, not stability. The per-step linear solve buys
freedom from the stiffness limit. That trade---a solve per step in exchange for big
steps---is why electromagnetics integrates implicitly.
\end{keyidea}

\subsection{Stability regions and A-stability}

The scalar test $\dot y = -\lambda y$ generalizes into a standard diagnostic. Let
$\lambda$ range over the \emph{complex} plane (a mode may oscillate as well as
decay, $\lambda = a + ib$), feed $z = \lambda\,\Delta t$ into the method's
amplification factor $g(z)$, and ask: for which $z$ is $|g(z)| \le 1$? That set is
the method's \emph{stability region}---the values of $z$ for which the method does
not amplify. A method is useful for a given step exactly when every mode's $z$
lands inside it.

\Cref{fig:stability} draws the two Euler regions in the $z$-plane. Forward Euler's
region is a single disk of radius $1$ centered at $z = -1$: a bounded set, so
large $\lambda\Delta t$ falls outside it and blows up---the limit
\cref{eq:fe-stability} read geometrically. Backward Euler's region is the
\emph{entire complex plane outside} the disk of radius $1$ centered at $+1$, which
contains the whole left half-plane $\Re(z) < 0$ where the physical decaying modes
live. A method whose stability region contains the entire left half-plane is called
\emph{A-stable}: it is stable for every decaying mode at every step size. Backward
Euler is A-stable; forward Euler is not. A-stability is the precise property the
stiffness argument was groping toward, and it is the property every integrator in
this chapter possesses.

\begin{figure}[t]
\centering
\begin{tikzpicture}[font=\footnotesize]
  % ===== LEFT: forward Euler region (disk |1+z|<=1) =====
  \begin{scope}
    \node[font=\small\bfseries, text=simRust] at (0,2.5) {forward Euler};
    \draw[->, simGray] (-2.6,0) -- (1.4,0) node[right, font=\scriptsize]{$\Re z$};
    \draw[->, simGray] (0,-1.9) -- (0,1.9) node[above, font=\scriptsize]{$\Im z$};
    \fill[simRust!20, draw=simRust, thick] (-1,0) circle (1);
    \node[font=\scriptsize, text=simRust] at (-1,0) {stable};
    \node[font=\scriptsize, text=simGray] at (-1,-1.35) {disk, $|1+z|\le 1$};
    \node[font=\scriptsize, text=simRust!80!black, align=center] at (0.55,1.45)
      {outside:\\blows up};
  \end{scope}
  \draw[simGray!50, thick] (2.0,-2.0) -- (2.0,2.7);
  % ===== RIGHT: backward Euler region (everything outside |1-z|<1) =====
  \begin{scope}[xshift=5.4cm]
    \node[font=\small\bfseries, text=simGreen] at (0,2.5) {backward Euler};
    % shade the whole left half (the important part) lightly green
    \fill[simBgGreen] (-2.6,-1.9) rectangle (0,1.9);
    \draw[->, simGray] (-2.6,0) -- (2.0,0) node[right, font=\scriptsize]{$\Re z$};
    \draw[->, simGray] (0,-1.9) -- (0,1.9) node[above, font=\scriptsize]{$\Im z$};
    % the UNSTABLE disk |1-z|<1 centered at +1
    \fill[white, draw=simGreen, thick, dashed] (1,0) circle (1);
    \node[font=\scriptsize, text=simGray, align=center] at (1.0,0) {un-\\stable};
    \node[font=\scriptsize, text=simGreen, align=center] at (-1.4,1.0)
      {stable\\(all of $\Re z<0$)};
    \node[font=\scriptsize, text=simGreen!60!black] at (-1.3,-1.55) {A-stable};
  \end{scope}
\end{tikzpicture}
\caption[Stability regions of the two Euler methods]{The stability regions in the
$z = \lambda\,\Delta t$ plane (\cref{sec:explicit-implicit}). \emph{Forward Euler}
(left) is stable only on the disk $|1+z|\le 1$---a \emph{bounded} set, so a stiff
mode with large $\lambda\Delta t$ falls outside and the step blows up. \emph{Backward
Euler} (right) is stable \emph{everywhere except} the small disk $|1-z|<1$; its
region contains the entire left half-plane $\Re z<0$ where every physically decaying
mode lives. A method whose region contains all of $\Re z<0$ is \emph{A-stable}.
Backward Euler is A-stable and forward Euler is not---the geometric statement of
why the implicit step has no step-size limit.}
\label{fig:stability}
\end{figure}

\subsection{Accuracy: order, and the truncation error}

Stability says the march does not blow up; \emph{accuracy} says it lands near the
true trajectory. The standard measure is the \emph{order} $p$ of a method: the
local truncation error---the error committed in a single step from exact data---is
$O(\Delta t^{p+1})$, and the global error accumulated over a fixed interval is
$O(\Delta t^{p})$. A first-order method ($p=1$, like both Eulers) halves its global
error when you halve the step; a second-order method ($p=2$) quarters it. Higher
order buys accuracy faster, at the cost of more work or more storage per step.

Order and stability are \emph{independent} axes, and a good integrator wants both:
high order for accuracy and A-stability for the stiff transient system. The plain
Eulers give one or the other---backward Euler is A-stable but only first order;
nothing here is high order. The integrators Palace actually uses, next, are the
ones engineered to be \emph{both} A-stable \emph{and} second order, with one extra
feature the plain methods lack.

\section{The integrators Palace uses}
\label{sec:integrators}

Palace offers a small menu of implicit integrators, selected at construction. We
present the \emph{idea} of each---what property it is engineered for---without
reproducing the Butcher tableaux that specify their coefficients; the tableaux are
in any numerical-methods reference, and the ideas are what matter for reading the
spec.

\subsection{Generalized-$\alpha$: tunable numerical damping}

The workhorse is the \emph{generalized-$\alpha$} method. It is second-order
accurate and A-stable---the two properties the plain Eulers could not combine---and
it adds a third, decisive feature: \emph{tunable, frequency-selective numerical
damping}. A pure second-order A-stable method (the trapezoidal rule) tracks every
mode without damping, which sounds ideal but is not: the highest-frequency modes a
mesh can represent are numerical junk---unresolved, aliased, physically
meaningless---and an undamped method lets them oscillate forever, polluting the
solution. Generalized-$\alpha$ damps them on purpose.

The point is best seen in its \emph{spectral-radius} curve, \cref{fig:genalpha}.
The method exposes a single dial, a high-frequency spectral radius $\rho_\infty \in
[0,1]$, that sets how aggressively the infinitely-fast modes are damped while
leaving the low, resolved modes essentially untouched. Set $\rho_\infty = 1$ and no
mode is damped (the trapezoidal limit); set $\rho_\infty = 0$ and the
fastest modes are annihilated in one step (the maximally dissipative limit); in
between, the low modes pass with amplification near $1$ and the high modes are
selectively bled away. The dial lets the analyst keep the physics and discard the
numerical noise, with one number, without touching accuracy on the resolved
band.\footnote{Palace constructs a \code{GeneralizedAlphaSolver} (alongside
\code{SDIRK23Solver} and the SUNDIALS \code{ARKStepSolver}) at
\code{palace/models/timeoperator.cpp:311-335}; the choice is a configuration token,
the integrator library-owned.}

\begin{figure}[t]
\centering
\begin{tikzpicture}
\begin{axis}[
    width=0.74\linewidth, height=52mm,
    xlabel={mode frequency $\;\omega\,\Delta t$ (low $\to$ high)},
    ylabel={amplification $|g|$},
    xmin=0, xmax=40, ymin=0, ymax=1.10,
    axis lines=left, ytick={0,0.5,1},
    xtick={0,10,20,30,40},
    yticklabel style={font=\scriptsize}, xticklabel style={font=\scriptsize},
    xlabel style={font=\scriptsize}, ylabel style={font=\scriptsize},
    legend style={font=\scriptsize, draw=simGray!50, at={(0.98,0.97)},
      anchor=north east, legend cell align=left},
    every axis plot/.append style={very thick, samples=120, domain=0:40},
  ]
  % rho_inf = 1: trapezoidal, no damping -> |g| stays 1
  \addplot[simTeal] {1};
  \addlegendentry{$\rho_\infty=1$: no damping (trapezoidal)}
  % rho_inf = 0.5: high modes damped toward 0.5
  \addplot[simBlue] {0.5 + 0.5/(1 + (x/5)^2)};
  \addlegendentry{$\rho_\infty=0.5$: moderate damping}
  % rho_inf = 0: high modes annihilated
  \addplot[simRust] {1/(1 + (x/5)^2)};
  \addlegendentry{$\rho_\infty=0$: maximal damping}
  % mark the resolved band
  \draw[simGray, dashed] (axis cs:7,0) -- (axis cs:7,1.10);
  \node[font=\scriptsize, text=simGray, anchor=south west]
    at (axis cs:7.5,0.05) {resolved $\to$ unresolved};
\end{axis}
\end{tikzpicture}
\caption[Numerical damping of generalized-$\alpha$]{The frequency-selective damping
of generalized-$\alpha$ (\cref{sec:integrators}), drawn as the per-step
amplification $|g|$ against mode frequency. The single dial $\rho_\infty$ sets the
high-frequency behavior: at $\rho_\infty = 1$ (teal) no mode is damped, the
undamped trapezoidal limit; at $\rho_\infty = 0$ (rust) the fastest modes are
annihilated in one step; intermediate values (blue) bleed the high modes away while
leaving the low, \emph{resolved} modes (left of the dashed line) essentially
untouched at $|g|\approx 1$. This is the feature the plain Eulers lack: the analyst
keeps the resolved physics and discards the unresolved numerical noise with one
number.}
\label{fig:genalpha}
\end{figure}

\subsection{SDIRK and ARKStep}

The other two menu items are Runge--Kutta-family methods. \emph{SDIRK}---singly
diagonally implicit Runge--Kutta---takes a small number of internal stages per
step, each stage an implicit solve, but arranges the stages so they share a
\emph{single} matrix factorization (the ``singly diagonal'' structure means every
stage solves a system with the same $\mat I - \gamma\Delta t\,\mat J$, so the
preconditioner is built once and reused across stages). SDIRK methods reach second
order and beyond while staying A-stable, at the cost of several solves per step
rather than one---the price of higher accuracy bought stage by stage.
\emph{ARKStep} is the additive Runge--Kutta integrator from the SUNDIALS library,
designed to split a system into a stiff part it treats implicitly and a non-stiff
part it treats explicitly---an implicit-explicit (IMEX) method---and to adapt its
step size automatically to a requested error tolerance.

For our purposes the three integrators share the structure that matters and differ
only in details we do not need:

\begin{description}[style=nextline, leftmargin=0pt, font=\normalfont\itshape]
  \item[All are implicit.] Every one solves at least one linear system per step (or
  per stage)---a \code{ksp\_solve} nested in the march. None escapes the per-step
  solve, because none escapes A-stability, and A-stability for these methods
  requires the implicit evaluation.
  \item[All are A-stable.] Each is stable on the whole stiff transient system at a
  step chosen for accuracy, not pinned by the largest eigenvalue. This is the
  non-negotiable property; it is why the menu has no explicit entry.
  \item[They differ in order, stage count, and damping control.]
  Generalized-$\alpha$ offers the damping dial; SDIRK trades stages for order;
  ARKStep adapts the step and splits stiff from non-stiff. These are knobs on a
  shared frame, not different algorithms.
\end{description}

\begin{notationbox}
The key fact to carry forward is structural, not numerical: \emph{whichever}
integrator is chosen, one step is an opaque routine that consumes the current
state and the step size and returns the next state, performing one or more
\code{ksp\_solve}s internally. The framework does not look inside that routine---it
quantifies over it. That opacity is exactly what lets the next section treat
``one step'' as a single black box and focus on how the steps \emph{chain}.
\end{notationbox}

\section{The march is a state-threaded fold}
\label{sec:fold-march}

We now reach the chapter's structural payoff, the reason the transient analysis was
promised as the canonical example back in \cref{ch:fold} and given a combinator
name in \cref{ch:combinators}. Strip away the integrator's internals---treat one
step as an opaque operator \code{time\_step\_op} that advances the state---and look
at how the march is assembled from steps. It is a \emph{fold}.

\subsection{Why a fold and not a map}

Recall the fundamental split of \cref{ch:combinators}. A \code{map} runs a family
of \emph{independent} computations: no carry between elements, so they reorder and
parallelize. A \code{fold} threads a \emph{carry} through a chain: each step
consumes what the last produced, so the steps are strictly ordered. The transient
march is unambiguously the second. The field-state at $t_{n+1}$ is computed
\emph{from} the field-state at $t_n$ by one time step; step $n+1$ cannot begin
until step $n$ has finished, because step $n$'s output is step $n+1$'s input. The
state is threaded forward along a single timeline. This is the
\emph{state-threaded fold}, the combinator \code{fold\_solve}.

Contrast it, as \cref{ch:combinators} did, with the \code{solve\_family} of
electrostatics. There the several solves shared one operator and depended on
nothing but their own right-hand side---independent, a map, embarrassingly
parallel. Here the steps share one operator \emph{and} a threaded state, and the
state makes them sequential. The presence or absence of the carry is the entire
difference, and it is visible in the type. \Cref{fig:fold-vs-map} sets the two side
by side on the same timeline: the map fans out, the fold lines up.

\begin{figure}[t]
\centering
\begin{tikzpicture}[font=\footnotesize]
  % ===== LEFT: the transient FOLD (carry threaded) =====
  \begin{scope}
    \node[font=\small\bfseries, text=simGreen] at (2.4,2.5) {transient: a \code{fold} (carry threaded)};
    \node[data, minimum width=9mm] (s0) at (-0.3,1.1) {$\vy_0$};
    \foreach \i/\x in {1/1.1,2/2.6,3/4.1}{
      \node[opbox, minimum width=11mm] (g\i) at (\x,1.1) {step};
    }
    \node[data, minimum width=9mm] (s3) at (5.4,1.1) {$\vy_3$};
    \draw[flow] (s0.east) -- (g1.west);
    \draw[flow] (g1.east) -- node[above,font=\scriptsize]{$\vy_1$} (g2.west);
    \draw[flow] (g2.east) -- node[above,font=\scriptsize]{$\vy_2$} (g3.west);
    \draw[flow] (g3.east) -- (s3.west);
    % each step hides a ksp_solve
    \foreach \i/\x in {1/1.1,2/2.6,3/4.1}{
      \node[extern, minimum width=11mm, minimum height=5mm, font=\scriptsize]
        (k\i) at (\x,-0.1) {\code{ksp}};
      \draw[refedge] (g\i.south) -- (k\i.north);
    }
    \node[font=\scriptsize, text=simRust, align=center] at (2.4,-0.95)
      {each step reads the prior $\vy$ \emph{and} hides an implicit solve:\\
       strictly ordered---cannot overlap timesteps};
  \end{scope}
  \draw[simGray!50, thick] (6.2,-1.3) -- (6.2,2.7);
  % ===== RIGHT: the electrostatic MAP (no carry) =====
  \begin{scope}[xshift=7.0cm]
    \node[font=\small\bfseries, text=simBlue] at (2.0,2.5) {electrostatic: a \code{map} (no carry)};
    \node[data, minimum width=11mm] (fam) at (2.0,1.6) {one $\Kop$};
    \foreach \i/\x in {1/0.4,2/2.0,3/3.6}{
      \node[opbox, minimum width=10mm] (m\i) at (\x,0.6) {solve};
      \node[product, minimum width=10mm, minimum height=5mm] (r\i) at (\x,-0.5) {$\vx_{\i}$};
      \draw[flow] (fam.south) -- (m\i.north);
      \draw[flow] (m\i.south) -- (r\i.north);
    }
    \node[font=\scriptsize, text=simGreen, align=center] at (2.0,-1.15)
      {no arrows \emph{between} solves:\\independent, reorderable, parallel};
  \end{scope}
\end{tikzpicture}
\caption[The transient fold versus the electrostatic map]{The transient march
(left) as a state-threaded \code{fold}, beside the electrostatic family (right) as
an independent \code{map}. In the fold the doubled state $\vy_0 \to \vy_1 \to \vy_2
\to \vy_3$ is threaded forward---each step consumes the prior state, so the steps
are strictly ordered and \emph{cannot} overlap---and each step additionally hides
an implicit \code{ksp\_solve} (\cref{sec:explicit-implicit}). In the map the
several solves share one operator $\Kop$ but carry nothing between them, so they
fan out and parallelize. The horizontal carry arrows, present on the left and
absent on the right, are the whole difference: \code{fold\_solve} versus
\code{solve\_family}.}
\label{fig:fold-vs-map}
\end{figure}

\subsection{\texttt{fold\_solve}, in the calculus}

In the vocabulary of \cref{ch:combinators}, the march is one expression. Capture
the operator once, seed the state once, and \code{foldl} the opaque per-step
operator over the schedule of times:
\begin{lstlisting}[style=l4style]
fold_solve :: OpParams -> TimeState -> [Time] -> TimeState
fold_solve op s0 schedule = foldl (\s t -> time_step_op op s t) s0 schedule

-- the opaque per-step operator the fold quantifies over:
-- advances the field-state one step; bottoms out in the library integrator
time_step_op :: OpParams -> TimeState -> Time -> TimeState
\end{lstlisting}
Read the signature and the whole structure of the march is there. The seed
\code{s0} and the result are both a \code{TimeState}---the field-state carried
forward. The step \lstinline[style=l4style]|time_step_op op s t| reads \code{s},
the prior carry; \emph{that read} is what makes this a fold and forbids reordering
at the type level. The operator \code{op} (the captured $\Mop,\Cop,\Kop$ and the
chosen integrator) is bound once, before the fold, and threaded unchanged into
every step---the construction hoist of \cref{sec:integrators} made structural.
And \code{time\_step\_op} is \emph{opaque}: it bottoms out in the library
integrator step of \cref{sec:integrators}, which the combinator only quantifies
over. \code{fold\_solve} owns the loop; the library owns one step.

This is one of the four solve corners of \cref{ch:situating,ch:combinators}---the
\emph{state-threaded fold} corner---and the same combinator drives the
adaptive-mesh-refinement loop of \cref{ch:amr}, where each refinement pass begins
from the mesh the last pass produced. The transient march is simply its purest
instance.

\begin{keyidea}
The time-march is a \emph{sequential fold}: each step begins where the last ended.
It is the \code{fold\_solve} combinator---\lstinline[style=l4style]|foldl time_step_op s0 schedule|---threading
a single field-state carry forward along one timeline, with the operator captured
once outside the loop and each opaque step hiding an implicit \code{ksp\_solve}.
This is the genuinely sequential corner of the simulator: unlike the electrostatic
map, the transient fold cannot overlap its timesteps, because step $n+1$'s input is
step $n$'s output. \emph{Implicit steps cost a linear solve; the carry makes the
march serial.}
\end{keyidea}

Two laws follow immediately from \code{fold\_solve} being a left fold, and they are
the opposite of the map's laws. The map of electrostatics \emph{distributes} over
concatenation, $f(a \mathbin{+\!\!+} b) = f\,a \mathbin{+\!\!+} f\,b$---the algebraic
license for splitting the family across machines. The fold instead
\emph{threads} through concatenation, which is the subject of the next section and
the cleanest gift the fold structure gives us. First, though, note the
\emph{non}-laws: the schedule does not commute (permuting the timesteps changes the
trajectory) and the steps do not fuse (each is an opaque integrator step with its
own internal state advance). A fold is not a map wearing a disguise; the carry is
real, and the type records it.

\section{Checkpoint and resume: the schedule-split law}
\label{sec:checkpoint}

The single most useful algebraic fact about \code{fold\_solve} is the law that a
left fold satisfies over a concatenated list. Split the schedule into a prefix
$a$ and a suffix $b$; then folding the whole is the same as folding the prefix and
folding the suffix \emph{from the state the prefix produced}:
\begin{equation}
  \boxed{\;\code{fold\_solve}\ \code{op}\ \vy_0\ (a \mathbin{+\!\!+} b)
  \;=\;
  \code{fold\_solve}\ \code{op}\ \bigl(\code{fold\_solve}\ \code{op}\ \vy_0\ a\bigr)\ b\;}
  \label{eq:schedule-split}
\end{equation}
This is the standard identity $\code{foldl}\,(a \mathbin{+\!\!+} b) = \code{foldl}\,b
\circ \code{foldl}\,a$, specialized to the captured-operator march. Read it
operationally and it says something an engineer cares about a great deal:
\textbf{a march can be checkpointed and resumed}. Run the prefix $a$, save the
final state $\vy_{|a|} = \code{fold\_solve}\ \code{op}\ \vy_0\ a$ to disk, and stop.
Later, load $\vy_{|a|}$ and run the suffix $b$ seeded from it. The law guarantees
the result is \emph{bit-for-bit} the same trajectory as the uninterrupted run---the
split changed nothing, because the only thing that crosses the split point is the
carry, and the carry is exactly what we saved.

\Cref{fig:checkpoint} draws the equality. The top row is the uninterrupted march;
the bottom row splits it, saving the carry at the cut and resuming from it. The two
produce the same final state because the carry threaded across the cut is the same
value either way.

\begin{figure}[t]
\centering
\begin{tikzpicture}[font=\footnotesize]
  % ===== TOP: uninterrupted march =====
  \node[font=\small\bfseries, text=simInk, anchor=west] at (-0.3,2.5)
    {one uninterrupted march};
  \node[data, minimum width=8mm] (u0) at (0,1.7) {$\vy_0$};
  \foreach \i/\x in {1/1.0,2/2.2,3/3.4,4/4.6}{
    \node[opbox, minimum width=8mm, minimum height=5mm] (us\i) at (\x,1.7) {};
  }
  \node[data, minimum width=8mm] (u4) at (5.6,1.7) {$\vy_4$};
  \draw[flow] (u0)--(us1); \draw[flow] (us1)--(us2); \draw[flow] (us2)--(us3);
  \draw[flow] (us3)--(us4); \draw[flow] (us4)--(u4);
  \node[font=\scriptsize, text=simGray] at (2.8,1.15) {schedule $a \mathbin{+\!\!+} b$};
  % ===== BOTTOM: split + checkpoint =====
  \node[font=\small\bfseries, text=simGreen, anchor=west] at (-0.3,0.3)
    {split: run $a$, save carry, resume $b$};
  \node[data, minimum width=8mm] (b0) at (0,-0.5) {$\vy_0$};
  \foreach \i/\x in {1/1.0,2/2.2}{
    \node[opbox, minimum width=8mm, minimum height=5mm] (bs\i) at (\x,-0.5) {};
  }
  % the checkpoint
  \node[product, minimum width=10mm, fill=simBgGold, draw=simGold] (ckpt) at (3.4,-0.5) {$\vy_2$};
  \foreach \i/\x in {3/4.6,4/5.6}{
    \node[opbox, minimum width=8mm, minimum height=5mm] (bs\i) at (\x,-0.5) {};
  }
  \node[data, minimum width=8mm] (b4) at (6.6,-0.5) {$\vy_4$};
  \draw[flow] (b0)--(bs1); \draw[flow] (bs1)--(bs2); \draw[flow] (bs2)--(ckpt);
  \draw[flow] (ckpt)--(bs3); \draw[flow] (bs3)--(bs4); \draw[flow] (bs4)--(b4);
  \node[font=\scriptsize, text=simGreen] at (1.6,-1.05) {fold $a$};
  \node[font=\scriptsize, text=simRust] at (5.1,-1.05) {fold $b$ from $\vy_2$};
  % save-to-disk callout
  \node[font=\scriptsize, text=simGold!75!black, align=center] at (3.4,-1.5)
    {save / load\\(checkpoint)};
  % equality brace
  \node[font=\Large, text=simGray] at (6.3,0.6) {$=$};
\end{tikzpicture}
\caption[Checkpoint and resume as a schedule split]{The schedule-split law
\cref{eq:schedule-split} drawn as checkpoint-and-resume. \emph{Top}: an
uninterrupted march folds the whole schedule $a \mathbin{+\!\!+} b$ from $\vy_0$ to
$\vy_4$. \emph{Bottom}: the same march split at the midpoint---fold the prefix $a$
to the carry $\vy_2$, \emph{save it} (the gold checkpoint), stop, then later load
$\vy_2$ and fold the suffix $b$ from it. The two produce the identical $\vy_4$,
because the only thing crossing the cut is the carry, and the carry is exactly what
is saved. This is why long runs can be checkpointed against faults and resumed
without changing the answer: it is a clean algebraic property of the fold, not an
engineering approximation.}
\label{fig:checkpoint}
\end{figure}

\begin{keyidea}
Splitting the time schedule and resuming from the saved state yields \emph{exactly}
the same trajectory: $\code{fold\_solve}\ \code{op}\ \vy_0\ (a\mathbin{+\!\!+}b) =
\code{fold\_solve}\ \code{op}\ (\code{fold\_solve}\ \code{op}\ \vy_0\ a)\ b$. This
is the fold's thread-through law---the carry is the only thing that crosses a split
point, so saving the carry \emph{is} a checkpoint. Long runs survive faults
(checkpoint, restart from disk), expensive runs split across job-scheduler windows,
and none of it perturbs the result. Restartability is not bolted on; it is what the
fold structure means.
\end{keyidea}

This is the practical reason the fold's sequentiality, which sounded like pure cost
in \cref{sec:fold-march}, comes with a genuine compensation. A map parallelizes but
has no notion of ``resume''---there is no carry to save. A fold serializes but is
\emph{restartable}, because the carry it threads is precisely the minimal state
needed to continue. The same property---a single carry threaded forward---that
forbids parallelism across timesteps is what makes a checkpoint a single saved
value rather than a replay of the whole history.

\section{Reading the results}
\label{sec:time-readout}

The fold produces a \emph{trajectory}: the field-state at each step. From it the
transient driver reads the physical products, and here the demand-driven pruning of
\cref{ch:fold} earns its keep. If a consumer reads only the final state, the
intermediate states are pruned and the fold degenerates to a bare carry-threading
loop; if a consumer reads the per-step outputs---and the transient driver does, at
every step---the trajectory materializes. What the driver reads at each step is:

\begin{description}[style=nextline, leftmargin=0pt, font=\normalfont\itshape]
  \item[The field trajectory.] From each step's state the electric and magnetic
  fields $\vE, \vB$ are recovered (the \code{GetE}/\code{GetB} of the time
  operator) and written out---the ``movie'' the transient analysis exists to
  produce, a frame per step.
  \item[Port voltages and currents.] Integrating the fields against the port modes
  gives a time-domain voltage and current at each terminal: the waveforms an
  engineer reads off an oscilloscope, computed step by step.
  \item[Energy versus time.] The field energy
  $\tfrac12(\vs^{\!\top}\Mop\,\vs + \dots)$, evaluated at each step, traces how
  energy enters with the pulse, sloshes between electric and magnetic forms, and
  bleeds away through loss and radiation. \Cref{fig:energy} draws a representative
  trajectory: an injected pulse rings up the stored energy, which then decays as
  the damping $\Cop$ and the ports carry it off.
\end{description}

\begin{figure}[t]
\centering
\begin{tikzpicture}
\begin{axis}[
    width=0.80\linewidth, height=50mm,
    xlabel={time $t$}, ylabel={},
    xmin=0, xmax=10, ymin=-1.15, ymax=1.25,
    axis lines=left, xtick={0,2,4,6,8,10}, ytick={-1,0,1},
    yticklabel style={font=\scriptsize}, xticklabel style={font=\scriptsize},
    xlabel style={font=\scriptsize},
    legend style={font=\scriptsize, draw=simGray!50, at={(0.98,0.97)},
      anchor=north east, legend cell align=left},
    every axis plot/.append style={very thick, samples=120, domain=0:10},
  ]
  % injected pulse: a modulated gaussian
  \addplot[simRust] {0.9*exp(-2*(x-1.5)^2)*cos(deg(6*(x-1.5)))};
  \addlegendentry{drive $\vJ(t)$ (a pulse)}
  % a field component: rings up then decays
  \addplot[simBlue] {0.8*(1-exp(-3*x))*exp(-0.35*x)*cos(deg(4*x-1))};
  \addlegendentry{field $E(t)$ (rings, decays)}
  % energy: rises then decays monotonically
  \addplot[simGreen, densely dashed] {0.95*(1-exp(-2.2*x))*exp(-0.5*x)};
  \addlegendentry{stored energy (rise then decay)}
\end{axis}
\end{tikzpicture}
\caption[A transient field and energy trajectory]{A representative readout of the
transient march (\cref{sec:time-readout}). The pulse drive $\vJ(t)$ (rust) injects
a short burst; a field component $E(t)$ (blue) rings up as the pulse arrives and
then decays as energy leaves the system; the stored energy (green, dashed) rises
with the injection and decays monotonically thereafter, bled away by the damping
$\Cop$ and carried off through the ports. Each curve is sampled once per fold step,
so the resolution of the readout is exactly the resolution of the march---one frame
per timestep.}
\label{fig:energy}
\end{figure}

\subsection{A bridge to the frequency domain}

One readout deserves a forward pointer. A transient run injects a \emph{broadband}
pulse---a superposition of many frequencies at once (\cref{sec:transient})---and
records the time-domain response. Taking the Fourier transform of an input
waveform and the corresponding output waveform, and dividing, recovers the
frequency response across the whole band of the pulse \emph{in a single run}. This
is the time-domain route to the same S-parameters the driven analysis
(\cref{ch:reductions}) computes frequency by frequency: where the driven analysis
sweeps a \code{frequency\_sweep} map, one frequency per solve, the transient
analysis folds one march and Fourier-transforms the result. The two analyses
compute the same physics by the two routes of the harmonic rule
(\cref{ch:reductions-pde})---a single frequency at a time, or all frequencies at
once through a pulse. The transient driver and its products are studied in full in
\cref{ch:transient}.

\section{Worked examples}

\begin{workedexample}{Stiffness made concrete: explicit blows up, implicit
survives}{stiff}
We make the stiffness argument of \cref{sec:explicit-implicit} numerical on the
scalar test equation $\dot y = -\lambda y$ with a \emph{stiff} rate $\lambda = 100$
and a modest step $\Delta t = 0.05$. The exact solution from $y_0 = 1$ is
$y(t) = e^{-100 t}$, which decays to nothing almost instantly.

\medskip
\emph{Explicit Euler.} The amplification factor \cref{eq:fe-scalar} is
\[
  g_{\text{exp}} = 1 - \lambda\,\Delta t = 1 - 100\cdot 0.05 = 1 - 5 = -4.
\]
Each step multiplies $y$ by $-4$. The stability bound \cref{eq:fe-stability} demands
$\Delta t < 2/\lambda = 0.02$, and our step $0.05$ violates it, so $|g| = 4 > 1$ and
the numerical solution \emph{grows} while the true one decays:
\[
  y_0 = 1,\quad y_1 = -4,\quad y_2 = 16,\quad y_3 = -64,\quad y_4 = 256,\ \dots
\]
an explosive oscillation, the exact opposite of $e^{-100t}\to 0$. To make explicit
Euler stable we would have to shrink $\Delta t$ below $0.02$---and on a real mesh
$\lambda_{\max}$ is far larger than $100$, forcing $\Delta t$ smaller still.

\medskip
\emph{Implicit (backward) Euler.} The amplification factor \cref{eq:be-scalar} is
\[
  g_{\text{imp}} = \frac{1}{1 + \lambda\,\Delta t} = \frac{1}{1 + 5} = \frac{1}{6}
  \approx 0.167.
\]
Now $|g| < 1$ for our step---indeed for \emph{any} step---so the solution decays as
it should:
\[
  y_0 = 1,\quad y_1 = \tfrac16 \approx 0.167,\quad y_2 = \tfrac1{36} \approx 0.028,
  \quad y_3 \approx 0.0046,\ \dots
\]
monotone toward zero, tracking the true decay. \Cref{fig:stiff} plots the two
against the exact solution: explicit Euler oscillates out of the frame, implicit
Euler decays cleanly. The implicit step paid for itself with one solve (here a
trivial scalar division; on the full system, a \code{ksp\_solve}) and in exchange
removed the step-size limit entirely. This single picture is the whole reason
electromagnetic transient solvers are implicit.
\end{workedexample}

\begin{figure}[t]
\centering
\begin{tikzpicture}
\begin{axis}[
    width=0.78\linewidth, height=52mm,
    xlabel={step $n$ \;($t = 0.05\,n$)}, ylabel={$y_n$},
    xmin=0, xmax=6, ymin=-9, ymax=9,
    axis lines=middle, xtick={1,2,3,4,5}, ytick={-8,-4,0,4,8},
    yticklabel style={font=\scriptsize}, xticklabel style={font=\scriptsize},
    xlabel style={font=\scriptsize, at={(axis description cs:0.5,-0.06)}},
    ylabel style={font=\scriptsize},
    legend style={font=\scriptsize, draw=simGray!50, at={(0.5,-0.22)},
      anchor=north, legend columns=3, column sep=5pt},
    every axis plot/.append style={very thick, mark=*, mark size=1.5pt},
  ]
  % explicit Euler: 1, -4, 16, -64(clip), ... oscillating blowup
  \addplot[simRust] coordinates {(0,1) (1,-4) (2,8.8) (3,-8.8)};
  \addlegendentry{explicit Euler ($g=-4$)}
  % implicit Euler: 1, 1/6, 1/36, ... decay
  \addplot[simBlue] coordinates {(0,1) (1,0.167) (2,0.028) (3,0.0046) (4,0.0008) (5,0.0001)};
  \addlegendentry{implicit Euler ($g=\tfrac16$)}
  % exact: e^{-100 t} = e^{-5 n}
  \addplot[simGreen, dashed, mark=none, samples=60, domain=0:6] {exp(-5*x)};
  \addlegendentry{exact $e^{-100t}$}
\end{axis}
\end{tikzpicture}
\caption[Explicit blow-up versus implicit stability]{The stiff test of
\cref{wex:stiff}: $\dot y = -\lambda y$, $\lambda = 100$, $\Delta t = 0.05$
(violating the explicit bound $\Delta t < 0.02$). \emph{Explicit Euler} (rust)
amplifies by $g = -4$ each step and oscillates explosively out of the frame---the
later points $\pm 64, \pm 256$ are clipped. \emph{Implicit Euler} (blue) contracts
by $g = \tfrac16$ each step and decays cleanly onto the exact solution (green,
dashed). The implicit step is unconditionally stable; the explicit step is useless
here without a step ten times smaller. This is the stiffness catastrophe, and its
cure, in one plot.}
\label{fig:stiff}
\end{figure}

\begin{workedexample}{Marching a two-degree-of-freedom oscillator}{twodof}
We now run the real machinery---the recast, the per-step implicit solve, and the
state threading---on the smallest non-trivial system: a two-degree-of-freedom
mass--spring--damper, $\Mop\ddot\vs + \Cop\dot\vs + \Kop\vs = \vJ(t)$ with
\[
  \Mop = \begin{psmallmatrix}1 & 0\\ 0 & 1\end{psmallmatrix},\quad
  \Cop = \begin{psmallmatrix}0.2 & 0\\ 0 & 0.2\end{psmallmatrix},\quad
  \Kop = \begin{psmallmatrix}2 & -1\\ -1 & 2\end{psmallmatrix},\quad
  \vJ(t) = \begin{psmallmatrix}1\\ 0\end{psmallmatrix}\ (\text{constant}),
\]
starting from rest, $\vs_0 = \vv_0 = (0,0)$. This is a hand-sized stand-in for the
semi-discrete transient system; $\Kop$ is the stiffness coupling the two degrees of
freedom.

\medskip
\emph{Recast.} The doubled state is $\vy = (\vs, \vv) \in \mathbb R^4$, with rate
$\vF(\vy) = (\vv,\ \Mop^{-1}(\vJ - \Cop\vv - \Kop\vs))$ as in
\cref{eq:time-recast}. Since $\Mop = \mat I$ here, $\dot\vv = \vJ - \Cop\vv -
\Kop\vs$.

\medskip
\emph{One implicit step.} Take backward Euler with $\Delta t = 0.1$. The step
\cref{eq:bwd-euler} couples the two halves: $\vs_{n+1} = \vs_n + \Delta t\,\vv_{n+1}$
and $\vv_{n+1} = \vv_n + \Delta t\,(\vJ - \Cop\vv_{n+1} - \Kop\vs_{n+1})$.
Substituting the first into the second eliminates $\vs_{n+1}$ and leaves a
\emph{linear system} for $\vv_{n+1}$ alone:
\[
  \underbrace{\bigl(\mat I + \Delta t\,\Cop + \Delta t^2\,\Kop\bigr)}_{\displaystyle \mat A}
  \,\vv_{n+1}
  = \vv_n + \Delta t\,\bigl(\vJ - \Kop\,\vs_n\bigr).
\]
This is the implicit step's linear solve, the $2\times 2$ stand-in for the
\code{ksp\_solve} (\cref{sec:explicit-implicit}). With our numbers,
\[
  \mat A = \mat I + 0.1\,\Cop + 0.01\,\Kop
  = \begin{psmallmatrix}1.04 & -0.01\\ -0.01 & 1.04\end{psmallmatrix}.
\]
The matrix $\mat A$ is \emph{fixed for the whole march}---it is built from the
captured $\Mop,\Cop,\Kop$ and $\Delta t$, none of which change---so its
factorization is computed once and reused every step. This is the
operator-capture-once hoist of \code{fold\_solve} (\cref{sec:fold-march}) in
miniature: \code{op} is bound before the fold; every step reads the same $\mat A$.

\medskip
\emph{Thread the state forward.} Starting from $\vy_0 = \mathbf 0$, solve for
$\vv_1$ then set $\vs_1 = \Delta t\,\vv_1$, feed $\vy_1 = (\vs_1, \vv_1)$ into the
next step, and repeat. Three steps:
\begin{center}
\setlength{\tabcolsep}{7pt}
\renewcommand{\arraystretch}{1.2}
\begin{tabular}{@{}clll@{}}
\toprule
step $n$ & RHS $\vv_n + \Delta t(\vJ - \Kop\vs_n)$ & solve $\vv_{n+1} = \mat A^{-1}(\cdot)$ & $\vs_{n+1} = \vs_n + \Delta t\,\vv_{n+1}$ \\
\midrule
0 & $(0.100,\ 0.000)$ & $(0.0961,\ 0.0009)$ & $(0.00961,\ 0.00009)$ \\
1 & $(0.1854,\ 0.0105)$ & $(0.1784,\ 0.0118)$ & $(0.02745,\ 0.00127)$ \\
2 & $(0.2589,\ 0.0264)$ & $(0.2492,\ 0.0278)$ & $(0.05237,\ 0.00405)$ \\
\bottomrule
\end{tabular}
\end{center}
Each row's input is the previous row's output---the carry $\vy_n \to \vy_{n+1}$
threaded forward, exactly the fold of \cref{sec:fold-march}. The driven coordinate
$s^{(1)}$ climbs steadily toward its static deflection while the coupled coordinate
$s^{(2)}$ follows through the off-diagonal stiffness; with the damping $\Cop$
present, the march settles rather than oscillating forever. Three points are the
whole chapter in miniature: \emph{(i)} the recast doubled $\mathbb R^2$ into $\mathbb
R^4$; \emph{(ii)} each step \emph{solved a linear system} $\mat A\vv_{n+1} =
\text{RHS}$, with $\mat A$ fixed and factored once; \emph{(iii)} the state was
\emph{threaded} from each step into the next---a \code{fold\_solve}, not a map. Were
we to checkpoint after step~1, saving $\vy_1$, and resume, the schedule-split law
\cref{eq:schedule-split} guarantees we would land on the identical $\vy_3$.
\end{workedexample}

\summarysection
The transient analysis marches the semi-discrete second-order system $\Mop\ddot\vs
+ \Cop\dot\vs + \Kop\vs = \vJ(t)$ forward in time. First it is recast as a
first-order system $\dot\vy = \vF(\vy, t)$ in the doubled state $\vy = (\vs,
\dot\vs)$, so standard integrators apply. Because the assembled spatial operator is
\emph{stiff}, explicit steps are limited to a tiny $\Delta t < 2/\lambda_{\max}$,
so electromagnetics integrates \emph{implicitly}: each implicit step is a linear
solve---a \code{ksp\_solve} (\cref{ch:krylov})---and the methods used
(generalized-$\alpha$ with its tunable numerical damping, SDIRK, ARKStep) are all
A-stable, trading a per-step solve for freedom from the stiffness limit.
Structurally, the whole march is the \code{fold\_solve} combinator of
\cref{ch:combinators}: a state-threaded \code{foldl} of an opaque per-step operator,
each step beginning where the last ended, the operator captured once outside the
loop. The fold's thread-through law,
$\code{fold\_solve}\,\code{op}\,\vy_0\,(a\mathbin{+\!\!+}b) =
\code{fold\_solve}\,\code{op}\,(\code{fold\_solve}\,\code{op}\,\vy_0\,a)\,b$, makes
checkpoint-and-restart exact: the carry is the only thing crossing a split point.
The march is genuinely sequential---it cannot overlap its timesteps, the price of
the carry the electrostatic map does not pay.

\furtherreading
The recast, stability regions, A-stability, order, and the stiffness phenomenon are
standard material in any numerical-ODE or matrix-computations
reference; Golub and Van Loan~\citep{golubvanloan2013} cover the linear-algebra
machinery of the per-step solve and the eigenvalue spread that makes a system
stiff. Jin's finite-element electromagnetics text~\citep{jin2014} treats the
time-domain finite-element formulation and the specific integrators used for
Maxwell's equations, including the generalized-$\alpha$ family and its numerical
damping. The per-step linear solve is a Krylov solve in the sense of
\cref{ch:krylov}, for which Saad~\citep{saad2003} is the standard reference; the
fold structure and its laws are developed in \cref{ch:fold,ch:combinators}, and the
transient driver that wraps the march is \cref{ch:transient}.

\exercisessection
\begin{enumerate}[nosep]
  \item \emph{The recast preserves the system.} Starting from the doubled
  first-order form \cref{eq:time-recast}, eliminate $\vv$ by differentiating the
  top equation and substituting, and recover the original second-order system
  \cref{eq:time-secondorder}. Conclude that the recast adds no physics---it only
  changes the shape of the unknowns.
  \item \emph{The explicit step-size bound.} For the scalar test $\dot y =
  -\lambda y$ with $\lambda = 50$, find the largest $\Delta t$ for which explicit
  Euler is stable. If a mesh refinement raises $\lambda_{\max}$ by a factor of $4$,
  by what factor must $\Delta t$ shrink? What does this say about explicit methods
  on fine meshes?
  \item \emph{Backward Euler is A-stable.} Show that for $\dot y = -\lambda y$ with
  \emph{any} complex $\lambda$ in the left half-plane ($\Re\lambda > 0$, decaying)
  and \emph{any} $\Delta t > 0$, the backward-Euler factor $g = 1/(1 +
  \lambda\Delta t)$ satisfies $|g| < 1$. (Hint: $|1 + \lambda\Delta t|^2 = (1 +
  \Delta t\,\Re\lambda)^2 + (\Delta t\,\Im\lambda)^2$.)
  \item \emph{Map versus fold, by the type.} Explain why \code{solve\_family}
  (electrostatics) can be evaluated on three machines in any order but
  \code{fold\_solve} (transient) cannot, citing the one feature of the
  \code{fold\_solve} signature that forbids reordering. Which line of the signature
  is it?
  \item \emph{Checkpoint exactness.} Using the schedule-split law
  \cref{eq:schedule-split}, argue that a transient run checkpointed after every
  $K$ steps and resumed produces \emph{exactly} the same final state as an
  uninterrupted run, for any $K$. What single quantity must the checkpoint save?
  Why does an electrostatic \code{solve\_family} have no analogous ``resume''?
  \item \emph{One implicit step by hand.} For the two-degree-of-freedom system of
  \cref{wex:twodof}, take \emph{one} backward-Euler step of size $\Delta t = 0.2$
  (rather than $0.1$) from rest. Form the step matrix $\mat A = \mat I + \Delta
  t\,\Cop + \Delta t^2\Kop$, write the right-hand side, and solve the $2\times 2$
  system for $\vv_1$. Confirm $\mat A$ depends only on $\Delta t$ and the operators,
  not on the step index---the operator-capture-once property.
  \item \emph{Numerical damping, qualitatively.} Generalized-$\alpha$ with
  $\rho_\infty = 0$ annihilates the highest-frequency modes in one step, while
  $\rho_\infty = 1$ preserves them. For a transient simulation polluted by ringing
  on the smallest mesh elements (unresolved high modes), which setting would you
  choose, and what is the risk of choosing $\rho_\infty$ too small (too much
  damping)?
  \item \emph{The two routes to S-parameters.} The driven analysis computes a
  frequency response by sweeping one frequency per solve (a \code{frequency\_sweep}
  map); the transient analysis computes the same response by one pulsed march
  followed by a Fourier transform. Give one reason to prefer each route. (Consider:
  how many frequencies you want, and whether the response is needed at a few sharp
  frequencies or across a broad band.)
\end{enumerate}

\chapter{Adaptive Mesh Refinement}
\label{ch:amr}

\chapterorient{\Cref{ch:lifecycle} showed the simulation lifecycle closing into a
loop: solve, estimate the error, mark the worst elements, refine the mesh, solve
again. We called that the adaptive estimate--mark--refine loop and named it as our
first \emph{fold}, but we left its three working stages as black boxes. This
chapter opens them. We ask three questions in turn. \emph{Where} is the error
worst---when we do not know the true solution and so cannot subtract it? \emph{Which}
elements should we refine---given a per-element error guess, how many, and which?
And \emph{how} does the refinement happen, and why does the whole loop terminate
with a good answer using far fewer unknowns than a uniform mesh would need? Two of
the three answers are short, sharp algorithms we can run by hand---a flux-recovery
estimate and a bulk-marking rule---and the third is a mesh operation we will treat
as a boundary. By the end you will be able to compute an error indicator on a tiny
mesh, mark elements by the D\"orfler rule on a list of numbers, and read a
convergence curve that shows adaptivity earning its keep.}

\section{Why adapt at all}
\label{sec:why-adapt}

Refining a mesh means cutting its elements into smaller ones, raising the number of
unknowns $N$ and---by C\'ea's lemma (\cref{ch:galerkin})---driving the discretization
error down. The simplest way to refine is \emph{uniformly}: cut every element. In
one step a 2-D mesh of $N$ elements becomes $4N$; in three dimensions, $8N$. Uniform
refinement is easy to implement and easy to reason about, and for a perfectly smooth
solution on a simple domain it is hard to beat. But the solutions an engineering
simulator actually computes are rarely smooth everywhere.

Real fields have \emph{features}. An electrostatic potential develops a singularity
at a sharp conductor corner---the field strength blows up like $r^{-\gamma}$ as the
distance $r$ to the corner goes to zero. A magnetic field crowds into the gap between
two poles. The current concentrates at a material interface where the permittivity
jumps. Near such a feature the solution bends violently and the piecewise-polynomial
basis of \cref{ch:galerkin} struggles to follow it, so the local error is large.
Away from the feature---across the broad, gently-varying interior---the same basis is
already more than accurate enough, and the local error is tiny.

\Cref{fig:cornermesh} is the canonical picture. A re-entrant (``L-shaped'') corner
pulls a singularity into the field; a good mesh is \emph{fine} in a small
neighbourhood of the corner and \emph{coarse} everywhere else. Uniform refinement,
by contrast, would spend the same effort everywhere---lavishing resolution on the
smooth interior, where it buys almost nothing, in order to resolve the corner.

\begin{figure}[t]
\centering
\begin{tikzpicture}[scale=1.0, every node/.style={font=\footnotesize}]
  % L-shaped domain outline (re-entrant corner at the origin)
  \def\a{4.0}
  \draw[simInk, thick, fill=simBgGray]
    (0,0) -- (\a,0) -- (\a,2.0) -- (2.0,2.0) -- (2.0,\a) -- (0,\a) -- cycle;
  % the re-entrant corner
  \fill[simRust] (2.0,2.0) circle (2.2pt);
  \node[simRust, above right=0pt and 1pt] at (2.0,2.0) {re-entrant corner};
  % graded mesh: coarse far cells
  \draw[simBlue!55, thin] (0,0) grid[step=1.0] (\a,2.0);
  \draw[simBlue!55, thin] (0,2.0) grid[step=1.0] (2.0,\a);
  % progressively finer cells toward the corner: a few nested rings
  \foreach \s in {0.5,0.25,0.125}{
    \draw[simBlue!75, thin] (2.0-\s,2.0-\s) rectangle (2.0,2.0);
    \draw[simBlue!75, thin] (2.0-\s,2.0) rectangle (2.0,2.0+\s);
    \draw[simBlue!75, thin] (2.0,2.0-\s) rectangle (2.0+\s,2.0);
  }
  % subdivide those finer rings further (just visual: cross lines)
  \draw[simBlue!75, thin] (1.5,2.0)--(2.0,2.0); % already there
  % field-strength annotation
  \draw[refedge, simRust] (3.4,3.0) to[bend right=20] (2.15,2.15);
  \node[simRust, font=\scriptsize, align=center] at (3.55,3.2)
    {$|\vE|\sim r^{-\gamma}$\\(error large)};
  \node[simGray, font=\scriptsize, align=center] at (1.0,1.0)
    {smooth interior\\(error tiny)};
\end{tikzpicture}
\caption[A mesh graded to a corner singularity]{The canonical adaptive mesh. A
re-entrant corner (here the inner corner of an L-shaped domain) forces the field to
blow up like $r^{-\gamma}$ as the distance $r$ to the corner shrinks, so the
discretization error concentrates there. A well-adapted mesh is heavily refined in a
small neighbourhood of the corner---the nested rings of ever-smaller cells---and
left coarse across the smooth interior, where the basis already resolves the field.
Uniform refinement would instead cut \emph{every} cell, spending most of its new
unknowns on the interior, where they buy almost nothing.}
\label{fig:cornermesh}
\end{figure}

The payoff is captured by a single principle. C\'ea's lemma told us the total error
is governed by how well the basis can approximate the solution; interpolation theory
adds that, for a fixed total number of unknowns, the approximation error is
\emph{minimised} when the local error contributions are \emph{equal} across
elements---when the error is \emph{equidistributed}. A uniform mesh equidistributes
the cell \emph{size}, not the error: it leaves large error at the corner and
negligible error in the interior, a wildly unequal distribution. An adaptive mesh
equidistributes the \emph{error} by making cells small exactly where the error per
cell is large. The result is dramatic: for a corner singularity in two dimensions, a
uniform mesh achieves error $\mathcal{O}(N^{-\gamma/2})$ with a stalled exponent set
by the singularity strength $\gamma$, while a properly graded adaptive mesh recovers
the full smooth-solution rate $\mathcal{O}(N^{-p/2})$---the same accuracy for
\emph{orders of magnitude} fewer unknowns. \Cref{sec:convergence} makes this
quantitative; \cref{fig:convergence} draws it.

\begin{keyidea}
Uniform refinement equidistributes \emph{cell size}; adaptive refinement
equidistributes \emph{error}. Because a solution's error concentrates at features
(corners, interfaces, field concentrations) and is negligible across smooth regions,
putting small cells only where the error is large reaches a target accuracy with far
fewer unknowns than cutting every cell. The recipe to do this automatically is the
loop of \cref{ch:lifecycle}: \emph{estimate} where the error is, \emph{mark} a bulk
fraction of it, \emph{refine} those elements, and \emph{repeat}. The rest of this
chapter is those four words.
\end{keyidea}

There is a circularity to defeat first. To put small cells where the error is large,
we must know where the error is large---but the error is $u - u_h$, and we do not
know the exact solution $u$ (if we did, we would not be running the simulation). The
whole edifice rests on \emph{estimating} the error from the computed solution alone,
without ever seeing the truth. That is the subject of the next section, and it is the
cleverest of the three stages.

\section{Estimate: Zienkiewicz--Zhu flux recovery}
\label{sec:estimate}

We have a computed solution $u_h$ and want a number $\eta_K \ge 0$ for each element
$K$---an \emph{error indicator}---large where the local error $\norm{u - u_h}_K$ is
large and small where it is small. We cannot compute the error directly. The
Zienkiewicz--Zhu (ZZ) idea is to find a \emph{computable quantity that is zero
exactly when the error is zero} and grows with it---a proxy we \emph{can} evaluate.

\subsection{The discontinuous flux and the smoothing trick}

The proxy is built from the \emph{flux}. In the electrostatic problem the field is
$\vE = -\Grad V$ and the physical flux is the displacement $\vD = \varepsilon\vE$;
in the magnetostatic problem the flux is $\vH = \mu^{-1}\vB$. Generically write the
flux as $\mat F = Q\,\diffop u_h$, a material weight $Q$ times a derivative of the
discrete solution. Here is the key observation about this flux as the finite-element
method computes it.

The discrete solution $u_h$ is continuous across element boundaries (it is built from
the conforming basis of \cref{ch:galerkin}), but its \emph{derivative} is not. A
piecewise-linear $u_h$ has a constant gradient on each element that jumps as you
cross into the next. So the computed flux $\mat F = Q\,\diffop u_h$ is
\emph{discontinuous across element faces}---it has jumps. The \emph{true} flux of the
exact solution, by contrast, is smooth (continuity of the normal flux is a physical
conservation law). The size of the jumps in the computed flux therefore measures how
far the discrete solution is from the truth: where $u_h$ resolves the solution well,
its flux is nearly continuous and the jumps are small; where $u_h$ struggles, the
flux is wildly discontinuous and the jumps are large.

ZZ turns this qualitative observation into a number with one move: \emph{recover a
smooth flux} $G$ from the discontinuous one and measure the gap. Concretely, $G$ is
the $L^2$ projection of the computed flux $\mat F$ onto a conforming
(continuous-where-it-should-be) finite-element space:
\begin{equation}
G = \arg\min_{g \in V_{\text{smooth}}} \int_\dom \norm{g - \mat F}^2 \dV,
\qquad\text{equivalently}\qquad
G = \mat M^{-1}\,\bigl(\text{Flux}\cdot u_h\bigr),
\label{eq:zz-recover}
\end{equation}
where $\mat M$ is the mass matrix of the smooth space and $\text{Flux}\cdot u_h$
assembles the projection's right-hand side---a single mass-matrix solve, exactly the
kind of system the solvers of \cref{ch:krylov} dispatch. The smoothed flux $G$ is the
code's best guess at the \emph{true}, continuous flux; the gap $\mat F - G$ between
the rough computed flux and its smoothed version stands in for the unknowable error.

\begin{figure}[t]
\centering
\begin{tikzpicture}[scale=1.0]
\begin{axis}[
  width=12.5cm, height=5.4cm,
  xmin=-0.05, xmax=4.05, ymin=-0.15, ymax=2.6,
  xtick={0,1,2,3,4},
  xticklabels={$x_0$,$x_1$,$x_2$,$x_3$,$x_4$},
  ytick=\empty, ylabel={flux}, axis lines=left,
  tick label style={font=\footnotesize},
  legend style={at={(0.97,0.97)}, anchor=north east, font=\footnotesize,
    draw=simGray!50},
  clip=false,
]
  % discontinuous element flux: constant per element, jumps at nodes
  \addplot[simRust, very thick] coordinates {(0,0.35)(1,0.35)};
  \addplot[simRust, very thick] coordinates {(1,0.75)(2,0.75)};
  \addplot[simRust, very thick] coordinates {(2,1.7)(3,1.7)};
  \addplot[simRust, very thick] coordinates {(3,2.15)(4,2.15)};
  \addlegendentry{$\mathbf{F}$ (computed, discontinuous)}
  % the jumps (dotted vertical risers)
  \addplot[simRust, densely dotted, thin] coordinates {(1,0.35)(1,0.75)};
  \addplot[simRust, densely dotted, thin] coordinates {(2,0.75)(2,1.7)};
  \addplot[simRust, densely dotted, thin] coordinates {(3,1.7)(3,2.15)};
  % recovered smooth flux: a continuous (piecewise-linear through midpoints) curve
  \addplot[simBlue, very thick, mark=*, mark size=1.4pt] coordinates
    {(0,0.2)(0.5,0.35)(1.5,0.75)(2.5,1.7)(3.5,2.15)(4,2.35)};
  \addlegendentry{$G$ (recovered, smooth)}
  % gap annotation on the big-jump element (between x2 and x3)
  \draw[<->, simGreen, thick] (axis cs:2.5,1.30) -- (axis cs:2.5,1.70);
  \node[simGreen, font=\scriptsize, anchor=west] at (axis cs:2.58,1.5)
    {gap $=\eta_K$};
  \node[simGray, font=\scriptsize, anchor=north] at (axis cs:0.5,0.30)
    {small gap};
\end{axis}
\end{tikzpicture}
\caption[Zienkiewicz--Zhu flux recovery]{The ZZ estimator in one dimension. The
computed flux $\mathbf F = Q\,\mathcal{D}u_h$ (rust) is \emph{piecewise constant} and
\emph{jumps} across each element boundary, because $u_h$'s derivative is
discontinuous. The recovered flux $G$ (blue) is the smooth $L^2$ projection of
$\mathbf F$ onto a conforming space---the code's best estimate of the true,
continuous flux. The per-element indicator $\eta_K$ is the size of the gap $\mathbf
F - G$ over element $K$ (\cref{eq:zz-indicator}). It is small where the flux is
nearly continuous (left, smooth region) and large where the flux jumps sharply
(centre, where the solution bends most). The estimator never sees the true error---it
infers it from the roughness the discretization left behind.}
\label{fig:zz}
\end{figure}

\subsection{The per-element indicator}

With a smooth flux $G$ in hand, the error indicator on element $K$ is the $L^2$ norm
of the gap, restricted to that element:
\begin{equation}
\boxed{\;
\eta_K^2 = \int_K \norm{\mat F - G}^2 \dV
\;}
\label{eq:zz-indicator}
\end{equation}
We work with the \emph{squared} indicator $\eta_K^2$ throughout, for a reason that
will matter in marking: squared error contributions \emph{add}, and the total squared
error $\sum_K \eta_K^2$ is the quantity the marking step apportions. The square root
that recovers $\eta_K$ itself (with a physical energy normalisation) is a thin
epilogue applied later, when the indicators are compared against a tolerance.

Three properties make \cref{eq:zz-indicator} a good indicator, and each is worth
naming because each is a fact you can check.

\emph{Local and cheap.} The integral in \cref{eq:zz-indicator} is over a single
element $K$, so $\eta_K^2$ is computed element-by-element from the local values of
$\mat F$ and $G$---one small quadrature per element, no global coupling in this step.
(The recovery \cref{eq:zz-recover} of $G$ is a global solve, but it is one solve,
shared by all elements; the per-element \emph{difference} is local.) The cost is a
small constant per element, negligible beside the solve that produced $u_h$.

\emph{Quadratic in the field.} The indicator is a squared norm of $\mat F - G$, and
both $\mat F$ and $G$ are \emph{linear} in $u_h$ (the flux map and the projection are
linear operators). So $\eta_K^2$ is a \emph{quadratic} functional of the solution: it
scales as $\eta_K^2(\alpha u_h) = \alpha^2\,\eta_K^2(u_h)$ but is \emph{not} additive,
$\eta_K^2(u_h + w_h) \ne \eta_K^2(u_h) + \eta_K^2(w_h)$. It is an energy-like
quantity, not a linear measurement.

\emph{Vanishes exactly when the flux is already smooth.} If the computed flux $\mat
F$ already lies in the smooth recovery space---no jumps to recover away---then its
$L^2$ projection $G$ equals $\mat F$ itself, and $\eta_K^2 = 0$ for every element.
The estimator fires precisely on the \emph{non-smoothness} the discretization
introduced, which is exactly the ZZ premise: a big jump means a big error, a smooth
flux means a converged solution.

\begin{pitfall}
The ZZ estimate is an \emph{indicator}, not the error. It is not a rigorous upper
bound on $\norm{u - u_h}$---it can over- or under-estimate the true error by a
problem-dependent constant. What it does reliably is rank the elements: it is large
where the error is large and small where it is small, so it tells you \emph{which}
elements to refine even when it does not tell you the error's exact size. For
\emph{driving refinement} that ranking is all we need (the marking step uses only the
relative sizes); for \emph{certifying} a final accuracy one wants a guaranteed bound,
which is a separate and harder theory. Do not read $\eta_K$ as ``the error on element
$K$''; read it as ``element $K$'s share of the work still to do.''
\end{pitfall}

\subsection{Why recover rather than just measure the jumps?}

One might ask: if the jumps in $\mat F$ are the signal, why project at all---why not
measure the jumps directly, face by face? Both strategies exist in the literature
(the direct one is the ``residual'' or ``jump'' estimator). The recovery approach has
two practical virtues that make it Palace's choice. First, it produces a smooth flux
$G$ that is itself a \emph{better} approximation to the true flux than $\mat F$
was---a useful by-product, sometimes used to improve the reported solution.\ Second,
the projection \cref{eq:zz-recover} is a single mass-matrix solve in machinery the
code already has (\cref{ch:cg}'s conjugate gradients on an SPD mass matrix), so the
estimator reuses the solver infrastructure rather than introducing face-by-face jump
bookkeeping. The estimator is, structurally, ``one extra solve and one cheap
per-element norm''---a small tax on each adaptive pass.

\section{Mark: D\"orfler bulk marking}
\label{sec:mark}

The estimate step hands us a list of nonnegative numbers, one per element:
$\{\eta_K^2\}_{K=1}^{N}$. The mark step must decide \emph{which} elements to refine.
This is a more subtle question than it looks, and the answer---D\"orfler's bulk
marking---is the one choice in the loop with a convergence theorem attached.

\subsection{Three marking strategies, one good one}

The naive strategies are tempting and wrong.

\emph{Refine the top $k$.} Mark the $k$ elements with the largest indicators, for a
fixed $k$. But $k$ is a magic number unrelated to the error: on a mesh where the
error is concentrated in two elements, marking the top $50$ wastes effort; on a mesh
where it is spread over hundreds, marking the top $50$ refines too little.

\emph{Refine the top fraction of elements.} Mark the worst-ranked $10\%$ of
\emph{elements}. This couples to the mesh size but still not to the error: $10\%$ of a
mesh might carry $90\%$ of the error or $5\%$ of it, depending on how concentrated the
features are. Marking by element \emph{count} ignores the very thing we are trying to
equidistribute.

\emph{Refine the elements carrying a fixed fraction of the error.} This is
D\"orfler's idea, and it is the right one. Mark the \emph{smallest} set of elements
whose summed squared error reaches a target fraction $\theta$ of the \emph{total}
squared error:
\begin{equation}
\boxed{\;
\text{find the smallest } \mathcal{S} \text{ with } \;
\sum_{K \in \mathcal{S}} \eta_K^2 \;\ge\; \theta \sum_{K} \eta_K^2.
\;}
\label{eq:dorfler}
\end{equation}
The marking is by \emph{error mass}, not element count. If the error is concentrated
in a few elements, $\mathcal{S}$ is tiny; if it is spread out, $\mathcal{S}$ is large.
Either way the marked set captures the same fraction $\theta$ of the work. The
default fraction in Palace is $\theta = 0.7$: mark the smallest set of elements
responsible for $70\%$ of the total squared error. This is called \emph{bulk}
marking because it marks the bulk of the error, however few or many elements that
takes.

\subsection{The algorithm: sort, accumulate, pivot}

\Cref{eq:dorfler} reads as a combinatorial optimisation---``the smallest set
with''---but it has a one-pass realisation, because greedily taking the largest
indicators first is optimal. The algorithm is:

\begin{enumerate}[nosep]
\item \textbf{Sort} the squared indicators in descending order.
\item \textbf{Accumulate} a running sum down the sorted list.
\item \textbf{Stop} (the \emph{pivot}) at the first point where the running sum
reaches $\theta$ of the total. Everything down to and including the pivot is marked.
\end{enumerate}

That ``largest first'' greedy is optimal for \cref{eq:dorfler}: to reach a fixed
error budget $\theta\sum_K\eta_K^2$ with the fewest elements, you must take the
biggest contributors first---swapping any marked element for a smaller unmarked one
could only require \emph{more} elements to reach the same budget. So the prefix of the
sorted list \emph{is} the smallest set. \Cref{alg:dorfler} writes it in the
framework's pseudo-language; the implementation realises the running sum as a
\emph{prefix sum} and the pivot as a binary search (a \code{lower\_bound}) for speed,
but the meaning is exactly the three steps above.

\begin{lstlisting}[style=l4style,caption={D\"orfler bulk marking: the smallest set covering fraction $\theta$ of the squared error.},label={alg:dorfler}]
dorfler_mark :: Float -> [Float] -> IndexSet
dorfler_mark theta etas2 =
  let total   = sum etas2                       -- total squared error
      ranked  = sortDescendingByIndicator etas2 -- (index, eta2) pairs, biggest first
      target  = theta * total                   -- the error budget to cover
      go acc marked ((i, e2):rest)
        | acc >= target = marked                -- pivot reached: stop
        | otherwise     = go (acc + e2) (i:marked) rest
      go acc marked []  = marked                -- (all elements; only if theta = 1)
  in  go 0.0 [] ranked
\end{lstlisting}

\Cref{fig:dorfler-hist} draws the algorithm as a picture: the squared indicators as
bars sorted tallest-first, and the marked prefix as the leading bars whose cumulative
area first crosses the $\theta$ line. It is the most vivid way to see what bulk
marking does---it walks in from the left, taking the biggest bars, until it has
captured fraction $\theta$ of the total area, then stops.

\begin{figure}[t]
\centering
\begin{tikzpicture}[scale=1.0]
\begin{axis}[
  width=12.5cm, height=6.0cm,
  ybar, bar width=8mm,
  xmin=0.3, xmax=8.7, ymin=0, ymax=46,
  xtick={1,2,3,4,5,6,7,8},
  xticklabels={$K_1$,$K_2$,$K_3$,$K_4$,$K_5$,$K_6$,$K_7$,$K_8$},
  xlabel={elements, sorted by indicator (largest first)},
  ylabel={$\eta_K^2$},
  axis lines=left, tick label style={font=\footnotesize},
  clip=false,
]
  % marked bars (first three: 40+25+14 = 79 >= 70 = theta*total), then unmarked
  \addplot[fill=simBgRust, draw=simRust, thick] coordinates
    {(1,40)(2,25)(3,14)};
  \addplot[fill=simBgGray, draw=simGray] coordinates
    {(4,9)(5,5)(6,3)(7,2.5)(8,1.5)};
  % pivot line between K3 (last marked) and K4 (first unmarked)
  \draw[simGreen, thick, densely dashed] (axis cs:3.5,0) -- (axis cs:3.5,44);
  \node[simGreen, font=\scriptsize, anchor=south] at (axis cs:3.5,44) {pivot};
  \node[simRust, font=\scriptsize\bfseries, align=center] at (axis cs:2,35)
    {\textsc{marked}\\(bulk of the error)};
  \node[simGray, font=\scriptsize, align=center] at (axis cs:6,18)
    {unmarked\\(refine later, if ever)};
\end{axis}
\end{tikzpicture}
\caption[D\"orfler bulk marking on a sorted histogram]{D\"orfler marking, drawn as a
sorted histogram. The per-element squared indicators $\eta_K^2$ are sorted
\emph{descending} and accumulated from the left; the marked set (rust) is the leading
prefix of bars whose cumulative area first reaches fraction $\theta$ of the total. The
pivot (green dashed line) is exactly where the running sum crosses $\theta\sum_K
\eta_K^2$. Here a few tall bars on the left carry most of the error, so a small marked
set suffices---the hallmark of a solution with concentrated features. A flat
histogram (no features) would push the pivot far to the right, marking nearly
everything, which is bulk marking correctly reporting that the error is spread out.
\Cref{wex:dorfler} traces the arithmetic for $\theta = 0.7$.}
\label{fig:dorfler-hist}
\end{figure}

\subsection{The four laws that make marking trustworthy}

The marking rule \cref{eq:dorfler} satisfies four properties that are worth stating
as laws, because the convergence theory of \cref{sec:loop} depends on them and because
they make the rule predictable.

\begin{description}[style=nextline, leftmargin=0pt, font=\normalfont\itshape]
\item[Coverage.] The marked set provably covers at least fraction $\theta$ of the
total squared error: $\sum_{K \in \mathcal S}\eta_K^2 \ge \theta\sum_K \eta_K^2$.
This is the defining inequality, and it is the hypothesis the convergence theorem
needs. The implementation verifies it as a post-condition---it never returns an
under-marking.
\item[Minimality.] Among all sets achieving coverage, $\mathcal S$ has the fewest
elements. This is the ``smallest set'' in \cref{eq:dorfler}, guaranteed by the
largest-first greedy. Minimality is what keeps the mesh from growing faster than it
must: we refine the least we can while still capturing the bulk.
\item[Monotonicity in $\theta$.] A larger target fraction marks a superset:
$\theta_1 \le \theta_2 \implies \mathcal S(\theta_1) \subseteq \mathcal S(\theta_2)$.
At $\theta \to 0^+$ the marked set shrinks to the single largest-indicator element;
at $\theta = 1$ it grows to the whole mesh. The fraction is a smooth dial between
``refine only the very worst'' and ``refine everything (uniform).''
\item[The over-mark tie-break.] The achievable cumulative fractions form a discrete
ladder---one rung per element---so the target $\theta$ usually falls \emph{between}
rungs. When it does, the rule takes the \emph{larger} set: it covers the smallest
fraction $\ge \theta$, never the largest fraction $< \theta$. This ``over-mark rather
than under-mark'' tie-break is load-bearing, not cosmetic---it is what guarantees the
coverage inequality holds with $\ge$, and the convergence theorem assumes exactly
that direction. Flipping it to under-mark would void the guarantee.
\end{description}

\begin{keyidea}
D\"orfler (\emph{bulk}) marking refines the smallest set of elements responsible for
a fixed fraction $\theta$ of the \emph{total squared error}---not a fixed count of
elements, and not a fixed fraction of elements. The algorithm is: sort the indicators
largest-first, accumulate, and stop at the pivot where the running sum reaches
$\theta$. Marking by error \emph{mass} (rather than element count) is what
equidistributes the error, and the ``over-mark on ties'' rule is what makes the
covered fraction provably $\ge \theta$---the single hypothesis behind the convergence
guarantee of \cref{sec:loop}.
\end{keyidea}

\section{Refine: subdividing the marked elements}
\label{sec:refine}

Marking produced a set of element indices. Refinement turns that set into a finer
mesh: each marked element is \emph{subdivided} into smaller children, and the global
finite-element space and operators are rebuilt on the new mesh before the next solve.

\subsection{How an element is split}

Two subdivision strategies dominate, and \cref{fig:refine} draws both for a triangle.

\emph{Red--green refinement.} A ``red'' refinement of a triangle joins the midpoints
of its three edges, splitting it into four congruent children (and a tetrahedron into
eight). This is clean and uniform, but it creates a problem at the boundary of the
refined region: a red-refined triangle has new nodes at its edge midpoints, and if the
neighbouring triangle was \emph{not} refined, those midpoint nodes sit in the middle
of the neighbour's edge with nothing to match them---a \emph{hanging node}. A hanging
node breaks conformity (the mesh is no longer a clean tiling) and must be dealt with.
The ``green'' part of red--green refinement is a closure step: neighbours adjacent to a
red element get a temporary bisection (a single cut to a vertex) that absorbs the
hanging node and restores a conforming mesh, at the cost of some sliver elements that
are undone before the next refinement pass.

\emph{Bisection.} An alternative cuts each marked element in two by a single edge
(the longest edge, or a designated ``refinement edge''), propagating the cut into
neighbours as needed to keep the mesh conforming. Repeated bisection, organised
carefully (newest-vertex bisection), provably keeps element shapes bounded---no cell
degenerates into a sliver however many times it is cut---which is what protects the
conditioning of the assembled matrix.

\begin{figure}[t]
\centering
\begin{tikzpicture}[scale=1.0, every node/.style={font=\footnotesize}]
  % ===== LEFT: red refinement (1 -> 4) =====
  \begin{scope}
    \node[font=\small\bfseries, text=simRust] at (1.0,2.5) {red: $1 \to 4$};
    \coordinate (A) at (0,0); \coordinate (B) at (2,0); \coordinate (C) at (1,1.8);
    \coordinate (mAB) at (1,0); \coordinate (mBC) at (1.5,0.9); \coordinate (mCA) at (0.5,0.9);
    \draw[simInk, thick] (A)--(B)--(C)--cycle;
    \draw[simRust, thick] (mAB)--(mBC)--(mCA)--cycle;
    \fill[simRust] (mAB) circle (1.3pt); \fill[simRust] (mBC) circle (1.3pt);
    \fill[simRust] (mCA) circle (1.3pt);
    \node[simGray, font=\scriptsize] at (1.0,-0.4) {four congruent children};
  \end{scope}
  % ===== MIDDLE: bisection (1 -> 2) =====
  \begin{scope}[xshift=4.0cm]
    \node[font=\small\bfseries, text=simGreen] at (1.0,2.5) {bisect: $1 \to 2$};
    \coordinate (A2) at (0,0); \coordinate (B2) at (2,0); \coordinate (C2) at (1,1.8);
    \coordinate (mid2) at (1,0); % midpoint of longest edge AB
    \draw[simInk, thick] (A2)--(B2)--(C2)--cycle;
    \draw[simGreen, thick] (C2)--(mid2);
    \fill[simGreen] (mid2) circle (1.3pt);
    \node[simGray, font=\scriptsize] at (1.0,-0.4) {cut the longest edge};
  \end{scope}
  % ===== RIGHT: hanging node =====
  \begin{scope}[xshift=8.2cm]
    \node[font=\small\bfseries, text=simBlue] at (1.1,2.5) {hanging node};
    % refined triangle on the left, coarse on the right sharing an edge
    \coordinate (P) at (0,0); \coordinate (Q) at (1.1,0); \coordinate (R) at (0.55,1.7);
    \coordinate (mPR) at (0.275,0.85); \coordinate (mQR) at (0.825,0.85);
    \coordinate (mPQ) at (0.55,0);
    \draw[simInk, thick] (P)--(Q)--(R)--cycle;
    \draw[simRust, thick] (mPQ)--(mQR)--(mPR)--cycle;
    % coarse neighbour to the right sharing edge QR
    \coordinate (S) at (2.0,0.85);
    \draw[simInk, thick] (Q)--(S)--(R);
    % the hanging node sits at mQR, mid of shared edge QR, unmatched in the neighbour
    \fill[simBlue] (mQR) circle (2.0pt);
    \draw[refedge, simBlue] (1.7,1.9) to[bend right=15] (mQR);
    \node[simBlue, font=\scriptsize, align=center] at (1.75,2.05)
      {unmatched\\midpoint};
    \node[simGray, font=\scriptsize, align=center] at (1.0,-0.45)
      {neighbour not refined\\$\Rightarrow$ constrain or close};
  \end{scope}
\end{tikzpicture}
\caption[Refinement of a marked element]{Subdividing a marked triangle. \emph{Left}:
\emph{red} refinement joins the three edge midpoints, producing four congruent
children (eight for a tetrahedron). \emph{Middle}: \emph{bisection} cuts a single
edge (here the longest) into two children, propagating into neighbours to stay
conforming. \emph{Right}: the difficulty either strategy must handle---a refined
element creates a new midpoint node on an edge it shares with an \emph{unrefined}
neighbour, a \emph{hanging node} unmatched on the coarse side. It is resolved either
by a \emph{green} closure (bisecting the neighbour to absorb it) or by a
\emph{constraint} that algebraically ties the hanging degree of freedom to its coarse
neighbours, keeping the global space conforming.}
\label{fig:refine}
\end{figure}

\subsection{Where our framework stops}

The subdivision geometry---which edges to cut, how to propagate cuts, how to generate
and constrain hanging nodes, how many levels of nonconformity to allow---is intricate,
and in Palace it is handled entirely by the underlying mesh library (MFEM). Our
framework treats \code{refine} as an \emph{opaque boundary}: a function
$\code{refine} : \text{Mesh} \to \text{IndexSet} \to \text{Mesh}$ that takes the
current mesh and the marked set and returns a finer conforming mesh, with the
subdivision arithmetic owned by the library. This is a deliberate scope decision, not
an oversight: the estimate and mark stages are short, well-defined algorithms we
specify fully and could re-implement; the refinement kernel is a large, mature piece
of mesh infrastructure we \emph{call} rather than \emph{re-derive}.

\begin{notationbox}
We write the refinement step as \lstinline[style=l4style]|refine mesh marked|,
returning a new mesh. Like every value in this book it is treated as
\emph{immutable}: refinement \emph{produces} a finer mesh rather than mutating the old
one in place (the underlying library does mutate a buffer, but that is below our
boundary, exactly as the in-place arithmetic of a BLAS kernel sits below
\code{linear\_combination}). The new mesh is the value the loop threads forward;
the old mesh is spent.
\end{notationbox}

What matters for the loop is only \code{refine}'s contract: given a marked set, it
returns a mesh that is finer \emph{exactly there}, conforming, and ready for the next
assemble--solve. With that contract we can close the loop.

\section{The loop as a state-generated fold}
\label{sec:loop}

We can now assemble the three stages into the loop of \cref{ch:lifecycle} and see
precisely what kind of computation it is. The shape is a \emph{fold}---the structure
\cref{ch:combinators} made the spine of the whole simulator---but a fold of a special
and important kind: one whose \emph{schedule is generated from its own state}.

\subsection{Recalling the two kinds of fold}

\Cref{ch:combinators} drew the map/fold split: a \code{fold} threads a carry through a
chain of steps, each step consuming what the last produced. The transient analysis of
\cref{ch:time} was the exemplar---it folds a time-step operator over a \emph{fixed,
pre-known} schedule of times $[t_0, t_1, \dots, t_M]$. You know the whole schedule
before you start; the fold simply walks it, threading the field state.

The AMR loop is a fold of the \emph{same combinator}, \code{fold\_solve}, but its
schedule is not fixed in advance. There is no list of meshes laid out before the loop
begins. Instead each iteration \emph{computes the next mesh from the current
solution's error}: solve on the current mesh, estimate the error, mark, refine---and
the refined mesh is the next ``schedule element,'' generated on the spot. The loop
also decides \emph{when to stop} from the state: it halts when the estimated error
falls below tolerance (or a resource budget---a maximum element count or iteration
count---is exhausted). This is a \emph{state-generated} (or \emph{data-dependent})
fold: the schedule is a function of the running carry, not a constant.

\begin{figure}[t]
\centering
\begin{tikzpicture}[font=\footnotesize]
  % ===== LEFT: fixed-schedule fold (time) =====
  \begin{scope}
    \node[font=\small\bfseries, text=simBlue] at (1.8,3.1) {fixed schedule (time-stepping)};
    \node[data, minimum width=9mm] (t0) at (-0.4,2.0) {$s_0$};
    \node[opbox, minimum width=10mm] (k1) at (1.0,2.0) {\code{step}};
    \node[opbox, minimum width=10mm] (k2) at (2.6,2.0) {\code{step}};
    \node[data, minimum width=9mm] (t3) at (4.0,2.0) {$s_M$};
    \draw[flow] (t0) -- (k1);
    \draw[flow] (k1) -- node[above,font=\scriptsize]{$s_1$} (k2);
    \draw[flow] (k2) -- (t3);
    % the pre-known schedule sitting above, feeding the steps
    \node[data, minimum width=30mm, fill=simBgBlue] (sched) at (1.8,0.9)
      {schedule $[t_0,\dots,t_M]$ (known up front)};
    \draw[refedge, simBlue] (1.0,1.55) -- (k1.south);
    \draw[refedge, simBlue] (2.6,1.55) -- (k2.south);
    \draw[flow, simBlue] (sched.north) -- (1.8,1.55);
  \end{scope}
  \draw[simGray!50, thick] (5.3,0.3) -- (5.3,3.4);
  % ===== RIGHT: state-generated fold (AMR) =====
  \begin{scope}[xshift=6.1cm]
    \node[font=\small\bfseries, text=simGreen] at (2.0,3.1) {state-generated (AMR)};
    \node[data, minimum width=11mm] (m0) at (-0.2,2.0) {mesh$_0$};
    \node[opbox, minimum width=13mm, fill=simBgGreen, draw=simGreen] (s1) at (1.6,2.0)
      {\scriptsize solve$\to$est$\to$mark$\to$refine};
    \node[data, minimum width=11mm] (m1) at (3.7,2.0) {mesh$_1$};
    \draw[flow] (m0) -- (s1);
    \draw[flow] (s1) -- (m1);
    % the back-edge: next mesh computed FROM the solve
    \draw[backflow, simGreen] (m1.south) -- ++(0,-9mm) -| node[pos=0.25,below,font=\scriptsize]{if $\eta > \text{tol}$: continue} (s1.south);
    % terminate
    \node[product, minimum width=12mm] (done) at (3.7,0.55) {output};
    \draw[flow] (m1.east) ++(0.0,0) .. controls +(0.6,0) and +(0.6,0) .. (done.east)
      node[midway, right, font=\scriptsize] {$\eta \le \text{tol}$};
    \node[simGreen, font=\scriptsize, align=center] at (1.6,0.95)
      {next mesh is\\\emph{computed}, not listed};
  \end{scope}
\end{tikzpicture}
\caption[Fixed-schedule versus state-generated fold]{Two folds of the same
combinator. \emph{Left}: time-stepping (\cref{ch:time}) folds over a schedule
$[t_0,\dots,t_M]$ that is \emph{known before the loop starts}---the steps read their
times from a pre-built list. \emph{Right}: the AMR loop folds the
solve--estimate--mark--refine step, but its schedule is \emph{generated from the
state}: the next mesh is computed from the current solution's error, and the loop
continues or halts depending on whether the estimated error $\eta$ has fallen below
tolerance. Same carry-threading \code{fold} structure; the difference is whether the
schedule is a constant (left) or a function of the carry (right). This is why the
adaptive loop is a fold with a \emph{data-dependent} stopping condition, not a
\code{map} over a fixed range.}
\label{fig:amrfold}
\end{figure}

\subsection{Writing the fold}

In the framework's vocabulary the loop is one \code{fold\_solve} whose step is the
three-stage body and whose carry bundles the mesh with the running error indicators.
\Cref{alg:amr-fold} writes it. Read it against \cref{fig:amrfold}: the carry is the
mesh (plus the indicators and the scalar error), the step is solve--estimate--mark--
refine, and the loop predicate tests the state.

\begin{lstlisting}[style=l4style,caption={The AMR loop as a state-generated \code{fold\_solve}. The schedule (the sequence of meshes) is generated from the carry, not supplied as a fixed list.},label={alg:amr-fold}]
-- the carry threaded through the fold: a mesh and its error state
type AmrCarry = { mesh : Mesh, etas : [Float], err : Float, ndof : Int }

amr_step :: Estimator -> Float -> AmrCarry -> AmrCarry
amr_step est theta carry =
  let (field, ndof) = solve carry.mesh                 -- assemble + solve (per-driver)
      etas          = flux_recovery_estimate est field -- ESTIMATE (Section 37.2)
      err           = nrm2 etas                         -- global error scalar
      marked        = dorfler_mark theta etas           -- MARK    (Section 37.3)
      mesh'         = refine carry.mesh marked           -- REFINE  (Section 37.4)
  in  { mesh = mesh', etas = etas, err = err, ndof = ndof }

-- the loop: iterate the step WHILE the state says "not yet accurate enough"
adaptive :: Estimator -> RefineConfig -> Mesh -> Product
adaptive est cfg mesh0 =
  let stop carry = carry.err <= cfg.tol            -- data-dependent stop ...
                || exhausted carry.ndof cfg         -- ... or budget hit
      final = iterate_while (not . stop) (amr_step est cfg.theta) (seed mesh0)
  in  product_of final
\end{lstlisting}

Two features of \cref{alg:amr-fold} are the whole point. First, the step \emph{reads
its carry}: \code{amr\_step} solves on \code{carry.mesh}, the mesh the previous step
produced. That read is what makes this a fold and not a \code{map}---the iterations do
not commute and cannot run in parallel or be reordered, because mesh $n+1$ is
\emph{computed from} the solution on mesh $n$. Second, the loop's \code{stop}
predicate is a function of the \emph{state} (the error \code{carry.err} and the
DOF count \code{carry.ndof}), not of a fixed iteration count. The number of
refinement passes is not known in advance; it is however many it takes to drive the
error below tolerance. This is the \code{iterate\_while} of \cref{ch:fold} at its most
characteristic---a loop whose length is discovered, not declared.

\begin{pitfall}
It is tempting to picture the AMR loop as a \code{map} over a list of meshes
``$[\text{mesh}_0, \text{mesh}_1, \dots]$.'' It is not, and the difference is not
pedantic. A \code{map}'s elements are independent and known up front; here neither
holds. Mesh$_{n+1}$ \emph{does not exist} until mesh$_n$ has been solved, estimated,
and refined---it is \emph{produced by} the step, from the carry. You cannot list the
schedule before running the loop, and you cannot run the iterations out of order or in
parallel. The AMR loop is intrinsically a sequential, state-threaded fold; reading it
as a map would promise a parallelism that the data dependence forbids.
\end{pitfall}

\subsection{Why it converges: D\"orfler's theorem}

The loop would be useless if refining did not reliably reduce the error. It does, and
the guarantee is the reason bulk marking, not the naive strategies of
\cref{sec:mark}, is used. The result, due to D\"orfler, is a \emph{contraction}
theorem.

\begin{theorem}[D\"orfler, contraction of adaptive refinement; informal]
\label{thm:dorfler}
For a model elliptic problem, with the error estimator equivalent to the true error
and elements marked by the bulk criterion \cref{eq:dorfler} with fraction $\theta$,
each adaptive pass reduces the error by at least a fixed factor: there is a constant
$\rho = \rho(\theta) < 1$, independent of the mesh, with
\begin{equation}
\norm{u - u_{h}^{(n+1)}} \;\le\; \rho \, \norm{u - u_{h}^{(n)}}
\label{eq:contraction}
\end{equation}
where $u_h^{(n)}$ is the discrete solution after $n$ adaptive passes. Consequently the
error decreases geometrically in the number of passes, and---combined with the
minimality of marking---the adaptive sequence achieves the optimal error-versus-DOF
rate.
\end{theorem}

The content of \cref{thm:dorfler} is that the loop is a genuine \emph{contraction}:
every pass shrinks the error by a guaranteed factor $\rho < 1$, so $n$ passes shrink
it by $\rho^n \to 0$. The bulk-marking coverage law of \cref{sec:mark} is exactly the
hypothesis: because the marked set carries at least fraction $\theta$ of the error,
refining it removes at least a corresponding fixed share, which is what forces
$\rho < 1$. Had we marked by the naive ``top $10\%$ of elements'' rule, no such
guarantee holds---one can construct meshes where the top $10\%$ of elements carry a
vanishing fraction of the error, and refinement stalls. Bulk marking buys the
theorem; the theorem is why the loop is trusted to terminate with a good answer.

\section{Reading the payoff: adaptive versus uniform convergence}
\label{sec:convergence}

The promise of \cref{sec:why-adapt} was that adaptivity reaches a target accuracy with
far fewer unknowns. \Cref{fig:convergence} is where that promise is cashed: a
log--log plot of error against the number of degrees of freedom $N$, for uniform and
adaptive refinement on a corner-singularity problem. It is the single most important
diagnostic an AMR practitioner reads.

On a log--log plot, a convergence rate $\text{error} = C\,N^{-r}$ appears as a
\emph{straight line of slope $-r$}: the steeper the line, the faster the error falls
per unknown added. Two lines tell the story.

\emph{Uniform refinement} (the shallow line) is rate-limited by the singularity. The
corner's $r^{-\gamma}$ blow-up caps the smooth-solution rate: no matter how finely you
cut \emph{every} cell, the under-resolved corner drags the global error down only as
$N^{-\gamma/2}$, a shallow slope. The interior is over-resolved and the corner is
under-resolved, and the corner wins.

\emph{Adaptive refinement} (the steep line) recovers the full rate. By concentrating
cells at the corner---equidistributing the error---it resolves the singularity with a
small, local investment of unknowns and lets the interior stay coarse. The recovered
slope is $N^{-p/2}$, the rate a smooth solution would enjoy. The gap between the lines
is the payoff: read horizontally at any fixed error level, the adaptive mesh needs
\emph{orders of magnitude} fewer unknowns to reach it.

\begin{figure}[t]
\centering
\begin{tikzpicture}[scale=1.0]
\begin{axis}[
  width=12cm, height=7.0cm,
  xmode=log, log basis x=10, ymode=log, log basis y=10,
  xmin=2e2, xmax=3e5, ymin=2e-4, ymax=2e-1,
  xlabel={degrees of freedom $N$}, ylabel={error $\|u-u_h\|$},
  axis lines=left, tick align=outside, grid=major, grid style={simGray!18},
  legend style={at={(0.03,0.03)}, anchor=south west, font=\footnotesize,
    draw=simGray!50},
  every axis plot/.append style={very thick},
  tick label style={font=\footnotesize},
]
  % uniform: slope ~ -1/3 (singularity-limited), shallow
  \addplot[simRust, mark=*, mark size=1.6pt] coordinates {
    (400,1.0e-1)(1600,6.3e-2)(6400,4.0e-2)(25600,2.5e-2)(102400,1.6e-2)};
  \addlegendentry{uniform: slope $-\gamma/2$ (singularity-limited)}
  % adaptive: slope ~ -1 (full rate), steep
  \addplot[simBlue, mark=square*, mark size=1.6pt] coordinates {
    (300,8.0e-2)(900,2.7e-2)(2700,9.0e-3)(8100,3.0e-3)(24300,1.0e-3)(72900,3.3e-4)};
  \addlegendentry{adaptive: slope $-p/2$ (full smooth rate)}
  % horizontal "same accuracy" guide at error = 3.3e-3
  \draw[simGreen, densely dashed, thick] (axis cs:2e2,3.3e-3) -- (axis cs:3e5,3.3e-3);
  \node[simGreen, font=\scriptsize, anchor=south west] at (axis cs:2.2e2,3.5e-3)
    {same accuracy};
  % the DOF gap at that level: adaptive ~7e3, uniform off-chart (extrapolated)
  \node[simGreen, font=\scriptsize, anchor=west] at (axis cs:1.1e4,1.8e-3)
    {$\sim\!10\times$ fewer DoFs};
\end{axis}
\end{tikzpicture}
\caption[Adaptive versus uniform convergence]{The payoff, on a log--log error-vs-DOF
plot for a corner-singularity problem. \emph{Uniform} refinement (rust) is limited by
the singularity to the shallow rate $\|u-u_h\| = \mathcal{O}(N^{-\gamma/2})$---cutting
every cell cannot outpace the under-resolved corner. \emph{Adaptive} refinement (blue)
concentrates unknowns at the corner and recovers the full smooth-solution rate
$\mathcal{O}(N^{-p/2})$, a steeper line. Reading horizontally at any fixed accuracy
(green dashed), the adaptive mesh reaches it with far fewer degrees of freedom---the
gap widens without bound as the target accuracy tightens. This separation, not any
single solve, is what adaptive mesh refinement buys.}
\label{fig:convergence}
\end{figure}

\begin{keyidea}
On a log--log error-versus-DOF plot, a convergence rate is a straight line and a
\emph{steeper line is better}. Uniform refinement on a singular problem is stuck on a
shallow line (the singularity caps the rate); adaptive refinement recovers the steep,
full-rate line by equidistributing the error. The horizontal gap between the lines---
how many fewer unknowns adaptivity needs for the same accuracy---widens as the target
accuracy tightens. That gap is the entire economic case for the
estimate--mark--refine loop.
\end{keyidea}

\section{Worked examples}
\label{sec:worked}

\begin{workedexample}{D\"orfler marking on eight elements, $\theta = 0.7$}{dorfler}
\label{wex:dorfler}
We run the bulk-marking algorithm of \cref{sec:mark} explicitly on a concrete list of
indicators, showing exactly which elements are marked and verifying the coverage law.

\emph{The data.} Eight elements with squared error indicators (in arbitrary units)
\[
\eta^2 = (\,4,\;\;36,\;\;9,\;\;1,\;\;25,\;\;2,\;\;16,\;\;7\,),
\qquad K = 1,\dots,8.
\]
The total squared error is
\[
T = \sum_K \eta_K^2 = 4+36+9+1+25+2+16+7 = 100,
\]
and with $\theta = 0.7$ the error budget to cover is $\theta T = 70$.

\emph{Step 1 --- sort descending.} Rank the elements largest-first, carrying their
original indices:
\[
\begin{array}{c|cccccccc}
\text{rank} & 1 & 2 & 3 & 4 & 5 & 6 & 7 & 8 \\\hline
\text{element } K & 2 & 5 & 7 & 3 & 8 & 1 & 6 & 4 \\
\eta_K^2 & 36 & 25 & 16 & 9 & 7 & 4 & 2 & 1
\end{array}
\]

\emph{Step 2 --- accumulate, and find the pivot.} Walk the sorted list, summing until
the running total reaches the budget $70$:
\[
\begin{array}{c|c|c|c}
\text{take element} & \eta^2 & \text{running sum} & \ge 70? \\\hline
K=2 & 36 & 36 & \text{no} \\
K=5 & 25 & 61 & \text{no} \\
K=7 & 16 & 77 & \textbf{yes} \;\leftarrow\; \text{pivot} \\
\end{array}
\]
The running sum first reaches $70$ when we add the third element, $K = 7$, bringing
the total to $77$. The pivot is here: we stop.

\emph{Step 3 --- the marked set.} The marked elements are the prefix taken up to and
including the pivot:
\[
\mathcal S = \{\,K_2,\;K_5,\;K_7\,\} = \{2, 5, 7\},
\qquad |\mathcal S| = 3.
\]

\emph{Verify the laws.} \textbf{Coverage:} the marked set carries
$36 + 25 + 16 = 77$ of the total $100$, so the covered fraction is $0.77 \ge 0.70 =
\theta$. The coverage law holds. \textbf{Minimality:} could a smaller set cover
$70$? A two-element set tops out at $36 + 25 = 61 < 70$ (the two largest), so no two
elements suffice---three is the minimum, and the greedy found it. \textbf{Over-mark
tie-break:} the achievable cumulative fractions are the ladder
$0.36,\,0.61,\,0.77,\,\dots$; the target $0.70$ falls \emph{between} the second rung
($0.61$) and the third ($0.77$), so the rule takes the larger set ($0.77 \ge 0.70$),
never the smaller ($0.61 < 0.70$). It over-marks, exactly as the convergence theory
requires.

Three of eight elements---the bulk of the error---are marked for refinement; the other
five stay coarse. Notice the marking ignored element \emph{count} entirely: it took
whatever number of elements was needed to reach $70\%$ of the error mass, which here
happened to be three. On a mesh where the error were more concentrated (say all $100$
in one element) it would mark just one; spread evenly over the eight (each
$\eta_K^2 = 12.5$) it would mark six. Bulk marking adapts the set size to the
error distribution.
\end{workedexample}

\begin{workedexample}{A ZZ flux indicator on a 1-D mesh}{zz}
\label{wex:zz}
We compute the Zienkiewicz--Zhu indicator of \cref{sec:estimate} on a tiny
one-dimensional example, showing it is largest where the solution bends most.

\emph{The setup.} Take the interval $[0,4]$ meshed into four unit elements
$K_1=[0,1],\,K_2=[1,2],\,K_3=[2,3],\,K_4=[3,4]$, with a piecewise-linear solution
$u_h$ whose nodal values are
\[
u_h(0)=0,\quad u_h(1)=1,\quad u_h(2)=2,\quad u_h(3)=4,\quad u_h(4)=8.
\]
The field is the (1-D) gradient $\diffop u_h = u_h'$; take material weight $Q=1$ so
the flux is $\mat F = u_h'$. On each element $u_h$ is linear, so $u_h'$ is the
constant slope:
\[
\mat F|_{K_1}=\frac{1-0}{1}=1,\quad
\mat F|_{K_2}=\frac{2-1}{1}=1,\quad
\mat F|_{K_3}=\frac{4-2}{1}=2,\quad
\mat F|_{K_4}=\frac{8-4}{1}=4.
\]
This is the \emph{discontinuous} computed flux: a constant on each element, jumping
$1\!\to\!1\!\to\!2\!\to\!4$ at the interior nodes. The jumps are $0$ (at $x=1$), $1$
(at $x=2$), $2$ (at $x=3$)---the flux is smooth on the left and increasingly broken
on the right, where the solution bends upward.

\emph{Recover a smooth flux.} A simple nodal recovery (the 1-D analogue of the $L^2$
projection \cref{eq:zz-recover}) assigns each interior node the average of the two
adjacent element fluxes, and each end node the one flux it touches:
\[
G(0)=1,\quad
G(1)=\tfrac{1+1}{2}=1,\quad
G(2)=\tfrac{1+2}{2}=1.5,\quad
G(3)=\tfrac{2+4}{2}=3,\quad
G(4)=4.
\]
The recovered flux $G$ is the continuous, piecewise-linear function through these
nodal values---a smooth stand-in for the true flux.

\emph{The per-element indicator.} On element $K$, the indicator
$\eta_K^2 = \int_K (\mat F - G)^2\,dx$ measures the gap between the constant element
flux $\mat F|_K$ and the linear recovered flux $G$. On a unit element from node values
$G_{\text{left}}, G_{\text{right}}$ against the constant $\mat F$, the integral of the
squared difference of a constant and a line evaluates to
\[
\eta_K^2 = \int_0^1\!\Bigl(\mat F - \bigl[G_{\text{left}} + (G_{\text{right}}-G_{\text{left}})\,s\bigr]\Bigr)^2 ds .
\]
Computing each:
\[
\begin{array}{c|c|c|c}
\text{element} & \mat F|_K & (G_{\text{left}},G_{\text{right}}) & \eta_K^2 \\\hline
K_1 & 1 & (1,\,1)   & 0 \\
K_2 & 1 & (1,\,1.5) & \int_0^1(1-(1+0.5s))^2\,ds = \int_0^1 (0.5s)^2 ds = \tfrac{1}{12}\approx 0.083 \\
K_3 & 2 & (1.5,\,3) & \int_0^1(2-(1.5+1.5s))^2 ds = \int_0^1 (0.5-1.5s)^2 ds = 0.25 \\
K_4 & 4 & (3,\,4)   & \int_0^1(4-(3+s))^2 ds = \int_0^1 (1-s)^2 ds = \tfrac{1}{3}\approx 0.333 \\
\end{array}
\]

\emph{Read the result.} The indicators are
\[
\eta^2 = (\,0,\;\;0.083,\;\;0.25,\;\;0.333\,),
\qquad K_1,\dots,K_4,
\]
\emph{increasing left to right}---exactly tracking where the solution bends most. On
$K_1$, where the flux is perfectly constant ($1\!\to\!1$, no jump), the indicator is
zero: the discretization already resolves this region, and the estimator correctly
reports nothing to do. On $K_4$, where the flux jumps hardest ($2\!\to\!4$ across the
node and the solution curves upward fastest), the indicator is largest. Feed these
four numbers to the D\"orfler marker of \cref{wex:dorfler}: with $\theta = 0.7$ and
total $0.666$, the budget is $0.466$; the sorted prefix $0.333 + 0.25 = 0.583 \ge
0.466$ marks $\{K_4, K_3\}$---the two right-hand elements, precisely the region that
needs refining. The estimator found the feature, and the marker selected it, with no
knowledge of the true solution at all.
\end{workedexample}

\section{Where this leaves us}

We have opened the three black boxes of the adaptive loop. \emph{Estimate} is
Zienkiewicz--Zhu flux recovery: smooth the discontinuous computed flux by an $L^2$
projection and take the per-element gap as the indicator---local, cheap, quadratic in
the field, and zero exactly when the flux is already smooth. \emph{Mark} is
D\"orfler bulk marking: sort the indicators, accumulate, and take the smallest prefix
carrying fraction $\theta$ of the total squared error---marking by error mass, with an
over-mark tie-break that makes the coverage provable. \emph{Refine} subdivides the
marked elements (red--green or bisection, keeping the mesh conforming) and rebuilds the
space---a mesh-library operation we treat as a boundary. And the whole loop is a
\code{fold\_solve} whose schedule is \emph{generated from its own state}: the next
mesh is computed from the current solution's error, and the fold runs until the error
falls below tolerance. D\"orfler's contraction theorem guarantees each pass shrinks
the error by a fixed factor, and the log--log convergence plot shows the payoff: the
full smooth-solution rate on a singular problem, at orders of magnitude fewer
unknowns than uniform refinement.

This closes the adaptive loop that \cref{ch:lifecycle} sketched and
\cref{ch:combinators} typed as a state-threaded fold. The error indicators the loop
consumes are themselves \emph{reductions}---a per-element norm of the flux gap, and a
global norm of the indicator vector for the stopping test. Those reductions, and the
broader family of fold-to-a-scalar operations the simulator runs, are the subject of
\cref{ch:reductions}, which we take up next.

\summarysection
Adaptive mesh refinement reaches a target accuracy with far fewer unknowns than
uniform refinement by \emph{equidistributing the error}: it puts small cells only
where the solution's features (corners, interfaces, field concentrations) make the
local error large, leaving smooth regions coarse. It does this automatically through
the estimate--mark--refine loop of \cref{ch:lifecycle}. \emph{Estimate} uses the
Zienkiewicz--Zhu flux-recovery indicator $\eta_K^2 = \int_K \norm{\mat F - G}^2$,
where $\mat F = Q\,\diffop u_h$ is the discontinuous computed flux and $G = \mat
M^{-1}(\text{Flux}\cdot u_h)$ is its smooth $L^2$ projection---a local, cheap,
quadratic indicator that vanishes exactly when the flux is already smooth and is large
where the solution bends most. \emph{Mark} uses D\"orfler bulk marking: the smallest
set of elements whose summed squared error covers a fraction $\theta$ (default $0.7$)
of the total, found by sorting the indicators largest-first and accumulating to a
pivot, with an over-mark-on-ties rule that makes coverage provably $\ge \theta$.
\emph{Refine} subdivides the marked elements (red--green or bisection, kept conforming)
and rebuilds the space---a mesh-library operation at our framework's boundary. The
whole loop is a \code{fold\_solve} whose schedule is generated from its own state: the
next mesh is computed from the current solution's error, and the loop halts when the
error drops below tolerance or a budget is hit---a state-generated fold, not a map
over a fixed schedule. D\"orfler's theorem guarantees geometric error contraction per
pass, and the log--log error-vs-DOF plot shows the payoff: the full smooth-solution
convergence rate on a singular problem, at a fraction of the unknowns uniform
refinement would demand.

\furtherreading
The bulk-marking strategy and its convergence theory are due to
D\"orfler~\citep{dorfler1996}, whose ``convergent adaptive algorithm for Poisson's
equation'' introduced the criterion \cref{eq:dorfler} and proved the contraction
\cref{eq:contraction}; it remains the standard reference and is well worth reading for
the elegance of the marking argument. The flux-recovery error estimator is due to
Zienkiewicz and Zhu~\citep{zienkiewiczzhu1987}, whose ``simple error estimator and
adaptive procedure'' is one of the most cited papers in computational engineering;
the recovered-gradient idea of \cref{sec:estimate} is theirs. For the mathematical
theory tying it all together---a-posteriori estimates, the role of bulk marking in
optimal convergence rates, and the approximation theory behind the rate recovery of
\cref{sec:convergence}---Brenner and Scott~\citep{brennerscott2008} is the standard
graduate reference, building directly on the C\'ea-lemma foundation of
\cref{ch:galerkin}. The mesh-refinement geometry of \cref{sec:refine}---red--green
refinement, newest-vertex bisection, and the hanging-node constraints---is treated in
the finite-element-implementation literature; readers wanting the algorithmic detail
of the boundary we declined to cross will find it there.

\exercisessection
\begin{enumerate}[nosep]
  \item \textbf{Marking by hand.} Six elements have squared indicators
  $\eta^2 = (10,\,2,\,30,\,5,\,18,\,35)$. With $\theta = 0.6$, run the D\"orfler
  algorithm: sort descending, accumulate, find the pivot, and write the marked set.
  Verify the coverage law and confirm minimality (no smaller set reaches the budget).

  \item \textbf{The dial $\theta$.} For the data of \cref{wex:dorfler}
  ($\eta^2 = (4,36,9,1,25,2,16,7)$, total $100$), find the marked set for
  $\theta = 0.3$, $\theta = 0.7$, and $\theta = 0.95$. Confirm the monotonicity law
  $\mathcal S(0.3) \subseteq \mathcal S(0.7) \subseteq \mathcal S(0.95)$. What is the
  marked set as $\theta \to 0^+$? As $\theta \to 1$?

  \item \textbf{Why not ``top $10\%$ of elements''?} Construct a list of $20$
  squared indicators in which the top $10\%$ of \emph{elements} (the two largest)
  carry only a tiny fraction of the total error. Explain, using D\"orfler's theorem
  \cref{thm:dorfler}, why refining those two would stall convergence, and how bulk
  marking avoids the trap. (Hint: make the error spread nearly uniformly.)

  \item \textbf{The ZZ indicator vanishes on smooth flux.} Suppose the computed flux
  $\mat F$ is already continuous and piecewise linear---so its $L^2$ projection $G$
  onto the smooth space equals $\mat F$ itself. Show every $\eta_K^2 = 0$. Then argue,
  conversely, that a single jump in $\mat F$ forces at least one $\eta_K^2 > 0$. Which
  property of the indicator (\cref{sec:estimate}) is this?

  \item \textbf{Quadratic, not linear.} Using $\eta_K^2 = \int_K \norm{\mat F - G}^2$
  and the linearity of the flux map and projection in $u_h$, show
  $\eta_K^2(\alpha u_h) = \alpha^2\,\eta_K^2(u_h)$ but give a concrete two-element
  example where $\eta_K^2(u_h + w_h) \ne \eta_K^2(u_h) + \eta_K^2(w_h)$. Why does this
  mean the estimator is \emph{not} an \code{apply\_linop}-style linear operator?

  \item \textbf{Map or fold?} Explain why the AMR loop of \cref{alg:amr-fold} is a
  \code{fold\_solve} and not a \code{map} over a list of meshes. Identify the carry,
  the step, and the data-dependent stopping condition. Which line of the pseudocode is
  the one that forbids reordering the iterations, and why?

  \item \textbf{Reading the convergence curve.} On the log--log plot of
  \cref{fig:convergence}, the uniform line has slope $-\tfrac13$ and the adaptive line
  slope $-1$. (a)~If a target error of $10^{-3}$ requires $N_{\text{ad}} = 8\times10^3$
  adaptive unknowns, estimate how many \emph{uniform} unknowns the shallow line would
  need to reach the same error by extrapolation. (b)~By what factor does the DOF
  advantage of adaptivity grow each time you tighten the target error by $10\times$?

  \item \textbf{Composite indicators.} A driven (frequency-domain) analysis computes
  \emph{two} ZZ indicators per element---an electric-flux one $\eta_{K,E}^2$ and a
  magnetic-flux one $\eta_{K,B}^2$---and combines them as
  $\eta_K^2 = \eta_{K,E}^2 + \eta_{K,B}^2$. Why is the combination done in the
  \emph{squared} domain rather than on the indicators $\eta_K$ themselves? (Hint: which
  quantity is it that the total error $\sum_K \eta_K^2$ must equal, and how does that
  decompose?) Relate this to the channel-additivity of \cref{ch:reductions}'s
  reduction verbs.
\end{enumerate}

\chapter{Output Reductions: From Fields to Physical Quantities}
\label{ch:reductions}

\chapterorient{Every analysis in this book ends the same way. A driver has
worked hard to produce a \emph{family of fields}---one potential per terminal,
one scattered field per frequency, one eigenvector per resonant mode---and now
that family must be boiled down to the single number, or small table, the
engineer actually asked for: a capacitance, an S-matrix, a list of resonant
frequencies and their quality factors. This last stage is the \texttt{reduce}
of the skeleton (\cref{ch:situating}), and it is where the physics finally
re-enters: a Gram matrix \emph{is} a capacitance, a projection \emph{is} an
S-parameter, an eigenvalue's imaginary part \emph{is} a damping rate. The
surprise of this chapter---the same surprise the whole book keeps finding---is
that the dozen physical outputs Palace can produce are built from only
\emph{three} reduction shapes. We meet all three in detail, compute two of them
by hand, and close Part~IV by handing the assembled machinery forward to
Part~V, where the modalities compose it.}

\section{The setting: a solution family, and the answer hidden in it}
\label{sec:red-setting}

By the time we reach the reduce stage, the hard work is done. \Cref{ch:situating}
drew the skeleton---\texttt{config} $\to$ \texttt{assemble} $\to$ \texttt{solve}
$\to$ \texttt{reduce} $\to$ \texttt{product}---and Parts~II through~IV have now
filled in every box but the last. Assembly turns the continuous physics into
operators (\cref{ch:assembly}); the solve stage, in one of its four corners
(\cref{ch:situating}), turns those operators into fields. What the solve hands
to the reduce stage is never a single field. It is a \emph{family}:

\begin{itemize}[nosep]
  \item the electrostatic and magnetostatic drivers produce one solved field
  \emph{per excitation}---per terminal held at unit voltage, per source carrying
  unit current (\cref{ch:electrostatics}, \cref{ch:magnetostatics});
  \item the driven driver produces one scattered field \emph{per frequency} in
  the sweep, and within each, one column \emph{per driven port}
  (\cref{ch:driven});
  \item the eigenmode and boundary-mode drivers produce one eigenvector
  \emph{per converged mode} (\cref{ch:eigenmodes}, \cref{ch:waveports}).
\end{itemize}

The reduce stage takes that family and collapses it to the product. The
physical answer is, in a precise sense, already \emph{present} in the family---it
is a pairing, a projection, or a per-item readout away. What the reduction does
is extract it (\cref{fig:red-pipeline}).

\begin{figure}[t]
\centering
\begin{tikzpicture}[node distance=8mm and 12mm]
  \node[stage] (solve) {\textbf{solve}};
  \node[data, right=14mm of solve, align=center, text width=26mm] (fam)
    {solution family\\[1pt]\footnotesize $\{\vx_i\}$, $\{\vE_j(\omega)\}$,
     $\{(\lambda_m,\vx_m)\}$};
  \node[stage, right=14mm of fam, fill=simBgGold, draw=simGold] (red)
    {\textbf{reduce}};
  \node[product, right=12mm of red, align=center, text width=24mm] (prod)
    {physical\\product\\[1pt]\footnotesize $C$, $S$, $(f,Q)$, $\dots$};
  \draw[flow] (solve) -- (fam);
  \draw[flow] (fam) -- node[above,font=\scriptsize]{map-then-collect} (red);
  \draw[flow] (red) -- (prod);
  \node[font=\scriptsize\itshape, simGray, below=2mm of red, align=center]
    {embarrassingly parallel:\\no inter-element state};
\end{tikzpicture}
\caption[The family-to-product pipeline]{The reduce stage in one picture. The
solve produces not a single field but a \emph{family}---one member per
excitation, per frequency, or per mode. The reduction maps a per-element
computation over that family and collects the results into the small physical
product. Because no element depends on another, every reduction in this chapter
is a \texttt{map}-then-collect (or \texttt{map}-then-reduce), the embarrassingly
parallel combinator of \cref{ch:combinators}: there is no state threaded between
elements, so the work is independent and reorderable.}
\label{fig:red-pipeline}
\end{figure}

That structural observation---\emph{every} reduction is a map over the family
with no inter-element state---is worth stating up front, because it is what
makes the reductions so uniform. None of them is a \texttt{fold}; none threads a
running accumulator from one family member to the next. Each is the
embarrassingly parallel \texttt{map} combinator of \cref{ch:combinators},
followed by a collect into a matrix or a table. In the framework's vocabulary
they are \emph{pure value-producing reductions}: no \code{Solve} monad, no
convergence predicate, no carried state---post-processing readouts that take a
solved family and return a value.

\begin{keyidea}
The reduce stage turns a \emph{solution family} into the user-facing physical
quantity. Every reduction is a \texttt{map} over the family with \emph{no
inter-element state}---embarrassingly parallel, never a fold
(\cref{ch:combinators}). What varies between reductions is only the per-element
computation and how the results are collected. There are exactly three
collection \emph{shapes}, and the rest of the chapter is those three.
\end{keyidea}

\Cref{ch:situating} named the three shapes when it first sketched the framework;
here we give each its full treatment. They are, in order: the \emph{rank-2
symmetric Gram} (capacitance, inductance); the \emph{rank-2 port-projection}
(S-parameters); and the \emph{rank-1 scalar-or-mode table} (energy and
participation, eigenfrequency and $Q$, waveguide modes). \Cref{fig:red-taxonomy}
is the hero figure of the chapter: the three shapes side by side, the taxonomy
that organizes everything that follows.

\begin{figure}[p]
\centering
\begin{tikzpicture}[font=\footnotesize]
  % ===== (a) symmetric Gram =====
  \begin{scope}[shift={(0,0)}]
    \foreach \i in {0,1,2}{
      \foreach \j in {0,1,2}{
        \node[draw=simGray, semithick,
          fill={\ifnum\i=\j simBgRust\else simBgBlue\fi},
          minimum size=7mm, inner sep=0pt]
          (g\i\j) at (\j*0.72, -\i*0.72) {};}}
    % mirror diagonal
    \draw[densely dashed, simRust, semithick] (-0.5,0.5) -- (1.94,-1.94);
    % an off-diagonal pair, mirrored
    \node[font=\scriptsize, simBlue] at (1.44,0.0) {$G_{02}$};
    \node[font=\scriptsize, simBlue] at (0.0,-1.44) {$G_{20}$};
    \node[font=\scriptsize\bfseries, align=center] at (0.72,-2.55)
      {\textsc{(a)} rank-2 Gram};
    \node[font=\scriptsize, align=center] at (0.72,-3.05)
      {$G_{ij}=w(i,j)\,\vx_j^{\!\top}\Kop\,\vx_i$};
    \node[font=\scriptsize, align=center] at (0.72,-3.5)
      {\emph{symmetric}: $G_{ij}=G_{ji}$};
    \node[font=\scriptsize\itshape, text=simRust, align=center] at (0.72,-4.0)
      {capacitance,\\inductance};
  \end{scope}

  % ===== (b) non-symmetric projection =====
  \begin{scope}[shift={(4.6,0)}]
    \foreach \i in {0,1,2}{
      \foreach \j in {0,1,2}{
        \node[draw=simGray, semithick, fill=simBgGreen,
          minimum size=7mm, inner sep=0pt]
          (p\i\j) at (\j*0.72, -\i*0.72) {};}}
    \node[font=\scriptsize, text=simGreen] at (0.72,0.6) {drive port $j\to$};
    \node[font=\scriptsize, rotate=90, text=simGreen] at (-0.72,-0.72)
      {receiver port $i\to$};
    \node[font=\scriptsize\bfseries, align=center] at (0.72,-2.55)
      {\textsc{(b)} rank-2 projection};
    \node[font=\scriptsize, align=center] at (0.72,-3.05)
      {$S_{ij}=\sigma_{ij}\,(\inner{\vs_i}{\vE_j}-\delta_{ij})$};
    \node[font=\scriptsize, align=center] at (0.72,-3.5)
      {\emph{not} symmetric, complex};
    \node[font=\scriptsize\itshape, text=simGreen, align=center] at (0.72,-4.0)
      {S-parameters};
  \end{scope}

  % ===== (c) rank-1 table =====
  \begin{scope}[shift={(9.2,0)}]
    \foreach \i in {0,1,2,3}{
      \node[draw=simGray, semithick, fill=simBgGold,
        minimum width=22mm, minimum height=5.5mm, inner sep=0pt]
        (r\i) at (0, -\i*0.66) {};}
    \node[font=\scriptsize] at (0,0)     {mode 1: $f_1,\,Q_1$};
    \node[font=\scriptsize] at (0,-0.66) {mode 2: $f_2,\,Q_2$};
    \node[font=\scriptsize] at (0,-1.32) {mode 3: $f_3,\,Q_3$};
    \node[font=\scriptsize] at (0,-1.98) {$\vdots$};
    \node[font=\scriptsize\bfseries, align=center] at (0,-2.55)
      {\textsc{(c)} rank-1 table};
    \node[font=\scriptsize, align=center] at (0,-3.05)
      {one row per item};
    \node[font=\scriptsize, align=center] at (0,-3.5)
      {a \texttt{map}, no pairing};
    \node[font=\scriptsize\itshape, text=simGold, align=center] at (0,-4.0)
      {energy$+p$, $f{+}Q$,\\waveguide modes};
  \end{scope}
\end{tikzpicture}
\caption[The three reduction shapes]{The three reduction shapes---the chapter's
organizing taxonomy, first sketched in \cref{ch:situating}. \textsc{(a)} A
\emph{rank-2 symmetric Gram} pairs a solution family with itself through an
operator, $G_{ij}=w(i,j)\,\vx_j^{\!\top}\Kop\,\vx_i$; its symmetry (dashed
mirror) is structural, so only the upper triangle is computed and the lower is
copied ($G_{02}=G_{20}$). \textsc{(b)} A \emph{rank-2 projection} also yields a
matrix, but it projects each driven field onto a basis of port modes; its two
indices play different roles, it is complex, and it is in general \emph{not}
symmetric. \textsc{(c)} A \emph{rank-1 table} is a one-dimensional list, one row
per item, with no pairing at all---a pure \texttt{map}. The single most
over-unified point in the design is the line between \textsc{(a)} and
\textsc{(b)}: two matrix-shaped results are \emph{not} the same reduction
(\cref{pit:gram-not-proj}).}
\label{fig:red-taxonomy}
\end{figure}

\section{Shape 1: the symmetric Gram reduction}
\label{sec:gram}

The first shape is a \emph{Gram matrix}: a symmetric matrix of operator-weighted
inner products formed by pairing a solution family with itself. It is the
reduction behind both capacitance and inductance, and it is the cleanest
instance in the book of ``one verb, two products.''

\subsection{The verb}

Given a solution family $\vx_0,\dots,\vx_{m-1}$, the assembled \emph{energy
operator} $\Kop$ (symmetric positive-definite), and a scalar \emph{weight}
function $w(i,j)$, the Gram reduction computes the $m\times m$ matrix
\begin{equation}
  \boxed{\;G_{ij} \;=\; w(i,j)\;\vx_j^{\!\top}\,\Kop\,\vx_i\;}
  \label{eq:gram-entry}
\end{equation}
The entry $G_{ij}$ is an \emph{operator-weighted bilinear form} of family
members $i$ and $j$, scaled by the weight. In the framework this verb is
\code{gram\_reduce}, with the signature and body
\begin{lstlisting}[style=l4style]
gram_reduce :: LinOp -> [Tensor] -> (Int -> Int -> Scalar) -> Matrix
gram_reduce k xs w =
  symmetric_from_upper                          -- mirror lower triangle from upper
    [ [ w i j * entry k xs i j | j <- [i .. m-1] ]   -- map over upper-triangle pairs
      | i <- [0 .. m-1] ]
  where
    m = length xs
    entry k xs i j
      | i == j    = matrix_weighted_norm (xs!!i) k   -- diagonal: x_i^T K x_i
      | otherwise = bilinear_form (xs!!j) k (xs!!i)  -- off-diag:  x_j^T K x_i
\end{lstlisting}
Two structural facts are baked into this body. First, each grid entry is
\emph{independent}: the upper-triangle comprehension is a \texttt{map} over
family pairs, with no state threaded between entries---the embarrassingly
parallel pattern of \cref{sec:red-setting}. Second, the matrix is \emph{symmetric
by construction}: the body computes only the upper triangle and mirrors the lower
from it (\code{symmetric\_from\_upper}). The symmetry is not checked after the
fact; it is guaranteed because $\Kop$ is symmetric, so $\vx_j^{\!\top}\Kop\vx_i
=\vx_i^{\!\top}\Kop\vx_j$, and the verb exploits it to halve the work.

\begin{notationbox}
The diagonal and off-diagonal entries call two slightly different lower-layer
verbs. On the diagonal, $G_{ii}=w(i,i)\,\vx_i^{\!\top}\Kop\vx_i$ is a
\emph{self}-pairing---an operator-weighted squared norm, $\code{matrix\_weighted\_norm}$
(its radicand). Off the diagonal, $G_{ij}=w(i,j)\,\vx_j^{\!\top}\Kop\vx_i$ is a
\emph{cross}-pairing of two different members, the bilinear form
$\code{bilinear\_form}$. They are the same expression $\vx_j^{\!\top}\Kop\vx_i$
specialized to $i=j$ and $i\neq j$; the verb names them separately only because
the self-case has a cheaper, symmetric implementation.
\end{notationbox}

\subsection{Two products, one weight apart}

Here is the payoff. \emph{Capacitance and inductance are the same reduction.}
They call the identical verb \code{gram\_reduce} on the identical operator-weighted
bilinear form; they differ in exactly one place---the scalar weight $w(i,j)$---and
in what the family members physically are.

\paragraph{Capacitance.} The electrostatic driver holds each conductor $i$ at
unit voltage and solves for the potential field $\vx_i=V_i$
(\cref{ch:electrostatics}). The operator $\Kop$ is the
$\varepsilon$-weighted energy operator (the discrete form of
$\tfrac12\varepsilon\abs{\vE}^2$, \cref{ch:maxwell}, eq.~\eqref{eq:energy-density}).
The weight is \emph{unity}:
\begin{equation}
  w(i,j) = 1,
  \qquad
  C_{ij} = V_j^{\!\top}\,\Kop\,V_i .
  \label{eq:cap-gram}
\end{equation}
The diagonal entry $C_{ii}=V_i^{\!\top}\Kop V_i$ is exactly twice the electric
field energy stored when conductor $i$ alone is at unit voltage---the energy form
$C_{ii}=2U_e/V^2$ with $V=1$ that \cref{ch:maxwell} promised
(eq.~\eqref{eq:energy-density}: $C=2U/V^2$). The off-diagonal $C_{ij}$ is the
\emph{mutual} capacitance, the charge induced on conductor $i$ by a unit voltage
on conductor $j$.

\paragraph{Inductance.} The magnetostatic driver drives each source $i$ with
unit current and solves for the vector potential field $\vx_i=\vA_i$
(\cref{ch:magnetostatics}). The operator $\Kop$ is now the $\nu$-weighted
curl--curl energy operator (the magnetic energy $\tfrac{1}{2\mu}\abs{\vB}^2$).
The verb, the bilinear form, the symmetry, the map---all identical. The only
change is the weight, which now \emph{normalizes by the currents}:
\begin{equation}
  w(i,j) = \frac{1}{I_i\,I_j},
  \qquad
  L_{ij} = \frac{1}{I_i I_j}\,\vA_j^{\!\top}\,\Kop\,\vA_i .
  \label{eq:ind-gram}
\end{equation}
The diagonal $L_{ii}$ is the self-inductance ($L=2U_m/I^2$, \cref{ch:maxwell});
the off-diagonal $L_{ij}$ is the mutual inductance between loops $i$ and $j$.

Why does inductance need the current normalization that capacitance does not? An
honest one-line answer: in electrostatics the excitation is a \emph{voltage} and
we choose it to be unity, so dividing by it does nothing; in magnetostatics the
natural excitation is a \emph{current}, and the energy bilinear form scales with
the square of the driving current, so to recover the geometry-only quantity
$L=2U_m/I^2$ we must divide out $I_iI_j$. The weight is exactly the
``per-unit-excitation'' rescaling, trivial in one case and nontrivial in the
other. \Cref{fig:gram-energy} draws the energy picture that underlies the
diagonal.

\begin{figure}[t]
\centering
\begin{tikzpicture}[font=\footnotesize, node distance=7mm]
  % the field-energy picture for the Gram diagonal
  \node[data, align=center, text width=22mm] (vi) {field $\vx_i$\\\footnotesize
    (unit excitation $i$)};
  \node[opbox, minimum width=20mm, right=10mm of vi] (Kx) {$\Kop\,\vx_i$};
  \node[opbox, fill=simBgRust, draw=simRust, minimum width=30mm, right=10mm of Kx]
    (form) {$\vx_i^{\!\top}\Kop\,\vx_i = 2U$};
  \node[product, align=center, text width=26mm, right=10mm of form] (out)
    {$G_{ii}=w(i,i)\cdot 2U$\\[1pt]\footnotesize $=C_{ii}$ or $L_{ii}$};
  \draw[flow] (vi) -- node[above,font=\scriptsize]{apply $\Kop$} (Kx);
  \draw[flow] (Kx) -- node[above,font=\scriptsize]{$\inner{\cdot}{\vx_i}$} (form);
  \draw[flow] (form) -- node[above,font=\scriptsize]{$\times\,w$} (out);
  \node[font=\scriptsize\itshape, simGray, below=2mm of form, align=center]
    {$U=\tfrac12\!\int\!\varepsilon\abs{\vE}^2$ (elec.)\\
     or $\tfrac{1}{2\mu}\!\int\!\abs{\vB}^2$ (mag.)};
  \node[font=\scriptsize\itshape, simRust, below=10mm of out, align=center]
    {$w=1$: capacitance\\$w=1/I_i^2$: inductance};
\end{tikzpicture}
\caption[The Gram diagonal is field energy]{The energy picture behind the Gram
diagonal. Applying the energy operator $\Kop$ to a solved field $\vx_i$ and
pairing the result back with $\vx_i$ yields $\vx_i^{\!\top}\Kop\vx_i$, which is
\emph{twice the stored field energy} $2U$ (\cref{ch:maxwell},
eq.~\eqref{eq:energy-density}). Scaling by the weight $w$ gives the physical
diagonal: with $w=1$ this is the self-capacitance $C_{ii}=2U_e/V^2$ at unit
voltage; with $w=1/I_i^2$ it is the self-inductance $L_{ii}=2U_m/I^2$ at unit
current. The off-diagonal repeats the picture with two \emph{different} fields,
$\vx_j^{\!\top}\Kop\vx_i$, giving mutual capacitance or inductance.}
\label{fig:gram-energy}
\end{figure}

\begin{keyidea}
The \textbf{symmetric Gram reduction} $G_{ij}=w(i,j)\,\vx_j^{\!\top}\Kop\vx_i$ is
one verb (\code{gram\_reduce}) that produces \emph{two} physical products.
Capacitance uses unit weight $w=1$ on the per-terminal potential family;
inductance uses the current-normalized weight $w=1/(I_iI_j)$ on the per-source
vector-potential family. Same operator-weighted bilinear form, same structural
symmetry ($G_{ij}=G_{ji}$, computed on the upper triangle and mirrored), same
embarrassingly parallel map---differing only in a scalar weight and in what the
family physically is. The diagonal is twice the stored field energy; the
off-diagonal is the mutual coupling.
\end{keyidea}

\subsection{Worked example: a two-terminal capacitance}

\begin{workedexample}{A $2\times 2$ capacitance from a symmetric Gram}{cap2x2}
A two-conductor electrostatic problem has been solved. The driver held conductor
$1$ at unit voltage (conductor $2$ grounded) and solved for the potential field
$V_1$; then it held conductor $2$ at unit voltage and solved for $V_2$. The
discrete fields and the assembled energy operator $\Kop$ are given; we are to
build the $2\times 2$ capacitance matrix.

Suppose evaluating the operator-weighted bilinear form $\vx_j^{\!\top}\Kop\vx_i$
on the three needed pairs gives the raw numbers (in femtofarads, with unit
voltage so $w\equiv 1$):
\[
  V_1^{\!\top}\Kop V_1 = 4.20,\qquad
  V_2^{\!\top}\Kop V_2 = 3.60,\qquad
  V_2^{\!\top}\Kop V_1 = -1.10 .
\]
We compute the upper triangle and \emph{mirror} the off-diagonal---we never
recompute $V_1^{\!\top}\Kop V_2$, because the symmetry of $\Kop$ guarantees it
equals $V_2^{\!\top}\Kop V_1$:
\[
  C_{11}=4.20,\qquad
  C_{22}=3.60,\qquad
  C_{12}=C_{21}=-1.10 .
\]
The capacitance matrix is therefore
\[
  \mat{C} \;=\;
  \begin{pmatrix} 4.20 & -1.10 \\ -1.10 & 3.60 \end{pmatrix}\ \si{\femto\farad}.
\]

\emph{Interpretation.} The diagonal entries are the \emph{self}-capacitances:
held alone at unit voltage, conductor $1$ stores $\tfrac12 C_{11}V^2$ of field
energy, conductor $2$ slightly less. The off-diagonal is the \emph{mutual}
capacitance, and its \emph{sign is negative}---a universal feature of the
Maxwell capacitance matrix. Raising conductor $2$'s voltage draws charge of the
\emph{opposite} sign onto conductor $1$ through their mutual field, so the
induced charge $Q_1=C_{11}V_1+C_{12}V_2$ falls as $V_2$ rises; the off-diagonal
coupling is an inducement of opposite charge, hence $C_{12}<0$. The matrix is
symmetric ($C_{12}=C_{21}$) because the underlying energy operator is symmetric:
the charge that conductor $2$ induces on conductor $1$ per volt equals the charge
conductor $1$ induces on conductor $2$ per volt---a reciprocity that the Gram
construction makes automatic, not coincidental.

\emph{The inductance contrast.} Had this been a magnetostatic problem with the
same raw bilinear values but driving currents $I_1=2$, $I_2=5$ (in amperes), the
weight $w(i,j)=1/(I_iI_j)$ would have rescaled each entry:
$L_{11}=4.20/4=1.05$, $L_{22}=3.60/25=0.144$, and
$L_{12}=L_{21}=-1.10/(2\cdot 5)=-0.110$ (in the corresponding units). Same verb,
same symmetry, same fields-in; one weight closure swapped, and the energy form
becomes an inductance instead of a capacitance.
\end{workedexample}

\section{Shape 2: the port-projection reduction}
\label{sec:proj}

The second shape is also a matrix indexed by two families---and it is tempting to
file it under the first. Resist. The S-parameter matrix is a \emph{projection},
not a Gram, and the difference is not cosmetic.

\subsection{The verb}

The driven driver solves the time-harmonic system at each frequency $\omega$,
once per driven port, producing a family of complex scattered fields
$\vE_j(\omega)$---one column per drive port $j$ (\cref{ch:driven}). The
S-parameter reduction projects each of those fields onto every \emph{receiver}
port's mode $\vs_i$, subtracts a self-reflection on the diagonal, and scales by
an impedance/de-embedding factor:
\begin{equation}
  \boxed{\;S_{ij} \;=\; \sigma_{ij}\,\bigl(\inner{\vs_i}{\vE_j(\omega)} - \delta_{ij}\bigr)\;}
  \label{eq:sparam-entry}
\end{equation}
where $\inner{\vs_i}{\vE_j}$ is the linear projection of the drive-$j$ field onto
receiver mode $i$, $\delta_{ij}$ is the Kronecker delta (the self-term, $1$ when
$i=j$ and $0$ otherwise), and $\sigma_{ij}$ is the port-kind scale (an impedance
ratio for lumped ports, a de-embedding phase for wave ports). The framework verb
is \code{sparameter\_reduce}:
\begin{lstlisting}[style=l4style]
sparameter_reduce :: [PortMode] -> [(Int, Tensor)] -> Matrix
sparameter_reduce ports family =
  matrix_from_columns
    [ [ scale ports i j * (project (ports!!i) e - selfterm i j)  -- entry S_ij
        | i <- [0 .. p-1] ]                                      -- over receivers i
      | (j, e) <- family ]                                       -- one column per drive port j
  where
    p             = length ports
    project s e   = port_dot s e            -- linear functional s_i . E
    selfterm i j  = if i == j then 1 else 0 -- the -1 self-reflection on the diagonal
    scale i j     = port_scale (ports!!i) (ports!!j)  -- impedance / de-embed
\end{lstlisting}
The reduction is applied once per swept frequency; the $\omega$ index rides as an
outer family axis that the driven composition root owns, so the verb itself is
the per-port reduction at a single $\omega$.

\subsection{Why this is rank-2 like a Gram but deliberately not a Gram}

Both \code{gram\_reduce} and \code{sparameter\_reduce} produce square matrices
indexed by a port-like family. That shared \emph{output rank} is the whole of
their resemblance, and reading more into it is the single most over-unified
mistake the design could make.

\begin{pitfall}
\label{pit:gram-not-proj}
\textbf{Two reductions producing matrices are not the same reduction.} It is
sorely tempting to unify \code{gram\_reduce} and \code{sparameter\_reduce} under
one ``matrix reduction'' verb, since both yield a $p\times p$ array indexed by
ports. Three structural facts forbid it.
\begin{itemize}[nosep]
  \item \emph{The Gram is a self-pairing; the projection is not.} The Gram pairs
  a family \emph{with itself} through an operator, $\vx_j^{\!\top}\Kop\vx_i$. The
  S-matrix projects one family (the solved drive fields $\vE_j$) onto a
  \emph{different} basis (the port modes $\vs_i$)---two families in different
  roles, not one family paired with itself.
  \item \emph{The Gram is symmetric; the projection is not.} The Gram's symmetry
  $G_{ij}=G_{ji}$ is structural (\cref{sec:gram}) and the verb \emph{computes
  only the upper triangle}. The S-matrix has no such guarantee---every entry is
  computed independently---and its near-symmetry $S_{ij}\approx S_{ji}$, where it
  holds at all, is reciprocity \emph{physics}, not a construction the reduction
  imposes. There is no \code{symmetric\_from\_upper} in the projection body.
  \item \emph{The Gram is real and energy-valued; the projection is complex and
  has an inhomogeneous diagonal.} S-parameters are complex (magnitude and phase),
  and the diagonal carries the $-\delta_{ij}$ self-reflection term, which has no
  analogue in the Gram. Forcing one verb would either smuggle a false symmetry
  into the S-matrix or burden the Gram with a self-term and complex arithmetic it
  does not need.
\end{itemize}
The right move---which the framework makes---is two verbs that happen to share an
output rank. \emph{Sharing a shape is not sharing a meaning.}
\end{pitfall}

The port modes $\vs_i$ themselves are not produced by the driven driver; they
come from the \emph{boundary-mode} analysis (\cref{ch:waveports}), which solves a
small two-dimensional eigenproblem on each port's cross-section to find its
propagation modes. The driven reduction consumes those modes as the basis it
projects onto. \Cref{fig:proj} draws the projection.

\begin{figure}[t]
\centering
\begin{tikzpicture}[font=\footnotesize]
  % the solved field as an arrow in "field space"
  \draw[->, simGray] (0,0) -- (0,3.2) node[above,font=\scriptsize]{};
  \draw[->, simGray] (0,0) -- (4.6,0) node[right,font=\scriptsize]{};
  % port mode basis vectors
  \draw[-{Stealth[length=2.4mm]}, very thick, simGreen] (0,0) -- (3.4,0.0)
    node[right,font=\scriptsize,simGreen]{$\vs_1$ (port 1 mode)};
  \draw[-{Stealth[length=2.4mm]}, very thick, simGreen] (0,0) -- (0.0,2.6)
    node[above,font=\scriptsize,simGreen]{$\vs_2$ (port 2 mode)};
  % the solved driven field E_j
  \draw[-{Stealth[length=2.6mm]}, very thick, simRust] (0,0) -- (2.7,1.7)
    node[above right,font=\scriptsize,simRust]{$\vE_j(\omega)$ (drive $j$)};
  % projections (dashed drops)
  \draw[densely dashed, simBlue] (2.7,1.7) -- (2.7,0)
    node[midway,right,font=\scriptsize,simBlue]{};
  \draw[densely dashed, simBlue] (2.7,1.7) -- (0,1.7);
  \node[font=\scriptsize, simBlue] at (2.7,-0.28) {$\inner{\vs_1}{\vE_j}$};
  \node[font=\scriptsize, simBlue] at (-0.6,1.7) {$\inner{\vs_2}{\vE_j}$};
  % annotation
  \node[font=\scriptsize\itshape, simGray, align=left] at (3.2,2.6)
    {each projection\\$=$ one $S_{\bullet j}$\\(before $-\delta$, scale)};
\end{tikzpicture}
\caption[Port projection of a driven field]{The S-parameter reduction as a
\emph{linear projection}. The driven driver's solved field $\vE_j(\omega)$ (rust)
for drive port $j$ is projected onto the basis of \emph{port modes} $\vs_i$
(green, supplied by the boundary-mode analysis, \cref{ch:waveports}). Each
projection $\inner{\vs_i}{\vE_j}$ (blue drop) becomes one column entry
$S_{ij}$---after subtracting the diagonal self-reflection $\delta_{ij}$ and
applying the impedance/de-embed scale $\sigma_{ij}$. Unlike the Gram of
\cref{fig:gram-energy}, the two index families are \emph{different} (drive fields
vs.\ port modes), the operation is a projection rather than an energy
self-pairing, and the result is complex and not symmetric.}
\label{fig:proj}
\end{figure}

\section{Shape 3: the rank-1 scalar and mode tables}
\label{sec:tables}

The third shape is the simplest in form and the most populous in membership: a
\emph{table}---a one-dimensional list with one row per item, each row a few
scalars (or, in one case, a few small fields). There is no matrix, no pairing of
a family with itself, just a \texttt{map} over a collection producing one tuple
per element. Three distinct products live here, and we take them in turn.

\subsection{Energy and participation, per domain}
\label{sec:energy-table}

The most driver-agnostic reduction in the book asks, of \emph{any} field-bearing
analysis: where does the field energy sit? Given a solved field and a configured
map of domain-restricted energy operators $\{i\mapsto \mat{M}_i\}$, the
\emph{domain-energy} reduction computes, per domain, the energy in that domain
and the fraction of the total it represents:
\begin{equation}
  \mathrm{energy}_i \;=\; \tfrac12\,\inner{\text{field}}{\mat{M}_i\,\text{field}},
  \qquad
  p_i \;=\; \frac{\mathrm{energy}_i}{e_{\text{total}}},
  \label{eq:energy-part}
\end{equation}
where $\mat{M}_i$ is the energy operator restricted to domain $i$ (an SPD form,
the domain-local piece of the global energy operator of \cref{ch:maxwell},
eq.~\eqref{eq:energy-density}) and $e_{\text{total}}$ is the whole-domain total.
The participation factor $p_i$ is the headline number---it says what fraction of
the electric (or magnetic) field energy lives in each material region, the
quantity a designer inspects to see whether energy is concentrating where it
should. The verb is \code{domain\_energy\_reduce}; it is a \texttt{map} over the
domain map with no inter-domain state.

This reduction is \emph{driver-agnostic}: it cares only that a field exists and
that domain operators are configured, so it runs identically as a post-process of
the electrostatic, eigenmode, driven, or transient analyses. It runs once per
field kind---electric ($\tfrac12\inner{\vE}{\mat{M}_i\vE}$) and magnetic
($\tfrac12\inner{\vB}{\mat{M}_i\vB}$)---producing one table per kind. The field
may be complex (in the driven and eigenmode cases), but the energy form is a real
nonnegative number, so the table is real. When the configured domains
\emph{partition} the field's support, the participations sum to one,
$\sum_i p_i = 1$; this is the natural sanity check on the table.
\Cref{fig:energy-bars} shows a typical participation table as a bar chart.

\begin{figure}[t]
\centering
\begin{tikzpicture}
\begin{axis}[
    width=0.82\linewidth, height=46mm,
    ybar, bar width=10mm,
    ymin=0, ymax=0.7,
    ylabel={participation $p_i$}, ylabel style={font=\scriptsize},
    xlabel={material domain}, xlabel style={font=\scriptsize},
    symbolic x coords={substrate, dielectric, air, metal},
    xtick=data, xticklabel style={font=\scriptsize},
    ytick={0,0.2,0.4,0.6}, yticklabel style={font=\scriptsize},
    nodes near coords, nodes near coords style={font=\scriptsize},
    every node near coord/.append style={/pgf/number format/fixed,
      /pgf/number format/precision=2},
    enlarge x limits=0.18,
    axis lines=left,
  ]
  \addplot[fill=simBgGold, draw=simGold] coordinates {
    (substrate,0.52) (dielectric,0.31) (air,0.13) (metal,0.04)};
\end{axis}
\end{tikzpicture}
\caption[Energy participation per domain]{An electric-energy participation table,
drawn as a bar chart. Each bar is one row of the rank-1 table: the fraction
$p_i=\mathrm{energy}_i/e_{\text{total}}$ of the total field energy stored in
material domain $i$ (\cref{eq:energy-part}). Here most of the energy sits in the
substrate and dielectric, very little in the surrounding air, and almost none in
the metal---a typical signature for a planar device. The four participations sum
to $1.00$ because the configured domains partition the field's support; that sum
is the table's natural sanity check. The reduction is \emph{driver-agnostic}: any
analysis that produces a field can be post-processed this way.}
\label{fig:energy-bars}
\end{figure}

\subsection{Eigenfrequency and quality factor, per mode}
\label{sec:fq-table}

The eigenmode driver hands back a family of converged eigenpairs
$(\lambda_m,\vx_m)$ (\cref{ch:eigenvalues}). The \emph{eigenfrequency--$Q$}
reduction turns each into a row $(f_m,Q_m)$ of a per-mode table. Two pieces of
work happen per row, and both \emph{undo} a transform the solve stage applied.

First, the \emph{eigenvalue un-transform}. \Cref{ch:eigenvalues} bent the
spectrum with a shift to make the wanted low eigenvalues easy to find; the
eigenvalue the library returns is therefore in transformed coordinates, and we
must map it back to a physical angular frequency $\omega_m$. For the lossless
\emph{linear} problem $\lambda=\omega^2$, so $\omega=\sqrt\lambda$; for the lossy
\emph{quadratic} problem $\lambda=i\omega$, so $\omega=\lambda/i=-i\lambda$. The
\emph{eigenfrequency} is the real part of the recovered $\omega$:
\begin{equation}
  f_m \;=\; \Re\,\omega_m .
  \label{eq:eigfreq}
\end{equation}

Second, the \emph{quality factor}. $Q$ measures how slowly a mode loses
energy---the ratio of energy stored to energy dissipated per cycle, equivalently
the ratio of the oscillation rate to the decay rate. In the lossy problem the
imaginary part of $\omega$ \emph{is} the decay rate (the loss rate $\kappa_m$),
so
\begin{equation}
  Q_m \;=\; \frac{\omega_m}{\kappa_m},
  \qquad
  \kappa_m = \text{(loss rate from }\Im\,\omega_m\text{)} .
  \label{eq:qfactor}
\end{equation}
A lossless mode has $\kappa_m=0$ and rings forever; its $Q$ is infinite. That is
not a numerical accident to be papered over but a physical fact, and the
reduction guards it explicitly: the verb \code{eigenfreq\_qfactor\_reduce}
returns $Q=\infty$ when $\kappa=0$ rather than dividing by zero.
\begin{lstlisting}[style=l4style]
eigenfreq_qfactor_reduce ptype kappa eigs =
  [ let omega = untransform ptype lambda        -- omega = sqrt(mu) | lambda/i
        f     = re omega                         -- eigenfrequency  f_m = Re omega
        k     = kappa mode                       -- loss rate kappa_m
        q     = if k == 0 then infinity else f / abs k  -- Q = omega/kappa, guarded
    in  (f, q)
  | (lambda, mode) <- eigs ]                      -- map over modes, no inter-mode state
\end{lstlisting}
The two halves are the two scalar projections of each mode; the mode-\emph{field}
recovery ($\vx_m$ itself, and $\vB=-\Curl\vx_m/(i\omega)$) is a separate field
readout, not part of this scalar table. \Cref{fig:omega-plane} shows where $f$
and $Q$ live in the complex $\omega$-plane.

\begin{figure}[t]
\centering
\begin{tikzpicture}[font=\footnotesize, scale=1.0]
  % complex omega-plane
  \draw[->, simGray] (-0.4,0) -- (5.2,0) node[right,font=\scriptsize]{$\Re\,\omega$ (frequency)};
  \draw[->, simGray] (0,-0.4) -- (0,2.7) node[above,font=\scriptsize]{$\Im\,\omega$ (loss rate $\kappa$)};
  % the eigenvalue omega_m
  \coordinate (w) at (3.8,1.1);
  \fill[simRust] (w) circle (2.2pt);
  \node[font=\scriptsize, simRust, anchor=south west] at (w) {$\omega_m$};
  % vector from origin
  \draw[-{Stealth[length=2.2mm]}, thick, simBlue] (0,0) -- (w);
  % drop to real axis: f
  \draw[densely dashed, simGreen, thick] (w) -- (3.8,0);
  \node[font=\scriptsize, simGreen] at (1.9,-0.32) {$f_m=\Re\,\omega_m$};
  % horizontal to imag axis: kappa
  \draw[densely dashed, simRust, thick] (w) -- (0,1.1);
  \node[font=\scriptsize, simRust, anchor=east] at (0,1.1) {$\kappa_m$};
  % distance-to-real-axis annotation
  \node[font=\scriptsize\itshape, simGray, align=left] at (4.2,2.1)
    {$Q_m=\dfrac{\Re\,\omega_m}{\Im\,\omega_m}$:\\high $Q$ $=$\\close to real axis};
  % an arc indicating small angle = high Q
  \draw[simGray, densely dotted] (1.2,0) arc (0:16:1.2);
\end{tikzpicture}
\caption[Eigenfrequency and $Q$ in the complex plane]{Reading $f$ and $Q$ from a
complex eigenvalue. The recovered angular frequency $\omega_m$ is a point in the
complex plane: its \emph{real} part is the oscillation frequency
$f_m=\Re\,\omega_m$ (\cref{eq:eigfreq}), its \emph{imaginary} part is the loss
rate $\kappa_m$ (the decay). The quality factor $Q_m=\Re\,\omega_m/\Im\,\omega_m$
is large precisely when $\omega_m$ sits \emph{close to the real axis}---a small
imaginary part means slow decay, a sharp long-ringing resonance. A purely real
$\omega_m$ (on the axis) is lossless: $\kappa_m=0$, $Q_m=\infty$, the guarded
case in the verb.}
\label{fig:omega-plane}
\end{figure}

\begin{workedexample}{Eigenfrequency and $Q$ from a complex eigenvalue}{fq}
A lossy eigenmode solve returns a converged eigenpair whose recovered angular
frequency is the complex number
\[
  \omega_m \;=\; \omega_r + i\,\omega_i
            \;=\; 2\pi\,(\SI{5.000}{\giga\hertz}) \;+\; i\,(2\pi\cdot\SI{2.5}{\mega\hertz}).
\]
That is, the real part corresponds to a \SI{5}{\giga\hertz} oscillation and the
imaginary part to a \SI{2.5}{\mega\hertz} decay rate. We want the eigenfrequency
$f_m$ and the quality factor $Q_m$.

\emph{Eigenfrequency.} The frequency is the real part, divided by $2\pi$ to go
from angular to ordinary frequency:
\[
  f_m \;=\; \frac{\Re\,\omega_m}{2\pi} \;=\; \SI{5.000}{\giga\hertz}.
\]

\emph{Quality factor.} With the energy/loss convention $Q=\Re\,\omega/(2\,\Im\,\omega)$
(equivalently $\omega_r/(2\omega_i)$, the form in which $Q$ counts radians of
oscillation per e-folding of amplitude),
\[
  Q_m \;=\; \frac{\omega_r}{2\,\omega_i}
        \;=\; \frac{2\pi\cdot 5.000\times 10^{9}}{2\cdot 2\pi\cdot 2.5\times 10^{6}}
        \;=\; \frac{5.000\times 10^{9}}{5.0\times 10^{6}}
        \;=\; 1000 .
\]
A $Q$ of $1000$ is a respectably sharp resonance: the mode rings for about $1000$
radians---some $160$ full cycles---before its amplitude decays by a factor of
$e$. The factor of $2$ in the denominator is the standard convention relating the
amplitude decay rate to the \emph{energy} decay rate (energy goes as amplitude
squared, so it decays twice as fast); different texts place the $2$ differently,
which is why the verb takes the loss-rate convention $\kappa$ as a parameter
rather than hard-coding it.

\emph{The lossless limit.} Now let the loss vanish, $\omega_i\to 0$. The mode
sits exactly on the real axis of \cref{fig:omega-plane}, the denominator goes to
zero, and $Q_m\to\infty$. A lossless cavity rings forever. The reduction does not
divide by zero and emit a \texttt{NaN}; it tests $\kappa=0$ and returns
$Q=\infty$, the honest physical answer. This guard is the one piece of
control flow in an otherwise-arithmetic reduction, and it exists because the
infinite-$Q$ case is real, not exceptional---an ideal lossless resonator is the
limit every low-loss design approaches.
\end{workedexample}

\subsection{Waveguide propagation modes, per mode}
\label{sec:wg-table}

The boundary-mode driver (\cref{ch:waveports}) solves a small two-dimensional
eigenproblem on a port's cross-section and hands back converged eigenpairs. The
\emph{waveguide-mode} reduction turns each into a table row---but this is the one
rank-1 table whose rows carry \emph{fields}, not just scalars. Per mode it
produces:
\begin{itemize}[nosep]
  \item the \emph{propagation constant} $k_n$, recovered from the eigenvalue by
  the same shift-invert un-transform the eigenmode reduction uses
  (\cref{ch:eigenvalues})---it says how fast the mode travels down the guide, and
  whether it travels at all;
  \item the \emph{effective index} $n_{\text{eff}}=k_n/\omega$, the propagation
  constant normalized by the operating frequency (a dimensionless ratio
  comparing the mode's phase velocity to the speed of light);
  \item the \emph{transverse fields} $(\vE_t,\vE_n)$, recovered from the
  eigenvector and \emph{power-normalized} so the mode carries unit Poynting power
  ($\abs{P}=1$); and, for propagating modes only, the longitudinal magnetic field
  \begin{equation}
    \vB_z \;=\; \frac{\Curl\vE_t}{i\,\omega},
    \label{eq:Bz}
  \end{equation}
  the curl of the transverse electric field divided by $i\omega$---a direct
  discrete reading of Faraday's law (\cref{ch:maxwell}, eq.~\eqref{eq:faraday}).
\end{itemize}
The verb is \code{waveguide\_mode\_reduce}; like its scalar-only siblings it is a
\texttt{map} over the converged modes with no inter-mode state, and the operating
frequency $\omega$ rides as a fixed parameter (the $n_{\text{eff}}$ divisor and
the $\vB_z$ scale), not a per-mode datum. The conditional $\vB_z$---present for
propagating modes, absent for evanescent ones---is the one place a table row's
shape varies by content; the framework records it as an optional field.
\Cref{fig:wg-table} sketches such a table, fields and all.

\begin{figure}[t]
\centering
\begin{tikzpicture}[font=\footnotesize]
  % header
  \node[font=\scriptsize\bfseries] at (0,0.55)   {mode};
  \node[font=\scriptsize\bfseries] at (1.4,0.55)  {$k_n$};
  \node[font=\scriptsize\bfseries] at (2.8,0.55)  {$n_{\text{eff}}$};
  \node[font=\scriptsize\bfseries] at (4.9,0.55)  {transverse $\vE_t$};
  \node[font=\scriptsize\bfseries] at (7.0,0.55)  {$\vB_z$};
  \draw[simGray, semithick] (-0.6,0.3) -- (7.7,0.3);
  % --- row 1 (propagating) ---
  \node[font=\scriptsize] at (0,0)   {$1$};
  \node[font=\scriptsize] at (1.4,0) {$31.4$};
  \node[font=\scriptsize] at (2.8,0) {$1.50$};
  \node[font=\scriptsize] at (7.0,0) {\checkmark};
  \draw[simGray, thin] (4.3,-0.25) rectangle (5.5,0.25);
  \draw[-{Stealth[length=1.3mm]}, simRust] (4.55,-0.13) -- (4.55,0.13);
  \draw[-{Stealth[length=1.3mm]}, simRust] (4.85,-0.13) -- (4.85,0.13);
  \draw[-{Stealth[length=1.3mm]}, simRust] (5.15,-0.13) -- (5.15,0.13);
  % --- row 2 (propagating) ---
  \node[font=\scriptsize] at (0,-1.05)   {$2$};
  \node[font=\scriptsize] at (1.4,-1.05) {$22.1$};
  \node[font=\scriptsize] at (2.8,-1.05) {$1.06$};
  \node[font=\scriptsize] at (7.0,-1.05) {\checkmark};
  \draw[simGray, thin] (4.3,-1.30) rectangle (5.5,-0.80);
  \draw[-{Stealth[length=1.3mm]}, simBlue] (4.55,-1.18) -- (4.70,-0.92);
  \draw[-{Stealth[length=1.3mm]}, simBlue] (4.90,-1.18) -- (5.05,-0.92);
  \draw[-{Stealth[length=1.3mm]}, simBlue] (5.20,-1.18) -- (5.30,-0.92);
  % --- row 3 (evanescent) ---
  \node[font=\scriptsize] at (0,-2.10)   {$3$};
  \node[font=\scriptsize] at (1.4,-2.10) {$i\,8.0$};
  \node[font=\scriptsize] at (2.8,-2.10) {---};
  \node[font=\scriptsize] at (7.0,-2.10) {---};
  \node[font=\scriptsize\itshape, simGray] at (4.9,-2.10) {(evanescent: no field)};
  % footnote
  \node[font=\scriptsize\itshape, simGreen, anchor=west] at (-0.6,-2.85)
    {modes 1--2 propagate ($k_n$ real); mode 3 is evanescent ($k_n$ imaginary)};
\end{tikzpicture}
\caption[A waveguide-mode table carrying fields]{The waveguide-mode reduction
produces a rank-1 table whose rows carry \emph{fields}. Each converged mode
yields a propagation constant $k_n$, an effective index
$n_{\text{eff}}=k_n/\omega$, the power-normalized transverse field pattern
$\vE_t$ (the little arrow cartoons), and---for propagating modes only---the
longitudinal field $\vB_z=\Curl\vE_t/(i\omega)$. Modes $1$ and $2$ propagate
($k_n$ real, $n_{\text{eff}}$ defined, $\vB_z$ present, marked \checkmark); mode
$3$ is \emph{evanescent} ($k_n$ imaginary), carries no propagating power, and has
no $\vB_z$. The conditional $\vB_z$ is the one place a table row's shape depends
on its content---recorded as an optional field. This is the only rank-1 table
whose entries are not purely scalar.}
\label{fig:wg-table}
\end{figure}

\section{The taxonomy, and what it buys}
\label{sec:taxonomy}

We have now met all three shapes in full, and the organizing payoff of the
chapter can be stated plainly. \emph{Every} user-facing physical output Palace
produces is one of three reduction shapes (\cref{fig:red-taxonomy}):

\begin{itemize}[nosep]
  \item the \textbf{rank-2 symmetric Gram} $G_{ij}=w(i,j)\,\vx_j^{\!\top}\Kop\vx_i$
  ---capacitance and inductance, one verb differing only in the weight $w$;
  \item the \textbf{rank-2 port-projection}
  $S_{ij}=\sigma_{ij}(\inner{\vs_i}{\vE_j}-\delta_{ij})$ ---S-parameters, a
  linear projection, complex, deliberately not a Gram (\cref{pit:gram-not-proj});
  \item the \textbf{rank-1 scalar-or-mode table}---energy and participation
  (per domain), eigenfrequency and $Q$ (per mode), and waveguide modes (per mode,
  carrying fields), each a pure \texttt{map} with no pairing.
\end{itemize}

\begin{table}[t]
\centering
\caption[The reduction shapes and their products]{The three reduction shapes and
the products each yields. The \emph{shape} column is the structural taxonomy; the
\emph{verb} column is the framework combinator; the \emph{per-item} column is the
computation mapped over the solution family. Reading down the \emph{shape} column,
two reductions produce matrices (Gram, projection) and four produce
tables---and the matrix pair is two genuinely different reductions, not one.}
\label{tab:red-shapes}
\small
\setlength{\tabcolsep}{4pt}
\renewcommand{\arraystretch}{1.25}
\begin{tabular}{@{}llll@{}}
\toprule
Shape & Verb & Per-item computation & Product \\
\midrule
rank-2 Gram & \code{gram\_reduce} & $w(i,j)\,\vx_j^{\!\top}\Kop\vx_i$ & capacitance \\
rank-2 Gram & \code{gram\_reduce} & $w{=}1/(I_iI_j)$, same form & inductance \\
rank-2 projection & \code{sparameter\_reduce} & $\sigma_{ij}(\inner{\vs_i}{\vE_j}-\delta_{ij})$ & S-parameters \\
rank-1 table & \code{domain\_energy\_reduce} & $(\tfrac12\inner{f}{\mat{M}_if},\,p_i)$ & energy, participation \\
rank-1 table & \code{eigenfreq\_qfactor\_reduce} & $(\Re\,\omega_m,\,\omega_m/\kappa_m)$ & eigenfrequency, $Q$ \\
rank-1 table & \code{waveguide\_mode\_reduce} & $(k_n,n_{\text{eff}},\vE_t,\vB_z)$ & waveguide modes \\
\bottomrule
\end{tabular}
\end{table}

\Cref{tab:red-shapes} collects them. The lesson is the book's recurring refrain,
now at the output stage: a small handful of shapes, specialized by a weight, a
basis, or a per-item formula, produce the entire menu of physical answers. The
taxonomy is not a tidy-minded over-unification---we were careful in
\cref{pit:gram-not-proj} to keep the Gram and the projection distinct---but a
genuine structural economy: there really are only three ways the machinery
collapses a field family into an answer.

\begin{keyidea}
The \textbf{three reduction shapes} produce \emph{every} physical output:
the rank-2 symmetric Gram (capacitance, inductance), the rank-2 port-projection
(S-parameters), and the rank-1 scalar-or-mode table (energy/participation,
eigenfrequency/$Q$, waveguide modes). Two shapes yield matrices and one yields
tables; the two matrix reductions are \emph{different} reductions that merely
share an output rank. Each is an embarrassingly parallel \texttt{map} over the
solution family---the last and simplest stage of every analysis, and the place
where the physics re-enters as a capacitance, a scattering coefficient, or a
resonant frequency.
\end{keyidea}

\section{Closing Part IV: the machinery is assembled}
\label{sec:bridge}

This chapter closes Part~IV, and with it the long climb up the machinery. It is
worth standing back to see how far we have come, and where Part~V takes us.

Part~II built the \emph{field theory}: Maxwell's equations, the weak form, the
finite-element function spaces, and the assembly that turns continuous physics
into discrete operators $\Kop,\Mop,\Cop$. Part~III built the \emph{calculus}: the
abstract vocabulary of tensors, records, folds, combinators, monads, and
operators in which the whole simulator can be written down precisely. Part~IV
built the \emph{numerical machinery}: the Krylov solvers (\cref{ch:krylov,ch:cg,ch:gmres}),
their orthogonalization and rotation kernels (\cref{ch:orthog,ch:givens}), the
preconditioners and multigrid that make them converge (\cref{ch:precond,ch:smoothers,ch:multigrid}),
the eigenvalue machinery and its spectral transforms (\cref{ch:eigenvalues,ch:krylovschur,ch:deflation}),
the time integrators (\cref{ch:time}), the adaptive-refinement outer loop
(\cref{ch:amr}), and now the output reductions that turn solved fields into
physical products.

Every box of the skeleton (\cref{ch:situating}) is now filled with real
machinery. What remains is to \emph{compose} it. That is Part~V---the
\emph{modalities}. Each modality chapter takes one of the six drivers and walks
its full path through the assembled machinery: which function space it assembles
on, which solve corner it occupies, which solver it calls, and---the subject of
this chapter---which reduction shape produces its answer. The electrostatic
driver (\cref{ch:electrostatics}) composes a fixed-operator family solve with a
unit-weight Gram; the magnetostatic (\cref{ch:magnetostatics}) swaps the space
and the Gram weight; the driven (\cref{ch:driven}) swaps in an operator-varying
sweep and the port-projection; the eigenmode (\cref{ch:eigenmodes}) calls the
black-box eigensolver and reduces to an $(f,Q)$ table; the boundary-mode
(\cref{ch:waveports}) does the same on a cross-section and reduces to a waveguide
table. The reductions of this chapter are the last stage of each of those paths.
The parts are all on the bench; Part~V assembles the machines.

\summarysection
The \texttt{reduce} stage turns a driver's solution \emph{family} into the
user-facing physical product. Every reduction is an embarrassingly parallel
\texttt{map} over the family with no inter-element state---never a fold
(\cref{ch:combinators}). There are exactly three reduction \emph{shapes}. The
\textbf{rank-2 symmetric Gram} $G_{ij}=w(i,j)\,\vx_j^{\!\top}\Kop\vx_i$
(\code{gram\_reduce}) produces capacitance ($w=1$, per-terminal potentials) and
inductance ($w=1/(I_iI_j)$, per-source vector potentials) from one verb, its
symmetry structural and its diagonal twice the stored field energy. The
\textbf{rank-2 port-projection}
$S_{ij}=\sigma_{ij}(\inner{\vs_i}{\vE_j}-\delta_{ij})$
(\code{sparameter\_reduce}) produces S-parameters by projecting each driven field
onto the port modes; it is complex, non-symmetric, and deliberately \emph{not} a
Gram---sharing an output rank is not sharing a meaning. The \textbf{rank-1
scalar-or-mode table} produces energy and participation per domain
(\code{domain\_energy\_reduce}, driver-agnostic), eigenfrequency and $Q$ per mode
(\code{eigenfreq\_qfactor\_reduce}, with the $\kappa{=}0\Rightarrow Q{=}\infty$
guard), and waveguide modes per mode (\code{waveguide\_mode\_reduce}, the one
table carrying fields). These three shapes, specialized by a weight, a basis, or
a per-item formula, produce the simulator's entire menu of physical answers---and
they are the last stage of each modality in Part~V.

\furtherreading
The physical meanings of the reduced quantities are standard electromagnetics:
Griffiths~\citep{griffiths2017} for capacitance, inductance, and stored field
energy (the Gram diagonal); Pozar~\citep{pozar2011} for S-parameters, quality
factors, and waveguide propagation modes (the microwave quantities, including the
de-embedding and impedance scalings of the port projection). For the
finite-element side---how the energy operators $\Kop,\mat{M}_i$ are assembled and
why the diagonal of the Gram is exactly twice the discrete field energy---Jin's
finite-element electromagnetics text~\citep{jin2014} is the natural companion,
and it treats the divergence-free projection and the spurious-mode cleanup that
keep the eigenfrequency table honest. The structural observation that all of
these are one of three \texttt{map}-shaped reductions is this book's own, distilled
from reading Palace's postprocessing side by side; \cref{ch:situating} states the
taxonomy and this chapter fills it in.

\exercisessection
\begin{enumerate}[nosep]
  \item \textbf{One verb, two products.} Write down the Gram reduction entry
  $G_{ij}$ and the weight $w(i,j)$ for capacitance and for inductance. In one
  sentence each, explain why capacitance's weight is trivial ($w=1$) but
  inductance's ($w=1/(I_iI_j)$) is not.
  \item \textbf{Mirror, don't recompute.} A three-terminal capacitance problem
  needs the full $3\times 3$ matrix $\mat{C}$. How many distinct bilinear-form
  evaluations $\vx_j^{\!\top}\Kop\vx_i$ does \code{gram\_reduce} actually perform,
  and which entries are filled by \code{symmetric\_from\_upper} instead? What
  property of $\Kop$ licenses the mirroring?
  \item \textbf{Sign of the off-diagonal.} In \cref{wex:cap2x2} the mutual
  capacitance $C_{12}$ came out negative. Explain physically why the off-diagonal
  of a Maxwell capacitance matrix is always negative, using the relation
  $Q_1=C_{11}V_1+C_{12}V_2$.
  \item \textbf{Not a Gram.} Give the three structural reasons (from
  \cref{pit:gram-not-proj}) that the S-parameter matrix is not a Gram matrix,
  even though both are square and port-indexed. For each reason, name the specific
  feature of \code{sparameter\_reduce}'s body that a unified ``matrix reduction''
  verb would have to violate.
  \item \textbf{The $Q$ guard.} Recompute $Q_m$ from \cref{wex:fq} if the loss
  rate were $\omega_i=2\pi\cdot\SI{5}{\mega\hertz}$ instead of \SI{2.5}{\mega\hertz}.
  Then take $\omega_i\to 0$: what does the reduction return, and why is that the
  physically correct answer rather than a numerical failure?
  \item \textbf{Two undone transforms.} The eigenfrequency--$Q$ reduction
  ``un-transforms'' the eigenvalue twice over. Name the two transforms it undoes:
  one applied by the eigenvalue \emph{solve} stage (\cref{ch:eigenvalues}) and one
  implicit in the lossy problem's eigenvalue convention $\lambda=i\omega$. What
  does each contribute to the final $(f,Q)$ row?
  \item \textbf{Driver-agnostic energy.} The domain-energy reduction is the one
  reduction that runs identically as a post-process of \emph{any} field-bearing
  analysis. State the minimal precondition it needs from the driver above it, and
  explain why the participation factors sum to one only when the configured
  domains partition the field's support.
  \item \textbf{Place the products.} For each of the six Palace analyses, name its
  product and the reduction shape (Gram / projection / table) that produces it.
  Which shape is used by the most analyses, and why is that unsurprising given
  that a table is just a \texttt{map} with no pairing?
\end{enumerate}


\part{The Constituent Modalities}\label{part:modalities}
\chapter{Electrostatics and Capacitance}
\label{ch:electrostatics}

\chapterorient{This chapter opens Part~V, and it is the one the rest of the part
is written against. Every modality chapter to come is a variation on a single
theme, and here we play that theme in full: we take the electrostatic
analysis---the simplest of the six---and trace it end to end, from a
configuration of conductors to the capacitance matrix that comes out, with no box
left closed. We will watch the skeleton of \cref{ch:situating} become an actual
computation: assemble the stiffness operator \emph{once}, map a family of
independent solves over it (the fixed-operator corner), and Gram-reduce the solved
fields into the answer. By the end you will have seen a complete analysis run, and
you will be able to read every later modality as \emph{this pipeline, with one
stage swapped}.}

\section{The physical question}
\label{sec:es-question}

An engineer hands us a device---a few metal conductors embedded in an insulating
dielectric---and asks a question as old as the Leyden jar: \emph{how much charge
does it store per volt?} For a single isolated conductor the answer is one number,
its capacitance. For $n$ conductors the answer is a \emph{matrix}, because the
charge on each conductor depends on the voltage of \emph{every} conductor, not
just its own. That $n\times n$ table is the \emph{capacitance matrix}
$\mat{C}$, and computing it is the entire job of the electrostatic analysis.

\begin{physicsbox}
Several conductors sit in a dielectric medium (\cref{fig:es-setup}). Each
conductor is a \emph{perfect equipotential}: a perfect conductor cannot sustain a
tangential electric field, so its whole surface sits at one voltage. We are free
to \emph{set} those voltages; the dielectric between the conductors responds with
an electric field $\vE$, and charge collects on the conductor surfaces. The
capacitance matrix relates the charges $\mathbf q=(q_1,\dots,q_n)$ to the applied
voltages $\mathbf v=(V_1,\dots,V_n)$ by $\mathbf q = \mat{C}\,\mathbf v$. Nothing
moves and nothing oscillates---this is a \emph{static} problem,
$\partial/\partial t=0$ (\cref{ch:maxwell}, \cref{sec:regimes}).
\end{physicsbox}

\begin{figure}[t]
\centering
\begin{tikzpicture}[scale=1.0]
  % dielectric background
  \fill[simBgBlue] (-3.6,-1.7) rectangle (3.6,1.7);
  \node[font=\scriptsize, text=simBlue] at (2.85,1.35) {dielectric $\varepsilon$};
  % conductor 1 (left plate, held at +V)
  \fill[simInk] (-2.7,0.55) rectangle (-0.4,0.95);
  \node[font=\scriptsize, text=white] at (-1.55,0.75) {conductor 1};
  \node[font=\scriptsize] at (-3.15,0.75) {$V_1$};
  % conductor 2 (right plate, held at 0)
  \fill[simInk] (-2.7,-0.95) rectangle (-0.4,-0.55);
  \node[font=\scriptsize, text=white] at (-1.55,-0.75) {conductor 2};
  \node[font=\scriptsize] at (-3.15,-0.75) {$V_2$};
  % field lines E = -grad V between the plates
  \foreach \x in {-2.4,-1.9,-1.4,-0.9}{
    \draw[flow, simRust, shorten <=1pt, shorten >=1pt] (\x,0.5)--(\x,-0.5);}
  \node[font=\scriptsize\itshape, text=simRust] at (-0.05,0.0) {$\vE=-\Grad V$};
  % a third, floating conductor on the right
  \fill[simInk] (1.3,-0.5) circle (0.45);
  \node[font=\scriptsize, text=white] at (1.3,-0.5) {3};
  \draw[flow, simRust, shorten >=1pt] (0.55,0.8) .. controls (1.0,0.4) .. (1.0,-0.2);
  \draw[flow, simRust, shorten >=1pt] (2.05,0.8) .. controls (1.6,0.4) .. (1.6,-0.2);
  % equipotential note
  \node[font=\scriptsize, text=simGray, align=center] at (2.5,-1.25)
    {each conductor:\\ an equipotential};
\end{tikzpicture}
\caption[The electrostatic setup]{The physical setup. A set of conductors
(dark) is embedded in a dielectric of permittivity $\varepsilon$ (blue). Each
conductor is a perfect equipotential held at a chosen voltage $V_i$; the
dielectric between them carries the electric field $\vE=-\Grad V$ (rust lines,
running from high potential to low and meeting every conductor at a right angle).
The simulation prescribes the voltages and computes how much charge collects on
each conductor per applied volt---the capacitance matrix. The field lines from the
upper plate that terminate on the floating conductor~3 are exactly the
\emph{mutual} coupling the off-diagonal entries of $\mat{C}$ measure.}
\label{fig:es-setup}
\end{figure}

The governing physics is the simplest in the book, and we derived it in full in
\cref{sec:electrostatic}. With $\vE=-\Grad V$ (a scalar potential, which makes the
static law $\Curl\vE=\mathbf0$ hold identically) and the dielectric constitutive
relation $\vD=\varepsilon\vE$, Gauss's law $\Div\vD=\rho$ collapses to the
generalized Poisson equation
\begin{equation}
  -\Div(\varepsilon\,\Grad V) = \rho .
  \label{eq:es-poisson-ch39}
\end{equation}
In the charge-free dielectric \emph{between} the conductors $\rho=0$, so the
equation Palace actually solves in the interior is the homogeneous Laplace
equation $-\Div(\varepsilon\Grad V)=0$; all the free charge lives on the conductor
surfaces, and it enters the problem through the boundary conditions, not the
right-hand side. The unknown $V$ is a scalar field, so it lives in the scalar
finite-element space $\Hone$ (continuous, nodal) of \cref{ch:functionspaces}---the
first slot of the de Rham complex, the natural home for a potential.

\subsection{The method of unit excitations}
\label{sec:unit-excitations}

The capacitance matrix has $n^2$ entries, but we do not solve $n^2$ problems to
find them. We solve $n$. The trick is the \emph{method of unit excitations}, and
it is the structural heart of the whole analysis.

The relation $\mathbf q=\mat{C}\,\mathbf v$ is linear. Its $j$th column,
$\mat{C}\,\mathbf e_j$, is the vector of charges produced when the voltage vector
is the $j$th standard basis vector $\mathbf e_j$---that is, when conductor $j$ is
held at \emph{one volt} and all the others are \emph{grounded} at zero. So to read
off column $j$ of $\mat{C}$ we set up exactly that excitation, solve the Laplace
problem for the resulting potential field $V_j$, and extract the charges. Repeat
for each conductor $j=1,\dots,n$ and we have all $n$ columns---the whole matrix
(\cref{fig:es-unit}). The $n$ solves produce a \emph{family} of potential fields
$\{V_1,\dots,V_n\}$, one per terminal, and that family is the raw material the
reduction stage turns into $\mat{C}$.

\begin{figure}[t]
\centering
\begin{tikzpicture}[scale=1.0, font=\footnotesize]
  % three little panels, one per unit excitation
  \foreach \k/\dx in {1/0, 2/4.2, 3/8.4}{
    \begin{scope}[shift={(\dx,0)}]
      \draw[simGray, semithick, rounded corners=2pt] (-0.55,-1.25) rectangle (3.0,1.35);
      \node[font=\scriptsize\bfseries] at (1.25,1.05) {excite conductor \k};
    \end{scope}}
  % panel 1: conductor 1 at 1V, others 0
  \begin{scope}[shift={(0,0)}]
    \fill[simRust] (0.0,0.45) rectangle (1.0,0.7);  \node[font=\scriptsize] at (-0.3,0.57){$1$};
    \fill[simInk] (0.0,-0.7) rectangle (1.0,-0.45); \node[font=\scriptsize] at (-0.3,-0.57){$0$};
    \fill[simInk] (1.7,-0.1) circle (0.22);          \node[font=\scriptsize] at (2.25,-0.1){$0$};
    \foreach \x in {0.2,0.5,0.8}{\draw[flow,simBlue,shorten <=1pt,shorten >=1pt](\x,0.4)--(\x,-0.4);}
    \node[font=\scriptsize\itshape] at (0.5,-1.02){$\to V_1$};
  \end{scope}
  % panel 2: conductor 2 at 1V
  \begin{scope}[shift={(4.2,0)}]
    \fill[simInk] (0.0,0.45) rectangle (1.0,0.7);   \node[font=\scriptsize] at (-0.3,0.57){$0$};
    \fill[simRust] (0.0,-0.7) rectangle (1.0,-0.45);\node[font=\scriptsize] at (-0.3,-0.57){$1$};
    \fill[simInk] (1.7,-0.1) circle (0.22);          \node[font=\scriptsize] at (2.25,-0.1){$0$};
    \foreach \x in {0.2,0.5,0.8}{\draw[flow,simBlue,shorten <=1pt,shorten >=1pt](\x,-0.4)--(\x,0.4);}
    \node[font=\scriptsize\itshape] at (0.5,-1.02){$\to V_2$};
  \end{scope}
  % panel 3: conductor 3 at 1V
  \begin{scope}[shift={(8.4,0)}]
    \fill[simInk] (0.0,0.45) rectangle (1.0,0.7);   \node[font=\scriptsize] at (-0.3,0.57){$0$};
    \fill[simInk] (0.0,-0.7) rectangle (1.0,-0.45); \node[font=\scriptsize] at (-0.3,-0.57){$0$};
    \fill[simRust] (1.5,-0.1) circle (0.24);         \node[font=\scriptsize] at (2.25,-0.1){$1$};
    \draw[flow,simBlue,shorten >=1pt](1.25,0.3) .. controls (0.9,0.0) .. (1.0,-0.05);
    \draw[flow,simBlue,shorten >=1pt](1.0,0.55) .. controls (0.7,0.2) .. (0.85,0.0);
    \node[font=\scriptsize\itshape] at (1.25,-1.02){$\to V_3$};
  \end{scope}
\end{tikzpicture}
\caption[The method of unit excitations]{The method of unit excitations on three
conductors. To recover column $j$ of the capacitance matrix we hold conductor $j$
at \emph{one volt} (rust) and ground all the others (dark, $0$), then solve the
Laplace problem $-\Div(\varepsilon\Grad V_j)=0$ for the resulting potential field
$V_j$ (blue field lines). Three excitations give three solved fields
$\{V_1,V_2,V_3\}$---the solution \emph{family}. Each solve has \emph{the same}
governing operator; only the prescribed boundary voltage moves. This is why the
analysis sits in the fixed-operator corner of \cref{ch:situating}.}
\label{fig:es-unit}
\end{figure}

\begin{keyidea}
The electrostatic analysis never assembles the capacitance matrix directly. It
exploits linearity: column $j$ of $\mat{C}$ is what you get by holding conductor
$j$ at unit voltage and the rest at ground, solving for the potential, and reading
off the charges. So $n$ conductors need $n$ solves, producing a \emph{family}
$\{V_1,\dots,V_n\}$ of potential fields. Each solve uses the \emph{same} discrete
operator and differs only in the prescribed boundary voltage---which is exactly the
shape that makes the next three sections possible.
\end{keyidea}

With the physical question and the method fixed, we can run the skeleton. The
electrostatic pipeline is
\[
  \texttt{config} \to \underbrace{\texttt{assemble}}_{\S\ref{sec:es-assemble}}
  \to \underbrace{\texttt{solve}}_{\S\ref{sec:es-solve}}
  \to \underbrace{\texttt{reduce}}_{\S\ref{sec:es-reduce}} \to \texttt{product},
\]
and the next three sections are those three stages, in order. We then post-process
the fields (\S\ref{sec:es-fields}) and step back to read the whole pipeline as the
anchor for Part~V (\S\ref{sec:es-anchor}).

\section{Assemble, once}
\label{sec:es-assemble}

The first stage turns the continuous physics of \cref{eq:es-poisson-ch39} into a
discrete operator. We did the mathematics in \cref{ch:weak} and \cref{ch:assembly};
here we only need its result and---crucially---\emph{when} it runs.

Multiplying the Laplace equation by a test function $v\in\Hone$ and integrating by
parts (the weak form, \cref{ch:weak}) turns the differential operator
$-\Div(\varepsilon\Grad\,\cdot\,)$ into the $\varepsilon$-weighted \emph{stiffness}
bilinear form
\begin{equation}
  a(u,v) \;=\; \int_\dom \varepsilon\,\Grad u \cdot \Grad v \dV ,
  \label{eq:es-bilinear-ch39}
\end{equation}
the single \code{weak\_form\_term} pairing the coefficient $\varepsilon$ with the
\emph{gradient} differential operator. Assembled over the mesh by the
assembly fold of \cref{ch:assembly}, this becomes a sparse, symmetric,
positive-definite matrix---the electrostatic stiffness operator $\Kop$. In the
framework's vocabulary the assembly is \code{fe\_assemble}, a fold that sums one
contribution per weak-form term; electrostatics is its simplest possible use,
because the term list has exactly one entry:
\begin{lstlisting}[style=l4style]
-- assemble the H1 stiffness operator from the single epsilon-weighted gradient term
space = h1_space cfg                          -- the H1 finite-element space (read-only)
terms = [ diffusion (permittivity cfg) ]      -- a one-term weak form: epsilon * grad
k     = fe_assemble space terms               -- fold the term list -> the stiffness K
\end{lstlisting}
The space and the term list are read-only construction-stratum data
(\cref{ch:strata}); $\Kop$ is a fresh value the fold returns. There is no mutation
and no loop we write---the loop is \emph{inside} \code{fe\_assemble}, over the term
list, and that list has length one.

\subsection{The boundary conditions enter by elimination}

The conductor voltages are essential (Dirichlet) boundary conditions
(\cref{ch:maxwell}, \S\ref{sec:bc}; \cref{ch:bc}): they constrain the value of $V$
on the conductor surfaces, so they must be built into the solution space rather
than added as a source. The mechanism is \emph{elimination}: the degrees of
freedom pinned by the boundary condition are removed from the unknown vector, and
their known values are moved to the right-hand side. The constrained operator
$\Kop$ is the stiffness matrix restricted to the free (interior) degrees of
freedom, and it is this reduced, still-symmetric-positive-definite matrix that the
solver receives. The free charge that Gauss's law puts on the conductor surfaces
never appears as a volume source; it is implicit in the eliminated boundary data.

\subsection{The load-bearing word is \emph{once}}

Here is the fact this section exists to make. The operator $\Kop$ is assembled
\emph{exactly once}, before any solving begins, and then held unchanged for the
entire analysis. Every one of the $n$ unit-excitation solves uses the same $\Kop$.
The driver makes this concrete: it builds the Krylov solver and hands it the
operator with a single \code{SetOperators} call, and that call sits \emph{outside}
the terminal loop.\footnote{In the Palace driver the operator is assembled by
\code{laplace\_op.GetStiffnessMatrix()} and captured by \code{ksp.SetOperators(*K,
*K)} \emph{before} the per-terminal \code{for} loop begins
(\code{palace/drivers/electrostaticsolver.cpp:30,36,59}). The placement of that one
call relative to that one loop is the entire structural signature of the
fixed-operator corner.}

\begin{pitfall}
The single most important line in the electrostatic driver is not a line of
arithmetic---it is the \emph{position} of the \code{SetOperators} call. It belongs
\emph{above} the terminal loop. An implementation that ``tidied'' it inside the
loop would reassemble (or at least re-factor and re-capture) the operator on every
terminal, doing $n$ times the setup work for no change in the answer. Worse, it
would erase the very property that distinguishes electrostatics from the driven
analysis of \cref{ch:driven}, where the operator genuinely \emph{does} change each
step and the call \emph{must} live inside the loop. Hoisting the operator is not an
optimization detail bolted onto the analysis; it \emph{is} the analysis's identity.
\end{pitfall}

\section{Solve: the fixed-operator map}
\label{sec:es-solve}

With $\Kop$ assembled and captured, the middle stage runs the family of solves.
This is the stage with all the structure, and it is the canonical occupant of the
first solve corner of \cref{ch:situating}: the \emph{fixed-operator map}.

\subsection{One operator, a fan of right-hand sides}

For each terminal $j$ the driver forms the right-hand side $\vb_j$ that encodes
``conductor $j$ at one volt, the rest grounded''---the eliminated boundary data of
the previous section, assembled into a vector. It then solves the linear system
\begin{equation}
  \Kop\,V_j \;=\; \vb_j
  \label{eq:es-system}
\end{equation}
for the potential field $V_j$. The operator $\Kop$ on the left is identical across
all $n$ systems; only the right-hand side $\vb_j$ changes. The $n$ systems are
therefore \emph{independent}---solving for $V_2$ needs nothing from $V_1$---so the
collection of solves is a \texttt{map} over the family of right-hand sides, not a
fold. They could be done in any order, or on $n$ separate machines at once
(\cref{fig:es-solvefamily}). This is the framework combinator \code{solve\_family},
and electrostatics is its defining witness:
\begin{lstlisting}[style=l4style]
-- the fixed-operator map: capture K ONCE, map the solver over the RHS family
solve_family :: LinOp -> [Tensor] -> [Tensor]
solve_family k bs = map (ksp_solve k) bs       -- k is captured once; reused per RHS

rhss = [ excitation cfg idx | idx <- terminal_sources cfg ]  -- one RHS per terminal
vs   = solve_family k rhss                       -- the solution family [V_1 .. V_n]
\end{lstlisting}
Read the body slowly. The operator \code{k} is a \emph{free variable} of the
mapped function \code{ksp\_solve k}; the \texttt{map} closes over it once and reuses
it for every element. That closure \emph{is} the hoist---the operator-capture made
structural by the type, so that no implementation can accidentally rebuild it
inside the loop. Each element solve \code{ksp\_solve k b\_j} is one preconditioned
conjugate-gradient solve of the SPD system \cref{eq:es-system}.

\begin{figure}[t]
\centering
\begin{tikzpicture}[font=\footnotesize, node distance=6mm]
  % the single fixed operator
  \node[opbox, minimum width=14mm, minimum height=12mm, fill=simBgBlue, draw=simBlue]
    (K) {$\Kop$};
  \node[font=\scriptsize\itshape, above=1.5mm of K] {assembled \emph{once}};
  % the fan of RHS and solves
  \foreach \j/\y in {1/2.2, 2/0.9, 3/-0.4, 4/-1.7}{
    \node[data, minimum width=9mm, right=20mm of K, yshift=(\y-0.25)*10mm]
      (b\j) {$\vb_{\j}$};
    \node[product, minimum width=10mm, right=14mm of b\j] (v\j) {$V_{\j}$};
    \draw[flow, shorten >=1pt] (K.east) .. controls +(0.9,0) .. (b\j.west);
    \draw[flow, shorten >=1pt] (b\j.east) -- node[above,font=\scriptsize]{\code{ksp\_solve}} (v\j.west);
  }
  \node[font=\scriptsize\itshape, text=simGreen, align=center]
    at ($(v4)+(0.1,-1.0)$) {independent solves:\\ embarrassingly parallel over terminals};
\end{tikzpicture}
\caption[\code{solve\_family} as a fixed-operator map]{The solve stage as a
\emph{fixed-operator map}. One operator $\Kop$ (blue, assembled once) feeds a fan
of independent solves, each with its own right-hand side $\vb_j$ (the $j$th unit
excitation) and producing its own potential field $V_j$. Because every solve
reuses the same $\Kop$ and none depends on another's result, the family is a
\texttt{map}: reorderable, parallelizable, ``embarrassingly parallel over
terminals.'' This is the combinator \code{solve\_family} of \cref{ch:combinators},
and the electrostatic terminal sweep is its defining witness. Contrast the driven
analysis (\cref{ch:driven}), where a \emph{fresh} operator feeds each solve.}
\label{fig:es-solvefamily}
\end{figure}

\subsection{The element solve is preconditioned CG}

Each \code{ksp\_solve} is a call to the conjugate-gradient method we built in
\cref{ch:cg}. CG is the right engine here for a specific reason: the constrained
stiffness $\Kop$ is symmetric positive-definite, and CG is the optimal Krylov
method for exactly that class (\cref{ch:cg}). In practice it is run with a
preconditioner---a cheap approximate inverse $\mat{P}^{-1}\approx\Kop^{-1}$ that
clusters the spectrum and cuts the iteration count, as developed in
\cref{ch:precond}. The convergence is the geometric decay of the CG error in the
$\Kop$-norm; \cref{fig:es-cg} shows a representative residual history for one
terminal solve, with and without preconditioning. The point for \emph{this} chapter
is only that the element solve is a self-contained, well-understood SPD solve---a
single piece of Part~IV machinery---and that the family runs $n$ of them over the
one operator.

\begin{figure}[t]
\centering
\begin{tikzpicture}
\begin{axis}[
    width=0.74\linewidth, height=52mm,
    xlabel={CG iteration $k$}, ylabel={relative residual $\norm{\vr_k}/\norm{\vb}$},
    xlabel style={font=\scriptsize}, ylabel style={font=\scriptsize},
    xmin=0, xmax=60, ymin=1e-8, ymax=2,
    ymode=log, ytick={1,1e-2,1e-4,1e-6,1e-8},
    xtick={0,20,40,60},
    yticklabel style={font=\scriptsize}, xticklabel style={font=\scriptsize},
    legend style={font=\scriptsize, draw=simGray!50, at={(0.98,0.98)},
      anchor=north east},
    every axis plot/.append style={very thick},
    axis lines=left,
  ]
  % unpreconditioned: slow geometric decay
  \addplot[simRust, domain=0:60, samples=61] {exp(-0.26*x)};
  \addlegendentry{plain CG}
  % preconditioned: fast geometric decay
  \addplot[simBlue, domain=0:18, samples=19] {exp(-0.95*x)};
  \addlegendentry{preconditioned CG}
  % convergence tolerance line
  \addplot[simGray, densely dashed, domain=0:60, samples=2] {1e-6};
  \addlegendentry{tolerance}
\end{axis}
\end{tikzpicture}
\caption[CG convergence for one terminal solve]{A representative convergence
history for a single unit-excitation solve $\Kop V_j=\vb_j$. The relative residual
falls geometrically (a straight line on the log axis) until it crosses the
convergence tolerance (dashed). A good preconditioner (\cref{ch:precond}, blue)
clusters the spectrum of $\Kop$ and reaches tolerance in a fraction of the
iterations of plain CG (\cref{ch:cg}, rust). Because the operator is fixed across
the whole family, the same preconditioner is built once and reused for every
terminal solve---another dividend of the hoist.}
\label{fig:es-cg}
\end{figure}

\begin{notationbox}
The fixed operator pays a second dividend beyond avoiding reassembly. Because
$\Kop$ does not change, any \emph{preconditioner} built from it---an incomplete
factorization, a multigrid hierarchy (\cref{ch:precond})---is also built once and
reused across all $n$ solves. The expensive setup is amortized over the whole
family. A direct factorization would amortize even more dramatically: factor
$\Kop=\mat{L}\mat{L}^{\!\top}$ once, then each terminal is a pair of cheap
triangular solves. ``Assemble once'' quietly licenses ``factor once'' and
``precondition once''---the entire economic case for the fixed-operator corner.
\end{notationbox}

\section{Reduce: the symmetric Gram}
\label{sec:es-reduce}

The solve stage has handed us the family $\{V_1,\dots,V_n\}$. The reduction stage
turns it into the capacitance matrix. We met this reduction in full generality in
\cref{ch:reductions}; here it appears in its simplest and most important instance.

\subsection{Capacitance is an operator-weighted Gram}

The $(i,j)$ entry of the capacitance matrix is the operator-weighted bilinear form
of two family members,
\begin{equation}
  \boxed{\;C_{ij} \;=\; V_j^{\!\top}\,\Kop\,V_i\;}
  \label{eq:cap-gram-ch39}
\end{equation}
where $\Kop$ is the $\varepsilon$-weighted energy operator (the discrete form of
the electric energy density $\tfrac12\varepsilon\abs{\vE}^2$ of \cref{ch:maxwell},
eq.~\eqref{eq:energy-density}). This is the \emph{symmetric Gram reduction}
$G_{ij}=w(i,j)\,V_j^{\!\top}\Kop V_i$ of \cref{ch:reductions}, with the weight set
to unity, $w(i,j)=1$. The framework verb is \code{gram\_reduce}, and capacitance is
its $w\equiv1$ specialization:
\begin{lstlisting}[style=l4style]
-- the symmetric Gram reduction; capacitance is the unit-weight (w = 1) member
gram_reduce :: LinOp -> [Tensor] -> (Int -> Int -> Scalar) -> Matrix
gram_reduce k vs w =
  symmetric_from_upper                              -- mirror lower triangle from upper
    [ [ w i j * entry k vs i j | j <- [i .. m-1] ]  -- map over upper-triangle pairs
      | i <- [0 .. m-1] ]
  where
    m = length vs
    entry k vs i j
      | i == j    = matrix_weighted_norm (vs!!i) k  -- diagonal: V_i^T K V_i
      | otherwise = bilinear_form (vs!!j) k (vs!!i) -- off-diag:  V_j^T K V_i

capacitance = gram_reduce k vs (\i j -> 1)          -- w = 1: unit voltage excitation
\end{lstlisting}
Two structural facts are built into this body, and both are load-bearing for
capacitance specifically. First, the grid of entries is a \texttt{map} over family
pairs with no state threaded between them---the embarrassingly parallel pattern of
\cref{ch:reductions}: no entry needs another's value. Second, the matrix is
\emph{symmetric by construction}. Because $\Kop$ is symmetric,
$V_j^{\!\top}\Kop V_i = V_i^{\!\top}\Kop V_j$, so the verb computes only the upper
triangle and mirrors the lower (\code{symmetric\_from\_upper}). For an $n$-conductor
problem this is $\binom{n}{2}+n$ bilinear evaluations rather than $n^2$---the
symmetry is exploited, not merely observed (\cref{fig:es-gram}).

\begin{figure}[t]
\centering
\begin{tikzpicture}[font=\footnotesize]
  % the symmetric Gram grid
  \foreach \i in {0,1,2}{
    \foreach \j in {0,1,2}{
      \node[draw=simGray, semithick,
        fill={\ifnum\i=\j simBgRust\else simBgBlue\fi},
        minimum size=8mm, inner sep=0pt]
        (g\i\j) at (\j*0.84, -\i*0.84) {};}}
  % labels on entries
  \node[font=\scriptsize] at (0,0)        {$C_{11}$};
  \node[font=\scriptsize] at (0.84,-0.84) {$C_{22}$};
  \node[font=\scriptsize] at (1.68,-1.68) {$C_{33}$};
  \node[font=\scriptsize, simBlue] at (0.84,0)    {$C_{12}$};
  \node[font=\scriptsize, simBlue] at (0,-0.84)   {$C_{21}$};
  % mirror line
  \draw[densely dashed, simRust, semithick] (-0.55,0.55) -- (2.23,-2.23);
  \node[font=\scriptsize, text=simRust] at (2.0,0.4) {mirror};
  % computed-half brace
  \node[font=\scriptsize, align=center, text=simBlue] at (3.6,-0.1)
    {upper triangle\\ \emph{computed}:\\ $V_j^{\!\top}\Kop V_i$};
  \node[font=\scriptsize, align=center, text=simRust] at (3.6,-1.7)
    {lower triangle\\ \emph{mirrored}:\\ $C_{ji}=C_{ij}$};
  % diagonal = energy callout
  \node[font=\scriptsize, align=center, text=simRust] at (1.0,-2.7)
    {diagonal $=$ twice the stored field energy:\ $C_{ii}=V_i^{\!\top}\Kop V_i = 2U_e$};
\end{tikzpicture}
\caption[The capacitance Gram]{The capacitance matrix as a symmetric Gram. Each
entry is the operator-weighted bilinear form $C_{ij}=V_j^{\!\top}\Kop V_i$ of two
solved potential fields. The diagonal (rust) is a \emph{self}-pairing,
$C_{ii}=V_i^{\!\top}\Kop V_i$, which equals \emph{twice} the field energy stored
when conductor $i$ alone is at unit voltage---the self-capacitance. The
off-diagonal (blue) is a \emph{cross}-pairing of two different fields, the mutual
capacitance. Because $\Kop$ is symmetric, $C_{ij}=C_{ji}$: only the upper triangle
is computed and the lower is mirrored across the dashed diagonal. The symmetry is
structural, a consequence of the symmetric energy operator, not a coincidence to be
checked afterward.}
\label{fig:es-gram}
\end{figure}

\subsection{Self versus mutual; the energy diagonal}

The two halves of the matrix carry two distinct physical meanings.

The \emph{diagonal} $C_{ii}=V_i^{\!\top}\Kop V_i$ is a self-pairing. Applying the
energy operator $\Kop$ to the field $V_i$ and pairing the result back with $V_i$
yields exactly twice the electric field energy stored when conductor $i$ alone is
held at unit voltage. This is the \emph{energy form} of capacitance,
\begin{equation}
  C_{ii} \;=\; \frac{2U_e}{V_i^2}\Bigg|_{V_i=1} \;=\; 2U_e ,
  \qquad U_e=\tfrac12\!\int_\dom \varepsilon\abs{\vE}^2\dV ,
  \label{eq:cap-energy}
\end{equation}
the relation \cref{ch:maxwell} promised (eq.~\eqref{eq:energy-density}, $C=2U/V^2$),
specialized to the unit voltage we excite with. It is the same number a
field-energy integration would produce---which is exactly the cross-check of
Worked example~\ref{wex:es-energy} below. (This energy formulation is the one
Palace's driver actually computes, following the COMSOL AC/DC manual's
$C_{ii}=2U_e/V_i^2$ recipe.)

The \emph{off-diagonal} $C_{ij}=V_j^{\!\top}\Kop V_i$ is a cross-pairing of two
different fields, the \emph{mutual} capacitance: the charge induced on conductor
$i$ by a unit voltage on conductor $j$. A robust and slightly surprising fact is
that the off-diagonal of a Maxwell capacitance matrix is always \emph{negative}.
Raising conductor $j$'s voltage draws charge of the \emph{opposite} sign onto a
neighbouring grounded conductor through their shared field, so in
$q_i=C_{ii}V_i+\sum_{j\neq i}C_{ij}V_j$ the coupling term $C_{ij}V_j$ must reduce
$q_i$ as $V_j$ rises---hence $C_{ij}<0$. The matrix is symmetric, $C_{ij}=C_{ji}$,
a reciprocity the Gram construction makes automatic: the charge conductor $j$
induces on conductor $i$ per volt equals the charge $i$ induces on $j$ per volt,
because the energy operator is symmetric.

\begin{figure}[t]
\centering
\begin{tikzpicture}[font=\footnotesize, scale=1.0]
  % a 3x3 matrix with self (diag) vs mutual (off-diag) highlighted
  \matrix[matrix of math nodes, left delimiter={(}, right delimiter={)},
    nodes={minimum size=8mm, anchor=center},
    column sep=2pt, row sep=2pt] (C) at (0,0)
  {
    {\color{simRust}C_{11}} & {\color{simBlue}C_{12}} & {\color{simBlue}C_{13}} \\
    {\color{simBlue}C_{21}} & {\color{simRust}C_{22}} & {\color{simBlue}C_{23}} \\
    {\color{simBlue}C_{31}} & {\color{simBlue}C_{32}} & {\color{simRust}C_{33}} \\
  };
  % annotations
  \node[font=\scriptsize, text=simRust, align=left, anchor=west] at (2.3,0.85)
    {\textbf{diagonal} $C_{ii}>0$:\\ \emph{self}-capacitance\\ ($=2U_e$, the stored energy)};
  \node[font=\scriptsize, text=simBlue, align=left, anchor=west] at (2.3,-0.7)
    {\textbf{off-diagonal} $C_{ij}<0$:\\ \emph{mutual} capacitance\\ ($C_{ij}=C_{ji}$, reciprocity)};
\end{tikzpicture}
\caption[Self versus mutual capacitance]{The structure of the Maxwell capacitance
matrix. The \emph{diagonal} entries (rust) are the self-capacitances; each is
positive and equals twice the field energy stored with that conductor alone at unit
voltage. The \emph{off-diagonal} entries (blue) are the mutual capacitances; each is
\emph{negative}---a unit voltage on conductor $j$ induces opposite-sign charge on
conductor $i$---and the matrix is symmetric, $C_{ij}=C_{ji}$, by the reciprocity the
symmetric Gram construction guarantees. Inverting $\mat{C}$ gives the alternate
\emph{Maxwell} form relating the conductor potentials to their charges,
$\mathbf v=\mat{C}^{-1}\mathbf q$.}
\label{fig:es-selfmutual}
\end{figure}

\begin{keyidea}
The capacitance matrix is the symmetric Gram reduction
$C_{ij}=V_j^{\!\top}\Kop V_i$ of the per-terminal potential family against the
energy operator $\Kop$, with unit weight. Its diagonal is twice the stored field
energy (self-capacitance, always positive); its off-diagonal is the mutual
capacitance (always negative); and it is symmetric by construction because $\Kop$
is. Capacitance and inductance (\cref{ch:magnetostatics}) call the \emph{same}
verb, \code{gram\_reduce}, and differ only in the weight $w$---one verb, two
products.
\end{keyidea}

\subsection{The product, and its inverse}

The verb returns the $n\times n$ matrix $\mat{C}$. This is the physical product the
user ran the analysis to compute. Palace additionally forms the inverse
$\mat{C}^{-1}$ by a dense LAPACK factorization---a small, cheap operation on the
$n\times n$ result, where $n$ is the handful of conductors, not the millions of
mesh degrees of freedom. The inverse is the \emph{other} Maxwell form: where
$\mathbf q=\mat{C}\,\mathbf v$ gives charges from voltages, $\mathbf
v=\mat{C}^{-1}\mathbf q$ gives the voltages that a prescribed set of charges would
induce. The inversion is a downstream consumer of the reduction's output, not part
of the Gram reduction itself.

\section{Field post-processing}
\label{sec:es-fields}

Before assembling the whole pipeline, we account for one more thing the analysis
produces: the fields themselves. The capacitance matrix is the headline number, but
each solved potential $V_j$ also yields a visualizable electric field and a stored
energy, and these are computed as cheap post-processes of the family.

The electric field of excitation $j$ is the (negative) gradient of its potential,
\begin{equation}
  \vE_j \;=\; -\Grad V_j ,
  \label{eq:es-Efield}
\end{equation}
computed by applying the discrete gradient operator to the solved nodal potential
(the same $\Grad$ matrix the assembly produced). The energy stored in that field is
\begin{equation}
  U_e^{(j)} \;=\; \tfrac12\!\int_\dom \varepsilon\,\abs{\vE_j}^2 \dV
            \;=\; \tfrac12\,V_j^{\!\top}\Kop V_j ,
  \label{eq:es-energy}
\end{equation}
and the second equality is the bridge that makes the diagonal of $\mat{C}$ an energy:
the field-energy integral is \emph{exactly} half the Gram diagonal entry $C_{jj}$.
This is no coincidence---the energy operator $\Kop$ is, by construction
(\cref{ch:maxwell}, eq.~\eqref{eq:energy-density}), the discrete quadratic form
whose value on a field is twice its stored energy. The same observation lets the
\emph{domain-energy} reduction of \cref{ch:reductions} report \emph{where} the energy
sits---what fraction lives in each dielectric region (the participation factors of
\cref{ch:maxwell}, \S\ref{sec:Q})---by restricting the energy operator to each
domain. The energy is also the thread back to Poynting's theorem (\cref{ch:maxwell},
\S\ref{sec:energy}): in the static case the Poynting flux is zero and all the energy
is the stored electric energy \cref{eq:es-energy}, the term that the capacitance
reports.

\section{The anchor: the whole pipeline, and what it anchors}
\label{sec:es-anchor}

We have now traced every stage. It is time to stand back and see the analysis
whole, because its \emph{shape}---not its electrostatics---is what the rest of
Part~V is built on.

\subsection{The pipeline in one line}

The three stages compose into a single expression. At the framework level the whole
electrostatic feature is
\begin{lstlisting}[style=l4style]
electrostatic = gram_reduce . solve_family . fe_assemble
\end{lstlisting}
read right to left: \code{fe\_assemble} builds the operator once, \code{solve\_family}
maps the family of solves over it, and \code{gram\_reduce} collapses the solved
family to the capacitance matrix. Spelled out with the data flowing through it,
\begin{lstlisting}[style=l4style]
-- inputs = config; output = the capacitance matrix (the physical product)
electrostatic :: ElectrostaticConfig -> CapacitanceMatrix
electrostatic cfg =
  let space = h1_space cfg                          -- the H1 finite-element space
      k     = fe_assemble space [ diffusion (permittivity cfg) ]  -- (1) assemble K ONCE
      rhss  = [ excitation cfg idx | idx <- terminal_sources cfg ] -- per-terminal RHS family
      vs    = solve_family k rhss                    -- (2) fixed-operator per-terminal map
  in  gram_reduce k vs (\i j -> 1)                   -- (3) C_ij = V_j^T K V_i, w = 1
\end{lstlisting}
\Cref{fig:es-pipeline} draws it as a labeled diagram. This figure is the one to
remember: every later modality chapter is, quite literally, a redrawing of it with
one or two boxes replaced.

\begin{figure}[p]
\centering
\resizebox{\linewidth}{!}{%
\begin{tikzpicture}[node distance=7mm and 13mm, font=\footnotesize]
  % the pipeline, left to right
  \node[data, align=center, text width=20mm] (cfg)
    {config\\[1pt]\footnotesize $\varepsilon$, mesh,\\terminals};
  \node[stage, right=of cfg, align=center, text width=22mm] (asm)
    {\textbf{assemble}\\[1pt]\footnotesize \code{fe\_assemble}\\(once)};
  \node[stage, right=of asm, align=center, text width=24mm, fill=simBgGreen, draw=simGreen]
    (solve) {\textbf{solve}\\[1pt]\footnotesize \code{solve\_family}\\(per terminal)};
  \node[stage, right=of solve, align=center, text width=22mm, fill=simBgGold, draw=simGold]
    (red) {\textbf{reduce}\\[1pt]\footnotesize \code{gram\_reduce}\\($w{=}1$)};
  \node[product, right=of red, align=center, text width=18mm] (prod)
    {$\mat{C}$\\[1pt]\footnotesize capacitance\\matrix};
  % flow edges with the carried data labelled
  \draw[flow] (cfg) -- (asm);
  \draw[flow] (asm) -- node[above,font=\scriptsize]{$\Kop$} (solve);
  \draw[flow] (solve) -- node[above,font=\scriptsize]{$\{V_j\}$} (red);
  \draw[flow] (red) -- (prod);
  % the shared-operator capture: K also feeds the reduction
  \draw[depedge, shorten >=1pt] (asm.south) ++(0,-1mm)
    .. controls +(0,-9mm) and +(0,-9mm) .. (red.south)
    node[midway, below, font=\scriptsize, text=simBlue]
      {the \emph{same} $\Kop$ weights the Gram};
  % corner / shape annotations under each stage
  \node[font=\scriptsize\itshape, text=simGreen, below=10mm of solve, align=center]
    {fixed-operator map\\(\cref{ch:situating}, corner 1)};
  \node[font=\scriptsize\itshape, text=simGold, above=2mm of red, align=center]
    {rank-2 symmetric Gram};
\end{tikzpicture}%
}
\caption[The electrostatic pipeline---the anchor diagram]{The complete
electrostatic pipeline, $\texttt{electrostatic} = \code{gram\_reduce}\circ
\code{solve\_family}\circ\code{fe\_assemble}$. A configuration of permittivity,
mesh, and terminals enters; \code{fe\_assemble} builds the stiffness operator
$\Kop$ \emph{once}; \code{solve\_family} maps the per-terminal solves over that
fixed operator (the fixed-operator corner of \cref{ch:situating}), producing the
potential family $\{V_j\}$; and \code{gram\_reduce} pairs that family against
$\Kop$ into the capacitance matrix $\mat{C}$ (the rank-2 symmetric Gram). The blue
back-edge marks a feature special to this analysis: the \emph{same} operator
$\Kop$ that drove the solves also weights the reduction. \textbf{This diagram is
the anchor.} Every other analysis in Part~V is this pipeline with the FE space, the
solve corner, and/or the reduction swapped---and the swaps are small and
well-localized.}
\label{fig:es-pipeline}
\end{figure}

\subsection{Every other analysis is this, with one stage swapped}

Here is the payoff that earns this chapter its length. The electrostatic pipeline
of \cref{fig:es-pipeline} is the \emph{template} for all six analyses. Each of the
others is obtained by swapping one or two boxes---and nothing else:

\begin{itemize}
  \item \textbf{Magnetostatics} (\cref{ch:magnetostatics}) swaps the
  \emph{function space} ($\Hone\to\Hcurl$, scalar potential $\to$ vector potential)
  and the \emph{Gram weight} ($w=1\to w=1/(I_iI_j)$, the current normalization). The
  solve corner is unchanged---it is still a fixed-operator map---and the reduction
  is still \code{gram\_reduce}. Two cells move; the skeleton is identical. It is the
  closest sibling, almost line for line the same analysis.
  \item \textbf{The driven analysis} (\cref{ch:driven}) keeps the map but swaps the
  \emph{solve corner} from fixed to operator-varying: the system operator
  $\mat{A}(\omega)=\Kop+i\omega\Cop-\omega^2\Mop$ depends on frequency, so it is
  rebuilt \emph{inside} the sweep---the \code{SetOperators} call moves \emph{into}
  the loop, the exact negation of \S\ref{sec:es-assemble}'s hoist. The reduction
  also swaps, from the symmetric Gram to the port-projection
  (\cref{ch:reductions}).
  \item \textbf{The eigenmode analysis} (\cref{ch:eigenmodes}) swaps the
  \emph{solve corner} for the opaque black box: there is no loop we write at all---the
  assembled pencil $(\Kop,\Mop)$ goes to a library eigensolver with one call. The
  reduction swaps to the rank-1 $(f,Q)$ table.
  \item \textbf{The transient analysis} (\cref{ch:transient}) swaps the
  \emph{map for a fold}: each time step begins from the state the previous step
  produced, so the independent family becomes a chained, sequential evolution.
  \item \textbf{The boundary-mode analysis} (\cref{ch:waveports}) is the eigenmode
  swap, performed on a two-dimensional cross-section instead of the full volume.
\end{itemize}

\begin{table}[t]
\centering
\caption[The modalities as swaps of the anchor]{Every Part~V modality as a swap of
the electrostatic anchor (\cref{fig:es-pipeline}). The columns are the stages of the
skeleton; a dash means ``unchanged from electrostatics.'' Reading down a column shows
which stage each analysis perturbs; reading across a row shows how few cells move.
Magnetostatics is the closest sibling (space and weight only); the others each swap
the solve corner, and most also the reduction. This table is the map of Part~V.}
\label{tab:es-swaps}
\small
\setlength{\tabcolsep}{4.5pt}
\renewcommand{\arraystretch}{1.3}
\begin{tabular}{@{}lllll@{}}
\toprule
Analysis & FE space & Assemble & Solve corner & Reduction \\
\midrule
\textbf{Electrostatic} (anchor) & $\Hone$ & $\Kop$ once & fixed-operator map & Gram, $w{=}1$ \\
Magnetostatic & $\Hcurl$ & --- & --- & Gram, $w{=}\tfrac{1}{I_iI_j}$ \\
Driven & $\Hcurl$ & $\mat{A}(\omega)$ in loop & operator-varying map & projection \\
Eigenmode & $\Hcurl$ & pencil once & opaque black box & table $(f,Q)$ \\
Transient & $\Hcurl$ & --- & state-threaded fold & table (history) \\
Boundary mode & $\Hcurl{\oplus}\Hone$ & 2-D pencil & opaque black box & table (modes) \\
\bottomrule
\end{tabular}
\end{table}

\Cref{tab:es-swaps} lays the swaps out as a table. The dashes are the point: most
cells of most rows are ``unchanged from electrostatics.'' This is why the book can
afford to teach electrostatics once, here, in complete detail, and then present
each remaining modality as a short, well-localized set of differences from it. The
full end-to-end trace---this chapter's worked example below, but on a real
three-dimensional device---is revisited as the capstone of Part~V in
\cref{ch:fulltrace}, where we run the entire pipeline against actual machinery and
check the answer.

\begin{keyidea}
The electrostatic pipeline $\code{gram\_reduce}\circ\code{solve\_family}\circ
\code{fe\_assemble}$ is the \textbf{anchor} for all of Part~V: assemble the operator
\emph{once}, \texttt{map} the solve over the excitation family (the fixed-operator
corner), and Gram-reduce the solved family to the product. Every other analysis is
this template with the FE space, the solve corner, and/or the reduction swapped---
magnetostatics swaps the space and the Gram weight, the driven analysis swaps to an
operator-varying map and a projection, eigenmode swaps the solve to a black box,
transient swaps the map for a fold. Learn this pipeline and you have learned the
spine of every analysis Palace runs.
\end{keyidea}

\section{Two worked examples}
\label{sec:es-worked}

The pipeline is general; we now run it end to end on a problem small enough to do by
hand. The first example traces the whole analysis---assemble, solve the two unit
excitations, Gram-reduce, mirror, interpret---and recovers the textbook
parallel-plate formula. The second checks that the Gram diagonal really is the field
energy, closing the loop on \cref{eq:es-energy}.

\begin{workedexample}{A two-conductor capacitance, end to end}{es-twocond}
Two conductors form a parallel-plate-like capacitor: an upper plate (conductor~1)
and a lower plate (conductor~2), with a uniform dielectric between them. We discretize
the gap so coarsely that the free potential reduces to a \emph{single} interior
degree of freedom $u$---the potential at the midplane---which is all we need to see
every stage of the pipeline turn.

\medskip
\emph{Assemble (once).} On this one-degree-of-freedom mesh the $\varepsilon$-weighted
stiffness form \cref{eq:es-bilinear-ch39} assembles, before elimination, into a
$3\times3$ matrix over the nodes (plate~1, midplane, plate~2). Writing
$g=\varepsilon A/(L/2)$ for the conductance of each half-gap (plate area $A$,
half-gap $L/2$), the unconstrained stiffness is
\[
  \Kop_{\text{full}} =
  \begin{pmatrix} g & -g & 0 \\ -g & 2g & -g \\ 0 & -g & g \end{pmatrix}.
\]
The plate nodes are Dirichlet (essential) degrees of freedom; eliminating them
leaves the single free interior unknown $u$ with the scalar constrained operator
$\Kop=2g$ and a right-hand side built from the prescribed plate voltages. This is
the assemble stage, run \emph{once}; the same $\Kop=2g$ serves both solves below.

\medskip
\emph{Solve the family (fixed-operator map).} Two unit excitations, two solves of
$\Kop u=b$ with the \emph{same} $\Kop=2g$:
\begin{itemize}[nosep]
  \item \textbf{Excitation 1:} plate~1 at $1$, plate~2 at $0$. Elimination gives the
  right-hand side $b_1=g\cdot1+g\cdot0=g$, so $u=b_1/\Kop=g/(2g)=\tfrac12$. The solved
  field $V_1$ has plate~1 at $1$, the midplane at $\tfrac12$, plate~2 at $0$---a linear
  ramp, the discrete image of the parallel-plate solution of \cref{wex:plates1d}.
  \item \textbf{Excitation 2:} plate~2 at $1$, plate~1 at $0$. By symmetry $u=\tfrac12$
  again, and $V_2$ ramps the other way: plate~2 at $1$, midplane at $\tfrac12$,
  plate~1 at $0$.
\end{itemize}
The two solves reused one operator; neither needed the other. That is the
fixed-operator map in miniature.

\medskip
\emph{Reduce (symmetric Gram).} Now form $C_{ij}=V_j^{\!\top}\Kop_{\text{full}}V_i$
on the full nodal fields (the energy operator pairs over all nodes). With
$V_1=(1,\tfrac12,0)$ and $V_2=(0,\tfrac12,1)$ we compute the upper triangle:
\[
  C_{11}=V_1^{\!\top}\Kop_{\text{full}}V_1
       = g\!\left[(1{-}\tfrac12)^2+(\tfrac12{-}0)^2\right] = g\cdot\tfrac12 = \tfrac{g}{2},
\]
and by the same geometry $C_{22}=\tfrac{g}{2}$. The off-diagonal is
\[
  C_{12}=V_2^{\!\top}\Kop_{\text{full}}V_1
       = g\bigl[(0{-}\tfrac12)(1{-}\tfrac12)+(\tfrac12{-}1)(\tfrac12{-}0)\bigr]
       = g\bigl[-\tfrac14-\tfrac14\bigr] = -\tfrac{g}{2}.
\]
We do \emph{not} recompute $C_{21}$: the symmetric Gram mirrors it,
$C_{21}=C_{12}=-\tfrac{g}{2}$. Assembling,
\[
  \mat{C} =
  \begin{pmatrix} \tfrac{g}{2} & -\tfrac{g}{2} \\[2pt] -\tfrac{g}{2} & \tfrac{g}{2} \end{pmatrix}
  = \frac{g}{2}\begin{pmatrix} 1 & -1 \\ -1 & 1 \end{pmatrix},
  \qquad g=\frac{\varepsilon A}{L/2}=\frac{2\varepsilon A}{L}.
\]

\medskip
\emph{Interpret.} The off-diagonal is negative, as every Maxwell capacitance matrix
demands: raising plate~2 induces opposite-sign charge on plate~1. The matrix is
symmetric, $C_{12}=C_{21}$, by reciprocity. And the physically meaningful
two-terminal capacitance---the charge that flows between the plates per volt of
\emph{difference}---is read from the matrix as $C_{\text{cap}}=-C_{12}=g/2=\varepsilon
A/L$, recovering the textbook parallel-plate result $C=\varepsilon A/L$ of
\cref{wex:plates1d}. The whole pipeline---assemble once, solve the two excitations,
Gram-reduce, mirror---has run in a page of arithmetic and produced the right answer.
On a fine three-dimensional mesh the only thing that changes is that the solves are
millions of degrees of freedom handled by CG (\cref{ch:cg}) instead of one division;
the \emph{structure} is exactly what we just did.
\end{workedexample}

\begin{workedexample}{The energy form agrees with the Gram diagonal}{es-energy}
We check the claim of \cref{eq:es-energy}: the diagonal Gram entry $C_{11}$ equals
twice the stored field energy of excitation~1, computed independently from the field.

\medskip
\emph{The Gram diagonal.} From \cref{wex:es-twocond}, $C_{11}=g/2$ with $g=2\varepsilon
A/L$, so $C_{11}=\varepsilon A/L$.

\emph{The field energy, computed directly.} Excitation~1 produced the field $V_1$
ramping linearly from $1$ at plate~1 to $0$ at plate~2 across the gap of width $L$,
so its electric field is uniform, $\vE_1=-\Grad V_1$ with magnitude
$\abs{\vE_1}=1/L$ (a unit potential drop over a distance $L$). The stored electric
energy is the energy density $\tfrac12\varepsilon\abs{\vE_1}^2$ integrated over the
gap volume $A\cdot L$:
\[
  U_e^{(1)} = \tfrac12\!\int_\dom \varepsilon\abs{\vE_1}^2\dV
            = \tfrac12\,\varepsilon\,\frac{1}{L^2}\,(A\,L)
            = \frac{\varepsilon A}{2L}.
\]

\emph{Compare.} Twice the energy is $2U_e^{(1)}=\varepsilon A/L$, which is
\emph{exactly} $C_{11}$. The Gram diagonal and the field-energy integral agree, as
\cref{eq:es-energy} promised:
\[
  C_{11} \;=\; V_1^{\!\top}\Kop V_1 \;=\; 2U_e^{(1)} \;=\; \frac{\varepsilon A}{L}.
\]
This is not a numerical coincidence. The energy operator $\Kop$ \emph{is}, by
construction, the discrete quadratic form whose value on a field is twice its stored
energy (\cref{ch:maxwell}, eq.~\eqref{eq:energy-density}); so the self-pairing
$V_i^{\!\top}\Kop V_i$ and the energy integral $2U_e$ are two evaluations of the same
quadratic form. The agreement is the reason the diagonal of the capacitance matrix
\emph{is} a capacitance: it measures the field energy a conductor stores per unit
voltage, which is exactly what capacitance means.
\end{workedexample}

\summarysection
The electrostatic analysis computes the capacitance matrix $\mat{C}$ of a system of
$n$ conductors in a dielectric. It exploits linearity through the \emph{method of
unit excitations}: hold conductor $j$ at one volt, ground the rest, solve the Laplace
problem $-\Div(\varepsilon\Grad V_j)=0$ for the potential, and repeat for each
conductor, producing a \emph{family} $\{V_1,\dots,V_n\}$ of potential fields. The
pipeline runs the skeleton of \cref{ch:situating} in its simplest form. \emph{Assemble}
($\S\ref{sec:es-assemble}$): the $\varepsilon$-weighted stiffness operator $\Kop$ on
the scalar space $\Hone$, with the terminal voltages eliminated as essential boundary
conditions, built \emph{once} and held---the load-bearing hoist. \emph{Solve}
($\S\ref{sec:es-solve}$): the fixed-operator map \code{solve\_family}, $n$ independent
preconditioned-CG solves $\Kop V_j=\vb_j$ reusing the one operator---the defining
witness of the first solve corner. \emph{Reduce} ($\S\ref{sec:es-reduce}$): the
symmetric Gram \code{gram\_reduce} with unit weight, $C_{ij}=V_j^{\!\top}\Kop V_i$,
whose diagonal is twice the stored field energy (self-capacitance, positive) and whose
off-diagonal is the mutual capacitance (negative), symmetric by construction. The whole
feature is the one-line composition
$\texttt{electrostatic}=\code{gram\_reduce}\circ\code{solve\_family}\circ
\code{fe\_assemble}$. This pipeline is the \textbf{anchor} of Part~V: every other
modality is this diagram (\cref{fig:es-pipeline}) with the FE space, the solve corner,
and/or the reduction swapped (\cref{tab:es-swaps}).

\furtherreading
For the electrostatics itself---capacitance, the energy method $C=2U/V^2$, and the
sign and reciprocity of the Maxwell capacitance matrix---Griffiths~\citep{griffiths2017}
is the standard undergraduate source, and Jin's finite-element electromagnetics
text~\citep{jin2014} treats the discretization, the $\Hone$ stiffness assembly, and the
energy-form capacitance extraction in the exact form Palace uses. Pozar~\citep{pozar2011}
places the capacitance matrix in its microwave-engineering context, alongside the
inductance matrix it pairs with (\cref{ch:magnetostatics}) and the S-parameters of the
driven analysis (\cref{ch:driven}). The fixed-operator solve corner and the symmetric
Gram reduction are this book's own framing, developed in \cref{ch:situating} and
\cref{ch:reductions}; the numerical engine behind each element solve is the conjugate
gradient method of \cref{ch:cg} with the preconditioners of \cref{ch:precond}.

\exercisessection
\begin{enumerate}[nosep]
  \item \textbf{Why $n$ solves, not $n^2$.} The capacitance matrix has $n^2$ entries.
  Explain, using the linearity of $\mathbf q=\mat{C}\,\mathbf v$, why the analysis
  needs only $n$ solves to fill them all. What is each solve's right-hand side, and
  which column of $\mat{C}$ does it produce?
  \item \textbf{The hoist.} Identify the single line of the pipeline whose position
  (outside versus inside the terminal loop) is the structural signature of the
  fixed-operator corner. What goes wrong---in \emph{cost}, and separately in
  \emph{correctness of the abstraction}---if it is moved inside? (Contrast with the
  driven analysis of \cref{ch:driven}, where it belongs inside.)
  \item \textbf{Mirror, don't recompute.} A four-conductor problem needs the full
  $4\times4$ matrix $\mat{C}$. How many distinct bilinear evaluations
  $V_j^{\!\top}\Kop V_i$ does \code{gram\_reduce} actually perform, and which entries
  does \code{symmetric\_from\_upper} fill instead? What property of $\Kop$ licenses the
  mirroring?
  \item \textbf{Sign of the off-diagonal.} Using $q_i=C_{ii}V_i+\sum_{j\neq
  i}C_{ij}V_j$ and the physical picture of induced charge, argue that every
  off-diagonal entry of a Maxwell capacitance matrix is negative. Check your argument
  against the $\mat{C}$ of \cref{wex:es-twocond}.
  \item \textbf{The energy diagonal.} Redo \cref{wex:es-energy} for excitation~2:
  compute the field energy $U_e^{(2)}$ of $V_2$ directly and confirm it equals
  $\tfrac12 C_{22}$. Why must $C_{11}=C_{22}$ for this symmetric two-plate geometry?
  \item \textbf{Recover $C=\varepsilon A/d$.} In \cref{wex:es-twocond} the
  two-terminal capacitance came out $-C_{12}=\varepsilon A/L$. Explain why the
  physical between-plates capacitance is $-C_{12}$ (the negative of the mutual entry)
  rather than a diagonal entry, and relate this to the difference between the
  $\mat{C}$ matrix and the elementary single-number capacitor formula.
  \item \textbf{One stage swapped.} Using \cref{tab:es-swaps}, describe the
  magnetostatic analysis (\cref{ch:magnetostatics}) as ``the electrostatic analysis
  with these stages changed.'' Exactly which cells move, and which stay? Then do the
  same for the eigenmode analysis (\cref{ch:eigenmodes})---which cell is the
  most radical swap, and why?
  \item \textbf{The composition.} Write out the meaning of
  $\texttt{electrostatic}=\code{gram\_reduce}\circ\code{solve\_family}\circ
  \code{fe\_assemble}$ as data flowing left to right: name the type of the value
  passed from each stage to the next (config, operator, field family, matrix), and
  identify the one value---the operator $\Kop$---that is consumed by \emph{two}
  stages rather than one. Why does that shared consumption appear as the back-edge in
  \cref{fig:es-pipeline}?
\end{enumerate}

\chapter{Magnetostatics and Inductance}
\label{ch:magnetostatics}

\chapterorient{The previous chapter took apart the electrostatic driver and found
a machine of disarmingly few parts: assemble one operator, solve it once per
conductor, and pair the solved family with itself through an energy operator to
read off a capacitance matrix. This chapter builds the magnetostatic driver---the
one that computes \emph{inductance}---and the surprise is how little there is to
build. It is the \emph{same machine}. The skeleton is identical, the solver is the
same conjugate-gradient solver, the reduction is the same Gram verb. Exactly three
things change, and we will name them at the outset and then spend the chapter
earning the right to say they are the \emph{only} three: the partial differential
equation and its field (a vector potential obeying the curl--curl equation, not a
scalar obeying Poisson), the finite-element space that field lives in
($\Hcurl$ edge elements, not $\Hone$ nodal elements), and one scalar weight in the
final reduction (a current normalization, not unity). Everything else---the
once-only assembly, the solve-per-source, the symmetric Gram, even the
positive-definiteness that lets conjugate gradients run---is shared verbatim with
\cref{ch:electrostatics}. By the end you will be able to point at the magnetostatic
pipeline, lay it beside the electrostatic one, and circle precisely the three boxes
that differ.}

\section{The physical question, and the machine we already have}
\label{sec:mag-question}

A handful of current-carrying loops sit in space---traces on a circuit board, turns
of a coil, the signal lines of a connector. Drive loop $i$ with a current $I_i$ and
it produces a magnetic field; that field threads the other loops, and a changing
current in loop $i$ induces a voltage in loop $j$. The bookkeeping of ``how much
flux loop $i$'s current links through loop $j$'' is the \emph{inductance matrix}
$\mat{L}$. Its diagonal $L_{ii}$ is the \emph{self}-inductance of loop $i$; its
off-diagonal $L_{ij}$ is the \emph{mutual} inductance between loops $i$ and $j$.
These are the magnetostatic analogue of the capacitances of
\cref{ch:electrostatics}, and an engineer wants them for the same reason: they are
the lumped circuit parameters that summarize a piece of geometry.

The temptation is to reach for the Biot--Savart law and integrate. We do not. The
whole point of \cref{ch:electrostatics} was that the simulator does not compute a
capacitance by a special-purpose formula; it runs a \emph{generic pipeline} and the
capacitance falls out of a generic reduction. Magnetostatics runs the \emph{same}
pipeline. Before we change anything, let us write down the skeleton we are
inheriting.

\begin{physicsbox}
\textbf{Setup.} $n$ current sources (loops, ports, coil terminals). Drive source
$i$ with unit current $I_i=1$; everything else carries zero current. Solve for the
resulting magnetic vector potential $\vA_i$, from which the magnetic field is
$\vB_i=\Curl\vA_i$. The stored magnetic energy when source $i$ alone is driven is
$U_m^{(i)}=\tfrac{1}{2}L_{ii}I_i^2$, and pairing the field of source $i$ against the
field of source $j$ gives the mutual inductance $L_{ij}$. \emph{We never integrate
Biot--Savart.} We solve a boundary-value problem for $\vA$ and read the inductances
off the solved family---exactly as electrostatics read capacitances off the solved
potential family.
\end{physicsbox}

\Cref{ch:electrostatics} distilled the electrostatic driver to a three-stage
skeleton---the \texttt{assemble}, \texttt{solve}, \texttt{reduce} stages of
\cref{ch:reductions}. Written as a composition, with the fixed-operator solver
corner that \cref{ch:electrostatics} identified:
\begin{lstlisting}[style=l4style]
-- the ELECTROSTATIC driver (Chapter 39), for reference
electrostatic = gram_reduce (w = 1) . solve_family . fe_assemble
  --             ^ reduce: symmetric    ^ one CG solve   ^ assemble
  --               Gram, unit weight       per terminal     ONCE
\end{lstlisting}
Read it right to left, the direction the data flows. \code{fe\_assemble} builds the
operator \emph{once}---the geometry does not change between excitations, so the
matrix is assembled a single time and reused. \code{solve\_family} solves that one
fixed operator against a family of right-hand sides, one per terminal; because the
operator is symmetric positive-definite and never changes, this is the
\emph{fixed-operator} corner of the solve taxonomy (\cref{ch:reductions}), and the
solver is plain conjugate gradients with the operator's preconditioner reused
across the family. \code{gram\_reduce} pairs the solved family with itself through
the energy operator and, with unit weight, returns the capacitance matrix
(\cref{ch:reductions}, eq.~\eqref{eq:cap-gram}).

\begin{keyidea}
The magnetostatic driver is the electrostatic driver with three substitutions and
\emph{nothing else}. The skeleton
\code{reduce} $\circ$ \code{solve\_family} $\circ$ \code{fe\_assemble} is identical;
the assemble-once / solve-per-source / Gram-reduce structure is identical; the
conjugate-gradient solver in the fixed-operator corner is identical. The three
swaps are: \textbf{(1)} the PDE and field---the curl--curl equation
$\Curl(\nu\Curl\vA)=\vJ$ for a vector potential $\vA$, not Poisson for a scalar
$V$; \textbf{(2)} the finite-element space---$\Hcurl$ edge elements, not $\Hone$
nodal elements; \textbf{(3)} the Gram weight---$w=1/(I_iI_j)$, not $w=1$. We spend
the chapter on those three, because they are all there is.
\end{keyidea}

\Cref{fig:mag-pipeline} draws the magnetostatic pipeline laid directly over the
electrostatic one, with the three swapped boxes called out in rust and everything
shared drawn in the muted shared palette. It is the hero figure of the chapter:
keep it in view as we take the swaps one at a time.

\begin{figure}[t]
\centering
\resizebox{\linewidth}{!}{%
\begin{tikzpicture}[node distance=6mm and 13mm, font=\footnotesize]
  % --- shared config input ---
  \node[data, align=center, text width=20mm] (cfg) {config\\\footnotesize geometry,
    sources};
  % --- assemble (SWAPPED: operator) ---
  \node[stage, fill=simBgRust, draw=simRust, right=of cfg, align=center,
    text width=24mm] (asm) {\textbf{fe\_assemble}\\\footnotesize \textsc{swap 1+2}};
  % --- solve (SHARED) ---
  \node[stage, fill=simBgGray, draw=simGray, right=of asm, align=center,
    text width=24mm] (slv) {\textbf{solve\_family}\\\footnotesize CG, fixed op
    \emph{(shared)}};
  % --- reduce (SWAPPED: weight) ---
  \node[stage, fill=simBgRust, draw=simRust, right=of slv, align=center,
    text width=24mm] (red) {\textbf{gram\_reduce}\\\footnotesize \textsc{swap 3}};
  % --- product ---
  \node[product, right=of red, align=center, text width=18mm] (prod)
    {$\mat{L}$\\\footnotesize inductance};
  \draw[flow] (cfg) -- (asm);
  \draw[flow] (asm) -- (slv);
  \draw[flow] (slv) -- (red);
  \draw[flow] (red) -- (prod);
  % --- swap annotations underneath ---
  \node[font=\scriptsize, simRust, align=center, below=8mm of asm]
    (s12) {\textsc{swap 1}: curl--curl PDE\\\textsc{swap 2}: $\Hcurl$ space\\
      $\Curl(\nu\Curl\vA)=\vJ$};
  \node[font=\scriptsize, simGray, align=center, below=8mm of slv]
    (ss) {identical to\\\cref{ch:electrostatics}:\\SPD operator, CG};
  \node[font=\scriptsize, simRust, align=center, below=8mm of red]
    (s3) {\textsc{swap 3}: weight\\$w=\dfrac{1}{I_iI_j}$\\(was $w=1$)};
  \draw[refedge] (asm) -- (s12);
  \draw[refedge] (slv) -- (ss);
  \draw[refedge] (red) -- (s3);
\end{tikzpicture}%
}
\caption[The magnetostatic pipeline, three swaps from electrostatics]{The
magnetostatic driver drawn over the electrostatic skeleton of
\cref{ch:electrostatics}. The three rust boxes are the \emph{only} things that
change. \textsc{Swap 1+2} live in \code{fe\_assemble}: the partial differential
equation becomes the curl--curl equation $\Curl(\nu\Curl\vA)=\vJ$ (\textsc{swap 1})
and the finite-element space becomes $\Hcurl$ edge elements (\textsc{swap 2}), so a
different stiffness operator is assembled. \textsc{Swap 3} lives in
\code{gram\_reduce}: the Gram weight becomes the current normalization
$w=1/(I_iI_j)$. The grey \code{solve\_family} box is byte-for-byte the electrostatic
solve---the same conjugate-gradient solver on a symmetric positive-definite,
fixed-across-the-family operator. Assemble once, solve per source, reduce by
pairing: the machine is shared; only the three rust boxes are magnetostatic.}
\label{fig:mag-pipeline}
\end{figure}

\section{Swap 1: the curl--curl equation and its field}
\label{sec:swap1}

Electrostatics solves Poisson's equation for a scalar potential. Magnetostatics
solves the \emph{curl--curl equation} for a vector potential. This is the first
swap, and it is forced by the physics: a magnetic field is a fundamentally
different kind of object from an electric one, and the de Rham complex
(\cref{ch:derham}) told us exactly which potential each admits.

\subsection{From Maxwell to curl--curl}

Two of Maxwell's laws govern magnetostatics. Amp\`ere's law (in the static,
source-driven regime) relates the magnetic field $\vH$ to the current density
$\vJ$, and the constitutive law relates $\vH$ to $\vB$ through the permeability:
\[
  \Curl\vH = \vJ, \qquad \vB = \mu\,\vH, \qquad \Div\vB = 0 .
\]
The last equation, $\Div\vB=0$---no magnetic monopoles---is the structural fact.
By exactness of the de Rham complex at the $\Hdiv$ rung (\cref{ch:derham},
\cref{thm:exact}), a divergence-free field is a curl: there exists a
\emph{magnetic vector potential} $\vA$ with
\[
  \boxed{\;\vB = \Curl\vA\;}
\]
Substitute this into Amp\`ere's law. Writing $\nu=1/\mu$ for the
\emph{reluctivity} (the magnetic analogue of conductivity---``how much a material
resists magnetic flux''), we have $\vH=\nu\vB=\nu\Curl\vA$, and Amp\`ere's law
becomes
\begin{equation}
  \boxed{\;\Curl\!\left(\nu\,\Curl\vA\right) = \vJ\;}
  \label{eq:curlcurl}
\end{equation}
the \emph{curl--curl equation} for the vector potential $\vA$, driven by the
prescribed source current $\vJ$. This is the magnetostatic governing equation
(\cref{ch:reductions-pde}). \Cref{fig:loop-field} shows the physical picture: a
current loop, the vector potential $\vA$ that circulates with the current, and the
field $\vB=\Curl\vA$ that loops through.

\begin{figure}[t]
\centering
\begin{tikzpicture}[scale=1.0, font=\footnotesize]
  % --- the current loop (an ellipse seen at an angle) ---
  \draw[simRust, very thick] (0,0) ellipse (1.5 and 0.55);
  % current direction arrows ON the loop
  \draw[-{Stealth[length=2mm]}, simRust, very thick] (1.5,0.05)--(1.5,-0.05);
  \draw[-{Stealth[length=2mm]}, simRust, very thick] (-1.5,-0.05)--(-1.5,0.05);
  \node[simRust] at (0,-0.85) {current loop, $I$};
  \node[simRust, font=\scriptsize] at (1.95,0.0) {$\vJ$};
  % --- B field lines threading the loop (vertical-ish loops) ---
  \draw[simBlue, opacity=0.4,  thick] (0,0) ellipse [x radius=0.45, y radius=0.675];
  \draw[simBlue, opacity=0.55, thick] (0,0) ellipse [x radius=0.85, y radius=1.275];
  \draw[simBlue, opacity=0.4,  thick] (0,0) ellipse [x radius=1.25, y radius=1.875];
  % B arrow up through the center
  \draw[-{Stealth[length=2.4mm]}, simBlue, thick] (0,-1.7)--(0,1.7);
  \node[simBlue] at (0.35,1.55) {$\vB$};
  \node[simBlue, font=\scriptsize, align=center] at (2.4,1.3)
    {$\vB=\Curl\vA$\\(field lines\\thread the loop)};
  % --- A field circulating WITH the current (faint, around the wire) ---
  \node[simGreen, font=\scriptsize, align=center] at (-2.5,0.9)
    {$\vA$ circulates\\\emph{with} the current\\(same direction as $\vJ$)};
  \draw[-{Stealth[length=1.8mm]}, simGreen, thick] (-1.95,0.35) arc (160:380:0.32 and 0.20);
\end{tikzpicture}
\caption[A current loop and its potential]{The physical picture behind the
curl--curl equation. A current $I$ runs around the loop (rust). The magnetic vector
potential $\vA$ (green) circulates \emph{with} the current---it lives one rung up
the de Rham complex from $\vB$, in $\Hcurl$, the space of circulations
(\cref{ch:derham}). Its curl is the magnetic field $\vB=\Curl\vA$ (blue), whose
field lines thread \emph{through} the loop. The simulator solves
$\Curl(\nu\Curl\vA)=\vJ$ for $\vA$ on the whole domain, then post-processes
$\vB=\Curl\vA$. Compare the electrostatic picture, where the scalar potential $V$
sits at \emph{vertices} and $\vE=-\Grad V$.}
\label{fig:loop-field}
\end{figure}

\subsection{The contrast with Poisson, term by term}

Lay \cref{eq:curlcurl} beside the electrostatic Poisson equation
$-\Div(\varepsilon\,\Grad V)=\rho$ of \cref{ch:electrostatics}. The two are the
\emph{same shape}---a second-order, self-adjoint, source-driven elliptic
equation---built from a different rung of the de Rham complex:
\[
  \underbrace{-\Div(\varepsilon\,\Grad V)=\rho}_{\text{electrostatics}}
  \qquad\longleftrightarrow\qquad
  \underbrace{\Curl(\nu\,\Curl\vA)=\vJ}_{\text{magnetostatics}} .
\]
The gradient is replaced by the curl, the divergence (its adjoint) by the curl
again (the curl is its own adjoint up to boundary terms), the scalar coefficient
$\varepsilon$ by the scalar coefficient $\nu$, the scalar unknown $V$ by the vector
unknown $\vA$, and the scalar source $\rho$ by the vector source $\vJ$. The weak
form (\cref{ch:weak}) follows the identical recipe: multiply by a test field,
integrate by parts once, and move one curl onto the test field. The bilinear form
is
\begin{equation}
  a(\vA,\vF) = \int_\dom \nu\,(\Curl\vA)\cdot(\Curl\vF)\dV,
  \qquad
  \ell(\vF) = \int_\dom \vJ\cdot\vF \dV,
  \label{eq:weak-curlcurl}
\end{equation}
to be compared with the electrostatic
$a(V,W)=\int\varepsilon\,\Grad V\cdot\Grad W$. Discretizing
\cref{eq:weak-curlcurl} by Galerkin projection (\cref{ch:galerkin}) produces the
\emph{curl--curl stiffness matrix} $\mat{K}$, the magnetostatic energy operator:
\begin{equation}
  \mat{K}_{ab} = \int_\dom \nu\,(\Curl\boldsymbol\phi_a)\cdot(\Curl\boldsymbol\phi_b)\dV,
  \label{eq:stiff-curlcurl}
\end{equation}
where $\boldsymbol\phi_a$ are the basis functions of the finite-element space---and
\emph{which} space those are is the second swap.

\begin{notationbox}
The same letter $\mat{K}$ denotes the energy/stiffness operator in both drivers,
because it plays the identical structural role: it is the symmetric
positive-(semi)definite operator that \code{solve\_family} inverts and that
\code{gram\_reduce} pairs the family through. In electrostatics
$\mat{K}_{ab}=\int\varepsilon\,\Grad\phi_a\cdot\Grad\phi_b$; in magnetostatics
$\mat{K}_{ab}=\int\nu\,\Curl\boldsymbol\phi_a\cdot\Curl\boldsymbol\phi_b$. One
symbol, one role, two integrands---the swap is entirely inside the integrand.
\end{notationbox}

\section{Swap 2: the field lives in \texorpdfstring{$\Hcurl$}{H(curl)}, on edge elements}
\label{sec:swap2}

The unknown is now a vector field $\vA$, and the single most consequential decision
in the whole chapter is which finite-element space to put it in. The naive choice
is to take three copies of the nodal $\Hone$ space---one per Cartesian component---
and store $\vA$'s components at vertices, as we stored $V$. \emph{This is wrong},
and the way it is wrong is the deepest lesson of \cref{ch:functionspaces} and
\cref{ch:derham}. The correct space is $\Hcurl$, realized by \emph{edge elements}
(N\'ed\'elec elements). This is the second swap.

\subsection{Why a vector potential wants edge elements}

The de Rham complex (\cref{ch:derham}) names the home of every electromagnetic
field, and it is unambiguous: $\vA$ generates $\vB$ by $\vB=\Curl\vA$, so $\vA$ sits
at the $\Hcurl$ rung---the space whose degrees of freedom are \emph{edge
circulations} $\int_e\vA\cdot d\vec\ell$, and whose defining continuity is that the
\emph{tangential component} of $\vA$ is continuous across element faces while the
normal component may jump. That tangential continuity is exactly what a vector
potential physically demands. The tangential component of $\vA$ carries its
circulation, hence its curl, hence the magnetic field $\vB$; the normal component
carries the gauge (\cref{ch:derham}, \cref{sec:potentials}) and is physically
free to jump. \Cref{fig:nodal-vs-edge} draws the two choices side by side.

\begin{figure}[t]
\centering
\begin{tikzpicture}[scale=1.0, font=\scriptsize]
% ===== LEFT: nodal vector (WRONG) =====
\begin{scope}
  \node[font=\footnotesize\bfseries, simRust] at (1.2,2.35) {nodal vector (wrong)};
  \coordinate (a) at (0,0); \coordinate (b) at (2.4,0); \coordinate (c) at (1.2,1.9);
  \fill[simBgRust] (a)--(b)--(c)--cycle;
  \draw[simInk, thick] (a)--(b)--(c)--cycle;
  % full vector at each vertex (all 3 components stored)
  \foreach \p in {a,b,c}{
    \fill[simRust] (\p) circle (2.6pt);
    \draw[-{Stealth[length=1.6mm]}, simRust, thick] (\p)++(0,0)--++(0.45,0.28);
    \draw[-{Stealth[length=1.6mm]}, simRust, thick] (\p)++(0,0)--++(-0.10,0.50);
  }
  \node[simRust, align=center] at (1.2,-0.55) {\emph{both} components\\at each vertex};
  \node[align=center, text=simInk] at (1.2,-1.35)
    {forces \textbf{full} continuity\\$\Rightarrow$ spurious gradient modes};
\end{scope}
% ===== RIGHT: edge (RIGHT) =====
\begin{scope}[shift={(5.6,0)}]
  \node[font=\footnotesize\bfseries, simGreen] at (1.2,2.35) {edge / N\'ed\'elec (right)};
  \coordinate (a) at (0,0); \coordinate (b) at (2.4,0); \coordinate (c) at (1.2,1.9);
  \fill[simBgGreen] (a)--(b)--(c)--cycle;
  \draw[simInk, thick] (a)--(b)--(c)--cycle;
  % one tangential circulation per edge
  \draw[-{Stealth[length=1.8mm]}, simGreen, very thick] ($(a)!0.32!(b)$)--($(a)!0.68!(b)$);
  \draw[-{Stealth[length=1.8mm]}, simGreen, very thick] ($(b)!0.32!(c)$)--($(b)!0.68!(c)$);
  \draw[-{Stealth[length=1.8mm]}, simGreen, very thick] ($(c)!0.32!(a)$)--($(c)!0.68!(a)$);
  \node[simGreen, align=center] at (1.2,-0.55) {\emph{one tangential}\\DoF per edge};
  \node[align=center, text=simInk] at (1.2,-1.35)
    {tangential continuity only\\$\Rightarrow$ exact discrete complex};
\end{scope}
\end{tikzpicture}
\caption[Nodal vector versus edge elements]{Two ways to discretize a vector field,
and why only one works. \emph{Left:} the naive nodal-vector element stores the full
vector (all components) at each \emph{vertex}, forcing \emph{full} continuity of
$\vA$ across faces---more continuity than the physics allows. The excess continuity
admits non-physical fields and (as the next section shows) pollutes the curl--curl
operator's null space with spurious modes. \emph{Right:} the edge (N\'ed\'elec)
element of $\Hcurl$ stores one \emph{tangential circulation} per \emph{edge},
imposing only tangential continuity---exactly what a vector potential requires. The
edge space sits in an \emph{exact} discrete de Rham complex (\cref{ch:derham},
\cref{thm:discrete-exact}), so its discrete curl has precisely the kernel the
physics demands, with no spurious extras. This is \emph{the} reason edge elements
exist (\cref{ch:bc}).}
\label{fig:nodal-vs-edge}
\end{figure}

The nodal-vector choice imposes \emph{full} continuity of $\vA$---both tangential
and normal components glued across every face. That is strictly more continuity
than $\vB=\Curl\vA$ requires, and the excess is not harmless. A field
discretization that is ``too continuous'' for its operator admits non-physical
solutions; in the curl--curl setting these surface as \emph{spurious modes}
(\cref{ch:functionspaces}, \cref{ch:bc})---finite-element fields with no physical
counterpart that nonetheless satisfy the discrete equations and contaminate the
spectrum. \Cref{ch:derham} diagnosed the cause precisely: nodal-vector elements do
not sit in an exact discrete de Rham complex, so their discrete curl has the
\emph{wrong} kernel, and the wrong kernel is the breeding ground of spurious modes.
Edge elements sit in an exact discrete complex (\cref{ch:derham},
\cref{thm:discrete-exact}); their discrete curl's kernel is \emph{exactly} the range
of the discrete gradient---the gauge fields, and nothing else.

\begin{pitfall}
Do \emph{not} discretize the magnetic vector potential with three copies of nodal
$\Hone$ elements ``one per component.'' It looks natural---it is how you would store
a vector in a general-purpose array---but it imposes more continuity than $\vA$
admits and breaks the exact discrete de Rham complex. The cost is spurious modes
(\cref{ch:bc}): non-physical fields polluting the curl--curl operator's null space,
which corrupt eigenvalue computations and stall iterative solvers. The vector
potential lives in $\Hcurl$, full stop, and $\Hcurl$ is realized by edge elements.
This is not an efficiency preference; it is a correctness requirement.
\end{pitfall}

\subsection{The curl--curl operator has a large null space}
\label{sec:nullspace}

Choosing edge elements fixes \emph{which} kernel the operator has, but it does not
make the kernel small. The curl--curl operator carries an enormous null space, and
acknowledging it is the price of working with vector potentials. The reason is the
gauge freedom of \cref{ch:derham}: for any scalar field $\psi$,
\[
  \Curl\!\left(\nu\,\Curl(\Grad\psi)\right) = \Curl\!\left(\nu\cdot\mathbf 0\right)
  = \mathbf 0,
  \qquad\text{since } \Curl\Grad\psi = \mathbf 0 .
\]
Every gradient field is annihilated by the curl--curl operator. Discretely, every
discrete-gradient vector $\mat{G}\psi_h$ lies in the null space of the curl--curl
stiffness $\mat{K}$ (\cref{eq:stiff-curlcurl}). So $\mat{K}$ is only \emph{positive
semi-definite}: it is singular, with a null space that is the entire range of the
discrete gradient $\mat{G}$ (\cref{ch:derham}, \cref{def:discrete-grad}). This is
the same gradient null space the eigenmode solver must project out
(\cref{ch:bc}) and the same null space a plain smoother cannot damp
(\cref{ch:smoothers}). \Cref{fig:curlcurl-null} draws it.

\begin{figure}[t]
\centering
\begin{tikzpicture}[font=\footnotesize]
  % the edge-dof space as a big rounded box
  \node[draw=simInk, thick, rounded corners=3pt, minimum width=66mm,
    minimum height=30mm, fill=simBgGray] (box) at (0,0) {};
  \node[font=\scriptsize\itshape, simInk, anchor=north west] at (box.north west)
    {\;edge-DoF space $\Hcurl_h$};
  % the gradient null space as an inner blob
  \node[draw=simBlue, thick, rounded corners=3pt, fill=simBgBlue,
    minimum width=30mm, minimum height=20mm, align=center] (null) at (-1.4,-0.3)
    {$\ker(\mat{K})$\\$=\operatorname{range}(\mat{G})$\\[1pt]
     \scriptsize gradient fields $\Grad\psi$\\\scriptsize ($\mat{K}\,\Grad\psi=0$)};
  % the physical / div-free complement
  \node[font=\scriptsize, simGreen, align=center, anchor=west] at (0.7,0.55)
    {physical part\\(divergence-free)\\$\mat{K}$ is SPD \emph{here}};
  \node[font=\scriptsize, simRust, align=center, below=1mm of box.south]
    {curl--curl $\mat{K}$ is only \emph{positive semi-definite}: singular on
     $\operatorname{range}(\mat{G})$};
\end{tikzpicture}
\caption[The curl--curl gradient null space]{The curl--curl stiffness operator
$\mat{K}$ on the edge-DoF space. Its null space (blue) is the \emph{entire} range of
the discrete gradient $\mat{G}$---every gradient field $\Grad\psi$ has zero curl, so
$\mat{K}\,\Grad\psi=0$ (the gauge freedom of \cref{ch:derham}). The operator is
therefore only positive \emph{semi}-definite, singular on a large subspace, and a
plain conjugate-gradient solve does not automatically converge. Two cures
(\cref{sec:solvecare}): make the right-hand side $\vJ$ \emph{consistent}
(divergence-free, so it has no component along the null space, and CG runs in the
SPD complement), or precondition with an auxiliary-space method (AMS,
\cref{ch:smoothers}) that handles the gradient block explicitly. Either way the
\emph{solver itself} is the shared CG of \cref{ch:electrostatics}; only its setup
accounts for the null space.}
\label{fig:curlcurl-null}
\end{figure}

\subsection{What this costs the solve---and what it does not}
\label{sec:solvecare}

A singular operator threatens \cref{fig:mag-pipeline}'s claim that the solve is
``identical to electrostatics.'' It is worth being precise about what changes. The
electrostatic $\mat{K}$ is positive \emph{definite} (after grounding a reference
node), so conjugate gradients converges with no further thought. The magnetostatic
$\mat{K}$ is positive \emph{semi}-definite, so a naive CG could wander in the null
space. Two standard cures restore convergence, and crucially \emph{neither changes
the solver}:

\begin{itemize}[nosep]
\item \textbf{Consistent source.} If the driving current $\vJ$ is
  divergence-free---$\Div\vJ=0$, the physical condition for a steady current that
  neither accumulates nor depletes charge---then $\vJ$ has \emph{no component along
  the null space} (the null space is gradients, and a divergence-free field is
  orthogonal to gradients). Conjugate gradients then runs entirely in the
  positive-definite complement, never excited along the singular directions, and
  converges to the (gauge-fixed) solution exactly as in electrostatics. This is the
  cheapest cure and the one Palace prefers: enforce a discretely-consistent source
  and the existing CG just works.
\item \textbf{Auxiliary-space preconditioning (AMS).} When robustness across mesh
  and coefficient is needed, the curl--curl solve is preconditioned by the
  \emph{auxiliary-space Maxwell} (AMS) preconditioner~\citep{hiptmair1998}, which
  splits the error into its divergence-free part and its gradient part and
  preconditions each on the appropriate auxiliary space (the nodal $\Hone$ space for
  the gradient block, via the discrete gradient $\mat{G}$). AMS is exactly the
  multigrid-for-edge-elements machinery of \cref{ch:smoothers}: it exists
  \emph{because} a plain smoother cannot see the gradient null space.
\end{itemize}

\begin{keyidea}
The curl--curl operator's gradient null space changes the solve's \emph{setup}, not
its \emph{solver}. Make the source consistent ($\Div\vJ=0$) so the null space is
never excited, or precondition with AMS (\cref{ch:smoothers}) so the gradient block
is handled explicitly---but in either case the iteration is the same
conjugate-gradient method on a symmetric, fixed-across-the-family operator that
electrostatics used. \code{solve\_family} in \cref{fig:mag-pipeline} is genuinely
shared; the null space is accounted for in how $\mat{K}$ and $\vJ$ are prepared, not
in a new algorithm.
\end{keyidea}

\section{Swap 3: the Gram weight is a current normalization}
\label{sec:swap3}

The solved family $\{\vA_0,\dots,\vA_{n-1}\}$ in hand, the reduce stage runs the
\emph{same} verb electrostatics ran---\code{gram\_reduce} of \cref{ch:reductions}---
and the third and final swap is a single scalar. This is the punchline of the
chapter, and it is the cleanest instance in the book of ``one verb, two physical
quantities.''

\subsection{One verb, both quantities}

Recall the Gram reduction (\cref{ch:reductions}, eq.~\eqref{eq:gram-entry}). Given a
solution family, the energy operator $\mat{K}$, and a scalar weight $w(i,j)$, it
computes the symmetric matrix $G_{ij}=w(i,j)\,\vx_j^{\!\top}\mat{K}\,\vx_i$:
\begin{lstlisting}[style=l4style]
gram_reduce :: LinOp -> [Tensor] -> (Int -> Int -> Scalar) -> Matrix
gram_reduce k xs w =
  symmetric_from_upper                              -- mirror lower triangle
    [ [ w i j * (xs!!j `bilinear` k $ xs!!i) | j <- [i .. m-1] ]
      | i <- [0 .. m-1] ]                           -- map over upper-triangle pairs
  where m = length xs
\end{lstlisting}
This is verbatim the reduction electrostatics used. The energy operator $\mat{K}$,
the symmetric pairing $\vx_j^{\!\top}\mat{K}\vx_i$, the upper-triangle map, the
mirror---all shared. Capacitance plugged in $w(i,j)=1$ and the per-terminal
potential family, getting $C_{ij}=V_j^{\!\top}\mat{K}\,V_i$
(\cref{ch:reductions}, eq.~\eqref{eq:cap-gram}). Inductance plugs in the
\emph{current-normalized} weight and the per-source vector-potential family:
\begin{equation}
  \boxed{\;w(i,j) = \frac{1}{I_i\,I_j},
  \qquad
  L_{ij} = \frac{1}{I_i I_j}\,\vA_j^{\!\top}\,\mat{K}\,\vA_i\;}
  \label{eq:ind-gram-local}
\end{equation}
with $\mat{K}$ the $\nu$-weighted curl--curl energy operator of
\cref{eq:stiff-curlcurl}. \Cref{fig:one-verb} sets the two side by side.

\begin{figure}[t]
\centering
\begin{tikzpicture}[font=\footnotesize]
% --- shared verb in the center ---
\node[stage, fill=simBgGold, draw=simGold, align=center, text width=34mm]
  (verb) at (0,0) {\textbf{gram\_reduce}\\[1pt]
  $G_{ij}=w(i,j)\,\vx_j^{\!\top}\mat{K}\,\vx_i$\\[1pt]\footnotesize
  symmetric, upper-triangle map};
% --- left branch: capacitance ---
\node[product, align=center, text width=30mm, above left=11mm and 14mm of verb]
  (cap) {\textbf{capacitance} $\mat{C}$\\[1pt]\footnotesize
  $w=1$\\family: potentials $V_i\in\Hone$};
% --- right branch: inductance ---
\node[product, fill=simBgBlue, draw=simBlue, align=center, text width=30mm,
  above right=11mm and 14mm of verb]
  (ind) {\textbf{inductance} $\mat{L}$\\[1pt]\footnotesize
  $w=\dfrac{1}{I_iI_j}$\\family: potentials $\vA_i\in\Hcurl$};
\draw[flow] (verb) -- node[left, font=\scriptsize, simGray]{$w\!=\!1$} (cap);
\draw[flow] (verb) -- node[right, font=\scriptsize, simBlue]
  {$w\!=\!\tfrac{1}{I_iI_j}$} (ind);
% --- the "everything else identical" band below ---
\node[font=\scriptsize\itshape, simGray, align=center, below=7mm of verb,
  text width=72mm]
  {identical: the operator-weighted bilinear form $\vx_j^{\!\top}\mat{K}\vx_i$, the
   structural symmetry $G_{ij}=G_{ji}$, the embarrassingly parallel map. \emph{The
   only differences are the scalar weight $w$ and which family is fed in.}};
\end{tikzpicture}
\caption[One verb, two physical quantities]{The single Gram verb
\code{gram\_reduce} (\cref{ch:reductions}) produces \emph{both} capacitance and
inductance. Feed it the per-terminal potential family with unit weight $w=1$ and it
returns the capacitance matrix $\mat{C}$ (left); feed it the per-source
vector-potential family with the current-normalized weight $w=1/(I_iI_j)$ and it
returns the inductance matrix $\mat{L}$ (right). The operator-weighted bilinear
form, the structural symmetry, and the parallel map are byte-for-byte identical---
the two physical quantities differ \emph{only} in the scalar weight and in which
solved family is supplied. This is \textsc{swap 3}, and it is the entire difference
in the reduce stage.}
\label{fig:one-verb}
\end{figure}

\subsection{Why inductance needs the normalization and capacitance does not}

The asymmetry is real but shallow, and it is worth saying plainly. In
electrostatics the excitation is a \emph{voltage}, and we choose it to be unity
($V_i=1$); the energy bilinear form $V_j^{\!\top}\mat{K}V_i$ is already
per-unit-voltage, so dividing by $V_iV_j=1$ does nothing and the weight is $1$. In
magnetostatics the natural excitation is a \emph{current} $I_i$, and the magnetic
energy bilinear form $\vA_j^{\!\top}\mat{K}\vA_i$ scales with the \emph{product of
the driving currents} (double the current, double $\vA$, quadruple the energy). To
recover the geometry-only inductance $L=2U_m/I^2$ (\cref{ch:maxwell}), we must
divide out that current scaling: $w=1/(I_iI_j)$. The weight is precisely the
``per-unit-excitation'' rescaling---invisible when the excitation is chosen to be
unity, explicit when it is a physical current we do not get to set to $1$. Same
energy form, same machine; one quantity happens to normalize trivially and the other
does not.

\begin{keyidea}
Capacitance and inductance are the \emph{same Gram reduction} differing in one
scalar weight: $w=1$ for capacitance, $w=1/(I_iI_j)$ for inductance. The diagonal
$L_{ii}=\vA_i^{\!\top}\mat{K}\vA_i/I_i^2$ is twice the stored magnetic energy per
unit current squared---the self-inductance $L=2U_m/I^2$. The off-diagonal $L_{ij}$
is the mutual inductance, the flux of source $j$'s field linked by loop $i$,
normalized by both currents. The current normalization is the \emph{only} thing the
reduce stage does differently from capacitance.
\end{keyidea}

\section{The composition, and reading the result}
\label{sec:compose}

We can now write the magnetostatic driver as a single composition and put it beside
the electrostatic one. They are the same expression with three substitutions:

\begin{lstlisting}[style=l4style]
-- ELECTROSTATIC (Chapter 39):
electrostatic = gram_reduce (w = 1)          . solve_family . fe_assemble{- Poisson, H^1 -}

-- MAGNETOSTATIC (this chapter), three swaps marked:
magnetostatic = gram_reduce (w = \i j -> 1/(I i * I j))   -- SWAP 3: current weight
              . solve_family                              --   (shared CG, fixed op)
              . fe_assemble{- curl-curl, H(curl) -}       -- SWAP 1+2: PDE + space

-- read B from the solved A:
postprocess A_i = curl A_i      -- B_i = Curl A_i  (one rung down the complex)
\end{lstlisting}
The structure is identical; the differences are exactly the curl--curl-on-$\Hcurl$
assembly (swaps 1 and 2, both inside \code{fe\_assemble}) and the current-normalized
Gram weight (swap 3, inside \code{gram\_reduce}). The forward reference to the
eigenmode driver (\cref{ch:eigenmodes}) is worth flagging now: it reuses this very
same curl--curl $\mat{K}$ on this same $\Hcurl$ space---the magnetostatic assembly
is the eigenmode assembly---and the gradient null space we met in
\cref{sec:nullspace} is precisely what the eigenmode solver must project out.

\subsection{The inductance matrix: self and mutual}

The product is the symmetric $n\times n$ inductance matrix $\mat{L}$. Reading it:

\begin{itemize}[nosep]
\item \textbf{Diagonal $L_{ii}$ (self-inductance).} The energy stored when source
  $i$ alone carries unit current, doubled: $L_{ii}=\vA_i^{\!\top}\mat{K}\vA_i/I_i^2
  =2U_m^{(i)}/I_i^2$. Always positive (the energy operator is positive
  semi-definite and $\vA_i$ is in its definite complement)---self-inductance is
  never negative.
\item \textbf{Off-diagonal $L_{ij}$ (mutual inductance).} The flux of source $j$'s
  field linked by loop $i$, per unit of each current:
  $L_{ij}=\vA_j^{\!\top}\mat{K}\vA_i/(I_iI_j)$. Its \emph{sign} encodes the relative
  orientation of the two loops---positive when their fields reinforce, negative when
  they oppose. Symmetric: $L_{ij}=L_{ji}$, the reciprocity that the Gram's structural
  symmetry guarantees for free.
\end{itemize}

\Cref{fig:Lmatrix} shows the matrix and its physical reading.

\begin{figure}[t]
\centering
\begin{tikzpicture}[font=\footnotesize]
  % --- the 3x3 inductance matrix with labeled cells ---
  \foreach \i in {0,1,2}{
    \foreach \j in {0,1,2}{
      \node[draw=simGray, semithick,
        fill={\ifnum\i=\j simBgBlue\else simBgGold\fi},
        minimum size=10mm, inner sep=0pt] (L\i\j) at (\j*1.0, -\i*1.0) {};}}
  \node at (0,0)    {$L_{11}$};
  \node at (1,-1)   {$L_{22}$};
  \node at (2,-2)   {$L_{33}$};
  \node at (1,0)    {$L_{12}$};
  \node at (2,0)    {$L_{13}$};
  \node at (0,-1)   {$L_{21}$};
  \node at (2,-1)   {$L_{23}$};
  \node at (0,-2)   {$L_{31}$};
  \node at (1,-2)   {$L_{32}$};
  % symmetry mirror
  \draw[densely dashed, simRust, semithick] (-0.55,0.55)--(2.55,-2.55);
  % annotations
  \node[simBlue, align=left, anchor=west, font=\scriptsize] at (3.0,-0.1)
    {\textbf{diagonal}: self-inductance\\$L_{ii}=2U_m^{(i)}/I_i^2>0$};
  \node[simGold, align=left, anchor=west, font=\scriptsize] at (3.0,-1.4)
    {\textbf{off-diagonal}: mutual\\$L_{ij}=\vA_j^{\!\top}\mat{K}\vA_i/(I_iI_j)$\\
     sign $=$ relative orientation};
  \node[simRust, align=left, anchor=west, font=\scriptsize] at (3.0,-2.4)
    {\textbf{symmetric}: $L_{ij}=L_{ji}$\\(reciprocity, free from the Gram)};
\end{tikzpicture}
\caption[The inductance matrix]{The inductance matrix $\mat{L}$, the magnetostatic
product. The diagonal entries (blue) are the self-inductances $L_{ii}=2U_m^{(i)}/
I_i^2$, always positive. The off-diagonal entries (gold) are the mutual inductances
$L_{ij}$, whose sign records whether loops $i$ and $j$'s fields reinforce
($+$) or oppose ($-$). The matrix is symmetric (dashed mirror): $L_{ij}=L_{ji}$, a
reciprocity that the Gram reduction's structural symmetry supplies at no cost---only
the upper triangle is computed (\cref{ch:reductions}). This is the same matrix
\emph{shape} as the capacitance matrix of \cref{ch:electrostatics}; the entries are
inductances because the family is vector potentials and the weight is
current-normalized.}
\label{fig:Lmatrix}
\end{figure}

\begin{remark}[Loop inductance, partial inductance, and what we measured]
We have computed the \emph{loop} inductance of a set of complete current loops: each
$\vA_i$ is the field of a full, closed, divergence-free current path, so $L_{ij}$ is
a well-defined, gauge-independent number. A related notion, \emph{partial}
inductance, ascribes inductance to a single open segment of wire (useful for
circuit extraction from layout); it is not gauge-independent in isolation and only
the loop sums it builds are physical. The magnetostatic driver as composed here
returns loop inductances directly. The partial-inductance bookkeeping is a
reduce-stage variation---a different weighting and pairing of the same solved
family---and we note it only to flag that ``inductance matrix'' can mean slightly
different things downstream; the field solve underneath is the one in this chapter.
\end{remark}

\section{Worked examples}

\subsection{A two-loop inductance, end to end}

We trace the magnetostatic pipeline through a two-loop problem, in deliberate
parallel with the two-terminal \emph{capacitance} computed in \cref{ch:reductions}
(\cref{wex:cap2x2}). Same machine, same numbers shape---inductances out instead of
capacitances, because of the three swaps.

\begin{workedexample}{A $2\times2$ inductance from a current-normalized Gram}{ind2x2}
\textbf{Setup.} Two current loops, driven one at a time at unit current
$I_1=I_2=1\,\si{\ampere}$. The magnetostatic driver assembles the curl--curl energy
operator $\mat{K}$ \emph{once} (swaps 1+2: curl--curl on $\Hcurl$), then solves it
twice (\code{solve\_family}), once per loop, for the vector-potential fields $\vA_1$
and $\vA_2$. We are handed the three operator-weighted pairings that the solve
produced:
\[
  \vA_1^{\!\top}\mat{K}\,\vA_1 = 8,\qquad
  \vA_2^{\!\top}\mat{K}\,\vA_2 = 6,\qquad
  \vA_2^{\!\top}\mat{K}\,\vA_1 = 2
\]
(in nanohenry-equivalent units; these are the magnetic-energy bilinear forms, twice
the stored energy in each configuration).

\textbf{The Gram reduction (swap 3).} The inductance is \code{gram\_reduce} with the
current-normalized weight $w(i,j)=1/(I_iI_j)$ (\cref{eq:ind-gram-local}). With unit
currents, $w\equiv 1$ \emph{numerically}---but the normalization is structurally
present, and would rescale every entry if the currents were not $1$. The verb
computes the upper triangle and mirrors:
\[
  L_{11} = \frac{\vA_1^{\!\top}\mat{K}\vA_1}{I_1^2} = \frac{8}{1} = 8,\qquad
  L_{22} = \frac{\vA_2^{\!\top}\mat{K}\vA_2}{I_2^2} = \frac{6}{1} = 6,
\]
\[
  L_{12} = \frac{\vA_2^{\!\top}\mat{K}\vA_1}{I_1 I_2} = \frac{2}{1} = 2,
  \qquad L_{21}=L_{12}=2 \ \text{(mirrored, not recomputed).}
\]
The inductance matrix is therefore
\[
  \mat{L} =
  \begin{pmatrix} L_{11} & L_{12}\\ L_{21} & L_{22} \end{pmatrix}
  = \begin{pmatrix} 8 & 2\\ 2 & 6 \end{pmatrix}\si{\nano\henry}.
\]

\textbf{Interpretation.} The diagonal entries are the \emph{self}-inductances:
$L_{11}=\SI{8}{\nano\henry}$ is twice the magnetic energy stored when loop $1$ alone
carries unit current; $L_{22}=\SI{6}{\nano\henry}$ likewise for loop $2$. Both are
positive, as self-inductance must be. The off-diagonal $L_{12}=\SI{2}{\nano\henry}$
is the \emph{mutual} inductance: the flux of loop $1$'s field linked by loop $2$ per
unit current. Here it is positive---the two loops' fields reinforce. Its
\emph{symmetry} $L_{12}=L_{21}$ is reciprocity, and we got it for free: the Gram
computes only the upper triangle and mirrors (\cref{ch:reductions}).

\textbf{The capacitance contrast.} Run \emph{exactly this arithmetic} for the
two-terminal capacitance of \cref{wex:cap2x2} and one thing changes: the weight is
$w=1$ outright (no currents to normalize), the family is potentials $V_i\in\Hone$
not vector potentials $\vA_i\in\Hcurl$, and the operator is the
$\varepsilon$-weighted Poisson stiffness not the $\nu$-weighted curl--curl
stiffness. The off-diagonal also flips \emph{sign} in the capacitance case (the
Maxwell capacitance matrix has negative off-diagonals; mutual inductance need not).
But the reduction \emph{verb} is the same, the symmetry is the same, the
upper-triangle map is the same. Two physical matrices, one Gram.
\end{workedexample}

\subsection{One change at a time: the electrostatic--magnetostatic ledger}

The second worked example is not a computation but an accounting: we lay the two
drivers side by side, line by line, and verify that exactly three lines differ. This
is the chapter's thesis stated as a table.

\begin{workedexample}{Capacitance versus inductance, one change at a time}{ledger}
We tabulate every stage of the two drivers and mark each line ``same'' or
``swap''. Reading down the table, the three swaps stand alone; every other line is
shared verbatim.

\begin{center}\small
\renewcommand{\arraystretch}{1.25}
\begin{tabular}{@{}p{0.20\linewidth}p{0.30\linewidth}p{0.30\linewidth}c@{}}
\toprule
\textbf{stage} & \textbf{electrostatic} & \textbf{magnetostatic} & \\
\midrule
governing PDE & $-\Div(\varepsilon\Grad V)=\rho$ & $\Curl(\nu\Curl\vA)=\vJ$
  & \textcolor{simRust}{\textbf{swap 1}} \\
field unknown & scalar potential $V$ & vector potential $\vA$
  & \textcolor{simRust}{\textbf{swap 1}} \\
post-process & $\vE=-\Grad V$ & $\vB=\Curl\vA$
  & \textcolor{simRust}{\textbf{swap 1}} \\
FE space & $\Hone$ nodal & $\Hcurl$ edge (N\'ed\'elec)
  & \textcolor{simRust}{\textbf{swap 2}} \\
energy operator $\mat{K}$ & $\int\varepsilon\,\Grad\phi_a\!\cdot\!\Grad\phi_b$
  & $\int\nu\,\Curl\boldsymbol\phi_a\!\cdot\!\Curl\boldsymbol\phi_b$
  & \textcolor{simRust}{(swap 2)} \\
\addlinespace
$\mat{K}$ definiteness & SPD (grounded) & SP\emph{S}D (gradient null space)
  & \textcolor{simGray}{setup} \\
solver & CG, fixed op, reused precond. & CG, fixed op, reused precond.
  & \textcolor{simGreen}{\textbf{same}} \\
assemble cadence & once, reused per excitation & once, reused per excitation
  & \textcolor{simGreen}{\textbf{same}} \\
solve cadence & one per terminal & one per source
  & \textcolor{simGreen}{\textbf{same}} \\
reduce verb & \code{gram\_reduce} & \code{gram\_reduce}
  & \textcolor{simGreen}{\textbf{same}} \\
Gram symmetry & $G_{ij}=G_{ji}$, upper triangle & $G_{ij}=G_{ji}$, upper triangle
  & \textcolor{simGreen}{\textbf{same}} \\
\addlinespace
Gram weight $w(i,j)$ & $1$ & $1/(I_iI_j)$
  & \textcolor{simRust}{\textbf{swap 3}} \\
product & capacitance $\mat{C}$ & inductance $\mat{L}$ & \\
\bottomrule
\end{tabular}
\end{center}

\textbf{The reading.} Three rust ``swap'' regions: the PDE-and-field block (swap 1,
which drags post-processing $\vE\to\vB$ along with it), the FE-space line (swap 2,
which determines the energy operator's integrand), and the Gram-weight line (swap
3). The definiteness line is marked grey ``setup'': it is a consequence of swap 2,
but it changes how $\mat{K}$ and $\vJ$ are \emph{prepared} (consistent source or AMS,
\cref{sec:solvecare}), not the solver itself---so the solver line is green
``same.'' Everything in green---the solver, both cadences, the reduce verb, the Gram
symmetry---is shared without modification. The magnetostatic driver is the
electrostatic driver, three swaps over.
\end{workedexample}

\summarysection

Magnetostatics computes an inductance matrix the way electrostatics computes a
capacitance matrix: \emph{the same way}. The driver is the composition
\code{gram\_reduce} $\circ$ \code{solve\_family} $\circ$ \code{fe\_assemble}, and it
differs from the electrostatic driver in exactly three places.

\begin{itemize}[nosep]
\item \textbf{Swap 1---the PDE and field.} Magnetostatics solves the curl--curl
  equation $\Curl(\nu\Curl\vA)=\vJ$ for a magnetic \emph{vector} potential $\vA$,
  not Poisson for a scalar $V$; it post-processes $\vB=\Curl\vA$. Both are
  second-order self-adjoint elliptic equations from adjacent rungs of the de Rham
  complex (\cref{ch:derham}).
\item \textbf{Swap 2---the FE space.} The vector potential lives in $\Hcurl$,
  realized by edge (N\'ed\'elec) elements, because $\vA$'s tangential component is
  continuous while its normal component is free. Nodal-vector elements would impose
  too much continuity and breed spurious modes (\cref{ch:bc}); edge elements sit in
  the exact discrete de Rham complex and avoid them. The curl--curl operator
  inherits a large gradient null space (the gauge fields), making it only positive
  semi-definite; the solve handles this through a consistent ($\Div\vJ=0$) source or
  AMS preconditioning (\cref{ch:smoothers})---a change of \emph{setup}, not of
  solver.
\item \textbf{Swap 3---the Gram weight.} Inductance is the same \code{gram\_reduce}
  as capacitance with the current-normalized weight $w=1/(I_iI_j)$ in place of
  $w=1$ (\cref{ch:reductions}). One verb, two physical quantities; the diagonal is
  twice the magnetic energy per unit current ($L_{ii}=2U_m/I^2$), the off-diagonal
  is the (signed, symmetric) mutual inductance.
\end{itemize}

The shared remainder---assemble once, solve per source with the same CG, reduce by a
symmetric Gram---is the deeper lesson: a simulator earns its generality by making
distinct physical analyses \emph{the same machine over different vocabulary}. The
next chapter (\cref{ch:eigenmodes}) reuses the very same curl--curl operator on the
very same $\Hcurl$ space, now as an eigenvalue problem, and the gradient null space
of \cref{sec:nullspace} returns as the subspace the eigensolver must project out.

\furtherreading

The curl--curl equation, the magnetic vector potential, and edge-element
discretization of magnetostatics are developed in Jin's computational
electromagnetics text~\citep{jin2014} and, with full finite-element-theoretic rigor
(N\'ed\'elec spaces, the discrete de Rham complex, and convergence), in
Monk~\citep{monk2003}. The physical content---vector potential, magnetic energy,
self- and mutual inductance---is the magnetostatics of any intermediate
electromagnetism course; Griffiths~\citep{griffiths2017} is the standard reference,
and connects the inductance matrix to the circuit parameters an engineer expects.
The auxiliary-space Maxwell (AMS) preconditioner for the curl--curl gradient null
space is due to Hiptmair~\citep{hiptmair1998}; \cref{ch:smoothers} treats it in the
multigrid context.

\exercisessection

\begin{enumerate}[label=\textbf{40.\arabic*}]

\item \textbf{The three swaps from memory.} Without looking back, write the
electrostatic driver as a composition \code{reduce} $\circ$ \code{solve\_family}
$\circ$ \code{fe\_assemble} and then derive the magnetostatic driver from it by
applying the three swaps. For each swap, name (i) the box it lives in and (ii) one
sentence of why the physics forces it.

\item \textbf{Curl--curl is self-adjoint.} Show that the bilinear form
$a(\vA,\vF)=\int_\dom\nu\,(\Curl\vA)\cdot(\Curl\vF)\dV$ is symmetric,
$a(\vA,\vF)=a(\vF,\vA)$, and hence that the assembled stiffness $\mat{K}$ of
\cref{eq:stiff-curlcurl} is a symmetric matrix. Why does this symmetry matter for
(a) the conjugate-gradient solver and (b) the Gram reduction's reciprocity
$L_{ij}=L_{ji}$?

\item \textbf{The null space is gradients.} Let $\psi_h$ be any discrete nodal field
and $\mat{G}$ the discrete gradient (\cref{ch:derham}, \cref{def:discrete-grad}).
Show $\mat{K}\,(\mat{G}\psi_h)=\mathbf 0$ directly from
\cref{eq:stiff-curlcurl}, using $\Curl\Grad=0$. Conclude that $\mat{K}$ is only
positive semi-definite, with null space $\supseteq\operatorname{range}(\mat{G})$.

\item \textbf{Why a consistent source rescues CG.} The curl--curl null space is the
range of the discrete gradient. Argue that a divergence-free source $\vJ$
($\Div\vJ=0$) is orthogonal to every gradient field, hence has no component along
the null space, hence conjugate gradients applied to $\mat{K}\vA=\vJ$ stays in the
positive-definite complement and converges. (Hint: integrate $\int\vJ\cdot\Grad\psi$
by parts.)

\item \textbf{The current normalization.} Suppose loop $1$ is driven at
$I_1=\SI{2}{\ampere}$ and loop $2$ at $I_2=\SI{3}{\ampere}$ (not unit currents), and
the solve returns $\vA_1^{\!\top}\mat{K}\vA_1=48$,
$\vA_2^{\!\top}\mat{K}\vA_2=54$, $\vA_2^{\!\top}\mat{K}\vA_1=18$ (energy-form units).
Compute the inductance matrix $\mat{L}$ using the weight $w(i,j)=1/(I_iI_j)$. Verify
$L_{12}=L_{21}$ and that the values are the same as if you had first rescaled each
field to unit current.

\item \textbf{Self-inductance is positive; mutual need not be.} Explain, from
$L_{ii}=\vA_i^{\!\top}\mat{K}\vA_i/I_i^2$ and the positive-semidefiniteness of
$\mat{K}$, why $L_{ii}\ge 0$ always. Then construct a physical scenario (two loops,
relative orientation) in which $L_{12}<0$, and contrast with the capacitance matrix,
whose off-diagonals are always negative (\cref{ch:reductions}, \cref{wex:cap2x2}).

\item \textbf{Nodal vector would break it.} Describe, in two or three sentences, what
goes wrong if one discretizes $\vA$ with three copies of nodal $\Hone$ elements
instead of $\Hcurl$ edge elements. Identify the property of the discrete de Rham
complex (\cref{ch:derham}) that fails and the symptom it produces in the spectrum
(\cref{ch:bc}).

\item \textbf{The eigenmode connection (forward to \cref{ch:eigenmodes}).} The
eigenmode driver assembles the \emph{same} curl--curl $\mat{K}$ on the \emph{same}
$\Hcurl$ space, but solves $\mat{K}\vx=\lambda\mat{M}\vx$ instead of
$\mat{K}\vA=\vJ$. Predict: which eigenvalues does the gradient null space of
\cref{sec:nullspace} produce, and why must the eigensolver project them out? (You
are reconstructing the spurious-zero-mode discussion of \cref{ch:bc} from the
null-space structure alone.)

\end{enumerate}

\chapter{Eigenmodes: Resonant Frequencies and Q}
\label{ch:eigenmodes}

\chapterorient{We have studied electrostatics (\cref{ch:electrostatics}) as the
anchor pipeline: \texttt{assemble} an operator once, run a family of independent
solves over it, and \texttt{reduce} the results to a physical product. This
chapter takes that same skeleton and swaps exactly one corner---the
\texttt{solve}. Instead of a linear solve we write ourselves, the eigenmode
analysis hands the whole assembled problem to an opaque eigenvalue library and
receives the answer. It is the fourth solve corner from \cref{ch:situating}: the
black box. We will see that almost everything else about the pipeline is
unchanged, that the one Palace-written loop is the harmless post-processing
readout, and that the price of the black box is a subtlety---spurious
non-physical modes---that we must handle \emph{before} the solve, not after.}

\section{The physical question: how does a cavity ring?}

Strike a wine glass and it sings at a definite pitch. Pluck a guitar string and
it vibrates at a fundamental frequency plus a ladder of overtones. A closed
electromagnetic cavity---a hollow metal box, a superconducting resonator, a
dielectric puck---does exactly the same thing with electromagnetic fields. With
\emph{no external drive at all}, the cavity supports a discrete set of standing
electromagnetic oscillations, its \emph{resonant modes}. Each mode rings at its
own \emph{resonant frequency} $f$, decays at its own rate (captured by a
\emph{quality factor} $Q$), and has its own spatial pattern of $\vE$ and $\vB$
(\cref{fig:cavity-setup}). The eigenmode analysis computes the lowest few of
these---typically the handful of modes nearest some target frequency---because
those are the ones that matter for a filter passband, a qubit's readout
resonator, or an accelerator cavity.

\begin{physicsbox}
A \emph{closed} cavity is a region of space bounded (mostly) by conducting
walls, with no source inside and no port feeding energy in. Maxwell's equations
on such a domain are \emph{homogeneous}: there is no right-hand side. A
homogeneous linear system has only the trivial solution $\vE=\mathbf 0$
\emph{except} at special frequencies, where a nonzero field can sustain itself.
Those special frequencies are the resonant frequencies; the self-sustaining
fields are the modes. Physically, the cavity is a trap: a mode is a field
pattern that bounces between the walls and reinforces itself, neither growing nor
(in the ideal lossless case) decaying. Finding the modes is therefore not a
matter of \emph{driving} the cavity and seeing what comes out---it is a matter of
asking \emph{at what frequencies can the cavity oscillate on its own?}
\end{physicsbox}

\begin{figure}[t]
\centering
\begin{tikzpicture}[font=\footnotesize]
  % cavity walls
  \draw[thick, simInk, fill=simBgGray] (0,0) rectangle (6.2,3.0);
  \node[anchor=south west, font=\scriptsize, text=simGray] at (0.05,3.02)
    {conducting walls (PEC): $\vn\times\vE=\mathbf 0$};
  % standing-wave E field pattern (vertical arrows, sin envelope)
  \foreach \x in {0.4,0.8,...,5.8}{
    \pgfmathsetmacro{\amp}{1.05*sin(deg(pi*\x/6.2))}
    \draw[->, simBlue, thick] (\x,1.5) -- (\x,{1.5+\amp});
    \draw[->, simBlue, thick] (\x,1.5) -- (\x,{1.5-\amp});
  }
  % envelope
  \draw[densely dashed, simRust, semithick, domain=0:6.2, samples=80, smooth]
    plot (\x,{1.5+1.15*sin(deg(pi*\x/6.2))});
  \draw[densely dashed, simRust, semithick, domain=0:6.2, samples=80, smooth]
    plot (\x,{1.5-1.15*sin(deg(pi*\x/6.2))});
  \node[text=simBlue] at (3.1,2.85) {$\vE$};
  \node[text=simRust, anchor=west, font=\scriptsize] at (6.25,2.4)
    {standing-wave envelope};
  % half-wavelength bracket
  \draw[<->, simGreen, thick] (0,0.25) -- (6.2,0.25)
    node[midway, fill=simBgGray, inner sep=1pt, text=simGreen]
      {$L=\tfrac{1}{2}\lambda$ (fundamental)};
\end{tikzpicture}
\caption[A cavity's lowest mode]{The fundamental mode of a closed cavity. With no
external drive, the field cannot escape the conducting walls (the perfect-electric-conductor
boundary forces the tangential $\vE$ to vanish there, \cref{ch:bc}), so it forms a
standing wave: an electric-field pattern (blue) whose amplitude follows a fixed
spatial envelope (dashed) and oscillates in time at the resonant frequency. The
fundamental fits exactly half a wavelength across the cavity, $L=\tfrac12\lambda$;
higher modes fit more half-wavelengths and ring at higher frequencies. The
eigenmode analysis computes these patterns and their frequencies directly, from
the geometry alone.}
\label{fig:cavity-setup}
\end{figure}

\subsection{The same skeleton, one corner swapped}

In \cref{ch:situating} we reduced every Palace analysis to one skeleton,
\[
  \texttt{config} \to \texttt{assemble} \to \texttt{solve}
  \to \texttt{reduce} \to \texttt{product},
\]
and we placed electrostatics (\cref{ch:electrostatics}) at its center as the
universal anchor. Eigenmode analysis is the \emph{same machine} with a single
part swapped---but the part it swaps is the most consequential one. In
electrostatics the \texttt{solve} stage is a \emph{fixed-operator map}: assemble
the stiffness $\Kop$ once, then call our own linear solver $\Kop\vx_i=\vb_i$ once
per conductor, in a loop we write and control (\code{solve\_family}). In the
eigenmode analysis the \texttt{solve} stage is not a linear solve at all, and
there is no loop we write: we assemble the operator pencil, hand it whole to an
eigenvalue library, and get the converged modes back from a single call.

\begin{keyidea}
The eigenmode pipeline is the anchor pipeline with the solve corner replaced.
Electrostatics solves $\Kop\vx_i=\vb_i$ in a Palace-written loop; the eigenmode
analysis solves the \emph{eigenproblem} $\Kop\vx=\lambda\Mop\vx$, and the entire
iteration that does so lives \emph{inside an opaque library}. Palace writes no
solve loop here at all. \emph{The only loop Palace writes in the eigenmode driver
is the post-processing readout map over the modes the library already
converged.} This is the fourth solve corner of \cref{ch:situating}: the opaque
black box.
\end{keyidea}

The rest of this chapter walks the three stages in order---assemble, the
black-box solve, reduce---then handles the one subtlety the black box forces on
us (spurious modes), and closes by reading out the mode fields and naming a few
applications.

\section{Assemble: the eigenproblem pencil, built once}
\label{sec:eig-assemble}

The assemble stage is almost unchanged from the anchor. We build operators by
folding weak-form terms over the finite-element space (\cref{ch:fem}); the only
difference is \emph{which} operators, and that there is more than one.

The fields live in the Nédélec edge-element space $\Hcurl$ (\cref{ch:functionspaces}),
the natural home for $\vE$: it builds tangential continuity across element faces
into the basis, exactly matching the physics of an electric field at a material
interface. On that space we assemble up to three operators:
\begin{itemize}[nosep]
  \item the \textbf{curl--curl stiffness} $\Kop$, the discretization of
  $\Curl(\nu\,\Curl\,\cdot\,)$ where $\nu=1/\mu$ is the reluctivity---this is the
  ``spring constant'' of the electromagnetic field, the operator that resists
  curvature;
  \item the \textbf{mass} $\Mop$, the discretization of $\varepsilon\,(\cdot)$
  weighted by permittivity---the ``inertia'' of the field;
  \item optionally the \textbf{damping} $\Cop$, present only when the cavity is
  \emph{lossy} (finite wall conductivity, lossy dielectric, or a resistive
  port)---the operator that drains energy.
\end{itemize}
All three come from the same assemble fold we have used since \cref{ch:fem}; in
the source they are three calls,
\lstinline[style=l4style]|GetStiffnessMatrix|,
\lstinline[style=l4style]|GetDampingMatrix|, and
\lstinline[style=l4style]|GetMassMatrix|, each run \emph{once}.

\subsection{Lossless: the generalized eigenproblem}

When the cavity is lossless, $\Cop$ is absent and Maxwell's source-free
time-harmonic equation, discretized, becomes a \emph{generalized eigenvalue
problem} (\cref{ch:reductions-pde}):
\begin{equation}
  \Kop\,\vx \;=\; \lambda\,\Mop\,\vx,
  \qquad \lambda = \omega^2,
  \label{eq:gevp}
\end{equation}
where $\vx$ is the vector of degrees of freedom of the mode field $\vE$, the
eigenvalue $\lambda$ is the squared angular frequency $\omega^2$, and the
eigenvector $\vx$ \emph{is} the discretized mode. This is exactly the generalized
EVP we met abstractly in \cref{ch:reductions-pde} and studied numerically in
\cref{ch:eigenvalues}; here it arrives with a concrete physical reading. Each
eigenpair $(\lambda,\vx)$ is one resonance: $\omega=\sqrt\lambda$ is its angular
frequency, $\vx$ is its field pattern.

\begin{notationbox}
The eigenvalue is the \emph{square} of the angular frequency, $\lambda=\omega^2$,
not the frequency itself. This is the single most important bookkeeping fact in
the chapter. It means the solver returns $\lambda$, and the readout must
\emph{un-transform} it---$\omega=\sqrt\lambda$---to get a frequency, then divide
by $2\pi$ to get the ordinary frequency $f=\omega/2\pi$ in hertz. Forgetting the
square root, or the $2\pi$, is the classic first-run-off-by-a-lot bug. We meet
the un-transform again in \cref{sec:eig-reduce}.
\end{notationbox}

\subsection{Lossy: the quadratic eigenproblem}

When the cavity loses energy, the damping operator $\Cop$ enters and the problem
becomes \emph{quadratic} in the eigenvalue (\cref{ch:reductions-pde}):
\begin{equation}
  \bigl(\Kop + \lambda\,\Cop + \lambda^2\,\Mop\bigr)\,\vx \;=\; \mathbf 0,
  \qquad \lambda = i\,\omega.
  \label{eq:qevp}
\end{equation}
Now $\lambda$ is not $\omega^2$ but $i\omega$, so a lossy mode's eigenvalue is
\emph{complex}: its imaginary part carries the frequency and its real part the
decay rate. The quadratic form is what lets a mode decay---the term linear in
$\lambda$ couples frequency to loss. We do not solve \cref{eq:qevp} by hand;
\cref{ch:eigenvalues} showed how a quadratic eigenproblem is linearized into a
larger ordinary one that the same library swallows. For our purposes the
quadratic case differs from the linear case only in (i) which operators the
pencil carries, and (ii) how the eigenvalue un-transforms to $\omega$. Both
differences are absorbed into the single \emph{problem-type} parameter the solve
takes; the shape of the pipeline does not change.

\begin{definition}[The operator pencil]
The \emph{pencil} is the bundle of assembled operators handed to the eigensolver:
the pair $(\Kop,\Mop)$ for the linear generalized eigenproblem \cref{eq:gevp}, or
the triple $(\Kop,\Cop,\Mop)$ for the quadratic eigenproblem \cref{eq:qevp},
together with a spectral-transform target $\sigma$ (\cref{ch:eigenvalues}) telling
the solver \emph{which} part of the spectrum we want. The pencil is assembled
exactly once; nothing in the solve mutates it.
\end{definition}

\begin{keyidea}
Assembly produces an operator \emph{pencil}, not a single matrix: $(\Kop,\Mop)$
for a lossless cavity (the generalized EVP $\Kop\vx=\lambda\Mop\vx$,
$\lambda=\omega^2$) or $(\Kop,\Cop,\Mop)$ for a lossy one (the quadratic EVP
$(\Kop+\lambda\Cop+\lambda^2\Mop)\vx=\mathbf 0$, $\lambda=i\omega$). It is built
\emph{once}, exactly like the single $\Kop$ of electrostatics---the assemble
corner is shared; only the count and the meaning of the operators differ.
\end{keyidea}

\section{Solve: the opaque black box}
\label{sec:eig-solve}

Here is where the pipeline diverges sharply from the anchor. We have the pencil.
What do we do with it?

In electrostatics we would now write a loop: for each right-hand side, call our
own conjugate-gradient or direct solver, collect the answer. We control that
loop; we can see every linear solve it performs; the iteration is
\emph{ours}. The eigenmode analysis does nothing of the kind. We make a single
call,
\[
  \texttt{eigs} \;=\; \texttt{eigsolve}\;\bigl(\text{pencil},\ \sigma,\
  n_{\text{modes}}\bigr),
\]
handing the entire assembled problem---operators, target, requested mode
count---to a specialized eigenvalue library (SLEPc or ARPACK in Palace's case,
\cref{ch:eigenvalues}). The library runs its own sophisticated iteration, a
Krylov--Schur or Arnoldi process with a shift-and-invert spectral transform
(\cref{ch:krylovschur}), and returns the converged eigenpairs. We do not write
that iteration. We do not see it. We see one call and its result
(\cref{fig:eig-pipeline}).

\begin{figure}[t]
\centering
\begin{tikzpicture}[node distance=7mm and 10mm, font=\small]
  \node[data] (cfg) {config};
  \node[stage, right=of cfg, align=center] (asm)
    {assemble\\pencil};
  \node[extern, right=of asm, minimum width=22mm, minimum height=11mm,
        align=center] (solve) {\textbf{eigsolve}\\\scriptsize SLEPc / ARPACK};
  \node[stage, right=of solve, align=center] (map) {map\\readout};
  \node[product, right=of map, align=center] (prod) {$(f,Q)$\\table};
  \draw[flow] (cfg) -- (asm);
  \draw[flow] (asm) -- node[above, font=\scriptsize]{$(\Kop,\Cop,\Mop),\sigma$} (solve);
  \draw[flow] (solve) -- node[above, font=\scriptsize]{eigenpairs} (map);
  \draw[flow] (map) -- (prod);
  % the library's hidden inner loop: callback for operator apply
  \draw[refedge] (solve.south) ++(0,-1mm)
    .. controls +(0,-7mm) and +(0,-7mm) .. ($(solve.south)+(0,-1mm)$)
    node[midway, below, font=\scriptsize, text=simViolet, align=center]
      {hidden inner loop:\\Krylov--Schur + shift-invert\\(calls back for $\mathbf{A}\vx$)};
  % the one Palace-written loop annotation
  \node[font=\scriptsize, text=simGreen, align=center, below=2mm of map]
    {the \emph{only} loop\\Palace writes};
\end{tikzpicture}
\caption[The eigenmode pipeline]{The eigenmode pipeline, drawn to contrast with
the electrostatic anchor (\cref{ch:electrostatics}). Assembly builds the operator
pencil once. The \texttt{eigsolve} box is dashed and violet---the convention of
\cref{ch:situating} for a corner where \emph{we write no loop}: the entire
eigen-iteration (Krylov--Schur with shift-and-invert, \cref{ch:krylovschur})
happens inside the library, which only calls back out to us for the operator
application $\mathbf A\vx$. The single Palace-written loop is the post-processing
\texttt{map readout} (green) that turns each converged eigenpair into a row of the
$(f,Q)$ table. Compare \cref{ch:electrostatics}, where the solve box is a solid
Palace-written \code{solve\_family} loop instead of a black box.}
\label{fig:eig-pipeline}
\end{figure}

\subsection{The black box still calls back}

``Opaque'' does not mean ``self-contained.'' The library does not know what
$\Kop$, $\Cop$, or $\Mop$ \emph{are}; it knows only that it can ask us to apply
them. The shift-and-invert spectral transform (\cref{ch:eigenvalues}) needs to
apply the operator $(\Kop-\sigma\Mop)^{-1}\Mop$ to a vector on every inner step,
and \emph{that inverse is itself a linear solve}. So the library, mid-iteration,
calls back into our code asking us to apply the operator---and to apply that
inverse we run one of our own linear solves (the very \code{ksp\_solve}
machinery from electrostatics, \cref{ch:electrostatics}). The picture is a black
box with a wire poking out: the library owns the outer eigen-iteration, but the
operator applications---including an inner linear solve buried inside the
shift-and-invert---are still ours. We see those applications; we do not see how
the library schedules them or decides it has converged.

\begin{keyidea}
The eigensolver is opaque but not isolated: it runs the outer eigen-iteration
itself, and \emph{calls back} into our code for each operator application
$\mathbf A\vx$. Because the spectral transform applies $(\Kop-\sigma\Mop)^{-1}$,
each callback contains an \emph{inner linear solve}---the same linear-solve
machinery electrostatics uses, now nested one level down, inside someone else's
loop. The opacity is precise: we own the operator applies, the library owns the
iteration and the convergence test.

\end{keyidea}

\subsection{Why hand it off at all?}

It is fair to ask why we surrender control here when we kept it for every other
solve corner. The answer is honest engineering, not laziness. A robust
eigenvalue iteration that finds a few interior eigenpairs of a large
non-symmetric (in the lossy case) pencil, restarts to control memory, locks
converged pairs, and deflates against them is a substantial piece of numerical
software---Krylov--Schur with all its restart and purging logic
(\cref{ch:krylovschur}). Palace's authors made the deliberate choice to call a
mature library (SLEPc, itself wrapping PETSc; or ARPACK) rather than rewrite it.
Our framework records that choice faithfully: the \code{eigsolve} combinator is
\emph{deliberately} opaque. It names the library boundary, states the contract
across it (in: a pencil and a target; out: converged eigenpairs and a count), and
stops.

\begin{pitfall}
The opacity of \code{eigsolve} is a property of \emph{Palace's} code, not of the
mathematics. It is tempting to read ``black box'' as ``we do not understand the
eigen-iteration.'' We do---\cref{ch:eigenvalues} and \cref{ch:krylovschur} build
Krylov--Schur from scratch. The black box marks where \emph{Palace} stops writing
its own loop and calls out, not where understanding stops. A reader who wants to
see inside the box opens those two chapters, which reconstruct exactly what the
library does, in our own vocabulary. The honest record is: \emph{Palace} treats
it as a black box; the book does not.
\end{pitfall}

\subsection{Contrast with the anchor's solve}

The contrast with electrostatics is worth drawing as a side-by-side
(\cref{fig:solve-vs-map}), because it is the whole point of the chapter. In
electrostatics the solve stage is a loop \emph{we} write, the
\code{solve\_family} fixed-operator map: visible, controllable, one linear solve
per right-hand side. In the eigenmode analysis the solve stage is a single
library call, and the only loop we write comes \emph{after} the solve---the
readout map over the modes the library has already produced. The difference is
not which equation we solve; it is who owns the iteration.

\begin{figure}[t]
\centering
\begin{tikzpicture}[font=\footnotesize]
  % --- left: electrostatic solve (a loop we write) ---
  \begin{scope}[shift={(0,0)}]
    \node[font=\small\bfseries, align=center] at (1.6,2.55)
      {Electrostatic: a loop \emph{we} write};
    \node[opbox, minimum width=8mm] (k) at (0.1,1.2) {$\Kop$};
    \foreach \i/\y in {1/1.85,2/1.2,3/0.55}{
      \node[product, minimum width=5mm, minimum height=4mm, scale=0.7]
        (es\i) at (2.1,\y) {};
      \draw[flow, shorten >=1pt] (k.east) -- (es\i.west);}
    \node[font=\scriptsize\ttfamily, text=simGreen] at (1.6,-0.05)
      {solve\_family (corner 1)};
    \node[font=\scriptsize\itshape, align=center] at (1.6,-0.55)
      {one linear solve per RHS,\\every solve visible};
  \end{scope}
  % divider
  \draw[simGray, densely dashed] (4.0,-0.9) -- (4.0,2.8);
  % --- right: eigenmode solve (a black box) ---
  \begin{scope}[shift={(4.6,0)}]
    \node[font=\small\bfseries, align=center] at (1.6,2.55)
      {Eigenmode: a black box};
    \node[opbox, minimum width=8mm, align=center] (pen) at (0.0,1.2)
      {$(\Kop,$\\$\Cop,\Mop)$};
    \node[extern, minimum width=13mm, minimum height=9mm] (bb) at (1.9,1.2)
      {\scriptsize eigsolve};
    \draw[flow, shorten >=1pt] (pen.east) -- (bb.west);
    \node[product, minimum width=8mm, minimum height=5mm, scale=0.75]
      (out) at (3.35,1.2) {\scriptsize eigs};
    \draw[flow, shorten >=1pt] (bb.east) -- (out.west);
    \node[font=\scriptsize\ttfamily, text=simViolet] at (1.6,-0.05)
      {eigsolve (corner 4)};
    \node[font=\scriptsize\itshape, align=center] at (1.6,-0.55)
      {one call, iteration hidden;\\our loop is the readout, after};
  \end{scope}
\end{tikzpicture}
\caption[Solve corner vs.\ map corner]{The two solve stages side by side. Left:
electrostatics runs a \code{solve\_family} map---a Palace-written loop of
independent linear solves over the fixed operator $\Kop$, every solve visible
(\cref{ch:electrostatics}). Right: the eigenmode analysis runs \code{eigsolve}---a
single call into an opaque library (dashed violet), whose internal iteration we do
not see. The eigenmode driver's \emph{own} loop is not in the solve at all; it is
the post-processing readout map that follows. This is the structural heart of the
chapter: same skeleton, the solve corner swapped from a loop-we-write to a
loop-we-do-not.}
\label{fig:solve-vs-map}
\end{figure}

\section{Keeping the modes physical: the gradient null-space}
\label{sec:eig-spurious}

There is a subtlety that the black box forces us to confront, and it is not
optional. The curl--curl stiffness operator $\Kop$ has a large \emph{null
space}: for \emph{any} scalar potential $\phi$, the gradient field $\Grad\phi$
satisfies $\Curl\Grad\phi=\mathbf 0$ identically, so $\Kop(\Grad\phi)=\mathbf 0$.
On the discrete space this means $\Kop$ has a whole subspace of eigenvectors with
eigenvalue \emph{zero}---one for every discrete gradient. These are not
resonances. They are an artifact of the curl--curl operator's structure, the
discrete shadow of the Helmholtz decomposition (\cref{ch:bc}): every $\Hcurl$
field splits into a curl part (the physics) and a gradient part (the null space),
and only the curl part carries a real mode.

If we hand the raw pencil to the eigensolver and ask for the modes nearest zero
frequency, it will dutifully return a cloud of these spurious zero and
near-zero ``modes''---numerically polluted gradients---and bury the real
low-frequency resonances among them (\cref{fig:spurious}). The black box cannot
tell physics from artifact; it returns eigenpairs, and these are genuine
eigenpairs of the matrix we gave it. The fix has to happen on \emph{our} side,
\emph{before} the solve.

\begin{figure}[t]
\centering
\begin{tikzpicture}[font=\footnotesize]
  % --- before: spectrum with a spurious cloud near 0 ---
  \begin{scope}[shift={(0,0)}]
    \node[font=\small\bfseries] at (3.0,1.9) {Before projection};
    \draw[->, thick, simInk] (0,0) -- (6.2,0) node[right]{$\omega$};
    \draw[thick, simInk] (0,-0.1) -- (0,0.1) node[above, font=\scriptsize]{$0$};
    % spurious cloud near 0
    \foreach \x in {0.05,0.12,0.2,0.28,0.35,0.45,0.55,0.62,0.7,0.8}{
      \fill[simRust] (\x,0) circle (1.6pt);}
    \node[text=simRust, font=\scriptsize, align=center] at (0.5,0.75)
      {spurious gradient\\null-space modes};
    \draw[simRust, ->] (0.5,0.5) -- (0.4,0.12);
    % real modes
    \foreach \x in {2.4,3.6,4.9}{
      \fill[simBlue] (\x,0) circle (2.4pt);}
    \node[text=simBlue, font=\scriptsize] at (3.6,0.55) {real resonances};
    \node[text=simRust, font=\scriptsize, align=center] at (3.0,-0.7)
      {the real modes are \emph{buried} in the cloud};
  \end{scope}
  % --- after: clean spectrum ---
  \begin{scope}[shift={(0,-2.6)}]
    \node[font=\small\bfseries] at (3.0,1.9) {After divergence-free projection};
    \draw[->, thick, simInk] (0,0) -- (6.2,0) node[right]{$\omega$};
    \draw[thick, simInk] (0,-0.1) -- (0,0.1) node[above, font=\scriptsize]{$0$};
    % only real modes
    \foreach \x in {2.4,3.6,4.9}{
      \fill[simBlue] (\x,0) circle (2.4pt);}
    \node[text=simBlue, font=\scriptsize] at (3.6,0.55) {real resonances};
    \node[text=simGreen, font=\scriptsize, align=center] at (0.5,0.7)
      {null space\\removed};
    \draw[simGreen, ->] (0.5,0.42) -- (0.18,0.1);
    \node[text=simGreen, font=\scriptsize, align=center] at (3.0,-0.7)
      {the solver now returns physical modes first};
  \end{scope}
\end{tikzpicture}
\caption[Spurious modes and the projection]{The gradient null-space and its
removal. Top: the curl--curl operator $\Kop$ annihilates every gradient field
($\Kop\Grad\phi=\mathbf 0$), so its raw spectrum carries a cloud of spurious
zero/near-zero ``modes'' (rust) that are numerical gradients, not resonances; the
real low-frequency resonances (blue) are buried among them, and the solver,
asked for the modes nearest zero, returns the artifacts first. Bottom: the
\emph{divergence-free projection} (\cref{ch:bc}) restricts the solve to the
gradient-free subspace, deleting the null-space cloud so the eigensolver returns
the physical resonances directly. The projection is wired in \emph{before} the
solve, because the black box cannot tell artifact from physics on its own.}
\label{fig:spurious}
\end{figure}

The fix is the \emph{divergence-free projection} of \cref{ch:bc}. We restrict the
solve to the subspace of fields orthogonal to all gradients---the divergence-free
(curl) part of the Helmholtz decomposition---so the null-space directions are
simply unavailable to the eigensolver. Concretely, Palace builds a projector $P$
onto the divergence-free subspace and wires it into the eigensolver: the starting
vector is projected, and the projector is applied so the iteration never strays
into the gradient subspace. With the null space projected out, the spectrum the
library sees has no spurious cloud, and the modes it returns are physical
resonances.

\begin{keyidea}
The curl--curl operator's gradient null space ($\Kop\Grad\phi=\mathbf 0$ for
every scalar $\phi$, \cref{ch:bc}) produces a cloud of spurious zero-frequency
``modes'' that are not resonances. The \emph{divergence-free projection} restricts
the solve to the gradient-free subspace, removing them, so the black box returns
physical modes. This projection is wired in \emph{before} the opaque solve---it
must be, because the black box cannot distinguish artifact from physics once it
is running. Handling the null space is the one place where the eigenmode pipeline
demands real physics knowledge the anchor never needed.
\end{keyidea}

\section{Reduce: the eigenfrequency-and-Q table}
\label{sec:eig-reduce}

The library returns converged eigenpairs $(\lambda_m,\vx_m)$. The reduce stage
turns them into the physical product the user asked for: a small \emph{table},
one row per mode (\cref{ch:reductions}, \cref{ch:situating})---the rank-1
``Shape~3'' reduction. Each row carries the mode's resonant frequency $f_m$ and
its quality factor $Q_m$. This stage \emph{is} the one Palace-written loop in the
driver: a \code{map} over the converged modes, with no coupling between rows
(mode~3's row does not depend on mode~2's), which is exactly what makes it a map
and not a fold.

\subsection{From eigenvalue to frequency: the un-transform}

The first column of the table is the resonant frequency, and getting it right is
pure bookkeeping---the un-transform of the eigenvalue back to a physical
frequency (\cref{ch:eigenvalues}). Two things must be undone. First, the
\emph{problem-type} transform from \cref{sec:eig-assemble}: for the linear EVP
the eigenvalue is $\lambda=\omega^2$, so $\omega=\sqrt\lambda$; for the quadratic
EVP it is $\lambda=i\omega$, so $\omega=\lambda/i$. Second, any \emph{spectral
shift} $\sigma$ used to target the interior of the spectrum
(\cref{ch:eigenvalues}) must be inverted---the solver works in shifted
coordinates, and the physical eigenvalue is recovered by undoing the shift,
exactly as in the shift-and-invert worked example of \cref{ch:eigenvalues}. With
$\omega$ in hand, the eigenfrequency is its real part over $2\pi$:
\begin{equation}
  f_m \;=\; \frac{\Re\,\omega_m}{2\pi}.
  \label{eq:freadout}
\end{equation}

\subsection{From the complex eigenvalue to Q}

The second column is the quality factor, and it lives entirely in the imaginary
part. A lossless cavity's $\omega$ is purely real and the mode rings forever; a
lossy cavity's $\omega$ has a positive imaginary part, the decay rate, and the
mode rings down. The quality factor measures how many radians of oscillation the
mode completes per e-folding of its amplitude. In the energy/loss convention
(\cref{ch:reductions}),
\begin{equation}
  Q_m \;=\; \frac{\omega_m}{\kappa_m}
        \;=\; \frac{\Re\,\omega_m}{2\,\Im\,\omega_m},
  \label{eq:qreadout}
\end{equation}
where $\kappa_m$ is the mode's loss rate. The factor of $2$ converts the
amplitude decay rate ($\Im\,\omega$) to the energy decay rate (energy is
amplitude squared, so it decays twice as fast); different texts place the $2$
differently, which is precisely why the reduction takes $\kappa$ as a parameter
rather than hard-coding the convention.

\begin{notationbox}
The loss rate $\kappa_m$ is not always read straight off $\Im\,\omega$. When the
loss enters through a specific physical channel---a resistive lumped port, say---Palace
computes $\kappa_m$ as a \emph{participation ratio}: the fraction of the mode's
energy dissipated in that channel, $\kappa_m=\tfrac12 R\,\abs{I_m}^2/E_m$ (resistor
self-energy over total mode energy), which then enters \cref{eq:qreadout} as
$Q_m=\omega_m/\kappa_m$. The two readings agree---$\Im\,\omega$ \emph{is} the
total loss rate---but the participation form attributes the loss to its source,
which is what a designer wants to know. Either way $Q$ is a scalar ratio, one
number per mode.
\end{notationbox}

\subsection{The lossless limit}

The honest edge case is the lossless cavity: $\Im\,\omega\to 0$, $\kappa\to 0$,
and \cref{eq:qreadout} divides by zero. The physical answer is $Q\to\infty$: an
ideal lossless resonator rings forever. The reduction does not produce a
\texttt{NaN}; it tests $\kappa=0$ and returns $Q=\infty$ directly, the one piece
of control flow in an otherwise purely arithmetic readout. The infinite-$Q$ case
is real, not exceptional: it is the limit every low-loss design---a
superconducting cavity, a high-$Q$ dielectric resonator---approaches, and the
table reports it faithfully.

\begin{keyidea}
The reduce stage is a \code{map} over the converged eigenpairs producing a
rank-1 table, one $(f,Q)$ row per mode. The frequency is the eigenvalue
\emph{un-transformed} ($\omega=\sqrt\lambda$ linear, $\lambda/i$ quadratic; shift
undone) and read as $f=\Re\,\omega/2\pi$. The quality factor is the
energy/loss ratio $Q=\Re\,\omega/(2\,\Im\,\omega)=\omega/\kappa$, with the
$\kappa\to0$ lossless limit handled as $Q=\infty$ rather than a division error.
This is the \emph{only} loop the eigenmode driver writes, and it is a pure
post-processing map---never a solve.
\end{keyidea}

\section{Worked examples}

\begin{workedexample}{The modes of a 1-D cavity, by hand}{cavity1d}
We compute the lowest resonances of the simplest cavity: a one-dimensional
segment of length $L$ with conducting ends, the electromagnetic analog of a
string fixed at both ends. The continuous problem is
\[
  -u''(x) = \lambda\, u(x), \qquad u(0)=u(L)=0, \qquad \lambda = \omega^2/c^2,
\]
whose exact modes are $u_n(x)=\sin(n\pi x/L)$ with $\lambda_n=(n\pi/L)^2$, giving
the familiar resonant-frequency ladder
\[
  f_n \;=\; \frac{c}{2\pi}\sqrt{\lambda_n} \;=\; \frac{n\,c}{2L},
  \qquad n=1,2,3,\dots
\]
The fundamental fits half a wavelength across the cavity ($f_1=c/2L$,
\cref{fig:cavity-setup}); each higher mode adds one more half-wavelength. We will
recover this ladder as the eigenvalues of a \emph{discrete} pencil
$\Kop\vx=\lambda\Mop\vx$, exactly the matrices the eigensolver would receive.

\emph{Discretize.} Put $N$ interior nodes at spacing $h=L/(N+1)$ and use the
standard second-difference stencil for $-u''$ and the (lumped) identity for the
mass. With $N=3$ ($h=L/4$), the stiffness and mass matrices are
\[
  \Kop = \frac{1}{h^2}\!\begin{pmatrix} 2 & -1 & 0\\ -1 & 2 & -1\\ 0 & -1 & 2\end{pmatrix},
  \qquad
  \Mop = \mathsf{I}_3,
\]
so the pencil reduces to the ordinary eigenproblem $\Kop\vx=\lambda\vx$ for the
tridiagonal $\Kop$.

\emph{Solve the $3\times3$ eigenproblem.} The eigenvalues of the symmetric
tridiagonal $\begin{psmallmatrix}2&-1&0\\-1&2&-1\\0&-1&2\end{psmallmatrix}$ are
known in closed form, $2-2\cos\!\bigl(\tfrac{k\pi}{4}\bigr)$ for $k=1,2,3$:
\[
  2-\sqrt2 \approx 0.586,\qquad 2,\qquad 2+\sqrt2 \approx 3.414.
\]
Dividing by $h^2=(L/4)^2=L^2/16$ gives the discrete $\lambda_k$:
\[
  \lambda_1 \approx \frac{9.37}{L^2},\quad
  \lambda_2 = \frac{32}{L^2},\quad
  \lambda_3 \approx \frac{54.6}{L^2}.
\]
\emph{Compare to the exact ladder.} The exact values are
$\lambda_n=(n\pi/L)^2 = \{9.87, 39.5, 88.8\}/L^2$. The lowest discrete eigenvalue,
$9.37/L^2$, is within $5\%$ of the exact $9.87/L^2$---a coarse three-node mesh
already pins the fundamental. The higher discrete eigenvalues sag below the exact
ones (the familiar dispersion error of a coarse mesh: short-wavelength modes are
underresolved), and they would tighten toward the ladder as $N$ grows. The
\emph{eigenvectors} of the $3\times3$ matrix are
$\vx_k=(\sin\tfrac{k\pi}{4},\sin\tfrac{2k\pi}{4},\sin\tfrac{3k\pi}{4})$---the
samples of $\sin(k\pi x/L)$ at the interior nodes, the discrete standing waves.

\emph{Read out the frequency.} Take the fundamental. Un-transform
$\lambda_1$ to a frequency: $\omega_1=c\sqrt{\lambda_1}=c\sqrt{9.37}/L=3.06\,c/L$,
so $f_1=\omega_1/2\pi=0.487\,c/L$, against the exact $f_1=c/2L=0.500\,c/L$. The
discrete eigenvalue, un-transformed exactly as \cref{sec:eig-reduce} prescribes,
lands within $3\%$ of the physical resonant frequency---and \emph{this little
$3\times3$ problem is, in miniature, precisely what the black box does at scale}:
build $(\Kop,\Mop)$, find $\lambda$, un-transform to $f$. The only difference at
full size is the matrix dimension and the fact that we hand the eigenproblem to a
library instead of reading the eigenvalues off a known formula.
\end{workedexample}

\begin{workedexample}{The Q of a lossy mode, and the lossless limit}{qfactor}
A lossy eigenmode solve returns a converged eigenpair whose recovered angular
frequency, after the un-transform of \cref{sec:eig-reduce}, is the complex number
\[
  \omega_m \;=\; \omega_r + i\,\omega_i
            \;=\; 2\pi\,(\SI{6.000}{\giga\hertz}) \;+\; i\,\bigl(2\pi\cdot\SI{1.5}{\mega\hertz}\bigr).
\]
The real part is a \SI{6}{\giga\hertz} oscillation; the imaginary part is the
amplitude decay rate. We want the table row $(f_m, Q_m)$.

\emph{Frequency.} By \cref{eq:freadout}, the frequency is the real part over
$2\pi$:
\[
  f_m \;=\; \frac{\Re\,\omega_m}{2\pi} \;=\; \SI{6.000}{\giga\hertz}.
\]

\emph{Quality factor.} By \cref{eq:qreadout} with the energy/loss convention
$Q=\Re\,\omega/(2\,\Im\,\omega)$,
\[
  Q_m \;=\; \frac{\omega_r}{2\,\omega_i}
        \;=\; \frac{2\pi\cdot 6.000\times 10^{9}}{2\cdot 2\pi\cdot 1.5\times 10^{6}}
        \;=\; \frac{6.000\times 10^{9}}{3.0\times 10^{6}}
        \;=\; 2000 .
\]
A $Q$ of $2000$ is a sharp resonance: the mode completes about $2000$
radians---roughly $320$ full cycles---before its amplitude falls by a factor of
$e$. The factor of $2$ in the denominator is the amplitude-to-energy conversion
(\cref{sec:eig-reduce}); had we used the bare $\omega_r/\omega_i$ convention we
would have reported $4000$, which is why the reduction carries $\kappa$ as a
parameter and never hard-codes the factor.

\emph{The lossless limit.} Now let the loss vanish, $\omega_i\to 0$: a perfect
cavity. The denominator of \cref{eq:qreadout} goes to zero and $Q_m\to\infty$.
The reduction does not divide by zero and emit a \texttt{NaN}; it tests
$\kappa=0$ and returns $Q=\infty$, the honest physical answer---a lossless cavity
rings forever. This single guard is the one branch in an otherwise straight-line
arithmetic readout, and it exists because the infinite-$Q$ case is the ideal that
every superconducting or low-loss design approaches, not an error to be trapped.
\end{workedexample}

\section{Reading out the mode fields, and where this goes}
\label{sec:eig-fields}

The $(f,Q)$ table is the headline product, but each converged eigenpair also
carries the mode \emph{field}, and the readout map recovers it too. The
eigenvector $\vx_m$ \emph{is} the electric mode field $\vE_m$, in the
$\Hcurl$ degrees of freedom; the magnetic field follows from Faraday's law as the
curl of $\vE$,
\begin{equation}
  \vB_m \;=\; -\frac{1}{i\,\omega_m}\,\Curl\,\vE_m,
  \label{eq:bfield}
\end{equation}
one curl and a complex scale per mode. (An eigenvector is defined only up to a
complex scalar, so Palace fixes the phase by a normalization before recording
the field---otherwise the same mode could come back rotated arbitrarily in the
complex plane.) These fields are what a designer plots to \emph{see} the mode:
where the electric field piles up (the high-field regions that set the breakdown
limit), where the magnetic field loops, which part of the structure a given
resonance lives in.

The whole pipeline, finally, is the clean composition the framework names
(\cref{lst:eigcompose}): assemble the pencil, one black-box solve, one readout
map.

\begin{figure}[t]
\begin{lstlisting}[style=l4style, caption={[The eigenmode composition]The
eigenmode feature as a composition of firm vocabulary. Assemble the
$(\Kop,\Cop,\Mop)$ pencil once (three folds), hand it to the opaque
\texttt{eigsolve} black box once, and \texttt{map} the per-mode readout over the
converged eigenpairs. The only Palace-written loop is the final \texttt{map};
the solve is a single opaque call.}, label={lst:eigcompose}]
eigenmode :: EigenmodeConfig -> EigenmodeResult
eigenmode cfg =
  let space  = nd_space cfg                          -- H(curl) edge space
      k      = fe_assemble space [ curl_curl  (reluctivity cfg) ]   -- K, once
      c      = fe_assemble space [ damping    (conductivity cfg) ]  -- C, once (may be empty)
      m      = fe_assemble space [ mass_term  (permittivity cfg) ]  -- M, once
      pencil = eig_pencil k c m (target cfg) (n_modes cfg)          -- bundle + target
      eigs   = eigsolve pencil (initial_space cfg)   -- ONE opaque black-box solve
  in  map (readout cfg) eigs                          -- the only loop we write: readout map
\end{lstlisting}
\end{figure}

\begin{remark}[Applications]
The eigenmode analysis is the workhorse behind a surprising range of devices.
Microwave \emph{cavity filters} are built by coupling several resonators, each a
controlled mode whose $f$ and $Q$ this analysis predicts before any metal is cut.
Superconducting-qubit \emph{readout resonators} and \emph{coupling buses} are
high-$Q$ modes of on-chip cavities; the eigenmode solve gives their frequencies
and participation ratios (\cref{sec:eig-reduce}), which set qubit--cavity
coupling and dictate where loss must be controlled. Particle-accelerator
\emph{RF cavities} are designed mode-by-mode for a target accelerating frequency.
In every case the deliverable is the same rank-1 table: a short list of
$(f,Q)$ rows, plus the field patterns that explain them.
\end{remark}

A close cousin of this analysis appears one chapter on. The \emph{boundary-mode}
(wave-port) analysis of \cref{ch:waveports} is, in the taxonomy of
\cref{ch:situating}, eigenmode on a two-dimensional cross-section: the same
opaque black-box solve corner, the same rank-1 table shape, applied to the
transverse modes of a waveguide cross-section rather than the volume modes of a
cavity. Having understood the volume eigenmode, the wave-port analysis will read
as a small, well-localized change of domain.

\summarysection
The eigenmode analysis is the anchor pipeline (\cref{ch:electrostatics}) with the
\texttt{solve} corner swapped for the fourth corner of \cref{ch:situating}: the
opaque black box. \emph{Assemble} builds an operator \emph{pencil} once---$(\Kop,\Mop)$
for a lossless cavity (the generalized EVP $\Kop\vx=\lambda\Mop\vx$,
$\lambda=\omega^2$) or $(\Kop,\Cop,\Mop)$ for a lossy one (the quadratic EVP,
$\lambda=i\omega$). \emph{Solve} is a single \code{eigsolve} call into a
specialized library (SLEPc/ARPACK, Krylov--Schur with shift-and-invert,
\cref{ch:eigenvalues}, \cref{ch:krylovschur}); Palace writes no eigen-iteration
loop, and the library only calls back for the operator applies (each containing
an inner linear solve from the shift-invert). Before the solve, the
\emph{divergence-free projection} (\cref{ch:bc}) removes the curl--curl operator's
spurious gradient null-space ``modes,'' so the black box returns physical
resonances. \emph{Reduce} is a \code{map} over the converged eigenpairs producing
the rank-1 $(f,Q)$ table: $f=\Re\,\omega/2\pi$ after the eigenvalue un-transform
(\cref{ch:eigenvalues}, \cref{ch:reductions}), $Q=\Re\,\omega/(2\,\Im\,\omega)$
with the lossless $\kappa\to0$ limit reported as $Q=\infty$. The mode fields read
out as $\vE=\vx$ and $\vB=-\tfrac{1}{i\omega}\Curl\vE$. The single Palace-written
loop in the whole driver is that final readout map---never a solve.

\furtherreading
Pozar's \emph{Microwave Engineering}~\citep{pozar2011} is the standard
device-level account of cavity resonators, resonant frequency, and the quality
factor, including the energy/loss definition of $Q$ and the loaded-versus-unloaded
distinction we glossed over. Jin's \emph{Finite Element Method in
Electromagnetics}~\citep{jin2014} develops the curl--curl eigenproblem on
$\Hcurl$ edge elements in full, including the gradient null-space and the
divergence-free constraint that tames it (\cref{sec:eig-spurious})---the FEM-level
companion to this chapter. For the eigensolver itself---why a few interior
eigenpairs of a large pencil are best found by Krylov--Schur with a
shift-and-invert spectral transform, and why that is the right thing to hand to a
library---Saad's \emph{Numerical Methods for Large Eigenvalue
Problems}~\citep{saad2011eig} is the reference; \cref{ch:eigenvalues} and
\cref{ch:krylovschur} build the relevant machinery on its foundation.

\exercisessection
\begin{enumerate}[nosep]
  \item \textbf{The swapped corner.} State precisely which stage of the
  \texttt{config}$\to$\texttt{assemble}$\to$\texttt{solve}$\to$\texttt{reduce}
  skeleton differs between electrostatics (\cref{ch:electrostatics}) and the
  eigenmode analysis, and which stages are unchanged. In one sentence, name the
  single structural property of the solve stage that makes the eigenmode corner
  ``opaque.''
  \item \textbf{The only loop.} The chapter insists the eigenmode driver writes
  exactly one loop, and that it is not a solve. Which loop is it, what does each
  of its iterations produce, and why is it a \code{map} (independent rows) rather
  than a \code{fold} (chained state)?
  \item \textbf{Un-transform.} A linear-EVP solve returns the eigenvalue
  $\lambda=39.5/L^2$ (in units where $c=1$). Recover the ordinary frequency $f$.
  Then a quadratic-EVP solve returns $\lambda=i\,(2\pi\cdot 4.0\times10^9)$;
  recover $f$ from \emph{it}. Why do the two cases use different un-transforms?
  \item \textbf{Q and the factor of two.} A mode has
  $\omega=2\pi(\SI{10}{\giga\hertz})+i\,2\pi(\SI{5}{\mega\hertz})$. Compute $Q$ in
  the energy/loss convention $Q=\Re\,\omega/(2\,\Im\,\omega)$. A colleague reports
  $Q=2000$ for the same mode; what convention did they use, and which factor did
  they drop?
  \item \textbf{The lossless limit.} Explain why the reduce stage tests
  $\kappa=0$ and returns $Q=\infty$ rather than letting \cref{eq:qreadout} divide
  by zero. Is $Q=\infty$ a numerical degeneracy to be avoided, or a physical
  answer to be reported? Give one real device whose $Q$ this limit describes.
  \item \textbf{Spurious modes.} The curl--curl operator satisfies
  $\Kop\Grad\phi=\mathbf 0$ for every scalar field $\phi$. Explain why this gives
  the raw eigenproblem a cloud of zero-eigenvalue ``modes,'' why those are not
  resonances, and why the divergence-free projection must be applied
  \emph{before} the black-box solve rather than filtered out of its output
  afterward.
  \item \textbf{The callback.} The eigensolver is opaque, yet the chapter says it
  ``calls back'' into Palace's code mid-iteration. What does it call back to
  compute, and why does that callback contain an \emph{inner} linear solve?
  Relate the inner solve to the shift-and-invert spectral transform of
  \cref{ch:eigenvalues}.
  \item \textbf{Refining the mesh.} In \cref{wex:cavity1d} the coarse $3\times3$
  mesh gave $\lambda_1\approx 9.37/L^2$ against the exact $9.87/L^2$. Without
  recomputing, predict qualitatively what happens to the three discrete
  eigenvalues as $N$ grows from $3$ to $30$: which converge fastest, and why does
  the highest mode lag? (Hint: how many nodes per wavelength does each mode get?)
\end{enumerate}

\chapter{Driven Simulations and S-Parameters}
\label{ch:driven}

\chapterorient{We have built the anchor: electrostatics assembles one operator,
\emph{hoists} it out of the loop, and maps a family of independent solves over it
(\cref{ch:electrostatics}). This chapter takes that anchor and moves a single
line of code---the line that captures the operator into the solver---from
\emph{outside} the loop to \emph{inside} it. That one move changes the analysis
from electrostatics to the driven frequency sweep, the workhorse of RF and
microwave engineering. We feed a time-harmonic signal into the ports of a device,
sweep the frequency across a band, and measure how much of each wave is
transmitted and how much reflected: the scattering matrix $S(\omega)$. Along the
way the operator is no longer fixed---it is rebuilt at every frequency from a
small fixed basis---and the reduction is no longer a Gram but a port-projection.
By the end you will see the driven analysis for exactly what it is: the anchor
with the operator rebuilt \emph{inside} the sweep, capped by a projection
readout.}

\section{The physical question: what a network analyzer measures}

Point a vector network analyzer at a microwave device---a filter, a coupler, a
length of transmission line, a resonator on a chip---and it does something very
concrete. It connects to the device's \emph{ports}: the places where signal
enters and leaves, typically coaxial connectors or on-chip waveguide
cross-sections. It then injects a pure sinusoid at one port, at one frequency
$\omega$, waits for the transients to die away so the device is in
\emph{time-harmonic steady state}, and measures two things at every port: how
much of the injected wave came \emph{back} out of the port it went in
(reflection), and how much came out of each \emph{other} port (transmission). It
repeats this across a whole band of frequencies. The result is a table of complex
numbers $S_{ij}(\omega)$---the \emph{scattering parameters}---and it is the
universal currency of the field. A filter's passband, a coupler's split ratio, an
antenna's match, a chip interconnect's loss: all of them are read directly off
$S(\omega)$.

\begin{physicsbox}
A two-port device (\cref{fig:setup}) sits between a source and a load. At a single
frequency, the analyzer launches an \emph{incident} wave $a_1$ into port~1. Part
of it \emph{reflects} straight back out of port~1 (the wave $b_1$), and part
\emph{transmits} through the device and emerges at port~2 (the wave $b_2$). The
scattering parameters are the ratios of outgoing to incoming wave amplitudes,
\[
  b_1 = S_{11}\,a_1 + S_{12}\,a_2, \qquad
  b_2 = S_{21}\,a_1 + S_{22}\,a_2,
\]
so that $S_{ij}$ is ``the wave that comes out of port~$i$ per unit wave injected
into port~$j$, with all other ports terminated.'' The diagonal entries are
\emph{reflection} coefficients (how much bounces back); the off-diagonals are
\emph{transmission} coefficients (how much gets through). Everything in this
chapter exists to compute that matrix, frequency by frequency.
\end{physicsbox}

\begin{figure}[t]
\centering
\begin{tikzpicture}[font=\footnotesize, >=Stealth]
  % the device
  \node[stage, minimum width=30mm, minimum height=18mm, fill=simBgGray,
        draw=simGray] (dev) at (0,0) {2-port device\\\scriptsize (filter / line)};
  % port 1 (left) feed line
  \draw[very thick, simBlue] (-3.6,0.35) -- (dev.west |- 0,0.35);
  \draw[very thick, simBlue] (-3.6,-0.35) -- (dev.west |- 0,-0.35);
  \node[font=\scriptsize, text=simBlue] at (-3.0,0.85) {port 1};
  % port 2 (right) feed line
  \draw[very thick, simGreen] (dev.east |- 0,0.35) -- (3.6,0.35);
  \draw[very thick, simGreen] (dev.east |- 0,-0.35) -- (3.6,-0.35);
  \node[font=\scriptsize, text=simGreen] at (3.0,0.85) {port 2};
  % incident wave a1
  \draw[->, thick, simRust] (-3.2,1.25) -- (-1.9,1.25)
     node[midway, above, font=\scriptsize]{incident $a_1$};
  % reflected wave b1
  \draw[->, thick, simRust] (-1.9,-1.25) -- (-3.2,-1.25)
     node[midway, below, font=\scriptsize]{reflected $b_1=S_{11}a_1$};
  % transmitted wave b2
  \draw[->, thick, simGreen!70!black] (1.9,1.25) -- (3.2,1.25)
     node[midway, above, font=\scriptsize]{transmitted $b_2=S_{21}a_1$};
  % the field inside
  \node[font=\scriptsize\itshape, text=simInk] at (0,0.0) {};
  \node[font=\scriptsize\itshape, align=center] at (0,-1.5)
    {drive port 1, terminate port 2:\\ read $S_{11}$ (reflect), $S_{21}$ (transmit)};
\end{tikzpicture}
\caption[A two-port device and its scattering waves]{The physical setup. A
time-harmonic wave $a_1$ is injected into port~1 of a two-port device. Part
reflects back out of port~1 (the wave $b_1 = S_{11}a_1$); part transmits through
and emerges at port~2 ($b_2 = S_{21}a_1$). Driving port~1 and reading both ports
gives the first column of the scattering matrix; driving port~2 gives the second.
A simulation reproduces this measurement: it solves the time-harmonic field
inside the device with port~1 excited, then projects the solved field onto each
port's mode to recover $S_{11}$ and $S_{21}$.}
\label{fig:setup}
\end{figure}

The analyzer's loop---inject at a frequency, measure the scattered waves, step to
the next frequency---is exactly the loop the simulation runs. The simulator does
not have a physical sinusoid and a physical analyzer; it has the time-harmonic
\emph{field equation} and a way to read scattered waves off the solved field. But
the shape is identical: a sweep over frequency, an independent steady-state
problem at each one, and a projection readout. That shape is the operator-varying
map of \cref{ch:situating}, and the rest of this chapter is its anatomy.

\section{Assemble the basis once: the frequency-dependent operator}
\label{sec:assemble}

Recall the time-harmonic reduction of Maxwell from \cref{ch:reductions-pde}. A
field oscillating at a single angular frequency $\omega$ is written
$\vE(\vx,t) = \Re\{\vE(\vx)\,e^{i\omega t}\}$, and the substitution
$\partial_t \mapsto i\omega$, $\partial_t^2 \mapsto -\omega^2$ turns Maxwell's
curl--curl equation with sources into a single complex linear system for the
phasor field $\vE(\vx)$. Discretized on a Nédélec $\Hcurl$ finite-element space
(\cref{ch:functionspaces}), that system reads
\begin{equation}
  \mat{A}(\omega)\,\vE \;=\; \vb(\omega),
  \qquad
  \mat{A}(\omega) \;=\; \Kop \;+\; i\omega\,\Cop \;-\; \omega^2\Mop
    \;\bigl[\,+\,\mat{A}_2(\omega)\,\bigr],
  \label{eq:driven-system}
\end{equation}
where the three operators are the same characters from every chapter of this
book: $\Kop$ is the curl--curl \emph{stiffness} (built from the reluctivity
$\nu=1/\mu$), $\Cop$ is the \emph{damping} (built from the surface impedance /
conductivity---the term that makes the system lossy and complex), and $\Mop$ is
the \emph{mass} (built from the permittivity $\varepsilon$). The optional
$\mat{A}_2(\omega)$ collects any genuinely frequency-dependent contribution that
is not a clean monomial in $\omega$---a dispersive material model, or a
wave-port term whose modal admittance varies with frequency.

Here is the structural fact this whole chapter turns on, and it is worth stating
in two halves.

\emph{The basis is fixed.} The three operators $\Kop$, $\Cop$, $\Mop$ do not
depend on $\omega$. They are built from the mesh, the function space, and the
materials---all of which are frozen the moment the configuration is parsed. So we
assemble them exactly \emph{once}, before the sweep ever begins, and store them as
a small fixed record (\cref{ch:records}):
\[
  \texttt{fam} \;=\; \{\,\Kop,\ \Cop,\ \Mop,\ \mat{A}_2\,\}.
\]
This is precisely the ``assemble once'' that electrostatics did with its single
$\Kop$ (\cref{ch:electrostatics}); the driven analysis just has three operators in
the basis instead of one.

\emph{The system operator is not fixed.} The operator that actually appears on the
left-hand side of the solve, $\mat{A}(\omega)$, is a frequency-dependent
\emph{combination} of the fixed basis. At each frequency it is a different
operator: the weights $1$, $i\omega$, $-\omega^2$ multiply the same three stored
operators, but those weights change with $\omega$, so $\mat{A}(\omega_1)$ and
$\mat{A}(\omega_2)$ are genuinely different matrices. We never store
$\mat{A}(\omega)$ permanently; we \emph{form} it from the basis at each frequency
we visit, use it for one solve, and discard it.

\begin{figure}[t]
\centering
\begin{tikzpicture}[font=\footnotesize, >=Stealth]
  % fixed basis on the left
  \node[opbox, minimum width=9mm] (K) at (0,1.4) {$\Kop$};
  \node[opbox, minimum width=9mm] (C) at (0,0.0) {$\Cop$};
  \node[opbox, minimum width=9mm] (M) at (0,-1.4) {$\Mop$};
  \node[font=\scriptsize\itshape, text=simGray, align=center] at (0,-2.3)
    {fixed basis\\(assembled once)};
  % weights
  \node[font=\footnotesize] (wK) at (2.7,1.4) {$1$};
  \node[font=\footnotesize] (wC) at (2.7,0.0) {$i\omega$};
  \node[font=\footnotesize] (wM) at (2.7,-1.4) {$-\omega^2$};
  \node[font=\scriptsize\itshape, text=simRust, align=center] at (2.7,-2.3)
    {weights\\(depend on $\omega$)};
  \draw[->, simGray] (K.east) -- (wK.west);
  \draw[->, simGray] (C.east) -- (wC.west);
  \draw[->, simGray] (M.east) -- (wM.west);
  % sum node
  \node[draw=simBlue, circle, thick, fill=simBgBlue, inner sep=1.5pt]
    (sum) at (5.0,0.0) {$\Sigma$};
  \draw[->, simInk] (wK.east) .. controls +(0.8,0) and +(-0.8,0.3) .. (sum.west);
  \draw[->, simInk] (wC.east) -- (sum.west);
  \draw[->, simInk] (wM.east) .. controls +(0.8,0) and +(-0.8,-0.3) .. (sum.west);
  % result
  \node[product, minimum width=20mm] (A) at (8.0,0.0) {$\mat{A}(\omega)$};
  \draw[->, thick, simInk] (sum.east) -- (A.west);
  \node[font=\scriptsize\itshape, align=center] at (8.0,-1.1)
    {one system operator\\\emph{per frequency}};
  \node[font=\scriptsize, text=simBlue, align=center] at (5.0,-2.3)
    {\code{assemble\_frequency\_operator}\\$=$ \code{linear\_combination}};
\end{tikzpicture}
\caption[The system operator as an operator-valued linear combination]{The
frequency-dependent system operator is an \emph{operator-valued linear
combination} (\cref{ch:combinators}) of the fixed basis. The three operators
$\Kop,\Cop,\Mop$ are assembled once and never change. The scalar weights
$1,\,i\omega,\,-\omega^2$ \emph{do} change with $\omega$. Summing
weight-times-operator yields the system operator $\mat{A}(\omega)$---a fresh
operator at every frequency, formed from a fixed basis. This is the same
\code{linear\_combination} combinator that sums tensors in BLAS-1; here its
operands are \emph{operators} and its scalars are \emph{complex}.}
\label{fig:opcombo}
\end{figure}

\subsection{The forming step is a linear combination}

Forming $\mat{A}(\omega)$ from the basis is not a new idea. It is the
\code{linear\_combination} combinator from \cref{ch:combinators}---the same one
that builds $\alpha\vx + \beta\vy$ in BLAS-1---specialized to the case where the
\emph{operands are operators} rather than tensors and the \emph{scalars are
complex} (\cref{fig:opcombo}). Spelled out as the framework verb
\code{assemble\_frequency\_operator}:
\begin{lstlisting}[style=l4style]
-- form the per-omega system operator from the fixed basis
assemble_frequency_operator :: FrequencyOperatorFamily -> Scalar -> LinOp
assemble_frequency_operator fam omega =
  linear_combination                  -- operator-operand corner; complex scalars
    [ (1,            fam.K)           --  K   with weight 1
    , (1i * omega,   fam.C)           --  C   with weight  i*omega
    , (-(omega^2),   fam.M)           --  M   with weight -omega^2
    , (1,            fam.A2 omega) ]  --  A2(omega) with weight 1
\end{lstlisting}
There is no separate fold algebra here---the body is literally
\code{linear\_combination} at term-list length four, with the operand monoid being
operator addition (commutative, identity the zero operator) instead of tensor
addition. The framework gives it its own name only because it is the named verb the
driven sweep reaches for; underneath, it is the BLAS-1 combinator the whole book
already trusts. The system operator $\mat{A}(\omega)$ is \emph{affine in $\omega$}
in the monomial part ($1, i\omega, -\omega^2$) modulo the extra term
$\mat{A}_2$; that affine structure is exactly what makes a reduced-order model
possible (\cref{sec:cost}), though the uniform sweep does not exploit it.

\begin{notationbox}
The system operator carries a complex element type because two of its three
weights are complex: $i\omega$ is imaginary and multiplies the damping $\Cop$,
$-\omega^2$ is real and multiplies the mass $\Mop$. So even when $\Kop,\Cop,\Mop$
are real operators, $\mat{A}(\omega)$ is a \emph{complex} operator, and the field
$\vE$ it solves for is a complex (phasor) field: its magnitude is the amplitude of
the oscillation and its argument is the phase. This is why the driven solve needs
GMRES over $\Complex^N$, not CG over $\Real^N$ (\cref{sec:solve}). The complex
arithmetic is not an implementation detail; it carries the phase information that
\emph{is} the answer.
\end{notationbox}

\section{Solve: the operator-varying map (the hoist, negated)}
\label{sec:solve}

Now the sweep. We are given a list of frequencies
$[\omega_0, \omega_1, \dots, \omega_{m-1}]$---the band, sampled---and at each one
we must solve \cref{eq:driven-system} for the field $\vE_k = \vE(\omega_k)$. The
solves are \emph{independent}: the field at $\omega_3$ does not depend on the
field at $\omega_2$, because each is the steady-state response to its own
frequency's excitation. Independence means the sweep is a \emph{map}, not a fold
(\cref{ch:situating}): the frequencies could be solved in any order, or on
separate machines, and the answers would be identical.

What makes it the operator-\emph{varying} map---corner~2 of \cref{ch:situating},
not corner~1---is that the operator on the left-hand side is not the same from one
iteration to the next. At each frequency we must:
\begin{enumerate}[nosep]
  \item \textbf{rebuild} the system operator $\mat{A}(\omega_k)$ from the fixed
  basis (\cref{sec:assemble});
  \item \textbf{re-capture} it into the solver---the call that tells the Krylov
  solver ``\emph{this} is the matrix you are about to invert'';
  \item \textbf{form} the right-hand side $\vb(\omega_k)$ from the port
  excitation at this frequency;
  \item \textbf{solve} $\mat{A}(\omega_k)\,\vE_k = \vb(\omega_k)$ by preconditioned
  GMRES (\cref{ch:gmres}).
\end{enumerate}
The framework names this combinator \code{frequency\_sweep}:
\begin{lstlisting}[style=l4style]
-- the operator-VARYING map: rebuild the operator INSIDE the loop
frequency_sweep :: FrequencyOperatorFamily -> [Scalar] -> [Tensor]
frequency_sweep fam omegas =
  map solve_at omegas               -- independent: a map, not a fold
  where
    solve_at omega =
      let a = assemble_frequency_operator fam omega   -- (1) rebuild A(omega)
          b = excitation fam omega                    -- (3) per-omega RHS
      in  ksp_solve a b                               -- (2,4) capture + GMRES solve
\end{lstlisting}

\subsection{The one line that is the whole difference}

Set the driven sweep beside electrostatics and almost everything is shared: both
assemble a basis once, both run a family of independent solves, both produce a
family of fields, both feed a rank-2 reduction. The single structural difference
is the \emph{position of the operator-capture call} relative to the loop
(\cref{fig:hoist}).

\begin{lstlisting}[style=l4style]
-- ELECTROSTATIC (corner 1): operator-capture HOISTED OUTSIDE the loop
solve_family k rhss =                    -- k assembled once, captured once
  let solver = ksp_setup k               -- <<-- SetOperators OUTSIDE
  in  map (\b -> ksp_run solver b) rhss   -- map: same operator every solve

-- DRIVEN (corner 2): operator-capture INSIDE the loop
frequency_sweep fam omegas =
  map (\omega ->
        let a      = assemble_frequency_operator fam omega
            solver = ksp_setup a          -- <<-- SetOperators INSIDE
        in  ksp_run solver (excitation fam omega))
      omegas                              -- map: fresh operator every solve
\end{lstlisting}

In electrostatics, \code{ksp\_setup}---the call that factors the preconditioner and
hands the operator to the solver---happens \emph{once}, before the map, because the
operator $\Kop$ is the same for every right-hand side. We \emph{hoist} it out of
the loop. In the driven sweep, \code{ksp\_setup} happens \emph{inside} the map,
once per frequency, because the operator $\mat{A}(\omega)$ is a function of the
loop index. We \emph{cannot} hoist it.

\begin{figure}[p]
\centering
\begin{tikzpicture}[font=\footnotesize, >=Stealth]
  % ============== LEFT: electrostatic, hoisted ==============
  \begin{scope}[shift={(0,0)}]
    \node[font=\small\bfseries, text=simBlue] at (1.5,4.3)
      {Electrostatic (corner 1)};
    \node[font=\scriptsize\itshape, text=simGray] at (1.5,3.9)
      {assemble once, \emph{hoist}};
    % the operator capture, outside
    \node[opbox, minimum width=30mm, fill=simBgBlue, draw=simBlue]
      (setup1) at (1.5,3.2) {\code{SetOperators}($\Kop$)};
    \node[font=\scriptsize, text=simBlue] at (1.5,2.65) {(once, OUTSIDE)};
    % the loop box
    \node[draw=simGray, densely dashed, rounded corners,
          minimum width=42mm, minimum height=24mm, fill=simBgGray]
      (loop1) at (1.5,0.6) {};
    \node[font=\scriptsize, text=simGray, anchor=north west]
      at (loop1.north west) {\ for each $\vb_i$:};
    \node[opbox, minimum width=26mm] (sv1) at (1.5,0.4)
      {\code{ksp\_run}($\vb_i$)};
    \node[font=\scriptsize\itshape] at (1.5,-0.4) {same $\Kop$ each pass};
    \draw[->, thick, simInk] (setup1.south) -- (loop1.north);
    % loop back arrow
    \draw[->, simGray, densely dashed]
      (loop1.east) ++(0,-0.4) .. controls +(0.9,0) and +(0.9,0) ..
      ($(loop1.east)+(0,0.6)$);
  \end{scope}

  % divider
  \draw[simGray, thin] (4.1,-1.0) -- (4.1,4.5);

  % ============== RIGHT: driven, inside ==============
  \begin{scope}[shift={(5.3,0)}]
    \node[font=\small\bfseries, text=simRust] at (1.7,4.3)
      {Driven (corner 2)};
    \node[font=\scriptsize\itshape, text=simGray] at (1.7,3.9)
      {basis once, rebuild \emph{inside}};
    % the basis, outside (only the basis is hoisted)
    \node[opbox, minimum width=34mm, fill=simBgBlue, draw=simBlue]
      (basis2) at (1.7,3.2) {assemble basis $\{\Kop,\Cop,\Mop\}$};
    \node[font=\scriptsize, text=simBlue] at (1.7,2.65) {(once, OUTSIDE)};
    % the loop box
    \node[draw=simRust, densely dashed, rounded corners,
          minimum width=46mm, minimum height=30mm, fill=simBgRust]
      (loop2) at (1.7,0.1) {};
    \node[font=\scriptsize, text=simRust, anchor=north west]
      at (loop2.north west) {\ for each $\omega_k$:};
    \node[opbox, minimum width=36mm, fill=white] (form2) at (1.7,0.95)
      {form $\mat{A}(\omega_k)=\Kop{+}i\omega_k\Cop{-}\omega_k^2\Mop$};
    \node[opbox, minimum width=36mm, fill=simBgRust, draw=simRust]
      (setup2) at (1.7,0.1) {\code{SetOperators}($\mat{A}(\omega_k)$)};
    \node[font=\scriptsize, text=simRust] at (1.7,-0.45) {(every pass, INSIDE)};
    \node[opbox, minimum width=36mm, fill=white] (sv2) at (1.7,-1.05)
      {\code{ksp\_run}($\vb(\omega_k)$)};
    \draw[->, simInk] (form2.south) -- (setup2.north);
    \draw[->, simInk] (setup2.south) -- (sv2.north);
    \draw[->, thick, simInk] (basis2.south) -- (loop2.north);
    \draw[->, simRust, densely dashed]
      (loop2.east) ++(0,-0.4) .. controls +(0.9,0) and +(0.9,0) ..
      ($(loop2.east)+(0,1.0)$);
  \end{scope}
\end{tikzpicture}
\caption[The hoist, and its negation]{The chapter's centerpiece: the operator
capture, hoisted versus not. \emph{Left}---electrostatics assembles $\Kop$ once
and captures it into the solver \emph{outside} the loop; every pass of the map
solves a new right-hand side against the \emph{same} operator. \emph{Right}---the
driven sweep assembles only the \emph{basis} $\{\Kop,\Cop,\Mop\}$ outside the
loop, then, \emph{inside} the loop, forms a fresh $\mat{A}(\omega_k)$ and captures
it (\code{SetOperators}) before each solve. The capture call moves from outside
the loop to inside it. That single move is the entire structural difference
between the two analyses---and moving it the wrong way does not slow the code
down, it computes the wrong answer.}
\label{fig:hoist}
\end{figure}

\begin{keyidea}
\textbf{Driven $=$ the anchor with the operator rebuilt \emph{inside} the sweep,
plus a port-projection reduction.} Electrostatics hoists the operator capture
\emph{outside} the loop (corner~1, the fixed-operator map). The driven analysis
moves that one call \emph{inside} the loop (corner~2, the operator-varying map),
because the system operator $\mat{A}(\omega) = \Kop + i\omega\Cop - \omega^2\Mop$
is a function of the loop index. The fixed \emph{basis} $\{\Kop,\Cop,\Mop\}$ is
still assembled once; only the \emph{combination} is rebuilt per frequency. Cap
the sweep with the S-parameter port-projection instead of the Gram, and you have
the entire driven driver. The position of one capture call relative to one loop is
load-bearing: hoist it and you solve every frequency against the operator of the
first one.
\end{keyidea}

\begin{pitfall}
A well-meaning optimizer sees the \code{SetOperators} call inside the loop, notices
it is ``expensive'' (it refactors the preconditioner), and hoists it out to ``save
work.'' In electrostatics that is correct and free. In the driven sweep it is a
\emph{silent wrong answer}: every frequency is now solved against
$\mat{A}(\omega_0)$, the operator of the first frequency, so the computed
$\vE_k$ is the steady-state response of the wrong physical system. Nothing crashes;
the numbers are simply wrong, and wrong in a way that looks plausible (a smooth
curve, just not the device's curve). The lesson of \cref{ch:situating} stated
operationally: the loop position of one setup call is part of the algorithm, not a
performance knob.
\end{pitfall}

\subsection{Why the solve is GMRES, not CG}

The driven system operator $\mat{A}(\omega) = \Kop + i\omega\Cop - \omega^2\Mop$ is
\emph{complex} and \emph{non-Hermitian}: the $i\omega\Cop$ term is purely
imaginary, and the $-\omega^2\Mop$ term subtracts mass from stiffness, so the
operator is indefinite as well. None of the structure CG relies on (\cref{ch:cg})
survives: there is no symmetric positive-definite form, hence no short
recurrence, hence no CG. This is exactly the regime GMRES was built for
(\cref{ch:gmres})---a general nonsymmetric, possibly complex non-Hermitian
operator---and the driven solve is its canonical application. Each frequency's
solve is a full preconditioned GMRES (or FGMRES, if the preconditioner itself
varies), and because the operator changes every frequency, the preconditioner must
be rebuilt or refreshed every frequency too (\cref{sec:cost}).

\section{The composition: \texttt{driven} end to end}

We can now write the whole driver as one composition, in the framework's
vocabulary. It is three stages: assemble the basis once, sweep with the operator
rebuilt inside, reduce to S-parameters (\cref{fig:pipeline}).
\begin{lstlisting}[style=l4style]
-- the whole driven driver: config in, frequency response out
driven :: DrivenConfig -> FrequencyResponse
driven cfg =
  let space  = nd_space cfg                                    -- H(curl) space
      fam    = { K  = fe_assemble space [ curl_curl (reluctivity cfg) ]  -- assemble
               , C  = fe_assemble space [ damping cfg ]                  --   basis
               , M  = fe_assemble space [ mass (permittivity cfg) ]      --   ONCE
               , A2 = \w -> extra_system_matrix space cfg w }
      omegas = sample_frequencies cfg                          -- the swept band
      es     = frequency_sweep fam omegas    -- operator-VARYING map: A(w) rebuilt inside
  in  sparameter_reduce (ports cfg) es       -- port-projection reduction -> S(w)
\end{lstlisting}
Read it as a pipeline composition:
\[
  \texttt{driven}
    \;=\;
  \texttt{sparameter\_reduce}
    \;\circ\;
  \texttt{frequency\_sweep}
    \;\circ\;
  \texttt{fe\_assemble}^{\times 3}.
\]
The first stage (\code{fe\_assemble}, three times) is the once-only basis; the
middle stage (\code{frequency\_sweep}) is the operator-varying solve map; the last
stage (\code{sparameter\_reduce}) is the projection readout we turn to next. Each
stage is firm, and each links down to its own chapter: the assemble fold to
\cref{ch:assembly}, the combination to \cref{ch:combinators}, the solve to
\cref{ch:gmres}, the reduction to \cref{ch:reductions}.

\begin{figure}[t]
\centering
\begin{tikzpicture}[node distance=7mm and 9mm, font=\footnotesize, >=Stealth]
  \node[data] (cfg) {config};
  \node[stage, right=of cfg, align=center] (asm)
    {assemble basis\\$\{\Kop,\Cop,\Mop\}$\\\scriptsize(once)};
  \node[stage, right=of asm, fill=simBgRust, draw=simRust, align=center] (sweep)
    {\textbf{frequency\_sweep}\\\scriptsize per $\omega$:};
  \node[stage, right=of sweep, align=center] (red)
    {sparameter\_\\reduce};
  \node[product, right=of red] (S) {$S(\omega)$};
  \draw[flow] (cfg) -- (asm);
  \draw[flow] (asm) -- node[above,font=\scriptsize]{\code{fam}} (sweep);
  \draw[flow] (sweep) -- node[above,font=\scriptsize]{$[\vE_k]$} (red);
  \draw[flow] (red) -- (S);
  % the per-omega inner body, drawn below the sweep box
  \node[opbox, below=10mm of sweep, fill=white, align=center] (inner)
    {form $\mat{A}(\omega)$ $\to$ \code{SetOperators} $\to$ GMRES};
  \draw[refedge] (sweep.south) -- (inner.north);
  \node[font=\scriptsize\itshape, text=simRust, below=1.5mm of inner]
    {operator rebuilt + re-captured \emph{inside}};
\end{tikzpicture}
\caption[The driven pipeline]{The driven pipeline. The basis
$\{\Kop,\Cop,\Mop\}$ is assembled once. \code{frequency\_sweep} maps over the
swept band; its per-$\omega$ body (drawn below) forms the system operator
$\mat{A}(\omega)$, re-captures it with \code{SetOperators}, and solves by GMRES.
The solved fields $[\vE_k]$ feed \code{sparameter\_reduce}, the port-projection
that yields the scattering matrix $S(\omega)$. The red box marks the one stage
whose operator changes each step---the entire difference from electrostatics.}
\label{fig:pipeline}
\end{figure}

\begin{remark}[Uniform sweep versus adaptive ROM]
The composition above is the \emph{uniform} sweep: a fixed list of frequencies,
each solved fully and independently---an honest map. Palace also offers an
\emph{adaptive} sweep, which builds a projection-based reduced-order model (a PROM)
and greedily adds the next sample frequency where the model's error is largest. The
adaptive sweep is \emph{not} a map: each new frequency choice depends on the model
built from all previous solves, so it is a state-generated \emph{fold}
(\cref{ch:situating}, \cref{ch:transient})---the running model is threaded through
the loop. We mention it only to place it: the uniform sweep is the composed shape
this chapter studies, and it is the operator-varying map. The affine-in-$\omega$
structure of $\mat{A}(\omega)$ (\cref{sec:assemble}) is exactly what makes the
reduced-order model cheap to evaluate, but the uniform sweep does not exploit it.
\end{remark}

\section{Reduce: the port-projection}
\label{sec:reduce}

The solve gives us a family of complex fields, one per swept frequency (and, when
several ports are driven, one column per driven port). The reduction collapses
each solved field into the scattering parameters the user asked for. This is the
\emph{rank-2 port-projection}---reduction shape~2 of \cref{ch:situating} and
\cref{ch:reductions}---and it is emphatically \emph{not} a Gram.

The port modes $\vs_i$ themselves do not come from the driven solve. They come
from the \emph{boundary-mode} analysis (\cref{ch:waveports}), which solves a small
two-dimensional eigenproblem on each port's cross-section to find the field
pattern that propagates there. The driven reduction \emph{consumes} those modes as
the basis it projects onto. Given the solved field $\vE_j$ for drive port $j$ and
the receiver modes $\vs_i$, each scattering entry is
\begin{equation}
  \boxed{\;S_{ij}
    \;=\;
  \sigma_{ij}\,\bigl(\,\inner{\vs_i}{\vE_j(\omega)} \;-\; \delta_{ij}\,\bigr)\;}
  \label{eq:sparam}
\end{equation}
where each piece earns its place:
\begin{itemize}[nosep]
  \item $\inner{\vs_i}{\vE_j(\omega)}$ is the \emph{linear projection} of the
  drive-$j$ field onto receiver mode $i$---how much of the solved field looks like
  the wave that propagates out of port $i$. For a lumped port this is the inner
  product $\vs_i\cdot\vE_j$; for a wave port it is the modal overlap
  $(\vE_j\times\vH_i^\star)\cdot\vn$ integrated over the port face. It is
  \emph{linear} in the field $\vE_j$.
  \item $\delta_{ij}$ is the Kronecker self-term: $1$ when $i=j$, $0$ otherwise.
  On the diagonal it subtracts the \emph{incident} wave from the total field,
  leaving only the \emph{reflected} part. Without it, $S_{11}$ would report the
  whole field at port~1 (incident plus reflected); with it, $S_{11}$ reports the
  reflection alone, which is what ``return loss'' means.
  \item $\sigma_{ij}$ is the port-kind \emph{scale}: an impedance ratio
  $\sqrt{R_j/R_i}$ for lumped ports (generalized scattering parameters, which
  normalize for unequal port impedances), or a de-embedding phase
  $e^{ik_{n,i}d_i}\,e^{ik_{n,j}d_j}$ for wave ports (which shift the reference
  plane from the mesh boundary to the physical port location).
\end{itemize}

\begin{figure}[t]
\centering
\begin{tikzpicture}[font=\footnotesize, >=Stealth]
  % field-space axes
  \draw[->, simGray] (0,0) -- (0,3.4);
  \draw[->, simGray] (0,0) -- (5.0,0);
  % port mode basis vectors
  \draw[->, very thick, simGreen] (0,0) -- (3.6,0)
    node[right, font=\scriptsize, text=simGreen]{$\vs_1$ (port 1 mode)};
  \draw[->, very thick, simGreen] (0,0) -- (0,2.8)
    node[above, font=\scriptsize, text=simGreen]{$\vs_2$ (port 2 mode)};
  % solved field
  \draw[->, very thick, simRust] (0,0) -- (2.6,2.0)
    node[midway, above left, font=\scriptsize, text=simRust]{$\vE_1(\omega)$};
  % projections (dashed)
  \draw[densely dashed, simBlue] (2.6,2.0) -- (2.6,0)
    node[midway, right, font=\scriptsize, text=simBlue]{};
  \draw[densely dashed, simBlue] (2.6,2.0) -- (0,2.0);
  \node[font=\scriptsize, text=simBlue] at (2.6,-0.32)
    {$\inner{\vs_1}{\vE_1}$};
  \node[font=\scriptsize, text=simBlue, anchor=east] at (-0.1,2.0)
    {$\inner{\vs_2}{\vE_1}$};
  % the entries
  \node[font=\footnotesize, align=left, text=simInk] at (6.9,2.2)
    {$S_{11}=\sigma_{11}(\inner{\vs_1}{\vE_1}-1)$};
  \node[font=\scriptsize, text=simGray] at (6.9,1.75)
    {(reflection: self-term subtracted)};
  \node[font=\footnotesize, align=left, text=simInk] at (6.9,1.0)
    {$S_{21}=\sigma_{21}\,\inner{\vs_2}{\vE_1}$};
  \node[font=\scriptsize, text=simGray] at (6.9,0.55)
    {(transmission: no self-term)};
\end{tikzpicture}
\caption[The port projection]{The port-projection reduction. The solved field
$\vE_1(\omega)$ (driving port~1) lives in field space; the port modes $\vs_1,\vs_2$
(from the boundary-mode analysis, \cref{ch:waveports}) are the directions we
project onto. Projecting $\vE_1$ onto $\vs_1$ and subtracting the incident
self-term gives the reflection $S_{11}$; projecting onto $\vs_2$ (no self-term,
since $2\neq1$) gives the transmission $S_{21}$. The projection is \emph{linear}
in the field and the two index families ($\vs_i$ versus $\vE_j$) play different
roles---this is why it is a projection, not a Gram.}
\label{fig:proj}
\end{figure}

In the framework this reduction is the verb \code{sparameter\_reduce}, applied once
per swept frequency:
\begin{lstlisting}[style=l4style]
-- project each solved field onto each receiver port mode -> the S-matrix
sparameter_reduce :: [PortMode] -> [(Int, Tensor)] -> Matrix
sparameter_reduce ports family =
  matrix_from_columns
    [ [ scale ports i j * (project (ports !! i) e - selfterm i j)  -- entry S_ij
        | i <- [0 .. p-1] ]            -- over receiver ports i
      | (j, e) <- family ]            -- one column per drive port j
  where
    p              = length ports
    project s e    = port_dot s e               -- linear functional  s_i . E
    selfterm i j   = if i == j then 1 else 0    -- the -1 self-reflection (diagonal)
    scale ports i j = port_scale (ports!!i) (ports!!j)  -- impedance / de-embed
\end{lstlisting}

\subsection{Why it is rank-2 like a Gram but deliberately not a Gram}

The S-matrix and the capacitance matrix (\cref{ch:reductions}) both come out as
$p\times p$ arrays indexed by a port-like family. That shared \emph{output rank}
is the whole of their resemblance, and \cref{ch:situating} already warned that
reading more into it is the single most over-unified mistake the design could
make. The driven reduction makes the distinction sharp.

\begin{pitfall}
\textbf{A matrix-shaped result is not a Gram.} The Gram reduction
(capacitance, inductance) pairs a family \emph{with itself} through an operator,
$\vx_j^{\!\top}\Kop\vx_i$---a \emph{bilinear} form, symmetric by construction, with
the verb computing only the upper triangle. The S-parameter reduction projects one
family (the solved fields $\vE_j$) onto a \emph{different} basis (the port modes
$\vs_i$)---a \emph{linear} functional, two families in different roles, every entry
computed independently. Three facts forbid merging them: (1) the Gram is a
self-pairing, the projection is not; (2) the Gram is symmetric by construction,
the projection's near-symmetry $S_{ij}\approx S_{ji}$ is reciprocity \emph{physics}
when the medium is reciprocal, never a construction the reduction imposes---there
is no \code{symmetric\_from\_upper} in its body; (3) the Gram is real and
energy-valued, the projection is complex and carries the inhomogeneous
$-\delta_{ij}$ self-term on its diagonal. Forcing one verb would smuggle a false
symmetry into the S-matrix or burden the Gram with machinery it does not need.
\emph{Sharing a shape is not sharing a meaning.}
\end{pitfall}

\section{Reading the answer: return loss and insertion loss}

The product is the scattering matrix as a function of frequency, $S(\omega)$, and
microwave engineers read it through two standard plots, both in decibels
(\cref{fig:splot}). The \emph{return loss} is $\abs{S_{11}}$ in dB---how much power
reflects off port~1; near a perfect match it dives toward $-\infty$ (a sharp dip),
because almost nothing reflects. The \emph{insertion loss} is $\abs{S_{21}}$ in
dB---how much power transmits from port~1 to port~2; in a passband it sits near
$0\,$dB (almost everything gets through), and in a stopband it plunges (the device
blocks the signal). A bandpass filter shows $\abs{S_{21}}$ flat and high across its
passband and low outside it, with $\abs{S_{11}}$ dipping sharply wherever the match
is good. The whole point of the simulation is to predict these curves before the
device is fabricated.

\begin{figure}[t]
\centering
\begin{tikzpicture}
\begin{axis}[
  width=0.84\linewidth, height=6.0cm,
  xlabel={frequency $f$ (GHz)}, ylabel={magnitude (dB)},
  xmin=3, xmax=7, ymin=-45, ymax=3,
  ytick={0,-10,-20,-30,-40},
  xtick={3,4,5,6,7},
  legend pos=south east, legend cell align=left,
  grid=both, grid style={simGray!20},
  tick label style={font=\footnotesize}, label style={font=\footnotesize},
  legend style={font=\footnotesize},
]
% |S21| insertion loss: passband near 0 dB around 5 GHz, rolling off outside
\addplot[very thick, simGreen, smooth, samples=120, domain=3:7]
  { -2 - 30*( (x-5)/1.0 )^2 / (1 + 0.55*((x-5))^2 ) - 0.5*(x-5)^2 };
\addlegendentry{$|S_{21}|$ (insertion loss)}
% |S11| return loss: deep dip near the match frequency 5 GHz
\addplot[very thick, simRust, smooth, samples=160, domain=3:7]
  { -3 - 32*exp( -((x-5)^2)/0.06 ) + 2.0*((x-5))^2 };
\addlegendentry{$|S_{11}|$ (return loss)}
% passband band marker
\draw[simBlue!50, densely dashed] (axis cs:4.4,3) -- (axis cs:4.4,-45);
\draw[simBlue!50, densely dashed] (axis cs:5.6,3) -- (axis cs:5.6,-45);
\node[font=\scriptsize, text=simBlue] at (axis cs:5.0,-43) {passband};
\end{axis}
\end{tikzpicture}
\caption[A swept S-parameter plot]{The product: a swept S-parameter plot for a
bandpass-like two-port. The insertion loss $\abs{S_{21}}$ (green) sits near
$0\,$dB across the passband---signal passes through---and rolls off on either
side. The return loss $\abs{S_{11}}$ (rust) dips sharply at the match frequency
near $5\,$GHz, where almost no power reflects (a good match), and rises away from
it. Each point on each curve is one converged GMRES solve at one frequency,
projected onto the port modes. Reading these two curves is reading the device.}
\label{fig:splot}
\end{figure}

\section{Cost: many independent solves, preconditioner reuse}
\label{sec:cost}

Because the sweep is a \emph{map}, the frequencies are embarrassingly parallel: a
hundred-frequency sweep is a hundred independent GMRES solves, and they can run on a
hundred machines with no communication. That is the good news, and it follows
directly from corner~2's independence.

The bad news is that each of those solves is a \emph{full} GMRES on a complex
indefinite operator, and---unlike electrostatics---we cannot amortize a single
factorization across the family, because the operator $\mat{A}(\omega)$ is
different every frequency. The preconditioner (\cref{ch:precond}) is what makes
each solve converge in a reasonable number of iterations, and rebuilding it from
scratch at every frequency is expensive. Two structural facts soften this:
\begin{itemize}[nosep]
  \item \textbf{Nearby frequencies have nearby operators.} Because
  $\mat{A}(\omega)$ depends smoothly on $\omega$, the operator at $\omega_{k+1}$ is
  close to the operator at $\omega_k$. A preconditioner built for $\omega_k$ is
  often a perfectly good preconditioner for $\omega_{k+1}$, so one can \emph{reuse}
  a preconditioner across a band of nearby frequencies and only rebuild it
  occasionally---trading a little extra GMRES iteration for a lot of saved
  factorization. When the preconditioner is allowed to drift this way, the inner
  solve becomes FGMRES (\cref{ch:gmres}), the flexible variant that tolerates a
  changing preconditioner.
  \item \textbf{The affine structure enables a reduced-order model.} The
  affine-in-$\omega$ form $\mat{A}(\omega) = \Kop + i\omega\Cop - \omega^2\Mop$
  means the solution $\vE(\omega)$ is a rational function of $\omega$, and a small
  number of full solves can fit a reduced-order model that predicts $\vE$ at the
  remaining frequencies cheaply. This is the adaptive sweep's payoff (the
  state-generated fold of the earlier remark); the uniform sweep leaves it on the
  table in exchange for simplicity and perfect parallelism.
\end{itemize}

\section{Two worked examples}

\begin{workedexample}{The operator-varying map, made concrete}{opvary}
We make the central claim---the operator changes every frequency---numerical, on a
deliberately tiny basis so the arithmetic is visible. Take the $2\times2$ real
basis operators
\[
  \Kop = \begin{pmatrix} 4 & -1 \\ -1 & 3 \end{pmatrix},\quad
  \Cop = \begin{pmatrix} 2 & 0 \\ 0 & 1 \end{pmatrix},\quad
  \Mop = \begin{pmatrix} 1 & 0 \\ 0 & 1 \end{pmatrix},
\]
assembled \emph{once}, and take $\mat{A}_2 \equiv 0$ for clarity. The system
operator is $\mat{A}(\omega) = \Kop + i\omega\Cop - \omega^2\Mop$. We form it at
three frequencies.

\medskip
\emph{At $\omega = 0$} (the static limit): the weights are $1,\,0,\,0$, so
\[
  \mat{A}(0) = \Kop = \begin{pmatrix} 4 & -1 \\ -1 & 3 \end{pmatrix}.
\]
At zero frequency the driven operator collapses to the electrostatic stiffness:
no damping, no mass. This is the anchor reappearing as a special case.

\medskip
\emph{At $\omega = 1$:} the weights are $1,\,i,\,-1$, so
\[
  \mat{A}(1)
    = \Kop + i\Cop - \Mop
    = \begin{pmatrix} 4-1 & -1 \\ -1 & 3-1 \end{pmatrix}
      + i\begin{pmatrix} 2 & 0 \\ 0 & 1 \end{pmatrix}
    = \begin{pmatrix} 3+2i & -1 \\ -1 & 2+i \end{pmatrix}.
\]

\emph{At $\omega = 2$:} the weights are $1,\,2i,\,-4$, so
\[
  \mat{A}(2)
    = \Kop + 2i\Cop - 4\Mop
    = \begin{pmatrix} 4-4 & -1 \\ -1 & 3-4 \end{pmatrix}
      + i\begin{pmatrix} 4 & 0 \\ 0 & 2 \end{pmatrix}
    = \begin{pmatrix} 0+4i & -1 \\ -1 & -1+2i \end{pmatrix}.
\]

Three frequencies, \emph{three different operators}---and all three complex, all
three formed from the same fixed $\{\Kop,\Cop,\Mop\}$ by changing only the scalar
weights. Contrast electrostatics, where the operator is $\Kop$ at \emph{every}
solve. This is corner~1 versus corner~2 in two lines of arithmetic.

\medskip
\emph{One solve.} Take the excitation $\vb = (1,\,0)^{\!\top}$ and solve
$\mat{A}(1)\,\vE = \vb$. With
$\det\mat{A}(1) = (3+2i)(2+i) - (-1)(-1) = (4 + 7i) - 1 = 3 + 7i$,
\[
  \vE
    = \frac{1}{3+7i}
      \begin{pmatrix} 2+i & 1 \\ 1 & 3+2i \end{pmatrix}
      \begin{pmatrix} 1 \\ 0 \end{pmatrix}
    = \frac{1}{3+7i}\begin{pmatrix} 2+i \\ 1 \end{pmatrix}.
\]
Multiplying numerator and denominator by the conjugate $3-7i$ (with
$\abs{3+7i}^2 = 58$),
\[
  \vE
    = \frac{1}{58}\begin{pmatrix} (2+i)(3-7i) \\ (3-7i) \end{pmatrix}
    = \frac{1}{58}\begin{pmatrix} 13 - 11i \\ 3 - 7i \end{pmatrix}
    \approx \begin{pmatrix} 0.224 - 0.190\,i \\ 0.052 - 0.121\,i \end{pmatrix}.
\]
The solved field is genuinely complex: it has magnitude (amplitude) and phase, and
the phase is the steady-state lag of the field behind the drive---information that
a real-valued static solve simply does not carry. In the real driver this $2\times2$
direct solve is replaced by preconditioned GMRES on a million-dimensional complex
operator, but the structure is exactly this: rebuild $\mat{A}(\omega)$, then solve.
\end{workedexample}

\begin{workedexample}{A two-port S-parameter readout}{sparam}
Now the reduction. Suppose we have driven port~1 of a two-port device at one
frequency and solved for the field $\vE_1$. The boundary-mode analysis
(\cref{ch:waveports}) gave us the two port modes $\vs_1,\vs_2$, and the projections
of the solved field onto them came out as
\[
  \inner{\vs_1}{\vE_1} = 0.20 + 0.10\,i,
  \qquad
  \inner{\vs_2}{\vE_1} = 0.70 - 0.30\,i.
\]
Take lumped ports of equal impedance, so the scale is trivial,
$\sigma_{ij} = \sqrt{R_j/R_i} = 1$. Apply \cref{eq:sparam}.

\medskip
\emph{Reflection $S_{11}$} (receiver $1$, drive $1$, so $i=j$---subtract the
self-term):
\[
  S_{11}
    = \sigma_{11}\bigl(\inner{\vs_1}{\vE_1} - \delta_{11}\bigr)
    = (0.20 + 0.10\,i) - 1
    = -0.80 + 0.10\,i.
\]
\emph{Transmission $S_{21}$} (receiver $2$, drive $1$, so $i\neq j$---no
self-term):
\[
  S_{21}
    = \sigma_{21}\bigl(\inner{\vs_2}{\vE_1} - \delta_{21}\bigr)
    = (0.70 - 0.30\,i) - 0
    = 0.70 - 0.30\,i.
\]

\medskip
\emph{Read them.} The magnitudes are
\[
  \abs{S_{11}} = \sqrt{0.80^2 + 0.10^2} \approx 0.806,
  \qquad
  \abs{S_{21}} = \sqrt{0.70^2 + 0.30^2} \approx 0.762.
\]
In decibels, the return loss is
$20\log_{10}\abs{S_{11}} \approx -1.87\,$dB and the insertion loss is
$20\log_{10}\abs{S_{21}} \approx -2.36\,$dB. A return loss this shallow means a
\emph{poor} match---about $65\%$ of the incident power
($\abs{S_{11}}^2 \approx 0.65$) reflects straight back---so this frequency is not
in a passband. Notice the load-bearing role of the self-term: without the
$-\delta_{11}$, $S_{11}$ would have been the raw projection $0.20+0.10\,i$, with
magnitude $\approx 0.22$, reporting a \emph{good} match (low reflection)---the
exact opposite conclusion. The $-1$ is what converts ``total field at the driven
port'' into ``reflected field at the driven port,'' and it is the difference
between a right answer and a plausible wrong one.

\medskip
Sweeping this readout across the band---one solved field, one projection, per
frequency---traces the $\abs{S_{11}}$ and $\abs{S_{21}}$ curves of
\cref{fig:splot}. That sweep is the whole driven driver.
\end{workedexample}

\summarysection
The driven analysis computes the frequency response of a microwave device: feed a
time-harmonic signal into the ports, sweep the frequency band, and read off the
scattering matrix $S(\omega)$. It is the anchor (\cref{ch:electrostatics}) with one
corner swapped. \emph{Assemble} builds a fixed basis $\{\Kop,\Cop,\Mop\}$ once, but
the system operator $\mat{A}(\omega) = \Kop + i\omega\Cop - \omega^2\Mop$
$[+\mat{A}_2(\omega)]$ is a frequency-dependent linear combination of that basis
(\code{assemble\_frequency\_operator}, a specialization of \code{linear\_combination}
with operator operands and complex scalars). \emph{Solve} is the operator-varying
map \code{frequency\_sweep}: at each frequency, rebuild $\mat{A}(\omega)$, re-capture
it into the solver \emph{inside} the loop (the exact negation of electrostatics'
hoist), and solve the complex non-Hermitian system by preconditioned GMRES
(\cref{ch:gmres}). The solves are independent (a map, hence parallel), but each is a
full GMRES, so preconditioner reuse across nearby frequencies matters. \emph{Reduce}
is the rank-2 port-projection \code{sparameter\_reduce}:
$S_{ij} = \sigma_{ij}(\inner{\vs_i}{\vE_j} - \delta_{ij})$---project each solved
field onto each port mode (from the boundary-mode analysis, \cref{ch:waveports}),
subtract the self-reflection on the diagonal, scale for impedance or de-embedding.
It produces a matrix like the Gram but is deliberately \emph{not} a Gram: a linear
projection of two distinct families, complex, asymmetric. The whole driver is
$\texttt{driven} = \texttt{sparameter\_reduce}\circ\texttt{frequency\_sweep}\circ
\texttt{fe\_assemble}^{\times3}$, and reading $\abs{S_{11}}$ (return loss) and
$\abs{S_{21}}$ (insertion loss) across the band is reading the device.

\furtherreading
The scattering-parameter formalism, ports, impedance normalization, and the
network-analyzer measurement this chapter simulates are the foundation of
microwave engineering; Pozar's text~\citep{pozar2011} is the standard treatment and
the place to deepen the physical reading of $S(\omega)$, return loss, insertion
loss, and de-embedding. The finite-element formulation of the time-harmonic
curl--curl system on $\Hcurl$, including wave-port boundary conditions and the modal
projection that defines the port modes $\vs_i$, is developed at length by
Jin~\citep{jin2014}. The complex, non-Hermitian, indefinite linear solve at each
frequency---and the GMRES/FGMRES machinery that handles it, together with
preconditioner reuse across a parameter sweep---is the subject of Saad's account of
iterative methods~\citep{saad2003}, on which \cref{ch:gmres} and \cref{ch:precond}
both draw. The boundary-mode analysis that produces the port modes is
\cref{ch:waveports}; the state-generated fold that the adaptive sweep becomes is the
shape of \cref{ch:transient}.

\exercisessection
\begin{enumerate}[nosep]
  \item \textbf{The one moved line.} In one sentence, name the single operation
  whose position---inside the frequency loop versus outside it---is the entire
  structural difference between the driven sweep and electrostatics. Explain why
  hoisting it out of the driven loop changes the \emph{answer}, not merely the
  speed.
  \item \textbf{Form the operator.} Using the basis of
  \cref{wex:opvary} ($\Kop,\Cop,\Mop$ as given, $\mat{A}_2\equiv0$), write down
  $\mat{A}(\omega)$ at $\omega = \tfrac12$ and at $\omega = 3$. Confirm that the two
  operators differ and that both are complex. Which weight is responsible for the
  imaginary part, and which for making the operator indefinite at large $\omega$?
  \item \textbf{Solve at a frequency.} With $\mat{A}(1)$ from \cref{wex:opvary} and
  the excitation $\vb = (0,\,1)^{\!\top}$, solve $\mat{A}(1)\vE = \vb$ by hand
  (the $2\times2$ inverse is given in the example). Report the magnitude and phase
  of each component of $\vE$.
  \item \textbf{The self-term.} In \cref{wex:sparam}, recompute $S_{11}$ and its
  magnitude \emph{without} the $-\delta_{11}$ self-term. Show that omitting it flips
  the physical conclusion from ``poor match'' to ``good match,'' and explain in one
  sentence what the $-1$ physically subtracts.
  \item \textbf{Projection, not Gram.} Give two structural reasons the S-parameter
  matrix is not a Gram matrix (\cref{ch:reductions}). For each reason, name the
  feature of \code{sparameter\_reduce}'s body that a Gram reduction would lack or
  forbid.
  \item \textbf{Why GMRES.} The driven system operator
  $\mat{A}(\omega) = \Kop + i\omega\Cop - \omega^2\Mop$ is solved by GMRES, not CG.
  Identify which two terms break the symmetric-positive-definite structure CG
  requires, and say what each one does to the operator (one makes it complex
  non-Hermitian; one makes it indefinite).
  \item \textbf{Parallelism and cost.} The frequency sweep is a map, so its solves
  are independent and parallel---yet the chapter calls each one expensive. Explain
  the tension: what can the driven sweep \emph{not} amortize across the family that
  electrostatics \emph{can}, and name the two mitigations the chapter offers.
  \item \textbf{Place it on the grid.} Using the driver $\times$ framework grid of
  \cref{ch:situating}, describe the driven analysis as ``electrostatics with these
  cells changed.'' Exactly which stage-cells (function space, assemble, solve
  corner, reduction) differ, and which stay the same?
\end{enumerate}

\chapter{Transient Simulations}
\label{ch:transient}

% local: velocity unknown in the first-order recast (no \vv collision; \vv is defined in simbook)

\chapterorient{We return to the anchor of \cref{ch:electrostatics} one last time,
and this time we swap the corner that the rest of Part~V has been circling. The
electrostatic and driven analyses solve a \emph{family} of problems---a fan of
independent solves, a \texttt{map}, parallel over terminals or frequencies. The
transient analysis does something the other four cannot: it \emph{marches}. Each
timestep begins exactly where the last one ended, so the solves form a chain, a
\texttt{fold}, that no machine on earth can run out of order. This is the one
genuinely sequential driver of the simulator. We assemble the three operators
$\Kop,\Cop,\Mop$ once, build an implicit ODE integrator once, and then thread a
single field-state forward through time, one opaque step at a time---and out
falls not a matrix but a \emph{movie}, the whole time history of the field. The
machinery of one step we built in \cref{ch:time}; here we compose it into the
driver and draw, as vividly as we can, the contrast between the map drivers and
this fold.}

\section{The physical question}
\label{sec:tr-question}

An engineer hands us a device and a \emph{pulse}: ``inject this current burst and
show me what happens.'' Not a steady voltage, not a single tone---a short,
broadband packet of current $\vJ(t)$ that arrives, excites the structure, and
rings down. The question is not ``what is the capacitance'' or ``what is the
response at \SI{5}{GHz}''; it is \emph{``what does the field do, frame by frame,
as time runs?''} The answer is a time history: the fields, the port voltages and
currents, the stored energy, each sampled at every instant of a simulated
interval. A movie, not a snapshot (\cref{fig:tr-setup}).

\begin{physicsbox}
A current pulse $\vJ(t)$---a modulated wave packet, a step, an impulse---is
injected into a device at $t=t_0$. The fields it launches propagate, reflect off
boundaries, couple into ports, and decay as energy bleeds away through loss and
radiation. At each instant $t_n$ the simulation holds a complete field-state: the
electric field $\vE(\vx,t_n)$ and magnetic field $\vB(\vx,t_n)$ everywhere in the
domain. The output is the whole sequence
$\vE(\cdot,t_0),\vE(\cdot,t_1),\vE(\cdot,t_2),\dots$---a time-march filmstrip, one
frame per timestep. Nothing here is frozen ($\partial_t\neq0$) and nothing is
assumed single-frequency: time is kept in full.
\end{physicsbox}

\begin{figure}[t]
\centering
\begin{tikzpicture}[font=\footnotesize, scale=1.0]
  % the drive pulse on the left
  \begin{scope}[shift={(-4.4,0)}]
    \draw[->, simGray] (-0.3,0)--(1.7,0) node[right,font=\scriptsize]{$t$};
    \draw[->, simGray] (0,-0.5)--(0,0.95);
    \draw[simRust, very thick, domain=-0.2:1.5, samples=60, smooth]
      plot (\x, {0.75*exp(-9*(\x-0.55)^2)*cos(deg(13*(\x-0.55)))});
    \node[font=\scriptsize, text=simRust] at (0.7,-0.85) {pulse $\vJ(t)$};
  \end{scope}
  \draw[flow] (-2.5,0) -- (-1.7,0);
  % three field snapshots t0, t1, t2 ...
  \foreach \k/\dx/\amp in {0/0/0.15, 1/2.4/0.55, 2/4.8/0.32}{
    \begin{scope}[shift={(\dx,0)}]
      \draw[simGray, semithick, rounded corners=2pt] (-0.85,-1.05) rectangle (0.85,1.05);
      % wavy field inside, amplitude grows then shrinks
      \draw[simBlue, thick, domain=-0.8:0.8, samples=40, smooth]
        plot (\x, {\amp*sin(deg(7*\x))});
      \draw[simBlue, thick, domain=-0.8:0.8, samples=40, smooth]
        plot (\x, {-\amp*sin(deg(7*\x))});
      \node[font=\scriptsize, anchor=north] at (0,-1.1) {$t_\k$};
    \end{scope}}
  \node[font=\scriptsize] at (6.1,0) {$\cdots$};
  \node[font=\scriptsize\itshape, text=simGreen] at (2.4,1.5)
    {the field evolves: a frame per timestep};
\end{tikzpicture}
\caption[The transient setup---a pulse and a filmstrip]{The transient analysis
in one picture. A current pulse $\vJ(t)$ (rust, left) is injected into the device;
the fields it launches evolve in time, and the simulation records a complete
field-state snapshot $(\vE,\vB)$ at every timestep $t_0,t_1,t_2,\dots$ (blue, the
field amplitude rising as the pulse arrives and decaying as energy leaves). The
product is not a single number but the \emph{whole filmstrip}---the time history
of the field. This is the analysis of choice for broadband excitation, for
nonlinear materials, or simply for watching a wave propagate.}
\label{fig:tr-setup}
\end{figure}

When do you reach for it? Three situations recommend the transient analysis over
its frequency-domain cousins. First, \emph{broadband excitation}: a single pulse
contains a whole band of frequencies at once (a narrow pulse in time is a wide
band in frequency), so one transient run can recover the response across an entire
band---where the driven analysis of \cref{ch:driven} sweeps one frequency per
solve. Second, \emph{nonlinearity}: a material whose properties depend on the
field amplitude has no single-frequency steady state to solve for, and only a
time-march can follow it. Third, simply \emph{watching a wave move}: when the
physical insight you want is spatial and temporal---a pulse crossing a junction, a
mode building up in a cavity---the movie \emph{is} the answer.

The transient analysis is the time-domain complement of the driven analysis. The
two compute the same physics by two routes of the harmonic rule
(\cref{ch:reductions-pde}): the driven analysis works one frequency at a time; the
transient analysis works all frequencies at once, through a pulse, and recovers
the spectrum afterward by a Fourier transform. We will close the chapter on
exactly that bridge (\S\ref{sec:tr-fourier}).

\section{Assemble, once: the three operators and the integrator}
\label{sec:tr-assemble}

The transient pipeline runs the same skeleton as the anchor
(\cref{ch:situating}),
\[
  \texttt{config} \to \underbrace{\texttt{assemble}}_{\S\ref{sec:tr-assemble}}
  \to \underbrace{\texttt{solve}}_{\S\ref{sec:tr-solve}}
  \to \underbrace{\texttt{reduce}}_{\S\ref{sec:tr-reduce}} \to \texttt{product},
\]
and the assemble stage, as always, runs \emph{once}, before any marching begins.
What differs from electrostatics is only \emph{how many} operators it builds and
\emph{what} it hands them to.

\subsection{Three operators on \texorpdfstring{$\Hcurl$}{H(curl)}}

The governing physics is the time-domain wave equation we derived in
\cref{ch:reductions-pde} (\S\ref{sec:transient}). Discretizing it in space only
(the method of lines) turns Maxwell's time-domain system into a \emph{second-order}
semi-discrete system of ordinary differential equations in the
degree-of-freedom vector $\vs(t)$:
\begin{equation}
  \boxed{\;\Mop\,\ddot\vs + \Cop\,\dot\vs + \Kop\,\vs = \vJ(t)\;}
  \label{eq:tr-secondorder}
\end{equation}
the damped-oscillator system of \cref{eq:transient-ode}. Here $\Mop$ is the
$\varepsilon$-weighted \emph{mass}, $\Cop$ the $\sigma$-weighted \emph{damping},
$\Kop$ the curl--curl \emph{stiffness}, and $\vJ(t)$ the discretized forcing
$-\partial_t\vJ_{\text{src}}$. The unknown is the electric field, a vector field,
so it lives in the Nédélec edge space $\Hcurl$ of \cref{ch:functionspaces}---the
same space as the driven, magnetostatic, and eigenmode analyses, and the swap from
electrostatics' scalar $\Hone$ that \cref{tab:es-swaps} recorded.

Where electrostatics assembled one operator, transient assembles \emph{three}, by
three calls to the same assembly fold \code{fe\_assemble} (\cref{ch:fem}), each
over its own one-term weak form:
\begin{lstlisting}[style=l4style]
-- assemble the three time-domain operators, each a one-term fe_assemble fold; ALL once
space = nd_space cfg                              -- the Nedelec H(curl) FE space (read-only)
k = fe_assemble space [ curl_curl  (reluctivity   cfg) ]  -- (1) stiffness K  (nu * curl)
c = fe_assemble space [ damping     (conductivity  cfg) ]  --     damping  C  (sigma * .)
m = fe_assemble space [ mass        (permittivity  cfg) ]  --     mass     M  (epsilon * .)
\end{lstlisting}
This is the $\Kop/\Mop/\Cop$ recurrence of \cref{ch:reductions-pde} made concrete:
the three operators are the fixed, reusable raw material, and every analysis is a
recipe for recombining them. The transient recipe is
\cref{eq:tr-secondorder}.\footnote{In Palace the three matrices are assembled once
inside the \code{TimeOperator} constructor:
\code{K = space\_op.GetStiffnessMatrix(...)},
\code{C = space\_op.GetDampingMatrix(...)},
\code{M = space\_op.GetMassMatrix(...)}
(\code{palace/models/timeoperator.cpp:65-67}). All three are built before the time
loop and held unchanged for the entire march.}

\subsection{Build the integrator once, and seed the state}

The assemble stage does two more things, both once. First it constructs the
\emph{ODE integrator}---the routine that will advance the state one step at a time.
As \cref{ch:time} explained at length, the semi-discrete system
\cref{eq:tr-secondorder} is stiff (its largest eigenvalue is set by the smallest
mesh element), so an explicit integrator would be pinned to absurdly tiny steps,
and Palace integrates \emph{implicitly}. It recasts \cref{eq:tr-secondorder} as a
first-order system in the doubled state $\vy=(\vs,\dot\vs)$ and hands it to a
library integrator---generalized-$\alpha$, SDIRK, or an ARKStep---chosen by a
config token and constructed \emph{once}, outside the loop
(\cref{sec:integrators}).\footnote{Palace builds a
\code{TimeDependentFirstOrderOperator} from $\Kop,\Cop,\Mop$ and constructs the
integrator at \code{palace/models/timeoperator.cpp:311-335}, before the driver's
time loop at \code{palace/drivers/transientsolver.cpp:77}.} Second it seeds the
initial field-state $\vy_0$---typically rest, $\vs_0=\dot\vs_0=\mathbf0$---the value
the fold will thread forward.

\begin{keyidea}
The transient assemble stage builds, \emph{once}, everything the march will reuse:
the three operators $\Kop,\Cop,\Mop$ on $\Hcurl$ (by three \code{fe\_assemble}
folds), the implicit ODE integrator (constructed from those operators, the step
machinery of \cref{ch:time}), and the seed field-state $\vy_0$. None of these
changes during the march---the operator capture is hoisted out of the time loop,
exactly as the electrostatic operator was hoisted out of the terminal loop
(\cref{ch:electrostatics}). The hoist is the same; what differs is the loop it sits
outside of.
\end{keyidea}

\section{Solve: the state-threaded fold}
\label{sec:tr-solve}

Here is the corner this chapter exists to teach. With the operators assembled, the
integrator built, and the state seeded, the solve stage marches. And the march is
\emph{not} a map. It is a \texttt{fold}: the one solve corner of \cref{ch:situating}
in which each solve begins from the previous solve's output. This is the structural
centerpiece of the chapter, and the rest of Part~V has been pointing at it.

\subsection{The crucial contrast: a fold, not a map}

Recall the fundamental split of \cref{ch:combinators}, which the anchor chapter
leaned on (\cref{fig:es-solvefamily}). A \texttt{map} runs a family of
\emph{independent} computations: no carry passes between elements, so they
reorder freely and run in parallel. A \texttt{fold} threads a \emph{carry} through
a chain: each step consumes what the last produced, so the steps are strictly
ordered. The electrostatic and driven solve corners are maps---a fan of
independent solves over a family of terminals or frequencies. The transient solve
corner is a fold, and the difference is the whole point.

In the transient march the field-state at time $t_{n+1}$ is computed \emph{from}
the field-state at time $t_n$ by one timestep. Step $n+1$ cannot begin until step
$n$ has finished, because step $n$'s output \emph{is} step $n+1$'s input. The state
is threaded forward along a single timeline. You cannot solve the timesteps in
parallel, in any order, or on separate machines, because there is no ``step 5'' to
hand a machine until ``step 4'' has produced the state it starts from. This is the
\emph{state-threaded fold}---the combinator \code{fold\_solve}---and the transient
march is its purest and primary witness. \Cref{fig:tr-foldvmap} sets it against the
electrostatic map, the two side by side.

\begin{figure}[p]
\centering
\begin{tikzpicture}[font=\footnotesize]
  % ===== LEFT: the map drivers (electrostatic / driven) =====
  \begin{scope}
    \node[font=\small\bfseries, text=simBlue] at (2.1,3.0)
      {\textsc{map} drivers: independent, parallel};
    \node[font=\scriptsize\itshape, text=simBlue] at (2.1,2.6)
      {electrostatic, driven};
    \node[opbox, minimum width=13mm, fill=simBgBlue, draw=simBlue] (K) at (2.1,1.7)
      {operator};
    \foreach \i/\x in {1/0.3,2/1.5,3/2.7,4/3.9}{
      \node[data, minimum width=7mm, minimum height=5mm] (b\i) at (\x,0.75) {$\vb_{\i}$};
      \node[product, minimum width=7mm, minimum height=5mm] (x\i) at (\x,-0.3) {$\vx_{\i}$};
      \draw[flow, shorten >=1pt] (K.south) -- (b\i.north);
      \draw[flow, shorten >=1pt] (b\i.south) -- (x\i.north);
    }
    \node[font=\scriptsize, text=simGreen, align=center] at (2.1,-1.1)
      {NO arrows between solves:\\ reorderable, embarrassingly parallel};
  \end{scope}
  \draw[simGray!55, very thick] (4.9,-1.5) -- (4.9,3.2);
  % ===== RIGHT: the transient fold =====
  \begin{scope}[xshift=5.6cm]
    \node[font=\small\bfseries, text=simRust] at (2.4,3.0)
      {\textsc{fold} driver: chained, sequential};
    \node[font=\scriptsize\itshape, text=simRust] at (2.4,2.6)
      {transient (the one fold driver)};
    \node[data, minimum width=8mm] (y0) at (-0.25,1.5) {$\vy_0$};
    \foreach \i/\x in {1/1.0,2/2.5,3/4.0}{
      \node[opbox, minimum width=11mm] (g\i) at (\x,1.5) {step};
    }
    \node[data, minimum width=8mm] (y3) at (5.25,1.5) {$\vy_3$};
    \draw[flow] (y0.east) -- (g1.west);
    \draw[flow] (g1.east) -- node[above,font=\scriptsize]{$\vy_1$} (g2.west);
    \draw[flow] (g2.east) -- node[above,font=\scriptsize]{$\vy_2$} (g3.west);
    \draw[flow] (g3.east) -- (y3.west);
    % each step hides a ksp_solve
    \foreach \i/\x in {1/1.0,2/2.5,3/4.0}{
      \node[extern, minimum width=11mm, minimum height=5mm, font=\scriptsize]
        (k\i) at (\x,0.3) {\code{ksp\_solve}};
      \draw[refedge] (g\i.south) -- (k\i.north);
    }
    \node[font=\scriptsize, text=simRust, align=center] at (2.4,-0.85)
      {horizontal carry $\vy_n\!\to\!\vy_{n+1}$:\\ strictly ordered---cannot overlap timesteps};
  \end{scope}
\end{tikzpicture}
\caption[The map drivers versus the transient fold]{The centerpiece contrast. On
the \emph{left}, the map drivers (electrostatic, driven): one operator feeds a fan
of independent solves $\vb_i\mapsto\vx_i$, with \emph{no} arrows between them---they
are reorderable, parallelizable, embarrassingly parallel over the family. On the
\emph{right}, the transient fold: the doubled state $\vy_0\to\vy_1\to\vy_2\to\vy_3$
is threaded \emph{forward}, each step consuming the prior state, so the steps are
strictly ordered and \emph{cannot} overlap. The horizontal carry arrows---present on
the right, absent on the left---are the whole difference: \code{fold\_solve} versus
\code{solve\_family}. Each fold step \emph{additionally} hides an implicit
\code{ksp\_solve} (\cref{ch:krylov}), but the implicit solve is not what makes it a
fold; the \emph{carry} is.}
\label{fig:tr-foldvmap}
\end{figure}

\begin{keyidea}
\textbf{Transient is the one sequential driver: a fold where each step begins where
the last ended---the opposite of the map drivers.} Electrostatic and driven solve a
\emph{family} of independent problems (a \texttt{map}: reorderable, parallel over
terminals or frequencies). Transient threads a \emph{single field-state} forward
through time (a \texttt{fold}: each step's input is the prior step's output). The
presence or absence of the carry between solves is the entire structural
difference, and it is the difference between a contact sheet of independent
snapshots and a movie in which every frame depends on the last. No amount of
hardware can parallelize the transient march across timesteps, because step $n+1$
literally has nothing to compute until step $n$ is done.
\end{keyidea}

\subsection{\texttt{fold\_solve}, in the calculus}

In the vocabulary of \cref{ch:combinators}, the march is one expression. Capture
the operator once, seed the state once, and \code{foldl} the opaque per-step
operator over the schedule of times:
\begin{lstlisting}[style=l4style]
-- the state-threaded march: foldl the per-step operator over the fixed schedule
fold_solve :: OpParams -> TimeState -> [Time] -> TimeState
fold_solve op s0 schedule = foldl (\s t -> time_step_op op s t) s0 schedule

-- the OPAQUE per-step operator the fold quantifies over:
-- advances the field-state one step; bottoms out in the library integrator
time_step_op :: OpParams -> TimeState -> Time -> TimeState
\end{lstlisting}
Read the signature and the whole structure of the march is in it. The seed
\code{s0} and the result are both a \code{TimeState}---the field-state carried
forward. The operator \code{op} (the captured $\Kop,\Cop,\Mop$ and the chosen
integrator) is bound \emph{once}, before the fold, and threaded unchanged into every
step: the construction hoist of \S\ref{sec:tr-assemble} made structural by the type.
And the step \lstinline[style=l4style]|time_step_op op s t| \emph{reads} \code{s},
the prior carry---\emph{that read} is what makes this a fold and forbids reordering
at the type level. A map's step would ignore the accumulator; a fold's reads it.

\code{time\_step\_op} is deliberately \emph{opaque}. It bottoms out in the library
integrator's one step---an MFEM \code{ODESolver::Step} call---which internally does
the implicit linear solve that \cref{ch:time} dissected
(\cref{fig:tr-innersolve}).\footnote{The fold body is \code{time\_op.Step(t,
delta\_t)} (\code{palace/drivers/transientsolver.cpp:93}), which dispatches to
\code{ode->Step(sol, t, dt)} (\code{palace/models/timeoperator.cpp:410})---the
opaque MFEM step advancing the persistent \code{sol} field-state in place. The prior
step's \code{sol} is the next step's input: the genuine \code{foldl}.} The
combinator only quantifies over that step; it does not look inside it.
\code{fold\_solve} owns the loop, and the library owns one step.

\begin{figure}[t]
\centering
\begin{tikzpicture}[font=\footnotesize]
  % the fold spine
  \node[data, minimum width=8mm] (sn) at (0,0) {$\vy_n$};
  \node[opbox, minimum width=26mm, minimum height=16mm, fill=simBgRust, draw=simRust,
    align=center] (step) at (3.0,0) {\code{time\_step\_op}\\[1pt]\scriptsize one implicit step};
  \node[data, minimum width=10mm] (sn1) at (6.2,0) {$\vy_{n+1}$};
  \draw[flow] (sn) -- (step);
  \draw[flow] (step) -- (sn1);
  % the inner ksp_solve, blown up below
  \node[extern, minimum width=30mm, minimum height=10mm, align=center]
    (ksp) at (3.0,-2.2) {\scriptsize MFEM \code{ODESolver::Step}\\[-1pt]\scriptsize $(\mat I-\Delta t\,\mat J)\,\vy_{n+1}=\text{rhs}$\\[-1pt]\scriptsize one \code{ksp\_solve}};
  \draw[refedge] (step.south) -- (ksp.north)
    node[midway, right, font=\scriptsize]{hides};
  \node[font=\scriptsize\itshape, text=simGray, align=left, anchor=west] at (-1.6,-2.2)
    {\emph{inside} each};
  \node[font=\scriptsize\itshape, text=simGray, align=left, anchor=west] at (-1.6,-2.55)
    {opaque step:};
\end{tikzpicture}
\caption[The implicit solve inside each fold step]{Each fold step is opaque, but it
is not empty: inside \code{time\_step\_op} the integrator solves one linear system.
Because the integrator is \emph{implicit} (the cure for the stiffness of
\cref{ch:time}), advancing $\vy_n\to\vy_{n+1}$ requires solving
$(\mat I-\Delta t\,\mat J)\,\vy_{n+1}=\text{rhs}$ for the unknown next state---a
single \code{ksp\_solve}, the Krylov solve of \cref{ch:krylov}, nested in the march.
The fold quantifies over this step; the library performs it. Note the two distinct
costs the fold carries: the \emph{carry} (which makes the march sequential) and the
\emph{implicit solve} (which makes each step expensive). They are independent---a
fold could be cheap per step, and a map could be expensive per step---but the
transient march has both.}
\label{fig:tr-innersolve}
\end{figure}

\subsection{The pipeline in one line}

The two stages compose into a single expression. At the framework level the whole
transient feature is
\begin{lstlisting}[style=l4style]
transient = fold_solve . fe_assemble
\end{lstlisting}
read right to left: \code{fe\_assemble} builds the operators once, and
\code{fold\_solve} threads the time-march fold over them. Spelled out with the data
flowing through it,
\begin{lstlisting}[style=l4style]
-- inputs = config; output = the time-domain field-state trajectory (the product)
transient :: TransientConfig -> FieldTrajectory
transient cfg =
  let space    = nd_space cfg                             -- the H(curl) FE space
      (k,c,m)  = ( fe_assemble space [ curl_curl (reluctivity cfg) ]   -- (1) K, C, M
                 , fe_assemble space [ damping    (conductivity cfg) ]  --     assembled
                 , fe_assemble space [ mass       (permittivity cfg) ] )--     ONCE
      op       = time_operator (k,c,m) (dJdt cfg)         -- captured ODE operator (read-only)
      s0       = init_state cfg                           -- seed field-state (rest)
      schedule = uniform_steps (delta_t cfg) (n_step cfg) -- the FIXED [Time] schedule
  in  fold_solve op s0 schedule                           -- (2) the state-threaded march FOLD
\end{lstlisting}
\Cref{fig:tr-pipeline} draws it. Set it beside the anchor diagram
(\cref{fig:es-pipeline}) and the swap is exactly the one \cref{tab:es-swaps}
promised: the FE space moves from $\Hone$ to $\Hcurl$, the solve corner moves from
the \code{solve\_family} \emph{map} to the \code{fold\_solve} \emph{fold}, and the
reduction moves from the Gram to a trajectory table. The skeleton is identical; one
corner is swapped.

\begin{figure}[p]
\centering
\begin{tikzpicture}[node distance=6mm and 6mm, font=\footnotesize]
  \node[data, align=center, text width=18mm] (cfg)
    {config\\[1pt]\footnotesize $\varepsilon,\sigma,\nu$,\\mesh, pulse};
  \node[stage, right=of cfg, align=center, text width=21mm] (asm)
    {\textbf{assemble}\\[1pt]\footnotesize \code{fe\_assemble}\\$\Kop,\Cop,\Mop$ (once)};
  \node[stage, right=of asm, align=center, text width=22mm, fill=simBgRust, draw=simRust]
    (solve) {\textbf{solve}\\[1pt]\footnotesize \code{fold\_solve}\\(time march)};
  \node[stage, right=of solve, align=center, text width=18mm, fill=simBgGold, draw=simGold]
    (red) {\textbf{reduce}\\[1pt]\footnotesize per-step\\readout};
  \node[product, right=of red, align=center, text width=17mm] (prod)
    {trajectory\\[1pt]\footnotesize field history};
  \draw[flow] (cfg) -- (asm);
  \draw[flow] (asm) -- node[above,font=\scriptsize]{$\Kop,\Cop,\Mop$} (solve);
  \draw[flow] (solve) -- node[above,font=\scriptsize]{$\{\vy_n\}$} (red);
  \draw[flow] (red) -- (prod);
  % the self-loop on the solve box: the carry
  \draw[backflow, shorten >=1pt] (solve.north) ++(-0.4,0)
    .. controls +(0,7mm) and +(0,7mm) .. ($(solve.north)+(0.4,0)$)
    node[midway, above, font=\scriptsize, text=simRust]
      {the carry $\vy_n\!\to\!\vy_{n+1}$};
  % corner annotations
  \node[font=\scriptsize\itshape, text=simRust, below=9mm of solve, align=center]
    {state-threaded fold\\(\cref{ch:situating}, corner 3)};
  \node[font=\scriptsize\itshape, text=simGold, above=8mm of red, align=center]
    {trajectory, not a matrix};
\end{tikzpicture}
\caption[The transient pipeline]{The complete transient pipeline, the composition
$\code{fold\_solve}\circ\code{fe\_assemble}$. A configuration of
materials, mesh, and excitation pulse enters; \code{fe\_assemble} builds the three
operators $\Kop,\Cop,\Mop$ \emph{once} (and, with them, the integrator and seed
state); \code{fold\_solve} threads the time-march fold, producing the field-state
trajectory $\{\vy_n\}$; and a per-step readout reduces each state to the physical
products. The rust self-loop on the solve box is the feature special to this
analysis: the \emph{carry}, $\vy_n\to\vy_{n+1}$, that makes the march sequential.
Compared with the anchor (\cref{fig:es-pipeline}), exactly one corner is swapped:
the \code{solve\_family} map becomes the \code{fold\_solve} fold.}
\label{fig:tr-pipeline}
\end{figure}

\subsection{Checkpoint and resume: the schedule-split property}

The fold's sequentiality sounds like pure cost, but it comes with a genuine
compensation that the map cannot offer: \emph{restartability}. A left fold over a
concatenated schedule satisfies a clean algebraic law (\cref{ch:time},
\cref{eq:schedule-split}). Split the schedule into a prefix $a$ and a suffix $b$;
then folding the whole is the same as folding the prefix and folding the suffix
\emph{from the state the prefix produced}:
\begin{equation}
  \code{fold\_solve}\ \code{op}\ \vy_0\ (a \mathbin{+\!\!+} b)
  \;=\;
  \code{fold\_solve}\ \code{op}\ \bigl(\code{fold\_solve}\ \code{op}\ \vy_0\ a\bigr)\ b .
  \label{eq:tr-split}
\end{equation}
Read operationally, this says: run the prefix, save the final state $\vy_{|a|}$ to
disk, stop; later, load it and run the suffix seeded from it. The law guarantees the
result is the \emph{identical} trajectory---bit for bit---as the uninterrupted run,
because the only thing that crosses the split point is the carry, and the carry is
exactly what we saved (\cref{fig:tr-checkpoint}). Long runs survive faults; expensive
runs split across job-scheduler windows; none of it perturbs the answer. The same
single threaded carry that \emph{forbids} parallelism across timesteps is what makes
a checkpoint a single saved value rather than a replay of the whole history. A map
parallelizes but has no notion of ``resume''---there is no carry to save.

\begin{figure}[t]
\centering
\begin{tikzpicture}[font=\footnotesize]
  % top: uninterrupted
  \node[font=\scriptsize\bfseries, text=simInk, anchor=west] at (-0.3,1.9)
    {one uninterrupted march};
  \node[data, minimum width=8mm] (u0) at (0,1.2) {$\vy_0$};
  \foreach \i/\x in {1/0.95,2/1.95,3/2.95,4/3.95}{
    \node[opbox, minimum width=7mm, minimum height=5mm] (us\i) at (\x,1.2) {};}
  \node[data, minimum width=8mm] (u4) at (4.95,1.2) {$\vy_4$};
  \draw[flow] (u0)--(us1);\draw[flow] (us1)--(us2);\draw[flow] (us2)--(us3);
  \draw[flow] (us3)--(us4);\draw[flow] (us4)--(u4);
  \node[font=\scriptsize, text=simGray] at (2.5,0.72) {schedule $a \mathbin{+\!\!+} b$};
  % bottom: split
  \node[font=\scriptsize\bfseries, text=simRust, anchor=west] at (-0.3,-0.1)
    {split: run $a$, save carry, resume $b$};
  \node[data, minimum width=8mm] (b0) at (0,-0.8) {$\vy_0$};
  \foreach \i/\x in {1/0.95,2/1.95}{
    \node[opbox, minimum width=7mm, minimum height=5mm] (bs\i) at (\x,-0.8) {};}
  \node[product, minimum width=9mm, fill=simBgGold, draw=simGold] (ckpt) at (2.95,-0.8) {$\vy_2$};
  \foreach \i/\x in {3/3.95,4/4.95}{
    \node[opbox, minimum width=7mm, minimum height=5mm] (bs\i) at (\x,-0.8) {};}
  \node[data, minimum width=8mm] (b4) at (5.95,-0.8) {$\vy_4$};
  \draw[flow] (b0)--(bs1);\draw[flow] (bs1)--(bs2);\draw[flow] (bs2)--(ckpt);
  \draw[flow] (ckpt)--(bs3);\draw[flow] (bs3)--(bs4);\draw[flow] (bs4)--(b4);
  \node[font=\scriptsize, text=simRust] at (1.4,-1.3) {fold $a$};
  \node[font=\scriptsize, text=simRust] at (4.4,-1.3) {fold $b$ from $\vy_2$};
  \node[font=\scriptsize, text=simGold!75!black, align=center] at (2.95,-1.75)
    {save / load (checkpoint)};
  \node[font=\Large, text=simGray] at (5.6,0.2) {$=$};
\end{tikzpicture}
\caption[Checkpoint and resume as a schedule split]{The schedule-split law
\cref{eq:tr-split} drawn as checkpoint-and-resume. \emph{Top}: the uninterrupted
march folds the whole schedule $a\mathbin{+\!\!+}b$ from $\vy_0$ to $\vy_4$.
\emph{Bottom}: the same march split at the midpoint---fold the prefix $a$ to the
carry $\vy_2$, \emph{save it} (the gold checkpoint), stop, then later load $\vy_2$
and fold the suffix $b$ from it. The two produce the identical $\vy_4$, because the
only thing crossing the cut is the carry. Restartability is not bolted onto the
fold; it \emph{is} what the fold structure means---a clean algebraic property, not
an engineering approximation.}
\label{fig:tr-checkpoint}
\end{figure}

\section{Reduce: reading the trajectory}
\label{sec:tr-reduce}

The solve stage has handed us not a family of fields but a \emph{trajectory}: the
field-state at every step of the march. The reduce stage reads the physical products
from it. Where the electrostatic reduction collapsed a family into a single
capacitance matrix (\cref{ch:electrostatics}), the transient ``reduction'' is
mostly a per-step \emph{readout}---it runs once per frame, and the output is the
whole time history, not one number.

\subsection{What the driver reads at each step}

At each step the transient driver recovers three kinds of product from the current
field-state:

\begin{description}[style=nextline, leftmargin=0pt, font=\normalfont\itshape]
  \item[The fields.] From each step's state the electric and magnetic fields
  $(\vE,\vB)$ are recovered (the \code{GetE}/\code{GetB} of the time operator) and
  written out---the ``movie'' frames, one per timestep.\footnote{The per-step readout
  is \code{time\_op.GetE()} / \code{time\_op.GetB()}
  (\code{palace/drivers/transientsolver.cpp:98-99}), consumed immediately by the
  per-step postprocess \code{post\_op.MeasureAndPrintAll(...)}
  (\code{:104}).}
  \item[Port voltages and currents.] Integrating the fields against the port modes
  gives a time-domain voltage and current at each terminal: the waveforms an
  engineer reads off an oscilloscope, computed step by step.
  \item[Energy versus time.] The field energy, evaluated at each step, traces how
  energy enters with the pulse, sloshes between electric and magnetic forms, and
  bleeds away through the damping $\Cop$ and the ports. This is the
  \emph{driver-agnostic} energy reduction of \cref{ch:reductions}---the same
  field-energy form the electrostatic diagonal used (\cref{ch:electrostatics}),
  here evaluated once per frame instead of once per terminal.
\end{description}

\Cref{fig:tr-energy} draws a representative trajectory: an injected pulse rings up
the stored energy, which then decays as loss and radiation carry it off. Because
each curve is sampled once per fold step, the resolution of the readout is exactly
the resolution of the march---one frame per timestep.

\begin{figure}[t]
\centering
\begin{tikzpicture}
\begin{axis}[
    width=0.80\linewidth, height=50mm,
    xlabel={time $t$}, ylabel={},
    xmin=0, xmax=10, ymin=-1.15, ymax=1.30,
    axis lines=left, xtick={0,2,4,6,8,10}, ytick={-1,0,1},
    yticklabel style={font=\scriptsize}, xticklabel style={font=\scriptsize},
    xlabel style={font=\scriptsize},
    legend style={font=\scriptsize, draw=simGray!50, at={(0.98,0.98)},
      anchor=north east, legend cell align=left},
    every axis plot/.append style={very thick, samples=120, domain=0:10},
  ]
  \addplot[simRust] {0.9*exp(-2*(x-1.5)^2)*cos(deg(6*(x-1.5)))};
  \addlegendentry{drive $\vJ(t)$ (a pulse)}
  \addplot[simBlue] {0.8*(1-exp(-3*x))*exp(-0.35*x)*cos(deg(4*x-1))};
  \addlegendentry{field $E(t)$ (rings, decays)}
  \addplot[simGreen, densely dashed] {0.95*(1-exp(-2.2*x))*exp(-0.5*x)};
  \addlegendentry{stored energy (rise then decay)}
\end{axis}
\end{tikzpicture}
\caption[A transient field and energy trajectory]{A representative readout of the
transient march. The pulse drive $\vJ(t)$ (rust) injects a short burst; a field
component $E(t)$ (blue) rings up as the pulse arrives and decays as energy leaves;
the stored energy (green, dashed) rises with the injection and decays monotonically
thereafter, bled away by the damping $\Cop$ and carried off through the ports. Each
curve is sampled once per fold step: the readout resolution is the march resolution,
one frame per timestep. The output is this \emph{whole history}, not a single
reduced number.}
\label{fig:tr-energy}
\end{figure}

\section{The bridge to the frequency domain}
\label{sec:tr-fourier}

One readout deserves its own section, because it ties the transient analysis back
to the driven analysis of \cref{ch:driven} and closes a loop the whole of Part~V has
left open. A transient run injects a \emph{broadband} pulse---a superposition of
many frequencies at once. Record the input waveform and the corresponding output
waveform in time, take the Fourier transform of each, and divide: the quotient is
the frequency response across the \emph{whole band} of the pulse, recovered \emph{in
a single run}.

This is the time-domain route to the same S-parameters the driven analysis computes
frequency by frequency. Where the driven analysis sweeps a \code{frequency\_sweep}
map---one frequency per solve, an independent family (\cref{ch:driven})---the
transient analysis folds one march and Fourier-transforms the result. The two compute
the same physics by the two routes of the harmonic rule
(\cref{ch:reductions-pde}): one frequency at a time, or all frequencies at once
through a pulse. \Cref{fig:tr-fourier} shows a sampled port-voltage signal beside its
transform, a resonance appearing as a peak in the spectrum.

\begin{figure}[t]
\centering
\begin{tikzpicture}
% LEFT: time-domain port signal
\begin{axis}[
    name=tplot,
    width=0.46\linewidth, height=44mm,
    xlabel={time $t$}, ylabel={port voltage $v(t)$},
    xmin=0, xmax=12, ymin=-1.1, ymax=1.1,
    axis lines=left, xtick={0,4,8,12}, ytick={-1,0,1},
    yticklabel style={font=\scriptsize}, xticklabel style={font=\scriptsize},
    xlabel style={font=\scriptsize}, ylabel style={font=\scriptsize},
    title style={font=\scriptsize\itshape}, title={time domain (the march)},
    every axis plot/.append style={very thick, samples=160, domain=0:12},
  ]
  % a decaying ring at a single frequency: the "resonance"
  \addplot[simBlue] {exp(-0.32*x)*cos(deg(5*x))};
\end{axis}
% RIGHT: frequency-domain magnitude with a peak
\begin{axis}[
    at={(tplot.right of east)}, anchor=left of west, xshift=10mm,
    width=0.46\linewidth, height=44mm,
    xlabel={frequency $f$}, ylabel={$|\hat v(f)|$},
    xmin=0, xmax=2, ymin=0, ymax=1.15,
    axis lines=left, xtick={0,0.5,1,1.5,2}, ytick={0,0.5,1},
    yticklabel style={font=\scriptsize}, xticklabel style={font=\scriptsize},
    xlabel style={font=\scriptsize}, ylabel style={font=\scriptsize},
    title style={font=\scriptsize\itshape}, title={frequency domain (FFT)},
    every axis plot/.append style={very thick, samples=160, domain=0:2},
  ]
  % a Lorentzian peak at f0 ~ 0.8 (= 5/(2 pi)), the spectrum of the decaying ring
  \addplot[simRust] {1/sqrt(1 + 80*(x-0.8)^2)};
  \draw[simGray, densely dashed] (axis cs:0.8,0) -- (axis cs:0.8,1.0);
  \node[font=\scriptsize, text=simGray, anchor=south] at (axis cs:0.8,1.0) {$f_0$};
\end{axis}
\end{tikzpicture}
\caption[A port signal and its Fourier transform]{The time$\to$frequency bridge.
\emph{Left}: a sampled port-voltage waveform $v(t)$ from one transient march---a
decaying ring at the structure's resonant frequency. \emph{Right}: its Fourier
transform $|\hat v(f)|$, a peak at the resonance $f_0$ (a slow decay in time is a
sharp peak in frequency---a high-$Q$ resonance). Dividing the transform of the
\emph{output} by that of the \emph{input} pulse recovers the frequency response
across the whole band, in one run. This is the same frequency-domain information the
driven analysis (\cref{ch:driven}) computes one frequency at a time; the transient
fold plus an FFT is its broadband, all-at-once complement.}
\label{fig:tr-fourier}
\end{figure}

\section{Cost and stability}
\label{sec:tr-cost}

Two practical facts about the march, both carried over from the machinery of
\cref{ch:time}, are worth collecting before the worked examples.

\paragraph{Implicit stepping, and the timestep trade.} Each fold step hides an
implicit linear solve precisely because the semi-discrete system is \emph{stiff}: an
explicit step would be pinned to $\Delta t<2/\lambda_{\max}$ by the fastest mesh
mode, a mode physically irrelevant to the answer (\cref{ch:time}). An implicit step
is unconditionally stable, so $\Delta t$ is chosen for \emph{accuracy on the modes we
care about}, not forced small by modes we do not. The trade is direct: a smaller
$\Delta t$ resolves faster physics and reduces truncation error but costs more steps
(more \code{ksp\_solve}s); a larger $\Delta t$ is cheaper but coarser. The integrator's
order (\cref{ch:time}) sets how fast that error falls as $\Delta t$ shrinks.

\paragraph{Reuse the factorization across steps.} For a \emph{uniform} schedule the
step operator $\mat I-\Delta t\,\mat J$ is built from the captured $\Kop,\Cop,\Mop$
and the fixed $\Delta t$, none of which change from step to step---so it is a
\emph{constructed operator}, fixed for the whole march (\cref{ch:operators}). Its
factorization (or its preconditioner) is therefore computed \emph{once} and reused
every step, exactly as the electrostatic operator's factorization was reused across
the terminal family (\cref{ch:electrostatics}). The operator-capture-once hoist of
\code{fold\_solve} is what licenses it: \code{op} is bound before the fold, and every
step reads the same operator. This is the fold's version of the dividend the
fixed-operator map paid---``assemble once'' quietly licenses ``factor once'' here too,
even though the solves are chained rather than independent.

\begin{notationbox}
The transient march carries \emph{two} costs that are easy to conflate but logically
independent. The \emph{carry} makes the march \emph{sequential}: timesteps cannot
overlap, regardless of how cheap each one is. The \emph{implicit solve} makes each
step \emph{expensive}: one \code{ksp\_solve} per step, regardless of whether the
steps are chained. A hypothetical explicit transient march would still be a fold
(sequential) but with cheap steps; an electrostatic map is parallel but with
expensive steps. Transient is the corner that pays both: sequential \emph{and}
per-step-expensive. Keeping the two costs separate is the key to reasoning about it.
\end{notationbox}

\section{Two worked examples}
\label{sec:tr-worked}

The pipeline is general; we now run it on problems small enough to do by hand. The
first marches a two-degree-of-freedom oscillator a few implicit steps, showing each
step solve a small linear system and the state thread forward---and contrasts it
explicitly with the independent map a \code{solve\_family} would run. The second
takes a sampled port signal and shows its DFT recovering a resonance, making the
time$\to$frequency bridge of \S\ref{sec:tr-fourier} numerical.

\begin{workedexample}{Marching a two-DoF oscillator---a fold, not a map}{tr-twodof}
We run the real machinery---the recast, the per-step implicit solve, and the state
threading---on the smallest non-trivial system: a two-degree-of-freedom
mass--spring--damper $\Mop\ddot\vs+\Cop\dot\vs+\Kop\vs=\vJ(t)$, a hand-sized
stand-in for the semi-discrete transient system, with
\[
  \Mop=\begin{psmallmatrix}1&0\\0&1\end{psmallmatrix},\quad
  \Cop=\begin{psmallmatrix}0.2&0\\0&0.2\end{psmallmatrix},\quad
  \Kop=\begin{psmallmatrix}2&-1\\-1&2\end{psmallmatrix},\quad
  \vJ=\begin{psmallmatrix}1\\0\end{psmallmatrix}\ (\text{constant}),
\]
starting from rest, $\vs_0=\vv_0=(0,0)$. The off-diagonal of $\Kop$ couples the two
degrees of freedom.

\medskip
\emph{Recast and the step matrix.} With $\vv=\dot\vs$ the doubled state is
$\vy=(\vs,\vv)$. Take backward Euler with $\Delta t=0.1$. Substituting
$\vs_{n+1}=\vs_n+\Delta t\,\vv_{n+1}$ into the velocity update eliminates
$\vs_{n+1}$ and leaves a linear system for $\vv_{n+1}$ alone:
\[
  \underbrace{\bigl(\mat I+\Delta t\,\Cop+\Delta t^2\,\Kop\bigr)}_{\displaystyle\mat A}
  \,\vv_{n+1}=\vv_n+\Delta t\,(\vJ-\Kop\,\vs_n).
\]
This is the implicit step's linear solve---the $2\times2$ stand-in for the
\code{ksp\_solve} (\cref{fig:tr-innersolve}). With our numbers,
\[
  \mat A=\mat I+0.1\,\Cop+0.01\,\Kop
        =\begin{psmallmatrix}1.04&-0.01\\-0.01&1.04\end{psmallmatrix}.
\]
The matrix $\mat A$ is \emph{fixed for the whole march}---built from the captured
operators and $\Delta t$, none of which change---so it is factored \emph{once} and
reused every step (\S\ref{sec:tr-cost}). This is the operator-capture-once hoist of
\code{fold\_solve} in miniature: \code{op} is bound before the fold, and every step
reads the same $\mat A$.

\medskip
\emph{Thread the state forward.} Starting from $\vy_0=\mathbf0$, solve for $\vv_1$,
set $\vs_1=\Delta t\,\vv_1$, feed $\vy_1=(\vs_1,\vv_1)$ into the next step, and
repeat. Three steps:
\begin{center}
\setlength{\tabcolsep}{6pt}
\renewcommand{\arraystretch}{1.2}
\begin{tabular}{@{}clll@{}}
\toprule
step $n$ & RHS $\vv_n+\Delta t(\vJ-\Kop\vs_n)$ & solve $\vv_{n+1}=\mat A^{-1}(\cdot)$ & $\vs_{n+1}=\vs_n+\Delta t\,\vv_{n+1}$ \\
\midrule
0 & $(0.100,\ 0.000)$ & $(0.0961,\ 0.0009)$ & $(0.00961,\ 0.00009)$ \\
1 & $(0.1854,\ 0.0105)$ & $(0.1784,\ 0.0118)$ & $(0.02745,\ 0.00127)$ \\
2 & $(0.2589,\ 0.0264)$ & $(0.2492,\ 0.0278)$ & $(0.05237,\ 0.00405)$ \\
\bottomrule
\end{tabular}
\end{center}
Each row's input is the previous row's output---the carry $\vy_n\to\vy_{n+1}$
threaded forward, exactly the fold. The driven coordinate $s^{(1)}$ climbs toward
its static deflection while the coupled coordinate $s^{(2)}$ follows through the
off-diagonal stiffness; with the damping present, the march settles rather than
oscillating forever.

\medskip
\emph{Contrast: what a \code{solve\_family} map would do instead.} Suppose we had,
instead, three \emph{independent} right-hand sides $\vb_1,\vb_2,\vb_3$ and the one
fixed operator $\mat A$---an electrostatic-style \code{solve\_family}. We would solve
\[
  \mat A\,\vx_1=\vb_1,\qquad \mat A\,\vx_2=\vb_2,\qquad \mat A\,\vx_3=\vb_3,
\]
and \emph{nothing in $\vb_2$ depends on $\vx_1$}. We could compute $\vx_2$ first, or
on a separate machine, or all three at once: the solves commute. That is the map. The
transient table above is its exact negation: the RHS of step~1 is built from
$\vs_1$, which came from step~0; the RHS of step~2 from $\vs_2$, which came from
step~1. The carry column \emph{is} the dependence chain. Same fixed operator $\mat A$
in both; the difference is whether the right-hand sides are a fixed independent
\emph{family} (map) or a sequence each \emph{generated by the previous solve} (fold).
Three points are the chapter in miniature: \emph{(i)} the recast doubled
$\mathbb R^2$ into $\mathbb R^4$; \emph{(ii)} each step \emph{solved a linear system}
with $\mat A$ fixed and factored once; \emph{(iii)} the state was \emph{threaded}
from each step into the next---a \code{fold\_solve}, not a \code{solve\_family}. Were
we to checkpoint after step~1, saving $\vy_1$, and resume, the schedule-split law
\cref{eq:tr-split} guarantees the identical $\vy_3$.
\end{workedexample}

\begin{workedexample}{From a port waveform to a resonance: the DFT bridge}{tr-dft}
We make the time$\to$frequency bridge of \S\ref{sec:tr-fourier} numerical. Suppose a
transient march produced a sampled port-voltage waveform---a decaying ring at the
structure's resonant frequency. To keep the arithmetic in hand we sample $N=8$
points of a clean cosine ring,
\[
  v_n=\cos\!\Bigl(\tfrac{2\pi}{8}\,2\,n\Bigr),\qquad n=0,1,\dots,7,
\]
that is, a cosine that completes exactly \emph{two} full cycles over the eight
samples: $v=(1,0,-1,0,1,0,-1,0)$.

\medskip
\emph{The discrete Fourier transform.} The DFT of the sampled signal is
\[
  \hat v_k=\sum_{n=0}^{7} v_n\,e^{-i\,2\pi kn/8},\qquad k=0,1,\dots,7,
\]
each $\hat v_k$ measuring how much of the signal oscillates at frequency-index $k$
(i.e.\ at $k$ cycles over the window). Our signal is, by construction, a pure cosine
of \emph{two} cycles, so we expect all its weight to land on $k=2$ (and its mirror
$k=6=N-2$, since a real cosine splits its energy between $+f$ and $-f$). Compute the
suspect bins. For $k=2$, the exponential is $e^{-i\pi n/2}=(1,-i,-1,i,\dots)$, whose
\emph{real} part is $(1,0,-1,0,1,0,-1,0)$---identical to $v_n$---so
\[
  \hat v_2=\sum_{n}v_n\,e^{-i\pi n/2}
         =\underbrace{(1+1+1+1)}_{\text{the four }\pm1\text{ aligned}}=4,
\]
and by the same alignment $\hat v_6=4$. For a generic other bin, say $k=1$, the
exponential is $e^{-i\pi n/4}$ and the four nonzero samples
$v_0,v_2,v_4,v_6=(+1,-1,+1,-1)$ multiply phases that sum to zero:
\[
  \hat v_1=(1)(1)+(-1)(e^{-i\pi/2})+(1)(e^{-i\pi})+(-1)(e^{-i3\pi/2})
         =1+i-1-i=0,
\]
and likewise $\hat v_0=\hat v_3=\hat v_4=\hat v_5=\hat v_7=0$.

\medskip
\emph{Read off the resonance.} The magnitude spectrum is
\[
  |\hat v_k|=(0,\,0,\,4,\,0,\,0,\,0,\,4,\,0),
\]
a single peak at $k=2$ (and its real-signal mirror at $k=6$). The transform has
located the ring's frequency---two cycles per window---as cleanly as
\cref{fig:tr-fourier} drew it. On a real port signal the ring is not a single pure
tone but a band, and the peak has finite width set by the decay rate (a slow decay,
a high-$Q$ resonance, gives a sharp peak); dividing the output spectrum by the input
pulse's spectrum then yields the S-parameter across the whole band. But the essential
move is exactly the one we just did by hand: a time-domain march, followed by a
Fourier transform, recovers the frequency-domain answer the driven analysis
(\cref{ch:driven}) computes one frequency at a time. One pulsed fold, one transform,
the whole band.
\end{workedexample}

\summarysection
The transient analysis computes the time-domain response of a device to an
excitation pulse $\vJ(t)$: not a snapshot but a \emph{movie}, the whole field
history. It runs the anchor skeleton of \cref{ch:situating} with one corner swapped.
\emph{Assemble} (\S\ref{sec:tr-assemble}): three operators $\Kop,\Cop,\Mop$ on
$\Hcurl$ by three \code{fe\_assemble} folds, plus an implicit ODE integrator and a
seed state---all built \emph{once}, before the march (\cref{ch:time}). \emph{Solve}
(\S\ref{sec:tr-solve}): the \code{fold\_solve} corner, the one genuinely sequential
driver. Where the electrostatic and driven solves are \texttt{map}s---independent
families, parallel over terminals or frequencies---the transient march is a
\texttt{fold}: \code{foldl} of an opaque per-step operator over the timestep
schedule, threading a single field-state carry forward, each step beginning where the
last ended, each step hiding an implicit \code{ksp\_solve} (\cref{ch:krylov}). The
whole feature is the one-line composition
$\texttt{transient}=\code{fold\_solve}\circ\code{fe\_assemble}$. The fold's
schedule-split law makes checkpoint-and-restart exact---the carry is the only thing
crossing a split point. \emph{Reduce} (\S\ref{sec:tr-reduce}): a per-step readout of
fields, port waveforms, and energy, producing the whole trajectory; and, by Fourier
transform of a port signal (\S\ref{sec:tr-fourier}), the broadband bridge to the
frequency-domain S-parameters of \cref{ch:driven}. \textbf{Transient is the one
sequential driver: a fold where each step begins where the last ended---the opposite
of the map drivers.}

\furtherreading
The finite-element time-domain formulation, the second-order semi-discrete system,
and the specific implicit integrators (the generalized-$\alpha$ family and its
numerical damping) are treated in Jin's finite-element electromagnetics
text~\citep{jin2014} in the exact form Palace uses. The time-integration machinery
this chapter composes---the stiffness phenomenon, A-stability, the per-step implicit
solve---is developed in \cref{ch:time}, building on the matrix-computation and
eigenvalue-spread material of Golub and Van Loan~\citep{golubvanloan2013}, who also
cover the dense and sparse linear algebra behind each step's solve. For the physics
of the time-domain response and the energy bookkeeping, Griffiths~\citep{griffiths2017}
is the standard undergraduate source. The map-versus-fold distinction that organizes
this chapter is the functional programmer's daily bread, built from scratch in
\cref{ch:combinators}; the fixed-operator map it is contrasted against is the anchor
of \cref{ch:electrostatics}, and the frequency-domain analysis it bridges to is
\cref{ch:driven}.

\exercisessection
\begin{enumerate}[nosep]
  \item \textbf{Map versus fold, by the picture.} Redraw \cref{fig:tr-foldvmap} from
  memory. Which arrows are present in the fold panel but absent in the map panel, and
  what do those arrows represent physically? State, in one sentence each, why the map
  can be run on three machines at once and why the fold cannot.
  \item \textbf{The one swapped corner.} Using \cref{tab:es-swaps}, describe the
  transient analysis as ``the electrostatic anchor with these stages changed.''
  Exactly which cells move (FE space, solve corner, reduction), and which stay? Which
  single change is the structural heart of the chapter?
  \item \textbf{Two independent costs.} The notation box of \S\ref{sec:tr-cost}
  distinguishes the \emph{carry} cost from the \emph{implicit-solve} cost. Classify
  each of the following as affected by one, the other, both, or neither: (a) doubling
  the number of CPU cores; (b) switching from an implicit to a (hypothetical stable)
  explicit integrator; (c) halving $\Delta t$; (d) refining the mesh so the largest
  eigenvalue grows.
  \item \textbf{One implicit step by hand.} For the two-DoF system of
  \cref{wex:tr-twodof}, take \emph{one} backward-Euler step of size $\Delta t=0.2$
  (rather than $0.1$) from rest. Form $\mat A=\mat I+\Delta t\,\Cop+\Delta t^2\Kop$,
  write the right-hand side, and solve the $2\times2$ system for $\vv_1$. Confirm
  $\mat A$ depends only on $\Delta t$ and the operators, not on the step index---the
  operator-capture-once property.
  \item \textbf{Checkpoint exactness.} Using the schedule-split law \cref{eq:tr-split},
  argue that a transient run checkpointed after every $K$ steps and resumed produces
  \emph{exactly} the same final state as an uninterrupted run, for any $K$. What
  single quantity must the checkpoint save? Why does an electrostatic
  \code{solve\_family} (\cref{ch:electrostatics}) have no analogous ``resume''?
  \item \textbf{The DFT bridge, varied.} Repeat \cref{wex:tr-dft} for the $N=8$ signal
  $v_n=\cos(2\pi\cdot1\cdot n/8)$ (one cycle per window instead of two). Predict which
  bins $|\hat v_k|$ are nonzero \emph{before} computing, then verify by evaluating
  $\hat v_1$ and $\hat v_3$. What does the shifted peak say about the resonant
  frequency the march recorded?
  \item \textbf{Why a pulse?} Explain, in terms of the harmonic rule
  (\cref{ch:reductions-pde}), why a single \emph{narrow} pulse in time lets one
  transient run recover a \emph{wide} band of the frequency response, whereas the
  driven analysis (\cref{ch:driven}) needs one solve per frequency. Give one situation
  where you would nonetheless prefer the driven sweep.
  \item \textbf{The opaque step.} The chapter treats \code{time\_step\_op} as opaque,
  quantifying over it rather than rendering it (\cref{fig:tr-innersolve}). What does
  the fold combinator actually \emph{need} to know about one step for the
  schedule-split law to hold, and what can it safely \emph{not} know? Relate your
  answer to why the same \code{fold\_solve} combinator can also drive the
  adaptive-mesh-refinement loop (\cref{ch:situating}), whose per-step body is an
  entire solve rather than an ODE step.
\end{enumerate}

\chapter{Wave Ports and Boundary Modes}
\label{ch:waveports}

\chapterorient{The eigenmode analysis (\cref{ch:eigenmodes}) found the resonances
of a three-dimensional cavity by handing an operator pencil to an opaque
eigenvalue library. This chapter runs \emph{the very same} black-box solve corner,
but on a two-dimensional cross-section sliced out of a three-dimensional boundary,
and it computes something different: not the resonances of a closed volume but the
\emph{propagation modes} of a waveguide---the field patterns that carry energy
along a pipe, a cable, an on-chip line. There is exactly one new idea, and it
comes \emph{before} the solve: a stage-0 extraction that cuts a 2D submesh from
named 3D boundary faces. Everything after it---assemble a pencil, one opaque
\texttt{eigsolve}, a per-mode readout---is the eigenmode machine you already know.
And the product closes a loop: these port modes are precisely the basis the driven
analysis (\cref{ch:driven}) projects its solved fields onto to define its
scattering parameters. By the end you will read the wave-port analysis as
``eigenmodes on an extracted 2D cross-section, feeding the driven solver's ports.''}

\section{The physical question: how does a pipe carry a wave?}

A microwave signal does not travel through free space to get from a generator to a
device; it travels down a \emph{guide}---a hollow rectangular pipe, a coaxial
cable, a microstrip trace on a circuit board, an on-chip coplanar waveguide. Each
of these guides has a characteristic way of carrying energy: the field does not
fill the cross-section arbitrarily but settles into one of a discrete set of
\emph{propagation modes}, each with its own transverse field pattern that simply
\emph{translates} down the guide's axis, multiplied by a travelling-wave factor
$e^{-i k_n z}$. The number $k_n$ is the mode's \emph{propagation constant}: how
many radians of phase the wave accumulates per unit length as it moves down the
guide (\cref{fig:wg-setup}). A mode with a real $k_n$ \emph{propagates}---it
carries energy down the guide without decaying; a mode with an imaginary $k_n$ is
\emph{evanescent}---it dies away exponentially and carries no net power.

\begin{physicsbox}
Picture a long uniform guide running along the $z$ axis. Because the guide is the
same at every cross-section, a field that ``fits'' the cross-section can ride along
the axis unchanged in shape:
\[
  \vE(x,y,z) \;=\; \vE_t(x,y)\,e^{-i k_n z},
\]
a fixed transverse pattern $\vE_t(x,y)$ multiplied by the travelling-wave phase
$e^{-i k_n z}$. The pattern $\vE_t$ is the \emph{mode}; the number $k_n$ is its
\emph{propagation constant}, the spatial frequency of the wave along the guide.
When $k_n$ is real the mode is a true travelling wave (it propagates with phase
velocity $\omega/k_n$ and carries power); when $k_n$ is imaginary the factor
$e^{-i k_n z}=e^{-\abs{k_n} z}$ is a pure decay (the mode is \emph{evanescent} and
stores energy near its source without transporting it). Each guide supports a
\emph{ladder} of such modes; which ones propagate depends on the operating
frequency. Finding the modes of a given cross-section is the whole job of this
chapter.
\end{physicsbox}

\begin{figure}[t]
\centering
\begin{tikzpicture}[font=\footnotesize, >=Stealth]
  % --- the guide drawn in perspective, axis to the right ---
  % back cross-section (2D port face)
  \draw[thick, simInk, fill=simBgGray] (0,0) rectangle (1.7,2.2);
  % front cross-section, offset for depth
  \draw[thick, simInk, fill=simBgGray, opacity=0.55]
    (3.4,0.55) rectangle (5.1,2.75);
  % connecting edges (the pipe length)
  \draw[simGray] (0,0) -- (3.4,0.55);
  \draw[simGray] (1.7,0) -- (5.1,0.55);
  \draw[simGray] (0,2.2) -- (3.4,2.75);
  \draw[simGray] (1.7,2.2) -- (5.1,2.75);
  % transverse field pattern Et on the near face (vertical arrows, sine envelope)
  \foreach \xx in {0.28,0.57,0.85,1.13,1.42}{
    \pgfmathsetmacro{\amp}{0.7*sin(deg(pi*\xx/1.7))}
    \draw[->, simBlue, thick] (\xx,1.1) -- (\xx,{1.1+\amp});
    \draw[->, simBlue, thick] (\xx,1.1) -- (\xx,{1.1-\amp});}
  \node[text=simBlue] at (0.85,2.45) {$\vE_t(x,y)$};
  \node[font=\scriptsize, text=simGray, anchor=north] at (0.85,-0.08)
    {2D cross-section (the port)};
  % propagation axis arrow
  \draw[->, very thick, simRust] (1.9,1.3) -- (3.2,1.55)
    node[midway, above, sloped, font=\scriptsize, text=simRust]
      {$e^{-i k_n z}$};
  \node[text=simRust, font=\scriptsize] at (3.0,0.95) {propagate along $z$};
  % phase wavelength bracket along the axis
  \draw[<->, simGreen, thick] (1.95,0.15) -- (3.35,0.42)
    node[midway, below, sloped, font=\scriptsize, text=simGreen]
      {$\lambda_g=2\pi/k_n$};
\end{tikzpicture}
\caption[A waveguide and its propagation mode]{A uniform waveguide carries energy
in propagation modes. A mode is a fixed transverse field pattern $\vE_t(x,y)$ on
the 2D cross-section (left face, blue), multiplied by the travelling-wave factor
$e^{-i k_n z}$ as it moves down the guide's axis (rust). The propagation constant
$k_n$ sets the guide wavelength $\lambda_g = 2\pi/k_n$ along the axis. A real $k_n$
is a propagating wave that carries power; an imaginary $k_n$ is an evanescent mode
that decays without transporting energy. The wave-port analysis computes these
transverse patterns and their $k_n$ directly from the cross-section geometry---and
the cross-section is a \emph{2D surface}, which is the first hint of what makes
this analysis structurally distinct.}
\label{fig:wg-setup}
\end{figure}

\subsection{Why we compute them: the driven analysis needs its ports}

The wave-port analysis is not usually run for its own sake. It is run because
\emph{something downstream needs the modes}---and that something is the driven
frequency sweep of \cref{ch:driven}. Recall how the driven analysis ends. It
solves the time-harmonic field $\vE_j(\omega)$ inside a device with port $j$
excited, then \emph{projects} that solved field onto each port's mode to read off
a scattering parameter,
\[
  S_{ij} \;=\; \sigma_{ij}\bigl(\inner{\vs_i}{\vE_j(\omega)} - \delta_{ij}\bigr),
\]
the central reduction of that chapter (\cref{ch:driven}). The projection basis
$\vs_i$---``the wave that propagates out of port $i$''---does \emph{not} come from
the driven solve. It is a \emph{precomputed input}: the propagation mode of port
$i$'s cross-section. The driven analysis cannot define its ports without first
knowing what waves those ports support, and \emph{that is what this chapter
computes}. Wave ports sit \emph{upstream} of the driven sweep; their output is the
driven reduction's input (\cref{fig:loop-close}, revisited at the end of the
chapter).

\begin{keyidea}
\textbf{Wave ports $=$ the eigenmode black box (\cref{ch:eigenmodes}) on an
extracted 2D cross-section, producing the port modes the driven analysis
(\cref{ch:driven}) projects onto.} The boundary-mode analysis runs the same opaque
\texttt{eigsolve} solve corner as the eigenmode analysis, with the same
assemble--solve--readout skeleton, but on a 2D submesh cut from named 3D boundary
faces and over a \emph{block} field space. Its product is a table of propagation
modes $\{k_n,\,n_\text{eff},\,(\vE_t,E_n,B_z)\}$, one row per mode---and those mode
fields $\vE_t$ are exactly the basis $\vs_i$ that the driven solver projects its
solved fields onto to compute $S_{ij}$. The chapter has one genuinely new stage
(the 2D extraction) bolted onto a solve corner you already understand.
\end{keyidea}

The rest of the chapter walks the pipeline in order. Stage~0 is the new
idea---extracting a 2D submesh from a 3D boundary (\cref{sec:extract}). Stage~1
assembles the block eigenproblem on that submesh (\cref{sec:assemble}). Stage~2 is
the opaque solve, identical to the eigenmode corner (\cref{sec:solve}). Stage~3
reads out the waveguide-mode table (\cref{sec:reduce}). Two worked examples make
the rectangular guide concrete, and the closing section feeds the modes back to
the driven solver.

\section{Stage 0: extracting the 2D submesh (the distinguishing preface)}
\label{sec:extract}

Every other analysis in this book---the five volume drivers---poses its problem on
the simulation mesh as given. The boundary-mode analysis is the one exception, and
the exception is its defining feature. \emph{Before} it can assemble anything, it
must first manufacture the mesh it will work on, by cutting a two-dimensional
cross-section out of the three-dimensional model.

The physical reason is simple: a port is a \emph{surface}, not a volume. When a
designer marks ``this face of the bounding box is port~1,'' they are naming a 2D
patch of the 3D boundary---the open end of a waveguide, the cross-section of a
coax, the footprint where a feed line meets the device. The modes that propagate
through that port live on that 2D surface. So the analysis must reduce the 3D
boundary description down to a self-contained 2D mesh on which a genuinely 2D
eigenproblem can be posed (\cref{fig:extract}).

\begin{figure}[t]
\centering
\begin{tikzpicture}[font=\footnotesize, >=Stealth]
  % --- LEFT: 3D box with a highlighted boundary face ---
  \begin{scope}[shift={(0,0)}]
    \node[font=\small\bfseries] at (1.5,3.1) {3D model};
    % box (simple cabinet projection)
    \coordinate (A) at (0,0);     \coordinate (B) at (2.6,0);
    \coordinate (C) at (2.6,1.9); \coordinate (D) at (0,1.9);
    \coordinate (dx) at (0.9,0.6);
    \draw[thick, simInk] (A) -- (B) -- (C) -- (D) -- cycle;
    \draw[thick, simInk] (B) -- ($(B)+(dx)$) -- ($(C)+(dx)$) -- (C);
    \draw[thick, simInk] (D) -- ($(D)+(dx)$) -- ($(C)+(dx)$);
    \draw[simGray, densely dashed] (A) -- ($(A)+(dx)$) -- ($(B)+(dx)$);
    \draw[simGray, densely dashed] ($(A)+(dx)$) -- ($(D)+(dx)$);
    % highlighted boundary face = the named port face (left face)
    \fill[simBlue, opacity=0.28] (A) -- (D) -- ($(D)+(dx)$) -- ($(A)+(dx)$) -- cycle;
    \node[text=simBlue, font=\scriptsize, align=center] at (0.55,-0.55)
      {named\\boundary face};
  \end{scope}
  % --- arrow: CreateFromBoundary ---
  \draw[->, very thick, simRust] (4.0,1.3) -- (5.4,1.3)
    node[midway, above, font=\scriptsize, text=simRust, align=center]
      {\code{CreateFromBoundary}}
    node[midway, below, font=\scriptsize, text=simRust, align=center]
      {project $\to$ 2D};
  % --- MIDDLE: extracted 2D mesh ---
  \begin{scope}[shift={(5.8,0)}]
    \node[font=\small\bfseries] at (1.1,3.1) {2D submesh};
    \draw[thick, simBlue, fill=simBgBlue] (0,0) rectangle (2.2,2.4);
    % a coarse triangulation
    \draw[simGray] (0,0) -- (2.2,2.4); \draw[simGray] (0,1.2) -- (1.1,2.4);
    \draw[simGray] (1.1,0) -- (2.2,1.2); \draw[simGray] (0,1.2) -- (2.2,1.2);
    \draw[simGray] (1.1,0) -- (1.1,2.4);
    \node[font=\scriptsize, text=simGray, anchor=north] at (1.1,-0.15)
      {tangent-plane mesh};
  \end{scope}
  % --- the tensor-rotation annotation below ---
  \node[font=\scriptsize, text=simGreen, align=center] at (4.5,-1.15)
    {and: rotate material tensors $\varepsilon,\mu$ into the local tangent frame
     $\;\Rightarrow\;$ a genuinely 2D problem};
\end{tikzpicture}
\caption[The 2D-submesh extraction]{Stage~0, the distinguishing preface. The
analysis selects the named boundary faces of the 3D model (left, blue patch),
extracts them with \code{CreateFromBoundary}, and projects the 3D coordinates of
those faces into the 2D tangent plane to produce a self-contained 2D submesh
(middle). It also rotates the material tensors $\varepsilon,\mu$ (and any path
coordinates) from the global 3D frame into the submesh's local tangent frame
(green), so that the mode problem posed on the 2D mesh is a faithful 2D reduction
of the 3D physics. This extraction is what no other driver does---it is why
wave-ports is a distinct analysis and not merely ``eigenmodes again.'' When the
configuration supplies a 2D mesh directly, this whole stage is the identity.}
\label{fig:extract}
\end{figure}

Concretely, stage~0 does three things, in order:
\begin{enumerate}[nosep]
  \item \textbf{Select and extract the faces.} The configuration names a set of
  boundary attributes---the faces that constitute the port. Those faces are pulled
  out of the parent 3D mesh as a standalone submesh (\code{CreateFromBoundary}),
  with the boundary attributes remapped and the internal boundary edges added so
  the submesh is a complete 2D mesh in its own right.
  \item \textbf{Project to the tangent plane.} The extracted faces still carry
  their 3D coordinates. A 3D$\to$2D projection drops them onto the face's tangent
  plane, giving each node an honest $(x,y)$ pair in the cross-section's own 2D
  coordinate system.
  \item \textbf{Rotate the material tensors.} Permittivity $\varepsilon$ and
  permeability $\mu$ are tensors in the global 3D frame; the 2D mode problem needs
  them expressed in the cross-section's local tangent frame. The extraction
  rotates them (and the port path coordinates) into that frame, so the assembled
  operators see the correct 2D material response.
\end{enumerate}

\begin{notationbox}
The extraction is a genuine \emph{mesh} operation, not a change of variables on the
fly. It produces a new, smaller, lower-dimensional mesh object---two dimensions
where the parent had three, a handful of elements where the parent had many---and
every later stage (assemble, solve, readout) runs on \emph{that} 2D mesh, knowing
nothing of the parent. This is exactly why the cost of the wave-port analysis is
tiny compared with a volume solve: a 2D eigenproblem on a port cross-section is
orders of magnitude smaller than the 3D volume problem the same device's driven
sweep solves. The port modes are cheap precisely because they live on a surface.
\end{notationbox}

With the 2D submesh in hand, the remaining three stages are the eigenmode pipeline
of \cref{ch:eigenmodes}, transplanted onto the cross-section.

\section{Stage 1: assemble the block ND$\,\oplus\,$H1 pencil}
\label{sec:assemble}

On the 2D submesh we now pose the eigenproblem whose eigenvalues are the
propagation constants. The subtlety that distinguishes it from the volume
eigenmode problem of \cref{ch:eigenmodes}---beyond living in 2D---is that the field
splits into \emph{two pieces of different character}, and the eigenproblem is a
\emph{block} system over a pair of function spaces.

\subsection{Why the field splits: transverse and longitudinal}

A waveguide mode's electric field has a part lying \emph{in} the cross-section
plane and a part pointing \emph{along} the propagation axis. These two parts obey
different continuity requirements at material interfaces, and so they want
different finite-element spaces (\cref{ch:functionspaces}, \cref{ch:derham}):
\begin{itemize}[nosep]
  \item the \textbf{transverse field} $\vE_t$---the in-plane part, the pattern
  that actually looks like \cref{fig:wg-setup}'s arrows---lives in the
  two-dimensional Nédélec edge space $\Hcurl$ on the submesh. As always for an
  electric field, $\Hcurl$ builds tangential continuity into the basis, the right
  physics for $\vE$ at an interface;
  \item the \textbf{longitudinal field} $E_n$---the component pointing down the
  guide axis, normal to the cross-section---is a scalar on the cross-section and
  lives in the ordinary $\Hone$ nodal space.
\end{itemize}
The full mode is the pair $(\vE_t, E_n)$, an element of the \emph{direct sum}
$\Hcurl \oplus \Hone$ (\cref{fig:block-space}). This block structure is forced by
the physics: you cannot describe a waveguide mode with a single space, because its
two field components genuinely belong to different members of the de Rham complex
(\cref{ch:derham}).

\begin{figure}[t]
\centering
\begin{tikzpicture}[font=\footnotesize, >=Stealth]
  % H(curl) block
  \node[space, minimum width=30mm, minimum height=15mm, align=center] (curl)
    at (0,0) {$\Hcurl$\\\scriptsize (Nédélec, 2D)};
  \node[font=\scriptsize, text=simGreen, align=center, below=1mm of curl]
    {transverse $\vE_t$\\(in-plane vector)};
  % oplus
  \node[font=\Large] (op) at (3.0,0) {$\oplus$};
  % H1 block
  \node[space, minimum width=30mm, minimum height=15mm, align=center, fill=simBgGold,
        draw=simGold] (h1) at (6.0,0) {$\Hone$\\\scriptsize (nodal scalar)};
  \node[font=\scriptsize, text=simGold!70!black, align=center, below=1mm of h1]
    {longitudinal $E_n$\\(axial scalar)};
  % the combined mode bracket
  \draw[decorate, decoration={brace, amplitude=5pt}, simInk]
    (-1.65,1.05) -- (7.65,1.05);
  \node[font=\small, text=simInk] at (3.0,1.5)
    {one waveguide mode $=(\vE_t,\,E_n)\in \Hcurl\oplus\Hone$};
  % the block operator shape
  \node[font=\footnotesize, align=center, text=simBlue] at (3.0,-2.6)
    {block pencil $(\mat{A}(\omega,\sigma),\,\mat{B})$ over the combined space:};
  \node[font=\small] at (3.0,-3.5)
    {$\mat{A} = \begin{pmatrix}\mat{A}_{tt} & \mat{A}_{tn}\\
                               \mat{A}_{nt} & \mat{A}_{nn}\end{pmatrix}$,
     \quad coupling $\vE_t$ and $E_n$};
\end{tikzpicture}
\caption[The block ND $\oplus$ H1 space]{The waveguide mode lives in a \emph{block}
function space. The transverse field $\vE_t$ is an in-plane vector and lives in the
2D Nédélec edge space $\Hcurl$ (green); the longitudinal field $E_n$ is an axial
scalar and lives in the nodal $\Hone$ space (gold). A single mode is the pair
$(\vE_t,E_n)$ in the direct sum $\Hcurl\oplus\Hone$. The eigenproblem's operators
are therefore \emph{block} operators, with diagonal blocks acting within each space
and off-diagonal blocks coupling the transverse and longitudinal parts---the
de Rham structure of \cref{ch:derham} made into a matrix.}
\label{fig:block-space}
\end{figure}

\subsection{The generalized eigenproblem, shift-inverted}

On this block space we assemble the generalized eigenproblem whose eigenvalues
carry the propagation constants. It is a \emph{generalized} eigenproblem
(\cref{ch:eigenvalues}, \cref{ch:reductions-pde}), a pencil $(\mat{A},\mat{B})$:
\begin{equation}
  \mat{A}(\omega,\sigma)\,\vx \;=\; \lambda\,\mat{B}\,\vx,
  \qquad \vx = (\vE_t, E_n) \in \Hcurl\oplus\Hone,
  \label{eq:wg-gevp}
\end{equation}
where $\mat{A}(\omega,\sigma)$ is the block system operator---it depends on the
operating frequency $\omega$ \emph{and} on the spectral shift $\sigma$, because the
boundary-mode formulation folds the shift into the operator---and $\mat{B}$ is the
block mass operator on the right-hand side. Both are built by the ordinary
assemble fold of \cref{ch:fem}, just with block-structured terms; in the
framework's vocabulary they are two \code{fe\_assemble} calls, run once each.

The eigenvalue $\lambda$ relates to the propagation constant. The formulation
\emph{shift-inverts} (\cref{ch:eigenvalues}, \cref{ch:eigenmodes}) around a target,
\begin{equation}
  \sigma \;=\; -\,k_{n,\text{target}}^2,
  \label{eq:wg-shift}
\end{equation}
the negative square of the propagation constant we are hunting for. Shift-invert is
the same trick the eigenmode chapter used to pull \emph{interior} eigenvalues out
of a large spectrum (\cref{ch:eigenvalues}): by centering the spectral transform on
$-k_{n,\text{target}}^2$ we ask the solver to return the modes whose propagation
constant is nearest the target, which are the modes of physical interest---the
propagating ones near the operating frequency. The target itself comes either from
a user-supplied effective index ($k_{n,\text{target}} = n_\text{eff}\cdot\omega$)
or, when none is given, from the material's light speed.

\begin{keyidea}
Stage~1 assembles a \emph{block} generalized eigenproblem
$\mat{A}(\omega,\sigma)\vx=\lambda\mat{B}\vx$ on the combined Nédélec$\,\oplus\,$nodal
space $\Hcurl\oplus\Hone$ of the 2D submesh: the transverse field $\vE_t$ in
$\Hcurl$, the longitudinal field $E_n$ in $\Hone$. It is shift-inverted around
$\sigma=-k_{n,\text{target}}^2$ to target the modes near a chosen propagation
constant. Two differences from the volume eigenmode pencil (\cref{ch:eigenmodes}):
the space is a \emph{block} $\Hcurl\oplus\Hone$ rather than a single $\Hcurl$, and
the mesh is the 2D submesh of stage~0 rather than the 3D domain. The assemble
\emph{shape}---fold weak-form terms into a pencil, once---is unchanged.
\end{keyidea}

\section{Stage 2: solve --- the same opaque black box}
\label{sec:solve}

Here is the payoff of the whole chapter's framing: stage~2 is \emph{nothing new}.
We have a generalized eigenproblem pencil $(\mat{A},\mat{B})$; we hand it whole to
the opaque eigenvalue library, with the shift target, and get the converged modes
back from a single call. This is the identical solve corner the eigenmode analysis
used (\cref{ch:eigenmodes})---the fourth corner of \cref{ch:situating}, the opaque
black box---transplanted onto the 2D block problem (\cref{fig:wg-pipeline}).
\[
  \texttt{eigs} \;=\; \texttt{eigsolve}\;\bigl((\mat{A},\mat{B}),\ \sigma,\
  n_{\text{modes}}\bigr).
\]
Palace writes \emph{no} eigen-iteration loop here, exactly as in the volume case.
The library (SLEPc or ARPACK, \cref{ch:eigenvalues}) runs its own Krylov--Schur or
Arnoldi process with the shift-and-invert spectral transform
(\cref{ch:krylovschur}), calling back into our code only to apply the operators---and,
because the shift-invert needs $(\mat{A}-\sigma\mat{B})^{-1}$, each callback hides
an inner linear solve, the same nested structure the eigenmode chapter dissected
(\cref{ch:eigenmodes}). We see one call and its result.

\begin{figure}[t]
\centering
\begin{tikzpicture}[node distance=6mm and 8mm, font=\small, >=Stealth]
  \node[data] (cfg) {config};
  \node[stage, right=of cfg, align=center] (ext)
    {extract\\2D submesh};
  \node[stage, right=of ext, align=center] (asm)
    {assemble\\block pencil};
  \node[extern, right=of asm, minimum width=20mm, minimum height=11mm,
        align=center] (solve) {\textbf{eigsolve}\\\scriptsize SLEPc / ARPACK};
  \node[stage, right=of solve, align=center] (map) {map\\readout};
  \node[product, right=of map, align=center] (prod)
    {waveguide-\\mode table};
  \draw[flow] (cfg) -- (ext);
  \draw[flow] (ext) -- node[above, font=\scriptsize]{2D mesh} (asm);
  \draw[flow] (asm) -- node[above, font=\scriptsize]{$(\mat{A},\mat{B}),\sigma$} (solve);
  \draw[flow] (solve) -- node[above, font=\scriptsize]{eigenpairs} (map);
  \draw[flow] (map) -- (prod);
  % the distinguishing-stage callout
  \node[font=\scriptsize, text=simBlue, align=center, below=4mm of ext]
    {the \emph{new} stage\\(3D boundary $\to$ 2D)};
  \draw[refedge] (ext.south) ++(0,-1mm) -- ++(0,-2.5mm);
  % the library's hidden inner loop
  \node[font=\scriptsize, text=simViolet, align=center, below=4mm of solve]
    {opaque eigen-iteration\\(SAME corner as\\\cref{ch:eigenmodes})};
  \draw[refedge] (solve.south) ++(0,-1mm) -- ++(0,-2.5mm);
  % the only Palace loop
  \node[font=\scriptsize, text=simGreen, align=center, below=4mm of map]
    {the only loop\\Palace writes};
  \draw[refedge] (map.south) ++(0,-1mm) -- ++(0,-2.5mm);
\end{tikzpicture}
\caption[The boundary-mode pipeline]{The boundary-mode pipeline. A stage~0
extraction (blue) cuts a 2D submesh from the named 3D boundary faces---the one new
idea. Assembly then builds the block $\Hcurl\oplus\Hone$ pencil
$(\mat{A},\mat{B})$ on that submesh. The \texttt{eigsolve} box is dashed and
violet, the \cref{ch:situating} convention for a corner where Palace writes no
loop: the entire eigen-iteration happens inside the library, exactly as in the
volume eigenmode analysis (\cref{ch:eigenmodes})---this is the \emph{same} black
box, just fed a 2D block problem. The single Palace-written loop (green) is the
per-mode readout map that turns each converged eigenpair into a row of the
waveguide-mode table. Compare \cref{ch:eigenmodes}: same skeleton, with one extra
stage at the front and a block space inside the pencil.}
\label{fig:wg-pipeline}
\end{figure}

\begin{figure}[t]
\centering
\begin{tikzpicture}[font=\footnotesize, >=Stealth]
  % --- LEFT: eigenmode (volume) corner ---
  \begin{scope}[shift={(0,0)}]
    \node[font=\small\bfseries, align=center] at (1.7,2.6)
      {Eigenmode (\cref{ch:eigenmodes})};
    \node[opbox, minimum width=18mm, align=center] (pen1) at (0.1,1.2)
      {$(\Kop,\Cop,\Mop)$\\\scriptsize 3D, single $\Hcurl$};
    \node[extern, minimum width=14mm, minimum height=9mm] (bb1) at (2.4,1.2)
      {\scriptsize eigsolve};
    \draw[flow, shorten >=1pt] (pen1.east) -- (bb1.west);
    \node[product, minimum width=9mm, minimum height=5mm, scale=0.75]
      (o1) at (3.9,1.2) {\scriptsize eigs};
    \draw[flow, shorten >=1pt] (bb1.east) -- (o1.west);
    \node[font=\scriptsize\itshape, align=center] at (1.9,0.2)
      {one opaque call,\\volume modes};
  \end{scope}
  % divider
  \draw[simGray, densely dashed] (4.7,-0.4) -- (4.7,2.9);
  % --- RIGHT: boundary-mode (2D block) corner ---
  \begin{scope}[shift={(5.3,0)}]
    \node[font=\small\bfseries, align=center] at (1.7,2.6)
      {Boundary-mode (this chapter)};
    \node[opbox, minimum width=18mm, align=center, fill=simBgBlue, draw=simBlue]
      (pen2) at (0.1,1.2) {$(\mat{A},\mat{B})$\\\scriptsize 2D, $\Hcurl\oplus\Hone$};
    \node[extern, minimum width=14mm, minimum height=9mm] (bb2) at (2.4,1.2)
      {\scriptsize eigsolve};
    \draw[flow, shorten >=1pt] (pen2.east) -- (bb2.west);
    \node[product, minimum width=9mm, minimum height=5mm, scale=0.75]
      (o2) at (3.9,1.2) {\scriptsize eigs};
    \draw[flow, shorten >=1pt] (bb2.east) -- (o2.west);
    \node[font=\scriptsize\itshape, align=center] at (1.9,0.2)
      {one opaque call,\\port modes};
  \end{scope}
  % the shared box
  \node[font=\scriptsize, text=simViolet] at (5.0,-0.95)
    {\emph{identical} \code{eigsolve} corner --- only the pencil and the mesh differ};
\end{tikzpicture}
\caption[The same black box, two pencils]{The solve corner is literally the same
combinator in both analyses. Left: the eigenmode analysis hands a 3D pencil over a
single $\Hcurl$ space to the opaque \texttt{eigsolve} (\cref{ch:eigenmodes}).
Right: the boundary-mode analysis hands a 2D pencil over the block
$\Hcurl\oplus\Hone$ space to the \emph{same} opaque \texttt{eigsolve}. The dashed
violet box---the library boundary where Palace writes no loop---is drawn identically
on both sides on purpose: only the operands change (the pencil's blocks and the
mesh's dimension), never the corner. This is why understanding
\cref{ch:eigenmodes} is most of the work of understanding this chapter.}
\label{fig:same-box}
\end{figure}

\begin{keyidea}
The solve is \emph{one} \code{eigsolve} call into the opaque eigenvalue library,
\emph{identical} to the eigenmode solve corner (\cref{ch:eigenmodes})---the fourth
corner of \cref{ch:situating}. Palace writes no eigen-iteration loop; the library
owns the Krylov--Schur shift-invert iteration (\cref{ch:krylovschur}) and only
calls back for operator applies, each containing an inner linear solve from the
shift-and-invert. There is no \code{solve\_family} map (no operator/RHS family) and
no state-threaded fold here: the single black-box call \emph{is} the entire solve.
The only thing that changed from \cref{ch:eigenmodes} is what we fed it---a 2D
block pencil instead of a 3D single-space one.
\end{keyidea}

\begin{remark}[Where the null-space subtlety went]
The volume eigenmode analysis had to fight the curl--curl operator's gradient null
space \emph{before} its solve, projecting it out so the black box would not return
a cloud of spurious zero-frequency modes (\cref{ch:eigenmodes}). The boundary-mode
problem carries the same $\Hcurl$ structure on its transverse block, so the same
concern is present in principle; the block $\Hcurl\oplus\Hone$ formulation and the
shift-invert targeting $\sigma=-k_{n,\text{target}}^2$ together keep the iteration
focused on the physical propagating modes near the target rather than the
gradient cloud near the origin. The lesson is identical to
\cref{ch:eigenmodes}: the black box cannot tell physics from artifact, so the
formulation must be arranged \emph{before} the solve to return the modes we want.
\end{remark}

\section{Stage 3: reduce --- the waveguide-mode table}
\label{sec:reduce}

The library returns converged eigenpairs $(\lambda_m,\vx_m)$. The reduce stage
turns them into the physical product: a \emph{table}, one row per mode---the rank-1
``Shape~3'' reduction of \cref{ch:reductions}, the same shape as the eigenmode
$(f,Q)$ table, but now each row carries \emph{fields}. This stage is the one
Palace-written loop in the driver: a \code{map} over the converged modes, with no
coupling between rows, which is exactly what makes it a map and not a fold. The
framework names it \code{waveguide\_mode\_reduce}. Each row carries four things.

\subsection{The propagation constant and effective index}

The headline number is the \emph{propagation constant} $k_n$, recovered from the
eigenvalue by the shift-invert un-transform (\cref{ch:eigenvalues},
\cref{ch:eigenmodes})---the solver works in shifted coordinates, and undoing the
shift \eqref{eq:wg-shift} returns the physical $k_n$. From it the
\emph{effective index} follows by a single divide,
\begin{equation}
  n_\text{eff} \;=\; \frac{k_n}{\omega},
  \label{eq:neff}
\end{equation}
the ratio of the mode's propagation constant to the free-space wavenumber---how
much more slowly (or quickly) the mode travels down the guide than a wave in free
space would. The effective index is the number a microwave engineer reads first: a
mode with $n_\text{eff}$ near $1$ is loosely guided, one with large
$n_\text{eff}$ is tightly bound to the structure.

\subsection{The mode fields, back-transformed and normalized}

Each eigenvector $\vx_m$ \emph{is} the mode, but in the raw block degrees of
freedom of the shift-inverted problem. The readout applies a \emph{back-transform}
to recover the physical transverse and longitudinal fields $(\vE_t, E_n)$, then
\emph{power-normalizes} the mode so that its Poynting power is unity, $\abs{P}=1$.
Normalization matters because an eigenvector is defined only up to a scalar; fixing
the power to $1$ gives every mode a canonical amplitude, which is what makes the
later projection $\inner{\vs_i}{\vE_j}$ in the driven analysis (\cref{ch:driven})
mean ``fraction of unit power in this mode'' rather than an arbitrary multiple.
(The normalization has a totality guard: if a mode's computed power is zero---a
degenerate or purely evanescent edge case---it is left unscaled rather than divided
by zero, the same defensive pattern as the eigenmode $Q=\infty$ lossless limit of
\cref{ch:eigenmodes}.)

Finally, for \emph{propagating} modes only, the readout forms the longitudinal
magnetic field from the transverse electric field by a discrete curl,
\begin{equation}
  B_z \;=\; \frac{\Curl\,\vE_t}{i\,\omega},
  \label{eq:bz}
\end{equation}
one curl and a complex scale per mode (the discrete-curl operator is the de Rham
interpolator of \cref{ch:derham}). For an evanescent mode this field is omitted---a
per-mode conditional, ``$B_z$ if propagating, nothing otherwise,'' keyed on whether
$k_n$ is real.

\subsection{Propagating versus evanescent, and the cutoff}

The single most important reading of the table is the split between
\emph{propagating} and \emph{evanescent} modes, and it lives entirely in whether
$k_n$ is real or imaginary (\cref{fig:dispersion}). A guide mode has a
\emph{cutoff frequency} $f_c$: below it the mode cannot propagate. The relationship
is the dispersion law
\begin{equation}
  k_n^2 \;=\; \Bigl(\frac{\omega}{c}\Bigr)^2 - k_c^2,
  \qquad k_c = \frac{\omega_c}{c} = \frac{2\pi f_c}{c},
  \label{eq:dispersion}
\end{equation}
where $k_c$ is the mode's \emph{cutoff wavenumber}, a fixed property of the
cross-section geometry. Read it directly:
\begin{itemize}[nosep]
  \item \textbf{Above cutoff} ($\omega > \omega_c$): the right-hand side is
  positive, so $k_n$ is \emph{real}. The mode \emph{propagates}---it carries power
  down the guide with phase velocity $\omega/k_n$, and the readout forms its $B_z$.
  \item \textbf{Below cutoff} ($\omega < \omega_c$): the right-hand side is
  negative, so $k_n$ is \emph{imaginary}. The mode is \emph{evanescent}---the
  factor $e^{-ik_n z}=e^{-\abs{k_n}z}$ is a pure decay, carrying no net power, and
  $B_z$ is omitted.
\end{itemize}
At cutoff exactly, $k_n=0$: the mode neither propagates nor decays, oscillating in
place. The cutoff is the gate, and the dispersion law \eqref{eq:dispersion} is the
single equation that decides, for every mode at every frequency, which side of the
gate it is on.

\begin{figure}[t]
\centering
\begin{tikzpicture}
\begin{axis}[
  width=0.84\linewidth, height=6.2cm,
  xlabel={angular frequency $\omega$}, ylabel={propagation constant},
  xmin=0, xmax=4, ymin=-2.2, ymax=2.6,
  xtick={1.2}, xticklabels={$\omega_c$},
  ytick={0}, yticklabels={$0$},
  axis lines=middle,
  legend pos=north west, legend cell align=left,
  tick label style={font=\footnotesize}, label style={font=\footnotesize},
  legend style={font=\footnotesize, draw=simGray!40},
  clip=false,
]
% propagating branch: k_n = sqrt(omega^2 - wc^2) for omega>wc  (c=1, wc=1.2)
\addplot[very thick, simBlue, smooth, samples=120, domain=1.2:4]
  { sqrt(x^2 - 1.44) };
\addlegendentry{$k_n$ real (propagating)}
% evanescent branch: |k_n| = sqrt(wc^2 - omega^2) drawn negative (imaginary magnitude)
\addplot[very thick, simRust, densely dashed, smooth, samples=120, domain=0:1.2]
  { -sqrt(1.44 - x^2) };
\addlegendentry{$k_n$ imaginary (evanescent)}
% cutoff marker
\draw[simGray, densely dotted] (axis cs:1.2,-2.2) -- (axis cs:1.2,2.6);
\node[font=\scriptsize, text=simGreen, anchor=west] at (axis cs:2.4,1.9)
  {above cutoff: carries power};
\node[font=\scriptsize, text=simRust, anchor=west] at (axis cs:0.05,-1.6)
  {below cutoff: decays};
\node[font=\scriptsize, text=simInk] at (axis cs:1.25,2.35) {cutoff};
\end{axis}
\end{tikzpicture}
\caption[Dispersion: propagating versus evanescent]{The dispersion relation
$k_n^2=(\omega/c)^2-k_c^2$ for a single guide mode, the gate between propagation
and decay. Above the cutoff frequency $\omega_c$ (right of the dotted line) the
right-hand side is positive: $k_n$ is real (blue), the mode is a true travelling
wave carrying power down the guide. Below cutoff (left) the right-hand side is
negative: $k_n$ is imaginary (rust, drawn as its magnitude below the axis), the
mode is evanescent and decays without transporting energy. Exactly at $\omega_c$
the mode has $k_n=0$. Each mode of a cross-section has its own cutoff; at a given
operating frequency, only the modes whose cutoff lies below it propagate. The
wave-port readout reports $k_n$ for every mode, and this sign of $k_n^2$ is how the
table is read.}
\label{fig:dispersion}
\end{figure}

\begin{keyidea}
The reduce stage is a \code{map} over the converged eigenpairs producing the
\emph{waveguide-mode table}, one row per mode, each carrying: the propagation
constant $k_n$ (the eigenvalue, shift-invert un-transformed); the effective index
$n_\text{eff}=k_n/\omega$; the back-transformed, power-normalized mode fields
$(\vE_t,E_n)$; and, for propagating modes only, $B_z=\Curl\vE_t/(i\omega)$. It is a
rank-1 ``Shape~3'' table like the eigenmode $(f,Q)$ table (\cref{ch:reductions}),
but its rows carry \emph{fields}, not just scalars---the defining difference from
its scalar-only sibling. A mode propagates when $k_n$ is real (above cutoff) and is
evanescent when $k_n$ is imaginary (below cutoff). This is the only loop the
boundary-mode driver writes, and it is a pure post-processing map---never a solve.
\end{keyidea}

\section{The composition, and closing the loop to the driven solver}
\label{sec:compose}

The whole boundary-mode driver is the clean four-stage composition the framework
names (\cref{lst:wgcompose}): extract the 2D submesh, assemble the block pencil,
one opaque eigensolve, one per-mode readout map.

\begin{figure}[t]
\begin{lstlisting}[style=l4style, caption={[The boundary-mode composition]The
boundary-mode feature as a composition of firm vocabulary. Stage~0 extracts the 2D
submesh from the named 3D boundary faces---the one new idea. Stages~1--3 are the
eigenmode pipeline of \cref{ch:eigenmodes}: assemble the block
$\Hcurl\oplus\Hone$ pencil, hand it to the opaque \texttt{eigsolve} once, and
\texttt{map} the per-mode readout over the converged modes. The only
Palace-written loop is the final \texttt{map}.}, label={lst:wgcompose}]
boundary_mode :: BoundaryModeConfig -> WaveguideModeTable
boundary_mode cfg =
  let mesh2d = extract_boundary_2d_submesh (parent_mesh cfg) (surface_attrs cfg) -- (0) 3D boundary -> 2D submesh
      space  = mode_space mesh2d cfg                  -- the block ND (+) H1 space on the 2D submesh
      opA    = fe_assemble space [ block_system (omega cfg) (sigma cfg) ]  -- (1a) block system A(w,s)
      opB    = fe_assemble space [ mass_block ]       -- (1b) generalized RHS block B
      pencil = eig_pencil opA opB (sigma cfg) (n_modes cfg)  -- (A, B) + shift-invert s = -kn_target^2
      eigs   = eigsolve pencil (initial_space cfg)    -- (2) ONE opaque eigensolve -- SAME corner as eigenmode
  in  map (waveguide_mode_readout cfg (omega cfg)) eigs  -- (3) per-mode readout map: kn, n_eff, (Et, En, Bz)
\end{lstlisting}
\end{figure}

Read as a pipeline composition,
\[
  \texttt{boundary\_mode}
    \;=\;
  (\texttt{map}\ \texttt{readout})
    \;\circ\;
  \texttt{eigsolve}
    \;\circ\;
  \texttt{fe\_assemble}^{\times 2}
    \;\circ\;
  \texttt{extract\_2d\_submesh}.
\]
Three of the four stages are shared verbatim with the eigenmode analysis
(\cref{ch:eigenmodes})---the assemble fold, the opaque solve, the readout
map---and the fourth, the extraction, is the prefix that makes this analysis
distinct.

\subsection{Where the modes go: the driven port projection}

Now we close the loop opened in \cref{sec:assemble}. The waveguide-mode table is
not the end of the story; it is an \emph{input} to the driven analysis
(\cref{ch:driven}). Each port mode field $\vE_t$ in the table \emph{is} a
projection basis vector $\vs_i$ in the driven reduction
(\cref{fig:loop-close}). When the driven sweep solves its 3D field $\vE_j(\omega)$
and needs the scattering parameter $S_{ij}$, it computes the modal overlap of that
solved field with port $i$'s mode---the very mode this chapter produced:
\[
  S_{ij}
    \;=\;
  \sigma_{ij}\bigl(\inner{\vs_i}{\vE_j(\omega)} - \delta_{ij}\bigr),
  \qquad
  \vs_i \;=\; \text{port $i$'s mode field } \vE_t \text{ (from this chapter)}.
\]
The power-normalization of \cref{sec:reduce} is what makes this projection a clean
fraction; the propagation constant $k_n$ is what supplies the de-embedding phase
$e^{ik_n d}$ that shifts the reference plane in $\sigma_{ij}$ (\cref{ch:driven}).
The boundary-mode analysis is, in this sense, a \emph{subroutine} of the driven
analysis: run once per port to characterize what waves it supports, before the
frequency sweep can interpret its solved fields.

\begin{figure}[t]
\centering
\begin{tikzpicture}[font=\footnotesize, node distance=8mm and 14mm, >=Stealth]
  % boundary-mode pipeline (compressed)
  \node[stage, align=center] (bm) {boundary-mode\\\scriptsize(this chapter)};
  \node[product, right=of bm, align=center] (modes)
    {port modes\\$\vs_i=\vE_t$, $k_n$};
  % driven pipeline
  \node[stage, right=of modes, align=center, fill=simBgRust, draw=simRust] (drv)
    {driven sweep\\\scriptsize(\cref{ch:driven})};
  \node[product, right=of drv, align=center] (S) {$S_{ij}(\omega)$};
  \draw[flow] (bm) -- node[above, font=\scriptsize]{eigsolve} (modes);
  \draw[flow] (modes) -- node[above, font=\scriptsize, align=center]
    {projection\\basis} (drv);
  \draw[flow] (drv) -- node[above, font=\scriptsize]{project} (S);
  % the projection equation below
  \node[font=\footnotesize, text=simInk, below=7mm of drv]
    {$S_{ij}=\sigma_{ij}(\inner{\vs_i}{\vE_j(\omega)}-\delta_{ij})$};
  \draw[refedge] (modes.south) .. controls +(0,-7mm) and +(-1.2,4mm) ..
    ($(drv.south)+(-0.3,-6mm)$);
  \node[font=\scriptsize\itshape, text=simGray, below=1mm of modes]
    {computed once per port};
\end{tikzpicture}
\caption[Closing the modality loop]{How the wave-port analysis feeds the driven
analysis, closing the modality loop. This chapter computes a port's propagation
modes (the table $\{\vs_i=\vE_t,\,k_n,\,\dots\}$) by the eigenmode black box on a 2D
cross-section. The driven analysis (\cref{ch:driven}) then \emph{consumes} those
modes: it solves its 3D field $\vE_j(\omega)$ with port $j$ excited and projects
the result onto each port mode $\vs_i$ to read off the scattering parameter
$S_{ij}$. The port modes are the projection basis the driven reduction cannot
define without; the boundary-mode analysis runs once per port, upstream of the
sweep. The two modalities are not independent---one produces exactly what the other
needs.}
\label{fig:loop-close}
\end{figure}

\subsection{Fundamental and higher modes; single-mode operation}

A guide's modes form a ladder ordered by cutoff frequency. The
\emph{fundamental}---the mode with the lowest cutoff---is the one a guide is
usually \emph{designed} to carry: between the fundamental's cutoff and the next
mode's cutoff lies the \emph{single-mode band}, the frequency range in which
exactly one mode propagates and all others are evanescent. Operating a guide in its
single-mode band is the norm in microwave engineering, because a single propagating
mode means the field at the port is unambiguous: one $\vs_i$ to project onto, one
$k_n$ to de-embed with, a clean port definition. The wave-port analysis typically
requests the lowest few modes (via the shift-invert target of \cref{sec:assemble}),
identifies which propagate at the operating frequency, and hands the fundamental
(and any other propagating modes) to the driven solver as the port basis.

\section{Worked examples}

\begin{workedexample}{The rectangular waveguide, by hand}{rectwg}
We compute the modes of the simplest guide---a hollow rectangular pipe of
cross-section $a\times b$ (with $a>b$), filled with a medium of light speed $c$,
its walls perfect conductors---and show how the discrete 2D eigenproblem of
\cref{sec:assemble} recovers the textbook answer. The continuous problem has
closed-form modes (the TE and TM families); we use them to anchor what the
black box returns.

\emph{The analytic cutoffs.} A rectangular guide's modes are indexed by two
integers $(m,n)$, counting half-wavelengths across the two transverse directions.
Each mode's cutoff \emph{wavenumber} is set by the geometry,
\[
  k_{c,mn} \;=\; \pi\sqrt{\Bigl(\frac{m}{a}\Bigr)^2 + \Bigl(\frac{n}{b}\Bigr)^2},
\]
and its cutoff \emph{frequency} follows as $f_{c,mn}=c\,k_{c,mn}/2\pi$, i.e.
\[
  f_{c,mn} \;=\; \frac{c}{2}\sqrt{\Bigl(\frac{m}{a}\Bigr)^2
                                 + \Bigl(\frac{n}{b}\Bigr)^2}.
\]
The \emph{fundamental} is the $\mathrm{TE}_{10}$ mode ($m=1,n=0$), the lowest
cutoff because $a$ is the larger dimension:
\[
  f_{c,10} \;=\; \frac{c}{2a},
\]
a mode that fits exactly half a transverse wavelength across the wide wall---the
2D analog of the half-wavelength cavity fundamental of \cref{ch:eigenmodes}.

\emph{A number.} Take the standard X-band guide WR-90, with
$a=\SI{22.86}{\milli\metre}$ and $b=\SI{10.16}{\milli\metre}$, air-filled
($c=\SI{3.0e8}{\metre\per\second}$). The fundamental cutoff is
\[
  f_{c,10} \;=\; \frac{c}{2a}
    \;=\; \frac{3.0\times10^8}{2\cdot 0.02286}\,\si{\hertz}
    \;\approx\; \SI{6.56}{\giga\hertz}.
\]
The next modes up are $\mathrm{TE}_{20}$ at $f_{c,20}=c/a\approx\SI{13.1}{\giga\hertz}$
and $\mathrm{TE}_{01}$ at $f_{c,01}=c/2b\approx\SI{14.8}{\giga\hertz}$. So between
$\SI{6.56}{}$ and $\SI{13.1}{\giga\hertz}$ exactly one mode propagates: the
single-mode band, which is why WR-90 is specified for roughly
$\SI{8.2}{}$--$\SI{12.4}{\giga\hertz}$ operation.

\emph{What the discrete solve does.} The wave-port analysis does not know these
formulas. It triangulates the $a\times b$ rectangle into the 2D submesh
(\cref{sec:extract}), assembles the block $\Hcurl\oplus\Hone$ pencil
$(\mat{A},\mat{B})$ on it (\cref{sec:assemble}), shift-inverts around a target near
the fundamental, and hands the pencil to \texttt{eigsolve} (\cref{sec:solve}). The
library returns eigenvalues $\lambda_m$; the readout un-transforms each to a
propagation constant $k_{n,m}$ (\cref{sec:reduce}). For the fundamental at an
operating frequency $\omega$ above cutoff, that recovered $k_n$ satisfies---to
discretization error---the dispersion law \eqref{eq:dispersion} with the analytic
cutoff wavenumber:
\[
  k_n^2 \;=\; \Bigl(\frac{\omega}{c}\Bigr)^2 - k_{c,10}^2,
  \qquad
  k_{c,10} = \frac{\pi}{a}.
\]
Exactly as the coarse $3\times3$ cavity mesh of \cref{ch:eigenmodes} pinned the
fundamental resonance to within a few percent, a modest 2D mesh of the WR-90
cross-section pins $k_{c,10}=\pi/a$---hence the fundamental cutoff---to within
discretization error, with the eigenvector reproducing the
$\mathrm{TE}_{10}$ transverse pattern ($\vE_t$ a single half-sine across the wide
wall). \emph{This little eigenproblem is, in miniature, exactly what the black box
does at scale}: build $(\mat{A},\mat{B})$ on the cross-section, find $\lambda$,
un-transform to $k_n$. The only difference at full size is the mesh resolution and
that we hand the pencil to a library instead of reading $k_c$ off a formula.
\end{workedexample}

\begin{workedexample}{Propagating versus evanescent at a fixed frequency}{propevan}
We run the WR-90 guide of \cref{wex:rectwg} at a single operating frequency and ask
the question the table is read for: which modes propagate and which decay? Take
$f=\SI{10}{\giga\hertz}$ (mid-band), so $\omega/c = 2\pi f/c$ with
$\omega/c \approx \SI{209.4}{\per\metre}$.

\emph{The dispersion test.} For each mode we compute $k_n^2 = (\omega/c)^2 - k_c^2$
from \eqref{eq:dispersion} and read the sign. The cutoff wavenumbers are
$k_{c,10}=\pi/a$, $k_{c,20}=2\pi/a$, $k_{c,01}=\pi/b$:
\[
  k_{c,10} = \frac{\pi}{0.02286} \approx \SI{137.4}{\per\metre},
  \quad
  k_{c,20} \approx \SI{274.9}{\per\metre},
  \quad
  k_{c,01} = \frac{\pi}{0.01016} \approx \SI{309.2}{\per\metre}.
\]
With $(\omega/c)^2 \approx (209.4)^2 = \SI{4.385e4}{\per\metre\squared}$:
\begin{itemize}[nosep]
  \item $\mathrm{TE}_{10}$: $k_n^2 = 4.385\times10^4 - (137.4)^2
    = 4.385\times10^4 - 1.888\times10^4 = +2.497\times10^4>0$, so
    $k_n=\SI{158.0}{\per\metre}$ is \emph{real}---\textbf{propagating}.
  \item $\mathrm{TE}_{20}$: $k_n^2 = 4.385\times10^4 - (274.9)^2
    = 4.385\times10^4 - 7.557\times10^4 = -3.172\times10^4<0$, so $k_n$ is
    \emph{imaginary}---\textbf{evanescent}.
  \item $\mathrm{TE}_{01}$: $k_n^2 = 4.385\times10^4 - (309.2)^2
    = 4.385\times10^4 - 9.560\times10^4 = -5.175\times10^4<0$, so $k_n$ is
    \emph{imaginary}---\textbf{evanescent}.
\end{itemize}
At $\SI{10}{\giga\hertz}$ exactly \emph{one} mode propagates---the fundamental
$\mathrm{TE}_{10}$---confirming the single-mode band of \cref{wex:rectwg}.

\emph{The effective index.} For the propagating fundamental, the effective index
\eqref{eq:neff} is
\[
  n_\text{eff} \;=\; \frac{k_n}{\omega/c}
    \;=\; \frac{158.0}{209.4}
    \;\approx\; 0.755.
\]
The mode travels with $n_\text{eff}<1$: its guide wavelength
$\lambda_g=2\pi/k_n\approx\SI{39.8}{\milli\metre}$ is \emph{longer} than the
free-space wavelength ($\SI{30}{\milli\metre}$ at $\SI{10}{\giga\hertz}$), the
familiar fact that a guided wave's phase races ahead of free space near cutoff.

\emph{What the readout records.} For $\mathrm{TE}_{10}$ the readout reports the real
$k_n$, the effective index $0.755$, the back-transformed power-normalized fields
$(\vE_t,E_n)$, and---because the mode propagates---the longitudinal magnetic field
$B_z=\Curl\vE_t/(i\omega)$ \eqref{eq:bz}. For $\mathrm{TE}_{20}$ and
$\mathrm{TE}_{01}$ the readout reports the imaginary $k_n$ and the fields but
\emph{omits} $B_z$ (the per-mode conditional of \cref{sec:reduce}), because an
evanescent mode carries no propagating longitudinal flux. This is the table the
driven analysis (\cref{ch:driven}) would receive: at $\SI{10}{\giga\hertz}$, drive
the single propagating $\mathrm{TE}_{10}$ mode and project onto it---one clean port
basis vector.
\end{workedexample}

\section{A gallery of transverse mode patterns}

Different guides support different transverse patterns, and the eigenvector $\vx_m$
is exactly that pattern (\cref{fig:patterns}). The fundamental rectangular mode is
a single transverse lobe; higher modes add nodal lines; a coaxial line's
fundamental is the radial TEM pattern that has no cutoff at all. The wave-port
analysis returns these as the field part of each table row, and a designer reads
them to understand \emph{where} on the port the field concentrates---which sets
where loss matters and how the port couples to the device.

\begin{figure}[t]
\centering
\begin{tikzpicture}[font=\footnotesize, >=Stealth]
  % --- TE10: single half-sine across width ---
  \begin{scope}[shift={(0,0)}]
    \draw[thick, simInk] (0,0) rectangle (2.6,1.6);
    \foreach \xx in {0.3,0.6,...,2.4}{
      \pgfmathsetmacro{\amp}{0.45*sin(deg(pi*\xx/2.6))}
      \draw[->, simBlue, semithick] (\xx,0.8) -- (\xx,{0.8+\amp});
      \draw[->, simBlue, semithick] (\xx,0.8) -- (\xx,{0.8-\amp});}
    \node[font=\scriptsize, anchor=north] at (1.3,-0.1)
      {$\mathrm{TE}_{10}$ (fundamental)};
  \end{scope}
  % --- TE20: two half-sines ---
  \begin{scope}[shift={(3.6,0)}]
    \draw[thick, simInk] (0,0) rectangle (2.6,1.6);
    \foreach \xx in {0.3,0.6,...,2.4}{
      \pgfmathsetmacro{\amp}{0.42*sin(deg(2*pi*\xx/2.6))}
      \draw[->, simRust, semithick] (\xx,0.8) -- (\xx,{0.8+\amp});
      \draw[->, simRust, semithick] (\xx,0.8) -- (\xx,{0.8-\amp});}
    \draw[densely dashed, simGray] (1.3,0) -- (1.3,1.6);
    \node[font=\scriptsize, anchor=north] at (1.3,-0.1)
      {$\mathrm{TE}_{20}$ (one node)};
  \end{scope}
  % --- coax TEM: radial arrows ---
  \begin{scope}[shift={(7.2,0)}]
    \draw[thick, simInk] (1.3,0.8) circle (0.78);
    \fill[simInk] (1.3,0.8) circle (0.22);
    \foreach \ang in {0,45,...,315}{
      \draw[->, simGreen, semithick]
        ({1.3+0.28*cos(\ang)},{0.8+0.28*sin(\ang)}) --
        ({1.3+0.72*cos(\ang)},{0.8+0.72*sin(\ang)});}
    \node[font=\scriptsize, anchor=north] at (1.3,-0.1)
      {coax TEM (no cutoff)};
  \end{scope}
\end{tikzpicture}
\caption[Transverse mode patterns]{A small gallery of transverse mode patterns
$\vE_t$, the field part of the waveguide-mode table. Left: the rectangular
fundamental $\mathrm{TE}_{10}$, a single transverse lobe (half a sine across the
wide wall). Middle: the higher $\mathrm{TE}_{20}$ mode, with one interior nodal line
(dashed) where the field reverses---it has a higher cutoff and propagates only at
higher frequency (\cref{wex:propevan}). Right: a coaxial line's fundamental, the
radial TEM mode, which uniquely has \emph{no} cutoff (it propagates at every
frequency, down to DC)---which is why coax is the medium of choice for broadband
connections. Each pattern is one eigenvector $\vx_m$ returned by the black box;
the designer reads it to see where the port's field lives.}
\label{fig:patterns}
\end{figure}

\begin{remark}[Why ``boundary'' mode]
The analysis is called \emph{boundary}-mode in Palace because its defining stage
extracts the mesh from the model's \emph{boundary} faces (\cref{sec:extract}),
where the ports physically are, rather than from the interior volume. The names
``wave-port mode,'' ``boundary mode,'' and ``port mode'' all denote the same
object: the propagation modes of a cross-section, computed so the driven analysis
can use them as ports. The dispatch is a sixth analysis branch alongside the five
volume drivers, sharing their solve infrastructure but distinguished by the 2D
extraction at its front.
\end{remark}

\summarysection
The wave-port (boundary-mode) analysis computes the propagation modes of a
waveguide cross-section---the field patterns $\vE_t$ that ride down a guide as
$\vE_t\,e^{-ik_n z}$, each with its propagation constant $k_n$ and effective index
$n_\text{eff}=k_n/\omega$. It is the eigenmode analysis (\cref{ch:eigenmodes}) with
one new stage at the front. \emph{Stage~0}, the distinguishing preface, extracts a
2D submesh from named 3D boundary faces (\code{CreateFromBoundary} $\to$ project to
the tangent plane $\to$ rotate the material tensors), because a port is a 2D
surface. \emph{Stage~1} assembles a generalized eigenproblem
$\mat{A}(\omega,\sigma)\vx=\lambda\mat{B}\vx$ on the \emph{block}
$\Hcurl\oplus\Hone$ space of that submesh---transverse $\vE_t$ in the 2D Nédélec
space, longitudinal $E_n$ in the nodal space (\cref{ch:functionspaces},
\cref{ch:derham})---shift-inverted around $\sigma=-k_{n,\text{target}}^2$.
\emph{Stage~2} is \emph{one} \code{eigsolve} call into the opaque eigenvalue
library---the \emph{identical} black-box solve corner as the eigenmode analysis
(\cref{ch:eigenmodes}, \cref{ch:eigenvalues}), Palace writes no eigen-iteration
loop---just fed a 2D block pencil instead of a 3D single-space one. \emph{Stage~3}
is a \code{map} producing the rank-1 waveguide-mode table (\cref{ch:reductions}):
per mode, the un-transformed $k_n$, the effective index, the back-transformed
power-normalized fields $(\vE_t,E_n)$, and---for propagating modes
($k_n$ real, above cutoff)---the longitudinal $B_z=\Curl\vE_t/(i\omega)$; evanescent
modes ($k_n$ imaginary, below cutoff) omit $B_z$. The whole driver is
$\texttt{boundary\_mode}=(\texttt{map}\ \texttt{readout})\circ\texttt{eigsolve}
\circ\texttt{fe\_assemble}^{\times2}\circ\texttt{extract\_2d\_submesh}$. Its product
closes the modality loop: the port mode fields $\vE_t$ are exactly the basis
$\vs_i$ the driven analysis (\cref{ch:driven}) projects its solved fields onto to
define the scattering parameters $S_{ij}=\sigma_{ij}(\inner{\vs_i}{\vE_j}-\delta_{ij})$.
Wave ports are eigenmodes on an extracted 2D cross-section, computed so the driven
sweep can have its ports.

\furtherreading
Pozar's \emph{Microwave Engineering}~\citep{pozar2011} is the standard
device-level account of waveguides, the TE/TM mode families, cutoff frequencies,
the propagation constant and guide wavelength, and the single-mode band---the
physical reading of everything in \cref{wex:rectwg} and \cref{wex:propevan}, and the
de-embedding and impedance normalization that turn a port mode into a scattering
reference. Jin's \emph{Finite Element Method in
Electromagnetics}~\citep{jin2014} develops the finite-element waveguide
eigenproblem in full---the block $\Hcurl\oplus\Hone$ formulation, the
transverse/longitudinal field split, and the modal projection that defines the
wave-port boundary condition feeding the driven solve (\cref{ch:driven})---and is
the FEM-level companion to this chapter. For the underlying mathematics of edge
elements, the de Rham complex on a 2D cross-section, and why the transverse field
belongs in $\Hcurl$ while the longitudinal field belongs in $\Hone$, Monk's
\emph{Finite Element Methods for Maxwell's Equations}~\citep{monk2003} is the
rigorous reference; \cref{ch:derham} builds the relevant structure. The opaque
eigensolver corner itself is the subject of \cref{ch:eigenmodes} and the
shift-and-invert machinery of \cref{ch:eigenvalues}.

\exercisessection
\begin{enumerate}[nosep]
  \item \textbf{The one new stage.} State precisely which stage of the
  \texttt{extract}$\to$\texttt{assemble}$\to$\texttt{solve}$\to$\texttt{reduce}
  pipeline distinguishes the boundary-mode analysis from the eigenmode analysis
  (\cref{ch:eigenmodes}), and which three stages the two analyses share. In one
  sentence, say \emph{why} a port needs that extra stage but a cavity does not.
  \item \textbf{The block space.} A waveguide mode is the pair $(\vE_t, E_n)$ in
  $\Hcurl\oplus\Hone$. Which physical field component lives in which space, and
  why---what continuity does each space build into its basis
  (\cref{ch:functionspaces}, \cref{ch:derham})? Why can a single function space not
  describe the whole mode?
  \item \textbf{The same black box.} The chapter insists stage~2 is identical to
  the eigenmode solve corner (\cref{ch:eigenmodes}). Name the two things that
  \emph{do} differ between the two solves, and the one thing---the structural
  property of the solve stage---that is the same. Why does Palace write no loop in
  either?
  \item \textbf{Cutoff by hand.} For an air-filled rectangular guide with
  $a=\SI{40}{\milli\metre}$, $b=\SI{20}{\milli\metre}$, compute the cutoff
  frequencies of the $\mathrm{TE}_{10}$, $\mathrm{TE}_{20}$, and $\mathrm{TE}_{01}$
  modes using $f_{c,mn}=\tfrac{c}{2}\sqrt{(m/a)^2+(n/b)^2}$. What is the single-mode
  band (the range in which exactly one mode propagates)?
  \item \textbf{Propagating or evanescent.} Using the guide of the previous
  exercise at $f=\SI{5}{\giga\hertz}$, apply the dispersion test
  $k_n^2=(\omega/c)^2-k_c^2$ to each of the three modes and state which propagate
  ($k_n$ real) and which are evanescent ($k_n$ imaginary). For each propagating
  mode, compute the effective index $n_\text{eff}=k_n/(\omega/c)$.
  \item \textbf{Why $B_z$ only sometimes.} The readout forms
  $B_z=\Curl\vE_t/(i\omega)$ for propagating modes but omits it for evanescent
  ones. Explain physically why the longitudinal magnetic field is recorded only
  when $k_n$ is real, and identify which per-mode conditional in the reduce stage
  (\cref{sec:reduce}) makes this choice.
  \item \textbf{Closing the loop.} The driven analysis (\cref{ch:driven}) computes
  $S_{ij}=\sigma_{ij}(\inner{\vs_i}{\vE_j(\omega)}-\delta_{ij})$. Identify exactly
  which output of \emph{this} chapter supplies $\vs_i$, and which supplies the
  de-embedding phase inside $\sigma_{ij}$. Why must the boundary-mode analysis run
  \emph{before} the driven sweep rather than after?
  \item \textbf{Power normalization.} The readout normalizes each mode to unit
  Poynting power $\abs{P}=1$, with a guard that leaves a mode unscaled when its
  computed power is zero. Explain (i) why normalization is needed at all---what
  ambiguity in an eigenvector it removes---and (ii) why the zero-power guard is the
  same kind of defensive branch as the eigenmode $Q=\infty$ lossless limit
  (\cref{ch:eigenmodes}).
\end{enumerate}

\chapter{Assembling a Driver: A Complete Worked Trace}
\label{ch:fulltrace}

\chapterorient{Every chapter so far has built one machine and set it on the shelf:
the mesh, the function spaces, the assembly fold, the weak forms, the Krylov
solvers, the multigrid preconditioners, the matrix-free applies, the reduction
verbs. This chapter takes the whole shelf down at once. We pick one concrete
simulation---the driven analysis of a two-port filter---and trace it from the
configuration file the engineer writes to the S-parameter curve the simulation
draws, naming, at every stage, the exact machine from the exact earlier chapter
that does the work. Nothing new is built here. Instead we watch \emph{all of it
run together}, see the whole simulator as a single composition of perhaps fifteen
named pieces, and---reading the same trace twice, once for the driven filter and
once for the simpler electrostatic capacitor---confirm the book's central claim
in motion: one simulation is one composition of the same handful of machines.
This is the synthesis of Parts~I through~V, and the bridge into Part~VI, where
that very composition is rendered as code.}

\section{One simulation, named all the way down}

We have spent five parts disassembling Palace. \Cref{part:modalities} took the
six analyses apart one at a time; before it, Part~IV built the numerical engines,
Part~III the calculus that names them, and Part~II the field theory they
discretize. Each chapter handed us a part and a name: \code{fe\_assemble},
\code{solve\_family}, \code{frequency\_sweep}, \code{ksp\_solve},
\code{gram\_reduce}, \code{sparameter\_reduce}, the V-cycle, the Chebyshev
smoother, the matrix-free apply. We have, by now, a complete vocabulary. What we
have \emph{not} yet done is run one real simulation slowly enough to point at
every part as it engages.

That is this chapter's only job. We will not derive a new equation or define a new
combinator. We will choose one problem, walk it through the lifecycle of
\cref{ch:lifecycle} stage by stage, and at each stage say: \emph{this is the
machine from that chapter, and here is what it does to the data right now}. The
value of the chapter is in its cross-references---it is the index of the whole
book's machinery caught in the act.

\begin{keyidea}
A complete Palace simulation is not a monolith. It is one \emph{composition} of a
small, fixed set of named machines---a dozen-and-a-half pieces, each built in an
earlier chapter, each with a precise type and a precise job. To ``understand a
driver'' is exactly to know which pieces it composes and in what order. By the end
of this chapter you will have seen the entire machine in a single expression, and
every box in that expression will be a chapter you have already read.
\end{keyidea}

\subsection{The problem we will trace}

We trace the analysis that exercises the most machinery: the \emph{driven
frequency sweep} of a two-port microwave \emph{filter} (\cref{ch:driven}). It is
the richest of the six because, end to end, it touches almost everything the book
built: it meshes a three-dimensional geometry, lays a Nédélec $\Hcurl$ space over
it, assembles three operators by the weak-form fold, forms a fresh complex
operator at every frequency, solves each by preconditioned GMRES wrapped around a
multigrid V-cycle whose smoother is Chebyshev/Hiptmair and whose every apply is
matrix-free, and finally projects the solved fields onto port modes borrowed from
a boundary-mode analysis. A single driven run is a guided tour of Parts~II--V.

Where the driven path grows too dense to read, we run the \emph{electrostatic
capacitance} extraction (\cref{ch:electrostatics}) beside it as the clean
contrast. It walks the identical lifecycle but lands in the simplest cells---the
fixed-operator map, the symmetric Gram---so it shows the same skeleton stripped to
its bones. Two analyses, one machine: the comparison is the chapter's spine.

\begin{physicsbox}
\textbf{The device.} An engineer hands us a two-port bandpass filter: a small
metal structure---coupled resonators on a substrate---with a coaxial or
waveguide \emph{port} at each end (\cref{fig:trace-pipeline}, far left). The
question is the network-analyzer question of \cref{ch:driven}: across a band of
frequencies $f \in [f_{\min}, f_{\max}]$, how much of a wave injected at port~1
is \emph{reflected} ($S_{11}$) and how much is \emph{transmitted} to port~2
($S_{21}$)? The answer is the scattering matrix $S(\omega)$ as a function of
frequency---the passband, the roll-off, the match. We will compute it by solving
the time-harmonic Maxwell field inside the filter at each frequency and reading
the scattered waves off the solved field.
\end{physicsbox}

\section{The lifecycle, walked once, naming every machine}
\label{sec:trace-walk}

Recall the lifecycle of \cref{ch:lifecycle}: parse, configure, dispatch, build the
mesh, solve, (optionally refine), output. We now walk those stages \emph{for this
one run}, and at each stage we name the component---and the chapter---doing the
work. \Cref{fig:trace-pipeline} is the whole walk in one annotated picture; keep
it in view, because the subsections below are simply its boxes, read left to right.

\begin{figure}[p]
\centering
\begin{tikzpicture}[node distance=6mm and 6mm, font=\footnotesize, >=Stealth]
  % the spine, left to right
  \node[data, align=center, text width=17mm] (cfg)
    {config file\\[1pt]\scriptsize geometry,\\$\varepsilon,\mu$, ports,\\freq.\ band};
  \node[stage, right=of cfg, align=center, text width=20mm] (io)
    {\textbf{parse}\\[1pt]\scriptsize \code{IoData}\\record};
  \node[stage, right=of io, align=center, text width=18mm] (mesh)
    {\textbf{mesh}\\[1pt]\scriptsize \code{build\_mesh}};
  \node[stage, right=of mesh, align=center, text width=18mm, fill=simBgViolet, draw=simViolet] (disp)
    {\textbf{dispatch}\\[1pt]\scriptsize driven driver};
  \node[stage, below=14mm of io, align=center, text width=21mm] (asm)
    {\textbf{assemble}\\[1pt]\scriptsize \code{fe\_assemble}$^{\times3}$\\$\{\Kop,\Cop,\Mop\}$};
  \node[stage, right=of asm, align=center, text width=24mm, fill=simBgRust, draw=simRust] (solve)
    {\textbf{solve}\\[1pt]\scriptsize \code{frequency\_sweep}\\(per $\omega$: GMRES)};
  \node[stage, right=of solve, align=center, text width=21mm, fill=simBgGold, draw=simGold] (red)
    {\textbf{reduce}\\[1pt]\scriptsize \code{sparameter\_}\\\code{reduce}};
  \node[product, right=of red, align=center, text width=15mm] (prod)
    {$S(\omega)$\\[1pt]\scriptsize S-param\\plot};
  % top-row forward flow
  \draw[flow] (cfg) -- (io);
  \draw[flow] (io) -- (mesh);
  \draw[flow] (mesh) -- (disp);
  % down into the solve band
  \draw[flow] (disp.south) |- (asm.north west);
  % the solve band, left to right
  \draw[flow] (asm) -- node[above,font=\scriptsize]{\code{fam}} (solve);
  \draw[flow] (solve) -- node[above,font=\scriptsize]{$[\vE_k]$} (red);
  \draw[flow] (red) -- (prod);
  % chapter tags under each box
  \node[cornertag, below=1mm of io]   {\cref{ch:records}};
  \node[cornertag, below=1mm of mesh] {\cref{ch:functionspaces}};
  \node[cornertag, below=1mm of disp] {\cref{ch:lifecycle}};
  \node[cornertag, below=1mm of asm]  {\cref{ch:assembly}, \cref{ch:weak}};
  \node[cornertag, below=1mm of solve]{\cref{ch:gmres}, \cref{ch:multigrid}};
  \node[cornertag, below=1mm of red]  {\cref{ch:reductions}, \cref{ch:waveports}};
  % the port modes feeding the reduction from the side
  \node[extern, above=8mm of red, align=center, text width=20mm] (ports)
    {port modes $\vs_i$\\[1pt]\scriptsize boundary-mode\\(\cref{ch:waveports})};
  \draw[refedge] (ports.south) -- (red.north);
\end{tikzpicture}
\caption[The complete driven pipeline, annotated by chapter]{The capstone
diagram: one complete driven simulation, from the configuration file to the
S-parameter plot, with every stage tagged by the chapter that built it. The
top row is the shared lifecycle set-up (parse the \code{IoData} record, build the
mesh, dispatch on the problem type, \cref{ch:lifecycle}); the bottom band is the
driver's \texttt{assemble--solve--reduce} body (\cref{ch:situating}). The
\emph{assemble} box folds three weak-form term lists into the fixed basis
$\{\Kop,\Cop,\Mop\}$ (\cref{ch:assembly}); the \emph{solve} box is the
operator-varying frequency sweep, each frequency a preconditioned GMRES solve
(\cref{ch:gmres}) over a multigrid V-cycle (\cref{ch:multigrid}); the
\emph{reduce} box projects the solved fields onto the port modes that the
boundary-mode analysis supplies from the side (\cref{ch:waveports}). Every box is
a machine from an earlier part; this chapter watches them all run together.}
\label{fig:trace-pipeline}
\end{figure}

\subsection{Parse the configuration: the \texorpdfstring{\code{IoData}}{IoData} record}
\label{sec:trace-parse}

The run begins with one file the engineer wrote: a configuration naming the mesh,
the polynomial order, the materials ($\varepsilon$ and $\mu$ per region), the two
ports and their excitations, the frequency band to sweep, and---the field that
decides everything downstream---the problem type, \code{Driven}. Parsing it
produces the single in-memory value the rest of the run consults: the
\code{config}, which is the \code{IoData} record of \cref{ch:records}.

The decisive property, established in \cref{ch:lifecycle} and \cref{ch:strata}, is
that this record is \emph{read-only}. It is built once and never modified; every
later stage reads it, none writes it. It is \emph{construction-stratum} data
(\cref{ch:strata}): fixed the instant parsing finishes, so any computation that
depends only on it is referentially stable for the whole run. That single fact is
what licenses every hoist, reuse, and reorder we are about to watch.

\begin{notationbox}
We write \code{config} for the parsed \code{IoData} value. It is a record
(\cref{ch:records})---a bundle of named fields read with a dot:
\lstinline[style=l4style]|config.problem_type|, the materials
\lstinline[style=l4style]|config.materials|, the swept band
\lstinline[style=l4style]|config.frequencies|, the ports
\lstinline[style=l4style]|config.ports|. The configuration's records get their
definitional home in \cref{ch:records}; here we simply read fields off it. The
two analyses we trace read the \emph{same} record type and differ only in which
fields they consult---the electrostatic path reads \code{config.terminals} where
the driven path reads \code{config.frequencies} and \code{config.ports}.
\end{notationbox}

\subsection{Configure the device, then dispatch}

Two short stages follow, both from \cref{ch:lifecycle}. \emph{Configure the
device} brings up the CPU or GPU backend the element kernels will run on---the
choice read from \code{config}, the computation downstream expressed against an
abstract device so the same description runs on either. Then \emph{dispatch}: the
run reads \code{config.problem\_type}, finds \code{Driven}, and on the strength of
that one value selects the \emph{driven driver}---the object that knows how to run
this one analysis end to end. This is the \emph{specialization seam} of
\cref{ch:lifecycle}: everything above it is shared by all six analyses; below it,
the driven driver fills in the one stage that differs. Our electrostatic contrast
takes the same seam to the \emph{electrostatic} driver instead---one field, one
\code{switch}, the entire difference of \emph{which} analysis runs.

\subsection{Build the mesh, and lay the \texorpdfstring{$\Hcurl$}{H(curl)} space}
\label{sec:trace-mesh}

With a driver chosen, the run builds the computational mesh: it loads the geometry
named in \code{config}, applies any preprocessing, partitions it, and performs any
a-priori refinement requested. The framework verb is \code{build\_mesh}, and the
value it returns---the \code{Mesh}---is the first value the lifecycle
\emph{produces} rather than reads (the config came from the user; the mesh is
computed). Mesh construction is driver-agnostic in shape: both our analyses build
the mesh identically.

Over that mesh the driver lays its \emph{finite-element space} (\cref{ch:functionspaces}).
Here the two analyses first diverge in physics. The driven analysis solves for a
vector field $\vE$ that must support tangential continuity and a curl, so it lives
in the Nédélec $\Hcurl$ space---the edge-element slot of the de Rham complex
(\cref{ch:derham}). The electrostatic analysis solves for a scalar potential $V$,
so it lives in the nodal $\Hone$ space---the vertex slot. The choice of space is
\emph{one cell} of the grid of \cref{ch:situating}, and it is the first cell our
two traces fill differently: $\Hcurl$ for driven, $\Hone$ for electrostatic.

\begin{remark}[The adaptive loop, deferred]
The lifecycle wraps the solve in the adaptive estimate--mark--refine loop of
\cref{ch:amr}---the \emph{outer fold} whose threaded value is the mesh. For our
trace we keep refinement off, so the fold runs its body exactly once and the
lifecycle is the straight line of \cref{ch:lifecycle}. Everything below is
\emph{one pass} of that fold. With refinement on, the entire
assemble--solve--reduce body we are about to walk simply runs again on a finer
mesh, each pass beginning from the mesh the last produced---the
\code{fold\_solve} of \cref{ch:fold} wrapped around all of it.
\end{remark}

\subsection{Assemble the basis, once: the weak-form fold}
\label{sec:trace-assemble}

Now the driver's body begins, and the first stage is \emph{assemble}. The driven
analysis needs three operators on $\Hcurl$, the same three characters from every
chapter: the curl--curl stiffness $\Kop$ (built from the reluctivity $\nu=1/\mu$),
the damping $\Cop$ (from surface impedance / conductivity), and the mass $\Mop$
(from permittivity $\varepsilon$). Each is the result of the assembly fold
\code{fe\_assemble} of \cref{ch:assembly}, summing one contribution per weak-form
term (\cref{ch:weak}) over the mesh, with the essential boundary degrees of
freedom eliminated (\cref{ch:bc}):
\begin{lstlisting}[style=l4style]
-- the fixed basis: three weak-form folds, run ONCE, before any sweep
space = nd_space cfg                                     -- H(curl) Nedelec space
fam   = { K  = fe_assemble space [ curl_curl (reluctivity cfg) ]   -- stiffness
        , C  = fe_assemble space [ damping cfg ]                   -- damping
        , M  = fe_assemble space [ mass (permittivity cfg) ] }     -- mass
\end{lstlisting}
Three structural facts, each a load-bearing reference. First, every
\code{fe\_assemble} is the \emph{fold} of \cref{ch:assembly}---a loop over the term
list that we do not write; the framework folds it. Second, each operator's matrix
action, when the solver later applies it, is \emph{matrix-free}
(\cref{ch:matrixfree}): the operator is never stored as an explicit sparse matrix
but applied as a sequence of element-local tensor contractions, the libCEED-style
quadrature kernel. Third---and this is the word from \cref{ch:electrostatics} that
governs the whole solve---the basis is assembled \emph{once}, before the sweep, and
held. The electrostatic contrast is the same fold run once over a \emph{single}
term, the $\varepsilon$-weighted gradient, producing one operator $\Kop$.

\begin{keyidea}
The \emph{assemble} stage is the weak-form fold \code{fe\_assemble}
(\cref{ch:assembly}), run once per operator before any solving. The driven
analysis folds three term lists into a fixed basis $\{\Kop,\Cop,\Mop\}$; the
electrostatic analysis folds one into $\Kop$. The operators' applies are
matrix-free (\cref{ch:matrixfree}); the essential boundary data is eliminated
(\cref{ch:bc}); the result is read-only and reused across the entire solve. ``Once''
is the structural signature shared by both analyses---it is the upstream half of
the hoist.
\end{keyidea}

\subsection{Solve: the operator-varying sweep, and its nested engine}
\label{sec:trace-solve}

The solve stage is where the driven analysis spends almost all of its time, and
where the most machinery engages at once. Recall from \cref{ch:driven} that the
system operator is \emph{not} fixed: at each frequency $\omega$ in the swept band
we form a fresh complex operator from the fixed basis,
\begin{equation}
  \mat{A}(\omega) \;=\; \Kop \;+\; i\omega\,\Cop \;-\; \omega^2\Mop,
  \label{eq:trace-Aomega}
\end{equation}
solve $\mat{A}(\omega)\,\vE = \vb(\omega)$ for the phasor field, and move on. The
solves are \emph{independent}---the field at one frequency needs nothing from
another---so the sweep is a \emph{map}, the operator-varying map (corner~2) of
\cref{ch:situating}. The framework verb is \code{frequency\_sweep}, and forming
\cref{eq:trace-Aomega} from the basis is the \code{linear\_combination} combinator
of \cref{ch:combinators}, specialized to operator operands and complex scalars.

What makes this stage worth a figure of its own is that each single
\code{ksp\_solve} inside the sweep is itself a \emph{stack} of the machines from
Part~IV, nested four deep (\cref{fig:trace-solverstack}):

\begin{enumerate}[nosep]
  \item \textbf{GMRES} (\cref{ch:gmres}) drives the outer iteration. The operator
  $\mat{A}(\omega)$ is complex, non-Hermitian, and indefinite---the $i\omega\Cop$
  term makes it complex, the $-\omega^2\Mop$ term makes it indefinite---so none of
  CG's structure survives and GMRES is the right Krylov engine (\cref{ch:cg}
  versus \cref{ch:gmres}).
  \item \textbf{A multigrid V-cycle} (\cref{ch:multigrid}) is the
  \emph{preconditioner} GMRES calls each iteration: restrict the residual down the
  mesh hierarchy, solve cheaply on the coarsest level, prolong the correction back
  up.
  \item \textbf{A Chebyshev / Hiptmair smoother} (\cref{ch:smoothers}) runs at
  each level of the V-cycle, damping the high-frequency error the coarse grid
  cannot see---the Hiptmair variant handling the $\Hcurl$ null-space that an
  ordinary smoother would miss.
  \item \textbf{A matrix-free apply} (\cref{ch:matrixfree}) is the innermost
  operation of all of them. Every GMRES Arnoldi step, every smoother sweep, every
  residual evaluation ultimately \emph{applies $\mat{A}(\omega)$ to a vector}, and
  that apply is the element-local tensor contraction of \cref{ch:matrixfree}---no
  assembled sparse matrix ever materializes.
\end{enumerate}

Threaded through all four levels is the solver \emph{state}---the Krylov basis, the
running residual, the iteration counters---carried by the solve monad of
\cref{ch:monad} as the \code{SimState} record (\cref{ch:records}), so that the
mutation the C++ performs in place is, in our reading, a pure state thread.

\begin{figure}[p]
\centering
\begin{tikzpicture}[font=\footnotesize, >=Stealth]
  % nested boxes: GMRES > V-cycle > smoother > matrix-free apply
  \node[stage, minimum width=72mm, minimum height=52mm, fill=simBgRust,
        draw=simRust, align=center] (g) at (0,0) {};
  \node[anchor=north, font=\small\bfseries, text=simRust] at (g.north)
    {\ GMRES (\cref{ch:gmres})};
  \node[font=\scriptsize\itshape, text=simRust, anchor=north]
    at ($(g.north)+(0,-0.5)$) {outer Krylov: complex, non-Hermitian $\mat{A}(\omega)$};

  \node[stage, minimum width=60mm, minimum height=36mm, fill=simBgBlue,
        draw=simBlue, align=center] (v) at (0,-0.55) {};
  \node[anchor=north, font=\footnotesize\bfseries, text=simBlue] at (v.north)
    {\ multigrid V-cycle (\cref{ch:multigrid})};
  \node[font=\scriptsize\itshape, text=simBlue, anchor=north]
    at ($(v.north)+(0,-0.42)$) {preconditioner: restrict $\to$ coarse-solve $\to$ prolong};

  \node[stage, minimum width=46mm, minimum height=20mm, fill=simBgGreen,
        draw=simGreen, align=center] (s) at (0,-1.1) {};
  \node[anchor=north, font=\footnotesize\bfseries, text=simGreen] at (s.north)
    {\ Chebyshev / Hiptmair smoother};
  \node[font=\scriptsize\itshape, text=simGreen, anchor=north]
    at ($(s.north)+(0,-0.4)$) {(\cref{ch:smoothers}): damp high-freq.\ error};

  \node[extern, minimum width=32mm, minimum height=7mm] (mf) at (0,-1.75)
    {matrix-free apply (\cref{ch:matrixfree})};

  % a side label: the threaded state
  \node[opbox, anchor=west, align=center, text width=24mm, fill=white]
    (state) at (4.6,1.4) {\code{SimState}\\\scriptsize threaded by the\\solve monad\\(\cref{ch:monad}, \cref{ch:records})};
  \draw[refedge] (state.south) .. controls +(0,-1.4) and +(0,1.2) .. (g.east);
  % apply call-down annotation
  \node[font=\scriptsize\itshape, text=simViolet, anchor=west] at (2.7,-1.75)
    {$\leftarrow$ every level};
  \node[font=\scriptsize\itshape, text=simViolet, anchor=west] at (2.7,-2.15)
    {\ \ bottoms out here};
\end{tikzpicture}
\caption[The solver stack for one frequency]{The solver stack engaged by a
\emph{single} frequency of the sweep---four machines from Part~IV, nested. The
outermost is preconditioned \textbf{GMRES} (\cref{ch:gmres}), the only Krylov
method that handles the complex, non-Hermitian, indefinite $\mat{A}(\omega)$. Its
preconditioner is a \textbf{multigrid V-cycle} (\cref{ch:multigrid}), which at each
mesh level calls a \textbf{Chebyshev/Hiptmair smoother} (\cref{ch:smoothers}) to
damp the error the coarse grid cannot resolve. Beneath all of them, every
operator application---every Arnoldi step, every smoother sweep, every residual---is
a \textbf{matrix-free apply} (\cref{ch:matrixfree}): an element-local tensor
contraction, never an assembled sparse matrix. The whole stack's mutable state (the
Krylov basis, residual, counters) is the \code{SimState} record threaded by the
solve monad (\cref{ch:monad}). And this entire stack runs \emph{once per
frequency}, fresh operator each time.}
\label{fig:trace-solverstack}
\end{figure}

The electrostatic contrast collapses this stack dramatically, and the collapse is
instructive. Its operator $\Kop$ is symmetric positive-definite and \emph{fixed},
so the outer Krylov method is \emph{CG}, not GMRES (\cref{ch:cg}); the
preconditioner is built \emph{once} and reused across the whole family (the hoist
of \cref{ch:electrostatics}); and the sweep is the fixed-operator map
\code{solve\_family} (corner~1) rather than \code{frequency\_sweep} (corner~2). The
inner machines---a V-cycle, a smoother, a matrix-free apply---are the \emph{same
machines}; only the outer Krylov method and the position of the operator capture
change. \emph{One stack, two configurations.}

\begin{pitfall}
The single line whose position separates the two analyses is the operator capture
inside \code{frequency\_sweep}. In electrostatics it is \emph{hoisted} above the
loop (the operator never changes); in the driven sweep it lives \emph{inside} the
loop (the operator changes every frequency). Moving it the wrong way in the driven
sweep does not slow the code---it solves every frequency against
$\mat{A}(\omega_0)$, the first frequency's operator, and returns a smooth,
plausible, \emph{wrong} curve (\cref{ch:driven}). The loop position of one capture
call is part of the algorithm, not a performance knob.
\end{pitfall}

\subsection{Reduce: the port-projection, with borrowed modes}
\label{sec:trace-reduce}

The sweep hands the reduction a family of complex fields $[\vE_k]$, one per swept
frequency. The reduction collapses each into the scattering parameters. This is
the rank-2 \emph{port-projection} of \cref{ch:reductions}---reduction shape~2 of
\cref{ch:situating}, and emphatically \emph{not} a Gram. The verb is
\code{sparameter\_reduce}, applied once per frequency:
\begin{equation}
  S_{ij}(\omega) \;=\; \sigma_{ij}\,\bigl(\inner{\vs_i}{\vE_j(\omega)} - \delta_{ij}\bigr),
  \label{eq:trace-sparam}
\end{equation}
the projection of the drive-$j$ field onto receiver mode $i$, minus the Kronecker
self-term on the diagonal (which converts ``total field at the driven port'' into
``reflected field''), scaled by the port-kind factor $\sigma_{ij}$.

The one component this stage \emph{consumes from elsewhere} is the set of port
modes $\vs_i$. They do not come from the driven solve; they come from the
\emph{boundary-mode analysis} of \cref{ch:waveports}, which solves a small
two-dimensional eigenproblem on each port's cross-section to find the field
pattern that propagates there. In the pipeline diagram
(\cref{fig:trace-pipeline}) this is the arrow entering the reduce box from the
side: the driven driver \emph{composes} the boundary-mode driver's product as one
of its inputs. The electrostatic contrast instead reduces with the symmetric Gram
\code{gram\_reduce} (\cref{ch:reductions}), $C_{ij}=V_j^{\!\top}\Kop V_i$, pairing
the solved family \emph{with itself} through the energy operator---no borrowed
basis, symmetric by construction.

\subsection{Output}

When the reduction finishes, the driver writes its product---the table of
$S_{ij}(\omega)$ across the band---to disk, and the top-level lifecycle writes the
run metadata around it (\cref{ch:lifecycle}). Plotted, $S(\omega)$ is the two
familiar microwave curves: the return loss $\abs{S_{11}}$ in dB and the insertion
loss $\abs{S_{21}}$ in dB (\cref{fig:trace-splot}). That plot is the answer the
engineer ran the simulation to obtain.

\section{The whole machine in one expression}
\label{sec:trace-composition}

We have walked the stages. Now collapse the walk into a single composition and
\emph{expand it}, so the reader sees the entire machine---every nested engine---in
one tree.

\subsection{The three-stage composition}

At the top level the driven driver is the three-stage composition of
\cref{ch:driven}:
\begin{lstlisting}[style=l4style]
driven = sparameter_reduce . frequency_sweep . fe_assemble
\end{lstlisting}
read right to left: assemble the basis once, sweep the operator-varying map over
the band, reduce by port-projection. Spelled out with its data, exactly the
composition of \cref{ch:driven}:
\begin{lstlisting}[style=l4style]
-- inputs = config; output = the frequency response S(omega)
driven :: DrivenConfig -> FrequencyResponse
driven cfg =
  let space  = nd_space cfg                                       -- H(curl) space
      fam    = { K = fe_assemble space [ curl_curl (reluctivity cfg) ]  -- (1) assemble
               , C = fe_assemble space [ damping cfg ]                  --     the basis
               , M = fe_assemble space [ mass (permittivity cfg) ] }    --     ONCE
      omegas = sample_frequencies cfg                             -- the swept band
      es     = frequency_sweep fam omegas       -- (2) operator-VARYING map: A(w) rebuilt inside
  in  sparameter_reduce (ports cfg) es          -- (3) port-projection -> S(omega)
\end{lstlisting}

\subsection{Expanded all the way down}

The three boxes hide the whole of Part~IV. Expand each piece and the entire call
tree appears---the assemble fold, the per-$\omega$ solve with its full
preconditioner stack, the reduction over borrowed modes. This is the hero figure
of the chapter: the whole simulator in one tree (\cref{fig:trace-tree}).

\begin{lstlisting}[style=l4style]
driven cfg
|
+- (1) fe_assemble x3              -- the basis, folded ONCE   [ch:assembly, ch:weak]
|     +- fold over weak-form terms                            -- [ch:assembly]
|     +- element-local quadrature  -- matrix-free contraction  [ch:matrixfree]
|     +- eliminate essential BC dofs                          -- [ch:bc]
|
+- (2) frequency_sweep fam omegas  -- map over the band (corner 2) [ch:driven]
|     +- FOR EACH omega:           -- independent: a map, not a fold
|         +- A(w) = linear_combination [(1,K),(iw,C),(-w^2,M)]  -- [ch:combinators]
|         +- b    = excitation fam omega                       -- per-omega RHS
|         +- ksp_solve A(w) b      -- one GMRES solve           [ch:gmres]
|             +- GMRES outer iteration                          -- [ch:gmres]
|             |   +- precondition: multigrid V-cycle            -- [ch:multigrid]
|             |   |   +- restrict / coarse-solve / prolong
|             |   |   +- smoother (Chebyshev / Hiptmair)        -- [ch:smoothers]
|             |   |       +- matrix-free apply of A(w)          -- [ch:matrixfree]
|             |   +- Arnoldi step: matrix-free apply of A(w)    -- [ch:matrixfree]
|             +- state threaded by the solve monad (SimState)  -- [ch:monad, ch:records]
|
+- (3) sparameter_reduce ports es  -- port-projection (shape 2) [ch:reductions]
      +- S_ij = scale * (project s_i E_j - selfterm i j)       -- [ch:reductions]
      +- port modes s_i FROM boundary-mode analysis            -- [ch:waveports]
\end{lstlisting}

\begin{figure}[p]
\centering
\begin{tikzpicture}[font=\scriptsize, >=Stealth,
    grow=down, level distance=8mm,
    edge from parent/.style={draw=simGray, semithick, -{Stealth[length=1.6mm]}},
    every node/.style={align=center}]
  \node[stage, fill=simBgViolet, draw=simViolet, text width=26mm]
    {\code{driven cfg}\\\scriptsize the whole machine}
    child {node[stage, text width=24mm] {\textbf{(1)} \code{fe\_assemble}$^{\times3}$\\\scriptsize basis, once}
      child {node[opbox, text width=22mm] {fold over\\weak-form terms\\\scriptsize\cref{ch:assembly}}}
      child {node[extern, text width=20mm] {matrix-free\\quadrature\\\scriptsize\cref{ch:matrixfree}}}
    }
    child {node[stage, fill=simBgRust, draw=simRust, text width=30mm] {\textbf{(2)} \code{frequency\_sweep}\\\scriptsize map over $\omega$ (corner 2)}
      child {node[opbox, text width=26mm] {$\mat{A}(\omega)=$\code{lin\_comb}\\\scriptsize\cref{ch:combinators}}}
      child {node[stage, fill=simBgRust, draw=simRust, text width=22mm] {\code{ksp\_solve}\\\scriptsize GMRES \cref{ch:gmres}}
        child {node[stage, fill=simBgBlue, draw=simBlue, text width=22mm] {V-cycle\\precond.\\\scriptsize\cref{ch:multigrid}}
          child {node[stage, fill=simBgGreen, draw=simGreen, text width=20mm] {smoother\\\scriptsize\cref{ch:smoothers}}
            child {node[extern, text width=20mm] {matrix-free\\apply\\\scriptsize\cref{ch:matrixfree}}}
          }
        }
      }
    }
    child {node[stage, fill=simBgGold, draw=simGold, text width=26mm] {\textbf{(3)} \code{sparameter\_reduce}\\\scriptsize projection (shape 2)}
      child {node[opbox, text width=22mm] {$S_{ij}$ projection\\\scriptsize\cref{ch:reductions}}}
      child {node[extern, text width=20mm] {port modes $\vs_i$\\\scriptsize\cref{ch:waveports}}}
    };
\end{tikzpicture}
\caption[The driven driver, expanded to its leaves]{The hero figure: the entire
driven driver expanded from the top-level composition down to its matrix-free
leaves. The three top-level stages---\code{fe\_assemble} (assemble the basis once),
\code{frequency\_sweep} (the operator-varying solve map), and
\code{sparameter\_reduce} (the port-projection)---each unfold into the machines
that implement them. The solve branch is the deepest: a GMRES outer iteration
(\cref{ch:gmres}) preconditioned by a multigrid V-cycle (\cref{ch:multigrid})
whose smoother (\cref{ch:smoothers}) bottoms out, like every operator application
in the tree, in a matrix-free apply (\cref{ch:matrixfree}). Read the whole tree
and you have read the whole book: every leaf is a chapter. This is exactly the
``synthesized library'' view that Part~VI (\cref{ch:synthlib}, \cref{ch:composing})
renders as code.}
\label{fig:trace-tree}
\end{figure}

\begin{keyidea}
The whole simulator is \emph{one expression}. \texttt{driven $=$
sparameter\_reduce $\circ$ frequency\_sweep $\circ$ fe\_assemble}, and each of
those three pieces expands into the machines of Parts~III--IV until every leaf is a
named component from an earlier chapter. There is no hidden glue, no special-case
plumbing: a driver \emph{is} its composition tree, and the tree is finite and
shallow. You have now seen the entire machine in one figure---and you have read
every box in it.
\end{keyidea}

\section{Two analyses, one machine}
\label{sec:trace-twoanalyses}

The driven trace is one path through the machinery. The chapter's claim is that the
\emph{same} machinery, recomposed, is every other analysis too. We make that
concrete by setting the electrostatic capacitance trace beside the driven filter
trace and reading off what is \emph{shared} and what is \emph{swapped}
(\cref{fig:trace-contrast}).

\begin{figure}[t]
\centering
\begin{tikzpicture}[font=\footnotesize, >=Stealth, node distance=5mm and 7mm]
  % ===== DRIVEN row =====
  \node[font=\small\bfseries, text=simRust, anchor=west] at (-0.3,2.0) {Driven filter};
  \node[stage, align=center, text width=15mm] (dasm) at (1.4,1.4) {assemble\\\scriptsize$\{\Kop,\Cop,\Mop\}$};
  \node[stage, fill=simBgRust, draw=simRust, align=center, text width=18mm] (dsolve) at (4.0,1.4)
    {\code{frequency\_}\\\code{sweep}\\\scriptsize GMRES, corner~2};
  \node[stage, fill=simBgGold, draw=simGold, align=center, text width=17mm] (dred) at (6.7,1.4)
    {\code{sparameter\_}\\\code{reduce}};
  \node[product, align=center, text width=11mm] (dprod) at (8.9,1.4) {$S(\omega)$};
  \draw[flow] (dasm) -- (dsolve); \draw[flow] (dsolve) -- (dred); \draw[flow] (dred) -- (dprod);
  % ===== ELECTROSTATIC row =====
  \node[font=\small\bfseries, text=simBlue, anchor=west] at (-0.3,-0.3) {Electrostatic capacitor};
  \node[stage, align=center, text width=15mm] (easm) at (1.4,-0.9) {assemble\\\scriptsize$\Kop$ only};
  \node[stage, fill=simBgBlue, draw=simBlue, align=center, text width=18mm] (esolve) at (4.0,-0.9)
    {\code{solve\_family}\\\scriptsize CG, corner~1};
  \node[stage, fill=simBgGold, draw=simGold, align=center, text width=17mm] (ered) at (6.7,-0.9)
    {\code{gram\_reduce}\\\scriptsize $w{=}1$};
  \node[product, align=center, text width=11mm] (eprod) at (8.9,-0.9) {$\mat{C}$};
  \draw[flow] (easm) -- (esolve); \draw[flow] (esolve) -- (ered); \draw[flow] (ered) -- (eprod);
  % ===== shared / swapped annotations =====
  \node[font=\scriptsize\itshape, text=simGreen, align=center] at (1.4,0.25)
    {\textbf{shared:}\\\code{fe\_assemble} fold\\(\cref{ch:assembly})};
  \node[font=\scriptsize\itshape, text=simRust, align=center] at (4.0,0.25)
    {\textbf{swapped:}\\corner~1 $\leftrightarrow$ 2\\CG $\leftrightarrow$ GMRES};
  \node[font=\scriptsize\itshape, text=simRust, align=center] at (6.7,0.25)
    {\textbf{swapped:}\\Gram $\leftrightarrow$ projection};
  % shared inner-stack note
  \node[font=\scriptsize\itshape, text=simGreen, align=center] at (4.0,-2.1)
    {\textbf{shared inner stack:} V-cycle (\cref{ch:multigrid}) $\cdot$ smoother (\cref{ch:smoothers}) $\cdot$ matrix-free apply (\cref{ch:matrixfree})};
\end{tikzpicture}
\caption[Two analyses, one machine]{The driven filter and the electrostatic
capacitor as two recompositions of one machine. \textbf{Shared} (green): both
build their operators with the \emph{same} assembly fold \code{fe\_assemble}
(\cref{ch:assembly}), and both solves bottom out in the \emph{same} inner
stack---a multigrid V-cycle (\cref{ch:multigrid}) over a smoother
(\cref{ch:smoothers}) over a matrix-free apply (\cref{ch:matrixfree}).
\textbf{Swapped} (rust): the solve \emph{corner} (fixed-operator map, corner~1,
versus operator-varying map, corner~2) and with it the outer Krylov method (CG
versus GMRES); and the \emph{reduction} (symmetric Gram versus port-projection).
Two boxes move; the skeleton, and most of the inner machinery, is identical. This
is the claim of \cref{ch:situating} caught in a single picture.}
\label{fig:trace-contrast}
\end{figure}

What is \emph{shared} is most of the machine. Both analyses parse the same
\code{IoData} record, build the mesh the same way, dispatch through the same seam,
assemble with the same fold \code{fe\_assemble}, and---this is the part worth
dwelling on---bottom out in the \emph{same} inner solver stack: a multigrid
V-cycle over a smoother over a matrix-free apply. The expensive numerical core is
identical; the analyses do not each carry their own private solver.

What is \emph{swapped} is small and well-localized: the function space ($\Hcurl$
versus $\Hone$), the solve corner (operator-varying map versus fixed-operator map,
and with it GMRES versus CG), and the reduction (port-projection versus symmetric
Gram). This is precisely the ``one stage swapped'' punchline of
\cref{ch:situating} and \cref{ch:electrostatics}, now visible as two real traces
side by side rather than a table of dashes.

\begin{keyidea}
The driven filter and the electrostatic capacitor are the \emph{same machine} in
two configurations. They share the lifecycle, the \code{IoData} record, the mesh
build, the assembly fold, and---crucially---the entire inner solver stack
(V-cycle, smoother, matrix-free apply). They differ in exactly three swapped
cells: the function space, the solve corner (with its Krylov method), and the
reduction. A different analysis again would swap a different few cells and reuse
everything else. The simulator is not six programs; it is one composition with a
handful of swappable parts.
\end{keyidea}

\section{Counting the corners and shapes this one run touched}
\label{sec:trace-count}

Step back to the taxonomy of \cref{ch:situating} and ask: of the four solve
corners and three reduction shapes, how many did this \emph{single} driven run
use? The answer makes the economy of the design tangible (\cref{fig:trace-grid}).

The driven run touched \emph{one} solve corner---the operator-varying map
(corner~2)---and \emph{one} reduction shape---the rank-2 port-projection (shape~2).
It used, inside that one corner, a stack of Krylov and multigrid layers, and it
\emph{borrowed} one product from another analysis (the port modes from the
black-box boundary-mode corner, corner~4). That is the entire taxonomic footprint
of a complete microwave simulation: one corner, one reduction, a shared inner
stack, one borrowed input.

\begin{figure}[t]
\centering
\begin{tikzpicture}[font=\footnotesize]
  % the 2x2 corners grid + the table, with THIS run's cells highlighted
  \node[cornertag, anchor=south] at (0.9,2.5) {operator \emph{fixed}};
  \node[cornertag, anchor=south] at (3.5,2.5) {operator \emph{varying}};
  \node[cornertag, rotate=90, anchor=south] at (-1.4,1.8) {\textbf{map}};
  \node[cornertag, rotate=90, anchor=south] at (-1.4,0.4) {\textbf{fold}};
  % top-left: fixed map (electrostatic uses; driven does NOT)
  \node[stage, minimum width=20mm, minimum height=12mm, fill=simBgGray, draw=simGray,
        align=center] (c1) at (0.9,1.7) {\scriptsize fixed-op map\\\scriptsize(electrostatic)};
  % top-right: operator-varying map -- THIS RUN
  \node[stage, minimum width=20mm, minimum height=12mm, fill=simBgRust, draw=simRust,
        very thick, align=center] (c2) at (3.5,1.7) {\scriptsize \textbf{op-varying map}\\\scriptsize\textbf{(driven: HERE)}};
  % bottom-left: fold
  \node[stage, minimum width=20mm, minimum height=12mm, fill=simBgGray, draw=simGray,
        align=center] (c3) at (0.9,0.3) {\scriptsize state fold\\\scriptsize(transient/AMR)};
  % bottom-right: black box
  \node[extern, minimum width=20mm, minimum height=12mm, align=center]
        (c4) at (3.5,0.3) {\scriptsize black-box\\\scriptsize(borrowed: ports)};
  % the borrow arrow
  \draw[refedge, simViolet] (c4.north) -- node[right, font=\scriptsize, text=simViolet]{modes} (c2.south);
  % the reductions column, to the right
  \node[font=\scriptsize\bfseries, anchor=south] at (7.2,2.5) {reduction shape};
  \node[stage, minimum width=24mm, minimum height=7mm, fill=simBgGray, draw=simGray]
        (r1) at (7.2,1.75) {\scriptsize Gram (electrostatic)};
  \node[stage, minimum width=24mm, minimum height=7mm, fill=simBgGold, draw=simGold,
        very thick] (r2) at (7.2,0.95) {\scriptsize \textbf{projection (driven)}};
  \node[stage, minimum width=24mm, minimum height=7mm, fill=simBgGray, draw=simGray]
        (r3) at (7.2,0.15) {\scriptsize table ($f,Q$/energy)};
\end{tikzpicture}
\caption[The taxonomic footprint of one run]{What one complete driven run touches
on the grid of \cref{ch:situating}. Of the four solve corners (left $2\times2$),
the driven filter occupies exactly \emph{one}: the operator-varying map (corner~2,
rust). Of the three reduction shapes (right column), it uses exactly \emph{one}:
the rank-2 port-projection (gold). It \emph{borrows} the port modes from the
black-box corner's boundary-mode analysis (violet arrow). The electrostatic
contrast (gray) occupies the fixed-operator map and the Gram instead. A complete
microwave simulation is one corner, one reduction, a shared inner stack, and one
borrowed input---and a different analysis swaps only which cells light up. The grid
is small because the simulator is: a few corners $\times$ a few reductions $\times$
the shared field theory and calculus.}
\label{fig:trace-grid}
\end{figure}

This is the payoff of the whole taxonomy. A simulator that looked, from outside, like
six elaborate programs is, from inside, a \emph{few corners} $\times$ a \emph{few
reductions}, all resting on one shared field theory (Part~II) and one shared
calculus (Part~III). The driven run lit up one cell of each axis. The electrostatic
contrast lit up a different cell of each. Every Palace analysis is a choice of one
cell per axis over the same substrate---which is exactly why
\cref{ch:situating}'s grid has so few columns, and why this book could be written
at all.

\section{Two worked traces}
\label{sec:trace-worked}

We close the disassembly with two end-to-end traces small enough to follow by hand:
the driven filter at one frequency, and the electrostatic capacitor, each rendered
as a \emph{trace table} that names, at every stage, the component, the data shape,
and the chapter of origin.

\begin{workedexample}{The driven filter, one frequency, by data shape}{trace-driven}
We follow the data through a \emph{single} frequency $\omega = \omega_k$ of the
sweep---config to mesh to assemble to one GMRES solve to port projection to one
S-matrix entry---tracking the \emph{shape} of the value at each stage rather than
its numerical contents (the numerics are \cref{wex:opvary} and \cref{wex:sparam}
of \cref{ch:driven}; here we trace the \emph{pipeline}). Write $N$ for the number
of $\Hcurl$ degrees of freedom (millions, in practice) and $p=2$ for the number
of ports.

\medskip
\begin{center}\small
\setlength{\tabcolsep}{4pt}
\renewcommand{\arraystretch}{1.25}
\begin{tabular}{@{}ll>{\raggedright\arraybackslash}p{42mm}l@{}}
\toprule
Stage & Component & Data in $\to$ out & From \\
\midrule
parse    & \code{IoData} record      & file $\to$ \code{config}                  & \cref{ch:records} \\
mesh     & \code{build\_mesh}         & \code{config} $\to$ \code{Mesh}            & \cref{ch:functionspaces} \\
space    & \code{nd\_space}           & \code{Mesh} $\to$ $\Hcurl$ space (dim $N$) & \cref{ch:functionspaces} \\
assemble & \code{fe\_assemble}$^{\times3}$ & space $\to$ basis $\{\Kop,\Cop,\Mop\}$, each $N{\times}N$ & \cref{ch:assembly} \\
form     & \code{linear\_combination} & basis, $\omega_k$ $\to$ $\mat{A}(\omega_k)$ ($N{\times}N$, complex) & \cref{ch:combinators} \\
excite   & \code{excitation}          & \code{config}, $\omega_k$ $\to$ $\vb \in \Complex^N$ & \cref{ch:driven} \\
solve    & \code{ksp\_solve} (GMRES)  & $\mat{A}(\omega_k),\vb$ $\to$ $\vE_k \in \Complex^N$ & \cref{ch:gmres} \\
 \ \ precond. & V-cycle              & residual $\to$ correction ($\Complex^N$)   & \cref{ch:multigrid} \\
 \ \ smooth   & Chebyshev/Hiptmair    & per-level error $\to$ damped error         & \cref{ch:smoothers} \\
 \ \ apply    & matrix-free           & $\vw \in \Complex^N$ $\to$ $\mat{A}\vw$    & \cref{ch:matrixfree} \\
modes    & boundary-mode             & port faces $\to$ $\{\vs_1,\vs_2\}$         & \cref{ch:waveports} \\
reduce   & \code{sparameter\_reduce}  & $\vE_k, \{\vs_i\}$ $\to$ $S(\omega_k)$ ($p{\times}p$, complex) & \cref{ch:reductions} \\
\bottomrule
\end{tabular}
\end{center}

\medskip
Read the rightmost column top to bottom: it is a tour of the book. Read the
\emph{shape} column and watch the data collapse---from a file, to an $N$-dimensional
complex field, to a $2\times2$ complex matrix at this frequency. The big dimension
$N$ (the field) lives only \emph{inside} the solve; it is born at \code{nd\_space}
and dies at \code{sparameter\_reduce}, which projects it down to the handful of
port numbers the engineer wanted. One S-matrix entry---say $S_{21}(\omega_k)$, the
transmission---is a single complex number distilled from one $N$-dimensional GMRES
solve by one inner product against a port mode. Sweep this whole table across the
band and you have the plot of \cref{fig:trace-splot}; that sweep is the entire
driven driver.
\end{workedexample}

\begin{workedexample}{The electrostatic capacitor: the same skeleton, different cells}{trace-es}
Now run the identical lifecycle for the electrostatic capacitance extraction (the
clean contrast), and lay its trace table beside the driven one. The problem is $n$
conductors; write $N$ for the number of $\Hone$ degrees of freedom and $n$ for the
number of conductors.

\medskip
\begin{center}\small
\setlength{\tabcolsep}{4pt}
\renewcommand{\arraystretch}{1.25}
\begin{tabular}{@{}ll>{\raggedright\arraybackslash}p{30mm}>{\raggedright\arraybackslash}p{32mm}@{}}
\toprule
Stage & Component & Data in $\to$ out & Same as driven? \\
\midrule
parse    & \code{IoData} record      & file $\to$ \code{config}                  & \textbf{same} (\cref{ch:records}) \\
mesh     & \code{build\_mesh}         & \code{config} $\to$ \code{Mesh}            & \textbf{same} (\cref{ch:functionspaces}) \\
space    & \code{h1\_space}           & \code{Mesh} $\to$ $\Hone$ (dim $N$)        & \emph{swapped}: $\Hone$ not $\Hcurl$ \\
assemble & \code{fe\_assemble}         & space $\to$ $\Kop$ ($N{\times}N$, SPD)     & \textbf{same fold}, one term \\
excite   & \code{excitation}          & \code{config}, $j$ $\to$ $\vb_j \in \Real^N$ & per terminal (real) \\
solve    & \code{solve\_family} (CG)  & $\Kop, [\vb_j]$ $\to$ $[V_j] \in \Real^N$  & \emph{swapped}: corner~1, CG \\
 \ \ precond. & V-cycle              & residual $\to$ correction                   & \textbf{same} (\cref{ch:multigrid}) \\
 \ \ apply    & matrix-free           & $\vw$ $\to$ $\Kop\vw$                       & \textbf{same} (\cref{ch:matrixfree}) \\
reduce   & \code{gram\_reduce}        & $[V_j], \Kop$ $\to$ $\mat{C}$ ($n{\times}n$, sym.) & \emph{swapped}: Gram not projection \\
\bottomrule
\end{tabular}
\end{center}

\medskip
Set the two tables side by side and the book's whole thesis is visible in one
glance. The first two rows are \emph{identical}---same record, same mesh build. The
inner solver rows (V-cycle, matrix-free apply) are \emph{identical}---the same
numerical core. Exactly three rows differ: the space row ($\Hone$ for $\Hcurl$),
the solve row (\code{solve\_family}/CG/corner~1 for \code{frequency\_sweep}/GMRES/
corner~2), and the reduce row (\code{gram\_reduce} for \code{sparameter\_reduce}).
And one stage \emph{drops out} entirely: electrostatics needs no operator-forming
\code{linear\_combination} step, because its operator is fixed---which is precisely
the hoist of \cref{ch:electrostatics}. The electrostatic operator is also real and
SPD, so CG replaces GMRES and the V-cycle's smoother is a simpler real one; but the
\emph{stack is the same stack}. Two analyses, one machine, three swapped cells: the
trace tables prove it row by row.
\end{workedexample}

\begin{figure}[t]
\centering
\begin{tikzpicture}
\begin{axis}[
  width=0.84\linewidth, height=5.6cm,
  xlabel={frequency $f$ (GHz)}, ylabel={magnitude (dB)},
  xmin=3, xmax=7, ymin=-45, ymax=4,
  ytick={0,-10,-20,-30,-40}, xtick={3,4,5,6,7},
  legend pos=south east, legend cell align=left,
  grid=both, grid style={simGray!20},
  tick label style={font=\footnotesize}, label style={font=\footnotesize},
  legend style={font=\footnotesize},
]
% |S21| insertion loss: bandpass near 5 GHz
\addplot[very thick, simGreen, smooth, samples=120, domain=3:7]
  { -2 - 30*( (x-5)/1.0 )^2 / (1 + 0.55*((x-5))^2 ) - 0.5*(x-5)^2 };
\addlegendentry{$|S_{21}|$ (insertion loss)}
% |S11| return loss: deep dip at the match
\addplot[very thick, simRust, smooth, samples=160, domain=3:7]
  { -3 - 32*exp( -((x-5)^2)/0.06 ) + 2.0*((x-5))^2 };
\addlegendentry{$|S_{11}|$ (return loss)}
\draw[simBlue!50, densely dashed] (axis cs:4.4,4) -- (axis cs:4.4,-45);
\draw[simBlue!50, densely dashed] (axis cs:5.6,4) -- (axis cs:5.6,-45);
\node[font=\scriptsize, text=simBlue] at (axis cs:5.0,-43) {passband};
% mark one swept frequency = one full solve
\node[circle, fill=simInk, inner sep=1.1pt, label={[font=\scriptsize,text=simInk]above:{$\omega_k$}}]
  at (axis cs:4.7,-1.3) {};
\end{axis}
\end{tikzpicture}
\caption[The product: the swept S-parameter plot]{The output of the whole trace:
the scattering parameters of the two-port filter across the band. The insertion
loss $\abs{S_{21}}$ (green) sits near $0\,$dB in the passband around $5\,$GHz---signal
passes through---and rolls off outside it; the return loss $\abs{S_{11}}$ (rust)
dips sharply at the match. \emph{Every point on every curve is one complete pass of
the trace} of \cref{wex:trace-driven}: one $\mat{A}(\omega)$ formed, one GMRES solve
over the full multigrid stack, one port projection. The single marked point
$\omega_k$ is the frequency whose trace table we walked. Reading these two curves is
reading the device---and producing them is running the entire machine of
\cref{fig:trace-tree} once per frequency.}
\label{fig:trace-splot}
\end{figure}

\section{The bridge into Part~VI}
\label{sec:trace-bridge}

We have now done what Part~V set out to do. We met the six analyses, anchored them
all on electrostatics, and---in this chapter---ran one of them, the driven filter,
all the way through, naming every machine the book built and watching the whole
simulator collapse into a single composition (\cref{fig:trace-tree}). We saw, in two
trace tables, that a second analysis is the \emph{same} composition with three cells
swapped. The disassembly is complete: we know every part, and we know how they
compose into any driver.

The natural next question is the one Part~VI answers. If every analysis is a
composition of the same named pieces, then those pieces---taken together, with their
types and their bodies written out---are a small \emph{library}, and each driver is a
few lines of code over it. That is exactly the \emph{synthesized library} view. The
composition tree of \cref{fig:trace-tree} is not a metaphor; it is, almost literally,
the call graph of a real program, and Part~VI writes that program down. \Cref{ch:synthlib}
lays out the library's architecture---the handful of modules (types, iteration,
data-algebra, coordination, drivers) that hold the machines we have been naming---and
\cref{ch:composing} composes a driver end to end in code, the rendered form of the
very tree we just expanded.

\begin{keyidea}
\textbf{You have now seen the whole machine.} One simulation is one composition of
the same handful of machines---a mesh build, an assembly fold, a solve corner with
its nested Krylov/multigrid/matrix-free stack, and a reduction---and a different
analysis swaps only a few of those boxes. Part~V disassembled the simulator into
those boxes and showed they compose; Part~VI reassembles them into a single small
library and renders the whole simulator as a few pages of code. The trace tree of
this chapter \emph{is} that program, drawn rather than typed.
\end{keyidea}

\summarysection
This capstone traces one complete simulation---the driven frequency sweep of a
two-port filter---from the configuration file to the S-parameter plot, naming at
every stage the exact machine, and chapter, that does the work. Walking the
lifecycle of \cref{ch:lifecycle}: \emph{parse} produces the read-only \code{IoData}
record (\cref{ch:records}, \cref{ch:strata}); \emph{build mesh} produces the
\code{Mesh} (\cref{ch:functionspaces}); \emph{dispatch} selects the driven driver at
the specialization seam (\cref{ch:lifecycle}); \emph{assemble} folds three weak-form
term lists (\cref{ch:weak}) into the fixed basis $\{\Kop,\Cop,\Mop\}$ by
\code{fe\_assemble} (\cref{ch:assembly}), with matrix-free applies
(\cref{ch:matrixfree}) and eliminated boundary conditions (\cref{ch:bc});
\emph{solve} is the operator-varying map \code{frequency\_sweep} (\cref{ch:driven}),
each frequency a preconditioned GMRES (\cref{ch:gmres}) over a multigrid V-cycle
(\cref{ch:multigrid}) whose smoother is Chebyshev/Hiptmair (\cref{ch:smoothers}) and
whose every apply is matrix-free, the state threaded by the solve monad
(\cref{ch:monad}); \emph{reduce} is the port-projection \code{sparameter\_reduce}
(\cref{ch:reductions}) over port modes borrowed from the boundary-mode analysis
(\cref{ch:waveports}). The whole driver is one composition, \texttt{driven $=$
sparameter\_reduce $\circ$ frequency\_sweep $\circ$ fe\_assemble}, which expands into
a single finite tree whose every leaf is an earlier chapter (\cref{fig:trace-tree}).
Set beside it, the electrostatic capacitance trace is the \emph{same} machine---same
record, mesh, assembly fold, inner solver stack---with exactly three cells swapped:
the function space ($\Hone$/$\Hcurl$), the solve corner (fixed-operator map/CG versus
operator-varying map/GMRES), and the reduction (symmetric Gram versus
port-projection). One run touches one solve corner and one reduction shape on the
grid of \cref{ch:situating}: the simulator is a few corners $\times$ a few reductions
over one shared field theory and calculus. Part~VI (\cref{ch:synthlib},
\cref{ch:composing}) renders that composition as code.

\furtherreading
The scattering-parameter physics this trace computes---ports, impedance
normalization, return and insertion loss, the network-analyzer measurement---is the
foundation of microwave engineering, and Pozar's text~\citep{pozar2011} is the
standard treatment; the finite-element formulation of the time-harmonic curl--curl
system on $\Hcurl$, the wave-port modes, and the modal projection are developed at
length by Jin~\citep{jin2014}. The numerical engine the trace nests four
deep---GMRES on the complex indefinite operator, the multigrid preconditioner with
its smoother, and the reuse economics across the family---is the subject of Saad's
account of iterative methods~\citep{saad2003}, on which \cref{ch:gmres},
\cref{ch:multigrid}, and \cref{ch:smoothers} all draw. The individual machines named
in this trace are each developed in their own chapters, cited inline; this chapter
adds no new content, only the composition---which Part~VI (\cref{ch:synthlib},
\cref{ch:composing}) then renders as a working library.

\exercisessection
\begin{enumerate}[nosep]
  \item \textbf{Name the machine at each stage.} Without looking back at
  \cref{fig:trace-pipeline}, list the lifecycle stages of the driven trace in order
  (parse, mesh, dispatch, assemble, solve, reduce, output) and name the framework
  component and the chapter responsible for each. Which single stage consumes a
  product computed by a \emph{different} analysis, and which analysis?
  \item \textbf{Expand one branch.} Take the \code{frequency\_sweep} box of the
  composition tree (\cref{fig:trace-tree}) and write out, as a nested list, every
  machine that engages inside a single \code{ksp\_solve}---from the outer Krylov
  method down to the innermost operation. At what operation does \emph{every} branch
  of the tree bottom out, and why is that the same operation for GMRES, the V-cycle,
  and the smoother alike?
  \item \textbf{The one line that swaps the analysis.} The driven and electrostatic
  traces share most of their stages. Identify the single operation whose
  \emph{position} (inside the solve loop versus hoisted outside it) flips the run
  from corner~1 to corner~2, and explain why electrostatics can hoist it while the
  driven sweep cannot. (Cross-check against \cref{ch:driven}.)
  \item \textbf{Trace the shape.} Following the trace table of
  \cref{wex:trace-driven}, state the data shape carried out of each stage for the
  driven filter, from the config file to one S-matrix entry. At which stage is the
  large dimension $N$ (the field) born, and at which does it die? What is the shape
  of the final product at one frequency, and how many ports does it have on a side?
  \item \textbf{Two analyses, three swaps.} Using the two trace tables
  (\cref{wex:trace-driven}, \cref{wex:trace-es}), list exactly which rows are
  identical between the driven and electrostatic traces and which three are swapped.
  Which stage present in the driven trace \emph{drops out} entirely for
  electrostatics, and why? (Hint: what does a fixed operator not need to do each
  iteration?)
  \item \textbf{Count the footprint.} Of the four solve corners and three reduction
  shapes of \cref{ch:situating}, how many of each did this one driven run use? Name
  them. The run also \emph{borrowed} a product from a corner it did not itself
  occupy---which corner, and what did it borrow? Repeat the count for the
  electrostatic trace.
  \item \textbf{A third analysis on the same machine.} Pick the \emph{transient}
  analysis (\cref{ch:transient}). Predict, by analogy with the two traces of this
  chapter, which cells of its trace table would be \emph{shared} with the driven
  filter (record, mesh, assembly fold, inner stack) and which would be
  \emph{swapped}. Which solve corner does it occupy, and why is its solve a
  \emph{fold} rather than a \emph{map}?
  \item \textbf{From tree to code.} The composition tree of \cref{fig:trace-tree} is
  described in the bridge (\S\ref{sec:trace-bridge}) as ``almost literally the call
  graph of a real program.'' In two or three sentences, explain what would have to
  be written down to turn the tree into running code: which leaves are already-built
  library functions, which boxes are the composition itself, and which single piece
  (named in \cref{ch:waveports}) must be supplied as an input. (Part~VI,
  \cref{ch:synthlib} and \cref{ch:composing}, does exactly this.)
\end{enumerate}


\part{The Collection}\label{part:collection}
\chapter{The Synthesized Library: Architecture}
\label{ch:synthlib}

\chapterorient{Five parts have disassembled Palace into named pieces and shown
that every analysis is a composition of the same handful of them. This part puts
the pieces back---not as the C++ they came from, but as a small, clean
\emph{library} written in the calculus of Part~III: the simulator rendered as
code. This opening chapter lays out that library's \emph{architecture}. We meet
the three \emph{views} of the machine the book has been building (top-down,
bottom-up, and---new here---the implementation view), then the five modules the
library splits into, the dependency graph that orders them, the load-bearing rule
that keeps the code a \emph{view} of the spec rather than a second copy of it, and
the \texttt{\#extern} boundary where three opaque kernels stop. By the end you
will know the shape of the whole library; the next two chapters tour its modules
(\cref{ch:fivelibs}) and trace one driver through it as code
(\cref{ch:composing}).}

\section{Three views of one machine}
\label{sec:threeviews}

A complete simulator is a large object, and a large object is best understood from
more than one side. This book has, by now, walked around the machine and looked at
it from three. It is worth naming all three at once, because this part is the
third of them, and it only makes sense in relation to the other two.

The first view is \emph{top-down}: \emph{what does the machine do?} You stand at
the user-facing entry points---the six analyses an engineer actually invokes---and
read each as a coherent feature: the electrostatic capacitance extraction, the
driven frequency sweep, the eigenmode solve. Part~V took this view, devoting one
chapter to each modality and closing, in \cref{ch:fulltrace}, with a single driver
traced from its configuration file to its output plot. A driver, seen top-down, is
a \emph{composition root}: inputs are configuration, the output is the physical
product, and the body composes already-built vocabulary. The top-down view answers
``what is this analysis, end to end?''

The second view is \emph{bottom-up}: \emph{what is each piece?} You start at the
ground---cited Palace source---and climb, re-expressing each operation in a
cleaner vocabulary one rotation at a time, until you reach the small calculus of
combinators at the top (\cref{ch:operators}, \cref{ch:rotations}). Parts~II
through~IV took this view: the function spaces, the assembly fold, the Krylov
solvers, the multigrid preconditioners, the reduction verbs, each defined in
itself, with its signature, its laws, its place in the vocabulary. The bottom-up
view answers ``what is \code{fe\_assemble}, exactly, and where is it used?''

The third view is the one this part adds. It is the \emph{implementation view}:
\emph{what is the code?} Not what the machine does (top-down), not what each piece
is (bottom-up), but the actual library a programmer would read if the spec were
realized as software---the def bodies, written out, in the pseudo-language of
\cref{ch:language}, organized into modules, with every cross-reference resolving to
a real function in a real file. The implementation view answers ``show me the code
that realizes all of this.''

\begin{figure}[t]
\centering
\begin{tikzpicture}[font=\footnotesize, >=Stealth]
  % the central machine
  \node[stage, minimum width=26mm, minimum height=14mm, fill=simBgGray, draw=simGray,
        align=center] (m) at (0,0) {\textbf{the simulator}\\[1pt]\scriptsize one machine};
  % --- top-down view (from above) ---
  \node[product, align=center, text width=30mm] (td) at (0,2.7)
    {\textbf{top-down view}\\[1pt]\scriptsize \emph{what does it do?}\\6 driver features\\(Part~V, \cref{ch:fulltrace})};
  \draw[flow, simRust] (td.south) -- node[right, font=\scriptsize]{compose} (m.north);
  % --- bottom-up view (from below) ---
  \node[opbox, align=center, text width=30mm] (bu) at (0,-2.7)
    {\textbf{bottom-up view}\\[1pt]\scriptsize \emph{what is each piece?}\\vocabulary L0$\to$L4\\(Parts~II--IV)};
  \draw[flow, simTeal] (bu.north) -- node[right, font=\scriptsize]{decompose} (m.south);
  % --- implementation view (from the side: THIS PART) ---
  \node[stage, fill=simBgViolet, draw=simViolet, very thick, align=center, text width=32mm]
    (iv) at (5.6,0) {\textbf{implementation view}\\[1pt]\scriptsize \emph{what is the code?}\\the synthesized library\\(\textbf{this Part})};
  \draw[flow, simViolet] (iv.west) -- node[above, font=\scriptsize]{render} (m.east);
  % the three angles annotation
  \node[font=\scriptsize\itshape, text=simGray, align=center] at (0,-4.15)
    {three views of the \emph{same} machine---the simulator does not change,\\only the side we read it from};
\end{tikzpicture}
\caption[Three views of one machine]{The same simulator, read from three sides.
The \emph{top-down} view (above, rust) enters at the user-facing driver features
and asks what each analysis \emph{does}, reading it as a composition of vocabulary;
Part~V took this view. The \emph{bottom-up} view (below, teal) enters at cited
source and asks what each \emph{piece is}, climbing the vocabulary from L0 to L4;
Parts~II--IV took this view. The \emph{implementation view} (right, violet) is this
Part: it renders the same vocabulary as concrete library \emph{code}---the def
bodies a programmer would read. The machine is one object; the three views are
three angles on it, not three machines.}
\label{fig:threeviews}
\end{figure}

\Cref{fig:threeviews} sets the three side by side. They are not three machines;
they are three angles on one machine. The top-down view recomposes the vocabulary
into features; the bottom-up view decomposes the features into vocabulary; the
implementation view renders that same vocabulary as the code that realizes it. The
last is the natural endpoint of the journey: \cref{ch:bigidea} chose the
pure-function picture precisely because it is ``the form the accelerator wants''---a
library of pure tensor-in, tensor-out functions, ready to lower onto a GPU. This
part writes that library down.

\begin{keyidea}
The simulator can be read from three sides. \emph{Top-down} (Part~V): what does each
analysis \emph{do}---a driver as a composition root over the vocabulary.
\emph{Bottom-up} (Parts~II--IV): what is each \emph{piece}---the vocabulary defined
in itself, L0 to L4. \emph{The implementation view} (this Part): what is the
\emph{code}---the same vocabulary rendered as a small library of pure functions, the
bridge to a real GPU or accelerator backend. The synthesized library \emph{is} the
simulator written out as code; it is the destination \cref{ch:bigidea} pointed at
when it chose the pure-function picture for being ``the form the accelerator wants.''
\end{keyidea}

\subsection{Why a synthesized \emph{library}, and not a port}

A word on what we are and are not building. We are not porting Palace---not
translating its C++ line by line into another language. We are \emph{synthesizing}
a library: writing down, in one clean place, the code that the spec describes,
factored the way the spec factored it (combinator-primary, \cref{ch:combinators}),
with the optimization tricks of the original erased (\cref{ch:bigidea}) and the
pure-function reading made literal. The result is shorter than the original by a
large factor, because the original's bulk is mostly performance machinery and
boilerplate that the spec has already stripped away. What remains is the algorithm.

This is why the library is small enough to print. The composition tree of
\cref{ch:fulltrace}---the whole driven driver expanded to its matrix-free
leaves---was finite and shallow; rendered as code, each box in that tree is a few
lines, and the whole library is a few hundred lines across five modules. That
smallness is not a trick of presentation. It is the payoff of five parts of
disassembly: once you have found the handful of combinators that the whole
simulator is built from, writing them down as code is a short exercise.

\section{The five-library partition}
\label{sec:fivelib}

The library splits into five modules. The split is not arbitrary: three of the
five \emph{mirror} the three groups the calculus of \cref{ch:combinators} already
fell into, and the other two \emph{bracket} those three---one below, one above---to
make the whole stack come out in a clean order. \Cref{fig:bracket} is the hero
figure of the chapter: the five-library bracket, drawn as a stack.

\begin{figure}[p]
\centering
\resizebox{\textwidth}{!}{%
\begin{tikzpicture}[font=\footnotesize, >=Stealth]
  % ===== drivers cap (top) =====
  \node[product, minimum width=92mm, minimum height=11mm, align=center] (drv) at (0,5.2)
    {\textbf{\code{drivers}} \;\; \emph{the entry-point surfaces}\\[1pt]
     \scriptsize 6 sim drivers \;$\cdot$\; 6 output products \;$\cdot$\; the \code{lifecycle} ROOT \;\;(composes everything below)};
  % ===== three middle columns =====
  \node[stage, fill=simBgGreen, draw=simGreen, minimum width=28mm, minimum height=22mm,
        align=center] (iter) at (-3.1,2.6)
    {\textbf{\code{iteration}}\\[2pt]\scriptsize \code{iterate\_while}\\\code{krylov\_step}\\\code{chebyshev}\\[2pt]\itshape iteration \&\\step combinators};
  \node[stage, fill=simBgGold, draw=simGold, minimum width=28mm, minimum height=22mm,
        align=center] (data) at (0,2.6)
    {\textbf{\code{data-algebra}}\\[2pt]\scriptsize \code{linear\_combination}\\\code{inner\_product}\\\code{fe\_assemble} \dots\\[2pt]\itshape data-algebra\\combinators \& verbs};
  \node[stage, fill=simBgRust, draw=simRust, minimum width=28mm, minimum height=22mm,
        align=center] (coord) at (3.1,2.6)
    {\textbf{\code{coordination}}\\[2pt]\scriptsize \code{ksp\_solve}\\\code{solve\_family}\\\code{fold\_solve} \dots\\[2pt]\itshape outer-driver caps\\\& coordination};
  % ===== types slab (bottom) =====
  \node[stage, fill=simBgBlue, draw=simBlue, minimum width=92mm, minimum height=11mm,
        align=center] (types) at (0,-0.3)
    {\textbf{\code{types}} \;\; \emph{the shared cross-cutting records}\\[1pt]
     \scriptsize \code{IoData} \;$\cdot$\; \code{OpParams} \;$\cdot$\; \code{SimState} \;\;(every consumer rests on these)};
  % ===== brackets / arrows =====
  % drivers composes the three middle
  \draw[depedge] (drv.south west) ++(8mm,0) -- (iter.north);
  \draw[depedge] (drv.south)              -- (data.north);
  \draw[depedge] (drv.south east) ++(-8mm,0) -- (coord.north);
  % three middle rest on types
  \draw[depedge] (iter.south)  -- (types.north west) ++(8mm,0);
  \draw[depedge] (data.south)  -- (types.north);
  \draw[depedge] (coord.south) -- (types.north east) ++(-8mm,0);
  % the "mirrors the calculus" tag
  \node[font=\scriptsize\itshape, text=simGray, align=center] at (0,0.95)
    {the three middle libraries \emph{mirror} the three calculus doc-groups of \cref{ch:combinators}};
  % bracket labels on the side
  \node[font=\scriptsize\bfseries, text=simRust, anchor=west, align=left] at (5.2,5.2)
    {top bracket:\\\scriptsize composes down};
  \node[font=\scriptsize\bfseries, text=simBlue, anchor=west, align=left] at (5.2,-0.3)
    {bottom bracket:\\\scriptsize everything rests on it};
  % topological-order arrow
  \draw[->, simGray, thick] (-5.4,-0.3) -- (-5.4,5.2)
    node[midway, rotate=90, anchor=south, font=\scriptsize\itshape, text=simGray]
    {topological order: a def appears after everything it uses};
\end{tikzpicture}}
\caption[The five-library bracket]{The architecture of the synthesized library: a
five-module stack. The three middle libraries (\code{iteration},
\code{data-algebra}, \code{coordination}) \emph{mirror} the three doc-groups the
calculus of \cref{ch:combinators} already fell into---iteration and step
combinators, data-combining verbs, and outer-driver coordination. They are
\emph{bracketed} below by \code{types}, the genuinely shared record definitions
(\code{IoData}, \code{OpParams}, \code{SimState}) that every consumer rests on, and
above by \code{drivers}, the entry-point surfaces (the 6 drivers, 6 output
products, and the \code{lifecycle} ROOT) that \emph{compose} everything beneath
them. The bottom slab must come first and the top cap last because the stack is in
\emph{topological order}: a definition appears only after everything it uses, so
shared types precede their consumers and composing drivers come last. Read the
stack bottom to top and you are reading the library in the order you could actually
compile it.}
\label{fig:bracket}
\end{figure}

\subsection{The three middle libraries mirror the calculus}

{\sloppy
Recall \cref{ch:combinators}. The combinator vocabulary fell naturally into groups:
the \emph{iteration-structural} combinators that thread a carry through a loop or a
family (\code{iterate\_while} and its relatives), and the \emph{data-algebra}
combinators that combine tensors into tensors or scalars (\code{linear\_combination}
and \code{inner\_product}). The synthesized library keeps that grouping, refining it
into three modules:\par}

\begingroup\sloppy
\begin{description}[style=nextline, leftmargin=0pt, font=\normalfont\itshape]
  \item[\code{iteration} --- iteration \& step combinators.] The value-threaded loop
  primitive \code{iterate\_while} (\cref{ch:fold}) and the step kernels it folds:
  the Krylov step \code{krylov\_step}, the Chebyshev smoother \code{chebyshev}.
  These are the machines that drive one iterative solve to convergence.
  \item[\code{data-algebra} --- data-combining verbs.] The tensor-sum combinator
  \code{linear\_combination} and the scalar-producing \code{inner\_product}
  (\cref{ch:combinators}), together with the named verbs that the steps and
  reductions call: \code{fe\_assemble}, the reduction verbs \code{gram\_reduce} /
  \code{sparameter\_reduce} / the rest, the operator builders. These are the
  machines that combine data.
  \item[\code{coordination} --- outer-driver caps \& coordination.] The
  \code{Solve}-monadic outer driver \code{ksp\_solve} (\cref{ch:monad}) and the
  family combinators that drive it over a collection: \code{solve\_family},
  \code{frequency\_sweep}, \code{fold\_solve}, and the eigensolve cap
  \code{eigsolve}. These are the machines that coordinate many solves.
\end{description}
\endgroup

These three are the heart of the library, and they are exactly the three questions
\cref{ch:bigidea} taught us to ask of any computation: \emph{what operations}
(\code{data-algebra}), \emph{who owns the state} (the \code{Solve} monad, in
\code{coordination}), \emph{how is evolution coordinated} (the combinators, split
between \code{iteration} and \code{coordination}). The library's middle is the
calculus, modularized.

\subsection{The two brackets: \texttt{types} below, \texttt{drivers} above}

The three middle libraries do not stand alone. They are bracketed.

Beneath them sits \code{types}, the foundation. It holds only the records that are
\emph{genuinely shared}---used by two or more of the API groups above it. There are
three: \code{IoData}, the parsed configuration every solve reads (\cref{ch:records},
\cref{ch:strata}); \code{OpParams}, the read-only operator-internal parameters the
Krylov kernel and the coordination caps both close over; and \code{SimState}, the
solve-state record the iteration kernel evolves and the coordination caps thread.
A record used by only \emph{one} group does not live here---it is placed with that
group (the rule has a name, below). \code{types} is the cross-cutting remainder,
and it sits at the bottom because every consumer rests on it.

Above the three sits \code{drivers}, the cap. It holds the entry-point
surfaces---the same composition roots Part~V studied top-down: the six simulation
drivers, the six output products, and the \code{lifecycle} ROOT that dispatches
among them. Each is rendered as code that \emph{composes} the three middle
libraries. The electrostatic driver, written out, is a few lines calling
\code{fe\_assemble}, \code{solve\_family}, and \code{gram\_reduce}---exactly the
composition of \cref{ch:fulltrace}, now as a function. \code{drivers} sits at the
top because it uses everything below it and nothing uses it.

\begin{notationbox}
The three middle libraries take their names---\code{iteration},
\code{data-algebra}, \code{coordination}---directly from the three calculus
doc-groups of \cref{ch:combinators}; reading them is reading that chapter's two
combinator branches as modules. The two brackets, \code{types} and \code{drivers},
are the \emph{new} names this Part introduces, and they are bookkeeping names: the
shared-record floor and the entry-point cap. Five modules, three of them the
calculus you already know and two of them the brackets that order it.
\end{notationbox}

\subsection{Why a bracket, and why this order}

Why must \code{types} be first and \code{drivers} last? Because the library is
written in \emph{topological order}: a definition appears only after everything it
mentions has already been defined. This is the ordinary discipline of any
single-pass module---you cannot call a function before you have written it---and it
forces the bracket shape. A shared record like \code{SimState} is mentioned in the
signature of \code{krylov\_step} (which evolves it) and \code{ksp\_solve} (which
threads it); so \code{SimState} must be defined before both, which puts \code{types}
at the bottom. A driver like \code{electrostatic} mentions \code{fe\_assemble},
\code{solve\_family}, and \code{gram\_reduce}; so those must all be defined before
it, which puts \code{drivers} at the top. The brackets are not a stylistic choice;
they are where topological order \emph{forces} the shared floor and the composing
cap to go.

\begin{figure}[t]
\centering
\resizebox{\textwidth}{!}{%
\begin{tikzpicture}[font=\footnotesize, >=Stealth, node distance=4mm]
  % a linear "reading order" of defs, left to right, with dependency arcs pointing backward
  \def\dy{0}
  \node[stage, fill=simBgBlue, draw=simBlue] (t1) at (0,\dy) {\code{SimState}};
  \node[stage, fill=simBgBlue, draw=simBlue, right=of t1] (t2) {\code{OpParams}};
  \node[opbox, fill=simBgGreen, draw=simGreen, right=of t2] (iw) {\code{iterate\_while}};
  \node[opbox, fill=simBgGreen, draw=simGreen, right=of iw] (ks) {\code{krylov\_step}};
  \node[opbox, fill=simBgRust, draw=simRust, right=of ks] (kspv) {\code{ksp\_solve}};
  \node[opbox, fill=simBgRust, draw=simRust, right=of kspv] (sf) {\code{solve\_family}};
  \node[product, right=of sf] (el) {\code{electrostatic}};
  % the reading-order baseline
  \draw[->, simGray, very thick] (-1.0,-0.85) -- (12.6,-0.85)
    node[right, font=\scriptsize\itshape, text=simGray] {reading / compile order};
  % drop ticks
  \foreach \n in {t1,t2,iw,ks,kspv,sf,el}{
    \draw[simGray!50] (\n.south) -- (\n.south |- 0,-0.85);}
  % dependency arcs (each points BACKWARD to something already defined): u -> v means u uses v
  \draw[depedge] (ks.north) to[out=120,in=60] (iw.north);          % krylov_step uses iterate_while
  \draw[depedge] (ks.north) to[out=140,in=40] (t1.north);          % ...and SimState
  \draw[depedge] (kspv.north) to[out=120,in=60] (ks.north);        % ksp_solve folds krylov_step
  \draw[depedge] (kspv.north) to[out=150,in=30] (t2.north);        % ...and OpParams
  \draw[depedge] (sf.north) to[out=120,in=60] (kspv.north);        % solve_family maps ksp_solve
  \draw[depedge] (el.north) to[out=120,in=60] (sf.north);          % electrostatic uses solve_family
  % annotation: every arc points LEFT (backward)
  \node[font=\scriptsize\itshape, text=simBlue, align=center] at (6.0,2.0)
    {every dependency arc points \emph{backward}: a def appears only after everything it uses};
\end{tikzpicture}}
\caption[Topological definition order]{Topological order, read along one line. The
library's definitions are laid out left to right in the order a single-pass reader
(or compiler) meets them; the lower arc is that reading order. Each curved arc above
is a \code{depends-on} edge---\code{krylov\_step} uses \code{iterate\_while} and
\code{SimState}; \code{ksp\_solve} folds \code{krylov\_step} and reads
\code{OpParams}; \code{solve\_family} maps \code{ksp\_solve}; the
\code{electrostatic} driver composes \code{solve\_family}. The defining property is
that \emph{every arc points backward}: a definition appears only after everything it
uses has already appeared. This is exactly why the shared \code{types}
(\code{SimState}, \code{OpParams}) come first and the composing driver comes last---a
forward-pointing arc would be a use-before-definition, which topological order
forbids. The five-library bracket of \cref{fig:bracket} is this one line, folded into
a stack.}
\label{fig:topo}
\end{figure}

\Cref{fig:topo} draws the order as a single line of defs with their dependency arcs,
and makes the defining property visible at a glance: every arc points
\emph{backward}. A forward-pointing arc would name something not yet defined---a
use-before-definition---and topological order is precisely the absence of such arcs.
The bracket of \cref{fig:bracket} is this line folded into a stack: the floor is its
left end, the cap its right.

\begin{keyidea}
The library is five modules: a \code{types} floor, three middle libraries
(\code{iteration}, \code{data-algebra}, \code{coordination}) that \emph{mirror} the
three calculus doc-groups of \cref{ch:combinators}, and a \code{drivers} cap that
\emph{composes} them into the entry-point surfaces of Part~V. The shape is a
bracket because the library is in \emph{topological order}---a def appears after
everything it uses---so the genuinely shared records must come first (the floor)
and the composing drivers must come last (the cap). Read bottom to top, the stack
is the order you could compile it in.
\end{keyidea}

\subsection{The type-placement rule}
\label{sec:typeplacement}

One subtlety of the floor deserves its own statement, because it is the rule that
keeps \code{types} small. \code{types} holds only records shared across \emph{two or
more} API groups. A record used by exactly \emph{one} group is not cross-cutting,
and putting it in the floor would force the reader to hold an unfamiliar type in
mind long before reaching the code that uses it. So a single-group record is placed
\emph{immediately before its own group}, bundled with its \emph{utility} API---its
constructors, field accessors, and trivial projections, the small namespace
intrinsic to the type. Its \emph{consumer} methods---the substantive operators that
read the record to do real work---stay in the group proper, after the type.

The Krylov workspace \code{Krylov} is used only by the iteration kernels, so it
lives at the head of \code{iteration}, not in the floor. The eigensolve state
\code{EigState} is used only by \code{eigsolve}, so it lives at the head of
\code{coordination}. The per-driver configuration records are thin projections of
\code{IoData} used only by their own driver, so they live at the head of
\code{drivers}. Only the truly cross-cutting three---\code{IoData}, \code{OpParams},
\code{SimState}---earn a place in the shared floor. Topological order still governs
throughout: a type and its utility API always precede their consumers, whether the
type sits in the floor or at the head of its group.

\section{The dependency graph and its two edge classes}
\label{sec:depgraph}

The bracket of \cref{fig:bracket} is a picture of \emph{dependencies}: who uses
whom. To make the picture precise, we draw the dependency graph explicitly and
distinguish its two kinds of edge, because the distinction is what makes the
ordering well-defined.

A \emph{dependency} between two defs comes in two flavors. A \textbf{\code{depends-on}}
edge is a real build dependency: def $u$ \code{depends-on} def $v$ when $u$'s body
\emph{calls} $v$, so $v$ must be defined before $u$ can be. These edges are
\emph{blocking}---they constrain the order, and they are exactly the edges
topological order must respect. \code{ksp\_solve} \code{depends-on}
\code{krylov\_step} (it folds it); \code{electrostatic} \code{depends-on}
\code{fe\_assemble}, \code{solve\_family}, and \code{gram\_reduce} (it composes
them); every consumer \code{depends-on} the \code{types} it reads.

A \textbf{\code{reference}} edge is navigational only: $u$ \code{reference}s $v$
when $u$ \emph{points the reader at} $v$ without calling it---a cross-link, a ``see
also,'' a ``this is the same idea as.'' These edges are \emph{free}: they do not
constrain the order and they impose no build dependency. When the \code{driven}
driver's documentation notes that its port modes ``come from the boundary-mode
analysis,'' that is a \code{reference}, not a \code{depends-on}: \code{driven} does
not \emph{call} \code{boundary\_mode}; it consumes a product that happens to be
produced elsewhere, and the cross-link is a courtesy to the reader, not a wire in
the build graph.

\begin{figure}[t]
\centering
\resizebox{\textwidth}{!}{%
\begin{tikzpicture}[font=\footnotesize, >=Stealth, node distance=10mm and 16mm]
  % the stack, as nodes
  \node[product, minimum width=24mm] (drv) at (0,3.6) {\code{drivers}};
  \node[stage, fill=simBgGreen, draw=simGreen, minimum width=20mm] (it) at (-3,1.8) {\code{iteration}};
  \node[stage, fill=simBgGold,  draw=simGold,  minimum width=22mm] (da) at (0,1.8) {\code{data-algebra}};
  \node[stage, fill=simBgRust,  draw=simRust,  minimum width=22mm] (co) at (3,1.8) {\code{coordination}};
  \node[stage, fill=simBgBlue,  draw=simBlue,  minimum width=24mm] (ty) at (0,0) {\code{types}};
  % depends-on edges (solid blue): drivers -> three middle -> types
  \draw[depedge] (drv) -- (it);
  \draw[depedge] (drv) -- (da);
  \draw[depedge] (drv) -- (co);
  \draw[depedge] (it) -- (ty);
  \draw[depedge] (da) -- (ty);
  \draw[depedge] (co) -- (ty);
  % an internal depends-on: coordination -> iteration (ksp_solve folds krylov_step)
  \draw[depedge] (co) -- node[above, font=\scriptsize, pos=0.45]{folds} (it);
  % coordination -> data-algebra (drives the verbs)
  % a reference edge (dotted gray): a driver cross-links a sibling driver's product
  \draw[refedge] (drv.east) .. controls +(2.4,-0.4) and +(1.6,1.4) ..
    node[right, font=\scriptsize, pos=0.5, text=simGray]{\emph{reference} (sibling product)} (co.east);
  % legend (sample edges drawn directly, with adjacent labels)
  \draw[depedge] (-6.6,3.6) -- (-6.0,3.6);
  \node[anchor=west, font=\scriptsize] at (-5.9,3.6) {\code{depends-on}};
  \node[anchor=west, font=\scriptsize, text=simGray] at (-6.6,3.15) {blocking; constrains order};
  \draw[refedge] (-6.6,2.4) -- (-6.0,2.4);
  \node[anchor=west, font=\scriptsize] at (-5.9,2.4) {\code{reference}};
  \node[anchor=west, font=\scriptsize, text=simGray] at (-6.6,1.95) {navigational; free};
\end{tikzpicture}}
\caption[The dependency graph and its two edge classes]{The library's dependency
graph, with its two edge classes distinguished. A \code{depends-on} edge (solid
blue) is a real build dependency---$u$ \code{depends-on} $v$ when $u$'s body calls
$v$---so it \emph{constrains the order}: the arrows all point downward, from
\code{drivers} through the three middle libraries to \code{types}, and an internal
arrow records that \code{coordination} folds \code{iteration}'s kernels. These are
the edges topological order respects. A \code{reference} edge (dotted gray) is
merely navigational---a cross-link to a sibling's product, here a driver pointing
at another driver's output it consumes---and is \emph{free}: it imposes no build
dependency and constrains nothing. The stack \code{types} $\leftarrow$
\{\code{iteration}, \code{data-algebra}, \code{coordination}\} $\leftarrow$
\code{drivers} is exactly the \code{depends-on} order.}
\label{fig:depgraph}
\end{figure}

\Cref{fig:depgraph} draws the graph. The \code{depends-on} edges (solid) all run
downward: \code{drivers} \code{depends-on} the three middle libraries, each of which
\code{depends-on} \code{types}, with the one internal cross-edge that
\code{coordination} folds \code{iteration}'s step kernels. Follow only the solid
edges and you recover the bracket stack of \cref{fig:bracket}; that is the whole
content of ``topological order''---the solid arrows never point up. The
\code{reference} edges (dotted) are the free cross-links the reader navigates by;
they are allowed to point any direction, because they bind nothing.

\begin{pitfall}
The distinction between \code{depends-on} and \code{reference} is not pedantry; it
is what keeps the dependency graph \emph{well-founded}. If we treated every
cross-link as a build dependency, the graph would tangle: the driven driver
``references'' the boundary-mode driver, the boundary-mode driver shares vocabulary
with the eigenmode driver, and soon every module appears to depend on every other
and no topological order exists. By reserving \code{depends-on} for actual calls and
demoting cross-links to free \code{reference} edges, the build graph stays a clean
downward stack while the reader still gets every helpful pointer. A cross-link that
is \emph{not} a call must be a \code{reference}, or the order breaks.
\end{pitfall}

\section{The implementation view versus the semantics}
\label{sec:viewvssem}

We come now to the load-bearing distinction of the whole part---the one rule that
makes the synthesized library a \emph{view} of the specification rather than a
second, rivalrous copy of it. State it plainly: \emph{the library renders the def
bodies; it does not restate the semantics.}

Every operator in this book has a definitional home---a chapter where its
signature, its algebraic laws, its meaning, and (for a record) its field schema are
stated \emph{once} and authoritatively. \code{linear\_combination}'s laws live in
\cref{ch:combinators}; \code{ksp\_solve}'s convergence semantics live in the Krylov
chapters; \code{SimState}'s field schema lives with the records of
\cref{ch:records}. The synthesized library does \emph{not} repeat any of that. It
renders the \emph{code}---the def body, the concrete expression that computes the
result---and for everything else it \emph{links} to the authoritative home. A
rendered \code{linear\_combination} in the library is a few lines of foldl; it does
not re-derive the concatenation-homomorphism law, it points at \cref{ch:combinators}
where that law is proved.

\begin{figure}[t]
\centering
\begin{tikzpicture}[font=\footnotesize, >=Stealth]
  % LEFT: the semantics home
  \node[stage, fill=simBgGold, draw=simGold, minimum width=34mm, minimum height=22mm,
        align=center] (sem) at (-3.4,0)
    {\textbf{the semantics home}\\[2pt]\scriptsize (the operator chapter,\\\cref{ch:combinators};\\the calculus, Part~III)\\[3pt]
     \footnotesize signature\\laws\\meaning\\record schema};
  \node[font=\scriptsize\itshape, text=simGold, anchor=north] at (-3.4,-1.7)
    {stated ONCE, authoritative};
  % RIGHT: the rendered code
  \node[stage, fill=simBgViolet, draw=simViolet, minimum width=34mm, minimum height=22mm,
        align=center] (code) at (3.4,0)
    {\textbf{the synthesized def}\\[2pt]\scriptsize (the library chapter,\\this Part)\\[4pt]
     \footnotesize the def \emph{body}:\\the concrete code\\that computes it};
  \node[font=\scriptsize\itshape, text=simViolet, anchor=north] at (3.4,-1.7)
    {the implementation VIEW};
  % the link (reference) from code to semantics
  \draw[refedge, very thick] (code.west) -- node[above, font=\scriptsize]{\emph{links to}}
    node[below, font=\scriptsize, text=simGray]{(does not restate)} (sem.east);
  % the "no back-copy" forbidden arrow
  \draw[-{Stealth[length=2mm]}, simRust, dashed, thick] (sem.north east) ++(0,0.2)
    .. controls +(1.5,1.0) and +(-1.5,1.0) .. (code.north west) ++(0,0.2);
  \node[font=\scriptsize\itshape, text=simRust] at (0,2.05) {$\times$ no second copy of the laws};
\end{tikzpicture}
\caption[Implementation view versus semantics]{The load-bearing distinction. The
\emph{semantics} of an operator---its signature, its algebraic laws, its meaning,
and (for a record) its field schema---live \emph{once}, authoritatively, in its
home chapter (left, gold): \code{linear\_combination}'s laws in
\cref{ch:combinators}, a record's schema with \cref{ch:records}. The synthesized
library (right, violet) renders only the def \emph{body}---the concrete code that
computes the result---and \emph{links} to the home for everything else (solid
arrow). It does \emph{not} copy the laws across (the forbidden dashed arrow): the
library is a \emph{view} of the spec, like generated code is a view of its source,
not a second authoritative statement of it. This is what keeps the code and the spec
from drifting into two rivalrous copies.}
\label{fig:viewvssem}
\end{figure}

\Cref{fig:viewvssem} draws the rule. Why insist on it? Because a specification and
its implementation, if both state the laws, will eventually \emph{disagree}---someone
edits one and forgets the other, and now there are two answers to ``what does this
operator mean.'' By making the library a pure \emph{view}---bodies plus links, no
restated laws---we guarantee there is exactly one authoritative statement of every
semantic fact, and the code is checked \emph{against} it rather than competing with
it. This is the same relationship generated code has to its source: the generated
form is real, runnable, and useful to read, but it is not where the meaning
\emph{lives}; the meaning lives in the source it was generated from. The synthesized
library is generated code in exactly this sense, and the calculus of Part~III is its
source.

\begin{keyidea}
The synthesized library is a \emph{view}, not a second copy. It renders the def
\emph{bodies}---the concrete code---and \emph{links} to each operator's home chapter
for the signature, laws, meaning, and record schema, which are stated \emph{once and
only once} there. It never restates the semantics. This is the discipline that keeps
the implementation and the specification from drifting into two rivalrous truths:
there is one authoritative statement of every fact, and the code is a checkable view
of it, exactly as generated code is a view of its source. The implementation view is
parallel to the top-down feature spine---another angle on the one machine, not a new
machine.
\end{keyidea}

\section{The \texttt{\#extern} boundary: where the library stops}
\label{sec:extern}

There is one more thing the library must be honest about: its \emph{edges}. Not
every operation the simulator performs is ours to write. Three of them bottom out in
an opaque external kernel---a specialized library Palace calls but does not
itself implement---and the synthesized code must mark those boundaries rather than
fabricate a body for them. The mark is \code{\#extern}.

When a def's implementation lives behind an opaque library boundary, the library
renders its \emph{type signature} and then, in place of a body, writes
\code{\#extern NAME}. This says, exactly: \emph{here is the contract---the input and
output types---and here is where our code stops and the external kernel takes over;
we do not fabricate what we do not own.} Three kernels in the whole simulator earn
an \code{\#extern}:

\begin{description}[style=nextline, leftmargin=0pt, font=\normalfont\itshape]
  \item[The libCEED element-quadrature kernel.] Inside \code{fe\_assemble}
  (\cref{ch:assembly}), the element-local quadrature contraction is performed by
  libCEED, the matrix-free kernel library (\cref{ch:matrixfree}). The assembly fold
  is ours; the per-element kernel it calls is \code{\#extern}.
  \item[The SLEPc eigensolve loop.] Inside \code{eigsolve} (\cref{ch:krylovschur}),
  the eigenvalue iteration itself runs inside SLEPc---Palace writes no loop, it
  hands the whole problem across and waits. The cap is ours; the iteration is
  \code{\#extern}.
  \item[The ODE time-step.] Inside \code{fold\_solve} (\cref{ch:time}), the transient
  analysis's per-step time advance bottoms out in an MFEM time-integrator (an
  implicit Runge--Kutta or generalized-$\alpha$ step). The fold that threads the
  state is ours; the step it folds is \code{\#extern}.
\end{description}

\begin{figure}[t]
\centering
\begin{tikzpicture}[font=\footnotesize, >=Stealth]
  % the def: type signature (ours) + extern body (opaque) + the impl link
  \node[stage, minimum width=40mm, minimum height=8mm, align=center] (sig) at (0,2.4)
    {\code{eigen\_iterate :: OpParams -> ...}\\[1pt]\scriptsize the type signature (ours: the contract)};
  \node[extern, minimum width=40mm, minimum height=9mm, align=center] (ext) at (0,1.0)
    {\code{\#extern eigen\_iterate}\\[1pt]\scriptsize the opaque kernel (SLEPc EPS loop);\\our code stops here};
  \draw[flow] (sig) -- (ext);
  % the kernel-API label
  \node[font=\scriptsize\bfseries, text=simViolet, anchor=west] at (3.2,1.0)
    {= the kernel-\emph{API}};
  \node[font=\scriptsize\itshape, text=simGray, anchor=west, align=left] at (3.2,0.55)
    {the boundary / contract\\(what the spine calls)};
  % the separate kernel-IMPL node, linked by realizes-kernel-api (dashed ref)
  \node[stage, fill=simBgGreen, draw=simGreen, minimum width=46mm, minimum height=12mm,
        align=center] (impl) at (0,-1.4)
    {\textbf{the kernel-\emph{implementation}}\\[1pt]\scriptsize a SEPARATE constructive realization\\in OUR vocabulary: Lanczos / Arnoldi /\\Krylov--Schur (\cref{ch:krylovschur})};
  \draw[refedge, very thick] (ext.south) --
    node[right, font=\scriptsize, pos=0.5]{\code{realizes-kernel-api}}
    node[left, font=\scriptsize, text=simGray, pos=0.5, align=right]{(reviewable;\\not a call)} (impl.north);
\end{tikzpicture}
\caption[The \texttt{\#extern} boundary and the kernel-API/impl split]{The
\code{\#extern} boundary, and the distinction it does \emph{not} erase. An opaque
kernel renders as its \emph{type signature} (top: the contract, ours to state)
followed by \code{\#extern NAME} (middle, violet dashed) in place of a body: our
code stops here and the external library---here SLEPc's eigenvalue loop---takes
over. This \code{\#extern} surface is the \emph{kernel-API}: the boundary and
contract the simulator calls across. It is \emph{not} the last word. Separately, the
book gives a \emph{kernel-implementation} (bottom, green): a constructive realization
of the same kernel in our own vocabulary---for the eigensolve, the Lanczos / Arnoldi
/ Krylov--Schur iteration of \cref{ch:krylovschur}. The two are linked by a
\code{realizes-kernel-api} edge (dashed): a \emph{reviewable correspondence}, not a
build call---a reader checks that the from-our-primitives version matches the opaque
contract. The \code{\#extern} marks where \emph{the library} stops; the kernel-impl
is where \emph{the book} keeps going.}
\label{fig:extern}
\end{figure}

\Cref{fig:extern} draws the boundary and the crucial thing it does \emph{not} do:
it does not abandon the question. An \code{\#extern} surface is the
\emph{kernel-API}---the contract the simulator calls across, the opaque library
boundary. But for each of these kernels the book \emph{separately} provides a
\emph{kernel-implementation}: a constructive realization of the same operation in
our own vocabulary, built from primitives we already own. The SLEPc eigensolve loop
is \code{\#extern} in the library, but \cref{ch:krylovschur} reconstructs what that
loop does---the Lanczos, Arnoldi, and Krylov--Schur iterations---in our calculus, so
the curious reader can read both the black-box contract and a glass-box realization
of it. The two are tied together by a \code{realizes-kernel-api} link: a
\emph{reviewable correspondence}, not a build call. The implementation does not
\code{depends-on} the API and the API does not call the implementation; rather, a
reader checks that the from-our-primitives version computes what the opaque contract
promises.

\begin{notationbox}
\code{\#extern NAME} is not a stub and not a failure---it is an honest edge. It
appears \emph{after} the type signature, never before: the library always states the
\emph{contract} (the input and output types) even where it does not own the
\emph{body}. Read \code{\#extern} as ``the contract is here; the implementation is
external; and---for these three---the book gives a constructive realization
separately, linked by \code{realizes-kernel-api}.'' Three kernels carry it: libCEED
quadrature, the SLEPc eigensolve loop, the ODE time-step.
\end{notationbox}

\begin{keyidea}
The library marks its edges with \code{\#extern}. Where an operation bottoms out in
an opaque external kernel---libCEED quadrature inside \code{fe\_assemble}, the SLEPc
eigensolve loop inside \code{eigsolve}, the ODE step inside \code{fold\_solve}---the
library renders the type signature and then writes \code{\#extern NAME} in place of
a body: it states the contract but does not fabricate what it does not own. The
\code{\#extern} surface is the \emph{kernel-API}. Distinct from it is the
\emph{kernel-implementation}: a separate, constructive realization of the same kernel
in our own vocabulary (e.g. the Krylov--Schur iteration of \cref{ch:krylovschur} for
the eigensolve), tied to the API by a \emph{reviewable} \code{realizes-kernel-api}
link, not a build call. The \code{\#extern} is where the \emph{library} stops; the
kernel-impl is where the \emph{book} keeps going.
\end{keyidea}

\section{Worked examples}

\begin{workedexample}{Reading the topological order, and placing \texttt{ksp\_solve}}{topo}
We make ``topological order'' concrete by deriving why \code{types} must come first
and \code{lifecycle} last, and then placing one specific operator,
\code{ksp\_solve}, in its module with its dependencies all pointing downward.

\emph{Why \code{types} first.} Pick any record in the floor, say \code{SimState}.
It is mentioned in the \emph{signature} of \code{krylov\_step} (whose monadic effect
\emph{is} the \code{SimState} transition) and in the signature of \code{ksp\_solve}
(which returns a \code{SimState}). Topological order forbids mentioning a name before
it is defined, so \code{SimState} must precede both \code{krylov\_step} (in
\code{iteration}) and \code{ksp\_solve} (in \code{coordination}). The same holds for
\code{OpParams} (read by the kernel and the caps) and \code{IoData} (read by every
driver). Three records, each used by two or more groups, each forced to precede its
users---so the cross-cutting records form a floor, and the floor is \code{types}.

\emph{Why \code{lifecycle} last.} The \code{lifecycle} ROOT \emph{dispatches} on the
problem type to one of the six drivers and calls it; so it mentions all six driver
names, which must already be defined. The drivers in turn mention the middle-library
verbs, which mention the floor. Nothing mentions \code{lifecycle}---it is the entry
point, the thing the user invokes, so no other def calls it. A def that uses
everything and is used by nothing is, by definition, topologically \emph{last}.

\emph{Placing \code{ksp\_solve}.} The preconditioned Krylov solve cap is an
outer-driver coordination machine: it runs the outer loop, classifies termination,
and projects the terminal state. So it belongs in \code{coordination}. Read its
dependencies:
\begin{lstlisting}[style=l4style]
ksp_solve :: OpParams -> Inputs -> SimState
ksp_solve op inp = execState (solve_loop op inp) (initial_state inp)
  where
    solve_loop op inp = do
      o <- one_cycle op inp
      unless (done o) (solve_loop op inp)
    one_cycle op inp = do
      s <- get
      let (kn, _) = iterate_while (krylov_step op (fresh_krylov op inp s)) cont
      modify (\s -> s { x = s.x `plus` (kn.basis `applyBasis` kn.y) })
      pure (classify kn op s)
\end{lstlisting}
Every name on the right-hand side is in a \emph{lower or equal} stratum:
\code{OpParams}, \code{SimState}, \code{Inputs} are in \code{types} (the floor);
\code{krylov\_step} and \code{iterate\_while} are in \code{iteration} (a middle
sibling, defined before \code{coordination} because \code{coordination} folds it);
\code{execState}, \code{get}, \code{modify} are the \code{Solve}-monad utilities
defined at the head of \code{coordination} itself, before \code{ksp\_solve}. No
dependency points \emph{up}: \code{ksp\_solve} never mentions a driver, never
mentions \code{lifecycle}. Its \code{depends-on} edges all run downward---\code{types}
and \code{iteration}---so it sits cleanly in \code{coordination}, above its
dependencies and below the \code{drivers} that will fold it into
\code{solve\_family}. That downward-only property is exactly what topological order
\emph{is}, checked at one node.
\end{workedexample}

\begin{workedexample}{The \texttt{\#extern} boundary inside \texttt{fe\_assemble}}{externassemble}
We render the assembly fold with its opaque kernel marked, and read off precisely
what the boundary means---contrasting it with the kernel-\emph{implementation} that
\cref{ch:matrixfree} provides.

The assembly verb \code{fe\_assemble} (\cref{ch:assembly}) folds a list of
weak-form terms over the mesh, accumulating each term's element-local contributions
into the operator. The \emph{fold} is ours---it is the ordinary
\code{linear\_combination}-shaped accumulation of \cref{ch:combinators}. But the
innermost step, evaluating one term's quadrature contraction on one element, is
performed by libCEED, the matrix-free kernel library. So the synthesized library
renders the fold and marks the kernel:
\begin{lstlisting}[style=l4style]
-- the assembly fold: ours. the per-element quadrature kernel: external.
fe_assemble :: FESpace -> [WeakTerm] -> LinOp
fe_assemble space terms =
  foldl (\op term -> op `plus_op` assemble_term space term) zero_op terms
  where
    -- one weak-form term -> its operator contribution, via the libCEED kernel.
    -- the CONTRACT (types) is ours; the BODY is the opaque matrix-free kernel.
    assemble_term :: FESpace -> WeakTerm -> LinOp
    #extern assemble_term
\end{lstlisting}
Read the two parts. The body of \code{fe\_assemble} is a plain \code{foldl}: it is
our code, the assembly fold of \cref{ch:assembly}, accumulating one operator
contribution per term. But \code{assemble\_term}---the per-element quadrature
contraction that turns one weak-form term into its matrix-free operator
action---renders as \code{\#extern} after its signature. That is the boundary: we
\emph{call into} libCEED's element-quadrature kernel, and the body is not ours to
fabricate. We state the contract (a \code{FESpace} and a \code{WeakTerm} in, a
\code{LinOp} out) and stop.

\emph{The contrast with the kernel-impl.} The \code{\#extern} is not the end of the
story for this kernel either. \Cref{ch:matrixfree} gives the
kernel-\emph{implementation}: a constructive realization of exactly this element
quadrature as a sequence of tensor contractions---gather the element DoFs, apply the
basis-evaluation tensor, contract against the quadrature weights and the geometry
factors, apply the basis-transpose, scatter back. That realization is built from
tensor primitives we already own, and it \emph{realizes-kernel-api} the
\code{\#extern assemble\_term} contract above: a reviewable correspondence, checked
by reading both, not a build call. So the same quadrature kernel appears twice in the
book, deliberately: once as the opaque contract the assembly fold calls
(\code{\#extern}, here), and once as the from-our-primitives tensor-contraction
implementation (\cref{ch:matrixfree}). The library renders the call site; the book
supplies the glass box beside it. That double rendering---API where the library
stops, impl where the book continues---is exactly the kernel-API/kernel-impl
distinction of \cref{sec:extern}, made concrete on the assembly kernel.
\end{workedexample}

\section{What the next two chapters do}
\label{sec:whatnext}

We have the architecture. \Cref{fig:bracket} is the floor plan: five modules, a
\code{types} floor, three middle libraries mirroring the calculus, a \code{drivers}
cap, all in topological order, with the \code{depends-on} edges running downward and
the \code{reference} cross-links free. We know the library is a \emph{view}---bodies
and links, never a second copy of the laws---and we know where it stops, at the three
\code{\#extern} kernels whose contracts it states but whose bodies it does not own.

The next two chapters fill the plan in. \Cref{ch:fivelibs} tours the five libraries
one at a time, opening each module and reading its rendered defs---the actual code,
module by module, from the \code{types} floor up through the \code{drivers} cap.
\Cref{ch:composing} then takes one driver---the electrostatic capacitance extraction,
the cleanest of the six---and composes it end to end \emph{in code}, the rendered
form of the very composition tree \cref{ch:fulltrace} expanded as a diagram. Together
they turn this architecture into the running library it describes, and
\cref{ch:wholemachine} closes the book by standing back from the whole machine and
asking where it goes next---onto the GPU and accelerator backends the pure-function
picture was chosen for in the first place.

\begin{keyidea}
\textbf{The whole simulator is one small library of composed functions---the
implementation view, mirroring the calculus.} Five modules: a \code{types} floor of
shared records, three middle libraries (\code{iteration}, \code{data-algebra},
\code{coordination}) that \emph{are} the three calculus doc-groups of
\cref{ch:combinators} rendered as code, and a \code{drivers} cap that composes them
into the entry-point surfaces of Part~V---all in topological order. It renders the
def \emph{bodies} and \emph{links} to the spec for the laws (a view, not a copy), and
marks its three opaque edges with \code{\#extern}. A hundred-thousand-line C++
simulator becomes a few hundred lines across five small modules, because once you
have found the handful of combinators it is built from, writing them as code is
short. That smallness is the whole book's thesis, finally rendered as the program it
always was.
\end{keyidea}

\summarysection
This chapter opens Part~VI, which renders the whole specification as a small
\emph{synthesized library}---the simulator written out as code in the calculus of
Part~III. We first named the three \emph{views} of the one machine: the
\emph{top-down} view of Part~V (what each analysis \emph{does}, a driver as a
composition root), the \emph{bottom-up} view of Parts~II--IV (what each \emph{piece
is}, the vocabulary L0$\to$L4), and---new here---the \emph{implementation view} (what
the \emph{code} is), the destination \cref{ch:bigidea} pointed at when it chose the
pure-function picture for matching the accelerator backend. The library splits into
\emph{five modules}: a \code{types} floor (the genuinely shared records
\code{IoData}, \code{OpParams}, \code{SimState}); three middle libraries
(\code{iteration}, \code{data-algebra}, \code{coordination}) that \emph{mirror} the
three calculus doc-groups of \cref{ch:combinators}; and a \code{drivers} cap that
\emph{composes} them into Part~V's entry-point surfaces. The bracket shape is forced
by \emph{topological order}---a def appears after everything it uses---so shared types
come first and composing drivers last; the type-placement rule keeps the floor small
by placing single-group records with their group. The dependency graph has two edge
classes: \code{depends-on} (a real call; blocking; constrains the order; the downward
arrows of the stack) and \code{reference} (navigational; free; the cross-links),
whose distinction keeps the graph well-founded. The load-bearing discipline is that
the library is a \emph{view}: it renders the def \emph{bodies} and \emph{links} to
each operator's home for the laws, meaning, and record schema---stated once
there---never restating them, exactly as generated code is a view of its source.
Finally, three opaque kernels (libCEED quadrature in \code{fe\_assemble}, the SLEPc
eigensolve loop in \code{eigsolve}, the ODE step in \code{fold\_solve}) render as
\code{\#extern NAME} after their type signature: the library states the contract but
does not fabricate the body. Each \code{\#extern} kernel-\emph{API} has a separate
kernel-\emph{implementation}---a constructive realization in our own vocabulary
(e.g. \cref{ch:krylovschur} for the eigensolve)---tied to it by a \emph{reviewable}
\code{realizes-kernel-api} link. The payoff: a hundred-thousand-line simulator becomes
a few hundred lines across five small modules; \cref{ch:fivelibs} tours them and
\cref{ch:composing} composes one driver as code.

\furtherreading
The pseudo-language the library is rendered in---Haskell-style type signatures over
pure functions, with do-notation for the state monad---is the notation of functional
programming; the standard reference for the language that crystallized it is the
\emph{Haskell} report and Peyton Jones's account of its design~\citep{peytonjones2003},
and the modular, type-directed style the five-library split follows is its everyday
practice. The \code{\#extern} libCEED quadrature kernel is documented in the CEED
project's matrix-free finite-element interface~\citep{ceed2021}, whose element-local
evaluation is exactly the contract \code{assemble\_term} renders as external and
\cref{ch:matrixfree} reconstructs as tensor contractions. The simulator whose surface
this library synthesizes is Palace~\citep{palace2023}; the synthesized form here is a
clean-room \emph{view} of that codebase's algorithm, not a transcription of its C++,
and the correspondence between the two---rendered body against cited source---is the
reviewable link the whole book has been building. The three-view framing and the
combinator-primary organization the middle libraries mirror were developed in
\cref{ch:bigidea} and \cref{ch:combinators}; this chapter only assembles them into an
architecture.

\exercisessection
\begin{enumerate}[nosep]
  \item \textbf{Name the three views.} State, in one sentence each, the question
  that the top-down view, the bottom-up view, and the implementation view answer
  about the simulator. Which Parts of the book take each view? Why is the
  implementation view the natural \emph{last} one, given the choice
  \cref{ch:bigidea} made about how to describe a computation?
  \item \textbf{Place a record.} The Krylov workspace \code{Krylov} is read only by
  the iteration kernels; the configuration record \code{IoData} is read by every
  driver and the coordination caps. Using the type-placement rule
  (\S\ref{sec:typeplacement}), say which module each belongs in and where within
  it, and explain why \code{IoData} earns a place in the \code{types} floor while
  \code{Krylov} does not.
  \item \textbf{Two edge classes.} The driven driver consumes port modes produced by
  the boundary-mode analysis but does not call \code{boundary\_mode}. Is the link
  from \code{driven} to \code{boundary\_mode} a \code{depends-on} edge or a
  \code{reference} edge? Justify your answer from the definitions, and explain what
  would go wrong with the topological order if every such cross-link were treated as
  a \code{depends-on}.
  \item \textbf{The view discipline.} The library's rendered \code{linear\_combination}
  is a few lines of \code{foldl} and contains \emph{no} statement of the
  concatenation-homomorphism law. Where does that law live, and why is it
  \emph{deliberately} absent from the rendered code? State the failure mode the
  ``view, not a copy'' discipline (\S\ref{sec:viewvssem}) is designed to prevent.
  \item \textbf{Mark the externs.} Of the three \code{\#extern} kernels (libCEED
  quadrature, the SLEPc eigensolve loop, the ODE time-step), say which module's def
  each appears inside, and what part of that def is \emph{ours} versus external. For
  one of them, explain in a sentence what the corresponding kernel-\emph{implementation}
  provides and which chapter gives it.
  \item \textbf{Topological check at a node.} Take \code{solve\_family}, which is
  \lstinline[style=l4style]|map (\inp -> ksp_solve op inp) rhss|. List every name on
  its right-hand side and assign each to a module; confirm that none of them sits in
  a module \emph{above} \code{coordination} (where \code{solve\_family} lives). What
  single fact about \code{solve\_family}'s dependencies makes it legal to place it in
  \code{coordination}?
  \item \textbf{Why the library is short.} The original Palace is roughly
  $10^5$ lines of C++; the synthesized library is a few hundred lines across five
  modules. Drawing on \cref{ch:bigidea}'s distinction between transparent
  optimization tricks and the base algebra, give two reasons the synthesized form is
  so much smaller, and name one thing that is \emph{deliberately not} reproduced in
  it. Is the library a \emph{port} of Palace? Explain the difference.
\end{enumerate}

\chapter{A Tour of the Five Libraries}
\label{ch:fivelibs}

\chapterorient{The previous chapter drew the architecture of the synthesized
library---five modules, a topological order, an \texttt{\#extern} boundary where
an opaque kernel sits. This chapter walks the shelves. We open four of the five
libraries one at a time---the shared \emph{types}, and the three calculus
libraries \emph{iteration}, \emph{data-algebra}, and \emph{coordination}---and
read the actual synthesized code each holds: the record definitions, the def
bodies, the \texttt{where}-local leaves, the kernel callouts. (The fifth library,
\emph{drivers}, composes these four into the five analyses and is the subject of
\cref{ch:composing}.) Two def bodies get a line-by-line reading. By the end you
will have seen the whole vocabulary of the simulator rendered as code, and you
will have noticed the one fact the tour exists to make vivid: across all four
libraries the same handful of shapes keeps reappearing, because the library is
small for the same reason the calculus is small.}

\section{Four libraries, one tour}

\Cref{ch:synthlib} laid out the synthesized library as five modules in a
topological stack: a foundational \emph{types} module of shared records, the
three calculus modules---\emph{iteration}, \emph{data-algebra},
\emph{coordination}---and a top \emph{drivers} module that composes them into the
five simulation analyses. A module appears in the stack \emph{after} everything it
uses: \emph{types} defines records the others thread; \emph{iteration} holds the
loop combinator the solvers fold; \emph{data-algebra} holds the tensor verbs the
steps call; \emph{coordination} holds the outer caps that drive the loops; and
\emph{drivers} sits on top, wiring the rest into electrostatics, the driven sweep,
the eigenmode solve, and so on.

This chapter tours the four lower libraries---the shared types and the three
calculus libraries---and shows what each \emph{holds}, in real synthesized code.
\Cref{fig:shelves} is the map for the tour: four labelled shelves, each stocked
with the operators and types we will read. Keep it in view; every section that
follows pulls items off one of these shelves and renders the def body. We save
\emph{drivers} for \cref{ch:composing}, because a driver is not a new kind of
thing---it is a \emph{composition} of the things on these four shelves, and you
cannot compose a vocabulary you have not yet met.

\begin{figure}[t]
\centering
\begin{tikzpicture}[font=\footnotesize, node distance=4mm]
  \def\shelfw{34mm}
  % ---- types ----
  \node[stage, minimum width=\shelfw, align=center, fill=simBgGray, draw=simGray]
    (tH) at (0,0) {\textbf{types}\\\scriptsize shared records};
  \node[opbox, minimum width=\shelfw, align=left, below=of tH] (tB)
    {\code{IoData} \scriptsize\itshape (config)\\
     \code{OpParams} \scriptsize\itshape (build-time)\\
     \code{SimState} \scriptsize\itshape (run-time)};
  % ---- iteration ----
  \node[stage, minimum width=\shelfw, align=center, fill=simBgGreen, draw=simGreen,
    right=9mm of tH] (iH) {\textbf{iteration}\\\scriptsize loops \& steps};
  \node[opbox, minimum width=\shelfw, align=left, below=of iH] (iB)
    {\code{iterate\_while}\\
     \code{iterate\_while\_with\_prev}\\
     \code{krylov\_step} \scriptsize\itshape (A/B)\\
     \code{chebyshev}\\
     \scriptsize\itshape \code{Krylov} \code{StepOutputs} \code{PrevCarry}};
  % ---- data-algebra ----
  \node[stage, minimum width=\shelfw, align=center, fill=simBgGold, draw=simGold,
    right=9mm of iH] (dH) {\textbf{data-algebra}\\\scriptsize tensor verbs};
  \node[opbox, minimum width=\shelfw, align=left, below=of dH] (dB)
    {\code{linear\_combination}\\
     \code{inner\_product} \code{dot} \code{nrm2}\\
     \code{mk\_matrix\_free\_operator}\\
     \code{fe\_assemble} \scriptsize\itshape (\#extern)\\
     \scriptsize\itshape gram / energy / S-param / \dots};
  % ---- coordination ----
  \node[stage, minimum width=\shelfw, align=center, fill=simBgBlue, draw=simBlue,
    right=9mm of dH] (cH) {\textbf{coordination}\\\scriptsize outer caps};
  \node[opbox, minimum width=\shelfw, align=left, below=of cH] (cB)
    {\code{Solve} \scriptsize\itshape monad\\
     \code{ksp\_solve} \code{eigsolve} \scriptsize\itshape(\#extern)\\
     \code{solve\_family}\\
     \code{frequency\_sweep} \code{fold\_solve}\\
     \scriptsize\itshape \code{Outcome} \code{EigState}};
\end{tikzpicture}
\caption[The four libraries side by side]{The four libraries this chapter tours,
each a shelf stocked with the operators and types we will render. \emph{types}
holds the three cross-cutting records every solve threads. \emph{iteration} holds
the loop combinator of \cref{ch:fold} and the step kernels it folds.
\emph{data-algebra} holds the tensor-combining verbs of \cref{ch:combinators}.
\emph{coordination} holds the \code{Solve}-monadic outer caps that drive the
loops over families. The fifth library, \emph{drivers}, composes all four into the
five analyses and is \cref{ch:composing}. Read left to right: the order is
topological---each shelf uses only the shelves to its left.}
\label{fig:shelves}
\end{figure}

One convention governs every code block below, and it is worth stating once. These
libraries are the \emph{implementation view} of the spec: they render the
synthesized code, but they do not restate the laws or the semantics, which live in
the calculus chapters. So each def we read here links \emph{up} to its authoritative
home---\code{iterate\_while}'s laws are in \cref{ch:fold}, \code{linear\_combination}'s
in \cref{ch:combinators}, and so on. The code is what is new; the meaning we have
already built. With that, we open the first shelf.

\section{The \texttt{types} library: three records, three lifetimes}
\label{sec:types}

The foundational shelf holds the records that \emph{every} other library threads.
There are exactly three that are genuinely cross-cutting---shared across two or
more of the calculus libraries---and they are the three lifetimes we named in
\cref{ch:strata}: a configuration parsed once and read forever, a bundle of
operator parameters fixed at construction, and a solve-state that evolves as the
loop runs. \Cref{fig:typecards} draws them as three cards, each tagged with its
stratum.

\begin{figure}[t]
\centering
\begin{tikzpicture}[font=\footnotesize, node distance=6mm]
  \def\cardw{40mm}
  % IoData card
  \node[draw=simGray, fill=simBgGray, rounded corners=3pt, thick,
    minimum width=\cardw, align=left, inner sep=5pt] (io) at (0,0)
    {\textbf{\code{IoData}}\\[1pt]\scriptsize\itshape parse-time, read-only\\[3pt]
     \scriptsize problem / model / domains\\
     \scriptsize boundaries / solver / units\\[2pt]
     \scriptsize\textbf{selects the driver}\\\scriptsize via \code{problem.type}};
  \node[band, fill=simGray, text=white, font=\scriptsize\bfseries, above=1mm of io]
    {stratum 0: parse once};
  % OpParams card
  \node[draw=simGreen, fill=simBgGreen, rounded corners=3pt, thick,
    minimum width=\cardw, align=left, inner sep=5pt, right=10mm of io] (op)
    {\textbf{\code{OpParams}}\\[1pt]\scriptsize\itshape build-time, read-only\\[3pt]
     \scriptsize \code{T} (apply surface)\\
     \scriptsize \code{eps} (stopping test)\\
     \scriptsize variant selectors,\\\scriptsize tolerances, \code{max\_it}};
  \node[band, fill=simGreen, text=white, font=\scriptsize\bfseries, above=1mm of op]
    {stratum 1: fixed at build};
  % SimState card
  \node[draw=simRust, fill=simBgRust, rounded corners=3pt, thick,
    minimum width=\cardw, align=left, inner sep=5pt, right=10mm of op] (ss)
    {\textbf{\code{SimState}}\\[1pt]\scriptsize\itshape run-time, evolving\\[3pt]
     \scriptsize \code{x} (the iterate)\\
     \scriptsize \code{it} \code{converged}\\
     \scriptsize \code{final\_res} \code{initial\_res}\\[2pt]
     \scriptsize\textbf{uniform across solvers}};
  \node[band, fill=simRust, text=white, font=\scriptsize\bfseries, above=1mm of ss]
    {stratum 2: evolves each step};
\end{tikzpicture}
\caption[The \texttt{types} records as stratum-tagged cards]{The three
cross-cutting records of the \emph{types} library, each tagged with its stratum
from \cref{ch:strata}. \code{IoData} is parsed once from the user's config and
read for the rest of the run; its \code{problem.type} field selects which driver
fires. \code{OpParams} is fixed at solver construction and read but never written
during the solve. \code{SimState} is the one record that \emph{evolves}---a fresh
value each step---and it has the \emph{same five fields} for every Krylov-shaped
solver in the book. Three records, three lifetimes; the type-placement rule keeps
only these cross-cutting three here.}
\label{fig:typecards}
\end{figure}

\subsection{\texttt{IoData}: the configuration parsed once}

The first record is the whole configuration, parsed once at startup from the
user's JSON file and then read---never written---for the rest of the run. It is an
aggregate of six sub-records: the problem type and solver-pipeline knobs, the mesh
and material model, the per-domain data, the boundary conditions, the solver
settings, and a units converter.

\begin{lstlisting}[style=l4style]
-- The aggregate parsed once at startup; read-only across the whole solve.
IoData = {
  problem    : ProblemConfig,   -- driver selector (problem.type) + pipeline knobs
  model      : ModelConfig,     -- mesh file + refinement + material assignment
  domains    : DomainConfig,    -- per-domain materials + postprocess energy regions
  boundaries : BoundaryConfig,  -- BC surfaces (PEC/PMC/impedance/lumped-/wave-port)
  solver     : SolverConfig,    -- linear/eigen/driven/transient settings + tolerances
  units      : Units            -- SI <-> nondimensional scale converter
}
-- utility API (parse-time only)
parseConfig :: FilePath -> IoData        -- parse the JSON config tree once
problemType :: IoData -> ProblemType     -- the driver-dispatch selector (a projection)
\end{lstlisting}
The load-bearing field is \code{problem.type}. It is read once, at the very top of
the run, by the dispatch that chooses which of the five drivers to invoke---the
single point where the ``six analyses, one machine'' of \cref{ch:targets} becomes
a concrete branch. The function \code{problemType} that reads it is a trivial
projection: \code{IoData} carries no behaviour, only data. Everything an analysis
does with the configuration, it does by reading fields out of this one immutable
object.

\subsection{\texttt{OpParams}: the operator's build-time parameters}

The second record holds the parameters fixed when the operator is
\emph{constructed}---the variant selectors, the tolerances, and the
constructed-operator surfaces the kernel will apply. It is read by the step kernel
(through its surfaces) and captured once by the outer caps; it never changes during
the solve.

\begin{lstlisting}[style=l4style]
-- read-only; captured once at solve construction, never re-inspected per step.
OpParams = {
  T          : ConstructedOp,    -- the apply surface (preconditioned operator)
  orthog?    : OrthogSurface,    -- GMRES/Arnoldi only; absent (no-op) for CG
  scalars?   : ScalarSurface,    -- Chebyshev polynomial-recurrence scalars; else absent
  eps        : Convergence,      -- the stopping-predicate surface
  -- variant selectors (closed over by the surfaces above; not read by the body)
  pc_side    : PreconditionerSide,
  gs_orthog  : Orthogonalization,
  flexible   : Bool,
  poly_kind? : PolynomialKind,
  restart    : RestartMode,
  -- termination knobs (close into eps)
  max_dim : Int, max_it : Int, rel_tol : Scalar, abs_tol : Scalar
}
\end{lstlisting}
Notice what \cref{ch:strata} promised: the variant selectors (\code{pc\_side},
\code{gs\_orthog}, \code{flexible}, \dots) are \emph{folded into the surfaces}
\code{T}, \code{orthog}, \code{eps} at construction. The kernel never branches on
``am I CG or GMRES?''---that question is answered once, when the surfaces are
built, and the body just applies \code{T} and \code{eps} as given. The
question-marks (\code{orthog?}, \code{scalars?}, \code{poly\_kind?}) mark fields
present for some solvers and absent for others: CG has no orthogonalization
surface, Chebyshev has no Krylov restart. A solver \emph{is} a choice of which of
these are present.

\subsection{\texttt{SimState}: the run-time solve state}

The third record is the only one that evolves. It carries the externally-visible
quantities a Krylov-shaped solve produces and reports, and---this is the fact worth
underlining---it has the \emph{same five fields for every solver in the book}.

\begin{lstlisting}[style=l4style]
-- the value threaded by `Solve a = StateT SimState Identity a`; every field run-time.
-- the iterate `x` is named with shape group S, not a fixed rank-1 axis.
SimState = {
  x           : Tensor[(S: ...)],  -- the current iterate (the solve's primary product)
  it          : Int,               -- iteration count
  converged   : Bool,              -- convergence flag
  final_res   : Scalar,            -- final (absolute) residual, possibly an estimate
  initial_res : Scalar             -- initial (absolute) residual, captured at entry
}
\end{lstlisting}

\begin{keyidea}
CG, GMRES, FGMRES, and Chebyshev all evolve the \emph{same} \code{SimState}: an
iterate \code{x}, a counter \code{it}, a converged flag, and two residuals. The
solvers differ in their \emph{ephemeral} workspace---CG's single search direction,
GMRES's growing basis---but the run-time state they expose to the outside world is
one five-field record. This is why \code{ksp\_solve} can be \emph{one} cap over
\emph{four} solvers (\cref{sec:coordination}): the state it threads is uniform, and
the differences live below it, in the step kernel's private \code{Krylov} bundle.
\end{keyidea}

\begin{notationbox}
Only these three records live in \emph{types}, because only these three are shared
across two or more libraries. A record used by just \emph{one} library---CG's
\code{Krylov} workspace, the eigensolver's \code{EigState}---is placed with that
library instead, bundled next to the operators that consume it. The
\emph{type-placement rule} (\cref{ch:synthlib}) is exactly this: cross-cutting
records go in \emph{types}; single-library records go with their library. We meet
the single-library records as we reach each shelf.
\end{notationbox}

\section{The \texttt{iteration} library: the loop and its kernels}
\label{sec:iteration}

The second shelf holds the heart of the framework: the one loop combinator of
\cref{ch:fold} and the step kernels it folds. Four operators in topological
order---\code{iterate\_while} first (everything below it folds it), then its
two-term-recurrence sibling, then the Krylov step, then the Chebyshev smoother---and
three single-library workspace types rendered just ahead of them. This is the
library whose first def we read line by line, in \cref{wex:readwhile}.

\subsection{The clustering types: \texttt{Krylov}, \texttt{StepOutputs}, \texttt{PrevCarry}}

Before the operators come the records they thread, the ones the type-placement rule
keeps \emph{here} rather than in \emph{types} because only \emph{iteration} uses
them. \code{Krylov} is the ephemeral per-restart workspace---born at restart entry,
discarded at restart exit, threaded as a \emph{plain value} (not a monadic effect,
because its lifetime is strictly inside one cycle). Its fields differ by solver:

\begin{lstlisting}[style=l4style]
-- Ephemeral workspace; born at restart, discarded at restart exit. Plain value.
type Krylov_CG = {
  r  : Tensor[(S: ...)],  -- residual
  p  : Tensor[(S: ...)],  -- search direction
  z? : Tensor[(S: ...)],  -- preconditioned residual (iff a preconditioner is set)
  alpha : Scalar,         -- step length
  beta  : Scalar          -- direction-update coeff (the convergence-test proxy)
}
type Krylov_GMRES = {
  V  : [Tensor[(S: ...)]],  -- Arnoldi basis (array of basis columns)
  Z? : [Tensor[(S: ...)]],  -- preconditioned basis (iff OpParams.flexible -- FGMRES)
  H  : DenseMatrix,         -- upper-Hessenberg (small-dense, scalar-stratum)
  s, cs, sn : [Scalar],     -- LS RHS / Givens cosines / Givens sines (small-dense)
  beta : Scalar, j : Int    -- residual proxy; inner index within the restart cycle
}
\end{lstlisting}
The contrast is the whole story of CG versus GMRES: CG's workspace is a fixed
handful of vectors, so its memory is bounded; GMRES's \code{V} \emph{grows} with
each inner iteration, which is exactly why GMRES restarts. \code{StepOutputs} is
the demand-prunable per-step readout (a residual norm, an optional
least-squares residual, an optional breakdown token)---the ``extras'' of
\cref{ch:fold}, rendered as a record. And \code{PrevCarry} is the one-slot
recurrence carry that Form~B of the Krylov step threads through the loop instead of
storing in the state---the first-iteration-unrolling of \cref{ch:fold} made
concrete.

\subsection{\texttt{iterate\_while}: the def body itself}

Here is the combinator the whole book turns on, rendered as synthesized code. We
read it in full in \cref{wex:readwhile}; for now, see it whole.

\begin{lstlisting}[style=l4style]
-- # Arguments
--   a    : a          -- the value-threaded carry (immutable); init
--   cont : a -> Bool  -- loop predicate, PURE on the carry; fires before each step
--   step : a -> Solve { state: a, ...e }  -- per-step body; Solve effect on SimState
-- # Returns
--   Solve { final_state: a, trajectory: [{ ...e }] }
--     trajectory -- per-step extras in order; demand-pruned when only final_state read
iterate_while :: a -> (a -> Bool) -> (a -> Solve { state: a, ...e })
              -> Solve { final_state: a, trajectory: [{ ...e }] }
iterate_while a cont step =
  if cont a
    then do
      { state: a', ...e } <- step a
      { final_state, trajectory } <- iterate_while a' cont step
      pure { final_state, trajectory: [{ ...e }] ++ trajectory }
    else
      pure { final_state: a, trajectory: [] }

-- the no-extras sugar (a bounded loop): e = (), body non-monadic.
iterate_while_pure :: a -> (a -> Bool) -> (a -> a) -> a
iterate_while_pure a cont f =
  (iterate_while a cont (\x -> pure { state: f x })).final_state
\end{lstlisting}
This is the small-step rule of \cref{ch:fold} transcribed verbatim into code: test
\code{cont a}; if true, run one \code{step}, recurse on the new carry, and prepend
this step's extras to the trajectory; if false, stop with the carry as
\code{final\_state} and an empty trajectory. The recursion sits in tail position, so
it is a loop, not a stack. The \code{iterate\_while\_pure} below it is the sugar for
loops with nothing to report---the carry threads, the trajectory is always empty,
and you read back only the final value. \Cref{fig:whilestructure} diagrams the
structure of the def; \cref{wex:readwhile} reads it line by line.

\begin{figure}[t]
\centering
\begin{tikzpicture}[font=\footnotesize, node distance=5mm]
  % the predicate gate
  \node[diamond, draw=simBlue, fill=simBgBlue, aspect=1.7, inner sep=1pt,
    font=\scriptsize] (gate) at (0,0) {\code{cont a}?};
  % FALSE branch (else)
  \node[product, minimum width=30mm, align=left, right=20mm of gate] (base)
    {\scriptsize\textbf{else (base case)}\\\scriptsize \code{final\_state: a}\\\scriptsize \code{trajectory: []}};
  \draw[flow] (gate.east) -- node[above,font=\scriptsize,text=simRust]{no} (base.west);
  % TRUE branch (then) -- three lines stacked below
  \node[opbox, minimum width=46mm, align=left, below=10mm of gate] (s1)
    {\scriptsize 1.\ \code{\{state: a', ...e\} <- step a}\\\scriptsize\itshape one step: next carry + extras};
  \node[opbox, minimum width=46mm, align=left, below=4mm of s1] (s2)
    {\scriptsize 2.\ \code{... <- iterate\_while a' cont step}\\\scriptsize\itshape tail recurse on the new carry};
  \node[opbox, minimum width=46mm, align=left, below=4mm of s2] (s3)
    {\scriptsize 3.\ \code{pure \{..., trajectory: [e] ++ ...\}}\\\scriptsize\itshape prepend this step's extras};
  \draw[flow] (gate.south) -- node[left,font=\scriptsize,text=simGreen]{yes} (s1.north);
  \draw[flow] (s1) -- (s2);
  \draw[flow] (s2) -- (s3);
  % the carry feedback to the gate (recursion)
  \draw[backflow] (s2.east) .. controls +(2.4,0) and +(2.4,0) ..
    node[right,font=\scriptsize,text=simRust,align=left,pos=0.5]{recurse:\\carry \code{a'}} (gate.east);
\end{tikzpicture}
\caption[The structure of the \texttt{iterate\_while} def]{The synthesized
\code{iterate\_while} def as a structure diagram. The predicate \code{cont a} is
tested \emph{first} (the \texttt{while} convention of \cref{ch:fold}). On
\textcolor{simRust}{no}, the \texttt{else} arm returns the base case---the carry as
\code{final\_state}, an empty trajectory. On \textcolor{simGreen}{yes}, the
\texttt{then} arm runs three lines: take one step (producing the next carry \code{a'}
and this step's extras \code{e}), tail-recurse on \code{a'} (the dashed feedback
edge), and prepend \code{e} to the trajectory the recursion returns. Three lines and
a gate; \cref{wex:readwhile} reads each in turn.}
\label{fig:whilestructure}
\end{figure}

\subsection{The step kernels: \texttt{krylov\_step} and \texttt{chebyshev}}

Above the loop sit the kernels it folds. \code{krylov\_step} is the per-step body
of a Krylov solve, embedded in the \code{Solve} monad. It comes in two
\emph{forms}, and seeing them side by side is the clearest picture of the
first-iteration-unrolling rotation from \cref{ch:fold}. \emph{Form~A}---the
default---keeps the ``is this the first iteration?'' branch \emph{inside} the body
and folds with plain \code{iterate\_while}:

\begin{lstlisting}[style=l4style]
-- Form A (branch-in-body, default): folded by iterate_while.
krylov_step :: OpParams -> Krylov -> (SimState -> Solve { krylov: Krylov, outputs: StepOutputs })
krylov_step op K = \s -> do
  let w     = apply_linop op.T K.input_field   -- operator apply (no SimState read)
  let K_aux = optionally_apply_auxiliary op K w -- GMRES orthog / Chebyshev scalars / CG no-op
  let K'    = krylov_update K_aux op w          -- bundle update (axpy/dot/nrm2, pure on K)
  let outputs = derived_views K' op             -- residual_norm; ls_residual; breakdown
  modify (\s -> s { it = s.it + 1 })            -- the sole monadic effect
  pure { krylov: K', outputs }
\end{lstlisting}
\emph{Form~B}---opt-in---splits the body into a \code{first\_step} that
manufactures the initial recurrence value and a \emph{branch-free}
\code{steady\_step} that consumes it, folded by
\code{iterate\_while\_with\_prev}. The canonical instance is CG, and reading the
two CG steps next to each other shows what unrolling buys: the steady step's
\code{p'} line is unconditional, where Form~A's would carry an \code{if it==0}
test.

\begin{lstlisting}[style=l4style]
-- Form B, CG: the steady step is branch-free; beta_prev is the closure carry.
cg_steady_step :: LinOp[(S: ...)] -> Scalar -> Scalar -> CgState
              -> { state: CgState, residual_norm: Scalar }
cg_steady_step opA eps beta_prev s =
  let p'    = axpby 1.0 s.r (s.beta / beta_prev) s.p in  -- unconditional: p <- r + (beta/beta_prev) p
  let Ap    = apply_linop opA p' in
  let alpha = s.beta / dot Ap p' in
  let x'    = axpy alpha p' s.x in                       -- x <- x + alpha p
  let r'    = axpy (negate alpha) Ap s.r in              -- r <- r - alpha Ap
  let beta' = dot r' r' in
  let res'  = sqrt (abs beta') in
  { state: { x: x', r: r', p: p', beta: beta', it: s.it + 1, converged: res' < eps },
    residual_norm: res' }
\end{lstlisting}
Every line of that step is a \emph{data-algebra} call---\code{axpby}, \code{axpy},
\code{dot}---which is the cue for the next shelf. But notice first that the step's
shape is exactly the carry-threading fold of \cref{ch:fold}: a record in, a record
out, the termination quantity (\code{beta}, \code{converged}) folded into the
carry so the predicate can read it. \code{chebyshev}, the fourth operator, is the
same story at two nesting levels: an outer Richardson sweep and an inner
polynomial recurrence, both written as bounded \code{iterate\_while\_pure} folds
with step-count predicates. We do not reproduce its body here---it is long, and
\cref{ch:matrixfree} reads it as a smoother---but its shape is no surprise: a fold
inside a fold, each a setting of the one combinator.

\begin{notationbox}
The \emph{iteration} library has \emph{no} \code{\#extern} kernel of its own. The
loop combinators are pure calculus, and the opaque library boundaries the steps
eventually hit---the libCEED quadrature inside \code{apply\_linop}, the SLEPc
eigensolve loop---belong to the operators that \emph{own} those applies, and render
\code{\#extern} on the \emph{data-algebra} and \emph{coordination} shelves
instead. Here, \code{apply\_linop} is an abstract surface the step folds; its
implementation lives where the operator is built.
\end{notationbox}

\section{The \texttt{data-algebra} library: the tensor verbs}
\label{sec:dataalgebra}

The third shelf holds the verbs that \emph{combine} tensors---the right-hand branch
of the combinator vocabulary in \cref{ch:combinators}. Every line of the CG step we
just read was one of these. There are two combinators at the top
(\code{linear\_combination} and \code{inner\_product}), their named leaves and
consumers, an operator constructor, the assemble fold with its one \code{\#extern}
kernel, and a family of reduce verbs. We render the headline def and survey the
rest.

\subsection{\texttt{linear\_combination} and its arity leaves}

The most-used verb in the book is \code{linear\_combination}, the scalar-weighted
tensor sum $\sum_i a_i t_i$. It is a \code{foldl} over a term list, and the four
BLAS routines hang off it as \code{where}-local arity leaves---specializations at a
fixed term-list length, exactly the picture of \cref{fig:lincomb}.

\begin{lstlisting}[style=l4style]
-- The scalar-weighted-tensor-sum combinator: a pure fold over a (Scalar, Tensor)
-- term list, producing the sum of a_i * t_i. No Solve monad, no carry.
-- Laws (concatenation-homomorphism, multilinearity): see ch:combinators.
linear_combination :: [(Scalar, Tensor[(S: ...)])] -> Tensor[(S: ...)]
linear_combination pairs = foldl (\acc (a, t) -> acc + scal a t) (zeros) pairs
  where
    -- the four BLAS-1 arity leaves: this combinator at a fixed term-list length
    -- (specialization notes, not co-equal ops -- accelerated kernels stopped low).
    scal     a x          = linear_combination [(a, x)]
    axpy     a x y        = linear_combination [(a, x), (1, y)]
    axpby    a x b y      = linear_combination [(a, x), (b, y)]
    axpbypcz a x b y c z  = linear_combination [(a, x), (b, y), (c, z)]
\end{lstlisting}
Read the four leaves top to bottom: \code{scal} is a one-term list, \code{axpy} a
two-term list with the second coefficient pinned to $1$, \code{axpby} the general
two-term list, \code{axpbypcz} the three-term list. They differ only in arity, and
the concatenation law makes the chain exact---\code{axpbypcz}'s list is
\code{axpby}'s list concatenated with \code{scal}'s. \Cref{fig:lincombleaves}
redraws this as the trunk-and-leaves the synthesized code makes literal: one def, a
\code{where}-block of four one-liners.

\begin{figure}[t]
\centering
\begin{tikzpicture}[font=\footnotesize, node distance=5mm]
  % trunk
  \node[stage, minimum width=66mm, align=center] (trunk) at (0,0)
    {\code{linear\_combination} \;\itshape\quad sum of $a_i\,t_i$, term list, any length};
  \def\ly{-1.7}
  \node[opbox, minimum width=24mm, align=center] (l1) at (-4.8,\ly)
    {\code{scal}\\\scriptsize arity 1\\\scriptsize \code{[(a,x)]}};
  \node[opbox, minimum width=24mm, align=center] (l2) at (-1.6,\ly)
    {\code{axpy}\\\scriptsize arity 2\\\scriptsize \code{[(a,x),(1,y)]}};
  \node[opbox, minimum width=24mm, align=center] (l3) at (1.6,\ly)
    {\code{axpby}\\\scriptsize arity 2\\\scriptsize \code{[(a,x),(b,y)]}};
  \node[opbox, minimum width=24mm, align=center] (l4) at (4.8,\ly)
    {\code{axpbypcz}\\\scriptsize arity 3\\\scriptsize \code{[..,(c,z)]}};
  \foreach \l in {l1,l2,l3,l4}{ \draw[refedge] (trunk.south) -- (\l.north); }
  \node[font=\scriptsize\itshape, text=simGray, align=center] at (0,-2.9)
    {rendered as a \code{where}-block of four one-liners beneath the one def:\\
     the combinator is the entry; the leaves are settings of it};
\end{tikzpicture}
\caption[\texttt{linear\_combination} with its \texttt{where}-local leaves]{The
synthesized \code{linear\_combination} and its four BLAS arity leaves, as the code
renders them: one def for the combinator (trunk), a \code{where}-block of four
one-liners beneath it (leaves). Each leaf is the combinator at a fixed term-list
length---\code{scal} at one term, \code{axpy}/\code{axpby} at two, \code{axpbypcz}
at three. This is the combinator-primary model of \cref{ch:combinators} made
literal in the library: the fused kernels are ``stopped low'' as notes, and the
combinator rises to be the entry. Compare \cref{fig:lincomb}.}
\label{fig:lincombleaves}
\end{figure}

\subsection{\texttt{inner\_product}, \texttt{dot}, and \texttt{nrm2}}

The scalar-producing sibling is \code{inner\_product}, $\langle x, y\rangle_M =
x^{\!\top} M y$, a reduction over the shape group. \code{dot} is its
$M = I$ \emph{leaf} (a specialization, hung below), and \code{nrm2} is a
\emph{consumer}---$\sqrt{\,\cdot\,}\circ\code{dot}$ at the diagonal, built
\emph{above} \code{dot}, not beneath it. The synthesized code keeps the distinction
visible:

\begin{lstlisting}[style=l4style]
inner_product :: Tensor[(S: ...)] -> Tensor[(S: ...)] -> Scalar
inner_product x y = reduce (+) zero (zipWith kernel x y)
  where kernel xi yi = conj_if_complex xi * yi   -- Hermitian: arg-1 conjugated

dot  x y = inner_product x y                       -- the M = I leaf
nrm2 :: Tensor[(S: ...)] -> Scalar
nrm2 x = sqrt (abs (inner_product x x))            -- a CONSUMER of inner_product
\end{lstlisting}
The \code{abs} inside \code{nrm2} is not decoration---it is the defensive
non-negativity guard that keeps a round-off-negative radicand from reaching the
square root, a load-bearing numerical detail preserved as an explicit part of the
scalar map (\cref{ch:bigidea}). The CG step's \code{dot Ap p'} and the residual's
\code{sqrt (abs beta')} we read in \cref{sec:iteration} are exactly these verbs.

\subsection{The operator constructor and the assemble fold}

Two larger items sit on this shelf. \code{mk\_matrix\_free\_operator} builds an
operator \emph{value}---a closure whose \code{apply} is the five-stage contraction
chain $G^{\!\top}\!\cdot B^{\!\top}\!\cdot D \cdot B \cdot G$ of \cref{ch:matrixfree},
rendered inline because that chain \emph{is} the implementation:

\begin{lstlisting}[style=l4style]
-- A = G^T . B_D^T . D . B_D . G  (the un-materialized contraction chain; ch:matrixfree)
mk_matrix_free_operator space term geom = mkOp (\v -> apply_chain v)
  where apply_chain v =
    element_restrict_T (basis_apply_T term
      (quad_point_contract geom (basis_apply term (element_restrict space v))))
\end{lstlisting}
And \code{fe\_assemble} folds the weak-form-term list into a global operator,
$K = \sum_t \code{assemble\_term}(\text{space}, t)$---a \code{foldr} over a
commutative-monoid sum. Its per-term leaf is the \emph{first} \code{\#extern} we
meet: \code{assemble\_term} is the libCEED-owned quadrature kernel, rendered as a
callout in place of a body.

\begin{lstlisting}[style=l4style]
fe_assemble space terms = foldr (\t acc -> assemble_term space t + acc) zero terms

-- The opaque per-term leaf: ONE weak-form term to its global-dof contribution.
-- libCEED-owned (the kernel-API boundary). Rendered #extern after its signature.
assemble_term :: FiniteElementSpace[N] -> WeakFormTerm -> LinOp[(N: ...)]
#extern assemble_term
\end{lstlisting}

\subsection{The reduce verbs}

The bottom of the shelf holds the reductions of \cref{ch:reductions}---the verbs
that turn a solved field into the physical product. There are five firm ones, each
a \code{map}-then-collect over a family with no inter-element state:
\code{gram\_reduce} (the symmetric Gram matrix behind capacitance and inductance),
\code{domain\_energy\_reduce} (the per-domain energy table), \code{sparameter\_reduce}
(the driven scattering matrix), \code{eigenfreq\_qfactor\_reduce} (the per-mode
frequency and quality factor), and \code{waveguide\_mode\_reduce} (the boundary-mode
field table). They share a shape but deliberately not a definition---a scalar table,
a matrix, and a field-carrying table are different products, and \cref{ch:reductions}
explains why forcing them together would be the wrong design. One example shows the
shape:

\begin{lstlisting}[style=l4style]
-- per-mode (f, Q) table over the converged eigenpairs: map-then-collect, no state.
eigenfreq_qfactor_reduce ptype kappa eigs =
  [ let omega = untransform ptype lambda            -- eigenvalue -> angular frequency
        f     = re omega                            -- eigenfrequency f = Re omega
        k     = kappa mode                          -- loss rate
        q     = if k == 0 then infinity else f / abs k  -- quality factor Q = omega / kappa
    in (f, q)
  | (lambda, _E) <- eigs ]                          -- map over the eigenpair family
\end{lstlisting}
The \code{if k == 0 then infinity} is the lossless-cavity guard---a cavity with no
loss has infinite $Q$---the same kind of explicit boundary case as \code{nrm2}'s
\code{abs}. Read past the physics and the shape is a \code{map} with a guard: the
combinator-primary model again, this time on the reduce side.

\section{The \texttt{coordination} library: the outer caps}
\label{sec:coordination}

The fourth shelf holds the \code{Solve}-monadic outer caps---the combinators that
drive the iteration kernels to convergence and run them over families. This is the
left-hand branch of \cref{ch:combinators}'s vocabulary and the home of the monad of
\cref{ch:monad}. Its headline def, \code{ksp\_solve}, is the second we read line by
line, in \cref{wex:kspsolve}.

\subsection{The \texttt{Solve} monad and the termination sums}

The coordination surface threads \code{SimState} through a state monad and
classifies termination into sum types. The monad is the thin one of \cref{ch:monad}:

\begin{lstlisting}[style=l4style]
-- The outer-driver state monad: the effect domain is exactly SimState.
type Solve a = StateT SimState Identity a

-- the discharge that turns a monadic cap into a pure function:
execState :: Solve a -> SimState -> SimState   -- run the action, project the state
\end{lstlisting}
And termination is a \emph{value}, a sum type the cap classifies once per cycle.
Krylov solves use a three-arm \code{Outcome}; the eigensolver uses a richer four-arm
\code{EigOutcome} that adds a first-class ``partially converged'' arm with no Krylov
analog:

\begin{lstlisting}[style=l4style]
data Outcome    = Continue | Done Bool                    -- the Krylov 3-arm sum
data EigStatus  = Converged | PartialConverged Int | MaxIterReached | LinearSolveFailed
data EigOutcome = Continue | Done EigStatus               -- the eigen 4-arm extension
\end{lstlisting}
The eigensolver's persistent state, \code{EigState}, is a single-library record
(only \code{eigsolve} names it), so the type-placement rule keeps it here, not in
\emph{types}. It does not collapse to \code{SimState} because it carries a whole
\emph{family} of converged eigenpairs, not one iterate.

\subsection{\texttt{ksp\_solve}: the Krylov cap}

Here is the cap, the def \cref{wex:kspsolve} reads. It seeds a \code{SimState},
runs the outer driver over it, and projects the terminal state with
\code{execState}---turning a monadic computation back into a pure
$\code{OpParams}\to\code{Inputs}\to\code{SimState}$ function.

\begin{lstlisting}[style=l4style]
-- # Returns SimState: terminal .x ~ A^-1 b, plus .it / .converged / .final_res
ksp_solve :: OpParams -> Inputs -> SimState
ksp_solve op inp = execState (solve_loop op inp) (initial_state inp)
  where
    -- outer driver: tail-recurse the per-cycle body until the Outcome says stop.
    solve_loop op inp = do
      o <- restart_cycle op inp
      unless (done o) (solve_loop op inp)

    -- per-cycle body: fresh Krylov, inner kernel-fold, fold the correction, classify once.
    restart_cycle op inp = do
      s <- get
      let k0        = fresh_krylov op inp s                 -- ephemeral bundle (plain value)
          (kn, _os) = iterate_while (krylov_step op k0) cont -- inner fold of the kernel
      modify (\s -> s { x = s.x `plus` (kn.basis `applyBasis` kn.y) })  -- the single per-cycle x write
      pure (classify kn op s)                               -- Done True / Done False / Continue
\end{lstlisting}
The inner \code{krylov\_step} and \code{iterate\_while} are not re-rendered
here---they live on the \emph{iteration} shelf, and the cap simply \emph{folds}
them. That deep link is the topological order paying off: \code{coordination} sits
above \code{iteration}, so it calls down to it.

\subsection{\texttt{eigsolve} and the second \texttt{\#extern}}

\code{eigsolve} is the cap that does \emph{not} own its loop. Where \code{ksp\_solve}
writes the tail-recursion itself, the eigen-iteration runs \emph{inside} a library
(SLEPc's \code{EPSSolve}), so the cap names the iteration by role and renders it
\code{\#extern}---the second kernel boundary in the library:

\begin{lstlisting}[style=l4style]
eigsolve op inp = execState (solve_loop op inp) (initial_eig_state inp)
  where solve_loop op inp = do
    o <- eigen_iterate op inp     -- OPAQUE library fold (no Palace loop)
    classify_into_state o

-- The opaque eigen-iteration: the entire SLEPc EPSSolve loop. The kernel-API boundary.
eigen_iterate :: OpParams -> Inputs -> Solve EigOutcome
#extern eigen_iterate
\end{lstlisting}
This is the deliberate non-loop of \cref{ch:combinators}: the fourth solve corner,
the one outside the map/fold grid because we write no loop at all. The
\code{\#extern} marks exactly where our vocabulary stops and the library's begins.
(\Cref{ch:krylov} separately reconstructs what that loop does internally, in our own
\code{lanczos\_step}/\code{krylov\_step} vocabulary---a glass-box version that
\emph{realizes} this API for the curious reader, without the cap depending on it.)

\subsection{\texttt{solve\_family}, \texttt{frequency\_sweep}, \texttt{fold\_solve}}

The top of the shelf holds the three family combinators of \cref{ch:combinators},
and each is a near one-liner over the caps below it. \Cref{fig:caps} lines all five
coordination caps up as the map/fold/black-box family they form.

\begin{lstlisting}[style=l4style]
-- fixed-operator map: capture op once, map ksp_solve over the RHS family.
solve_family op rhss = map (\inp -> ksp_solve op inp) rhss

-- operator-VARYING map: rebuild A(omega) INSIDE the map, then solve.
frequency_sweep fam omegas =
  map (\w -> ksp_solve (assemble_frequency_operator fam w) (rhs_at fam w)) omegas

-- state-threaded fold: thread the field-state through a schedule by foldl.
fold_solve op s0 schedule = foldl (\s t -> time_step_op op s t) s0 schedule
time_step_op :: OpParams -> TimeState -> Time -> TimeState
#extern time_step_op                              -- the opaque per-step integrator
\end{lstlisting}
The three differ by exactly the axes \cref{ch:combinators} named:
\code{solve\_family} hoists the operator out of the map; \code{frequency\_sweep}
rebuilds it inside; \code{fold\_solve} drops the independence and threads a carry,
folding a third \code{\#extern} per-step integrator. Three lines, three corners.

\begin{figure}[t]
\centering
\begin{tikzpicture}[font=\footnotesize]
  % the family header
  \node[stage, minimum width=64mm, align=center] (hdr) at (0,2.0)
    {the \emph{coordination} caps: \code{map} / \code{fold} / black-box};
  % map row
  \node[opbox, minimum width=30mm, fill=simBgBlue, draw=simBlue, align=center] (sf)
    at (-3.6,0.3) {\code{solve\_family}\\\scriptsize fixed-op \textbf{map}};
  \node[opbox, minimum width=30mm, fill=simBgBlue, draw=simBlue, align=center] (fs)
    at (-0.0,0.3) {\code{frequency\_sweep}\\\scriptsize op-varying \textbf{map}};
  \node[opbox, minimum width=30mm, fill=simBgGreen, draw=simGreen, align=center] (fo)
    at (3.6,0.3) {\code{fold\_solve}\\\scriptsize state-threaded \textbf{fold}};
  % bottom row: the two caps the families call / the black box
  \node[opbox, minimum width=30mm, align=center] (ksp) at (-1.8,-1.6)
    {\code{ksp\_solve}\\\scriptsize the Krylov cap (our loop)};
  \node[extern, minimum width=30mm, align=center] (eig) at (1.8,-1.6)
    {\code{eigsolve}\\\scriptsize black-box (\#extern loop)};
  \draw[depedge] (hdr) -- (sf); \draw[depedge] (hdr) -- (fs); \draw[depedge] (hdr) -- (fo);
  \draw[refedge] (sf.south) -- (ksp.north);
  \draw[refedge] (fs.south) -- (ksp.north);
  \node[font=\scriptsize\itshape, text=simViolet] at (1.8,-2.6) {(no map/fold: a single call)};
\end{tikzpicture}
\caption[The coordination caps as the map/fold/black-box family]{The five
coordination caps as the family of \cref{ch:combinators}. \code{solve\_family} and
\code{frequency\_sweep} are \code{map}s---independent solves, the first hoisting its
operator, the second rebuilding it per member---and both call the Krylov cap
\code{ksp\_solve} (our loop) per element. \code{fold\_solve} is the
\code{fold}---a threaded carry, intrinsically sequential. \code{eigsolve} is the
black-box: a single \code{\#extern} library call, outside the map/fold grid because
no loop is ours to write. One header, three structural shapes.}
\label{fig:caps}
\end{figure}

\section{The same shapes, again and again}
\label{sec:samesahapes}

Step back from the four shelves and look at what we actually rendered. The
\emph{iteration} library was a fold and two kernels that are folds. The
\emph{data-algebra} library was two folds (\code{linear\_combination},
\code{inner\_product}) with named leaves, an operator that is a contraction chain,
and five reductions that are \code{map}s-then-collect. The \emph{coordination}
library was a state monad, three caps that are \code{map}s or a \code{fold}, and
two \code{\#extern} boundaries. And the \emph{types} library was three records,
distinguished only by lifetime. \Cref{fig:samesahapes} tallies the recurring
shapes across the libraries; the count is strikingly small.

\begin{figure}[t]
\centering
\setlength{\tabcolsep}{5pt}
\renewcommand{\arraystretch}{1.25}
\begin{tabular}{@{}lcccc@{}}
\toprule
& \textbf{types} & \textbf{iteration} & \textbf{data-algebra} & \textbf{coordination} \\
\midrule
a \emph{fold}        & --- & \code{iterate\_while} & \code{linear\_combination} & \code{fold\_solve} \\
                     &     & \code{krylov\_step}   & \code{inner\_product}      &                    \\
a \emph{map}         & --- & ---                   & the reduce verbs           & \code{solve\_family} \\
                     &     &                       &                            & \code{frequency\_sweep} \\
a record by lifetime & all 3 & \code{Krylov} etc.  & \code{DofSet} etc.         & \code{EigState} etc. \\
an \code{\#extern}   & --- & ---                   & \code{assemble\_term}      & \code{eigen\_iterate} \\
                     &     &                       &                            & \code{time\_step\_op} \\
\bottomrule
\end{tabular}
\caption[The recurring shapes, tallied across the libraries]{The four libraries,
tallied by the shape of what they hold. Every operator is a \emph{fold}, a
\emph{map}, or one of three \code{\#extern} boundaries; every type is a record
distinguished by its lifetime. The table is mostly the same four rows repeating
across the columns---which is the chapter's whole point. There is no fifth kind of
thing. The library is small because the calculus is small: a few combinators, a few
records threaded by lifetime, and a marked boundary where an opaque kernel sits.}
\label{fig:samesahapes}
\end{figure}

\begin{keyidea}
Across all four libraries you keep meeting the \emph{same handful of shapes}: a
fold (\code{iterate\_while}, \code{linear\_combination}, \code{fold\_solve}), a map
(\code{solve\_family}, the reduce verbs), a record distinguished only by its
lifetime (\code{IoData}/\code{OpParams}/\code{SimState} and their single-library
cousins), and a marked \code{\#extern} boundary where an opaque kernel lives
(\code{assemble\_term}, \code{eigen\_iterate}, \code{time\_step\_op}). There is no
fifth kind of thing. \textbf{The library is small because the calculus is small}:
the entire synthesized implementation of a hundred-thousand-line simulator is a few
combinators, a few records, and three kernel callouts---the vocabulary of
\cref{part:framework}, rendered as code.
\end{keyidea}

This is the dividend the whole book has been compounding. \Cref{part:framework}
argued that the simulator is built from a small calculus; \cref{part:machinery} and
\cref{part:modalities} spent the bulk of the text showing that the argument holds
component by component. Here, with the spec rendered as a library, we can simply
\emph{count}: four shelves, and on them a fold, a map, a record, and an
\code{\#extern}, repeated. The next chapter, \cref{ch:composing}, opens the fifth
shelf---\emph{drivers}---and watches a single analysis composed end to end out of
the items we have just toured. Nothing new appears there either; a driver is these
four libraries, wired together.

\section{Worked examples}

\begin{workedexample}{Reading the \texttt{iterate\_while} def body line by line}{readwhile}
We read the synthesized \code{iterate\_while} of \cref{sec:iteration} one line at a
time, connecting each to a concept from \cref{ch:fold}. The body is:
\begin{lstlisting}[style=l4style]
iterate_while a cont step =
  if cont a                                              -- (1)
    then do
      { state: a', ...e } <- step a                      -- (2)
      { final_state, trajectory } <- iterate_while a' cont step  -- (3)
      pure { final_state, trajectory: [{ ...e }] ++ trajectory }  -- (4)
    else
      pure { final_state: a, trajectory: [] }            -- (5)
\end{lstlisting}

\emph{Line (1): the predicate fires first.} \code{if cont a} tests the predicate
\emph{before} any step runs---this is the \texttt{while} convention of
\cref{ch:fold}, Rule~1. If \code{cont a} is already false on the initial carry, the
body never executes and we fall straight to line~(5). A solver handed an
already-converged input does zero work, exactly as promised.

\emph{Line (2): one step, producing the next carry and the extras.}
\code{\{ state: a', ...e \} <- step a} runs the step once and \emph{destructures}
its output record: the \code{state} field becomes the next carry \code{a'}, and the
remaining fields---the open \code{...e}---are this step's extras. The \code{<-} is
the monadic bind: \code{step} runs in the \code{Solve} monad, so this line also
threads the \code{SimState} effect (the counter increment) under the hood, while
the carry \code{a'} stays a pure value. This is the carry/extras split of
\cref{ch:fold}: the carry threads forward, the extras peel off.

\emph{Line (3): tail-recurse on the new carry.}
\code{iterate\_while a' cont step} calls the combinator again with the
\emph{new} carry \code{a'} and the \emph{same} predicate and step. Because this call
is the last thing the \texttt{then} arm does before line~(4) assembles the result,
it is in tail position---the implementation realizes it as a loop, not a growing
stack. The predicate will be re-tested on \code{a'} at the top of this recursive
call; we may \emph{not} hoist that test out (\cref{ch:fold}'s non-law), or the loop
never stops.

\emph{Line (4): prepend this step's extras to the trajectory.}
\code{[\{ ...e \}] ++ trajectory} puts this step's extras at the \emph{front} of
whatever trajectory the recursion returned. Unrolled, the extras line up in
iteration order: step~1's first, step~2's next, and so on---the trajectory of
\cref{fig:unroll}. And because the trajectory is a \emph{value} the consumer may or
may not read, the demand-driven pruning of \cref{ch:fold} applies: if nothing
downstream reads \code{trajectory}, the \code{e}-producing work in line~(2) is
eliminated at every step, and what remains is the bare carry spine of
\code{iterate\_while\_pure}.

\emph{Line (5): the base case.} \code{\{ final\_state: a, trajectory: [] \}}
returns the current carry as the answer with an empty trajectory. This is the
\texttt{else} arm, reached the first time \code{cont a} is false. It is the formal
content of \cref{ch:fold}'s ``base case'' law: for \emph{any} step, an initial
carry that fails the predicate yields itself and nothing else.

Five lines: a gate, three lines of recursive work, one base case. Every iterative
algorithm in the simulator---CG, GMRES, time-stepping, the adaptive
loop---specializes \emph{this} body by supplying its own carry type, predicate, and
step. The loop is the same five lines every time.
\end{workedexample}

\begin{workedexample}{Reading the \texttt{ksp\_solve} cap}{kspsolve}
We read the \code{ksp\_solve} of \cref{sec:coordination}, connecting it to the monad
of \cref{ch:monad} and the Krylov machinery of \cref{ch:krylov}. The cap is:
\begin{lstlisting}[style=l4style]
ksp_solve op inp = execState (solve_loop op inp) (initial_state inp)  -- (1)
  where
    solve_loop op inp = do
      o <- restart_cycle op inp                            -- (2)
      unless (done o) (solve_loop op inp)                  -- (3)
    restart_cycle op inp = do
      s <- get                                             -- (4)
      let k0        = fresh_krylov op inp s                -- (5)
          (kn, _os) = iterate_while (krylov_step op k0) cont  -- (6)
      modify (\s -> s { x = s.x `plus` (kn.basis `applyBasis` kn.y) })  -- (7)
      pure (classify kn op s)                              -- (8)
\end{lstlisting}

\emph{Line (1): the \code{execState} discharge.} The cap's type is
$\code{OpParams}\to\code{Inputs}\to\code{SimState}$---a \emph{pure} function---yet
its body is monadic. \code{execState} bridges the two: it runs the
\code{Solve}-action \code{solve\_loop op inp} starting from a seeded state
\code{initial\_state inp}, threads the \code{SimState} through every \code{modify}
inside, and projects out the final state, dropping the action's unit return value.
This is exactly the monad-to-pure-function discharge of \cref{ch:monad}: the thin
\code{StateT} wrapper exists only inside the cap, and \code{execState} peels it
away at the boundary so the outside world sees a plain function.

\emph{Lines (2)--(3): the outer restart loop.} \code{solve\_loop} runs one
\code{restart\_cycle}, binds its \code{Outcome} to \code{o}, and---via
\code{unless (done o)}---recurses only if the outcome is not \code{Done}. This is
the GMRES \emph{restart} structure of \cref{ch:krylov}: an \emph{outer} loop around
the inner Krylov fold, deliberately not flattened (the two stop on different
conditions, the non-fusion non-law of \cref{ch:fold}). For CG and Chebyshev, which
do not restart, this outer loop runs exactly once.

\emph{Line (4): read the state.} \code{s <- get} reads the current \code{SimState}
out of the monad---the iterate, the counter, the residuals as they stand at the
start of this cycle.

\emph{Line (5): a fresh ephemeral workspace.} \code{fresh\_krylov op inp s} builds
the \code{Krylov} bundle of \cref{sec:iteration}---born here, at cycle entry, as a
\emph{plain value}, not a monadic effect. Its lifetime is exactly this one cycle;
it will be discarded when the cycle ends. This is the ephemeral stratum of
\cref{ch:strata} in action.

\emph{Line (6): the inner Krylov fold.} \code{iterate\_while (krylov\_step op k0) cont}
is the heart---the very combinator we read in \cref{wex:readwhile}, now folding the
\code{krylov\_step} kernel of \cref{sec:iteration} over the working state until the
convergence predicate \code{cont} fails. Every CG/GMRES iteration of
\cref{ch:krylov} happens inside this one call. The result \code{kn} is the final
Krylov bundle; the discarded \code{\_os} would be its trajectory, pruned here
because the cap does not read it.

\emph{Line (7): the single per-cycle \code{x} write.} \code{modify (\textbackslash s -> s \{ x = ... \})}
is the cap's \emph{only} write to the iterate. After the inner fold finishes, the
accumulated correction \code{kn.basis `applyBasis` kn.y} (the back-substituted
solution increment) is added to \code{x} once. Not per inner iteration---\emph{once
per restart cycle}. This is the deliberate write discipline of \cref{ch:monad}: the
externally-visible iterate changes at cycle boundaries, while the inner fold works
in the ephemeral \code{Krylov} space and never touches \code{SimState.x}.

\emph{Line (8): classify the outcome once.} \code{classify kn op s} inspects the
final bundle (its residual proxy, its inner-iteration count) and the state, and
returns \code{Done True} (converged), \code{Done False} (hit \code{max\_it}), or
\code{Continue} (restart). The classification happens \emph{once}, at the cycle
boundary---the \code{Outcome} sum of \cref{sec:coordination}, the
classify-once law of \cref{ch:monad}.

Read whole: \code{ksp\_solve} is an outer restart loop (lines 2--3) around an inner
Krylov fold (line 6), with one state read (4), one ephemeral workspace (5), one
iterate write per cycle (7), and one termination classification (8), all discharged
to a pure function by \code{execState} (1). Every Krylov solver in the book---CG,
GMRES, FGMRES---is this same cap; what changes is only the \code{krylov\_step}
kernel it folds.
\end{workedexample}

\summarysection
The synthesized library is five modules in topological order; this chapter toured
four of them. The \emph{types} library holds three cross-cutting records
distinguished by lifetime: \code{IoData} (parsed once, selects the driver via
\code{problem.type}), \code{OpParams} (fixed at build time, variant selectors
folded into surfaces), and \code{SimState} (the run-time state, uniform across all
Krylov solvers). The \emph{iteration} library holds the \code{iterate\_while}
combinator of \cref{ch:fold}---rendered here as five lines of code---plus the
\code{krylov\_step} kernel (Form~A branch-in-body, Form~B unrolled) and the
\code{chebyshev} smoother, with \code{Krylov}/\code{StepOutputs}/\code{PrevCarry} as
single-library types. The \emph{data-algebra} library holds
\code{linear\_combination} with its \code{where}-local BLAS arity leaves,
\code{inner\_product}/\code{dot}/\code{nrm2}, the matrix-free operator constructor,
the \code{fe\_assemble} fold with its libCEED \code{\#extern}, and the five reduce
verbs. The \emph{coordination} library holds the \code{Solve} monad, the
\code{ksp\_solve} cap (read line by line), the \code{eigsolve} black-box with its
SLEPc \code{\#extern}, and the three family combinators
\code{solve\_family}/\code{frequency\_sweep}/\code{fold\_solve}. Across all four,
the same handful of shapes recurs---a fold, a map, a record threaded by lifetime, a
marked \code{\#extern} boundary---because the library is small for the same reason
the calculus is small. The fifth library, \emph{drivers}, composes these four and
is \cref{ch:composing}.

\furtherreading
The implementation idioms rendered here---the state monad as a thin wrapper
discharged by \code{execState}, \code{where}-local helper clusters, the records
threaded by lifetime---are standard functional-programming practice; Peyton Jones's
account of Haskell's design~\citep{peytonjones2003} is the canonical reference for
the monadic and record machinery. The Krylov kernels the \emph{iteration} and
\emph{coordination} libraries fold (CG, GMRES, and the Arnoldi/Lanczos relatives
behind \code{eigsolve}) are developed in full by Saad~\citep{saad2003}. The libCEED
quadrature kernel behind \code{fe\_assemble}'s \code{\#extern} boundary---the
matrix-free element-local contraction that \code{mk\_matrix\_free\_operator} renders
inline---is documented by the CEED project~\citep{ceed2021}.

\exercisessection
\begin{enumerate}[nosep]
  \item For each of the three \emph{types} records, name the stratum it belongs to
  (\cref{ch:strata}) and one field whose value is fixed for the whole run. Which of
  the three is the \emph{only} one written during the solve?
  \item \code{SimState} has the same five fields for CG and for GMRES, yet the two
  solvers clearly differ. Where does the difference live? (Hint: which library, and
  which record, holds the search direction for CG versus the growing basis for
  GMRES?)
  \item In the \code{iterate\_while} body of \cref{wex:readwhile}, suppose a
  consumer reads only \code{final\_state} and never the \code{trajectory}. Which
  line's work is pruned, and what combinator does the def collapse to? (Cite the
  relevant law from \cref{ch:fold}.)
  \item The four BLAS leaves of \code{linear\_combination} are written as a
  \code{where}-block of one-liners. Write \code{axpbypcz} in terms of \code{axpby}
  and \code{scal} using the concatenation law, and explain why this makes them ``one
  combinator, four leaves'' rather than four separate operators (\cref{ch:combinators}).
  \item \code{nrm2} calls \code{inner\_product} but is \emph{not} a leaf of it,
  while \code{dot} \emph{is}. Explain the difference using the
  specialization-versus-consumer distinction of \cref{ch:combinators}, and say what
  the \code{abs} inside \code{nrm2} guards against (\cref{ch:bigidea}).
  \item Three operators across the four libraries render \code{\#extern}:
  \code{assemble\_term} (\emph{data-algebra}), and \code{eigen\_iterate} /
  \code{time\_step\_op} (\emph{coordination}). For each, name the opaque library
  that owns it and say in one sentence why the synthesized code marks a boundary
  there rather than rendering a body.
  \item Reading \code{ksp\_solve} (\cref{wex:kspsolve}): the cap writes
  \code{SimState.x} exactly once per restart cycle (line~7), not once per inner
  iteration. Where do the \emph{inner} iterations do their work instead, and why is
  it correct for the externally-visible iterate to change only at cycle boundaries?
  Relate your answer to the strata of \cref{ch:strata}.
  \item \code{solve\_family}, \code{frequency\_sweep}, and \code{fold\_solve} are
  each near one-liners. Which two are \code{map}s and which is a \code{fold}? For the
  two maps, what single difference distinguishes them (\cref{ch:combinators})? Which
  of the three forbids reordering its elements, and why?
\end{enumerate}

\chapter{Composing a Driver, End to End}
\label{ch:composing}

\chapterorient{Part~V disassembled the simulator into named machines and showed,
in a drawn trace (\cref{ch:fulltrace}), that one analysis is one composition of
them. Part~VI takes that picture seriously: if every analysis is a composition of
the same library functions, then the compositions are \emph{code}, and we can
write them down. This chapter does exactly that. We open the \code{drivers}
library---the topmost shelf of the synthesized library (\cref{ch:fivelibs})---and
read its six driver definitions, each a handful of lines over the four calculus
libraries beneath it. We render one driver, electrostatics, in full and walk it
line by line, naming the library every call comes from. Then we descend into the
single most important step in the book, the conjugate-gradient \code{krylov\_step},
and watch the abstract calculus of Part~III produce real, runnable-looking
algorithm code. The payoff is the chapter's thesis: every driver is a few lines of
\emph{map}, \emph{fold}, or \emph{single call} over \texttt{assemble}$\to$%
\texttt{reduce}, and the calculus we built really does synthesize working
algorithms.}

\section{The \texttt{drivers} library: six compositions over four}

We have, by now, the whole library. \Cref{ch:synthlib} laid out its
architecture---five modules---and \cref{ch:fivelibs} toured the four that hold the
\emph{vocabulary}: \code{types} (the records every solve threads), \code{iteration}
(the loop combinators of \cref{ch:fold} and the step kernels), \code{data-algebra}
(the assembly fold \code{fe\_assemble} and the reduction verbs \code{gram\_reduce},
\code{sparameter\_reduce}, and their siblings), and \code{coordination} (the
outer-driver shapes \code{solve\_family}, \code{frequency\_sweep},
\code{fold\_solve}, and the opaque \code{eigsolve} cap). This chapter is about the
\emph{fifth} module, the one that sits on top of all of them:

\begin{quote}\itshape
The \code{drivers} library holds the six entry-point analyses---electrostatic,
magnetostatic, driven, transient, eigenmode, boundary-mode---each written as a short
composition of the four libraries below it. It composes the vocabulary; it does not
add to it.
\end{quote}

\noindent That is the whole nature of the module, and it is worth stating sharply
before we read any code. A driver definition introduces no new operator, no new
combinator, no new record. Every name in its body has already been built and named
in an earlier chapter; the driver's only job is to \emph{wire those names together}
in the right shape. Reading the \code{drivers} library is therefore reading the
book's table of contents executed: each line points down at a machine you already
own (\cref{fig:lib-stack}).

\begin{figure}[t]
\centering
\begin{tikzpicture}[font=\footnotesize, >=Stealth]
  % top: the drivers module
  \node[stage, fill=simBgViolet, draw=simViolet, minimum width=92mm,
        minimum height=11mm, align=center] (drv) at (0,3.2)
    {\textbf{\code{drivers}}\ \scriptsize --- the 6 entry-point compositions
     (this chapter)};
  % the four calculus libraries below
  \node[stage, fill=simBgGold, draw=simGold, minimum width=20mm,
        minimum height=15mm, align=center] (da) at (-3.45,1.1)
    {\code{data-}\\\code{algebra}\\[1pt]\scriptsize assemble\\+ reduce};
  \node[stage, fill=simBgRust, draw=simRust, minimum width=20mm,
        minimum height=15mm, align=center] (co) at (-1.15,1.1)
    {\code{coord-}\\\code{ination}\\[1pt]\scriptsize the solve\\shapes};
  \node[stage, fill=simBgBlue, draw=simBlue, minimum width=20mm,
        minimum height=15mm, align=center] (it) at (1.15,1.1)
    {\code{iter-}\\\code{ation}\\[1pt]\scriptsize loop +\\step};
  \node[stage, fill=simBgGray, draw=simGray, minimum width=20mm,
        minimum height=15mm, align=center] (ty) at (3.45,1.1)
    {\code{types}\\[1pt]\scriptsize the\\records};
  % composition edges
  \foreach \n in {da,co,it,ty}{ \draw[depedge] (drv.south) -- (\n.north); }
  \node[font=\scriptsize\itshape, text=simGray] at (0,0.05)
    {(each library composes the ones to its right)};
\end{tikzpicture}
\caption[The library stack]{The synthesized library as a stack. The
\code{drivers} module (violet, this chapter) sits on top and \emph{composes} the
four calculus libraries of \cref{ch:fivelibs}: \code{data-algebra} (the assembly
fold and the reduction verbs), \code{coordination} (the four solve shapes plus the
opaque eigensolver cap), \code{iteration} (the loop combinators and step kernels),
and \code{types} (the records). A driver definition adds no new vocabulary; every
name in its body is built below it. Reading a driver is reading those four
libraries wired into the shape one analysis needs.}
\label{fig:lib-stack}
\end{figure}

\subsection{The unifying algebra, in one line each}

Here are all six drivers, each collapsed to the composition that defines it. Read
each right to left, the way function composition reads---\emph{assemble} on the
right, \emph{reduce} on the left, the solve shape in the middle:

\begin{lstlisting}[style=l4style]
electrostatic = gram_reduce(w=1)          . solve_family    . fe_assemble
magnetostatic = gram_reduce(w=1/(Ii*Ij))  . solve_family    . fe_assemble
driven        = sparameter_reduce         . frequency_sweep . fe_assemble x3
transient     = fold_solve                                  . fe_assemble x3
eigenmode     = map readout . eigsolve    . eig_pencil      . fe_assemble x3
boundary_mode = map readout . eigsolve    . eig_pencil      . fe_assemble . extract_2d
\end{lstlisting}

\noindent Six analyses, six one-liners, and one shape. Every driver is
\texttt{assemble} on the right, a \emph{solve shape} in the middle, and a
\texttt{reduce} on the left. What differs from row to row is small and named: how
many operators get assembled, which of the four \emph{coordination shapes} runs the
solve, and which reduction verb caps it. \Cref{ch:situating} taught us to read those
middle terms as a small grid of \emph{solve corners}; the \code{drivers} library is
that grid made executable. \Cref{fig:six-trees} draws the six as composition trees,
grouped by the corner each occupies.

\begin{figure}[p]
\centering
\begin{tikzpicture}[font=\scriptsize, >=Stealth,
    every node/.style={align=center}]
% ===================== GROUP 1: the maps =====================
\node[cornertag, anchor=west, font=\footnotesize\bfseries] at (-0.3,7.7)
  {\textsc{the maps} --- a family of independent solves (\cref{ch:situating}, corners 1--2)};

% electrostatic
\node[stage, fill=simBgGold, draw=simGold, text width=20mm] (e0) at (1.4,6.4)
  {\code{electrostatic}};
\node[opbox, text width=15mm] (e1) at (-0.4,5.2) {\code{fe\_assemble}\\\scriptsize K once};
\node[stage, fill=simBgBlue, draw=simBlue, text width=15mm] (e2) at (1.4,5.2)
  {\code{solve\_family}\\\scriptsize fixed-op map};
\node[opbox, text width=15mm] (e3) at (3.2,5.2) {\code{gram\_reduce}\\\scriptsize $w{=}1$};
\draw[depedge] (e0)--(e1); \draw[depedge] (e0)--(e2); \draw[depedge] (e0)--(e3);

% magnetostatic
\node[stage, fill=simBgGold, draw=simGold, text width=20mm] (m0) at (6.6,6.4)
  {\code{magnetostatic}};
\node[opbox, text width=15mm] (m1) at (4.8,5.2) {\code{fe\_assemble}\\\scriptsize K once};
\node[stage, fill=simBgBlue, draw=simBlue, text width=15mm] (m2) at (6.6,5.2)
  {\code{solve\_family}\\\scriptsize fixed-op map};
\node[opbox, text width=16mm] (m3) at (8.5,5.2) {\code{gram\_reduce}\\\scriptsize $w{=}\frac{1}{I_iI_j}$};
\draw[depedge] (m0)--(m1); \draw[depedge] (m0)--(m2); \draw[depedge] (m0)--(m3);

% driven
\node[stage, fill=simBgGold, draw=simGold, text width=18mm] (d0) at (11.4,6.4)
  {\code{driven}};
\node[opbox, text width=15mm] (d1) at (9.7,5.2) {\code{fe\_assemble}\\\scriptsize $\times3$ basis};
\node[stage, fill=simBgRust, draw=simRust, text width=16mm] (d2) at (11.4,5.2)
  {\code{frequency\_}\\\code{sweep}\\\scriptsize op-varying map};
\node[opbox, text width=16mm] (d3) at (13.3,5.2) {\code{sparameter\_}\\\code{reduce}};
\draw[depedge] (d0)--(d1); \draw[depedge] (d0)--(d2); \draw[depedge] (d0)--(d3);

% ===================== GROUP 2: the fold =====================
\node[cornertag, anchor=west, font=\footnotesize\bfseries] at (-0.3,3.9)
  {\textsc{the fold} --- a state-threaded chain (\cref{ch:situating}, corner 3)};
\node[stage, fill=simBgGold, draw=simGold, text width=18mm] (t0) at (2.0,2.7)
  {\code{transient}};
\node[opbox, text width=16mm] (t1) at (0.0,1.5) {\code{fe\_assemble}\\\scriptsize $\times3$ K,C,M};
\node[stage, fill=simBgGreen, draw=simGreen, text width=18mm] (t2) at (2.6,1.5)
  {\code{fold\_solve}\\\scriptsize state-threaded\\time march};
\draw[depedge] (t0)--(t1); \draw[depedge] (t0)--(t2);

% ===================== GROUP 3: the black box =====================
\node[cornertag, anchor=west, font=\footnotesize\bfseries] at (6.2,3.9)
  {\textsc{the black box} --- one opaque call (\cref{ch:situating}, corner 4)};

% eigenmode
\node[stage, fill=simBgGold, draw=simGold, text width=16mm] (g0) at (8.0,2.7)
  {\code{eigenmode}};
\node[opbox, text width=14mm] (g1) at (6.6,1.5) {\code{fe\_assemble}\\\scriptsize $\times3$ pencil};
\node[extern, text width=12mm] (g2) at (8.1,1.5) {\code{eigsolve}\\\scriptsize 1 call};
\node[opbox, text width=12mm] (g3) at (9.5,1.5) {\code{map}\\\code{readout}};
\draw[depedge] (g0)--(g1); \draw[depedge] (g0)--(g2); \draw[depedge] (g0)--(g3);

% boundary mode
\node[stage, fill=simBgGold, draw=simGold, text width=18mm] (b0) at (12.4,2.7)
  {\code{boundary\_mode}};
\node[opbox, text width=14mm] (b1) at (10.7,1.5) {\code{extract\_2d}\\\scriptsize submesh};
\node[opbox, text width=13mm] (b2) at (12.1,1.5) {\code{fe\_assemble}\\\scriptsize $+$ pencil};
\node[extern, text width=11mm] (b3) at (13.5,1.5) {\code{eigsolve}\\\scriptsize 1 call};
\draw[depedge] (b0)--(b1); \draw[depedge] (b0)--(b2); \draw[depedge] (b0)--(b3);

\end{tikzpicture}
\caption[The six drivers as composition trees]{The hero figure: all six drivers as
composition trees, grouped by the solve corner each occupies (\cref{ch:situating}).
\textsc{The maps} (top): electrostatic and magnetostatic run \code{solve\_family},
a fixed-operator map (corner~1); the driven sweep runs \code{frequency\_sweep}, an
operator-varying map (corner~2). \textsc{The fold} (lower left): transient threads
its state through \code{fold\_solve}, a sequential chain (corner~3). \textsc{The
black box} (lower right): eigenmode and boundary-mode hand the assembled pencil to a
single opaque \code{eigsolve} call (corner~4), boundary-mode prefacing it with a
2-D submesh extraction. Every tree has the same skeleton---an assemble on one side,
a reduce/readout on the other, a solve shape between---and differs only in which
shape and which reduction it picks. This grid \emph{is} the \code{drivers} library.}
\label{fig:six-trees}
\end{figure}

\subsection{The three coordination shapes are the deep unifier}

Look again at the middle column of the six one-liners and a smaller pattern
emerges---smaller even than the four corners of \cref{ch:situating}. Underneath the
corners there are really only \emph{three coordination shapes}, the three ways a
driver can run its solves (\cref{ch:combinators}):

\begin{description}[style=nextline, leftmargin=0pt, font=\normalfont\itshape]
  \item[map] a family of \emph{independent} solves, run in any order---no carry.
  Electrostatics and magnetostatics map over a fixed operator (\code{solve\_family});
  the driven sweep maps over a per-frequency operator (\code{frequency\_sweep}). The
  shape is the same \texttt{map}; the corners differ only in whether the operator is
  reused or rebuilt.
  \item[fold] a \emph{sequential} chain, each step consuming the state the last
  produced. Transient time-stepping is the canonical fold (\code{fold\_solve}); so is
  the adaptive mesh loop that the lifecycle wraps around all of it.
  \item[single call] no loop \emph{we} write at all---the assembled pencil goes to an
  opaque library routine (\code{eigsolve}) and the modes come back. Eigenmode and
  boundary-mode are the two black-box drivers.
\end{description}

\noindent Map, fold, single call. That is the entire vocabulary of \emph{how a
driver coordinates its solves}, and every analysis Palace runs is one of the three
wrapped around \texttt{assemble} and capped by a \texttt{reduce}. \Cref{fig:three-shapes}
sets them side by side; the rest of the chapter is, in effect, three readings of
this one figure---one map (electrostatics, in full), one fold (the CG step that
drives every solve inside any of them), and a tour of the single call.

\begin{figure}[t]
\centering
\begin{tikzpicture}[font=\footnotesize, >=Stealth]
  % ---- MAP ----
  \begin{scope}[shift={(0,0)}]
    \node[font=\small\bfseries, text=simBlue] at (1.3,2.2) {\code{map}};
    \node[opbox, minimum width=10mm] (mk) at (1.3,1.55) {op (fixed)};
    \foreach \j/\y in {1/0.7,2/0.0,3/-0.7}{
      \node[data, minimum width=7mm] (mb\j) at (0.3,\y) {$\vb_{\j}$};
      \node[product, minimum width=8mm] (mv\j) at (2.3,\y) {$\vx_{\j}$};
      \draw[flow, shorten >=1pt] (mb\j) -- (mv\j);
    }
    \node[font=\scriptsize\itshape, text=simBlue, align=center] at (1.3,-1.4)
      {independent;\\reorderable, parallel};
  \end{scope}
  \draw[simGray!45, thick] (3.6,-1.7) -- (3.6,2.4);
  % ---- FOLD ----
  \begin{scope}[shift={(4.6,0)}]
    \node[font=\small\bfseries, text=simGreen] at (1.6,2.2) {\code{fold}};
    \node[data, minimum width=8mm] (s0) at (-0.1,0.7) {$\vs_0$};
    \node[opbox, minimum width=7mm] (st1) at (0.9,0.7) {step};
    \node[data, minimum width=8mm] (s1) at (1.85,0.7) {$\vs_1$};
    \node[opbox, minimum width=7mm] (st2) at (2.8,0.7) {step};
    \node[data, minimum width=8mm] (s2) at (3.75,0.7) {$\vs_2$};
    \draw[flow] (s0)--(st1); \draw[flow] (st1)--(s1);
    \draw[flow] (s1)--(st2); \draw[flow] (st2)--(s2);
    \node[font=\scriptsize\itshape, text=simGreen, align=center] at (1.8,-1.4)
      {sequential;\\each step needs the last};
  \end{scope}
  \draw[simGray!45, thick] (9.0,-1.7) -- (9.0,2.4);
  % ---- SINGLE CALL ----
  \begin{scope}[shift={(9.9,0)}]
    \node[font=\small\bfseries, text=simViolet] at (1.5,2.2) {single call};
    \node[data, minimum width=12mm] (pen) at (0.2,0.7) {pencil};
    \node[extern, minimum width=15mm] (eig) at (2.2,0.7) {\code{eigsolve}};
    \node[product, minimum width=12mm] (modes) at (2.2,-0.4) {modes};
    \draw[flow] (pen) -- (eig);
    \draw[flow] (eig) -- (modes);
    \node[font=\scriptsize\itshape, text=simViolet, align=center] at (1.5,-1.4)
      {opaque; no loop\\we write};
  \end{scope}
\end{tikzpicture}
\caption[The three coordination shapes]{The three ways a driver coordinates its
solves---the deep unifier beneath the six drivers. A \code{map} (blue) runs a family
of independent solves over an operator; nothing is carried, so the solves are
reorderable and parallel (electrostatic, magnetostatic, driven). A \code{fold}
(green) threads a state through a sequential chain, each step consuming the last
(transient, and the adaptive mesh loop). A \code{single call} (violet) hands the
assembled pencil to an opaque library routine and reads the modes back (eigenmode,
boundary-mode). Every driver in the library is one of these three, wrapped around
\code{fe\_assemble} and capped by a reduction.}
\label{fig:three-shapes}
\end{figure}

\begin{keyidea}
Every driver is a few lines: a \texttt{map}, a \texttt{fold}, or a single opaque
call, wrapped around \texttt{assemble} and capped by \texttt{reduce}. The six
analyses differ only in \emph{which} of the three coordination shapes runs the solve
and \emph{which} reduction verb caps it---and the calculus of Part~III produces real
algorithm code that fills those shapes in. The \code{drivers} library is the grid of
\cref{ch:situating} written out as composition; reading it is reading the book's
machinery wired together.
\end{keyidea}

\subsection{Per-driver config records are projections, not new shapes}
\label{sec:config-projections}

Before we render a driver in full, one piece of plumbing deserves naming, because it
recurs in every signature. Each driver's body opens by reading fields off a
configuration value typed \code{ElectrostaticConfig}, \code{DrivenConfig}, and so on.
These look like six different record types. They are not. There is exactly
\emph{one} configuration record---the \code{IoData} parsed from the user's file
(\cref{ch:records}, \cref{ch:fulltrace})---and the per-driver config types are
\emph{projection-views} of it: type aliases that name which sub-records a given
driver reads.

\begin{lstlisting}[style=l4style]
-- the per-driver config records are VIEWS of the one IoData, not new shapes
type ElectrostaticConfig = IoData   -- reads: model, domains.epsilon, terminals, linear solver
type DrivenConfig        = IoData   -- reads: model, domains.{nu,sigma,eps}, ports, omega-sweep
type EigenmodeConfig     = IoData   -- reads: model, domains.{nu,sigma,eps}, n-modes + target
-- ... and likewise for magnetostatic, transient, boundary_mode
\end{lstlisting}

\noindent Writing \code{ElectrostaticConfig} rather than the bare \code{IoData} in a
driver's type is documentation, not a new data shape: it says ``this driver consults
\emph{this} slice of the one config.'' The accessors a driver uses---%
\lstinline[style=l4style]|permittivity cfg|, \lstinline[style=l4style]|terminal_sources cfg|,
\lstinline[style=l4style]|ports cfg|---are the projection's utility API, trivial
field selectors into \code{IoData}. Keeping these as views and not as separate
records is the record-definition discipline of \cref{ch:records} applied at the top
of the stack: the data shape is defined once, in \code{types}, and the drivers
\emph{view} it rather than redefining it.

\begin{notationbox}
The shared \code{IoData} is why the lifecycle (\cref{ch:fulltrace}) can parse one
file and dispatch to any of the six drivers: the parse produces a single
\code{IoData}, and \emph{dispatch} reads one field of it---%
\lstinline[style=l4style]|problemType cfg|---to choose which driver to run. The
chosen driver then projects the slice it needs. One parse, one record, six views, a
one-field switch: that is the entire shared front end, and it is why
\S\ref{sec:lifecycle} can write the spine root as a dispatch over the six driver
defs.
\end{notationbox}

\section{One driver, rendered and walked: \texttt{electrostatic}}
\label{sec:es-rendered}

We now render a complete driver as synthesized code and read every line of it. We
choose \code{electrostatic}---the cleanest of the six, the anchor of all of Part~V
(\cref{ch:electrostatics})---because its shape is the bare skeleton: assemble one
operator, map a family of solves over it, Gram-reduce the family to the answer. Once
we have walked this one, the other five are short variations, and we contrast the
richest of them, \code{driven}, in \S\ref{sec:driven-contrast}.

\subsection{The whole driver}

Here is the synthesized \code{electrostatic} definition in full---the rendered form
of the composition $\code{gram\_reduce}\circ\code{solve\_family}\circ
\code{fe\_assemble}$ that \cref{ch:electrostatics} anchored Part~V on:

\begin{lstlisting}[style=l4style]
-- # Arguments
--   cfg : ElectrostaticConfig   -- the IoData electrostatic view: H1 space, epsilon,
--                                  terminal sources, linear solver config
-- # Returns
--   CapacitanceMatrix           -- the n x n Maxwell capacitance matrix (the product)
electrostatic :: ElectrostaticConfig -> CapacitanceMatrix
electrostatic cfg =
  let space = h1_space cfg                                  -- (0) the H1 FE space
      k     = fe_assemble space [ diffusion (permittivity cfg) ]   -- (1) assemble K ONCE
      rhss  = [ excitation cfg idx | idx <- terminal_sources cfg ] -- (2) per-terminal RHS family
      vs    = solve_family k rhss                           -- (3) fixed-operator solve map
  in  gram_reduce k vs (\i j -> 1)                          -- (4) C_ij = V_j^T K V_i, w = 1
\end{lstlisting}

\noindent Five lines of body. That is a complete electromagnetic analysis---the same
one that, on a real mesh, fills a capacitance matrix for a chip interconnect. The
brevity is the point: the driver is \emph{nothing but} the composition. All the
weight lives in the four library calls, each built in an earlier chapter. We walk
them in order, naming the library each comes from (\cref{fig:es-dataflow}).

\subsection{Line by line, naming the library}
\label{sec:es-linewalk}

\paragraph{(0) \code{space = h1\_space cfg}---from \code{types}/\code{data-algebra}.}
The first line builds the finite-element space the analysis lives in: the scalar
nodal $\Hone$ space (\cref{ch:electrostatics}), read off the mesh and polynomial
order in \code{cfg}. It is construction-stratum data (\cref{ch:strata})---fixed the
instant the mesh is built, read-only thereafter. Everything downstream is expressed
against this space.

\paragraph{(1) \code{k = fe\_assemble space [\,diffusion (permittivity cfg)\,]}---from
\code{data-algebra}.} The second line assembles the operator, and it is the
load-bearing line of the whole fixed-operator corner. \code{fe\_assemble} is the
assembly fold of \cref{ch:assembly}: it walks a list of weak-form terms and sums one
contribution per term into a sparse operator. For electrostatics the list has
\emph{exactly one} term---the $\varepsilon$-weighted gradient,
\code{diffusion (permittivity cfg)}, the discrete form of
$\int_\dom\varepsilon\Grad u\cdot\Grad v\dV$ (\cref{ch:electrostatics},
eq.~\eqref{eq:es-bilinear-ch39})---so the fold runs once and returns the stiffness
operator $\Kop$. The word that matters is the one the type makes invisible:
\emph{once}. This line sits \emph{outside} the solve map below it, so $\Kop$ is built
a single time and reused for every terminal. Hoisting it here, structurally, is the
fixed-operator corner (\cref{ch:electrostatics}, and the pitfall it warned of).

\paragraph{(2) \code{rhss = [\,excitation cfg idx | idx <- terminal\_sources cfg\,]}---from
the config view.} The third line builds the family of right-hand sides, one per
terminal. The list comprehension ranges over \code{terminal\_sources cfg}---the
conductors named in the config---and for each builds the vector
\code{excitation cfg idx} that encodes ``conductor \code{idx} at one volt, the rest
grounded'' (the eliminated boundary data, \cref{ch:electrostatics}). This is the
\emph{method of unit excitations} (\cref{ch:electrostatics}) rendered as a
comprehension: $n$ conductors, $n$ right-hand sides, no solving yet.

\paragraph{(3) \code{vs = solve\_family k rhss}---from \code{coordination}.} The
fourth line runs the solves, and it is the coordination shape of the driver. The
\code{solve\_family} combinator (\cref{ch:combinators}) is the \emph{fixed-operator
map}: it captures the operator \code{k} once and maps a solver over the RHS family,
\lstinline[style=l4style]|map (ksp_solve k) rhss|. Because \code{k} is a free
variable of the mapped function, the \texttt{map} closes over it once and reuses it
for every element---the hoist made structural by the type, so no implementation can
rebuild the operator inside the loop. Each element solve \code{ksp\_solve k b\_j} is
one preconditioned conjugate-gradient solve of the SPD system $\Kop V_j=\vb_j$, and
\code{ksp\_solve} itself bottoms out in the \code{krylov\_step} of
\S\ref{sec:cg-flagship}---the step we render in full next. The result \code{vs} is
the solution family $[V_1,\dots,V_n]$.

\paragraph{(4) \code{gram\_reduce k vs (\textbackslash i j -> 1)}---from
\code{data-algebra}.} The final line reduces the family to the answer.
\code{gram\_reduce} (\cref{ch:reductions}) is the symmetric Gram verb: it pairs the
family with itself through the operator \code{k}, computing
$C_{ij}=V_j^{\!\top}\Kop V_i$ over the upper triangle and mirroring the lower (since
$\Kop$ is symmetric). The weight closure \lstinline[style=l4style]|\i j -> 1| is the
unit-voltage specialization that makes this reduction \emph{capacitance}
specifically; magnetostatics will pass a different weight to the same verb
(\S\ref{sec:driven-contrast}). Notice that \code{k} appears \emph{twice}---in the
solve (line 3) and again in the reduce (line 4): the same operator that drives the
solves also weights the Gram. That double use is the back-edge of the anchor diagram
(\cref{ch:electrostatics}, \cref{fig:es-pipeline}), now visible in the code as a
shared \code{let}-binding.

\begin{figure}[t]
\centering
\begin{tikzpicture}[node distance=7mm and 11mm, font=\footnotesize, >=Stealth]
  \node[data, align=center, text width=17mm] (cfg)
    {\code{cfg}\\[1pt]\scriptsize \code{IoData} view};
  \node[stage, right=of cfg, align=center, text width=20mm, fill=simBgGold, draw=simGold]
    (asm) {\code{fe\_assemble}\\[1pt]\scriptsize (1) $\Kop$ once};
  \node[stage, right=of asm, align=center, text width=22mm, fill=simBgBlue, draw=simBlue]
    (solve) {\code{solve\_family}\\[1pt]\scriptsize (3) fixed-op map};
  \node[stage, right=of solve, align=center, text width=20mm, fill=simBgGold, draw=simGold]
    (red) {\code{gram\_reduce}\\[1pt]\scriptsize (4) $w{=}1$};
  \node[product, right=of red, align=center, text width=15mm] (prod)
    {$\mat{C}$\\[1pt]\scriptsize capacitance};
  \draw[flow] (cfg) -- (asm);
  \draw[flow] (asm) -- node[above,font=\scriptsize]{$\Kop$} (solve);
  % the RHS family entering the solve from below
  \node[data, below=8mm of solve, align=center, text width=20mm] (rhs)
    {\code{rhss} $=[\vb_j]$\\[1pt]\scriptsize (2) per terminal};
  \draw[flow] (rhs.north) -- (solve.south);
  \draw[flow] (solve) -- node[above,font=\scriptsize]{$\{V_j\}$} (red);
  \draw[flow] (red) -- (prod);
  % the shared K back-edge
  \draw[depedge, shorten >=1pt] (asm.south) ++(0,-1mm)
    .. controls +(0,-12mm) and +(0,-12mm) .. (red.south)
    node[midway, below, font=\scriptsize, text=simBlue]
      {the \emph{same} $\Kop$ weights the Gram (\code{k} used twice)};
  % library tags
  \node[cornertag, above=0.5mm of asm]   {\code{data-algebra}};
  \node[cornertag, above=0.5mm of solve] {\code{coordination}};
  \node[cornertag, above=0.5mm of red]   {\code{data-algebra}};
\end{tikzpicture}
\caption[The \texttt{electrostatic} driver's data flow]{The data flow of the
rendered \code{electrostatic} driver (\S\ref{sec:es-rendered}), each stage tagged by
the library its call comes from. The config view \code{cfg} feeds
\code{fe\_assemble} (\code{data-algebra}), which builds the stiffness $\Kop$
\emph{once}; the per-terminal RHS family \code{rhss} enters from below;
\code{solve\_family} (\code{coordination}) maps the fixed-operator solve over the
family, yielding $\{V_j\}$; and \code{gram\_reduce} (\code{data-algebra}) pairs that
family against $\Kop$ into the capacitance matrix $\mat{C}$. The blue back-edge is
the same operator $\Kop$ consumed twice---by the solve and again by the reduce---which
is the \code{k} that appears in both line~(3) and line~(4) of the code.}
\label{fig:es-dataflow}
\end{figure}

\subsection{What the rendering shows that the diagram could not}

The trace of \cref{ch:fulltrace} \emph{drew} this pipeline; the rendering
\emph{types} it. Three things become concrete only in the code. First, the operators
flow as ordinary \code{let}-bound values---$\Kop$ is a value named \code{k}, passed
to two later lines---so ``the same operator weights the Gram'' is not a diagram
annotation but a variable used twice. Second, the per-terminal family is a list
comprehension, so ``$n$ solves'' is a structural consequence of
\code{terminal\_sources cfg} having $n$ elements, not a loop bound written by hand.
Third, the whole driver has the type \code{ElectrostaticConfig -> CapacitanceMatrix}:
config in, product out, no monad, no mutation, no mention of the iteration machinery
that the solve hides. The driver is pure all the way across; the only place the
state monad of \cref{ch:monad} appears is buried four levels down, inside
\code{ksp\_solve}'s \code{krylov\_step}---which is exactly where we go next.

\subsection{The richest contrast: \texttt{driven} and its per-\texorpdfstring{$\omega$}{omega} rebuild}
\label{sec:driven-contrast}

Set \code{driven} beside \code{electrostatic} and the family resemblance is total
except for one move. We render it more briefly, to surface the single difference:

\begin{lstlisting}[style=l4style]
driven :: DrivenConfig -> FrequencyResponse
driven cfg =
  let space  = nd_space cfg                                       -- H(curl) space
      fam    = { K  = fe_assemble space [ curl_curl (reluctivity cfg) ]  -- (1) assemble
               , C  = fe_assemble space [ damping cfg ]                  --     the basis
               , M  = fe_assemble space [ mass (permittivity cfg) ] }    --     ONCE (x3)
      omegas = sample_frequencies cfg                             -- the swept band
      es     = frequency_sweep fam omegas    -- (2) operator-VARYING map: A(w) rebuilt INSIDE
  in  sparameter_reduce (ports cfg) es       -- (3) port-projection -> S(omega)
\end{lstlisting}

\noindent The skeleton is identical: assemble, map, reduce. But three cells have
moved, exactly as \cref{ch:driven} promised. The \emph{assemble} now folds
\emph{three} term lists into a fixed \emph{basis} $\{\Kop,\Cop,\Mop\}$ rather than
one operator. The \emph{reduce} swaps \code{gram\_reduce} for
\code{sparameter\_reduce}, the port-projection (not a Gram). And the load-bearing
swap is the middle: \code{solve\_family} becomes \code{frequency\_sweep}. Both are
\texttt{map}s---independent solves, reorderable and parallel---but \code{solve\_family}
captures \emph{one fixed operator} outside the loop, while \code{frequency\_sweep}
rebuilds a \emph{fresh} operator $\mat{A}(\omega)=\Kop+i\omega\Cop-\omega^2\Mop$
\emph{inside} the loop, once per frequency, from the fixed basis. The rebuild lives in
\code{frequency\_sweep}'s body (\code{coordination}); the driver only composes it.
\Cref{fig:driven-dataflow} draws the per-$\omega$ rebuild loop.

\begin{figure}[t]
\centering
\begin{tikzpicture}[font=\footnotesize, >=Stealth, node distance=7mm and 10mm]
  \node[data, align=center, text width=15mm] (cfg) {\code{cfg}};
  \node[stage, right=of cfg, align=center, text width=22mm, fill=simBgGold, draw=simGold]
    (asm) {\code{fe\_assemble}$^{\times3}$\\[1pt]\scriptsize basis $\{\Kop,\Cop,\Mop\}$ once};
  \node[stage, right=of asm, align=center, text width=24mm, fill=simBgRust, draw=simRust]
    (sweep) {\code{frequency\_sweep}\\[1pt]\scriptsize per $\omega$:};
  \node[stage, right=of sweep, align=center, text width=20mm, fill=simBgGold, draw=simGold]
    (red) {\code{sparameter\_}\\\code{reduce}};
  \node[product, right=of red, text width=14mm, align=center] (S) {$S(\omega)$};
  \draw[flow] (cfg) -- (asm);
  \draw[flow] (asm) -- node[above,font=\scriptsize]{\code{fam}} (sweep);
  \draw[flow] (sweep) -- node[above,font=\scriptsize]{$[\vE_k]$} (red);
  \draw[flow] (red) -- (S);
  % the per-omega inner loop body
  \node[opbox, below=12mm of sweep, fill=white, align=center, text width=44mm] (inner)
    {\textbf{for each $\omega_k$:}\ form $\mat{A}(\omega_k)=\Kop{+}i\omega_k\Cop{-}\omega_k^2\Mop$\\
     $\to$ \code{ksp\_solve} (GMRES) $\to \vE_k$};
  \draw[refedge] (sweep.south) -- (inner.north);
  % the rebuild loop arrow
  \draw[backflow] (inner.east) .. controls +(1.4,0) and +(1.4,0) ..
    (inner.north east) node[midway, right=1mm, font=\scriptsize, text=simRust, align=center]
    {operator\\rebuilt\\\emph{inside}};
  \node[cornertag, above=0.5mm of asm]   {\code{data-algebra}};
  \node[cornertag, above=0.5mm of sweep] {\code{coordination}};
\end{tikzpicture}
\caption[The \texttt{driven} driver's per-\texorpdfstring{$\omega$}{omega} rebuild]{%
The \code{driven} driver's data flow, contrasting \code{electrostatic}
(\cref{fig:es-dataflow}). The \emph{basis} $\{\Kop,\Cop,\Mop\}$ is assembled once
(three \code{fe\_assemble} folds). But the system operator is \emph{not} fixed:
\code{frequency\_sweep} (the rust box) maps over the swept band, and its per-$\omega$
body (drawn below) \emph{rebuilds} $\mat{A}(\omega_k)$ from the basis and solves it
by GMRES---the operator is re-formed and re-captured \emph{inside} the loop (the rust
back-edge), the exact negation of the electrostatic hoist. The solved fields
$[\vE_k]$ feed \code{sparameter\_reduce}, the port-projection that yields $S(\omega)$.
One coordination shape (\code{solve\_family}) swapped for another
(\code{frequency\_sweep}); the rest of the skeleton is unchanged.}
\label{fig:driven-dataflow}
\end{figure}

\begin{notationbox}
\textbf{Magnetostatics: the same verb, a different weight.} The closest sibling of
all is \code{magnetostatic}. Its rendered body is \code{electrostatic} with the
$\Hone$ space changed to $\Hcurl$, the diffusion term changed to curl--curl, and
\emph{one closure changed}: the Gram weight becomes
\lstinline[style=l4style]|\i j -> 1 / (is!!i * is!!j)|, the current normalization
$1/(I_iI_j)$, where \code{is = currents cfg}. The solve corner is identical
(\code{solve\_family}, fixed-operator map); the reduction verb is identical
(\code{gram\_reduce}); only the weight closure differs. ``One verb, two products''
(\cref{ch:electrostatics}) is, in the rendering, one function \code{gram\_reduce}
called with two different weight lambdas.
\end{notationbox}

\section{The flagship: \texttt{krylov\_step} and CG as code}
\label{sec:cg-flagship}

We have seen a driver compose its solves; now we open one solve and watch the
abstract calculus produce a real algorithm. The \code{ksp\_solve} that
\code{solve\_family} maps---and that lives, in some form, inside \emph{every} driver
above---is at bottom a single step kernel, \code{krylov\_step}, folded by the
iteration combinator of \cref{ch:fold}. This is the showcase of the book: the place
where high-order combinators (\cref{ch:fold}), the state monad (\cref{ch:monad}),
and the strata of \cref{ch:strata} converge to synthesize the conjugate-gradient
method (\cref{ch:cg}) as runnable-looking code.

\subsection{Form A: the branch-in-body step}

The first rendering of \code{krylov\_step} is the direct one---a single step body
that does everything, with the first-iteration special case handled by a branch
inside it. It is the cleanest illustration of the monadic discipline of
\cref{ch:monad}: a sea of pure \code{let}s and exactly one \code{modify}.

\begin{lstlisting}[style=l4style]
-- Form A: branch-in-body. One step of CG, embedded in the Solve monad.
krylov_step :: OpParams -> Krylov -> (SimState -> Solve { krylov: Krylov, outputs: StepOutputs })
krylov_step op K = \s -> do
  -- ===== OUTSIDE the monad: pure let-bindings, the actual numerical work =====
  let Ap     = apply_linop op.A K.p     in   -- operator apply        (pure)
  let alpha  = K.beta / dot K.p Ap      in   -- step length           (pure dot)
  let x'     = axpy alpha K.p s.x       in   -- new iterate           (pure axpy)
  let r'     = axpy (negate alpha) Ap K.r in -- new residual          (pure axpy)
  let rdot'  = dot r' r'                in   -- residual energy       (pure dot)
  let beta'  = rdot' / K.beta           in   -- direction coefficient (pure)
  let p'     = axpy beta' K.p r'        in   -- new search direction   (pure axpy)
  let K'     = { ...K, r: r', p: p', beta: rdot' }  -- fresh workspace (plain value)
  let outputs = { residual_norm: sqrt (abs rdot') } -- demand-prunable readout
  -- ===== INSIDE the monad: the ONE state write =====
  modify (\s -> s { it: s.it + 1, x: x', converged: rdot' < op.tol })
  pure { krylov: K', outputs }
\end{lstlisting}

\noindent Count the monadic operations: there is exactly \emph{one}, the closing
\code{modify}. Everything above it---the operator apply \code{Ap}, the dot products
\code{alpha} and \code{rdot'}, the new iterate \code{x'}, residual \code{r'}, and
direction \code{p'}---is a pure \code{let}-binding, the algorithm's actual arithmetic
(\cref{ch:monad}). The three strata of \cref{ch:strata} are visible by eye:
\code{op} is read-only construction-stratum config (the operator, the tolerance);
\code{K} is the ephemeral per-restart workspace, threaded as a \emph{plain value} and
returned fresh; and the \code{SimState} \code{s} is the persistent run-time state,
touched only by the lone \code{modify}. The step \emph{returns} its readout in
\code{outputs} as a value, never logging it---so when no consumer reads the residual
history, the \code{residual\_norm} computation is pruned away (\cref{ch:fold}).

\subsection{Form B: first-iteration unrolling into bootstrap and steady}

Form~A carries a hidden cost. The conjugate-gradient recurrence builds each search
direction from the current residual \emph{and} a coefficient $\beta$ computed from
the \emph{previous} residual (\cref{ch:cg}). On the very first iteration there is no
previous residual, so the search direction is just $\vp\gets\vr$. Form~A would
handle that with an \lstinline[style=l4style]|if it == 0| branch inside the step---a
conditional that runs (and is mispredicted) on every iteration but is meaningful on
only the first.

The framework removes the branch by the \emph{first-iteration unrolling} rotation of
\cref{ch:fold}: split the step into a one-time \code{cg\_first\_step} that
manufactures the first $\beta$, followed by a branch-free \code{cg\_steady\_step}
that always has a real previous $\beta$ to work with, threaded by the combinator as a
\emph{closure argument} rather than a state field. This is
\code{iterate\_while\_with\_prev} (\cref{ch:fold}) in its defining use.
\Cref{fig:formAB} draws the rotation.

\begin{figure}[t]
\centering
\begin{tikzpicture}[font=\footnotesize, >=Stealth]
  % ===== LEFT: Form A, branch in body =====
  \begin{scope}
    \node[font=\small\bfseries, text=simRust] at (1.6,2.7) {Form A: branch in body};
    \node[stage, minimum width=34mm, minimum height=20mm, fill=simBgRust, draw=simRust,
          align=center] (stepA) at (1.6,1.0) {};
    \node[font=\scriptsize, anchor=north] at ($(stepA.north)+(0,-1mm)$) {\code{krylov\_step}};
    \node[diamond, draw=simRust, aspect=1.7, inner sep=1pt, font=\scriptsize,
          fill=white] (br) at (1.6,1.05) {\code{if it==0}};
    \node[font=\scriptsize\itshape, text=simRust, below=0.5mm of br] {tested every pass};
    \draw[backflow] (stepA.east) .. controls +(0.8,0) and +(0.8,0) .. (stepA.north east)
      node[midway, right=1mm, font=\scriptsize, text=simRust] {loop};
    \node[font=\scriptsize\itshape, text=simGray, align=center] at (1.6,-0.55)
      {one step, one\\extra slot \code{beta\_prev}};
  \end{scope}
  % rotation arrow
  \draw[depedge, line width=1pt] (3.7,1.0) -- node[above, font=\scriptsize, text=simBlue]
    {unroll} node[below, font=\scriptsize, text=simBlue] {iter.~1} (5.1,1.0);
  % ===== RIGHT: Form B, bootstrap + steady =====
  \begin{scope}[shift={(5.6,0)}]
    \node[font=\small\bfseries, text=simGreen] at (2.4,2.7) {Form B: bootstrap then steady};
    \node[opbox, minimum width=24mm, fill=simBgBlue, draw=simBlue, align=center]
      (boot) at (1.1,1.5) {\code{cg\_first\_step}\\\scriptsize runs once; makes $\beta_0$};
    \node[stage, minimum width=26mm, fill=simBgGreen, draw=simGreen, align=center]
      (steady) at (3.4,0.5) {\code{cg\_steady\_step}\\\scriptsize branch-free};
    \draw[flow] (boot.south) .. controls +(0,-0.5) and +(-1.2,0.4) ..
      node[below left=-1mm, font=\scriptsize, text=simBlue] {$\beta_0$} (steady.west);
    \draw[backflow] (steady.east) .. controls +(0.9,0) and +(0.9,0) .. (steady.north east)
      node[midway, right=1mm, font=\scriptsize, text=simGreen] {loop};
    \node[font=\scriptsize\itshape, text=simGreen, align=center] at (2.4,-0.7)
      {$\beta$ threaded as a CLOSURE arg\\(\code{iterate\_while\_with\_prev});\\carry one slot lighter};
  \end{scope}
\end{tikzpicture}
\caption[The Form A \texorpdfstring{$\to$}{->} Form B rotation]{The
first-iteration-unrolling rotation (\cref{ch:fold}) applied to \code{krylov\_step}.
\emph{Left, Form A}: a single step body carries an \lstinline[style=l4style]|if it==0|
branch (the $\beta_{k-1}$ has no value on iteration~0), tested on every pass though
meaningful only on the first, and a \code{beta\_prev} slot in the carry. \emph{Right,
Form B}: the loop is unrolled once into a \code{cg\_first\_step} (runs exactly once,
manufactures the first $\beta_0$) and a \emph{branch-free} \code{cg\_steady\_step};
the previous $\beta$ is threaded by \code{iterate\_while\_with\_prev} as a
\emph{closure argument}, so the steady carry is one slot lighter and the steady step
has no conditional. Same algorithm, rotated for a cleaner steady state.}
\label{fig:formAB}
\end{figure}

Here are the two unrolled steps and the driver that folds them, rendered in full.
This is the synthesized conjugate-gradient method:

\begin{lstlisting}[style=l4style]
-- the v0.5 CgState: ONE scalar lighter than Form A -- beta_prev is gone, threaded
-- instead as the PrevCarry closure parameter of iterate_while_with_prev.
type CgState = { x: Tensor[$S], r: Tensor[$S], p: Tensor[$S],
                 beta: Scalar, it: Int, converged: Bool }

-- the FIRST step: it == 0, so p <- r unconditionally (the Form-A branch, hoisted out)
cg_first_step :: LinOp[$S] -> Scalar -> CgState -> { state: CgState, residual_norm: Scalar }
cg_first_step opA eps s =
  let p'    = s.r                  in   -- it == 0  =>  p <- r   (no branch: this IS the first step)
  let Ap    = apply_linop opA p'   in
  let alpha = s.beta / dot Ap p'   in
  let x'    = axpy alpha p' s.x    in
  let r'    = axpy (negate alpha) Ap s.r in
  let beta' = dot r' r'            in
  let res'  = sqrt (abs beta')     in
  { state: { x: x', r: r', p: p', beta: beta', it: 1, converged: res' < eps },
    residual_norm: res' }

-- the STEADY step: branch-free; beta_prev is the closure carry (NOT a state field)
cg_steady_step :: LinOp[$S] -> Scalar -> Scalar -> CgState -> { state: CgState, residual_norm: Scalar }
cg_steady_step opA eps beta_prev s =
  let p'    = axpby 1.0 s.r (s.beta / beta_prev) s.p in   -- p <- r + (beta/beta_prev) p
  let Ap    = apply_linop opA p'   in
  let alpha = s.beta / dot Ap p'   in
  let x'    = axpy alpha p' s.x    in
  let r'    = axpy (negate alpha) Ap s.r in
  let beta' = dot r' r'            in
  let res'  = sqrt (abs beta')     in
  { state: { x: x', r: r', p: p', beta: beta', it: s.it + 1, converged: res' < eps },
    residual_norm: res' }

-- the DRIVER: first step, then fold cg_steady_step with iterate_while_with_prev,
-- threading the prior step's beta as the next step's beta_prev (closure, not a field).
cg_solve :: CgConfig -> LinOp[$S] -> Tensor[$S] -> Tensor[$S] -> Bool
         -> { final_state: CgState, residual_history: [Scalar] }
cg_solve config opA b x0 initial_guess =
  let { state: s0, initial_res } = cg_init opA b x0 initial_guess in
  let eps = max (config.rel_tol * initial_res) config.abs_tol in
  if sqrt (abs s0.beta) < eps                              -- already converged? do zero work
    then { final_state: { ...s0, converged: True }, residual_history: [] }
    else
      let { state: s1, residual_norm: res1 } = cg_first_step opA eps s0 in   -- the bootstrap
      if s1.converged || s1.it >= config.max_it
        then { final_state: s1, residual_history: [res1] }
        else
          let { final_state, trajectory } =
            iterate_while_with_prev
              (\_ -> pure { state: s1, prev: s0.beta })   -- bootstrap: seed (s1, beta_prev = s0.beta)
              s1                                          -- initial carry
              (\(s, beta_prev) ->                         -- steady body: (carry, prev)
                 let r = cg_steady_step opA eps beta_prev s in
                 pure { state: r.state, prev: s.beta, residual_norm: r.residual_norm })
              (\s -> s.it < config.max_it && not s.converged)   -- predicate: PURE on the carry
          in { final_state, residual_history: [res1] ++ trajectory.map (\t -> t.residual_norm) }
\end{lstlisting}

\noindent Read \code{cg\_solve} as the chapter's punchline. It is the
conjugate-gradient method of \cref{ch:cg}---a real, convergent linear solver---written
as a few lines over the combinator of \cref{ch:fold}. The \texttt{while} convention
of \cref{ch:fold} is visible in the very first \code{if}: an already-converged input
does \emph{zero} work. The bootstrap-then-steady split is
\code{iterate\_while\_with\_prev}. The predicate
\lstinline[style=l4style]|\s -> s.it < config.max_it && not s.converged| is pure on
the carry, reading the convergence flag and counter the step folded into the
state---never the residual the step just produced (\cref{ch:fold}'s predicate rule).
And the \code{residual\_history} is assembled from the trajectory the combinator
collected---demand-pruned away when nobody reads it. Everything those chapters built
abstractly is here, doing work.

\begin{keyidea}
The calculus produces real algorithm code. \code{cg\_solve} is the conjugate-gradient
method (\cref{ch:cg}) written as a short composition over
\code{iterate\_while\_with\_prev} (\cref{ch:fold}): a bootstrap step manufactures the
first $\beta$, a branch-free steady step folds the recurrence, the previous $\beta$
rides as a closure argument (\cref{ch:strata}), the predicate is pure on the carry,
the sole mutation is one \code{modify} of the counter (\cref{ch:monad}), and the
residual history is a demand-pruned trajectory. No part of it is pseudocode dressed
up---it is the abstract combinators of Part~III, instantiated, running a real solver.
\end{keyidea}

\section{Two worked examples}
\label{sec:worked}

The first worked example renders and walks the \code{electrostatic} driver against a
problem small enough to do by hand---from \code{IoData} projection to the capacitance
matrix---naming each library call as it fires. The second traces one CG step through
the synthesized \code{cg\_steady\_step}, tying every line to the conjugate-gradient
math of \cref{ch:cg}.

\begin{workedexample}{The \texttt{electrostatic} driver, run by hand}{es-driver}
We run the rendered driver of \S\ref{sec:es-rendered} on the two-conductor capacitor
of \cref{ch:electrostatics} (\cref{wex:es-twocond}), discretized to a single interior
degree of freedom, watching each library call do its job.

\medskip
\emph{(0) Project the config.} The parsed \code{IoData} has
\lstinline[style=l4style]|problemType cfg = ELECTROSTATIC|, so the lifecycle
(\S\ref{sec:lifecycle}) dispatched to \code{electrostatic} and handed it the config.
The driver reads its electrostatic view: \code{permittivity cfg} $=\varepsilon$, and
\lstinline[style=l4style]|terminal_sources cfg = [1, 2]| (two conductors). Line~(0)
builds \code{space = h1\_space cfg}, the scalar $\Hone$ space; on this coarse mesh it
has one free interior node.

\emph{(1) \code{fe\_assemble}, once.} Line~(1) folds the one-term list
\code{[\,diffusion}~$\varepsilon$\code{\,]} over the mesh. With $g=2\varepsilon A/L$
the half-gap conductance, the assembled-then-constrained stiffness is the scalar
$\Kop=2g$ (\cref{wex:es-twocond}). This is \code{k}, built a single time.

\emph{(2) The RHS family.} Line~(2)'s comprehension ranges over the two terminals.
\code{excitation cfg 1} encodes ``conductor~1 at one volt, conductor~2 grounded,''
giving $b_1=g$; \code{excitation cfg 2} gives $b_2=g$ by symmetry. So
\lstinline[style=l4style]|rhss = [g, g]|.

\emph{(3) \code{solve\_family k rhss}.} Line~(3) maps \code{ksp\_solve k} over the
family, reusing the \emph{one} operator \code{k} $=2g$ for both. Each solve is
$\Kop u=b$: $u_1=b_1/\Kop=g/(2g)=\tfrac12$ and $u_2=\tfrac12$. Reconstructing the full
nodal fields gives $V_1=(1,\tfrac12,0)$ and $V_2=(0,\tfrac12,1)$. So
\lstinline[style=l4style]|vs = [V1, V2]|---the solution family, produced by a map that
captured \code{k} once.

\emph{(4) \code{gram\_reduce k vs (\textbackslash i j -> 1)}.} Line~(4) pairs the
family against \code{k} with unit weight. The upper triangle is
$C_{11}=V_1^{\!\top}\Kop V_1=g/2$, $C_{22}=g/2$, and
$C_{12}=V_2^{\!\top}\Kop V_1=-g/2$; the verb mirrors $C_{21}=C_{12}$ (because $\Kop$
is symmetric). The returned product is
\[
  \mat{C}=\frac{g}{2}\begin{pmatrix}1&-1\\-1&1\end{pmatrix},
  \qquad -C_{12}=\frac{g}{2}=\frac{\varepsilon A}{L},
\]
the textbook parallel-plate capacitance. \tcblower
\emph{What the rendering added.} The arithmetic is identical to
\cref{wex:es-twocond}---but now we watched it run as \emph{code}. The same \code{k}
appeared in line~(3) (the solve) and line~(4) (the reduce): one value, two
consumers, the back-edge of \cref{fig:es-dataflow} made literal. The ``two solves''
were the two elements of \code{rhss}, not a hand-written loop. And the whole thing had
type \code{ElectrostaticConfig -> CapacitanceMatrix}: \code{IoData} in, $\mat{C}$ out,
nothing mutable crossing the boundary. On a real mesh only line~(3) changes
character---the two scalar divisions become two million-dimensional CG solves, each
the \code{cg\_solve} of \S\ref{sec:cg-flagship}---but the composition is exactly the
five lines we ran.
\end{workedexample}

\begin{workedexample}{One CG step through \texttt{cg\_steady\_step}}{cg-step}
We trace a single iteration of the synthesized \code{cg\_steady\_step}
(\S\ref{sec:cg-flagship}) and tie each line to the conjugate-gradient math of
\cref{ch:cg}. Take the tiny SPD system
\[
  \mat{A}=\begin{pmatrix}4&1\\1&3\end{pmatrix},\quad \vb=\begin{pmatrix}1\\2\end{pmatrix},
\]
and suppose the loop has reached a steady state with carry
\lstinline[style=l4style]|s = { x, r, p, beta, it, ... }| where
$\vr=(\tfrac{1}{2},\tfrac{1}{2})^{\!\top}$, $\vp=(1,0)^{\!\top}$, the carry's
$\code{beta}=\inner{\vr}{\vr}=\tfrac12$, and the closure carry
$\code{beta\_prev}=2$ (the previous step's residual energy). We run
\code{cg\_steady\_step opA eps beta\_prev s} line by line.

\medskip
\emph{Line 1 --- the direction update \code{p' = axpby 1.0 s.r (s.beta/beta\_prev) s.p}.}
The coefficient is $\beta/\beta_{\text{prev}}=\tfrac12/2=\tfrac14$. This is the CG
$\beta_k=\inner{\vr_k}{\vr_k}/\inner{\vr_{k-1}}{\vr_{k-1}}$ of \cref{ch:cg}, and the
threaded \code{beta\_prev} is exactly $\inner{\vr_{k-1}}{\vr_{k-1}}$---the
\emph{previous} residual energy, supplied as a closure argument, not read from a state
field. The new direction is
$\vp'=\vr+\tfrac14\vp=(\tfrac12+\tfrac14,\ \tfrac12)=(\tfrac34,\tfrac12)^{\!\top}$.

\emph{Line 2 --- the apply \code{Ap = apply\_linop opA p'}.} One operator application,
$\mat{A}\vp'=(4\cdot\tfrac34+\tfrac12,\ \tfrac34+3\cdot\tfrac12)=(\tfrac72,\tfrac94)^{\!\top}$.
This is the single matrix--vector product that is the cost of a CG step (\cref{ch:cg});
in the driver it is one \code{krylov\_step}'s operator apply, the matrix-free leaf of
\cref{ch:fulltrace}.

\emph{Line 3 --- the step length \code{alpha = s.beta / dot Ap p'}.} The denominator
is $\inner{\mat{A}\vp'}{\vp'}=\tfrac72\cdot\tfrac34+\tfrac94\cdot\tfrac12=\tfrac{21}{8}+\tfrac98=\tfrac{30}{8}=\tfrac{15}{4}$,
so $\alpha=\beta/\inner{\mat{A}\vp'}{\vp'}=\tfrac12/\tfrac{15}{4}=\tfrac{2}{15}$.
This is the CG line-search $\alpha_k=\inner{\vr_k}{\vr_k}/\inner{\vp_k}{\mat{A}\vp_k}$
(\cref{ch:cg}), the step that minimizes the energy norm along $\vp'$.

\emph{Lines 4--5 --- the updates \code{x' = axpy alpha p' s.x},
\code{r' = axpy (-alpha) Ap s.r}.} The iterate moves $\vx'=\vx+\alpha\vp'$ and the
residual updates $\vr'=\vr-\alpha\mat{A}\vp'$---the two \code{axpy} fused-multiply-adds
that advance CG (\cref{ch:cg}). Numerically
$\vr'=(\tfrac12,\tfrac12)-\tfrac{2}{15}(\tfrac72,\tfrac94)=(\tfrac12-\tfrac{7}{15},\ \tfrac12-\tfrac{3}{10})=(\tfrac{1}{30},\tfrac{1}{5})^{\!\top}$.

\emph{Lines 6--7 --- the new residual energy and convergence.}
$\code{beta'}=\inner{\vr'}{\vr'}=(\tfrac{1}{30})^2+(\tfrac15)^2\approx0.0411$, and
$\code{res'}=\sqrt{\code{beta'}}\approx0.203$---a smaller residual than the
$\sqrt{\code{beta}}=\sqrt{\tfrac12}\approx0.707$ we started with. The step writes
\code{converged: res' < eps} into the state.

\tcblower
\emph{Reading the trace as the calculus.} Every CG quantity of \cref{ch:cg} appears
exactly once and in the expected place: $\beta_k$ on line~1, the apply on line~2,
$\alpha_k$ on line~3, the iterate and residual updates on lines~4--5. The step is
\emph{branch-free}---there is no \lstinline[style=l4style]|if it==0|---because the
unrolling of \S\ref{sec:cg-flagship} moved that case into \code{cg\_first\_step}; and
the previous residual energy reached this step as the closure argument
\code{beta\_prev}, threaded by \code{iterate\_while\_with\_prev} (\cref{ch:fold}), not
stored in \code{CgState}. The one effect on the persistent \code{SimState}---the
counter bump and convergence flag---is the lone \code{modify} of Form~A
(\cref{ch:monad}); the iterate, residual, and direction are pure \code{let}s. This is
the conjugate-gradient method of \cref{ch:cg}, running, with each line a calculus
construct we built earlier in the book.
\end{workedexample}

\section{Writing a new analysis, and the spine that runs them}
\label{sec:lifecycle}

We close by stepping back from the six existing drivers to the shape they share, and
asking what it would take to write a seventh.

\subsection{How a new analysis is written}

Because a driver is nothing but a composition, adding an analysis is a matter of
\emph{three choices}, not a new program (\cref{fig:new-analysis}):

\begin{enumerate}[nosep]
  \item \textbf{Pick the corner.} Decide how the solves coordinate: a fixed-operator
  map (\code{solve\_family}), an operator-varying map (\code{frequency\_sweep}), a
  state-threaded fold (\code{fold\_solve}), or a single opaque call (\code{eigsolve}).
  This is the choice of coordination shape from \code{coordination}.
  \item \textbf{Pick the reduction.} Decide what physical product the solved family
  collapses to: a symmetric Gram (\code{gram\_reduce}), a port-projection
  (\code{sparameter\_reduce}), a per-mode table (\code{eigenfreq\_qfactor\_reduce}),
  an energy table (\code{domain\_energy\_reduce}). This is a verb from
  \code{data-algebra}.
  \item \textbf{Compose.} Wire \code{fe\_assemble} (the right term list) into the
  chosen corner into the chosen reduction. That composition \emph{is} the driver.
\end{enumerate}

\noindent No new combinator, no new record, no new loop. The expressive economy is
the whole argument of the book cashed out: the simulator is a small library, and an
analysis is a short composition over it. The four existing reduction verbs and four
coordination shapes already span the six analyses Palace ships; a genuinely new
analysis would reuse them and differ only in its term list and its choice of corner
and reduction.

\begin{figure}[t]
\centering
\begin{tikzpicture}[font=\footnotesize, >=Stealth, node distance=5mm]
  \node[data, align=center, text width=20mm] (cfg) at (0,0) {term list +\\config view};
  \node[stage, fill=simBgGold, draw=simGold, right=10mm of cfg, text width=18mm, align=center]
    (asm) {\code{fe\_assemble}\\\scriptsize \code{data-algebra}};
  % the corner choice
  \node[stage, fill=simBgRust, draw=simRust, right=9mm of asm, text width=24mm, align=center]
    (corner) {\textbf{pick the corner}\\\scriptsize map / map / fold / call\\\code{coordination}};
  \node[stage, fill=simBgBlue, draw=simBlue, right=9mm of corner, text width=22mm, align=center]
    (red) {\textbf{pick the reduce}\\\scriptsize Gram / proj / table\\\code{data-algebra}};
  \node[product, right=7mm of red, text width=14mm, align=center] (prod) {new\\product};
  \draw[flow] (cfg)--(asm); \draw[flow] (asm)--(corner);
  \draw[flow] (corner)--(red); \draw[flow] (red)--(prod);
  \node[cornertag, below=1mm of corner] {choice 1};
  \node[cornertag, below=1mm of red] {choice 2};
  \node[font=\scriptsize\itshape, text=simGray] at (6.3,-1.5)
    {choice 3: compose the three into one driver def};
\end{tikzpicture}
\caption[Writing a new analysis]{Writing a new analysis is three choices, not a new
program. Start from a weak-form term list and a config view; assemble with
\code{fe\_assemble} (\code{data-algebra}); \textbf{pick the corner}---one of the four
coordination shapes (\code{coordination}); \textbf{pick the reduction}---one of the
verbs (\code{data-algebra}); and \textbf{compose} the three into a driver definition.
No new combinator or record is written; the new analysis reuses the same library and
differs only in its term list and its two choices. This is why the six drivers are
each a handful of lines.}
\label{fig:new-analysis}
\end{figure}

\subsection{The spine root: \texttt{lifecycle} dispatches over the six}

One definition sits above all six drivers and runs them: the \code{lifecycle} root
(\cref{ch:fulltrace}). It is the topologically last definition in the library---it
\emph{composes the driver defs themselves}---and it is short, because dispatch is one
field read:

\begin{lstlisting}[style=l4style]
lifecycle :: IoData -> Product
lifecycle cfg =
  let mesh0 = build_mesh cfg                          -- shared front end: load + partition mesh
      drv   = dispatch (problemType cfg)              -- the specialization seam: pick ONE driver
      step  = \m -> estimate_mark_refine cfg (drv cfg m)  -- run the driver, then mark/refine
  in  fold_solve_amr step mesh0 (refinement cfg)      -- adaptive estimate-mark-refine OUTER fold
  where
    dispatch ELECTROSTATIC = electrostatic            -- the six driver defs, by problem type
    dispatch MAGNETOSTATIC = magnetostatic
    dispatch DRIVEN        = driven
    dispatch TRANSIENT     = transient
    dispatch EIGENMODE     = eigenmode
    dispatch BOUNDARYMODE  = boundary_mode
\end{lstlisting}

\noindent Three things to read. First, \code{build\_mesh} and the parse that produced
\code{cfg} are the \emph{shared front end} (\cref{ch:fulltrace})---driver-agnostic,
run identically whichever analysis follows. Second, \code{dispatch (problemType cfg)}
is the \emph{specialization seam}: one field of the one \code{IoData} selects one of
the six driver defs above. Everything before the seam is shared; the seam picks the
one stage that differs. Third, the whole run is wrapped in \code{fold\_solve}---the
adaptive estimate--mark--refine loop is itself a fold (\cref{ch:fold}), whose carry is
the mesh, degenerating to a single solve when refinement is off. The spine root is the
\emph{fourth} use of the same coordination vocabulary: a fold wrapped around a
dispatch wrapped around the six compositions.

\begin{keyidea}
The whole simulator is one library and a one-field dispatch. \code{lifecycle} parses
\code{IoData}, builds the mesh (the shared front end), reads \code{problemType} to
select one of the six driver compositions (the specialization seam), and threads the
adaptive loop through \code{fold\_solve}. Above the six few-line drivers sits one
few-line root that runs any of them. From parse to product, the entire machine is the
five-module library of \cref{ch:fivelibs} composed---no glue beyond the compositions
themselves.
\end{keyidea}

\subsection{Where this leaves us}

We set out, in \cref{ch:fulltrace}, to \emph{draw} the simulator as one composition;
in this chapter we \emph{wrote} it. The six drivers are six short compositions over
four libraries; the conjugate-gradient solver inside every one of them is a few lines
over the combinator of \cref{ch:fold}; and a seventh analysis would be three choices
and a composition. The synthesized library is not a metaphor for the simulator---it is
the simulator, with its compositions made into code. What remains is to step back one
final time and see the whole machine---library, drivers, and lifecycle---as a single
object, which is the closing chapter, \cref{ch:wholemachine}.

\summarysection
The \code{drivers} library is the top module of the synthesized library
(\cref{ch:fivelibs}): it holds the six entry-point analyses, each a short composition
over the four calculus libraries (\code{types}, \code{iteration}, \code{data-algebra},
\code{coordination}) and adding no new vocabulary. The six are one-liners in a single
algebra---\code{electrostatic} $=$ \code{gram\_reduce}$(w{=}1)\circ$%
\code{solve\_family}$\circ$\code{fe\_assemble}; \code{magnetostatic} the same with the
weight $1/(I_iI_j)$; \code{driven} $=$ \code{sparameter\_reduce}$\circ$%
\code{frequency\_sweep}$\circ$\code{fe\_assemble}$^{\times3}$; \code{transient} $=$
\code{fold\_solve}$\circ$\code{fe\_assemble}$^{\times3}$; \code{eigenmode} and
\code{boundary\_mode} $=$ \code{map readout}$\circ$\code{eigsolve}$\circ$%
\code{eig\_pencil}$\circ\cdots$. Beneath the four solve corners of \cref{ch:situating}
lie just three coordination shapes---\texttt{map} (independent solves), \texttt{fold}
(sequential chain), and \texttt{single call} (opaque eigensolve)---and every driver is
one of the three wrapped around \texttt{assemble} and capped by \texttt{reduce}. The
per-driver config types are projection-views of the one \code{IoData}, not new shapes.
We rendered \code{electrostatic} in full and walked its five lines, naming the library
each call comes from---\code{fe\_assemble} assembles $\Kop$ once (\code{data-algebra});
a comprehension builds the per-terminal RHS family; \code{solve\_family} maps the
fixed-operator solve (\code{coordination}); \code{gram\_reduce} pairs the family
against the same $\Kop$ (\code{data-algebra})---and contrasted \code{driven}, whose only
real difference is rebuilding $\mat{A}(\omega)$ \emph{inside} \code{frequency\_sweep}.
The flagship is \code{krylov\_step}/CG: Form~A (branch-in-body, one \code{modify} amid
pure \code{let}s) rotated by first-iteration unrolling into Form~B
(\code{cg\_first\_step} $+$ branch-free \code{cg\_steady\_step}), folded by
\code{iterate\_while\_with\_prev} in \code{cg\_solve}---the conjugate-gradient method of
\cref{ch:cg} as runnable code, with $\beta$ threaded as a closure argument
(\cref{ch:strata}), the predicate pure on the carry (\cref{ch:fold}), and the residual
history demand-pruned. A new analysis is three choices---pick the corner, pick the
reduction, compose---and the \code{lifecycle} root parses \code{IoData}, builds the
mesh, dispatches on one field to one of the six drivers, and wraps the adaptive loop in
\code{fold\_solve}. Every driver is a few lines of map/fold/single-call over
\texttt{assemble}$\to$\texttt{reduce}, and the abstract calculus of Part~III produces
real algorithm code.

\furtherreading
The conjugate-gradient method rendered here as \code{cg\_solve}, and the broader family
of Krylov solvers the other drivers invoke, are developed in full by
Saad~\citep{saad2003}, whose treatment of CG, its short recurrence, and the role of the
$\beta$ coefficient is the standard reference behind \cref{ch:cg} and this chapter's
flagship. The functional-programming machinery that makes the synthesized rendering
possible---the state monad threading \code{SimState}, the \code{do}-notation, the
higher-order combinators that \code{iterate\_while\_with\_prev} specializes---is
documented by Peyton Jones~\citep{peytonjones2003}, the source behind \cref{ch:monad}
and \cref{ch:fold}. The simulator whose six drivers we have rendered is Palace itself,
documented at the project~\citep{palace2023}; readers comparing the synthesized
compositions to the original C++ will find the driver dispatch and the per-analysis
solvers there, in the imperative form this book has spent six parts rotating into the
pure compositions of this chapter. The closing chapter, \cref{ch:wholemachine},
assembles library, drivers, and lifecycle into one final view of the whole machine.

\exercisessection
\begin{enumerate}[nosep]
  \item \textbf{Read the one-liners.} For each of the six driver one-liners (the list
  at the opening of the chapter), name (a) how many operators \code{fe\_assemble}
  builds, (b) which coordination shape runs the solve (\code{solve\_family} /
  \code{frequency\_sweep} / \code{fold\_solve} / \code{eigsolve}), and (c) which
  reduction caps it. Which two drivers share the \emph{same} coordination shape but
  differ only in a weight?
  \item \textbf{Walk the rendering.} Take the rendered \code{electrostatic} driver of
  \S\ref{sec:es-rendered} and, for each of its five body lines, name the library the
  call comes from and the one fact the line establishes. Which value is \code{let}-bound
  on one line and consumed on \emph{two} later lines, and why does that double use
  appear as a back-edge in \cref{fig:es-dataflow}?
  \item \textbf{The one moved call.} Comparing \code{electrostatic} and \code{driven},
  identify the single coordination call that swaps (\code{solve\_family} for
  \code{frequency\_sweep}) and explain, in terms of \emph{where} the operator capture
  sits relative to the loop, why one rebuilds the operator inside the map and the other
  does not. Cross-check against \cref{ch:driven}.
  \item \textbf{Magnetostatics as one closure.} Magnetostatics calls the \emph{same}
  \code{gram\_reduce} as electrostatics. Write the weight closure each passes, and
  explain why the inductance matrix and the capacitance matrix come from one verb with
  two lambdas rather than two reduction routines. (See the notation box in
  \S\ref{sec:driven-contrast}.)
  \item \textbf{Count the monad.} In Form~A of \code{krylov\_step}
  (\S\ref{sec:cg-flagship}), classify every line as ``in the monad'' or ``outside the
  monad'' by the rule of thumb of \cref{ch:monad}. How many are in the monad? Name the
  three strata (\cref{ch:strata}) the step touches and which variable carries each.
  \item \textbf{Why unroll.} Explain what the first-iteration unrolling
  (\cref{fig:formAB}) buys: what branch does \code{cg\_steady\_step} avoid that Form~A
  carries, where did that case go, and how is the previous $\beta$ supplied to the
  steady step if it is no longer a field of \code{CgState}? Which combinator threads it?
  \item \textbf{Trace a CG step.} Repeat \cref{wex:cg-step} for the \emph{next} step:
  use the carry it produced (the new $\vr'$, $\vp'$, and $\beta'=\inner{\vr'}{\vr'}$),
  take \code{beta\_prev} to be the previous step's $\beta=\tfrac12$, and run
  \code{cg\_steady\_step} line by line to compute the new direction, step length,
  iterate, and residual. Confirm the residual norm continues to fall, and name the CG
  formula each line implements.
  \item \textbf{Write the seventh analysis.} Suppose you wanted a driver that computes
  the \emph{energy} stored by each of several independent excitations (no cross terms).
  Using the three-choice recipe of \S\ref{sec:lifecycle}: which coordination shape would
  you pick, which reduction verb, and what term list would \code{fe\_assemble} take?
  Sketch the five-line driver by analogy with \code{electrostatic}, and say which one
  line differs from it.
\end{enumerate}

\chapter{The Whole Machine, and Where It Goes Next}
\label{ch:wholemachine}

\chapterorient{This is the last chapter of the main text, and its job is to stand
back far enough to see everything at once. We have built, named, and composed every
part: the field theory, the calculus, the numerical engines, the six analyses, and
finally the synthesized library that renders them all as code. Here we tie the
single thread that runs through them---the \code{lifecycle} root from which every
other piece of the book is reachable---and draw the whole machine on one page. Then
we close two loops. The first is the book's opening promise (\cref{ch:bigidea}):
``fields as tensors, steps as pure functions'' was never an aesthetic preference; it
was the form in which the algorithm could be handed, unchanged, to an accelerator.
The second is the book's thesis itself, now that you have watched it proven component
by component: a large physics simulator is the repeated composition of a small set of
mathematical building blocks. We end honestly---naming the frontiers this book did
\emph{not} cross---and warmly, with what you can now do that you could not before.}

\section{One root, and everything below it}
\label{sec:wm-root}

Across five parts we disassembled a real electromagnetic simulator and, in Part~VI,
reassembled it as a small library: a handful of modules---\code{types},
\code{iteration}, \code{data-algebra}, \code{coordination}, and \code{drivers}---each
holding a few of the named machines we built (\cref{ch:synthlib},
\cref{ch:fivelibs}). \Cref{ch:composing} composed individual drivers from those
modules. One question remains, and it is the question that turns a \emph{library}
into a \emph{program}: what is the single entry point that ties all of it together?
What is the \code{main} of the synthesized library?

It has a name. It is \code{lifecycle}, the topmost composition in the \code{drivers}
module, and it is short enough to state in one line:
\begin{equation}
  \code{lifecycle} \;=\; \code{fold\_solve}\,\bigl(\code{dispatch}\,(\code{problem\_type}\;\code{cfg})\bigr)
                        \;\circ\; \code{build\_mesh}.
  \label{eq:wm-lifecycle}
\end{equation}
Read right to left: build the mesh from the parsed configuration; dispatch on the
configuration's \code{problem.type} to one of the six drivers; and wrap the chosen
driver in the adaptive estimate--mark--refine fold of \cref{ch:amr}, which threads
the mesh through repeated solves and---when adaptive refinement is switched
off---degenerates to a single solve. That is the entire top of the machine. Every
other piece of this book hangs beneath it.

\begin{keyidea}
The synthesized library has exactly one root: \code{lifecycle}. It is the \code{main}
of the whole simulator, and \cref{eq:wm-lifecycle} is its body in full---build the
mesh, dispatch on the problem type, fold the driver through the adaptive loop. From
this single function, \emph{every} component the book built is reachable by following
composition edges: the drivers, the solve corners, the assembly fold, the
matrix-free applies, the records. To hold \code{lifecycle} in your head is to hold
the whole machine.
\end{keyidea}

\subsection{The root, rendered as code}
\label{sec:wm-root-code}

\Cref{eq:wm-lifecycle} is the skeleton; here is the synthesized body, the way the
\code{drivers} module actually writes it. It is worth reading slowly, because every
identifier in it is a chapter.

\begin{lstlisting}[style=l4style]
-- The synthesized library's main. inputs = the parsed config; output = the
-- driver-selected physical product (capacitance | inductance | S-params |
-- (f,Q) | trajectory | fields | modes).
lifecycle :: IoData -> Product
lifecycle cfg =
  let mesh0 = build_mesh cfg                    -- (1) load + partition + a-priori refine   [ch:functionspaces]
      drv   = dispatch (problem_type cfg)       -- (2) select ONE driver: the specialization seam  [ch:lifecycle]
      step  = \m -> estimate_mark_refine cfg (drv cfg m)  -- one adaptive iterate: solve, then refine
  in  fold_solve step mesh0 (refinement cfg)    -- (3) adaptive estimate-mark-refine FOLD  [ch:amr]
  where
    -- the dispatch over the six driver defs, by problem type
    dispatch Electrostatic = electrostatic      -- [ch:electrostatics]
    dispatch Magnetostatic = magnetostatic      -- [ch:magnetostatics]
    dispatch Driven        = driven             -- [ch:driven]
    dispatch Transient     = transient          -- [ch:transient]
    dispatch Eigenmode     = eigenmode          -- [ch:eigenmodes]
    dispatch BoundaryMode  = boundary_mode      -- [ch:waveports]
\end{lstlisting}

\noindent Three lines of body and a six-way \code{dispatch}. Line~(1) is the mesh
build of \cref{ch:functionspaces}---the only stage every analysis runs identically.
Line~(2) is the \emph{specialization seam} of \cref{ch:lifecycle}: a single
field, \code{problem\_type}, read off the configuration record, selects which of the
six drivers fills in the one stage that differs. Line~(3) is the adaptive loop of
\cref{ch:amr}, written as exactly the \code{fold\_solve} combinator of \cref{ch:fold}:
its threaded value is the mesh, its body solves and then refines, and when no
refinement is requested it runs that body once. Nothing else is at the top. The
machine's entire control structure---which we spent Part~V tracing by hand---is these
three lines.

\begin{notationbox}
\code{dispatch} is an ordinary function from a \code{ProblemType} tag to a driver. It
is the synthesized form of the \texttt{switch} that \cref{ch:lifecycle} called the
specialization seam: above it, everything is shared (parse, mesh, the fold); below
it, the six drivers diverge. Because \code{dispatch} returns a \emph{function} (a
driver), \cref{eq:wm-lifecycle}'s middle factor is higher-order---the closure-returning
shape the calculus uses throughout (\cref{ch:language}). The whole of the
simulator's variety is funneled through this one tag.
\end{notationbox}

\subsection{Reachability: the root sees the whole book}
\label{sec:wm-reach}

The claim that ``everything is reachable from \code{lifecycle}'' is not rhetorical.
The synthesized library is a dependency graph: each function names the functions it
composes, and those edges form a finite tree rooted at \code{lifecycle}. Follow the
edges and you visit, in order, every machine the book built (\cref{fig:wm-root}).
\code{lifecycle} composes \code{build\_mesh} and \code{dispatch}; \code{dispatch}
reaches each of the six drivers; each driver composes an \emph{assemble} stage
(\code{fe\_assemble}, \cref{ch:assembly}), a \emph{solve} stage (one of four
coordination combinators, \cref{ch:situating}), and a \emph{reduce} stage (one of
three reduction verbs, \cref{ch:reductions}); the solve stages reach the Krylov
methods (\cref{ch:cg}, \cref{ch:gmres}), the multigrid preconditioner
(\cref{ch:multigrid}), and the smoothers (\cref{ch:smoothers}); and \emph{every}
operator application in any of them bottoms out in the same matrix-free apply
(\cref{ch:matrixfree}). There is no part of the book that this tree does not touch,
and no node in the tree that is not a chapter you have read.

\begin{figure}[p]
\centering
\begin{tikzpicture}[font=\footnotesize, >=Stealth, node distance=5mm and 4mm]
  % --- the root ---
  \node[stage, fill=simBgViolet, draw=simViolet, very thick, text width=30mm]
    (root) {\code{lifecycle}\\[1pt]\scriptsize the library's \code{main}};
  % parse / mesh feeding in
  \node[data, left=10mm of root, text width=16mm] (cfg)
    {\code{IoData}\\\scriptsize parsed config};
  \draw[flow] (cfg) -- node[above,font=\scriptsize]{parse} (root);
  % the AMR fold wrapping
  \node[opbox, above=7mm of root, text width=30mm] (amr)
    {\code{fold\_solve}: AMR loop\\\scriptsize estimate--mark--refine (\cref{ch:amr})};
  \draw[refedge] (amr.south) -- (root.north);
  \node[opbox, left=5mm of amr, text width=18mm] (mesh)
    {\code{build\_mesh}\\\scriptsize\cref{ch:functionspaces}};
  \draw[refedge] (mesh.south) to[out=-90,in=150] (root.north west);
  % --- dispatch fan-out to six drivers ---
  \node[stage, fill=simBgGray, draw=simGray, below=8mm of root, text width=24mm]
    (disp) {\code{dispatch} (\code{problem\_type})\\\scriptsize the specialization seam};
  \draw[flow] (root) -- (disp);
  % six driver boxes in a row
  \node[stage, below=8mm of disp, xshift=-44mm, text width=15mm] (d1)
    {\code{electro-}\\\code{static}};
  \node[stage, right=of d1, text width=15mm] (d2) {\code{magneto-}\\\code{static}};
  \node[stage, fill=simBgRust, draw=simRust, right=of d2, text width=13mm] (d3) {\code{driven}};
  \node[stage, right=of d3, text width=14mm] (d4) {\code{transient}};
  \node[stage, right=of d4, text width=15mm] (d5) {\code{eigen-}\\\code{mode}};
  \node[stage, right=of d5, text width=15mm] (d6) {\code{boundary-}\\\code{mode}};
  \foreach \d in {d1,d2,d3,d4,d5,d6}{\draw[flow] (disp) -- (\d);}
  % chapter tags under drivers
  \node[cornertag, below=0.5mm of d1] {\cref{ch:electrostatics}};
  \node[cornertag, below=0.5mm of d3] {\cref{ch:driven}};
  \node[cornertag, below=0.5mm of d5] {\cref{ch:eigenmodes}};
  % --- the shared body each driver composes ---
  \node[stage, fill=simBgBlue, draw=simBlue, below=10mm of d3, xshift=-14mm, text width=20mm]
    (asm) {\textbf{assemble}\\\code{fe\_assemble}\\\scriptsize\cref{ch:assembly}};
  \node[stage, fill=simBgRust, draw=simRust, right=7mm of asm, text width=24mm]
    (solve) {\textbf{solve}\\\scriptsize 4 corners (\cref{ch:situating})\\Krylov + multigrid};
  \node[stage, fill=simBgGold, draw=simGold, right=7mm of solve, text width=22mm]
    (red) {\textbf{reduce}\\\scriptsize 3 shapes (\cref{ch:reductions})};
  \foreach \d in {d1,d2,d3,d4,d5,d6}{
    \draw[refedge, simGray!60] (\d.south) to[out=-90,in=90] (asm.north);}
  \draw[flow] (asm) -- (solve);
  \draw[flow] (solve) -- (red);
  % the solve stack reaching down
  \node[opbox, below=6mm of solve, text width=26mm] (stack)
    {Krylov (\cref{ch:cg},\,\cref{ch:gmres})\\$\cdot$ V-cycle (\cref{ch:multigrid})\\$\cdot$ smoother (\cref{ch:smoothers})};
  \draw[refedge] (solve.south) -- (stack.north);
  \node[extern, below=5mm of stack, text width=30mm] (mf)
    {\textbf{matrix-free apply} (\cref{ch:matrixfree})\\\scriptsize every operator application bottoms out here};
  \draw[refedge] (stack.south) -- (mf.north);
  % records thread to the side
  \node[opbox, fill=white, right=8mm of stack, text width=20mm] (rec)
    {records: \code{SimState},\\\code{IoData}\\\scriptsize\cref{ch:records}, monad \cref{ch:monad}};
  \draw[refedge] (rec.west) -- (stack.east);
\end{tikzpicture}
\caption[The \code{lifecycle} root and its reachable machine]{The synthesized
library rooted at \code{lifecycle}, its \code{main}. The parsed \code{IoData} enters
on the left; \code{lifecycle} builds the mesh (\cref{ch:functionspaces}), wraps
everything in the adaptive \code{fold\_solve} loop (\cref{ch:amr}), and dispatches on
\code{problem\_type} (the specialization seam, \cref{ch:lifecycle}) to one of the six
drivers. Every driver composes the same three-stage body---\code{fe\_assemble}
(\cref{ch:assembly}), one of four solve corners (\cref{ch:situating}), one of three
reduction shapes (\cref{ch:reductions})---and every solve reaches the same nested
Krylov/multigrid/smoother stack, which bottoms out, like everything else, in the
matrix-free apply (\cref{ch:matrixfree}). The threaded state lives in the records
(\cref{ch:records}) carried by the solve monad (\cref{ch:monad}). Follow any edge and
you reach a chapter; this single tree \emph{is} the book.}
\label{fig:wm-root}
\end{figure}

\section{The whole machine on one page}
\label{sec:wm-onepage}

\Cref{fig:wm-root} drew the root and its first few layers. We now draw the
\emph{entire} machine---root to leaves, every driver, every combinator, every record,
every external-kernel boundary---in a single capstone figure
(\cref{fig:wm-hero}). It is dense by design: the point is not to read each box in
isolation but to \emph{see that everything connects}, that the simulator which looked,
from \cref{ch:targets}, like six elaborate and separate programs is in fact one finite
tree with a handful of shared leaves.

\begin{figure}[p]
\centering
\begin{tikzpicture}[font=\scriptsize, >=Stealth,
    grow=down, level distance=11mm, sibling distance=0mm,
    edge from parent/.style={draw=simGray, semithick, -{Stealth[length=1.4mm]}},
    every node/.style={align=center},
    level 1/.style={sibling distance=46mm},
    level 2/.style={sibling distance=23mm},
    level 3/.style={sibling distance=23mm}]
  \node[stage, fill=simBgViolet, draw=simViolet, very thick, text width=34mm]
    {\code{lifecycle}\\\scriptsize the whole machine (\S\ref{sec:wm-root})}
    child {node[stage, text width=22mm] {\code{build\_mesh}\\\scriptsize\cref{ch:functionspaces}}}
    child {node[opbox, text width=24mm] {\code{fold\_solve} (AMR)\\\scriptsize\cref{ch:amr}, \cref{ch:fold}}}
    child {node[stage, fill=simBgGray, draw=simGray, text width=26mm] {\code{dispatch}\\\scriptsize 6 drivers (\cref{ch:situating})}
      child {node[stage, text width=20mm] {\textbf{assemble}\\\code{fe\_assemble}\\\scriptsize\cref{ch:assembly}, \cref{ch:weak}}
        child {node[extern, text width=20mm] {matrix-free\\quadrature\\\scriptsize\cref{ch:matrixfree}}}
      }
      child {node[stage, fill=simBgRust, draw=simRust, text width=22mm] {\textbf{solve}\\\scriptsize 4 corners\\\code{solve\_family} /\\\code{frequency\_sweep} /\\\code{fold\_solve} / \code{eigsolve}}
        child {node[stage, fill=simBgBlue, draw=simBlue, text width=20mm] {Krylov + V-cycle\\\scriptsize\cref{ch:gmres}, \cref{ch:multigrid}}
          child {node[stage, fill=simBgGreen, draw=simGreen, text width=18mm] {smoother\\\scriptsize\cref{ch:smoothers}}
            child {node[extern, text width=18mm] {matrix-free\\apply\\\scriptsize\cref{ch:matrixfree}}}
          }
        }
        child {node[extern, text width=18mm] {\code{eigsolve}\\\scriptsize SLEPc kernel\\\cref{ch:krylovschur}}}
      }
      child {node[stage, fill=simBgGold, draw=simGold, text width=22mm] {\textbf{reduce}\\\scriptsize 3 shapes\\\code{gram\_reduce} /\\projection / table\\\cref{ch:reductions}}}
    };
  % the threaded-records side note
  \node[opbox, fill=white, text width=26mm, anchor=west] (rec) at (5.2,-1.1)
    {\textbf{records} (\cref{ch:records})\\\scriptsize \code{IoData} (read-only),\\\code{SimState} (threaded by\\the solve monad, \cref{ch:monad})};
\end{tikzpicture}
\caption[The whole machine on one page]{The capstone hero figure: the entire
synthesized simulator as a single finite tree, rooted at \code{lifecycle} and
expanded down to its leaves. Reading top to bottom: \code{lifecycle}
(\S\ref{sec:wm-root}) builds the mesh, wraps the run in the adaptive
\code{fold\_solve} loop (\cref{ch:amr}), and dispatches to one of six drivers
(\cref{ch:situating}). Each driver is the three-stage body
\emph{assemble--solve--reduce}: the assembly fold \code{fe\_assemble}
(\cref{ch:assembly}); one of the \emph{four} solve corners
(\code{solve\_family}, \code{frequency\_sweep}, \code{fold\_solve},
\code{eigsolve}), whose iterative cases nest a Krylov/multigrid/smoother stack
(\cref{ch:gmres}--\cref{ch:smoothers}) and whose black-box case calls the
\code{eigsolve} kernel (\cref{ch:krylovschur}); and one of the \emph{three} reduction
shapes (\cref{ch:reductions}). Every operator application in the iterative branch
bottoms out in the matrix-free apply (\cref{ch:matrixfree}, shown twice---it is the
universal leaf). The run's mutable state is the records (\cref{ch:records}) threaded
by the solve monad (\cref{ch:monad}). The figure is dense on purpose: it shows, in
one view, that the whole book is one connected machine.}
\label{fig:wm-hero}
\end{figure}

Two features of \cref{fig:wm-hero} deserve a word, because they are the whole point.
First, the tree is \emph{shallow and finite}. There is no recursion of structure, no
open-ended plumbing; from the root to the deepest leaf is six or seven composition
steps, and the same leaves are shared across all six drivers. A simulator is a small
tree, not a sprawl. Second, the leaves \emph{repeat}. The matrix-free apply appears
under assembly and again under every level of the solve stack; the reduction verbs are
shared across drivers that look physically unrelated; the records and the monad thread
through everything. The economy is not an accident of this book's exposition---it is a
property of the simulator, which is exactly why the book could be written at the length
it was rather than six times over.

\begin{keyidea}
The whole simulator fits on one page because it \emph{is} one finite, shallow tree of
shared parts. From \code{lifecycle} down to the matrix-free leaf is a handful of
composition steps, and the deepest, most-reused leaves---the matrix-free apply, the
reduction verbs, the records---are shared across all six drivers. The variety the
reader met in \cref{ch:targets} lives almost entirely in \emph{which few boxes a
driver swaps}, not in the structure of the tree.
\end{keyidea}

\section{Why this shape: the implementation view as a backend target}
\label{sec:wm-backend}

We can now close the book's oldest promise. Back in \cref{ch:bigidea} we made a
choice that, at the time, looked like a matter of taste: describe a simulation not as
in-place mutation of arrays but as pure functions over immutable tensors---``fields as
tensors, steps as pure functions.'' We gave three reasons (an explicit dependency
graph, clean composition, and a passing remark about accelerators) and promised the
third was the destination. This is that destination.

The synthesized library of \cref{fig:wm-hero} is written in a calculus whose semantics
were chosen to \emph{match} an external tensor backend
(\cref{ch:matrixfree}). Its values are immutable tensors. Its operations are pure
functions: an operator apply takes an input tensor and returns an output tensor, with
no aliasing, no hidden global, and no ordering hazard to police. Its combinators---the
fold, the map, the solve corners---express coordination without ever naming a memory
cell. This is precisely the form a GPU, or any massively parallel backend, wants to
consume: ``here is an immutable input tensor; produce an output tensor,'' applied to
all entries at once (\cref{fig:wm-backend}). The matrix-free apply at the bottom of
the tree is a sequence of element-local tensor contractions
(\cref{ch:matrixfree})---a pure, whole-tensor transformation that lowers onto an
accelerator kernel without rewriting.

\begin{figure}[t]
\centering
\begin{tikzpicture}[font=\footnotesize, >=Stealth, node distance=8mm and 12mm]
  % left: the synthesized library (pure tensor code)
  \node[stage, fill=simBgViolet, draw=simViolet, text width=34mm, align=left] (lib)
    {\textbf{synthesized library}\\[2pt]\scriptsize
     \lstinline[style=l4style]|apply :: Tensor[N] -> Tensor[N]|\\
     \scriptsize immutable tensors,\\pure functions, combinators\\(\cref{ch:bigidea}, \cref{fig:wm-hero})};
  % middle: the lowering arrow
  \node[opbox, right=of lib, text width=22mm] (lower)
    {\textbf{lower}\\\scriptsize no rewrite:\\the form already matches};
  % right: the accelerator
  \node[extern, right=of lower, text width=30mm, minimum height=18mm] (gpu)
    {\textbf{GPU / tensor backend}\\[2pt]\scriptsize whole-tensor ops,\\no aliasing,\\all entries at once\\(\cref{ch:matrixfree})};
  \draw[flow] (lib) -- (lower);
  \draw[flow] (lower) -- (gpu);
  % the closed-loop callback to ch:bigidea
  \node[font=\scriptsize\itshape, text=simGray, below=3mm of lower, text width=44mm, align=center]
    {the promise of \cref{ch:bigidea}, kept:\\``fields as tensors, steps as pure functions''\\was always heading here};
  % grid of cells on the GPU to suggest parallelism
  \begin{scope}[shift={($(gpu.south)+(0,-1.0)$)}]
    \foreach \i in {0,1,2,3,4,5}{
      \fill[simBgViolet, draw=simViolet] (\i*0.42-1.25,0) rectangle (\i*0.42-0.9,0.32);}
    \node[font=\scriptsize, text=simViolet] at (0,-0.32) {one kernel, all cells at once};
  \end{scope}
\end{tikzpicture}
\caption[The implementation view as a backend-lowering target]{Why the pure picture
was the right one. The synthesized library (left) is written in tensor-in,
tensor-out pure functions---\lstinline[style=l4style]|apply :: Tensor[N] -> Tensor[N]|,
no mutation, no aliasing. That is, almost word for word, the form a GPU or other
tensor backend consumes (right): a whole-tensor operation applied to every entry at
once. So lowering the library onto an accelerator is not a translation---the
description \emph{already matches} the hardware's model. This closes the promise of
\cref{ch:bigidea}: ``fields as tensors, steps as pure functions'' was never an
aesthetic choice; it was the form in which the algorithm could be handed to an
accelerator and run unchanged.}
\label{fig:wm-backend}
\end{figure}

This is the deep reason the implementation view of Part~VI is worth having as a
distinct rendering at all. It is not merely ``the same spec, written as code.'' It is
the spec written in \emph{the form that runs on the hardware the whole effort
targets}. The book never produces a port; it produces this layered description. But the
top layer of that description---the synthesized library---is deliberately shaped so
that the final step, lowering it onto burn's device backends or any equivalent
tensor library, is the short one. The hard work of saying \emph{what the algorithm
is}, free of any one machine's cost model, is what the five parts before this one did.

\begin{keyidea}
The functional rendering matters because it is the \emph{backend-lowering target}. The
synthesized library's semantics---immutable tensors, pure functions, whole-tensor
combinators---were chosen to match an external GPU/tensor backend, so the algorithm can
be handed to an accelerator and run \emph{unchanged}. ``Fields as tensors, steps as
pure functions'' (\cref{ch:bigidea}) was never decoration; it was a commitment to the
one description from which the accelerator form is a short step rather than a rewrite.
\end{keyidea}

\subsection{A worked dispatch: the final \emph{everything-connects} trace}
\label{sec:wm-worked}

We have one worked example left, and it is the whole book in miniature: take the
\code{lifecycle} root, feed it one real configuration, and trace the single dispatch
all the way down until it bottoms out in the machinery of every part.

\begin{workedexample}{One dispatch, root to leaf}{wm-dispatch}
An engineer hands the simulator a configuration file for a capacitance extraction: a
mesh of two conductors, the permittivity $\varepsilon$ of the dielectric between them,
the terminals to drive, and the field \code{problem\_type:~Electrostatic}. We trace
\code{lifecycle cfg} from the top.

\medskip
\textbf{(1) Build the mesh.} \code{lifecycle} calls \code{build\_mesh cfg}
(\cref{ch:functionspaces}), which loads the geometry, partitions it, and returns the
\code{Mesh}---the first value the run \emph{produces} rather than reads. This stage is
identical for all six analyses.

\textbf{(2) Dispatch.} \code{lifecycle} reads \code{problem\_type cfg}, finds
\code{Electrostatic}, and \code{dispatch} returns the \code{electrostatic} driver
(\cref{ch:electrostatics}). One field, one branch, the entire choice of \emph{which}
analysis runs. Had the field read \code{Driven}, the same line would have returned the
\code{driven} driver of \cref{wex:trace-driven}; the root does not change.

\textbf{(3) The adaptive fold.} \code{lifecycle} wraps the chosen driver in
\code{fold\_solve} (\cref{ch:amr}, \cref{ch:fold}). The configuration requests no
refinement, so the fold runs its body \emph{once}---the straight-line lifecycle of
\cref{ch:lifecycle}. Everything below is one pass.

\textbf{(4) Inside the driver: assemble.} \code{electrostatic cfg} lays the nodal
$\Hone$ space (\cref{ch:functionspaces}) and assembles the stiffness operator $\Kop$
by the weak-form fold \code{fe\_assemble} (\cref{ch:assembly}, \cref{ch:weak})---one
term, the $\varepsilon$-weighted gradient---with the essential boundary degrees of
freedom eliminated (\cref{ch:bc}). $\Kop$ is symmetric positive-definite and built
\emph{once}.

\textbf{(5) Inside the driver: solve.} The driver maps the fixed-operator solve over
the per-terminal right-hand sides: \code{solve\_family} (corner~1 of
\cref{ch:situating}). Each solve is CG (\cref{ch:cg}, because $\Kop$ is SPD)
preconditioned by a multigrid V-cycle (\cref{ch:multigrid}) over a smoother
(\cref{ch:smoothers}), and \emph{every} operator application inside all of them is the
matrix-free apply (\cref{ch:matrixfree})---the same element-local tensor contraction
that the driven filter's GMRES bottomed out in. The solver's running state lives in the
\code{SimState} record (\cref{ch:records}) threaded by the solve monad
(\cref{ch:monad}).

\textbf{(6) Inside the driver: reduce.} The driver reduces the solved family by the
symmetric Gram \code{gram\_reduce} (\cref{ch:reductions}), $C_{ij}=V_j^{\!\top}\Kop
V_i$, pairing the fields with themselves through the energy operator to produce the
$n\times n$ capacitance matrix.

\textbf{(7) Output.} \code{lifecycle} writes the product and the run metadata
(\cref{ch:lifecycle}).

\medskip
Read the chapter tags down the trace: \cref{ch:functionspaces}, \cref{ch:lifecycle},
\cref{ch:amr}, \cref{ch:assembly}, \cref{ch:bc}, \cref{ch:situating}, \cref{ch:cg},
\cref{ch:multigrid}, \cref{ch:smoothers}, \cref{ch:matrixfree}, \cref{ch:records},
\cref{ch:monad}, \cref{ch:reductions}. \emph{One dispatch touched every part of the
book.} And the only lines that distinguish this run from the driven filter of
\cref{wex:trace-driven} are the three swapped cells of \cref{ch:fulltrace}---the
function space, the solve corner, and the reduction. The root is the same; the path
down from it is the whole machine; and which machine you get is a single field in a
file.
\end{workedexample}

\section{Where it goes next: the deferred frontiers}
\label{sec:wm-frontiers}

A closing chapter owes the reader an honest accounting of what the book did
\emph{not} do. There are three frontiers the synthesized library stops at, and it is
worth being precise about each---not as failures, but as the well-marked edges of a
deliberately scoped effort, each with a known direction onward (\cref{fig:wm-next}).

\paragraph{Distributed execution and sharding.}
The entire book targets a \emph{single machine}: one CPU lifting onto one GPU through a
device backend (\cref{ch:targets}). Real Palace runs across many machines, partitioning
the mesh and the fields into shards and exchanging data between them by message passing.
We deliberately read every distributed type as its single-rank equivalent and set the
parallel machinery aside, because the sharding theory assumes a memory and lifetime
structure that the pure, immutable-tensor spine of this book has rewritten---lifting it
naively could destabilize the very abstractions that made the spec clean. The
\emph{mathematics} of decomposition---splitting a tensor-field problem into coupled
sub-problems on sub-domains---is a genuine and tractable next direction, and it composes
with the pure picture rather than fighting it: a sharded solve is the same combinators
over a tensor that happens to be distributed. But it is future work, gated behind a
careful check that it does not disturb the single-machine spine, and we have marked it
as such throughout.

\paragraph{The opaque kernels, fully realized.}
Three places in the synthesized library are \code{\#extern} boundaries: a named
type signature standing in for an implementation we describe but do not fully render.
Each is an opaque library kernel that the spine \emph{depends on}, and each has a
known implementation in the vocabulary we built. The libCEED element-quadrature kernel
inside \code{fe\_assemble} is matrix-free operator application as element-local tensor
contractions (\cref{ch:matrixfree})---the implementation is understood; only its full
rendering is deferred. The SLEPc eigenvalue loop behind \code{eigsolve} is a
constructive Krylov--Schur iteration in our own \code{krylov\_step} and
\code{lanczos\_step} vocabulary (\cref{ch:krylovschur}). The MFEM time-step behind the
transient \code{fold\_solve} is one ODE-integrator step (\cref{ch:time}). For each, the
book gives \emph{both} surfaces---the opaque contract the spine calls, and the
in-our-vocabulary implementation it corresponds to---and leaves the final reconciliation
of the two as marked, well-understood work.

\paragraph{Deeper modularization.}
The five-module partition of the synthesized library (\cref{ch:fivelibs}) is a sensible
default, not a theorem. As the library is used---as components are reused across
drivers, as new analyses are added---some boundaries will want to move: a type may
migrate to cluster with the API that consumes it; a combinator family may earn its own
module. This is the ordinary refinement-by-use of any growing library, and the book
leaves the partition explicitly open to it.

\begin{figure}[t]
\centering
\begin{tikzpicture}[font=\footnotesize, >=Stealth, node distance=6mm and 10mm]
  % the firm core
  \node[stage, fill=simBgViolet, draw=simViolet, very thick, text width=30mm, minimum height=16mm]
    (core) {\textbf{this book}\\[1pt]\scriptsize the single-machine\\synthesized library\\(\cref{fig:wm-hero})};
  % three frontiers radiating out
  \node[extern, above right=4mm and 16mm of core, text width=34mm] (shard)
    {\textbf{distributed / sharding}\\\scriptsize decomposition as math;\\gated on spine stability\\(\cref{ch:targets})};
  \node[extern, right=22mm of core, text width=34mm] (kern)
    {\textbf{opaque kernels, realized}\\\scriptsize libCEED quadrature (\cref{ch:matrixfree}),\\\code{eigsolve} (\cref{ch:krylovschur}),\\ODE step (\cref{ch:time})};
  \node[extern, below right=4mm and 16mm of core, text width=34mm] (mod)
    {\textbf{deeper modularization}\\\scriptsize refine the 5-module\\partition by use (\cref{ch:fivelibs})};
  \draw[refedge] (core) -- (shard);
  \draw[refedge] (core) -- (kern);
  \draw[refedge] (core) -- (mod);
  \node[font=\scriptsize\itshape, text=simGray, below=2mm of core] {where it goes next};
\end{tikzpicture}
\caption[The deferred frontiers]{Where the book goes next. The firm single-machine
synthesized library (left) is bordered by three marked frontiers, each a direction
onward rather than a gap. \emph{Distributed/sharding}: the mathematics of domain
decomposition, composing with the pure picture, gated behind a check that it does not
destabilize the single-machine spine. \emph{The opaque kernels realized}: the three
\code{\#extern} boundaries---libCEED quadrature (\cref{ch:matrixfree}), the
\code{eigsolve} loop (\cref{ch:krylovschur}), the ODE step (\cref{ch:time})---each with
a known in-our-vocabulary implementation to be fully rendered. \emph{Deeper
modularization}: the five-module partition refined by use. None is a failure; each is a
well-understood next step from a deliberately scoped foundation.}
\label{fig:wm-next}
\end{figure}

\begin{pitfall}
It would be a misreading to take the \code{\#extern} boundaries as holes in the
argument. They are the opposite: a \code{\#extern} is a place where the book knows
\emph{exactly} what the implementation is---it has a type, a contract, and a
described realization in the calculus---and has chosen to defer the full rendering, not
the understanding. The frontier is one of \emph{effort remaining}, not of
\emph{comprehension missing}. That distinction is the difference between an honest
boundary and a gap.
\end{pitfall}

\section{The reflective close: three views, one machine}
\label{sec:wm-reflect}

Step all the way back. This book showed you the \emph{same} simulator three times.

Part~V showed it \emph{top-down}: the six analyses as the engineer meets them, each a
physical question with a physical product, walked through the lifecycle stage by stage
until the whole thing collapsed into one composition (\cref{ch:fulltrace}). Parts~II--IV
showed it \emph{bottom-up}: the vocabulary, built from the ground---the tensor and its
shapes, the de Rham complex, the weak forms, the calculus of pure transformations, the
Krylov and multigrid engines---each piece firm in itself before it was composed.
Part~VI showed it as the \emph{synthesized library}: the same machines, written out
with their types and bodies, rendered as the code an accelerator would run. Three
views---top-down modalities, bottom-up vocabulary, the implementation
library---and the central claim of the book is that they are \emph{one machine seen
three ways} (\cref{fig:wm-triality}).

\begin{figure}[t]
\centering
\begin{tikzpicture}[font=\footnotesize, >=Stealth]
  % the central machine
  \node[stage, fill=simBgViolet, draw=simViolet, very thick, circle, minimum size=22mm,
        text width=18mm, align=center] (m) at (0,0) {\textbf{one}\\\textbf{machine}};
  % three views around it
  \node[stage, fill=simBgGold, draw=simGold, text width=30mm, align=center] (top) at (0,2.7)
    {\textbf{top-down: modalities}\\\scriptsize 6 analyses, walked\\(Part~V, \cref{ch:fulltrace})};
  \node[stage, fill=simBgBlue, draw=simBlue, text width=30mm, align=center] (bot) at (-4.4,-2.2)
    {\textbf{bottom-up: vocabulary}\\\scriptsize tensors, de Rham, calculus,\\engines (Parts~II--IV)};
  \node[stage, fill=simBgGreen, draw=simGreen, text width=30mm, align=center] (lib) at (4.4,-2.2)
    {\textbf{the library: code}\\\scriptsize synthesized, accelerator-ready\\(Part~VI)};
  \draw[depedge] (top) -- (m);
  \draw[depedge] (bot) -- (m);
  \draw[depedge] (lib) -- (m);
  % the "they agree" annotations on the connecting edges
  \node[font=\scriptsize\itshape, text=simGray] at (1.9,1.5) {same composition};
  \node[font=\scriptsize\itshape, text=simGray] at (-2.9,-0.6) {same parts};
  \node[font=\scriptsize\itshape, text=simGray] at (2.9,-0.6) {same algorithm};
\end{tikzpicture}
\caption[The three views as one machine]{The triality. The book renders one
simulator three ways: \emph{top-down} as the six analyses walked through the lifecycle
(Part~V); \emph{bottom-up} as the vocabulary built from tensors and the de Rham complex
up through the calculus and the numerical engines (Parts~II--IV); and as the
\emph{synthesized library}, the same machines written as accelerator-ready code
(Part~VI). The book's central claim is that the three are not three descriptions of
three things but three views of \emph{one machine}---they compose the same parts into the
same algorithm, and \cref{ch:fulltrace} caught all three agreeing in a single trace.}
\label{fig:wm-triality}
\end{figure}

And here is the deep claim those three views were built to establish, now that you have
seen it proven component by component rather than asserted: \emph{a large physics
simulator is the repeated composition of a small set of mathematical building blocks.}
Not a sprawl of special cases. A small kit, composed (\cref{fig:wm-kit}). The kit is
short enough to list:

\begin{itemize}[leftmargin=*]
  \item \textbf{Four solve corners} (\cref{ch:situating})---the fixed-operator map, the
  operator-varying map, the state fold, and the black-box eigensolve. Every solve in
  every analysis is one of these four ways of threading the step.
  \item \textbf{Three reduction shapes} (\cref{ch:reductions})---the symmetric Gram, the
  port-projection, and the per-mode table. Every physical product is one of these three
  ways of collapsing a solution family to numbers.
  \item \textbf{The de Rham complex} (\cref{ch:derham})---the four function spaces
  ($\Hone$, $\Hcurl$, $\Hdiv$, $\Ltwo$) and the gradient, curl, and divergence that link
  them, from which every operator in the book is assembled.
  \item \textbf{The shared substrate}---the fold (\cref{ch:fold}), the state monad
  (\cref{ch:monad}), and the matrix-free operator (\cref{ch:matrixfree}). These thread
  through everything: the fold coordinates, the monad owns the state, the matrix-free
  apply is the universal leaf.
\end{itemize}

\begin{figure}[t]
\centering
\begin{tikzpicture}[font=\footnotesize, >=Stealth, node distance=4mm and 6mm]
  % four corners (2x2)
  \node[font=\small\bfseries, text=simRust] at (-3.2,2.4) {4 solve corners};
  \node[stage, fill=simBgRust, draw=simRust, minimum width=16mm, minimum height=8mm] (c1) at (-4.0,1.5) {\scriptsize fixed map};
  \node[stage, fill=simBgRust, draw=simRust, minimum width=16mm, minimum height=8mm] (c2) at (-2.3,1.5) {\scriptsize op-vary map};
  \node[stage, fill=simBgRust, draw=simRust, minimum width=16mm, minimum height=8mm] (c3) at (-4.0,0.6) {\scriptsize state fold};
  \node[extern, minimum width=16mm, minimum height=8mm] (c4) at (-2.3,0.6) {\scriptsize black-box};
  \node[cornertag] at (-3.15,0.0) {\cref{ch:situating}};
  % three reductions
  \node[font=\small\bfseries, text=simGold] at (0.9,2.4) {3 reductions};
  \node[stage, fill=simBgGold, draw=simGold, minimum width=20mm, minimum height=6mm] (r1) at (0.9,1.6) {\scriptsize Gram};
  \node[stage, fill=simBgGold, draw=simGold, minimum width=20mm, minimum height=6mm] (r2) at (0.9,0.9) {\scriptsize projection};
  \node[stage, fill=simBgGold, draw=simGold, minimum width=20mm, minimum height=6mm] (r3) at (0.9,0.2) {\scriptsize per-mode table};
  \node[cornertag] at (0.9,-0.35) {\cref{ch:reductions}};
  % shared substrate
  \node[font=\small\bfseries, text=simBlue] at (4.4,2.4) {shared substrate};
  \node[opbox, minimum width=24mm] (s1) at (4.4,1.6) {\scriptsize the fold (\cref{ch:fold})};
  \node[opbox, minimum width=24mm] (s2) at (4.4,0.9) {\scriptsize the monad (\cref{ch:monad})};
  \node[extern, minimum width=24mm] (s3) at (4.4,0.2) {\scriptsize matrix-free (\cref{ch:matrixfree})};
  % de Rham underneath
  \node[space, minimum width=86mm, minimum height=8mm] (dr) at (0.5,-1.3)
    {\footnotesize \textbf{the de Rham complex}: $\Hone \xrightarrow{\Grad} \Hcurl \xrightarrow{\Curl} \Hdiv \xrightarrow{\Div} \Ltwo$ \quad(\cref{ch:derham})};
  % composition arrows down to a product
  \draw[flow] (c4.south) -- (dr.north);
  \draw[flow] (r3.south) -- (dr.north);
  \draw[flow] (s3.south) -- (dr.north);
\end{tikzpicture}
\caption[The small kit of building blocks]{The book's thesis in one figure: a large
simulator is the repeated composition of a small kit. \emph{Four solve corners}
(\cref{ch:situating}) $\times$ \emph{three reduction shapes} (\cref{ch:reductions}),
resting on a \emph{shared substrate}---the fold (\cref{ch:fold}), the state monad
(\cref{ch:monad}), and the matrix-free operator (\cref{ch:matrixfree})---all assembled
over the \emph{de Rham complex} (\cref{ch:derham}). Six physically distinct analyses are
choices of one corner and one reduction over this same kit. The kit is small; the
simulator is its composition.}
\label{fig:wm-kit}
\end{figure}

That is the entire kit. Four corners, three reductions, one de Rham complex, and a
shared substrate of fold, monad, and matrix-free operator. Six physically distinct
simulations---capacitance, inductance, resonances, S-parameters, time-domain fields,
waveguide modes---are each a choice of \emph{one corner and one reduction} over this
kit, assembled on \emph{one} function-space complex. The thesis we held only as a
promise in \cref{ch:targets}---``six settings of one machine''---is now a fact you have
watched assemble itself, block by block, across forty-eight chapters.

\begin{keyidea}
A large physics simulator is the repeated composition of a small set of mathematical
building blocks---four solve corners, three reduction shapes, the de Rham complex, the
fold, the monad, and the matrix-free operator. You have now seen the whole machine: one
root, one finite tree, a handful of shared leaves. And you have seen \emph{why} it is
shaped the way it is---pure functions over immutable tensors---because that is the form
in which it will run on the accelerators it was always meant for. The variety is in the
composition; the kit is small.
\end{keyidea}

\subsection{What you can now do}
\label{sec:wm-cando}

The point of building this understanding was never the disassembly of one simulator. It
was a way of reading. With the kit in hand, you can do three things you could not before.

You can \emph{read any simulator's algorithms this way}. Open an unfamiliar
finite-element or finite-difference code and the questions are no longer ``what is all
this?'' but the three of \cref{ch:bigidea}: what operations happen, who owns the state,
how is the evolution coordinated? The in-place loops resolve into pure transformations;
the tangle of arrays resolves into threaded values; the control flow resolves into folds
and maps.

You can \emph{recognize the corners}. Faced with a new analysis, you can ask which of
the four solve corners its middle stage occupies---is the operator fixed or varying, is
the solve a map or a fold or a black box?---and which of the three reduction shapes
collapses its answer. Most of a new simulator is already familiar before you read a line
of its solver, because most of it is the shared substrate.

And you can \emph{compose a new analysis}. To add an analysis the book did not
cover---a thermal solve, a coupled multiphysics problem---is to choose a function space
from the de Rham complex, assemble its operators by the weak-form fold, pick the solve
corner its physics demands, and reduce its solution to the product an engineer wants. The
machine is a kit; you now know the pieces and how they snap together.

\section{A pointer to the appendices}
\label{sec:wm-appendix}

The figures of this book are deliberately pictures---block diagrams, composition trees,
trace tables---because pictures are how the \emph{shape} of the machine is best seen. A
reader who wants the exact rules behind the pictures will find them in the appendices,
which carry the deep computation semantics the main text kept at arm's length: the formal
calculus and its grammar (the precise syntax of the language sketched in
\cref{ch:language}), the shape algebra and its named-axis congruence rules, the algebraic
laws each combinator obeys, the numerical conventions that distinguish a transparent
performance trick from a load-bearing one, and the full matrix-free kernel detail behind
the universal leaf of \cref{ch:matrixfree}. The main text is the machine drawn; the
appendices are the machine specified to the last rule. Read them when a picture leaves you
wanting the proof.

\begin{keyidea}
\textbf{You have now seen the whole machine}---and why it is shaped for the accelerators it
will run on. A large simulator is not a monolith but the repeated composition of a small
kit of mathematical building blocks, rooted in a single \code{lifecycle} function from which
every part is reachable. The pure, immutable-tensor form was the right choice all along,
because it is the form in which the algorithm becomes a backend-lowering target rather than
a rewrite. That is the end of the disassembly, and the beginning of a way of reading every
simulator you meet from here on.
\end{keyidea}

\medskip
\noindent And that is the whole machine. We met it as six black boxes in
\cref{ch:targets}; we leave it as one small, connected, accelerator-ready
composition---understood, on one page, root to leaf. Thank you for building it with us.

\summarysection
This closing chapter views the whole simulator at once. The synthesized library has a
single root, \code{lifecycle}---its \code{main}---whose body is
\code{lifecycle $=$ fold\_solve (dispatch (problem\_type cfg)) $\circ$ build\_mesh}: build
the mesh (\cref{ch:functionspaces}), dispatch on the problem type at the specialization
seam (\cref{ch:lifecycle}) to one of six drivers, and fold the chosen driver through the
adaptive estimate--mark--refine loop (\cref{ch:amr}, degenerating to a single solve when
refinement is off). Every component of the book is reachable from this root by following
composition edges, and the whole machine fits on one page as a finite, shallow tree whose
deepest, most-reused leaf is the matrix-free apply (\cref{ch:matrixfree}). The functional
rendering is the \emph{backend-lowering target}: its immutable-tensor, pure-function,
whole-tensor-combinator form matches an external GPU/tensor backend, so the algorithm
lowers onto an accelerator unchanged---closing the promise of \cref{ch:bigidea}. Three
frontiers remain, marked as directions onward rather than failures: distributed/sharding
(deferred, gated on spine stability); the three \code{\#extern} opaque kernels (libCEED
quadrature, the \code{eigsolve} loop, the ODE step---each with a known in-vocabulary
implementation); and deeper modularization. Reflectively, the book showed one simulator
three ways---top-down modalities (Part~V), bottom-up vocabulary (Parts~II--IV), and the
synthesized library (Part~VI)---and proved they are one machine. The deep thesis, now
demonstrated component by component: a large physics simulator is the repeated composition
of a small kit---four solve corners (\cref{ch:situating}), three reduction shapes
(\cref{ch:reductions}), the de Rham complex (\cref{ch:derham}), and the shared substrate of
the fold (\cref{ch:fold}), the monad (\cref{ch:monad}), and the matrix-free operator
(\cref{ch:matrixfree}). With that kit, the reader can read any simulator's algorithms this
way, recognize its solve corners, and compose a new analysis. The appendices carry the exact
computation semantics behind the pictures.

\furtherreading
The simulator dissected throughout this book is AWS Labs' Palace~\citep{palace2023}, whose
source is the ground truth every layer of the spec cites. The matrix-free, element-local
operator-evaluation model that the synthesized library lowers onto an accelerator---the
universal leaf of \cref{fig:wm-hero} and the libCEED \code{\#extern} of
\S\ref{sec:wm-frontiers}---is developed by the CEED project~\citep{ceed2021}, the standard
reference for high-order matrix-free finite-element kernels. The iterative-solver
machinery the solve corners nest---the Krylov methods, the multigrid preconditioners, and
the reuse economics across an operator family---is treated comprehensively by
Saad~\citep{saad2003}, on which \cref{ch:gmres}, \cref{ch:multigrid}, and
\cref{ch:krylovschur} all draw. For the frontier of distributed execution and domain
decomposition (\S\ref{sec:wm-frontiers}), Saad's later chapters and the parallel-solver
literature are the natural next reading; the mathematics composes with the pure picture
this book built, and is the clearest single direction onward.

\exercisessection
\begin{enumerate}[nosep]
  \item \textbf{Read the root.} Write out the body of \code{lifecycle} from memory
  (\cref{eq:wm-lifecycle}), then name the chapter responsible for each of its three stages
  and for the six-way \code{dispatch}. Which single stage is identical for all six
  analyses, and which single field selects among them?
  \item \textbf{Trace a different dispatch.} Repeat the worked trace of
  \cref{wex:wm-dispatch}, but for a configuration with \code{problem\_type:~Eigenmode}.
  Which line of \code{lifecycle} changes its result, which driver does \code{dispatch}
  return, and at what point does this trace \emph{differ} from the electrostatic one---the
  solve corner, the reduction, or both? (Cross-check against \cref{ch:eigenmodes}.)
  \item \textbf{The universal leaf.} \Cref{fig:wm-hero} shows the matrix-free apply
  appearing in more than one place in the tree. List every distinct path from the root
  \code{lifecycle} down to a matrix-free apply, and explain why \emph{every} operator
  application in an iterative solve---the Krylov step, the smoother sweep, the residual---
  bottoms out in the same leaf.
  \item \textbf{Why pure, again.} In one paragraph, explain to a skeptic who has written
  fast in-place C why the synthesized library is written in pure functions over immutable
  tensors rather than mutating arrays. Your answer should connect the choice to the
  accelerator-lowering claim of \S\ref{sec:wm-backend} and the dependency-graph and
  composition payoffs of \cref{ch:bigidea}---and should address the objection that purity
  ``copies the array on every step.'' (Hint: revisit the pitfall of \cref{ch:bigidea}.)
  \item \textbf{The kit, counted.} \Cref{fig:wm-kit} lists the kit: four solve corners,
  three reduction shapes, the de Rham complex, and a shared substrate. For each of the six
  analyses of \cref{ch:targets}, state which solve corner and which reduction shape it uses.
  How many of the twelve corner-$\times$-reduction combinations does Palace actually
  occupy, and what does the answer say about the economy of the design?
  \item \textbf{Where would sharding go?} (Open-ended.) Suppose you were to add distributed
  execution to the synthesized library (\S\ref{sec:wm-frontiers}). Identify the \emph{one}
  node in the tree of \cref{fig:wm-hero} where a domain decomposition would most naturally
  enter, and argue whether it would change the \emph{structure} of the tree or only the
  implementation of a leaf. Why does the book claim sharding ``composes with the pure
  picture rather than fighting it''? What is the risk that gates it as future work?
  \item \textbf{Compose a new analysis.} (Open-ended.) Sketch how you would add a
  steady-state \emph{thermal} analysis (solve $-\Div(\kappa\Grad T)=q$ for a temperature
  field $T$, reduce to a thermal-resistance matrix) to the synthesized library. Which
  function space of the de Rham complex (\cref{ch:derham}) does $T$ live in? Which existing
  solve corner (\cref{ch:situating}) does the steady linear solve occupy? Which reduction
  shape (\cref{ch:reductions}) produces the resistance matrix, and which existing driver is
  your new one most nearly a copy of? Which pieces of the kit would you reuse unchanged, and
  what is the smallest set of new pieces you would have to write?
\end{enumerate}


\appendix
\part{Appendices --- Deep Computation Semantics}\label{part:appendices}
\chapter{The Formal Calculus: Grammar and Reduction Rules}
\label{app:calculus}

\chapterorient{Part~III taught the calculus the way one learns a spoken
language---by use, with intuition leading and precision following. This appendix
is the grammar book. It states, in one place and with no hand-waving, the two
pieces that \emph{define} the language: a \emph{grammar} (the legal types and
terms---what you are allowed to write down) and a \emph{small-step reduction
relation} (how an expression computes, one rewrite at a time---what those terms
\emph{mean}). Everything the main text said informally about values, records,
the state monad, \texttt{iterate\_while}, and demand-driven pruning is here as a
formal rule. You can read this appendix as the precise companion to
\cref{ch:language} (values and functions), \cref{ch:monad} (the state monad),
and \cref{ch:fold} (the fold and its pruning); each rule below carries a
back-reference to the chapter that uses it. The reward for the formality is
concrete: the whole language is a handful of reduction rules, and its one
genuinely unusual rule---demand-driven pruning---is what makes ``compute the
diagnostics'' and ``skip them for speed'' the \emph{same program}.}

\section{What a formal calculus is, and why this one is small}
\label{app:calc:orient}

A \emph{calculus}, in the sense used here, is a tiny formal language given by two
ingredients. The first is a \emph{grammar}: a finite set of rules, written in
Backus--Naur form (BNF), that says exactly which strings of symbols are legal
\emph{types} and legal \emph{terms}. The grammar is purely syntactic---it tells
you what you may write, not what it does. The second is an \emph{operational
semantics}: a \emph{reduction relation} $e \to e'$, read ``the expression $e$
steps to the expression $e'$,'' that says how a term computes by rewriting,
one small step at a time, until no rule applies and we are left with a
\emph{value}---a fully reduced answer.

This is the same machinery that defines the lambda calculus and, through it, every
typed functional language; our notation is a deliberately small dialect of it.
``Small'' is a design goal, not an accident. The calculus admits exactly what
Palace's algorithms need to be written cleanly---immutable values, records,
first-class functions, a state monad for orchestration, and one iteration
combinator---and nothing more. There is no inheritance, no exception machinery,
no mutable reference cell, no general recursion beyond the one loop combinator.
The smaller the calculus, the shorter the argument that an L3 form and its L4
re-expression compute the same thing (\cref{ch:rotations}), and the less there is
for a reader---or a checking tool---to get wrong.

\begin{keyidea}
The entire computational content of the calculus is a \emph{handful of reduction
rules}: $\beta$-reduction and \code{let} for functions, projection and
destructuring for records, three monad laws and the state effects for
orchestration, one rule for operator application, one recursive rule for
\code{iterate\_while}, and---the signature move---one rule for demand-driven
pruning. Master those and you have the formal meaning of every program in the
book. The grammar says what you may write; the reduction rules say what it means.
\end{keyidea}

\Cref{fig:appA:layered} shows the shape of the whole appendix: the grammar
defines two layers (the type language and the term language), and the reduction
relation acts on the term layer, rewriting terms to values while the type
language stands guard over what is well-formed.

\begin{figure}[t]
\centering
\begin{tikzpicture}[font=\footnotesize, node distance=5mm]
  % ===== top band: the TYPE language =====
  \node[band, fill=simBgBlue, draw=simBlue, minimum width=118mm, thick] (typesband)
    at (0,2.6) {};
  \node[font=\small\bfseries, simBlue, anchor=west] at (-5.7,2.6)
    {\textsc{type language}~~$\tau$};
  \node[font=\scriptsize, simInk, anchor=east, align=right] at (5.7,2.6)
    {\code{Scalar} \;\code{Tensor[$\sigma$]} \;\code{\{l:$\tau$\}} \;$\tau_1{\to}\tau_2$
     \;\code{Sim $\tau$} \;\code{Op[$\cdot$]} \;\code{!$\tau$}};
  % ===== middle band: the TERM language =====
  \node[band, fill=simBgGray, draw=simGray, minimum width=118mm, thick] (termsband)
    at (0,1.3) {};
  \node[font=\small\bfseries, simInk, anchor=west] at (-5.7,1.3)
    {\textsc{term language}~~$e$};
  \node[font=\scriptsize, simInk, anchor=east, align=right] at (5.7,1.3)
    {$x$ \;$\lambda x.e$ \;$e_1\,e_2$ \;\code{let} \;\code{\{l:e\}} \;\code{e.l}
     \;\code{if} \;\code{op($\bar e$)} \;\code{apply} \;\code{do} \;\code{iterate\_while}};
  % types classify terms
  \draw[depedge] (typesband.south) -- node[right, font=\scriptsize, simBlue]
    {classifies (\;$\Gamma \vdash e : \tau$\;)} (termsband.north);
  % ===== reduction acts on terms, producing values =====
  \node[product, minimum width=26mm] (val) at (0,-0.9) {a \emph{value} $v$};
  \draw[flow] (termsband.south) -- node[right, font=\scriptsize]
    {reduction $e \to e'$ (\cref{app:calc:reduction})} (val);
  \node[align=left, font=\scriptsize\itshape, simGray, text width=40mm]
    at (4.9,-0.4) {a term reduces, one small step at a time, until no rule
    applies---then it is a value.};
\end{tikzpicture}
\caption[The layered structure of the calculus]{The calculus has two grammatical
layers and one dynamics. The \emph{type language} $\tau$
(\cref{app:calc:types}) names the kinds of thing that exist---scalars, tensors,
records, functions, the simulator monad, operators. The \emph{term language} $e$
(\cref{app:calc:terms}) names the expressions you may write. A typing judgment
$\Gamma \vdash e : \tau$ says a term is well-formed and assigns it a type
(\cref{app:calc:typing}). The \emph{reduction relation} $e \to e'$
(\cref{app:calc:reduction}) is the dynamics: it rewrites a well-typed term, one
step at a time, down to a value. This appendix gives all three.}
\label{fig:appA:layered}
\end{figure}

\section{The grammar: types}
\label{app:calc:types}

We give the grammar in two halves, types first. A \emph{type} classifies values:
it is the static description of what kind of thing an expression denotes, checked
before any reduction happens. The type language is the following BNF, with each
production a way of forming a legal type $\tau$. (Throughout, $\tau$ ranges over
types, $\sigma$ over shape expressions, $l$ over field labels.)

\begin{lstlisting}[style=l4style, language=]
tau ::= Scalar | Bool | Int | Float         -- ground scalar types
      | Tensor[sigma]                        -- tensor with symbolic shape sigma
      | { l1: tau1, ..., ln: taun }          -- record (TypeScript-flavored)
      | tau1 x ... x taun                    -- tuple
      | tau1 -> tau2                          -- function
      | Sim tau                              -- simulator monad
      | Op[tau_in -> tau_out]                -- operator carrying baked-in params
      | !tau                                 -- shareable (non-linear) annotation
\end{lstlisting}

\noindent Each production earns its place. We name what it is \emph{for}, with a
one-line example, taking them in the order written.

\begin{description}[style=nextline, leftmargin=1.4em, font=\normalfont\itshape]
  \item[Ground scalars \code{Scalar}, \code{Bool}, \code{Int}, \code{Float}.] The
  atoms: a single number, a truth value, an integer, a floating-point number.
  \code{Scalar} is the generic numeric atom (a coefficient $\alpha$ in
  \code{axpy}); \code{Bool} types a predicate's answer; \code{Int} an iteration
  counter; \code{Float} an explicit floating value. \emph{Example:}
  \lstinline[style=l4style]|tol : Float|, \lstinline[style=l4style]|converged : Bool|.
  \item[Tensors \code{Tensor[$\sigma$]}.] An immutable multidimensional array
  whose shape is the symbolic expression $\sigma$ (the shape grammar is
  \cref{app:shapes}; we treat $\sigma$ as opaque here). This is the
  workhorse---every field, every vector, every operator's argument is a tensor.
  \emph{Example:} \lstinline[style=l4style]|Tensor[N]| is a length-$N$ vector;
  \lstinline[style=l4style]|Tensor[m, n]| a matrix.
  \item[Records \code{\{$l_1$: $\tau_1$, \dots, $l_n$: $\tau_n$\}}.] A bundle of
  named fields, each with its own type---the calculus's only aggregate
  (\cref{ch:records}). Solver state, configuration, and every step's output are
  records. \emph{Example:}
  \lstinline[style=l4style]|{ x: Float, residual: Float }|.
  \item[Tuples \code{$\tau_1 \times \cdots \times \tau_n$}.] A positional bundle:
  a fixed-length, heterogeneous sequence accessed by position rather than by
  name. Used where a couple of results travel together and naming them would be
  ceremony. \emph{Example:} \lstinline[style=l4style]|(Tensor[N], Scalar)|.
  \item[Functions \code{$\tau_1 \to \tau_2$}.] The type of a value that maps a
  $\tau_1$ to a $\tau_2$. The arrow associates to the right, so
  \lstinline[style=l4style]|A -> B -> C| is \lstinline[style=l4style]|A -> (B -> C)|
  (\cref{ch:language}). \emph{Example:}
  \lstinline[style=l4style]|Tensor[N] -> Scalar| is the type of \code{nrm2}.
  \item[The simulator monad \code{Sim $\tau$}.] The type of an \emph{action} that,
  when run against the simulator's state, produces a $\tau$ while threading that
  state (\cref{ch:monad}). It marks orchestration---file I/O, multi-stage
  workflows, RNG threading---as distinct from pure per-step computation.
  \emph{Example:} \lstinline[style=l4style]|get : Sim S|.
  \item[Operators \code{Op[$\tau_{\text{in}} \to \tau_{\text{out}}$]}.] A function
  value that \emph{also carries baked-in parameters}---matrix entries, a
  factorization, basis tables (\cref{ch:operators}). It is the closure ``configure
  once, apply many times'' given a first-class spelling; the brackets group the
  in/out arrow. \emph{Example:}
  \lstinline[style=l4style]|Op[Tensor[N] -> Tensor[N]]|, an assembled linear
  operator.
  \item[The shareable annotation \code{!$\tau$}.] A tag marking a value as
  \emph{shareable} (non-linear): it may be read repeatedly without being consumed.
  It distinguishes an operator's read-only baked-in parameters from the linearly
  threaded simulator state (\cref{app:calc:ownership}). \emph{Example:}
  \lstinline[style=l4style]|!LbmTables|, the velocity/weight tables closed into an
  operator.
\end{description}

\begin{notationbox}
The shape expression $\sigma$ inside \code{Tensor[$\sigma$]} has its own grammar
and its own algebra (literal axes, named axes, congruence groups
\lstinline[style=l4style]|Tensor[(S: ...)]|, rank wildcards). That grammar is the
subject of \cref{app:shapes} and is treated as opaque here: as far as the
\emph{term} calculus is concerned, \code{Tensor[$\sigma$]} is one ground type,
and shape checking is a side condition resolved by a separate shape solver. The
two appendices compose---this one gives the term dynamics, that one gives the
shape discipline---but neither depends on the internals of the other.
\end{notationbox}

\section{The grammar: terms}
\label{app:calc:terms}

The term language is what you actually write. Each production is an expression
form $e$; the reduction rules of \cref{app:calc:reduction} act on exactly these.

\begin{lstlisting}[style=l4style, language=]
e ::= x                                      -- variable
    | c                                      -- constant / literal
    | \(x: tau). e                           -- lambda abstraction
    | e1 e2                                  -- application
    | let x = e1 in e2                       -- let-binding
    | let { l1, ..., ln } = e1 in e2         -- destructuring let
    | let { l1, ..., lk, ...r } = e1 in e2   -- destructure with rest
    | { l1: e1, ..., ln: en }                -- record literal
    | { ...e0, li: ei, ..., lj: ej }         -- record spread / update
    | e.l                                    -- field access
    | (e1, ..., en)                          -- tuple literal
    | proj_i e                               -- tuple projection
    | if e0 then e1 else e2                  -- conditional
    | op(e1, ..., en)                        -- primitive operator application
    | op-with-params { p = e ; \(x: t). e }  -- operator-VALUE introducer
    | apply A e                              -- operator application (eliminator)
    | return e                               -- monadic return (Sim)
    | e1 >>= e2                              -- monadic bind
    | do { s1; ...; sn; e }                  -- do-notation block
    | get | put e | modify e                 -- state effects (Sim)
    | while_loop init cond step              -- pure tail-recursive loop
\end{lstlisting}

\noindent We group the productions by role and explain what each is for. The
order mirrors the reduction rules to come.

\paragraph{The lambda core.} A \emph{variable} \code{x} stands for a value bound
in scope; a \emph{constant} \code{c} is a literal (\code{7}, \code{true}). A
\emph{lambda} \lstinline[style=l4style]|\(x: tau). e| is an anonymous function of
one parameter \code{x} of type \code{tau} with body \code{e}---the only way to
introduce a function value. \emph{Application} \code{e1 e2} supplies an argument
to a function by juxtaposition. \emph{Let} \lstinline[style=l4style]|let x = e1 in e2|
names an intermediate: bind \code{x} to the value of \code{e1}, then evaluate
\code{e2} with \code{x} in scope. These four are the lambda calculus, and they
do everything the main text's ``values and functions'' chapter teaches
(\cref{ch:language}). \emph{Example:}
\lstinline[style=l4style]|let sq = \(x: Float). x * x in sq 5|.

\paragraph{Records.} A \emph{record literal}
\lstinline[style=l4style]|{ l1: e1, ..., ln: en }| builds a bundle of named
fields. \emph{Field access} \code{e.l} reads one field out. The two
\emph{destructuring lets} pull several fields at once---the plain form names a
fixed set of fields, the rest form
\lstinline[style=l4style]|{ l1, ..., lk, ...r }| binds some fields and collects
the remainder into \code{r}. \emph{Record spread / update}
\lstinline[style=l4style]|{ ...e0, li: ei }| builds a new record equal to
\code{e0} but with the listed fields replaced---never a mutation, always a fresh
record (\cref{ch:records}). \emph{Example:}
\lstinline[style=l4style]|{ ...s, dist: dist_next, step: s.step + 1 }|.

\paragraph{Tuples and conditionals.} A \emph{tuple literal}
\code{(e1, ..., en)} bundles values positionally; \emph{projection}
\lstinline[style=l4style]|proj_i e| reads the $i$-th component. The
\emph{conditional} \lstinline[style=l4style]|if e0 then e1 else e2| chooses a
branch on a \code{Bool}. \emph{Example:}
\lstinline[style=l4style]|if s.residual > tol then keep else stop|.

\paragraph{Operators.} \emph{Primitive operator application}
\lstinline[style=l4style]|op(e1, ..., en)| applies a built-in primitive
(\code{axpy}, \code{dot}, \code{matvec}, \dots), each governed by its own
$\delta$-rule (\cref{app:laws}). The \emph{operator-value introducer}
\lstinline[style=l4style]|op-with-params { p = e ; \(x: t). e_body }| builds a
first-class operator value: a record of closed-over \code{!}-shareable parameters
\code{p}, paired with a body lambda over the run-time argument \code{x}, of type
\lstinline[style=l4style]|Op[t -> t_out]| (\cref{ch:operators}). Its eliminator is \emph{operator
application} \code{apply A e}, which supplies the run-time argument to an operator
value \code{A}. \emph{Example:}
\lstinline[style=l4style]|apply bgkOp streamed|.

\paragraph{The monad and state effects.} \emph{Return}
\lstinline[style=l4style]|return e| injects a pure value into a \code{Sim}
action; \emph{bind} \lstinline[style=l4style]|e1 >>= e2| sequences two actions,
feeding the first's result into the second; a \emph{do-block}
\lstinline[style=l4style]|do { s1; ...; sn; e }| is the readable sugar for a chain
of binds (\cref{ch:monad}). A \code{do}-statement \code{s} is one of
\lstinline[style=l4style]|x <- e| (bind), \lstinline[style=l4style]|let x = e|
(pure binding), or a bare \code{e} (an effect of type \code{Sim ()}). The
\emph{state effects} \code{get} (read the current state), \lstinline[style=l4style]|put e|
(replace it), and \lstinline[style=l4style]|modify e| (apply a function to it) are
the only way state is touched. \emph{Example:}
\lstinline[style=l4style]|do { s <- get; put (step s) }|.

\paragraph{The loop.} \lstinline[style=l4style]|while_loop init cond step| is the
one primitive iteration form---a pure tail-recursive loop. The book's
\code{iterate\_while} (\cref{ch:fold}) is the named, trajectory-collecting
specialization of it, and its small-step rule is \cref{app:calc:iterate} below.
This is the \emph{only} unbounded control construct in the calculus; everything
else terminates structurally.

\begin{keyidea}
The grammar is closed: there is no production for a mutable variable, no
assignment, no \code{goto}, no general recursion outside \code{while\_loop}.
Every term is built from this finite list, and every term's behavior is fixed by
the reduction rules of the next section. The expressive reach---six simulators
from one vocabulary---comes from \emph{first-class functions and operators}
(lambda, \code{apply}, \code{Op}) combined with the single loop, not from a large
language.
\end{keyidea}

\section{Ownership: three categories the syntax enforces}
\label{app:calc:ownership}

Before the dynamics, one static distinction the grammar quietly encodes, because
several reduction rules below depend on it. The calculus sorts every value into
exactly one of three \emph{ownership categories}, and the sort is syntactic---you
can read it off the form.

\begin{description}[style=nextline, leftmargin=1.4em, font=\normalfont\itshape]
  \item[Simulator state.] Values threaded through the \code{Sim} monad, touched
  only via \code{get} / \code{put} / \code{modify}. Treated as \emph{linear}:
  a \code{Sim} action consumes the current state and produces the next, so the old
  state is shadowed and aliasing is structurally impossible. There is exactly one
  state slot in scope at any point in a \code{do} block.
  \item[Operator internal parameters.] Values closed over inside an
  \code{Op[$\cdot$]} and tagged \code{!} (shareable). Read-only across a solve;
  they never evolve. Matrix entries, factorizations, basis and lattice tables,
  time-step constants.
  \item[Ephemeral intermediates.] Values bound by a \code{let} or a \code{do}
  binding. Local in scope; they participate in neither state evolution nor
  operator closure. The residual computed mid-iteration, the CG search direction,
  a predictor in a multi-step integrator.
\end{description}

\noindent Linearity is tracked only to enforce the first-versus-second
distinction---which slot owns the canonical current value. It is \emph{not}
tracked at the tensor level: every tensor is immutable, so two names for one
tensor can never surprise each other, and \code{clone()} is meaningless
(\cref{ch:language}). The \code{!} annotation is the syntactic marker; inside an
\code{Op[$\cdot$]}, all closed-over data is \code{!}-tagged by construction. We
will see this category distinction surface in the state-effect rules
(\cref{app:calc:state}) and the operator-application rule
(\cref{app:calc:apply}).

\section{Small-step operational semantics}
\label{app:calc:reduction}

We now give the dynamics: the reduction relation $e \to e'$, presented as a set
of \emph{inference rules}. An inference rule is read as ``\emph{if} the premises
above the line hold, \emph{then} the conclusion below the line holds.'' For most
of our reduction rules there are no premises---they are \emph{axioms}, an
unconditional rewrite, drawn as a conclusion with an empty line above it (or
simply stated as an equation $e \to e'$). The pruning rule of
\cref{app:calc:pruning} is the one rule with a genuine premise.

Throughout, $v$ ranges over \emph{values}---fully reduced terms that no rule
rewrites further (a literal, a lambda, a fully-built record of values, an operator
value). The notation $e[v/x]$ means ``the term $e$ with every free occurrence of
the variable $x$ replaced by $v$''---capture-avoiding substitution, the single
operation underlying nearly every rule.

\Cref{fig:appA:reduction} traces a small expression all the way down to a value
by these rules, to fix the idea of ``reduce one step at a time'' before we
catalogue the steps.

\begin{figure}[t]
\centering
\begin{tikzpicture}[font=\footnotesize, node distance=3mm]
  \node[opbox, align=left, minimum width=104mm] (e0) at (0,0)
    {\lstinline[style=l4style]|let r = axpy 2 x y in r.|... \quad\textit{(start: a let over an operator app)}};
  \node[font=\scriptsize, simBlue, right=24mm of e0.east, anchor=west] (l0) {};
  \node[opbox, align=left, minimum width=104mm, below=8mm of e0] (e1)
    {\lstinline[style=l4style]|let r = (2 .* x .+ y) in r.|... \quad\textit{($\delta$-rule for \code{axpy}: unfold the primitive)}};
  \node[opbox, align=left, minimum width=104mm, below=8mm of e1] (e2)
    {\lstinline[style=l4style]|let r = v_r in r.|... \quad\textit{(the tensor arithmetic evaluates to a value $v_r$)}};
  \node[opbox, align=left, minimum width=104mm, below=8mm of e2] (e3)
    {\lstinline[style=l4style]|v_r.|... \quad\textit{(\textbf{let} rule, \S\ref{app:calc:lambda}: substitute $v_r$ for \code{r})}};
  \node[product, align=left, minimum width=104mm, below=8mm of e3] (e4)
    {a \emph{value} --- no rule applies; reduction halts.};
  % step arrows with rule labels
  \draw[flow] (e0) -- node[right, font=\scriptsize, simBlue]{$\to$ \;$\delta$} (e1);
  \draw[flow] (e1) -- node[right, font=\scriptsize, simBlue]{$\to^{*}$ \;(arith)} (e2);
  \draw[flow] (e2) -- node[right, font=\scriptsize, simBlue]{$\to$ \;\textsf{let}} (e3);
  \draw[flow] (e3) -- node[right, font=\scriptsize, simBlue]{$\to^{*}$ \;\textsf{proj}} (e4);
\end{tikzpicture}
\caption[A term reducing step by step]{Reduction is rewriting, one small step at
a time. An expression is repeatedly matched against the reduction rules; each
match rewrites it to a simpler term ($\to$ is one step, $\to^{*}$ several), until
no rule applies and the term is a \emph{value}. Here a \code{let} over an
operator application unfolds the \code{axpy} $\delta$-rule (\cref{app:laws}),
evaluates the tensor arithmetic, then fires the \code{let} rule
(\cref{app:calc:lambda}) to substitute the result. The rule fired at each step is
named on the arrow. The order is forced only by data dependency
(\cref{app:calc:confluence}); the answer does not depend on which legal step we
take first.}
\label{fig:appA:reduction}
\end{figure}

\subsection{The lambda core: \texorpdfstring{$\beta$}{beta} and \texttt{let}}
\label{app:calc:lambda}

The two foundational rules. \emph{$\beta$-reduction} is function
application---supply an argument to a lambda by substituting it into the body.
\emph{The \code{let} rule} names a value---bind a value to a name by substituting
it into the body. They are the same operation (substitution) wearing two faces.

\[
\begin{array}{rcll}
(\lambda(x:\tau).\, e)\ v &\;\to\;& e[v/x] & \text{($\beta$)} \\[2pt]
\textsf{let}\ x = v\ \textsf{in}\ e &\;\to\;& e[v/x] & \text{(\textsf{let})}
\end{array}
\]

\noindent \emph{Intuition.} $\beta$ says: to apply
\lstinline[style=l4style]|(\(x: Float). x * x)| to \code{5}, erase the lambda and
replace \code{x} by \code{5} in the body, giving \lstinline[style=l4style]|5 * 5|.
The \code{let} rule says the same for naming: \lstinline[style=l4style]|let y = 6 in y + 1|
becomes \lstinline[style=l4style]|6 + 1|. There are no boxes and no overwrites
(\cref{ch:language}): a name is a stand-in for a value, and reduction simply
substitutes the value back in. Currying and partial application
(\cref{ch:language}) are nothing more than $\beta$ applied one argument at a time:
\lstinline[style=l4style]|add 3 4| is \lstinline[style=l4style]|(add 3) 4|, two
$\beta$-steps.

\subsection{Records: projection, destructuring, spread}
\label{app:calc:records}

Records have three rules: \emph{projection} reads one field; \emph{destructuring}
binds several fields at once; \emph{spread} builds an updated copy. All operate on
a record of \emph{values} $\{l_1\!:\!v_1, \dots, l_n\!:\!v_n\}$.

\[
\begin{array}{rcll}
\{l_1\!:\!v_1,\dots,l_n\!:\!v_n\}.l_k &\;\to\;& v_k & \text{(proj)} \\[3pt]
\textsf{let}\ \{l_1,\dots,l_n\} = \{l_1\!:\!v_1,\dots,l_n\!:\!v_n\}\ \textsf{in}\ e
  &\;\to\;& e[v_1/l_1,\dots,v_n/l_n] & \text{(destr)} \\[3pt]
\{\dots\{l_1\!:\!v_1,\dots,l_n\!:\!v_n\},\ l_k\!:\!w\}
  &\;\to\;& \{l_1\!:\!v_1,\dots,l_k\!:\!w,\dots,l_n\!:\!v_n\} & \text{(spread)}
\end{array}
\]

\noindent \emph{Intuition.} Projection \code{r.x} pulls the named field's value
straight out. Destructuring is projection done in bulk: it binds each field name
to its value in one stroke, so
\lstinline[style=l4style]|let { x, residual } = s in ...| makes both fields
available without writing \code{s.x} and \code{s.residual} everywhere. Spread is
the immutable update: \lstinline[style=l4style]|{ ...s, x: x' }| produces a
\emph{new} record identical to \code{s} except in field \code{x}---the original
\code{s} is untouched, exactly as immutability requires (\cref{ch:records}). The
``destructure with rest'' form
\lstinline[style=l4style]|let { l1, ..., lk, ...r } = ...| is the same destr rule
with \code{r} bound to the record of remaining fields; we elide its rule as a
routine variant. The spread rule is the formal content of every
``\lstinline[style=l4style]|{ ...s, dist: dist_next }|'' you have seen: it is how
a step ``updates'' state without mutating it.

\subsection{The monad: \texttt{return}, \texttt{>>=}, and \texttt{do}}
\label{app:calc:monad}

Orchestration is sequenced by the three \emph{monad laws}. These govern how
\code{Sim} actions compose; they are the precise meaning of a \code{do}-block
(\cref{ch:monad}).

\[
\begin{array}{rcll}
\textsf{return}\ v \,\mathbin{\texttt{>>=}}\, f &\;\to\;& f\ v & \text{(left identity)} \\[2pt]
m \,\mathbin{\texttt{>>=}}\, \textsf{return} &\;\to\;& m & \text{(right identity)} \\[2pt]
(m \,\mathbin{\texttt{>>=}}\, f) \,\mathbin{\texttt{>>=}}\, g
  &\;\to\;& m \,\mathbin{\texttt{>>=}}\, (\lambda x.\, f\ x \,\mathbin{\texttt{>>=}}\, g)
  & \text{(associativity)}
\end{array}
\]

\noindent \emph{Intuition.} \emph{Left identity}: injecting a value with
\code{return} and then binding it is just calling the continuation on the
value---\code{return} adds no effect. \emph{Right identity}: binding an action to
\code{return} changes nothing---returning the result is a no-op. \emph{Associativity}:
how you parenthesize a chain of binds does not matter, which is exactly why the
flat \code{do}-block notation is well-defined---\lstinline[style=l4style]|do { x <- m; y <- n; k }|
is unambiguous because the binds re-associate freely. A \code{do}-block desugars
to a right-nested chain of \code{>>=}: \lstinline[style=l4style]|do { x <- m; rest }|
is \lstinline[style=l4style]|m >>= (\x -> do { rest })|, and a pure
\lstinline[style=l4style]|let x = e; rest| desugars to
\lstinline[style=l4style]|(\x -> do { rest }) e|. The three laws guarantee these
desugarings all mean the same thing.

\subsection{State effects}
\label{app:calc:state}

The state effects \code{get} / \code{put} / \code{modify} are the only way the
threaded simulator state (\cref{app:calc:ownership}, category one) is touched.
Two rewrite rules relate them, capturing that \code{modify} is exactly
``get, transform, put'':

\[
\begin{array}{rcl}
\textsf{do}\ \{\, x \leftarrow \textsf{get};\ \textsf{put}\ f(x);\ \textsf{return}\ ()\,\}
  &\;\to\;& \textsf{modify}\ f \\[3pt]
\textsf{modify}\ (f \circ g) &\;\to\;& \textsf{do}\ \{\,\textsf{modify}\ g;\ \textsf{modify}\ f\,\}
\end{array}
\]

\noindent \emph{Intuition.} The first rule says \code{modify f} is shorthand for
reading the state, applying \code{f}, and writing it back---there is no separate
``read-modify-write'' hazard because the state slot is linear and the old value is
shadowed. The second rule says modifying by a composed function $f \circ g$ is the
same as modifying by $g$ then by $f$: state updates compose in application order.
Together they let the orchestration layer be rewritten freely between the explicit
\code{get}/\code{put} form and the compact \code{modify} form without changing
meaning. Crucially, per-algorithm state evolution does \emph{not} need these---a
solver's state threads through \code{iterate\_while} as a pure carry
(\cref{app:calc:iterate}); \code{Sim} recedes to genuine orchestration (I/O,
multi-stage workflows).

\subsection{Operator application}
\label{app:calc:apply}

An operator value pairs closed-over parameters \code{p} with a body lambda
$\lambda x.\, e$ (\cref{app:calc:terms}, the operator-value introducer). Applying
it substitutes \emph{both} the parameters and the run-time argument into the body:

\[
\textsf{apply}\ (\textsf{op-with-params}\ p,\ \lambda x.\, e)\ v
  \;\to\; e[p/\textsf{params},\ v/x]
\]

\noindent \emph{Intuition.} This is $\beta$-reduction with a twist: an ordinary
lambda has only its parameter to substitute, but an operator value also carries a
record of baked-in parameters (the matrix, the factorization, the basis tables),
and applying it substitutes those \emph{and} the run-time argument
simultaneously. The notation $e[p/\textsf{params}]$ means the body's references to
its closed-over parameters are replaced by the actual closed-over values \code{p}.
This is the formal content of ``build once with configuration baked in, apply many
times'' (\cref{ch:operators}): the \code{op-with-params} introducer is the
``build'' (it captures \code{p} once), and each \code{apply A v} is one ``apply.''
The parameters \code{p} are \code{!}-shareable
(\cref{app:calc:ownership}, category two), so the same operator value may be
applied to many arguments with no copying and no consumption.

\subsection{Primitive \texorpdfstring{$\delta$}{delta}-rules}
\label{app:calc:delta}

Each primitive operator---\code{axpy}, \code{dot}, \code{matvec}, \code{conv},
\code{reduce}, \dots---comes with its own \emph{$\delta$-rule}, an axiom giving
its action on value arguments. We do not enumerate them here: the primitive
semantics live in the concept library and the algebraic-law appendix
(\cref{app:laws}), where each $\delta$-rule can be checked against a reference
implementation. The calculus assumes only that every primitive \code{op} comes
\emph{with} such a rule, so that \lstinline[style=l4style]|op(v1, ..., vn)|
reduces to a value. \Cref{fig:appA:reduction} fired the \code{axpy} $\delta$-rule
as its first step. The point of structure here is the division of labor: the
\emph{shape} of computation (functions, records, loops, pruning) is fixed by this
appendix; the \emph{arithmetic} of each primitive is fixed by its $\delta$-rule
elsewhere.

\subsection{The \texttt{iterate\_while} rule}
\label{app:calc:iterate}

Here is the heart of the dynamics: the small-step rule for the iteration
combinator the whole book turns on (\cref{ch:fold}). Given a carry $a$, a
predicate $p : \alpha \to \code{Bool}$, and a step
$f : \alpha \to \{\code{state}\!:\!\alpha, \dots e\}$, one unrolling step of
\code{iterate\_while} is:

\[
\begin{aligned}
\textsf{iterate\_while}\ a\ p\ f \;\to\;\ &\textsf{if}\ p(a)\\
  &\quad\textsf{then}\ \textsf{let}\ \{\textsf{state}\!:\!a',\ \dots e\} = f(a)\ \textsf{in}\\
  &\qquad \textsf{let}\ \{\textsf{final\_state},\ \textsf{trajectory}\}
        = \textsf{iterate\_while}\ a'\ p\ f\ \textsf{in}\\
  &\qquad \{\textsf{final\_state},\ \textsf{trajectory}\!:\
        [\{\dots e\}] \mathbin{+\!\!+} \textsf{trajectory}\}\\
  &\quad\textsf{else}\ \{\textsf{final\_state}\!:\ a,\ \textsf{trajectory}\!:\ [\,]\,\}
\end{aligned}
\]

\noindent \emph{Intuition.} Read it as a recursion that unrolls one loop step per
firing. Test the predicate on the \emph{current} carry $a$. If it is false, stop:
the answer is $a$ as \code{final\_state} with an empty \code{trajectory}. If it is
true, run one step $f(a)$, which yields the next carry $a'$ together with this
step's per-step extras $e$; recurse on $a'$; and prepend $e$ to whatever
trajectory the recursion returns. Two structural facts are baked into the rule and
are load-bearing in \cref{ch:fold}. First, the predicate fires \emph{before} any
step (the \code{while}, not \code{do}/\code{while}, convention): an
already-converged input does zero work. Second, the predicate reads only the carry
$a$---its type is $\alpha \to \code{Bool}$---so any termination quantity must be
\emph{folded into the carry}; letting the predicate read the step's extras would
be a circular dependency, and the type forbids it. The recursion sits in tail
position, so an implementation realizes it as an ordinary loop.

\Cref{fig:appA:unroll} draws one firing of the rule as the operation it really
is: a single step splits into a \emph{carry update} (the next $a'$, threaded
forward) and a \emph{trajectory extension} (the extras $e$, prepended to the
list).

\begin{figure}[t]
\centering
\begin{tikzpicture}[font=\footnotesize, node distance=7mm]
  % the rule as: one application splits into carry-update + trajectory-extend
  \node[stage, minimum width=40mm, align=center] (lhs) at (0,0)
    {\lstinline[style=l4style]|iterate_while a p f|};
  \node[font=\scriptsize, simBlue] (arrow) at (3.3,0) {$\to$ \;(one unrolling)};
  % predicate gate
  \node[diamond, draw=simBlue, fill=simBgBlue, aspect=1.7, inner sep=1pt,
        font=\scriptsize, text=simInk] (p) at (6.7,0) {$p(a)$?};
  \draw[flow] (lhs.east) -- (p.west);
  % FALSE branch
  \node[product, minimum width=34mm, align=center] (stop) at (6.7,-2.2)
    {\lstinline[style=l4style]|{ final_state: a, trajectory: [] }|};
  \draw[backflow] (p.south) -- node[right, font=\scriptsize, simRust]{no} (stop.north);
  % TRUE branch: run the step
  \node[opbox, minimum width=20mm] (step) at (10.4,0) {\lstinline[style=l4style]|f(a)|};
  \draw[flow] (p.east) -- node[above, font=\scriptsize, simGreen]{yes} (step.west);
  % step splits into carry-update and trajectory-extend
  \node[data, fill=white, draw=simGray, minimum width=20mm] (carry) at (10.4,-1.5)
    {carry update: $a'$};
  \node[data, fill=simBgGold, draw=simGold, minimum width=26mm] (traj) at (10.4,-2.9)
    {trajectory ext.: $[\{\dots e\}]\mathbin{+\!\!+}\cdots$};
  \draw[flow] (step.south) -- node[right, font=\scriptsize]{\code{.state}} (carry.north);
  \draw[refedge, simGold!80!black] (step.south west) to[out=240,in=80] (traj.north);
  % the recursion arrow: carry update feeds the recursive call
  \draw[flow, simBlue, densely dashed] (carry.west) to[out=180,in=270]
    node[below, font=\scriptsize, simBlue, pos=0.4]{recurse on $a'$} (lhs.south);
\end{tikzpicture}
\caption[One unrolling step of \texttt{iterate\_while}]{The
\code{iterate\_while} small-step rule (\cref{app:calc:iterate}), drawn as the
single operation it performs. The predicate $p(a)$ gates the firing. On
\textcolor{simRust}{no}, the loop stops, surfacing the current carry as
\code{final\_state} with an empty trajectory. On \textcolor{simGreen}{yes}, one
step \code{f(a)} runs and \emph{splits in two}: its \code{.state} field becomes
the next carry $a'$ (threaded forward, and fed back into the recursive call,
dashed), while its extras $e$ are prepended to the growing
\code{trajectory}. This is the formal version of the carry-update-plus-append
picture of \cref{ch:fold}: every iterative algorithm in the book is this rule
fired repeatedly, with ``step'' replaced by the algorithm's per-step kernel.}
\label{fig:appA:unroll}
\end{figure}

The combinator has a sugared special case,
\lstinline[style=l4style]|iterate_while_pure init cont step|, for steps with no
extras: it is \lstinline[style=l4style]|(iterate_while init cont (\a -> { state: step a })).final_state|,
the general rule with $\code{extras} = ()$ and the trajectory always empty
(\cref{ch:fold}). Total correctness---that the loop halts---is an obligation on
the \emph{algorithm} that uses the rule (a bounded \code{max\_it} folded into the
carry, or a convergence proof), not a guarantee of the rule itself; the type
promises nothing about termination.

\subsection{Demand-driven pruning: the signature move}
\label{app:calc:pruning}

The calculus has one rule unlike anything in an ordinary operational
semantics, and it is the most important rule in the appendix. It formalizes the
demand-driven pruning of \cref{ch:fold}: an output that no consumer reads is never
computed. It is the only rule here with a real premise.

The calculus is a \emph{graph-evaluation language}: operators and aggregate
operations produce \emph{records} of outputs, and at solve time only the outputs
that some root consumer transitively reaches are computed. Observation propagates
structurally through bindings---if a body reads field $l_k$, then $l_k$ is
observed in the producing record; if the body never references $l_k$ (uses a
\code{\_} placeholder, or simply omits it), $l_k$ is unobserved. The pruning rule
is then an equivalence:

\[
\dfrac{\;\text{output } l_k \text{ of } \textsf{op} \text{ is not observed in } \textsf{body}\;}
{\;
  \textsf{let}\ \{l_1\!:\!x_1, \dots, l_k\!:\!\_, \dots, l_n\!:\!x_n\}
    = \textsf{op}(\bar{e})\ \textsf{in}\ \textsf{body}
  \;\equiv\;
  \textsf{let}\ \{l_i\!:\!x_i\}_{i \neq k} = \textsf{op}_{\neg k}(\bar{e})\ \textsf{in}\ \textsf{body}
\;}
\]

\noindent where $\textsf{op}_{\neg k}$ is the subgraph of $\textsf{op}$ that does
\emph{not} compute the $l_k$-output---formally, the dependency-closure of the
remaining outputs over $\textsf{op}$'s body. \emph{Intuition.} The premise (above
the line) is a fact about \emph{use}: nothing downstream reads field $l_k$. Given
that fact, the conclusion (below the line) licenses rewriting the producer to a
strictly smaller operator that omits $l_k$ entirely, along with every
sub-computation that fed \emph{only} into $l_k$. This is dead-code elimination,
lifted from a compiler's view of a function body to the evaluation of the whole
computation graph. The rule reads as a subgraph-derivation: from the full graph
and a non-observation premise, derive the pruned subgraph.

\Cref{fig:appA:pruning} draws the derivation as the graph transformation it is:
the full dependency graph on the left, and on the right the subgraph that survives
when one output is unread, with the dead nodes and edges erased.

\begin{figure}[t]
\centering
\begin{tikzpicture}[font=\footnotesize, node distance=6mm]
  % ============ LEFT: full graph (all outputs observed) ============
  \begin{scope}
    \node[font=\small\bfseries, simBlue] at (0,2.7) {full graph: every output observed};
    \node[data, minimum width=14mm] (in) at (0,1.8) {inputs $\bar e$};
    \node[opbox, minimum width=12mm] (core) at (-1.2,0.7) {core};
    \node[opbox, minimum width=14mm] (extra) at (1.4,0.7) {extra calc};
    \draw[flow] (in.south) -- (core.north);
    \draw[flow] (in.south) -- (extra.north);
    \node[product, minimum width=12mm] (o1) at (-1.2,-0.7) {$l_1$};
    \node[product, fill=simBgGold, draw=simGold, minimum width=12mm] (ok) at (1.4,-0.7) {$l_k$};
    \draw[flow] (core.south) -- (o1.north);
    \draw[flow] (extra.south) -- (ok.north);
    % both consumed
    \draw[refedge] (o1.south) -- ++(0,-0.55) node[below, font=\scriptsize]{read};
    \draw[refedge] (ok.south) -- ++(0,-0.55) node[below, font=\scriptsize, simGold!80!black]{read};
  \end{scope}
  \draw[simGray!50, thick] (3.4,-1.6) -- (3.4,2.9);
  % the rule arrow with the premise
  \node[align=center, font=\scriptsize, simBlue] at (4.5,1.0)
    {premise:\\$l_k$ \emph{not}\\observed\\[2pt]$\Longrightarrow$};
  % ============ RIGHT: pruned subgraph ============
  \begin{scope}[xshift=8.4cm]
    \node[font=\small\bfseries, simGreen] at (0,2.7) {pruned subgraph: \code{op}$_{\neg k}$};
    \node[data, minimum width=14mm] (in2) at (0,1.8) {inputs $\bar e$};
    \node[opbox, minimum width=12mm] (core2) at (-1.2,0.7) {core};
    % the extra calc is dead (greyed)
    \node[opbox, draw=simGray!40, fill=white, text=simGray!55, minimum width=14mm]
      (extra2) at (1.4,0.7) {\textcolor{simGray!55}{extra calc}};
    \draw[flow] (in2.south) -- (core2.north);
    \draw[flow, simGray!35] (in2.south) -- (extra2.north);
    \node[product, minimum width=12mm] (o12) at (-1.2,-0.7) {$l_1$};
    \node[product, draw=simGray!40, fill=white, text=simGray!55, minimum width=12mm]
      (ok2) at (1.4,-0.7) {\textcolor{simGray!55}{$l_k$}};
    \draw[flow] (core2.south) -- (o12.north);
    \draw[flow, simGray!35] (extra2.south) -- (ok2.north);
    \draw[refedge] (o12.south) -- ++(0,-0.55) node[below, font=\scriptsize]{read};
    \node[font=\scriptsize, simGray!60] at (1.4,-1.45) {pruned};
  \end{scope}
\end{tikzpicture}
\caption[Demand-driven pruning as graph dead-code elimination]{The pruning rule
(\cref{app:calc:pruning}) as a subgraph derivation. \emph{Left}: the full
operator graph---two outputs $l_1$ and $l_k$, each produced by its own
sub-computation, both read by a consumer. \emph{Right}: the premise ``$l_k$ is
not observed'' fires the rule, deriving $\textsf{op}_{\neg k}$---the subgraph that
omits $l_k$ \emph{and} the \emph{extra calc} node that fed only into it (greyed,
erased). The surviving graph computes exactly the observed outputs. This is the
calculus's signature move: no verbosity flag, no conditional branch, no Writer
monad---one operator definition, and what materializes is decided by what the
consumer reads. Applied inside \code{iterate\_while}, it is why an unread
\code{trajectory} costs nothing (\cref{ch:fold}).}
\label{fig:appA:pruning}
\end{figure}

This single rule subsumes three mechanisms other systems need separately: a
\emph{Writer monad} (logging/metric accumulation), a \emph{verbosity flag}
(``if verbose, also compute the residual''), and a \emph{conditional branch}
around each optional output. In the calculus none of those exist as separate
concepts: every operator unconditionally declares all its outputs as record
fields, every consumer reads what it needs, and pruning removes the rest. An
algorithm is written \emph{once} and is correct whether or not its optional
outputs are consumed---which is why no solver chapter in this book presents a
``debug version'' beside a ``fast version'' (\cref{ch:fold}).

\begin{keyidea}
Demand-driven pruning is the calculus's one genuinely unusual rule, and its whole
payoff. ``Compute the residual history for monitoring'' and ``skip it for speed''
are not two algorithms but \emph{one}, distinguished only by what a downstream
consumer reads. The rule is graph dead-code elimination given as a formal
equivalence: a non-observation premise licenses replacing a producer by the
subgraph of its observed outputs. It is what lets the framework write each solver
once and use it for both diagnostics and production.
\end{keyidea}

\section{Typing in one paragraph}
\label{app:calc:typing}

The grammar and reduction rules are the load-bearing content; the type system is
standard typed-$\lambda$ and we state only its shape. A typing judgment
$\Gamma \vdash e : \tau$ reads ``in context $\Gamma$ (a set of variable-to-type
assumptions), the term $e$ has type $\tau$.'' The selected rules are the usual
ones---a variable has the type the context assigns it; an application
$e_1\,e_2$ requires $e_1 : \tau_1 \to \tau_2$ and $e_2 : \tau_1$ and yields
$\tau_2$; a record literal is typed field-by-field; a spread preserves the record
type; \code{return} lifts a $\tau$ to \code{Sim $\tau$}; and \code{>>=} sequences
\code{Sim $\alpha$} with $\alpha \to \code{Sim $\beta$}$ to give \code{Sim $\beta$}:

\begin{gather*}
\dfrac{}{\Gamma, x{:}\tau \vdash x : \tau}\;(\textsf{var})
\qquad
\dfrac{\Gamma \vdash e_1 : \tau_1 \to \tau_2 \quad \Gamma \vdash e_2 : \tau_1}
      {\Gamma \vdash e_1\,e_2 : \tau_2}\;(\textsf{app})
\\[6pt]
\dfrac{\Gamma \vdash m : \textsf{Sim}\,\alpha \quad \Gamma, x{:}\alpha \vdash k : \textsf{Sim}\,\beta}
      {\Gamma \vdash m \mathbin{\texttt{>>=}} \lambda x.k : \textsf{Sim}\,\beta}\;(\textsf{bind})
\end{gather*}

\noindent The one place the type system does real, distinctive work is
\emph{shapes}: a primitive's signature carries symbolic tensor shapes
(\lstinline[style=l4style]|axpy : Scalar -> Tensor[s] -> Tensor[s] -> Tensor[s]|),
and the requirement that two shapes ``match'' is a side condition discharged by a
separate shape solver, not by these rules. That discipline---congruence groups,
rank wildcards, the named-axis algebra---is the entire subject of
\cref{app:shapes}. Linearity, the second distinctive ingredient, is tracked only
to enforce the simulator-state-versus-operator-parameter ownership split
(\cref{app:calc:ownership}); it rides on the \code{Sim} monad's \code{get}/\code{put}
discipline and the \code{!} annotation rather than on per-tensor annotations.

\section{Evaluation order and confluence}
\label{app:calc:confluence}

A final, reassuring property. The reduction rules above can often fire in more
than one place in a term---in \cref{fig:appA:reduction}, we could have evaluated a
sub-expression before or after another. Does the order matter? \emph{No, beyond
what data dependency forces.} Because every function is pure---its output depends
only on its inputs, with no side effects (\cref{ch:language})---the value of an
expression depends only on the values of its sub-expressions, never on the order
in which those sub-expressions were reduced. Two redexes that do not depend on
each other may be fired in either order, or simultaneously, and reach the same
value; this is the calculus's \emph{confluence}. The only ordering the dynamics
respects is the one the data edges demand: a step that consumes the residual must
wait for the step that produces it.

Two qualifications keep this honest. First, demand-driven pruning
(\cref{app:calc:pruning}) means some redexes are never fired at all---an unobserved
output's sub-computation is pruned rather than reduced---but this does not disturb
the observed result, since a pruned node by definition feeds nothing observed.
Second, the calculus is deterministic only \emph{up to floating-point
reassociation}: when a load-bearing numerical trick fixes a particular reduction
or accumulation order for a property it buys (determinism, a condition number, IEEE
compliance), that order is part of the algorithm and is preserved as an explicit
algebraic claim, not freely reassociated (the load-bearing-versus-transparent
distinction of the main text). With that caveat, the program text \emph{is} the
dependency graph: there is no hidden sequencing, only data flowing along visible
edges, exactly as \cref{ch:language} promised.

\summarysection
This appendix is the formal companion to the main text's calculus chapters. The
\emph{grammar} fixes the legal syntax in two layers: a \emph{type language} $\tau$
(\cref{app:calc:types}: ground scalars, \code{Tensor[$\sigma$]}, records, tuples,
functions, \code{Sim $\tau$}, \code{Op[$\cdot$]}, and the \code{!} shareable
annotation) and a \emph{term language} $e$ (\cref{app:calc:terms}: the lambda core,
records, tuples, conditionals, primitive and operator application, the monad and
state effects, and the single loop \code{while\_loop}). The \emph{small-step
operational semantics} (\cref{app:calc:reduction}) is the dynamics, a reduction
relation $e \to e'$ presented as inference rules: $\beta$ and \code{let} for the
lambda core; projection, destructuring, and spread for records; the three monad
laws and the state-effect rewrites for orchestration; one substitution rule for
operator \code{apply}; the recursive unrolling rule for \code{iterate\_while}; and
the demand-driven pruning rule. The grammar says what you may write; the reduction
rules say what it means; and the whole computational content is the handful of
rules collected here. The signature move is \emph{pruning}: a non-observation
premise licenses replacing a producer by the subgraph of its observed outputs, so
diagnostics and production are one program. Because every function is pure, the
calculus is confluent (\cref{app:calc:confluence})---evaluation order beyond data
dependency is unobservable, deterministic up to deliberate floating-point
reassociation---so the program text is exactly the dependency graph. The ownership
split (\cref{app:calc:ownership}) and the standard typing judgment
(\cref{app:calc:typing}) round out the static side; tensor shapes are handed off
to \cref{app:shapes} and primitive $\delta$-rules and algebraic laws to
\cref{app:laws}.

\furtherreading
The calculus is a small dialect of the typed lambda calculus extended with a state
monad and one iteration combinator; the two most readable entry points are the
standard reference for the language whose grammar and reduction conventions we
borrow~\citep{peytonjones2003}, and Wadler's exposition of how a pure calculus
nonetheless models stateful computation through monads~\citep{wadler1995monads},
which underwrites both the \code{Sim} effects of \cref{app:calc:state} and the
\code{do}-notation desugaring of \cref{app:calc:monad}. Readers who want the
operational-semantics machinery itself---inference rules, $\beta$-reduction,
substitution, confluence---in its original setting will find it in any text on the
lambda calculus or programming-language theory. The two companion appendices carry
the pieces this one hands off: \cref{app:shapes} formalizes the shape grammar
$\sigma$ and its congruence algebra, and \cref{app:laws} collects the primitive
$\delta$-rules and the algebraic equational laws the reduction chains of
\cref{ch:rotations} appeal to.

\chapter{Types, Shapes, and the Shape Algebra}
\label{app:shapes}

\chapterorient{\Cref{ch:tensors} fixed the vocabulary of tensors and shapes, and
\cref{ch:shapes} taught the reading discipline: a shape is a type, named groups
assert congruence, scalar promotion lets a real number quietly join a complex
computation. Those chapters argued by example and picture. This appendix supplies
the \emph{formal} object underneath them. We write the shape language as a
grammar; we state named-group resolution as a single unification rule; we give the
type-and-shape judgment $\Gamma \vdash e : \tau$ as a small set of inference rules;
we collect the primitive shape contracts into one reference table; and we cast
scalar promotion as a subtyping rule \lstinline[style=l4style]|real <: complex|.
Nothing here is new physics or new notation---it is the precise version of the
rules the main text reads informally, the companion to \cref{app:calculus}'s
grammar and reduction rules and the home for the equational laws of
\cref{app:laws}. Read it as reference: the rule you reach for when a shape
contract or a typing question must be settled exactly rather than by eye.}

\section{What this appendix formalizes}
\label{sec:b-scope}

\Cref{ch:tensors,ch:shapes} carry the conceptual load; this appendix carries the
formal load behind them. The split is deliberate. The main text earns its rules on
worked vectors and matrices and draws each contract as a picture, so that a reader
learns to \emph{read} a shape the way they read a type. Here we write those same
rules as grammar productions, inference rules, and a reference table, so that a
reader can \emph{settle} a shape question exactly---decide whether two shapes are
congruent, whether a call type-checks, which group binds where---without appeal to
intuition.

Four objects organize the appendix, each the formal counterpart of a main-text
idea:

\begin{itemize}[leftmargin=1.4em, nosep]
\item \textbf{The shape language} (\cref{sec:b-shapelang})---a BNF grammar for the
shape expression $\sigma$, with its literals, axis variables, algebraic
combinations, the rank-wildcard \lstinline[style=l4style]|...|, and the named-group
binding and use forms. This is the formal version of \cref{ch:tensors}'s bracket
notation and \cref{ch:shapes}'s named groups.
\item \textbf{Named-group resolution} (\cref{sec:b-resolve})---the precise
unification rule by which a use \lstinline[style=l4style]|$S| is matched against a
binding \lstinline[style=l4style]|(S: ...)|, the formal version of
\cref{ch:shapes}'s rank-agnostic congruence.
\item \textbf{The type-and-shape judgment} (\cref{sec:b-judgment})---the inference
rules of $\Gamma \vdash e : \tau$, including the primitive-op shape-contract rule
that ties shapes into typing.
\item \textbf{The primitive shape contracts} (\cref{sec:b-contracts}) and
\textbf{scalar promotion} (\cref{sec:b-promotion})---a reference table of the
primitive signatures, and the subtyping rule
\lstinline[style=l4style]|real <: complex|.
\end{itemize}

\begin{keyidea}
\textbf{Shapes are types; congruence is unification.} A tensor's shape is a
type-level expression, checked at every function boundary exactly as a type is. A
named shape group \lstinline[style=l4style]|(S: ...)| is a unification variable
over a run of axes of \emph{unknown rank}; a use \lstinline[style=l4style]|$S|
resolves by \emph{unifying} the labelled shape against the binding, modulo a small
equational theory on axis extents (commutative, associative $+$ and $\times$). The
whole shape discipline of this book is therefore one rule applied everywhere:
\emph{unify the named groups of a signature against the argument shapes, modulo
that equational theory, rank first.}
\end{keyidea}

\section{The shape language}
\label{sec:b-shapelang}

A tensor type is \lstinline[style=l4style]|Tensor[s]| for a shape expression
$\sigma$. A shape expression is a tuple of \emph{items}, each item an axis
extent---a literal, a named axis, an algebraic combination---or one of the two
rank-flexible constructs: the wildcard and the named group. We give the grammar,
then read each production.

\subsection{Grammar}
\label{sec:b-grammar}

A shape $\sigma$ is a bracketed sequence of \emph{group items} $\gamma$; a group
item is either a single axis extent or one of the rank-flexible forms.

\begin{lstlisting}[style=l4style]
sigma ::= n                      -- literal natural (a pinned extent, e.g. 3)
        | s                      -- axis variable (a named axis: N, m, H, ...)
        | sigma + sigma          -- axis arithmetic
        | sigma - sigma
        | sigma * sigma
        | sigma ^ n
        | [ gamma_1, ..., gamma_k ]  -- a shape: a sequence of group items

gamma ::= sigma                  -- one axis (literal, variable, or algebraic)
        | ...                    -- rank wildcard: zero or more unconstrained axes
        | (S: gamma_1, ..., gamma_j)  -- named group BINDING: binds S to this run
        | $S                     -- named group USE: a back-reference to S
\end{lstlisting}

\noindent The grammar is taken verbatim from the source calculus's
\lstinline[style=l4style]|DimExpr| algebra; \cref{ch:tensors,ch:shapes} use exactly
its shapes. Reading the productions:

\begin{description}[leftmargin=1.6em, style=nextline, nosep]
\item[Literal $n$.] A \emph{pinned} extent: an axis whose length is fixed in the
signature, like the $3$ in \lstinline[style=l4style]|Tensor[H, W, 3]| (a 3-vector
per cell). A call supplying any other extent at that position is rejected.
\item[Axis variable $s$.] A \emph{named} axis: a single letter ($N$, $m$, $n$,
$H$, $W$) standing for a size not fixed in the signature but required to
\emph{match} wherever the same name recurs. An axis variable is \emph{one} axis,
hence rank-bearing---this is the load-bearing fact behind the
\lstinline[style=l4style]|Tensor[N]| pitfall of \cref{ch:shapes}.
\item[Algebraic combination.] An extent built from others by $+,-,\times,{}^\wedge$.
These arise from reshapes and factorizations: flattening a height axis $h$ and a
width axis $w$ into one produces the extent \lstinline[style=l4style]|h*w|. The
matcher reasons about these \emph{modulo} a small algebra (\cref{sec:b-dimexpr}).
\item[{Shape tuple $[\gamma_1,\dots,\gamma_k]$.}] The shape itself: an ordered
sequence of group items. The \emph{number} of axes the tuple denotes is the
tensor's \emph{rank}---but a single \lstinline[style=l4style]|...| or
\lstinline[style=l4style]|(S: ...)| item may itself stand for \emph{several} axes,
so the count of \emph{items} is not in general the rank.
\item[Rank wildcard \code{...}.] Matches \emph{zero or more}
axes of any extent. It is the device that makes a signature rank-agnostic: it
declines to commit to how many axes fill its position.
\item[Named group, binding \code{(S: ...)}.] Binds the name
$S$ to the run of axes its body matches. The binding occurs \emph{exactly once},
at $S$'s first appearance.
\item[{Named group, use \code{\$S}.}] A \emph{back-reference}: it
asserts that the run of axes at this position is \emph{congruent} to the run bound
to $S$. The \lstinline[style=l4style]|$| sigil is what visually distinguishes a use
from a fresh binding or a single-axis variable.
\end{description}

\Cref{fig:b-shapealgebra} draws the shape algebra as labelled axis boxes: a row of
boxes for a tensor's axes, a spanning bracket for a named group's binding, and the
\lstinline[style=l4style]|$S| back-arrow for a use---the formal-side companion to
the hero figure of \cref{ch:shapes}.

\begin{figure}[t]
\centering
\begin{tikzpicture}[scale=1.0]
  % LEGEND-style: literal, variable, group, wildcard
  % Row of axis boxes with mixed item kinds: [ (S: a, ...), k ]
  \begin{scope}[shift={(0,1.3)}]
    % the named-group run: a leading pinned 'a' then a wildcard tail
    \draw[draw=simBlue, fill=simBgBlue, thick] (0,-0.4) rectangle (0.8,0.4);
    \node[font=\small, simInk] at (0.4,0) {$a$};
    \draw[draw=simBlue, fill=simBgBlue, thick] (1.0,-0.4) rectangle (1.8,0.4);
    \node[font=\footnotesize, simGray] at (1.4,0) {$\cdot$};
    \node[font=\scriptsize, simGray] at (2.15,0) {$\cdots$};
    \draw[draw=simBlue, fill=simBgBlue, thick] (2.5,-0.4) rectangle (3.3,0.4);
    \node[font=\footnotesize, simGray] at (2.9,0) {$\cdot$};
    % the free trailing axis k
    \draw[draw=simRust, fill=simBgRust, thick] (3.9,-0.4) rectangle (4.7,0.4);
    \node[font=\small, simInk] at (4.3,0) {$k$};
    % bracket binding the group S over the leading run
    \draw[decorate, decoration={brace, amplitude=4pt}, simBlue]
      (-0.02,0.55) -- (3.32,0.55)
      node[midway, above=4pt, font=\scriptsize\bfseries, simBlue]
      {binds \texttt{S} (a leading run: pinned \texttt{a}, then \texttt{...})};
    \node[font=\scriptsize\bfseries, simRust, above=4pt] at (4.3,0.55) {free \texttt{k}};
    \node[font=\scriptsize\ttfamily, simInk] at (-2.0,0)
      {\lstinline[style=l4style]|Tensor[(S: a, ...), k]|};
  \end{scope}
  % a second operand: $S then a free m  -- the USE row
  \begin{scope}[shift={(0,-0.6)}]
    \draw[draw=simTeal, fill=simBgGreen, thick] (0,-0.4) rectangle (0.8,0.4);
    \node[font=\small, simInk] at (0.4,0) {$a$};
    \draw[draw=simTeal, fill=simBgGreen, thick] (1.0,-0.4) rectangle (1.8,0.4);
    \node[font=\footnotesize, simGray] at (1.4,0) {$\cdot$};
    \node[font=\scriptsize, simGray] at (2.15,0) {$\cdots$};
    \draw[draw=simTeal, fill=simBgGreen, thick] (2.5,-0.4) rectangle (3.3,0.4);
    \node[font=\footnotesize, simGray] at (2.9,0) {$\cdot$};
    \draw[draw=simGold, fill=simBgGold, thick] (3.9,-0.4) rectangle (4.7,0.4);
    \node[font=\small, simInk] at (4.3,0) {$m$};
    \draw[decorate, decoration={brace, amplitude=4pt, mirror}, simTeal]
      (-0.02,-0.55) -- (3.32,-0.55)
      node[midway, below=4pt, font=\scriptsize\bfseries, simTeal]
      {uses \texttt{\$S} (must be congruent)};
    \node[font=\scriptsize\bfseries, simGold, below=4pt] at (4.3,-0.55) {free \texttt{m}};
    \node[font=\scriptsize\ttfamily, simInk] at (-2.0,0)
      {\lstinline[style=l4style]|Tensor[$S, m]|};
  \end{scope}
  % the congruence back-arrow
  \draw[-{Stealth[length=2.6mm]}, simRust, thick]
    (3.55,-0.85) .. controls (4.5,-0.1) and (4.5,0.7) .. (3.55,1.45)
    node[midway, right=1pt, font=\scriptsize\itshape, simRust, align=left]
    {congruent\\(same run)};
\end{tikzpicture}
\caption[The shape algebra: items, groups, and the congruence back-arrow]{The
shape language as labelled axis boxes. \emph{Top (blue/rust):} the binding shape
\lstinline[style=l4style]|Tensor[(S: a, ...), k]|---a named group \texttt{S}
spanning a leading run (a pinned axis \texttt{a} followed by a wildcard
\texttt{...} of any number of further axes), then a single \emph{free} trailing
axis \texttt{k} outside the group. \emph{Bottom (teal/gold):} a second operand
\lstinline[style=l4style]|Tensor[$S, m]|, whose leading run is a \emph{use}
\lstinline[style=l4style]|$S| and whose own trailing axis \texttt{m} is free and
independent of \texttt{k}. The rust back-arrow is the congruence assertion the
\lstinline[style=l4style]|$S| carries: \emph{my leading run is the same run that
bound \texttt{S}}. Group items may be literals, variables, algebraic extents,
wildcards, or nested groups; a group splits a shape into a named congruent part
and a free part.}
\label{fig:b-shapealgebra}
\end{figure}

\subsection{The DimExpr equational theory}
\label{sec:b-dimexpr}

Two shapes that differ on the page may denote the same run of extents. A reshape
that merges axes $h$ and $w$ yields the extent \lstinline[style=l4style]|h*w|; an
operation consuming the flattened tensor must recognize
\lstinline[style=l4style]|h*w| and \lstinline[style=l4style]|w*h| as the same
length. Shape matching is therefore \emph{not} raw character comparison; it is
comparison modulo a small equational theory on extent expressions---the
\lstinline[style=l4style]|DimExpr| theory.

\begin{definition}[DimExpr equational theory]
\label{def:b-dimexpr}
Two axis extents are \emph{equal} when they are equal as elements of the free
commutative semiring on the axis variables, with $+$ and $\times$ the semiring
operations: that is, modulo the axioms

\begin{center}
\begin{tabular}{@{}ll@{}}
\toprule
\emph{commutativity} & $a + b = b + a$,\qquad $a \times b = b \times a$ \\
\emph{associativity} & $(a+b)+c = a+(b+c)$,\qquad $(a\,b)\,c = a\,(b\,c)$ \\
\emph{distributivity} & $a\,(b+c) = a\,b + a\,c$ \\
\emph{identities} & $a + 0 = a$,\qquad $a \times 1 = a$,\qquad $a \times 0 = 0$ \\
\emph{powers} & $a^{m}\,a^{n} = a^{m+n}$,\qquad $a^{1}=a$ \\
\bottomrule
\end{tabular}
\end{center}

\noindent Subtraction $a - b$ is admitted as the formal inverse of addition
(extents arising from a reshape are assumed nonnegative; the theory does not
attempt to prove nonnegativity). Two shapes $[\gamma_1,\dots,\gamma_k]$ and
$[\delta_1,\dots,\delta_k]$ of \emph{equal rank} are equal when $\gamma_i = \delta_i$
in this theory for every $i$.
\end{definition}

In practice the matcher normalizes each extent expression to a canonical
sum-of-monomials form and compares; \lstinline[style=l4style]|h*w| and
\lstinline[style=l4style]|w*h| both normalize to the monomial $hw$, so they match,
while \lstinline[style=l4style]|h*w| and \lstinline[style=l4style]|h+w| do not.
The theory is deliberately weak---it decides the \emph{algebraic} coincidences a
reshape genuinely produces, and nothing subtler.

\begin{notationbox}
The equational theory is a fact the matcher \emph{verifies}; congruence is a
contract you \emph{assert}. When two operands must share a shape, say so with a
group (\lstinline[style=l4style]|(S: ...)| / \lstinline[style=l4style]|$S|) and let
the matcher enforce it; reserve the extent algebra
(\lstinline[style=l4style]|h*w| and friends) for the positions where a genuine
reshape \emph{produced} a composite extent. Writing two operands as
\lstinline[style=l4style]|Tensor[h*w]| and hoping the algebra makes them agree is
backwards: the algebra equates \emph{extents at matching positions}, it does not
discover that two whole shapes were meant to be congruent.
\end{notationbox}

\section{Named-group resolution}
\label{sec:b-resolve}

We can now state, precisely, what \cref{ch:shapes} drew as a back-arrow: how a use
\lstinline[style=l4style]|$S| is resolved against a binding
\lstinline[style=l4style]|(S: ...)|. The resolution is a \emph{unification} of
axis runs, with one wrinkle---a named group spans a \emph{run} of axes of unknown
length, so the unifier must split a concrete shape into the part the group claims
and the part left free.

\subsection{The resolution rule}
\label{sec:b-resolve-rule}

A signature is checked left to right against the argument shapes. The first
position labelled by a group binds it; later uses are tested.

\begin{definition}[Named-group resolution]
\label{def:b-resolve}
Let a signature contain a binding \lstinline[style=l4style]|(S: p)| (where the
\emph{pattern} $p$ is a sequence of group items, possibly containing a wildcard
\lstinline[style=l4style]|...|) and one or more uses \lstinline[style=l4style]|$S|.
Resolution proceeds in two steps.

\begin{description}[leftmargin=1.6em, style=nextline, nosep]
\item[Bind.] Match the pattern $p$ against the corresponding run of the
\emph{first} argument shape that labels $S$. If $p$ is a pure wildcard
\lstinline[style=l4style]|(S: ...)|, the run is the \emph{entire} shape at that
position; if $p$ pins leading or trailing axes (e.g.
\lstinline[style=l4style]|(S: a, ...)|), those pinned axes must match by the
\lstinline[style=l4style]|DimExpr| theory (\cref{def:b-dimexpr}) and the wildcard
absorbs the remainder. The matched run---a concrete sequence of extents
$\hat{S} = [\hat{e}_1,\dots,\hat{e}_r]$ of definite rank $r$---becomes the
\emph{value} of $S$.
\item[Use.] At each later \lstinline[style=l4style]|$S|, require the run of axes at
that position to be \emph{congruent} to $\hat{S}$: \textbf{(i)} the same rank $r$,
and \textbf{(ii)} extent $\hat{e}_i$ at axis $i$ for every $i$, equality decided by
the \lstinline[style=l4style]|DimExpr| theory. Rank is checked \emph{first}: a rank
mismatch fails immediately, before any extent is compared.
\end{description}

\noindent A group with both a pinned part and a free part splits each labelled
shape into ``the run $S$ claims'' and ``the free axes outside it''; only the
claimed run participates in congruence, while the free axes (\texttt{k},
\texttt{m} in \cref{fig:b-shapealgebra}) are unconstrained across operands.
\end{definition}

Read as a unification rule, the binding/use machinery is one inference:

\[
\resizebox{\linewidth}{!}{$\displaystyle
\dfrac{
  \begin{array}{c}
  \text{$\hat{S}$ is the run matched by \lstinline[style=l4style]|(S: p)| in the binding position}
  \end{array}
  \qquad
  \text{rank}(\tau) = \text{rank}(\hat{S})
  \qquad
  \tau \;=_{\text{DimExpr}}\; \hat{S}
}{
  \text{a shape $\tau$ at a use position \lstinline[style=l4style]|$S| resolves (is congruent)}
}
\quad(\textsc{congr})
$}
\]

\noindent and the dual---a rank mismatch---is the rejection rule:

\[
\dfrac{\text{rank}(\tau) \neq \text{rank}(\hat{S})}
      {\text{$\tau$ at \lstinline[style=l4style]|$S| is rejected (rank mismatch, no extent compared)}}
\quad(\textsc{congr-fail})
\]

This is exactly why \lstinline[style=l4style]|Tensor[N]| (\cref{ch:shapes}) can
never match a rank-$2$ argument: a single axis variable $N$ is a run of rank $1$,
so \textsc{congr-fail} fires on the rank check before $N$'s extent is ever
considered. \Cref{fig:b-unify} draws the two-step resolution as a matched
binding-against-use, the formal version of the boundary gate of
\cref{ch:shapes}.

\begin{figure}[t]
\centering
\begin{tikzpicture}[node distance=8mm and 12mm]
  % the binding shape on the left
  \node[data, text width=26mm] (bind)
    {\textbf{binding}\\\lstinline[style=l4style]|(S: ...)|\\matched run\\$\hat{S}=[3,4]$};
  % the unify gate
  \node[stage, right=16mm of bind, minimum height=24mm, text width=20mm] (gate)
    {\textbf{unify}\\$\$S \stackrel{?}{=} \hat{S}$};
  % two use candidates
  \node[data, right=16mm of gate, yshift=11mm, text width=30mm] (ok)
    {use \lstinline[style=l4style]|$S| $=[3,4]$\\rank $2$, extents agree};
  \node[data, right=16mm of gate, yshift=-11mm, text width=30mm] (bad)
    {use \lstinline[style=l4style]|$S| $=[12]$\\rank $1\neq 2$};
  \draw[flow] (bind) -- node[above, font=\scriptsize, simGray]{bind $\hat{S}$} (gate);
  \draw[-{Stealth[length=2.4mm]}, simGreen, thick] (gate.east|-ok) -- (ok)
    node[midway, above, font=\scriptsize, simGreen]{congruent};
  \draw[-{Stealth[length=2.4mm]}, simRust, thick, densely dashed] (gate.east|-bad) -- (bad)
    node[midway, below, font=\scriptsize, simRust]{reject};
  % annotation: rank first
  \node[font=\scriptsize\itshape, simGray, below=7mm of gate, text width=64mm,
        align=center]
    {\textbf{step 1} bind $\hat{S}$ from the first labelled run;\quad
     \textbf{step 2} at each use, check rank \emph{first}, then extents modulo
     \lstinline[style=l4style]|DimExpr|. The lower candidate fails on rank
     ($1\neq 2$) before any extent is compared---even though $3\cdot4=12$.};
\end{tikzpicture}
\caption[Named-group resolution as unification]{Resolving a use
\lstinline[style=l4style]|$S| against a binding \lstinline[style=l4style]|(S: ...)|.
The binding's first labelled run fixes $\hat{S}$ (here $[3,4]$, rank $2$). Each
later use is \emph{unified} against $\hat{S}$: rank is checked first (a mismatch is
rejected on the spot, lower path), then extents axis-by-axis modulo the
\lstinline[style=l4style]|DimExpr| equational theory (\cref{def:b-dimexpr}, upper
path). The flattened candidate $[12]$ is rejected on \emph{rank} even though its
single extent equals the product $3\cdot4$ of the bound run: congruence demands the
same number of axes, not merely the same total size. This is the precise content of
\cref{ch:shapes}'s back-arrow and boundary gate.}
\label{fig:b-unify}
\end{figure}

\subsection{Operator shapes: range-first}
\label{sec:b-operators}

An operator whose domain and range shapes differ names \emph{two} groups, a range
group $R$ and a domain group $D$, written \emph{range-first} (matching the
$m \times n$ matrix convention, rows then columns):

\begin{lstlisting}[style=l4style]
LinOp[(R: ...), (D: ...)]                                   -- general operator
apply_linop :: LinOp[(R: ...), (D: ...)] -> Tensor[$D] -> Tensor[$R]
\end{lstlisting}

\noindent The operand is a use of the \emph{domain} group
(\lstinline[style=l4style]|$D|); the result is a use of the \emph{range} group
(\lstinline[style=l4style]|$R|). Resolution is \cref{def:b-resolve} applied to two
groups at once: the operator value binds both $R$ and $D$, the operand is checked
against $\hat{D}$, the result is congruent to $\hat{R}$. An \emph{endomorphism}
(domain $\equiv$ range) binds one group and uses it for both:

\begin{lstlisting}[style=l4style]
LinOp[(S: ...), $S]                                        -- square / endomorphic
apply_linop :: LinOp[(S: ...), $S] -> Tensor[$S] -> Tensor[$S]
\end{lstlisting}

\noindent which is the signature a Krylov solver wants: feed it a vector of shape
$S$, get back a vector of shape $S$ you can feed it again
(\cref{ch:matrixfree}). At the lowest layers, where Palace's operators act on
genuine flat dof-vectors of a definite length, the concrete rank-$1$ spelling
\lstinline[style=l4style]|LinearOperator[M, N]| over
\lstinline[style=l4style]|Tensor[N]| is faithful and is kept; the group form is the
rank-agnostic calculus rendering used at the higher layers.

\subsection{Congruence groups versus named axes of fixed meaning}
\label{sec:b-fixedaxes}

The resolution rule applies only to \emph{congruence groups}---the
\lstinline[style=l4style]|(S: ...)| / \lstinline[style=l4style]|$S| variables. A
shape bracket holds a second, entirely different kind of name, and the two must not
be confused. A \emph{named axis of fixed meaning} is a single axis whose letter
denotes one specific, fixed quantity---it is \emph{not} a unification variable and
\emph{not} a back-reference. The book's recurring family is the element-local
vocabulary of the matrix-free kernel (\cref{ch:matrixfree}):

\begin{description}[leftmargin=1.6em, style=nextline, nosep]
\item[\texttt{E}] element count---how many mesh elements the operator runs over.
\item[\texttt{L}] local dofs per element, $(p+1)^d$ for a degree-$p$ element.
\item[\texttt{P}] quadrature points per element.
\item[\texttt{C}] value (or derivative) components carried at a quadrature point.
\item[\texttt{G}] geometry-factor components stored per quadrature point.
\end{description}

\noindent with the three rank-structured shapes
\lstinline[style=l4style]|Tensor[(E, L)]|, \lstinline[style=l4style]|Tensor[(E, P, C)]|,
and \lstinline[style=l4style]|Tensor[(E, P, G)]|. Each letter has a definite role
fixed by the finite-element construction; none is rank-agnostic, none is a
back-reference. A fixed-meaning axis answers \emph{what is this axis}; a congruence
group answers \emph{does this shape match that one}. The two are governed by
different rules: a fixed-meaning axis is matched like an ordinary axis variable
(its letter must agree wherever it recurs, at rank $1$), whereas a congruence group
is resolved by \cref{def:b-resolve} over a run of unknown rank. The flat global
dof-axis $N$ is a third species again---a genuine rank-$1$
\lstinline[style=l4style]|Tensor[N]|, neither fixed-meaning element-local nor a
congruence group---and the gather/scatter pair $G,G^{\mathsf T}$ of
\cref{ch:matrixfree} is precisely the boundary between the flat-$N$ world and the
fixed-meaning \lstinline[style=l4style]|(E, L)| world.

\section{The type-and-shape judgment}
\label{sec:b-judgment}

Shapes do not float free; they ride on \emph{types}, and a tensor type
\lstinline[style=l4style]|Tensor[s]| is a type former carrying a shape. The
calculus's typing judgment is the standard typed-$\lambda$ form,
\[
  \Gamma \vdash e : \tau,
\]
read ``under context $\Gamma$ (a finite map from variables to types), the
expression $e$ has type $\tau$.'' The type $\tau$ ranges over the ground scalar
types, tensor types, records, tuples, functions, the simulator monad
\lstinline[style=l4style]|Sim t|, and operator types
\lstinline[style=l4style]|Op[t_in -> t_out]| (the full grammar is
\cref{app:calculus}). The shape discipline of the previous sections enters typing
at exactly one place: the rule for primitive-operator application, where the op's
declared shape contract must match the argument shapes modulo the equational theory.

\subsection{The core rules}
\label{sec:b-coretyping}

The structural rules are the textbook ones. We give variable, application,
abstraction, record formation and spread, and the two monadic rules; the rest are
mechanical.

\begin{description}[leftmargin=1.4em, style=nextline]
\item[Variable.] A variable has the type the context assigns it.
\[
  \dfrac{x : \tau \,\in\, \Gamma}{\Gamma \vdash x : \tau}\quad(\textsc{var})
\]

\item[Application and abstraction.] Function application demands the argument's
type meet the function's domain; abstraction discharges a hypothesis.
\[
  \dfrac{\Gamma \vdash e_1 : \tau_1 \to \tau_2 \qquad \Gamma \vdash e_2 : \tau_1}
        {\Gamma \vdash e_1\,e_2 : \tau_2}\quad(\textsc{app})
  \qquad\qquad
  \dfrac{\Gamma,\, x : \tau_1 \vdash e : \tau_2}
        {\Gamma \vdash \lambda(x{:}\tau_1).\,e \;:\; \tau_1 \to \tau_2}\quad(\textsc{abs})
\]

\item[Records and spread.] A record literal is typed field-by-field; a spread
(functional update) preserves the record type while replacing one field's value
with one of the same type.
\[
  \dfrac{\Gamma \vdash e_i : \tau_i \ \ (\text{each } i)}
        {\Gamma \vdash \{\, l_1{:}\,e_1,\dots,l_n{:}\,e_n \,\} : \{\, l_1{:}\,\tau_1,\dots,l_n{:}\,\tau_n \,\}}
  \quad(\textsc{record})
\]
\[
  \dfrac{\Gamma \vdash e : \{\, l_1{:}\,\tau_1,\dots,l_n{:}\,\tau_n \,\}
         \qquad \Gamma \vdash w : \tau_k}
        {\Gamma \vdash \{\, \dots e,\ l_k{:}\,w \,\} : \{\, l_1{:}\,\tau_1,\dots,l_n{:}\,\tau_n \,\}}
  \quad(\textsc{spread})
\]

\item[Monad: return and bind.] \lstinline[style=l4style]|return| injects a pure
value into the simulator monad; bind sequences two \lstinline[style=l4style]|Sim|
actions, threading the state (the ownership discipline of \cref{ch:monad}).
\[
  \dfrac{\Gamma \vdash e : \tau}{\Gamma \vdash \texttt{return}\ e : \texttt{Sim}\ \tau}
  \quad(\textsc{return})
  \qquad\quad
  \dfrac{\Gamma \vdash m : \texttt{Sim}\ \alpha \qquad \Gamma,\, x{:}\,\alpha \vdash k : \texttt{Sim}\ \beta}
        {\Gamma \vdash m \mathbin{>\!\!>\!\!=} \lambda x.\,k \;:\; \texttt{Sim}\ \beta}
  \quad(\textsc{bind})
\]
\end{description}

\noindent These rules are shape-agnostic: a shape appears only \emph{inside} a
tensor type and is carried along opaquely by \textsc{var}, \textsc{app}, and the
rest. Where shapes acquire force is the next rule.

\subsection{The primitive shape-contract rule}
\label{sec:b-primrule}

Each primitive operator $\mathit{op}$ carries a declared signature---a sequence of
argument tensor (or scalar) types and a result type, with shapes drawn from
\cref{sec:b-shapelang}, possibly mentioning groups. Applying $\mathit{op}$ to
arguments type-checks when each argument's shape matches the contract's
corresponding shape \emph{modulo the equational theory and the group resolution of
\cref{def:b-resolve}}, after which the result wears the contract's result shape with
the resolved groups substituted in.

Write the contract of $\mathit{op}$ as
\lstinline[style=l4style]|op :: Tensor[s_1] -> ... -> Tensor[s_k] -> Tensor[s_r]|,
possibly with scalar arguments interspersed and with groups shared among the
$\sigma$'s. Let $\theta$ be the substitution that resolves every group
(\cref{def:b-resolve}) against the argument shapes. The rule is:

\[
\dfrac{
  \begin{array}{c}
  \mathit{op} :: \texttt{Tensor[}\sigma_1\texttt{]} \to \cdots \to \texttt{Tensor[}\sigma_k\texttt{]} \to \texttt{Tensor[}\sigma_r\texttt{]}
  \qquad
  \Gamma \vdash e_i : \texttt{Tensor[}\tau_i\texttt{]} \ \ (\text{each } i) \\[2pt]
  \tau_i \;=_{\text{DimExpr}}\; \theta(\sigma_i) \ \ (\text{each } i)
  \qquad
  \theta \text{ resolves every group of the contract (\cref{def:b-resolve})}
  \end{array}
}{
  \Gamma \vdash \mathit{op}(e_1,\dots,e_k) : \texttt{Tensor[}\theta(\sigma_r)\texttt{]}
}\quad(\textsc{prim})
\]

\noindent In words: the op's declared shape contract must match the argument shapes
modulo the equational theory (with groups resolved consistently), and the result
shape is the contract's result with the resolved group values plugged in. This is
the single rule that turns ``a shape is a type'' into a checkable obligation---and
it is the formal content of every shape-matching picture in
\cref{ch:tensors,ch:shapes}. A mismatch in any $\tau_i =_{\text{DimExpr}}
\theta(\sigma_i)$, or a failed group resolution, is a type error caught at the
boundary, before any arithmetic runs.

\begin{example}[\code{matvec} under \textsc{prim}]
\label{ex:b-matvec}
The contract is \lstinline[style=l4style]|matvec :: Tensor[m, n] -> Tensor[n] -> Tensor[m]|.
Apply it to a matrix of shape $[3,4]$ and a vector of shape $[4]$. Resolution
binds $m \mapsto 3$, $n \mapsto 4$ from the first argument; the second argument's
shape $[4]$ must equal $\theta([n]) = [4]$---it does. By \textsc{prim} the result
type is $\theta([m]) = [3]$, a length-$3$ vector. Hand it a length-$3$ vector
instead and the premise $[3] =_{\text{DimExpr}} [4]$ fails: a type error at the
boundary, exactly the shape-mismatch picture of \cref{ch:tensors}.
\end{example}

\Cref{fig:b-deriv} shows \textsc{prim} composing with the structural rules in a
small derivation tree: typing the composite expression
\lstinline[style=l4style]|axpy a x (matvec K x)|, where the inner
\lstinline[style=l4style]|matvec| produces the second \lstinline[style=l4style]|axpy|
operand.

\begin{figure}[t]
\centering
\begin{tikzpicture}[scale=1.0,
   every node/.style={font=\footnotesize}]
  % leaves
  \node (g1) at (-5.4,2.4) {$\Gamma \vdash \mathsf{K} : \texttt{Tensor[3,3]}$};
  \node (g2) at (-2.6,2.4) {$\Gamma \vdash \vx : \texttt{Tensor[3]}$};
  \node (g3) at (-2.6,0.9) {$\Gamma \vdash \vx : \texttt{Tensor[3]}$};
  \node (ga) at (-5.0,0.9) {$\Gamma \vdash a : \texttt{Scalar}$};
  % matvec node (PRIM)
  \node (mv) at (-3.9,1.55)
    {$\Gamma \vdash \texttt{matvec}\ \mathsf{K}\ \vx : \texttt{Tensor[3]}$};
  \draw[simGray] (-5.4,2.2) -- (-2.0,2.2);
  \node[font=\scriptsize, simRust] at (-1.5,2.2) {(\textsc{prim})};
  % axpy node (PRIM), top conclusion
  \node (ax) at (-3.4,-0.3)
    {$\Gamma \vdash \texttt{axpy}\ a\ \vx\ (\texttt{matvec}\ \mathsf{K}\ \vx) : \texttt{Tensor[3]}$};
  \draw[simGray] (-5.4,0.65) -- (-1.6,0.65);
  \node[font=\scriptsize, simRust] at (-1.1,0.65) {(\textsc{prim})};
  % connectors (informal tree edges)
  \draw[simGray, -] (mv.south) -- (-3.9,0.7);
  \node[font=\scriptsize\itshape, simGray, align=center] at (1.9,1.55)
    {\textbf{inner} \texttt{matvec}\\binds $m,n\mapsto 3$;\\result \texttt{Tensor[3]}};
  \node[font=\scriptsize\itshape, simGray, align=center] at (1.9,-0.3)
    {\textbf{outer} \texttt{axpy}\\binds $\$S\mapsto[3]$ from $\vx$;\\
     the \texttt{matvec} result, also\\\texttt{[3]}, is congruent};
  \draw[refedge, simBlue] (0.3,1.55) -- (mv.east);
  \draw[refedge, simBlue] (0.3,-0.3) -- (ax.east);
\end{tikzpicture}
\caption[A typing derivation for a composite expression]{A derivation of
$\Gamma \vdash \texttt{axpy}\ a\ \vx\ (\texttt{matvec}\ \mathsf{K}\ \vx) :
\texttt{Tensor[3]}$, showing \textsc{prim} (the primitive shape-contract rule)
applied twice and composed through the result of one op feeding the next. The
\emph{inner} \lstinline[style=l4style]|matvec| (top) takes the $[3,3]$ operator
$\mathsf{K}$ and the $[3]$ vector $\vx$, binds $m,n\mapsto 3$, and yields a $[3]$
result by \textsc{prim} (\cref{ex:b-matvec}). That $[3]$ result becomes the third
operand of the \emph{outer} \lstinline[style=l4style]|axpy| (bottom), whose group
\lstinline[style=l4style]|$S| is bound to $[3]$ by $\vx$; the
\lstinline[style=l4style]|matvec| result, also $[3]$, is congruent, so the whole
expression types at \lstinline[style=l4style]|Tensor[3]|. Each horizontal bar is one
rule instance; the tree is read bottom-up as ``to conclude the root, discharge these
premises.''}
\label{fig:b-deriv}
\end{figure}

\section{The primitive shape contracts: a reference}
\label{sec:b-contracts}

We collect the primitive operators' shape contracts in one place. Each is a pure
function (\cref{ch:operators} gives their algebra and laws; \cref{app:laws}
collects the equational laws). The table is the formal companion to the
shape-matching pictures of \cref{ch:tensors}: read each row as a signature plus the
side-condition its groups and axis names impose under \textsc{prim}.

\begin{figure}[t]
\centering
\small
\renewcommand{\arraystretch}{1.35}
\resizebox{\linewidth}{!}{%
\begin{tabular}{@{}>{\ttfamily}l l l@{}}
\toprule
\normalfont\textbf{op} & \textbf{signature} & \textbf{shape side-condition} \\
\midrule
axpy &
  \lstinline[style=l4style]|Scalar -> Tensor[$S] -> Tensor[$S] -> Tensor[$S]| &
  all three congruent ($S$) \\
dot &
  \lstinline[style=l4style]|Tensor[$S] -> Tensor[$S] -> Scalar| &
  both congruent; collapses \emph{all} axes \\
matvec &
  \lstinline[style=l4style]|Tensor[m, n] -> Tensor[n] -> Tensor[m]| &
  inner $n$ must match, cancels; $m$ survives \\
outer &
  \lstinline[style=l4style]|Tensor[m] -> Tensor[n] -> Tensor[m, n]| &
  no axis matches; rank \emph{raised} ($1{+}1$) \\
reduce &
  \lstinline[style=l4style]|Tensor[..., k] -> Tensor[...]| &
  labelled axis $k$ collapsed; rest survives \\
unfold &
  \lstinline[style=l4style]!Tensor[$S] -> (axis,size,stride) -> Tensor[..., size]! &
  windows one axis into a new trailing \texttt{size} \\
broadcast &
  \lstinline[style=l4style]!(target: Shape) -> Tensor[$Sa] -> Tensor[target]! &
  \texttt{target} a broadcast extension of $S_a$ \\
conv &
  \lstinline[style=l4style]|Tensor[N,Ci,Si] -> Tensor[Co,Ci,Sk] -> Tensor[N,Co,So]| &
  channels $C_i$ match; $S_o$ from $S_i,S_k$ \\
\bottomrule
\end{tabular}%
}
\caption[Primitive shape contracts]{The primitive operators and their shape
contracts, the formal reference behind \cref{ch:tensors}'s picture-per-op
treatment. \emph{Same-shape} ops (\lstinline[style=l4style]|axpy|,
\lstinline[style=l4style]|dot|) demand their operands be congruent over a group
$S$; \emph{contraction} (\lstinline[style=l4style]|matvec|) and
\emph{rank-raising} (\lstinline[style=l4style]|outer|) ops state the input/output
axis relation explicitly; \emph{rank-moving} ops
(\lstinline[style=l4style]|reduce| drops an axis,
\lstinline[style=l4style]|broadcast| / \lstinline[style=l4style]|unfold| add one);
and \lstinline[style=l4style]|conv| is the structured kernel contract (matching
input channels $C_i$, output spatial extent $S_o$ a function of input and kernel
extents $S_i$, $S_k$). Each side-condition is what the \textsc{prim} rule
(\cref{sec:b-primrule}) checks at the boundary. In running text we abbreviate the
broadcast/unfold subscript names; the source signatures use
\lstinline[style=l4style]!Tensor[$S_a]! and the explicit \texttt{target\_shape}
tuple.}
\label{fig:b-contracts}
\end{figure}

A few rows reward a second look, each a representative \emph{kind} of contract:

\begin{description}[leftmargin=1.4em, style=nextline, nosep]
\item[Same-shape: \code{axpy}, \code{dot}.] One group $S$ shared across all tensor
positions. \lstinline[style=l4style]|axpy| returns shape $S$;
\lstinline[style=l4style]|dot| collapses \emph{every} axis to a \code{Scalar} (the
$S=[\,]$ case of a reduction). The contract is the simplest: the shapes must be
congruent, decided by \cref{def:b-resolve}.
\item[Contraction: \code{matvec}.] Two distinct axis variables $m,n$; the inner
$n$ appears on both the matrix's second axis and the vector's only axis, so they
must agree and are summed away; the outer $m$ survives as the result.
\code{matvec} is the rank-$1$/rank-$2$ instance of the general contraction whose
group-level form is the operator-apply of \cref{sec:b-operators}.
\item[Rank-raising: \code{outer}.] No axis is required to match; the result rank is
the \emph{sum} of the inputs'. The opposite move to \code{matvec}.
\item[Rank-moving: \code{reduce}, \code{broadcast}, \code{unfold}.]
\lstinline[style=l4style]|reduce| drops the labelled axis $k$ (\code{dot} is its
collapse-everything special case); \lstinline[style=l4style]|broadcast| stretches a
tensor to a larger \texttt{target} shape, with the matching side-condition
(\texttt{target} must be a broadcast extension of the source, decided by the
\lstinline[style=l4style]|DimExpr| solver); \lstinline[style=l4style]|unfold|
windows one axis into a new trailing \texttt{size} axis, the rank-raising mechanism
behind the matrix-free element gather.
\item[Structured kernel: \code{conv}.] The convolution contract is the most
structured: batch $N$ and input channels $C_i$ must agree between data and kernel,
the kernel's output-channel axis $C_o$ becomes the result's channel axis, and the
output spatial extent $S_o$ is a \lstinline[style=l4style]|DimExpr| function of the
input and kernel spatial extents $S_i$, $S_k$.
\end{description}

\section{Scalar promotion as subtyping}
\label{sec:b-promotion}

Everything so far concerned \emph{shapes}. One rule about \emph{element types}
threads through the contracts and deserves its own formal statement: scalar
promotion, the rule that lets a \emph{real} scalar join a \emph{complex} operation
without minting a second operator. Cast formally, it is a one-step \emph{subtyping}
relation \lstinline[style=l4style]|real <: complex|.

\begin{definition}[Scalar promotion subtyping]
\label{def:b-promotion}
The scalar element types carry the one-step subtype order
\[
  \texttt{real} \;\sqsubseteq\; \texttt{complex},
\]
witnessed by the coercion $\alpha \mapsto \alpha + 0i$ (attach a zero imaginary
part). The coercion is \emph{exact}: zero is representable exactly, so $\alpha$ and
$\alpha + 0i$ are the same number to the last bit, and every algebraic law of an
operator holds identically before and after promotion. The order is one-way: there
is no demotion \lstinline[style=l4style]|complex <: real|.
\end{definition}

As an inference rule, promotion is the standard \emph{subsumption} rule
specialized to the scalar lattice: a real-typed expression may be used where a
complex one is expected.

\[
  \dfrac{\Gamma \vdash e : \texttt{real}}
        {\Gamma \vdash e : \texttt{complex}}\quad(\textsc{sub-scalar})
  \qquad\qquad
  \dfrac{}{\texttt{real} \sqsubseteq \texttt{complex}}\quad(\textsc{sub-base})
\]

\noindent and it lifts to tensors elementwise: a
\lstinline[style=l4style]|Tensor[s]| over \texttt{real} subsumes to a
\lstinline[style=l4style]|Tensor[s]| over \texttt{complex} of the \emph{same}
shape---promotion changes only the element type, never the shape. The practical
payoff is the collapse of overloads: an operator like \code{axpy} need not come in
separate real-scalar and complex-scalar forms. When its tensors are \texttt{complex}
and a scalar argument arrives \texttt{real}, \textsc{sub-scalar} promotes the scalar
and the one complex \code{axpy} runs; its laws are written once.
\Cref{fig:b-lattice} draws the lattice.

\begin{figure}[t]
\centering
\begin{tikzpicture}[scale=1.0]
  \node[draw=simViolet, fill=simBgViolet, rounded corners=3pt, thick,
        minimum width=22mm, minimum height=8mm, font=\small] (c) at (0,1.5)
    {\texttt{complex}};
  \node[draw=simBlue, fill=simBgBlue, rounded corners=3pt, thick,
        minimum width=22mm, minimum height=8mm, font=\small] (r) at (0,0)
    {\texttt{real}};
  \draw[-{Stealth[length=2.6mm]}, simRust, thick] (r) -- (c)
    node[midway, right=2mm, font=\scriptsize\itshape, simRust, align=left]
    {\textsc{sub-scalar}:\\ $\alpha \mapsto \alpha + 0i$\\ (exact, no rounding)};
  \node[font=\scriptsize\itshape, simGray, align=right, text width=28mm]
    at (-2.4,0.75)
    {the subtype lattice\\ \lstinline[style=l4style]!real <: complex!};
  \draw[-{Stealth[length=2.2mm]}, simGray, dotted] (c.west) to[bend right=42]
    node[midway, left=1pt, font=\scriptsize, simGray] {\,no} (r.west);
\end{tikzpicture}
\caption[The scalar-promotion subtype lattice]{Scalar promotion is the one-step
subtype order \lstinline[style=l4style]!real <: complex!. Upward (rust): a
\texttt{real} scalar is \emph{promoted}---used where \texttt{complex} is
expected---by the exact coercion $\alpha\mapsto\alpha+0i$, the witness of
\textsc{sub-scalar}. The promotion attaches a zero imaginary part with no rounding,
so the operator's algebraic laws (\cref{app:laws}) are invariant. Downward (grey,
dotted): \emph{no} demotion---a \texttt{complex} scalar against a \texttt{real}
operation is a type error, because the imaginary part has nowhere to go. The
lattice runs one way, collapsing the real/complex operator overloads into one.}
\label{fig:b-lattice}
\end{figure}

\begin{pitfall}
Promotion lifts a \texttt{real} scalar \emph{into} a complex operation; it never
demotes, and it never touches a \emph{reduction's output} scalar. A
\lstinline[style=l4style]|dot :: Tensor[$S] -> Tensor[$S] -> Scalar| has no
\emph{incoming} scalar to promote---its scalar is the \emph{output}, fixed by the
operands' element type (complex operands give a complex result); likewise a norm's
result is real regardless. \textsc{sub-scalar} is about an incoming scalar
argument, not a result. If you find yourself wanting to ``promote'' a complex
coefficient down to a real result, the type is genuinely wrong: there is no
information-preserving demotion, and the calculus rejects it rather than silently
discard the imaginary part.
\end{pitfall}

\section{Tying back: where these rules are used}
\label{sec:b-tieback}

The rules of this appendix are not free-standing formalism; each is the precise
version of a discipline the main text relies on continuously.

\begin{itemize}[leftmargin=1.4em, nosep]
\item \textbf{To \cref{ch:tensors}.} The shape language (\cref{sec:b-shapelang})
and the primitive-contract table (\cref{fig:b-contracts}) formalize that chapter's
bracket notation and its picture-per-op treatment of \code{axpy}, \code{dot},
\code{matvec}, \code{outer}, \code{reduce}, \code{broadcast}. The \textsc{prim}
rule (\cref{sec:b-primrule}) is what its ``checked at every function boundary''
slogan means exactly.
\item \textbf{To \cref{ch:shapes}.} Named-group resolution (\cref{def:b-resolve})
is the formal version of that chapter's rank-agnostic congruence and its
boundary-gate figure; the range-first operator shapes (\cref{sec:b-operators}) and
the congruence-group-versus-fixed-axis distinction (\cref{sec:b-fixedaxes}) are its
\lstinline[style=l4style]|LinOp[(R: ...), (D: ...)]| and \texttt{E}/\texttt{L}/%
\texttt{P}/\texttt{C}/\texttt{G} material made formal; and scalar promotion
(\cref{sec:b-promotion}) is its \lstinline[style=l4style]!real <: complex! rule
cast as subtyping.
\item \textbf{To \cref{ch:matrixfree}.} The fixed-meaning element-local axes
(\cref{sec:b-fixedaxes}) are the shape vocabulary of the matrix-free kernel; the
square operator form \lstinline[style=l4style]|LinOp[(S: ...), $S]| is the shape a
Krylov solver demands of its operator, re-feedable at every power.
\item \textbf{To \cref{app:calculus} and \cref{app:laws}.} The typing judgment
$\Gamma \vdash e : \tau$ (\cref{sec:b-judgment}) shares its term grammar and
reduction rules with \cref{app:calculus}; the algebraic laws the \textsc{prim} rule
presupposes (and which scalar promotion leaves invariant) are collected in
\cref{app:laws}.
\end{itemize}

Read together, the appendix says one thing in four registers: a shape is a type, a
named group is a unification variable over axes of unknown rank, the
shape-contract rule \textsc{prim} is where shapes acquire force, and scalar
promotion is a one-step subtyping that keeps the algebra single-voiced across real
and complex.

\summarysection
This appendix supplies the formal object underneath \cref{ch:tensors,ch:shapes}. The
\emph{shape language} (\cref{sec:b-shapelang}) is a BNF grammar for a shape
expression $\sigma$: a bracketed tuple of items, each a literal extent, a named axis
variable, an algebraic combination ($+,-,\times,{}^\wedge$), the rank-wildcard
\lstinline[style=l4style]|...|, or a named group (binding
\lstinline[style=l4style]|(S: ...)|, use \lstinline[style=l4style]|$S|). Shapes are
compared modulo the \lstinline[style=l4style]|DimExpr| equational theory
(\cref{def:b-dimexpr})---commutative, associative $+$ and $\times$---so that
\lstinline[style=l4style]|h*w| matches \lstinline[style=l4style]|w*h|.
\emph{Named-group resolution} (\cref{def:b-resolve}) is a two-step unification:
\emph{bind} the group from the first labelled run, then at each use check rank
\emph{first} and extents modulo \lstinline[style=l4style]|DimExpr|---the precise
content of the congruence back-arrow, and the reason \lstinline[style=l4style]|Tensor[N]|
cannot match a rank-$2$ argument. Operator shapes name two groups range-first,
\lstinline[style=l4style]|LinOp[(R: ...), (D: ...)]|, collapsing to the square form
for endomorphisms; congruence groups are kept distinct from named axes of fixed
meaning (\texttt{E}/\texttt{L}/\texttt{P}/\texttt{C}/\texttt{G}). The
\emph{type-and-shape judgment} $\Gamma \vdash e : \tau$ (\cref{sec:b-judgment}) is
the standard typed-$\lambda$ form; shapes acquire force at exactly one rule, the
primitive shape-contract rule \textsc{prim}, which requires each argument's shape to
match the op's contract modulo the equational theory and group resolution, then
substitutes the resolved groups into the result. The \emph{primitive shape
contracts} (\cref{fig:b-contracts}) collect \code{axpy}, \code{dot}, \code{matvec},
\code{outer}, \code{reduce}, \code{unfold}, \code{broadcast}, and \code{conv} as a
reference table. Finally, \emph{scalar promotion}
(\lstinline[style=l4style]!real <: complex!, \cref{def:b-promotion}) is a one-step
subtyping witnessed by the exact coercion $\alpha\mapsto\alpha+0i$: a real scalar
joins a complex operation, never the reverse, and never on a reduction's output
scalar.

\furtherreading
The view of a matrix as a typed map between shaped spaces, rather than a grid of
numbers, is the organizing idea of Trefethen and Bau~\citep{trefethenbau1997}; their
treatment of matrix--vector products and the four fundamental subspaces is the
linear-algebra backdrop for the range-first operator-shape convention
(\cref{sec:b-operators}) and the contraction contract of \code{matvec}. The
inference-rule presentation of the typing judgment $\Gamma\vdash e:\tau$
(\cref{sec:b-judgment}), the subsumption rule behind scalar promotion
(\cref{sec:b-promotion}), and the general discipline of reading a type system as a
set of judgment rules follow the standard implementation-of-functional-languages
account of Peyton Jones~\citep{peytonjones2003}, whose treatment of type inference,
unification, and the typed-$\lambda$ core is the closest analogue to the
shape-as-type / congruence-as-unification reading developed here. The grammar and
reduction rules this appendix's typing judgment shares are collected in
\cref{app:calculus}; the algebraic laws the \textsc{prim} rule presupposes are
collected in \cref{app:laws}.

\chapter{Monad and Algebraic Laws}
\label{app:laws}

\chapterorient{The combinators of \cref{ch:combinators} and the state monad of
\cref{ch:monad} are only as trustworthy as the rewrites they license. This appendix
is the reference catalogue of those rewrites: the monad laws that let you refactor a
\code{do}-block without changing what it computes, the state-effect algebra of
\code{get}/\code{put}/\code{modify}, and the equational laws of the tensor
primitives. But the catalogue has a second half that is every bit as important as the
first, and far easier to forget---the laws that do \emph{not} hold. A refactor that
\emph{looks} valid because it resembles a law you half-remember can silently change
the algorithm, or, more insidiously, change only its last bits. Those silent traps
are pinned here, each with the reason it fails and the harm it does. Read the
right-hand column as carefully as the left.}

\section{How to read this appendix}
\label{sec:cread}

The chapters of Parts~III and~IV state laws in passing, as they need them. This
appendix gathers them into one place and adds the precision a reference demands: each
law is an \emph{equation} (in math mode) or a \emph{rewrite} (a pseudocode
\code{==}), together with one sentence on what it \emph{guarantees}---what you are
allowed to do because it holds. The companion appendix on the formal calculus
(\cref{app:calculus}) gives the small-step reduction rules these laws sit on top of;
where a law is just a reduction rule renamed, we say so and point there.

Three kinds of statement appear, and the typography distinguishes them.

\begin{description}[style=nextline, leftmargin=0pt, font=\normalfont\itshape]
  \item[A law.] An equation that holds and may be used as a rewrite in either
  direction. Stated in a \code{theorem}, \code{lemma}, or a numbered list.
  \item[A non-law.] An equation that \emph{looks} like a law but is \emph{false}, set
  off explicitly. Each carries the reason it fails and the concrete damage assuming
  it would do. These are collected in \cref{sec:nonlaws} and are the most valuable
  content here.
  \item[A floating-point caveat.] A law that holds over an idealized exact field but
  only \emph{approximately} under IEEE-754 arithmetic. These are a special, sharp
  class of non-law---the equation is mathematically true and numerically false---and
  they are why \cref{app:numerics} exists. They are flagged inline and revisited in
  \cref{sec:cfp}.
\end{description}

\begin{keyidea}
The non-laws are as important as the laws. A law tells you a rewrite is \emph{safe};
a non-law tells you a rewrite that \emph{looks} safe is not. The dangerous case is
not the rewrite that obviously breaks---the types catch that---but the one that
resembles a law you half-remember: hoisting a loop predicate, merging two folds,
reordering a reduction. Such a rewrite can leave the program type-correct, still
running, and computing the wrong thing---or, worse, computing a \emph{bit-different}
thing that diverges only after thousands of iterations. \textbf{This is why a
refactor that looks valid can silently change the last bits.} The right-hand column
of this appendix---the laws that do \emph{not} hold---is the one you consult before
touching working solver code.
\end{keyidea}

\section{The monad laws}
\label{sec:cmonad}

We begin with the three laws that govern \code{>>=} (``bind'') and \code{return}, the
two operations underneath every \code{do}-block of \cref{ch:monad}. A \code{do}-block
is sugar: each line is desugared into a \code{>>=} that threads the state from one
line into the next (\cref{app:calculus} gives the desugaring in full). The three
monad laws are what make that sugar \emph{sound}---they are the guarantee that
rearranging the block in the ways the laws permit changes nothing about what it
computes. They hold for the \code{Solve} monad of \cref{ch:monad} exactly as they
hold for every monad.

\begin{theorem}[The three monad laws]
\label{thm:monad}
For all values \code{a}, recipes \code{m}, and recipe-valued functions \code{f},
\code{g}:
\begin{align}
  \textsf{return}\ a \,\texttt{>>=}\, f
    &\;=\; f\ a
    & \text{(left identity)} \label{eq:leftid}\\[2pt]
  m \,\texttt{>>=}\, \textsf{return}
    &\;=\; m
    & \text{(right identity)} \label{eq:rightid}\\[2pt]
  (m \,\texttt{>>=}\, f) \,\texttt{>>=}\, g
    &\;=\; m \,\texttt{>>=}\, (\lambda x.\, f\ x \,\texttt{>>=}\, g)
    & \text{(associativity)} \label{eq:assoc}
\end{align}
\end{theorem}

\noindent Take them one at a time, each with the refactor it licenses.

\paragraph{Left identity \eqref{eq:leftid}: \code{return} then bind is just
substitution.} If you lift a pure value into a recipe with \code{return} and
immediately feed it to \code{f}, you may as well have called \code{f} on the value
directly. In \code{do}-notation: a line
\lstinline[style=l4style]|x <- return a| followed by a use of \code{x} is identical
to binding \code{x} with a plain \lstinline[style=l4style]|let x = a|. The guarantee:
\emph{a \code{return} immediately consumed adds nothing}---you may always erase it and
substitute. This is the law that lets the over-monadifying mistake of
\cref{ch:monad}'s pitfall be \emph{undone}: a value needlessly wrapped in
\code{return} and then bound back out collapses to the pure binding it should have
been.

\paragraph{Right identity \eqref{eq:rightid}: binding into \code{return} is a
no-op.} If a recipe \code{m} runs and you immediately \code{return} whatever it
yielded, changing nothing, the trailing \code{return} can be dropped: the result is
just \code{m}. The guarantee: \emph{a recipe that ends by returning its own result is
that recipe}. This is what lets a \code{do}-block's final
\lstinline[style=l4style]|x <- step; return x| be shortened to \code{step}---the
redundant last line is provably inert.

\paragraph{Associativity \eqref{eq:assoc}: grouping of binds does not matter.} This
is the load-bearing one. It says that when three recipes run in sequence, it makes no
difference whether you parenthesize the sequencing as ``(first two) then third'' or
``first then (last two).'' In \code{do}-notation this is exactly the statement that
\emph{flattening and splitting \code{do}-blocks is safe}: a block

\begin{lstlisting}[style=l4style]
do { a <- m; do { b <- f a; g b } }   -- inner do-block nested
\end{lstlisting}

\noindent computes the same final state and value as the flattened

\begin{lstlisting}[style=l4style]
do { a <- m; b <- f a; g b }          -- flattened
\end{lstlisting}

\noindent because the two are \eqref{eq:assoc} read left-to-right and right-to-left.
The guarantee: \emph{you may regroup, factor out, and inline sub-blocks freely}. When
\cref{ch:monad} factors a Krylov step's body into a helper recipe and splices it back
into the loop, associativity is the warrant that the splice changed nothing.
\Cref{fig:monadlaw} draws this rewrite as a before/after.

\begin{figure}[t]
\centering
\begin{tikzpicture}[font=\footnotesize]
  % ===== LEFT: nested do-block =====
  \begin{scope}
    \node[font=\small\bfseries, text=simBlue] at (0,2.55) {before: nested};
    \node[draw=simGray!55, fill=simBgGray, rounded corners=4pt, thick,
          minimum width=42mm, minimum height=34mm] (outer) at (0,0.4) {};
    \node[anchor=north west, font=\ttfamily\scriptsize] at ($(outer.north west)+(2mm,-2mm)$)
      {do \{};
    \node[anchor=west, font=\ttfamily\scriptsize] at ($(outer.north west)+(4mm,-7mm)$)
      {a <- m;};
    \node[draw=simBlue!60, fill=simBgBlue, rounded corners=3pt, thick,
          minimum width=30mm, minimum height=15mm] (inner) at (0.6,-0.55) {};
    \node[anchor=north west, font=\ttfamily\scriptsize] at ($(inner.north west)+(1.5mm,-1.5mm)$)
      {do \{};
    \node[anchor=west, font=\ttfamily\scriptsize] at ($(inner.north west)+(3.5mm,-6mm)$)
      {b <- f a;};
    \node[anchor=west, font=\ttfamily\scriptsize] at ($(inner.north west)+(3.5mm,-10mm)$)
      {g b \} \}};
  \end{scope}
  % ===== arrow =====
  \node[font=\large, text=simRust] (arr) at (3.3,0.4) {$\equiv$};
  \node[font=\scriptsize\itshape, text=simRust, align=center] at (3.3,-0.55)
    {assoc.\\\eqref{eq:assoc}};
  % ===== RIGHT: flattened =====
  \begin{scope}[xshift=6.6cm]
    \node[font=\small\bfseries, text=simGreen] at (0,2.55) {after: flattened};
    \node[draw=simGreen!60, fill=simBgGreen, rounded corners=4pt, thick,
          minimum width=42mm, minimum height=34mm] (flat) at (0,0.4) {};
    \node[anchor=north west, font=\ttfamily\scriptsize] at ($(flat.north west)+(2mm,-2mm)$)
      {do \{};
    \node[anchor=west, font=\ttfamily\scriptsize] at ($(flat.north west)+(4mm,-8mm)$)
      {a <- m;};
    \node[anchor=west, font=\ttfamily\scriptsize] at ($(flat.north west)+(4mm,-14mm)$)
      {b <- f a;};
    \node[anchor=west, font=\ttfamily\scriptsize] at ($(flat.north west)+(4mm,-20mm)$)
      {g b \}};
  \end{scope}
  \node[font=\scriptsize\itshape, text=simGray, align=center] at (3.3,-2.4)
    {same threaded state, same yielded value};
\end{tikzpicture}
\caption[A monad law as a \texttt{do}-block rewrite]{Associativity \eqref{eq:assoc}
read as a \code{do}-block refactor. The nested block on the left---an inner
\code{do} spliced into an outer one---computes exactly the same final state and
yielded value as the flattened block on the right. The two are the two
parenthesizations of \eqref{eq:assoc}: \lstinline[style=l4style]|(m >>= f) >>= g|
versus \lstinline[style=l4style]|m >>= (\x -> f x >>= g)|. Because the law holds, you
may flatten, factor, and inline sub-blocks at will; the threading the
\code{do}-notation hides is rearranged but its net effect is invariant. This is the
formal license behind every ``pull this step out into a helper'' refactor in
\cref{ch:monad}.}
\label{fig:monadlaw}
\end{figure}

\begin{notationbox}
The three laws are stated as \emph{equalities} here, but in \cref{app:calculus} they
appear as left-to-right \emph{reduction rules} (the arrow $\to$). The direction
matters for an evaluator---reduction always moves left to right---but for
\emph{refactoring} both directions are sound, because the relation is an equality.
When this appendix writes a law with \code{=} or $=$ you may rewrite either way; when
\cref{app:calculus} writes the same fact with $\to$ it is fixing the evaluation
direction. Same fact, two uses.
\end{notationbox}

\section{The state-effect laws}
\label{sec:cstate}

The monad laws of \cref{sec:cmonad} hold for \emph{every} monad; they know nothing of
state. The \code{Solve} monad adds three state primitives---\code{get}, \code{put},
\code{modify}---and these obey their own small algebra, the rules for how reads and
writes of the threaded \code{SimState} interact. This algebra is what lets you
simplify a chain of state operations into fewer, clearer ones, and it is the formal
content behind \cref{ch:monad}'s claim that \code{modify} is the one honest verb for
a state change.

\begin{lemma}[The \texttt{get}/\texttt{put}/\texttt{modify} algebra]
\label{lem:state}
Writing \lstinline[style=l4style]|do {...}| for sequencing and $f, g$ for pure
\code{SimState -> SimState} functions:
\begin{align}
  \textsf{modify}\ f
    &\;=\; \textsf{do}\ \{\, s \leftarrow \textsf{get};\ \textsf{put}\ (f\ s) \,\}
    & \text{(\code{modify} \emph{is} get-apply-put)} \label{eq:modifydef}\\[2pt]
  \textsf{modify}\ (f \circ g)
    &\;=\; \textsf{do}\ \{\, \textsf{modify}\ g;\ \textsf{modify}\ f \,\}
    & \text{(\code{modify} fusion)} \label{eq:modifyfuse}\\[2pt]
  \textsf{do}\ \{\, s \leftarrow \textsf{get};\ \textsf{put}\ s \,\}
    &\;=\; \textsf{return}\ ()
    & \text{(get-then-put is inert)} \label{eq:getput}\\[2pt]
  \textsf{do}\ \{\, \textsf{put}\ s_1;\ \textsf{put}\ s_2 \,\}
    &\;=\; \textsf{put}\ s_2
    & \text{(a later \code{put} wins)} \label{eq:putput}\\[2pt]
  \textsf{do}\ \{\, \textsf{put}\ s;\ \textsf{get} \,\}
    &\;=\; \textsf{do}\ \{\, \textsf{put}\ s;\ \textsf{return}\ s \,\}
    & \text{(get after put sees the put)} \label{eq:putget}
\end{align}
\end{lemma}

\noindent These five are the complete working algebra; everything else follows by
combining them with the monad laws. Read each as a rewrite you are permitted.

\paragraph{\code{modify} is get-apply-put \eqref{eq:modifydef}.} This is the
\emph{definition} of \code{modify}, promoted to a law because it justifies both
expansions: any \code{modify} may be unfolded into an explicit get/apply/put, and---
read backwards---any get/apply/put triple may be \emph{folded} into a single
\code{modify}. The backward direction is the useful one: it is how three lines of
state plumbing collapse into the one named verb that, per \cref{ch:monad}, makes the
change stand out. \Cref{ch:monad}'s preference for \code{modify} over \code{put} is
this law's left-hand side being more honest than its right.

\paragraph{\code{modify} fusion \eqref{eq:modifyfuse}.} Two consecutive
\code{modify}s are one \code{modify} by the composed function. The guarantee:
\emph{adjacent state updates merge}. If a step bumps the counter and then, later in
the same block, records a residual---two \code{modify}s---they may be fused into a
single \code{modify} that does both, or, read the other way, one fat \code{modify}
may be split into stages for clarity. Note the order: \eqref{eq:modifyfuse} fuses to
$f \circ g$ with \code{g} applied \emph{first} (it ran first), so the composition
reads right-to-left exactly as function composition does. Getting this order backward
is a classic slip---which is why we write it out.

\paragraph{Get-then-put is inert \eqref{eq:getput}.} Reading the state and writing
back exactly what you read changes nothing: it equals \code{return ()}. The
guarantee: \emph{a read followed by an identity write can be deleted}. This is the
state-effect analog of the monad's right identity \eqref{eq:rightid}, and it is the
law that catches a particular dead-code shape---a refactor that left behind a
\lstinline[style=l4style]|s <- get; put s| with no change in between is doing nothing
and may be removed.

\paragraph{A later \code{put} wins \eqref{eq:putput}.} Two \code{put}s in a row: the
first is overwritten wholesale, so only the second survives. The guarantee:
\emph{a write with no intervening read is dead}. If two \code{put}s sit adjacent with
nothing reading the state between them, the first is pure overhead. (Contrast
\code{modify}: two \code{modify}s do \emph{not} collapse to the last one---each reads
the state the previous wrote, so they compose \eqref{eq:modifyfuse} rather than
overwrite. The difference is exactly that \code{modify} reads and \code{put} does
not, which is the algebra's whole texture.)

\paragraph{Get after put sees the put \eqref{eq:putget}.} If you \code{put} a state
and then \code{get}, the \code{get} yields exactly what you just \code{put}---reads
observe the most recent write. The guarantee: \emph{state is consistent within a
block}; there is no stale-read hazard once the monad threads for you. This is the law
that formalizes the line-by-line trace of \cref{ch:monad}'s walked \code{do}-block,
where the second \code{get} saw the \code{modify}'d state, not the original.

\Cref{fig:state} draws the four interaction laws \eqref{eq:modifyfuse}--\eqref{eq:putget}
as a small interaction table: what each adjacent pair of primitives simplifies to.

\begin{figure}[t]
\centering
\begin{tikzpicture}[font=\footnotesize]
  % header row
  \node[font=\scriptsize\bfseries, text=simGray] at (0,0) {\emph{then} $\downarrow$ / \emph{first} $\rightarrow$};
  \node[opbox, minimum width=16mm] (hg) at (2.2,0) {\code{get}};
  \node[opbox, minimum width=16mm] (hp) at (4.4,0) {\code{put s'}};
  \node[opbox, minimum width=16mm] (hm) at (6.6,0) {\code{modify g}};
  % row 1: get
  \node[opbox, minimum width=16mm] (rg) at (0,-1.0) {\code{get}};
  \node[font=\scriptsize, align=center] at (2.2,-1.0) {(idempotent\\read)};
  \node[font=\scriptsize, text=simRust, align=center] at (4.4,-1.0) {yields \code{s'}\\\eqref{eq:putget}};
  \node[font=\scriptsize, align=center] at (6.6,-1.0) {yields \code{g s}};
  % row 2: put s
  \node[opbox, minimum width=16mm] (rp) at (0,-2.2) {\code{put s}};
  \node[font=\scriptsize, align=center] at (2.2,-2.2) {read then\\overwrite};
  \node[font=\scriptsize, text=simRust, align=center] at (4.4,-2.2) {\code{put s}\\\eqref{eq:putput}};
  \node[font=\scriptsize, align=center] at (6.6,-2.2) {\code{put (g s)}};
  % row 3: modify f
  \node[opbox, minimum width=16mm] (rm) at (0,-3.4) {\code{modify f}};
  \node[font=\scriptsize, align=center] at (2.2,-3.4) {read then\\\code{f}};
  \node[font=\scriptsize, align=center] at (4.4,-3.4) {ignores \code{s'}};
  \node[font=\scriptsize, text=simRust, align=center] at (6.6,-3.4) {\code{modify}\\$(f\!\circ\! g)$ \eqref{eq:modifyfuse}};
  % separators
  \draw[simGray!40] (-1.4,-0.5) -- (7.6,-0.5);
  \draw[simGray!40] (1.1,0.4) -- (1.1,-4.0);
  \node[font=\scriptsize\itshape, text=simGray, align=center] at (3,-4.4)
    {rust cells are the simplification laws of \cref{lem:state};
     \emph{first} primitive runs, \emph{then} the second};
\end{tikzpicture}
\caption[The state-effect algebra]{The \code{get}/\code{put}/\code{modify}
interaction table. Each cell shows what the pair ``run the column primitive, then the
row primitive'' simplifies to. The rust cells are the named laws of
\cref{lem:state}: a \code{get} after a \code{put} yields the put value
\eqref{eq:putget}; a \code{put} after a \code{put} discards the first \eqref{eq:putput};
a \code{modify} after a \code{modify} fuses to the composed function
\eqref{eq:modifyfuse}. The asymmetry between \code{put} (overwrites, so an earlier
one dies) and \code{modify} (reads first, so they compose) is the whole point of the
algebra---and the reason \cref{ch:monad} prefers \code{modify}, the verb that names
its change.}
\label{fig:state}
\end{figure}

\begin{pitfall}
The \code{put}/\code{put} law \eqref{eq:putput} and the \code{modify}/\code{modify}
law \eqref{eq:modifyfuse} look symmetric and are not. Two \code{put}s collapse to the
\emph{last} one---the first is dead. Two \code{modify}s collapse to a
\emph{composition}---neither is dead, because the second \code{modify} reads the state
the first produced. Mistaking one for the other is a real bug: ``simplifying'' two
adjacent \code{modify}s by keeping only the second \emph{drops the first's effect
entirely}. Before collapsing adjacent state writes, ask whether the second
\emph{reads} the state---if it is a \code{modify}, it does, and you must compose, not
discard.
\end{pitfall}

\section{The algebraic operator laws}
\label{sec:coplaws}

We turn from the monad to the tensor primitives it threads. These are the equational
laws of the data-algebra verbs of \cref{ch:combinators}---the BLAS-style operations a
solver step calls---and of the matrix-vector apply. They are the rewrites the
synthesis surface appeals to when arguing that one form of a step computes the same
tensor as another. Unlike the monad laws, several of these carry a floating-point
caveat; those are flagged here and indicted fully in \cref{sec:nonlaws}.

\subsection{The scaling and saxpy laws}

\begin{lemma}[\texttt{axpy} corner laws]
\label{lem:axpy}
For \lstinline[style=l4style]|axpy a x y = a*x + y| (the two-term
\code{linear\_combination} of \cref{ch:combinators}):
\begin{align}
  \textsf{axpy}\ \alpha\ \vv\ \mathbf{0} &\;=\; \alpha \cdot \vv
    & \text{(zero accumulator $\Rightarrow$ scale)} \label{eq:axpyzeroacc}\\
  \textsf{axpy}\ 0\ \vv\ \vy &\;=\; \vy
    & \text{(zero coefficient $\Rightarrow$ identity)} \label{eq:axpyzerocoeff}
\end{align}
\end{lemma}

\noindent These are the two boundary cases of the saxpy. With a zero accumulator the
add drops away and \code{axpy} is a pure scale \eqref{eq:axpyzeroacc}; with a zero
coefficient the scaled term drops away and \code{axpy} returns its accumulator
untouched \eqref{eq:axpyzerocoeff}. The second is the algebraic content of the
implementation's ``skip the term if its coefficient is zero'' branch, lifted to
\code{linear\_combination}'s zero-coefficient term-drop. They license the obvious
peephole simplifications and nothing more.

\subsection{Inner-product laws}

\begin{lemma}[\texttt{dot} and \texttt{inner\_product} laws]
\label{lem:dot}
For the Euclidean inner product \lstinline[style=l4style]|dot v w| and its
operator-weighted generalization
$\inner{\vx}{\vy}_M = \vx^{\!\top} M \vy$:
\begin{align}
  \textsf{dot}\ \vv\ \vw &\;=\; \textsf{dot}\ \vw\ \vv
    & \text{(symmetry, real case)} \label{eq:dotsym}\\
  \textsf{dot}\ (\textsf{axpy}\ \alpha\ \vv\ \vw)\ \vx
    &\;=\; \alpha \cdot \textsf{dot}\ \vv\ \vx + \textsf{dot}\ \vw\ \vx
    & \text{(bilinearity)} \label{eq:dotbilin}\\
  \inner{\vx}{\vx}_M &\;\ge\; 0,\quad
    \inner{\vx}{\vx}_M = 0 \iff \vx = \mathbf{0}
    & \text{(positive-definiteness, SPD $M$)} \label{eq:dotpd}\\
  \inner{\vx}{\vy} &\;=\; \overline{\inner{\vy}{\vx}}
    & \text{(Hermitian symmetry, complex)} \label{eq:dotherm}
\end{align}
\end{lemma}

\noindent Symmetry \eqref{eq:dotsym} lets a solver write \code{dot p Ap} or
\code{dot Ap p} interchangeably. Bilinearity \eqref{eq:dotbilin} is the law that lets
a step-length computation distribute the inner product over a scaled sum---the
workhorse of every Krylov derivation. Positive-definiteness \eqref{eq:dotpd} on an
SPD weight $M$ is the law the norm \code{nrm2} rests its square root on, and the law
that makes the conjugate-gradient energy norm well-defined (\cref{ch:cg}). In the
complex case symmetry weakens to Hermitian symmetry \eqref{eq:dotherm}---the inner
product equals the conjugate of its swap---and the unconjugated bilinear variant
\code{tdot} loses positive-definiteness entirely (it is \emph{not} PSD; see
\cref{sec:nonlaws}).

\begin{pitfall}
Symmetry \eqref{eq:dotsym} is exact in \emph{exact} arithmetic. In floating point,
\code{dot v w} and \code{dot w v} may differ in the last bit if the two
accumulations visit operands in a different order---the multiply is commutative but
the running sum's rounding is order-sensitive. For the symmetric Euclidean dot the
two orders usually coincide, but do not \emph{rely} on bit-identity across a swap;
the safe statement is mathematical equality, not bit-equality. This is the
floating-point caveat of \cref{sec:cfp} in miniature.
\end{pitfall}

\subsection{Matrix-vector apply laws}

\begin{lemma}[\texttt{matvec} laws]
\label{lem:matvec}
For \lstinline[style=l4style]|matvec A v|, the application of a linear operator $A$:
\begin{align}
  \textsf{matvec}\ A\ (\alpha\,\vv)
    &\;=\; \alpha \cdot \textsf{matvec}\ A\ \vv
    & \text{(homogeneity in the vector)} \label{eq:mvhom}\\
  \textsf{matvec}\ (A + B)\ \vv
    &\;=\; \textsf{matvec}\ A\ \vv + \textsf{matvec}\ B\ \vv
    & \text{(additivity in the operator)} \label{eq:mvadd}\\
  (\textsf{matvec}\ A) \circ (\textsf{matvec}\ B)
    &\;=\; \textsf{matvec}\ (A \circ B)
    & \text{(composition)} \label{eq:mvcomp}\\
  \textsf{matvec}\ A\ (\alpha\vv + \beta\vw)
    &\;=\; \alpha\,\textsf{matvec}\ A\ \vv + \beta\,\textsf{matvec}\ A\ \vw
    & \text{(linearity)} \label{eq:mvlin}
\end{align}
\end{lemma}

\noindent Homogeneity \eqref{eq:mvhom} and additivity-in-the-operator \eqref{eq:mvadd}
combine into full linearity \eqref{eq:mvlin}: a \code{matvec} distributes over a
linear combination of its vector argument and pulls scalars out. The
composition law \eqref{eq:mvcomp} is the one with teeth in this book---it says that
applying $B$ then $A$ equals applying the single fused operator $A \circ B$, which is
exactly the algebraic license for the assembled-operator-product forms of the driven
analysis (\cref{ch:combinators} writes the frequency operator
$\mathsf{A}(\omega) = \mathsf{K} + i\omega\,\mathsf{C} - \omega^2\mathsf{M}$ as one
\code{linear\_combination} of operators, justified by \eqref{eq:mvadd}). Additivity
in the operator is also what lets the driven operator be \emph{assembled} as a sum
rather than applied term by term.

\subsection{The combinator structural laws}

The data-algebra combinators of \cref{ch:combinators} carry structural laws that
organize their whole family. These are restated here as the reference catalogue; the
chapter derives them.

\begin{theorem}[\texttt{linear\_combination} is a monoid homomorphism]
\label{thm:lincomb}
\code{linear\_combination} carries the monoid of term-lists under concatenation to
the monoid of tensors under addition:
\begin{align}
  \textsf{linear\_combination}\ [\,] &\;=\; \mathbf{0}
    & \text{(empty list $\mapsto$ unit)} \label{eq:lcempty}\\
  \textsf{linear\_combination}\ (a \mathbin{+\!\!+} b)
    &\;=\; \textsf{linear\_combination}\ a + \textsf{linear\_combination}\ b
    & \text{(concatenation-homomorphism)} \label{eq:lcconcat}\\
  \textsf{linear\_combination}\ ((\kappa a, t):\mathit{rest})
    &\;=\; \textsf{linear\_combination}\ ((a, \kappa t):\mathit{rest})
    & \text{(scalar absorption)} \label{eq:lcabsorb}\\
  \textsf{linear\_combination}\ ((0, t):\mathit{rest})
    &\;=\; \textsf{linear\_combination}\ \mathit{rest}
    & \text{(zero-term drop)} \label{eq:lcdrop}
\end{align}
\end{theorem}

\noindent The concatenation-homomorphism \eqref{eq:lcconcat} is the structural fact
of \cref{ch:combinators}: it is what makes the four fixed-arity BLAS routines
(\code{scal}, \code{axpy}, \code{axpby}, \code{axpbypcz}) \emph{one} combinator,
since a longer term-list is the concatenation of shorter ones. The empty list maps to
the zero tensor \eqref{eq:lcempty}; a scalar slides freely between coefficient and
tensor \eqref{eq:lcabsorb}; a zero-coefficient term drops out \eqref{eq:lcdrop}. Over
an \emph{exact} field the sum is also permutation-invariant---but that last one is a
floating-point caveat, deferred to \cref{sec:nonlaws} as the load-bearing non-law it
becomes under IEEE-754.

\begin{theorem}[\texttt{inner\_product} structural laws]
\label{thm:ip}
\code{inner\_product} is the scalar-producing dual: a monoid homomorphism from
shape-concatenated tensors to scalar addition, bilinear in its arguments.
\begin{align}
  \textsf{inner\_product}\ (\vx_1 \mathbin{+\!\!+} \vx_2)\ (\vy_1 \mathbin{+\!\!+} \vy_2)
    &\;=\; \textsf{inner\_product}\ \vx_1\ \vy_1 + \textsf{inner\_product}\ \vx_2\ \vy_2
    \label{eq:ipsplit}\\
  \textsf{inner\_product}\ \mathbf{0}\ \vy
    &\;=\; \textsf{inner\_product}\ \vx\ \mathbf{0} \;=\; 0
    \label{eq:ipzero}
\end{align}
\end{theorem}

\noindent The split-additivity \eqref{eq:ipsplit} is the law that licenses blocked
and parallel evaluation of an inner product: split the operands, reduce the pieces
independently, add the partial sums. It is the formal twin of
\code{linear\_combination}'s concatenation law, only collapsing to a scalar rather
than a tensor. \Cref{fig:catalogue} gathers all of \cref{sec:coplaws} into a single
reference table.

\begin{figure}[t]
\centering
\renewcommand{\arraystretch}{1.35}
\setlength{\tabcolsep}{6pt}
\begin{tabular}{@{}p{0.26\textwidth} p{0.40\textwidth} p{0.24\textwidth}@{}}
\toprule
\textbf{Operator} & \textbf{Law} & \textbf{Guarantees} \\
\midrule
\code{axpy} &
  $\textsf{axpy}\ \alpha\ \vv\ \mathbf{0} = \alpha\vv$;\ \
  $\textsf{axpy}\ 0\ \vv\ \vy = \vy$ &
  peephole drop of a dead scale/add \\
\code{dot} &
  $\textsf{dot}\ \vv\ \vw = \textsf{dot}\ \vw\ \vv$ (symmetry) &
  operand order free (math, not bits) \\
\code{dot} &
  $\textsf{dot}\,(\textsf{axpy}\,\alpha\,\vv\,\vw)\,\vx
   = \alpha\,\textsf{dot}\,\vv\,\vx + \textsf{dot}\,\vw\,\vx$ &
  distribute over scaled sums \\
\code{matvec} &
  $\textsf{matvec}\,A\,(\alpha\vv) = \alpha\,\textsf{matvec}\,A\,\vv$ &
  pull scalars through apply \\
\code{matvec} &
  $\textsf{matvec}\,(A{+}B)\,\vv = \textsf{matvec}\,A\,\vv + \textsf{matvec}\,B\,\vv$ &
  assemble operator as a sum \\
\code{matvec} &
  $(\textsf{matvec}\,A)\!\circ\!(\textsf{matvec}\,B) = \textsf{matvec}\,(A\!\circ\! B)$ &
  fuse two applies into one \\
\code{linear\_combination} &
  $\textsf{lc}\,(a {+}\mkern-3mu{+}\, b) = \textsf{lc}\,a + \textsf{lc}\,b$ &
  split / chunk / reorder the sum (exact) \\
\code{inner\_product} &
  $\textsf{ip}\,(\vx_1{+}\mkern-3mu{+}\vx_2)\,(\vy_1{+}\mkern-3mu{+}\vy_2)
   = \textsf{ip}\,\vx_1\,\vy_1 + \textsf{ip}\,\vx_2\,\vy_2$ &
  blocked / parallel reduction (exact) \\
\bottomrule
\end{tabular}
\caption[The operator-law catalogue]{The algebraic operator laws of
\cref{sec:coplaws} as a reference table. Each row gives an operator, an equational
law, and the rewrite it permits. The two combinator homomorphism laws
(\code{linear\_combination}, \code{inner\_product}) are marked ``exact''---they hold
over an idealized field but become floating-point caveats once the reduction order
is fixed, the subject of \cref{sec:nonlaws} and \cref{app:numerics}. The
\code{matvec} composition law is the one with the most downstream leverage: it is the
algebraic license for assembling and fusing operators in the driven analysis.}
\label{fig:catalogue}
\end{figure}

\section{The laws that do \emph{not} hold}
\label{sec:nonlaws}

We arrive at the appendix's most valuable content. The laws above tell you which
rewrites are safe. This section tells you which \emph{plausible} rewrites are
\emph{unsafe}---the ones that resemble a law closely enough to tempt a refactor, but
fail, and fail silently. Each is stated as the false equation it is, with the reason
it fails and the concrete damage assuming it would do. They are grouped by the
operator they tempt you to mishandle.

\subsection{Non-laws of \texttt{iterate\_while}}
\label{sec:nlloop}

The loop combinator of \cref{ch:fold} carries four traps, all variations on ``treat
the loop as if it were an ordinary algebraic object.''

\begin{description}[style=nextline, leftmargin=0pt, font=\normalfont\itshape]
  \item[Predicate hoisting is invalid.] One might hope to hoist the loop predicate
  out of the iteration---test it once, then run the body unconditionally:
  \begin{lstlisting}[style=l4style]
iterate_while a p f  =/=  if p a then (repeat f forever) else stop   -- WRONG
  \end{lstlisting}
  This is false. The predicate is re-evaluated on \emph{every} iteration, not once at
  entry. Hoisting it converts a bounded loop into an unbounded one driven by the
  \emph{initial} predicate value---if \code{p a} is true at entry, the hoisted form
  runs forever, because nothing re-tests it. \textbf{What breaks:} a solver that was
  converging in fifty steps now never terminates; the bug is a hang, not a wrong
  number, and it surfaces only on inputs where the predicate was initially true.
  \item[Loop fusion across iterations is invalid.] Two \code{iterate\_while}s in
  sequence are \emph{not} one \code{iterate\_while} with a merged predicate:
  \begin{lstlisting}[style=l4style]
(iterate_while a p1 f) then (iterate_while _ p2 f)  =/=  iterate_while a (p1 || p2) f
  \end{lstlisting}
  except in the special case where the two predicates agree at the handoff state.
  The two-phase fold's trajectory is the concatenation \code{traj1 ++ traj2} with the
  second phase \emph{restarting} from the first phase's final state; the flattened
  single fold is one continuous trajectory under one predicate. \textbf{What breaks:}
  this is precisely why GMRES is written as an \emph{outer} restart loop wrapped
  around an \emph{inner} \code{iterate\_while} (\cref{ch:gmres}), not as one flat
  loop---the restart \emph{is} the phase boundary, and fusing it away discards the
  restart's basis reset. Merging the two loops silently changes the algorithm from
  restarted GMRES to a single un-restarted Krylov method with a different convergence
  and memory profile.
  \item[\code{while} and \code{do-while} are not interchangeable.] The book's
  \code{iterate\_while} tests the predicate \emph{before} the body (a \code{while});
  the variant that tests \emph{after} (a \code{do-while}, which always runs the body
  at least once) is a \emph{different combinator}:
  \begin{lstlisting}[style=l4style]
iterate_while a p f  =/=  iterate_while_post a p f   -- differ when p a is false at entry
  \end{lstlisting}
  They disagree exactly when the predicate is false at entry: the \code{while} runs
  zero steps and returns \code{a}; the \code{do-while} runs one step anyway.
  \textbf{What breaks:} swapping one for the other off-by-ones the iteration count and,
  for a loop entered already-converged, performs one spurious step on a state that
  should have been returned untouched.
  \item[No identity element.] There is no carry \code{a\_id} for which
  \code{iterate\_while a\_id p f} is a unit for all \code{p}, \code{f}. The loop is a
  \emph{fold}, not a monoid operation; folds do not have a unit in the carry.
  \textbf{What breaks:} any refactor that tries to ``initialize the loop with a
  neutral state'' and expects it to be inert is resting on a structure that is not
  there.
\end{description}

\subsection{Non-laws of \texttt{krylov\_step}}
\label{sec:nlkrylov}

The Krylov step is the richest source of non-laws, because it \emph{looks} like a
composition of well-behaved linear primitives and is not---its scalar updates involve
divisions and its convergence flag a comparison, and those destroy the linearity and
composability the primitives individually enjoy. Each non-law below is a property the
\emph{primitives} have that the \emph{step} does not.

\begin{description}[style=nextline, leftmargin=0pt, font=\normalfont\itshape]
  \item[The primitive sequence does not commute.] The five primitive groups of a
  step---operator apply, auxiliary, iterate-update, scalar-update, output-readout---
  cannot be reordered. The dataflow is rigid: \code{apply} produces $\vw$, the
  \code{axpy} reads it, the \code{dot} reads it. \textbf{What breaks:} reorder the
  apply past the \code{dot} that consumes its output and you read uninitialized or
  stale data; the result is not a different valid ordering, it is wrong.
  \cref{fig:noncommute} traces two orderings of a reduction to bit-different results.
  \item[No fold-merge / associativity.] Folding a Krylov step over predicate $p_1$
  and then over $p_2$ is \emph{not} folding over a merged predicate
  \lstinline[style=l4style]|p1 || p2|. The inner fold's per-step outputs are invisible
  to the outer fold, and convergence predicates that depend on monotonic-loss
  properties do not compose. \textbf{What breaks:} the same restart-flattening hazard
  as the loop's, now at the step level---it is why restart logic is an outer loop
  (\cref{ch:gmres}), not a flattened single fold.
  \item[No identity element.] There is no Krylov state \code{K\_id} with
  \lstinline[style=l4style]|krylov_step op K_id = pure { krylov: K_id, ... }|. The
  tempting candidate---a zero step length $\alpha = 0$---is not identity but
  \emph{breakdown}: in CG, $\alpha = 0$ is a division-by-zero-adjacent degeneracy, not
  a no-op. \textbf{What breaks:} treating a stalled step as an identity hides a
  breakdown the solver must detect and report.
  \item[No linearity in any argument.] The step is \emph{not} linear in its state:
  \[
    \textsf{krylov\_step}\ op\ (\alpha K_1 + \beta K_2)
      \;\neq\; \alpha\,\textsf{krylov\_step}\ op\ K_1
             + \beta\,\textsf{krylov\_step}\ op\ K_2.
  \]
  It is built from linear primitives, but their composition with \code{dot} and the
  scalar \emph{divisions} of the coefficient updates ($\alpha = r\!\cdot\! r / p\!\cdot\! Ap$)
  is non-linear, and the convergence comparison is not even continuous.
  \textbf{What breaks:} any attempt to ``superpose'' two solves into one by linearity
  computes nonsense; the step is an affine-ish iteration, not a linear map.
  \item[Bit-non-determinism across orthogonalization variants.] Switching the
  orthogonalization scheme---modified Gram--Schmidt (MGS), classical (CGS), or
  twice-iterated (CGS2)---produces \emph{mathematically equivalent} but
  \emph{bit-distinct} trajectories. \textbf{What breaks:} a regression test that pins
  exact bits will fail across a variant switch even though the math is unchanged; this
  is a load-bearing numerical choice (\cref{ch:bigidea}'s transparent-versus-load-bearing
  distinction), not a free refactor.
  \item[Form-A / Form-B equivalence is \emph{not} a monad-law rewrite.] The
  first-iteration-unrolled Form B of a Krylov step produces
  iteration-for-iteration-identical trajectories to the rolled Form A---but the two
  are \emph{not} related by the monad laws of \cref{sec:cmonad}. The rotation between
  them is a \emph{structural} rewrite (it drops a state field and threads a closure
  argument), not a $\beta$-reduction or a \code{>>=}-associativity step.
  \textbf{What breaks:} citing ``the monad laws'' to justify the unrolling is a
  category error; the equivalence holds, but it must be argued structurally, and a
  tool that only knows the monad laws cannot verify it.
\end{description}

\begin{figure}[t]
\centering
\begin{tikzpicture}[font=\footnotesize]
  % ===== ordering A =====
  \begin{scope}
    \node[font=\small\bfseries, text=simBlue] at (0,2.45) {ordering A};
    \node[data, minimum width=11mm] (a0) at (-1.4,1.5) {$c_1$};
    \node[opbox, minimum width=9mm] (aA) at (0.2,1.5) {$+$};
    \node[data, minimum width=11mm] (a1) at (1.8,1.5) {$c_2$};
    \draw[flow] (a0)--(aA); \draw[flow] (a1)--(aA);
    \node[opbox, minimum width=9mm] (aB) at (0.2,0.4) {$+$};
    \node[data, minimum width=11mm] (a2) at (1.8,0.4) {$c_3$};
    \draw[flow] (aA)--(aB); \draw[flow] (a2)--(aB);
    \node[product, minimum width=22mm] (ar) at (0.2,-0.7)
      {$(c_1\!\oplus\! c_2)\!\oplus\! c_3$};
    \draw[flow] (aB)--(ar);
    \node[font=\scriptsize\ttfamily, text=simRust] at (0.2,-1.5) {= 1.0000003};
  \end{scope}
  \draw[simGray!50, thick] (3.3,-1.8) -- (3.3,2.6);
  % ===== ordering B =====
  \begin{scope}[xshift=6.6cm]
    \node[font=\small\bfseries, text=simGreen] at (0,2.45) {ordering B};
    \node[data, minimum width=11mm] (b0) at (-1.4,1.5) {$c_2$};
    \node[opbox, minimum width=9mm] (bA) at (0.2,1.5) {$+$};
    \node[data, minimum width=11mm] (b1) at (1.8,1.5) {$c_3$};
    \draw[flow] (b0)--(bA); \draw[flow] (b1)--(bA);
    \node[opbox, minimum width=9mm] (bB) at (0.2,0.4) {$+$};
    \node[data, minimum width=11mm] (b2) at (1.8,0.4) {$c_1$};
    \draw[flow] (bA)--(bB); \draw[flow] (b2)--(bB);
    \node[product, minimum width=22mm] (br) at (0.2,-0.7)
      {$c_1\!\oplus\!(c_2\!\oplus\! c_3)$};
    \draw[flow] (bB)--(br);
    \node[font=\scriptsize\ttfamily, text=simRust] at (0.2,-1.5) {= 1.0000002};
  \end{scope}
  \node[font=\scriptsize\itshape, text=simGray, align=center] at (3.3,-2.3)
    {same operands, same exact sum --- different rounding, different last bit};
\end{tikzpicture}
\caption[Non-commutativity of a reduction: two orderings, different last bits]{Two
reduction orderings of the same three floating-point operands $c_1, c_2, c_3$, with
$\oplus$ the rounding floating-point add. Ordering A accumulates
$(c_1 \oplus c_2) \oplus c_3$; ordering B accumulates $c_1 \oplus (c_2 \oplus c_3)$.
The \emph{exact} sum is identical and order-independent---real addition is
associative---but each rounding step discards low bits differently, so the two orders
return results that differ in their last bit (the schematic values
\code{1.0000003} versus \code{1.0000002}). This is why \code{dot}, \code{nrm2}, and
\code{linear\_combination} are \emph{not} bit-associative: the reduction tree is part
of the algorithm, not a free choice (\cref{sec:cfp}, \cref{app:numerics}).}
\label{fig:noncommute}
\end{figure}

\subsection{Floating-point non-laws of the reductions}
\label{sec:cfp}

The most consequential non-laws in the entire book are the ones that hold in
mathematics and fail in floating point. The reduction verbs---\code{dot},
\code{nrm2}, \code{inner\_product}, and the summation inside
\code{linear\_combination}---are over an \emph{exact} field associative and
commutative, so the catalogue of \cref{sec:coplaws} stated their split and
permutation laws as exact. Under IEEE-754 those laws become \emph{caveats}.

\begin{description}[style=nextline, leftmargin=0pt, font=\normalfont\itshape]
  \item[Reductions are not bit-associative.] Floating-point $\oplus$ is not
  associative: $(a \oplus b) \oplus c \neq a \oplus (b \oplus c)$ in general, because
  each step rounds. So the value of \code{dot v w} or
  \lstinline[style=l4style]|linear_combination pairs| \emph{depends on the order in
  which the partial sums are accumulated}---the reduction tree. The mathematical value
  is order-agnostic; the computed value is not. \cref{fig:noncommute} is this fact in
  three operands.
  \item[Permutation-invariance of \code{linear\_combination} \eqref{eq:lcconcat} is
  exact-only.] Reordering the term list leaves the \emph{exact} sum unchanged but
  perturbs the rounding. The framework therefore pins the left-to-right
  \code{foldl} order as \emph{canonical} and treats any reordering as a load-bearing
  numerical change, not a free rewrite.
  \item[Split-additivity of \code{inner\_product} \eqref{eq:ipsplit} licenses blocked
  evaluation only up to rounding.] Splitting the operands and adding the partial
  inner products gives the \emph{exact} answer but a \emph{bit-different} one from the
  unsplit reduction, because the partial-sum tree differs.
  \item[Cauchy--Schwarz can fail by ULPs.] $\abs{\inner{\vx}{\vy}} \le \norm{\vx}\norm{\vy}$
  holds mathematically but can be violated by unit-in-last-place amounts in floating
  point, because the two sides are reduced by different trees.
\end{description}

\textbf{What breaks if you assume these are exact laws:} bit-reproducibility. A
solver refactored to ``reorder the reduction for cache locality'' or ``split the dot
product across two threads and add the halves'' computes a \emph{mathematically
equal} but \emph{bit-different} residual. For a single step that is invisible; over
thousands of iterations the bit-difference compounds through the convergence
trajectory, and two builds that should be identical now disagree on iteration counts
and final digits. This is the precise mechanism by which \emph{a refactor that looks
valid silently changes the last bits}, and it is the reason the book classifies
reduction order as a load-bearing numerical trick (\cref{ch:bigidea}) rather than a
transparent one. The full treatment---error bounds, compensated summation, the cost
of pinning a tree---is \cref{app:numerics}.

\begin{keyidea}
The floating-point non-laws are the sharpest edge in this appendix. An exact-field
law (associativity of $+$, permutation-invariance of a sum, Cauchy--Schwarz) is
\emph{true in mathematics and false in the machine}. Assuming the exact law when
refactoring a reduction does not produce a wrong answer you can see---it produces a
\emph{bit-different} answer that is mathematically correct and reproducibly wrong, and
the divergence only becomes visible after it has compounded through a long iteration.
Reduction order is part of the algorithm. Treat \code{dot}, \code{nrm2}, and the
\code{linear\_combination} sum as having a \emph{pinned tree}, not a free one, whenever
bit-reproducibility matters (\cref{app:numerics}).
\end{keyidea}

\section{Law versus non-law: the reference card}
\label{sec:ccard}

\Cref{fig:lawcard} is the two-column summary the rest of this appendix exists to
support: on the left, the rewrites you \emph{can} apply; on the right, the
look-alikes you \emph{cannot}. Consult it before any refactor of solver code.

\begin{figure}[t]
\centering
\renewcommand{\arraystretch}{1.3}
\setlength{\tabcolsep}{6pt}
\begin{tabular}{@{}p{0.46\textwidth} p{0.46\textwidth}@{}}
\toprule
\textbf{You CAN rewrite\ \ (laws)} & \textbf{You CANNOT rewrite\ \ (non-laws)} \\
\midrule
flatten / split / inline \code{do}-blocks \eqref{eq:assoc} &
  hoist a loop predicate out of \code{iterate\_while} \\
drop a \code{return} immediately consumed \eqref{eq:leftid} &
  fuse two sequential loops into one merged-predicate loop \\
fold get-apply-put into one \code{modify} \eqref{eq:modifydef} &
  swap \code{while} for \code{do-while} \\
fuse two adjacent \code{modify}s to $f\!\circ\! g$ \eqref{eq:modifyfuse} &
  collapse two \code{modify}s to the last one (that is a \code{put} law) \\
delete an inert get-then-put \eqref{eq:getput} &
  reorder the primitive sequence inside a \code{krylov\_step} \\
distribute \code{dot} over a scaled sum \eqref{eq:dotbilin} &
  treat $\alpha=0$ as a step identity (it is breakdown) \\
fuse two \code{matvec}s into one \eqref{eq:mvcomp} &
  superpose two Krylov solves by ``linearity'' \\
split a \code{linear\_combination} / \code{inner\_product} \emph{exactly} &
  reorder a reduction and expect \emph{bit}-identity \\
\bottomrule
\end{tabular}
\caption[Law versus non-law reference card]{The two-column reference card. The
left column lists the safe rewrites of \cref{sec:cmonad}--\cref{sec:coplaws}, each
keyed to its law; the right column lists the look-alike rewrites of
\cref{sec:nonlaws} that are \emph{not} laws. The pairing is deliberate: each
right-column entry resembles a left-column one closely enough to tempt the
substitution. The distinction between fusing two \code{modify}s (a law,
left) and collapsing two writes to the last (a \code{put} law misapplied to
\code{modify}, right) is the kind of near-miss this card exists to catch. When a
refactor is not on the left, assume it is on the right until proven otherwise.}
\label{fig:lawcard}
\end{figure}

\summarysection
This appendix is the reference catalogue of the rewrites the calculus licenses, in
two halves. The \emph{monad laws}---left identity \eqref{eq:leftid}, right identity
\eqref{eq:rightid}, associativity \eqref{eq:assoc}---govern \code{>>=} and
\code{return} and make \code{do}-block refactoring sound: you may erase a consumed
\code{return}, drop a trailing one, and flatten, split, or inline sub-blocks freely
(\cref{sec:cmonad}). The \emph{state-effect laws} \eqref{eq:modifydef}--\eqref{eq:putget}
are the algebra of \code{get}/\code{put}/\code{modify}: \code{modify} is get-apply-put,
adjacent \code{modify}s fuse to a composition, a get-then-put is inert, a later
\code{put} wins, and a get after a put sees the put---with the load-bearing asymmetry
that \code{put}s overwrite (so an earlier one dies) while \code{modify}s read first
(so they compose, not collapse) (\cref{sec:cstate}). The \emph{operator laws}
(\cref{sec:coplaws}) are the equational reference for the tensor primitives:
\code{axpy} corner cases, \code{dot} symmetry and bilinearity, \code{matvec}
linearity and composition, and the combinator homomorphisms that make
\code{linear\_combination} and \code{inner\_product} their families' entries. The
second half---the \emph{non-laws} (\cref{sec:nonlaws})---is the load-bearing content:
predicate hoisting and loop fusion are invalid; \code{while} and \code{do-while}
differ; a Krylov step does not commute, merge, have an identity, or stay linear, and
its Form-A/Form-B equivalence is structural, not a monad-law rewrite; and---sharpest
of all---the reductions are \emph{not} bit-associative, so reordering a \code{dot} or
a \code{linear\_combination} sum is a load-bearing numerical change, not a free one
(\cref{sec:cfp}). The reference card of \cref{sec:ccard} pairs each safe rewrite with
its dangerous look-alike. The governing lesson is the one the \cref{sec:cread} key
idea states and \cref{sec:cfp} sharpens: \textbf{the non-laws are as important as the
laws, because they are why a refactor that looks valid can silently change the last
bits.}

\furtherreading
The monad laws and the \code{get}/\code{put}/\code{modify} algebra are standard
functional-programming material; Wadler's tutorial derives them from first principles
and is the readable canonical source for why they make stateful-looking pure code
sound to refactor~\citep{wadler1995monads}, while the precise desugaring of
\code{do}-notation into \code{>>=}---on which the associativity law's
\code{do}-block reading rests---is documented in the language report Peyton Jones
edits~\citep{peytonjones2003}. The floating-point non-laws of \cref{sec:cfp}---the
non-associativity of summation and its consequences for iterative solvers---are
treated at length in Saad's standard reference on iterative methods, which is also
the home of the Krylov convergence theory the step's non-linearity and
non-composability (\cref{sec:nlkrylov}) ultimately rest on~\citep{saad2003}. The
reduction-order error analysis those non-laws demand is the subject of the companion
\cref{app:numerics}; the formal reduction rules these equational laws sit atop are in
\cref{app:calculus}.

\chapter{Numerical and Floating-Point Semantics}
\label{app:numerics}

\chapterorient{Every algorithm in this book has been presented in exact
arithmetic---inner products are sums, the residual is $\vb-\mat{A}\vx$, a
polynomial is its value. A real machine does none of this exactly. It rounds, and
the order in which it rounds is visible in the last bits of the answer. This
appendix is about which of those rounding choices the reader may safely
\emph{ignore} and which are \emph{part of the algorithm}. The distinction is the
one \cref{ch:bigidea} drew between performance tricks and base algebra, sharpened
to the floating-point level: a \emph{transparent} trick computes the same result
faster and is invisible to the mathematics; a \emph{load-bearing} trick changes a
property the algorithm depends on---determinism, conditioning, IEEE
compliance---and must be stated as an explicit claim. We first give a working
floating-point primer (rounding, machine epsilon, non-associativity, cancellation,
conditioning), then catalogue the transparent tricks the book has quietly skipped
over, then treat the load-bearing ones in depth: non-associative reduction
orderings, fast-math reassociation, mixed precision, the Chebyshev
recurrence-versus-monomial choice (\cref{ch:smoothers}), and modified-versus-classical
Gram--Schmidt (\cref{ch:orthog}). \Cref{app:laws} catalogues which algebraic
\emph{laws} survive the rotation to pure functions; this appendix explains why
several of them hold only \emph{up to rounding}, and why that ``up to'' is itself a
specification.}

\section{A floating-point primer}
\label{sec:primer}

The reader who has written a numerical program has met the surprises of floating
point---a sum that depends on the order of its terms, a subtraction that loses all
its digits---without necessarily having a vocabulary for them. This section
supplies the vocabulary, at the level needed to follow the rest of the appendix.
Nothing here is deep; all of it is load-bearing for what follows.

\subsection{Floating-point numbers and rounding}

A binary floating-point number is a sign, a fraction (the \emph{significand} or
\emph{mantissa}), and an exponent: $x = \pm\, (1.b_1 b_2 \cdots b_{p-1})_2 \times
2^{e}$, with $p$ significand bits. IEEE-754 double precision (``binary64'',
the workhorse of every solver in this book) has $p = 53$ significand bits and an
11-bit exponent; single precision (``binary32'') has $p = 24$ and an 8-bit
exponent. The finite significand is the entire story: only numbers expressible in
$p$ bits exist, and every arithmetic result must be \emph{rounded} to the nearest
representable one.

\begin{definition}[Unit roundoff]
\label{def:eps}
The \emph{unit roundoff} (or \emph{machine epsilon}) $u$ is the largest relative
error incurred by rounding a real number to the nearest floating-point number:
under round-to-nearest, $\mathrm{fl}(x) = x(1+\delta)$ with $\abs{\delta} \le u$,
where $u = 2^{-p}$. For double precision $u = 2^{-53} \approx 1.11\times10^{-16}$;
for single precision $u = 2^{-24} \approx 5.96\times10^{-8}$.
\end{definition}

The fundamental model of floating-point arithmetic follows: each elementary
operation $\circ \in \{+,-,\times,\div\}$ returns the exactly-rounded result,
\begin{equation}
\label{eq:fpmodel}
  \mathrm{fl}(a \circ b) \;=\; (a \circ b)(1 + \delta), \qquad \abs{\delta} \le u .
\end{equation}
Each \emph{single} operation is as accurate as the format allows---a relative
error of at most $u$. The trouble is never one operation; it is the way errors
from many operations combine, and \eqref{eq:fpmodel} is the seed from which every
phenomenon below grows.

\begin{notationbox}
We write $\mathrm{fl}(\cdot)$ for ``the floating-point value of,'' $u$ for the unit
roundoff (\cref{def:eps}), and $\kappa$ for a condition number (\cref{def:cond}
below, and \cref{def:kappa} in \cref{ch:cg} for the matrix case). A relative error
``of order $u$'' means $O(u)$; ``of order $\kappa u$'' means the rounding has been
amplified by the conditioning. These three symbols---$u$, $\kappa$, and their
product $\kappa u$---carry the whole quantitative story of the appendix.
\end{notationbox}

\subsection{Addition is not associative}
\label{sec:nonassoc}

The single fact that drives most of this appendix is that floating-point addition,
though commutative, is \emph{not associative}:
\begin{equation}
\label{eq:nonassoc}
  \mathrm{fl}\bigl(\mathrm{fl}(a + b) + c\bigr)
  \;\ne\;
  \mathrm{fl}\bigl(a + \mathrm{fl}(b + c)\bigr)
  \quad\text{in general.}
\end{equation}
Each inner sum is rounded before the outer sum sees it, and the two roundings
discard different bits. A three-number example makes it concrete. In a toy decimal
format carrying three significant digits, take $a = 1.00$, $b = c = 5.00\times
10^{-3}$. Left-associating, $\mathrm{fl}(a+b) = \mathrm{fl}(1.005) = 1.00$ (the
$0.005$ falls off the third digit), and then $\mathrm{fl}(1.00 + 0.005) = 1.00$:
both small terms are lost. Right-associating, $\mathrm{fl}(b+c) =
\mathrm{fl}(0.010) = 0.010$ (no rounding---it fits), and then $\mathrm{fl}(1.00 +
0.010) = 1.01$: the two small terms, summed \emph{first}, survive their encounter
with the large one. Same three numbers, two bracketings, different answers in the
last digit. \Cref{fig:nonassoc} draws the two summation trees and the bits each
one keeps.

\begin{figure}[t]
\centering
\begin{tikzpicture}[font=\footnotesize]
  % ===== LEFT: left-leaning tree =====
  \begin{scope}
    \node[font=\small\bfseries, text=simBlue] at (1.3,3.1) {left-associated $((a{+}b){+}c)$};
    \node[data] (a) at (0,2.2) {$a$};
    \node[data] (b) at (1.0,2.2) {$b$};
    \node[data] (c) at (2.6,2.2) {$c$};
    \node[opbox, minimum width=9mm] (s1) at (0.5,1.2) {$+$};
    \node[opbox, minimum width=9mm] (s2) at (1.55,0.2) {$+$};
    \node[product, minimum width=14mm] (r1) at (1.55,-0.9) {sum$_L$};
    \draw[flow] (a) -- (s1); \draw[flow] (b) -- (s1);
    \draw[flow] (s1) -- (s2); \draw[flow] (c) -- (s2);
    \draw[flow] (s2) -- (r1);
    \node[font=\scriptsize, text=simRust, align=center] at (1.55,-1.7)
      {$a{+}b$ rounds first:\\small bits of $b$ lost early};
  \end{scope}
  \draw[simGray!50, thick] (3.8,-2.0) -- (3.8,3.3);
  % ===== RIGHT: right-leaning tree =====
  \begin{scope}[xshift=5.2cm]
    \node[font=\small\bfseries, text=simGreen] at (1.5,3.1) {right-associated $(a{+}(b{+}c))$};
    \node[data] (a2) at (0,2.2) {$a$};
    \node[data] (b2) at (1.6,2.2) {$b$};
    \node[data] (c2) at (2.6,2.2) {$c$};
    \node[opbox, minimum width=9mm] (t1) at (2.1,1.2) {$+$};
    \node[opbox, minimum width=9mm] (t2) at (1.05,0.2) {$+$};
    \node[product, minimum width=14mm] (r2) at (1.05,-0.9) {sum$_R$};
    \draw[flow] (b2) -- (t1); \draw[flow] (c2) -- (t1);
    \draw[flow] (a2) -- (t2); \draw[flow] (t1) -- (t2);
    \draw[flow] (t2) -- (r2);
    \node[font=\scriptsize, text=simGreen, align=center] at (1.05,-1.7)
      {$b{+}c$ summed first:\\small terms survive together};
  \end{scope}
  % the verdict band across the bottom
  \node[font=\scriptsize\itshape, text=simInk, align=center] at (3.8,-2.6)
    {$a=1.00,\ b=c=5{\times}10^{-3}$ (3-digit decimal):\quad
     sum$_L = 1.00 \ne 1.01 = $ sum$_R$};
\end{tikzpicture}
\caption[Two summation trees, two answers]{The non-associativity of
\eqref{eq:nonassoc} made visible. The same three numbers are summed by two
\emph{reduction trees}. \emph{Left:} left-association adds $a+b$ first; in a coarse
format the small term $b$ is swallowed by the large $a$ before $c$ is ever added,
and its contribution is lost. \emph{Right:} right-association adds the two small
terms $b+c$ first, so they accumulate to a magnitude that survives the encounter
with $a$. The two trees compute the identical mathematical sum $a+b+c$ but disagree
in the last digit, because rounding happens \emph{between} the additions and each
tree rounds at different places. A reduction is a tree, and the tree is observable
in the result.}
\label{fig:nonassoc}
\end{figure}

This is not pathology; it is the rule. A sum of $N$ numbers can be evaluated by any
of exponentially many trees---sequential left-to-right, pairwise (a balanced
binary tree), blocked, or whatever a parallel reduction happens to produce---and
the trees agree only in exact arithmetic. In floating point they agree to within a
bound of order $N u$ (sequential) or $u\log_2 N$ (pairwise) times the sum of
magnitudes, but they do not agree \emph{bit for bit}. \Cref{sec:reduction} returns
to this with the global reductions of CG and GMRES, where it becomes the central
load-bearing fact.

\subsection{Catastrophic cancellation}
\label{sec:cancel}

The second hazard is \emph{subtraction of nearly equal quantities}. When $a$ and
$b$ agree in their leading digits, the difference $a - b$ \emph{cancels} those
leading digits, and the result is built entirely from the low-order bits---the very
bits already corrupted by the roundings that produced $a$ and $b$. The relative
error in $a - b$ is magnified by roughly $\abs{a}/\abs{a-b}$, which is enormous
when the operands are close. \Cref{fig:cancel} shows the bookkeeping.

\begin{figure}[t]
\centering
\begin{tikzpicture}[font=\footnotesize, x=0.34cm, y=0.5cm]
  % a row of "digit cells": green = correct, gray = garbage
  \def\cell#1#2#3{\draw[fill=#3] (#1,#2) rectangle ++(0.9,0.8);}
  % a = 1.2345678  (all 7 digits correct, last is rounding-uncertain)
  \node[anchor=east, font=\scriptsize] at (-0.4,3.4) {$a = 1.2345678\underline{9}$};
  \foreach \i in {0,...,6}{\cell{\i*1}{3}{simBgGreen}}
  \cell{7}{3}{simBgGray}
  % b = 1.2345677  (agree in first 7)
  \node[anchor=east, font=\scriptsize] at (-0.4,2.2) {$b = 1.2345677\underline{2}$};
  \foreach \i in {0,...,6}{\cell{\i*1}{2}{simBgGreen}}
  \cell{7}{2}{simBgGray}
  % a-b = 0.0000001.7 : leading 7 cancel, only the last (garbage) survives, shifted up
  \node[anchor=east, font=\scriptsize] at (-0.4,0.9) {$a - b = 0.000000\underline{17}$};
  \foreach \i in {0,...,5}{\cell{\i*1}{0.6}{white}}
  \cell{6}{0.6}{simBgRust}
  \cell{7}{0.6}{simBgRust}
  % brace over the cancelled region
  \draw[decorate, decoration={brace, amplitude=4pt}, simGray]
    (6.9,3.95) -- (-0.05,3.95);
  \node[font=\scriptsize, text=simGray] at (3.4,4.45) {7 leading digits cancel};
  \draw[decorate, decoration={brace, amplitude=4pt, mirror}, simRust]
    (5.95,0.45) -- (7.85,0.45);
  \node[font=\scriptsize, text=simRust, align=center] at (9.6,0.9)
    {result built only from\\the \emph{uncertain} trailing bits};
  % legend
  \node[anchor=west, font=\scriptsize] at (0,-1.1)
    {\tikz\fill[simBgGreen](0,0)rectangle(0.3,0.3); correct\quad
     \tikz\fill[simBgGray](0,0)rectangle(0.3,0.3); rounding-uncertain\quad
     \tikz\fill[simBgRust](0,0)rectangle(0.3,0.3); amplified garbage};
\end{tikzpicture}
\caption[Catastrophic cancellation]{Subtracting two nearly equal numbers destroys
significant digits. The operands $a$ and $b$ agree in their first seven digits
(green); only their trailing digit (gray) carries rounding uncertainty. Their
difference $a-b$ \emph{cancels} the seven agreeing digits exactly, leaving a result
whose significant figures are built entirely from the previously-uncertain trailing
bits (rust)---which are then re-normalized up to the leading position, magnifying
their relative error by roughly $\abs{a}/\abs{a-b}$. The subtraction itself is
exactly rounded (its own $\delta$ is below $u$); the damage was done \emph{earlier},
when $a$ and $b$ were rounded, and cancellation merely exposes it. The cure is
always algebraic---reformulate to avoid the near-equal subtraction---never a more
accurate subtractor.}
\label{fig:cancel}
\end{figure}

\begin{pitfall}
Cancellation is widely \emph{mis}-diagnosed. The subtraction $a-b$ is itself
exactly rounded---its relative error is below $u$ like any other operation. What is
catastrophic is that it \emph{reveals} errors already present in $a$ and $b$ from
earlier roundings. Replacing the subtraction with a ``more accurate'' one buys
nothing; the fix is always to reformulate the computation so the near-equal
subtraction never occurs---summing positive terms before subtracting, using a
stable recurrence instead of an explicit difference (the Chebyshev story of
\cref{sec:cheby}), or rationalizing $\sqrt{x+1}-\sqrt{x}$ into
$1/(\sqrt{x+1}+\sqrt{x})$. The monomial form of a polynomial
(\cref{sec:cheby}) is catastrophic cancellation industrialized: large
alternating-sign terms that should sum to something small.
\end{pitfall}

\subsection{Conditioning: the problem versus the algorithm}
\label{sec:conditioning}

The last piece of vocabulary separates two sources of error that are constantly
conflated. The \emph{condition number} of a problem measures how much its
\emph{output} can change in response to a small change in its \emph{input}---a
property of the problem itself, before any algorithm is chosen.

\begin{definition}[Condition number of a problem]
\label{def:cond}
For a problem $y = f(x)$, the (relative) condition number is the smallest $\kappa$
such that a relative perturbation $\abs{\Delta x}/\abs{x}$ produces a relative
output change $\abs{\Delta y}/\abs{y} \lesssim \kappa\,\abs{\Delta x}/\abs{x}$. For
solving $\mat{A}\vx = \vb$ the relevant condition number is
$\kappa(\mat{A}) = \norm{\mat{A}}\,\norm{\mat{A}^{-1}}$, which for SPD $\mat{A}$ is
$\lambda_{\max}/\lambda_{\min}$ (\cref{def:kappa}).
\end{definition}

An algorithm is \emph{stable} (more precisely, \emph{backward stable}) if the
answer it computes is the exact answer to a slightly perturbed input---a
perturbation of relative size $O(u)$. Stability and conditioning then multiply:
\begin{equation}
\label{eq:errorrule}
  \underbrace{\frac{\norm{\hat{\vx} - \vx}}{\norm{\vx}}}_{\text{forward error}}
  \;\lesssim\;
  \underbrace{\kappa(\mat{A})}_{\text{conditioning of problem}}
  \times
  \underbrace{O(u)}_{\text{stability of algorithm}} .
\end{equation}
This single inequality is the organizing rule of numerical computing. A stable
algorithm on an ill-conditioned problem still returns a large forward error---not
because the algorithm misbehaved but because the problem amplifies any
perturbation, including the unavoidable $O(u)$ rounding of the input. An unstable
algorithm wastes accuracy even on a well-conditioned problem. The two failure modes
are independent, and \eqref{eq:errorrule} tells which is which: if $\kappa$ is huge,
no algorithm will save you (precondition the problem---\cref{ch:cg}); if $\kappa$ is
modest but the answer is still wrong, the algorithm is unstable (fix the algorithm,
as MGS fixes CGS in \cref{sec:orthog}).

\begin{keyidea}
Conditioning is a property of the \emph{problem}; stability is a property of the
\emph{algorithm}. Forward error $\lesssim \kappa \cdot O(u)$: the conditioning of
the problem times the (in)stability of the algorithm. The ill-conditioned operators
Palace assembles (fine meshes, high material contrast) make $\kappa$ large, which
is exactly why preconditioning (\cref{ch:cg}) and stable orthogonalization
(\cref{ch:orthog}) earn their keep---they attack the two factors separately.
\end{keyidea}

\section{Transparent tricks: invisible to the mathematics}
\label{sec:transparent}

With the primer in hand we can state the appendix's central distinction precisely
and then dispatch the easy half of it. A \emph{transparent} performance trick is a
rewrite that computes the \emph{same result}---bit for bit, or up to a benign
reordering whose error is of the same harmless order as the computation
already carries---but faster. The book has presented the unfolded, mathematically
transparent form throughout and relegated each such trick to a one-line note,
because by definition it changes nothing the mathematics can see.

\begin{definition}[Transparent performance trick]
\label{def:transparent}
A code transformation is \emph{transparent} if it is algebraically equivalent to
the unfolded form it replaces and does not change any property the algorithm
depends on---not the value (up to benign reordering), not the conditioning, not the
rounding behavior in any consumed way. Its only effect is speed or memory traffic.
A transparent trick may be \emph{ignored in the mathematical presentation} and
recorded as a footnote on the operation it accelerates.
\end{definition}

The catalogue of transparent tricks Palace uses, each algebraically equivalent to
its unfolded form:

\begin{itemize}[leftmargin=1.4em]
\item \textbf{Operator fusion.} Two consecutive elementwise passes (say a scale
  then an add) are fused into one loop that reads the input once and writes the
  output once, halving memory traffic. The arithmetic performed is identical; only
  the number of trips to memory changes. The \code{axpy} family of \cref{ch:cg} is
  a fused scale-and-add presented as one operator for exactly this reason.

\item \textbf{Tiling and blocking.} A large matrix operation is partitioned into
  cache-sized tiles processed one at a time, so the working set fits in fast memory.
  The set of multiply-adds is unchanged; only their \emph{schedule} changes. (When
  the schedule changes the \emph{reduction order} of a sum, the trick is no longer
  purely transparent---see \cref{sec:reduction}; blocking that keeps each reduction
  intact is transparent, blocking that re-trees a reduction is load-bearing.)

\item \textbf{Packing and batching.} Many small elementwise or per-element
  operations are gathered into one large contiguous kernel launch (the
  element-batched assembly of \cref{ch:matrixfree}). The per-element arithmetic is
  untouched; batching only amortizes launch and indexing overhead.

\item \textbf{Memory layout.} Choosing array-of-structs versus struct-of-arrays,
  or a blocked sparse format, changes which bytes are adjacent but not which
  numbers are computed.

\item \textbf{Recomputation versus lookup.} A cheap quantity (a Jacobian
  determinant at a quadrature point) may be recomputed on the fly rather than
  stored and reloaded, trading arithmetic for memory bandwidth. The recomputed
  value is bit-identical to the stored one (it is the same formula), so the choice
  is transparent.

\item \textbf{Sum-factorization} (\cref{ch:matrixfree}). High-order matrix-free
  operator application replaces the dense $O(P\cdot L)$ basis contraction by a
  sequence of one-dimensional sweeps, $O(d\,P^{1/d}L)$, exploiting the
  tensor-product structure of the basis. As \cref{ch:matrixfree} establishes at
  length, the two compute \emph{bit-for-bit the same quadrature-point values}---the
  product $\phi_i(\xi)\phi_j(\eta)$ is simply contracted one axis at a time---so
  sum-factorization is the canonical transparent trick: decisive for performance,
  invisible to the algebra.
\end{itemize}

\begin{keyidea}
A transparent trick is a \emph{factoring of the schedule}, not of the algorithm.
Fusion, tiling, packing, layout, recomputation, and sum-factorization all leave the
computed numbers unchanged (up to benign reordering) and change only how fast or
with how much memory traffic they are produced. The mathematics of the book
therefore presents the unfolded form and notes the trick in one line. The reader
who wants to understand \emph{what is computed} may ignore every transparent trick;
the reader who wants to understand \emph{why it is fast enough to run} should read
\cref{ch:matrixfree}.
\end{keyidea}

There is one boundary case worth flagging explicitly, because it is the bridge to
the next section. Tiling, blocking, and batching are transparent \emph{only so long
as they preserve the reduction order of every sum}. The moment a tiling re-brackets
a global reduction---computes partial sums per tile and then combines them in a
different tree than the sequential sweep would---it has crossed from transparent to
load-bearing, because by \cref{sec:nonassoc} it changes the last bits of the sum.
The same code transformation can be transparent or load-bearing depending on
whether a reduction tree hides inside it. This is precisely why the classification
is a matter of \emph{judgment}, flagged to the human when unclear (\cref{ch:bigidea}),
rather than a mechanical property of the code shape.

\section{Load-bearing tricks: part of the algorithm}
\label{sec:loadbearing}

A \emph{load-bearing} numerical trick is one that changes a property the algorithm
depends on---and therefore cannot be ignored in the mathematics, because ignoring
it would mis-state what the algorithm computes. Where a transparent trick gets a
footnote, a load-bearing trick gets an \emph{explicit algebraic claim}, naming the
property it buys.

\begin{definition}[Load-bearing numerical trick]
\label{def:loadbearing}
A numerical choice is \emph{load-bearing} if changing it changes a property of the
result that some consumer depends on: its \emph{value} in a way that matters (the
exact iteration count, the last bits when reproducibility is required), its
\emph{conditioning} or attained accuracy, its \emph{determinism}, or its
\emph{IEEE compliance}. A load-bearing trick must be stated in the specification as
an explicit claim, together with the property it secures, and may not be silently
``optimized'' into an algebraically-equivalent-in-exact-arithmetic alternative.
\end{definition}

\Cref{fig:taxonomy} sets the two kinds side by side as a decision the spec makes for
every numerical choice, and the rest of the section works through the load-bearing
column case by case.

\begin{figure}[t]
\centering
\begin{tikzpicture}[font=\footnotesize]
  % the input question
  \node[stage, minimum width=46mm] (q) at (0,2.4)
    {A numerical choice (reorder, reassociate,\\reduce precision, expand a polynomial)};
  % left branch: transparent
  \node[opbox, fill=simBgGreen, draw=simGreen, text width=42mm, align=left] (T) at (-3.6,0)
    {\textbf{Transparent}\\[2pt]
     same result (up to benign reorder),\\just faster: fusion, tiling, packing,\\
     layout, recompute, sum-factorization};
  \node[product, fill=simBgGreen, draw=simGreen, text width=34mm] (Tout) at (-3.6,-2.3)
    {\textbf{ignore in the math};\\ one-line note on the operation};
  % right branch: load-bearing
  \node[opbox, fill=simBgRust, draw=simRust, text width=44mm, align=left] (L) at (3.6,0)
    {\textbf{Load-bearing}\\[2pt]
     changes value / conditioning /\\determinism / IEEE compliance:\\
     reduction order, fast-math, mixed\\precision, recurrence-vs-monomial, MGS-vs-CGS};
  \node[product, fill=simBgRust, draw=simRust, text width=36mm] (Lout) at (3.6,-2.3)
    {\textbf{state as an algebraic claim};\\ name the property it buys};
  % edges
  \draw[flow] (q.south) -- (T.north)
    node[midway, above left, font=\scriptsize, text=simGreen]
      {same answer?};
  \draw[flow] (q.south) -- (L.north)
    node[midway, above right, font=\scriptsize, text=simRust]
      {changes a property?};
  \draw[flow] (T) -- (Tout);
  \draw[flow] (L) -- (Lout);
  % unclear -> human
  \node[extern, text width=30mm] (H) at (0,-3.7)
    {\emph{unclear}: flag to the human (\cref{ch:bigidea})};
  \draw[refedge] (T.south) .. controls +(0,-0.9) and +(-0.7,0.3) .. (H.west);
  \draw[refedge] (L.south) .. controls +(0,-0.9) and +(0.7,0.3) .. (H.east);
\end{tikzpicture}
\caption[Transparent versus load-bearing: the spec's decision]{The classification
every numerical choice passes through. \emph{Left (green):} if the choice computes
the same result up to a benign reordering and changes no depended-on property, it is
\emph{transparent}---ignored in the mathematical presentation, recorded as a one-line
note on the operation it accelerates. \emph{Right (rust):} if it changes the value
in a consumed way, the conditioning, the determinism, or the IEEE compliance, it is
\emph{load-bearing}---stated as an explicit algebraic claim naming the property it
secures. The hazard is mis-routing a load-bearing trick down the transparent path,
which silently changes the algorithm; when the routing is unclear the spec flags it
to the human (\cref{ch:bigidea}) rather than guessing. The five load-bearing cases
named on the right are the subject of \crefrange{sec:reduction}{sec:orthog}.}
\label{fig:taxonomy}
\end{figure}

\subsection{Non-associative reduction orderings}
\label{sec:reduction}

The first and most pervasive load-bearing trick is the one \cref{sec:numerics} of
\cref{ch:cg} named and deferred to here: the \emph{order} in which a global
reduction is summed. Every Krylov coefficient---the CG step length $\alpha$, the
direction coefficient $\beta$, the GMRES least-squares residual, every convergence
test---is built from inner products $\inner{\vu}{\vv} = \sum_i u_i v_i$, and an
inner product is a global reduction over all $N$ entries. By the non-associativity
of \cref{sec:nonassoc}, the value of that reduction depends on the summation tree.

The summation tree is not a fixed property of the program; it is chosen by the
runtime. A sequential loop sums left-to-right; a multithreaded reduction splits the
range into per-thread partial sums and combines them; a GPU reduction uses a
tree of warps and blocks; a different thread count or a different device produces a
different tree. Each tree gives a bit-different but equally valid value of the inner
product. \Cref{fig:redtree} contrasts two trees over the same eight terms.

\begin{figure}[t]
\centering
\begin{tikzpicture}[font=\footnotesize, scale=0.95]
  % ===== LEFT: sequential (left-leaning) reduction =====
  \begin{scope}
    \node[font=\small\bfseries, text=simBlue] at (1.6,3.5) {sequential: one long chain};
    \foreach \i [count=\j from 0] in {0,1,2,3,4,5,6,7}{
      \node[data, minimum width=4mm, inner sep=1.5pt] (x\i) at (\j*0.62,2.7)
        {\scriptsize$\i$};
    }
    % left-leaning accumulator chain
    \node[opbox, minimum width=4mm, inner sep=1.5pt] (acc) at (0.0,1.0) {\scriptsize$+$};
    \foreach \i in {1,...,7}{
      \pgfmathsetmacro\xx{\i*0.45}
      \node[opbox, minimum width=4mm, inner sep=1.5pt] (a\i) at (\xx,{1.0-\i*0.22}) {\scriptsize$+$};
    }
    \draw[flow] (x0)--(acc); \draw[flow] (x1)--(acc);
    \draw[flow] (acc)--(a2);
    \foreach \i [evaluate=\i as \p using int(\i-1)] in {2,...,7}{
      \ifnum\i>2 \draw[flow] (a\p)--(a\i); \fi
      \draw[flow] (x\i)--(a\i);
    }
    \node[product, minimum width=12mm] (rL) at (3.15,-0.9) {sum$_{\text{seq}}$};
    \draw[flow] (a7)--(rL);
    \node[font=\scriptsize, text=simRust, align=center] at (1.6,-1.7)
      {error $\sim N u$;\\ depth $N$, not parallel};
  \end{scope}
  \draw[simGray!50, thick] (4.6,-2.0) -- (4.6,3.7);
  % ===== RIGHT: pairwise (balanced tree) reduction =====
  \begin{scope}[xshift=5.6cm]
    \node[font=\small\bfseries, text=simGreen] at (2.2,3.5) {pairwise: balanced tree};
    \foreach \i [count=\j from 0] in {0,1,2,3,4,5,6,7}{
      \node[data, minimum width=4mm, inner sep=1.5pt] (y\i) at (\j*0.62,2.7)
        {\scriptsize$\i$};
    }
    % level 1: 4 adds
    \foreach \k [count=\m from 0] in {0,2,4,6}{
      \node[opbox, minimum width=4mm, inner sep=1.5pt] (l1\m) at ({(\k+0.5)*0.62},1.7) {\scriptsize$+$};
    }
    \draw[flow] (y0)--(l10); \draw[flow] (y1)--(l10);
    \draw[flow] (y2)--(l11); \draw[flow] (y3)--(l11);
    \draw[flow] (y4)--(l12); \draw[flow] (y5)--(l12);
    \draw[flow] (y6)--(l13); \draw[flow] (y7)--(l13);
    % level 2: 2 adds
    \node[opbox, minimum width=4mm, inner sep=1.5pt] (l20) at (1.24,0.7) {\scriptsize$+$};
    \node[opbox, minimum width=4mm, inner sep=1.5pt] (l21) at (3.10,0.7) {\scriptsize$+$};
    \draw[flow] (l10)--(l20); \draw[flow] (l11)--(l20);
    \draw[flow] (l12)--(l21); \draw[flow] (l13)--(l21);
    % level 3: 1 add
    \node[product, minimum width=12mm] (rR) at (2.17,-0.9) {sum$_{\text{tree}}$};
    \draw[flow] (l20)--(rR); \draw[flow] (l21)--(rR);
    \node[font=\scriptsize, text=simGreen, align=center] at (2.2,-1.7)
      {error $\sim u\log_2 N$;\\ depth $\log_2 N$, parallel};
  \end{scope}
\end{tikzpicture}
\caption[Two reduction trees over the same data]{The same eight-term reduction
$\sum_{i=0}^{7} z_i$ computed by two trees. \emph{Left:} a sequential accumulator
chain---the natural left-to-right loop---has depth $N$ and a worst-case error bound
of order $N u$ times the sum of magnitudes; it cannot be parallelized. \emph{Right:}
a balanced pairwise tree has depth $\log_2 N$, a tighter error bound of order
$u\log_2 N$, and maps onto parallel hardware. Both are correct reductions of the
same data; they round at different places and so generally produce \emph{different
last bits}. A multithreaded or GPU run picks a tree determined by its thread and
block geometry, so ``the same'' inner product computed on two devices---or with two
thread counts---returns bit-different values, and the iterative solver downstream
takes a bit-different (but equally valid) trajectory.}
\label{fig:redtree}
\end{figure}

The consequence for an iterative solver is precise and, at first, unsettling. Two
runs of the ``same'' CG solve on two machines---or on one machine with two thread
counts---compute bit-different inner products, hence bit-different $\alpha$ and
$\beta$, hence iterates differing in their last bits, and may even cross the
convergence tolerance at \emph{different iteration numbers}. The converged answer is
correct to tolerance in both cases; what is \emph{not} reproducible bit-for-bit is
the trajectory and the exact step count. This is the canonical load-bearing trick:
not a bug, but a property of the computation that a consumer may or may not depend
on, and that the specification must therefore \emph{name}.

\begin{keyidea}
An inner product is a reduction; a reduction is a tree; the tree determines the last
bits. Reduction order is \emph{load-bearing} precisely because it is observable: it
changes the bit-level trajectory and the exact iteration count of every Krylov
solver. The mathematics is identical across trees---each is a correct sum---so the
spec does not pick a tree, but it must \emph{state} that the trajectory is
tree-dependent, so that a downstream consumer knows reproducibility is not free.
\end{keyidea}

\begin{pitfall}
\textbf{Bit-reproducibility is not free, and is not on by default.} A consumer that
demands the identical bit pattern from run to run---a regression test asserting on
the last digits, a certification requiring reproducible output---must \emph{pin the
reduction order}: fix the thread count, fix the device, use a deterministic
(fixed-tree or compensated) reduction. Each of these costs performance: a
deterministic reduction forgoes the fastest tree the hardware offers, and a
fixed thread count forgoes load balancing. CG does not, and cannot, promise
bit-reproducibility for free---the same property that lets it run fast on any device
geometry is the property that makes its last bits device-dependent. The trade is
explicit: \emph{determinism costs throughput}, and the spec records reduction order
as the place that bill comes due.
\end{pitfall}

\subsection{Fast-math and reassociation flags}
\label{sec:fastmath}

A compiler in its default mode honors \eqref{eq:fpmodel} operation by operation and
preserves the parenthesization the source code wrote---it will not turn
$(a+b)+c$ into $a+(b+c)$, because by \cref{sec:nonassoc} that would change the
result. The various \emph{fast-math} flags (\code{-ffast-math},
\code{-funsafe-math-optimizations}, and their relatives) lift this restriction:
they license the compiler to treat floating-point arithmetic as if it were
associative, distributive, and exact, enabling reassociation, vectorization across
reduction boundaries, contraction of $a\cdot b + c$ into a fused multiply-add,
flushing denormals to zero, and assuming no NaNs or infinities appear.

\begin{table}[t]
\centering
\caption[What fast-math buys and breaks]{The fast-math bargain. Each license the
flag grants buys a performance opportunity and breaks a guarantee. The flags are
\emph{load-bearing}: enabling them changes the computed last bits and can change
which branch a NaN-guard takes, so they are a specification decision, not a free
optimization. Palace and the solvers of this book are compiled \emph{without}
unsafe fast-math precisely because the reduction-order and recurrence-stability
properties of \crefrange{sec:reduction}{sec:orthog} depend on the standard model
holding.}
\label{tab:fastmath}
\small
\setlength{\tabcolsep}{6pt}
\renewcommand{\arraystretch}{1.25}
\begin{tabular}{@{}p{0.30\linewidth}p{0.30\linewidth}p{0.32\linewidth}@{}}
\toprule
License granted & Buys & Breaks \\
\midrule
Reassociation ($+,\times$ associative) & vectorized / tree reductions, loop
  fusion & bit-reproducibility; a pinned reduction order (\cref{sec:reduction}) \\
Fused multiply-add contraction & one rounding instead of two; throughput &
  the exact $\mathrm{fl}(a\cdot b)$ a stable formula may rely on \\
Flush denormals to zero & no slow-path denormal handling & gradual underflow; tiny
  residuals snap to $0$ \\
Assume no NaN / Inf & cheaper comparisons, no guard code & NaN-propagation
  diagnostics; divergence detection \\
\bottomrule
\end{tabular}
\end{table}

\Cref{tab:fastmath} lays out the trade. The point for the specification is that
fast-math reassociation is \emph{exactly} the reduction-order freedom of
\cref{sec:reduction}, granted to the compiler rather than the runtime: it lets the
compiler pick a summation tree the source did not write. That makes it load-bearing
for the same reason---it changes last bits and reproducibility---and additionally
hazardous because the ``assume no NaN/Inf'' license can silently delete the very
guard code that an algorithm uses to detect divergence or a singular operator. A
solver that relies on a NaN propagating to its convergence check (a crude but real
divergence detector) will, under fast-math, march past the divergence unnoticed.
The disposition in this book: the standard model \eqref{eq:fpmodel} is assumed to
hold, unsafe fast-math is off, and any place that genuinely wants a fused
multiply-add or a relaxed reduction states it explicitly as the load-bearing choice
it is.

\subsection{Mixed precision: single apply, double accumulate}
\label{sec:mixed}

The third load-bearing trick trades precision for speed deliberately and in a
controlled place. Many operator applies and elementwise kernels can be carried in
\emph{single} precision---halving memory traffic and, on GPUs, often doubling or
quadrupling arithmetic throughput---provided the \emph{accumulation} that follows is
carried in \emph{double}. The pattern is ``single-precision apply, double-precision
accumulate,'' and it is the workhorse of modern mixed-precision iterative solvers.

\begin{figure}[t]
\centering
\begin{tikzpicture}[font=\footnotesize, node distance=8mm and 13mm]
  % single-precision lane (top)
  \node[data, fill=simBgGray, minimum width=14mm] (vin) {$\vv$\\\scriptsize fp32};
  \node[opbox, fill=simBgRust, draw=simRust, right=of vin, minimum width=22mm]
    (apply) {operator apply\\$\mat{A}\vv$\\\scriptsize \textbf{fp32}, fast};
  \node[data, fill=simBgGray, minimum width=16mm, right=of apply] (prod)
    {products $u_i v_i$\\\scriptsize fp32};
  \draw[flow] (vin) -- (apply);
  \draw[flow] (apply) -- (prod);
  % the precision-promotion boundary
  \draw[densely dashed, simViolet, thick] (prod.north east)++(0.35,0.7)
    -- ($(prod.south east)+(0.35,-0.7)$);
  \node[font=\scriptsize, text=simViolet, anchor=south, rotate=90] at
    ($(prod.east)+(0.55,0)$) {promote fp32$\to$fp64};
  % double-precision accumulate (right)
  \node[opbox, fill=simBgBlue, draw=simBlue, right=18mm of prod, minimum width=24mm]
    (acc) {reduction\\$\sum_i u_i v_i$\\\scriptsize \textbf{fp64}, accurate};
  \node[product, right=of acc, minimum width=16mm] (out) {$\alpha,\beta,\norm{\vr}$\\\scriptsize fp64};
  \draw[flow] (prod) -- (acc);
  \draw[flow] (acc) -- (out);
  % annotations
  \node[font=\scriptsize, text=simRust, align=center, below=2mm of apply]
    {bulk work:\\cheap, low precision};
  \node[font=\scriptsize, text=simBlue, align=center, below=2mm of acc]
    {the reduction:\\full precision};
\end{tikzpicture}
\caption[Single-precision apply, double-precision accumulate]{The mixed-precision
dataflow. The expensive bulk work---the operator apply $\mat{A}\vv$ and the
elementwise products $u_i v_i$ (rust)---is carried in single precision (fp32),
where it is cheap in both memory bandwidth and throughput. At the
\emph{accumulation} boundary (violet, dashed) the products are promoted to double
precision (fp64), and the global reduction $\sum_i u_i v_i$ that forms each Krylov
coefficient is summed in double (blue). This is load-bearing, not transparent: the
result is \emph{not} the bit-equal double-precision answer, and the soundness of the
trade rests on the conditioning argument of \cref{sec:conditioning}---the
single-precision apply injects a relative error of order $u_{32}\sim 6\times10^{-8}$
that the conditioning of the problem amplifies, so the trick is safe only while
$\kappa\,u_{32}$ stays below the target tolerance.}
\label{fig:mixed}
\end{figure}

Why accumulate in double? Because the accumulation is the global reduction of
\cref{sec:reduction}, and a reduction of $N$ single-precision terms accumulates an
error of order $N u_{32}$---which, for $N$ in the millions and $u_{32}\approx
6\times10^{-8}$, overwhelms the answer. Promoting the running sum to double drops
the accumulation error to order $N u_{64} \approx N\cdot10^{-16}$, negligible again,
while the bulk arithmetic stays cheap. The single-precision \emph{apply} still
injects a relative error of order $u_{32}$ into each matrix-vector product, and here
the conditioning rule \eqref{eq:errorrule} governs: that error is amplified by
$\kappa(\mat{A})$, so the mixed-precision solve attains a forward accuracy of order
$\kappa\,u_{32}$ rather than $\kappa\,u_{64}$. The trade is sound \emph{exactly
when} $\kappa\,u_{32}$ is comfortably below the solve tolerance---a well-conditioned
or well-preconditioned system can run its applies in single and lose nothing
visible; an ill-conditioned one cannot.

\begin{keyidea}
Mixed precision is the conditioning rule \eqref{eq:errorrule} turned into an
engineering knob. Do the \emph{bulk} arithmetic (the apply, the products) in single
precision for speed, but carry the \emph{accumulation} (the reduction) in double so
the global sum does not lose its low bits. The result is load-bearing---it is not
the double-precision answer---and it is sound only while $\kappa\,u_{32}$ stays below
tolerance. This is why mixed-precision iterative refinement and single/double Krylov
solvers pair so naturally with preconditioning (\cref{ch:cg}): the preconditioner
shrinks $\kappa$, which is exactly what buys the headroom to drop the apply to
single precision.
\end{keyidea}

\subsection{The Chebyshev recurrence versus the monomial form}
\label{sec:cheby}

The fourth load-bearing trick is a case where two algebraically identical
evaluations of the \emph{same polynomial} differ by everything that matters
numerically. The Chebyshev smoother of \cref{ch:smoothers} applies a fixed-degree
polynomial $p_k(\mat{D}^{-1}\mat{A})$ to the residual. There are two ways to
evaluate it, and the choice between them is the algorithm.

The \emph{monomial form} writes the polynomial out in powers and sums them:
\begin{equation}
\label{eq:monomial}
  p_k(\mat{B})\,\vr \;=\; \sum_{j=0}^{k} c_j\,\mat{B}^{j}\vr,
  \qquad \mat{B} = \mat{D}^{-1}\mat{A},
\end{equation}
forming $\mat{B}\vr, \mat{B}^2\vr, \dots$ and accumulating with the coefficients
$c_j$. The \emph{three-term recurrence} (a Clenshaw/Horner-style evaluation) never
forms the powers or the coefficients; it threads three vectors through a short loop,
as \cref{ch:smoothers} derives. In exact arithmetic the two compute the identical
vector. In floating point they could not be more different.

The monomial form \eqref{eq:monomial} is numerically catastrophic for the degrees in
use. The high powers $\mat{B}^{j}\vr$ grow enormous (their size is governed by the
largest eigenvalue raised to the $j$), the coefficients $c_j$ alternate in sign and
also grow large, and the sum of large alternating-sign terms collapsing to a modest
result is catastrophic cancellation (\cref{sec:cancel}) at industrial scale: the
answer is built from the differences of huge numbers whose low bits are pure
rounding noise. The three-term recurrence sidesteps this entirely---it evaluates the
same polynomial through a sequence of moderate, well-scaled intermediate vectors,
never forming a large power or a large coefficient, so no cancellation occurs. This
is the stability that Phillips and Fischer analyze and that Palace follows
directly~\citep[\S2]{phillipsfischer2022}. \Cref{fig:cheby} plots the divergence.

\begin{figure}[t]
\centering
\begin{tikzpicture}
\begin{semilogyaxis}[
  width=12.0cm, height=6.6cm,
  xlabel={polynomial degree $k$},
  ylabel={relative error $\norm{\hat{p}_k(\mat{B})\vr - p_k(\mat{B})\vr}/\norm{p_k(\mat{B})\vr}$},
  ymode=log, ymin=1e-17, ymax=1e6,
  xmin=1, xmax=20,
  grid=both, grid style={simGray!18},
  legend pos=north west, legend cell align=left,
  legend style={font=\footnotesize, draw=simGray!40},
  label style={font=\footnotesize}, tick label style={font=\scriptsize},
  ytick={1e-16,1e-12,1e-8,1e-4,1e0,1e4},
  every axis plot/.append style={very thick}]
  % monomial: error blows up roughly like (cond)^k * u -> exponential growth
  \addplot[simRust, mark=*, mark size=1.2pt] coordinates {
    (1,2e-16)(2,1e-15)(3,8e-15)(4,9e-14)(5,1e-12)(6,2e-11)(7,3e-10)(8,5e-9)
    (9,8e-8)(10,1e-6)(11,2e-5)(12,3e-4)(13,5e-3)(14,8e-2)(15,1.0)(16,1e1)
    (17,2e2)(18,4e3)(19,7e4)(20,1e6)};
  \addlegendentry{monomial $\sum_j c_j \mat{B}^j\vr$ (unstable)}
  % recurrence: stays flat at machine precision (grows only mildly with k)
  \addplot[simGreen, mark=triangle*, mark size=1.4pt] coordinates {
    (1,2e-16)(2,2e-16)(3,3e-16)(4,3e-16)(5,4e-16)(6,4e-16)(7,5e-16)(8,5e-16)
    (9,6e-16)(10,6e-16)(11,7e-16)(12,7e-16)(13,8e-16)(14,8e-16)(15,9e-16)
    (16,9e-16)(17,1.0e-15)(18,1.1e-15)(19,1.2e-15)(20,1.3e-15)};
  \addlegendentry{three-term recurrence (stable)}
  % machine-epsilon reference + the "answer is garbage" line at 1.0
  \addplot[simGray, densely dashed, domain=1:20, samples=2] {2.2e-16};
  \node[font=\scriptsize, text=simGray, anchor=west] at (axis cs:1.2,5e-16) {machine $u$};
  \addplot[simInk, densely dotted, domain=1:20, samples=2] {1.0};
  \node[font=\scriptsize, text=simInk, anchor=west] at (axis cs:1.2,2.2) {no digits left};
\end{semilogyaxis}
\end{tikzpicture}
\caption[Chebyshev recurrence versus monomial evaluation]{Relative error of the two
evaluations of the same degree-$k$ Chebyshev polynomial, as the degree grows, on a
log scale (schematic of the Phillips--Fischer experiment). The \emph{monomial form}
(rust) accumulates error that grows roughly geometrically in $k$---the high powers
$\mat{B}^j\vr$ and the alternating coefficients $c_j$ produce catastrophic
cancellation (\cref{sec:cancel}), and by a modest degree the result has \emph{no
correct digits left} (the dotted line at relative error $1$). The \emph{three-term
recurrence} (green) evaluates the identical polynomial through well-scaled
intermediate vectors and stays pinned near machine precision $u$ for every degree.
The two forms are equal in exact arithmetic and not interchangeable in floating
point: the sequentiality of the recurrence is load-bearing, not a convenience, and
choosing the monomial form ``to vectorize the powers'' silently destroys the
smoother~\citep{phillipsfischer2022}.}
\label{fig:cheby}
\end{figure}

\begin{pitfall}
The recurrence form looks like a \emph{sequential bottleneck} begging to be
parallelized: each step reads the previous step's vectors. The temptation is to
``unroll'' it into the monomial form \eqref{eq:monomial}, where the powers
$\mat{B}^j\vr$ can be precomputed independently. \emph{Do not}---this is the same
trap as replacing MGS with CGS (\cref{sec:orthog}). The sequentiality of the
recurrence is not an obstacle to the algorithm; it \emph{is} the algorithm's
stability. Expanding to monomials trades a true sequential dependency for
catastrophic cancellation, computing the same polynomial in exact arithmetic and
garbage in floating point. The spec records the recurrence form as a load-bearing
algebraic claim (``evaluate $p_k$ by the three-term recurrence, not the monomial
expansion''), with the property it buys (floating-point stability) named
explicitly.
\end{pitfall}

\subsection{Modified versus classical Gram--Schmidt}
\label{sec:orthog}

The fifth load-bearing trick closes the loop with \cref{ch:orthog}, where it was
derived in full; here we file it in the taxonomy. Orthogonalizing a vector against a
basis can subtract the projections \emph{all at once against the original}
(classical Gram--Schmidt, CGS) or \emph{one at a time against the running residual}
(modified Gram--Schmidt, MGS). The two compute the identical orthogonal residual in
exact arithmetic---both are the projection $(\mat{I} - \mat{Q}\mat{Q}^{\!\top})\vw$.
In floating point they are not interchangeable: CGS loses orthogonality
\emph{quadratically} in the basis condition number, $\eta\sim\kappa^2 u$, while MGS
loses it only \emph{linearly}, $\eta\sim\kappa u$ (\cref{ch:orthog},
\cref{def:cond} here for $\kappa$).

The structural fact, drawn in full in \cref{ch:orthog}, is that CGS's independence
(all projections read the same input, so they fan out in parallel) is exactly what
makes it unstable, and MGS's sequential dependency (each projection reads the
freshly-cleaned residual) is exactly what makes it stable. The sequential dependency
is \emph{load-bearing}: it cannot be optimized into the parallel CGS form without
silently degrading the basis, and the degradation surfaces only later as a solver
that stalls or a spurious eigenvalue. The ``twice is enough'' repair (CGS2:
re-run the parallel classical pass a second time) recovers full orthogonality
$\eta\sim u$ at double the cost, recovering MGS-grade stability while keeping
batched, parallel-friendly reductions---the sweet spot of \cref{ch:orthog}'s
trade-off table. The cross-reference is the point: this is the same
transparent-versus-load-bearing decision, and MGS's column order lands firmly in the
load-bearing column, which is why \cref{ch:orthog} treats the order as a sequential
obstruction rather than an inefficiency to be removed.

\section{How the book has handled this throughout}
\label{sec:throughout}

The distinction of this appendix has been operating, quietly, in every chapter. The
governing rule was stated in \cref{ch:bigidea}: a significant fraction of Palace's
C++ exists because it was tuned for CPU cache and SIMD, a cost model that is not the
target's, and most of that code shape is a \emph{performance trick} to be separated
from the \emph{base algebra}. This appendix is the floating-point refinement of that
separation.

The mechanism throughout has been uniform. The pure-function form the book presents
(the ``L1 form,'' in the layered vocabulary) is always the \emph{unfolded} form---the
mathematics with every transparent trick removed. Sum-factorization is presented as a
dense contraction and footnoted (\cref{ch:matrixfree}); fusion is presented as
separate scale-and-add operators conceptually and fused only in the note; tiling and
layout never appear in the algebra at all. A reader who follows only the unfolded
forms computes the right answer, just slower---which is the definition of transparent
(\cref{def:transparent}).

The load-bearing tricks, by contrast, each earned an \emph{explicit note} at the
point of use, stating the property bought:
\begin{itemize}[leftmargin=1.4em, nosep]
\item CG's inner products carry the reduction-order note (\cref{ch:cg},
  \cref{sec:numerics}): the trajectory is tree-dependent; reproducibility must be
  bought.
\item The Chebyshev smoother carries the recurrence-versus-monomial note
  (\cref{ch:smoothers}): the three-term recurrence is the stable evaluation; the
  sequentiality is load-bearing.
\item The orthogonalizer carries the MGS-versus-CGS note (\cref{ch:orthog}): the
  column order is a sequential dependency, not an inefficiency.
\item The preconditioned and mixed-precision solvers carry the conditioning note
  (this appendix, \cref{sec:mixed}): the attainable accuracy is $\kappa$ times the
  apply precision.
\end{itemize}
That asymmetry---transparent tricks footnoted, load-bearing tricks claimed---is the
book's consistent discipline, and \cref{app:laws} relies on it: several of the
algebraic ``laws'' catalogued there (commutativity of a reduction, associativity of
an accumulation) hold \emph{only up to rounding}, and the reason a law is qualified
``up to reduction order'' rather than asserted outright is exactly the load-bearing
status of \cref{sec:reduction}. The non-law is the specification.

\section{Practical guidance: when to care}
\label{sec:guidance}

The appendix closes with the engineering question a reader actually faces: of all
this, when does it matter? The honest answer is that for most of a simulation it does
not, and knowing \emph{where} it does is the skill.

\begin{itemize}[leftmargin=1.4em]
\item \textbf{Iterative-solver convergence.} Reduction order, fast-math, and mixed
  precision all perturb the Krylov trajectory. The \emph{converged} answer is robust
  (it is correct to tolerance regardless), but the \emph{iteration count} and the
  \emph{last bits} are not. Care when the iteration count is itself a measured
  quantity, when a tolerance is set near machine precision (where the trajectory
  noise approaches the signal), or when a solve is near the edge of converging at
  all---there the choice of reduction order can decide convergence versus stagnation.

\item \textbf{Eigenvalue accuracy.} Orthogonalization quality is load-bearing for
  eigensolvers in a way it is not for linear solves. A linear solve tolerates a
  somewhat degraded Krylov basis (it re-converges); an eigensolver reads the basis
  directly, and a lost-orthogonality basis manufactures \emph{spurious eigenvalues}
  (\cref{ch:orthog}). Care: use CGS2 or re-orthogonalized MGS whenever eigenvalues
  are the product, never bare CGS.

\item \textbf{Reproducibility requirements.} If a downstream consumer asserts on
  exact bits---a regression test, a certification, a bit-exact replay---reduction
  order must be pinned and fast-math disabled, at a throughput cost
  (\cref{sec:reduction}). Care: decide reproducibility as a requirement \emph{up
  front}, because retrofitting determinism onto a solver tuned for the fastest tree
  is expensive.

\item \textbf{When \emph{not} to care.} For a single well-conditioned solve whose
  output is a field or a scalar reported to engineering precision, with no
  bit-exact reproducibility requirement, every transparent trick may be taken freely
  and the load-bearing ones tuned for speed: single-precision applies, the fastest
  reduction tree, aggressive fusion. The vast majority of a simulation lives here.
  The discipline is not paranoia everywhere; it is knowing the handful of places---near
  convergence, at eigenvalues, under reproducibility mandates---where the last bits
  stop being noise and start being the answer.
\end{itemize}

\begin{keyidea}
\textbf{The one rule to carry away.} A \emph{transparent} trick is invisible to the
mathematics---it computes the same result faster, so the book presents the unfolded
form and ignores the trick. A \emph{load-bearing} trick \emph{is} the
mathematics---it changes determinism, conditioning, or IEEE compliance, so the book
states it as an explicit algebraic claim naming the property it buys. Mis-classifying
a load-bearing trick as transparent (replacing MGS by CGS, the recurrence by
monomials, double accumulation by single, a pinned reduction by the fastest tree)
silently changes the algorithm. When in doubt, the choice is flagged to the human,
because the cost of guessing wrong is an answer that is plausibly, invisibly, wrong.
\end{keyidea}

\summarysection

A real machine rounds, and the order in which it rounds is observable in the last
bits. This appendix separated the rounding choices that the mathematics may
\emph{ignore} from those that \emph{are} the mathematics. The primer
(\cref{sec:primer}) established the working facts: the unit roundoff $u$
(\cref{def:eps}), the non-associativity of floating-point addition
(\cref{eq:nonassoc}), catastrophic cancellation (\cref{sec:cancel}), and the
conditioning rule, forward error $\lesssim \kappa\cdot O(u)$ (\cref{eq:errorrule}),
which separates a property of the \emph{problem} from a property of the
\emph{algorithm}. \emph{Transparent} tricks (\cref{sec:transparent})---fusion,
tiling, packing, layout, recomputation, sum-factorization---compute the same result
faster and are footnoted, not stated. \emph{Load-bearing} tricks
(\cref{sec:loadbearing}) change a depended-on property and are stated as claims:
non-associative reduction order makes the Krylov trajectory device-dependent and
bit-reproducibility a paid-for property (\cref{sec:reduction}); fast-math grants the
compiler the same reordering freedom and breaks reproducibility and NaN-guards
(\cref{sec:fastmath}); mixed precision does the bulk work in single and the
accumulation in double, sound only while $\kappa\,u_{32}$ stays below tolerance
(\cref{sec:mixed}); the Chebyshev three-term recurrence is the stable evaluation and
the monomial form is catastrophic cancellation (\cref{sec:cheby}); and the MGS
column order is a load-bearing sequential dependency, not an inefficiency
(\cref{sec:orthog}). The book has applied this discipline throughout
(\cref{sec:throughout})---unfolded forms for the transparent, explicit notes for the
load-bearing---and \cref{app:laws} depends on it, since several algebraic laws hold
only ``up to reduction order,'' a qualification that is itself a specification. The
practical lesson (\cref{sec:guidance}): care near convergence, at eigenvalues, and
under reproducibility mandates; elsewhere, take the speed.

\furtherreading

The standard reference for the analysis of every phenomenon in this appendix---the
floating-point model \eqref{eq:fpmodel}, error bounds for summation and inner
products, the stability of Gram--Schmidt and of polynomial evaluation---is Higham's
treatment as developed in the numerical-linear-algebra texts; the conditioning and
stability framing of \cref{sec:conditioning}, including the forward-error rule
\eqref{eq:errorrule} and the backward-stability viewpoint, follows
Trefethen and Bau~\citep{trefethenbau1997}, whose first lectures are the cleanest
short account of conditioning, stability, and floating point in print. The
summation-tree error bounds, the quantitative loss-of-orthogonality results for CGS
versus MGS ($\kappa^2 u$ versus $\kappa u$), the ``twice is enough''
re-orthogonalization theorem, and the matrix condition number are developed in full
in Golub and Van Loan~\citep{golubvanloan2013}, the definitive matrix-computations
reference and the source for the orthogonalization trade-off of \cref{ch:orthog}.
The load-bearing stability of the Chebyshev three-term recurrence over its monomial
expansion, the coefficient formulas, and the analysis Palace follows directly are due
to Phillips and Fischer~\citep{phillipsfischer2022}; their \S2 is the precise
statement that the sequentiality of the recurrence is part of the algorithm. For the
broader principle---separating performance tricks from base algebra, of which this
appendix is the floating-point instance---see \cref{ch:bigidea}; for the algebraic
laws that this appendix qualifies ``up to rounding,'' see \cref{app:laws}.

\chapter{The Matrix-Free Kernel in Detail}
\label{app:kernel}

\chapterorient{\Cref{ch:matrixfree} factored the finite-element operator into five
element-local stages, $\mat{A}\vv = G^{\mathsf T}B_{\diffop}^{\mathsf T}\,D\,
B_{\diffop}\,G\,\vv$, and promised this appendix as the deep version. Here we
deliver it. We take each of the five stages apart down to its exact tensor shape,
write out the index maps that $G$ and $G^{\mathsf T}$ actually perform, show how
the differential operator $\diffop$ selects the basis mode and shapes the
component axis $C$, and build the geometry factor that the diagonal $D$ consumes
from the Jacobian metric up. We separate, in full, the build-time stratum (the
geometry computed once per mesh) from the run-time stratum (the action recomputed
on every apply); we count the operations sum-factorization saves and present it as
a contraction reordering; and we draw the libCEED library boundary precisely,
distinguishing the kernel \emph{API} we call from the kernel \emph{implementation}
---the contraction chain of this appendix---that realizes it. By the end the whole
operator stands exposed as a composition of tensor contractions over named axes:
the GPU-native form, with no sparse matrix anywhere, that the book's
implementation view (\cref{ch:synthlib}) consumes. This is where the mathematics
meets the silicon.}

\section{Scope: the chain this appendix opens}
\label{sec:appk-scope}

\Cref{ch:matrixfree} established the central factorization and verified it against
an assembled matrix on a two-element 1-D mesh. We restate it once, because every
section below is a magnification of one of its factors:
\begin{equation}
  \mat{A}\,\vv
   \;=\;
   \underbrace{G^{\mathsf T}}_{\text{scatter}}\;
   \underbrace{B_{\diffop}^{\mathsf T}}_{\text{integrate}}\;
   \underbrace{D}_{\text{weight}}\;
   \underbrace{B_{\diffop}}_{\text{evaluate}}\;
   \underbrace{G}_{\text{gather}}\;
   \vv .
  \label{eq:appk-chain}
\end{equation}
The five factors carry a tensor through a sequence of named shapes,
\begin{equation}
\begin{aligned}
  \texttt{[(N: \dots)]}
  &\xrightarrow{\,G\,}
  \texttt{[(E, L)]}
  \xrightarrow{\,B_{\diffop}\,}
  \texttt{[(E, P, C)]}
  \xrightarrow{\,D\,}
  \texttt{[(E, P, C')]} \\
  &\xrightarrow{\,B_{\diffop}^{\mathsf T}\,}
  \texttt{[(E, L)]}
  \xrightarrow{\,G^{\mathsf T}\,}
  \texttt{[(N: \dots)]},
\end{aligned}
  \label{eq:appk-shapes}
\end{equation}
where the five axes are the element-local family fixed in \cref{ch:refelement} and
catalogued in \cref{app:shapes}: $N$ the global true-degree-of-freedom count, $E$
the element count, $L$ the local dofs per element, $P$ the quadrature points per
element, $C$ the value/derivative components carried at a quadrature point, and
$G$ the geometry-metric components stored per point. \Cref{tab:appk-axes} is the
working legend for this appendix; \cref{app:shapes} carries the full algebra.

\begin{table}[t]
\centering
\caption[The five element-local axes, with extents]{The element-local tensor axes
used throughout this appendix, with their typical extents on a degree-$p$,
$d$-dimensional tensor-product element. The stratum column---built once per
$(\text{mesh},\text{order})$ versus produced afresh on every apply---is the subject
of \cref{sec:appk-strata}. The global axis $N$ is \emph{not} an element-local axis;
it is the flat solver vector, and $G$ is the boundary between it and the
$(E,L)$ vocabulary.}
\label{tab:appk-axes}
\small
\setlength{\tabcolsep}{4.5pt}
\renewcommand{\arraystretch}{1.25}
\begin{tabular}{@{}cllcl@{}}
\toprule
axis & name & meaning & extent & stratum \\
\midrule
$N$ & global dofs & true dofs of the global vector & --- & solver vector \\
$E$ & element count & mesh elements assembled over & --- & build-time \\
$L$ & local dofs & dofs local to one element & $(p{+}1)^d$ & build-time \\
$P$ & quad points & quadrature points per element & $(p{+}1)^d$ & build-time \\
$C$ & components & value/deriv.\ components per point & $1$ or $d$ & run-time \\
$G$ & geom comps & stored metric comps per point & $2+d^2$ & build-time \\
\bottomrule
\end{tabular}
\end{table}

\Cref{fig:appk-pipeline} is the deep version of \cref{ch:matrixfree}'s hero
diagram: the same five stages, but with the \emph{exact} shape on every wire,
including the build-time inputs each stage draws from below. The sections that
follow walk it stage by stage.

\begin{figure}[p]
\centering
\resizebox{\linewidth}{!}{%
\begin{tikzpicture}[node distance=7mm and 8mm]
  % input
  \node[data, text width=20mm] (in) {global vector\\\code{[(N: \dots)]}};
  \node[stage, right=of in, text width=19mm] (G)  {$G$\\\footnotesize gather\\element restrict};
  \node[stage, right=of G,  text width=19mm] (B)  {$B_{\diffop}$\\\footnotesize evaluate\\basis apply};
  \node[extern, right=of B, text width=19mm] (D)  {$D$\\\footnotesize weight\\\code{\#extern}};
  \node[stage, right=of D,  text width=19mm] (Bt) {$B_{\diffop}^{\mathsf T}$\\\footnotesize integrate};
  \node[stage, right=of Bt, text width=19mm] (Gt) {$G^{\mathsf T}$\\\footnotesize scatter-add};
  \node[data, right=of Gt, text width=20mm] (out) {global vector\\\code{[(N: \dots)]}};

  \draw[flow] (in) -- (G);  \draw[flow] (G) -- (B); \draw[flow] (B) -- (D);
  \draw[flow] (D) -- (Bt);  \draw[flow] (Bt) -- (Gt); \draw[flow] (Gt) -- (out);

  % exact shapes on every wire
  \node[font=\scriptsize\ttfamily, simGray, below=2.4mm of G.south]  {[(E, L)]};
  \node[font=\scriptsize\ttfamily, simGray, below=2.4mm of B.south]  {[(E, P, C)]};
  \node[font=\scriptsize\ttfamily, simGray, below=2.4mm of D.south]  {[(E, P, C')]};
  \node[font=\scriptsize\ttfamily, simGray, below=2.4mm of Bt.south] {[(E, L)]};
  \node[font=\scriptsize\ttfamily, simBlue, below=3mm of in.south, align=center] {action\\input};
  \node[font=\scriptsize\ttfamily, simBlue, below=3mm of out.south, align=center] {action\\output};

  % build-time inputs fed from below
  \node[data, fill=simBgGold, draw=simGold, below=15mm of G, text width=22mm]
    (restr) {\footnotesize\itshape restriction\\index map (build)};
  \node[data, fill=simBgGold, draw=simGold, below=15mm of B, text width=22mm]
    (basis) {\footnotesize\itshape tabulated basis\\$\hat B,\,\hat G$ (build)};
  \node[data, fill=simBgGold, draw=simGold, below=15mm of D, text width=24mm]
    (geom) {\footnotesize\itshape geom\_data \code{[(E,P,G)]}\\(build)};
  \draw[flow, simGold] (restr) -- (G);
  \draw[flow, simGold] (basis) -- (B);
  \draw[flow, simGold] (geom)  -- (D);
  % Bt and Gt reuse the same build-time tables (dotted back-refs)
  \draw[refedge, simGold] (basis.north) .. controls +(0,0.9) and +(0,-1.0) .. (Bt.south);
  \draw[refedge, simGold] (restr.north) .. controls +(0,1.1) and +(0,-1.2) .. (Gt.south);

  % brackets: element-local middle
  \draw[decorate, decoration={brace, amplitude=4pt}, simRust]
    ([yshift=7mm]G.north west) -- ([yshift=7mm]Gt.north east)
    node[midway, above=4pt, font=\scriptsize, simRust]
    {element-local: block-diagonal in $E$, fully parallel};
\end{tikzpicture}%
}
\caption[The five-stage pipeline with exact tensor shapes]{The matrix-free
operator of \cref{eq:appk-chain}, drawn with the \emph{exact} element-local shape
on every wire (the deep version of \cref{ch:matrixfree}'s pipeline figure). A flat
global vector \code{[(N: \dots)]} enters; the gather $G$ produces per-element local
blocks \code{[(E, L)]}; basis evaluation $B_{\diffop}$ lifts them to
\code{[(E, P, C)]} values at the quadrature points; the center stage $D$, the
\code{\#extern} libCEED boundary (violet, \cref{sec:appk-extern}), weights each
point against the precomputed \code{geom\_data}, possibly changing the component
count to $C'$; $B_{\diffop}^{\mathsf T}$ integrates back to \code{[(E, L)]}; and
$G^{\mathsf T}$ scatters into a fresh global vector. The three build-time tables
(gold, fed from below)---the restriction index map, the tabulated basis $\hat B$
and reference-gradient $\hat G$, and the geometry factor---are computed once and
reused by both the forward stages and their transposes (dotted back-references).
The element-local middle (rust brace) is block-diagonal in $E$.}
\label{fig:appk-pipeline}
\end{figure}

\section{The element-local tensor, axis by axis}
\label{sec:appk-anatomy}

Before opening the stages we fix the three tensors they move between, because
their axes are the entire vocabulary. The flat solver vector is rank-1,
\code{[(N: \dots)]}; the moment the gather runs we acquire an explicit element
axis $E$ and the per-element structure under it. \Cref{fig:appk-anatomy} is the
anatomy chart.

\begin{figure}[t]
\centering
\begin{tikzpicture}[scale=1.0]
  % --- [E, L] ---
  \begin{scope}[shift={(-4.6,0)}]
    \foreach \k in {0,1,2}{
      \begin{scope}[shift={(0.12*\k,0.09*\k)}]
        \fill[simBgBlue, draw=simBlue] (-0.7,-0.5) rectangle (0.7,0.5);
        \foreach \y in {-0.25,0,0.25}{\draw[simBlue!40] (-0.7,\y)--(0.7,\y);}
      \end{scope}}
    \node[font=\scriptsize, simInk] at (-1.05,0.0) {$L$};
    \node[font=\scriptsize, simRust] at (0.55,0.78) {$E$};
    \node[font=\scriptsize\ttfamily, simInk] at (0.1,-1.0) {[(E, L)]};
    \node[font=\scriptsize\itshape, simGray, text width=26mm, align=center]
      at (0.1,-1.52) {per-element local dofs\\(restricted field)};
  \end{scope}
  \node[font=\large] at (-2.7,0) {$\xrightarrow{\,B_{\diffop}\,}$};
  % --- [E, P, C] ---
  \begin{scope}[shift={(-0.6,0)}]
    \foreach \k in {0,1,2}{
      \begin{scope}[shift={(0.12*\k,0.09*\k)}]
        \fill[simBgGreen, draw=simGreen] (-0.78,-0.5) rectangle (0.78,0.5);
        \foreach \y in {-0.17,0.17}{\draw[simGreen!40] (-0.78,\y)--(0.78,\y);}
        \foreach \x in {-0.26,0.26}{\draw[simGreen!40] (\x,-0.5)--(\x,0.5);}
      \end{scope}}
    \node[font=\scriptsize, simInk] at (-1.13,0.0) {$P$};
    \node[font=\scriptsize, simInk] at (0.0,-0.83) {$C$};
    \node[font=\scriptsize, simRust] at (0.62,0.78) {$E$};
    \node[font=\scriptsize\ttfamily, simInk] at (0.05,-1.0) {[(E, P, C)]};
    \node[font=\scriptsize\itshape, simGray, text width=28mm, align=center]
      at (0.05,-1.52) {$C$ values at\\$P$ quad points};
  \end{scope}
  \node[font=\large] at (1.55,0) {$\odot$};
  % --- [E, P, G]  geom_data ---
  \begin{scope}[shift={(3.7,0)}]
    \foreach \k in {0,1,2}{
      \begin{scope}[shift={(0.12*\k,0.09*\k)}]
        \fill[simBgGold, draw=simGold] (-0.78,-0.5) rectangle (0.78,0.5);
        \foreach \y in {-0.17,0.17}{\draw[simGold!50] (-0.78,\y)--(0.78,\y);}
        \foreach \x in {-0.26,0.26}{\draw[simGold!50] (\x,-0.5)--(\x,0.5);}
      \end{scope}}
    \node[font=\scriptsize, simInk] at (-1.13,0.0) {$P$};
    \node[font=\scriptsize, simInk] at (0.0,-0.83) {$G$};
    \node[font=\scriptsize, simRust] at (0.62,0.78) {$E$};
    \node[font=\scriptsize\ttfamily, simInk] at (0.05,-1.0) {[(E, P, G)]};
    \node[font=\scriptsize\itshape, simGray, text width=28mm, align=center]
      at (0.05,-1.52) {geom\_data\\(build-time)};
  \end{scope}
\end{tikzpicture}
\caption[Anatomy of the element-local tensor family]{The three element-local
tensors and the axes they share. Gather produces \code{[(E, L)]} (blue): an
element axis $E$ over per-element local-dof blocks of length $L$. Basis evaluation
$B_{\diffop}$ maps it to \code{[(E, P, C)]} (green): $C$ components at each of $P$
quadrature points, per element. The weight $D$ multiplies it pointwise ($\odot$)
against the precomputed geometry carrier \code{[(E, P, G)]} (gold). The element
axis $E$ (rust) threads through all three: every stage between gather and scatter
is block-diagonal in $E$. The $G$ axis is purely build-time storage; the $C$ axis
is the only run-time-shaped axis, and even it is fixed once the differential
operator $\diffop$ is chosen.}
\label{fig:appk-anatomy}
\end{figure}

\begin{notationbox}
The libCEED library, which owns the device kernel (\cref{sec:appk-extern}), names
two of these axes \emph{oppositely} to this book. Its basis constructor's
parameter \code{P} counts \emph{nodes per element}---our $L$---and its \code{Q}
counts \emph{quadrature points}---our $P$. This appendix uses $L$ for local dofs
and $P$ for quadrature points throughout; if you read the library source expect the
letters to swap. The binding is exact; only the labels differ.
\end{notationbox}

The $G$ axis deserves an early word because it is the one whose extent we have not
yet pinned. The geometry carrier stores $G = 2 + d^2$ components per quadrature
point: in three dimensions, $2 + 9 = 11$. The two scalar slots hold the
quadrature-weighted volume factor and an attribute (the material-region tag); the
$d^2$ slots hold the full metric matrix $J^{-1}J^{-\mathsf T}|\det J|$ that a
grad--grad term needs (\cref{sec:appk-geom}). A pure mass term uses only the volume
scalar, but the storage size is fixed by the largest term, so $G=2+d^2$ is the
allocation. This is a positive fact read off the library's build pass, not an
estimate: the assembler verifies \code{geom\_data\_size == 2 + space\_dim*dim}
before it runs.

\section{Stage 1: \texorpdfstring{$G$}{G}, element restriction (the gather)}
\label{sec:appk-gather}

The gather $G$ is the boundary between two vocabularies. On its left sits the
flat global vector \code{[(N: \dots)]}, a rank-1 list the solver owns; on its right
sit the per-element blocks \code{[(E, L)]} the rest of the pipeline runs on. $G$
is the map that crosses it.

\subsection{The index map}

$G$ performs no arithmetic. It is a pure \emph{indexed copy}, defined by a fixed
table---the \emph{element-to-dof map}---that records, for each element $e$ and each
of its $L$ local dof slots $\ell$, which global dof index $n = \iota(e,\ell)$ it
reads from. The action is simply
\begin{equation}
  (G\vv)[e,\ell] \;=\; \vv\bigl[\iota(e,\ell)\bigr],
  \qquad e\in\{0,\dots,E-1\},\ \ell\in\{0,\dots,L-1\}.
  \label{eq:appk-gather}
\end{equation}
For each element, $G$ walks that element's $L$ slots and copies in the global
values they point at. A global dof that several elements share---a vertex,
edge, or face dof on an inter-element boundary---is read by \emph{every} element
that touches it, so its value is \emph{copied out} into several local blocks. That
duplication is deliberate; it is what lets the middle stages run element by element
in ignorance of the sharing, and it is undone, by summation, at the far end of the
pipeline. \Cref{fig:appk-gatherscatter} draws the copy-out and the matching
sum-back.

\begin{figure}[t]
\centering
\begin{tikzpicture}[scale=1.0, every node/.style={font=\scriptsize}]
  % ---- gather (top): one shared global dof fans out ----
  \begin{scope}
    \node[font=\footnotesize\bfseries, simBlue] at (-2.4,1.1) {$G$ (gather)};
    % global slots
    \foreach \i/\lab in {0/$v_0$,1/$v_1$,2/$v_2$,3/$v_3$}{
      \node[draw=simGray, fill=simBgGray, minimum size=5.5mm] (g\i) at (\i*0.9,1.1) {\lab};}
    \node[simGray] at (1.35,1.7) {global \code{[(N)]}};
    % element blocks e0=(v0,v1) e1=(v1,v2) e2=(v2,v3)
    \foreach \k/\a/\b/\la/\lb in {0/0/0.9/$v_0$/$v_1$, 1/2.2/3.1/$v_1$/$v_2$, 2/4.4/5.3/$v_2$/$v_3$}{
      \node[draw=simBlue, fill=simBgBlue, minimum size=5.0mm] (e\k a) at (\a,-0.2) {\la};
      \node[draw=simBlue, fill=simBgBlue, minimum size=5.0mm] (e\k b) at (\b,-0.2) {\lb};
      \node[simRust] at ({(\a+\b)/2},-0.75) {$e_\k$};}
    % copy arrows for the shared dof v1 and v2 (fan-out)
    \draw[flow, simBlue] (g0.south) -- (e0a.north);
    \draw[flow, simBlue] (g1.south) -- (e0b.north);
    \draw[flow, simRust, thick] (g1.south) -- (e1a.north);   % v1 fans to e1 too
    \draw[flow, simBlue] (g2.south) -- (e1b.north);
    \draw[flow, simRust, thick] (g2.south) -- (e2a.north);   % v2 fans to e2 too
    \draw[flow, simBlue] (g3.south) -- (e2b.north);
    \node[simRust, font=\scriptsize\itshape] at (6.6,0.4) {shared dof};
    \node[simRust, font=\scriptsize\itshape] at (6.6,0.05) {copied out};
  \end{scope}
  % ---- scatter (bottom): contributions sum back ----
  \begin{scope}[shift={(0,-3.4)}]
    \node[font=\footnotesize\bfseries, simGreen] at (-2.4,-0.2) {$G^{\mathsf T}$ (scatter-add)};
    \foreach \k/\a/\b/\la/\lb in {0/0/0.9/$c^0_0$/$c^0_1$, 1/2.2/3.1/$c^1_0$/$c^1_1$, 2/4.4/5.3/$c^2_0$/$c^2_1$}{
      \node[draw=simBlue, fill=simBgBlue, minimum size=5.0mm] (s\k a) at (\a,0.5) {\la};
      \node[draw=simBlue, fill=simBgBlue, minimum size=5.0mm] (s\k b) at (\b,0.5) {\lb};}
    \foreach \i/\lab in {0/{$c^0_0$},1/{$c^0_1{+}c^1_0$},2/{$c^1_1{+}c^2_0$},3/{$c^2_1$}}{
      \node[draw=simGreen, fill=simBgGreen, minimum width=9mm, minimum height=5.5mm] (o\i) at (\i*1.55+0.0,-0.9) {\lab};}
    \draw[flow, simGreen] (s0a.south) -- (o0.north);
    \draw[flow, simGreen] (s0b.south) -- (o1.north);
    \draw[flow, simGreen] (s1a.south) -- (o1.north);
    \draw[flow, simGreen] (s1b.south) -- (o2.north);
    \draw[flow, simGreen] (s2a.south) -- (o2.north);
    \draw[flow, simGreen] (s2b.south) -- (o3.north);
    \node[simGreen, font=\scriptsize\itshape] at (6.6,-0.5) {summed};
    \node[simGreen, font=\scriptsize\itshape] at (6.6,-0.85) {at shared dofs};
  \end{scope}
\end{tikzpicture}
\caption[Gather copies out; scatter sums back]{The $G$/$G^{\mathsf T}$ transpose
pair on a 1-D mesh of three elements sharing two interior nodes. \emph{Top:} the
gather $G$ copies each global value into every element that references it; the
shared dofs $v_1,v_2$ (rust) fan out to two elements each. No arithmetic---pure
copying. \emph{Bottom:} the scatter-add $G^{\mathsf T}$ runs the table in reverse,
\emph{summing} the per-element contributions $c^e_\ell$ back into each global slot;
at a shared node the two elements' contributions add (e.g.\ $c^0_1+c^1_0$). The
copy-out/sum-back pair is the entire content of finite-element assembly, here
applied to a vector. Because $G^{\mathsf T}$ is the literal transpose of $G$,
$\inner{G^{\mathsf T}\vw}{\vv}=\inner{\vw}{G\vv}$ holds exactly---which is why the
assembled value is a \emph{sum} and not an average.}
\label{fig:appk-gatherscatter}
\end{figure}

\subsection{$G^{\mathsf T}$ is the scatter-add; $G^{\mathsf T}G$ is the multiplicity}

The transpose $G^{\mathsf T}$ runs the same index table backwards. Where $G$ read
$n=\iota(e,\ell)$ and copied $\vv[n]$ into slot $(e,\ell)$, the transpose takes a
per-element tensor $\vc\in\code{[(E, L)]}$ and \emph{adds} each local contribution
into the global slot it came from:
\begin{equation}
  (G^{\mathsf T}\vc)[n] \;=\; \sum_{(e,\ell)\,:\,\iota(e,\ell)=n} \vc[e,\ell].
  \label{eq:appk-scatter}
\end{equation}
This is exactly a \emph{scatter-add}: indexed writes that accumulate where indices
collide. It is the only stage in the pipeline that reduces---every other stage is a
map or a per-block contraction---and on a parallel machine it is the one stage that
needs an atomic or a coloring to make the colliding writes safe. The summation is
finite-element \emph{assembly} (\cref{ch:assembly}), here applied to a vector
rather than stamped into a matrix.

The composition $G^{\mathsf T}G$ is worth isolating, because it has a clean
meaning. Gather then immediately scatter-add: each global dof is copied out to
every element that touches it, then those copies are summed straight back with no
intervening computation. The result is the global dof scaled by how many elements
share it:
\begin{equation}
  (G^{\mathsf T}G\,\vv)[n] \;=\; m_n\,\vv[n],
  \qquad m_n = \#\{(e,\ell): \iota(e,\ell)=n\},
  \label{eq:appk-mult}
\end{equation}
the \emph{dof-multiplicity diagonal} $\mat{M}_{\mathrm{mult}}=\operatorname{diag}(m_n)$.
An interior vertex of a hex mesh has $m_n=8$; a face dof has $m_n=2$; a dof on the
domain boundary touched by one element has $m_n=1$. The multiplicity diagonal is
cheap, build-time, and useful: it is exactly the operator that converts a
\emph{summed} (assembled) quantity into an \emph{averaged} one, and it appears
whenever the pipeline must reconcile the redundant element-local representation
with the unique global one---for instance in building a diagonal for a smoother
(\cref{ch:smoothers}) or projecting between the two representations.

\begin{keyidea}
The gather $G$ and scatter $G^{\mathsf T}$ are a transpose pair built from one
index table $\iota(e,\ell)$: $G$ copies global dofs out to elements (duplicating
shared dofs), $G^{\mathsf T}$ sums element contributions back (assembling shared
dofs). All inter-element coupling in the whole operator lives in these two ends;
everything between them is block-diagonal in $E$. Their composition
$G^{\mathsf T}G=\operatorname{diag}(m_n)$ is the dof-multiplicity diagonal that
reconciles the duplicated element-local view with the unique global one.
\end{keyidea}

\section{Stage 2: \texorpdfstring{$B_{\diffop}$}{B}, basis evaluation}
\label{sec:appk-basis}

With the field restricted to elements, the weak form needs it---and its
derivative---sampled at the quadrature points. That is what $B_{\diffop}$ delivers:
it contracts each element's $L$ local coefficients against the tabulated reference
basis to produce $C$ components at each of $P$ quadrature points.

\subsection{The tabulated basis tensor}

The reference shape functions of \cref{ch:refelement} are sampled \emph{once} at
the quadrature nodes and stored. For the value (interpolation) mode this is a dense
matrix
\[
  \hat B \in \Real^{P\times L},
  \qquad \hat B_{q\ell} = \hat\phi_\ell(\hat\vx_q),
\]
the $\ell$-th reference shape function evaluated at the $q$-th quadrature point. For
a derivative mode the table is the reference \emph{gradient}
$\hat G\in\Real^{P\times L\times d}$ with $\hat G_{q\ell k}=\partial_k\hat\phi_\ell
(\hat\vx_q)$. Both are pure build-time data: tabulated once, shared by every element
of the mesh, independent of the input vector. The value-mode evaluation, per
element $e$, is the contraction over the local-dof axis $L$:
\begin{equation}
  (B\,\vc)[e,q,c] \;=\; \sum_{\ell=0}^{L-1} \hat B_{q\ell}\;\vc[e,\ell],
  \qquad C=1\ \text{for a scalar value},
  \label{eq:appk-basis}
\end{equation}
producing \code{[(E, P, C)]} from \code{[(E, L)]}. The same small $\hat B$ is
applied independently inside each of the $E$ elements---a batched matrix-vector
product, block-diagonal in $E$.

\subsection{The mode is selected by \texorpdfstring{$\diffop$}{D}}

The differential operator $\diffop$ of the weak-form term selects which tabulated
table $B_{\diffop}$ uses, and with it the meaning and extent of the component axis
$C$. \Cref{tab:appk-modes} lists the four de Rham modes (\cref{ch:refelement}).

\begin{table}[t]
\centering
\caption[Basis-evaluation modes and the component axis]{The four basis-evaluation
modes selected by the term's differential operator $\diffop$, with the
component-axis extent $C$ each produces. The interpolation mode evaluates basis
\emph{values} ($C=1$ scalar); the gradient, curl, and divergence modes evaluate
derivatives, with $C$ set by the de Rham slot. The same five-stage skeleton serves
all of them; only the tabulated table and the geometry-factor shape change.}
\label{tab:appk-modes}
\small
\setlength{\tabcolsep}{7pt}
\renewcommand{\arraystretch}{1.25}
\begin{tabular}{@{}llll@{}}
\toprule
$\diffop$ & mode & table & components $C$ \\
\midrule
$\mathrm{Id}$ & interp (values) & $\hat B$, $P\times L$ & $1$ (scalar) \\
$\Grad$ & grad & $\hat G$, $P\times L\times d$ & $d$ \\
$\Curl$ & curl & $\hat G$, curl-combined & $d$ (3-D), $1$ (2-D) \\
$\Div$ & div & $\hat G$, trace-combined & $1$ \\
\bottomrule
\end{tabular}
\end{table}

For a gradient term the contraction carries the extra $k$ axis,
\[
  (B_{\Grad}\vc)[e,q,c] \;=\; \sum_{\ell} \hat G_{q\ell c}\;\vc[e,\ell],
  \qquad c\in\{0,\dots,d-1\},
\]
so $C=d$: the reference gradient of the field at each quadrature point. (The
\emph{physical} gradient---the one the weak form actually integrates---needs the
inverse-transpose Jacobian $J^{-\mathsf T}$ applied to this reference gradient; that
multiplication is folded into the geometry factor and applied in stage 3, not here,
which keeps $B$ a pure reference-table contraction. We return to this in
\cref{sec:appk-geom}.) The curl and divergence modes combine the reference
derivatives the way the differential operator prescribes; the structural shape is
identical, only the table differs.

\subsection{The transpose $B_{\diffop}^{\mathsf T}$ integrates}

Stage 4 is the exact transpose. Where $B$ contracted the local-dof axis $L$ against
the table to produce quadrature-point values, $B_{\diffop}^{\mathsf T}$ contracts
the quadrature-point axis $P$ against the \emph{same} table to produce local-dof
contributions:
\begin{equation}
  (B^{\mathsf T}\vu)[e,\ell] \;=\; \sum_{q=0}^{P-1} \hat B_{q\ell}\;\vu[e,q,c],
  \label{eq:appk-basist}
\end{equation}
mapping \code{[(E, P, C)]} back to \code{[(E, L)]}. This is the discrete integral
$\int(\cdot)\,\phi_\ell$ against each test basis function: the quadrature sum
$\sum_q$ of \cref{ch:refelement} \emph{is} this transpose contraction, with the
quadrature weight already folded into the weighted values $\vu$ that reach it from
stage 3. The same tabulated $\hat B$ serves both directions---one of the reasons
the build-time stratum is so small.

\section{Stage 4 in passing, Stage 3 in full: the diagonal weight \texorpdfstring{$D$}{D}}
\label{sec:appk-diagonal}

We have already described stage 4 as the transpose of stage 2. The center stage
$D$---logically stage 3, the weight---is the simplest object in the pipeline and
the one that carries all the physics, so we give it its own treatment.

After basis evaluation the field is a tensor of values at quadrature points,
\code{[(E, P, C)]}. The weight stage multiplies each point's value, independently,
by a single precomputed factor read from \code{geom\_data}:
\begin{equation}
  (D\,\vu)[e,p,c'] \;=\; \sum_{c} \texttt{geom\_data}[e,p]\,{}_{c'c}\;\vu[e,p,c],
  \label{eq:appk-diagonal}
\end{equation}
where for a scalar (mass) term the factor is a single number and \cref{eq:appk-diagonal}
is a plain scaling; for a vector (grad--grad) term the factor is a small $d\times d$
block---the metric---and the sum over $c$ is a tiny per-point matrix-vector product
that may change the component count to $C'$. \Cref{fig:appk-diagonal} draws the
structure: the output at index $(e,p)$ depends only on the input at $(e,p)$ and on
\code{geom\_data}$[e,p]$, never on any other point.

\begin{figure}[t]
\centering
\begin{tikzpicture}[scale=1.0]
  \foreach \i in {0,...,4}{
    \foreach \j in {0,...,2}{
      \node[draw=simGreen, fill=simBgGreen, rounded corners=1pt,
            minimum size=7mm, font=\scriptsize] (c\i\j) at (\i*1.05, \j*0.9) {$u_{\i\j}$};
    }
  }
  \foreach \i in {0,...,4}{
    \foreach \j in {0,...,2}{
      \node[font=\scriptsize, simGold] at (\i*1.05+0.35, \j*0.9+0.35) {$\odot g$};
    }
  }
  \node[font=\scriptsize\itshape, simRust] at (-0.95, 0.9) {$(e,p)$};
  \node[font=\scriptsize, simInk] at (2.1,-0.85) {$P$ points $\times$ $E$ elements (flattened)};
  \draw[decorate, decoration={brace, amplitude=5pt}, simBlue]
    (-0.5, 2.2) -- (4.7, 2.2)
    node[midway, above=5pt, font=\footnotesize, simBlue]
    {every cell independent --- no coupling, any order, fully parallel};
\end{tikzpicture}
\caption[The weight stage as an embarrassingly parallel diagonal]{The weight stage
$D$ over the flattened $(e,p)$ index. Each cell is one quadrature point's evaluated
value $u_{ep}$, multiplied ($\odot$) by its own precomputed geometry block $g_{ep}$
from \code{geom\_data}. No cell depends on any other: $D$ is \emph{diagonal} in the
$(E,P)$ index---a per-point lift, mapped across the whole point set at once. There
is no reduction, no recurrence, no order-dependence; the output at $(e,p)$ is a
function of the input at $(e,p)$ alone. This is the embarrassingly parallel core of
the operator and the form a tensor accelerator consumes most directly.}
\label{fig:appk-diagonal}
\end{figure}

This is why we call $D$ the \emph{diagonal} of the operator. The basis stages
couple the $L$ local dofs to the $P$ quadrature points (a dense small contraction);
the gather and scatter couple elements through shared dofs. But $D$, sandwiched
between them, couples nothing: it is a diagonal scaling in the quadrature-point
basis, block-diagonal not merely in $E$ but in the full $(E,P)$ index. Everything
that distinguishes a mass term from a stiffness term, this material from that one,
a curved element from a straight one, is packed into the contents of
\code{geom\_data}---and that packing is the build-time pass we turn to now.

\section{Building the geometry factor: \texttt{geom\_factor\_build}}
\label{sec:appk-geom}

The tensor \code{geom\_data} of shape \code{[(E, P, G)]} is the one piece of the
build-time stratum that costs real arithmetic. It is produced by a setup pass,
\code{geom\_factor\_build}, whose signature is
\begin{lstlisting}[style=l4style]
-- the geometry-factor build pass: a pure setup-stratum function
geom_factor_build :: MeshNodes -> QuadWeights -> Tensor[(E, P, G)]
  -- per (e, p):  geom_data[e,p] = metric( J(mesh_nodes)[e,p] ) * w[e,p] * Q[e,p]
  --   E = elements;  P = quad points per element;  G = 2 + d*d metric components
\end{lstlisting}
At each quadrature point of each element it forms the Jacobian, combines it into
the metric the weak-form term needs, scales by the quadrature weight, and folds in
the material coefficient.

\subsection{From the Jacobian to the metric}

Recall from \cref{ch:refelement} that the geometric map $F_K$ carries the reference
element onto the physical element, and its derivative is the Jacobian $J=\partial
F_K/\partial\hat\vx\in\Real^{d\times d}$, whose columns are the physical edge
vectors. Two objects derived from $J$ are all the geometry a term needs:
\begin{equation}
\begin{aligned}
  \text{mass term ($\diffop=\mathrm{Id}$):}\quad & g = \abs{\det J}\,; \\
  \text{grad--grad ($\diffop=\Grad$):}\quad & \mat{G}_{\mathrm{metric}} = J^{-1}J^{-\mathsf T}\,\abs{\det J}.
\end{aligned}
  \label{eq:appk-metric}
\end{equation}
The mass factor is a scalar; the grad--grad metric is a symmetric $d\times d$
matrix. The metric arises by pulling both gradients in $\Grad\phi_i\cdot\Grad\phi_j$
back through $J^{-\mathsf T}$ and collecting the volume factor (\cref{ch:refelement}
derives it). Storing $J^{-1}J^{-\mathsf T}|\det J|$ is exactly the device that lets
the run-time $D$ apply the physical-gradient transformation as a single per-point
matrix multiply, with no Jacobian inversion at apply time.

\subsection{Folding in the weight and the coefficient}

The build pass does not stop at the bare metric. It pre-multiplies two more
factors, so that the run-time stage $D$ is a single multiply with nothing left to
combine:
\begin{equation}
  \texttt{geom\_data}[e,p] \;=\;
    \underbrace{\bigl(\text{metric from } J[e,p]\bigr)}_{\abs{\det J}\ \text{or}\ J^{-1}J^{-\mathsf T}\abs{\det J}}
    \;\cdot\;
    \underbrace{w_p}_{\text{quad weight}}
    \;\cdot\;
    \underbrace{Q[e,p]}_{\text{material coeff}}.
  \label{eq:appk-geomdata}
\end{equation}
The quadrature weight $w_p$ (\cref{ch:refelement}) rides inside \code{geom\_data}
rather than appearing as a separate stage---the weighted sum
$\sum_q w_q(\cdot)$ of the discrete integral is split across $D$ (the $w_q$ factor)
and $B^{\mathsf T}$ (the $\sum_q$ contraction). The material coefficient $Q$---the
$\varepsilon$, $\mu^{-1}$, or conductivity of the term, looked up by the element's
material-region attribute---is folded in the same way. The result is one number
(mass) or one small matrix (grad--grad) per quadrature point, $G=2+d^2$ stored
components, computed once and read by every apply.

\begin{remark}
On a \emph{straight-sided} (affine) element the Jacobian $J$ is constant over the
element, so the metric is constant in $p$ and \code{geom\_data}$[e,\cdot]$ repeats
the same block at every quadrature point---a degenerate case the build pass may
exploit by storing one block per element instead of $P$. On a \emph{curved}
(isoparametric high-order) element $J$ varies point to point and the full
\code{[(E, P, G)]} storage is needed. The pipeline is identical either way; only
the storage compresses.
\end{remark}

\section{Build-time versus run-time stratification}
\label{sec:appk-strata}

The single most consequential structural fact about the operator is that its data
divides into two strata by \emph{lifetime}. \Cref{ch:strata} treats this division
across the whole simulator; here we make it concrete for the kernel.

\begin{description}[leftmargin=1.4em, style=nextline]
\item[Build-time (setup) stratum.] Data fixed once per
$(\text{mesh},\text{FE order},\text{quadrature rule})$ and reused across every
application of the operator: the restriction index map $\iota(e,\ell)$, the
tabulated basis tables $\hat B$ and $\hat G$, and---the one with real
arithmetic---\code{geom\_data}. Their axis extents $E$, $L$, $P$, $G$ are all
build-time.
\item[Run-time (apply) stratum.] Data produced afresh on every action
$\vv\mapsto\mat{A}\vv$: the gathered \code{[(E, L)]} block, the evaluated and
weighted \code{[(E, P, C)]} values, the integrated contributions. These are
transient, born inside one apply and discarded when it returns. The only
run-time-\emph{shaped} axis is $C$, and even it is fixed once $\diffop$ is chosen;
what truly varies per apply is the field \emph{value} at each $(e,p,c)$ index.
\end{description}

\Cref{fig:appk-strata} draws the split. A setup pass runs the build \emph{once} and
produces the three tables; thereafter every operator application---and a Krylov
solve performs hundreds---reads them without rebuilding. The expensive geometry
work is amortized across the entire solve. Only when the mesh changes---adaptive
refinement (\cref{ch:amr}) being the usual cause---is the setup pass rerun, a
setup-stratum invalidation, not a per-apply cost.

\begin{figure}[t]
\centering
\begin{tikzpicture}[node distance=7mm and 12mm]
  % --- build-time row (gold) ---
  \node[data, fill=simBgGold, draw=simGold, text width=20mm] (mesh)
    {mesh nodes\\$+$ quad rule};
  \node[stage, fill=simBgGold, draw=simGold, right=of mesh, text width=27mm]
    (build) {\textbf{setup pass}\\\footnotesize $J,\;\abs{\det J},\;Q,\;w$\\(\code{geom\_factor\_build})};
  \node[data, fill=simBgGold, draw=simGold, right=of build, text width=23mm]
    (gd) {\code{geom\_data}\\\code{[(E, P, G)]}\\$+$ basis $+$ restr};
  \draw[flow, simGold] (mesh) -- (build);
  \draw[flow, simGold] (build) -- (gd);
  \node[font=\scriptsize\bfseries, simGold, left=1mm of mesh, rotate=90, anchor=south]
    {build-time};
  \node[font=\scriptsize\itshape, simGray, below=1mm of build]
    {runs \emph{once} per (mesh, order)};

  % --- run-time row (blue): many applies reading the tables ---
  \node[stage, below=18mm of build, text width=30mm] (apply1)
    {apply $\mat{A}\vv_1$\\\footnotesize $G^{\mathsf T}B^{\mathsf T}\,D\,B\,G$};
  \node[stage, right=5mm of apply1, text width=13mm] (apply2)
    {apply\\$\mat{A}\vv_2$};
  \node[font=\large, right=2mm of apply2] (dots) {$\cdots$};
  \node[stage, right=2mm of dots, text width=13mm] (applyk)
    {apply\\$\mat{A}\vv_k$};
  \node[font=\scriptsize\bfseries, simBlue, left=1mm of apply1, rotate=90, anchor=south]
    {run-time};
  \node[font=\scriptsize\itshape, simGray, below=1mm of apply1, xshift=14mm]
    {runs \emph{every} Krylov iteration --- reads the tables, never rebuilds them};
  \draw[flow, simGold] (gd.south) -- ++(0,-0.4) -| (apply1.north);
  \draw[flow, simGold] (gd.south) -- ++(0,-0.4) -| (apply2.north);
  \draw[flow, simGold] (gd.south) -- ++(0,-0.4) -| (applyk.north);
\end{tikzpicture}
\caption[Build-time versus run-time stratification]{The geometry factor and the two
index/basis tables are computed \emph{once}. A setup pass (gold, top) consumes the
mesh nodes and the quadrature rule and produces \code{geom\_data} together with the
restriction map and tabulated basis. Every subsequent operator application (blue,
bottom)---and an iterative solve performs hundreds---simply \emph{reads} these
build-time tables as it runs the five stages. The expensive geometry work is
amortized across the whole solve; only a mesh change (e.g.\ adaptive refinement,
\cref{ch:amr}) reruns the setup pass. This setup/apply split is the kernel's
instance of the stratification that organizes the whole simulator
(\cref{ch:strata}).}
\label{fig:appk-strata}
\end{figure}

\begin{keyidea}
The kernel carries two strata. The \emph{build-time} stratum---restriction map,
tabulated basis $\hat B,\hat G$, and the geometry factor \code{geom\_data}---is
computed once per $(\text{mesh},\text{order})$ and reused. The \emph{run-time}
stratum---the gathered, evaluated, weighted, integrated tensors---is rebuilt on
every application. Folding the Jacobian metric, quadrature weight, and material
coefficient into one precomputed \code{geom\_data} is precisely what reduces the
run-time weight stage $D$ to a single pointwise multiply.
\end{keyidea}

\section{Sum-factorization: the basis stage, reordered}
\label{sec:appk-sumfac}

The one stage with non-trivial arithmetic per apply is the basis evaluation $B$ and
its transpose. Taken naively, $B$ is a dense $P\times L$ contraction per element,
and at high order that is the dominant cost. On tensor-product elements
\emph{sum-factorization} cuts it dramatically by exploiting the product structure
of the basis---and it does so as a pure contraction reordering, computing
bit-for-bit the same result.

\subsection{The product structure}

On a tensor-product ($Q_p$) element the local dofs and the quadrature points each
lie on a $d$-dimensional grid of side $p+1$, so $L=P=(p+1)^d$, and the
$d$-dimensional basis factors into a product of one-dimensional bases:
\[
  \hat\phi_{i_1\cdots i_d}(\hat\xi_1,\dots,\hat\xi_d)
  \;=\; \hat\phi_{i_1}(\hat\xi_1)\,\cdots\,\hat\phi_{i_d}(\hat\xi_d).
\]
The naive evaluation ignores this and contracts against the full dense table:
\begin{equation}
  \text{naive cost} \;=\; P\cdot L \;=\; (p+1)^d\cdot(p+1)^d \;=\; (p+1)^{2d}
  \quad\text{multiply-adds per element.}
  \label{eq:appk-naive}
\end{equation}
Sum-factorization instead applies the \emph{one}-dimensional basis along each axis
in turn. Contract axis $1$ (a $(p+1)\times(p+1)$ matrix applied along one grid
direction, the other $d-1$ directions carried along as batch), then axis $2$, and
so on. Each of the $d$ sweeps costs $(p+1)$---the 1-D contraction---times
$(p+1)^d$---the rest of the grid carried along:
\begin{equation}
  \text{sum-factored cost} \;=\; d\cdot(p+1)\cdot(p+1)^d \;=\; d\,(p+1)^{d+1}.
  \label{eq:appk-factored}
\end{equation}
In the operator-count notation of \cref{ch:matrixfree} this is the move from
$O(P\cdot L)$ to $O(d\cdot P^{1/d}\cdot L)$, since one grid axis has length
$P^{1/d}=(p+1)$. \Cref{fig:appk-sumfac} draws the reordering for $d=3$ as a chain
of three batched 1-D contractions.

\begin{figure}[t]
\centering
\begin{tikzpicture}[node distance=5mm and 9mm]
  % naive: one dense contraction
  \node[data, text width=16mm] (din) {dofs\\$(p{+}1)^3$};
  \node[opbox, fill=simBgRust, draw=simRust, right=of din, text width=26mm]
    (naive) {dense $P\times L$\\\footnotesize $(p{+}1)^6$ ops};
  \node[data, right=of naive, text width=16mm] (dpts) {pts\\$(p{+}1)^3$};
  \draw[flow] (din) -- (naive);
  \draw[flow] (naive) -- (dpts);
  \node[font=\scriptsize\bfseries, simRust, left=1mm of din, rotate=90, anchor=south]
    {naive};

  % factored: three 1-D sweeps (d=3)
  \node[data, below=15mm of din, text width=16mm] (fin) {dofs\\$(p{+}1)^3$};
  \node[opbox, right=of fin, text width=17mm] (s1)
    {sweep $\xi_3$\\\footnotesize $(p{+}1)^4$};
  \node[opbox, right=of s1, text width=17mm] (s2)
    {sweep $\xi_2$\\\footnotesize $(p{+}1)^4$};
  \node[opbox, right=of s2, text width=17mm] (s3)
    {sweep $\xi_1$\\\footnotesize $(p{+}1)^4$};
  \node[data, right=of s3, text width=16mm] (fpts) {pts\\$(p{+}1)^3$};
  \draw[flow] (fin) -- (s1);
  \draw[flow] (s1) -- node[above, font=\scriptsize, simGray] {partial} (s2);
  \draw[flow] (s2) -- node[above, font=\scriptsize, simGray] {partial} (s3);
  \draw[flow] (s3) -- (fpts);
  \node[font=\scriptsize\bfseries, simBlue, left=1mm of fin, rotate=90, anchor=south]
    {factored};
\end{tikzpicture}
\caption[Sum-factorization as a contraction reordering ($d=3$)]{Basis evaluation on
a 3-D tensor-product element, two ways. \emph{Top (rust):} the naive contraction
treats the basis as one dense $P\times L$ matrix, costing $(p+1)^6$ multiply-adds
per element. \emph{Bottom (blue):} sum-factorization applies the one-dimensional
basis along each of the three axes in turn---a sweep in $\xi_3$, then $\xi_2$, then
$\xi_1$, each carrying the other two grid directions as batch and costing
$(p+1)^4$---for $3(p+1)^4$ total. The two produce \emph{identical} values; only the
contraction order differs. The saving is the move from $O(P\cdot L)$ to
$O(d\,P^{1/d}L)$, and the exponent gap ($2d$ versus $d+1$) makes it decisive at
high order.}
\label{fig:appk-sumfac}
\end{figure}

\begin{table}[t]
\centering
\caption[Sum-factorization operation counts, $d=3$]{Multiply-adds per element for
the basis-evaluation stage in three dimensions, naive $(p+1)^{2d}=(p+1)^6$ versus
sum-factored $d\,(p+1)^{d+1}=3(p+1)^4$. The speedup is unbounded in $p$: the
exponent drops from $2d=6$ to $d+1=4$, so the ratio grows like $(p+1)^{d-1}/d$.}
\label{tab:appk-sumfac}
\small
\setlength{\tabcolsep}{9pt}
\renewcommand{\arraystretch}{1.2}
\begin{tabular}{@{}cccc@{}}
\toprule
order $p$ & naive $(p{+}1)^6$ & factored $3(p{+}1)^4$ & speedup \\
\midrule
1 & 64               & 48                & $1.3\times$ \\
2 & 729              & 243               & $3.0\times$ \\
4 & \num{15625}      & \num{1875}        & $8.3\times$ \\
8 & \num{531441}     & \num{19683}       & $27\times$ \\
\bottomrule
\end{tabular}
\end{table}

\Cref{tab:appk-sumfac} tabulates the 3-D savings. At $p=8$ it is the difference
between $9^6=\num{531441}$ and $3\cdot 9^4=\num{19683}$ operations per element---a
factor of $27$, widening with every increment of $p$. Without sum-factorization the
basis stage would reintroduce the very $O(p^{2d})$ cost the matrix-free form set out
to escape; with it, high-order matrix-free is not merely affordable but the fastest
form available.

\begin{pitfall}
Sum-factorization is a \emph{transparent} performance trick (\cref{ch:bigidea}):
same result, faster, with no change to the answer, the conditioning, or the
rounding behavior that the algorithm depends on. It is recorded as a one-line note
on the basis stage, never as a separate algebraic claim. Do not mistake it for a
\emph{load-bearing} numerical trick---it reorders the sums of a single linear map
and nothing more. The reason to know it is that it is \emph{why} the high-order
kernel is fast, not a change to what the kernel computes.
\end{pitfall}

\section{The libCEED boundary: API versus implementation}
\label{sec:appk-extern}

We have now specified all five stages as tensor contractions. One honest
qualification remains, and it is the subject of this section: in the real Palace
build the innermost element-quadrature kernel is not hand-written from these
contractions but handed to a specialized accelerator library, libCEED, which owns
the device-level details---which loops to fuse, how to tile the quadrature points,
how to map them onto warps or wavefronts. The book draws this as a clean boundary
with two sides, and the distinction between them is load-bearing.

\subsection{The two surfaces}

The boundary has a \emph{kernel API} and a \emph{kernel implementation}, and the
spec keeps both, linked for review (\cref{ch:synthlib}).

\paragraph{The kernel API---what we call.} The API surface is the opaque library
contract: the shape signature and the semantics of the innermost kernel, with
\emph{no} claim about how it is realized. In the calculus it is written with the
\code{\#extern} marker---a node whose contract we state precisely but whose body we
delegate:
\begin{lstlisting}[style=l4style]
-- the kernel API: the contract is ours, the device body is libCEED's
quad_kernel :: GeomData -> Tensor[(E, P, C)] -> Tensor[(E, P, C')]
#extern quad_kernel   -- shape + semantics specified; body delegated
\end{lstlisting}
This is the violet, dashed \code{\#extern} node at the center of
\cref{fig:appk-pipeline}. Its contract is exactly the pointwise diagonal of
\cref{sec:appk-diagonal}: a shape-preserving (up to $C\to C'$) elementwise multiply
against \code{geom\_data}. The library owns the machine; we own the contract.

\paragraph{The kernel implementation---the contraction chain.} The implementation
surface is the constructive realization in our own tensor vocabulary: the five-stage
chain $G^{\mathsf T}B_{\diffop}^{\mathsf T}D B_{\diffop}G$ of this entire appendix,
built from the firm substrate operators \code{element\_restrict}, \code{basis\_apply},
\code{quad\_point\_contract}, and \code{geom\_factor\_build}. It makes a positive
claim and carries a normal rank: it \emph{is} the kernel, expressed from our
primitives.

\begin{lstlisting}[style=l4style]
-- the kernel implementation: the contraction chain that realizes quad_kernel,
-- written from our own firm substrate operators.  apply A v = scatter . Bt . D . B . gather
apply_kernel :: Operator -> Tensor[(N: ..)] -> Tensor[(N: ..)]
apply_kernel op v =
  scatter   op.restr           $   -- G^T  : [(E, L)]      -> [(N: ..)]   element_restrict^T
  integrate op.mode op.basis   $   -- B_D^T: [(E, P, C')]  -> [(E, L)]    basis_apply^T
  weight    op.geom            $   -- D    : [(E, P, C)]   -> [(E, P, C')] quad_point_contract  (= #extern quad_kernel)
  evaluate  op.mode op.basis   $   -- B_D  : [(E, L)]      -> [(E, P, C)] basis_apply
  gather    op.restr v             -- G    : [(N: ..)]     -> [(E, L)]    element_restrict
  where
    -- op.geom is the build-time geom_data, computed once:
    -- op.geom = geom_factor_build mesh_nodes quad_weights
\end{lstlisting}

\subsection{The correspondence}

The two surfaces are tied by a typed \emph{\code{realizes-kernel-api}}
correspondence: an edge from the implementation node to the API node asserting
``this contraction chain realizes that opaque contract.'' \Cref{fig:appk-extern}
draws it. The edge is a \emph{reference}, not a dependency: the implementation does
not build on the opaque API---it is an \emph{alternative} to it---so the edge
carries no rank or liveness constraint. A reviewer reads both sides and checks they
match: the opaque \code{\#extern} contract against the from-our-primitives chain.
That review is the discipline the boundary buys; the lowering-verifier audits the
correspondence.

\begin{figure}[t]
\centering
\resizebox{\linewidth}{!}{%
\begin{tikzpicture}[node distance=9mm and 22mm]
  % API node (extern, violet)
  \node[extern, text width=42mm] (api)
    {\textbf{kernel API}\;\code{\#extern}\\[2pt]
     \footnotesize \code{quad\_kernel ::}\\
     \footnotesize \code{GeomData -> [(E,P,C)] -> [(E,P,C')]}\\[2pt]
     \footnotesize\itshape opaque libCEED contract\\(shape $+$ semantics; body delegated)};
  % impl node (stage, blue)
  \node[stage, below=of api, text width=58mm] (impl)
    {\textbf{kernel implementation}\\[2pt]
     \footnotesize $G^{\mathsf T}\,B_{\diffop}^{\mathsf T}\,D\,B_{\diffop}\,G$\\[2pt]
     \footnotesize from firm substrate:\\
     \footnotesize \code{element\_restrict}, \code{basis\_apply},\\
     \footnotesize \code{quad\_point\_contract}, \code{geom\_factor\_build}};
  % the realizes-kernel-api reference edge (free, navigational)
  \draw[refedge] (impl.north) -- node[right, font=\scriptsize, simGray, align=left]
     {\code{realizes-}\\\code{kernel-api}\\(reference:\\free, reviewed)} (api.south);
  % the substrate, to the side
  \node[opbox, right=34mm of impl, text width=30mm] (sub)
    {\footnotesize firm \code{[(E,L)]}/\allowbreak\code{[(E,P,C)]}/\allowbreak\code{[(E,P,G)]}\\substrate ops};
  \draw[depedge] (impl.east) -- node[above, font=\scriptsize, simBlue] {composes} (sub.west);
\end{tikzpicture}%
}
\caption[The kernel API and the kernel implementation]{The two surfaces of the
libCEED boundary, kept distinct and linked for review (\cref{ch:synthlib}).
\emph{Top (violet, \code{\#extern}):} the kernel \emph{API}---the opaque libCEED
contract, stating the shape signature and pointwise semantics of the innermost
kernel with no claim about its realization. \emph{Bottom (blue):} the kernel
\emph{implementation}---the five-stage contraction chain of this appendix, built by
composing the firm \code{[(E,L)]}/\code{[(E,P,C)]}/\code{[(E,P,G)]} substrate
operators. The \code{realizes-kernel-api} edge (dotted) is a \emph{reference}, not a
dependency: the implementation is an alternative to the opaque API, not a consumer
of it, so the edge carries no rank or liveness---it is a correspondence a reviewer
checks. The same \code{\#extern} pattern recurs at the SLEPc eigensolve
(\cref{ch:eigenvalues}) and the triangular solves inside a smoother
(\cref{ch:smoothers}).}
\label{fig:appk-extern}
\end{figure}

\begin{keyidea}
The matrix-free operator is a chain of element-local tensor contractions over named
axes---the GPU-native form, and the kernel \emph{implementation} behind the libCEED
boundary. The boundary has two surfaces: the kernel \emph{API}, an opaque
\code{\#extern} contract stating shape and semantics; and the kernel
\emph{implementation}, the chain $G^{\mathsf T}B_{\diffop}^{\mathsf T}D B_{\diffop}G$
built from firm substrate operators. A \code{realizes-kernel-api} reference edge
links them for review. We own the mathematics; the library owns the machine; and
the two meet at a narrow, well-typed, reviewable interface.
\end{keyidea}

\section{Why this is the accelerator target}
\label{sec:appk-gpu}

Step back from the stages and look at what the whole operator turned out to be. The
finite-element operator $\mat{A}$---the discretization of Maxwell's weak form, the
object every solver in the book applies---is a composition of tensor contractions
over named axes. Gather is an indexed copy. Evaluate and integrate are dense
contractions of an \code{[(E, L)]} tensor against a tabulated basis, yielding
\code{[(E, P, C)]}, and on tensor-product elements they factor further into batched
1-D contractions. Weight is an elementwise multiply of \code{[(E, P, C)]} against
\code{[(E, P, G)]}. Scatter is an indexed sum. There is no sparse-matrix data
structure anywhere---no row pointers, no column indices, no stored entries. There
are only tensors and contractions.

This is precisely the shape a tensor accelerator is built to run. Every stage is a
pure function from tensors to tensors with no aliasing and no ordering hazard; the
element axis $E$ is a large batch dimension the hardware maps directly onto its
parallel units; the weight stage is a perfectly diagonal map; and the only stages
that couple data---gather and scatter---are pure index movement at the two ends,
not arithmetic in the hot loop. The assembled-matrix form, by contrast, is a
CPU-era data structure: irregular, pointer-chasing, bandwidth-bound, something a
tensor backend must be coaxed into. The matrix-free kernel is the form the
accelerator was built for, and it is exactly what the book's implementation view
(\cref{ch:synthlib}) expands the ``\,\lstinline[style=l4style]|apply A v|\,'' of
every solver into.

\summarysection
This appendix opened the five factors of $\mat{A}\vv=
G^{\mathsf T}B_{\diffop}^{\mathsf T}\,D\,B_{\diffop}\,G\,\vv$ down to their exact
tensor shapes. The gather $G$ is a pure indexed copy
$\texttt{[(N: \dots)]}\to\texttt{[(E, L)]}$ defined by the element-to-dof index map
$\iota(e,\ell)$; its transpose $G^{\mathsf T}$ is the scatter-add (assembly), the
operator's only reduction, and the composition $G^{\mathsf T}G=\operatorname{diag}
(m_n)$ is the dof-multiplicity diagonal. Basis evaluation $B_{\diffop}$ contracts
the local-dof axis $L$ against the tabulated reference basis to produce
\code{[(E, P, C)]}, with the mode---interp, grad, curl, div---and the component
count $C$ selected by the term's differential operator $\diffop$; $B^{\mathsf T}$
contracts the quadrature axis $P$ back to local dofs, carrying the discrete
integral. The center stage $D$ is the embarrassingly parallel diagonal: a pointwise
multiply against \code{geom\_data}, block-diagonal in the full $(E,P)$ index. That
geometry factor, of shape \code{[(E, P, G)]} with $G=2+d^2$ components, is built
once by \code{geom\_factor\_build} from the Jacobian metric
($\abs{\det J}$ for mass, $J^{-1}J^{-\mathsf T}\abs{\det J}$ for grad--grad), with
the quadrature weight and material coefficient pre-multiplied in. This is the
build-time stratum---computed once per $(\text{mesh},\text{order})$ and reused
across the hundreds of applies a solve performs---as against the transient run-time
tensors rebuilt on every action. On tensor-product elements sum-factorization
reorders the basis contraction from $O(P\cdot L)=O((p+1)^{2d})$ to
$O(d\,P^{1/d}L)=O(d(p+1)^{d+1})$---a transparent reordering, same answer, $27\times$
at $p=8$ in 3-D. The innermost kernel is sealed behind the libCEED boundary, which
the spec gives two surfaces: an opaque \code{\#extern} kernel \emph{API} stating
shape and semantics, and a kernel \emph{implementation}---the contraction chain of
this appendix from firm substrate operators---tied to it by a
\code{realizes-kernel-api} reference edge for review. Because every stage is a
contraction over named axes with no sparse-matrix structure anywhere, the
matrix-free kernel \emph{is} the GPU-native implementation view (\cref{ch:synthlib})
that the book has been building toward.

\furtherreading
The matrix-free, sum-factorized, tensor-contraction view of high-order
finite-element operators, and the kernel-API/implementation split realized through
libCEED's \code{\#extern} boundary, are the design philosophy of the CEED project;
its co-design survey~\citep{ceed2021} is the entry point to the libraries (libCEED,
MFEM) that realize the five-stage decomposition and the operation-count argument for
sum-factorization. For the finite-element foundations this appendix builds on---the
reference element, the tabulated bases and quadrature, the Jacobian metric, and the
element matrices the operator factors---Jin's text on the finite-element method in
electromagnetics~\citep{jin2014} gives the Maxwell-specific grounding, and the
mathematical theory of the spaces and their approximation properties is developed in
Brenner and Scott~\citep{brennerscott2008}. The build-time/run-time stratification
the geometry factor relies on is treated in its own right in \cref{ch:strata}; the
implementation view the contraction chain feeds is \cref{ch:synthlib}.

\chapter{The Layered Stack: A Reference Map}
\label{app:stack}

\chapterorient{This book taught one description of the simulator---the $L_4$
calculus of Part~III---grounded in the cited source and pure-function operators of
Parts~II and~IV, and composed into the modalities of Part~V. But the calculus is
only the top of a five-layer stack, and the full specification the book draws on is
built from all five layers and the rotations between them. This appendix is the
reference map of that stack. It is not a tutorial---\cref{ch:rotations} already
taught the rotation idea, and you should read that chapter first. It is a
\emph{legend}: a one-page-per-layer account of what each layer expresses and what
impedance it has turned; a reference table of the four rotations; a catalogue of the
obstructions that do not rotate; and a chart of how the book's parts map onto the
layers. Keep it beside you when you go to read the underlying spec.}

This appendix maps the layered representation that underlies the entire book. Where
\cref{ch:rotations} \emph{taught} the rotation idea---what a rotation is, the three
tests that detect one, the five justification kinds, and how an obstruction is an
honest failure---this appendix \emph{tabulates} the structure for reference. The
reader who has finished the book and now wants to open the full specification will
find here the map: the five layers in order, the four edges between them, the
negative results, and the correspondence to the parts already read.

\begin{keyidea}
The stack is \textbf{five layers, four rotations, and a catalogue of obstructions}.
The five layers ($L_0$ source, $L_1$ pure functions, $L_2$ algebra, $L_3$
tensor-field operations, $L_4$ calculus) are each a complete, correct account of the
simulator in its own vocabulary. The four rotations between adjacent layers each turn
exactly one impedance---mutation, fusion, iteration, coordination---without changing
the numbers. And the obstructions---the loops that will \emph{not} rotate to a
whole-field form---are findings, not failures: they mark precisely where the
abstract tensor picture meets a wall, and why. The book taught $L_4$; this map shows
the rest.
\end{keyidea}

\section{The five layers}
\label{sec:fivelayers}

We number the layers from the ground up: $L_0$ is the literal source, $L_4$ the
calculus. \Cref{fig:appstack} is the hero figure of this appendix---the polished
reference version of the stack first drawn in \cref{ch:rotations}. It shows all five
bands and the four labeled rotation arrows between them, each tagged with the
impedance it turns. The subsections that follow give each layer its own paragraph:
what it expresses, and what impedance it has rotated relative to the layer below.

\begin{figure}[p]
\centering
\begin{tikzpicture}[node distance=0mm]
  \def\lw{96mm}    % layer width
  \def\lh{14mm}    % layer height
  \def\gap{12mm}   % vertical gap between bands (room for rotation arrow + 2 labels)
  % five stacked layer bands, L0 at bottom -> L4 at top
  \foreach \i/\name/\desc/\col/\bg in {%
      0/$L_0$/the cited source {\textbullet} Palace C++ {\&} MFEM {\textbullet} reference notes/simInk/simBgGray,
      1/$L_1$/pure functions over immutable tensors {\textbullet} no in-place mutation/simGreen/simBgGreen,
      2/$L_2$/algebraic decompositions {\textbullet} HPC fusions erased/simTeal/simBgGray,
      3/$L_3$/global tensor-field operations {\textbullet} element loops lifted/simBlue/simBgBlue,
      4/$L_4$/the calculus {\textbullet} combinators {\textbullet} state monad {\textbullet} strata/simViolet/simBgViolet}{
    \node[draw=\col, fill=\bg, very thick, rounded corners=3pt,
          minimum width=\lw, minimum height=\lh, align=center, font=\small]
      (L\i) at (0, \i*\dimexpr\lh+\gap\relax) {%
        \textbf{\name}\quad\textcolor{simGray}{\footnotesize\desc}};
  }
  % rotation arrows between adjacent bands
  \foreach \lo/\hi in {0/1,1/2,2/3,3/4}{
    \draw[depedge, line width=1.3pt] (L\lo.north) -- (L\hi.south);
  }
  % rotation labels (name + impedance turned), placed to the RIGHT of each arrow gap
  \def\lx{0.52*\lw+3mm}
  \foreach \i/\rot/\imp in {%
      0/mutation rotation/{mutation $\to$ purity},
      1/fusion rotation/{fused kernel $\to$ base algebra},
      2/iteration rotation/{element loop $\to$ whole-field op},
      3/calculus rotation/{control flow $\to$ named combinators}}{
    \node[font=\scriptsize\bfseries\itshape, simRust, anchor=west]
      at (\lx, \i*\dimexpr\lh+\gap\relax+0.5*\lh+0.5*\gap+1.6mm) {\rot};
    \node[font=\scriptsize, simGray, anchor=west]
      at (\lx, \i*\dimexpr\lh+\gap\relax+0.5*\lh+0.5*\gap-1.6mm) {turns: \imp};
  }
  % side bracket: "complete & correct in itself"
  \draw[decorate, decoration={brace, amplitude=6pt}, simGray, thick]
    (-0.52*\lw-3mm, -0.5*\lh) -- (-0.52*\lw-3mm, 4*\dimexpr\lh+\gap\relax+0.5*\lh);
  \node[rotate=90, font=\scriptsize\itshape, simGray, anchor=south]
    at (-0.52*\lw-9mm, 2*\dimexpr\lh+\gap\relax) {each layer complete \& correct in its own vocabulary};
  % top / bottom annotations
  \node[font=\scriptsize\itshape, simViolet, anchor=south] at (0, 4*\dimexpr\lh+\gap\relax+0.5*\lh+2mm)
    {\textbf{top:} the accelerator-ready surface (what the backend consumes)};
  \node[font=\scriptsize\itshape, simInk, anchor=north] at (0, -0.5*\lh-2mm)
    {\textbf{ground:} every claim in the stack rests on an $L_0$ citation};
\end{tikzpicture}
\caption[The five-layer impedance-matching stack]{The impedance-matching stack, in
full. Read it from the ground up. $L_0$ is the cited Palace/MFEM source---the only
layer that \emph{is} the machine rather than a description of it. Each rotation above
turns exactly one impedance while computing the same numbers: \emph{mutation} into
purity ($L_0\!\to\!L_1$), a \emph{fused kernel} into the base algebra it implements
($L_1\!\to\!L_2$), an explicit \emph{element loop} into a whole-field operation
($L_2\!\to\!L_3$), and hand-written \emph{control flow} into named combinators
($L_3\!\to\!L_4$). Every band is a complete, correct account of the simulator in its
own vocabulary; the arrows are translations between vocabularies, never mere
renamings (\cref{ch:rotations}). The book has been teaching the $L_4$ view at the top,
with Parts~II and~IV supplying the $L_0$--$L_1$ grounding at the bottom.}
\label{fig:appstack}
\end{figure}

\subsection*{$L_0$ --- the cited source}

$L_0$ is the ground truth: ranges of Palace's C++ (and the MFEM library beneath it),
quoted by file and line, together with a small set of cross-cutting \emph{reference
notes} that explain the conventions those ranges use---output-argument versus
receiver mutation, the \code{Vector}/\code{ComplexVector} element-type duality,
free-function versus method-form symbols, transparent versus load-bearing
optimization tricks. $L_0$ is not ``ours''; it is the machine as written. Nothing is
interpreted here, only read and cited. It has turned \emph{no} impedance---there is no
layer below it---and every claim anywhere higher in the stack ultimately rests on an
$L_0$ citation. The reference notes are deliberately lean (a few paragraphs of
interpretation plus representative citations each), so that the higher layers can
\emph{point} at a convention rather than re-state it. $L_0$ is where the grep-and-IDE
workflow lives; the citation format is the plain-text \code{path/file.ext:start-end}
the whole book uses.

\subsection*{$L_1$ --- the mutation rotation}

$L_1$ re-expresses each $L_0$ operation as a \emph{pure function over immutable
tensors}. The destination array that $L_0$ overwrote in place becomes a returned
value; \lstinline[style=l4style]|y.Add(alpha, x)| (a mutating member method) becomes
\lstinline[style=l4style]|y_new = axpy(alpha, x, y_old)|, with the $L_0$ destination
buffer gone from the signature. An operator application \lstinline[style=l4style]|A.Mult(x, y)|
that wrote into \code{y} becomes \lstinline[style=l4style]|y = apply_linop(A, x)|, no
destination mentioned. This is the \emph{mutation rotation}, applied operator by
operator: it has turned the impedance \emph{mutation into purity}. $L_1$ is the layer
the book mostly worked in for its operator definitions (Part~IV)---it is the closest
pure-functional layer to the source, its structure still following the source loop,
with workspace scratch buffers erased and only semantically-meaningful aliasing kept.
Reduction-tree non-associativity (the order in which \code{dot} and \code{nrm2} sum)
is preserved here as a \emph{load-bearing} algebraic claim, not quietly discarded.

\subsection*{$L_2$ --- the fusion rotation}

$L_2$ unfolds each $L_1$ function back into a \emph{composition of base algebraic
primitives}, with the performance tricks erased. A fused kernel that computed
$a x + b y + c z$ in one cache-friendly sweep is decomposed into the scaled sums it
stands for; a packed, blocked, SIMD-tuned matrix apply becomes the plain linear map it
implements; a batched specialized BLAS call becomes a composition of base operations.
This is the \emph{fusion rotation}: it has turned the impedance \emph{a fused kernel
into its base algebra}. The discipline of \cref{ch:bigidea} governs the erasure---a
\emph{transparent} performance trick (fusion, tiling, packing) is unfolded and gets a
one-line note, while a \emph{load-bearing} numerical trick (a non-associative reduction
ordering, a deterministic-versus-atomic accumulation) is \emph{preserved} as an
explicit algebraic claim with the property it buys called out. $L_2$ is where the
fold cohorts live: \code{linear\_combination} folds the term axis to a tensor,
\code{inner\_product} folds the length axis to a scalar, and the once-standalone leaves
(\code{axpy}, \code{dot}, \code{scal}, \dots) become thin specializations of their
combinator.

\subsection*{$L_3$ --- the iteration rotation}

$L_3$ lifts the explicit per-degree-of-freedom loops of $L_2$, where the algebra
permits, to \emph{whole-field operations}: ``add these two fields,'' ``apply this
operator to this field,'' with no surviving element loop in the spec form. This is the
\emph{iteration rotation}: it has turned the impedance \emph{an explicit element loop
into a whole-field operation}. It is the rotation that does \emph{not} always go
through---and where it fails, the failure is recorded as a first-class
\emph{obstruction} (\cref{sec:appobstruction}). A clean BLAS-1 update lifts trivially
(every entry independent); a Gauss--Seidel sweep does not (each entry reads the
previous one). $L_3$ therefore arranges its operators on an obstruction-profile
spectrum, from the obstruction-free leaves (\code{apply\_linop}, \code{dot},
\code{normalize}) through operators that carry an obstruction \emph{by reference}
(the divergence-free projector, whose inner solve iterates) to the four
\emph{partial-obstruction} operators whose body lifts cleanly but whose loop does not
(\code{chebyshev}, \code{eigsolve}, \code{orthogonalize}, \code{fold\_solve}).

\subsection*{$L_4$ --- the calculus}

$L_4$ is the small, formally defined graph-evaluation calculus the book taught in
Part~III: high-order combinators (\cref{ch:fold,ch:combinators}), the state monad
(\cref{ch:monad}), the three-stratum state stratification (\cref{ch:strata}), and
operators as values (\cref{ch:operators}). It is \emph{vocabulary, not architecture}---it
says what the algorithm \emph{is}, abstractly, in a form an accelerator backend can
consume, without committing to eager-versus-traced execution or to any particular
container or monad implementation. The \emph{calculus rotation} from $L_3$ replaces
hand-written control flow with named combinators: it has turned the impedance
\emph{hand-written coordination into a named combinator}. A solver's restart-and-converge
loop, threaded by hand at $L_3$, becomes \code{ksp\_solve} folding \code{krylov\_step}
through \code{iterate\_while}; a time march becomes \code{fold\_solve} over a schedule;
an RHS family becomes \code{solve\_family} mapping a fixed operator. $L_4$ is the
backend-lowering target---the whole-tensor, no-aliasing surface the GPU wants---and the
implementation rendered \emph{as} that surface is the synthesized library of
\cref{ch:synthlib}.

\begin{notationbox}
The layer numbers run \emph{opposite} to the direction of construction. The
\emph{stack} is numbered ground-up ($L_0$ source at the bottom, $L_4$ calculus at the
top), but the \emph{rotations} are performed and read \emph{up} the stack: a rotation
``$L_n\!\to\!L_{n+1}$'' takes the lower form to the higher one. A
\emph{lowering}---the spec's name for the formal correspondence---points the other
way, ``$L_{n+1}>L_n$,'' and narrates how the higher form \emph{lowers} into the lower
one. Both describe the same edge; the rotation is the act, the lowering is the
recorded correspondence. The spec stores lowerings in per-adjacent-edge directories
(\code{L4-L3/}, \code{L3-L2/}, \code{L2-L1/}, \code{L1-L0/}), one chapter per theme.
\end{notationbox}

\section{The four rotations between adjacent layers}
\label{sec:approtations}

The four edges of the stack are the four rotations. Each is a \emph{translation across
vocabularies}, not a one-to-one renaming (\cref{def:rotation}): the lower form is
re-expressed in a genuinely different idiom in which one impedance has been turned. A
move that merely pastes new names on the old operations---same state threaded, same
grain, same admitted algorithm in every slot---turns no impedance and is a
\emph{renaming}, the smell \cref{sec:tests} teaches you to catch. \Cref{tab:rotations}
is the reference table: each layer pair, the impedance it rotates, what changes in the
vocabulary, and a worked example drawn from the book.

\begin{table}[t]
\centering
\caption[The four rotations: a reference table]{The four rotations between adjacent
layers, as a reference table. Each row names the edge, the impedance it turns, the
vocabulary shift it performs (the genuine translation, not a rename), and a concrete
example met in the book. Read up the stack ($L_0$ at the bottom): mutation becomes
purity, fusion becomes base algebra, the element loop becomes a whole-field operation,
and hand-written coordination becomes a named combinator.}
\label{tab:rotations}
\small
\renewcommand{\arraystretch}{1.35}
\begin{tabular}{@{}>{\raggedright\arraybackslash}p{0.115\textwidth}
                   >{\raggedright\arraybackslash}p{0.155\textwidth}
                   >{\raggedright\arraybackslash}p{0.30\textwidth}
                   >{\raggedright\arraybackslash}p{0.30\textwidth}@{}}
\toprule
\textbf{Edge} & \textbf{Impedance turned} & \textbf{What rotates (the translation)} &
\textbf{Example} \\
\midrule
\textbf{$L_4 > L_3$} \newline calculus rotation &
hand-written control flow $\to$ named combinator &
The monad / strata / \code{readonly} wrappers dissolve; coordination loops become
\code{iterate\_while}, \code{fold\_solve}, \code{solve\_family}. State that was
threaded by hand is hidden in a combinator's contract. &
\code{ksp\_solve} folding \code{krylov\_step} through \code{iterate\_while} replaces a
hand-threaded restart loop; the time march becomes \code{fold\_solve} over a schedule
(\cref{ch:transient}). \\
\textbf{$L_3 > L_2$} \newline iteration rotation &
explicit element loop $\to$ whole-field op &
The per-DoF loop is erased and the operation re-stated on whole tensors---\emph{or},
where it cannot be, an \code{obstruction} marker is named first-class and shadows to an
$L_2$ non-law. &
\lstinline[style=l4style]|for i: y[i]=a*x[i]+y[i]| becomes the single whole-field
\code{axpy}; the Gauss--Seidel sweep refuses and records an obstruction
(\cref{sec:appobstruction}). \\
\textbf{$L_2 > L_1$} \newline fusion rotation &
fused kernel $\to$ base algebra &
A fused / packed / batched kernel is unfolded into a composition of base primitives;
transparent tricks get a note, load-bearing numerical tricks are preserved as explicit
claims. &
A fused $ax+by+cz$ sweep unfolds into scaled sums; the leaves \code{axpy}/\code{dot}
become specializations of the \code{linear\_combination} / \code{inner\_product}
folds. \\
\textbf{$L_1 > L_0$} \newline mutation rotation &
mutation $\to$ purity &
The in-place destination buffer disappears from the signature; a mutating method
becomes a function returning a fresh value. The lowering reintroduces the buffer when
reading down to source. &
\lstinline[style=l4style]|y.Add(alpha, x)| (writes \code{y}) becomes
\lstinline[style=l4style]|axpy alpha x y| (returns fresh);
\lstinline[style=l4style]|A.Mult(x, y)| becomes \code{apply\_linop}
(\cref{ch:bigidea}). \\
\bottomrule
\end{tabular}
\end{table}

A few structural facts about the table are worth stating once. First, not every edge
is substantive for every operator: a leaf that is already at the right grain---an
\code{axpy} that is whole-field by construction---``carries through'' an edge unchanged,
and the spec records that as an in-line identity note rather than a separate lowering
theme (a mirrored entry with a thin theme would be the renaming smell). The genuine
content of an edge lives in the operators that \emph{do} rotate: the solver caps at
$L_4>L_3$, the fold cohorts at $L_2>L_1$, the constructed-operator gates and the
sequential obstructions at $L_3>L_2$.

Second, an edge can be substantive in one direction's vocabulary and identity-in-form
in the other. The Krylov step's \emph{body} is identity-in-form across $L_4>L_3>L_2$
(the same primitive sequence, only the wrappers dissolving), while its \emph{driver}---the
outer restart-and-converge loop---rotates substantively at every edge (the iteration
view is rendered at $L_3$, erased at $L_2$, named as a combinator at $L_4$). The spec
captures this as a kernel-identity-plus-driver-non-identity pair of themes.

Third, identity rotations that span \emph{non-adjacent} layers are annotated in-line,
not given their own directory. There is no \code{L3-L1/}; when an operator's $L_3$ body
is value-thread-isomorphic to its $L_1$ form through an identity-like $L_2$ absorption,
that relationship is the transitive composition of the adjacent-edge themes, noted in
the $L_n$ entry.

\section{Obstructions: the things that do not rotate}
\label{sec:appobstruction}

Not everything rotates. Some work, faithfully described at $L_2$, genuinely cannot be
lifted to a clean $L_3$ whole-field operation, because its very meaning depends on doing
things in order. When that happens the spec does \emph{not} fake a rotation and does
\emph{not} treat the failure as a defect in the method. It records an
\emph{obstruction}: a precise statement of what blocks the lift and why. An obstruction
is a first-class negative result---a finding, not a failure (\cref{ch:rotations}). The
obstructions are some of the most valuable output in the stack, because they say
exactly where the abstract whole-tensor picture meets a wall.

There are two \emph{kinds} of obstruction, distinguished by how the blocking loop
relates to Palace, and the kind determines whether the obstruction can ever be
discharged. \Cref{tab:obstructions} catalogues both, with the book chapters where each
was met; \cref{fig:obstructionkinds} draws the distinction.

\begin{table}[t]
\centering
\caption[The obstruction catalogue]{The catalogue of non-rotation. Two kinds of
obstruction block the iteration rotation, distinguished by \emph{whose} loop carries
the blockage. A \emph{sequential-obstruction} is a Palace-authored loop whose
$i$-th step reads the $(i\!-\!1)$-th step's output (a loop-carried dependency); Palace
\emph{renders} it, so it survives as ``these reordering laws do not hold,'' and it
could in principle be re-described. An \emph{opaque-library-ownership} obstruction is a
loop that lives entirely in third-party code (SLEPc, MFEM); Palace authors no loop to
render, only a boundary marker, and it cannot be discharged unless the upstream
architecture changes.}
\label{tab:obstructions}
\small
\renewcommand{\arraystretch}{1.3}
\begin{tabular}{@{}>{\raggedright\arraybackslash}p{0.20\textwidth}
                   >{\raggedright\arraybackslash}p{0.34\textwidth}
                   >{\raggedright\arraybackslash}p{0.34\textwidth}@{}}
\toprule
\textbf{Kind} & \textbf{What blocks the lift} & \textbf{Where the book met it} \\
\midrule
\textbf{sequential-} \newline \textbf{obstruction} \newline \emph{(Palace-authored
loop)} &
A loop-carried dependency: step $i$ reads step $i\!-\!1$'s output, so there is no
``all entries at once'' form. Palace \emph{renders} the loop (an explicit tail
recursion); the lift erases the iteration view and the obstruction survives as a lost
reordering law. Root is usually numerical stability. &
The MGS arm of orthogonalization---the serial rank-1-projector chain that buys
roundoff-orthogonality (\cref{ch:orthog}); the Chebyshev smoother's inner $k$-recurrence
and outer Richardson sweep; the Gauss--Seidel-flavored multigrid smoother sweep; the
transient time march's carry-threading (\cref{ch:transient}). \\
\textbf{opaque-} \newline \textbf{library-} \newline \textbf{ownership} \newline
\emph{(third-party loop)} &
The iteration lives entirely inside a library Palace calls but does not implement.
Palace authors \emph{no} loop, so there is nothing to render---only a marker at the
library boundary. No promotion route unless Palace re-architects; these are the
\code{\#extern} edges of \cref{app:kernel}. &
The SLEPc / ARPACK eigensolve loop, handed the operators and run to convergence inside
the library (\cref{ch:eigenmodes}); the MFEM ODE time-integrator step inside
\code{fold\_solve} (\cref{ch:transient}); the libCEED element-quadrature kernel inside
\code{fe\_assemble}. \\
\bottomrule
\end{tabular}
\end{table}

\begin{figure}[t]
\centering
\begin{tikzpicture}[node distance=6mm]
  % ---- LEFT: liftable, independent ----
  \begin{scope}
    \node[font=\bfseries\small, simGreen] at (1.4,1.85) {Liftable: independent updates};
    \foreach \i in {0,1,2,3}{
      \node[opbox, minimum width=8mm, minimum height=6mm] (a\i) at (\i*0.95,0.7) {$y_\i$};
      \node[data, minimum width=8mm, minimum height=5mm] (x\i) at (\i*0.95,-0.4) {$x_\i$};
      \draw[flow] (x\i) -- (a\i);
    }
    \node[align=center, font=\scriptsize\itshape, simGray, text width=42mm] at (1.4,-1.3)
      {each $y_i = a x_i + y_i$ reads only its own input $\Rightarrow$ one whole-field op};
    \node[draw=simGreen, thick, rounded corners=2pt, fit=(a0)(a3), inner sep=3pt,
          fill=simGreen!8] {};
    \node[font=\scriptsize\bfseries, simGreen] at (1.4,1.2) {$L_3$ whole-field op};
  \end{scope}
  % ---- MIDDLE: sequential-obstruction (Palace-authored) ----
  \begin{scope}[shift={(5.4,0)}]
    \node[font=\bfseries\small, simRust, align=center, text width=42mm] at (1.4,1.85)
      {Sequential-obstruction\\[-1pt]\scriptsize(Palace authors the loop)};
    \foreach \i in {0,1,2,3}{
      \node[opbox, draw=simRust, minimum width=8mm, minimum height=6mm] (b\i) at (\i*0.95,0.7) {$y_\i$};
      \node[data, minimum width=8mm, minimum height=5mm] (u\i) at (\i*0.95,-0.4) {$x_\i$};
      \draw[flow] (u\i) -- (b\i);
    }
    \foreach \i/\j in {0/1,1/2,2/3}{
      \draw[backflow] (b\i.east) -- (b\j.west);}
    \node[align=center, font=\scriptsize\itshape, simGray, text width=42mm] at (1.4,-1.3)
      {$y_i$ reads the \emph{already-updated} $y_{i-1}$ $\Rightarrow$ loop renders, does not lift};
    \node[font=\scriptsize\bfseries, simRust] at (1.4,1.2) {loop-carried dependency};
  \end{scope}
  % ---- RIGHT: opaque-library-ownership ----
  \begin{scope}[shift={(10.8,0)}]
    \node[font=\bfseries\small, simViolet, align=center, text width=40mm] at (1.1,1.85)
      {Opaque-library-ownership\\[-1pt]\scriptsize(Palace authors no loop)};
    \node[extern, minimum width=34mm, minimum height=18mm, align=center] (lib) at (1.1,0.15)
      {\textbf{\#extern}\\[2pt]\scriptsize the loop lives entirely\\in SLEPc / MFEM / libCEED;\\only a boundary marker};
    \node[data, minimum width=10mm, above=4mm of lib] (in2) {operators in};
    \node[data, minimum width=10mm, below=4mm of lib] (out2) {result out};
    \draw[flow] (in2) -- (lib);
    \draw[flow] (lib) -- (out2);
    \node[align=center, font=\scriptsize\itshape, simGray, text width=40mm] at (1.1,-1.95)
      {nothing to render $\Rightarrow$ a marker, not a recurrence};
  \end{scope}
\end{tikzpicture}
\caption[The two kinds of obstruction]{Why some loops rotate to $L_3$ and some do not,
and the two kinds of those that do not. \emph{Left (green):} an \code{axpy}-style
update writes $y_i = a x_i + y_i$ reading only its own input; the entries are
independent, so the loop lifts to one whole-field operation. \emph{Middle (rust):} a
\emph{sequential-obstruction}---a Gauss--Seidel-like sweep where $y_i$ reads the
already-updated $y_{i-1}$ (the dashed loop-carried dependency). Palace authors this
loop, so $L_3$ \emph{renders} it as an explicit tail recursion; the lift to $L_2$
erases the iteration view and the obstruction survives as a lost reordering law. It
could, in principle, be re-described. \emph{Right (violet):} an
\emph{opaque-library-ownership} obstruction---the loop lives entirely in third-party
code (the eigensolver, the ODE step, the quadrature kernel). Palace authors no loop, so
there is nothing to render; only a boundary marker, the \code{\#extern} edge of
\cref{app:kernel}, and no discharge route unless the upstream changes.}
\label{fig:obstructionkinds}
\end{figure}

The distinction is not pedantry; it changes how an obstruction can ever be removed. A
\emph{rendered} recurrence (the sequential kind) records a loop Palace chose to
surface and could in principle re-describe---and indeed, for the three opaque kernels
that are spine dependencies, the book provides \emph{both} an opaque
\code{\#extern} contract and a separate, constructive realization in our own
vocabulary, linked by a reviewable correspondence (\cref{app:kernel}'s
kernel-API/kernel-implementation split; the eigensolve's Krylov--Schur reconstruction,
\cref{ch:eigenmodes}). An \emph{opaque-library} marker, by contrast, records a boundary
Palace never sees inside and will not discharge unless the upstream architecture
changes.

A finer reading---the \emph{erasure scope} of \cref{ch:rotations}---asks how much of
the iteration view a single layer-hop erases: a single sequential loop (a linear
solve's outer iteration), a nested double loop (a polynomial smoother's inner
recurrence inside an outer sweep), or an opaque library call with no loop to render at
all. The rung names what kind of obstruction you are looking at; \cref{ch:rotations}'s
erasure-scope ladder draws all three.

\begin{pitfall}
An obstruction is never an excuse to stop, and never a whole-operator verdict when
only the loop resists. The honest record is precise about the boundary: the per-step
\emph{body} usually rotates cleanly (it is identity-in-form to the layer below), and
only the \emph{iteration} fails. That is exactly the \emph{partial-obstruction} status
the spec uses for \code{chebyshev}, \code{eigsolve}, \code{orthogonalize}, and
\code{fold\_solve}---``the body lifts; the loop does not, because of this cited
dependency.'' Reading one of these as a total failure, or forcing a false whole-field
form to avoid recording the obstruction, both throw away the finding the obstruction
\emph{is}.
\end{pitfall}

\section{How the book navigated the stack}
\label{sec:appnavigation}

The book did not present the five layers in order. It taught the \emph{top}---the $L_4$
calculus---as the main object (Part~III), grounded it in the \emph{bottom}---the cited
source and the pure-function operators (Parts~II and~IV)---and then \emph{composed} the
top into the user-facing modalities (Part~V). The synthesized library (Part~VI)
rendered that same $L_4$ surface as code. \Cref{fig:partsmap} charts which part teaches
which layer, so the reader who now wants the full spec knows where each piece of the
book already sits in the stack.

\begin{figure}[t]
\centering
\begin{tikzpicture}[font=\footnotesize, >=Stealth, node distance=0mm]
  \def\lw{58mm}
  \def\lh{9mm}
  \def\gap{3.0mm}
  % the five layer bands (left column)
  \foreach \i/\name/\col/\bg in {%
      0/$L_0$ cited source/simInk/simBgGray,
      1/$L_1$ pure functions/simGreen/simBgGreen,
      2/$L_2$ algebra/simTeal/simBgGray,
      3/$L_3$ tensor-field ops/simBlue/simBgBlue,
      4/$L_4$ the calculus/simViolet/simBgViolet}{
    \node[draw=\col, fill=\bg, thick, rounded corners=2pt,
          minimum width=\lw, minimum height=\lh, align=center, font=\small]
      (B\i) at (0, \i*\dimexpr\lh+\gap\relax) {\textbf{\name}};
  }
  % the part annotations (right column), aligned to layers, with connector braces
  \def\px{0.5*\lw+6mm}
  % Parts II/IV -> L0/L1  (grounding)
  \draw[decorate, decoration={brace, amplitude=5pt, mirror}, simGreen, thick]
    ($(B0.east)+(3mm,-0.5*\lh+1mm)$) -- ($(B1.east)+(3mm,0.5*\lh-1mm)$);
  \node[product, anchor=west, text width=44mm, align=left] at (\px, 0.5*\dimexpr\lh+\gap\relax)
    {\textbf{Parts II \& IV} $\to$ $L_0$/$L_1$\\[1pt]\scriptsize the field theory + the
     pure-function operators: the bottom-up grounding};
  % L2 -> intervening
  \node[opbox, anchor=west, text width=44mm, align=left] at (\px, 2*\dimexpr\lh+\gap\relax)
    {\textbf{the intervening algebra}\\[1pt]\scriptsize $L_2$/$L_3$: the stepping-stones
     the rotations pass through (the spec's middle)};
  \draw[refedge] (B2.east) -- (\px-1mm, 2*\dimexpr\lh+\gap\relax);
  % Part III -> L4 (the taught surface)
  \draw[decorate, decoration={brace, amplitude=5pt, mirror}, simViolet, thick]
    ($(B4.east)+(3mm,-0.5*\lh+1mm)$) -- ($(B4.east)+(3mm,0.5*\lh-1mm)$);
  \node[stage, fill=simBgViolet, draw=simViolet, very thick, anchor=west,
        text width=44mm, align=left] at (\px, 4*\dimexpr\lh+\gap\relax)
    {\textbf{Part III} $\to$ $L_4$ \emph{(the taught surface)}\\[1pt]\scriptsize the calculus:
     fold, combinators, monad, strata, operators-as-values};
  % Part V / VI annotations along the top
  \node[font=\scriptsize\itshape, simRust, align=center, text width=0.95\linewidth]
    at (0.30\linewidth, 4*\dimexpr\lh+\gap\relax+0.5*\lh+5mm)
    {\textbf{Part V} reads the modalities as $L_4$ compositions (drivers as composition
     roots); \quad \textbf{Part VI} renders the $L_4$ surface as the synthesized library
     (\cref{ch:synthlib}).};
\end{tikzpicture}
\caption[How the book's parts map onto the layers]{The correspondence between the
book's parts and the stack's layers. The book taught the \emph{top} of the stack as its
primary object---\textbf{Part~III} is the $L_4$ calculus (violet, the taught
surface)---and grounded it in the \emph{bottom}: \textbf{Parts~II and~IV} supply the
$L_0$ cited source and the $L_1$ pure-function operators (green, the bottom-up
grounding). The intervening $L_2$/$L_3$ algebra (teal/blue) is the spec's middle---the
stepping-stones the rotations pass through, which the book touched mainly when an
operator's iteration rotation or fusion rotation was load-bearing. \textbf{Part~V} then
reads the five modalities as $L_4$ compositions (each driver a composition root over
the calculus), and \textbf{Part~VI} renders that same $L_4$ surface as the synthesized
library. The reader who wants the full specification now has the map: the book is the
$L_4$-and-grounding view; the spec is all five layers and the rotations between them.}
\label{fig:partsmap}
\end{figure}

The choice to teach $L_4$ first, rather than climb the stack from $L_0$, was
deliberate. A direct $L_0\!\to\!L_4$ leap would be a leap of faith; the stack replaces
it with a sequence of short, auditable rotations, but a \emph{reader} does not need to
re-perform that sequence to understand the result. The result---what the algorithm
\emph{is}, abstractly---is the $L_4$ calculus, and that is what the book taught. The
grounding in Parts~II and~IV is what lets you trust the calculus is faithful; the
rotations of \cref{ch:rotations} are what connect the two; and this appendix is the map
that lets you go read the rest. The payoff of the whole layered construction is three
things at once: \emph{auditability} (each short move is checkable), \emph{honesty} (the
obstructions are recorded, not papered over), and a \emph{target} (the $L_4$ surface an
accelerator backend consumes, rendered as the library of \cref{ch:synthlib}).

\summarysection
The simulator is described not once but at a stack of five \emph{layers}, each complete
and correct in its own vocabulary: $L_0$, the cited Palace/MFEM source plus its
reference notes; $L_1$, pure functions over immutable tensors; $L_2$, algebraic
decompositions with HPC fusions erased; $L_3$, global tensor-field operations; and
$L_4$, the calculus of Part~III. Adjacent layers are joined by one of \emph{four
rotations}, each turning exactly one impedance while computing the same numbers: the
\emph{mutation rotation} ($L_1>L_0$, mutation into purity), the \emph{fusion rotation}
($L_2>L_1$, a fused kernel into base algebra), the \emph{iteration rotation} ($L_3>L_2$,
an element loop into a whole-field operation), and the \emph{calculus rotation}
($L_4>L_3$, hand-written control flow into a named combinator). Every rotation is a
translation across vocabularies, never a renaming. Where the iteration rotation cannot
go through, the spec records a first-class \emph{obstruction}---of two kinds: a
\emph{sequential-obstruction} (a Palace-authored loop-carried dependency, rendered and
re-describable in principle---Gauss--Seidel smoothing, MGS orthogonalization, the
Chebyshev recurrence) and an \emph{opaque-library-ownership} obstruction (a loop owned
by SLEPc, MFEM, or libCEED, marked but not rendered, the \code{\#extern} edges of
\cref{app:kernel})---each tied to the chapter where the book met it
(\cref{ch:orthog,ch:eigenmodes,ch:transient}). The book navigated the stack by teaching
the top ($L_4$, Part~III), grounding it in the bottom ($L_0$/$L_1$, Parts~II and~IV),
and composing it into the modalities (Part~V) and the synthesized library (Part~VI).
The reader who wants the full specification now holds the map: five layers, four
rotations, and the obstructions that, by not rotating, are findings rather than
failures.

\furtherreading
The layered, rotation-based method this appendix maps is purpose-built for this project,
but the obstructions it catalogues are the standard hard cases of numerical linear
algebra. The sequential reductions that resist the iteration rotation---Gauss--Seidel
and SSOR smoothing, the modified-Gram--Schmidt orthogonalization whose serial
rank-one projector chain buys the roundoff-orthogonality the batched form loses---are
developed in Saad's \emph{Iterative Methods for Sparse Linear
Systems}~\citep{saad2003}, the standard reference for the Krylov and multigrid loops
behind the single-loop and nested-loop obstructions of \cref{sec:appobstruction}.
Trefethen and Bau's \emph{Numerical Linear Algebra}~\citep{trefethenbau1997} develops
orthogonalization, conditioning, and the eigenvalue problem lecture by lecture, and is
the ideal companion for seeing \emph{why} the eigensolve loop is the canonical
opaque-library obstruction---why the iteration is best left to a specialized library
and marked at the boundary rather than re-derived. The simulator whose five-layer
dissection this map summarizes is Palace~\citep{palace2023}; reading its source against
this stack, with the four rotations of \cref{sec:approtations} and the obstruction
catalogue of \cref{sec:appobstruction} in hand, is exactly the auditable
ground-to-calculus walk the layered method was built to make possible.

\chapter{Symbol Glossary and Guide to the Literature}
\label{app:glossary}

\chapterorient{This closing appendix is a reference, not a narrative. Its first
half is a glossary: every symbol the book uses often, grouped by the world it
lives in---fields, operators, function spaces, numerical methods, the
pseudo-language, and the layered methodology---with a pointer to the chapter that
introduced it. Its second half is a guide to the literature, a topic-organized
map of the sources the book leans on, each with a one-line note on what to read
it for. The notation primer at the front of the book introduced these symbols a
few at a time; here they are collected so you can look one up. When a symbol is
unfamiliar, start here.}

This appendix expands the front-matter notation primer (and its
\cref{tab:symbols}) into a full reference. The mathematical symbols are set the
usual way; the pseudo-language symbols obey the one standing convention of the
book---they belong in a shaded box, never inside a formula, because the
shape-group sigil \texttt{\$} clashes with mathematics typesetting. In the
tables below, pseudo-language identifiers are set in \texttt{typewriter}; where a
shape carries the sigil we write it as \texttt{Tensor[(S:\,\dots)]} and discuss
the \texttt{\$S} back-reference form in prose.

\section{Symbol and notation glossary}

The glossary is split into six categories, one cluster of tables each: physical
fields and quantities; discrete operators and matrices; function spaces and the
finite-element vocabulary; the symbols of the numerical methods; the
pseudo-language and calculus; and the layered methodology. Each row gives the
symbol, its meaning, and a cross-reference to the chapter where it is introduced
and used.

\subsection{Fields and physical quantities}

These are the quantities Maxwell's equations relate, the materials they move
through, and the scalar readouts a simulation finally reports. Vector fields are
upright bold; material coefficients and scalars are italic.

\begin{table}[ht]
\centering
\small
\caption[Fields and material coefficients]{Electromagnetic fields, potentials,
and the constitutive coefficients of the medium.}
\label{tab:gl-fields}
\begin{tabular}{@{}lll@{}}
\toprule
Symbol & Meaning & Where introduced \\
\midrule
$\vE$ & electric field & \cref{ch:maxwell} \\
$\vB$ & magnetic flux density & \cref{ch:maxwell} \\
$\vH$ & magnetic field intensity, $\vH=\nu\vB$ & \cref{ch:maxwell} \\
$\vD$ & electric displacement, $\vD=\varepsilon\vE$ & \cref{ch:maxwell} \\
$\vJ$ & (free) current density & \cref{ch:maxwell} \\
$\vA$ & magnetic vector potential, $\vB=\Curl\vA$ & \cref{ch:magnetostatics} \\
$V$ & electrostatic scalar potential, $\vE=-\Grad V$ & \cref{ch:electrostatics} \\
$\varepsilon$ & electric permittivity & \cref{ch:maxwell} \\
$\mu$ & magnetic permeability & \cref{ch:maxwell} \\
$\nu$ & magnetic reluctivity, $\nu=\mu^{-1}$ & \cref{ch:weak} \\
$\sigma$ & electrical conductivity (loss term) & \cref{ch:maxwell} \\
\bottomrule
\end{tabular}
\end{table}

\begin{table}[ht]
\centering
\small
\caption[Scalar outputs and modal quantities]{Frequencies, modal quantities, and
the scalar products a driver reports.}
\label{tab:gl-outputs}
\begin{tabular}{@{}lll@{}}
\toprule
Symbol & Meaning & Where introduced \\
\midrule
$\omega$ & angular frequency, $\omega=2\pi f$ & \cref{ch:maxwell} \\
$f$ & (ordinary) frequency & \cref{ch:driven} \\
$\lambda$ & eigenvalue (here $\lambda=\omega^2$, or a shifted variant) & \cref{ch:eigenvalues} \\
$k_n$ & wavenumber / propagation constant of mode $n$ & \cref{ch:waveports} \\
$Q$ & quality factor of a resonant mode, $Q=\omega/(2\,\Im\omega)$ & \cref{ch:eigenmodes} \\
$S_{ij}$ & scattering-matrix entry (port $j$ in, port $i$ out) & \cref{ch:driven} \\
$C$ & capacitance (matrix) & \cref{ch:electrostatics} \\
$M$ & inductance (matrix) & \cref{ch:magnetostatics} \\
$U$ & stored field energy, electric or magnetic & \cref{ch:reductions-pde} \\
\bottomrule
\end{tabular}
\end{table}

\begin{notationbox}
The letter $\sigma$ does double duty: it is the conductivity in
\cref{tab:gl-fields} (a physical material property) and the spectral
\emph{shift} of a shift-and-invert eigensolve in \cref{tab:gl-numerics} (a purely
algebraic device). The context---a constitutive law versus an eigenvalue
transform---always disambiguates, and the two never appear in the same equation.
\end{notationbox}

\subsection{Discrete operators and matrices}

Assembling a weak form over a mesh turns each continuous bilinear form into a
sparse matrix. These are set sans-serif to mark them as discrete operators acting
on vectors of degrees of freedom, as opposed to the continuous differential
operators of \cref{tab:gl-spaces}.

\begin{table}[ht]
\centering
\small
\caption[Discrete operators]{The assembled matrices and the maps between
discretizations. Sans-serif marks a discrete (matrix) operator.}
\label{tab:gl-operators}
\begin{tabular}{@{}lll@{}}
\toprule
Symbol & Meaning & Where introduced \\
\midrule
$\mat{K}$ & stiffness operator (curl--curl or grad--grad) & \cref{ch:weak} \\
$\mat{M}$ & mass operator & \cref{ch:weak} \\
$\mat{C}$ & damping / conductivity operator & \cref{ch:weak} \\
$\mat{A}$ & a generic system matrix, $\mat{A}\vx=\vb$ & \cref{ch:cg} \\
$\mat{A}(\omega)$ & frequency-dependent system matrix of a driven solve & \cref{ch:driven} \\
$\mat{P}$ & prolongation (coarse$\to$fine) or a projector & \cref{ch:multigrid} \\
$\mat{G}$ & discrete gradient (the $\Hone\to\Hcurl$ map); a restriction & \cref{ch:smoothers} \\
$a(u,v)$ & a bilinear (weak) form, the source of $\mat{K},\mat{M},\mat{C}$ & \cref{ch:weak} \\
\bottomrule
\end{tabular}
\end{table}

The driven system matrix is the recurring composite
$\mat{A}(\omega)=\mat{K}+i\omega\mat{C}-\omega^2\mat{M}$: one matrix-valued
function of frequency, built once from the three assembled operators and
re-evaluated across a sweep (\cref{ch:driven}). The discrete gradient $\mat{G}$
is the algebraic shadow of $\Grad$---it is what makes the Hiptmair smoother of
\cref{ch:smoothers} possible, by carrying a residual from the edge space into the
nodal space where the near-null-space lives.

\subsection{Function spaces and the finite-element vocabulary}

The four de Rham spaces and the operators that connect them are the backbone of
\cref{part:fieldtheory}; the discretization symbols ($N$, $\eta_K$, $\theta$) and
the element-local axes come from \cref{part:machinery}.

\begin{table}[ht]
\centering
\small
\caption[Function spaces and continuous operators]{The de Rham complex and its
differential operators.}
\label{tab:gl-spaces}
\begin{tabular}{@{}lll@{}}
\toprule
Symbol & Meaning & Where introduced \\
\midrule
$\Hone$ & scalar potentials; nodal elements & \cref{ch:derham} \\
$\Hcurl$ & curl-conforming (edge) elements; tangential continuity & \cref{ch:derham} \\
$\Hdiv$ & div-conforming (face) elements; normal continuity & \cref{ch:derham} \\
$\Ltwo$ & square-integrable fields; no continuity imposed & \cref{ch:derham} \\
$\Grad$ & gradient, $\Hone\to\Hcurl$ & \cref{ch:derham} \\
$\Curl$ & curl, $\Hcurl\to\Hdiv$ & \cref{ch:derham} \\
$\Div$ & divergence, $\Hdiv\to\Ltwo$ & \cref{ch:derham} \\
$\diffop$ & a generic first-order differential operator in a weak form & \cref{ch:weak} \\
\bottomrule
\end{tabular}
\end{table}

\begin{table}[ht]
\centering
\small
\caption[Discretization and adaptivity symbols]{Degree-of-freedom counts, error
indicators, and the element-local axes.}
\label{tab:gl-fem}
\begin{tabular}{@{}lll@{}}
\toprule
Symbol & Meaning & Where introduced \\
\midrule
$N$ & number of (true) degrees of freedom & \cref{ch:fem} \\
$\eta_K$ & local error indicator on element $K$ & \cref{ch:amr} \\
$\theta$ & D\"orfler marking fraction & \cref{ch:amr} \\
E & element axis (which element) & \cref{ch:refelement} \\
L & local-node axis (dof within an element) & \cref{ch:refelement} \\
P & quadrature-point axis & \cref{ch:matrixfree} \\
C & component axis (vector-field component) & \cref{ch:refelement} \\
G & global-dof axis (after gather/scatter) & \cref{ch:assembly} \\
\bottomrule
\end{tabular}
\end{table}

The five element-local axes \texttt{E/L/P/C/G} are the named shape run of the
matrix-free operator: \cref{ch:matrixfree} writes a quadrature contraction as a
sequence of moves along these axes, and \cref{ch:assembly} shows the gather and
scatter that fold the local axis \texttt{L} into the global axis \texttt{G}.

\subsection{Numerical-method symbols}

The Krylov, eigensolver, and multigrid chapters of \cref{part:machinery} share a
compact algebraic vocabulary. The same letters recur---$\alpha$ and $\beta$ for
step coefficients, $\mat{H}$ and $\mat{T}$ for the small projected operators,
$\sigma$ for a spectral shift---so they are collected here.

\begin{table}[ht]
\centering
\small
\caption[Krylov and eigensolver symbols]{Subspaces, coefficients, and projected
operators of the iterative methods.}
\label{tab:gl-numerics}
\begin{tabular}{@{}lll@{}}
\toprule
Symbol & Meaning & Where introduced \\
\midrule
$\mathcal{K}_m$ & Krylov subspace of dimension $m$ & \cref{ch:krylov} \\
$\alpha$ & CG step length / Lanczos diagonal coefficient & \cref{ch:cg} \\
$\beta$ & CG conjugacy factor / Lanczos off-diagonal coefficient & \cref{ch:cg} \\
$\bar{\mat{H}}$ & upper-Hessenberg matrix (Arnoldi / GMRES) & \cref{ch:gmres} \\
$\mat{T}$ & tridiagonal matrix (Lanczos) & \cref{ch:krylovschur} \\
$\sigma$ & spectral shift in shift-and-invert & \cref{ch:eigenvalues} \\
$\kappa$ & condition number, $\kappa(\mat{A})$ & \cref{ch:cg} \\
$c,s$ & cosine / sine of a Givens rotation & \cref{ch:givens} \\
$\mat{V}$ & matrix of (orthonormal) basis vectors & \cref{ch:krylov} \\
\bottomrule
\end{tabular}
\end{table}

The Hessenberg $\bar{\mat{H}}$ and the tridiagonal $\mat{T}$ are the same object
seen from two angles: $\mat{T}$ is what $\bar{\mat{H}}$ collapses to when the
operator is symmetric, which is exactly why Lanczos is the cheap special case of
Arnoldi (\cref{ch:krylovschur}). The Givens pair $(c,s)$ is the rotation that
GMRES applies to keep $\bar{\mat{H}}$ triangular incrementally
(\cref{ch:givens}); the \emph{V-cycle} (\cref{ch:multigrid}) is the recursive
coarsen--smooth--correct schedule that the smoothers of \cref{ch:smoothers}
plug into.

\subsection{The pseudo-language and calculus}

These are the symbols of the typed pseudo-language in which the book writes its
algorithms. They always appear in a shaded box. The combinators and named
operators are the verbs of the calculus; the records are its nouns; the
\texttt{Solve} monad threads the one piece of mutable state.

\begin{table}[ht]
\centering
\small
\caption[Calculus types and shapes]{The type and shape vocabulary of the
pseudo-language.}
\label{tab:gl-calc-types}
\begin{tabular}{@{}lll@{}}
\toprule
Identifier & Meaning & Where introduced \\
\midrule
\texttt{Tensor[N]} & flat vector of $N$ dofs & \cref{ch:tensors} \\
\texttt{Tensor[(S:\,\dots)]} & tensor with a named, rank-agnostic shape run \texttt{S} & \cref{ch:shapes} \\
\texttt{Tensor[\$S]} & a back-reference: ``the same shape \texttt{S} again'' & \cref{ch:shapes} \\
\texttt{Op[\dots]} & a (matrix-free) linear-operator value & \cref{ch:matrixfree} \\
\texttt{Scalar} & a rank-0 tensor (a number with shape) & \cref{ch:tensors} \\
\bottomrule
\end{tabular}
\end{table}

The shape run is the heart of the shape system (\cref{ch:shapes}): a function
typed \texttt{Tensor[(S:\,\dots)] -> Tensor[\$S]} works on a vector, a matrix, or
a higher array at once, and the \texttt{\$S} back-reference promises the output
carries the same axes as the input. We write the sigil form only inside the
shaded box; in prose we name it ``the shape group \texttt{S}.''

\begin{table}[ht]
\centering
\small
\caption[Combinators and named operators]{The loop combinator and the named solve
verbs of the synthesis surface.}
\label{tab:gl-calc-ops}
\begin{tabular}{@{}lll@{}}
\toprule
Identifier & Meaning & Where introduced \\
\midrule
\texttt{iterate\_while} & the one loop combinator (predicate + step) & \cref{ch:fold} \\
\texttt{fold\_solve} & a linear solve expressed as a fold over Krylov steps & \cref{ch:krylov} \\
\texttt{ksp\_solve} & the preconditioned Krylov solve verb & \cref{ch:precond} \\
\texttt{solve\_family} & one operator, many right-hand sides (reused factorization) & \cref{ch:electrostatics} \\
\texttt{frequency\_sweep} & a driven solve evaluated across a band of $\omega$ & \cref{ch:driven} \\
\texttt{eigsolve} & the eigenvalue-problem solve verb & \cref{ch:eigenvalues} \\
\bottomrule
\end{tabular}
\end{table}

\begin{table}[ht]
\centering
\small
\caption[Records and effects]{The state records, the \texttt{Solve} monad, and
the two edge kinds of the dependency graph.}
\label{tab:gl-calc-records}
\begin{tabular}{@{}lll@{}}
\toprule
Identifier & Meaning & Where introduced \\
\midrule
\texttt{IoData} & the parsed configuration record (the run's inputs) & \cref{ch:records} \\
\texttt{OpParams} & the assembled-operator parameter bundle & \cref{ch:records} \\
\texttt{SimState} & the evolving simulation state the monad threads & \cref{ch:monad} \\
\texttt{SolveResult} & a solve's output record (solution + diagnostics) & \cref{ch:records} \\
\texttt{Solve} & the state monad threading \texttt{SimState} & \cref{ch:monad} \\
\texttt{\#extern} & an opaque external kernel (no body; type only) & \cref{ch:matrixfree} \\
\texttt{depends-on} & a blocking edge: constrains resolution and liveness & \cref{ch:wholemachine} \\
\texttt{reference} & a navigational (free) edge: a cross-link, not a dependency & \cref{ch:wholemachine} \\
\bottomrule
\end{tabular}
\end{table}

The \texttt{Solve} monad is the one place state changes: a \texttt{do} block in
the \texttt{Solve} monad threads a \texttt{SimState} through a sequence of steps,
and \texttt{modify} is the single state-mutating verb (\cref{ch:monad}). An
\texttt{\#extern} marks a leaf the synthesis surface does not implement---a call
out to an opaque library kernel (libCEED quadrature, a vendor eigensolve), shown
by its type signature only (\cref{ch:matrixfree}).

\subsection{The layered stack and methodology}

The book's organizing idea is a stack of representations connected by rotations.
These tokens name its levels and its recurring structural patterns.

\begin{table}[ht]
\centering
\small
\caption[The layered stack]{Levels, rotations, and the structural categories of
the methodology.}
\label{tab:gl-method}
\begin{tabular}{@{}lll@{}}
\toprule
Token & Meaning & Where introduced \\
\midrule
L0 & ground-truth source level (cited ranges) & \cref{ch:strata} \\
L1 & pure-function level (mutation rotated out) & \cref{ch:strata} \\
L2 & fused-algebra level (HPC tricks erased) & \cref{ch:strata} \\
L3 & iteration level (global tensor-field ops) & \cref{ch:strata} \\
L4 & calculus level (combinators + monad + tensors) & \cref{ch:strata} \\
rotation & a representation-preserving rewrite between levels & \cref{ch:rotations} \\
obstruction & a place a rotation cannot complete cleanly & \cref{ch:rotations} \\
\bottomrule
\end{tabular}
\end{table}

A \emph{rotation} re-expresses a level in a new vocabulary while preserving what
it computes; an \emph{obstruction} is an honest record of a place that resists
the rotation (a sequential recurrence that will not become a tensor-field op, an
opaque library boundary). Two small taxonomies recur. The \emph{four solve
corners} (\cref{ch:composing}) are the cross product of two axes---one
right-hand side versus a family, and a fixed operator versus a frequency-varying
one---and they classify the five drivers: electrostatic and magnetostatic are
family-of-RHS / fixed-operator corners, the driven solve is the
varying-operator corner, and so on. The \emph{three reduction shapes}
(\cref{ch:reductions-pde}) are the three ways a field is collapsed to a
reported number---a global scalar (energy), a boundary integral (a port flux),
and a pointwise probe (a field sample at a point)---one shape per family of
output product.

\begin{keyidea}
Every symbol in this glossary lives at a definite level of the stack. The
physical fields ($\vE$, $\vB$) and the function spaces ($\Hcurl$, $\Ltwo$) name
the continuous problem; the assembled operators ($\mat{K}$, $\mat{M}$) and the
Krylov symbols ($\alpha$, $\bar{\mat{H}}$) name the discrete machinery; the
pseudo-language identifiers (\texttt{iterate\_while}, \texttt{Solve}) name the
calculus that rewrites that machinery into a synthesizable surface. Reading a
symbol's level is half of reading the symbol.
\end{keyidea}

\section{A guide to the literature}
\label{sec:furtherreading-app}

The book is a synthesis, and it stands on a small, deliberately chosen
bibliography. This section groups those sources by topic and says, in a line,
what to read each for. It is a map for going deeper, not a survey.

\paragraph{Finite elements and Maxwell's equations.}
For the finite-element treatment of electromagnetics, \citet{monk2003} is the
standard reference and the closest in spirit to the field-theory part of this
book---read it for the rigorous theory of edge elements and $\Hcurl$
conformity. \citet{jin2014} is the more applied companion, rich in engineering
detail and worked formulations; read it for how the weak forms become real
solvers. \citet{brennerscott2008} is the mathematical foundation of the finite
element method in general (approximation theory, error estimates); read it when
you want the analysis under the discretization. \citet{arnold2006feec} is the
deep structural account: finite element exterior calculus explains \emph{why}
the de Rham complex of \cref{ch:derham} is the right organizing principle, and
why the four spaces must connect as they do.

\paragraph{Krylov methods and iterative linear solvers.}
\citet{saad2003} is the encyclopedic reference for sparse iterative
methods---read it for CG, GMRES, preconditioning, and the Krylov-subspace theory
behind \cref{ch:krylov,ch:cg,ch:gmres}. \citet{trefethenbau1997} is the elegant
short course; read it for the cleanest possible derivations of QR, Arnoldi, and
the conditioning ideas ($\kappa$). \citet{golubvanloan2013} is the matrix-computations
handbook; read it for Givens rotations (\cref{ch:givens}), the dense linear
algebra under the small projected problems, and as a definitive cross-reference.

\paragraph{Eigensolvers.}
\citet{saad2011eig} is the companion volume to Saad's linear-systems book and the
reference for large-scale eigenvalue computation---read it for Lanczos,
Arnoldi, and the shift-and-invert ($\sigma$) transforms of
\cref{ch:eigenvalues}. \citet{stewart2001krylovschur} introduced the
Krylov--Schur restarting scheme of \cref{ch:krylovschur}; read it for the
stable thick-restart mechanism. \citet{effenberger2013} treats the nonlinear
eigenvalue problem and the successive-deflation strategy behind
\cref{ch:deflation}; read it when the eigenproblem's operator depends on its own
eigenvalue.

\paragraph{Multigrid and preconditioning.}
\citet{hiptmair1998} is the foundational paper on multigrid for Maxwell's
equations and the source of the Hiptmair smoother (\cref{ch:smoothers})---read
it for why curl-conforming problems need a special near-null-space correction.
\citet{hiptmairxu2007} generalizes this into nodal auxiliary-space
preconditioning (the AMS idea); read it for the auxiliary-space construction
that makes $\Hcurl$ and $\Hdiv$ solves scalable. \citet{phillipsfischer2022}
gives the optimal Chebyshev smoothers and one-sided V-cycle schedule of
\cref{ch:multigrid}; read it for the practical smoother tuning.

\paragraph{Adaptivity and error estimation.}
\citet{dorfler1996} introduced the marking strategy ($\theta$) that drives the
adaptive loop of \cref{ch:amr}---read it for the convergence guarantee behind
``refine the smallest set of elements holding a fixed fraction of the error.''
\citet{zienkiewiczzhu1987} is the classic recovery-based error estimator
($\eta_K$); read it for the practical, inexpensive indicator that tells the loop
\emph{where} to refine.

\paragraph{Microwave engineering and physics.}
\citet{pozar2011} is the microwave-engineering reference and the source of the
S-parameter and wave-port conventions ($S_{ij}$, $k_n$) of
\cref{ch:driven,ch:waveports}---read it for what the driven solve's outputs
mean physically. \citet{griffiths2017} is the undergraduate electrodynamics
text; read it to refresh Maxwell's equations, the potentials, and the energy
integrals of \cref{ch:maxwell} from first principles.

\paragraph{Functional programming and the calculus.}
\citet{peytonjones2003} is the Haskell language report and the reference for the
typed-functional notation the pseudo-language borrows (currying, records, type
signatures)---read it for the precise meaning of the syntax in the shaded
boxes. \citet{wadler1995monads} is the classic exposition of monads for
functional programming; read it for the idea that makes the \texttt{Solve} monad
of \cref{ch:monad} legitimate---state threaded purely, with one mutation point.

\paragraph{The tools.}
\citet{palace2023} is Palace itself, the AWS Labs finite-element electromagnetics
solver this book dissects; it is the L0 ground truth that every citation in the
book ultimately points at. \citet{ceed2021} is libCEED, the matrix-free
high-order element library behind the \texttt{\#extern} quadrature kernel of
\cref{ch:matrixfree}; read it for the element-local algebra (the E/L/P/C/G axes)
the synthesis surface lifts into tensor contractions.

\summarysection
You now have the whole book in three forms. You have the \emph{map}---the
layered stack L0 through L4, the rotations that connect the levels, and the
honest obstructions where a rotation cannot finish. You have the
\emph{machinery}---the weak forms and de Rham spaces of the field theory, the
Krylov, eigensolver, and multigrid algorithms of the numerical core, and the
five drivers that compose them into electrostatic, magnetostatic, eigenmode,
driven, and transient simulations. And you have the \emph{vocabulary}---the
symbols of this glossary and the typed pseudo-language they belong to, in which
the synthesis surface re-expresses a quarter-million lines of tuned C++ as a few
combinators over immutable tensors. The simulator is no longer a black box: it is
a small calculus, a handful of named operators, and a clear account of what each
one computes and why. That account was the goal; you are now holding it.

% \part{The Abstract Framework}      ...
% \part{The Numerical Machinery}     ...
% \part{The Constituent Modalities}  ...
% \part{The Collection}              ...
% \appendix \part{Appendices}        ...

\backmatter
\bibliographystyle{plainnat}
\bibliography{references}

\end{document}
